\chapter{Overview}
This is the \emph{full} declaration-level dependency graph of \texttt{pphi2}
together with its upstream dependencies in our own libraries
(\texttt{gaussian-field}, \texttt{gaussian-hilbert}, \dots), covering 2833
declarations and their real dependency edges. Edges are \emph{extracted} from
the compiled Lean environment (\texttt{scripts/BlueprintGraph.lean}), not
hand-drawn. Informal statements are reused from \texttt{summary/**/*.md} where
available (otherwise the node shows its Lean name). Declarations flagged
\textbf{(axiom)} are assumed, not proved.

\chapter{gaussian-field (upstream)}
\section{\texttt{SmoothCircle.Restriction}}
\begin{definition}[\texttt{circleRestrictionLM}]\label{GaussianField-circleRestrictionLM}
  \lean{GaussianField.circleRestrictionLM}\leanok
  \uses{GaussianField-SmoothMap-Circle-instFunLike, GaussianField-SmoothMap-Circle-instAddCommGroup, GaussianField-SmoothMap-Circle-instModule, GaussianField-SmoothMap-Circle, GaussianField-circleSpacing, GaussianField-circlePoint}
  \texttt{GaussianField.circleRestrictionLM}
\end{definition}
\begin{theorem}[\texttt{circleSpacing\_pos}]\label{GaussianField-circleSpacing-pos}
  \lean{GaussianField.circleSpacing_pos}\leanok
  \uses{GaussianField-circleSpacing}
  \texttt{GaussianField.circleSpacing\_pos}
\end{theorem}
\begin{definition}[\texttt{circleSpacing}]\label{GaussianField-circleSpacing}
  \lean{GaussianField.circleSpacing}\leanok
  \texttt{GaussianField.circleSpacing}
\end{definition}
\begin{theorem}[\texttt{circleRestriction\_sq\_sum}]\label{GaussianField-circleRestriction-sq-sum}
  \lean{GaussianField.circleRestriction_sq_sum}\leanok
  \uses{GaussianField-SmoothMap-Circle-instFunLike, GaussianField-SmoothMap-Circle-instAddCommGroup, GaussianField-circleRestriction, GaussianField-SmoothMap-Circle-instModule, GaussianField-SmoothMap-Circle, GaussianField-SmoothMap-Circle-instTopologicalSpace, GaussianField-circlePoint}
  \texttt{GaussianField.circleRestriction\_sq\_sum}
\end{theorem}
\begin{theorem}[\texttt{circleRestrictionGJ\_le}]\label{GaussianField-circleRestrictionGJ-le}
  \lean{GaussianField.circleRestrictionGJ_le}\leanok
  \uses{GaussianField-SmoothMap-Circle-instFunLike, GaussianField-SmoothMap-Circle-instAddCommGroup, GaussianField-SmoothMap-Circle-norm-iteratedDeriv-le-sobolevSeminorm-, GaussianField-SmoothMap-Circle-instSMul, GaussianField-circleRestrictionGJ, GaussianField-SmoothMap-Circle-instModule, GaussianField-SmoothMap-Circle, GaussianField-circleSpacing, GaussianField-SmoothMap-Circle-sobolevSeminorm, GaussianField-circleRestrictionGJ-apply, GaussianField-SmoothMap-Circle-instTopologicalSpace, GaussianField-circlePoint, Lattice-toSemilatticeInf, GaussianField-circleSpacing-pos}
  \texttt{GaussianField.circleRestrictionGJ\_le}
\end{theorem}
\begin{theorem}[\texttt{circleRestrictionGJ\_eq}]\label{GaussianField-circleRestrictionGJ-eq}
  \lean{GaussianField.circleRestrictionGJ_eq}\leanok
  \uses{GaussianField-SmoothMap-Circle-instFunLike, GaussianField-SmoothMap-Circle-instAddCommGroup, GaussianField-circleRestrictionGJ, GaussianField-SmoothMap-Circle-instModule, GaussianField-SmoothMap-Circle, GaussianField-circleSpacing, GaussianField-circleRestrictionGJ-apply, GaussianField-SmoothMap-Circle-instTopologicalSpace, GaussianField-circlePoint}
  \texttt{GaussianField.circleRestrictionGJ\_eq}
\end{theorem}
\begin{theorem}[\texttt{circleRestriction\_eq}]\label{GaussianField-circleRestriction-eq}
  \lean{GaussianField.circleRestriction_eq}\leanok
  \uses{GaussianField-SmoothMap-Circle-instFunLike, GaussianField-SmoothMap-Circle-instAddCommGroup, GaussianField-circleRestriction, GaussianField-SmoothMap-Circle-instModule, GaussianField-SmoothMap-Circle, GaussianField-circleSpacing, GaussianField-SmoothMap-Circle-instTopologicalSpace, GaussianField-circlePoint}
  \texttt{GaussianField.circleRestriction\_eq}
\end{theorem}
\begin{theorem}[\texttt{congr\_simp}]\label{GaussianField-circlePoint-congr-simp}
  \lean{GaussianField.circlePoint.congr_simp}\leanok
  \uses{GaussianField-circlePoint}
  \texttt{GaussianField.circlePoint.congr\_simp}
\end{theorem}
\begin{theorem}[\texttt{continuous\_eval\_at}]\label{GaussianField-continuous-eval-at}
  \lean{GaussianField.continuous_eval_at}\leanok
  \uses{GaussianField-SmoothMap-Circle-instFunLike, GaussianField-SmoothMap-Circle-smoothCircle-withSeminorms, GaussianField-SmoothMap-Circle-smul-apply, GaussianField-SmoothMap-Circle-instAddCommGroup, GaussianField-SmoothMap-Circle-add-apply, GaussianField-SmoothMap-Circle-norm-iteratedDeriv-le-sobolevSeminorm-, GaussianField-SmoothMap-Circle-instSMul, GaussianField-SmoothMap-Circle-instModule, GaussianField-SmoothMap-Circle, GaussianField-SmoothMap-Circle-sobolevSeminorm, GaussianField-SmoothMap-Circle-instTopologicalSpace}
  \texttt{GaussianField.continuous\_eval\_at}
\end{theorem}
\begin{theorem}[\texttt{congr\_simp}]\label{GaussianField-circleRestrictionGJ-congr-simp}
  \lean{GaussianField.circleRestrictionGJ.congr_simp}\leanok
  \uses{GaussianField-SmoothMap-Circle-instAddCommGroup, GaussianField-circleRestrictionGJ, GaussianField-SmoothMap-Circle-instModule, GaussianField-SmoothMap-Circle, GaussianField-SmoothMap-Circle-instTopologicalSpace}
  \texttt{GaussianField.circleRestrictionGJ.congr\_simp}
\end{theorem}
\begin{theorem}[\texttt{circleRestrictionGJ\_apply}]\label{GaussianField-circleRestrictionGJ-apply}
  \lean{GaussianField.circleRestrictionGJ_apply}\leanok
  \uses{GaussianField-SmoothMap-Circle-instFunLike, GaussianField-SmoothMap-Circle-instAddCommGroup, GaussianField-circleRestriction, GaussianField-circleRestrictionGJ, GaussianField-circleRestriction-apply, GaussianField-SmoothMap-Circle-instModule, GaussianField-SmoothMap-Circle, GaussianField-circleSpacing, GaussianField-SmoothMap-Circle-instTopologicalSpace, GaussianField-circlePoint, GaussianField-circleSpacing-pos}
  \texttt{GaussianField.circleRestrictionGJ\_apply}
\end{theorem}
\begin{theorem}[\texttt{congr\_simp}]\label{GaussianField-circleSpacing-congr-simp}
  \lean{GaussianField.circleSpacing.congr_simp}\leanok
  \uses{GaussianField-circleSpacing}
  \texttt{GaussianField.circleSpacing.congr\_simp}
\end{theorem}
\begin{theorem}[\texttt{circleSpacing\_eq}]\label{GaussianField-circleSpacing-eq}
  \lean{GaussianField.circleSpacing_eq}\leanok
  \uses{GaussianField-circleSpacing}
  \texttt{GaussianField.circleSpacing\_eq}
\end{theorem}
\begin{definition}[\texttt{circlePoint}]\label{GaussianField-circlePoint}
  \lean{GaussianField.circlePoint}\leanok
  \texttt{GaussianField.circlePoint}
\end{definition}
\begin{definition}[\texttt{circleRestriction}]\label{GaussianField-circleRestriction}
  \lean{GaussianField.circleRestriction}\leanok
  \uses{GaussianField-SmoothMap-Circle-instAddCommGroup, GaussianField-SmoothMap-Circle-instModule, GaussianField-SmoothMap-Circle, GaussianField-circleRestrictionLM, GaussianField-SmoothMap-Circle-instTopologicalSpace}
  \texttt{GaussianField.circleRestriction}
\end{definition}
\begin{theorem}[\texttt{circleRestriction\_apply}]\label{GaussianField-circleRestriction-apply}
  \lean{GaussianField.circleRestriction_apply}\leanok
  \uses{GaussianField-SmoothMap-Circle-instFunLike, GaussianField-SmoothMap-Circle-instAddCommGroup, GaussianField-circleRestriction, GaussianField-SmoothMap-Circle-instModule, GaussianField-SmoothMap-Circle, GaussianField-circleSpacing, GaussianField-SmoothMap-Circle-instTopologicalSpace, GaussianField-circlePoint}
  \texttt{GaussianField.circleRestriction\_apply}
\end{theorem}
\begin{definition}[\texttt{circleRestrictionGJ}]\label{GaussianField-circleRestrictionGJ}
  \lean{GaussianField.circleRestrictionGJ}\leanok
  \uses{GaussianField-SmoothMap-Circle-instAddCommGroup, GaussianField-circleRestriction, GaussianField-SmoothMap-Circle-instModule, GaussianField-SmoothMap-Circle, GaussianField-circleSpacing, GaussianField-SmoothMap-Circle-instTopologicalSpace}
  \texttt{GaussianField.circleRestrictionGJ}
\end{definition}
\begin{theorem}[\texttt{congr\_simp}]\label{GaussianField-circleRestriction-congr-simp}
  \lean{GaussianField.circleRestriction.congr_simp}\leanok
  \uses{GaussianField-SmoothMap-Circle-instAddCommGroup, GaussianField-circleRestriction, GaussianField-SmoothMap-Circle-instModule, GaussianField-SmoothMap-Circle, GaussianField-SmoothMap-Circle-instTopologicalSpace}
  \texttt{GaussianField.circleRestriction.congr\_simp}
\end{theorem}
\section{\texttt{Nuclear.NuclearTensorProduct}}
\begin{definition}[\texttt{pure}]\label{GaussianField-NuclearTensorProduct-pure}
  \lean{GaussianField.NuclearTensorProduct.pure}\leanok
  \uses{GaussianField-DyninMityaginSpace, GaussianField-DyninMityaginSpace-coeff, GaussianField-NuclearTensorProduct}
  \texttt{GaussianField.NuclearTensorProduct.pure}
\end{definition}
\begin{theorem}[\texttt{toPairIndex\_fromPairIndex}]\label{GaussianField-NuclearTensorProduct-toPairIndex-fromPairIndex}
  \lean{GaussianField.NuclearTensorProduct.toPairIndex_fromPairIndex}\leanok
  \uses{GaussianField-NuclearTensorProduct-toPairIndex, GaussianField-NuclearTensorProduct-fromPairIndex}
  \texttt{GaussianField.NuclearTensorProduct.toPairIndex\_fromPairIndex}
\end{theorem}
\begin{definition}[\texttt{instModule}]\label{GaussianField-RapidDecaySeq-instModule}
  \lean{GaussianField.RapidDecaySeq.instModule}\leanok
  \uses{GaussianField-RapidDecaySeq, GaussianField-RapidDecaySeq-instSMul, GaussianField-RapidDecaySeq-instAddCommGroup}
  \texttt{GaussianField.RapidDecaySeq.instModule}
\end{definition}
\begin{theorem}[\texttt{rapid\_decay}]\label{GaussianField-RapidDecaySeq-rapid-decay}
  \lean{GaussianField.RapidDecaySeq.rapid_decay}\leanok
  \uses{GaussianField-RapidDecaySeq, GaussianField-RapidDecaySeq-val}
  \texttt{GaussianField.RapidDecaySeq.rapid\_decay}
\end{theorem}
\begin{theorem}[\texttt{congr\_simp}]\label{GaussianField-RapidDecaySeq-mk-congr-simp}
  \lean{GaussianField.RapidDecaySeq.mk.congr_simp}\leanok
  \uses{GaussianField-RapidDecaySeq}
  \texttt{GaussianField.RapidDecaySeq.mk.congr\_simp}
\end{theorem}
\begin{definition}[\texttt{instModuleReal}]\label{GaussianField-NuclearTensorProduct-instModuleReal}
  \lean{GaussianField.NuclearTensorProduct.instModuleReal}\leanok
  \uses{GaussianField-NuclearTensorProduct, GaussianField-NuclearTensorProduct-instAddCommGroup}
  \texttt{GaussianField.NuclearTensorProduct.instModuleReal}
\end{definition}
\begin{definition}[\texttt{lift}]\label{GaussianField-NuclearTensorProduct-lift}
  \lean{GaussianField.NuclearTensorProduct.lift}\leanok
  \uses{GaussianField-DyninMityaginSpace, GaussianField-NuclearTensorProduct-instModuleReal, GaussianField-NuclearTensorProduct-instTopologicalSpace, GaussianField-DyninMityaginSpace--, GaussianField-NuclearTensorProduct, GaussianField-NuclearTensorProduct-instAddCommGroup, GaussianField-DyninMityaginSpace-p}
  \texttt{GaussianField.NuclearTensorProduct.lift}
\end{definition}
\begin{definition}[\texttt{RapidDecaySeq}]\label{GaussianField-RapidDecaySeq}
  \lean{GaussianField.RapidDecaySeq}\leanok
  \texttt{GaussianField.RapidDecaySeq}
\end{definition}
\begin{definition}[\texttt{instAddCommGroup}]\label{GaussianField-RapidDecaySeq-instAddCommGroup}
  \lean{GaussianField.RapidDecaySeq.instAddCommGroup}\leanok
  \uses{GaussianField-RapidDecaySeq-instAdd, GaussianField-RapidDecaySeq-instZero, GaussianField-RapidDecaySeq, GaussianField-RapidDecaySeq-instNeg}
  \texttt{GaussianField.RapidDecaySeq.instAddCommGroup}
\end{definition}
\begin{definition}[\texttt{val}]\label{GaussianField-RapidDecaySeq-val}
  \lean{GaussianField.RapidDecaySeq.val}\leanok
  \uses{GaussianField-RapidDecaySeq}
  \texttt{GaussianField.RapidDecaySeq.val}
\end{definition}
\begin{definition}[\texttt{evalCLM}]\label{GaussianField-NuclearTensorProduct-evalCLM}
  \lean{GaussianField.NuclearTensorProduct.evalCLM}\leanok
  \uses{GaussianField-DyninMityaginSpace, GaussianField-NuclearTensorProduct-instModuleReal, GaussianField-NuclearTensorProduct-instTopologicalSpace, GaussianField-DyninMityaginSpace--, GaussianField-NuclearTensorProduct, GaussianField-NuclearTensorProduct-instAddCommGroup, GaussianField-NuclearTensorProduct-lift, GaussianField-DyninMityaginSpace-p}
  \texttt{GaussianField.NuclearTensorProduct.evalCLM}
\end{definition}
\begin{theorem}[\texttt{summable\_inv\_one\_add\_sq}]\label{GaussianField-NuclearTensorProduct-summable-inv-one-add-sq}
  \lean{GaussianField.NuclearTensorProduct.summable_inv_one_add_sq}\leanok
  \texttt{GaussianField.NuclearTensorProduct.summable\_inv\_one\_add\_sq}
\end{theorem}
\begin{theorem}[\texttt{nat\_unpair\_le}]\label{GaussianField-nat-unpair-le}
  \lean{GaussianField.nat_unpair_le}\leanok
  \texttt{GaussianField.nat\_unpair\_le}
\end{theorem}
\begin{theorem}[\texttt{ext}]\label{GaussianField-NuclearTensorProduct-ext}
  \lean{GaussianField.NuclearTensorProduct.ext}\leanok
  \uses{GaussianField-RapidDecaySeq-val, GaussianField-NuclearTensorProduct, GaussianField-RapidDecaySeq-ext}
  \texttt{GaussianField.NuclearTensorProduct.ext}
\end{theorem}
\begin{definition}[\texttt{ofRapidDecayEquiv\_hasBiorthogonalBasis}]\label{GaussianField-DyninMityaginSpace-ofRapidDecayEquiv-hasBiorthogonalBasis}
  \lean{GaussianField.DyninMityaginSpace.ofRapidDecayEquiv_hasBiorthogonalBasis}\leanok
  \uses{GaussianField-RapidDecaySeq-mk-congr-simp, GaussianField-DyninMityaginSpace-coeff, GaussianField-DyninMityaginSpace-ofRapidDecayEquiv, GaussianField-RapidDecaySeq, GaussianField-RapidDecaySeq-val, GaussianField-RapidDecaySeq-instModule, GaussianField-RapidDecaySeq-basisVec, GaussianField-DyninMityaginSpace-basis, GaussianField-RapidDecaySeq-instAddCommGroup, GaussianField-RapidDecaySeq-instTopologicalSpace, GaussianField-DyninMityaginSpace-HasBiorthogonalBasis, GaussianField-RapidDecaySeq-coeffCLM}
  \texttt{GaussianField.DyninMityaginSpace.ofRapidDecayEquiv\_hasBiorthogonalBasis}
\end{definition}
\begin{definition}[\texttt{NuclearTensorProduct}]\label{GaussianField-NuclearTensorProduct}
  \lean{GaussianField.NuclearTensorProduct}\leanok
  \uses{GaussianField-RapidDecaySeq}
  \texttt{GaussianField.NuclearTensorProduct}
\end{definition}
\begin{definition}[\texttt{coeffCLM}]\label{GaussianField-RapidDecaySeq-coeffCLM}
  \lean{GaussianField.RapidDecaySeq.coeffCLM}\leanok
  \uses{GaussianField-RapidDecaySeq, GaussianField-RapidDecaySeq-instModule, GaussianField-RapidDecaySeq-coeffLM, GaussianField-RapidDecaySeq-instAddCommGroup, GaussianField-RapidDecaySeq-instTopologicalSpace}
  \texttt{GaussianField.RapidDecaySeq.coeffCLM}
\end{definition}
\begin{theorem}[\texttt{evalCLM\_pure}]\label{GaussianField-NuclearTensorProduct-evalCLM-pure}
  \lean{GaussianField.NuclearTensorProduct.evalCLM_pure}\leanok
  \uses{GaussianField-DyninMityaginSpace, GaussianField-NuclearTensorProduct-instModuleReal, GaussianField-NuclearTensorProduct-instTopologicalSpace, GaussianField-NuclearTensorProduct-lift-pure, GaussianField-DyninMityaginSpace--, GaussianField-NuclearTensorProduct, GaussianField-NuclearTensorProduct-instAddCommGroup, GaussianField-NuclearTensorProduct-pure, GaussianField-NuclearTensorProduct-lift, GaussianField-DyninMityaginSpace-p, GaussianField-NuclearTensorProduct-evalCLM}
  \texttt{GaussianField.NuclearTensorProduct.evalCLM\_pure}
\end{theorem}
\begin{definition}[\texttt{rapidDecaySeminorm}]\label{GaussianField-RapidDecaySeq-rapidDecaySeminorm}
  \lean{GaussianField.RapidDecaySeq.rapidDecaySeminorm}\leanok
  \uses{GaussianField-RapidDecaySeq, GaussianField-RapidDecaySeq-val, GaussianField-RapidDecaySeq-instSMul, GaussianField-RapidDecaySeq-instAddCommGroup}
  \texttt{GaussianField.RapidDecaySeq.rapidDecaySeminorm}
\end{definition}
\begin{theorem}[\texttt{rapidDecay\_t1Space}]\label{GaussianField-RapidDecaySeq-rapidDecay-t1Space}
  \lean{GaussianField.RapidDecaySeq.rapidDecay_t1Space}\leanok
  \uses{GaussianField-RapidDecaySeq-rapidDecaySeminorm, GaussianField-RapidDecaySeq, GaussianField-RapidDecaySeq-val, GaussianField-RapidDecaySeq-instSMul, GaussianField-RapidDecaySeq-instModule, GaussianField-RapidDecaySeq-ext, GaussianField-RapidDecaySeq-instAddCommGroup, GaussianField-RapidDecaySeq-instTopologicalSpace, GaussianField-RapidDecaySeq-weight-nonneg, GaussianField-RapidDecaySeq-rapidDecay-withSeminorms, GaussianField-RapidDecaySeq-rapid-decay, Lattice-toSemilatticeInf}
  \texttt{GaussianField.RapidDecaySeq.rapidDecay\_t1Space}
\end{theorem}
\begin{theorem}[\texttt{instIsTopologicalAddGroup}]\label{GaussianField-RapidDecaySeq-instIsTopologicalAddGroup}
  \lean{GaussianField.RapidDecaySeq.instIsTopologicalAddGroup}\leanok
  \uses{GaussianField-RapidDecaySeq-rapidDecaySeminorm, GaussianField-RapidDecaySeq, GaussianField-RapidDecaySeq-instModule, GaussianField-RapidDecaySeq-instAddCommGroup, GaussianField-RapidDecaySeq-instTopologicalSpace, GaussianField-RapidDecaySeq-rapidDecay-withSeminorms}
  \texttt{GaussianField.RapidDecaySeq.instIsTopologicalAddGroup}
\end{theorem}
\begin{definition}[\texttt{instAddCommGroup}]\label{GaussianField-NuclearTensorProduct-instAddCommGroup}
  \lean{GaussianField.NuclearTensorProduct.instAddCommGroup}\leanok
  \uses{GaussianField-NuclearTensorProduct}
  \texttt{GaussianField.NuclearTensorProduct.instAddCommGroup}
\end{definition}
\begin{definition}[\texttt{instZero}]\label{GaussianField-RapidDecaySeq-instZero}
  \lean{GaussianField.RapidDecaySeq.instZero}\leanok
  \uses{GaussianField-RapidDecaySeq}
  \texttt{GaussianField.RapidDecaySeq.instZero}
\end{definition}
\begin{definition}[\texttt{instTopologicalSpace}]\label{GaussianField-NuclearTensorProduct-instTopologicalSpace}
  \lean{GaussianField.NuclearTensorProduct.instTopologicalSpace}\leanok
  \uses{GaussianField-NuclearTensorProduct}
  \texttt{GaussianField.NuclearTensorProduct.instTopologicalSpace}
\end{definition}
\begin{definition}[\texttt{instSMul}]\label{GaussianField-RapidDecaySeq-instSMul}
  \lean{GaussianField.RapidDecaySeq.instSMul}\leanok
  \uses{GaussianField-RapidDecaySeq, GaussianField-RapidDecaySeq-val}
  \texttt{GaussianField.RapidDecaySeq.instSMul}
\end{definition}
\begin{theorem}[\texttt{zero\_val}]\label{GaussianField-RapidDecaySeq-zero-val}
  \lean{GaussianField.RapidDecaySeq.zero_val}\leanok
  \uses{GaussianField-RapidDecaySeq-instZero, GaussianField-RapidDecaySeq, GaussianField-RapidDecaySeq-val}
  \texttt{GaussianField.RapidDecaySeq.zero\_val}
\end{theorem}
\begin{theorem}[\texttt{hasSum\_basisVec}]\label{GaussianField-RapidDecaySeq-hasSum-basisVec}
  \lean{GaussianField.RapidDecaySeq.hasSum_basisVec}\leanok
  \uses{GaussianField-RapidDecaySeq-rapidDecaySeminorm, GaussianField-RapidDecaySeq, GaussianField-RapidDecaySeq-val, GaussianField-RapidDecaySeq-instSMul, GaussianField-RapidDecaySeq-instModule, GaussianField-RapidDecaySeq-basisVec, GaussianField-RapidDecaySeq-instAddCommGroup, GaussianField-RapidDecaySeq-instTopologicalSpace, GaussianField-RapidDecaySeq-weight-nonneg, GaussianField-RapidDecaySeq-rapidDecay-withSeminorms, GaussianField-RapidDecaySeq-rapid-decay, Lattice-toSemilatticeInf}
  \texttt{GaussianField.RapidDecaySeq.hasSum\_basisVec}
\end{theorem}
\begin{theorem}[\texttt{pure\_val}]\label{GaussianField-NuclearTensorProduct-pure-val}
  \lean{GaussianField.NuclearTensorProduct.pure_val}\leanok
  \uses{GaussianField-DyninMityaginSpace, GaussianField-DyninMityaginSpace-coeff, GaussianField-RapidDecaySeq-val, GaussianField-NuclearTensorProduct-pure}
  \texttt{GaussianField.NuclearTensorProduct.pure\_val}
\end{theorem}
\begin{theorem}[\texttt{nat\_pair\_bound}]\label{GaussianField-nat-pair-bound}
  \lean{GaussianField.nat_pair_bound}\leanok
  \texttt{GaussianField.nat\_pair\_bound}
\end{theorem}
\begin{theorem}[\texttt{basisVec\_val\_ne}]\label{GaussianField-RapidDecaySeq-basisVec-val-ne}
  \lean{GaussianField.RapidDecaySeq.basisVec_val_ne}\leanok
  \uses{GaussianField-RapidDecaySeq-val, GaussianField-RapidDecaySeq-basisVec}
  \texttt{GaussianField.RapidDecaySeq.basisVec\_val\_ne}
\end{theorem}
\begin{theorem}[\texttt{add\_val}]\label{GaussianField-NuclearTensorProduct-add-val}
  \lean{GaussianField.NuclearTensorProduct.add_val}\leanok
  \uses{GaussianField-RapidDecaySeq-val, GaussianField-NuclearTensorProduct, GaussianField-NuclearTensorProduct-instAddCommGroup}
  \texttt{GaussianField.NuclearTensorProduct.add\_val}
\end{theorem}
\begin{theorem}[\texttt{rapidDecay\_expansion}]\label{GaussianField-RapidDecaySeq-rapidDecay-expansion}
  \lean{GaussianField.RapidDecaySeq.rapidDecay_expansion}\leanok
  \uses{GaussianField-RapidDecaySeq, GaussianField-RapidDecaySeq-val, GaussianField-RapidDecaySeq-hasSum-basisVec, GaussianField-RapidDecaySeq-instSMul, GaussianField-RapidDecaySeq-instModule, GaussianField-RapidDecaySeq-basisVec, GaussianField-RapidDecaySeq-instAddCommGroup, GaussianField-RapidDecaySeq-instTopologicalSpace}
  \texttt{GaussianField.RapidDecaySeq.rapidDecay\_expansion}
\end{theorem}
\begin{definition}[\texttt{instAdd}]\label{GaussianField-RapidDecaySeq-instAdd}
  \lean{GaussianField.RapidDecaySeq.instAdd}\leanok
  \uses{GaussianField-RapidDecaySeq, GaussianField-RapidDecaySeq-val}
  \texttt{GaussianField.RapidDecaySeq.instAdd}
\end{definition}
\begin{definition}[\texttt{rapidDecay\_dyninMityaginSpace}]\label{GaussianField-RapidDecaySeq-rapidDecay-dyninMityaginSpace}
  \lean{GaussianField.RapidDecaySeq.rapidDecay_dyninMityaginSpace}\leanok
  \uses{GaussianField-DyninMityaginSpace, GaussianField-RapidDecaySeq-rapidDecaySeminorm, GaussianField-RapidDecaySeq-instCompleteSpace, GaussianField-RapidDecaySeq, GaussianField-RapidDecaySeq-instModule, GaussianField-RapidDecaySeq-basisVec, GaussianField-RapidDecaySeq-rapidDecay-t1Space, GaussianField-RapidDecaySeq-rapidDecay-expansion, GaussianField-RapidDecaySeq-instAddCommGroup, GaussianField-RapidDecaySeq-instIsTopologicalAddGroup, GaussianField-RapidDecaySeq-instTopologicalSpace, GaussianField-RapidDecaySeq-instContinuousSMul, GaussianField-RapidDecaySeq-rapidDecay-withSeminorms, GaussianField-RapidDecaySeq-coeffCLM}
  \texttt{GaussianField.RapidDecaySeq.rapidDecay\_dyninMityaginSpace}
\end{definition}
\begin{theorem}[\texttt{rapidDecay\_withSeminorms}]\label{GaussianField-RapidDecaySeq-rapidDecay-withSeminorms}
  \lean{GaussianField.RapidDecaySeq.rapidDecay_withSeminorms}\leanok
  \uses{GaussianField-RapidDecaySeq-rapidDecaySeminorm, GaussianField-RapidDecaySeq, GaussianField-RapidDecaySeq-instModule, GaussianField-RapidDecaySeq-instAddCommGroup, GaussianField-RapidDecaySeq-instTopologicalSpace}
  \texttt{GaussianField.RapidDecaySeq.rapidDecay\_withSeminorms}
\end{theorem}
\begin{definition}[\texttt{fromPairIndex}]\label{GaussianField-NuclearTensorProduct-fromPairIndex}
  \lean{GaussianField.NuclearTensorProduct.fromPairIndex}\leanok
  \texttt{GaussianField.NuclearTensorProduct.fromPairIndex}
\end{definition}
\begin{definition}[\texttt{dyninMityaginSpace}]\label{GaussianField-NuclearTensorProduct-dyninMityaginSpace}
  \lean{GaussianField.NuclearTensorProduct.dyninMityaginSpace}\leanok
  \uses{GaussianField-DyninMityaginSpace, GaussianField-NuclearTensorProduct-instModuleReal, GaussianField-NuclearTensorProduct-instTopologicalSpace, GaussianField-NuclearTensorProduct-instIsTopologicalAddGroup, GaussianField-NuclearTensorProduct, GaussianField-NuclearTensorProduct-instAddCommGroup, GaussianField-NuclearTensorProduct-instContinuousSMulReal, GaussianField-RapidDecaySeq-rapidDecay-dyninMityaginSpace}
  \texttt{GaussianField.NuclearTensorProduct.dyninMityaginSpace}
\end{definition}
\begin{theorem}[\texttt{pure\_continuous}]\label{GaussianField-NuclearTensorProduct-pure-continuous}
  \lean{GaussianField.NuclearTensorProduct.pure_continuous}\leanok
  \uses{GaussianField-DyninMityaginSpace, GaussianField-NuclearTensorProduct-ext, GaussianField-NuclearTensorProduct-instModuleReal, GaussianField-NuclearTensorProduct-instTopologicalSpace, GaussianField-DyninMityaginSpace-coeff, GaussianField-RapidDecaySeq-instZero, GaussianField-RapidDecaySeq-rapidDecaySeminorm, GaussianField-RapidDecaySeq, GaussianField-RapidDecaySeq-val, GaussianField-NuclearTensorProduct-instIsTopologicalAddGroup, GaussianField-NuclearTensorProduct-pureCLM-right, GaussianField-RapidDecaySeq-instSMul, GaussianField-RapidDecaySeq-instModule, GaussianField-DyninMityaginSpace--, GaussianField-NuclearTensorProduct, GaussianField-NuclearTensorProduct-instAddCommGroup, GaussianField-DyninMityaginSpace-h-with, GaussianField-NuclearTensorProduct-pure, GaussianField-RapidDecaySeq-instAddCommGroup, GaussianField-RapidDecaySeq-instTopologicalSpace, GaussianField-DyninMityaginSpace-p, GaussianField-NuclearTensorProduct-pureLin, GaussianField-RapidDecaySeq-rapidDecay-withSeminorms, GaussianField-NuclearTensorProduct-pure-seminorm-bound, GaussianField-NuclearTensorProduct-pure-continuous-left}
  \texttt{GaussianField.NuclearTensorProduct.pure\_continuous}
\end{theorem}
\begin{theorem}[\texttt{basisVec\_eq\_pure}]\label{GaussianField-NuclearTensorProduct-basisVec-eq-pure}
  \lean{GaussianField.NuclearTensorProduct.basisVec_eq_pure}\leanok
  \uses{GaussianField-DyninMityaginSpace, GaussianField-DyninMityaginSpace-coeff, GaussianField-RapidDecaySeq, GaussianField-RapidDecaySeq-val, GaussianField-RapidDecaySeq-basisVec, GaussianField-RapidDecaySeq-ext, GaussianField-DyninMityaginSpace-basis, GaussianField-NuclearTensorProduct-pure}
  \texttt{GaussianField.NuclearTensorProduct.basisVec\_eq\_pure}
\end{theorem}
\begin{theorem}[\texttt{weight\_nonneg}]\label{GaussianField-RapidDecaySeq-weight-nonneg}
  \lean{GaussianField.RapidDecaySeq.weight_nonneg}\leanok
  \uses{GaussianField-RapidDecaySeq-weight-pos}
  \texttt{GaussianField.RapidDecaySeq.weight\_nonneg}
\end{theorem}
\begin{definition}[\texttt{instTopologicalSpace}]\label{GaussianField-RapidDecaySeq-instTopologicalSpace}
  \lean{GaussianField.RapidDecaySeq.instTopologicalSpace}\leanok
  \uses{GaussianField-RapidDecaySeq-rapidDecaySeminorm, GaussianField-RapidDecaySeq, GaussianField-RapidDecaySeq-instModule, GaussianField-RapidDecaySeq-instAddCommGroup}
  \texttt{GaussianField.RapidDecaySeq.instTopologicalSpace}
\end{definition}
\begin{definition}[\texttt{pureCLM\_right}]\label{GaussianField-NuclearTensorProduct-pureCLM-right}
  \lean{GaussianField.NuclearTensorProduct.pureCLM_right}\leanok
  \uses{GaussianField-DyninMityaginSpace, GaussianField-NuclearTensorProduct-instModuleReal, GaussianField-NuclearTensorProduct-instTopologicalSpace, GaussianField-NuclearTensorProduct, GaussianField-NuclearTensorProduct-instAddCommGroup, GaussianField-NuclearTensorProduct-pureLin}
  \texttt{GaussianField.NuclearTensorProduct.pureCLM\_right}
\end{definition}
\begin{definition}[\texttt{instNeg}]\label{GaussianField-RapidDecaySeq-instNeg}
  \lean{GaussianField.RapidDecaySeq.instNeg}\leanok
  \uses{GaussianField-RapidDecaySeq, GaussianField-RapidDecaySeq-val}
  \texttt{GaussianField.RapidDecaySeq.instNeg}
\end{definition}
\begin{theorem}[\texttt{ext}]\label{GaussianField-RapidDecaySeq-ext}
  \lean{GaussianField.RapidDecaySeq.ext}\leanok
  \uses{GaussianField-RapidDecaySeq, GaussianField-RapidDecaySeq-val}
  \texttt{GaussianField.RapidDecaySeq.ext}
\end{theorem}
\begin{definition}[\texttt{basisVec}]\label{GaussianField-RapidDecaySeq-basisVec}
  \lean{GaussianField.RapidDecaySeq.basisVec}\leanok
  \uses{GaussianField-RapidDecaySeq}
  \texttt{GaussianField.RapidDecaySeq.basisVec}
\end{definition}
\begin{theorem}[\texttt{instContinuousSMulReal}]\label{GaussianField-NuclearTensorProduct-instContinuousSMulReal}
  \lean{GaussianField.NuclearTensorProduct.instContinuousSMulReal}\leanok
  \uses{GaussianField-NuclearTensorProduct-instModuleReal, GaussianField-NuclearTensorProduct-instTopologicalSpace, GaussianField-NuclearTensorProduct, GaussianField-NuclearTensorProduct-instAddCommGroup}
  \texttt{GaussianField.NuclearTensorProduct.instContinuousSMulReal}
\end{theorem}
\begin{definition}[\texttt{ctorIdx}]\label{GaussianField-RapidDecaySeq-ctorIdx}
  \lean{GaussianField.RapidDecaySeq.ctorIdx}\leanok
  \uses{GaussianField-RapidDecaySeq}
  \texttt{GaussianField.RapidDecaySeq.ctorIdx}
\end{definition}
\begin{theorem}[\texttt{finset\_sup\_rapidDecaySeminorm\_le}]\label{GaussianField-RapidDecaySeq-finset-sup-rapidDecaySeminorm-le}
  \lean{GaussianField.RapidDecaySeq.finset_sup_rapidDecaySeminorm_le}\leanok
  \uses{GaussianField-RapidDecaySeq-rapidDecaySeminorm-mono, Lattice-toSemilatticeSup, GaussianField-RapidDecaySeq-rapidDecaySeminorm, GaussianField-RapidDecaySeq, GaussianField-RapidDecaySeq-instSMul, GaussianField-RapidDecaySeq-instModule, GaussianField-RapidDecaySeq-instAddCommGroup}
  \texttt{GaussianField.RapidDecaySeq.finset\_sup\_rapidDecaySeminorm\_le}
\end{theorem}
\begin{theorem}[\texttt{neg\_val}]\label{GaussianField-RapidDecaySeq-neg-val}
  \lean{GaussianField.RapidDecaySeq.neg_val}\leanok
  \uses{GaussianField-RapidDecaySeq, GaussianField-RapidDecaySeq-val, GaussianField-RapidDecaySeq-instNeg}
  \texttt{GaussianField.RapidDecaySeq.neg\_val}
\end{theorem}
\begin{theorem}[\texttt{instCompleteSpace}]\label{GaussianField-RapidDecaySeq-instCompleteSpace}
  \lean{GaussianField.RapidDecaySeq.instCompleteSpace}\leanok
  \uses{GaussianField-RapidDecaySeq-instAdd, Lattice-toSemilatticeSup, GaussianField-RapidDecaySeq-instZero, GaussianField-RapidDecaySeq-rapidDecaySeminorm, GaussianField-RapidDecaySeq, GaussianField-RapidDecaySeq-val, GaussianField-RapidDecaySeq-instSMul, GaussianField-RapidDecaySeq-instModule, GaussianField-RapidDecaySeq-instAddCommGroup, GaussianField-RapidDecaySeq-instIsTopologicalAddGroup, GaussianField-RapidDecaySeq-instTopologicalSpace, GaussianField-RapidDecaySeq-weight-nonneg, GaussianField-RapidDecaySeq-rapidDecay-withSeminorms, Lattice-toSemilatticeInf}
  \texttt{GaussianField.RapidDecaySeq.instCompleteSpace}
\end{theorem}
\begin{theorem}[\texttt{ext\_iff}]\label{GaussianField-NuclearTensorProduct-ext-iff}
  \lean{GaussianField.NuclearTensorProduct.ext_iff}\leanok
  \uses{GaussianField-NuclearTensorProduct-ext, GaussianField-RapidDecaySeq-val, GaussianField-NuclearTensorProduct}
  \texttt{GaussianField.NuclearTensorProduct.ext\_iff}
\end{theorem}
\begin{theorem}[\texttt{rapidDecaySeminorm\_basisVec}]\label{GaussianField-RapidDecaySeq-rapidDecaySeminorm-basisVec}
  \lean{GaussianField.RapidDecaySeq.rapidDecaySeminorm_basisVec}\leanok
  \uses{GaussianField-RapidDecaySeq-rapidDecaySeminorm, GaussianField-RapidDecaySeq, GaussianField-RapidDecaySeq-val, GaussianField-RapidDecaySeq-instSMul, GaussianField-RapidDecaySeq-basisVec, GaussianField-RapidDecaySeq-instAddCommGroup}
  \texttt{GaussianField.RapidDecaySeq.rapidDecaySeminorm\_basisVec}
\end{theorem}
\begin{theorem}[\texttt{lift\_pure}]\label{GaussianField-NuclearTensorProduct-lift-pure}
  \lean{GaussianField.NuclearTensorProduct.lift_pure}\leanok
  \uses{GaussianField-DyninMityaginSpace, GaussianField-NuclearTensorProduct-instModuleReal, GaussianField-NuclearTensorProduct-instTopologicalSpace, GaussianField-DyninMityaginSpace-coeff, GaussianField-RapidDecaySeq-val, GaussianField-DyninMityaginSpace--, GaussianField-NuclearTensorProduct, GaussianField-DyninMityaginSpace-hasSum-basis, GaussianField-NuclearTensorProduct-instAddCommGroup, GaussianField-DyninMityaginSpace-h-with, GaussianField-DyninMityaginSpace-basis, GaussianField-NuclearTensorProduct-pure, GaussianField-NuclearTensorProduct-lift, GaussianField-DyninMityaginSpace-p}
  \texttt{GaussianField.NuclearTensorProduct.lift\_pure}
\end{theorem}
\begin{theorem}[\texttt{instContinuousSMul}]\label{GaussianField-RapidDecaySeq-instContinuousSMul}
  \lean{GaussianField.RapidDecaySeq.instContinuousSMul}\leanok
  \uses{GaussianField-RapidDecaySeq-rapidDecaySeminorm, GaussianField-RapidDecaySeq, GaussianField-RapidDecaySeq-instSMul, GaussianField-RapidDecaySeq-instModule, GaussianField-RapidDecaySeq-instAddCommGroup, GaussianField-RapidDecaySeq-instTopologicalSpace, GaussianField-RapidDecaySeq-rapidDecay-withSeminorms}
  \texttt{GaussianField.RapidDecaySeq.instContinuousSMul}
\end{theorem}
\begin{theorem}[\texttt{pure\_continuous\_left}]\label{GaussianField-NuclearTensorProduct-pure-continuous-left}
  \lean{GaussianField.NuclearTensorProduct.pure_continuous_left}\leanok
  \uses{GaussianField-DyninMityaginSpace, GaussianField-NuclearTensorProduct-ext, GaussianField-NuclearTensorProduct-instModuleReal, GaussianField-NuclearTensorProduct-instTopologicalSpace, GaussianField-DyninMityaginSpace-coeff, GaussianField-RapidDecaySeq-rapidDecaySeminorm, GaussianField-RapidDecaySeq, GaussianField-RapidDecaySeq-instSMul, GaussianField-RapidDecaySeq-instModule, GaussianField-DyninMityaginSpace--, GaussianField-NuclearTensorProduct, GaussianField-NuclearTensorProduct-instAddCommGroup, GaussianField-DyninMityaginSpace-h-with, GaussianField-NuclearTensorProduct-pure, GaussianField-RapidDecaySeq-instAddCommGroup, GaussianField-RapidDecaySeq-instTopologicalSpace, GaussianField-DyninMityaginSpace-p, GaussianField-NuclearTensorProduct-pureLin, GaussianField-RapidDecaySeq-rapidDecay-withSeminorms, GaussianField-NuclearTensorProduct-pure-seminorm-bound}
  \texttt{GaussianField.NuclearTensorProduct.pure\_continuous\_left}
\end{theorem}
\begin{theorem}[\texttt{pure\_seminorm\_bound}]\label{GaussianField-NuclearTensorProduct-pure-seminorm-bound}
  \lean{GaussianField.NuclearTensorProduct.pure_seminorm_bound}\leanok
  \uses{GaussianField-DyninMityaginSpace, GaussianField-DyninMityaginSpace-coeff, GaussianField-RapidDecaySeq-rapidDecaySeminorm, GaussianField-RapidDecaySeq, GaussianField-DyninMityaginSpace-coeff-decay, GaussianField-RapidDecaySeq-val, GaussianField-RapidDecaySeq-instSMul, GaussianField-DyninMityaginSpace--, GaussianField-NuclearTensorProduct-pure, GaussianField-RapidDecaySeq-instAddCommGroup, GaussianField-DyninMityaginSpace-p, GaussianField-RapidDecaySeq-rapid-decay, Lattice-toSemilatticeInf}
  \texttt{GaussianField.NuclearTensorProduct.pure\_seminorm\_bound}
\end{theorem}
\begin{theorem}[\texttt{smul\_val}]\label{GaussianField-NuclearTensorProduct-smul-val}
  \lean{GaussianField.NuclearTensorProduct.smul_val}\leanok
  \uses{GaussianField-NuclearTensorProduct-instModuleReal, GaussianField-RapidDecaySeq-val, GaussianField-NuclearTensorProduct, GaussianField-NuclearTensorProduct-instAddCommGroup}
  \texttt{GaussianField.NuclearTensorProduct.smul\_val}
\end{theorem}
\begin{theorem}[\texttt{fromPairIndex\_toPairIndex}]\label{GaussianField-NuclearTensorProduct-fromPairIndex-toPairIndex}
  \lean{GaussianField.NuclearTensorProduct.fromPairIndex_toPairIndex}\leanok
  \uses{GaussianField-NuclearTensorProduct-toPairIndex, GaussianField-NuclearTensorProduct-fromPairIndex}
  \texttt{GaussianField.NuclearTensorProduct.fromPairIndex\_toPairIndex}
\end{theorem}
\begin{theorem}[\texttt{rapidDecay\_hasBiorthogonalBasis}]\label{GaussianField-RapidDecaySeq-rapidDecay-hasBiorthogonalBasis}
  \lean{GaussianField.RapidDecaySeq.rapidDecay_hasBiorthogonalBasis}\leanok
  \uses{GaussianField-DyninMityaginSpace-coeff, GaussianField-RapidDecaySeq, GaussianField-RapidDecaySeq-instModule, GaussianField-DyninMityaginSpace-basis, GaussianField-RapidDecaySeq-instAddCommGroup, GaussianField-RapidDecaySeq-instIsTopologicalAddGroup, GaussianField-RapidDecaySeq-instTopologicalSpace, GaussianField-RapidDecaySeq-instContinuousSMul, GaussianField-DyninMityaginSpace-HasBiorthogonalBasis, GaussianField-RapidDecaySeq-rapidDecay-dyninMityaginSpace}
  \texttt{GaussianField.RapidDecaySeq.rapidDecay\_hasBiorthogonalBasis}
\end{theorem}
\begin{theorem}[\texttt{add\_val}]\label{GaussianField-RapidDecaySeq-add-val}
  \lean{GaussianField.RapidDecaySeq.add_val}\leanok
  \uses{GaussianField-RapidDecaySeq-instAdd, GaussianField-RapidDecaySeq, GaussianField-RapidDecaySeq-val}
  \texttt{GaussianField.RapidDecaySeq.add\_val}
\end{theorem}
\begin{theorem}[\texttt{ext\_iff}]\label{GaussianField-RapidDecaySeq-ext-iff}
  \lean{GaussianField.RapidDecaySeq.ext_iff}\leanok
  \uses{GaussianField-RapidDecaySeq, GaussianField-RapidDecaySeq-val, GaussianField-RapidDecaySeq-ext}
  \texttt{GaussianField.RapidDecaySeq.ext\_iff}
\end{theorem}
\begin{definition}[\texttt{ofRapidDecayEquiv}]\label{GaussianField-DyninMityaginSpace-ofRapidDecayEquiv}
  \lean{GaussianField.DyninMityaginSpace.ofRapidDecayEquiv}\leanok
  \uses{GaussianField-DyninMityaginSpace, GaussianField-RapidDecaySeq, GaussianField-RapidDecaySeq-instModule, GaussianField-RapidDecaySeq-basisVec, GaussianField-RapidDecaySeq-instAddCommGroup, GaussianField-RapidDecaySeq-instTopologicalSpace, GaussianField-RapidDecaySeq-coeffCLM}
  \texttt{GaussianField.DyninMityaginSpace.ofRapidDecayEquiv}
\end{definition}
\begin{definition}[\texttt{toPairIndex}]\label{GaussianField-NuclearTensorProduct-toPairIndex}
  \lean{GaussianField.NuclearTensorProduct.toPairIndex}\leanok
  \texttt{GaussianField.NuclearTensorProduct.toPairIndex}
\end{definition}
\begin{definition}[\texttt{pureLin}]\label{GaussianField-NuclearTensorProduct-pureLin}
  \lean{GaussianField.NuclearTensorProduct.pureLin}\leanok
  \uses{GaussianField-DyninMityaginSpace, GaussianField-NuclearTensorProduct-instModuleReal, GaussianField-NuclearTensorProduct, GaussianField-NuclearTensorProduct-instAddCommGroup, GaussianField-NuclearTensorProduct-pure}
  \texttt{GaussianField.NuclearTensorProduct.pureLin}
\end{definition}
\begin{theorem}[\texttt{rapidDecaySeminorm\_mono}]\label{GaussianField-RapidDecaySeq-rapidDecaySeminorm-mono}
  \lean{GaussianField.RapidDecaySeq.rapidDecaySeminorm_mono}\leanok
  \uses{GaussianField-RapidDecaySeq-rapidDecaySeminorm, GaussianField-RapidDecaySeq, GaussianField-RapidDecaySeq-val, GaussianField-RapidDecaySeq-instSMul, GaussianField-RapidDecaySeq-instAddCommGroup, GaussianField-RapidDecaySeq-rapid-decay, Lattice-toSemilatticeInf}
  \texttt{GaussianField.RapidDecaySeq.rapidDecaySeminorm\_mono}
\end{theorem}
\begin{theorem}[\texttt{instIsTopologicalAddGroup}]\label{GaussianField-NuclearTensorProduct-instIsTopologicalAddGroup}
  \lean{GaussianField.NuclearTensorProduct.instIsTopologicalAddGroup}\leanok
  \uses{GaussianField-NuclearTensorProduct-instTopologicalSpace, GaussianField-NuclearTensorProduct, GaussianField-NuclearTensorProduct-instAddCommGroup}
  \texttt{GaussianField.NuclearTensorProduct.instIsTopologicalAddGroup}
\end{theorem}
\begin{theorem}[\texttt{smul\_val}]\label{GaussianField-RapidDecaySeq-smul-val}
  \lean{GaussianField.RapidDecaySeq.smul_val}\leanok
  \uses{GaussianField-RapidDecaySeq, GaussianField-RapidDecaySeq-val, GaussianField-RapidDecaySeq-instSMul}
  \texttt{GaussianField.RapidDecaySeq.smul\_val}
\end{theorem}
\begin{theorem}[\texttt{basisVec\_val\_self}]\label{GaussianField-RapidDecaySeq-basisVec-val-self}
  \lean{GaussianField.RapidDecaySeq.basisVec_val_self}\leanok
  \uses{GaussianField-RapidDecaySeq-val, GaussianField-RapidDecaySeq-basisVec}
  \texttt{GaussianField.RapidDecaySeq.basisVec\_val\_self}
\end{theorem}
\begin{definition}[\texttt{coeffLM}]\label{GaussianField-RapidDecaySeq-coeffLM}
  \lean{GaussianField.RapidDecaySeq.coeffLM}\leanok
  \uses{GaussianField-RapidDecaySeq, GaussianField-RapidDecaySeq-val, GaussianField-RapidDecaySeq-instModule, GaussianField-RapidDecaySeq-instAddCommGroup}
  \texttt{GaussianField.RapidDecaySeq.coeffLM}
\end{definition}
\begin{theorem}[\texttt{weight\_pos}]\label{GaussianField-RapidDecaySeq-weight-pos}
  \lean{GaussianField.RapidDecaySeq.weight_pos}\leanok
  \texttt{GaussianField.RapidDecaySeq.weight\_pos}
\end{theorem}
\begin{theorem}[\texttt{finset\_sup\_rapidDecaySeminorm\_basisVec\_le}]\label{GaussianField-RapidDecaySeq-finset-sup-rapidDecaySeminorm-basisVec-le}
  \lean{GaussianField.RapidDecaySeq.finset_sup_rapidDecaySeminorm_basisVec_le}\leanok
  \uses{Lattice-toSemilatticeSup, GaussianField-RapidDecaySeq-rapidDecaySeminorm, GaussianField-RapidDecaySeq, GaussianField-RapidDecaySeq-rapidDecaySeminorm-basisVec, GaussianField-RapidDecaySeq-instSMul, GaussianField-RapidDecaySeq-instModule, GaussianField-RapidDecaySeq-basisVec, GaussianField-RapidDecaySeq-instAddCommGroup, GaussianField-RapidDecaySeq-finset-sup-rapidDecaySeminorm-le}
  \texttt{GaussianField.RapidDecaySeq.finset\_sup\_rapidDecaySeminorm\_basisVec\_le}
\end{theorem}
\section{\texttt{GaussianField.Density}}
\begin{theorem}[\texttt{gaussianDensity\_measurable}]\label{GaussianField-gaussianDensity-measurable}
  \lean{GaussianField.gaussianDensity_measurable}\leanok
  \uses{GaussianField-FinLatticeSites, GaussianField-massOperator, GaussianField-gaussianDensity, GaussianField-FinLatticeField}
  \texttt{GaussianField.gaussianDensity\_measurable}
\end{theorem}
\begin{definition}[\texttt{latticeGaussianFieldLaw}]\label{GaussianField-latticeGaussianFieldLaw}
  \lean{GaussianField.latticeGaussianFieldLaw}\leanok
  \uses{GaussianField-evalMap, GaussianField-FinLatticeSites, GaussianField-latticeGaussianMeasure, GaussianField-Configuration, GaussianField-instMeasurableSpaceConfiguration, GaussianField-FinLatticeField}
  \texttt{GaussianField.latticeGaussianFieldLaw}
\end{definition}
\begin{theorem}[\texttt{integral\_massEigenbasis\_cexp}]\label{GaussianField-integral-massEigenbasis-cexp}
  \lean{GaussianField.integral_massEigenbasis_cexp}\leanok
  \uses{GaussianField-FinLatticeSites, GaussianField-integral-cexp-neg-half-sum-mul-sq-add-linear, GaussianField-massEigenbasis-quadratic-sum-reprSymm-ofLp, GaussianField-massEigenbasis-linear-sum-reprSymm-ofLp, GaussianField-massEigenvalues, GaussianField-massEigenvectorBasis, GaussianField-massOperatorMatrix-eigenvalues-pos}
  \texttt{GaussianField.integral\_massEigenbasis\_cexp}
\end{theorem}
\begin{theorem}[\texttt{evalMap\_evalMapInv}]\label{GaussianField-evalMap-evalMapInv}
  \lean{GaussianField.evalMap_evalMapInv}\leanok
  \uses{GaussianField-evalMapInv-apply, GaussianField-finLatticeDelta, GaussianField-evalMap, GaussianField-FinLatticeSites, GaussianField-evalMapInv, GaussianField-Configuration, GaussianField-FinLatticeField}
  \texttt{GaussianField.evalMap\_evalMapInv}
\end{theorem}
\begin{theorem}[\texttt{congr\_simp}]\label{GaussianField-evalMapInv-congr-simp}
  \lean{GaussianField.evalMapInv.congr_simp}\leanok
  \uses{GaussianField-FinLatticeSites, GaussianField-evalMapInv, GaussianField-Configuration, GaussianField-FinLatticeField}
  \texttt{GaussianField.evalMapInv.congr\_simp}
\end{theorem}
\begin{definition}[\texttt{evalMap}]\label{GaussianField-evalMap}
  \lean{GaussianField.evalMap}\leanok
  \uses{GaussianField-finLatticeDelta, GaussianField-FinLatticeSites, GaussianField-Configuration, GaussianField-FinLatticeField}
  \texttt{GaussianField.evalMap}
\end{definition}
\begin{definition}[\texttt{strongDualToField}]\label{GaussianField-strongDualToField}
  \lean{GaussianField.strongDualToField}\leanok
  \uses{GaussianField-finLatticeDelta, GaussianField-FinLatticeSites, GaussianField-FinLatticeField}
  \texttt{GaussianField.strongDualToField}
\end{definition}
\begin{definition}[\texttt{normalizedGaussianDensityMeasure}]\label{GaussianField-normalizedGaussianDensityMeasure}
  \lean{GaussianField.normalizedGaussianDensityMeasure}\leanok
  \uses{GaussianField-FinLatticeSites, GaussianField-gaussianDensityMeasure, GaussianField-gaussianDensityNormConst, GaussianField-FinLatticeField}
  \texttt{GaussianField.normalizedGaussianDensityMeasure}
\end{definition}
\begin{theorem}[\texttt{latticeGaussianFieldLaw\_univ}]\label{GaussianField-latticeGaussianFieldLaw-univ}
  \lean{GaussianField.latticeGaussianFieldLaw_univ}\leanok
  \uses{GaussianField-measurable-evalMap, GaussianField-evalMap, GaussianField-FinLatticeSites, GaussianField-latticeGaussianFieldLaw, GaussianField-latticeGaussianMeasure-isProbability, GaussianField-latticeGaussianMeasure, GaussianField-Configuration, GaussianField-instMeasurableSpaceConfiguration, GaussianField-FinLatticeField}
  \texttt{GaussianField.latticeGaussianFieldLaw\_univ}
\end{theorem}
\begin{theorem}[\texttt{congr\_simp}]\label{GaussianField-gaussianDensityNormConst-congr-simp}
  \lean{GaussianField.gaussianDensityNormConst.congr_simp}\leanok
  \uses{GaussianField-gaussianDensityNormConst}
  \texttt{GaussianField.gaussianDensityNormConst.congr\_simp}
\end{theorem}
\begin{theorem}[\texttt{integrable\_mul\_gaussianDensity\_of\_fieldLaw}]\label{GaussianField-integrable-mul-gaussianDensity-of-fieldLaw}
  \lean{GaussianField.integrable_mul_gaussianDensity_of_fieldLaw}\leanok
  \uses{GaussianField-FinLatticeSites, GaussianField-gaussianDensityMeasure, GaussianField-gaussianDensityNormConst, GaussianField-latticeGaussianFieldLaw, GaussianField-latticeGaussianFieldLaw-eq-normalizedGaussianDensityMeasure, GaussianField-gaussianDensityNormConst-ne-zero, GaussianField-gaussianDensityWeight-measurable, GaussianField-gaussianDensityWeight, GaussianField-gaussianDensityWeight-toReal, GaussianField-gaussianDensityNormConst-ne-top, GaussianField-normalizedGaussianDensityMeasure, GaussianField-gaussianDensity, GaussianField-FinLatticeField}
  \texttt{GaussianField.integrable\_mul\_gaussianDensity\_of\_fieldLaw}
\end{theorem}
\begin{theorem}[\texttt{latticeGaussianFieldLaw\_fourier}]\label{GaussianField-latticeGaussianFieldLaw-fourier}
  \lean{GaussianField.latticeGaussianFieldLaw_fourier}\leanok
  \uses{GaussianField-measurable-evalMap, GaussianField-charFun, GaussianField-finLatticeField-dyninMityaginSpace, GaussianField-measure, GaussianField-ell2-separableSpace, GaussianField-config-apply-eq-sum-evalMap, GaussianField-evalMap, GaussianField-FinLatticeSites, GaussianField-latticeCovarianceGJ, GaussianField-latticeGaussianFieldLaw, GaussianField-latticeGaussianMeasure, GaussianField-Configuration, GaussianField-instMeasurableSpaceConfiguration, GaussianField-ell2-, GaussianField-FinLatticeField}
  \texttt{GaussianField.latticeGaussianFieldLaw\_fourier}
\end{theorem}
\begin{theorem}[\texttt{gaussianDensity\_eq\_exp\_massEigenbasis}]\label{GaussianField-gaussianDensity-eq-exp-massEigenbasis}
  \lean{GaussianField.gaussianDensity_eq_exp_massEigenbasis}\leanok
  \uses{GaussianField-gaussianDensity-eq-exp-spectral, GaussianField-FinLatticeSites, GaussianField-massEigenvalues, GaussianField-massEigenvectorBasis, GaussianField-gaussianDensity, GaussianField-FinLatticeField}
  \texttt{GaussianField.gaussianDensity\_eq\_exp\_massEigenbasis}
\end{theorem}
\begin{theorem}[\texttt{integrable\_gaussianDensity}]\label{GaussianField-integrable-gaussianDensity}
  \lean{GaussianField.integrable_gaussianDensity}\leanok
  \uses{GaussianField-measurable-evalMap, GaussianField-evalMap, GaussianField-FinLatticeSites, GaussianField-gaussianDensityMeasure, GaussianField-gaussianDensityNormConst, GaussianField-latticeGaussianFieldLaw, GaussianField-latticeGaussianMeasure-isProbability, GaussianField-latticeGaussianFieldLaw-eq-normalizedGaussianDensityMeasure, GaussianField-latticeGaussianMeasure, GaussianField-gaussianDensityNormConst-ne-zero, GaussianField-gaussianDensityWeight-measurable, GaussianField-Configuration, GaussianField-gaussianDensityWeight, GaussianField-gaussianDensityWeight-toReal, GaussianField-gaussianDensityNormConst-ne-top, GaussianField-instMeasurableSpaceConfiguration, GaussianField-normalizedGaussianDensityMeasure, GaussianField-gaussianDensity, GaussianField-FinLatticeField}
  \texttt{GaussianField.integrable\_gaussianDensity}
\end{theorem}
\begin{theorem}[\texttt{integral\_cexp\_neg\_half\_sum\_mul\_sq\_add\_linear}]\label{GaussianField-integral-cexp-neg-half-sum-mul-sq-add-linear}
  \lean{GaussianField.integral_cexp_neg_half_sum_mul_sq_add_linear}\leanok
  \texttt{GaussianField.integral\_cexp\_neg\_half\_sum\_mul\_sq\_add\_linear}
\end{theorem}
\begin{theorem}[\texttt{strongDual\_apply\_eq\_site\_sum}]\label{GaussianField-strongDual-apply-eq-site-sum}
  \lean{GaussianField.strongDual_apply_eq_site_sum}\leanok
  \uses{GaussianField-finLatticeDelta, GaussianField-FinLatticeSites, GaussianField-strongDualToField, GaussianField-field-basis-decomposition-density, GaussianField-FinLatticeField}
  \texttt{GaussianField.strongDual\_apply\_eq\_site\_sum}
\end{theorem}
\begin{theorem}[\texttt{congr\_simp}]\label{GaussianField-latticeGaussianFieldLaw-congr-simp}
  \lean{GaussianField.latticeGaussianFieldLaw.congr_simp}\leanok
  \uses{GaussianField-FinLatticeSites, GaussianField-latticeGaussianFieldLaw, GaussianField-FinLatticeField}
  \texttt{GaussianField.latticeGaussianFieldLaw.congr\_simp}
\end{theorem}
\begin{theorem}[\texttt{integrable\_normalizedGaussianDensityMeasure\_iff}]\label{GaussianField-integrable-normalizedGaussianDensityMeasure-iff}
  \lean{GaussianField.integrable_normalizedGaussianDensityMeasure_iff}\leanok
  \uses{GaussianField-FinLatticeSites, GaussianField-latticeGaussianFieldLaw, GaussianField-latticeGaussianFieldLaw-eq-normalizedGaussianDensityMeasure, GaussianField-normalizedGaussianDensityMeasure, GaussianField-FinLatticeField}
  \texttt{GaussianField.integrable\_normalizedGaussianDensityMeasure\_iff}
\end{theorem}
\begin{theorem}[\texttt{latticeGaussianMeasure\_density\_integral'}]\label{GaussianField-latticeGaussianMeasure-density-integral-}
  \lean{GaussianField.latticeGaussianMeasure_density_integral'}\leanok
  \uses{GaussianField-evalMap, GaussianField-FinLatticeSites, GaussianField-evalMapMeasurableEquiv, GaussianField-latticeGaussianFieldLaw, GaussianField-latticeGaussianFieldLaw-density-integral, GaussianField-latticeGaussianMeasure, GaussianField-Configuration, GaussianField-instMeasurableSpaceConfiguration, GaussianField-gaussianDensity, GaussianField-FinLatticeField}
  \texttt{GaussianField.latticeGaussianMeasure\_density\_integral'}
\end{theorem}
\begin{theorem}[\texttt{gaussianDensity\_integral\_pos}]\label{GaussianField-gaussianDensity-integral-pos}
  \lean{GaussianField.gaussianDensity_integral_pos}\leanok
  \uses{GaussianField-integrable-gaussianDensity, GaussianField-FinLatticeSites, GaussianField-massOperator, GaussianField-gaussianDensity, GaussianField-FinLatticeField}
  \texttt{GaussianField.gaussianDensity\_integral\_pos}
\end{theorem}
\begin{theorem}[\texttt{latticeGaussianFieldLaw\_pairing\_is\_gaussian}]\label{GaussianField-latticeGaussianFieldLaw-pairing-is-gaussian}
  \lean{GaussianField.latticeGaussianFieldLaw_pairing_is_gaussian}\leanok
  \uses{GaussianField-measurable-evalMap, GaussianField-finLatticeField-dyninMityaginSpace, GaussianField-measure, GaussianField-ell2-separableSpace, GaussianField-config-apply-eq-sum-evalMap, GaussianField-evalMap, GaussianField-FinLatticeSites, GaussianField-latticeCovarianceGJ, GaussianField-latticeGaussianFieldLaw, GaussianField-measurable-sitePairing, GaussianField-pairing-is-gaussian, GaussianField-latticeGaussianMeasure, GaussianField-Configuration, GaussianField-instMeasurableSpaceConfiguration, GaussianField-ell2-, GaussianField-FinLatticeField}
  \texttt{GaussianField.latticeGaussianFieldLaw\_pairing\_is\_gaussian}
\end{theorem}
\begin{theorem}[\texttt{evalMapInv\_evalMap}]\label{GaussianField-evalMapInv-evalMap}
  \lean{GaussianField.evalMapInv_evalMap}\leanok
  \uses{GaussianField-config-apply-eq-sum-evalMap, GaussianField-evalMap, GaussianField-FinLatticeSites, GaussianField-evalMapInv, GaussianField-Configuration, GaussianField-FinLatticeField}
  \texttt{GaussianField.evalMapInv\_evalMap}
\end{theorem}
\begin{theorem}[\texttt{latticeGaussianFieldLaw\_density\_integral}]\label{GaussianField-latticeGaussianFieldLaw-density-integral}
  \lean{GaussianField.latticeGaussianFieldLaw_density_integral}\leanok
  \uses{GaussianField-gaussianDensity-integral-pos, GaussianField-FinLatticeSites, GaussianField-gaussianDensityNormConst-eq-ofReal-integral, GaussianField-gaussianDensityMeasure, GaussianField-gaussianDensityNormConst, GaussianField-latticeGaussianFieldLaw, GaussianField-latticeGaussianFieldLaw-eq-normalizedGaussianDensityMeasure, GaussianField-gaussianDensityWeight-measurable, GaussianField-gaussianDensity-nonneg, GaussianField-gaussianDensityWeight, GaussianField-gaussianDensityWeight-toReal, GaussianField-normalizedGaussianDensityMeasure, GaussianField-gaussianDensity, GaussianField-FinLatticeField}
  \texttt{GaussianField.latticeGaussianFieldLaw\_density\_integral}
\end{theorem}
\begin{theorem}[\texttt{evalMapInv\_measurable}]\label{GaussianField-evalMapInv-measurable}
  \lean{GaussianField.evalMapInv_measurable}\leanok
  \uses{GaussianField-evalMapInv-apply, GaussianField-FinLatticeSites, GaussianField-evalMapInv, GaussianField-configuration-measurable-of-eval-measurable, GaussianField-Configuration, GaussianField-instMeasurableSpaceConfiguration, GaussianField-FinLatticeField}
  \texttt{GaussianField.evalMapInv\_measurable}
\end{theorem}
\begin{definition}[\texttt{evalMapInv}]\label{GaussianField-evalMapInv}
  \lean{GaussianField.evalMapInv}\leanok
  \uses{GaussianField-FinLatticeSites, GaussianField-Configuration, GaussianField-FinLatticeField}
  \texttt{GaussianField.evalMapInv}
\end{definition}
\begin{theorem}[\texttt{integrable\_mul\_gaussianDensity\_of\_comp\_eval}]\label{GaussianField-integrable-mul-gaussianDensity-of-comp-eval}
  \lean{GaussianField.integrable_mul_gaussianDensity_of_comp_eval}\leanok
  \uses{GaussianField-measurable-evalMap, GaussianField-integrable-mul-gaussianDensity-of-fieldLaw, GaussianField-evalMap, GaussianField-FinLatticeSites, GaussianField-latticeGaussianFieldLaw, GaussianField-latticeGaussianMeasure, GaussianField-Configuration, GaussianField-instMeasurableSpaceConfiguration, GaussianField-gaussianDensity, GaussianField-FinLatticeField}
  \texttt{GaussianField.integrable\_mul\_gaussianDensity\_of\_comp\_eval}
\end{theorem}
\begin{theorem}[\texttt{field\_basis\_decomposition\_density}]\label{GaussianField-field-basis-decomposition-density}
  \lean{GaussianField.field_basis_decomposition_density}\leanok
  \uses{GaussianField-finLatticeDelta, GaussianField-FinLatticeSites, GaussianField-FinLatticeField}
  \texttt{GaussianField.field\_basis\_decomposition\_density}
\end{theorem}
\begin{theorem}[\texttt{congr\_simp}]\label{GaussianField-gaussianDensityWeight-congr-simp}
  \lean{GaussianField.gaussianDensityWeight.congr_simp}\leanok
  \uses{GaussianField-gaussianDensityWeight, GaussianField-FinLatticeField}
  \texttt{GaussianField.gaussianDensityWeight.congr\_simp}
\end{theorem}
\begin{theorem}[\texttt{sitePairing\_eq\_massEigenbasis\_sum}]\label{GaussianField-sitePairing-eq-massEigenbasis-sum}
  \lean{GaussianField.sitePairing_eq_massEigenbasis_sum}\leanok
  \uses{GaussianField-FinLatticeSites, GaussianField-massEigenbasis-sum-mul-sum-eq-site-inner, GaussianField-massEigenvectorBasis, GaussianField-FinLatticeField}
  \texttt{GaussianField.sitePairing\_eq\_massEigenbasis\_sum}
\end{theorem}
\begin{theorem}[\texttt{gaussianDensityWeight\_measurable}]\label{GaussianField-gaussianDensityWeight-measurable}
  \lean{GaussianField.gaussianDensityWeight_measurable}\leanok
  \uses{GaussianField-FinLatticeSites, GaussianField-gaussianDensity-measurable, GaussianField-gaussianDensityWeight, GaussianField-gaussianDensity, GaussianField-FinLatticeField}
  \texttt{GaussianField.gaussianDensityWeight\_measurable}
\end{theorem}
\begin{theorem}[\texttt{integral\_latticeGaussianFieldLaw\_eq\_integral\_normalizedGaussianDensityMeasure}]\label{GaussianField-integral-latticeGaussianFieldLaw-eq-integral-normalizedGaussianDensityMeasure}
  \lean{GaussianField.integral_latticeGaussianFieldLaw_eq_integral_normalizedGaussianDensityMeasure}\leanok
  \uses{GaussianField-FinLatticeSites, GaussianField-latticeGaussianFieldLaw, GaussianField-latticeGaussianFieldLaw-eq-normalizedGaussianDensityMeasure, GaussianField-normalizedGaussianDensityMeasure, GaussianField-FinLatticeField}
  \texttt{GaussianField.integral\_latticeGaussianFieldLaw\_eq\_integral\_normalizedGaussianDensityMeasure}
\end{theorem}
\begin{theorem}[\texttt{normalizedGaussianDensityMeasure\_charFunDual\_eq\_latticeGaussianFieldLaw}]\label{GaussianField-normalizedGaussianDensityMeasure-charFunDual-eq-latticeGaussianFieldLaw}
  \lean{GaussianField.normalizedGaussianDensityMeasure_charFunDual_eq_latticeGaussianFieldLaw}\leanok
  \uses{GaussianField-latticeGaussianFieldLaw-fourier, GaussianField-FinLatticeSites, GaussianField-latticeCovarianceGJ, GaussianField-latticeGaussianFieldLaw, GaussianField-strongDualToField, GaussianField-normalizedGaussianDensityMeasure, GaussianField-charFunDual-eq-site-integral, GaussianField-normalizedGaussianDensityMeasure-linearFourier, GaussianField-ell2-, GaussianField-FinLatticeField}
  \texttt{GaussianField.normalizedGaussianDensityMeasure\_charFunDual\_eq\_latticeGaussianFieldLaw}
\end{theorem}
\begin{definition}[\texttt{quadraticGaussianMeasure}]\label{GaussianField-quadraticGaussianMeasure}
  \lean{GaussianField.quadraticGaussianMeasure}\leanok
  \uses{GaussianField-quadraticGaussianDensity}
  \texttt{GaussianField.quadraticGaussianMeasure}
\end{definition}
\begin{theorem}[\texttt{measurable\_evalMap}]\label{GaussianField-measurable-evalMap}
  \lean{GaussianField.measurable_evalMap}\leanok
  \uses{GaussianField-finLatticeDelta, GaussianField-evalMap, GaussianField-FinLatticeSites, GaussianField-configuration-eval-measurable, GaussianField-Configuration, GaussianField-instMeasurableSpaceConfiguration, GaussianField-FinLatticeField}
  \texttt{GaussianField.measurable\_evalMap}
\end{theorem}
\begin{definition}[\texttt{quadraticGaussianDensity}]\label{GaussianField-quadraticGaussianDensity}
  \lean{GaussianField.quadraticGaussianDensity}\leanok
  \texttt{GaussianField.quadraticGaussianDensity}
\end{definition}
\begin{theorem}[\texttt{normalizedGaussianDensityMeasure\_linearFourier}]\label{GaussianField-normalizedGaussianDensityMeasure-linearFourier}
  \lean{GaussianField.normalizedGaussianDensityMeasure_linearFourier}\leanok
  \uses{GaussianField-covariance, GaussianField-FinLatticeSites, GaussianField-latticeCovarianceGJ, GaussianField-lattice-covariance-GJ-eq-spectral, GaussianField-gaussianDensityMeasure, GaussianField-gaussianDensityNormConst, GaussianField-gaussianDensity-measurable, GaussianField-integral-massEigenbasis-cexp-GJ, GaussianField-massEigenvalues, GaussianField-gaussianDensity-nonneg, GaussianField-sitePairing-eq-massEigenbasis-sum, GaussianField-gaussianDensityWeight, GaussianField-massEigenvectorBasis, GaussianField-gaussianDensity-eq-exp-massEigenbasis, GaussianField-normalizedGaussianDensityMeasure, GaussianField-gaussianDensity, GaussianField-massOperatorMatrix-eigenvalues-pos, GaussianField-ell2-, GaussianField-FinLatticeField}
  \texttt{GaussianField.normalizedGaussianDensityMeasure\_linearFourier}
\end{theorem}
\begin{definition}[\texttt{evalMapMeasurableEquiv}]\label{GaussianField-evalMapMeasurableEquiv}
  \lean{GaussianField.evalMapMeasurableEquiv}\leanok
  \uses{GaussianField-measurable-evalMap, GaussianField-evalMap, GaussianField-FinLatticeSites, GaussianField-evalMap-evalMapInv, GaussianField-evalMapInv, GaussianField-Configuration, GaussianField-evalMapInv-evalMap, GaussianField-evalMapInv-measurable, GaussianField-instMeasurableSpaceConfiguration, GaussianField-FinLatticeField}
  \texttt{GaussianField.evalMapMeasurableEquiv}
\end{definition}
\begin{theorem}[\texttt{config\_apply\_eq\_sum\_delta}]\label{GaussianField-config-apply-eq-sum-delta}
  \lean{GaussianField.config_apply_eq_sum_delta}\leanok
  \uses{GaussianField-finLatticeDelta, GaussianField-FinLatticeSites, GaussianField-Configuration, GaussianField-field-basis-decomposition-density, GaussianField-FinLatticeField}
  \texttt{GaussianField.config\_apply\_eq\_sum\_delta}
\end{theorem}
\begin{theorem}[\texttt{latticeGaussianFieldLaw\_eq\_normalizedGaussianDensityMeasure}]\label{GaussianField-latticeGaussianFieldLaw-eq-normalizedGaussianDensityMeasure}
  \lean{GaussianField.latticeGaussianFieldLaw_eq_normalizedGaussianDensityMeasure}\leanok
  \uses{GaussianField-measurable-evalMap, GaussianField-evalMap, GaussianField-FinLatticeSites, GaussianField-latticeGaussianFieldLaw, GaussianField-latticeGaussianMeasure-isProbability, GaussianField-latticeGaussianMeasure, GaussianField-Configuration, GaussianField-normalizedGaussianDensityMeasure-charFunDual-eq-latticeGaussianFieldLaw, GaussianField-instMeasurableSpaceConfiguration, GaussianField-normalizedGaussianDensityMeasure, GaussianField-FinLatticeField, GaussianField-normalizedGaussianDensityMeasure-isFinite}
  \texttt{GaussianField.latticeGaussianFieldLaw\_eq\_normalizedGaussianDensityMeasure}
\end{theorem}
\begin{definition}[\texttt{gaussianDensityMeasure}]\label{GaussianField-gaussianDensityMeasure}
  \lean{GaussianField.gaussianDensityMeasure}\leanok
  \uses{GaussianField-FinLatticeSites, GaussianField-gaussianDensityWeight, GaussianField-FinLatticeField}
  \texttt{GaussianField.gaussianDensityMeasure}
\end{definition}
\begin{theorem}[\texttt{evalMapInv\_apply}]\label{GaussianField-evalMapInv-apply}
  \lean{GaussianField.evalMapInv_apply}\leanok
  \uses{GaussianField-FinLatticeSites, GaussianField-evalMapInv, GaussianField-Configuration, GaussianField-FinLatticeField}
  \texttt{GaussianField.evalMapInv\_apply}
\end{theorem}
\begin{theorem}[\texttt{gaussianDensityWeight\_toReal}]\label{GaussianField-gaussianDensityWeight-toReal}
  \lean{GaussianField.gaussianDensityWeight_toReal}\leanok
  \uses{GaussianField-gaussianDensity-nonneg, GaussianField-gaussianDensityWeight, GaussianField-gaussianDensity, GaussianField-FinLatticeField}
  \texttt{GaussianField.gaussianDensityWeight\_toReal}
\end{theorem}
\begin{definition}[\texttt{normalizedQuadraticGaussianMeasure}]\label{GaussianField-normalizedQuadraticGaussianMeasure}
  \lean{GaussianField.normalizedQuadraticGaussianMeasure}\leanok
  \uses{GaussianField-quadraticGaussianMeasure}
  \texttt{GaussianField.normalizedQuadraticGaussianMeasure}
\end{definition}
\begin{theorem}[\texttt{congr\_simp}]\label{GaussianField-gaussianDensity-congr-simp}
  \lean{GaussianField.gaussianDensity.congr_simp}\leanok
  \uses{GaussianField-gaussianDensity, GaussianField-FinLatticeField}
  \texttt{GaussianField.gaussianDensity.congr\_simp}
\end{theorem}
\begin{theorem}[\texttt{charFunDual\_eq\_site\_integral}]\label{GaussianField-charFunDual-eq-site-integral}
  \lean{GaussianField.charFunDual_eq_site_integral}\leanok
  \uses{GaussianField-FinLatticeSites, GaussianField-strongDual-apply-eq-site-sum, GaussianField-strongDualToField, GaussianField-FinLatticeField}
  \texttt{GaussianField.charFunDual\_eq\_site\_integral}
\end{theorem}
\begin{theorem}[\texttt{gaussianDensityNormConst\_ne\_top}]\label{GaussianField-gaussianDensityNormConst-ne-top}
  \lean{GaussianField.gaussianDensityNormConst_ne_top}\leanok
  \uses{GaussianField-latticeGaussianFieldLaw-univ, GaussianField-FinLatticeSites, GaussianField-gaussianDensityMeasure, GaussianField-gaussianDensityNormConst, GaussianField-latticeGaussianFieldLaw, GaussianField-latticeGaussianFieldLaw-eq-normalizedGaussianDensityMeasure, GaussianField-normalizedGaussianDensityMeasure, GaussianField-FinLatticeField}
  \texttt{GaussianField.gaussianDensityNormConst\_ne\_top}
\end{theorem}
\begin{theorem}[\texttt{latticeGaussianMeasure\_density\_integral}]\label{GaussianField-latticeGaussianMeasure-density-integral}
  \lean{GaussianField.latticeGaussianMeasure_density_integral}\leanok
  \uses{GaussianField-finLatticeDelta, GaussianField-FinLatticeSites, GaussianField-latticeGaussianMeasure, GaussianField-Configuration, GaussianField-instMeasurableSpaceConfiguration, GaussianField-gaussianDensity, GaussianField-latticeGaussianMeasure-density-integral-of-fieldLaw, GaussianField-FinLatticeField}
  \texttt{GaussianField.latticeGaussianMeasure\_density\_integral}
\end{theorem}
\begin{theorem}[\texttt{measurable\_sitePairing}]\label{GaussianField-measurable-sitePairing}
  \lean{GaussianField.measurable_sitePairing}\leanok
  \uses{GaussianField-FinLatticeSites, GaussianField-FinLatticeField}
  \texttt{GaussianField.measurable\_sitePairing}
\end{theorem}
\begin{theorem}[\texttt{congr\_simp}]\label{GaussianField-gaussianDensityMeasure-congr-simp}
  \lean{GaussianField.gaussianDensityMeasure.congr_simp}\leanok
  \uses{GaussianField-FinLatticeSites, GaussianField-gaussianDensityMeasure, GaussianField-FinLatticeField}
  \texttt{GaussianField.gaussianDensityMeasure.congr\_simp}
\end{theorem}
\begin{theorem}[\texttt{congr\_simp}]\label{GaussianField-Configuration-congr-simp}
  \lean{GaussianField.Configuration.congr_simp}\leanok
  \uses{GaussianField-Configuration}
  \texttt{GaussianField.Configuration.congr\_simp}
\end{theorem}
\begin{theorem}[\texttt{config\_apply\_eq\_sum\_evalMap}]\label{GaussianField-config-apply-eq-sum-evalMap}
  \lean{GaussianField.config_apply_eq_sum_evalMap}\leanok
  \uses{GaussianField-finLatticeDelta, GaussianField-evalMap, GaussianField-FinLatticeSites, GaussianField-Configuration, GaussianField-config-apply-eq-sum-delta, GaussianField-FinLatticeField}
  \texttt{GaussianField.config\_apply\_eq\_sum\_evalMap}
\end{theorem}
\begin{theorem}[\texttt{integrable\_mul\_gaussianDensity}]\label{GaussianField-integrable-mul-gaussianDensity}
  \lean{GaussianField.integrable_mul_gaussianDensity}\leanok
  \uses{GaussianField-measurable-evalMap, GaussianField-integrable-mul-gaussianDensity-of-fieldLaw, GaussianField-finLatticeDelta, GaussianField-evalMap, GaussianField-FinLatticeSites, GaussianField-latticeGaussianFieldLaw, GaussianField-latticeGaussianMeasure, GaussianField-Configuration, GaussianField-instMeasurableSpaceConfiguration, GaussianField-gaussianDensity, GaussianField-FinLatticeField}
  \texttt{GaussianField.integrable\_mul\_gaussianDensity}
\end{theorem}
\begin{theorem}[\texttt{normalizedGaussianDensityMeasure\_isFinite}]\label{GaussianField-normalizedGaussianDensityMeasure-isFinite}
  \lean{GaussianField.normalizedGaussianDensityMeasure_isFinite}\leanok
  \uses{GaussianField-FinLatticeSites, GaussianField-gaussianDensityMeasure, GaussianField-normalizedGaussianDensityMeasure, GaussianField-FinLatticeField}
  \texttt{GaussianField.normalizedGaussianDensityMeasure\_isFinite}
\end{theorem}
\begin{theorem}[\texttt{gaussianDensityNormConst\_eq\_ofReal\_integral}]\label{GaussianField-gaussianDensityNormConst-eq-ofReal-integral}
  \lean{GaussianField.gaussianDensityNormConst_eq_ofReal_integral}\leanok
  \uses{GaussianField-gaussianDensityNormConst-eq-lintegral, GaussianField-integrable-gaussianDensity, GaussianField-FinLatticeSites, GaussianField-gaussianDensityNormConst, GaussianField-gaussianDensity-nonneg, GaussianField-gaussianDensityWeight, GaussianField-gaussianDensity, GaussianField-FinLatticeField}
  \texttt{GaussianField.gaussianDensityNormConst\_eq\_ofReal\_integral}
\end{theorem}
\begin{theorem}[\texttt{congr\_simp}]\label{GaussianField-latticeCovarianceGJ-congr-simp}
  \lean{GaussianField.latticeCovarianceGJ.congr_simp}\leanok
  \uses{GaussianField-FinLatticeSites, GaussianField-latticeCovarianceGJ, GaussianField-ell2-, GaussianField-FinLatticeField}
  \texttt{GaussianField.latticeCovarianceGJ.congr\_simp}
\end{theorem}
\begin{definition}[\texttt{gaussianDensityNormConst}]\label{GaussianField-gaussianDensityNormConst}
  \lean{GaussianField.gaussianDensityNormConst}\leanok
  \uses{GaussianField-FinLatticeSites, GaussianField-gaussianDensityMeasure, GaussianField-FinLatticeField}
  \texttt{GaussianField.gaussianDensityNormConst}
\end{definition}
\begin{definition}[\texttt{gaussianDensityWeight}]\label{GaussianField-gaussianDensityWeight}
  \lean{GaussianField.gaussianDensityWeight}\leanok
  \uses{GaussianField-gaussianDensity, GaussianField-FinLatticeField}
  \texttt{GaussianField.gaussianDensityWeight}
\end{definition}
\begin{theorem}[\texttt{gaussianDensityNormConst\_eq\_lintegral}]\label{GaussianField-gaussianDensityNormConst-eq-lintegral}
  \lean{GaussianField.gaussianDensityNormConst_eq_lintegral}\leanok
  \uses{GaussianField-FinLatticeSites, GaussianField-gaussianDensityNormConst, GaussianField-gaussianDensityWeight, GaussianField-FinLatticeField}
  \texttt{GaussianField.gaussianDensityNormConst\_eq\_lintegral}
\end{theorem}
\begin{theorem}[\texttt{normalizedGaussianDensityMeasure\_eq\_normalizedQuadraticGaussianMeasure}]\label{GaussianField-normalizedGaussianDensityMeasure-eq-normalizedQuadraticGaussianMeasure}
  \lean{GaussianField.normalizedGaussianDensityMeasure_eq_normalizedQuadraticGaussianMeasure}\leanok
  \uses{GaussianField-quadraticGaussianDensity, GaussianField-normalizedQuadraticGaussianMeasure, GaussianField-quadraticGaussianMeasure, GaussianField-FinLatticeSites, GaussianField-gaussianDensityMeasure, GaussianField-massOperator, GaussianField-gaussianDensityNormConst, GaussianField-gaussianDensityWeight, GaussianField-normalizedGaussianDensityMeasure, GaussianField-gaussianDensity, GaussianField-FinLatticeField}
  \texttt{GaussianField.normalizedGaussianDensityMeasure\_eq\_normalizedQuadraticGaussianMeasure}
\end{theorem}
\begin{theorem}[\texttt{congr\_simp}]\label{GaussianField-normalizedGaussianDensityMeasure-congr-simp}
  \lean{GaussianField.normalizedGaussianDensityMeasure.congr_simp}\leanok
  \uses{GaussianField-FinLatticeSites, GaussianField-normalizedGaussianDensityMeasure, GaussianField-FinLatticeField}
  \texttt{GaussianField.normalizedGaussianDensityMeasure.congr\_simp}
\end{theorem}
\begin{theorem}[\texttt{integral\_massEigenbasis\_cexp\_GJ}]\label{GaussianField-integral-massEigenbasis-cexp-GJ}
  \lean{GaussianField.integral_massEigenbasis_cexp_GJ}\leanok
  \uses{GaussianField-FinLatticeSites, GaussianField-integral-cexp-neg-half-sum-mul-sq-add-linear, GaussianField-massEigenbasis-quadratic-sum-reprSymm-ofLp, GaussianField-massEigenbasis-linear-sum-reprSymm-ofLp, GaussianField-massEigenvalues, GaussianField-massEigenvectorBasis, GaussianField-massOperatorMatrix-eigenvalues-pos}
  \texttt{GaussianField.integral\_massEigenbasis\_cexp\_GJ}
\end{theorem}
\begin{theorem}[\texttt{gaussianDensityNormConst\_ne\_zero}]\label{GaussianField-gaussianDensityNormConst-ne-zero}
  \lean{GaussianField.gaussianDensityNormConst_ne_zero}\leanok
  \uses{GaussianField-latticeGaussianFieldLaw-univ, GaussianField-FinLatticeSites, GaussianField-gaussianDensityMeasure, GaussianField-gaussianDensityNormConst, GaussianField-latticeGaussianFieldLaw, GaussianField-latticeGaussianFieldLaw-eq-normalizedGaussianDensityMeasure, GaussianField-normalizedGaussianDensityMeasure, GaussianField-FinLatticeField}
  \texttt{GaussianField.gaussianDensityNormConst\_ne\_zero}
\end{theorem}
\begin{theorem}[\texttt{latticeGaussianMeasure\_density\_integral\_of\_fieldLaw}]\label{GaussianField-latticeGaussianMeasure-density-integral-of-fieldLaw}
  \lean{GaussianField.latticeGaussianMeasure_density_integral_of_fieldLaw}\leanok
  \uses{GaussianField-measurable-evalMap, GaussianField-evalMap, GaussianField-FinLatticeSites, GaussianField-latticeGaussianFieldLaw, GaussianField-latticeGaussianFieldLaw-density-integral, GaussianField-latticeGaussianMeasure, GaussianField-Configuration, GaussianField-instMeasurableSpaceConfiguration, GaussianField-gaussianDensity, GaussianField-FinLatticeField}
  \texttt{GaussianField.latticeGaussianMeasure\_density\_integral\_of\_fieldLaw}
\end{theorem}
\section{\texttt{SmoothCircle.Basic}}
\begin{theorem}[\texttt{coe\_mk}]\label{GaussianField-SmoothMap-Circle-coe-mk}
  \lean{GaussianField.SmoothMap_Circle.coe_mk}\leanok
  \uses{GaussianField-SmoothMap-Circle-instFunLike, GaussianField-SmoothMap-Circle}
  \texttt{GaussianField.SmoothMap\_Circle.coe\_mk}
\end{theorem}
\begin{theorem}[\texttt{sobolevSeminorm\_nonneg}]\label{GaussianField-SmoothMap-Circle-sobolevSeminorm-nonneg}
  \lean{GaussianField.SmoothMap_Circle.sobolevSeminorm_nonneg}\leanok
  \uses{GaussianField-SmoothMap-Circle-instFunLike, GaussianField-SmoothMap-Circle-instAddCommGroup, GaussianField-SmoothMap-Circle-instSMul, GaussianField-SmoothMap-Circle, GaussianField-SmoothMap-Circle-sobolevSeminorm, GaussianField-SmoothMap-Circle-bddAbove-norm-iteratedDeriv-image}
  \texttt{GaussianField.SmoothMap\_Circle.sobolevSeminorm\_nonneg}
\end{theorem}
\begin{theorem}[\texttt{smul\_apply}]\label{GaussianField-SmoothMap-Circle-smul-apply}
  \lean{GaussianField.SmoothMap_Circle.smul_apply}\leanok
  \uses{GaussianField-SmoothMap-Circle-instFunLike, GaussianField-SmoothMap-Circle-instSMul, GaussianField-SmoothMap-Circle}
  \texttt{GaussianField.SmoothMap\_Circle.smul\_apply}
\end{theorem}
\begin{theorem}[\texttt{continuous\_iteratedDeriv}]\label{GaussianField-SmoothMap-Circle-continuous-iteratedDeriv}
  \lean{GaussianField.SmoothMap_Circle.continuous_iteratedDeriv}\leanok
  \uses{GaussianField-SmoothMap-Circle-instFunLike, GaussianField-SmoothMap-Circle-smooth, GaussianField-SmoothMap-Circle}
  \texttt{GaussianField.SmoothMap\_Circle.continuous\_iteratedDeriv}
\end{theorem}
\begin{definition}[\texttt{fourierCoeffLM}]\label{GaussianField-SmoothMap-Circle-fourierCoeffLM}
  \lean{GaussianField.SmoothMap_Circle.fourierCoeffLM}\leanok
  \uses{GaussianField-SmoothMap-Circle-instAddCommGroup, GaussianField-SmoothMap-Circle-fourierCoeffReal, GaussianField-SmoothMap-Circle-instModule, GaussianField-SmoothMap-Circle, GaussianField-SmoothMap-Circle-fourierCoeffReal-add}
  \texttt{GaussianField.SmoothMap\_Circle.fourierCoeffLM}
\end{definition}
\begin{theorem}[\texttt{fourierBasis\_apply}]\label{GaussianField-SmoothMap-Circle-fourierBasis-apply}
  \lean{GaussianField.SmoothMap_Circle.fourierBasis_apply}\leanok
  \uses{GaussianField-SmoothMap-Circle-instFunLike, GaussianField-SmoothMap-Circle-fourierBasis, GaussianField-SmoothMap-Circle-fourierBasisFun, GaussianField-SmoothMap-Circle}
  \texttt{GaussianField.SmoothMap\_Circle.fourierBasis\_apply}
\end{theorem}
\begin{definition}[\texttt{instNeg}]\label{GaussianField-SmoothMap-Circle-instNeg}
  \lean{GaussianField.SmoothMap_Circle.instNeg}\leanok
  \uses{GaussianField-SmoothMap-Circle-instFunLike, GaussianField-SmoothMap-Circle}
  \texttt{GaussianField.SmoothMap\_Circle.instNeg}
\end{definition}
\begin{definition}[\texttt{fourierCoeffReal}]\label{GaussianField-SmoothMap-Circle-fourierCoeffReal}
  \lean{GaussianField.SmoothMap_Circle.fourierCoeffReal}\leanok
  \uses{GaussianField-SmoothMap-Circle-instFunLike, GaussianField-SmoothMap-Circle-fourierBasisFun, GaussianField-SmoothMap-Circle}
  \texttt{GaussianField.SmoothMap\_Circle.fourierCoeffReal}
\end{definition}
\begin{definition}[\texttt{instSub}]\label{GaussianField-SmoothMap-Circle-instSub}
  \lean{GaussianField.SmoothMap_Circle.instSub}\leanok
  \uses{GaussianField-SmoothMap-Circle-instFunLike, GaussianField-SmoothMap-Circle}
  \texttt{GaussianField.SmoothMap\_Circle.instSub}
\end{definition}
\begin{theorem}[\texttt{periodic}]\label{GaussianField-SmoothMap-Circle-periodic}
  \lean{GaussianField.SmoothMap_Circle.periodic}\leanok
  \uses{GaussianField-SmoothMap-Circle-instFunLike, GaussianField-SmoothMap-Circle-periodic-, GaussianField-SmoothMap-Circle}
  \texttt{GaussianField.SmoothMap\_Circle.periodic}
\end{theorem}
\begin{theorem}[\texttt{instContinuousSMul}]\label{GaussianField-SmoothMap-Circle-instContinuousSMul}
  \lean{GaussianField.SmoothMap_Circle.instContinuousSMul}\leanok
  \uses{GaussianField-SmoothMap-Circle-smoothCircle-withSeminorms, GaussianField-SmoothMap-Circle-instAddCommGroup, GaussianField-SmoothMap-Circle-instSMul, GaussianField-SmoothMap-Circle-instModule, GaussianField-SmoothMap-Circle, GaussianField-SmoothMap-Circle-sobolevSeminorm, GaussianField-SmoothMap-Circle-instTopologicalSpace}
  \texttt{GaussianField.SmoothMap\_Circle.instContinuousSMul}
\end{theorem}
\begin{definition}[\texttt{instAdd}]\label{GaussianField-SmoothMap-Circle-instAdd}
  \lean{GaussianField.SmoothMap_Circle.instAdd}\leanok
  \uses{GaussianField-SmoothMap-Circle-instFunLike, GaussianField-SmoothMap-Circle}
  \texttt{GaussianField.SmoothMap\_Circle.instAdd}
\end{definition}
\begin{theorem}[\texttt{fourierCoeffReal\_bound}]\label{GaussianField-SmoothMap-Circle-fourierCoeffReal-bound}
  \lean{GaussianField.SmoothMap_Circle.fourierCoeffReal_bound}\leanok
  \uses{GaussianField-SmoothMap-Circle-instFunLike, GaussianField-SmoothMap-Circle-fourierBasisFun-abs-le, GaussianField-SmoothMap-Circle-norm-iteratedDeriv-le-sobolevSeminorm, GaussianField-SmoothMap-Circle-instAddCommGroup, GaussianField-SmoothMap-Circle-fourierCoeffReal, GaussianField-SmoothMap-Circle-fourierBasisFun, GaussianField-SmoothMap-Circle-instSMul, GaussianField-SmoothMap-Circle, GaussianField-SmoothMap-Circle-sobolevSeminorm-nonneg, GaussianField-SmoothMap-Circle-sobolevSeminorm}
  \texttt{GaussianField.SmoothMap\_Circle.fourierCoeffReal\_bound}
\end{theorem}
\begin{definition}[\texttt{instModule}]\label{GaussianField-SmoothMap-Circle-instModule}
  \lean{GaussianField.SmoothMap_Circle.instModule}\leanok
  \uses{GaussianField-SmoothMap-Circle-instAddCommGroup, GaussianField-SmoothMap-Circle-instSMul, GaussianField-SmoothMap-Circle}
  \texttt{GaussianField.SmoothMap\_Circle.instModule}
\end{definition}
\begin{definition}[\texttt{fourierBasisFun}]\label{GaussianField-SmoothMap-Circle-fourierBasisFun}
  \lean{GaussianField.SmoothMap_Circle.fourierBasisFun}\leanok
  \texttt{GaussianField.SmoothMap\_Circle.fourierBasisFun}
\end{definition}
\begin{theorem}[\texttt{bddAbove\_norm\_iteratedDeriv\_image}]\label{GaussianField-SmoothMap-Circle-bddAbove-norm-iteratedDeriv-image}
  \lean{GaussianField.SmoothMap_Circle.bddAbove_norm_iteratedDeriv_image}\leanok
  \uses{GaussianField-SmoothMap-Circle-instFunLike, GaussianField-SmoothMap-Circle, GaussianField-SmoothMap-Circle-continuous-iteratedDeriv, Lattice-toSemilatticeInf}
  \texttt{GaussianField.SmoothMap\_Circle.bddAbove\_norm\_iteratedDeriv\_image}
\end{theorem}
\begin{theorem}[\texttt{periodic\_iteratedDeriv}]\label{GaussianField-SmoothMap-Circle-periodic-iteratedDeriv}
  \lean{GaussianField.SmoothMap_Circle.periodic_iteratedDeriv}\leanok
  \uses{GaussianField-SmoothMap-Circle-instFunLike, GaussianField-SmoothMap-Circle-periodic, GaussianField-SmoothMap-Circle-smooth, GaussianField-SmoothMap-Circle}
  \texttt{GaussianField.SmoothMap\_Circle.periodic\_iteratedDeriv}
\end{theorem}
\begin{theorem}[\texttt{fourierCoeffReal\_add}]\label{GaussianField-SmoothMap-Circle-fourierCoeffReal-add}
  \lean{GaussianField.SmoothMap_Circle.fourierCoeffReal_add}\leanok
  \uses{GaussianField-SmoothMap-Circle-instFunLike, GaussianField-SmoothMap-Circle-continuous, GaussianField-SmoothMap-Circle-instAdd, GaussianField-SmoothMap-Circle-fourierBasisFun-smooth, GaussianField-SmoothMap-Circle-fourierCoeffReal, GaussianField-SmoothMap-Circle-fourierBasisFun, GaussianField-SmoothMap-Circle}
  \texttt{GaussianField.SmoothMap\_Circle.fourierCoeffReal\_add}
\end{theorem}
\begin{theorem}[\texttt{instIsTopologicalAddGroup}]\label{GaussianField-SmoothMap-Circle-instIsTopologicalAddGroup}
  \lean{GaussianField.SmoothMap_Circle.instIsTopologicalAddGroup}\leanok
  \uses{GaussianField-SmoothMap-Circle-smoothCircle-withSeminorms, GaussianField-SmoothMap-Circle-instAddCommGroup, GaussianField-SmoothMap-Circle-instModule, GaussianField-SmoothMap-Circle, GaussianField-SmoothMap-Circle-sobolevSeminorm, GaussianField-SmoothMap-Circle-instTopologicalSpace}
  \texttt{GaussianField.SmoothMap\_Circle.instIsTopologicalAddGroup}
\end{theorem}
\begin{definition}[\texttt{SmoothMap\_Circle}]\label{GaussianField-SmoothMap-Circle}
  \lean{GaussianField.SmoothMap_Circle}\leanok
  \texttt{GaussianField.SmoothMap\_Circle}
\end{definition}
\begin{theorem}[\texttt{smooth}]\label{GaussianField-SmoothMap-Circle-smooth}
  \lean{GaussianField.SmoothMap_Circle.smooth}\leanok
  \uses{GaussianField-SmoothMap-Circle-instFunLike, GaussianField-SmoothMap-Circle-smooth-, GaussianField-SmoothMap-Circle}
  \texttt{GaussianField.SmoothMap\_Circle.smooth}
\end{theorem}
\begin{theorem}[\texttt{zero\_apply}]\label{GaussianField-SmoothMap-Circle-zero-apply}
  \lean{GaussianField.SmoothMap_Circle.zero_apply}\leanok
  \uses{GaussianField-SmoothMap-Circle-instFunLike, GaussianField-SmoothMap-Circle-instZero, GaussianField-SmoothMap-Circle}
  \texttt{GaussianField.SmoothMap\_Circle.zero\_apply}
\end{theorem}
\begin{theorem}[\texttt{coe\_zero}]\label{GaussianField-SmoothMap-Circle-coe-zero}
  \lean{GaussianField.SmoothMap_Circle.coe_zero}\leanok
  \uses{GaussianField-SmoothMap-Circle-instFunLike, GaussianField-SmoothMap-Circle-instZero, GaussianField-SmoothMap-Circle}
  \texttt{GaussianField.SmoothMap\_Circle.coe\_zero}
\end{theorem}
\begin{theorem}[\texttt{fourierBasisFun\_periodic}]\label{GaussianField-SmoothMap-Circle-fourierBasisFun-periodic}
  \lean{GaussianField.SmoothMap_Circle.fourierBasisFun_periodic}\leanok
  \uses{GaussianField-SmoothMap-Circle-fourierBasisFun}
  \texttt{GaussianField.SmoothMap\_Circle.fourierBasisFun\_periodic}
\end{theorem}
\begin{theorem}[\texttt{sobolevSeminorm\_fourierBasis\_le}]\label{GaussianField-SmoothMap-Circle-sobolevSeminorm-fourierBasis-le}
  \lean{GaussianField.SmoothMap_Circle.sobolevSeminorm_fourierBasis_le}\leanok
  \uses{GaussianField-SmoothMap-Circle-instFunLike, GaussianField-SmoothMap-Circle-fourierBasis, GaussianField-SmoothMap-Circle-Icc-nonempty, GaussianField-SmoothMap-Circle-instAddCommGroup, GaussianField-SmoothMap-Circle-fourierBasisFun, GaussianField-SmoothMap-Circle-instSMul, GaussianField-SmoothMap-Circle, GaussianField-SmoothMap-Circle-sobolevSeminorm, Lattice-toSemilatticeInf}
  \texttt{GaussianField.SmoothMap\_Circle.sobolevSeminorm\_fourierBasis\_le}
\end{theorem}
\begin{theorem}[\texttt{sobolevSeminorm\_affine\_precomp\_le}]\label{GaussianField-sobolevSeminorm-affine-precomp-le}
  \lean{GaussianField.sobolevSeminorm_affine_precomp_le}\leanok
  \uses{GaussianField-SmoothMap-Circle-instFunLike, GaussianField-SmoothMap-Circle-Icc-nonempty, GaussianField-SmoothMap-Circle-instAddCommGroup, GaussianField-SmoothMap-Circle-norm-iteratedDeriv-le-sobolevSeminorm-, GaussianField-SmoothMap-Circle-smooth, GaussianField-SmoothMap-Circle-instSMul, GaussianField-SmoothMap-Circle, GaussianField-SmoothMap-Circle-sobolevSeminorm, Lattice-toSemilatticeInf}
  \texttt{GaussianField.sobolevSeminorm\_affine\_precomp\_le}
\end{theorem}
\begin{theorem}[\texttt{fourierCoeffReal\_smul}]\label{GaussianField-SmoothMap-Circle-fourierCoeffReal-smul}
  \lean{GaussianField.SmoothMap_Circle.fourierCoeffReal_smul}\leanok
  \uses{GaussianField-SmoothMap-Circle-instFunLike, GaussianField-SmoothMap-Circle-fourierCoeffReal, GaussianField-SmoothMap-Circle-fourierBasisFun, GaussianField-SmoothMap-Circle-instSMul, GaussianField-SmoothMap-Circle}
  \texttt{GaussianField.SmoothMap\_Circle.fourierCoeffReal\_smul}
\end{theorem}
\begin{theorem}[\texttt{fourierBasisFun\_abs\_le}]\label{GaussianField-SmoothMap-Circle-fourierBasisFun-abs-le}
  \lean{GaussianField.SmoothMap_Circle.fourierBasisFun_abs_le}\leanok
  \uses{GaussianField-SmoothMap-Circle-fourierBasisFun}
  \texttt{GaussianField.SmoothMap\_Circle.fourierBasisFun\_abs\_le}
\end{theorem}
\begin{theorem}[\texttt{Icc\_nonempty}]\label{GaussianField-SmoothMap-Circle-Icc-nonempty}
  \lean{GaussianField.SmoothMap_Circle.Icc_nonempty}\leanok
  \texttt{GaussianField.SmoothMap\_Circle.Icc\_nonempty}
\end{theorem}
\begin{theorem}[\texttt{smoothCircle\_withSeminorms}]\label{GaussianField-SmoothMap-Circle-smoothCircle-withSeminorms}
  \lean{GaussianField.SmoothMap_Circle.smoothCircle_withSeminorms}\leanok
  \uses{GaussianField-SmoothMap-Circle-instAddCommGroup, GaussianField-SmoothMap-Circle-instModule, GaussianField-SmoothMap-Circle, GaussianField-SmoothMap-Circle-sobolevSeminorm, GaussianField-SmoothMap-Circle-instTopologicalSpace}
  \texttt{GaussianField.SmoothMap\_Circle.smoothCircle\_withSeminorms}
\end{theorem}
\begin{definition}[\texttt{instAddCommGroup}]\label{GaussianField-SmoothMap-Circle-instAddCommGroup}
  \lean{GaussianField.SmoothMap_Circle.instAddCommGroup}\leanok
  \uses{GaussianField-SmoothMap-Circle-instAdd, GaussianField-SmoothMap-Circle-instZero, GaussianField-SmoothMap-Circle-instSub, GaussianField-SmoothMap-Circle, GaussianField-SmoothMap-Circle-instNeg}
  \texttt{GaussianField.SmoothMap\_Circle.instAddCommGroup}
\end{definition}
\begin{definition}[\texttt{sobolevSeminorm}]\label{GaussianField-SmoothMap-Circle-sobolevSeminorm}
  \lean{GaussianField.SmoothMap_Circle.sobolevSeminorm}\leanok
  \uses{GaussianField-SmoothMap-Circle-instFunLike, GaussianField-SmoothMap-Circle-instAddCommGroup, GaussianField-SmoothMap-Circle-instSMul, GaussianField-SmoothMap-Circle}
  \texttt{GaussianField.SmoothMap\_Circle.sobolevSeminorm}
\end{definition}
\begin{definition}[\texttt{fourierFreq}]\label{GaussianField-SmoothMap-Circle-fourierFreq}
  \lean{GaussianField.SmoothMap_Circle.fourierFreq}\leanok
  \texttt{GaussianField.SmoothMap\_Circle.fourierFreq}
\end{definition}
\begin{definition}[\texttt{fourierBasis}]\label{GaussianField-SmoothMap-Circle-fourierBasis}
  \lean{GaussianField.SmoothMap_Circle.fourierBasis}\leanok
  \uses{GaussianField-SmoothMap-Circle-fourierBasisFun-smooth, GaussianField-SmoothMap-Circle-fourierBasisFun, GaussianField-SmoothMap-Circle, GaussianField-SmoothMap-Circle-fourierBasisFun-periodic}
  \texttt{GaussianField.SmoothMap\_Circle.fourierBasis}
\end{definition}
\begin{theorem}[\texttt{norm\_iteratedDeriv\_le\_sobolevSeminorm'}]\label{GaussianField-SmoothMap-Circle-norm-iteratedDeriv-le-sobolevSeminorm-}
  \lean{GaussianField.SmoothMap_Circle.norm_iteratedDeriv_le_sobolevSeminorm'}\leanok
  \uses{GaussianField-SmoothMap-Circle-periodic-iteratedDeriv, GaussianField-SmoothMap-Circle-instFunLike, GaussianField-SmoothMap-Circle-norm-iteratedDeriv-le-sobolevSeminorm, GaussianField-SmoothMap-Circle-instAddCommGroup, GaussianField-SmoothMap-Circle-instSMul, GaussianField-SmoothMap-Circle, GaussianField-SmoothMap-Circle-sobolevSeminorm, Lattice-toSemilatticeInf}
  \texttt{GaussianField.SmoothMap\_Circle.norm\_iteratedDeriv\_le\_sobolevSeminorm'}
\end{theorem}
\begin{theorem}[\texttt{norm\_iteratedDeriv\_le\_sobolevSeminorm}]\label{GaussianField-SmoothMap-Circle-norm-iteratedDeriv-le-sobolevSeminorm}
  \lean{GaussianField.SmoothMap_Circle.norm_iteratedDeriv_le_sobolevSeminorm}\leanok
  \uses{GaussianField-SmoothMap-Circle-instFunLike, GaussianField-SmoothMap-Circle-instAddCommGroup, GaussianField-SmoothMap-Circle-instSMul, GaussianField-SmoothMap-Circle, GaussianField-SmoothMap-Circle-sobolevSeminorm, GaussianField-SmoothMap-Circle-bddAbove-norm-iteratedDeriv-image}
  \texttt{GaussianField.SmoothMap\_Circle.norm\_iteratedDeriv\_le\_sobolevSeminorm}
\end{theorem}
\begin{theorem}[\texttt{fourierCoeffReal\_fourierBasis}]\label{GaussianField-SmoothMap-Circle-fourierCoeffReal-fourierBasis}
  \lean{GaussianField.SmoothMap_Circle.fourierCoeffReal_fourierBasis}\leanok
  \uses{GaussianField-SmoothMap-Circle-instFunLike, GaussianField-SmoothMap-Circle-fourierBasis, GaussianField-SmoothMap-Circle-fourierBasisFun-smooth, GaussianField-SmoothMap-Circle-fourierCoeffReal, GaussianField-SmoothMap-Circle-fourierBasisFun, GaussianField-SmoothMap-Circle, Lattice-toSemilatticeInf}
  \texttt{GaussianField.SmoothMap\_Circle.fourierCoeffReal\_fourierBasis}
\end{theorem}
\begin{theorem}[\texttt{smooth'}]\label{GaussianField-SmoothMap-Circle-smooth-}
  \lean{GaussianField.SmoothMap_Circle.smooth'}\leanok
  \uses{GaussianField-SmoothMap-Circle-toFun, GaussianField-SmoothMap-Circle}
  \texttt{GaussianField.SmoothMap\_Circle.smooth'}
\end{theorem}
\begin{definition}[\texttt{ctorIdx}]\label{GaussianField-SmoothMap-Circle-ctorIdx}
  \lean{GaussianField.SmoothMap_Circle.ctorIdx}\leanok
  \uses{GaussianField-SmoothMap-Circle}
  \texttt{GaussianField.SmoothMap\_Circle.ctorIdx}
\end{definition}
\begin{definition}[\texttt{instSMul}]\label{GaussianField-SmoothMap-Circle-instSMul}
  \lean{GaussianField.SmoothMap_Circle.instSMul}\leanok
  \uses{GaussianField-SmoothMap-Circle-instFunLike, GaussianField-SmoothMap-Circle}
  \texttt{GaussianField.SmoothMap\_Circle.instSMul}
\end{definition}
\begin{definition}[\texttt{instZero}]\label{GaussianField-SmoothMap-Circle-instZero}
  \lean{GaussianField.SmoothMap_Circle.instZero}\leanok
  \uses{GaussianField-SmoothMap-Circle}
  \texttt{GaussianField.SmoothMap\_Circle.instZero}
\end{definition}
\begin{definition}[\texttt{toFun}]\label{GaussianField-SmoothMap-Circle-toFun}
  \lean{GaussianField.SmoothMap_Circle.toFun}\leanok
  \uses{GaussianField-SmoothMap-Circle}
  \texttt{GaussianField.SmoothMap\_Circle.toFun}
\end{definition}
\begin{definition}[\texttt{instFunLike}]\label{GaussianField-SmoothMap-Circle-instFunLike}
  \lean{GaussianField.SmoothMap_Circle.instFunLike}\leanok
  \uses{GaussianField-SmoothMap-Circle-toFun, GaussianField-SmoothMap-Circle}
  \texttt{GaussianField.SmoothMap\_Circle.instFunLike}
\end{definition}
\begin{theorem}[\texttt{contDiffAt\_of\_smooth}]\label{GaussianField-SmoothMap-Circle-contDiffAt-of-smooth}
  \lean{GaussianField.SmoothMap_Circle.contDiffAt_of_smooth}\leanok
  \uses{GaussianField-SmoothMap-Circle-instFunLike, GaussianField-SmoothMap-Circle-smooth, GaussianField-SmoothMap-Circle}
  \texttt{GaussianField.SmoothMap\_Circle.contDiffAt\_of\_smooth}
\end{theorem}
\begin{theorem}[\texttt{neg\_apply}]\label{GaussianField-SmoothMap-Circle-neg-apply}
  \lean{GaussianField.SmoothMap_Circle.neg_apply}\leanok
  \uses{GaussianField-SmoothMap-Circle-instFunLike, GaussianField-SmoothMap-Circle, GaussianField-SmoothMap-Circle-instNeg}
  \texttt{GaussianField.SmoothMap\_Circle.neg\_apply}
\end{theorem}
\begin{theorem}[\texttt{fourierBasisFun\_smooth}]\label{GaussianField-SmoothMap-Circle-fourierBasisFun-smooth}
  \lean{GaussianField.SmoothMap_Circle.fourierBasisFun_smooth}\leanok
  \uses{GaussianField-SmoothMap-Circle-fourierBasisFun}
  \texttt{GaussianField.SmoothMap\_Circle.fourierBasisFun\_smooth}
\end{theorem}
\begin{theorem}[\texttt{periodic'}]\label{GaussianField-SmoothMap-Circle-periodic-}
  \lean{GaussianField.SmoothMap_Circle.periodic'}\leanok
  \uses{GaussianField-SmoothMap-Circle-toFun, GaussianField-SmoothMap-Circle}
  \texttt{GaussianField.SmoothMap\_Circle.periodic'}
\end{theorem}
\begin{theorem}[\texttt{congr\_simp}]\label{GaussianField-SmoothMap-Circle-mk-congr-simp}
  \lean{GaussianField.SmoothMap_Circle.mk.congr_simp}\leanok
  \uses{GaussianField-SmoothMap-Circle}
  \texttt{GaussianField.SmoothMap\_Circle.mk.congr\_simp}
\end{theorem}
\begin{theorem}[\texttt{continuous}]\label{GaussianField-SmoothMap-Circle-continuous}
  \lean{GaussianField.SmoothMap_Circle.continuous}\leanok
  \uses{GaussianField-SmoothMap-Circle-instFunLike, GaussianField-SmoothMap-Circle-smooth, GaussianField-SmoothMap-Circle}
  \texttt{GaussianField.SmoothMap\_Circle.continuous}
\end{theorem}
\begin{theorem}[\texttt{ext\_iff}]\label{GaussianField-SmoothMap-Circle-ext-iff}
  \lean{GaussianField.SmoothMap_Circle.ext_iff}\leanok
  \uses{GaussianField-SmoothMap-Circle-instFunLike, GaussianField-SmoothMap-Circle-ext, GaussianField-SmoothMap-Circle}
  \texttt{GaussianField.SmoothMap\_Circle.ext\_iff}
\end{theorem}
\begin{theorem}[\texttt{add\_apply}]\label{GaussianField-SmoothMap-Circle-add-apply}
  \lean{GaussianField.SmoothMap_Circle.add_apply}\leanok
  \uses{GaussianField-SmoothMap-Circle-instFunLike, GaussianField-SmoothMap-Circle-instAdd, GaussianField-SmoothMap-Circle}
  \texttt{GaussianField.SmoothMap\_Circle.add\_apply}
\end{theorem}
\begin{definition}[\texttt{fourierCoeffCLM}]\label{GaussianField-SmoothMap-Circle-fourierCoeffCLM}
  \lean{GaussianField.SmoothMap_Circle.fourierCoeffCLM}\leanok
  \uses{GaussianField-SmoothMap-Circle-instAddCommGroup, GaussianField-SmoothMap-Circle-fourierCoeffLM, GaussianField-SmoothMap-Circle-instModule, GaussianField-SmoothMap-Circle, GaussianField-SmoothMap-Circle-instTopologicalSpace}
  \texttt{GaussianField.SmoothMap\_Circle.fourierCoeffCLM}
\end{definition}
\begin{theorem}[\texttt{ext}]\label{GaussianField-SmoothMap-Circle-ext}
  \lean{GaussianField.SmoothMap_Circle.ext}\leanok
  \uses{GaussianField-SmoothMap-Circle-instFunLike, GaussianField-SmoothMap-Circle}
  \texttt{GaussianField.SmoothMap\_Circle.ext}
\end{theorem}
\begin{definition}[\texttt{instTopologicalSpace}]\label{GaussianField-SmoothMap-Circle-instTopologicalSpace}
  \lean{GaussianField.SmoothMap_Circle.instTopologicalSpace}\leanok
  \uses{GaussianField-SmoothMap-Circle-instAddCommGroup, GaussianField-SmoothMap-Circle-instModule, GaussianField-SmoothMap-Circle, GaussianField-SmoothMap-Circle-sobolevSeminorm}
  \texttt{GaussianField.SmoothMap\_Circle.instTopologicalSpace}
\end{definition}
\section{\texttt{Pphi2.TorusContinuumLimit.TorusEmbedding}}
\begin{theorem}[\texttt{congr\_simp}]\label{GaussianField-torusEmbedCLMGJ-congr-simp}
  \lean{GaussianField.torusEmbedCLMGJ.congr_simp}\leanok
  \uses{GaussianField-NuclearTensorProduct-instModuleReal, GaussianField-NuclearTensorProduct-instTopologicalSpace, GaussianField-NuclearTensorProduct-instIsTopologicalAddGroup, GaussianField-NuclearTensorProduct-instAddCommGroup, GaussianField-Configuration, GaussianField-NuclearTensorProduct-instContinuousSMulReal, GaussianField-SmoothMap-Circle, GaussianField-torusEmbedCLMGJ, GaussianField-FinLatticeField, GaussianField-TorusTestFunction}
  \texttt{GaussianField.torusEmbedCLMGJ.congr\_simp}
\end{theorem}
\section{\texttt{Nuclear.NuclearSpace}}
\begin{theorem}[\texttt{summable\_coeff\_seminorm\_basis}]\label{GaussianField-DyninMityaginSpace-summable-coeff-seminorm-basis}
  \lean{GaussianField.DyninMityaginSpace.summable_coeff_seminorm_basis}\leanok
  \uses{GaussianField-DyninMityaginSpace, GaussianField-DyninMityaginSpace-coeff, GaussianField-DyninMityaginSpace-coeff-decay, GaussianField-DyninMityaginSpace--, GaussianField-DyninMityaginSpace-basis, GaussianField-DyninMityaginSpace-basis-growth, GaussianField-DyninMityaginSpace-p, Lattice-toSemilatticeInf}
  \texttt{GaussianField.DyninMityaginSpace.summable\_coeff\_seminorm\_basis}
\end{theorem}
\begin{theorem}[\texttt{exists\_CLF\_le\_seminorm}]\label{GaussianField-exists-CLF-le-seminorm}
  \lean{GaussianField.exists_CLF_le_seminorm}\leanok
  \uses{Lattice-toSemilatticeSup, Lattice-toSemilatticeInf}
  \texttt{GaussianField.exists\_CLF\_le\_seminorm}
\end{theorem}
\begin{theorem}[\texttt{toNuclearSpace}]\label{GaussianField-DyninMityaginSpace-toNuclearSpace}
  \lean{GaussianField.DyninMityaginSpace.toNuclearSpace}\leanok
  \uses{GaussianField-DyninMityaginSpace, GaussianField-DyninMityaginSpace-coeff, GaussianField-DyninMityaginSpace-coeff-decay, GaussianField-DyninMityaginSpace--, GaussianField-seminorm-le-nuclear-expansion, GaussianField-NuclearSpace, GaussianField-finset-sup-poly-bound, GaussianField-DyninMityaginSpace-h-with, GaussianField-DyninMityaginSpace-basis, GaussianField-DyninMityaginSpace-basis-growth, GaussianField-DyninMityaginSpace-p}
  \texttt{GaussianField.DyninMityaginSpace.toNuclearSpace}
\end{theorem}
\begin{theorem}[\texttt{hasSum\_basis}]\label{GaussianField-DyninMityaginSpace-hasSum-basis}
  \lean{GaussianField.DyninMityaginSpace.hasSum_basis}\leanok
  \uses{GaussianField-DyninMityaginSpace, GaussianField-DyninMityaginSpace-coeff, GaussianField-DyninMityaginSpace--, GaussianField-DyninMityaginSpace-h-with, GaussianField-DyninMityaginSpace-summable-coeff-seminorm-basis, GaussianField-DyninMityaginSpace-basis, GaussianField-DyninMityaginSpace-p, Lattice-toSemilatticeInf}
  \texttt{GaussianField.DyninMityaginSpace.hasSum\_basis}
\end{theorem}
\begin{definition}[\texttt{NuclearSpace}]\label{GaussianField-NuclearSpace}
  \lean{GaussianField.NuclearSpace}\leanok
  \texttt{GaussianField.NuclearSpace}
\end{definition}
\begin{theorem}[\texttt{finset\_sup\_poly\_bound}]\label{GaussianField-finset-sup-poly-bound}
  \lean{GaussianField.finset_sup_poly_bound}\leanok
  \texttt{GaussianField.finset\_sup\_poly\_bound}
\end{theorem}
\begin{theorem}[\texttt{nuclear\_dominance}]\label{GaussianField-NuclearSpace-nuclear-dominance}
  \lean{GaussianField.NuclearSpace.nuclear_dominance}\leanok
  \uses{GaussianField-NuclearSpace}
  \texttt{GaussianField.NuclearSpace.nuclear\_dominance}
\end{theorem}
\begin{theorem}[\texttt{seminorm\_le\_nuclear\_expansion}]\label{GaussianField-seminorm-le-nuclear-expansion}
  \lean{GaussianField.seminorm_le_nuclear_expansion}\leanok
  \uses{GaussianField-DyninMityaginSpace, GaussianField-DyninMityaginSpace-coeff, GaussianField-DyninMityaginSpace-coeff-decay, GaussianField-DyninMityaginSpace-expansion, GaussianField-DyninMityaginSpace--, GaussianField-finset-sup-poly-bound, GaussianField-DyninMityaginSpace-h-with, GaussianField-DyninMityaginSpace-basis, GaussianField-DyninMityaginSpace-basis-growth, GaussianField-exists-CLF-le-seminorm, GaussianField-DyninMityaginSpace-p, Lattice-toSemilatticeInf}
  \texttt{GaussianField.seminorm\_le\_nuclear\_expansion}
\end{theorem}
\section{\texttt{Cylinder.FourierMultiplier}}
\begin{theorem}[\texttt{resolventMultiplierCLM\_comp\_inverseResolventMultiplierCLM}]\label{GaussianField-resolventMultiplierCLM-comp-inverseResolventMultiplierCLM}
  \lean{GaussianField.resolventMultiplierCLM_comp_inverseResolventMultiplierCLM}\leanok
  \uses{GaussianField-realFourierMultiplierCLM-congr-simp, GaussianField-inverseResolventSymbol-hasTemperateGrowth, GaussianField-inverseResolventSymbol-even, GaussianField-inverseResolventSymbol, GaussianField-resolventSymbol-hasTemperateGrowth, GaussianField-realFourierMultiplierCLM, GaussianField-resolventMultiplierCLM, GaussianField-realFourierMultiplierCLM-comp, GaussianField-resolventSymbol, GaussianField-inverseResolventMultiplierCLM, GaussianField-realFourierMultiplierCLM-one, GaussianField-resolventSymbol-mul-inverse}
  \texttt{GaussianField.resolventMultiplierCLM\_comp\_inverseResolventMultiplierCLM}
\end{theorem}
\begin{theorem}[\texttt{schwartzToReal\_comp\_schwartzToComplex}]\label{GaussianField-schwartzToReal-comp-schwartzToComplex}
  \lean{GaussianField.schwartzToReal_comp_schwartzToComplex}\leanok
  \uses{GaussianField-schwartzToComplex, GaussianField-schwartzToReal}
  \texttt{GaussianField.schwartzToReal\_comp\_schwartzToComplex}
\end{theorem}
\begin{theorem}[\texttt{inverseResolventMultiplierCLM\_comp\_resolventMultiplierCLM}]\label{GaussianField-inverseResolventMultiplierCLM-comp-resolventMultiplierCLM}
  \lean{GaussianField.inverseResolventMultiplierCLM_comp_resolventMultiplierCLM}\leanok
  \uses{GaussianField-realFourierMultiplierCLM-congr-simp, GaussianField-resolventSymbol-even, GaussianField-inverseResolventSymbol-hasTemperateGrowth, GaussianField-inverseResolventSymbol, GaussianField-resolventSymbol-hasTemperateGrowth, GaussianField-realFourierMultiplierCLM, GaussianField-resolventMultiplierCLM, GaussianField-realFourierMultiplierCLM-comp, GaussianField-resolventSymbol, GaussianField-inverseResolventMultiplierCLM, GaussianField-realFourierMultiplierCLM-one, GaussianField-resolventSymbol-mul-inverse}
  \texttt{GaussianField.inverseResolventMultiplierCLM\_comp\_resolventMultiplierCLM}
\end{theorem}
\begin{theorem}[\texttt{heatMultiplierCLM\_translation\_comm}]\label{GaussianField-heatMultiplierCLM-translation-comm}
  \lean{GaussianField.heatMultiplierCLM_translation_comm}\leanok
  \uses{GaussianField-realFourierMultiplierCLM-translation-comm, GaussianField-heatSymbol, GaussianField-heatSymbol-hasTemperateGrowth, GaussianField-heatMultiplierCLM, GaussianField-schwartzTranslation}
  \texttt{GaussianField.heatMultiplierCLM\_translation\_comm}
\end{theorem}
\begin{definition}[\texttt{heatSymbol}]\label{GaussianField-heatSymbol}
  \lean{GaussianField.heatSymbol}\leanok
  \texttt{GaussianField.heatSymbol}
\end{definition}
\begin{theorem}[\texttt{resolventSymbol\_mul\_inverse}]\label{GaussianField-resolventSymbol-mul-inverse}
  \lean{GaussianField.resolventSymbol_mul_inverse}\leanok
  \uses{GaussianField-inverseResolventSymbol, GaussianField-resolventSymbol, Lattice-toSemilatticeInf}
  \texttt{GaussianField.resolventSymbol\_mul\_inverse}
\end{theorem}
\begin{theorem}[\texttt{freeHeatSemigroup\_reflection\_comm}]\label{GaussianField-freeHeatSemigroup-reflection-comm}
  \lean{GaussianField.freeHeatSemigroup_reflection_comm}\leanok
  \uses{GaussianField-freeHeatSemigroup, GaussianField-heatMultiplierCLM-reflection-comm, GaussianField-schwartzReflection}
  \texttt{GaussianField.freeHeatSemigroup\_reflection\_comm}
\end{theorem}
\begin{definition}[\texttt{heatMultiplierCLM}]\label{GaussianField-heatMultiplierCLM}
  \lean{GaussianField.heatMultiplierCLM}\leanok
  \uses{GaussianField-heatSymbol, GaussianField-heatSymbol-hasTemperateGrowth, GaussianField-realFourierMultiplierCLM}
  \texttt{GaussianField.heatMultiplierCLM}
\end{definition}
\begin{definition}[\texttt{resolventMultiplierCLM}]\label{GaussianField-resolventMultiplierCLM}
  \lean{GaussianField.resolventMultiplierCLM}\leanok
  \uses{GaussianField-resolventSymbol-hasTemperateGrowth, GaussianField-realFourierMultiplierCLM, GaussianField-resolventSymbol}
  \texttt{GaussianField.resolventMultiplierCLM}
\end{definition}
\begin{theorem}[\texttt{realFourierMultiplierCLM\_apply\_eq}]\label{GaussianField-realFourierMultiplierCLM-apply-eq}
  \lean{GaussianField.realFourierMultiplierCLM_apply_eq}\leanok
  \uses{GaussianField-realFourierMultiplierCLM}
  \texttt{GaussianField.realFourierMultiplierCLM\_apply\_eq}
\end{theorem}
\begin{definition}[\texttt{freeHeatSemigroup}]\label{GaussianField-freeHeatSemigroup}
  \lean{GaussianField.freeHeatSemigroup}\leanok
  \uses{GaussianField-heatMultiplierCLM}
  \texttt{GaussianField.freeHeatSemigroup}
\end{definition}
\begin{theorem}[\texttt{heatSymbol\_even}]\label{GaussianField-heatSymbol-even}
  \lean{GaussianField.heatSymbol_even}\leanok
  \uses{GaussianField-heatSymbol}
  \texttt{GaussianField.heatSymbol\_even}
\end{theorem}
\begin{theorem}[\texttt{realFourierMultiplierCLM\_comp}]\label{GaussianField-realFourierMultiplierCLM-comp}
  \lean{GaussianField.realFourierMultiplierCLM_comp}\leanok
  \uses{GaussianField-realFourierMultiplierCLM, GaussianField-fourierMultiplier-preserves-real, GaussianField-schwartzToComplex, GaussianField-schwartzToReal}
  \texttt{GaussianField.realFourierMultiplierCLM\_comp}
\end{theorem}
\begin{definition}[\texttt{resolventSymbol}]\label{GaussianField-resolventSymbol}
  \lean{GaussianField.resolventSymbol}\leanok
  \texttt{GaussianField.resolventSymbol}
\end{definition}
\begin{theorem}[\texttt{resolventSymbol\_hasTemperateGrowth}]\label{GaussianField-resolventSymbol-hasTemperateGrowth}
  \lean{GaussianField.resolventSymbol_hasTemperateGrowth}\leanok
  \uses{GaussianField-resolventSymbol, Lattice-toSemilatticeInf}
  \texttt{GaussianField.resolventSymbol\_hasTemperateGrowth}
\end{theorem}
\begin{theorem}[\texttt{resolventMultiplierCLM\_translation\_comm}]\label{GaussianField-resolventMultiplierCLM-translation-comm}
  \lean{GaussianField.resolventMultiplierCLM_translation_comm}\leanok
  \uses{GaussianField-realFourierMultiplierCLM-translation-comm, GaussianField-resolventSymbol-hasTemperateGrowth, GaussianField-resolventMultiplierCLM, GaussianField-schwartzTranslation, GaussianField-resolventSymbol}
  \texttt{GaussianField.resolventMultiplierCLM\_translation\_comm}
\end{theorem}
\begin{theorem}[\texttt{inverseResolventSymbol\_hasTemperateGrowth}]\label{GaussianField-inverseResolventSymbol-hasTemperateGrowth}
  \lean{GaussianField.inverseResolventSymbol_hasTemperateGrowth}\leanok
  \uses{GaussianField-inverseResolventSymbol, Lattice-toSemilatticeInf}
  \texttt{GaussianField.inverseResolventSymbol\_hasTemperateGrowth}
\end{theorem}
\begin{theorem}[\texttt{freeHeatSemigroup\_translation\_comm}]\label{GaussianField-freeHeatSemigroup-translation-comm}
  \lean{GaussianField.freeHeatSemigroup_translation_comm}\leanok
  \uses{GaussianField-heatMultiplierCLM-translation-comm, GaussianField-freeHeatSemigroup, GaussianField-schwartzTranslation}
  \texttt{GaussianField.freeHeatSemigroup\_translation\_comm}
\end{theorem}
\begin{theorem}[\texttt{freeHeatSemigroup\_zero}]\label{GaussianField-freeHeatSemigroup-zero}
  \lean{GaussianField.freeHeatSemigroup_zero}\leanok
  \uses{GaussianField-heatMultiplierCLM-zero, GaussianField-freeHeatSemigroup}
  \texttt{GaussianField.freeHeatSemigroup\_zero}
\end{theorem}
\begin{definition}[\texttt{realFourierMultiplierCLM}]\label{GaussianField-realFourierMultiplierCLM}
  \lean{GaussianField.realFourierMultiplierCLM}\leanok
  \uses{GaussianField-schwartzToComplex, GaussianField-schwartzToReal}
  \texttt{GaussianField.realFourierMultiplierCLM}
\end{definition}
\begin{theorem}[\texttt{fourierMultiplier\_preserves\_real}]\label{GaussianField-fourierMultiplier-preserves-real}
  \lean{GaussianField.fourierMultiplier_preserves_real}\leanok
  \uses{GaussianField-schwartzToComplex, GaussianField-schwartzToReal}
  \texttt{GaussianField.fourierMultiplier\_preserves\_real}
\end{theorem}
\begin{theorem}[\texttt{freeHeatSemigroup\_comp}]\label{GaussianField-freeHeatSemigroup-comp}
  \lean{GaussianField.freeHeatSemigroup_comp}\leanok
  \uses{GaussianField-heatMultiplierCLM-comp, GaussianField-freeHeatSemigroup}
  \texttt{GaussianField.freeHeatSemigroup\_comp}
\end{theorem}
\begin{theorem}[\texttt{heatSymbol\_mul}]\label{GaussianField-heatSymbol-mul}
  \lean{GaussianField.heatSymbol_mul}\leanok
  \uses{GaussianField-heatSymbol}
  \texttt{GaussianField.heatSymbol\_mul}
\end{theorem}
\begin{theorem}[\texttt{heatMultiplierCLM\_zero}]\label{GaussianField-heatMultiplierCLM-zero}
  \lean{GaussianField.heatMultiplierCLM_zero}\leanok
  \uses{GaussianField-realFourierMultiplierCLM-congr-simp, GaussianField-heatSymbol, GaussianField-heatSymbol-hasTemperateGrowth, GaussianField-heatMultiplierCLM, GaussianField-realFourierMultiplierCLM, GaussianField-schwartzToComplex, GaussianField-schwartzToReal, GaussianField-heatSymbol-zero}
  \texttt{GaussianField.heatMultiplierCLM\_zero}
\end{theorem}
\begin{definition}[\texttt{schwartzToComplex}]\label{GaussianField-schwartzToComplex}
  \lean{GaussianField.schwartzToComplex}\leanok
  \texttt{GaussianField.schwartzToComplex}
\end{definition}
\begin{theorem}[\texttt{realFourierMultiplierCLM\_even\_reflection\_comm}]\label{GaussianField-realFourierMultiplierCLM-even-reflection-comm}
  \lean{GaussianField.realFourierMultiplierCLM_even_reflection_comm}\leanok
  \uses{GaussianField-fourierMultiplierCLM-even-reflection-comm, GaussianField-realFourierMultiplierCLM, GaussianField-schwartzToComplex, GaussianField-schwartzReflection, GaussianField-schwartzToReal}
  \texttt{GaussianField.realFourierMultiplierCLM\_even\_reflection\_comm}
\end{theorem}
\begin{theorem}[\texttt{resolventMultiplierCLM\_reflection\_comm}]\label{GaussianField-resolventMultiplierCLM-reflection-comm}
  \lean{GaussianField.resolventMultiplierCLM_reflection_comm}\leanok
  \uses{GaussianField-resolventSymbol-even, GaussianField-resolventSymbol-hasTemperateGrowth, GaussianField-realFourierMultiplierCLM-even-reflection-comm, GaussianField-resolventMultiplierCLM, GaussianField-resolventSymbol, GaussianField-schwartzReflection}
  \texttt{GaussianField.resolventMultiplierCLM\_reflection\_comm}
\end{theorem}
\begin{theorem}[\texttt{realFourierMultiplierCLM\_translation\_comm}]\label{GaussianField-realFourierMultiplierCLM-translation-comm}
  \lean{GaussianField.realFourierMultiplierCLM_translation_comm}\leanok
  \uses{GaussianField-fourierMultiplierCLM-translation-comm, GaussianField-realFourierMultiplierCLM, GaussianField-schwartzToComplex, GaussianField-schwartzTranslation, GaussianField-schwartzToReal}
  \texttt{GaussianField.realFourierMultiplierCLM\_translation\_comm}
\end{theorem}
\begin{definition}[\texttt{inverseResolventMultiplierCLM}]\label{GaussianField-inverseResolventMultiplierCLM}
  \lean{GaussianField.inverseResolventMultiplierCLM}\leanok
  \uses{GaussianField-inverseResolventSymbol-hasTemperateGrowth, GaussianField-inverseResolventSymbol, GaussianField-realFourierMultiplierCLM}
  \texttt{GaussianField.inverseResolventMultiplierCLM}
\end{definition}
\begin{theorem}[\texttt{resolventMultiplierCLM\_injective}]\label{GaussianField-resolventMultiplierCLM-injective}
  \lean{GaussianField.resolventMultiplierCLM_injective}\leanok
  \uses{GaussianField-resolventSymbol-even, GaussianField-resolventSymbol-pos, GaussianField-resolventSymbol-hasTemperateGrowth, GaussianField-fourierMultiplier-preserves-real, GaussianField-schwartzToComplex, GaussianField-resolventMultiplierCLM, GaussianField-resolventSymbol, GaussianField-schwartzToReal}
  \texttt{GaussianField.resolventMultiplierCLM\_injective}
\end{theorem}
\begin{theorem}[\texttt{resolventSymbol\_pos}]\label{GaussianField-resolventSymbol-pos}
  \lean{GaussianField.resolventSymbol_pos}\leanok
  \uses{GaussianField-resolventSymbol, Lattice-toSemilatticeInf}
  \texttt{GaussianField.resolventSymbol\_pos}
\end{theorem}
\begin{definition}[\texttt{schwartzToReal}]\label{GaussianField-schwartzToReal}
  \lean{GaussianField.schwartzToReal}\leanok
  \texttt{GaussianField.schwartzToReal}
\end{definition}
\begin{theorem}[\texttt{congr\_simp}]\label{GaussianField-realFourierMultiplierCLM-congr-simp}
  \lean{GaussianField.realFourierMultiplierCLM.congr_simp}\leanok
  \uses{GaussianField-realFourierMultiplierCLM}
  \texttt{GaussianField.realFourierMultiplierCLM.congr\_simp}
\end{theorem}
\begin{theorem}[\texttt{heatMultiplierCLM\_reflection\_comm}]\label{GaussianField-heatMultiplierCLM-reflection-comm}
  \lean{GaussianField.heatMultiplierCLM_reflection_comm}\leanok
  \uses{GaussianField-heatSymbol, GaussianField-heatSymbol-hasTemperateGrowth, GaussianField-heatMultiplierCLM, GaussianField-realFourierMultiplierCLM-even-reflection-comm, GaussianField-schwartzReflection, GaussianField-heatSymbol-even}
  \texttt{GaussianField.heatMultiplierCLM\_reflection\_comm}
\end{theorem}
\begin{theorem}[\texttt{heatMultiplierCLM\_comp}]\label{GaussianField-heatMultiplierCLM-comp}
  \lean{GaussianField.heatMultiplierCLM_comp}\leanok
  \uses{GaussianField-heatSymbol-mul, GaussianField-heatSymbol, GaussianField-heatSymbol-hasTemperateGrowth, GaussianField-heatMultiplierCLM, GaussianField-realFourierMultiplierCLM, GaussianField-realFourierMultiplierCLM-comp, GaussianField-heatSymbol-even}
  \texttt{GaussianField.heatMultiplierCLM\_comp}
\end{theorem}
\begin{definition}[\texttt{inverseResolventSymbol}]\label{GaussianField-inverseResolventSymbol}
  \lean{GaussianField.inverseResolventSymbol}\leanok
  \texttt{GaussianField.inverseResolventSymbol}
\end{definition}
\begin{theorem}[\texttt{fourierMultiplierCLM\_translation\_comm}]\label{GaussianField-fourierMultiplierCLM-translation-comm}
  \lean{GaussianField.fourierMultiplierCLM_translation_comm}\leanok
  \uses{Lattice-toSemilatticeInf}
  \texttt{GaussianField.fourierMultiplierCLM\_translation\_comm}
\end{theorem}
\begin{definition}[\texttt{resolventMultiplierCLE}]\label{GaussianField-resolventMultiplierCLE}
  \lean{GaussianField.resolventMultiplierCLE}\leanok
  \uses{GaussianField-resolventMultiplierCLM, GaussianField-inverseResolventMultiplierCLM}
  \texttt{GaussianField.resolventMultiplierCLE}
\end{definition}
\begin{theorem}[\texttt{fourierMultiplierCLM\_even\_reflection\_comm}]\label{GaussianField-fourierMultiplierCLM-even-reflection-comm}
  \lean{GaussianField.fourierMultiplierCLM_even_reflection_comm}\leanok
  \texttt{GaussianField.fourierMultiplierCLM\_even\_reflection\_comm}
\end{theorem}
\begin{theorem}[\texttt{inverseResolventSymbol\_pos}]\label{GaussianField-inverseResolventSymbol-pos}
  \lean{GaussianField.inverseResolventSymbol_pos}\leanok
  \uses{GaussianField-inverseResolventSymbol, Lattice-toSemilatticeInf}
  \texttt{GaussianField.inverseResolventSymbol\_pos}
\end{theorem}
\begin{theorem}[\texttt{heatSymbol\_zero}]\label{GaussianField-heatSymbol-zero}
  \lean{GaussianField.heatSymbol_zero}\leanok
  \uses{GaussianField-heatSymbol}
  \texttt{GaussianField.heatSymbol\_zero}
\end{theorem}
\begin{theorem}[\texttt{inverseResolventSymbol\_even}]\label{GaussianField-inverseResolventSymbol-even}
  \lean{GaussianField.inverseResolventSymbol_even}\leanok
  \uses{GaussianField-inverseResolventSymbol}
  \texttt{GaussianField.inverseResolventSymbol\_even}
\end{theorem}
\begin{theorem}[\texttt{realFourierMultiplierCLM\_one}]\label{GaussianField-realFourierMultiplierCLM-one}
  \lean{GaussianField.realFourierMultiplierCLM_one}\leanok
  \uses{GaussianField-realFourierMultiplierCLM, GaussianField-schwartzToComplex, GaussianField-schwartzToReal}
  \texttt{GaussianField.realFourierMultiplierCLM\_one}
\end{theorem}
\begin{theorem}[\texttt{resolventSymbol\_even}]\label{GaussianField-resolventSymbol-even}
  \lean{GaussianField.resolventSymbol_even}\leanok
  \uses{GaussianField-resolventSymbol}
  \texttt{GaussianField.resolventSymbol\_even}
\end{theorem}
\begin{theorem}[\texttt{heatSymbol\_hasTemperateGrowth}]\label{GaussianField-heatSymbol-hasTemperateGrowth}
  \lean{GaussianField.heatSymbol_hasTemperateGrowth}\leanok
  \uses{Lattice-toSemilatticeSup, GaussianField-heatSymbol, Lattice-toSemilatticeInf}
  \texttt{GaussianField.heatSymbol\_hasTemperateGrowth}
\end{theorem}
\section{\texttt{Lattice.FKG}}
\begin{theorem}[\texttt{abs\_max\_neg\_le}]\label{GaussianField-abs-max-neg-le}
  \lean{GaussianField.abs_max_neg_le}\leanok
  \uses{Lattice-toSemilatticeInf}
  \texttt{GaussianField.abs\_max\_neg\_le}
\end{theorem}
\begin{theorem}[\texttt{ad\_marginal\_preservation\_ae\_from\_fibers\_lintegral}]\label{GaussianField-ad-marginal-preservation-ae-from-fibers-lintegral}
  \lean{GaussianField.ad_marginal_preservation_ae_from_fibers_lintegral}\leanok
  \uses{GaussianField-ahlswede-daykin-one-dim-ae-lintegral}
  \texttt{GaussianField.ad\_marginal\_preservation\_ae\_from\_fibers\_lintegral}
\end{theorem}
\begin{theorem}[\texttt{integrable\_truncation\_mul}]\label{GaussianField-integrable-truncation-mul}
  \lean{GaussianField.integrable_truncation_mul}\leanok
  \uses{GaussianField-abs-sup-neg-nat-le-abs, Lattice-toSemilatticeInf}
  \texttt{GaussianField.integrable\_truncation\_mul}
\end{theorem}
\begin{theorem}[\texttt{integrable\_fiber\_ae}]\label{GaussianField-integrable-fiber-ae}
  \lean{GaussianField.integrable_fiber_ae}\leanok
  \texttt{GaussianField.integrable\_fiber\_ae}
\end{theorem}
\begin{theorem}[\texttt{ahlswede\_daykin\_one\_dim\_ae}]\label{GaussianField-ahlswede-daykin-one-dim-ae}
  \lean{GaussianField.ahlswede_daykin_one_dim_ae}\leanok
  \uses{GaussianField-algebraic-four-functions}
  \texttt{GaussianField.ahlswede\_daykin\_one\_dim\_ae}
\end{theorem}
\begin{theorem}[\texttt{fkg\_lattice\_condition\_mul}]\label{GaussianField-fkg-lattice-condition-mul}
  \lean{GaussianField.fkg_lattice_condition_mul}\leanok
  \uses{GaussianField-FKGLatticeCondition}
  \texttt{GaussianField.fkg\_lattice\_condition\_mul}
\end{theorem}
\begin{theorem}[\texttt{algebraic\_four\_functions\_ennreal}]\label{GaussianField-algebraic-four-functions-ennreal}
  \lean{GaussianField.algebraic_four_functions_ennreal}\leanok
  \texttt{GaussianField.algebraic\_four\_functions\_ennreal}
\end{theorem}
\begin{definition}[\texttt{FKGLatticeCondition}]\label{GaussianField-FKGLatticeCondition}
  \lean{GaussianField.FKGLatticeCondition}\leanok
  \texttt{GaussianField.FKGLatticeCondition}
\end{definition}
\begin{theorem}[\texttt{inf\_dite\_eq}]\label{GaussianField-inf-dite-eq}
  \lean{GaussianField.inf_dite_eq}\leanok
  \texttt{GaussianField.inf\_dite\_eq}
\end{theorem}
\begin{theorem}[\texttt{gaussianDensity\_integrable}]\label{GaussianField-gaussianDensity-integrable}
  \lean{GaussianField.gaussianDensity_integrable}\leanok
  \uses{GaussianField-FinLatticeSites, GaussianField-massOperator, GaussianField-gaussianDensity-nonneg, GaussianField-finiteLaplacian-neg-semidefinite, GaussianField-gaussianDensity, GaussianField-finiteLaplacian, GaussianField-FinLatticeField}
  \texttt{GaussianField.gaussianDensity\_integrable}
\end{theorem}
\begin{theorem}[\texttt{fkg\_from\_lattice\_condition\_nonneg}]\label{GaussianField-fkg-from-lattice-condition-nonneg}
  \lean{GaussianField.fkg_from_lattice_condition_nonneg}\leanok
  \uses{GaussianField-FKGLatticeCondition, GaussianField-ahlswede-daykin}
  \texttt{GaussianField.fkg\_from\_lattice\_condition\_nonneg}
\end{theorem}
\begin{theorem}[\texttt{integral\_empty\_pi}]\label{GaussianField-integral-empty-pi}
  \lean{GaussianField.integral_empty_pi}\leanok
  \texttt{GaussianField.integral\_empty\_pi}
\end{theorem}
\begin{theorem}[\texttt{sup\_dite\_eq}]\label{GaussianField-sup-dite-eq}
  \lean{GaussianField.sup_dite_eq}\leanok
  \texttt{GaussianField.sup\_dite\_eq}
\end{theorem}
\begin{theorem}[\texttt{ad\_marginal\_preservation\_lintegral}]\label{GaussianField-ad-marginal-preservation-lintegral}
  \lean{GaussianField.ad_marginal_preservation_lintegral}\leanok
  \uses{GaussianField-ahlswede-daykin-one-dim-ae-lintegral, GaussianField-inf-dite-eq, GaussianField-sup-dite-eq}
  \texttt{GaussianField.ad\_marginal\_preservation\_lintegral}
\end{theorem}
\begin{theorem}[\texttt{ad\_marginal\_preservation\_ennreal}]\label{GaussianField-ad-marginal-preservation-ennreal}
  \lean{GaussianField.ad_marginal_preservation_ennreal}\leanok
  \uses{GaussianField-inf-dite-eq, GaussianField-sup-dite-eq, GaussianField-ahlswede-daykin-one-dim-ennreal}
  \texttt{GaussianField.ad\_marginal\_preservation\_ennreal}
\end{theorem}
\begin{theorem}[\texttt{monotone\_max\_neg}]\label{GaussianField-monotone-max-neg}
  \lean{GaussianField.monotone_max_neg}\leanok
  \texttt{GaussianField.monotone\_max\_neg}
\end{theorem}
\begin{theorem}[\texttt{abs\_sup\_neg\_nat\_le\_abs}]\label{GaussianField-abs-sup-neg-nat-le-abs}
  \lean{GaussianField.abs_sup_neg_nat_le_abs}\leanok
  \texttt{GaussianField.abs\_sup\_neg\_nat\_le\_abs}
\end{theorem}
\begin{theorem}[\texttt{ahlswede\_daykin\_one\_dim\_ae\_lintegral}]\label{GaussianField-ahlswede-daykin-one-dim-ae-lintegral}
  \lean{GaussianField.ahlswede_daykin_one_dim_ae_lintegral}\leanok
  \uses{GaussianField-algebraic-four-functions}
  \texttt{GaussianField.ahlswede\_daykin\_one\_dim\_ae\_lintegral}
\end{theorem}
\begin{theorem}[\texttt{ahlswede\_daykin}]\label{GaussianField-ahlswede-daykin}
  \lean{GaussianField.ahlswede_daykin}\leanok
  \uses{GaussianField-ahlswede-daykin-ennreal}
  \texttt{GaussianField.ahlswede\_daykin}
\end{theorem}
\begin{theorem}[\texttt{quadraticForm\_submodular\_of\_nonpos\_offDiag}]\label{GaussianField-quadraticForm-submodular-of-nonpos-offDiag}
  \lean{GaussianField.quadraticForm_submodular_of_nonpos_offDiag}\leanok
  \uses{GaussianField-sup-inf-mul-add-le, GaussianField-IsSubmodular}
  \texttt{GaussianField.quadraticForm\_submodular\_of\_nonpos\_offDiag}
\end{theorem}
\begin{theorem}[\texttt{ad\_marginal\_preservation\_ae\_from\_fibers}]\label{GaussianField-ad-marginal-preservation-ae-from-fibers}
  \lean{GaussianField.ad_marginal_preservation_ae_from_fibers}\leanok
  \uses{GaussianField-ahlswede-daykin-one-dim-ae}
  \texttt{GaussianField.ad\_marginal\_preservation\_ae\_from\_fibers}
\end{theorem}
\begin{definition}[\texttt{IsSingleSite}]\label{GaussianField-IsSingleSite}
  \lean{GaussianField.IsSingleSite}\leanok
  \texttt{GaussianField.IsSingleSite}
\end{definition}
\begin{theorem}[\texttt{ahlswede\_daykin\_one\_dim\_ennreal}]\label{GaussianField-ahlswede-daykin-one-dim-ennreal}
  \lean{GaussianField.ahlswede_daykin_one_dim_ennreal}\leanok
  \uses{GaussianField-ennreal-cancel-two, GaussianField-algebraic-four-functions-ennreal}
  \texttt{GaussianField.ahlswede\_daykin\_one\_dim\_ennreal}
\end{theorem}
\begin{theorem}[\texttt{algebraic\_four\_functions}]\label{GaussianField-algebraic-four-functions}
  \lean{GaussianField.algebraic_four_functions}\leanok
  \texttt{GaussianField.algebraic\_four\_functions}
\end{theorem}
\begin{theorem}[\texttt{integrable\_marginal}]\label{GaussianField-integrable-marginal}
  \lean{GaussianField.integrable_marginal}\leanok
  \texttt{GaussianField.integrable\_marginal}
\end{theorem}
\begin{theorem}[\texttt{fubini\_pi\_decomp}]\label{GaussianField-fubini-pi-decomp}
  \lean{GaussianField.fubini_pi_decomp}\leanok
  \texttt{GaussianField.fubini\_pi\_decomp}
\end{theorem}
\begin{theorem}[\texttt{fkg\_lattice\_condition\_of\_submodular}]\label{GaussianField-fkg-lattice-condition-of-submodular}
  \lean{GaussianField.fkg_lattice_condition_of_submodular}\leanok
  \uses{GaussianField-FKGLatticeCondition, GaussianField-IsSubmodular}
  \texttt{GaussianField.fkg\_lattice\_condition\_of\_submodular}
\end{theorem}
\begin{theorem}[\texttt{measurable\_marginal\_lintegral}]\label{GaussianField-measurable-marginal-lintegral}
  \lean{GaussianField.measurable_marginal_lintegral}\leanok
  \texttt{GaussianField.measurable\_marginal\_lintegral}
\end{theorem}
\begin{theorem}[\texttt{fubini\_pi\_decomp\_lintegral}]\label{GaussianField-fubini-pi-decomp-lintegral}
  \lean{GaussianField.fubini_pi_decomp_lintegral}\leanok
  \texttt{GaussianField.fubini\_pi\_decomp\_lintegral}
\end{theorem}
\begin{theorem}[\texttt{fkg\_lattice\_condition\_single\_site}]\label{GaussianField-fkg-lattice-condition-single-site}
  \lean{GaussianField.fkg_lattice_condition_single_site}\leanok
  \uses{GaussianField-single-coord-sup-inf-sum, GaussianField-IsSingleSite, GaussianField-FKGLatticeCondition, Lattice-toSemilatticeInf}
  \texttt{GaussianField.fkg\_lattice\_condition\_single\_site}
\end{theorem}
\begin{theorem}[\texttt{chebyshev\_integral\_inequality}]\label{GaussianField-chebyshev-integral-inequality}
  \lean{GaussianField.chebyshev_integral_inequality}\leanok
  \texttt{GaussianField.chebyshev\_integral\_inequality}
\end{theorem}
\begin{theorem}[\texttt{massOperator\_offDiag\_nonpos}]\label{GaussianField-massOperator-offDiag-nonpos}
  \lean{GaussianField.massOperator_offDiag_nonpos}\leanok
  \uses{GaussianField-finLatticeDelta, GaussianField-FinLatticeSites, GaussianField-finiteLaplacianFun, GaussianField-finiteLaplacian, GaussianField-FinLatticeField, GaussianField-massOperatorEntry, Lattice-toSemilatticeInf}
  \texttt{GaussianField.massOperator\_offDiag\_nonpos}
\end{theorem}
\begin{definition}[\texttt{insertDecompEquiv}]\label{GaussianField-insertDecompEquiv}
  \lean{GaussianField.insertDecompEquiv}\leanok
  \texttt{GaussianField.insertDecompEquiv}
\end{definition}
\begin{definition}[\texttt{IsFieldMonotone}]\label{GaussianField-IsFieldMonotone}
  \lean{GaussianField.IsFieldMonotone}\leanok
  \uses{GaussianField-finLatticeDelta, GaussianField-FinLatticeSites, GaussianField-Configuration, GaussianField-FinLatticeField}
  \texttt{GaussianField.IsFieldMonotone}
\end{definition}
\begin{theorem}[\texttt{ahlswede\_daykin\_one\_dim}]\label{GaussianField-ahlswede-daykin-one-dim}
  \lean{GaussianField.ahlswede_daykin_one_dim}\leanok
  \uses{GaussianField-algebraic-four-functions}
  \texttt{GaussianField.ahlswede\_daykin\_one\_dim}
\end{theorem}
\begin{definition}[\texttt{IsSubmodular}]\label{GaussianField-IsSubmodular}
  \lean{GaussianField.IsSubmodular}\leanok
  \texttt{GaussianField.IsSubmodular}
\end{definition}
\begin{theorem}[\texttt{ahlswede\_daykin\_ennreal}]\label{GaussianField-ahlswede-daykin-ennreal}
  \lean{GaussianField.ahlswede_daykin_ennreal}\leanok
  \uses{GaussianField-ad-marginal-preservation-ennreal, GaussianField-measurable-marginal-lintegral, GaussianField-fubini-pi-decomp-lintegral}
  \texttt{GaussianField.ahlswede\_daykin\_ennreal}
\end{theorem}
\begin{theorem}[\texttt{single\_coord\_sup\_inf\_sum}]\label{GaussianField-single-coord-sup-inf-sum}
  \lean{GaussianField.single_coord_sup_inf_sum}\leanok
  \texttt{GaussianField.single\_coord\_sup\_inf\_sum}
\end{theorem}
\begin{theorem}[\texttt{fkg\_truncation\_dct\_prod}]\label{GaussianField-fkg-truncation-dct-prod}
  \lean{GaussianField.fkg_truncation_dct_prod}\leanok
  \uses{GaussianField-abs-sup-neg-nat-le-abs, Lattice-toSemilatticeInf}
  \texttt{GaussianField.fkg\_truncation\_dct\_prod}
\end{theorem}
\begin{theorem}[\texttt{fkg\_truncation\_dct}]\label{GaussianField-fkg-truncation-dct}
  \lean{GaussianField.fkg_truncation_dct}\leanok
  \uses{GaussianField-abs-sup-neg-nat-le-abs, Lattice-toSemilatticeInf}
  \texttt{GaussianField.fkg\_truncation\_dct}
\end{theorem}
\begin{theorem}[\texttt{gaussian\_fkg\_lattice\_condition}]\label{GaussianField-gaussian-fkg-lattice-condition}
  \lean{GaussianField.gaussian_fkg_lattice_condition}\leanok
  \uses{GaussianField-gaussianDensity-integrable, GaussianField-gaussianDensity-integral-pos, GaussianField-finLatticeDelta, GaussianField-FinLatticeSites, GaussianField-gaussianDensity-fkg-lattice-condition, GaussianField-gaussianDensity-measurable, GaussianField-integrable-mul-gaussianDensity-of-comp-eval, GaussianField-latticeGaussianMeasure, GaussianField-IsFieldMonotone, GaussianField-Configuration, GaussianField-gaussianDensity-nonneg, GaussianField-fkg-from-lattice-condition, GaussianField-instMeasurableSpaceConfiguration, GaussianField-gaussianDensity, GaussianField-latticeGaussianMeasure-density-integral-of-fieldLaw, GaussianField-FinLatticeField}
  \texttt{GaussianField.gaussian\_fkg\_lattice\_condition}
\end{theorem}
\begin{theorem}[\texttt{fkg\_from\_lattice\_condition}]\label{GaussianField-fkg-from-lattice-condition}
  \lean{GaussianField.fkg_from_lattice_condition}\leanok
  \uses{GaussianField-fkg-truncation-dct, GaussianField-fkg-truncation-dct-prod, GaussianField-integrable-truncation-prod-mul, GaussianField-FKGLatticeCondition, GaussianField-integrable-truncation-mul, GaussianField-fkg-from-lattice-condition-nonneg}
  \texttt{GaussianField.fkg\_from\_lattice\_condition}
\end{theorem}
\begin{theorem}[\texttt{fkg\_perturbed}]\label{GaussianField-fkg-perturbed}
  \lean{GaussianField.fkg_perturbed}\leanok
  \uses{GaussianField-fkg-lattice-condition-mul, GaussianField-gaussianDensity-integral-pos, GaussianField-finLatticeDelta, GaussianField-FinLatticeSites, GaussianField-gaussianDensity-fkg-lattice-condition, GaussianField-IsSingleSite, GaussianField-gaussianDensity-measurable, GaussianField-integrable-mul-gaussianDensity-of-comp-eval, GaussianField-latticeGaussianMeasure, GaussianField-fkg-lattice-condition-single-site, GaussianField-IsFieldMonotone, GaussianField-Configuration, GaussianField-gaussianDensity-nonneg, GaussianField-FKGLatticeCondition, GaussianField-fkg-from-lattice-condition, GaussianField-instMeasurableSpaceConfiguration, GaussianField-gaussianDensity, GaussianField-latticeGaussianMeasure-density-integral-of-fieldLaw, GaussianField-FinLatticeField, Lattice-toSemilatticeInf}
  \texttt{GaussianField.fkg\_perturbed}
\end{theorem}
\begin{theorem}[\texttt{gaussianDensity\_fkg\_lattice\_condition}]\label{GaussianField-gaussianDensity-fkg-lattice-condition}
  \lean{GaussianField.gaussianDensity_fkg_lattice_condition}\leanok
  \uses{GaussianField-massOperator-offDiag-nonpos, GaussianField-FinLatticeSites, GaussianField-quadraticForm-submodular-of-nonpos-offDiag, GaussianField-massOperator, GaussianField-FKGLatticeCondition, GaussianField-gaussianDensity, GaussianField-FinLatticeField, GaussianField-massOperatorEntry}
  \texttt{GaussianField.gaussianDensity\_fkg\_lattice\_condition}
\end{theorem}
\begin{theorem}[\texttt{fkg\_lattice\_gaussian}]\label{GaussianField-fkg-lattice-gaussian}
  \lean{GaussianField.fkg_lattice_gaussian}\leanok
  \uses{GaussianField-gaussian-fkg-lattice-condition, GaussianField-FinLatticeSites, GaussianField-latticeGaussianMeasure, GaussianField-IsFieldMonotone, GaussianField-Configuration, GaussianField-instMeasurableSpaceConfiguration, GaussianField-FinLatticeField}
  \texttt{GaussianField.fkg\_lattice\_gaussian}
\end{theorem}
\begin{theorem}[\texttt{ennreal\_cancel\_two}]\label{GaussianField-ennreal-cancel-two}
  \lean{GaussianField.ennreal_cancel_two}\leanok
  \texttt{GaussianField.ennreal\_cancel\_two}
\end{theorem}
\begin{theorem}[\texttt{integrable\_truncation\_prod\_mul}]\label{GaussianField-integrable-truncation-prod-mul}
  \lean{GaussianField.integrable_truncation_prod_mul}\leanok
  \uses{GaussianField-abs-sup-neg-nat-le-abs, Lattice-toSemilatticeInf}
  \texttt{GaussianField.integrable\_truncation\_prod\_mul}
\end{theorem}
\begin{theorem}[\texttt{sup\_inf\_mul\_add\_le}]\label{GaussianField-sup-inf-mul-add-le}
  \lean{GaussianField.sup_inf_mul_add_le}\leanok
  \texttt{GaussianField.sup\_inf\_mul\_add\_le}
\end{theorem}
\section{\texttt{SmoothCircle.Laplacian}}
\begin{definition}[\texttt{derivSCCLM}]\label{GaussianField-derivSCCLM}
  \lean{GaussianField.derivSCCLM}\leanok
  \uses{GaussianField-SmoothMap-Circle-instAddCommGroup, GaussianField-derivSC-continuous, GaussianField-SmoothMap-Circle-instModule, GaussianField-SmoothMap-Circle, GaussianField-derivSCLM, GaussianField-SmoothMap-Circle-instTopologicalSpace}
  \texttt{GaussianField.derivSCCLM}
\end{definition}
\begin{theorem}[\texttt{circleLaplacian\_fourierBasis}]\label{GaussianField-circleLaplacian-fourierBasis}
  \lean{GaussianField.circleLaplacian_fourierBasis}\leanok
  \uses{GaussianField-SmoothMap-Circle-instFunLike, GaussianField-SmoothMap-Circle-ext, GaussianField-circleLaplacian-apply, GaussianField-SmoothMap-Circle-fourierBasis, GaussianField-circleHasLaplacianEigenvalues, GaussianField-SmoothMap-Circle-instAddCommGroup, GaussianField-SmoothMap-Circle-fourierBasisFun, GaussianField-SmoothMap-Circle-instSMul, GaussianField-HasLaplacianEigenvalues-eigenvalue, GaussianField-SmoothMap-Circle-instModule, GaussianField-circleLaplacian, GaussianField-SmoothMap-Circle-instContinuousSMul, GaussianField-smoothCircle-dyninMityaginSpace, GaussianField-SmoothMap-Circle-instIsTopologicalAddGroup, GaussianField-SmoothMap-Circle, GaussianField-SmoothMap-Circle-instTopologicalSpace, GaussianField-SmoothMap-Circle-fourierFreq}
  \texttt{GaussianField.circleLaplacian\_fourierBasis}
\end{theorem}
\begin{definition}[\texttt{derivSC}]\label{GaussianField-derivSC}
  \lean{GaussianField.derivSC}\leanok
  \uses{GaussianField-SmoothMap-Circle-instFunLike, GaussianField-SmoothMap-Circle}
  \texttt{GaussianField.derivSC}
\end{definition}
\begin{theorem}[\texttt{derivSC\_apply}]\label{GaussianField-derivSC-apply}
  \lean{GaussianField.derivSC_apply}\leanok
  \uses{GaussianField-SmoothMap-Circle-instFunLike, GaussianField-derivSC, GaussianField-SmoothMap-Circle}
  \texttt{GaussianField.derivSC\_apply}
\end{theorem}
\begin{theorem}[\texttt{circleLaplacian\_apply}]\label{GaussianField-circleLaplacian-apply}
  \lean{GaussianField.circleLaplacian_apply}\leanok
  \uses{GaussianField-SmoothMap-Circle-instFunLike, GaussianField-SmoothMap-Circle-instAddCommGroup, GaussianField-SmoothMap-Circle-instModule, GaussianField-circleLaplacian, GaussianField-SmoothMap-Circle, GaussianField-SmoothMap-Circle-instTopologicalSpace}
  \texttt{GaussianField.circleLaplacian\_apply}
\end{theorem}
\begin{theorem}[\texttt{derivSC\_smul}]\label{GaussianField-derivSC-smul}
  \lean{GaussianField.derivSC_smul}\leanok
  \uses{GaussianField-SmoothMap-Circle-instFunLike, GaussianField-SmoothMap-Circle-ext, GaussianField-SmoothMap-Circle-instSMul, GaussianField-derivSC, GaussianField-SmoothMap-Circle-contDiffAt-of-smooth, GaussianField-SmoothMap-Circle}
  \texttt{GaussianField.derivSC\_smul}
\end{theorem}
\begin{definition}[\texttt{circleLaplacian}]\label{GaussianField-circleLaplacian}
  \lean{GaussianField.circleLaplacian}\leanok
  \uses{GaussianField-derivSCCLM, GaussianField-SmoothMap-Circle-instAddCommGroup, GaussianField-SmoothMap-Circle-instModule, GaussianField-SmoothMap-Circle-instIsTopologicalAddGroup, GaussianField-SmoothMap-Circle, GaussianField-SmoothMap-Circle-instTopologicalSpace}
  \texttt{GaussianField.circleLaplacian}
\end{definition}
\begin{theorem}[\texttt{derivSC\_add}]\label{GaussianField-derivSC-add}
  \lean{GaussianField.derivSC_add}\leanok
  \uses{GaussianField-SmoothMap-Circle-instFunLike, GaussianField-SmoothMap-Circle-ext, GaussianField-SmoothMap-Circle-instAdd, GaussianField-derivSC, GaussianField-SmoothMap-Circle-contDiffAt-of-smooth, GaussianField-SmoothMap-Circle}
  \texttt{GaussianField.derivSC\_add}
\end{theorem}
\begin{definition}[\texttt{derivSCLM}]\label{GaussianField-derivSCLM}
  \lean{GaussianField.derivSCLM}\leanok
  \uses{GaussianField-SmoothMap-Circle-instAddCommGroup, GaussianField-derivSC, GaussianField-SmoothMap-Circle-instModule, GaussianField-SmoothMap-Circle, GaussianField-derivSC-add, GaussianField-derivSC-smul}
  \texttt{GaussianField.derivSCLM}
\end{definition}
\begin{theorem}[\texttt{derivSC\_continuous}]\label{GaussianField-derivSC-continuous}
  \lean{GaussianField.derivSC_continuous}\leanok
  \uses{GaussianField-SmoothMap-Circle-instFunLike, GaussianField-SmoothMap-Circle-smoothCircle-withSeminorms, GaussianField-SmoothMap-Circle-norm-iteratedDeriv-le-sobolevSeminorm, GaussianField-SmoothMap-Circle-Icc-nonempty, GaussianField-SmoothMap-Circle-instAddCommGroup, GaussianField-SmoothMap-Circle-instSMul, GaussianField-derivSC, GaussianField-SmoothMap-Circle-instModule, GaussianField-SmoothMap-Circle, GaussianField-derivSCLM, GaussianField-SmoothMap-Circle-sobolevSeminorm, GaussianField-derivSC-add, GaussianField-SmoothMap-Circle-instTopologicalSpace, Lattice-toSemilatticeInf, GaussianField-derivSC-smul}
  \texttt{GaussianField.derivSC\_continuous}
\end{theorem}
\section{\texttt{GaussianField.Hypercontractive}}
\begin{theorem}[\texttt{integral\_pow4\_gaussianReal}]\label{GaussianField-integral-pow4-gaussianReal}
  \lean{GaussianField.integral_pow4_gaussianReal}\leanok
  \uses{GaussianField-fourth-moment-standard-gaussian, Lattice-toSemilatticeInf}
  \texttt{GaussianField.integral\_pow4\_gaussianReal}
\end{theorem}
\begin{theorem}[\texttt{gross\_log\_sobolev}]\label{GaussianField-gross-log-sobolev}
  \lean{GaussianField.gross_log_sobolev}\leanok
  \uses{GaussianField-DyninMityaginSpace, GaussianField-log-sobolev-1d, GaussianField-measure, GaussianField-second-moment-eq-covariance, GaussianField-pairing-memLp, GaussianField-pairing-is-gaussian, GaussianField-configuration-eval-measurable, GaussianField-Configuration, GaussianField-instMeasurableSpaceConfiguration, Lattice-toSemilatticeInf}
  \texttt{GaussianField.gross\_log\_sobolev}
\end{theorem}
\begin{theorem}[\texttt{integral\_sq\_gaussianReal}]\label{GaussianField-integral-sq-gaussianReal}
  \lean{GaussianField.integral_sq_gaussianReal}\leanok
  \uses{Lattice-toSemilatticeInf}
  \texttt{GaussianField.integral\_sq\_gaussianReal}
\end{theorem}
\begin{theorem}[\texttt{sq\_log\_div\_le}]\label{GaussianField-sq-log-div-le}
  \lean{GaussianField.sq_log_div_le}\leanok
  \uses{Lattice-toSemilatticeInf}
  \texttt{GaussianField.sq\_log\_div\_le}
\end{theorem}
\begin{theorem}[\texttt{fourth\_moment\_standard\_gaussian}]\label{GaussianField-fourth-moment-standard-gaussian}
  \lean{GaussianField.fourth_moment_standard_gaussian}\leanok
  \texttt{GaussianField.fourth\_moment\_standard\_gaussian}
\end{theorem}
\begin{theorem}[\texttt{log\_sobolev\_1d}]\label{GaussianField-log-sobolev-1d}
  \lean{GaussianField.log_sobolev_1d}\leanok
  \uses{GaussianField-integral-pow4-gaussianReal, GaussianField-integral-sq-gaussianReal, GaussianField-sq-log-div-le, Lattice-toSemilatticeInf}
  \texttt{GaussianField.log\_sobolev\_1d}
\end{theorem}
\begin{theorem}[\texttt{gaussian\_abs\_moment\_std}]\label{GaussianField-gaussian-abs-moment-std}
  \lean{GaussianField.gaussian_abs_moment_std}\leanok
  \texttt{GaussianField.gaussian\_abs\_moment\_std}
\end{theorem}
\section{\texttt{Lattice.RapidDecayLattice}}
\begin{definition}[\texttt{instAddCommGroup}]\label{GaussianField-RapidDecayLattice-instAddCommGroup}
  \lean{GaussianField.RapidDecayLattice.instAddCommGroup}\leanok
  \uses{GaussianField-RapidDecayLattice-instNeg, GaussianField-RapidDecayLattice-instZero, GaussianField-RapidDecayLattice, GaussianField-RapidDecayLattice-instAdd}
  \texttt{GaussianField.RapidDecayLattice.instAddCommGroup}
\end{definition}
\begin{definition}[\texttt{val}]\label{GaussianField-RapidDecayLattice-val}
  \lean{GaussianField.RapidDecayLattice.val}\leanok
  \uses{GaussianField-RapidDecayLattice}
  \texttt{GaussianField.RapidDecayLattice.val}
\end{definition}
\begin{theorem}[\texttt{latticeEnum\_index\_bound}]\label{GaussianField-RapidDecayLattice-latticeEnum-index-bound}
  \lean{GaussianField.RapidDecayLattice.latticeEnum_index_bound}\leanok
  \uses{GaussianField-RapidDecayLattice-latticeEnum, GaussianField-latticeNorm}
  \texttt{GaussianField.RapidDecayLattice.latticeEnum\_index\_bound}
\end{theorem}
\begin{definition}[\texttt{rapidDecayLattice\_hasPointEval}]\label{GaussianField-rapidDecayLattice-hasPointEval}
  \lean{GaussianField.rapidDecayLattice_hasPointEval}\leanok
  \uses{GaussianField-RapidDecayLattice, GaussianField-RapidDecayLattice-val, GaussianField-HasPointEval}
  \texttt{GaussianField.rapidDecayLattice\_hasPointEval}
\end{definition}
\begin{definition}[\texttt{lattice\_dyninMityaginSpace}]\label{GaussianField-RapidDecayLattice-lattice-dyninMityaginSpace}
  \lean{GaussianField.RapidDecayLattice.lattice_dyninMityaginSpace}\leanok
  \uses{GaussianField-DyninMityaginSpace, GaussianField-RapidDecayLattice-latticeRapidDecaySeminorm, GaussianField-DyninMityaginSpace-ofRapidDecayEquiv, GaussianField-RapidDecayLattice-instAddCommGroup, GaussianField-RapidDecayLattice-latticeRapidDecayEquiv, GaussianField-RapidDecayLattice, GaussianField-RapidDecayLattice-lattice-withSeminorms, GaussianField-RapidDecayLattice-instTopologicalSpace, GaussianField-RapidDecayLattice-instModule, GaussianField-RapidDecayLattice-instIsTopologicalAddGroup, GaussianField-RapidDecayLattice-instContinuousSMul}
  \texttt{GaussianField.RapidDecayLattice.lattice\_dyninMityaginSpace}
\end{definition}
\begin{definition}[\texttt{instSMul}]\label{GaussianField-RapidDecayLattice-instSMul}
  \lean{GaussianField.RapidDecayLattice.instSMul}\leanok
  \uses{GaussianField-RapidDecayLattice, GaussianField-RapidDecayLattice-val}
  \texttt{GaussianField.RapidDecayLattice.instSMul}
\end{definition}
\begin{theorem}[\texttt{instContinuousSMul}]\label{GaussianField-RapidDecayLattice-instContinuousSMul}
  \lean{GaussianField.RapidDecayLattice.instContinuousSMul}\leanok
  \uses{GaussianField-RapidDecayLattice-latticeRapidDecaySeminorm, GaussianField-RapidDecayLattice-instAddCommGroup, GaussianField-RapidDecayLattice, GaussianField-RapidDecayLattice-lattice-withSeminorms, GaussianField-RapidDecayLattice-instTopologicalSpace, GaussianField-RapidDecayLattice-instSMul, GaussianField-RapidDecayLattice-instModule}
  \texttt{GaussianField.RapidDecayLattice.instContinuousSMul}
\end{theorem}
\begin{theorem}[\texttt{smul\_val}]\label{GaussianField-RapidDecayLattice-smul-val}
  \lean{GaussianField.RapidDecayLattice.smul_val}\leanok
  \uses{GaussianField-RapidDecayLattice, GaussianField-RapidDecayLattice-val, GaussianField-RapidDecayLattice-instSMul}
  \texttt{GaussianField.RapidDecayLattice.smul\_val}
\end{theorem}
\begin{definition}[\texttt{latticeCoeffCLM}]\label{GaussianField-RapidDecayLattice-latticeCoeffCLM}
  \lean{GaussianField.RapidDecayLattice.latticeCoeffCLM}\leanok
  \uses{GaussianField-RapidDecayLattice-instAddCommGroup, GaussianField-RapidDecayLattice-latticeEnum, GaussianField-RapidDecayLattice, GaussianField-RapidDecayLattice-val, GaussianField-RapidDecayLattice-instTopologicalSpace, GaussianField-RapidDecayLattice-instModule}
  \texttt{GaussianField.RapidDecayLattice.latticeCoeffCLM}
\end{definition}
\begin{theorem}[\texttt{ext\_iff}]\label{GaussianField-RapidDecayLattice-ext-iff}
  \lean{GaussianField.RapidDecayLattice.ext_iff}\leanok
  \uses{GaussianField-RapidDecayLattice, GaussianField-RapidDecayLattice-val, GaussianField-RapidDecayLattice-ext}
  \texttt{GaussianField.RapidDecayLattice.ext\_iff}
\end{theorem}
\begin{definition}[\texttt{instZero}]\label{GaussianField-RapidDecayLattice-instZero}
  \lean{GaussianField.RapidDecayLattice.instZero}\leanok
  \uses{GaussianField-RapidDecayLattice}
  \texttt{GaussianField.RapidDecayLattice.instZero}
\end{definition}
\begin{theorem}[\texttt{congr\_simp}]\label{GaussianField-RapidDecayLattice-mk-congr-simp}
  \lean{GaussianField.RapidDecayLattice.mk.congr_simp}\leanok
  \uses{GaussianField-RapidDecayLattice, GaussianField-latticeNorm}
  \texttt{GaussianField.RapidDecayLattice.mk.congr\_simp}
\end{theorem}
\begin{theorem}[\texttt{zero\_val}]\label{GaussianField-RapidDecayLattice-zero-val}
  \lean{GaussianField.RapidDecayLattice.zero_val}\leanok
  \uses{GaussianField-RapidDecayLattice-instZero, GaussianField-RapidDecayLattice, GaussianField-RapidDecayLattice-val}
  \texttt{GaussianField.RapidDecayLattice.zero\_val}
\end{theorem}
\begin{definition}[\texttt{latticeRapidDecayEquiv}]\label{GaussianField-RapidDecayLattice-latticeRapidDecayEquiv}
  \lean{GaussianField.RapidDecayLattice.latticeRapidDecayEquiv}\leanok
  \uses{GaussianField-RapidDecayLattice-latticeRapidDecaySeminorm, GaussianField-RapidDecaySeq-rapidDecaySeminorm, GaussianField-RapidDecayLattice-instAddCommGroup, GaussianField-RapidDecaySeq, GaussianField-RapidDecayLattice-latticeEnum, GaussianField-RapidDecaySeq-val, GaussianField-RapidDecayLattice, GaussianField-RapidDecaySeq-instModule, GaussianField-RapidDecayLattice-val, GaussianField-RapidDecayLattice-instTopologicalSpace, GaussianField-RapidDecayLattice-instModule, GaussianField-RapidDecaySeq-instAddCommGroup, GaussianField-RapidDecaySeq-instTopologicalSpace}
  \texttt{GaussianField.RapidDecayLattice.latticeRapidDecayEquiv}
\end{definition}
\begin{theorem}[\texttt{latticeEnum\_norm\_bound}]\label{GaussianField-RapidDecayLattice-latticeEnum-norm-bound}
  \lean{GaussianField.RapidDecayLattice.latticeEnum_norm_bound}\leanok
  \uses{GaussianField-RapidDecayLattice-latticeEnum, GaussianField-latticeNorm}
  \texttt{GaussianField.RapidDecayLattice.latticeEnum\_norm\_bound}
\end{theorem}
\begin{theorem}[\texttt{weight\_pos}]\label{GaussianField-RapidDecayLattice-weight-pos}
  \lean{GaussianField.RapidDecayLattice.weight_pos}\leanok
  \uses{GaussianField-latticeNorm, GaussianField-latticeNorm-nonneg}
  \texttt{GaussianField.RapidDecayLattice.weight\_pos}
\end{theorem}
\begin{theorem}[\texttt{lattice\_withSeminorms}]\label{GaussianField-RapidDecayLattice-lattice-withSeminorms}
  \lean{GaussianField.RapidDecayLattice.lattice_withSeminorms}\leanok
  \uses{GaussianField-RapidDecayLattice-latticeRapidDecaySeminorm, GaussianField-RapidDecayLattice-instAddCommGroup, GaussianField-RapidDecayLattice, GaussianField-RapidDecayLattice-instTopologicalSpace, GaussianField-RapidDecayLattice-instModule}
  \texttt{GaussianField.RapidDecayLattice.lattice\_withSeminorms}
\end{theorem}
\begin{definition}[\texttt{instAdd}]\label{GaussianField-RapidDecayLattice-instAdd}
  \lean{GaussianField.RapidDecayLattice.instAdd}\leanok
  \uses{GaussianField-RapidDecayLattice, GaussianField-RapidDecayLattice-val}
  \texttt{GaussianField.RapidDecayLattice.instAdd}
\end{definition}
\begin{definition}[\texttt{RapidDecayLattice}]\label{GaussianField-RapidDecayLattice}
  \lean{GaussianField.RapidDecayLattice}\leanok
  \texttt{GaussianField.RapidDecayLattice}
\end{definition}
\begin{definition}[\texttt{latticeRapidDecaySeminorm}]\label{GaussianField-RapidDecayLattice-latticeRapidDecaySeminorm}
  \lean{GaussianField.RapidDecayLattice.latticeRapidDecaySeminorm}\leanok
  \uses{GaussianField-RapidDecayLattice-instAddCommGroup, GaussianField-RapidDecayLattice, GaussianField-latticeNorm, GaussianField-RapidDecayLattice-val, GaussianField-RapidDecayLattice-instSMul}
  \texttt{GaussianField.RapidDecayLattice.latticeRapidDecaySeminorm}
\end{definition}
\begin{definition}[\texttt{latticeEnum}]\label{GaussianField-RapidDecayLattice-latticeEnum}
  \lean{GaussianField.RapidDecayLattice.latticeEnum}\leanok
  \texttt{GaussianField.RapidDecayLattice.latticeEnum}
\end{definition}
\begin{definition}[\texttt{instTopologicalSpace}]\label{GaussianField-RapidDecayLattice-instTopologicalSpace}
  \lean{GaussianField.RapidDecayLattice.instTopologicalSpace}\leanok
  \uses{GaussianField-RapidDecayLattice-latticeRapidDecaySeminorm, GaussianField-RapidDecayLattice-instAddCommGroup, GaussianField-RapidDecayLattice, GaussianField-RapidDecayLattice-instModule}
  \texttt{GaussianField.RapidDecayLattice.instTopologicalSpace}
\end{definition}
\begin{theorem}[\texttt{neg\_val}]\label{GaussianField-RapidDecayLattice-neg-val}
  \lean{GaussianField.RapidDecayLattice.neg_val}\leanok
  \uses{GaussianField-RapidDecayLattice-instNeg, GaussianField-RapidDecayLattice, GaussianField-RapidDecayLattice-val}
  \texttt{GaussianField.RapidDecayLattice.neg\_val}
\end{theorem}
\begin{definition}[\texttt{instNeg}]\label{GaussianField-RapidDecayLattice-instNeg}
  \lean{GaussianField.RapidDecayLattice.instNeg}\leanok
  \uses{GaussianField-RapidDecayLattice, GaussianField-RapidDecayLattice-val}
  \texttt{GaussianField.RapidDecayLattice.instNeg}
\end{definition}
\begin{theorem}[\texttt{ext}]\label{GaussianField-RapidDecayLattice-ext}
  \lean{GaussianField.RapidDecayLattice.ext}\leanok
  \uses{GaussianField-RapidDecayLattice, GaussianField-latticeNorm, GaussianField-RapidDecayLattice-val}
  \texttt{GaussianField.RapidDecayLattice.ext}
\end{theorem}
\begin{definition}[\texttt{instModule}]\label{GaussianField-RapidDecayLattice-instModule}
  \lean{GaussianField.RapidDecayLattice.instModule}\leanok
  \uses{GaussianField-RapidDecayLattice-instAddCommGroup, GaussianField-RapidDecayLattice, GaussianField-RapidDecayLattice-instSMul}
  \texttt{GaussianField.RapidDecayLattice.instModule}
\end{definition}
\begin{theorem}[\texttt{instIsTopologicalAddGroup}]\label{GaussianField-RapidDecayLattice-instIsTopologicalAddGroup}
  \lean{GaussianField.RapidDecayLattice.instIsTopologicalAddGroup}\leanok
  \uses{GaussianField-RapidDecayLattice-latticeRapidDecaySeminorm, GaussianField-RapidDecayLattice-instAddCommGroup, GaussianField-RapidDecayLattice, GaussianField-RapidDecayLattice-lattice-withSeminorms, GaussianField-RapidDecayLattice-instTopologicalSpace, GaussianField-RapidDecayLattice-instModule}
  \texttt{GaussianField.RapidDecayLattice.instIsTopologicalAddGroup}
\end{theorem}
\begin{definition}[\texttt{latticeBasisAt}]\label{GaussianField-RapidDecayLattice-latticeBasisAt}
  \lean{GaussianField.RapidDecayLattice.latticeBasisAt}\leanok
  \uses{GaussianField-RapidDecayLattice}
  \texttt{GaussianField.RapidDecayLattice.latticeBasisAt}
\end{definition}
\begin{theorem}[\texttt{rapid\_decay}]\label{GaussianField-RapidDecayLattice-rapid-decay}
  \lean{GaussianField.RapidDecayLattice.rapid_decay}\leanok
  \uses{GaussianField-RapidDecayLattice, GaussianField-latticeNorm, GaussianField-RapidDecayLattice-val}
  \texttt{GaussianField.RapidDecayLattice.rapid\_decay}
\end{theorem}
\begin{theorem}[\texttt{weight\_nonneg}]\label{GaussianField-RapidDecayLattice-weight-nonneg}
  \lean{GaussianField.RapidDecayLattice.weight_nonneg}\leanok
  \uses{GaussianField-latticeNorm, GaussianField-RapidDecayLattice-weight-pos}
  \texttt{GaussianField.RapidDecayLattice.weight\_nonneg}
\end{theorem}
\begin{theorem}[\texttt{congr\_simp}]\label{GaussianField-RapidDecayLattice-latticeEnum-congr-simp}
  \lean{GaussianField.RapidDecayLattice.latticeEnum.congr_simp}\leanok
  \uses{GaussianField-RapidDecayLattice-latticeEnum}
  \texttt{GaussianField.RapidDecayLattice.latticeEnum.congr\_simp}
\end{theorem}
\begin{theorem}[\texttt{add\_val}]\label{GaussianField-RapidDecayLattice-add-val}
  \lean{GaussianField.RapidDecayLattice.add_val}\leanok
  \uses{GaussianField-RapidDecayLattice, GaussianField-RapidDecayLattice-val, GaussianField-RapidDecayLattice-instAdd}
  \texttt{GaussianField.RapidDecayLattice.add\_val}
\end{theorem}
\begin{definition}[\texttt{latticeBasisVec}]\label{GaussianField-RapidDecayLattice-latticeBasisVec}
  \lean{GaussianField.RapidDecayLattice.latticeBasisVec}\leanok
  \uses{GaussianField-RapidDecayLattice-latticeEnum, GaussianField-RapidDecayLattice, GaussianField-RapidDecayLattice-latticeBasisAt}
  \texttt{GaussianField.RapidDecayLattice.latticeBasisVec}
\end{definition}
\begin{definition}[\texttt{ctorIdx}]\label{GaussianField-RapidDecayLattice-ctorIdx}
  \lean{GaussianField.RapidDecayLattice.ctorIdx}\leanok
  \uses{GaussianField-RapidDecayLattice}
  \texttt{GaussianField.RapidDecayLattice.ctorIdx}
\end{definition}
\section{\texttt{GaussianField.ConfigurationEmbedding}}
\begin{theorem}[\texttt{configBasisEval\_continuous}]\label{GaussianField-configBasisEval-continuous}
  \lean{GaussianField.configBasisEval_continuous}\leanok
  \uses{GaussianField-DyninMityaginSpace, GaussianField-Configuration, GaussianField-DyninMityaginSpace-basis, GaussianField-configBasisEval}
  \texttt{GaussianField.configBasisEval\_continuous}
\end{theorem}
\begin{theorem}[\texttt{prokhorov\_sequential}]\label{GaussianField-prokhorov-sequential}
  \lean{GaussianField.prokhorov_sequential}\leanok
  On a Polish space $X$, if a sequence of probability measures $\{\mu_n\}$ is tight, then it has a weakly convergent subsequence. Fully proved from Mathlib.
\end{theorem}
\begin{theorem}[\texttt{configBasisEval\_injective}]\label{GaussianField-configBasisEval-injective}
  \lean{GaussianField.configBasisEval_injective}\leanok
  \uses{GaussianField-DyninMityaginSpace, GaussianField-DyninMityaginSpace-coeff, GaussianField-DyninMityaginSpace-expansion, GaussianField-Configuration, GaussianField-DyninMityaginSpace-basis, GaussianField-configBasisEval}
  \texttt{GaussianField.configBasisEval\_injective}
\end{theorem}
\begin{theorem}[\texttt{configBasisEval\_measurable}]\label{GaussianField-configBasisEval-measurable}
  \lean{GaussianField.configBasisEval_measurable}\leanok
  \uses{GaussianField-DyninMityaginSpace, GaussianField-configuration-eval-measurable, GaussianField-Configuration, GaussianField-DyninMityaginSpace-basis, GaussianField-instMeasurableSpaceConfiguration, GaussianField-configBasisEval}
  \texttt{GaussianField.configBasisEval\_measurable}
\end{theorem}
\begin{definition}[\texttt{configBasisEval}]\label{GaussianField-configBasisEval}
  \lean{GaussianField.configBasisEval}\leanok
  \uses{GaussianField-DyninMityaginSpace, GaussianField-Configuration, GaussianField-DyninMityaginSpace-basis}
  \texttt{GaussianField.configBasisEval}
\end{definition}
\begin{theorem}[\texttt{instMeasurableSpaceConfiguration\_eq\_comap}]\label{GaussianField-instMeasurableSpaceConfiguration-eq-comap}
  \lean{GaussianField.instMeasurableSpaceConfiguration_eq_comap}\leanok
  \uses{GaussianField-DyninMityaginSpace, GaussianField-DyninMityaginSpace-coeff, GaussianField-DyninMityaginSpace-hasSum-basis, GaussianField-configBasisEval-measurable, GaussianField-Configuration, GaussianField-DyninMityaginSpace-basis, GaussianField-instMeasurableSpaceConfiguration, GaussianField-configBasisEval}
  \texttt{GaussianField.instMeasurableSpaceConfiguration\_eq\_comap}
\end{theorem}
\begin{theorem}[\texttt{prokhorov\_configuration\_pushforward}]\label{GaussianField-prokhorov-configuration-pushforward}
  \lean{GaussianField.prokhorov_configuration_pushforward}\leanok
  \uses{GaussianField-DyninMityaginSpace, GaussianField-prokhorov-sequential, GaussianField-configBasisEval-measurable, GaussianField-Configuration, GaussianField-instMeasurableSpaceConfiguration, GaussianField-configBasisEval, GaussianField-configBasisEval-continuous}
  \texttt{GaussianField.prokhorov\_configuration\_pushforward}
\end{theorem}
\begin{theorem}[\texttt{prokhorov\_configuration}]\label{GaussianField-prokhorov-configuration}
  \lean{GaussianField.prokhorov_configuration}\leanok
  \uses{GaussianField-DyninMityaginSpace, Lattice-toSemilatticeSup, GaussianField-DyninMityaginSpace-coeff, GaussianField-prokhorov-configuration-pushforward, GaussianField-DyninMityaginSpace-hasSum-basis, GaussianField-configBasisEval-injective, GaussianField-configBasisEval-measurable, GaussianField-Configuration, GaussianField-DyninMityaginSpace-basis, GaussianField-instMeasurableSpaceConfiguration, GaussianField-instMeasurableSpaceConfiguration-eq-comap, GaussianField-configBasisEval, GaussianField-configBasisEval-continuous, Lattice-toSemilatticeInf}
  \texttt{GaussianField.prokhorov\_configuration}
\end{theorem}
\section{\texttt{GaussianField.NuclearSVD}}
\begin{definition}[\texttt{ell2'}]\label{GaussianField-ell2-}
  \lean{GaussianField.ell2'}\leanok
  Lean signature ``\texttt{lean abbrev ell2' := lp (fun \_ : ℕ => ℝ) 2 }``  Informal: The $\ell^2$ space $\ell^2(\mathbb{N}, \mathbb{R})$ used as the domain of the nuclear operator.
\end{definition}
\begin{theorem}[\texttt{gram\_isCompact'}]\label{GaussianField-gram-isCompact-}
  \lean{GaussianField.gram_isCompact'}\leanok
  \texttt{GaussianField.gram\_isCompact'}
\end{theorem}
\begin{theorem}[\texttt{ell2\_basis\_orthonormal}]\label{GaussianField-ell2-basis-orthonormal}
  \lean{GaussianField.ell2_basis_orthonormal}\leanok
  \uses{GaussianField-ell2-basis, GaussianField-ell2-}
  \texttt{GaussianField.ell2\_basis\_orthonormal}
\end{theorem}
\begin{theorem}[\texttt{congr\_simp}]\label{GaussianField-nuclear-clm-congr-simp}
  \lean{GaussianField.nuclear_clm.congr_simp}\leanok
  \uses{GaussianField-nuclear-clm, GaussianField-ell2-}
  \texttt{GaussianField.nuclear\_clm.congr\_simp}
\end{theorem}
\begin{theorem}[\texttt{smulRight\_isCompactOperator}]\label{GaussianField-smulRight-isCompactOperator}
  \lean{GaussianField.smulRight_isCompactOperator}\leanok
  \uses{GaussianField-ell2-}
  \texttt{GaussianField.smulRight\_isCompactOperator}
\end{theorem}
\begin{theorem}[\texttt{gram\_isSelfAdjoint'}]\label{GaussianField-gram-isSelfAdjoint-}
  \lean{GaussianField.gram_isSelfAdjoint'}\leanok
  \texttt{GaussianField.gram\_isSelfAdjoint'}
\end{theorem}
\begin{definition}[\texttt{nuclear\_linearMap}]\label{GaussianField-nuclear-linearMap}
  \lean{GaussianField.nuclear_linearMap}\leanok
  \uses{GaussianField-ell2-}
  Lean signature ``\texttt{lean def nuclear\_linearMap (y : ℕ → K) (hy : Summable (fun m => ‖y m‖)) : ell2' →ₗ[ℝ] K }``  Informal: The linear map underlying the nuclear operator: $A(x) = \sum_m x_m y_m$.
\end{definition}
\begin{theorem}[\texttt{summable\_smul\_of\_ell2}]\label{GaussianField-summable-smul-of-ell2}
  \lean{GaussianField.summable_smul_of_ell2}\leanok
  \uses{GaussianField-summable-norm-sq-of-summable-norm, GaussianField-ell2-, Lattice-toSemilatticeInf}
  \texttt{GaussianField.summable\_smul\_of\_ell2}
\end{theorem}
\begin{theorem}[\texttt{orthonormal\_countable'}]\label{GaussianField-orthonormal-countable-}
  \lean{GaussianField.orthonormal_countable'}\leanok
  \uses{Lattice-toSemilatticeInf}
  \texttt{GaussianField.orthonormal\_countable'}
\end{theorem}
\begin{theorem}[\texttt{nuclear\_sequence\_svd}]\label{GaussianField-nuclear-sequence-svd}
  \lean{GaussianField.nuclear_sequence_svd}\leanok
  \uses{GaussianField-compact-selfAdjoint-spectral-nat, GaussianField-nuclear-clm, GaussianField-nuclear-clm-basis, GaussianField-ell2-basis, GaussianField-summable-sqrt-eigenvalues, GaussianField-ell2-, GaussianField-nuclear-clm-isCompact, GaussianField-inner-ell2-basis-eq-coord, Lattice-toSemilatticeInf}
  \texttt{GaussianField.nuclear\_sequence\_svd}
\end{theorem}
\begin{theorem}[\texttt{ell2\_separableSpace}]\label{GaussianField-ell2-separableSpace}
  \lean{GaussianField.ell2_separableSpace}\leanok
  \uses{GaussianField-ell2-basis, GaussianField-ell2-}
  \texttt{GaussianField.ell2\_separableSpace}
\end{theorem}
\begin{theorem}[\texttt{finset\_sum\_isCompactOperator}]\label{GaussianField-finset-sum-isCompactOperator}
  \lean{GaussianField.finset_sum_isCompactOperator}\leanok
  \uses{GaussianField-ell2-}
  \texttt{GaussianField.finset\_sum\_isCompactOperator}
\end{theorem}
\begin{theorem}[\texttt{nuclear\_clm\_isCompact}]\label{GaussianField-nuclear-clm-isCompact}
  \lean{GaussianField.nuclear_clm_isCompact}\leanok
  \uses{GaussianField-rank1-norm-le, GaussianField-finset-sum-isCompactOperator, GaussianField-nuclear-clm, GaussianField-ell2-basis, GaussianField-ell2-, GaussianField-smulRight-isCompactOperator, GaussianField-inner-ell2-basis-eq-coord}
  \texttt{GaussianField.nuclear\_clm\_isCompact}
\end{theorem}
\begin{theorem}[\texttt{hilbertBasis\_fintype\_finiteDimensional}]\label{GaussianField-hilbertBasis-fintype-finiteDimensional}
  \lean{GaussianField.hilbertBasis_fintype_finiteDimensional}\leanok
  \texttt{GaussianField.hilbertBasis\_fintype\_finiteDimensional}
\end{theorem}
\begin{theorem}[\texttt{ell2\_not\_finiteDimensional}]\label{GaussianField-ell2-not-finiteDimensional}
  \lean{GaussianField.ell2_not_finiteDimensional}\leanok
  \uses{GaussianField-ell2-basis, GaussianField-ell2-basis-orthonormal, GaussianField-ell2-}
  \texttt{GaussianField.ell2\_not\_finiteDimensional}
\end{theorem}
\begin{definition}[\texttt{ell2\_basis}]\label{GaussianField-ell2-basis}
  \lean{GaussianField.ell2_basis}\leanok
  \uses{GaussianField-ell2-}
  Lean signature ``\texttt{lean def ell2\_basis (m : ℕ) : ell2' }``  Informal: The $m$-th standard basis vector $e_m \in \ell^2$, defined as \texttt{lp.single 2 m 1}.
\end{definition}
\begin{theorem}[\texttt{compact\_selfAdjoint\_spectral\_nat}]\label{GaussianField-compact-selfAdjoint-spectral-nat}
  \lean{GaussianField.compact_selfAdjoint_spectral_nat}\leanok
  \uses{GaussianField-hilbertBasis-fintype-finiteDimensional, GaussianField-compact-selfAdjoint-spectral, GaussianField-orthonormal-countable-}
  \texttt{GaussianField.compact\_selfAdjoint\_spectral\_nat}
\end{theorem}
\begin{theorem}[\texttt{nuclear\_clm\_basis}]\label{GaussianField-nuclear-clm-basis}
  \lean{GaussianField.nuclear_clm_basis}\leanok
  \uses{GaussianField-nuclear-clm, GaussianField-ell2-basis, GaussianField-ell2-}
  \texttt{GaussianField.nuclear\_clm\_basis}
\end{theorem}
\begin{theorem}[\texttt{rank1\_norm\_le}]\label{GaussianField-rank1-norm-le}
  \lean{GaussianField.rank1_norm_le}\leanok
  \uses{GaussianField-ell2-basis, GaussianField-ell2-}
  \texttt{GaussianField.rank1\_norm\_le}
\end{theorem}
\begin{theorem}[\texttt{nuclear\_operator\_bound}]\label{GaussianField-nuclear-operator-bound}
  \lean{GaussianField.nuclear_operator_bound}\leanok
  \uses{GaussianField-nuclear-linearMap, GaussianField-ell2-}
  \texttt{GaussianField.nuclear\_operator\_bound}
\end{theorem}
\begin{theorem}[\texttt{inner\_ell2\_basis\_eq\_coord}]\label{GaussianField-inner-ell2-basis-eq-coord}
  \lean{GaussianField.inner_ell2_basis_eq_coord}\leanok
  \uses{GaussianField-ell2-basis, GaussianField-ell2-}
  \texttt{GaussianField.inner\_ell2\_basis\_eq\_coord}
\end{theorem}
\begin{definition}[\texttt{nuclear\_clm}]\label{GaussianField-nuclear-clm}
  \lean{GaussianField.nuclear_clm}\leanok
  \uses{GaussianField-nuclear-linearMap, GaussianField-nuclear-operator-bound, GaussianField-ell2-}
  Lean signature ``\texttt{lean def nuclear\_clm (y : ℕ → K) (hy : Summable (fun m => ‖y m‖)) : ell2' →L[ℝ] K }``  Informal: The nuclear operator $A : \ell^2 \to K$ as a continuous linear map, with $A(e_m) = y_m$.
\end{definition}
\begin{theorem}[\texttt{summable\_sqrt\_eigenvalues}]\label{GaussianField-summable-sqrt-eigenvalues}
  \lean{GaussianField.summable_sqrt_eigenvalues}\leanok
  \uses{GaussianField-ell2-basis, GaussianField-ell2-, Lattice-toSemilatticeInf}
  \texttt{GaussianField.summable\_sqrt\_eigenvalues}
\end{theorem}
\begin{theorem}[\texttt{summable\_norm\_sq\_of\_summable\_norm}]\label{GaussianField-summable-norm-sq-of-summable-norm}
  \lean{GaussianField.summable_norm_sq_of_summable_norm}\leanok
  \texttt{GaussianField.summable\_norm\_sq\_of\_summable\_norm}
\end{theorem}
\section{\texttt{SmoothCircle.FourierTranslation}}
\begin{theorem}[\texttt{fourierCoeffReal\_circleTranslation\_sin}]\label{GaussianField-fourierCoeffReal-circleTranslation-sin}
  \lean{GaussianField.fourierCoeffReal_circleTranslation_sin}\leanok
  \uses{GaussianField-circleTranslation, GaussianField-SmoothMap-Circle-instFunLike, GaussianField-SmoothMap-Circle-continuous, GaussianField-SmoothMap-Circle-periodic, GaussianField-SmoothMap-Circle-instAddCommGroup, GaussianField-SmoothMap-Circle-fourierBasisFun-smooth, GaussianField-SmoothMap-Circle-fourierCoeffReal, GaussianField-SmoothMap-Circle-fourierBasisFun, GaussianField-SmoothMap-Circle-instModule, GaussianField-SmoothMap-Circle, GaussianField-SmoothMap-Circle-fourierBasisFun-periodic, GaussianField-SmoothMap-Circle-instTopologicalSpace}
  \texttt{GaussianField.fourierCoeffReal\_circleTranslation\_sin}
\end{theorem}
\begin{theorem}[\texttt{fourierCoeffReal\_circleTranslation\_zero}]\label{GaussianField-fourierCoeffReal-circleTranslation-zero}
  \lean{GaussianField.fourierCoeffReal_circleTranslation_zero}\leanok
  \uses{GaussianField-circleTranslation, GaussianField-SmoothMap-Circle-instFunLike, GaussianField-SmoothMap-Circle-continuous, GaussianField-SmoothMap-Circle-periodic, GaussianField-SmoothMap-Circle-instAddCommGroup, GaussianField-SmoothMap-Circle-fourierBasisFun-smooth, GaussianField-SmoothMap-Circle-fourierCoeffReal, GaussianField-SmoothMap-Circle-fourierBasisFun, GaussianField-SmoothMap-Circle-instModule, GaussianField-SmoothMap-Circle, GaussianField-SmoothMap-Circle-fourierBasisFun-periodic, GaussianField-SmoothMap-Circle-instTopologicalSpace}
  \texttt{GaussianField.fourierCoeffReal\_circleTranslation\_zero}
\end{theorem}
\begin{theorem}[\texttt{fourierCoeffReal\_circleReflection\_zero}]\label{GaussianField-fourierCoeffReal-circleReflection-zero}
  \lean{GaussianField.fourierCoeffReal_circleReflection_zero}\leanok
  \uses{GaussianField-SmoothMap-Circle-instFunLike, GaussianField-SmoothMap-Circle-continuous, GaussianField-SmoothMap-Circle-periodic, GaussianField-SmoothMap-Circle-instAddCommGroup, GaussianField-SmoothMap-Circle-fourierBasisFun-smooth, GaussianField-SmoothMap-Circle-fourierCoeffReal, GaussianField-SmoothMap-Circle-fourierBasisFun, GaussianField-circleReflection, GaussianField-SmoothMap-Circle-instModule, GaussianField-SmoothMap-Circle, GaussianField-SmoothMap-Circle-fourierBasisFun-periodic, GaussianField-SmoothMap-Circle-instTopologicalSpace}
  \texttt{GaussianField.fourierCoeffReal\_circleReflection\_zero}
\end{theorem}
\begin{theorem}[\texttt{paired\_coeff\_product\_circleReflection}]\label{GaussianField-paired-coeff-product-circleReflection}
  \lean{GaussianField.paired_coeff_product_circleReflection}\leanok
  \uses{GaussianField-SmoothMap-Circle-instAddCommGroup, GaussianField-SmoothMap-Circle-fourierCoeffReal, GaussianField-circleReflection, GaussianField-SmoothMap-Circle-instModule, GaussianField-SmoothMap-Circle, GaussianField-fourierCoeffReal-circleReflection-cos, GaussianField-fourierCoeffReal-circleReflection-sin, GaussianField-SmoothMap-Circle-instTopologicalSpace}
  \texttt{GaussianField.paired\_coeff\_product\_circleReflection}
\end{theorem}
\begin{theorem}[\texttt{paired\_coeff\_product\_circleTranslation}]\label{GaussianField-paired-coeff-product-circleTranslation}
  \lean{GaussianField.paired_coeff_product_circleTranslation}\leanok
  \uses{GaussianField-circleTranslation, GaussianField-SmoothMap-Circle-instAddCommGroup, GaussianField-SmoothMap-Circle-fourierCoeffReal, GaussianField-fourierCoeffReal-circleTranslation-sin, GaussianField-SmoothMap-Circle-instModule, GaussianField-SmoothMap-Circle, GaussianField-fourierCoeffReal-circleTranslation-cos, GaussianField-SmoothMap-Circle-instTopologicalSpace}
  \texttt{GaussianField.paired\_coeff\_product\_circleTranslation}
\end{theorem}
\begin{theorem}[\texttt{fourierCoeffReal\_circleReflection\_sin}]\label{GaussianField-fourierCoeffReal-circleReflection-sin}
  \lean{GaussianField.fourierCoeffReal_circleReflection_sin}\leanok
  \uses{GaussianField-SmoothMap-Circle-instFunLike, GaussianField-SmoothMap-Circle-continuous, GaussianField-SmoothMap-Circle-periodic, GaussianField-SmoothMap-Circle-instAddCommGroup, GaussianField-SmoothMap-Circle-fourierBasisFun-smooth, GaussianField-SmoothMap-Circle-fourierCoeffReal, GaussianField-SmoothMap-Circle-fourierBasisFun, GaussianField-circleReflection, GaussianField-SmoothMap-Circle-instModule, GaussianField-SmoothMap-Circle, GaussianField-SmoothMap-Circle-fourierBasisFun-periodic, GaussianField-SmoothMap-Circle-instTopologicalSpace}
  \texttt{GaussianField.fourierCoeffReal\_circleReflection\_sin}
\end{theorem}
\begin{theorem}[\texttt{fourierCoeffReal\_circleReflection\_cos}]\label{GaussianField-fourierCoeffReal-circleReflection-cos}
  \lean{GaussianField.fourierCoeffReal_circleReflection_cos}\leanok
  \uses{GaussianField-SmoothMap-Circle-instFunLike, GaussianField-SmoothMap-Circle-continuous, GaussianField-SmoothMap-Circle-periodic, GaussianField-SmoothMap-Circle-instAddCommGroup, GaussianField-SmoothMap-Circle-fourierBasisFun-smooth, GaussianField-SmoothMap-Circle-fourierCoeffReal, GaussianField-SmoothMap-Circle-fourierBasisFun, GaussianField-circleReflection, GaussianField-SmoothMap-Circle-instModule, GaussianField-SmoothMap-Circle, GaussianField-SmoothMap-Circle-fourierBasisFun-periodic, GaussianField-SmoothMap-Circle-instTopologicalSpace}
  \texttt{GaussianField.fourierCoeffReal\_circleReflection\_cos}
\end{theorem}
\begin{theorem}[\texttt{fourierCoeffReal\_circleTranslation\_cos}]\label{GaussianField-fourierCoeffReal-circleTranslation-cos}
  \lean{GaussianField.fourierCoeffReal_circleTranslation_cos}\leanok
  \uses{GaussianField-circleTranslation, GaussianField-SmoothMap-Circle-instFunLike, GaussianField-SmoothMap-Circle-continuous, GaussianField-SmoothMap-Circle-periodic, GaussianField-SmoothMap-Circle-instAddCommGroup, GaussianField-SmoothMap-Circle-fourierBasisFun-smooth, GaussianField-SmoothMap-Circle-fourierCoeffReal, GaussianField-SmoothMap-Circle-fourierBasisFun, GaussianField-SmoothMap-Circle-instModule, GaussianField-SmoothMap-Circle, GaussianField-SmoothMap-Circle-fourierBasisFun-periodic, GaussianField-SmoothMap-Circle-instTopologicalSpace}
  \texttt{GaussianField.fourierCoeffReal\_circleTranslation\_cos}
\end{theorem}
\section{\texttt{Lattice.AsymCovariance}}
\begin{theorem}[\texttt{massOperator\_product\_eigenvector\_asym}]\label{GaussianField-massOperator-product-eigenvector-asym}
  \lean{GaussianField.massOperator_product_eigenvector_asym}\leanok
  \uses{GaussianField-latticeEigenvalue1d, GaussianField-AsymLatticeSites, GaussianField-latticeFourierBasisFun, GaussianField-finiteLaplacianAsym, GaussianField-AsymLatticeField, GaussianField-dft-1d-eigenvalue-pointwise, GaussianField-massOperatorAsym}
  \texttt{GaussianField.massOperator\_product\_eigenvector\_asym}
\end{theorem}
\begin{definition}[\texttt{massOperatorEntryAsym}]\label{GaussianField-massOperatorEntryAsym}
  \lean{GaussianField.massOperatorEntryAsym}\leanok
  \uses{GaussianField-AsymLatticeSites, GaussianField-AsymLatticeField, GaussianField-massOperatorAsym, GaussianField-asymLatticeDelta}
  \texttt{GaussianField.massOperatorEntryAsym}
\end{definition}
\begin{theorem}[\texttt{abstract\_spectral\_eq\_dft\_spectral\_2d\_asym}]\label{GaussianField-abstract-spectral-eq-dft-spectral-2d-asym}
  \lean{GaussianField.abstract_spectral_eq_dft_spectral_2d_asym}\leanok
  \uses{GaussianField-spectralLatticeCovarianceAsym, GaussianField-latticeEigenvalue1d, GaussianField-AsymLatticeSites, GaussianField-latticeFourierBasisFun, GaussianField-covariance, GaussianField-massEigenvectorBasisAsym, GaussianField-massOperatorAsym-eigenCoeff-eq-eigenvalues-mul-eigenCoeff, GaussianField-latticeEigenvalue1d-nonneg, GaussianField-massEigenbasisAsym-sum-mul-sum-eq-site-inner, GaussianField-massEigenvaluesAsym, GaussianField-covariance-spectralLatticeCovarianceAsym-eq, GaussianField-latticeFourierNormSq, GaussianField-massOperatorMatrixAsym-eigenvalues-pos, GaussianField-AsymLatticeField, GaussianField-instFintypeAsymLatticeSitesOfNeZeroNat, GaussianField-dft-parseval-2d-asym, GaussianField-massOperatorAsym-surjective, GaussianField-massOperatorAsym, GaussianField-dft-eigencoeff-massOperatorAsym, GaussianField-ell2-}
  \texttt{GaussianField.abstract\_spectral\_eq\_dft\_spectral\_2d\_asym}
\end{theorem}
\begin{theorem}[\texttt{congr\_simp}]\label{GaussianField-spectralLatticeCovarianceAsym-congr-simp}
  \lean{GaussianField.spectralLatticeCovarianceAsym.congr_simp}\leanok
  \uses{GaussianField-spectralLatticeCovarianceAsym, GaussianField-AsymLatticeSites, GaussianField-AsymLatticeField, GaussianField-ell2-}
  \texttt{GaussianField.spectralLatticeCovarianceAsym.congr\_simp}
\end{theorem}
\begin{theorem}[\texttt{congr\_simp}]\label{GaussianField-massOperatorMatrixAsym-congr-simp}
  \lean{GaussianField.massOperatorMatrixAsym.congr_simp}\leanok
  \uses{GaussianField-AsymLatticeSites, GaussianField-massOperatorMatrixAsym}
  \texttt{GaussianField.massOperatorMatrixAsym.congr\_simp}
\end{theorem}
\begin{definition}[\texttt{finiteLaplacianAsymLM}]\label{GaussianField-finiteLaplacianAsymLM}
  \lean{GaussianField.finiteLaplacianAsymLM}\leanok
  \uses{GaussianField-AsymLatticeSites, GaussianField-finiteLaplacianAsymFun, GaussianField-AsymLatticeField}
  \texttt{GaussianField.finiteLaplacianAsymLM}
\end{definition}
\begin{theorem}[\texttt{dft\_eigencoeff\_massOperatorAsym}]\label{GaussianField-dft-eigencoeff-massOperatorAsym}
  \lean{GaussianField.dft_eigencoeff_massOperatorAsym}\leanok
  \uses{GaussianField-latticeEigenvalue1d, GaussianField-AsymLatticeSites, GaussianField-latticeFourierBasisFun, GaussianField-massOperator-product-eigenvector-asym, GaussianField-AsymLatticeField, GaussianField-instFintypeAsymLatticeSitesOfNeZeroNat, GaussianField-massOperatorAsym, GaussianField-massOperatorAsym-selfAdjoint}
  \texttt{GaussianField.dft\_eigencoeff\_massOperatorAsym}
\end{theorem}
\begin{theorem}[\texttt{massOperatorMatrixAsym\_isHermitian}]\label{GaussianField-massOperatorMatrixAsym-isHermitian}
  \lean{GaussianField.massOperatorMatrixAsym_isHermitian}\leanok
  \uses{GaussianField-AsymLatticeSites, GaussianField-massOperatorMatrixAsym, GaussianField-AsymLatticeField, GaussianField-instFintypeAsymLatticeSitesOfNeZeroNat, GaussianField-massOperatorAsym, GaussianField-asymLatticeDelta, GaussianField-massOperatorAsym-selfAdjoint}
  \texttt{GaussianField.massOperatorMatrixAsym\_isHermitian}
\end{theorem}
\begin{theorem}[\texttt{sum\_factor\_asym}]\label{GaussianField-sum-factor-asym}
  \lean{GaussianField.sum_factor_asym}\leanok
  \uses{GaussianField-AsymLatticeSites, GaussianField-instFintypeAsymLatticeSitesOfNeZeroNat}
  \texttt{GaussianField.sum\_factor\_asym}
\end{theorem}
\begin{theorem}[\texttt{congr\_simp}]\label{GaussianField-finiteLaplacianAsymFun-congr-simp}
  \lean{GaussianField.finiteLaplacianAsymFun.congr_simp}\leanok
  \uses{GaussianField-AsymLatticeSites, GaussianField-finiteLaplacianAsymFun, GaussianField-AsymLatticeField}
  \texttt{GaussianField.finiteLaplacianAsymFun.congr\_simp}
\end{theorem}
\begin{definition}[\texttt{spectralLatticeCovarianceAsym}]\label{GaussianField-spectralLatticeCovarianceAsym}
  \lean{GaussianField.spectralLatticeCovarianceAsym}\leanok
  \uses{GaussianField-AsymLatticeSites, GaussianField-massEigenvectorBasisAsym, GaussianField-massEigenvaluesAsym, GaussianField-AsymLatticeField, GaussianField-instFintypeAsymLatticeSitesOfNeZeroNat, GaussianField-ell2-}
  \texttt{GaussianField.spectralLatticeCovarianceAsym}
\end{definition}
\begin{theorem}[\texttt{congr\_simp}]\label{GaussianField-finiteLaplacianAsym-congr-simp}
  \lean{GaussianField.finiteLaplacianAsym.congr_simp}\leanok
  \uses{GaussianField-AsymLatticeSites, GaussianField-finiteLaplacianAsym, GaussianField-AsymLatticeField}
  \texttt{GaussianField.finiteLaplacianAsym.congr\_simp}
\end{theorem}
\begin{theorem}[\texttt{massOperatorAsym\_eq\_matrix\_mulVec}]\label{GaussianField-massOperatorAsym-eq-matrix-mulVec}
  \lean{GaussianField.massOperatorAsym_eq_matrix_mulVec}\leanok
  \uses{GaussianField-AsymLatticeSites, GaussianField-massOperatorEntryAsym, GaussianField-massOperatorMatrixAsym, GaussianField-AsymLatticeField, GaussianField-instFintypeAsymLatticeSitesOfNeZeroNat, GaussianField-massOperatorAsym, GaussianField-asymLatticeDelta}
  \texttt{GaussianField.massOperatorAsym\_eq\_matrix\_mulVec}
\end{theorem}
\begin{theorem}[\texttt{massOperatorMatrixAsym\_eigenvalues\_pos}]\label{GaussianField-massOperatorMatrixAsym-eigenvalues-pos}
  \lean{GaussianField.massOperatorMatrixAsym_eigenvalues_pos}\leanok
  \uses{GaussianField-AsymLatticeSites, GaussianField-massOperatorMatrixAsym, GaussianField-massEigenvaluesAsym, GaussianField-massOperatorMatrixAsym-isHermitian, GaussianField-AsymLatticeField, GaussianField-massOperatorAsym-pos-def, GaussianField-massMatrixAsymHerm, GaussianField-instFintypeAsymLatticeSitesOfNeZeroNat, GaussianField-massOperatorAsym, GaussianField-massOperatorAsym-eq-matrix-mulVec}
  \texttt{GaussianField.massOperatorMatrixAsym\_eigenvalues\_pos}
\end{theorem}
\begin{theorem}[\texttt{massOperatorAsym\_pos\_def}]\label{GaussianField-massOperatorAsym-pos-def}
  \lean{GaussianField.massOperatorAsym_pos_def}\leanok
  \uses{GaussianField-AsymLatticeSites, GaussianField-finiteLaplacianAsym-neg-semidefinite, GaussianField-finiteLaplacianAsym, GaussianField-AsymLatticeField, GaussianField-instFintypeAsymLatticeSitesOfNeZeroNat, GaussianField-massOperatorAsym, Lattice-toSemilatticeInf}
  \texttt{GaussianField.massOperatorAsym\_pos\_def}
\end{theorem}
\begin{theorem}[\texttt{covariance\_spectralLatticeCovarianceAsym\_eq}]\label{GaussianField-covariance-spectralLatticeCovarianceAsym-eq}
  \lean{GaussianField.covariance_spectralLatticeCovarianceAsym_eq}\leanok
  \uses{GaussianField-spectralLatticeCovarianceAsym, GaussianField-AsymLatticeSites, GaussianField-covariance, GaussianField-massEigenvectorBasisAsym, GaussianField-massEigenvaluesAsym, GaussianField-AsymLatticeField, GaussianField-instFintypeAsymLatticeSitesOfNeZeroNat, GaussianField-spectralLatticeCovarianceAsym-inner, GaussianField-ell2-}
  \texttt{GaussianField.covariance\_spectralLatticeCovarianceAsym\_eq}
\end{theorem}
\begin{definition}[\texttt{asymLatticeDelta}]\label{GaussianField-asymLatticeDelta}
  \lean{GaussianField.asymLatticeDelta}\leanok
  \uses{GaussianField-AsymLatticeSites, GaussianField-AsymLatticeField}
  \texttt{GaussianField.asymLatticeDelta}
\end{definition}
\begin{theorem}[\texttt{congr\_simp}]\label{GaussianField-massOperatorAsym-congr-simp}
  \lean{GaussianField.massOperatorAsym.congr_simp}\leanok
  \uses{GaussianField-AsymLatticeSites, GaussianField-AsymLatticeField, GaussianField-massOperatorAsym}
  \texttt{GaussianField.massOperatorAsym.congr\_simp}
\end{theorem}
\begin{theorem}[\texttt{massOperatorAsym\_selfAdjoint}]\label{GaussianField-massOperatorAsym-selfAdjoint}
  \lean{GaussianField.massOperatorAsym_selfAdjoint}\leanok
  \uses{GaussianField-AsymLatticeSites, GaussianField-finiteLaplacianAsym-selfAdjoint, GaussianField-finiteLaplacianAsym, GaussianField-AsymLatticeField, GaussianField-instFintypeAsymLatticeSitesOfNeZeroNat, GaussianField-massOperatorAsym}
  \texttt{GaussianField.massOperatorAsym\_selfAdjoint}
\end{theorem}
\begin{theorem}[\texttt{dft\_parseval\_2d\_asym}]\label{GaussianField-dft-parseval-2d-asym}
  \lean{GaussianField.dft_parseval_2d_asym}\leanok
  \uses{GaussianField-AsymLatticeSites, GaussianField-latticeFourierBasisFun, GaussianField-dft-parseval-1d, GaussianField-sum-factor-asym, GaussianField-latticeFourierNormSq, GaussianField-AsymLatticeField, GaussianField-instFintypeAsymLatticeSitesOfNeZeroNat}
  \texttt{GaussianField.dft\_parseval\_2d\_asym}
\end{theorem}
\begin{definition}[\texttt{massEigenvaluesAsym}]\label{GaussianField-massEigenvaluesAsym}
  \lean{GaussianField.massEigenvaluesAsym}\leanok
  \uses{GaussianField-AsymLatticeSites, GaussianField-massOperatorMatrixAsym, GaussianField-massMatrixAsymHerm, GaussianField-instFintypeAsymLatticeSitesOfNeZeroNat}
  \texttt{GaussianField.massEigenvaluesAsym}
\end{definition}
\begin{theorem}[\texttt{massEigenbasisAsym\_sum\_mul\_sum\_eq\_site\_inner}]\label{GaussianField-massEigenbasisAsym-sum-mul-sum-eq-site-inner}
  \lean{GaussianField.massEigenbasisAsym_sum_mul_sum_eq_site_inner}\leanok
  \uses{GaussianField-AsymLatticeSites, GaussianField-massEigenvectorBasisAsym, GaussianField-AsymLatticeField, GaussianField-instFintypeAsymLatticeSitesOfNeZeroNat}
  \texttt{GaussianField.massEigenbasisAsym\_sum\_mul\_sum\_eq\_site\_inner}
\end{theorem}
\begin{definition}[\texttt{latticeCovarianceAsymGJ}]\label{GaussianField-latticeCovarianceAsymGJ}
  \lean{GaussianField.latticeCovarianceAsymGJ}\leanok
  \uses{GaussianField-spectralLatticeCovarianceAsym, GaussianField-AsymLatticeSites, GaussianField-AsymLatticeField, GaussianField-ell2-}
  \texttt{GaussianField.latticeCovarianceAsymGJ}
\end{definition}
\begin{definition}[\texttt{finiteLaplacianAsym}]\label{GaussianField-finiteLaplacianAsym}
  \lean{GaussianField.finiteLaplacianAsym}\leanok
  \uses{GaussianField-AsymLatticeSites, GaussianField-finiteLaplacianAsymLM, GaussianField-AsymLatticeField}
  \texttt{GaussianField.finiteLaplacianAsym}
\end{definition}
\begin{theorem}[\texttt{spectralLatticeCovarianceAsym\_inner}]\label{GaussianField-spectralLatticeCovarianceAsym-inner}
  \lean{GaussianField.spectralLatticeCovarianceAsym_inner}\leanok
  \uses{GaussianField-spectralLatticeCovarianceAsym, GaussianField-AsymLatticeSites, GaussianField-massEigenvectorBasisAsym, GaussianField-massEigenvaluesAsym, GaussianField-massOperatorMatrixAsym-eigenvalues-pos, GaussianField-AsymLatticeField, GaussianField-instFintypeAsymLatticeSitesOfNeZeroNat, GaussianField-ell2-, Lattice-toSemilatticeInf}
  \texttt{GaussianField.spectralLatticeCovarianceAsym\_inner}
\end{theorem}
\begin{definition}[\texttt{massMatrixAsymHerm}]\label{GaussianField-massMatrixAsymHerm}
  \lean{GaussianField.massMatrixAsymHerm}\leanok
  \uses{GaussianField-AsymLatticeSites, GaussianField-massOperatorMatrixAsym, GaussianField-massOperatorMatrixAsym-isHermitian}
  \texttt{GaussianField.massMatrixAsymHerm}
\end{definition}
\begin{theorem}[\texttt{massOperatorAsym\_eigenCoeff\_eq\_eigenvalues\_mul\_eigenCoeff}]\label{GaussianField-massOperatorAsym-eigenCoeff-eq-eigenvalues-mul-eigenCoeff}
  \lean{GaussianField.massOperatorAsym_eigenCoeff_eq_eigenvalues_mul_eigenCoeff}\leanok
  \uses{GaussianField-AsymLatticeSites, GaussianField-massEigenvectorBasisAsym, GaussianField-massOperatorMatrixAsym, GaussianField-massEigenvaluesAsym, GaussianField-massOperatorMatrixAsym-isHermitian, GaussianField-AsymLatticeField, GaussianField-instFintypeAsymLatticeSitesOfNeZeroNat, GaussianField-massOperatorAsym, GaussianField-massOperatorAsym-selfAdjoint, GaussianField-massOperatorAsym-eq-matrix-mulVec}
  \texttt{GaussianField.massOperatorAsym\_eigenCoeff\_eq\_eigenvalues\_mul\_eigenCoeff}
\end{theorem}
\begin{definition}[\texttt{massEigenvectorBasisAsym}]\label{GaussianField-massEigenvectorBasisAsym}
  \lean{GaussianField.massEigenvectorBasisAsym}\leanok
  \uses{GaussianField-AsymLatticeSites, GaussianField-massOperatorMatrixAsym, GaussianField-massMatrixAsymHerm, GaussianField-instFintypeAsymLatticeSitesOfNeZeroNat}
  \texttt{GaussianField.massEigenvectorBasisAsym}
\end{definition}
\begin{theorem}[\texttt{congr\_simp}]\label{GaussianField-asymLatticeDelta-congr-simp}
  \lean{GaussianField.asymLatticeDelta.congr_simp}\leanok
  \uses{GaussianField-AsymLatticeSites, GaussianField-asymLatticeDelta}
  \texttt{GaussianField.asymLatticeDelta.congr\_simp}
\end{theorem}
\begin{theorem}[\texttt{finiteLaplacianAsym\_neg\_semidefinite}]\label{GaussianField-finiteLaplacianAsym-neg-semidefinite}
  \lean{GaussianField.finiteLaplacianAsym_neg_semidefinite}\leanok
  \uses{GaussianField-AsymLatticeSites, GaussianField-finiteLaplacianAsym, GaussianField-finiteLaplacianAsymFun, GaussianField-AsymLatticeField, GaussianField-instFintypeAsymLatticeSitesOfNeZeroNat, Lattice-toSemilatticeInf}
  \texttt{GaussianField.finiteLaplacianAsym\_neg\_semidefinite}
\end{theorem}
\begin{definition}[\texttt{massOperatorMatrixAsym}]\label{GaussianField-massOperatorMatrixAsym}
  \lean{GaussianField.massOperatorMatrixAsym}\leanok
  \uses{GaussianField-AsymLatticeSites, GaussianField-massOperatorEntryAsym}
  \texttt{GaussianField.massOperatorMatrixAsym}
\end{definition}
\begin{theorem}[\texttt{finiteLaplacianAsym\_selfAdjoint}]\label{GaussianField-finiteLaplacianAsym-selfAdjoint}
  \lean{GaussianField.finiteLaplacianAsym_selfAdjoint}\leanok
  \uses{GaussianField-AsymLatticeSites, GaussianField-finiteLaplacianAsym, GaussianField-finiteLaplacianAsymFun, GaussianField-AsymLatticeField, GaussianField-instFintypeAsymLatticeSitesOfNeZeroNat}
  \texttt{GaussianField.finiteLaplacianAsym\_selfAdjoint}
\end{theorem}
\begin{definition}[\texttt{massOperatorAsym}]\label{GaussianField-massOperatorAsym}
  \lean{GaussianField.massOperatorAsym}\leanok
  \uses{GaussianField-AsymLatticeSites, GaussianField-finiteLaplacianAsym, GaussianField-AsymLatticeField}
  \texttt{GaussianField.massOperatorAsym}
\end{definition}
\begin{theorem}[\texttt{lattice\_green\_tendsto\_continuum\_asym}]\label{GaussianField-lattice-green-tendsto-continuum-asym}
  \lean{GaussianField.lattice_green_tendsto_continuum_asym}\leanok
  \uses{GaussianField-spectralLatticeCovarianceAsym, GaussianField-NuclearTensorProduct-instModuleReal, GaussianField-NuclearTensorProduct-instTopologicalSpace, GaussianField-DyninMityaginSpace-coeff, GaussianField-AsymLatticeSites, GaussianField-DyninMityaginSpace-coeff-decay, GaussianField-covariance, GaussianField-circleHasLaplacianEigenvalues, GaussianField-NuclearTensorProduct-dyninMityaginSpace, GaussianField-SmoothMap-Circle-instAddCommGroup, GaussianField-NuclearTensorProduct-instIsTopologicalAddGroup, GaussianField-tensorProductHasLaplacianEigenvalues, GaussianField-DyninMityaginSpace-expansion, GaussianField-DyninMityaginSpace--, GaussianField-evalAsymTorusAtSite, GaussianField-NuclearTensorProduct-instAddCommGroup, GaussianField-SmoothMap-Circle-instModule, GaussianField-SmoothMap-Circle-instContinuousSMul, GaussianField-smoothCircle-dyninMityaginSpace, GaussianField-SmoothMap-Circle-instIsTopologicalAddGroup, GaussianField-DyninMityaginSpace-basis, GaussianField-NuclearTensorProduct-instContinuousSMulReal, GaussianField-NuclearTensorProduct-pure, GaussianField-AsymLatticeField, GaussianField-SmoothMap-Circle, GaussianField-DyninMityaginSpace-p, GaussianField-AsymTorusTestFunction, GaussianField-ell2-, GaussianField-SmoothMap-Circle-instTopologicalSpace, GaussianField-greenFunctionBilinear, Lattice-toSemilatticeInf, GaussianField-greenFunctionBilinear-symm}
  \texttt{GaussianField.lattice\_green\_tendsto\_continuum\_asym}
\end{theorem}
\begin{theorem}[\texttt{massOperatorAsym\_surjective}]\label{GaussianField-massOperatorAsym-surjective}
  \lean{GaussianField.massOperatorAsym_surjective}\leanok
  \uses{GaussianField-AsymLatticeSites, GaussianField-AsymLatticeField, GaussianField-massOperatorAsym-pos-def, GaussianField-instFintypeAsymLatticeSitesOfNeZeroNat, GaussianField-massOperatorAsym}
  \texttt{GaussianField.massOperatorAsym\_surjective}
\end{theorem}
\begin{definition}[\texttt{finiteLaplacianAsymFun}]\label{GaussianField-finiteLaplacianAsymFun}
  \lean{GaussianField.finiteLaplacianAsymFun}\leanok
  \uses{GaussianField-AsymLatticeSites, GaussianField-AsymLatticeField}
  \texttt{GaussianField.finiteLaplacianAsymFun}
\end{definition}
\begin{theorem}[\texttt{congr\_simp}]\label{GaussianField-massOperatorEntryAsym-congr-simp}
  \lean{GaussianField.massOperatorEntryAsym.congr_simp}\leanok
  \uses{GaussianField-AsymLatticeSites, GaussianField-massOperatorEntryAsym}
  \texttt{GaussianField.massOperatorEntryAsym.congr\_simp}
\end{theorem}
\begin{theorem}[\texttt{congr\_simp}]\label{GaussianField-massEigenvaluesAsym-congr-simp}
  \lean{GaussianField.massEigenvaluesAsym.congr_simp}\leanok
  \uses{GaussianField-AsymLatticeSites, GaussianField-massEigenvaluesAsym}
  \texttt{GaussianField.massEigenvaluesAsym.congr\_simp}
\end{theorem}
\section{\texttt{Lattice.LatticeFourierProduct}}
\begin{theorem}[\texttt{massOperator\_product\_eigenvector\_family}]\label{GaussianField-massOperator-product-eigenvector-family}
  \lean{GaussianField.massOperator_product_eigenvector_family}\leanok
  \uses{GaussianField-latticeEigenvalue1d, GaussianField-latticeFourierBasisFun, GaussianField-FinLatticeSites, GaussianField-massOperator, GaussianField-dft-1d-eigenvalue-pointwise, GaussianField-latticeFourierProductBasisFun, GaussianField-finiteLaplacian, GaussianField-FinLatticeField}
  \texttt{GaussianField.massOperator\_product\_eigenvector\_family}
\end{theorem}
\begin{theorem}[\texttt{congr\_simp}]\label{GaussianField-latticeFourierProductCoeff-congr-simp}
  \lean{GaussianField.latticeFourierProductCoeff.congr_simp}\leanok
  \uses{GaussianField-FinLatticeSites, GaussianField-latticeFourierProductCoeff}
  \texttt{GaussianField.latticeFourierProductCoeff.congr\_simp}
\end{theorem}
\begin{theorem}[\texttt{abstract\_spectral\_eq\_dft\_spectral\_family}]\label{GaussianField-abstract-spectral-eq-dft-spectral-family}
  \lean{GaussianField.abstract_spectral_eq_dft_spectral_family}\leanok
  \uses{GaussianField-latticeEigenvalue1d, GaussianField-dft-parseval-family, GaussianField-lattice-covariance-eq-spectral, GaussianField-covariance, GaussianField-latticeCovariance, GaussianField-FinLatticeSites, GaussianField-dft-eigencoeff-massOperator-family, GaussianField-massEigenbasis-sum-mul-sum-eq-site-inner, GaussianField-massOperator, GaussianField-latticeEigenvalue1d-nonneg, GaussianField-massEigenvalues, GaussianField-massEigenvectorBasis, GaussianField-massOperator-surjective, GaussianField-latticeFourierProductCoeff, GaussianField-massOperatorMatrix-eigenvalues-pos, GaussianField-ell2-, GaussianField-latticeFourierProductNormSq-pos, GaussianField-FinLatticeField, GaussianField-massOperator-eigenCoeff-eq-eigenvalues-mul-eigenCoeff, GaussianField-latticeFourierProductNormSq}
  \texttt{GaussianField.abstract\_spectral\_eq\_dft\_spectral\_family}
\end{theorem}
\begin{definition}[\texttt{latticeFourierProductNormSq}]\label{GaussianField-latticeFourierProductNormSq}
  \lean{GaussianField.latticeFourierProductNormSq}\leanok
  \uses{GaussianField-latticeFourierNormSq}
  \texttt{GaussianField.latticeFourierProductNormSq}
\end{definition}
\begin{theorem}[\texttt{latticeFourierProductNormSq\_pos}]\label{GaussianField-latticeFourierProductNormSq-pos}
  \lean{GaussianField.latticeFourierProductNormSq_pos}\leanok
  \uses{GaussianField-latticeFourierProductNormSq}
  \texttt{GaussianField.latticeFourierProductNormSq\_pos}
\end{theorem}
\begin{theorem}[\texttt{dft\_eigencoeff\_massOperator\_family}]\label{GaussianField-dft-eigencoeff-massOperator-family}
  \lean{GaussianField.dft_eigencoeff_massOperator_family}\leanok
  \uses{GaussianField-massOperator-selfAdjoint, GaussianField-latticeEigenvalue1d, GaussianField-massOperator-product-eigenvector-family, GaussianField-FinLatticeSites, GaussianField-massOperator, GaussianField-latticeFourierProductBasisFun, GaussianField-latticeFourierProductCoeff, GaussianField-FinLatticeField}
  \texttt{GaussianField.dft\_eigencoeff\_massOperator\_family}
\end{theorem}
\begin{theorem}[\texttt{dft\_parseval\_family}]\label{GaussianField-dft-parseval-family}
  \lean{GaussianField.dft_parseval_family}\leanok
  \uses{GaussianField-FinLatticeSites, GaussianField-latticeFourierProductCoeff, GaussianField-latticeFourierProductNormSq}
  \texttt{GaussianField.dft\_parseval\_family}
\end{theorem}
\begin{definition}[\texttt{latticeFourierProductCoeff}]\label{GaussianField-latticeFourierProductCoeff}
  \lean{GaussianField.latticeFourierProductCoeff}\leanok
  \uses{GaussianField-FinLatticeSites, GaussianField-latticeFourierProductBasisFun}
  \texttt{GaussianField.latticeFourierProductCoeff}
\end{definition}
\begin{theorem}[\texttt{congr\_simp}]\label{GaussianField-latticeFourierProductNormSq-congr-simp}
  \lean{GaussianField.latticeFourierProductNormSq.congr_simp}\leanok
  \uses{GaussianField-latticeFourierProductNormSq}
  \texttt{GaussianField.latticeFourierProductNormSq.congr\_simp}
\end{theorem}
\begin{theorem}[\texttt{latticeFourierProductNormSq\_zero}]\label{GaussianField-latticeFourierProductNormSq-zero}
  \lean{GaussianField.latticeFourierProductNormSq_zero}\leanok
  \uses{GaussianField-latticeFourierNormSq, GaussianField-latticeFourierProductNormSq}
  \texttt{GaussianField.latticeFourierProductNormSq\_zero}
\end{theorem}
\begin{definition}[\texttt{latticeFourierProductBasisFun}]\label{GaussianField-latticeFourierProductBasisFun}
  \lean{GaussianField.latticeFourierProductBasisFun}\leanok
  \uses{GaussianField-latticeFourierBasisFun, GaussianField-FinLatticeSites}
  \texttt{GaussianField.latticeFourierProductBasisFun}
\end{definition}
\begin{theorem}[\texttt{latticeFourierProductNormSq\_succ}]\label{GaussianField-latticeFourierProductNormSq-succ}
  \lean{GaussianField.latticeFourierProductNormSq_succ}\leanok
  \uses{GaussianField-latticeFourierNormSq-congr-simp, GaussianField-latticeFourierNormSq, GaussianField-latticeFourierProductNormSq}
  \texttt{GaussianField.latticeFourierProductNormSq\_succ}
\end{theorem}
\begin{theorem}[\texttt{latticeFourierProductCoeff\_succ}]\label{GaussianField-latticeFourierProductCoeff-succ}
  \lean{GaussianField.latticeFourierProductCoeff_succ}\leanok
  \uses{GaussianField-latticeFourierBasisFun, GaussianField-FinLatticeSites, GaussianField-sum-finLatticeSites-succ, GaussianField-latticeFourierProductBasisFun, GaussianField-latticeFourierProductCoeff}
  \texttt{GaussianField.latticeFourierProductCoeff\_succ}
\end{theorem}
\begin{theorem}[\texttt{massOperator\_surjective}]\label{GaussianField-massOperator-surjective}
  \lean{GaussianField.massOperator_surjective}\leanok
  \uses{GaussianField-FinLatticeSites, GaussianField-massOperator, GaussianField-massOperator-pos-def, GaussianField-FinLatticeField}
  \texttt{GaussianField.massOperator\_surjective}
\end{theorem}
\begin{theorem}[\texttt{latticeFourierProductBasis\_sq\_sum}]\label{GaussianField-latticeFourierProductBasis-sq-sum}
  \lean{GaussianField.latticeFourierProductBasis_sq_sum}\leanok
  \uses{GaussianField-FinLatticeSites, GaussianField-latticeFourierProductBasisFun, GaussianField-latticeFourierProductNormSq}
  \texttt{GaussianField.latticeFourierProductBasis\_sq\_sum}
\end{theorem}
\begin{theorem}[\texttt{sum\_finLatticeSites\_succ}]\label{GaussianField-sum-finLatticeSites-succ}
  \lean{GaussianField.sum_finLatticeSites_succ}\leanok
  \uses{GaussianField-FinLatticeSites}
  \texttt{GaussianField.sum\_finLatticeSites\_succ}
\end{theorem}
\section{\texttt{Lattice.SpectralCovariance}}
\begin{theorem}[\texttt{massOperatorMatrix\_eigenvalues\_pos}]\label{GaussianField-massOperatorMatrix-eigenvalues-pos}
  \lean{GaussianField.massOperatorMatrix_eigenvalues_pos}\leanok
  \uses{GaussianField-massOperatorMatrix-isHermitian, GaussianField-massOperatorMatrix, GaussianField-FinLatticeSites, GaussianField-massOperator, GaussianField-massMatrixHerm, GaussianField-massOperator-eq-matrix-mulVec, GaussianField-massEigenvalues, GaussianField-massOperator-pos-def, GaussianField-FinLatticeField}
  \texttt{GaussianField.massOperatorMatrix\_eigenvalues\_pos}
\end{theorem}
\begin{theorem}[\texttt{massEigenbasis\_linear\_sum\_reprSymm\_ofLp}]\label{GaussianField-massEigenbasis-linear-sum-reprSymm-ofLp}
  \lean{GaussianField.massEigenbasis_linear_sum_reprSymm_ofLp}\leanok
  \uses{GaussianField-massEigenbasis-coeff-reprSymm-ofLp, GaussianField-FinLatticeSites, GaussianField-massEigenvectorBasis, GaussianField-FinLatticeField}
  \texttt{GaussianField.massEigenbasis\_linear\_sum\_reprSymm\_ofLp}
\end{theorem}
\begin{theorem}[\texttt{massEigenbasis\_quadratic\_sum\_reprSymm\_ofLp}]\label{GaussianField-massEigenbasis-quadratic-sum-reprSymm-ofLp}
  \lean{GaussianField.massEigenbasis_quadratic_sum_reprSymm_ofLp}\leanok
  \uses{GaussianField-massEigenbasis-coeff-reprSymm-ofLp, GaussianField-FinLatticeSites, GaussianField-massEigenvalues, GaussianField-massEigenvectorBasis}
  \texttt{GaussianField.massEigenbasis\_quadratic\_sum\_reprSymm\_ofLp}
\end{theorem}
\begin{definition}[\texttt{massOperatorEntry}]\label{GaussianField-massOperatorEntry}
  \lean{GaussianField.massOperatorEntry}\leanok
  \uses{GaussianField-finLatticeDelta, GaussianField-FinLatticeSites, GaussianField-massOperator, GaussianField-FinLatticeField}
  \texttt{GaussianField.massOperatorEntry}
\end{definition}
\begin{theorem}[\texttt{spectralLatticeCovariance\_inner}]\label{GaussianField-spectralLatticeCovariance-inner}
  \lean{GaussianField.spectralLatticeCovariance_inner}\leanok
  \uses{GaussianField-FinLatticeSites, GaussianField-massEigenvalues, GaussianField-spectralLatticeCovariance, GaussianField-massEigenvectorBasis, GaussianField-massOperatorMatrix-eigenvalues-pos, GaussianField-ell2-, GaussianField-FinLatticeField, Lattice-toSemilatticeInf}
  \texttt{GaussianField.spectralLatticeCovariance\_inner}
\end{theorem}
\begin{theorem}[\texttt{spectralLatticeCovariance\_norm\_sq}]\label{GaussianField-spectralLatticeCovariance-norm-sq}
  \lean{GaussianField.spectralLatticeCovariance_norm_sq}\leanok
  \uses{GaussianField-FinLatticeSites, GaussianField-spectralLatticeCovariance-inner, GaussianField-massEigenvalues, GaussianField-spectralLatticeCovariance, GaussianField-massEigenvectorBasis, GaussianField-ell2-, GaussianField-FinLatticeField}
  \texttt{GaussianField.spectralLatticeCovariance\_norm\_sq}
\end{theorem}
\begin{theorem}[\texttt{massOperator\_eigenCoeff\_eq\_eigenvalues\_mul\_eigenCoeff}]\label{GaussianField-massOperator-eigenCoeff-eq-eigenvalues-mul-eigenCoeff}
  \lean{GaussianField.massOperator_eigenCoeff_eq_eigenvalues_mul_eigenCoeff}\leanok
  \uses{GaussianField-massOperatorMatrix-isHermitian, GaussianField-massOperator-selfAdjoint, GaussianField-massOperatorMatrix, GaussianField-FinLatticeSites, GaussianField-massOperator, GaussianField-massOperator-eq-matrix-mulVec, GaussianField-massEigenvalues, GaussianField-massEigenvectorBasis, GaussianField-FinLatticeField}
  \texttt{GaussianField.massOperator\_eigenCoeff\_eq\_eigenvalues\_mul\_eigenCoeff}
\end{theorem}
\begin{definition}[\texttt{massEigenvectorBasis}]\label{GaussianField-massEigenvectorBasis}
  \lean{GaussianField.massEigenvectorBasis}\leanok
  \uses{GaussianField-massOperatorMatrix, GaussianField-FinLatticeSites, GaussianField-massMatrixHerm}
  \texttt{GaussianField.massEigenvectorBasis}
\end{definition}
\begin{definition}[\texttt{gaussianDensity}]\label{GaussianField-gaussianDensity}
  \lean{GaussianField.gaussianDensity}\leanok
  \uses{GaussianField-FinLatticeSites, GaussianField-massOperator, GaussianField-FinLatticeField}
  \texttt{GaussianField.gaussianDensity}
\end{definition}
\begin{definition}[\texttt{spectralLatticeCovariance}]\label{GaussianField-spectralLatticeCovariance}
  \lean{GaussianField.spectralLatticeCovariance}\leanok
  \uses{GaussianField-FinLatticeSites, GaussianField-massEigenvalues, GaussianField-massEigenvectorBasis, GaussianField-ell2-, GaussianField-FinLatticeField}
  \texttt{GaussianField.spectralLatticeCovariance}
\end{definition}
\begin{theorem}[\texttt{massOperatorMatrix\_isHermitian}]\label{GaussianField-massOperatorMatrix-isHermitian}
  \lean{GaussianField.massOperatorMatrix_isHermitian}\leanok
  \uses{GaussianField-massOperator-selfAdjoint, GaussianField-massOperatorMatrix, GaussianField-finLatticeDelta, GaussianField-FinLatticeSites, GaussianField-massOperator, GaussianField-FinLatticeField}
  \texttt{GaussianField.massOperatorMatrix\_isHermitian}
\end{theorem}
\begin{theorem}[\texttt{massOperator\_selfAdjoint}]\label{GaussianField-massOperator-selfAdjoint}
  \lean{GaussianField.massOperator_selfAdjoint}\leanok
  \uses{GaussianField-FinLatticeSites, GaussianField-massOperator, GaussianField-finiteLaplacian-selfAdjoint, GaussianField-finiteLaplacian, GaussianField-FinLatticeField}
  \texttt{GaussianField.massOperator\_selfAdjoint}
\end{theorem}
\begin{theorem}[\texttt{congr\_simp}]\label{GaussianField-massEigenvalues-congr-simp}
  \lean{GaussianField.massEigenvalues.congr_simp}\leanok
  \uses{GaussianField-FinLatticeSites, GaussianField-massEigenvalues}
  \texttt{GaussianField.massEigenvalues.congr\_simp}
\end{theorem}
\begin{definition}[\texttt{massEigenvalues}]\label{GaussianField-massEigenvalues}
  \lean{GaussianField.massEigenvalues}\leanok
  \uses{GaussianField-massOperatorMatrix, GaussianField-FinLatticeSites, GaussianField-massMatrixHerm}
  \texttt{GaussianField.massEigenvalues}
\end{definition}
\begin{definition}[\texttt{massOperatorMatrix}]\label{GaussianField-massOperatorMatrix}
  \lean{GaussianField.massOperatorMatrix}\leanok
  \uses{GaussianField-FinLatticeSites, GaussianField-massOperatorEntry}
  \texttt{GaussianField.massOperatorMatrix}
\end{definition}
\begin{theorem}[\texttt{gaussianDensity\_eq\_exp\_spectral}]\label{GaussianField-gaussianDensity-eq-exp-spectral}
  \lean{GaussianField.gaussianDensity_eq_exp_spectral}\leanok
  \uses{GaussianField-massOperator-quadratic-eq-spectral, GaussianField-FinLatticeSites, GaussianField-massOperator, GaussianField-massEigenvalues, GaussianField-massEigenvectorBasis, GaussianField-gaussianDensity, GaussianField-FinLatticeField}
  \texttt{GaussianField.gaussianDensity\_eq\_exp\_spectral}
\end{theorem}
\begin{theorem}[\texttt{spectral\_eq\_site\_bilinear}]\label{GaussianField-spectral-eq-site-bilinear}
  \lean{GaussianField.spectral_eq_site_bilinear}\leanok
  \uses{GaussianField-FinLatticeSites, GaussianField-massEigenvalues, GaussianField-massEigenvectorBasis, GaussianField-FinLatticeField}
  \texttt{GaussianField.spectral\_eq\_site\_bilinear}
\end{theorem}
\begin{definition}[\texttt{massMatrixHerm}]\label{GaussianField-massMatrixHerm}
  \lean{GaussianField.massMatrixHerm}\leanok
  \uses{GaussianField-massOperatorMatrix-isHermitian, GaussianField-massOperatorMatrix, GaussianField-FinLatticeSites}
  \texttt{GaussianField.massMatrixHerm}
\end{definition}
\begin{theorem}[\texttt{congr\_simp}]\label{GaussianField-massOperatorEntry-congr-simp}
  \lean{GaussianField.massOperatorEntry.congr_simp}\leanok
  \uses{GaussianField-FinLatticeSites, GaussianField-massOperatorEntry}
  \texttt{GaussianField.massOperatorEntry.congr\_simp}
\end{theorem}
\begin{theorem}[\texttt{congr\_simp}]\label{GaussianField-massOperatorMatrix-congr-simp}
  \lean{GaussianField.massOperatorMatrix.congr_simp}\leanok
  \uses{GaussianField-massOperatorMatrix, GaussianField-FinLatticeSites}
  \texttt{GaussianField.massOperatorMatrix.congr\_simp}
\end{theorem}
\begin{theorem}[\texttt{congr\_simp}]\label{GaussianField-massOperator-congr-simp}
  \lean{GaussianField.massOperator.congr_simp}\leanok
  \uses{GaussianField-FinLatticeSites, GaussianField-massOperator, GaussianField-FinLatticeField}
  \texttt{GaussianField.massOperator.congr\_simp}
\end{theorem}
\begin{theorem}[\texttt{massEigenbasis\_quadratic\_sum\_reprSymm}]\label{GaussianField-massEigenbasis-quadratic-sum-reprSymm}
  \lean{GaussianField.massEigenbasis_quadratic_sum_reprSymm}\leanok
  \uses{GaussianField-FinLatticeSites, GaussianField-massEigenvalues, GaussianField-massEigenvectorBasis, GaussianField-massEigenbasis-coeff-reprSymm}
  \texttt{GaussianField.massEigenbasis\_quadratic\_sum\_reprSymm}
\end{theorem}
\begin{theorem}[\texttt{massOperator\_eq\_matrix\_mulVec}]\label{GaussianField-massOperator-eq-matrix-mulVec}
  \lean{GaussianField.massOperator_eq_matrix_mulVec}\leanok
  \uses{GaussianField-massOperatorMatrix, GaussianField-finLatticeDelta, GaussianField-FinLatticeSites, GaussianField-massOperator, GaussianField-FinLatticeField, GaussianField-massOperatorEntry}
  \texttt{GaussianField.massOperator\_eq\_matrix\_mulVec}
\end{theorem}
\begin{theorem}[\texttt{spectral\_singular\_values\_bounded}]\label{GaussianField-spectral-singular-values-bounded}
  \lean{GaussianField.spectral_singular_values_bounded}\leanok
  \uses{GaussianField-FinLatticeSites, GaussianField-massEigenvalues, Lattice-toSemilatticeInf, GaussianField-IsBoundedSeq}
  \texttt{GaussianField.spectral\_singular\_values\_bounded}
\end{theorem}
\begin{theorem}[\texttt{massEigenbasis\_sum\_mul\_sum\_eq\_site\_inner}]\label{GaussianField-massEigenbasis-sum-mul-sum-eq-site-inner}
  \lean{GaussianField.massEigenbasis_sum_mul_sum_eq_site_inner}\leanok
  \uses{GaussianField-FinLatticeSites, GaussianField-massEigenvectorBasis, GaussianField-FinLatticeField}
  \texttt{GaussianField.massEigenbasis\_sum\_mul\_sum\_eq\_site\_inner}
\end{theorem}
\begin{theorem}[\texttt{gaussianDensity\_nonneg}]\label{GaussianField-gaussianDensity-nonneg}
  \lean{GaussianField.gaussianDensity_nonneg}\leanok
  \uses{GaussianField-FinLatticeSites, GaussianField-massOperator, GaussianField-gaussianDensity, GaussianField-FinLatticeField}
  \texttt{GaussianField.gaussianDensity\_nonneg}
\end{theorem}
\begin{theorem}[\texttt{massEigenbasis\_coeff\_reprSymm\_ofLp}]\label{GaussianField-massEigenbasis-coeff-reprSymm-ofLp}
  \lean{GaussianField.massEigenbasis_coeff_reprSymm_ofLp}\leanok
  \uses{GaussianField-FinLatticeSites, GaussianField-massEigenvectorBasis, GaussianField-massEigenbasis-coeff-reprSymm}
  \texttt{GaussianField.massEigenbasis\_coeff\_reprSymm\_ofLp}
\end{theorem}
\begin{theorem}[\texttt{massEigenbasis\_linear\_sum\_reprSymm}]\label{GaussianField-massEigenbasis-linear-sum-reprSymm}
  \lean{GaussianField.massEigenbasis_linear_sum_reprSymm}\leanok
  \uses{GaussianField-FinLatticeSites, GaussianField-massEigenvectorBasis, GaussianField-massEigenbasis-coeff-reprSymm, GaussianField-FinLatticeField}
  \texttt{GaussianField.massEigenbasis\_linear\_sum\_reprSymm}
\end{theorem}
\begin{theorem}[\texttt{massEigenbasis\_coeff\_reprSymm}]\label{GaussianField-massEigenbasis-coeff-reprSymm}
  \lean{GaussianField.massEigenbasis_coeff_reprSymm}\leanok
  \uses{GaussianField-FinLatticeSites, GaussianField-massEigenvectorBasis}
  \texttt{GaussianField.massEigenbasis\_coeff\_reprSymm}
\end{theorem}
\begin{theorem}[\texttt{massOperator\_quadratic\_eq\_spectral}]\label{GaussianField-massOperator-quadratic-eq-spectral}
  \lean{GaussianField.massOperator_quadratic_eq_spectral}\leanok
  \uses{GaussianField-FinLatticeSites, GaussianField-massEigenbasis-sum-mul-sum-eq-site-inner, GaussianField-massOperator, GaussianField-massEigenvalues, GaussianField-massEigenvectorBasis, GaussianField-FinLatticeField, GaussianField-massOperator-eigenCoeff-eq-eigenvalues-mul-eigenCoeff}
  \texttt{GaussianField.massOperator\_quadratic\_eq\_spectral}
\end{theorem}
\section{\texttt{Cylinder.Symmetry}}
\begin{theorem}[\texttt{schwartzTranslation\_zero}]\label{GaussianField-schwartzTranslation-zero}
  \lean{GaussianField.schwartzTranslation_zero}\leanok
  \uses{GaussianField-schwartzTranslation}
  \texttt{GaussianField.schwartzTranslation\_zero}
\end{theorem}
\begin{theorem}[\texttt{schwartzPositiveTimeSubmodule\_isClosed}]\label{GaussianField-schwartzPositiveTimeSubmodule-isClosed}
  \lean{GaussianField.schwartzPositiveTimeSubmodule_isClosed}\leanok
  \uses{GaussianField-schwartzPositiveTimeSubmodule, GaussianField-schwartzEvalCLM}
  \texttt{GaussianField.schwartzPositiveTimeSubmodule\_isClosed}
\end{theorem}
\begin{theorem}[\texttt{schwartzTranslation\_apply}]\label{GaussianField-schwartzTranslation-apply}
  \lean{GaussianField.schwartzTranslation_apply}\leanok
  \uses{GaussianField-schwartzTranslation}
  \texttt{GaussianField.schwartzTranslation\_apply}
\end{theorem}
\begin{definition}[\texttt{schwartzReflection}]\label{GaussianField-schwartzReflection}
  \lean{GaussianField.schwartzReflection}\leanok
  \texttt{GaussianField.schwartzReflection}
\end{definition}
\begin{theorem}[\texttt{schwartzNegativeTimeSubmodule\_isClosed}]\label{GaussianField-schwartzNegativeTimeSubmodule-isClosed}
  \lean{GaussianField.schwartzNegativeTimeSubmodule_isClosed}\leanok
  \uses{GaussianField-schwartzEvalCLM, GaussianField-schwartzNegativeTimeSubmodule}
  \texttt{GaussianField.schwartzNegativeTimeSubmodule\_isClosed}
\end{theorem}
\begin{theorem}[\texttt{cylinderConfigReflection\_involution}]\label{GaussianField-cylinderConfigReflection-involution}
  \lean{GaussianField.cylinderConfigReflection_involution}\leanok
  \uses{GaussianField-NuclearTensorProduct-instModuleReal, GaussianField-NuclearTensorProduct-instTopologicalSpace, GaussianField-CylinderTestFunction, GaussianField-NuclearTensorProduct-instIsTopologicalAddGroup, GaussianField-NuclearTensorProduct-instAddCommGroup, GaussianField-cylinderTimeReflection, GaussianField-cylinderTimeReflection-involution, GaussianField-Configuration, GaussianField-NuclearTensorProduct-instContinuousSMulReal, GaussianField-SmoothMap-Circle, GaussianField-cylinderConfigReflection}
  \texttt{GaussianField.cylinderConfigReflection\_involution}
\end{theorem}
\begin{theorem}[\texttt{schwartzEvalCLM\_apply}]\label{GaussianField-schwartzEvalCLM-apply}
  \lean{GaussianField.schwartzEvalCLM_apply}\leanok
  \uses{GaussianField-schwartzEvalCLM}
  \texttt{GaussianField.schwartzEvalCLM\_apply}
\end{theorem}
\begin{definition}[\texttt{cylinderConfigTimeTranslation}]\label{GaussianField-cylinderConfigTimeTranslation}
  \lean{GaussianField.cylinderConfigTimeTranslation}\leanok
  \uses{GaussianField-NuclearTensorProduct-instModuleReal, GaussianField-NuclearTensorProduct-instTopologicalSpace, GaussianField-CylinderTestFunction, GaussianField-NuclearTensorProduct-instIsTopologicalAddGroup, GaussianField-NuclearTensorProduct-instAddCommGroup, GaussianField-Configuration, GaussianField-NuclearTensorProduct-instContinuousSMulReal, GaussianField-SmoothMap-Circle, GaussianField-cylinderTimeTranslation}
  \texttt{GaussianField.cylinderConfigTimeTranslation}
\end{definition}
\begin{theorem}[\texttt{cylinderTimeReflection\_involution}]\label{GaussianField-cylinderTimeReflection-involution}
  \lean{GaussianField.cylinderTimeReflection_involution}\leanok
  \uses{GaussianField-NuclearTensorProduct-instModuleReal, GaussianField-NuclearTensorProduct-instTopologicalSpace, GaussianField-nuclearTensorProduct-mapCLM-comp, GaussianField-nuclearTensorProduct-mapCLM-id, GaussianField-nuclearTensorProduct-mapCLM, GaussianField-schwartz-dyninMityaginSpace-1D, GaussianField-SmoothMap-Circle-instAddCommGroup, GaussianField-CylinderTestFunction, GaussianField-NuclearTensorProduct, GaussianField-NuclearTensorProduct-instAddCommGroup, GaussianField-smoothCircle-hasBiorthogonalBasis, GaussianField-cylinderTimeReflection, GaussianField-SmoothMap-Circle-instModule, GaussianField-SmoothMap-Circle-instContinuousSMul, GaussianField-smoothCircle-dyninMityaginSpace, GaussianField-SmoothMap-Circle-instIsTopologicalAddGroup, GaussianField-SmoothMap-Circle, GaussianField-schwartzReflection-involution, GaussianField-schwartzReflection, GaussianField-SmoothMap-Circle-instTopologicalSpace, GaussianField-schwartz-hasBiorthogonalBasis-1D}
  \texttt{GaussianField.cylinderTimeReflection\_involution}
\end{theorem}
\begin{definition}[\texttt{cylinderConfigSpatialTranslation}]\label{GaussianField-cylinderConfigSpatialTranslation}
  \lean{GaussianField.cylinderConfigSpatialTranslation}\leanok
  \uses{GaussianField-NuclearTensorProduct-instModuleReal, GaussianField-NuclearTensorProduct-instTopologicalSpace, GaussianField-cylinderSpatialTranslation, GaussianField-CylinderTestFunction, GaussianField-NuclearTensorProduct-instIsTopologicalAddGroup, GaussianField-NuclearTensorProduct-instAddCommGroup, GaussianField-Configuration, GaussianField-NuclearTensorProduct-instContinuousSMulReal, GaussianField-SmoothMap-Circle}
  \texttt{GaussianField.cylinderConfigSpatialTranslation}
\end{definition}
\begin{theorem}[\texttt{schwartzTranslation\_add}]\label{GaussianField-schwartzTranslation-add}
  \lean{GaussianField.schwartzTranslation_add}\leanok
  \uses{GaussianField-schwartzTranslation}
  \texttt{GaussianField.schwartzTranslation\_add}
\end{theorem}
\begin{theorem}[\texttt{cylinderConfigSpatialTranslation\_continuous}]\label{GaussianField-cylinderConfigSpatialTranslation-continuous}
  \lean{GaussianField.cylinderConfigSpatialTranslation_continuous}\leanok
  \uses{GaussianField-NuclearTensorProduct-instModuleReal, GaussianField-NuclearTensorProduct-instTopologicalSpace, GaussianField-cylinderSpatialTranslation, GaussianField-CylinderTestFunction, GaussianField-NuclearTensorProduct-instIsTopologicalAddGroup, GaussianField-NuclearTensorProduct-instAddCommGroup, GaussianField-cylinderConfigSpatialTranslation, GaussianField-Configuration, GaussianField-NuclearTensorProduct-instContinuousSMulReal, GaussianField-SmoothMap-Circle}
  \texttt{GaussianField.cylinderConfigSpatialTranslation\_continuous}
\end{theorem}
\begin{theorem}[\texttt{schwartzPositiveTime\_disjoint\_reflected}]\label{GaussianField-schwartzPositiveTime-disjoint-reflected}
  \lean{GaussianField.schwartzPositiveTime_disjoint_reflected}\leanok
  \uses{GaussianField-schwartzPositiveTimeSubmodule, GaussianField-schwartzReflection}
  \texttt{GaussianField.schwartzPositiveTime\_disjoint\_reflected}
\end{theorem}
\begin{definition}[\texttt{schwartzNegativeTimeSubmodule}]\label{GaussianField-schwartzNegativeTimeSubmodule}
  \lean{GaussianField.schwartzNegativeTimeSubmodule}\leanok
  \texttt{GaussianField.schwartzNegativeTimeSubmodule}
\end{definition}
\begin{definition}[\texttt{schwartzEvalCLM}]\label{GaussianField-schwartzEvalCLM}
  \lean{GaussianField.schwartzEvalCLM}\leanok
  \texttt{GaussianField.schwartzEvalCLM}
\end{definition}
\begin{definition}[\texttt{cylinderConfigTranslation}]\label{GaussianField-cylinderConfigTranslation}
  \lean{GaussianField.cylinderConfigTranslation}\leanok
  \uses{GaussianField-NuclearTensorProduct-instModuleReal, GaussianField-NuclearTensorProduct-instTopologicalSpace, GaussianField-CylinderTestFunction, GaussianField-NuclearTensorProduct-instIsTopologicalAddGroup, GaussianField-NuclearTensorProduct-instAddCommGroup, GaussianField-Configuration, GaussianField-NuclearTensorProduct-instContinuousSMulReal, GaussianField-SmoothMap-Circle, GaussianField-cylinderTranslation}
  \texttt{GaussianField.cylinderConfigTranslation}
\end{definition}
\begin{theorem}[\texttt{schwartzTranslation\_preserves\_positiveTime}]\label{GaussianField-schwartzTranslation-preserves-positiveTime}
  \lean{GaussianField.schwartzTranslation_preserves_positiveTime}\leanok
  \uses{GaussianField-schwartzPositiveTimeSubmodule, GaussianField-schwartzTranslation}
  \texttt{GaussianField.schwartzTranslation\_preserves\_positiveTime}
\end{theorem}
\begin{definition}[\texttt{cylinderSpatialTranslation}]\label{GaussianField-cylinderSpatialTranslation}
  \lean{GaussianField.cylinderSpatialTranslation}\leanok
  \uses{GaussianField-NuclearTensorProduct-instModuleReal, GaussianField-circleTranslation, GaussianField-NuclearTensorProduct-instTopologicalSpace, GaussianField-nuclearTensorProduct-mapCLM, GaussianField-schwartz-dyninMityaginSpace-1D, GaussianField-SmoothMap-Circle-instAddCommGroup, GaussianField-CylinderTestFunction, GaussianField-NuclearTensorProduct-instAddCommGroup, GaussianField-SmoothMap-Circle-instModule, GaussianField-SmoothMap-Circle-instContinuousSMul, GaussianField-smoothCircle-dyninMityaginSpace, GaussianField-SmoothMap-Circle-instIsTopologicalAddGroup, GaussianField-SmoothMap-Circle, GaussianField-SmoothMap-Circle-instTopologicalSpace}
  \texttt{GaussianField.cylinderSpatialTranslation}
\end{definition}
\begin{theorem}[\texttt{cylinderConfigTimeTranslation\_continuous}]\label{GaussianField-cylinderConfigTimeTranslation-continuous}
  \lean{GaussianField.cylinderConfigTimeTranslation_continuous}\leanok
  \uses{GaussianField-NuclearTensorProduct-instModuleReal, GaussianField-NuclearTensorProduct-instTopologicalSpace, GaussianField-CylinderTestFunction, GaussianField-NuclearTensorProduct-instIsTopologicalAddGroup, GaussianField-NuclearTensorProduct-instAddCommGroup, GaussianField-Configuration, GaussianField-NuclearTensorProduct-instContinuousSMulReal, GaussianField-SmoothMap-Circle, GaussianField-cylinderTimeTranslation, GaussianField-cylinderConfigTimeTranslation}
  \texttt{GaussianField.cylinderConfigTimeTranslation\_continuous}
\end{theorem}
\begin{theorem}[\texttt{schwartzReflection\_positive\_to\_negative}]\label{GaussianField-schwartzReflection-positive-to-negative}
  \lean{GaussianField.schwartzReflection_positive_to_negative}\leanok
  \uses{GaussianField-schwartzPositiveTimeSubmodule, GaussianField-schwartzReflection}
  \texttt{GaussianField.schwartzReflection\_positive\_to\_negative}
\end{theorem}
\begin{theorem}[\texttt{cylinderConfigReflection\_continuous}]\label{GaussianField-cylinderConfigReflection-continuous}
  \lean{GaussianField.cylinderConfigReflection_continuous}\leanok
  \uses{GaussianField-NuclearTensorProduct-instModuleReal, GaussianField-NuclearTensorProduct-instTopologicalSpace, GaussianField-CylinderTestFunction, GaussianField-NuclearTensorProduct-instIsTopologicalAddGroup, GaussianField-NuclearTensorProduct-instAddCommGroup, GaussianField-cylinderTimeReflection, GaussianField-Configuration, GaussianField-NuclearTensorProduct-instContinuousSMulReal, GaussianField-SmoothMap-Circle, GaussianField-cylinderConfigReflection}
  \texttt{GaussianField.cylinderConfigReflection\_continuous}
\end{theorem}
\begin{definition}[\texttt{cylinderTranslation}]\label{GaussianField-cylinderTranslation}
  \lean{GaussianField.cylinderTranslation}\leanok
  \uses{GaussianField-NuclearTensorProduct-instModuleReal, GaussianField-circleTranslation, GaussianField-NuclearTensorProduct-instTopologicalSpace, GaussianField-nuclearTensorProduct-mapCLM, GaussianField-schwartz-dyninMityaginSpace-1D, GaussianField-SmoothMap-Circle-instAddCommGroup, GaussianField-CylinderTestFunction, GaussianField-NuclearTensorProduct-instAddCommGroup, GaussianField-SmoothMap-Circle-instModule, GaussianField-SmoothMap-Circle-instContinuousSMul, GaussianField-smoothCircle-dyninMityaginSpace, GaussianField-SmoothMap-Circle-instIsTopologicalAddGroup, GaussianField-SmoothMap-Circle, GaussianField-schwartzTranslation, GaussianField-SmoothMap-Circle-instTopologicalSpace}
  \texttt{GaussianField.cylinderTranslation}
\end{definition}
\begin{definition}[\texttt{cylinderTimeReflection}]\label{GaussianField-cylinderTimeReflection}
  \lean{GaussianField.cylinderTimeReflection}\leanok
  \uses{GaussianField-NuclearTensorProduct-instModuleReal, GaussianField-NuclearTensorProduct-instTopologicalSpace, GaussianField-nuclearTensorProduct-mapCLM, GaussianField-schwartz-dyninMityaginSpace-1D, GaussianField-SmoothMap-Circle-instAddCommGroup, GaussianField-CylinderTestFunction, GaussianField-NuclearTensorProduct-instAddCommGroup, GaussianField-SmoothMap-Circle-instModule, GaussianField-SmoothMap-Circle-instContinuousSMul, GaussianField-smoothCircle-dyninMityaginSpace, GaussianField-SmoothMap-Circle-instIsTopologicalAddGroup, GaussianField-SmoothMap-Circle, GaussianField-schwartzReflection, GaussianField-SmoothMap-Circle-instTopologicalSpace}
  \texttt{GaussianField.cylinderTimeReflection}
\end{definition}
\begin{theorem}[\texttt{schwartzReflection\_apply}]\label{GaussianField-schwartzReflection-apply}
  \lean{GaussianField.schwartzReflection_apply}\leanok
  \uses{GaussianField-schwartzReflection}
  \texttt{GaussianField.schwartzReflection\_apply}
\end{theorem}
\begin{definition}[\texttt{schwartzTranslation}]\label{GaussianField-schwartzTranslation}
  \lean{GaussianField.schwartzTranslation}\leanok
  \texttt{GaussianField.schwartzTranslation}
\end{definition}
\begin{definition}[\texttt{schwartzPositiveTimeSubmodule}]\label{GaussianField-schwartzPositiveTimeSubmodule}
  \lean{GaussianField.schwartzPositiveTimeSubmodule}\leanok
  \texttt{GaussianField.schwartzPositiveTimeSubmodule}
\end{definition}
\begin{definition}[\texttt{cylinderConfigReflection}]\label{GaussianField-cylinderConfigReflection}
  \lean{GaussianField.cylinderConfigReflection}\leanok
  \uses{GaussianField-NuclearTensorProduct-instModuleReal, GaussianField-NuclearTensorProduct-instTopologicalSpace, GaussianField-CylinderTestFunction, GaussianField-NuclearTensorProduct-instIsTopologicalAddGroup, GaussianField-NuclearTensorProduct-instAddCommGroup, GaussianField-cylinderTimeReflection, GaussianField-Configuration, GaussianField-NuclearTensorProduct-instContinuousSMulReal, GaussianField-SmoothMap-Circle}
  \texttt{GaussianField.cylinderConfigReflection}
\end{definition}
\begin{definition}[\texttt{cylinderTimeTranslation}]\label{GaussianField-cylinderTimeTranslation}
  \lean{GaussianField.cylinderTimeTranslation}\leanok
  \uses{GaussianField-NuclearTensorProduct-instModuleReal, GaussianField-NuclearTensorProduct-instTopologicalSpace, GaussianField-nuclearTensorProduct-mapCLM, GaussianField-schwartz-dyninMityaginSpace-1D, GaussianField-SmoothMap-Circle-instAddCommGroup, GaussianField-CylinderTestFunction, GaussianField-NuclearTensorProduct-instAddCommGroup, GaussianField-SmoothMap-Circle-instModule, GaussianField-SmoothMap-Circle-instContinuousSMul, GaussianField-smoothCircle-dyninMityaginSpace, GaussianField-SmoothMap-Circle-instIsTopologicalAddGroup, GaussianField-SmoothMap-Circle, GaussianField-schwartzTranslation, GaussianField-SmoothMap-Circle-instTopologicalSpace}
  \texttt{GaussianField.cylinderTimeTranslation}
\end{definition}
\begin{theorem}[\texttt{schwartzReflection\_involution}]\label{GaussianField-schwartzReflection-involution}
  \lean{GaussianField.schwartzReflection_involution}\leanok
  \uses{GaussianField-schwartzReflection}
  \texttt{GaussianField.schwartzReflection\_involution}
\end{theorem}
\section{\texttt{Lattice.Covariance}}
\begin{theorem}[\texttt{latticeSingularValue\_nonneg}]\label{GaussianField-latticeSingularValue-nonneg}
  \lean{GaussianField.latticeSingularValue_nonneg}\leanok
  \uses{GaussianField-latticeSingularValue, GaussianField-latticeEigenvalue-pos, GaussianField-latticeEigenvalue}
  \texttt{GaussianField.latticeSingularValue\_nonneg}
\end{theorem}
\begin{theorem}[\texttt{latticeGaussianMeasure\_isProbability}]\label{GaussianField-latticeGaussianMeasure-isProbability}
  \lean{GaussianField.latticeGaussianMeasure_isProbability}\leanok
  \uses{GaussianField-finLatticeField-dyninMityaginSpace, GaussianField-ell2-separableSpace, GaussianField-FinLatticeSites, GaussianField-latticeCovarianceGJ, GaussianField-latticeGaussianMeasure, GaussianField-Configuration, GaussianField-instMeasurableSpaceConfiguration, GaussianField-ell2-, GaussianField-FinLatticeField, GaussianField-measure-isProbability}
  \texttt{GaussianField.latticeGaussianMeasure\_isProbability}
\end{theorem}
\begin{definition}[\texttt{latticeGaussianMeasure}]\label{GaussianField-latticeGaussianMeasure}
  \lean{GaussianField.latticeGaussianMeasure}\leanok
  \uses{GaussianField-finLatticeField-dyninMityaginSpace, GaussianField-measure, GaussianField-ell2-separableSpace, GaussianField-FinLatticeSites, GaussianField-latticeCovarianceGJ, GaussianField-Configuration, GaussianField-instMeasurableSpaceConfiguration, GaussianField-ell2-, GaussianField-FinLatticeField}
  \texttt{GaussianField.latticeGaussianMeasure}
\end{definition}
\begin{theorem}[\texttt{congr\_simp}]\label{GaussianField-latticeEigenvalue-congr-simp}
  \lean{GaussianField.latticeEigenvalue.congr_simp}\leanok
  \uses{GaussianField-latticeEigenvalue}
  \texttt{GaussianField.latticeEigenvalue.congr\_simp}
\end{theorem}
\begin{definition}[\texttt{latticeCovariance}]\label{GaussianField-latticeCovariance}
  \lean{GaussianField.latticeCovariance}\leanok
  \uses{GaussianField-FinLatticeSites, GaussianField-spectralLatticeCovariance, GaussianField-ell2-, GaussianField-FinLatticeField}
  \texttt{GaussianField.latticeCovariance}
\end{definition}
\begin{theorem}[\texttt{latticeCovariance\_GJ\_eq\_inv\_smul\_bare}]\label{GaussianField-latticeCovariance-GJ-eq-inv-smul-bare}
  \lean{GaussianField.latticeCovariance_GJ_eq_inv_smul_bare}\leanok
  \uses{GaussianField-covariance, GaussianField-latticeCovariance, GaussianField-FinLatticeSites, GaussianField-latticeCovarianceGJ, GaussianField-ell2-, GaussianField-FinLatticeField}
  \texttt{GaussianField.latticeCovariance\_GJ\_eq\_inv\_smul\_bare}
\end{theorem}
\begin{theorem}[\texttt{lattice\_cross\_moment}]\label{GaussianField-lattice-cross-moment}
  \lean{GaussianField.lattice_cross_moment}\leanok
  \uses{GaussianField-cross-moment-eq-covariance, GaussianField-finLatticeField-dyninMityaginSpace, GaussianField-ell2-separableSpace, GaussianField-covariance, GaussianField-FinLatticeSites, GaussianField-latticeCovarianceGJ, GaussianField-latticeGaussianMeasure, GaussianField-Configuration, GaussianField-instMeasurableSpaceConfiguration, GaussianField-ell2-, GaussianField-FinLatticeField}
  \texttt{GaussianField.lattice\_cross\_moment}
\end{theorem}
\begin{definition}[\texttt{latticeCovarianceGJ}]\label{GaussianField-latticeCovarianceGJ}
  \lean{GaussianField.latticeCovarianceGJ}\leanok
  \uses{GaussianField-latticeCovariance, GaussianField-FinLatticeSites, GaussianField-ell2-, GaussianField-FinLatticeField}
  \texttt{GaussianField.latticeCovarianceGJ}
\end{definition}
\begin{theorem}[\texttt{congr\_simp}]\label{GaussianField-latticeCovariance-congr-simp}
  \lean{GaussianField.latticeCovariance.congr_simp}\leanok
  \uses{GaussianField-latticeCovariance, GaussianField-FinLatticeSites, GaussianField-ell2-, GaussianField-FinLatticeField}
  \texttt{GaussianField.latticeCovariance.congr\_simp}
\end{theorem}
\begin{theorem}[\texttt{lattice\_singular\_values\_bounded}]\label{GaussianField-lattice-singular-values-bounded}
  \lean{GaussianField.lattice_singular_values_bounded}\leanok
  \uses{GaussianField-latticeSingularValue, GaussianField-latticeEigenvalue-eq, GaussianField-latticeLaplacianEigenvalue, GaussianField-latticeLaplacianEigenvalue-nonneg, GaussianField-latticeEigenvalue-pos, GaussianField-latticeSingularValue-nonneg, GaussianField-latticeEigenvalue, GaussianField-IsBoundedSeq}
  \texttt{GaussianField.lattice\_singular\_values\_bounded}
\end{theorem}
\begin{definition}[\texttt{latticeSingularValue}]\label{GaussianField-latticeSingularValue}
  \lean{GaussianField.latticeSingularValue}\leanok
  \uses{GaussianField-latticeEigenvalue}
  \texttt{GaussianField.latticeSingularValue}
\end{definition}
\begin{theorem}[\texttt{lattice\_covariance\_eq\_spectral}]\label{GaussianField-lattice-covariance-eq-spectral}
  \lean{GaussianField.lattice_covariance_eq_spectral}\leanok
  \uses{GaussianField-covariance, GaussianField-latticeCovariance, GaussianField-FinLatticeSites, GaussianField-spectralLatticeCovariance-inner, GaussianField-massEigenvalues, GaussianField-massEigenvectorBasis, GaussianField-ell2-, GaussianField-FinLatticeField}
  \texttt{GaussianField.lattice\_covariance\_eq\_spectral}
\end{theorem}
\begin{theorem}[\texttt{lattice\_covariance\_GJ\_eq\_spectral}]\label{GaussianField-lattice-covariance-GJ-eq-spectral}
  \lean{GaussianField.lattice_covariance_GJ_eq_spectral}\leanok
  \uses{GaussianField-latticeCovariance-GJ-eq-inv-smul-bare, GaussianField-lattice-covariance-eq-spectral, GaussianField-covariance, GaussianField-latticeCovariance, GaussianField-FinLatticeSites, GaussianField-latticeCovarianceGJ, GaussianField-massEigenvalues, GaussianField-massEigenvectorBasis, GaussianField-ell2-, GaussianField-FinLatticeField}
  \texttt{GaussianField.lattice\_covariance\_GJ\_eq\_spectral}
\end{theorem}
\section{\texttt{HeatKernel.Bilinear}}
\begin{theorem}[\texttt{greenFunctionBilinear\_add\_right}]\label{GaussianField-greenFunctionBilinear-add-right}
  \lean{GaussianField.greenFunctionBilinear_add_right}\leanok
  \uses{GaussianField-DyninMityaginSpace, GaussianField-HasLaplacianEigenvalues, GaussianField-greenFunctionBilinear-add-left, GaussianField-greenFunctionBilinear, GaussianField-greenFunctionBilinear-symm}
  \texttt{GaussianField.greenFunctionBilinear\_add\_right}
\end{theorem}
\begin{theorem}[\texttt{congr\_simp}]\label{GaussianField-greenFunctionBilinear-congr-simp}
  \lean{GaussianField.greenFunctionBilinear.congr_simp}\leanok
  \uses{GaussianField-DyninMityaginSpace, GaussianField-HasLaplacianEigenvalues, GaussianField-greenFunctionBilinear}
  \texttt{GaussianField.greenFunctionBilinear.congr\_simp}
\end{theorem}
\begin{theorem}[\texttt{greenFunctionBilinear\_finset\_sum}]\label{GaussianField-greenFunctionBilinear-finset-sum}
  \lean{GaussianField.greenFunctionBilinear_finset_sum}\leanok
  \uses{GaussianField-DyninMityaginSpace, GaussianField-HasLaplacianEigenvalues, GaussianField-greenFunctionBilinear-finset-sum-left, GaussianField-greenFunctionBilinear, GaussianField-greenFunctionBilinear-symm}
  \texttt{GaussianField.greenFunctionBilinear\_finset\_sum}
\end{theorem}
\begin{theorem}[\texttt{greenFunctionBilinear\_add\_left}]\label{GaussianField-greenFunctionBilinear-add-left}
  \lean{GaussianField.greenFunctionBilinear_add_left}\leanok
  \uses{GaussianField-DyninMityaginSpace, GaussianField-HasLaplacianEigenvalues, GaussianField-DyninMityaginSpace-coeff, GaussianField-greenFunctionBilinear-summable, GaussianField-HasLaplacianEigenvalues-eigenvalue, GaussianField-greenFunctionBilinear}
  \texttt{GaussianField.greenFunctionBilinear\_add\_left}
\end{theorem}
\begin{definition}[\texttt{ctorIdx}]\label{GaussianField-HasLaplacianEigenvalues-ctorIdx}
  \lean{GaussianField.HasLaplacianEigenvalues.ctorIdx}\leanok
  \uses{GaussianField-DyninMityaginSpace, GaussianField-HasLaplacianEigenvalues}
  \texttt{GaussianField.HasLaplacianEigenvalues.ctorIdx}
\end{definition}
\begin{definition}[\texttt{l2InnerProduct}]\label{GaussianField-l2InnerProduct}
  \lean{GaussianField.l2InnerProduct}\leanok
  \uses{GaussianField-DyninMityaginSpace}
  \texttt{GaussianField.l2InnerProduct}
\end{definition}
\begin{definition}[\texttt{tensorProductHasLaplacianEigenvalues}]\label{GaussianField-tensorProductHasLaplacianEigenvalues}
  \lean{GaussianField.tensorProductHasLaplacianEigenvalues}\leanok
  \uses{GaussianField-DyninMityaginSpace, GaussianField-NuclearTensorProduct-instModuleReal, GaussianField-HasLaplacianEigenvalues, GaussianField-NuclearTensorProduct-instTopologicalSpace, GaussianField-NuclearTensorProduct-dyninMityaginSpace, GaussianField-NuclearTensorProduct-instIsTopologicalAddGroup, GaussianField-NuclearTensorProduct, GaussianField-NuclearTensorProduct-instAddCommGroup, GaussianField-HasLaplacianEigenvalues-eigenvalue, GaussianField-NuclearTensorProduct-instContinuousSMulReal}
  \texttt{GaussianField.tensorProductHasLaplacianEigenvalues}
\end{definition}
\begin{theorem}[\texttt{greenFunctionBilinear\_summable}]\label{GaussianField-greenFunctionBilinear-summable}
  \lean{GaussianField.greenFunctionBilinear_summable}\leanok
  \uses{GaussianField-DyninMityaginSpace, GaussianField-HasLaplacianEigenvalues, GaussianField-DyninMityaginSpace-coeff, GaussianField-HasLaplacianEigenvalues-eigenvalue, GaussianField-HasLaplacianEigenvalues-eigenvalue-nonneg, GaussianField-l2InnerProduct-summable, Lattice-toSemilatticeInf}
  \texttt{GaussianField.greenFunctionBilinear\_summable}
\end{theorem}
\begin{theorem}[\texttt{greenFunctionBilinear\_sub\_left}]\label{GaussianField-greenFunctionBilinear-sub-left}
  \lean{GaussianField.greenFunctionBilinear_sub_left}\leanok
  \uses{GaussianField-DyninMityaginSpace, GaussianField-HasLaplacianEigenvalues, GaussianField-greenFunctionBilinear-add-left, GaussianField-greenFunctionBilinear-neg-left, GaussianField-greenFunctionBilinear}
  \texttt{GaussianField.greenFunctionBilinear\_sub\_left}
\end{theorem}
\begin{theorem}[\texttt{greenFunctionBilinear\_zero\_left}]\label{GaussianField-greenFunctionBilinear-zero-left}
  \lean{GaussianField.greenFunctionBilinear_zero_left}\leanok
  \uses{GaussianField-DyninMityaginSpace, GaussianField-HasLaplacianEigenvalues, GaussianField-greenFunctionBilinear-congr-simp, GaussianField-greenFunctionBilinear-smul-left, GaussianField-greenFunctionBilinear}
  \texttt{GaussianField.greenFunctionBilinear\_zero\_left}
\end{theorem}
\begin{theorem}[\texttt{greenFunctionBilinear\_neg\_left}]\label{GaussianField-greenFunctionBilinear-neg-left}
  \lean{GaussianField.greenFunctionBilinear_neg_left}\leanok
  \uses{GaussianField-DyninMityaginSpace, GaussianField-HasLaplacianEigenvalues, GaussianField-greenFunctionBilinear-smul-left, GaussianField-greenFunctionBilinear}
  \texttt{GaussianField.greenFunctionBilinear\_neg\_left}
\end{theorem}
\begin{definition}[\texttt{heatKernelBilinear}]\label{GaussianField-heatKernelBilinear}
  \lean{GaussianField.heatKernelBilinear}\leanok
  \uses{GaussianField-DyninMityaginSpace, GaussianField-HasLaplacianEigenvalues}
  \texttt{GaussianField.heatKernelBilinear}
\end{definition}
\begin{definition}[\texttt{HasLaplacianEigenvalues}]\label{GaussianField-HasLaplacianEigenvalues}
  \lean{GaussianField.HasLaplacianEigenvalues}\leanok
  \uses{GaussianField-DyninMityaginSpace}
  \texttt{GaussianField.HasLaplacianEigenvalues}
\end{definition}
\begin{definition}[\texttt{eigenvalue}]\label{GaussianField-HasLaplacianEigenvalues-eigenvalue}
  \lean{GaussianField.HasLaplacianEigenvalues.eigenvalue}\leanok
  \uses{GaussianField-DyninMityaginSpace, GaussianField-HasLaplacianEigenvalues}
  \texttt{GaussianField.HasLaplacianEigenvalues.eigenvalue}
\end{definition}
\begin{theorem}[\texttt{heatKernelBilinear\_nonneg}]\label{GaussianField-heatKernelBilinear-nonneg}
  \lean{GaussianField.heatKernelBilinear_nonneg}\leanok
  \uses{GaussianField-DyninMityaginSpace, GaussianField-HasLaplacianEigenvalues, GaussianField-DyninMityaginSpace-coeff, GaussianField-HasLaplacianEigenvalues-eigenvalue, GaussianField-heatKernelBilinear, Lattice-toSemilatticeInf}
  \texttt{GaussianField.heatKernelBilinear\_nonneg}
\end{theorem}
\begin{theorem}[\texttt{eigenvalue\_nonneg}]\label{GaussianField-HasLaplacianEigenvalues-eigenvalue-nonneg}
  \lean{GaussianField.HasLaplacianEigenvalues.eigenvalue_nonneg}\leanok
  \uses{GaussianField-DyninMityaginSpace, GaussianField-HasLaplacianEigenvalues, GaussianField-HasLaplacianEigenvalues-eigenvalue}
  \texttt{GaussianField.HasLaplacianEigenvalues.eigenvalue\_nonneg}
\end{theorem}
\begin{theorem}[\texttt{greenFunctionBilinear\_le}]\label{GaussianField-greenFunctionBilinear-le}
  \lean{GaussianField.greenFunctionBilinear_le}\leanok
  \uses{GaussianField-DyninMityaginSpace, GaussianField-l2InnerProduct, GaussianField-HasLaplacianEigenvalues, GaussianField-DyninMityaginSpace-coeff, GaussianField-greenFunctionBilinear-summable, GaussianField-HasLaplacianEigenvalues-eigenvalue, GaussianField-HasLaplacianEigenvalues-eigenvalue-nonneg, GaussianField-l2InnerProduct-summable, GaussianField-greenFunctionBilinear, Lattice-toSemilatticeInf}
  \texttt{GaussianField.greenFunctionBilinear\_le}
\end{theorem}
\begin{theorem}[\texttt{greenFunctionBilinear\_smul\_left}]\label{GaussianField-greenFunctionBilinear-smul-left}
  \lean{GaussianField.greenFunctionBilinear_smul_left}\leanok
  \uses{GaussianField-DyninMityaginSpace, GaussianField-HasLaplacianEigenvalues, GaussianField-DyninMityaginSpace-coeff, GaussianField-HasLaplacianEigenvalues-eigenvalue, GaussianField-greenFunctionBilinear}
  \texttt{GaussianField.greenFunctionBilinear\_smul\_left}
\end{theorem}
\begin{definition}[\texttt{greenFunctionBilinear}]\label{GaussianField-greenFunctionBilinear}
  \lean{GaussianField.greenFunctionBilinear}\leanok
  \uses{GaussianField-DyninMityaginSpace, GaussianField-HasLaplacianEigenvalues}
  \texttt{GaussianField.greenFunctionBilinear}
\end{definition}
\begin{theorem}[\texttt{heatKernelBilinear\_tensorProduct}]\label{GaussianField-heatKernelBilinear-tensorProduct}
  \lean{GaussianField.heatKernelBilinear_tensorProduct}\leanok
  \uses{GaussianField-DyninMityaginSpace, GaussianField-NuclearTensorProduct-instModuleReal, GaussianField-HasLaplacianEigenvalues, GaussianField-NuclearTensorProduct-instTopologicalSpace, GaussianField-NuclearTensorProduct-dyninMityaginSpace, GaussianField-NuclearTensorProduct-instIsTopologicalAddGroup, GaussianField-tensorProductHasLaplacianEigenvalues, GaussianField-NuclearTensorProduct, GaussianField-NuclearTensorProduct-instAddCommGroup, GaussianField-NuclearTensorProduct-instContinuousSMulReal, GaussianField-NuclearTensorProduct-pure, GaussianField-heatKernelBilinear}
  \texttt{GaussianField.heatKernelBilinear\_tensorProduct}
\end{theorem}
\begin{theorem}[\texttt{greenFunctionBilinear\_pos}]\label{GaussianField-greenFunctionBilinear-pos}
  \lean{GaussianField.greenFunctionBilinear_pos}\leanok
  \uses{GaussianField-DyninMityaginSpace, GaussianField-HasLaplacianEigenvalues, GaussianField-DyninMityaginSpace-coeff, GaussianField-greenFunctionBilinear-summable, GaussianField-DyninMityaginSpace-expansion, GaussianField-DyninMityaginSpace--, GaussianField-HasLaplacianEigenvalues-eigenvalue, GaussianField-DyninMityaginSpace-h-with, GaussianField-HasLaplacianEigenvalues-eigenvalue-nonneg, GaussianField-DyninMityaginSpace-basis, GaussianField-DyninMityaginSpace-p, GaussianField-greenFunctionBilinear, Lattice-toSemilatticeInf}
  \texttt{GaussianField.greenFunctionBilinear\_pos}
\end{theorem}
\begin{theorem}[\texttt{heatKernelBilinear\_le\_l2}]\label{GaussianField-heatKernelBilinear-le-l2}
  \lean{GaussianField.heatKernelBilinear_le_l2}\leanok
  \uses{GaussianField-DyninMityaginSpace, GaussianField-l2InnerProduct, GaussianField-HasLaplacianEigenvalues, GaussianField-DyninMityaginSpace-coeff, GaussianField-HasLaplacianEigenvalues-eigenvalue, GaussianField-HasLaplacianEigenvalues-eigenvalue-nonneg, GaussianField-heatKernelBilinear-summable, GaussianField-heatKernelBilinear, GaussianField-l2InnerProduct-summable, Lattice-toSemilatticeInf}
  \texttt{GaussianField.heatKernelBilinear\_le\_l2}
\end{theorem}
\begin{theorem}[\texttt{greenFunctionBilinear\_eq\_heatKernel}]\label{GaussianField-greenFunctionBilinear-eq-heatKernel}
  \lean{GaussianField.greenFunctionBilinear_eq_heatKernel}\leanok
  \uses{GaussianField-DyninMityaginSpace, GaussianField-HasLaplacianEigenvalues, GaussianField-DyninMityaginSpace-coeff, GaussianField-HasLaplacianEigenvalues-eigenvalue}
  \texttt{GaussianField.greenFunctionBilinear\_eq\_heatKernel}
\end{theorem}
\begin{theorem}[\texttt{greenFunctionBilinear\_symm}]\label{GaussianField-greenFunctionBilinear-symm}
  \lean{GaussianField.greenFunctionBilinear_symm}\leanok
  \uses{GaussianField-DyninMityaginSpace, GaussianField-HasLaplacianEigenvalues, GaussianField-DyninMityaginSpace-coeff, GaussianField-HasLaplacianEigenvalues-eigenvalue, GaussianField-greenFunctionBilinear}
  \texttt{GaussianField.greenFunctionBilinear\_symm}
\end{theorem}
\begin{theorem}[\texttt{heatKernelBilinear\_tendsto\_l2}]\label{GaussianField-heatKernelBilinear-tendsto-l2}
  \lean{GaussianField.heatKernelBilinear_tendsto_l2}\leanok
  \uses{GaussianField-DyninMityaginSpace, GaussianField-l2InnerProduct, GaussianField-HasLaplacianEigenvalues, GaussianField-DyninMityaginSpace-coeff, GaussianField-HasLaplacianEigenvalues-eigenvalue, GaussianField-HasLaplacianEigenvalues-eigenvalue-nonneg, GaussianField-heatKernelBilinear, GaussianField-l2InnerProduct-summable, Lattice-toSemilatticeInf}
  \texttt{GaussianField.heatKernelBilinear\_tendsto\_l2}
\end{theorem}
\begin{theorem}[\texttt{heatKernelBilinear\_summable}]\label{GaussianField-heatKernelBilinear-summable}
  \lean{GaussianField.heatKernelBilinear_summable}\leanok
  \uses{GaussianField-DyninMityaginSpace, GaussianField-HasLaplacianEigenvalues, GaussianField-DyninMityaginSpace-coeff, GaussianField-HasLaplacianEigenvalues-eigenvalue, GaussianField-HasLaplacianEigenvalues-eigenvalue-nonneg, GaussianField-l2InnerProduct-summable, Lattice-toSemilatticeInf}
  \texttt{GaussianField.heatKernelBilinear\_summable}
\end{theorem}
\begin{theorem}[\texttt{greenFunctionBilinear\_smul\_right}]\label{GaussianField-greenFunctionBilinear-smul-right}
  \lean{GaussianField.greenFunctionBilinear_smul_right}\leanok
  \uses{GaussianField-DyninMityaginSpace, GaussianField-HasLaplacianEigenvalues, GaussianField-greenFunctionBilinear-smul-left, GaussianField-greenFunctionBilinear, GaussianField-greenFunctionBilinear-symm}
  \texttt{GaussianField.greenFunctionBilinear\_smul\_right}
\end{theorem}
\begin{theorem}[\texttt{greenFunctionBilinear\_finset\_sum\_left}]\label{GaussianField-greenFunctionBilinear-finset-sum-left}
  \lean{GaussianField.greenFunctionBilinear_finset_sum_left}\leanok
  \uses{GaussianField-DyninMityaginSpace, GaussianField-HasLaplacianEigenvalues, GaussianField-greenFunctionBilinear-add-left, GaussianField-greenFunctionBilinear-smul-left, GaussianField-greenFunctionBilinear-zero-left, GaussianField-greenFunctionBilinear}
  \texttt{GaussianField.greenFunctionBilinear\_finset\_sum\_left}
\end{theorem}
\begin{theorem}[\texttt{l2InnerProduct\_summable}]\label{GaussianField-l2InnerProduct-summable}
  \lean{GaussianField.l2InnerProduct_summable}\leanok
  \uses{GaussianField-DyninMityaginSpace, GaussianField-DyninMityaginSpace-coeff, GaussianField-DyninMityaginSpace-coeff-decay, GaussianField-DyninMityaginSpace--, GaussianField-DyninMityaginSpace-p, Lattice-toSemilatticeInf}
  \texttt{GaussianField.l2InnerProduct\_summable}
\end{theorem}
\begin{theorem}[\texttt{greenFunctionBilinear\_nonneg}]\label{GaussianField-greenFunctionBilinear-nonneg}
  \lean{GaussianField.greenFunctionBilinear_nonneg}\leanok
  \uses{GaussianField-DyninMityaginSpace, GaussianField-HasLaplacianEigenvalues, GaussianField-DyninMityaginSpace-coeff, GaussianField-HasLaplacianEigenvalues-eigenvalue, GaussianField-HasLaplacianEigenvalues-eigenvalue-nonneg, GaussianField-greenFunctionBilinear, Lattice-toSemilatticeInf}
  \texttt{GaussianField.greenFunctionBilinear\_nonneg}
\end{theorem}
\begin{theorem}[\texttt{greenFunctionBilinear\_continuous\_diag}]\label{GaussianField-greenFunctionBilinear-continuous-diag}
  \lean{GaussianField.greenFunctionBilinear_continuous_diag}\leanok
  \uses{GaussianField-DyninMityaginSpace, GaussianField-HasLaplacianEigenvalues, Lattice-toSemilatticeSup, GaussianField-DyninMityaginSpace-coeff, GaussianField-DyninMityaginSpace-coeff-decay, GaussianField-DyninMityaginSpace--, GaussianField-HasLaplacianEigenvalues-eigenvalue, GaussianField-DyninMityaginSpace-h-with, GaussianField-HasLaplacianEigenvalues-eigenvalue-nonneg, GaussianField-DyninMityaginSpace-p, GaussianField-greenFunctionBilinear, Lattice-toSemilatticeInf}
  \texttt{GaussianField.greenFunctionBilinear\_continuous\_diag}
\end{theorem}
\section{\texttt{Lattice.CirculantDFT2d}}
\begin{theorem}[\texttt{abstract\_spectral\_eq\_dft\_spectral\_2d}]\label{GaussianField-abstract-spectral-eq-dft-spectral-2d}
  \lean{GaussianField.abstract_spectral_eq_dft_spectral_2d}\leanok
  \uses{GaussianField-latticeEigenvalue1d, GaussianField-latticeFourierBasisFun, GaussianField-lattice-covariance-eq-spectral, GaussianField-covariance, GaussianField-latticeCovariance, GaussianField-FinLatticeSites, GaussianField-massEigenbasis-sum-mul-sum-eq-site-inner, GaussianField-massOperator, GaussianField-latticeEigenvalue1d-nonneg, GaussianField-latticeFourierNormSq, GaussianField-dft-parseval-2d, GaussianField-massEigenvalues, GaussianField-massEigenvectorBasis, GaussianField-dft-eigencoeff-massOperator-2d, GaussianField-massOperator-surjective-2d, GaussianField-massOperatorMatrix-eigenvalues-pos, GaussianField-ell2-, GaussianField-FinLatticeField, GaussianField-massOperator-eigenCoeff-eq-eigenvalues-mul-eigenCoeff}
  \texttt{GaussianField.abstract\_spectral\_eq\_dft\_spectral\_2d}
\end{theorem}
\begin{theorem}[\texttt{dft\_parseval\_2d}]\label{GaussianField-dft-parseval-2d}
  \lean{GaussianField.dft_parseval_2d}\leanok
  \uses{GaussianField-latticeFourierBasisFun, GaussianField-dft-parseval-1d, GaussianField-FinLatticeSites, GaussianField-sum-factor, GaussianField-latticeFourierNormSq}
  \texttt{GaussianField.dft\_parseval\_2d}
\end{theorem}
\begin{theorem}[\texttt{sum\_factor}]\label{GaussianField-sum-factor}
  \lean{GaussianField.sum_factor}\leanok
  \uses{GaussianField-FinLatticeSites}
  \texttt{GaussianField.sum\_factor}
\end{theorem}
\begin{theorem}[\texttt{dft\_eigencoeff\_massOperator\_2d}]\label{GaussianField-dft-eigencoeff-massOperator-2d}
  \lean{GaussianField.dft_eigencoeff_massOperator_2d}\leanok
  \uses{GaussianField-massOperator-product-eigenvector, GaussianField-massOperator-selfAdjoint, GaussianField-latticeEigenvalue1d, GaussianField-latticeFourierBasisFun, GaussianField-FinLatticeSites, GaussianField-massOperator, GaussianField-FinLatticeField}
  \texttt{GaussianField.dft\_eigencoeff\_massOperator\_2d}
\end{theorem}
\begin{theorem}[\texttt{dft\_parseval\_1d}]\label{GaussianField-dft-parseval-1d}
  \lean{GaussianField.dft_parseval_1d}\leanok
  \uses{GaussianField-latticeFourierBasisFun, GaussianField-latticeFourierNormSq, GaussianField-dft-expansion}
  \texttt{GaussianField.dft\_parseval\_1d}
\end{theorem}
\begin{theorem}[\texttt{massOperator\_product\_eigenvector}]\label{GaussianField-massOperator-product-eigenvector}
  \lean{GaussianField.massOperator_product_eigenvector}\leanok
  \uses{GaussianField-latticeEigenvalue1d, GaussianField-latticeFourierBasisFun, GaussianField-FinLatticeSites, GaussianField-massOperator, GaussianField-dft-1d-eigenvalue-pointwise, GaussianField-finiteLaplacian, GaussianField-FinLatticeField}
  \texttt{GaussianField.massOperator\_product\_eigenvector}
\end{theorem}
\begin{theorem}[\texttt{dft\_1d\_eigenvalue\_pointwise}]\label{GaussianField-dft-1d-eigenvalue-pointwise}
  \lean{GaussianField.dft_1d_eigenvalue_pointwise}\leanok
  \uses{GaussianField-latticeEigenvalue1d, GaussianField-massOperatorMatrix, GaussianField-latticeFourierBasisFun, GaussianField-FinLatticeSites, GaussianField-massOperator, GaussianField-massOperator-eq-matrix-mulVec, GaussianField-negLaplacian1d-dft-eigenvector, GaussianField-finiteLaplacian, GaussianField-FinLatticeField, Lattice-toSemilatticeInf, GaussianField-negLaplacianMatrix}
  \texttt{GaussianField.dft\_1d\_eigenvalue\_pointwise}
\end{theorem}
\begin{theorem}[\texttt{massOperator\_surjective\_2d}]\label{GaussianField-massOperator-surjective-2d}
  \lean{GaussianField.massOperator_surjective_2d}\leanok
  \uses{GaussianField-FinLatticeSites, GaussianField-massOperator, GaussianField-massOperator-pos-def, GaussianField-FinLatticeField}
  \texttt{GaussianField.massOperator\_surjective\_2d}
\end{theorem}
\section{\texttt{GaussianField.WickMultivariate}}
\begin{definition}[\texttt{gffPositionCovariance}]\label{GaussianField-gffPositionCovariance}
  \lean{GaussianField.gffPositionCovariance}\leanok
  \uses{GaussianField-FinLatticeSites}
  \texttt{GaussianField.gffPositionCovariance}
\end{definition}
\begin{theorem}[\texttt{siteWickMonomial\_eigenbasis\_expansion}]\label{GaussianField-siteWickMonomial-eigenbasis-expansion}
  \lean{GaussianField.siteWickMonomial_eigenbasis_expansion}\leanok
  \uses{GaussianField-gffMultiWickMonomial, GaussianField-FinLatticeSites, GaussianField-gffOrthonormalCoord, GaussianField-gffSiteVariance, GaussianField-Configuration, GaussianField-multiIndicesOfDegree, GaussianField-FinLatticeField, GaussianField-MultiIndexLattice-totalDegree}
  \texttt{GaussianField.siteWickMonomial\_eigenbasis\_expansion}
\end{theorem}
\begin{theorem}[\texttt{gff\_wickPower\_two\_smeared\_inner}]\label{GaussianField-gff-wickPower-two-smeared-inner}
  \lean{GaussianField.gff_wickPower_two_smeared_inner}\leanok
  \uses{GaussianField-gffMultiWickMonomial, GaussianField-FinLatticeSites, GaussianField-gffSmearedCoeff, GaussianField-latticeGaussianMeasure, GaussianField-Configuration, GaussianField-gffMultiWickMonomial-orthogonality, GaussianField-instMeasurableSpaceConfiguration, GaussianField-FinLatticeField, GaussianField-gffSmearedCovariance}
  \texttt{GaussianField.gff\_wickPower\_two\_smeared\_inner}
\end{theorem}
\begin{theorem}[\texttt{congr\_simp}]\label{GaussianField-gffMultiWickMonomial-congr-simp}
  \lean{GaussianField.gffMultiWickMonomial.congr_simp}\leanok
  \uses{GaussianField-gffMultiWickMonomial, GaussianField-FinLatticeSites, GaussianField-Configuration, GaussianField-FinLatticeField}
  \texttt{GaussianField.gffMultiWickMonomial.congr\_simp}
\end{theorem}
\begin{definition}[\texttt{gffSmearedCovariance}]\label{GaussianField-gffSmearedCovariance}
  \lean{GaussianField.gffSmearedCovariance}\leanok
  \uses{GaussianField-FinLatticeSites, GaussianField-gffSmearedCoeff, GaussianField-FinLatticeField}
  \texttt{GaussianField.gffSmearedCovariance}
\end{definition}
\begin{theorem}[\texttt{gffMultiWickMonomial\_orthogonality}]\label{GaussianField-gffMultiWickMonomial-orthogonality}
  \lean{GaussianField.gffMultiWickMonomial_orthogonality}\leanok
  \uses{GaussianField-gffMultiWickMonomial, GaussianField-wickMonomial-inner-gaussianReal-one, GaussianField-FinLatticeSites, GaussianField-gffOrthonormalProj-measurable, GaussianField-gffOrthonormalProj-pushforward-eq-stdGaussian, GaussianField-gffOrthonormalProj, GaussianField-gffOrthonormalCoord, GaussianField-latticeGaussianMeasure, GaussianField-Configuration, GaussianField-instMeasurableSpaceConfiguration, GaussianField-FinLatticeField}
  \texttt{GaussianField.gffMultiWickMonomial\_orthogonality}
\end{theorem}
\begin{theorem}[\texttt{congr\_simp}]\label{GaussianField-gffOrthonormalProj-congr-simp}
  \lean{GaussianField.gffOrthonormalProj.congr_simp}\leanok
  \uses{GaussianField-FinLatticeSites, GaussianField-gffOrthonormalProj, GaussianField-Configuration, GaussianField-FinLatticeField}
  \texttt{GaussianField.gffOrthonormalProj.congr\_simp}
\end{theorem}
\begin{theorem}[\texttt{gffSiteVariance\_eq\_covarianceGJ}]\label{GaussianField-gffSiteVariance-eq-covarianceGJ}
  \lean{GaussianField.gffSiteVariance_eq_covarianceGJ}\leanok
  \uses{GaussianField-covariance, GaussianField-FinLatticeSites, GaussianField-latticeCovarianceGJ, GaussianField-gffPositionCovariance-self, GaussianField-gffPositionCovariance, GaussianField-gffSiteVariance, GaussianField-gffPositionCovariance-eq-covarianceGJ, GaussianField-ell2-, GaussianField-FinLatticeField}
  \texttt{GaussianField.gffSiteVariance\_eq\_covarianceGJ}
\end{theorem}
\begin{theorem}[\texttt{multiWickMonomial\_pi\_gaussianReal\_inner}]\label{GaussianField-multiWickMonomial-pi-gaussianReal-inner}
  \lean{GaussianField.multiWickMonomial_pi_gaussianReal_inner}\leanok
  \uses{GaussianField-wickMonomial-inner-gaussianReal-one}
  \texttt{GaussianField.multiWickMonomial\_pi\_gaussianReal\_inner}
\end{theorem}
\begin{theorem}[\texttt{congr\_simp}]\label{GaussianField-gffPositionCovariance-congr-simp}
  \lean{GaussianField.gffPositionCovariance.congr_simp}\leanok
  \uses{GaussianField-FinLatticeSites, GaussianField-gffPositionCovariance}
  \texttt{GaussianField.gffPositionCovariance.congr\_simp}
\end{theorem}
\begin{theorem}[\texttt{congr\_simp}]\label{GaussianField-gffSmearedCovariance-congr-simp}
  \lean{GaussianField.gffSmearedCovariance.congr_simp}\leanok
  \uses{GaussianField-FinLatticeField, GaussianField-gffSmearedCovariance}
  \texttt{GaussianField.gffSmearedCovariance.congr\_simp}
\end{theorem}
\begin{definition}[\texttt{gffSiteVariance}]\label{GaussianField-gffSiteVariance}
  \lean{GaussianField.gffSiteVariance}\leanok
  \uses{GaussianField-FinLatticeSites, GaussianField-massEigenvalues, GaussianField-massEigenvectorBasis}
  \texttt{GaussianField.gffSiteVariance}
\end{definition}
\begin{theorem}[\texttt{gffSmearedCovariance\_self}]\label{GaussianField-gffSmearedCovariance-self}
  \lean{GaussianField.gffSmearedCovariance_self}\leanok
  \uses{GaussianField-FinLatticeSites, GaussianField-gffSmearedCoeff, GaussianField-FinLatticeField, GaussianField-gffSmearedCovariance}
  \texttt{GaussianField.gffSmearedCovariance\_self}
\end{theorem}
\begin{theorem}[\texttt{congr\_simp}]\label{GaussianField-gffSmearedCoeff-congr-simp}
  \lean{GaussianField.gffSmearedCoeff.congr_simp}\leanok
  \uses{GaussianField-FinLatticeSites, GaussianField-gffSmearedCoeff, GaussianField-FinLatticeField}
  \texttt{GaussianField.gffSmearedCoeff.congr\_simp}
\end{theorem}
\begin{definition}[\texttt{gffMultiWickMonomial}]\label{GaussianField-gffMultiWickMonomial}
  \lean{GaussianField.gffMultiWickMonomial}\leanok
  \uses{GaussianField-FinLatticeSites, GaussianField-gffOrthonormalCoord, GaussianField-Configuration, GaussianField-FinLatticeField}
  \texttt{GaussianField.gffMultiWickMonomial}
\end{definition}
\begin{theorem}[\texttt{gffPositionCovariance\_self}]\label{GaussianField-gffPositionCovariance-self}
  \lean{GaussianField.gffPositionCovariance_self}\leanok
  \uses{GaussianField-FinLatticeSites, GaussianField-gffPositionCovariance, GaussianField-gffSiteVariance}
  \texttt{GaussianField.gffPositionCovariance\_self}
\end{theorem}
\begin{theorem}[\texttt{gff\_wickPower\_two\_site\_inner}]\label{GaussianField-gff-wickPower-two-site-inner}
  \lean{GaussianField.gff_wickPower_two_site_inner}\leanok
  \uses{GaussianField-gffMultiWickMonomial, GaussianField-FinLatticeSites, GaussianField-gffPositionCovariance, GaussianField-gffSiteVariance, GaussianField-latticeGaussianMeasure, GaussianField-Configuration, GaussianField-gffMultiWickMonomial-orthogonality, GaussianField-instMeasurableSpaceConfiguration, GaussianField-FinLatticeField}
  \texttt{GaussianField.gff\_wickPower\_two\_site\_inner}
\end{theorem}
\begin{theorem}[\texttt{eigenbasis\_completeness}]\label{GaussianField-eigenbasis-completeness}
  \lean{GaussianField.eigenbasis_completeness}\leanok
  \uses{GaussianField-FinLatticeSites, GaussianField-massEigenvectorBasis}
  \texttt{GaussianField.eigenbasis\_completeness}
\end{theorem}
\begin{theorem}[\texttt{congr\_simp}]\label{GaussianField-gffSiteVariance-congr-simp}
  \lean{GaussianField.gffSiteVariance.congr_simp}\leanok
  \uses{GaussianField-FinLatticeSites, GaussianField-gffSiteVariance}
  \texttt{GaussianField.gffSiteVariance.congr\_simp}
\end{theorem}
\begin{theorem}[\texttt{integrable\_wickMonomial\_smeared\_mul}]\label{GaussianField-integrable-wickMonomial-smeared-mul}
  \lean{GaussianField.integrable_wickMonomial_smeared_mul}\leanok
  \uses{GaussianField-gffMultiWickMonomial, GaussianField-FinLatticeSites, GaussianField-gffSmearedCoeff, GaussianField-latticeGaussianMeasure, GaussianField-Configuration, GaussianField-instMeasurableSpaceConfiguration, GaussianField-FinLatticeField}
  \texttt{GaussianField.integrable\_wickMonomial\_smeared\_mul}
\end{theorem}
\begin{definition}[\texttt{multiIndicesUpToDegree}]\label{GaussianField-multiIndicesUpToDegree}
  \lean{GaussianField.multiIndicesUpToDegree}\leanok
  \uses{GaussianField-FinLatticeSites}
  \texttt{GaussianField.multiIndicesUpToDegree}
\end{definition}
\begin{theorem}[\texttt{wickMonomial\_inner\_gaussianReal\_one}]\label{GaussianField-wickMonomial-inner-gaussianReal-one}
  \lean{GaussianField.wickMonomial_inner_gaussianReal_one}\leanok
  \texttt{GaussianField.wickMonomial\_inner\_gaussianReal\_one}
\end{theorem}
\begin{theorem}[\texttt{gffPositionCovariance\_eq\_covarianceGJ}]\label{GaussianField-gffPositionCovariance-eq-covarianceGJ}
  \lean{GaussianField.gffPositionCovariance_eq_covarianceGJ}\leanok
  \uses{GaussianField-covariance, GaussianField-FinLatticeSites, GaussianField-latticeCovarianceGJ, GaussianField-lattice-covariance-GJ-eq-spectral, GaussianField-gffPositionCovariance, GaussianField-massEigenvalues, GaussianField-massEigenvectorBasis, GaussianField-massOperatorMatrix-eigenvalues-pos, GaussianField-ell2-, GaussianField-FinLatticeField}
  \texttt{GaussianField.gffPositionCovariance\_eq\_covarianceGJ}
\end{theorem}
\begin{definition}[\texttt{multiIndicesOfDegree}]\label{GaussianField-multiIndicesOfDegree}
  \lean{GaussianField.multiIndicesOfDegree}\leanok
  \uses{GaussianField-FinLatticeSites, GaussianField-multiIndicesUpToDegree, GaussianField-MultiIndexLattice-totalDegree}
  \texttt{GaussianField.multiIndicesOfDegree}
\end{definition}
\begin{theorem}[\texttt{gffSmearedCovariance\_single\_right}]\label{GaussianField-gffSmearedCovariance-single-right}
  \lean{GaussianField.gffSmearedCovariance_single_right}\leanok
  \uses{GaussianField-FinLatticeSites, GaussianField-gffPositionCovariance, GaussianField-gffSmearedCovariance-eq-sum-position, GaussianField-FinLatticeField, GaussianField-gffSmearedCovariance}
  \texttt{GaussianField.gffSmearedCovariance\_single\_right}
\end{theorem}
\begin{theorem}[\texttt{gffPositionCovariance\_symm}]\label{GaussianField-gffPositionCovariance-symm}
  \lean{GaussianField.gffPositionCovariance_symm}\leanok
  \uses{GaussianField-FinLatticeSites, GaussianField-gffPositionCovariance}
  \texttt{GaussianField.gffPositionCovariance\_symm}
\end{theorem}
\begin{definition}[\texttt{totalDegree}]\label{GaussianField-MultiIndexLattice-totalDegree}
  \lean{GaussianField.MultiIndexLattice.totalDegree}\leanok
  \uses{GaussianField-FinLatticeSites}
  \texttt{GaussianField.MultiIndexLattice.totalDegree}
\end{definition}
\begin{theorem}[\texttt{gffMultiWickMonomial\_eq\_hermite\_product}]\label{GaussianField-gffMultiWickMonomial-eq-hermite-product}
  \lean{GaussianField.gffMultiWickMonomial_eq_hermite_product}\leanok
  \uses{GaussianField-gffMultiWickMonomial, GaussianField-FinLatticeSites, GaussianField-gffOrthonormalCoord, GaussianField-Configuration, GaussianField-FinLatticeField}
  \texttt{GaussianField.gffMultiWickMonomial\_eq\_hermite\_product}
\end{theorem}
\begin{theorem}[\texttt{congr\_simp}]\label{GaussianField-gffOrthonormalCoord-congr-simp}
  \lean{GaussianField.gffOrthonormalCoord.congr_simp}\leanok
  \uses{GaussianField-FinLatticeSites, GaussianField-gffOrthonormalCoord, GaussianField-Configuration, GaussianField-FinLatticeField}
  \texttt{GaussianField.gffOrthonormalCoord.congr\_simp}
\end{theorem}
\begin{theorem}[\texttt{gffSmearedCovariance\_eq\_sum\_position}]\label{GaussianField-gffSmearedCovariance-eq-sum-position}
  \lean{GaussianField.gffSmearedCovariance_eq_sum_position}\leanok
  \uses{GaussianField-FinLatticeSites, GaussianField-gffSmearedCoeff, GaussianField-gffPositionCovariance, GaussianField-FinLatticeField, GaussianField-gffSmearedCovariance}
  \texttt{GaussianField.gffSmearedCovariance\_eq\_sum\_position}
\end{theorem}
\begin{definition}[\texttt{gffSmearedCoeff}]\label{GaussianField-gffSmearedCoeff}
  \lean{GaussianField.gffSmearedCoeff}\leanok
  \uses{GaussianField-FinLatticeSites, GaussianField-FinLatticeField}
  \texttt{GaussianField.gffSmearedCoeff}
\end{definition}
\section{\texttt{Lattice.CirculantDFT}}
\begin{theorem}[\texttt{sum\_sin\_zmod\_eq\_zero}]\label{GaussianField-sum-sin-zmod-eq-zero}
  \lean{GaussianField.sum_sin_zmod_eq_zero}\leanok
  \uses{GaussianField-sin-zmod-succ, GaussianField-sum-cos-zmod-eq-zero, Lattice-toSemilatticeInf}
  \texttt{GaussianField.sum\_sin\_zmod\_eq\_zero}
\end{theorem}
\begin{theorem}[\texttt{latticeFourierNormSq\_eventually\_one}]\label{GaussianField-latticeFourierNormSq-eventually-one}
  \lean{GaussianField.latticeFourierNormSq_eventually_one}\leanok
  \uses{GaussianField-latticeFourierNormSq-zero, GaussianField-latticeFourierNormSq, GaussianField-latticeFourierNormSq-eq-one, GaussianField-SmoothMap-Circle-fourierFreq}
  \texttt{GaussianField.latticeFourierNormSq\_eventually\_one}
\end{theorem}
\begin{theorem}[\texttt{eigenvector\_transfer}]\label{GaussianField-eigenvector-transfer}
  \lean{GaussianField.eigenvector_transfer}\leanok
  \uses{GaussianField-latticeEigenvalue1d, GaussianField-latticeFourierBasisFun, GaussianField-SmoothMap-Circle-instAddCommGroup, GaussianField-circleRestriction, GaussianField-eigenvalue-formula, GaussianField-negLaplacian1d-sin-eigenvector, GaussianField-cos-zmod-pred, GaussianField-SmoothMap-Circle-instModule, GaussianField-sin-zmod-succ, GaussianField-SmoothMap-Circle, GaussianField-latticeDFTCoeff1d, GaussianField-circleSpacing, GaussianField-negLaplacian1d-cos-eigenvector, GaussianField-sin-zmod-pred, GaussianField-cos-zmod-succ, GaussianField-SmoothMap-Circle-instTopologicalSpace, GaussianField-SmoothMap-Circle-fourierFreq, Lattice-toSemilatticeInf}
  \texttt{GaussianField.eigenvector\_transfer}
\end{theorem}
\begin{theorem}[\texttt{latticeEigenvalue1d\_ge\_quadratic}]\label{GaussianField-latticeEigenvalue1d-ge-quadratic}
  \lean{GaussianField.latticeEigenvalue1d_ge_quadratic}\leanok
  \uses{GaussianField-latticeEigenvalue1d, GaussianField-circleSpacing, GaussianField-SmoothMap-Circle-fourierFreq, Lattice-toSemilatticeInf}
  \texttt{GaussianField.latticeEigenvalue1d\_ge\_quadratic}
\end{theorem}
\begin{theorem}[\texttt{eigenvalue\_formula}]\label{GaussianField-eigenvalue-formula}
  \lean{GaussianField.eigenvalue_formula}\leanok
  \texttt{GaussianField.eigenvalue\_formula}
\end{theorem}
\begin{theorem}[\texttt{latticeFourierNormSq\_ge\_one}]\label{GaussianField-latticeFourierNormSq-ge-one}
  \lean{GaussianField.latticeFourierNormSq_ge_one}\leanok
  \uses{GaussianField-latticeFourierNormSq-zero, GaussianField-latticeFourierNormSq, GaussianField-latticeFourierNormSq-eq-one, GaussianField-SmoothMap-Circle-fourierFreq}
  \texttt{GaussianField.latticeFourierNormSq\_ge\_one}
\end{theorem}
\begin{theorem}[\texttt{negLaplacian\_restriction\_bound}]\label{GaussianField-negLaplacian-restriction-bound}
  \lean{GaussianField.negLaplacian_restriction_bound}\leanok
  \uses{GaussianField-SmoothMap-Circle-instFunLike, GaussianField-SmoothMap-Circle-periodic, GaussianField-symmetric-second-diff-bound, GaussianField-SmoothMap-Circle-instAddCommGroup, GaussianField-circleRestriction, GaussianField-SmoothMap-Circle-instSMul, GaussianField-SmoothMap-Circle-instModule, GaussianField-SmoothMap-Circle, GaussianField-circleSpacing, GaussianField-SmoothMap-Circle-sobolevSeminorm, GaussianField-SmoothMap-Circle-instTopologicalSpace, GaussianField-circlePoint, Lattice-toSemilatticeInf, GaussianField-circleSpacing-pos}
  \texttt{GaussianField.negLaplacian\_restriction\_bound}
\end{theorem}
\begin{theorem}[\texttt{latticeFourierNormSq\_pos}]\label{GaussianField-latticeFourierNormSq-pos}
  \lean{GaussianField.latticeFourierNormSq_pos}\leanok
  \uses{GaussianField-latticeFourierBasisFun, GaussianField-latticeFourierNormSq-zero, GaussianField-latticeFourierNormSq, GaussianField-latticeFourierNormSq-eq-one, GaussianField-SmoothMap-Circle-fourierFreq, Lattice-toSemilatticeInf}
  \texttt{GaussianField.latticeFourierNormSq\_pos}
\end{theorem}
\begin{theorem}[\texttt{negLaplacian1d\_dft\_eigenvector}]\label{GaussianField-negLaplacian1d-dft-eigenvector}
  \lean{GaussianField.negLaplacian1d_dft_eigenvector}\leanok
  \uses{GaussianField-latticeEigenvalue1d, GaussianField-latticeFourierBasisFun, GaussianField-negLaplacian1d-cos-eigenvalue, GaussianField-FinLatticeSites, GaussianField-eigenvalue-formula, GaussianField-negLaplacian1d-sin-eigenvalue, GaussianField-latticeEigenvalue1d-zero, GaussianField-SmoothMap-Circle-fourierFreq, Lattice-toSemilatticeInf, GaussianField-negLaplacianMatrix}
  \texttt{GaussianField.negLaplacian1d\_dft\_eigenvector}
\end{theorem}
\begin{definition}[\texttt{latticeFourierBasis}]\label{GaussianField-latticeFourierBasis}
  \lean{GaussianField.latticeFourierBasis}\leanok
  \uses{GaussianField-latticeFourierBasisFun, GaussianField-latticeFourierBasisFun-linearIndependent}
  \texttt{GaussianField.latticeFourierBasis}
\end{definition}
\begin{theorem}[\texttt{sin\_zmod\_pred}]\label{GaussianField-sin-zmod-pred}
  \lean{GaussianField.sin_zmod_pred}\leanok
  \uses{Lattice-toSemilatticeInf}
  \texttt{GaussianField.sin\_zmod\_pred}
\end{theorem}
\begin{theorem}[\texttt{latticeFourierBasisFun\_span\_eq\_top}]\label{GaussianField-latticeFourierBasisFun-span-eq-top}
  \lean{GaussianField.latticeFourierBasisFun_span_eq_top}\leanok
  \uses{GaussianField-latticeFourierBasisFun, GaussianField-latticeFourierBasisFun-linearIndependent}
  \texttt{GaussianField.latticeFourierBasisFun\_span\_eq\_top}
\end{theorem}
\begin{theorem}[\texttt{sin\_zmod\_succ}]\label{GaussianField-sin-zmod-succ}
  \lean{GaussianField.sin_zmod_succ}\leanok
  \uses{Lattice-toSemilatticeInf}
  \texttt{GaussianField.sin\_zmod\_succ}
\end{theorem}
\begin{theorem}[\texttt{symmetric\_second\_diff\_bound}]\label{GaussianField-symmetric-second-diff-bound}
  \lean{GaussianField.symmetric_second_diff_bound}\leanok
  \uses{Lattice-toSemilatticeSup, GaussianField-SmoothMap-Circle-instFunLike, GaussianField-SmoothMap-Circle-instAddCommGroup, GaussianField-SmoothMap-Circle-norm-iteratedDeriv-le-sobolevSeminorm-, GaussianField-SmoothMap-Circle-smooth, GaussianField-SmoothMap-Circle-instSMul, GaussianField-SmoothMap-Circle, GaussianField-SmoothMap-Circle-sobolevSeminorm, Lattice-toSemilatticeInf}
  \texttt{GaussianField.symmetric\_second\_diff\_bound}
\end{theorem}
\begin{theorem}[\texttt{sum\_sin\_mul\_sin\_zmod\_eq\_zero}]\label{GaussianField-sum-sin-mul-sin-zmod-eq-zero}
  \lean{GaussianField.sum_sin_mul_sin_zmod_eq_zero}\leanok
  \uses{GaussianField-sum-cos-zmod-eq-zero-of-pos-lt}
  \texttt{GaussianField.sum\_sin\_mul\_sin\_zmod\_eq\_zero}
\end{theorem}
\begin{theorem}[\texttt{sum\_cos\_zmod\_eq\_zero\_of\_pos\_lt}]\label{GaussianField-sum-cos-zmod-eq-zero-of-pos-lt}
  \lean{GaussianField.sum_cos_zmod_eq_zero_of_pos_lt}\leanok
  \uses{GaussianField-cos-ne-one-of-pos-lt, GaussianField-sum-cos-zmod-eq-zero}
  \texttt{GaussianField.sum\_cos\_zmod\_eq\_zero\_of\_pos\_lt}
\end{theorem}
\begin{theorem}[\texttt{cos\_zmod\_pred}]\label{GaussianField-cos-zmod-pred}
  \lean{GaussianField.cos_zmod_pred}\leanok
  \uses{Lattice-toSemilatticeInf}
  \texttt{GaussianField.cos\_zmod\_pred}
\end{theorem}
\begin{theorem}[\texttt{latticeHeatKernelBilinear1d\_eq\_spectral'}]\label{GaussianField-latticeHeatKernelBilinear1d-eq-spectral-}
  \lean{GaussianField.latticeHeatKernelBilinear1d_eq_spectral'}\leanok
  \uses{GaussianField-latticeEigenvalue1d, GaussianField-sum-fin1-eq-sum-zmod, GaussianField-latticeFourierBasisFun, GaussianField-SmoothMap-Circle-instAddCommGroup, GaussianField-FinLatticeSites, GaussianField-circleRestriction, GaussianField-Matrix-IsHermitian-mulVec-exp-of-eigenvector, GaussianField-latticeFourierNormSq, GaussianField-SmoothMap-Circle-instModule, GaussianField-SmoothMap-Circle, GaussianField-dft-expansion, GaussianField-latticeDFTCoeff1d, GaussianField-negLaplacian1d-dft-eigenvector, GaussianField-circleSpacing, GaussianField-latticeHeatKernelBilinear1d, GaussianField-SmoothMap-Circle-instTopologicalSpace, GaussianField-latticeHeatKernelMatrix, GaussianField-negLaplacianMatrix-isHermitian, GaussianField-circleSpacing-pos, GaussianField-negLaplacianMatrix}
  \texttt{GaussianField.latticeHeatKernelBilinear1d\_eq\_spectral'}
\end{theorem}
\begin{theorem}[\texttt{latticeFourierBasisFun\_orthogonal}]\label{GaussianField-latticeFourierBasisFun-orthogonal}
  \lean{GaussianField.latticeFourierBasisFun_orthogonal}\leanok
  \uses{GaussianField-sum-cos-zmod-eq-zero-of-pos-lt, GaussianField-sum-sin-mul-sin-zmod-eq-zero, GaussianField-sum-cos-mul-cos-zmod-eq-zero, GaussianField-latticeFourierBasisFun, GaussianField-sum-cos-mul-sin-zmod-eq-zero, GaussianField-sum-sin-zmod-eq-zero}
  \texttt{GaussianField.latticeFourierBasisFun\_orthogonal}
\end{theorem}
\begin{theorem}[\texttt{latticeFourierNormSq\_zero}]\label{GaussianField-latticeFourierNormSq-zero}
  \lean{GaussianField.latticeFourierNormSq_zero}\leanok
  \uses{GaussianField-latticeFourierBasisFun, GaussianField-latticeFourierNormSq}
  \texttt{GaussianField.latticeFourierNormSq\_zero}
\end{theorem}
\begin{theorem}[\texttt{lattice\_heatKernel\_tendsto\_continuum\_1d}]\label{GaussianField-lattice-heatKernel-tendsto-continuum-1d}
  \lean{GaussianField.lattice_heatKernel_tendsto_continuum_1d}\leanok
  \uses{GaussianField-latticeFourierNormSq-eventually-one, GaussianField-latticeDFTCoeff1d-zero-of-ge, GaussianField-latticeEigenvalue1d-tendsto-continuum, GaussianField-latticeEigenvalue1d, GaussianField-DyninMityaginSpace-coeff, GaussianField-latticeDFTCoeff1d-flat-bound, GaussianField-circleHasLaplacianEigenvalues, GaussianField-SmoothMap-Circle-instAddCommGroup, GaussianField-SmoothMap-Circle-instSMul, GaussianField-latticeDFTCoeff1d-tendsto, GaussianField-latticeFourierNormSq-zero, GaussianField-HasLaplacianEigenvalues-eigenvalue, GaussianField-latticeFourierNormSq-ge-one, GaussianField-latticeFourierNormSq, GaussianField-SmoothMap-Circle-instModule, GaussianField-latticeFourierNormSq-pos, GaussianField-SmoothMap-Circle-instContinuousSMul, GaussianField-smoothCircle-dyninMityaginSpace, GaussianField-SmoothMap-Circle-instIsTopologicalAddGroup, GaussianField-latticeEigenvalue1d-zero, GaussianField-SmoothMap-Circle, GaussianField-latticeDFTCoeff1d, GaussianField-heatKernelBilinear, GaussianField-circleSpacing, GaussianField-SmoothMap-Circle-sobolevSeminorm-nonneg, GaussianField-latticeHeatKernelBilinear1d, GaussianField-SmoothMap-Circle-sobolevSeminorm, GaussianField-latticeEigenvalue1d-ge-8m, GaussianField-SmoothMap-Circle-instTopologicalSpace, Lattice-toSemilatticeInf}
  \texttt{GaussianField.lattice\_heatKernel\_tendsto\_continuum\_1d}
\end{theorem}
\begin{theorem}[\texttt{latticeDFTCoeff1d\_quadratic\_bound'}]\label{GaussianField-latticeDFTCoeff1d-quadratic-bound-}
  \lean{GaussianField.latticeDFTCoeff1d_quadratic_bound'}\leanok
  \uses{GaussianField-latticeDFTCoeff1d-zero-of-ge, GaussianField-latticeEigenvalue1d, GaussianField-latticeFourierBasisFun, GaussianField-laplacianDFT-flat-bound, GaussianField-latticeDFTCoeff1d-flat-bound, GaussianField-SmoothMap-Circle-instAddCommGroup, GaussianField-circleRestriction, GaussianField-SmoothMap-Circle-instSMul, GaussianField-latticeEigenvalue1d-ge-quadratic, GaussianField-eigenvector-transfer, GaussianField-SmoothMap-Circle-instModule, GaussianField-SmoothMap-Circle, GaussianField-latticeDFTCoeff1d, GaussianField-circleSpacing, GaussianField-SmoothMap-Circle-sobolevSeminorm-nonneg, GaussianField-SmoothMap-Circle-sobolevSeminorm, GaussianField-SmoothMap-Circle-instTopologicalSpace}
  \texttt{GaussianField.latticeDFTCoeff1d\_quadratic\_bound'}
\end{theorem}
\begin{theorem}[\texttt{laplacianDFT\_flat\_bound}]\label{GaussianField-laplacianDFT-flat-bound}
  \lean{GaussianField.laplacianDFT_flat_bound}\leanok
  \uses{GaussianField-latticeFourierBasisFun, GaussianField-SmoothMap-Circle-instAddCommGroup, GaussianField-circleRestriction, GaussianField-SmoothMap-Circle-instSMul, GaussianField-circleSpacing-eq, GaussianField-SmoothMap-Circle-instModule, GaussianField-SmoothMap-Circle, GaussianField-circleSpacing, GaussianField-SmoothMap-Circle-sobolevSeminorm-nonneg, GaussianField-SmoothMap-Circle-sobolevSeminorm, GaussianField-negLaplacian-restriction-bound, GaussianField-SmoothMap-Circle-instTopologicalSpace, Lattice-toSemilatticeInf, GaussianField-circleSpacing-pos}
  \texttt{GaussianField.laplacianDFT\_flat\_bound}
\end{theorem}
\begin{theorem}[\texttt{sum\_cos\_zmod\_eq\_N}]\label{GaussianField-sum-cos-zmod-eq-N}
  \lean{GaussianField.sum_cos_zmod_eq_N}\leanok
  \texttt{GaussianField.sum\_cos\_zmod\_eq\_N}
\end{theorem}
\begin{theorem}[\texttt{negLaplacian1d\_sin\_eigenvector}]\label{GaussianField-negLaplacian1d-sin-eigenvector}
  \lean{GaussianField.negLaplacian1d_sin_eigenvector}\leanok
  \uses{Lattice-toSemilatticeInf}
  \texttt{GaussianField.negLaplacian1d\_sin\_eigenvector}
\end{theorem}
\begin{theorem}[\texttt{sum\_sin\_zmod\_int\_eq\_zero}]\label{GaussianField-sum-sin-zmod-int-eq-zero}
  \lean{GaussianField.sum_sin_zmod_int_eq_zero}\leanok
  \uses{GaussianField-sum-sin-zmod-eq-zero}
  \texttt{GaussianField.sum\_sin\_zmod\_int\_eq\_zero}
\end{theorem}
\begin{theorem}[\texttt{dft\_expansion}]\label{GaussianField-dft-expansion}
  \lean{GaussianField.dft_expansion}\leanok
  \uses{GaussianField-latticeFourierBasisFun, GaussianField-latticeFourierBasisFun-orthogonal, GaussianField-latticeFourierNormSq, GaussianField-latticeFourierNormSq-pos, GaussianField-latticeFourierBasisFun-linearIndependent, GaussianField-latticeFourierBasisFun-span-eq-top}
  \texttt{GaussianField.dft\_expansion}
\end{theorem}
\begin{theorem}[\texttt{negLaplacian1d\_cos\_eigenvalue}]\label{GaussianField-negLaplacian1d-cos-eigenvalue}
  \lean{GaussianField.negLaplacian1d_cos_eigenvalue}\leanok
  \uses{GaussianField-massOperatorMatrix, GaussianField-FinLatticeSites, GaussianField-massOperator, GaussianField-massOperator-eq-matrix-mulVec, GaussianField-finiteLaplacian, GaussianField-FinLatticeField, GaussianField-negLaplacianMatrix}
  \texttt{GaussianField.negLaplacian1d\_cos\_eigenvalue}
\end{theorem}
\begin{theorem}[\texttt{latticeFourierBasisFun\_linearIndependent}]\label{GaussianField-latticeFourierBasisFun-linearIndependent}
  \lean{GaussianField.latticeFourierBasisFun_linearIndependent}\leanok
  \uses{GaussianField-latticeFourierBasisFun, GaussianField-latticeFourierBasisFun-orthogonal, GaussianField-latticeFourierNormSq, GaussianField-latticeFourierNormSq-pos}
  \texttt{GaussianField.latticeFourierBasisFun\_linearIndependent}
\end{theorem}
\begin{theorem}[\texttt{cos\_ne\_one\_of\_pos\_lt}]\label{GaussianField-cos-ne-one-of-pos-lt}
  \lean{GaussianField.cos_ne_one_of_pos_lt}\leanok
  \texttt{GaussianField.cos\_ne\_one\_of\_pos\_lt}
\end{theorem}
\begin{theorem}[\texttt{sum\_cos\_zmod\_eq\_zero}]\label{GaussianField-sum-cos-zmod-eq-zero}
  \lean{GaussianField.sum_cos_zmod_eq_zero}\leanok
  \uses{GaussianField-sin-zmod-succ, GaussianField-cos-zmod-succ, Lattice-toSemilatticeInf}
  \texttt{GaussianField.sum\_cos\_zmod\_eq\_zero}
\end{theorem}
\begin{definition}[\texttt{zmodToFin1}]\label{GaussianField-zmodToFin1}
  \lean{GaussianField.zmodToFin1}\leanok
  \uses{GaussianField-FinLatticeSites}
  \texttt{GaussianField.zmodToFin1}
\end{definition}
\begin{theorem}[\texttt{latticeFourierNormSq\_eq\_one}]\label{GaussianField-latticeFourierNormSq-eq-one}
  \lean{GaussianField.latticeFourierNormSq_eq_one}\leanok
  \uses{GaussianField-latticeFourierBasisFun, GaussianField-latticeFourierNormSq, GaussianField-sum-cos-zmod-eq-zero, GaussianField-SmoothMap-Circle-fourierFreq}
  \texttt{GaussianField.latticeFourierNormSq\_eq\_one}
\end{theorem}
\begin{theorem}[\texttt{sum\_cos\_mul\_sin\_zmod\_eq\_zero}]\label{GaussianField-sum-cos-mul-sin-zmod-eq-zero}
  \lean{GaussianField.sum_cos_mul_sin_zmod_eq_zero}\leanok
  \uses{GaussianField-sum-sin-zmod-int-eq-zero}
  \texttt{GaussianField.sum\_cos\_mul\_sin\_zmod\_eq\_zero}
\end{theorem}
\begin{theorem}[\texttt{negLaplacian1d\_sin\_eigenvalue}]\label{GaussianField-negLaplacian1d-sin-eigenvalue}
  \lean{GaussianField.negLaplacian1d_sin_eigenvalue}\leanok
  \uses{GaussianField-massOperatorMatrix, GaussianField-FinLatticeSites, GaussianField-massOperator, GaussianField-massOperator-eq-matrix-mulVec, GaussianField-finiteLaplacian, GaussianField-FinLatticeField, GaussianField-negLaplacianMatrix}
  \texttt{GaussianField.negLaplacian1d\_sin\_eigenvalue}
\end{theorem}
\begin{theorem}[\texttt{congr\_simp}]\label{GaussianField-latticeFourierNormSq-congr-simp}
  \lean{GaussianField.latticeFourierNormSq.congr_simp}\leanok
  \uses{GaussianField-latticeFourierNormSq}
  \texttt{GaussianField.latticeFourierNormSq.congr\_simp}
\end{theorem}
\begin{definition}[\texttt{latticeFourierNormSq}]\label{GaussianField-latticeFourierNormSq}
  \lean{GaussianField.latticeFourierNormSq}\leanok
  \uses{GaussianField-latticeFourierBasisFun}
  \texttt{GaussianField.latticeFourierNormSq}
\end{definition}
\begin{theorem}[\texttt{congr\_simp}]\label{GaussianField-latticeDFTCoeff1d-congr-simp}
  \lean{GaussianField.latticeDFTCoeff1d.congr_simp}\leanok
  \uses{GaussianField-SmoothMap-Circle, GaussianField-latticeDFTCoeff1d}
  \texttt{GaussianField.latticeDFTCoeff1d.congr\_simp}
\end{theorem}
\begin{theorem}[\texttt{negLaplacian1d\_cos\_eigenvector}]\label{GaussianField-negLaplacian1d-cos-eigenvector}
  \lean{GaussianField.negLaplacian1d_cos_eigenvector}\leanok
  \uses{Lattice-toSemilatticeInf}
  \texttt{GaussianField.negLaplacian1d\_cos\_eigenvector}
\end{theorem}
\begin{theorem}[\texttt{sum\_cos\_mul\_cos\_zmod\_eq\_zero}]\label{GaussianField-sum-cos-mul-cos-zmod-eq-zero}
  \lean{GaussianField.sum_cos_mul_cos_zmod_eq_zero}\leanok
  \uses{GaussianField-sum-cos-zmod-eq-zero-of-pos-lt}
  \texttt{GaussianField.sum\_cos\_mul\_cos\_zmod\_eq\_zero}
\end{theorem}
\begin{theorem}[\texttt{congr\_simp}]\label{GaussianField-latticeHeatKernelBilinear1d-congr-simp}
  \lean{GaussianField.latticeHeatKernelBilinear1d.congr_simp}\leanok
  \uses{GaussianField-SmoothMap-Circle, GaussianField-latticeHeatKernelBilinear1d}
  \texttt{GaussianField.latticeHeatKernelBilinear1d.congr\_simp}
\end{theorem}
\begin{theorem}[\texttt{sum\_fin1\_eq\_sum\_zmod}]\label{GaussianField-sum-fin1-eq-sum-zmod}
  \lean{GaussianField.sum_fin1_eq_sum_zmod}\leanok
  \uses{GaussianField-FinLatticeSites, GaussianField-zmodToFin1}
  \texttt{GaussianField.sum\_fin1\_eq\_sum\_zmod}
\end{theorem}
\begin{theorem}[\texttt{cos\_zmod\_succ}]\label{GaussianField-cos-zmod-succ}
  \lean{GaussianField.cos_zmod_succ}\leanok
  \uses{Lattice-toSemilatticeInf}
  \texttt{GaussianField.cos\_zmod\_succ}
\end{theorem}
\section{\texttt{Torus.Symmetry}}
\begin{theorem}[\texttt{circleTranslation\_add}]\label{GaussianField-circleTranslation-add}
  \lean{GaussianField.circleTranslation_add}\leanok
  \uses{GaussianField-circleTranslation, GaussianField-SmoothMap-Circle-instFunLike, GaussianField-SmoothMap-Circle-ext, GaussianField-SmoothMap-Circle-mk-congr-simp, GaussianField-SmoothMap-Circle-instAddCommGroup, GaussianField-SmoothMap-Circle-instModule, GaussianField-SmoothMap-Circle, GaussianField-SmoothMap-Circle-instTopologicalSpace}
  \texttt{GaussianField.circleTranslation\_add}
\end{theorem}
\begin{theorem}[\texttt{torusConfigReflection\_continuous}]\label{GaussianField-torusConfigReflection-continuous}
  \lean{GaussianField.torusConfigReflection_continuous}\leanok
  \uses{GaussianField-torusTimeReflection, GaussianField-NuclearTensorProduct-instModuleReal, GaussianField-NuclearTensorProduct-instTopologicalSpace, GaussianField-NuclearTensorProduct-instIsTopologicalAddGroup, GaussianField-NuclearTensorProduct-instAddCommGroup, GaussianField-Configuration, GaussianField-NuclearTensorProduct-instContinuousSMulReal, GaussianField-SmoothMap-Circle, GaussianField-torusConfigReflection, GaussianField-TorusTestFunction}
  \texttt{GaussianField.torusConfigReflection\_continuous}
\end{theorem}
\begin{definition}[\texttt{torusConfigTranslation}]\label{GaussianField-torusConfigTranslation}
  \lean{GaussianField.torusConfigTranslation}\leanok
  \uses{GaussianField-NuclearTensorProduct-instModuleReal, GaussianField-NuclearTensorProduct-instTopologicalSpace, GaussianField-NuclearTensorProduct-instIsTopologicalAddGroup, GaussianField-NuclearTensorProduct-instAddCommGroup, GaussianField-torusTranslation, GaussianField-Configuration, GaussianField-NuclearTensorProduct-instContinuousSMulReal, GaussianField-SmoothMap-Circle, GaussianField-TorusTestFunction}
  \texttt{GaussianField.torusConfigTranslation}
\end{definition}
\begin{definition}[\texttt{torusSwap}]\label{GaussianField-torusSwap}
  \lean{GaussianField.torusSwap}\leanok
  \uses{GaussianField-NuclearTensorProduct-instModuleReal, GaussianField-NuclearTensorProduct-instTopologicalSpace, GaussianField-SmoothMap-Circle-instAddCommGroup, GaussianField-nuclearTensorProduct-swapCLM, GaussianField-NuclearTensorProduct-instAddCommGroup, GaussianField-SmoothMap-Circle-instModule, GaussianField-SmoothMap-Circle-instContinuousSMul, GaussianField-smoothCircle-dyninMityaginSpace, GaussianField-SmoothMap-Circle-instIsTopologicalAddGroup, GaussianField-SmoothMap-Circle, GaussianField-TorusTestFunction, GaussianField-SmoothMap-Circle-instTopologicalSpace}
  \texttt{GaussianField.torusSwap}
\end{definition}
\begin{theorem}[\texttt{torusTimeReflection\_involution}]\label{GaussianField-torusTimeReflection-involution}
  \lean{GaussianField.torusTimeReflection_involution}\leanok
  \uses{GaussianField-torusTimeReflection, GaussianField-NuclearTensorProduct-instModuleReal, GaussianField-NuclearTensorProduct-instTopologicalSpace, GaussianField-nuclearTensorProduct-mapCLM-comp, GaussianField-nuclearTensorProduct-mapCLM-id, GaussianField-nuclearTensorProduct-mapCLM, GaussianField-SmoothMap-Circle-instAddCommGroup, GaussianField-NuclearTensorProduct, GaussianField-NuclearTensorProduct-instAddCommGroup, GaussianField-smoothCircle-hasBiorthogonalBasis, GaussianField-circleReflection, GaussianField-SmoothMap-Circle-instModule, GaussianField-SmoothMap-Circle-instContinuousSMul, GaussianField-smoothCircle-dyninMityaginSpace, GaussianField-SmoothMap-Circle-instIsTopologicalAddGroup, GaussianField-SmoothMap-Circle, GaussianField-circleReflection-involution, GaussianField-TorusTestFunction, GaussianField-SmoothMap-Circle-instTopologicalSpace}
  \texttt{GaussianField.torusTimeReflection\_involution}
\end{theorem}
\begin{definition}[\texttt{circleReflection}]\label{GaussianField-circleReflection}
  \lean{GaussianField.circleReflection}\leanok
  \uses{GaussianField-SmoothMap-Circle-instFunLike, GaussianField-SmoothMap-Circle-instAddCommGroup, GaussianField-SmoothMap-Circle-instModule, GaussianField-SmoothMap-Circle, GaussianField-SmoothMap-Circle-instTopologicalSpace}
  \texttt{GaussianField.circleReflection}
\end{definition}
\begin{theorem}[\texttt{circleReflection\_involution}]\label{GaussianField-circleReflection-involution}
  \lean{GaussianField.circleReflection_involution}\leanok
  \uses{GaussianField-SmoothMap-Circle-instFunLike, GaussianField-SmoothMap-Circle-ext, GaussianField-SmoothMap-Circle-mk-congr-simp, GaussianField-SmoothMap-Circle-instAddCommGroup, GaussianField-circleReflection, GaussianField-SmoothMap-Circle-instModule, GaussianField-SmoothMap-Circle, GaussianField-SmoothMap-Circle-instTopologicalSpace}
  \texttt{GaussianField.circleReflection\_involution}
\end{theorem}
\begin{theorem}[\texttt{circleTranslation\_zero}]\label{GaussianField-circleTranslation-zero}
  \lean{GaussianField.circleTranslation_zero}\leanok
  \uses{GaussianField-circleTranslation, GaussianField-SmoothMap-Circle-instFunLike, GaussianField-SmoothMap-Circle-ext, GaussianField-SmoothMap-Circle-mk-congr-simp, GaussianField-SmoothMap-Circle-instAddCommGroup, GaussianField-SmoothMap-Circle-instModule, GaussianField-SmoothMap-Circle, GaussianField-SmoothMap-Circle-instTopologicalSpace}
  \texttt{GaussianField.circleTranslation\_zero}
\end{theorem}
\begin{theorem}[\texttt{torusConfigReflection\_involution}]\label{GaussianField-torusConfigReflection-involution}
  \lean{GaussianField.torusConfigReflection_involution}\leanok
  \uses{GaussianField-torusTimeReflection, GaussianField-NuclearTensorProduct-instModuleReal, GaussianField-NuclearTensorProduct-instTopologicalSpace, GaussianField-torusTimeReflection-involution, GaussianField-NuclearTensorProduct-instIsTopologicalAddGroup, GaussianField-NuclearTensorProduct-instAddCommGroup, GaussianField-Configuration, GaussianField-NuclearTensorProduct-instContinuousSMulReal, GaussianField-SmoothMap-Circle, GaussianField-torusConfigReflection, GaussianField-TorusTestFunction}
  \texttt{GaussianField.torusConfigReflection\_involution}
\end{theorem}
\begin{theorem}[\texttt{torusConfigTranslation\_continuous}]\label{GaussianField-torusConfigTranslation-continuous}
  \lean{GaussianField.torusConfigTranslation_continuous}\leanok
  \uses{GaussianField-NuclearTensorProduct-instModuleReal, GaussianField-NuclearTensorProduct-instTopologicalSpace, GaussianField-NuclearTensorProduct-instIsTopologicalAddGroup, GaussianField-NuclearTensorProduct-instAddCommGroup, GaussianField-torusConfigTranslation, GaussianField-torusTranslation, GaussianField-Configuration, GaussianField-NuclearTensorProduct-instContinuousSMulReal, GaussianField-SmoothMap-Circle, GaussianField-TorusTestFunction}
  \texttt{GaussianField.torusConfigTranslation\_continuous}
\end{theorem}
\begin{definition}[\texttt{circleTranslation}]\label{GaussianField-circleTranslation}
  \lean{GaussianField.circleTranslation}\leanok
  \uses{GaussianField-SmoothMap-Circle-instFunLike, GaussianField-SmoothMap-Circle-instAddCommGroup, GaussianField-SmoothMap-Circle-instModule, GaussianField-SmoothMap-Circle, GaussianField-SmoothMap-Circle-instTopologicalSpace}
  \texttt{GaussianField.circleTranslation}
\end{definition}
\begin{definition}[\texttt{torusTranslation}]\label{GaussianField-torusTranslation}
  \lean{GaussianField.torusTranslation}\leanok
  \uses{GaussianField-NuclearTensorProduct-instModuleReal, GaussianField-circleTranslation, GaussianField-NuclearTensorProduct-instTopologicalSpace, GaussianField-nuclearTensorProduct-mapCLM, GaussianField-SmoothMap-Circle-instAddCommGroup, GaussianField-NuclearTensorProduct-instAddCommGroup, GaussianField-SmoothMap-Circle-instModule, GaussianField-SmoothMap-Circle-instContinuousSMul, GaussianField-smoothCircle-dyninMityaginSpace, GaussianField-SmoothMap-Circle-instIsTopologicalAddGroup, GaussianField-SmoothMap-Circle, GaussianField-TorusTestFunction, GaussianField-SmoothMap-Circle-instTopologicalSpace}
  \texttt{GaussianField.torusTranslation}
\end{definition}
\begin{definition}[\texttt{torusConfigReflection}]\label{GaussianField-torusConfigReflection}
  \lean{GaussianField.torusConfigReflection}\leanok
  \uses{GaussianField-torusTimeReflection, GaussianField-NuclearTensorProduct-instModuleReal, GaussianField-NuclearTensorProduct-instTopologicalSpace, GaussianField-NuclearTensorProduct-instIsTopologicalAddGroup, GaussianField-NuclearTensorProduct-instAddCommGroup, GaussianField-Configuration, GaussianField-NuclearTensorProduct-instContinuousSMulReal, GaussianField-SmoothMap-Circle, GaussianField-TorusTestFunction}
  \texttt{GaussianField.torusConfigReflection}
\end{definition}
\begin{definition}[\texttt{torusTimeReflection}]\label{GaussianField-torusTimeReflection}
  \lean{GaussianField.torusTimeReflection}\leanok
  \uses{GaussianField-NuclearTensorProduct-instModuleReal, GaussianField-NuclearTensorProduct-instTopologicalSpace, GaussianField-nuclearTensorProduct-mapCLM, GaussianField-SmoothMap-Circle-instAddCommGroup, GaussianField-NuclearTensorProduct-instAddCommGroup, GaussianField-circleReflection, GaussianField-SmoothMap-Circle-instModule, GaussianField-SmoothMap-Circle-instContinuousSMul, GaussianField-smoothCircle-dyninMityaginSpace, GaussianField-SmoothMap-Circle-instIsTopologicalAddGroup, GaussianField-SmoothMap-Circle, GaussianField-TorusTestFunction, GaussianField-SmoothMap-Circle-instTopologicalSpace}
  \texttt{GaussianField.torusTimeReflection}
\end{definition}
\section{\texttt{Lattice.AsymFiniteField}}
\begin{theorem}[\texttt{congr\_simp}]\label{GaussianField-asymLatticeCoeffCLM-congr-simp}
  \lean{GaussianField.asymLatticeCoeffCLM.congr_simp}\leanok
  \uses{GaussianField-AsymLatticeSites, GaussianField-asymLatticeCoeffCLM, GaussianField-AsymLatticeField}
  \texttt{GaussianField.asymLatticeCoeffCLM.congr\_simp}
\end{theorem}
\begin{theorem}[\texttt{congr\_simp}]\label{GaussianField-asymLatticeSeminorm-congr-simp}
  \lean{GaussianField.asymLatticeSeminorm.congr_simp}\leanok
  \uses{GaussianField-AsymLatticeSites, GaussianField-AsymLatticeField, GaussianField-asymLatticeSeminorm}
  \texttt{GaussianField.asymLatticeSeminorm.congr\_simp}
\end{theorem}
\begin{theorem}[\texttt{congr\_simp}]\label{GaussianField-asymLatticeBasisVec-congr-simp}
  \lean{GaussianField.asymLatticeBasisVec.congr_simp}\leanok
  \uses{GaussianField-AsymLatticeSites, GaussianField-asymLatticeBasisVec}
  \texttt{GaussianField.asymLatticeBasisVec.congr\_simp}
\end{theorem}
\begin{definition}[\texttt{asymLatticeCoeffCLM}]\label{GaussianField-asymLatticeCoeffCLM}
  \lean{GaussianField.asymLatticeCoeffCLM}\leanok
  \uses{GaussianField-AsymLatticeSites, GaussianField-AsymLatticeField, GaussianField-instFintypeAsymLatticeSitesOfNeZeroNat}
  \texttt{GaussianField.asymLatticeCoeffCLM}
\end{definition}
\begin{definition}[\texttt{asymLatticeSeminorm}]\label{GaussianField-asymLatticeSeminorm}
  \lean{GaussianField.asymLatticeSeminorm}\leanok
  \uses{GaussianField-AsymLatticeSites, GaussianField-AsymLatticeField, GaussianField-instFintypeAsymLatticeSitesOfNeZeroNat}
  \texttt{GaussianField.asymLatticeSeminorm}
\end{definition}
\begin{definition}[\texttt{asymLatticeBasisVec}]\label{GaussianField-asymLatticeBasisVec}
  \lean{GaussianField.asymLatticeBasisVec}\leanok
  \uses{GaussianField-AsymLatticeSites, GaussianField-AsymLatticeField, GaussianField-instFintypeAsymLatticeSitesOfNeZeroNat, GaussianField-asymLatticeDelta}
  \texttt{GaussianField.asymLatticeBasisVec}
\end{definition}
\begin{definition}[\texttt{asymLatticeField\_dyninMityaginSpace}]\label{GaussianField-asymLatticeField-dyninMityaginSpace}
  \lean{GaussianField.asymLatticeField_dyninMityaginSpace}\leanok
  \uses{GaussianField-DyninMityaginSpace, GaussianField-AsymLatticeSites, GaussianField-asymLatticeCoeffCLM, GaussianField-asymLatticeBasisVec, GaussianField-AsymLatticeField, GaussianField-asymLatticeSeminorm}
  \texttt{GaussianField.asymLatticeField\_dyninMityaginSpace}
\end{definition}
\section{\texttt{Pphi2.TorusContinuumLimit.MeasureUniqueness}}
\begin{theorem}[\texttt{measure\_eq\_of\_moments}]\label{GaussianField-measure-eq-of-moments}
  \lean{GaussianField.measure_eq_of_moments}\leanok
  \uses{GaussianField-DyninMityaginSpace, GaussianField-eval-map-eq-of-moments, GaussianField-pushforward-eq-of-eval-eq, GaussianField-configBasisEval-measurable, GaussianField-Configuration, GaussianField-instMeasurableSpaceConfiguration, GaussianField-instMeasurableSpaceConfiguration-eq-comap, GaussianField-configBasisEval}
  \texttt{GaussianField.measure\_eq\_of\_moments}
\end{theorem}
\begin{theorem}[\texttt{weakLimit\_centered\_gaussianReal}]\label{GaussianField-weakLimit-centered-gaussianReal}
  \lean{GaussianField.weakLimit_centered_gaussianReal}\leanok
  \uses{Lattice-toSemilatticeSup, Lattice-toSemilatticeInf}
  The 1D marginals of the weak limit equal $N(0, G_L(f,f))$ for each test function $f$.
\end{theorem}
\begin{theorem}[\texttt{gaussian\_measure\_unique\_of\_covariance}]\label{GaussianField-gaussian-measure-unique-of-covariance}
  \lean{GaussianField.gaussian_measure_unique_of_covariance}\leanok
  \uses{GaussianField-DyninMityaginSpace, GaussianField-pushforward-eq-of-eval-eq, GaussianField-eval-map-eq-of-covariance, GaussianField-configBasisEval-measurable, GaussianField-Configuration, GaussianField-instMeasurableSpaceConfiguration, GaussianField-instMeasurableSpaceConfiguration-eq-comap, GaussianField-configBasisEval}
  Main theorem: Two centered Gaussian probability measures on Configuration($E$) with the same covariance are equal. Proved by pulling back from $\mathbb{R}^{\mathbb{N}}$ using \texttt{instMeasurableSpaceConfiguration\_eq\_comap}.
\end{theorem}
\begin{theorem}[\texttt{mgf\_at\_scaled}]\label{GaussianField-mgf-at-scaled}
  \lean{GaussianField.mgf_at_scaled}\leanok
  \uses{GaussianField-Configuration, GaussianField-instMeasurableSpaceConfiguration}
  For a Gaussian measure, $\int e^{t \omega(f)}\, d\mu = \exp(t^2 \sigma^2(f)/2)$ where $\sigma^2(f) = \int (\omega f)^2\, d\mu$.
\end{theorem}
\begin{theorem}[\texttt{eval\_map\_eq\_of\_covariance}]\label{GaussianField-eval-map-eq-of-covariance}
  \lean{GaussianField.eval_map_eq_of_covariance}\leanok
  \uses{GaussianField-configuration-eval-measurable, GaussianField-Configuration, GaussianField-mgf-at-scaled, GaussianField-instMeasurableSpaceConfiguration, Lattice-toSemilatticeInf}
  Same covariance implies same 1D marginals for all test functions. Proved by matching real MGFs to \texttt{gaussianReal}, analytically continuing to complex MGFs via \texttt{eqOn\_complexMGF\_of\_mgf}, and applying \texttt{Measure.ext\_of\_complexMGF\_eq}.
\end{theorem}
\begin{theorem}[\texttt{integrable\_exp\_smul\_of\_integrable\_exp\_sq}]\label{GaussianField-integrable-exp-smul-of-integrable-exp-sq}
  \lean{GaussianField.integrable_exp_smul_of_integrable_exp_sq}\leanok
  \uses{GaussianField-configuration-eval-measurable, GaussianField-Configuration, GaussianField-instMeasurableSpaceConfiguration, Lattice-toSemilatticeInf}
  \texttt{GaussianField.integrable\_exp\_smul\_of\_integrable\_exp\_sq}
\end{theorem}
\begin{theorem}[\texttt{pushforward\_eq\_of\_eval\_eq}]\label{GaussianField-pushforward-eq-of-eval-eq}
  \lean{GaussianField.pushforward_eq_of_eval_eq}\leanok
  \uses{GaussianField-DyninMityaginSpace, GaussianField-configBasisEval-measurable, GaussianField-configuration-eval-measurable, GaussianField-Configuration, GaussianField-DyninMityaginSpace-basis, GaussianField-instMeasurableSpaceConfiguration, GaussianField-configBasisEval}
  Equal 1D marginals for all $f : E$ imply equal pushforward measures on $\mathbb{R}^{\mathbb{N}}$ via \texttt{configBasisEval}. Proved by Cramer-Wold (\texttt{ext\_of\_charFunDual} for finite-dim marginals) + Kolmogorov uniqueness (\texttt{IsProjectiveLimit.unique}).
\end{theorem}
\begin{theorem}[\texttt{mgf\_eq\_of\_covariance}]\label{GaussianField-mgf-eq-of-covariance}
  \lean{GaussianField.mgf_eq_of_covariance}\leanok
  \uses{GaussianField-Configuration, GaussianField-instMeasurableSpaceConfiguration}
  If two Gaussian measures have the same covariance, their moment generating functions agree on all test functions: $\int e^{\omega(f)}\, d\mu_1 = \int e^{\omega(f)}\, d\mu_2$.
\end{theorem}
\begin{theorem}[\texttt{eval\_map\_eq\_gaussianReal}]\label{GaussianField-eval-map-eq-gaussianReal}
  \lean{GaussianField.eval_map_eq_gaussianReal}\leanok
  \uses{GaussianField-configuration-eval-measurable, GaussianField-Configuration, GaussianField-mgf-at-scaled, GaussianField-instMeasurableSpaceConfiguration, Lattice-toSemilatticeInf}
  \texttt{GaussianField.eval\_map\_eq\_gaussianReal}
\end{theorem}
\begin{theorem}[\texttt{eval\_map\_eq\_of\_moments}]\label{GaussianField-eval-map-eq-of-moments}
  \lean{GaussianField.eval_map_eq_of_moments}\leanok
  \uses{GaussianField-configuration-eval-measurable, GaussianField-Configuration, GaussianField-instMeasurableSpaceConfiguration}
  \texttt{GaussianField.eval\_map\_eq\_of\_moments}
\end{theorem}
\section{\texttt{SchwartzNuclear.HermiteTensorProduct}}
\begin{theorem}[\texttt{hermiteFunctionNd\_unpair}]\label{GaussianField-hermiteFunctionNd-unpair}
  \lean{GaussianField.hermiteFunctionNd_unpair}\leanok
  \uses{GaussianField-euclideanSnoc-init-last, GaussianField-hermiteFunctionNd, GaussianField-MultiIndex, GaussianField-euclideanSnoc, GaussianField-euclideanInit, GaussianField-multiIndexEquiv}
  \texttt{GaussianField.hermiteFunctionNd\_unpair}
\end{theorem}
\begin{theorem}[\texttt{multiIndexEquiv\_symm\_growth}]\label{GaussianField-multiIndexEquiv-symm-growth}
  \lean{GaussianField.multiIndexEquiv_symm_growth}\leanok
  \uses{GaussianField-MultiIndex, GaussianField-MultiIndex-abs, GaussianField-multiIndexEquiv}
  \texttt{GaussianField.multiIndexEquiv\_symm\_growth}
\end{theorem}
\begin{definition}[\texttt{fromRapidDecayNdCLM}]\label{GaussianField-fromRapidDecayNdCLM}
  \lean{GaussianField.fromRapidDecayNdCLM}\leanok
  \uses{GaussianField-RapidDecaySeq, GaussianField-RapidDecaySeq-instModule, GaussianField-RapidDecaySeq-instAddCommGroup, GaussianField-RapidDecaySeq-instTopologicalSpace}
  \texttt{GaussianField.fromRapidDecayNdCLM}
\end{definition}
\begin{theorem}[\texttt{euclideanFin1MeasEquiv\_mp}]\label{GaussianField-euclideanFin1MeasEquiv-mp}
  \lean{GaussianField.euclideanFin1MeasEquiv_mp}\leanok
  \uses{GaussianField-euclideanFin1MeasEquiv}
  \texttt{GaussianField.euclideanFin1MeasEquiv\_mp}
\end{theorem}
\begin{definition}[\texttt{schwartzRapidDecayEquiv}]\label{GaussianField-schwartzRapidDecayEquiv}
  \lean{GaussianField.schwartzRapidDecayEquiv}\leanok
  \uses{GaussianField-schwartzDomCongr, GaussianField-RapidDecaySeq, GaussianField-schwartzRapidDecayEquivFin, GaussianField-RapidDecaySeq-instModule, GaussianField-RapidDecaySeq-instAddCommGroup, GaussianField-RapidDecaySeq-instTopologicalSpace}
  \texttt{GaussianField.schwartzRapidDecayEquiv}
\end{definition}
\begin{definition}[\texttt{euclideanFin1Equiv}]\label{GaussianField-euclideanFin1Equiv}
  \lean{GaussianField.euclideanFin1Equiv}\leanok
  \texttt{GaussianField.euclideanFin1Equiv}
\end{definition}
\begin{definition}[\texttt{abs}]\label{GaussianField-MultiIndex-abs}
  \lean{GaussianField.MultiIndex.abs}\leanok
  \uses{GaussianField-MultiIndex}
  \texttt{GaussianField.MultiIndex.abs}
\end{definition}
\begin{theorem}[\texttt{euclideanFin1MeasEquiv\_apply}]\label{GaussianField-euclideanFin1MeasEquiv-apply}
  \lean{GaussianField.euclideanFin1MeasEquiv_apply}\leanok
  \uses{GaussianField-euclideanFin1MeasEquiv}
  \texttt{GaussianField.euclideanFin1MeasEquiv\_apply}
\end{theorem}
\begin{theorem}[\texttt{schwartzRapidDecayEquivNd\_summable}]\label{GaussianField-schwartzRapidDecayEquivNd-summable}
  \lean{GaussianField.schwartzRapidDecayEquivNd_summable}\leanok
  \uses{GaussianField-RapidDecaySeq, GaussianField-RapidDecaySeq-val, GaussianField-hermiteFunctionNd, GaussianField-MultiIndex, GaussianField-multiIndexEquiv}
  \texttt{GaussianField.schwartzRapidDecayEquivNd\_summable}
\end{theorem}
\begin{theorem}[\texttt{schwartzRapidDecayEquivNd\_symm\_apply}]\label{GaussianField-schwartzRapidDecayEquivNd-symm-apply}
  \lean{GaussianField.schwartzRapidDecayEquivNd_symm_apply}\leanok
  \uses{GaussianField-RapidDecaySeq, GaussianField-RapidDecaySeq-val, GaussianField-schwartzRapidDecayEquivNd, GaussianField-hermiteFunctionNd, GaussianField-RapidDecaySeq-instModule, GaussianField-MultiIndex, GaussianField-fromRapidDecayNdCLM, GaussianField-RapidDecaySeq-instAddCommGroup, GaussianField-RapidDecaySeq-instTopologicalSpace, GaussianField-multiIndexEquiv}
  \texttt{GaussianField.schwartzRapidDecayEquivNd\_symm\_apply}
\end{theorem}
\begin{definition}[\texttt{schwartzRapidDecayEquiv1D}]\label{GaussianField-schwartzRapidDecayEquiv1D}
  \lean{GaussianField.schwartzRapidDecayEquiv1D}\leanok
  \uses{GaussianField-RapidDecaySeq, GaussianField-RapidDecaySeq-instModule, GaussianField-RapidDecaySeq-instAddCommGroup, GaussianField-RapidDecaySeq-instTopologicalSpace}
  \texttt{GaussianField.schwartzRapidDecayEquiv1D}
\end{definition}
\begin{definition}[\texttt{MultiIndex}]\label{GaussianField-MultiIndex}
  \lean{GaussianField.MultiIndex}\leanok
  \texttt{GaussianField.MultiIndex}
\end{definition}
\begin{definition}[\texttt{schwartzRapidDecayEquivNd}]\label{GaussianField-schwartzRapidDecayEquivNd}
  \lean{GaussianField.schwartzRapidDecayEquivNd}\leanok
  \uses{GaussianField-RapidDecaySeq, GaussianField-RapidDecaySeq-instModule, GaussianField-fromRapidDecayNdCLM, GaussianField-RapidDecaySeq-instAddCommGroup, GaussianField-RapidDecaySeq-instTopologicalSpace, GaussianField-toRapidDecayNdCLM}
  \texttt{GaussianField.schwartzRapidDecayEquivNd}
\end{definition}
\begin{theorem}[\texttt{schwartzMap\_T2Space}]\label{GaussianField-schwartzMap-T2Space}
  \lean{GaussianField.schwartzMap_T2Space}\leanok
  \texttt{GaussianField.schwartzMap\_T2Space}
\end{theorem}
\begin{definition}[\texttt{hermiteFunctionNd}]\label{GaussianField-hermiteFunctionNd}
  \lean{GaussianField.hermiteFunctionNd}\leanok
  \uses{GaussianField-MultiIndex}
  \texttt{GaussianField.hermiteFunctionNd}
\end{definition}
\begin{definition}[\texttt{multiIndexEquiv}]\label{GaussianField-multiIndexEquiv}
  \lean{GaussianField.multiIndexEquiv}\leanok
  \uses{GaussianField-MultiIndex}
  \texttt{GaussianField.multiIndexEquiv}
\end{definition}
\begin{theorem}[\texttt{schwartzRapidDecayEquiv1D\_symm\_apply}]\label{GaussianField-schwartzRapidDecayEquiv1D-symm-apply}
  \lean{GaussianField.schwartzRapidDecayEquiv1D_symm_apply}\leanok
  \uses{GaussianField-RapidDecaySeq, GaussianField-RapidDecaySeq-val, GaussianField-RapidDecaySeq-instModule, GaussianField-RapidDecaySeq-instAddCommGroup, GaussianField-RapidDecaySeq-instTopologicalSpace, GaussianField-schwartzRapidDecayEquiv1D}
  \texttt{GaussianField.schwartzRapidDecayEquiv1D\_symm\_apply}
\end{theorem}
\begin{definition}[\texttt{toRapidDecayNdCLM}]\label{GaussianField-toRapidDecayNdCLM}
  \lean{GaussianField.toRapidDecayNdCLM}\leanok
  \uses{GaussianField-RapidDecaySeq, GaussianField-RapidDecaySeq-instModule, GaussianField-RapidDecaySeq-instAddCommGroup, GaussianField-RapidDecaySeq-instTopologicalSpace}
  \texttt{GaussianField.toRapidDecayNdCLM}
\end{definition}
\begin{definition}[\texttt{schwartzDomCongr}]\label{GaussianField-schwartzDomCongr}
  \lean{GaussianField.schwartzDomCongr}\leanok
  \texttt{GaussianField.schwartzDomCongr}
\end{definition}
\begin{definition}[\texttt{schwartzRapidDecayEquivFin}]\label{GaussianField-schwartzRapidDecayEquivFin}
  \lean{GaussianField.schwartzRapidDecayEquivFin}\leanok
  \uses{GaussianField-schwartzDomCongr, GaussianField-RapidDecaySeq, GaussianField-schwartzRapidDecayEquivNd, GaussianField-RapidDecaySeq-instModule, GaussianField-RapidDecaySeq-instAddCommGroup, GaussianField-RapidDecaySeq-instTopologicalSpace, GaussianField-euclideanFin1Equiv, GaussianField-schwartzRapidDecayEquiv1D}
  \texttt{GaussianField.schwartzRapidDecayEquivFin}
\end{definition}
\begin{theorem}[\texttt{schwartzRapidDecayEquivNd\_norm\_summable}]\label{GaussianField-schwartzRapidDecayEquivNd-norm-summable}
  \lean{GaussianField.schwartzRapidDecayEquivNd_norm_summable}\leanok
  \uses{GaussianField-schwartzRapidDecayEquivNd-summable, GaussianField-RapidDecaySeq, GaussianField-RapidDecaySeq-val, GaussianField-hermiteFunctionNd, GaussianField-MultiIndex, GaussianField-multiIndexEquiv}
  \texttt{GaussianField.schwartzRapidDecayEquivNd\_norm\_summable}
\end{theorem}
\begin{theorem}[\texttt{hermiteCoeff1D\_basis\_kronecker}]\label{GaussianField-hermiteCoeff1D-basis-kronecker}
  \lean{GaussianField.hermiteCoeff1D_basis_kronecker}\leanok
  \texttt{GaussianField.hermiteCoeff1D\_basis\_kronecker}
\end{theorem}
\begin{theorem}[\texttt{schwartzRapidDecayEquiv1D\_norm\_summable}]\label{GaussianField-schwartzRapidDecayEquiv1D-norm-summable}
  \lean{GaussianField.schwartzRapidDecayEquiv1D_norm_summable}\leanok
  \uses{GaussianField-RapidDecaySeq, GaussianField-RapidDecaySeq-val, GaussianField-schwartzRapidDecayEquiv1D-summable}
  \texttt{GaussianField.schwartzRapidDecayEquiv1D\_norm\_summable}
\end{theorem}
\begin{definition}[\texttt{hermiteCoeffNd}]\label{GaussianField-hermiteCoeffNd}
  \lean{GaussianField.hermiteCoeffNd}\leanok
  \uses{GaussianField-hermiteFunctionNd, GaussianField-MultiIndex}
  \texttt{GaussianField.hermiteCoeffNd}
\end{definition}
\begin{theorem}[\texttt{schwartzRapidDecayEquiv1D\_summable}]\label{GaussianField-schwartzRapidDecayEquiv1D-summable}
  \lean{GaussianField.schwartzRapidDecayEquiv1D_summable}\leanok
  \uses{GaussianField-RapidDecaySeq, GaussianField-RapidDecaySeq-val, Lattice-toSemilatticeInf}
  \texttt{GaussianField.schwartzRapidDecayEquiv1D\_summable}
\end{theorem}
\begin{theorem}[\texttt{hermiteCoeffNd\_eq\_hermiteCoeff1D}]\label{GaussianField-hermiteCoeffNd-eq-hermiteCoeff1D}
  \lean{GaussianField.hermiteCoeffNd_eq_hermiteCoeff1D}\leanok
  \uses{GaussianField-schwartzDomCongr, GaussianField-hermiteFunctionNd, GaussianField-hermiteCoeffNd, GaussianField-euclideanFin1MeasEquiv-mp, GaussianField-euclideanFin1MeasEquiv, GaussianField-euclideanFin1MeasEquiv-apply, GaussianField-euclideanFin1Equiv}
  \texttt{GaussianField.hermiteCoeffNd\_eq\_hermiteCoeff1D}
\end{theorem}
\begin{theorem}[\texttt{multiIndexEquiv\_growth}]\label{GaussianField-multiIndexEquiv-growth}
  \lean{GaussianField.multiIndexEquiv_growth}\leanok
  \uses{GaussianField-MultiIndex, GaussianField-MultiIndex-abs, GaussianField-multiIndexEquiv}
  \texttt{GaussianField.multiIndexEquiv\_growth}
\end{theorem}
\begin{definition}[\texttt{euclideanFin1MeasEquiv}]\label{GaussianField-euclideanFin1MeasEquiv}
  \lean{GaussianField.euclideanFin1MeasEquiv}\leanok
  \texttt{GaussianField.euclideanFin1MeasEquiv}
\end{definition}
\begin{definition}[\texttt{schwartzHermiteBasisNd}]\label{GaussianField-schwartzHermiteBasisNd}
  \lean{GaussianField.schwartzHermiteBasisNd}\leanok
  \uses{GaussianField-hermiteFunctionNd, GaussianField-MultiIndex}
  \texttt{GaussianField.schwartzHermiteBasisNd}
\end{definition}
\section{\texttt{GaussianField.DensityAsym}}
\begin{theorem}[\texttt{normalizedGaussianDensityMeasureAsym\_charFunDual\_eq\_latticeGaussianFieldLawAsym}]\label{GaussianField-normalizedGaussianDensityMeasureAsym-charFunDual-eq-latticeGaussianFieldLawAsym}
  \lean{GaussianField.normalizedGaussianDensityMeasureAsym_charFunDual_eq_latticeGaussianFieldLawAsym}\leanok
  \uses{GaussianField-AsymLatticeSites, GaussianField-normalizedGaussianDensityMeasureAsym, GaussianField-latticeGaussianFieldLawAsym-fourier, GaussianField-latticeGaussianFieldLawAsym, GaussianField-strongDualToFieldAsym, GaussianField-charFunDualAsym-eq-site-integral, GaussianField-AsymLatticeField, GaussianField-instFintypeAsymLatticeSitesOfNeZeroNat, GaussianField-latticeCovarianceAsymGJ, GaussianField-ell2-, GaussianField-normalizedGaussianDensityMeasureAsym-linearFourier}
  \texttt{GaussianField.normalizedGaussianDensityMeasureAsym\_charFunDual\_eq\_latticeGaussianFieldLawAsym}
\end{theorem}
\begin{theorem}[\texttt{massEigenbasisAsym\_coeff\_reprSymm}]\label{GaussianField-massEigenbasisAsym-coeff-reprSymm}
  \lean{GaussianField.massEigenbasisAsym_coeff_reprSymm}\leanok
  \uses{GaussianField-AsymLatticeSites, GaussianField-massEigenvectorBasisAsym, GaussianField-instFintypeAsymLatticeSitesOfNeZeroNat}
  \texttt{GaussianField.massEigenbasisAsym\_coeff\_reprSymm}
\end{theorem}
\begin{theorem}[\texttt{measurable\_sitePairingAsym}]\label{GaussianField-measurable-sitePairingAsym}
  \lean{GaussianField.measurable_sitePairingAsym}\leanok
  \uses{GaussianField-AsymLatticeSites, GaussianField-AsymLatticeField, GaussianField-instFintypeAsymLatticeSitesOfNeZeroNat}
  \texttt{GaussianField.measurable\_sitePairingAsym}
\end{theorem}
\begin{theorem}[\texttt{gaussianDensityAsym\_measurable}]\label{GaussianField-gaussianDensityAsym-measurable}
  \lean{GaussianField.gaussianDensityAsym_measurable}\leanok
  \uses{GaussianField-AsymLatticeSites, GaussianField-gaussianDensityAsym, GaussianField-AsymLatticeField, GaussianField-instFintypeAsymLatticeSitesOfNeZeroNat, GaussianField-massOperatorAsym}
  \texttt{GaussianField.gaussianDensityAsym\_measurable}
\end{theorem}
\begin{theorem}[\texttt{congr\_simp}]\label{GaussianField-strongDualToFieldAsym-congr-simp}
  \lean{GaussianField.strongDualToFieldAsym.congr_simp}\leanok
  \uses{GaussianField-AsymLatticeSites, GaussianField-strongDualToFieldAsym, GaussianField-AsymLatticeField}
  \texttt{GaussianField.strongDualToFieldAsym.congr\_simp}
\end{theorem}
\begin{theorem}[\texttt{evalMap\_evalMapInvAsym}]\label{GaussianField-evalMap-evalMapInvAsym}
  \lean{GaussianField.evalMap_evalMapInvAsym}\leanok
  \uses{GaussianField-evalMapAsymInv-apply, GaussianField-AsymLatticeSites, GaussianField-Configuration, GaussianField-AsymLatticeField, GaussianField-evalMapAsym, GaussianField-instFintypeAsymLatticeSitesOfNeZeroNat, GaussianField-evalMapAsymInv, GaussianField-asymLatticeDelta}
  \texttt{GaussianField.evalMap\_evalMapInvAsym}
\end{theorem}
\begin{theorem}[\texttt{strongDualAsym\_apply\_eq\_site\_sum}]\label{GaussianField-strongDualAsym-apply-eq-site-sum}
  \lean{GaussianField.strongDualAsym_apply_eq_site_sum}\leanok
  \uses{GaussianField-AsymLatticeSites, GaussianField-asym-field-basis-decomp-density, GaussianField-strongDualToFieldAsym, GaussianField-AsymLatticeField, GaussianField-instFintypeAsymLatticeSitesOfNeZeroNat, GaussianField-asymLatticeDelta}
  \texttt{GaussianField.strongDualAsym\_apply\_eq\_site\_sum}
\end{theorem}
\begin{theorem}[\texttt{congr\_simp}]\label{GaussianField-gaussianDensityMeasureAsym-congr-simp}
  \lean{GaussianField.gaussianDensityMeasureAsym.congr_simp}\leanok
  \uses{GaussianField-AsymLatticeSites, GaussianField-gaussianDensityMeasureAsym, GaussianField-AsymLatticeField}
  \texttt{GaussianField.gaussianDensityMeasureAsym.congr\_simp}
\end{theorem}
\begin{theorem}[\texttt{config\_apply\_eq\_sum\_deltaAsym}]\label{GaussianField-config-apply-eq-sum-deltaAsym}
  \lean{GaussianField.config_apply_eq_sum_deltaAsym}\leanok
  \uses{GaussianField-AsymLatticeSites, GaussianField-asym-field-basis-decomp-density, GaussianField-Configuration, GaussianField-AsymLatticeField, GaussianField-instFintypeAsymLatticeSitesOfNeZeroNat, GaussianField-asymLatticeDelta}
  \texttt{GaussianField.config\_apply\_eq\_sum\_deltaAsym}
\end{theorem}
\begin{theorem}[\texttt{massEigenbasisAsym\_coeff\_reprSymm\_ofLp}]\label{GaussianField-massEigenbasisAsym-coeff-reprSymm-ofLp}
  \lean{GaussianField.massEigenbasisAsym_coeff_reprSymm_ofLp}\leanok
  \uses{GaussianField-massEigenbasisAsym-coeff-reprSymm, GaussianField-AsymLatticeSites, GaussianField-massEigenvectorBasisAsym, GaussianField-instFintypeAsymLatticeSitesOfNeZeroNat}
  \texttt{GaussianField.massEigenbasisAsym\_coeff\_reprSymm\_ofLp}
\end{theorem}
\begin{definition}[\texttt{strongDualToFieldAsym}]\label{GaussianField-strongDualToFieldAsym}
  \lean{GaussianField.strongDualToFieldAsym}\leanok
  \uses{GaussianField-AsymLatticeSites, GaussianField-AsymLatticeField, GaussianField-asymLatticeDelta}
  \texttt{GaussianField.strongDualToFieldAsym}
\end{definition}
\begin{theorem}[\texttt{evalMapAsymInv\_apply}]\label{GaussianField-evalMapAsymInv-apply}
  \lean{GaussianField.evalMapAsymInv_apply}\leanok
  \uses{GaussianField-AsymLatticeSites, GaussianField-Configuration, GaussianField-AsymLatticeField, GaussianField-instFintypeAsymLatticeSitesOfNeZeroNat, GaussianField-evalMapAsymInv}
  \texttt{GaussianField.evalMapAsymInv\_apply}
\end{theorem}
\begin{theorem}[\texttt{massEigenbasisAsym\_quadratic\_sum\_reprSymm\_ofLp}]\label{GaussianField-massEigenbasisAsym-quadratic-sum-reprSymm-ofLp}
  \lean{GaussianField.massEigenbasisAsym_quadratic_sum_reprSymm_ofLp}\leanok
  \uses{GaussianField-massEigenbasisAsym-coeff-reprSymm-ofLp, GaussianField-AsymLatticeSites, GaussianField-massEigenvectorBasisAsym, GaussianField-massEigenvaluesAsym, GaussianField-instFintypeAsymLatticeSitesOfNeZeroNat}
  \texttt{GaussianField.massEigenbasisAsym\_quadratic\_sum\_reprSymm\_ofLp}
\end{theorem}
\begin{theorem}[\texttt{gaussianDensityAsym\_eq\_exp\_spectral}]\label{GaussianField-gaussianDensityAsym-eq-exp-spectral}
  \lean{GaussianField.gaussianDensityAsym_eq_exp_spectral}\leanok
  \uses{GaussianField-AsymLatticeSites, GaussianField-massEigenvectorBasisAsym, GaussianField-massOperatorAsym-eigenCoeff-eq-eigenvalues-mul-eigenCoeff, GaussianField-massEigenbasisAsym-sum-mul-sum-eq-site-inner, GaussianField-gaussianDensityAsym, GaussianField-massEigenvaluesAsym, GaussianField-AsymLatticeField, GaussianField-instFintypeAsymLatticeSitesOfNeZeroNat, GaussianField-massOperatorAsym}
  \texttt{GaussianField.gaussianDensityAsym\_eq\_exp\_spectral}
\end{theorem}
\begin{theorem}[\texttt{sitePairingAsym\_eq\_massEigenbasis\_sum}]\label{GaussianField-sitePairingAsym-eq-massEigenbasis-sum}
  \lean{GaussianField.sitePairingAsym_eq_massEigenbasis_sum}\leanok
  \uses{GaussianField-AsymLatticeSites, GaussianField-massEigenvectorBasisAsym, GaussianField-massEigenbasisAsym-sum-mul-sum-eq-site-inner, GaussianField-AsymLatticeField, GaussianField-instFintypeAsymLatticeSitesOfNeZeroNat}
  \texttt{GaussianField.sitePairingAsym\_eq\_massEigenbasis\_sum}
\end{theorem}
\begin{theorem}[\texttt{integral\_massEigenbasisAsym\_cexp\_GJ}]\label{GaussianField-integral-massEigenbasisAsym-cexp-GJ}
  \lean{GaussianField.integral_massEigenbasisAsym_cexp_GJ}\leanok
  \uses{GaussianField-AsymLatticeSites, GaussianField-massEigenbasisAsym-linear-sum-reprSymm-ofLp, GaussianField-massEigenvectorBasisAsym, GaussianField-integral-cexp-neg-half-sum-mul-sq-add-linear, GaussianField-massEigenvaluesAsym, GaussianField-massEigenbasisAsym-quadratic-sum-reprSymm-ofLp, GaussianField-massOperatorMatrixAsym-eigenvalues-pos, GaussianField-instFintypeAsymLatticeSitesOfNeZeroNat}
  \texttt{GaussianField.integral\_massEigenbasisAsym\_cexp\_GJ}
\end{theorem}
\begin{definition}[\texttt{latticeGaussianFieldLawAsym}]\label{GaussianField-latticeGaussianFieldLawAsym}
  \lean{GaussianField.latticeGaussianFieldLawAsym}\leanok
  \uses{GaussianField-AsymLatticeSites, GaussianField-measure, GaussianField-ell2-separableSpace, GaussianField-asymLatticeField-dyninMityaginSpace, GaussianField-Configuration, GaussianField-AsymLatticeField, GaussianField-evalMapAsym, GaussianField-instMeasurableSpaceConfiguration, GaussianField-latticeCovarianceAsymGJ, GaussianField-ell2-}
  \texttt{GaussianField.latticeGaussianFieldLawAsym}
\end{definition}
\begin{theorem}[\texttt{measurable\_evalMapAsym}]\label{GaussianField-measurable-evalMapAsym}
  \lean{GaussianField.measurable_evalMapAsym}\leanok
  \uses{GaussianField-AsymLatticeSites, GaussianField-configuration-eval-measurable, GaussianField-Configuration, GaussianField-AsymLatticeField, GaussianField-evalMapAsym, GaussianField-instMeasurableSpaceConfiguration, GaussianField-asymLatticeDelta}
  \texttt{GaussianField.measurable\_evalMapAsym}
\end{theorem}
\begin{theorem}[\texttt{evalMapInv\_evalMapAsym}]\label{GaussianField-evalMapInv-evalMapAsym}
  \lean{GaussianField.evalMapInv_evalMapAsym}\leanok
  \uses{GaussianField-AsymLatticeSites, GaussianField-config-apply-eq-sum-evalMapAsym, GaussianField-Configuration, GaussianField-AsymLatticeField, GaussianField-evalMapAsym, GaussianField-instFintypeAsymLatticeSitesOfNeZeroNat, GaussianField-evalMapAsymInv}
  \texttt{GaussianField.evalMapInv\_evalMapAsym}
\end{theorem}
\begin{theorem}[\texttt{asym\_field\_basis\_decomp\_density}]\label{GaussianField-asym-field-basis-decomp-density}
  \lean{GaussianField.asym_field_basis_decomp_density}\leanok
  \uses{GaussianField-AsymLatticeSites, GaussianField-AsymLatticeField, GaussianField-instFintypeAsymLatticeSitesOfNeZeroNat, GaussianField-asymLatticeDelta}
  \texttt{GaussianField.asym\_field\_basis\_decomp\_density}
\end{theorem}
\begin{definition}[\texttt{evalMapAsym}]\label{GaussianField-evalMapAsym}
  \lean{GaussianField.evalMapAsym}\leanok
  \uses{GaussianField-AsymLatticeSites, GaussianField-Configuration, GaussianField-AsymLatticeField, GaussianField-asymLatticeDelta}
  \texttt{GaussianField.evalMapAsym}
\end{definition}
\begin{theorem}[\texttt{config\_apply\_eq\_sum\_evalMapAsym}]\label{GaussianField-config-apply-eq-sum-evalMapAsym}
  \lean{GaussianField.config_apply_eq_sum_evalMapAsym}\leanok
  \uses{GaussianField-AsymLatticeSites, GaussianField-Configuration, GaussianField-AsymLatticeField, GaussianField-evalMapAsym, GaussianField-instFintypeAsymLatticeSitesOfNeZeroNat, GaussianField-asymLatticeDelta, GaussianField-config-apply-eq-sum-deltaAsym}
  \texttt{GaussianField.config\_apply\_eq\_sum\_evalMapAsym}
\end{theorem}
\begin{theorem}[\texttt{lattice\_covariance\_AsymGJ\_eq\_spectral}]\label{GaussianField-lattice-covariance-AsymGJ-eq-spectral}
  \lean{GaussianField.lattice_covariance_AsymGJ_eq_spectral}\leanok
  \uses{GaussianField-spectralLatticeCovarianceAsym, GaussianField-AsymLatticeSites, GaussianField-covariance, GaussianField-massEigenvectorBasisAsym, GaussianField-massEigenvaluesAsym, GaussianField-covariance-spectralLatticeCovarianceAsym-eq, GaussianField-AsymLatticeField, GaussianField-instFintypeAsymLatticeSitesOfNeZeroNat, GaussianField-latticeCovarianceAsymGJ, GaussianField-ell2-}
  \texttt{GaussianField.lattice\_covariance\_AsymGJ\_eq\_spectral}
\end{theorem}
\begin{definition}[\texttt{gaussianDensityMeasureAsym}]\label{GaussianField-gaussianDensityMeasureAsym}
  \lean{GaussianField.gaussianDensityMeasureAsym}\leanok
  \uses{GaussianField-AsymLatticeSites, GaussianField-gaussianDensityWeightAsym, GaussianField-AsymLatticeField, GaussianField-instFintypeAsymLatticeSitesOfNeZeroNat}
  \texttt{GaussianField.gaussianDensityMeasureAsym}
\end{definition}
\begin{theorem}[\texttt{massEigenbasisAsym\_linear\_sum\_reprSymm\_ofLp}]\label{GaussianField-massEigenbasisAsym-linear-sum-reprSymm-ofLp}
  \lean{GaussianField.massEigenbasisAsym_linear_sum_reprSymm_ofLp}\leanok
  \uses{GaussianField-massEigenbasisAsym-coeff-reprSymm-ofLp, GaussianField-AsymLatticeSites, GaussianField-massEigenvectorBasisAsym, GaussianField-instFintypeAsymLatticeSitesOfNeZeroNat}
  \texttt{GaussianField.massEigenbasisAsym\_linear\_sum\_reprSymm\_ofLp}
\end{theorem}
\begin{theorem}[\texttt{latticeGaussianFieldLawAsym\_eq\_normalizedQuadraticGaussianMeasure}]\label{GaussianField-latticeGaussianFieldLawAsym-eq-normalizedQuadraticGaussianMeasure}
  \lean{GaussianField.latticeGaussianFieldLawAsym_eq_normalizedQuadraticGaussianMeasure}\leanok
  \uses{GaussianField-normalizedQuadraticGaussianMeasure, GaussianField-AsymLatticeSites, GaussianField-measure, GaussianField-ell2-separableSpace, GaussianField-normalizedGaussianDensityMeasureAsym, GaussianField-asymLatticeField-dyninMityaginSpace, GaussianField-measurable-evalMapAsym, GaussianField-latticeGaussianFieldLawAsym, GaussianField-normalizedGaussianDensityMeasureAsym-eq-normalizedQuadraticGaussianMeasure, GaussianField-normalizedGaussianDensityMeasureAsym-isFinite, GaussianField-Configuration, GaussianField-AsymLatticeField, GaussianField-evalMapAsym, GaussianField-instMeasurableSpaceConfiguration, GaussianField-instFintypeAsymLatticeSitesOfNeZeroNat, GaussianField-massOperatorAsym, GaussianField-latticeCovarianceAsymGJ, GaussianField-ell2-, GaussianField-measure-isProbability, GaussianField-normalizedGaussianDensityMeasureAsym-charFunDual-eq-latticeGaussianFieldLawAsym}
  \texttt{GaussianField.latticeGaussianFieldLawAsym\_eq\_normalizedQuadraticGaussianMeasure}
\end{theorem}
\begin{theorem}[\texttt{normalizedGaussianDensityMeasureAsym\_linearFourier}]\label{GaussianField-normalizedGaussianDensityMeasureAsym-linearFourier}
  \lean{GaussianField.normalizedGaussianDensityMeasureAsym_linearFourier}\leanok
  \uses{GaussianField-AsymLatticeSites, GaussianField-gaussianDensityWeightAsym, GaussianField-gaussianDensityNormConstAsym, GaussianField-normalizedGaussianDensityMeasureAsym, GaussianField-covariance, GaussianField-massEigenvectorBasisAsym, GaussianField-gaussianDensityAsym, GaussianField-integral-massEigenbasisAsym-cexp-GJ, GaussianField-sitePairingAsym-eq-massEigenbasis-sum, GaussianField-gaussianDensityAsym-measurable, GaussianField-massEigenvaluesAsym, GaussianField-gaussianDensityMeasureAsym, GaussianField-massOperatorMatrixAsym-eigenvalues-pos, GaussianField-AsymLatticeField, GaussianField-instFintypeAsymLatticeSitesOfNeZeroNat, GaussianField-latticeCovarianceAsymGJ, GaussianField-gaussianDensityAsym-nonneg, GaussianField-gaussianDensityAsym-eq-exp-spectral, GaussianField-ell2-, GaussianField-lattice-covariance-AsymGJ-eq-spectral}
  \texttt{GaussianField.normalizedGaussianDensityMeasureAsym\_linearFourier}
\end{theorem}
\begin{theorem}[\texttt{charFunDualAsym\_eq\_site\_integral}]\label{GaussianField-charFunDualAsym-eq-site-integral}
  \lean{GaussianField.charFunDualAsym_eq_site_integral}\leanok
  \uses{GaussianField-AsymLatticeSites, GaussianField-strongDualAsym-apply-eq-site-sum, GaussianField-strongDualToFieldAsym, GaussianField-AsymLatticeField, GaussianField-instFintypeAsymLatticeSitesOfNeZeroNat}
  \texttt{GaussianField.charFunDualAsym\_eq\_site\_integral}
\end{theorem}
\begin{theorem}[\texttt{congr\_simp}]\label{GaussianField-gaussianDensityNormConstAsym-congr-simp}
  \lean{GaussianField.gaussianDensityNormConstAsym.congr_simp}\leanok
  \uses{GaussianField-gaussianDensityNormConstAsym}
  \texttt{GaussianField.gaussianDensityNormConstAsym.congr\_simp}
\end{theorem}
\begin{theorem}[\texttt{normalizedGaussianDensityMeasureAsym\_eq\_normalizedQuadraticGaussianMeasure}]\label{GaussianField-normalizedGaussianDensityMeasureAsym-eq-normalizedQuadraticGaussianMeasure}
  \lean{GaussianField.normalizedGaussianDensityMeasureAsym_eq_normalizedQuadraticGaussianMeasure}\leanok
  \uses{GaussianField-quadraticGaussianDensity, GaussianField-normalizedQuadraticGaussianMeasure, GaussianField-AsymLatticeSites, GaussianField-quadraticGaussianMeasure, GaussianField-gaussianDensityWeightAsym, GaussianField-gaussianDensityNormConstAsym, GaussianField-normalizedGaussianDensityMeasureAsym, GaussianField-gaussianDensityAsym, GaussianField-gaussianDensityMeasureAsym, GaussianField-AsymLatticeField, GaussianField-instFintypeAsymLatticeSitesOfNeZeroNat, GaussianField-massOperatorAsym}
  \texttt{GaussianField.normalizedGaussianDensityMeasureAsym\_eq\_normalizedQuadraticGaussianMeasure}
\end{theorem}
\begin{theorem}[\texttt{congr\_simp}]\label{GaussianField-evalMapAsymInv-congr-simp}
  \lean{GaussianField.evalMapAsymInv.congr_simp}\leanok
  \uses{GaussianField-AsymLatticeSites, GaussianField-Configuration, GaussianField-AsymLatticeField, GaussianField-evalMapAsymInv}
  \texttt{GaussianField.evalMapAsymInv.congr\_simp}
\end{theorem}
\begin{theorem}[\texttt{latticeGaussianFieldLawAsym\_fourier}]\label{GaussianField-latticeGaussianFieldLawAsym-fourier}
  \lean{GaussianField.latticeGaussianFieldLawAsym_fourier}\leanok
  \uses{GaussianField-charFun, GaussianField-AsymLatticeSites, GaussianField-measure, GaussianField-ell2-separableSpace, GaussianField-asymLatticeField-dyninMityaginSpace, GaussianField-measurable-evalMapAsym, GaussianField-latticeGaussianFieldLawAsym, GaussianField-config-apply-eq-sum-evalMapAsym, GaussianField-Configuration, GaussianField-AsymLatticeField, GaussianField-evalMapAsym, GaussianField-instMeasurableSpaceConfiguration, GaussianField-instFintypeAsymLatticeSitesOfNeZeroNat, GaussianField-latticeCovarianceAsymGJ, GaussianField-ell2-}
  \texttt{GaussianField.latticeGaussianFieldLawAsym\_fourier}
\end{theorem}
\begin{definition}[\texttt{evalMapAsymInv}]\label{GaussianField-evalMapAsymInv}
  \lean{GaussianField.evalMapAsymInv}\leanok
  \uses{GaussianField-AsymLatticeSites, GaussianField-Configuration, GaussianField-AsymLatticeField, GaussianField-instFintypeAsymLatticeSitesOfNeZeroNat}
  \texttt{GaussianField.evalMapAsymInv}
\end{definition}
\begin{theorem}[\texttt{congr\_simp}]\label{GaussianField-gaussianDensityWeightAsym-congr-simp}
  \lean{GaussianField.gaussianDensityWeightAsym.congr_simp}\leanok
  \uses{GaussianField-gaussianDensityWeightAsym, GaussianField-AsymLatticeField}
  \texttt{GaussianField.gaussianDensityWeightAsym.congr\_simp}
\end{theorem}
\begin{theorem}[\texttt{evalMapAsymInv\_measurable}]\label{GaussianField-evalMapAsymInv-measurable}
  \lean{GaussianField.evalMapAsymInv_measurable}\leanok
  \uses{GaussianField-evalMapAsymInv-apply, GaussianField-AsymLatticeSites, GaussianField-configuration-measurable-of-eval-measurable, GaussianField-Configuration, GaussianField-AsymLatticeField, GaussianField-instMeasurableSpaceConfiguration, GaussianField-instFintypeAsymLatticeSitesOfNeZeroNat, GaussianField-evalMapAsymInv}
  \texttt{GaussianField.evalMapAsymInv\_measurable}
\end{theorem}
\begin{definition}[\texttt{normalizedGaussianDensityMeasureAsym}]\label{GaussianField-normalizedGaussianDensityMeasureAsym}
  \lean{GaussianField.normalizedGaussianDensityMeasureAsym}\leanok
  \uses{GaussianField-AsymLatticeSites, GaussianField-gaussianDensityNormConstAsym, GaussianField-gaussianDensityMeasureAsym, GaussianField-AsymLatticeField}
  \texttt{GaussianField.normalizedGaussianDensityMeasureAsym}
\end{definition}
\begin{theorem}[\texttt{congr\_simp}]\label{GaussianField-gaussianDensityAsym-congr-simp}
  \lean{GaussianField.gaussianDensityAsym.congr_simp}\leanok
  \uses{GaussianField-gaussianDensityAsym, GaussianField-AsymLatticeField}
  \texttt{GaussianField.gaussianDensityAsym.congr\_simp}
\end{theorem}
\begin{theorem}[\texttt{gaussianDensityAsym\_nonneg}]\label{GaussianField-gaussianDensityAsym-nonneg}
  \lean{GaussianField.gaussianDensityAsym_nonneg}\leanok
  \uses{GaussianField-AsymLatticeSites, GaussianField-gaussianDensityAsym, GaussianField-AsymLatticeField, GaussianField-instFintypeAsymLatticeSitesOfNeZeroNat, GaussianField-massOperatorAsym}
  \texttt{GaussianField.gaussianDensityAsym\_nonneg}
\end{theorem}
\begin{definition}[\texttt{gaussianDensityAsym}]\label{GaussianField-gaussianDensityAsym}
  \lean{GaussianField.gaussianDensityAsym}\leanok
  \uses{GaussianField-AsymLatticeSites, GaussianField-AsymLatticeField, GaussianField-instFintypeAsymLatticeSitesOfNeZeroNat, GaussianField-massOperatorAsym}
  \texttt{GaussianField.gaussianDensityAsym}
\end{definition}
\begin{definition}[\texttt{evalMapAsymMeasurableEquiv}]\label{GaussianField-evalMapAsymMeasurableEquiv}
  \lean{GaussianField.evalMapAsymMeasurableEquiv}\leanok
  \uses{GaussianField-AsymLatticeSites, GaussianField-evalMapAsymInv-measurable, GaussianField-measurable-evalMapAsym, GaussianField-evalMapInv-evalMapAsym, GaussianField-Configuration, GaussianField-AsymLatticeField, GaussianField-evalMapAsym, GaussianField-instMeasurableSpaceConfiguration, GaussianField-evalMap-evalMapInvAsym, GaussianField-evalMapAsymInv}
  \texttt{GaussianField.evalMapAsymMeasurableEquiv}
\end{definition}
\begin{theorem}[\texttt{congr\_simp}]\label{GaussianField-evalMapAsym-congr-simp}
  \lean{GaussianField.evalMapAsym.congr_simp}\leanok
  \uses{GaussianField-AsymLatticeSites, GaussianField-Configuration, GaussianField-AsymLatticeField, GaussianField-evalMapAsym}
  \texttt{GaussianField.evalMapAsym.congr\_simp}
\end{theorem}
\begin{definition}[\texttt{gaussianDensityWeightAsym}]\label{GaussianField-gaussianDensityWeightAsym}
  \lean{GaussianField.gaussianDensityWeightAsym}\leanok
  \uses{GaussianField-gaussianDensityAsym, GaussianField-AsymLatticeField}
  \texttt{GaussianField.gaussianDensityWeightAsym}
\end{definition}
\begin{theorem}[\texttt{normalizedGaussianDensityMeasureAsym\_isFinite}]\label{GaussianField-normalizedGaussianDensityMeasureAsym-isFinite}
  \lean{GaussianField.normalizedGaussianDensityMeasureAsym_isFinite}\leanok
  \uses{GaussianField-AsymLatticeSites, GaussianField-normalizedGaussianDensityMeasureAsym, GaussianField-gaussianDensityMeasureAsym, GaussianField-AsymLatticeField}
  \texttt{GaussianField.normalizedGaussianDensityMeasureAsym\_isFinite}
\end{theorem}
\begin{definition}[\texttt{gaussianDensityNormConstAsym}]\label{GaussianField-gaussianDensityNormConstAsym}
  \lean{GaussianField.gaussianDensityNormConstAsym}\leanok
  \uses{GaussianField-AsymLatticeSites, GaussianField-gaussianDensityMeasureAsym, GaussianField-AsymLatticeField}
  \texttt{GaussianField.gaussianDensityNormConstAsym}
\end{definition}
\section{\texttt{Lattice.Sites}}
\begin{definition}[\texttt{FinLatticeSites}]\label{GaussianField-FinLatticeSites}
  \lean{GaussianField.FinLatticeSites}\leanok
  \texttt{GaussianField.FinLatticeSites}
\end{definition}
\begin{definition}[\texttt{stdBasisInt}]\label{GaussianField-stdBasisInt}
  \lean{GaussianField.stdBasisInt}\leanok
  \texttt{GaussianField.stdBasisInt}
\end{definition}
\begin{theorem}[\texttt{latticeNorm\_zero}]\label{GaussianField-latticeNorm-zero}
  \lean{GaussianField.latticeNorm_zero}\leanok
  \uses{GaussianField-latticeNorm}
  \texttt{GaussianField.latticeNorm\_zero}
\end{theorem}
\begin{theorem}[\texttt{latticeNorm\_shift\_le}]\label{GaussianField-latticeNorm-shift-le}
  \lean{GaussianField.latticeNorm_shift_le}\leanok
  \uses{GaussianField-stdBasisInt, GaussianField-latticeNorm, GaussianField-latticeNorm-stdBasis, GaussianField-latticeNorm-triangle}
  \texttt{GaussianField.latticeNorm\_shift\_le}
\end{theorem}
\begin{definition}[\texttt{latticeNorm}]\label{GaussianField-latticeNorm}
  \lean{GaussianField.latticeNorm}\leanok
  \texttt{GaussianField.latticeNorm}
\end{definition}
\begin{theorem}[\texttt{latticeNorm\_stdBasis}]\label{GaussianField-latticeNorm-stdBasis}
  \lean{GaussianField.latticeNorm_stdBasis}\leanok
  \uses{GaussianField-stdBasisInt, GaussianField-latticeNorm}
  \texttt{GaussianField.latticeNorm\_stdBasis}
\end{theorem}
\begin{definition}[\texttt{neighborsFinite}]\label{GaussianField-neighborsFinite}
  \lean{GaussianField.neighborsFinite}\leanok
  \uses{GaussianField-FinLatticeSites}
  \texttt{GaussianField.neighborsFinite}
\end{definition}
\begin{theorem}[\texttt{latticeNorm\_nonneg}]\label{GaussianField-latticeNorm-nonneg}
  \lean{GaussianField.latticeNorm_nonneg}\leanok
  \uses{GaussianField-latticeNorm, Lattice-toSemilatticeInf}
  \texttt{GaussianField.latticeNorm\_nonneg}
\end{theorem}
\begin{definition}[\texttt{neighbors}]\label{GaussianField-neighbors}
  \lean{GaussianField.neighbors}\leanok
  \uses{GaussianField-stdBasisInt}
  \texttt{GaussianField.neighbors}
\end{definition}
\begin{theorem}[\texttt{latticeNorm\_shift\_sub\_le}]\label{GaussianField-latticeNorm-shift-sub-le}
  \lean{GaussianField.latticeNorm_shift_sub_le}\leanok
  \uses{GaussianField-stdBasisInt, GaussianField-latticeNorm, GaussianField-latticeNorm-stdBasis, GaussianField-latticeNorm-triangle}
  \texttt{GaussianField.latticeNorm\_shift\_sub\_le}
\end{theorem}
\begin{definition}[\texttt{InfLatticeSites}]\label{GaussianField-InfLatticeSites}
  \lean{GaussianField.InfLatticeSites}\leanok
  \texttt{GaussianField.InfLatticeSites}
\end{definition}
\begin{theorem}[\texttt{neighbors\_card\_le}]\label{GaussianField-neighbors-card-le}
  \lean{GaussianField.neighbors_card_le}\leanok
  \uses{GaussianField-stdBasisInt, GaussianField-neighbors}
  \texttt{GaussianField.neighbors\_card\_le}
\end{theorem}
\begin{theorem}[\texttt{latticeNorm\_triangle}]\label{GaussianField-latticeNorm-triangle}
  \lean{GaussianField.latticeNorm_triangle}\leanok
  \uses{GaussianField-latticeNorm}
  \texttt{GaussianField.latticeNorm\_triangle}
\end{theorem}
\section{\texttt{SchwartzNuclear.HermiteNuclear}}
\begin{definition}[\texttt{schwartz\_dyninMityaginSpace\_euclidean}]\label{GaussianField-schwartz-dyninMityaginSpace-euclidean}
  \lean{GaussianField.schwartz_dyninMityaginSpace_euclidean}\leanok
  \uses{GaussianField-DyninMityaginSpace, GaussianField-DyninMityaginSpace-ofRapidDecayEquiv, GaussianField-schwartzRapidDecayEquivNd}
  \texttt{GaussianField.schwartz\_dyninMityaginSpace\_euclidean}
\end{definition}
\begin{definition}[\texttt{schwartz\_dyninMityaginSpace}]\label{GaussianField-schwartz-dyninMityaginSpace}
  \lean{GaussianField.schwartz_dyninMityaginSpace}\leanok
  \uses{GaussianField-DyninMityaginSpace, GaussianField-DyninMityaginSpace-ofRapidDecayEquiv, GaussianField-schwartzRapidDecayEquiv}
  \texttt{GaussianField.schwartz\_dyninMityaginSpace}
\end{definition}
\begin{theorem}[\texttt{schwartz\_separableSpace}]\label{GaussianField-schwartz-separableSpace}
  \lean{GaussianField.schwartz_separableSpace}\leanok
  \uses{GaussianField-RapidDecaySeq, GaussianField-RapidDecaySeq-instModule, GaussianField-rapidDecaySeq-separableSpace, GaussianField-schwartzRapidDecayEquiv, GaussianField-RapidDecaySeq-instAddCommGroup, GaussianField-RapidDecaySeq-instTopologicalSpace}
  [\texttt{schwartz\_separableSpace}](../../Test/WhiteNoise.lean\#L64)  Statement: Schwartz space $\mathcal{S}(\mathbb{R})$ is separable.  ``\texttt{ axiom schwartz\_separableSpace : SeparableSpace (SchwartzMap ℝ ℝ) }``  The Hermite basis provides a countable dense set.
\end{theorem}
\begin{definition}[\texttt{schwartz\_dyninMityaginSpace\_1D}]\label{GaussianField-schwartz-dyninMityaginSpace-1D}
  \lean{GaussianField.schwartz_dyninMityaginSpace_1D}\leanok
  \uses{GaussianField-DyninMityaginSpace, GaussianField-DyninMityaginSpace-ofRapidDecayEquiv, GaussianField-schwartzRapidDecayEquiv1D}
  \texttt{GaussianField.schwartz\_dyninMityaginSpace\_1D}
\end{definition}
\begin{theorem}[\texttt{rapidDecaySeq\_separableSpace}]\label{GaussianField-rapidDecaySeq-separableSpace}
  \lean{GaussianField.rapidDecaySeq_separableSpace}\leanok
  \uses{GaussianField-RapidDecaySeq, GaussianField-RapidDecaySeq-val, GaussianField-RapidDecaySeq-hasSum-basisVec, GaussianField-RapidDecaySeq-instSMul, GaussianField-RapidDecaySeq-instModule, GaussianField-RapidDecaySeq-basisVec, GaussianField-RapidDecaySeq-instAddCommGroup, GaussianField-RapidDecaySeq-instIsTopologicalAddGroup, GaussianField-RapidDecaySeq-instTopologicalSpace, GaussianField-RapidDecaySeq-instContinuousSMul}
  \texttt{GaussianField.rapidDecaySeq\_separableSpace}
\end{theorem}
\begin{theorem}[\texttt{schwartz\_hasBiorthogonalBasis\_1D}]\label{GaussianField-schwartz-hasBiorthogonalBasis-1D}
  \lean{GaussianField.schwartz_hasBiorthogonalBasis_1D}\leanok
  \uses{GaussianField-schwartz-dyninMityaginSpace-1D, GaussianField-DyninMityaginSpace-ofRapidDecayEquiv-hasBiorthogonalBasis, GaussianField-DyninMityaginSpace-HasBiorthogonalBasis, GaussianField-schwartzRapidDecayEquiv1D}
  \texttt{GaussianField.schwartz\_hasBiorthogonalBasis\_1D}
\end{theorem}
\section{\texttt{HeatKernel.Axioms}}
\begin{theorem}[\texttt{spectralCLM\_zero}]\label{GaussianField-spectralCLM-zero}
  \lean{GaussianField.spectralCLM_zero}\leanok
  \uses{GaussianField-DyninMityaginSpace, GaussianField-DyninMityaginSpace-coeff, GaussianField-spectralCLM-coord, GaussianField-spectralCLM, GaussianField-ell2-}
  \texttt{GaussianField.spectralCLM\_zero}
\end{theorem}
\begin{definition}[\texttt{qftEigenvalue}]\label{GaussianField-qftEigenvalue}
  \lean{GaussianField.qftEigenvalue}\leanok
  \uses{GaussianField-SmoothMap-Circle-fourierFreq}
  \texttt{GaussianField.qftEigenvalue}
\end{definition}
\begin{definition}[\texttt{heatSingularValue}]\label{GaussianField-heatSingularValue}
  \lean{GaussianField.heatSingularValue}\leanok
  \uses{GaussianField-qftEigenvalue}
  \texttt{GaussianField.heatSingularValue}
\end{definition}
\begin{theorem}[\texttt{isBoundedSeq\_const}]\label{GaussianField-isBoundedSeq-const}
  \lean{GaussianField.isBoundedSeq_const}\leanok
  \uses{GaussianField-IsBoundedSeq}
  \texttt{GaussianField.isBoundedSeq\_const}
\end{theorem}
\begin{theorem}[\texttt{congr\_simp}]\label{GaussianField-spectralCLM-congr-simp}
  \lean{GaussianField.spectralCLM.congr_simp}\leanok
  \uses{GaussianField-DyninMityaginSpace, GaussianField-spectralCLM, GaussianField-ell2-, GaussianField-IsBoundedSeq}
  \texttt{GaussianField.spectralCLM.congr\_simp}
\end{theorem}
\begin{theorem}[\texttt{qft\_singular\_values\_bounded}]\label{GaussianField-qft-singular-values-bounded}
  \lean{GaussianField.qft_singular_values_bounded}\leanok
  \uses{GaussianField-qftSingularValue, GaussianField-qftEigenvalue, GaussianField-qftSingularValue-nonneg, GaussianField-qftEigenvalue-pos, GaussianField-SmoothMap-Circle-fourierFreq, Lattice-toSemilatticeInf, GaussianField-IsBoundedSeq}
  \texttt{GaussianField.qft\_singular\_values\_bounded}
\end{theorem}
\begin{theorem}[\texttt{const\_mul}]\label{GaussianField-IsBoundedSeq-const-mul}
  \lean{GaussianField.IsBoundedSeq.const_mul}\leanok
  \uses{GaussianField-isBoundedSeq-const, GaussianField-IsBoundedSeq-mul, GaussianField-IsBoundedSeq}
  \texttt{GaussianField.IsBoundedSeq.const\_mul}
\end{theorem}
\begin{definition}[\texttt{spectralCLM}]\label{GaussianField-spectralCLM}
  \lean{GaussianField.spectralCLM}\leanok
  \uses{GaussianField-DyninMityaginSpace, GaussianField-DyninMityaginSpace-coeff, GaussianField-ell2-, GaussianField-IsBoundedSeq}
  \texttt{GaussianField.spectralCLM}
\end{definition}
\begin{theorem}[\texttt{isBoundedSeq\_of\_le}]\label{GaussianField-isBoundedSeq-of-le}
  \lean{GaussianField.isBoundedSeq_of_le}\leanok
  \uses{GaussianField-IsBoundedSeq}
  \texttt{GaussianField.isBoundedSeq\_of\_le}
\end{theorem}
\begin{theorem}[\texttt{qftSingularValue\_nonneg}]\label{GaussianField-qftSingularValue-nonneg}
  \lean{GaussianField.qftSingularValue_nonneg}\leanok
  \uses{GaussianField-qftSingularValue, GaussianField-qftEigenvalue, GaussianField-qftEigenvalue-pos}
  \texttt{GaussianField.qftSingularValue\_nonneg}
\end{theorem}
\begin{theorem}[\texttt{heatSingularValue\_factors}]\label{GaussianField-heatSingularValue-factors}
  \lean{GaussianField.heatSingularValue_factors}\leanok
  \uses{GaussianField-qftEigenvalue, GaussianField-heatSingularValue, GaussianField-SmoothMap-Circle-fourierFreq}
  \texttt{GaussianField.heatSingularValue\_factors}
\end{theorem}
\begin{theorem}[\texttt{heat\_singular\_values\_bounded}]\label{GaussianField-heat-singular-values-bounded}
  \lean{GaussianField.heat_singular_values_bounded}\leanok
  \uses{GaussianField-heatSingularValue-le-one, GaussianField-heatSingularValue, GaussianField-heatSingularValue-pos, GaussianField-IsBoundedSeq}
  \texttt{GaussianField.heat\_singular\_values\_bounded}
\end{theorem}
\begin{theorem}[\texttt{qftEigenvalue\_pos}]\label{GaussianField-qftEigenvalue-pos}
  \lean{GaussianField.qftEigenvalue_pos}\leanok
  \uses{GaussianField-qftEigenvalue, GaussianField-SmoothMap-Circle-fourierFreq, Lattice-toSemilatticeInf}
  \texttt{GaussianField.qftEigenvalue\_pos}
\end{theorem}
\begin{definition}[\texttt{IsBoundedSeq}]\label{GaussianField-IsBoundedSeq}
  \lean{GaussianField.IsBoundedSeq}\leanok
  \texttt{GaussianField.IsBoundedSeq}
\end{definition}
\begin{theorem}[\texttt{mul}]\label{GaussianField-IsBoundedSeq-mul}
  \lean{GaussianField.IsBoundedSeq.mul}\leanok
  \uses{Lattice-toSemilatticeInf, GaussianField-IsBoundedSeq}
  \texttt{GaussianField.IsBoundedSeq.mul}
\end{theorem}
\begin{theorem}[\texttt{spectralCLM\_coord}]\label{GaussianField-spectralCLM-coord}
  \lean{GaussianField.spectralCLM_coord}\leanok
  \uses{GaussianField-DyninMityaginSpace, GaussianField-DyninMityaginSpace-coeff, GaussianField-DyninMityaginSpace--, GaussianField-spectralCLM, GaussianField-DyninMityaginSpace-p, GaussianField-ell2-, GaussianField-nuclear-ell2-embedding-from-decay, GaussianField-IsBoundedSeq}
  \texttt{GaussianField.spectralCLM\_coord}
\end{theorem}
\begin{theorem}[\texttt{spectralCLM\_smul}]\label{GaussianField-spectralCLM-smul}
  \lean{GaussianField.spectralCLM_smul}\leanok
  \uses{GaussianField-DyninMityaginSpace, GaussianField-DyninMityaginSpace-coeff, GaussianField-spectralCLM-coord, GaussianField-spectralCLM, GaussianField-ell2-, GaussianField-IsBoundedSeq}
  \texttt{GaussianField.spectralCLM\_smul}
\end{theorem}
\begin{theorem}[\texttt{heatSingularValue\_le\_one}]\label{GaussianField-heatSingularValue-le-one}
  \lean{GaussianField.heatSingularValue_le_one}\leanok
  \uses{GaussianField-qftEigenvalue, GaussianField-heatSingularValue, GaussianField-qftEigenvalue-pos}
  \texttt{GaussianField.heatSingularValue\_le\_one}
\end{theorem}
\begin{definition}[\texttt{qftSingularValue}]\label{GaussianField-qftSingularValue}
  \lean{GaussianField.qftSingularValue}\leanok
  \uses{GaussianField-qftEigenvalue}
  \texttt{GaussianField.qftSingularValue}
\end{definition}
\begin{theorem}[\texttt{heatSingularValue\_pos}]\label{GaussianField-heatSingularValue-pos}
  \lean{GaussianField.heatSingularValue_pos}\leanok
  \uses{GaussianField-qftEigenvalue, GaussianField-heatSingularValue}
  \texttt{GaussianField.heatSingularValue\_pos}
\end{theorem}
\section{\texttt{Nuclear.DyninMityagin}}
\begin{theorem}[\texttt{expansion\_H}]\label{GaussianField-DyninMityaginSpace-expansion-H}
  \lean{GaussianField.DyninMityaginSpace.expansion_H}\leanok
  \uses{GaussianField-DyninMityaginSpace, GaussianField-DyninMityaginSpace-coeff, GaussianField-DyninMityaginSpace-expansion, GaussianField-DyninMityaginSpace-basis}
  \texttt{GaussianField.DyninMityaginSpace.expansion\_H}
\end{theorem}
\begin{definition}[\texttt{basis}]\label{GaussianField-DyninMityaginSpace-basis}
  \lean{GaussianField.DyninMityaginSpace.basis}\leanok
  \uses{GaussianField-DyninMityaginSpace}
  \texttt{GaussianField.DyninMityaginSpace.basis}
\end{definition}
\begin{theorem}[\texttt{expansion}]\label{GaussianField-DyninMityaginSpace-expansion}
  \lean{GaussianField.DyninMityaginSpace.expansion}\leanok
  \uses{GaussianField-DyninMityaginSpace, GaussianField-DyninMityaginSpace-coeff, GaussianField-DyninMityaginSpace-basis}
  \texttt{GaussianField.DyninMityaginSpace.expansion}
\end{theorem}
\begin{definition}[\texttt{p}]\label{GaussianField-DyninMityaginSpace-p}
  \lean{GaussianField.DyninMityaginSpace.p}\leanok
  \uses{GaussianField-DyninMityaginSpace, GaussianField-DyninMityaginSpace--}
  \texttt{GaussianField.DyninMityaginSpace.p}
\end{definition}
\begin{definition}[\texttt{DyninMityaginSpace}]\label{GaussianField-DyninMityaginSpace}
  \lean{GaussianField.DyninMityaginSpace}\leanok
  \texttt{GaussianField.DyninMityaginSpace}
\end{definition}
\begin{definition}[\texttt{coeff}]\label{GaussianField-DyninMityaginSpace-coeff}
  \lean{GaussianField.DyninMityaginSpace.coeff}\leanok
  \uses{GaussianField-DyninMityaginSpace}
  Lean signature ``\texttt{lean private noncomputable def coeff (T : E →L[ℝ] H) (h\_inf : ¬ FiniteDimensional ℝ H) (n : ℕ) (f : E) : ℝ }``  Informal: The $n$-th Fourier coefficient of $T(f)$ in the adapted ONB: $c_n(f) = \langle e_n, T(f) \rangle_H$.
\end{definition}
\begin{theorem}[\texttt{coeff\_decay}]\label{GaussianField-DyninMityaginSpace-coeff-decay}
  \lean{GaussianField.DyninMityaginSpace.coeff_decay}\leanok
  \uses{GaussianField-DyninMityaginSpace, GaussianField-DyninMityaginSpace-coeff, GaussianField-DyninMityaginSpace--, GaussianField-DyninMityaginSpace-p}
  \texttt{GaussianField.DyninMityaginSpace.coeff\_decay}
\end{theorem}
\begin{theorem}[\texttt{h\_with}]\label{GaussianField-DyninMityaginSpace-h-with}
  \lean{GaussianField.DyninMityaginSpace.h_with}\leanok
  \uses{GaussianField-DyninMityaginSpace, GaussianField-DyninMityaginSpace--, GaussianField-DyninMityaginSpace-p}
  \texttt{GaussianField.DyninMityaginSpace.h\_with}
\end{theorem}
\begin{definition}[\texttt{ι}]\label{GaussianField-DyninMityaginSpace--}
  \lean{GaussianField.DyninMityaginSpace.ι}\leanok
  \uses{GaussianField-DyninMityaginSpace}
  \texttt{GaussianField.DyninMityaginSpace.ι}
\end{definition}
\begin{theorem}[\texttt{basis\_growth}]\label{GaussianField-DyninMityaginSpace-basis-growth}
  \lean{GaussianField.DyninMityaginSpace.basis_growth}\leanok
  \uses{GaussianField-DyninMityaginSpace, GaussianField-DyninMityaginSpace--, GaussianField-DyninMityaginSpace-basis, GaussianField-DyninMityaginSpace-p}
  \texttt{GaussianField.DyninMityaginSpace.basis\_growth}
\end{theorem}
\begin{theorem}[\texttt{h\_countable}]\label{GaussianField-DyninMityaginSpace-h-countable}
  \lean{GaussianField.DyninMityaginSpace.h_countable}\leanok
  \uses{GaussianField-DyninMityaginSpace, GaussianField-DyninMityaginSpace--}
  \texttt{GaussianField.DyninMityaginSpace.h\_countable}
\end{theorem}
\begin{theorem}[\texttt{h\_completeSpace}]\label{GaussianField-DyninMityaginSpace-h-completeSpace}
  \lean{GaussianField.DyninMityaginSpace.h_completeSpace}\leanok
  \uses{GaussianField-DyninMityaginSpace}
  \texttt{GaussianField.DyninMityaginSpace.h\_completeSpace}
\end{theorem}
\begin{theorem}[\texttt{coeff\_basis}]\label{GaussianField-DyninMityaginSpace-HasBiorthogonalBasis-coeff-basis}
  \lean{GaussianField.DyninMityaginSpace.HasBiorthogonalBasis.coeff_basis}\leanok
  \uses{GaussianField-DyninMityaginSpace, GaussianField-DyninMityaginSpace-coeff, GaussianField-DyninMityaginSpace-basis, GaussianField-DyninMityaginSpace-HasBiorthogonalBasis}
  \texttt{GaussianField.DyninMityaginSpace.HasBiorthogonalBasis.coeff\_basis}
\end{theorem}
\begin{theorem}[\texttt{toT1Space}]\label{GaussianField-DyninMityaginSpace-toT1Space}
  \lean{GaussianField.DyninMityaginSpace.toT1Space}\leanok
  \uses{GaussianField-DyninMityaginSpace}
  \texttt{GaussianField.DyninMityaginSpace.toT1Space}
\end{theorem}
\begin{definition}[\texttt{ctorIdx}]\label{GaussianField-DyninMityaginSpace-ctorIdx}
  \lean{GaussianField.DyninMityaginSpace.ctorIdx}\leanok
  \uses{GaussianField-DyninMityaginSpace}
  \texttt{GaussianField.DyninMityaginSpace.ctorIdx}
\end{definition}
\begin{definition}[\texttt{HasBiorthogonalBasis}]\label{GaussianField-DyninMityaginSpace-HasBiorthogonalBasis}
  \lean{GaussianField.DyninMityaginSpace.HasBiorthogonalBasis}\leanok
  \uses{GaussianField-DyninMityaginSpace}
  \texttt{GaussianField.DyninMityaginSpace.HasBiorthogonalBasis}
\end{definition}
\section{\texttt{Nuclear.NuclearExtension}}
\begin{theorem}[\texttt{nuclearLift\_pure}]\label{GaussianField-nuclearLift-pure}
  \lean{GaussianField.nuclearLift_pure}\leanok
  \uses{GaussianField-DyninMityaginSpace, GaussianField-NuclearTensorProduct-instModuleReal, GaussianField-NuclearTensorProduct-instTopologicalSpace, GaussianField-NuclearTensorProduct-lift-pure, GaussianField-nuclearLift, GaussianField-DyninMityaginSpace--, GaussianField-NuclearTensorProduct, GaussianField-NuclearTensorProduct-instAddCommGroup, GaussianField-NuclearTensorProduct-pure, GaussianField-DyninMityaginSpace-p}
  \texttt{GaussianField.nuclearLift\_pure}
\end{theorem}
\begin{definition}[\texttt{nuclearLift}]\label{GaussianField-nuclearLift}
  \lean{GaussianField.nuclearLift}\leanok
  \uses{GaussianField-DyninMityaginSpace, GaussianField-NuclearTensorProduct-instModuleReal, GaussianField-NuclearTensorProduct-instTopologicalSpace, GaussianField-DyninMityaginSpace--, GaussianField-NuclearTensorProduct, GaussianField-NuclearTensorProduct-instAddCommGroup, GaussianField-NuclearTensorProduct-lift, GaussianField-DyninMityaginSpace-p}
  \texttt{GaussianField.nuclearLift}
\end{definition}
\begin{theorem}[\texttt{nuclear\_bilinear\_extension\_unique}]\label{GaussianField-nuclear-bilinear-extension-unique}
  \lean{GaussianField.nuclear_bilinear_extension_unique}\leanok
  \uses{GaussianField-DyninMityaginSpace, GaussianField-NuclearTensorProduct-instModuleReal, GaussianField-NuclearTensorProduct-instTopologicalSpace, GaussianField-RapidDecaySeq-val, GaussianField-NuclearTensorProduct, GaussianField-NuclearTensorProduct-instAddCommGroup, GaussianField-DyninMityaginSpace-basis, GaussianField-NuclearTensorProduct-pure, GaussianField-DyninMityaginSpace-HasBiorthogonalBasis, GaussianField-hasSum-pure}
  \texttt{GaussianField.nuclear\_bilinear\_extension\_unique}
\end{theorem}
\begin{theorem}[\texttt{nuclear\_bilinear\_extension\_existsUnique}]\label{GaussianField-nuclear-bilinear-extension-existsUnique}
  \lean{GaussianField.nuclear_bilinear_extension_existsUnique}\leanok
  \uses{GaussianField-DyninMityaginSpace, GaussianField-NuclearTensorProduct-instModuleReal, GaussianField-NuclearTensorProduct-instTopologicalSpace, GaussianField-NuclearTensorProduct-lift-pure, GaussianField-nuclear-bilinear-extension-unique, GaussianField-DyninMityaginSpace--, GaussianField-NuclearTensorProduct, GaussianField-NuclearTensorProduct-instAddCommGroup, GaussianField-NuclearTensorProduct-pure, GaussianField-NuclearTensorProduct-lift, GaussianField-DyninMityaginSpace-p, GaussianField-DyninMityaginSpace-HasBiorthogonalBasis}
  \texttt{GaussianField.nuclear\_bilinear\_extension\_existsUnique}
\end{theorem}
\begin{theorem}[\texttt{basisVec\_eq\_pure}]\label{GaussianField-basisVec-eq-pure}
  \lean{GaussianField.basisVec_eq_pure}\leanok
  \uses{GaussianField-DyninMityaginSpace, GaussianField-DyninMityaginSpace-coeff, GaussianField-RapidDecaySeq, GaussianField-RapidDecaySeq-val, GaussianField-DyninMityaginSpace-HasBiorthogonalBasis-coeff-basis, GaussianField-RapidDecaySeq-basisVec, GaussianField-RapidDecaySeq-ext, GaussianField-DyninMityaginSpace-basis, GaussianField-NuclearTensorProduct-pure, GaussianField-DyninMityaginSpace-HasBiorthogonalBasis}
  \texttt{GaussianField.basisVec\_eq\_pure}
\end{theorem}
\begin{theorem}[\texttt{hasSum\_pure}]\label{GaussianField-hasSum-pure}
  \lean{GaussianField.hasSum_pure}\leanok
  \uses{GaussianField-DyninMityaginSpace, GaussianField-NuclearTensorProduct-instModuleReal, GaussianField-NuclearTensorProduct-instTopologicalSpace, GaussianField-RapidDecaySeq, GaussianField-RapidDecaySeq-val, GaussianField-RapidDecaySeq-hasSum-basisVec, GaussianField-RapidDecaySeq-instSMul, GaussianField-NuclearTensorProduct, GaussianField-RapidDecaySeq-basisVec, GaussianField-NuclearTensorProduct-instAddCommGroup, GaussianField-DyninMityaginSpace-basis, GaussianField-NuclearTensorProduct-pure, GaussianField-RapidDecaySeq-instAddCommGroup, GaussianField-RapidDecaySeq-instTopologicalSpace, GaussianField-DyninMityaginSpace-HasBiorthogonalBasis, GaussianField-basisVec-eq-pure}
  \texttt{GaussianField.hasSum\_pure}
\end{theorem}
\begin{theorem}[\texttt{nuclear\_bilinear\_extension}]\label{GaussianField-nuclear-bilinear-extension}
  \lean{GaussianField.nuclear_bilinear_extension}\leanok
  \uses{GaussianField-DyninMityaginSpace, GaussianField-NuclearTensorProduct-instModuleReal, GaussianField-NuclearTensorProduct-instTopologicalSpace, GaussianField-NuclearTensorProduct-lift-pure, GaussianField-DyninMityaginSpace--, GaussianField-NuclearTensorProduct, GaussianField-NuclearTensorProduct-instAddCommGroup, GaussianField-NuclearTensorProduct-pure, GaussianField-NuclearTensorProduct-lift, GaussianField-DyninMityaginSpace-p}
  \texttt{GaussianField.nuclear\_bilinear\_extension}
\end{theorem}
\section{\texttt{Cylinder.FreeHeatSemigroup}}
\begin{theorem}[\texttt{congr\_simp}]\label{GaussianField-freeHeatSemigroup-congr-simp}
  \lean{GaussianField.freeHeatSemigroup.congr_simp}\leanok
  \uses{GaussianField-freeHeatSemigroup}
  \texttt{GaussianField.freeHeatSemigroup.congr\_simp}
\end{theorem}
\begin{theorem}[\texttt{cylinderHeatSemigroup\_timeTranslation\_comm}]\label{GaussianField-cylinderHeatSemigroup-timeTranslation-comm}
  \lean{GaussianField.cylinderHeatSemigroup_timeTranslation_comm}\leanok
  \uses{GaussianField-NuclearTensorProduct-instModuleReal, GaussianField-NuclearTensorProduct-instTopologicalSpace, GaussianField-nuclearTensorProduct-mapCLM-comp, GaussianField-nuclearTensorProduct-mapCLM, GaussianField-schwartz-dyninMityaginSpace-1D, GaussianField-SmoothMap-Circle-instAddCommGroup, GaussianField-CylinderTestFunction, GaussianField-NuclearTensorProduct, GaussianField-freeHeatSemigroup-translation-comm, GaussianField-cylinderHeatSemigroup, GaussianField-NuclearTensorProduct-instAddCommGroup, GaussianField-smoothCircle-hasBiorthogonalBasis, GaussianField-SmoothMap-Circle-instModule, GaussianField-SmoothMap-Circle-instContinuousSMul, GaussianField-smoothCircle-dyninMityaginSpace, GaussianField-SmoothMap-Circle-instIsTopologicalAddGroup, GaussianField-SmoothMap-Circle, GaussianField-freeHeatSemigroup, GaussianField-schwartzTranslation, GaussianField-circleHeatSemigroup, GaussianField-SmoothMap-Circle-instTopologicalSpace, GaussianField-schwartz-hasBiorthogonalBasis-1D, GaussianField-cylinderTimeTranslation}
  \texttt{GaussianField.cylinderHeatSemigroup\_timeTranslation\_comm}
\end{theorem}
\begin{theorem}[\texttt{congr\_simp}]\label{GaussianField-circleHeatSemigroup-congr-simp}
  \lean{GaussianField.circleHeatSemigroup.congr_simp}\leanok
  \uses{GaussianField-SmoothMap-Circle-instAddCommGroup, GaussianField-SmoothMap-Circle-instModule, GaussianField-SmoothMap-Circle, GaussianField-circleHeatSemigroup, GaussianField-SmoothMap-Circle-instTopologicalSpace}
  \texttt{GaussianField.circleHeatSemigroup.congr\_simp}
\end{theorem}
\begin{definition}[\texttt{cylinderHeatSemigroup}]\label{GaussianField-cylinderHeatSemigroup}
  \lean{GaussianField.cylinderHeatSemigroup}\leanok
  \uses{GaussianField-NuclearTensorProduct-instModuleReal, GaussianField-NuclearTensorProduct-instTopologicalSpace, GaussianField-nuclearTensorProduct-mapCLM, GaussianField-schwartz-dyninMityaginSpace-1D, GaussianField-SmoothMap-Circle-instAddCommGroup, GaussianField-CylinderTestFunction, GaussianField-NuclearTensorProduct, GaussianField-NuclearTensorProduct-instAddCommGroup, GaussianField-SmoothMap-Circle-instModule, GaussianField-SmoothMap-Circle-instContinuousSMul, GaussianField-smoothCircle-dyninMityaginSpace, GaussianField-SmoothMap-Circle-instIsTopologicalAddGroup, GaussianField-SmoothMap-Circle, GaussianField-freeHeatSemigroup, GaussianField-circleHeatSemigroup, GaussianField-SmoothMap-Circle-instTopologicalSpace}
  \texttt{GaussianField.cylinderHeatSemigroup}
\end{definition}
\begin{theorem}[\texttt{cylinderHeatSemigroup\_timeReflection\_comm}]\label{GaussianField-cylinderHeatSemigroup-timeReflection-comm}
  \lean{GaussianField.cylinderHeatSemigroup_timeReflection_comm}\leanok
  \uses{GaussianField-NuclearTensorProduct-instModuleReal, GaussianField-NuclearTensorProduct-instTopologicalSpace, GaussianField-nuclearTensorProduct-mapCLM-comp, GaussianField-nuclearTensorProduct-mapCLM, GaussianField-schwartz-dyninMityaginSpace-1D, GaussianField-SmoothMap-Circle-instAddCommGroup, GaussianField-CylinderTestFunction, GaussianField-NuclearTensorProduct, GaussianField-cylinderHeatSemigroup, GaussianField-NuclearTensorProduct-instAddCommGroup, GaussianField-smoothCircle-hasBiorthogonalBasis, GaussianField-freeHeatSemigroup-reflection-comm, GaussianField-cylinderTimeReflection, GaussianField-SmoothMap-Circle-instModule, GaussianField-SmoothMap-Circle-instContinuousSMul, GaussianField-smoothCircle-dyninMityaginSpace, GaussianField-SmoothMap-Circle-instIsTopologicalAddGroup, GaussianField-SmoothMap-Circle, GaussianField-freeHeatSemigroup, GaussianField-circleHeatSemigroup, GaussianField-schwartzReflection, GaussianField-SmoothMap-Circle-instTopologicalSpace, GaussianField-schwartz-hasBiorthogonalBasis-1D}
  \texttt{GaussianField.cylinderHeatSemigroup\_timeReflection\_comm}
\end{theorem}
\begin{theorem}[\texttt{cylinderHeatSemigroup\_spatialTranslation\_comm}]\label{GaussianField-cylinderHeatSemigroup-spatialTranslation-comm}
  \lean{GaussianField.cylinderHeatSemigroup_spatialTranslation_comm}\leanok
  \uses{GaussianField-NuclearTensorProduct-instModuleReal, GaussianField-circleTranslation, GaussianField-NuclearTensorProduct-instTopologicalSpace, GaussianField-nuclearTensorProduct-mapCLM-comp, GaussianField-cylinderSpatialTranslation, GaussianField-nuclearTensorProduct-mapCLM, GaussianField-schwartz-dyninMityaginSpace-1D, GaussianField-SmoothMap-Circle-instAddCommGroup, GaussianField-CylinderTestFunction, GaussianField-NuclearTensorProduct, GaussianField-cylinderHeatSemigroup, GaussianField-NuclearTensorProduct-instAddCommGroup, GaussianField-smoothCircle-hasBiorthogonalBasis, GaussianField-SmoothMap-Circle-instModule, GaussianField-SmoothMap-Circle-instContinuousSMul, GaussianField-smoothCircle-dyninMityaginSpace, GaussianField-SmoothMap-Circle-instIsTopologicalAddGroup, GaussianField-SmoothMap-Circle, GaussianField-freeHeatSemigroup, GaussianField-circleHeatSemigroup, GaussianField-SmoothMap-Circle-instTopologicalSpace, GaussianField-circleHeatSemigroup-translation-comm, GaussianField-schwartz-hasBiorthogonalBasis-1D}
  \texttt{GaussianField.cylinderHeatSemigroup\_spatialTranslation\_comm}
\end{theorem}
\section{\texttt{Cylinder.GreenFunction}}
\begin{theorem}[\texttt{cylinderMassOperator\_inner}]\label{GaussianField-cylinderMassOperator-inner}
  \lean{GaussianField.cylinderMassOperator_inner}\leanok
  \uses{GaussianField-NuclearTensorProduct-instModuleReal, GaussianField-NuclearTensorProduct-instTopologicalSpace, GaussianField-CylinderTestFunction, GaussianField-cylinderGreen, GaussianField-NuclearTensorProduct-instAddCommGroup, GaussianField-cylinderMassOperator, GaussianField-SmoothMap-Circle, GaussianField-ell2-}
  \texttt{GaussianField.cylinderMassOperator\_inner}
\end{theorem}
\begin{definition}[\texttt{cylinderTimeReflectionSym}]\label{GaussianField-cylinderTimeReflectionSym}
  \lean{GaussianField.cylinderTimeReflectionSym}\leanok
  \uses{GaussianField-CylinderSpacetimeSymmetry, GaussianField-CylinderTestFunction, GaussianField-cylinderMassOperator-timeReflection-norm-eq, GaussianField-cylinderTimeReflection, GaussianField-cylinderHeatSemigroup-timeReflection-comm}
  \texttt{GaussianField.cylinderTimeReflectionSym}
\end{definition}
\begin{theorem}[\texttt{cylinderGreen\_symm}]\label{GaussianField-cylinderGreen-symm}
  \lean{GaussianField.cylinderGreen_symm}\leanok
  \uses{GaussianField-NuclearTensorProduct-instModuleReal, GaussianField-NuclearTensorProduct-instTopologicalSpace, GaussianField-CylinderTestFunction, GaussianField-cylinderGreen, GaussianField-NuclearTensorProduct-instAddCommGroup, GaussianField-cylinderMassOperator, GaussianField-SmoothMap-Circle, GaussianField-ell2-}
  \texttt{GaussianField.cylinderGreen\_symm}
\end{theorem}
\begin{theorem}[\texttt{cylinderMassOperator\_timeReflection\_norm\_eq}]\label{GaussianField-cylinderMassOperator-timeReflection-norm-eq}
  \lean{GaussianField.cylinderMassOperator_timeReflection_norm_eq}\leanok
  \uses{GaussianField-NuclearTensorProduct-instModuleReal, GaussianField-NuclearTensorProduct-instTopologicalSpace, GaussianField-CylinderTestFunction, GaussianField-NuclearTensorProduct-instAddCommGroup, GaussianField-cylinderMassOperator-timeReflection-norm-eq-proved, GaussianField-cylinderTimeReflection, GaussianField-cylinderMassOperator, GaussianField-SmoothMap-Circle, GaussianField-ell2-}
  \texttt{GaussianField.cylinderMassOperator\_timeReflection\_norm\_eq}
\end{theorem}
\begin{theorem}[\texttt{cylinderMassOperator\_timeReflection\_equivariant}]\label{GaussianField-cylinderMassOperator-timeReflection-equivariant}
  \lean{GaussianField.cylinderMassOperator_timeReflection_equivariant}\leanok
  \uses{GaussianField-NuclearTensorProduct-instModuleReal, GaussianField-NuclearTensorProduct-instTopologicalSpace, GaussianField-cylinderTimeReflectionSym, GaussianField-CylinderTestFunction, GaussianField-NuclearTensorProduct-instAddCommGroup, GaussianField-cylinderTimeReflection, GaussianField-cylinderMassOperator, GaussianField-SmoothMap-Circle, GaussianField-cylinderMassOperator-equivariant, GaussianField-ell2-}
  \texttt{GaussianField.cylinderMassOperator\_timeReflection\_equivariant}
\end{theorem}
\begin{definition}[\texttt{inverse}]\label{GaussianField-CylinderSpacetimeSymmetry-inverse}
  \lean{GaussianField.CylinderSpacetimeSymmetry.inverse}\leanok
  \uses{GaussianField-NuclearTensorProduct-instModuleReal, GaussianField-NuclearTensorProduct-instTopologicalSpace, GaussianField-CylinderSpacetimeSymmetry, GaussianField-CylinderTestFunction, GaussianField-NuclearTensorProduct-instAddCommGroup, GaussianField-SmoothMap-Circle}
  \texttt{GaussianField.CylinderSpacetimeSymmetry.inverse}
\end{definition}
\begin{theorem}[\texttt{cylinderGreen\_timeReflection\_invariant}]\label{GaussianField-cylinderGreen-timeReflection-invariant}
  \lean{GaussianField.cylinderGreen_timeReflection_invariant}\leanok
  \uses{GaussianField-NuclearTensorProduct-instModuleReal, GaussianField-NuclearTensorProduct-instTopologicalSpace, GaussianField-CylinderTestFunction, GaussianField-cylinderGreen, GaussianField-NuclearTensorProduct-instAddCommGroup, GaussianField-cylinderTimeReflection, GaussianField-cylinderMassOperator, GaussianField-SmoothMap-Circle, GaussianField-cylinderMassOperator-timeReflection-equivariant, GaussianField-ell2-}
  \texttt{GaussianField.cylinderGreen\_timeReflection\_invariant}
\end{theorem}
\begin{theorem}[\texttt{cylinderGreen\_bilinear}]\label{GaussianField-cylinderGreen-bilinear}
  \lean{GaussianField.cylinderGreen_bilinear}\leanok
  \uses{GaussianField-NuclearTensorProduct-instModuleReal, GaussianField-NuclearTensorProduct-instTopologicalSpace, GaussianField-CylinderTestFunction, GaussianField-cylinderGreen, GaussianField-NuclearTensorProduct-instAddCommGroup, GaussianField-cylinderMassOperator, GaussianField-SmoothMap-Circle, GaussianField-ell2-}
  \texttt{GaussianField.cylinderGreen\_bilinear}
\end{theorem}
\begin{theorem}[\texttt{cylinderMassOperator\_spatialTranslation\_equivariant}]\label{GaussianField-cylinderMassOperator-spatialTranslation-equivariant}
  \lean{GaussianField.cylinderMassOperator_spatialTranslation_equivariant}\leanok
  \uses{GaussianField-NuclearTensorProduct-instModuleReal, GaussianField-NuclearTensorProduct-instTopologicalSpace, GaussianField-cylinderSpatialTranslation, GaussianField-CylinderTestFunction, GaussianField-cylinderSpatialTranslationSym, GaussianField-NuclearTensorProduct-instAddCommGroup, GaussianField-cylinderMassOperator, GaussianField-SmoothMap-Circle, GaussianField-cylinderMassOperator-equivariant, GaussianField-ell2-}
  \texttt{GaussianField.cylinderMassOperator\_spatialTranslation\_equivariant}
\end{theorem}
\begin{definition}[\texttt{toCLM}]\label{GaussianField-CylinderSpacetimeSymmetry-toCLM}
  \lean{GaussianField.CylinderSpacetimeSymmetry.toCLM}\leanok
  \uses{GaussianField-NuclearTensorProduct-instModuleReal, GaussianField-NuclearTensorProduct-instTopologicalSpace, GaussianField-CylinderSpacetimeSymmetry, GaussianField-CylinderTestFunction, GaussianField-NuclearTensorProduct-instAddCommGroup, GaussianField-SmoothMap-Circle}
  \texttt{GaussianField.CylinderSpacetimeSymmetry.toCLM}
\end{definition}
\begin{theorem}[\texttt{preserves\_T\_norm\_inv}]\label{GaussianField-CylinderSpacetimeSymmetry-preserves-T-norm-inv}
  \lean{GaussianField.CylinderSpacetimeSymmetry.preserves_T_norm_inv}\leanok
  \uses{GaussianField-NuclearTensorProduct-instModuleReal, GaussianField-NuclearTensorProduct-instTopologicalSpace, GaussianField-CylinderSpacetimeSymmetry, GaussianField-CylinderTestFunction, GaussianField-NuclearTensorProduct-instAddCommGroup, GaussianField-cylinderMassOperator, GaussianField-SmoothMap-Circle, GaussianField-CylinderSpacetimeSymmetry-inverse, GaussianField-ell2-}
  \texttt{GaussianField.CylinderSpacetimeSymmetry.preserves\_T\_norm\_inv}
\end{theorem}
\begin{theorem}[\texttt{cylinderGreen\_nonneg}]\label{GaussianField-cylinderGreen-nonneg}
  \lean{GaussianField.cylinderGreen_nonneg}\leanok
  \uses{GaussianField-NuclearTensorProduct-instModuleReal, GaussianField-NuclearTensorProduct-instTopologicalSpace, GaussianField-CylinderTestFunction, GaussianField-cylinderGreen, GaussianField-NuclearTensorProduct-instAddCommGroup, GaussianField-cylinderMassOperator, GaussianField-SmoothMap-Circle, GaussianField-ell2-}
  \texttt{GaussianField.cylinderGreen\_nonneg}
\end{theorem}
\begin{definition}[\texttt{cylinderTimeTranslationSym}]\label{GaussianField-cylinderTimeTranslationSym}
  \lean{GaussianField.cylinderTimeTranslationSym}\leanok
  \uses{GaussianField-cylinderMassOperator-timeTranslation-norm-eq, GaussianField-CylinderSpacetimeSymmetry, GaussianField-cylinderHeatSemigroup-timeTranslation-comm, GaussianField-CylinderTestFunction, GaussianField-cylinderTimeTranslation}
  \texttt{GaussianField.cylinderTimeTranslationSym}
\end{definition}
\begin{definition}[\texttt{CylinderSpacetimeSymmetry}]\label{GaussianField-CylinderSpacetimeSymmetry}
  \lean{GaussianField.CylinderSpacetimeSymmetry}\leanok
  \texttt{GaussianField.CylinderSpacetimeSymmetry}
\end{definition}
\begin{theorem}[\texttt{heat\_comm}]\label{GaussianField-CylinderSpacetimeSymmetry-heat-comm}
  \lean{GaussianField.CylinderSpacetimeSymmetry.heat_comm}\leanok
  \uses{GaussianField-NuclearTensorProduct-instModuleReal, GaussianField-NuclearTensorProduct-instTopologicalSpace, GaussianField-CylinderSpacetimeSymmetry-toCLM, GaussianField-CylinderSpacetimeSymmetry, GaussianField-CylinderTestFunction, GaussianField-cylinderHeatSemigroup, GaussianField-NuclearTensorProduct-instAddCommGroup, GaussianField-SmoothMap-Circle}
  \texttt{GaussianField.CylinderSpacetimeSymmetry.heat\_comm}
\end{theorem}
\begin{theorem}[\texttt{cylinderMassOperator\_timeTranslation\_norm\_eq}]\label{GaussianField-cylinderMassOperator-timeTranslation-norm-eq}
  \lean{GaussianField.cylinderMassOperator_timeTranslation_norm_eq}\leanok
  \uses{GaussianField-NuclearTensorProduct-instModuleReal, GaussianField-NuclearTensorProduct-instTopologicalSpace, GaussianField-CylinderTestFunction, GaussianField-NuclearTensorProduct-instAddCommGroup, GaussianField-cylinderMassOperator-timeTranslation-norm-eq-proved, GaussianField-cylinderMassOperator, GaussianField-SmoothMap-Circle, GaussianField-ell2-, GaussianField-cylinderTimeTranslation}
  \texttt{GaussianField.cylinderMassOperator\_timeTranslation\_norm\_eq}
\end{theorem}
\begin{theorem}[\texttt{cylinderGreen\_spatialTranslation\_invariant}]\label{GaussianField-cylinderGreen-spatialTranslation-invariant}
  \lean{GaussianField.cylinderGreen_spatialTranslation_invariant}\leanok
  \uses{GaussianField-NuclearTensorProduct-instModuleReal, GaussianField-NuclearTensorProduct-instTopologicalSpace, GaussianField-cylinderSpatialTranslation, GaussianField-CylinderTestFunction, GaussianField-cylinderGreen, GaussianField-cylinderMassOperator-spatialTranslation-equivariant, GaussianField-NuclearTensorProduct-instAddCommGroup, GaussianField-cylinderMassOperator, GaussianField-SmoothMap-Circle, GaussianField-ell2-}
  \texttt{GaussianField.cylinderGreen\_spatialTranslation\_invariant}
\end{theorem}
\begin{theorem}[\texttt{inverse\_comp\_toCLM}]\label{GaussianField-CylinderSpacetimeSymmetry-inverse-comp-toCLM}
  \lean{GaussianField.CylinderSpacetimeSymmetry.inverse_comp_toCLM}\leanok
  \uses{GaussianField-NuclearTensorProduct-instModuleReal, GaussianField-NuclearTensorProduct-instTopologicalSpace, GaussianField-CylinderSpacetimeSymmetry-toCLM, GaussianField-CylinderSpacetimeSymmetry, GaussianField-CylinderTestFunction, GaussianField-NuclearTensorProduct-instAddCommGroup, GaussianField-SmoothMap-Circle, GaussianField-CylinderSpacetimeSymmetry-inverse}
  \texttt{GaussianField.CylinderSpacetimeSymmetry.inverse\_comp\_toCLM}
\end{theorem}
\begin{theorem}[\texttt{cylinderGreen\_pos}]\label{GaussianField-cylinderGreen-pos}
  \lean{GaussianField.cylinderGreen_pos}\leanok
  \uses{GaussianField-NuclearTensorProduct-instModuleReal, GaussianField-NuclearTensorProduct-instTopologicalSpace, GaussianField-cylinderMassOperator-injective, GaussianField-CylinderTestFunction, GaussianField-cylinderGreen, GaussianField-NuclearTensorProduct-instAddCommGroup, GaussianField-cylinderMassOperator, GaussianField-SmoothMap-Circle, GaussianField-ell2-}
  \texttt{GaussianField.cylinderGreen\_pos}
\end{theorem}
\begin{theorem}[\texttt{toCLM\_comp\_inverse}]\label{GaussianField-CylinderSpacetimeSymmetry-toCLM-comp-inverse}
  \lean{GaussianField.CylinderSpacetimeSymmetry.toCLM_comp_inverse}\leanok
  \uses{GaussianField-NuclearTensorProduct-instModuleReal, GaussianField-NuclearTensorProduct-instTopologicalSpace, GaussianField-CylinderSpacetimeSymmetry-toCLM, GaussianField-CylinderSpacetimeSymmetry, GaussianField-CylinderTestFunction, GaussianField-NuclearTensorProduct-instAddCommGroup, GaussianField-SmoothMap-Circle, GaussianField-CylinderSpacetimeSymmetry-inverse}
  \texttt{GaussianField.CylinderSpacetimeSymmetry.toCLM\_comp\_inverse}
\end{theorem}
\begin{theorem}[\texttt{cylinderMassOperator\_spatialTranslation\_norm\_eq}]\label{GaussianField-cylinderMassOperator-spatialTranslation-norm-eq}
  \lean{GaussianField.cylinderMassOperator_spatialTranslation_norm_eq}\leanok
  \uses{GaussianField-NuclearTensorProduct-instModuleReal, GaussianField-NuclearTensorProduct-instTopologicalSpace, GaussianField-cylinderSpatialTranslation, GaussianField-CylinderTestFunction, GaussianField-NuclearTensorProduct-instAddCommGroup, GaussianField-cylinderMassOperator-spatialTranslation-norm-eq-proved, GaussianField-cylinderMassOperator, GaussianField-SmoothMap-Circle, GaussianField-ell2-}
  \texttt{GaussianField.cylinderMassOperator\_spatialTranslation\_norm\_eq}
\end{theorem}
\begin{theorem}[\texttt{cylinderMassOperator\_timeTranslation\_equivariant}]\label{GaussianField-cylinderMassOperator-timeTranslation-equivariant}
  \lean{GaussianField.cylinderMassOperator_timeTranslation_equivariant}\leanok
  \uses{GaussianField-NuclearTensorProduct-instModuleReal, GaussianField-NuclearTensorProduct-instTopologicalSpace, GaussianField-cylinderTimeTranslationSym, GaussianField-CylinderTestFunction, GaussianField-NuclearTensorProduct-instAddCommGroup, GaussianField-cylinderMassOperator, GaussianField-SmoothMap-Circle, GaussianField-cylinderMassOperator-equivariant, GaussianField-ell2-, GaussianField-cylinderTimeTranslation}
  \texttt{GaussianField.cylinderMassOperator\_timeTranslation\_equivariant}
\end{theorem}
\begin{definition}[\texttt{ctorIdx}]\label{GaussianField-CylinderSpacetimeSymmetry-ctorIdx}
  \lean{GaussianField.CylinderSpacetimeSymmetry.ctorIdx}\leanok
  \uses{GaussianField-CylinderSpacetimeSymmetry}
  \texttt{GaussianField.CylinderSpacetimeSymmetry.ctorIdx}
\end{definition}
\begin{theorem}[\texttt{preserves\_T\_norm}]\label{GaussianField-CylinderSpacetimeSymmetry-preserves-T-norm}
  \lean{GaussianField.CylinderSpacetimeSymmetry.preserves_T_norm}\leanok
  \uses{GaussianField-NuclearTensorProduct-instModuleReal, GaussianField-NuclearTensorProduct-instTopologicalSpace, GaussianField-CylinderSpacetimeSymmetry-toCLM, GaussianField-CylinderSpacetimeSymmetry, GaussianField-CylinderTestFunction, GaussianField-NuclearTensorProduct-instAddCommGroup, GaussianField-cylinderMassOperator, GaussianField-SmoothMap-Circle, GaussianField-ell2-}
  \texttt{GaussianField.CylinderSpacetimeSymmetry.preserves\_T\_norm}
\end{theorem}
\begin{definition}[\texttt{cylinderGreen}]\label{GaussianField-cylinderGreen}
  \lean{GaussianField.cylinderGreen}\leanok
  \uses{GaussianField-NuclearTensorProduct-instModuleReal, GaussianField-NuclearTensorProduct-instTopologicalSpace, GaussianField-CylinderTestFunction, GaussianField-NuclearTensorProduct-instAddCommGroup, GaussianField-cylinderMassOperator, GaussianField-SmoothMap-Circle, GaussianField-ell2-}
  \texttt{GaussianField.cylinderGreen}
\end{definition}
\begin{theorem}[\texttt{cylinderGreen\_continuous\_diag}]\label{GaussianField-cylinderGreen-continuous-diag}
  \lean{GaussianField.cylinderGreen_continuous_diag}\leanok
  \uses{GaussianField-NuclearTensorProduct-instModuleReal, GaussianField-NuclearTensorProduct-instTopologicalSpace, GaussianField-CylinderTestFunction, GaussianField-cylinderGreen, GaussianField-NuclearTensorProduct-instAddCommGroup, GaussianField-cylinderMassOperator, GaussianField-SmoothMap-Circle, GaussianField-ell2-}
  \texttt{GaussianField.cylinderGreen\_continuous\_diag}
\end{theorem}
\begin{theorem}[\texttt{cylinderGreen\_timeTranslation\_invariant}]\label{GaussianField-cylinderGreen-timeTranslation-invariant}
  \lean{GaussianField.cylinderGreen_timeTranslation_invariant}\leanok
  \uses{GaussianField-NuclearTensorProduct-instModuleReal, GaussianField-NuclearTensorProduct-instTopologicalSpace, GaussianField-CylinderTestFunction, GaussianField-cylinderGreen, GaussianField-cylinderMassOperator-timeTranslation-equivariant, GaussianField-NuclearTensorProduct-instAddCommGroup, GaussianField-cylinderMassOperator, GaussianField-SmoothMap-Circle, GaussianField-ell2-, GaussianField-cylinderTimeTranslation}
  \texttt{GaussianField.cylinderGreen\_timeTranslation\_invariant}
\end{theorem}
\begin{theorem}[\texttt{cylinderMassOperator\_injective}]\label{GaussianField-cylinderMassOperator-injective}
  \lean{GaussianField.cylinderMassOperator_injective}\leanok
  \uses{GaussianField-NuclearTensorProduct-ext, GaussianField-NuclearTensorProduct-instModuleReal, GaussianField-NuclearTensorProduct-instTopologicalSpace, GaussianField-cylinderMassOperator-formula, GaussianField-DyninMityaginSpace-coeff, GaussianField-RapidDecaySeq-instZero, GaussianField-ntpExtractSlice, GaussianField-RapidDecaySeq, GaussianField-schwartz-dyninMityaginSpace-1D, GaussianField-RapidDecaySeq-val, GaussianField-CylinderTestFunction, GaussianField-resolventMultiplierCLM-injective, GaussianField-RapidDecaySeq-instModule, GaussianField-resolventMultiplierCLM-congr-simp, GaussianField-NuclearTensorProduct-instAddCommGroup, GaussianField-RapidDecaySeq-ext, GaussianField-ntpExtractSlice-val, GaussianField-resolventFreq-pos, GaussianField-resolventFreq, GaussianField-cylinderMassOperator, GaussianField-RapidDecaySeq-instAddCommGroup, GaussianField-SmoothMap-Circle, GaussianField-ntpSliceSchwartz, GaussianField-RapidDecaySeq-instTopologicalSpace, GaussianField-resolventMultiplierCLM, GaussianField-ell2-, GaussianField-schwartzRapidDecayEquiv1D}
  \texttt{GaussianField.cylinderMassOperator\_injective}
\end{theorem}
\begin{definition}[\texttt{cylinderSpatialTranslationSym}]\label{GaussianField-cylinderSpatialTranslationSym}
  \lean{GaussianField.cylinderSpatialTranslationSym}\leanok
  \uses{GaussianField-cylinderHeatSemigroup-spatialTranslation-comm, GaussianField-cylinderSpatialTranslation, GaussianField-cylinderMassOperator-spatialTranslation-norm-eq, GaussianField-CylinderSpacetimeSymmetry, GaussianField-CylinderTestFunction}
  \texttt{GaussianField.cylinderSpatialTranslationSym}
\end{definition}
\begin{theorem}[\texttt{cylinderMassOperator\_equivariant}]\label{GaussianField-cylinderMassOperator-equivariant}
  \lean{GaussianField.cylinderMassOperator_equivariant}\leanok
  \uses{GaussianField-NuclearTensorProduct-instModuleReal, GaussianField-CylinderSpacetimeSymmetry-inverse-comp-toCLM, GaussianField-NuclearTensorProduct-instTopologicalSpace, GaussianField-CylinderSpacetimeSymmetry-toCLM, GaussianField-CylinderSpacetimeSymmetry, GaussianField-cylinderMassOperator-range-dense, GaussianField-CylinderTestFunction, GaussianField-NuclearTensorProduct-instAddCommGroup, GaussianField-CylinderSpacetimeSymmetry-toCLM-comp-inverse, GaussianField-cylinderMassOperator, GaussianField-SmoothMap-Circle, GaussianField-CylinderSpacetimeSymmetry-preserves-T-norm, GaussianField-CylinderSpacetimeSymmetry-inverse, GaussianField-ell2-, GaussianField-CylinderSpacetimeSymmetry-preserves-T-norm-inv}
  \texttt{GaussianField.cylinderMassOperator\_equivariant}
\end{theorem}
\section{\texttt{SmoothCircle.HeatEquivariance}}
\begin{theorem}[\texttt{circleHeatSemigroup\_reflection\_comm}]\label{GaussianField-circleHeatSemigroup-reflection-comm}
  \lean{GaussianField.circleHeatSemigroup_reflection_comm}\leanok
  \uses{GaussianField-fourierCoeffReal-circleReflection-zero, GaussianField-circleHeatSemigroup-fourierCoeff, GaussianField-RapidDecaySeq, GaussianField-RapidDecaySeq-val, GaussianField-SmoothMap-Circle-smoothCircleRapidDecayEquiv, GaussianField-SmoothMap-Circle-instAddCommGroup, GaussianField-heatWeight, GaussianField-SmoothMap-Circle-fourierCoeffReal, GaussianField-RapidDecaySeq-instModule, GaussianField-RapidDecaySeq-ext, GaussianField-circleReflection, GaussianField-SmoothMap-Circle-instModule, GaussianField-RapidDecaySeq-instAddCommGroup, GaussianField-SmoothMap-Circle, GaussianField-RapidDecaySeq-instTopologicalSpace, GaussianField-fourierCoeffReal-circleReflection-cos, GaussianField-circleHeatSemigroup, GaussianField-fourierCoeffReal-circleReflection-sin, GaussianField-SmoothMap-Circle-instTopologicalSpace}
  \texttt{GaussianField.circleHeatSemigroup\_reflection\_comm}
\end{theorem}
\begin{theorem}[\texttt{eigenvalue\_paired}]\label{GaussianField-eigenvalue-paired}
  \lean{GaussianField.eigenvalue_paired}\leanok
  \uses{GaussianField-fourierFreq-paired, GaussianField-circleHasLaplacianEigenvalues, GaussianField-SmoothMap-Circle-instAddCommGroup, GaussianField-HasLaplacianEigenvalues-eigenvalue, GaussianField-SmoothMap-Circle-instModule, GaussianField-SmoothMap-Circle-instContinuousSMul, GaussianField-smoothCircle-dyninMityaginSpace, GaussianField-SmoothMap-Circle-instIsTopologicalAddGroup, GaussianField-SmoothMap-Circle, GaussianField-SmoothMap-Circle-instTopologicalSpace, GaussianField-SmoothMap-Circle-fourierFreq}
  \texttt{GaussianField.eigenvalue\_paired}
\end{theorem}
\begin{theorem}[\texttt{circleHeatSemigroup\_translation\_comm}]\label{GaussianField-circleHeatSemigroup-translation-comm}
  \lean{GaussianField.circleHeatSemigroup_translation_comm}\leanok
  \uses{GaussianField-circleTranslation, GaussianField-circleHeatSemigroup-fourierCoeff, GaussianField-RapidDecaySeq, GaussianField-RapidDecaySeq-val, GaussianField-SmoothMap-Circle-smoothCircleRapidDecayEquiv, GaussianField-SmoothMap-Circle-instAddCommGroup, GaussianField-heatWeight, GaussianField-SmoothMap-Circle-fourierCoeffReal, GaussianField-RapidDecaySeq-instModule, GaussianField-fourierCoeffReal-circleTranslation-sin, GaussianField-RapidDecaySeq-ext, GaussianField-fourierCoeffReal-circleTranslation-zero, GaussianField-heatWeight-paired, GaussianField-SmoothMap-Circle-instModule, GaussianField-RapidDecaySeq-instAddCommGroup, GaussianField-SmoothMap-Circle, GaussianField-RapidDecaySeq-instTopologicalSpace, GaussianField-fourierCoeffReal-circleTranslation-cos, GaussianField-circleHeatSemigroup, GaussianField-SmoothMap-Circle-instTopologicalSpace}
  \texttt{GaussianField.circleHeatSemigroup\_translation\_comm}
\end{theorem}
\begin{theorem}[\texttt{heatWeight\_paired}]\label{GaussianField-heatWeight-paired}
  \lean{GaussianField.heatWeight_paired}\leanok
  \uses{GaussianField-eigenvalue-paired, GaussianField-circleHasLaplacianEigenvalues, GaussianField-SmoothMap-Circle-instAddCommGroup, GaussianField-heatWeight, GaussianField-HasLaplacianEigenvalues-eigenvalue, GaussianField-SmoothMap-Circle-instModule, GaussianField-SmoothMap-Circle-instContinuousSMul, GaussianField-smoothCircle-dyninMityaginSpace, GaussianField-SmoothMap-Circle-instIsTopologicalAddGroup, GaussianField-SmoothMap-Circle, GaussianField-SmoothMap-Circle-instTopologicalSpace}
  \texttt{GaussianField.heatWeight\_paired}
\end{theorem}
\begin{theorem}[\texttt{fourierFreq\_paired}]\label{GaussianField-fourierFreq-paired}
  \lean{GaussianField.fourierFreq_paired}\leanok
  \uses{GaussianField-SmoothMap-Circle-fourierFreq}
  \texttt{GaussianField.fourierFreq\_paired}
\end{theorem}
\section{\texttt{GaussianField.Properties}}
\begin{theorem}[\texttt{congr\_simp}]\label{GaussianField-measure-congr-simp}
  \lean{GaussianField.measure.congr_simp}\leanok
  \uses{GaussianField-DyninMityaginSpace, GaussianField-measure, GaussianField-Configuration, GaussianField-instMeasurableSpaceConfiguration}
  \texttt{GaussianField.measure.congr\_simp}
\end{theorem}
\begin{theorem}[\texttt{cross\_moment\_eq\_covariance}]\label{GaussianField-cross-moment-eq-covariance}
  \lean{GaussianField.cross_moment_eq_covariance}\leanok
  \uses{GaussianField-DyninMityaginSpace, GaussianField-measure, GaussianField-second-moment-eq-covariance, GaussianField-pairing-memLp, GaussianField-Configuration, GaussianField-instMeasurableSpaceConfiguration, GaussianField-pairing-product-integrable}
  \texttt{GaussianField.cross\_moment\_eq\_covariance}
\end{theorem}
\begin{theorem}[\texttt{pairing\_is\_gaussian}]\label{GaussianField-pairing-is-gaussian}
  \lean{GaussianField.pairing_is_gaussian}\leanok
  \uses{GaussianField-DyninMityaginSpace, GaussianField-charFun, GaussianField-measure, GaussianField-configuration-eval-measurable, GaussianField-Configuration, GaussianField-instMeasurableSpaceConfiguration, GaussianField-measure-isProbability}
  \texttt{GaussianField.pairing\_is\_gaussian}
\end{theorem}
\begin{theorem}[\texttt{measure\_centered}]\label{GaussianField-measure-centered}
  \lean{GaussianField.measure_centered}\leanok
  \uses{GaussianField-DyninMityaginSpace, GaussianField-measure, GaussianField-pairing-is-gaussian, GaussianField-configuration-eval-measurable, GaussianField-Configuration, GaussianField-instMeasurableSpaceConfiguration}
  \texttt{GaussianField.measure\_centered}
\end{theorem}
\begin{theorem}[\texttt{pairing\_memLp}]\label{GaussianField-pairing-memLp}
  \lean{GaussianField.pairing_memLp}\leanok
  \uses{GaussianField-DyninMityaginSpace, GaussianField-measure, GaussianField-pairing-is-gaussian, GaussianField-configuration-eval-measurable, GaussianField-Configuration, GaussianField-instMeasurableSpaceConfiguration}
  \texttt{GaussianField.pairing\_memLp}
\end{theorem}
\begin{theorem}[\texttt{pairing\_integrable}]\label{GaussianField-pairing-integrable}
  \lean{GaussianField.pairing_integrable}\leanok
  \uses{GaussianField-DyninMityaginSpace, GaussianField-measure, GaussianField-pairing-memLp, GaussianField-Configuration, GaussianField-instMeasurableSpaceConfiguration}
  \texttt{GaussianField.pairing\_integrable}
\end{theorem}
\begin{theorem}[\texttt{pairing\_product\_integrable}]\label{GaussianField-pairing-product-integrable}
  \lean{GaussianField.pairing_product_integrable}\leanok
  \uses{GaussianField-DyninMityaginSpace, GaussianField-measure, GaussianField-pairing-memLp, GaussianField-Configuration, GaussianField-instMeasurableSpaceConfiguration}
  \texttt{GaussianField.pairing\_product\_integrable}
\end{theorem}
\begin{theorem}[\texttt{second\_moment\_eq\_covariance}]\label{GaussianField-second-moment-eq-covariance}
  \lean{GaussianField.second_moment_eq_covariance}\leanok
  \uses{GaussianField-DyninMityaginSpace, GaussianField-measure, GaussianField-measure-centered, GaussianField-pairing-is-gaussian, GaussianField-configuration-eval-measurable, GaussianField-Configuration, GaussianField-instMeasurableSpaceConfiguration}
  \texttt{GaussianField.second\_moment\_eq\_covariance}
\end{theorem}
\section{\texttt{SmoothCircle.Nuclear}}
\begin{theorem}[\texttt{fromRapidDecay\_toRapidDecay}]\label{GaussianField-SmoothMap-Circle-fromRapidDecay-toRapidDecay}
  \lean{GaussianField.SmoothMap_Circle.fromRapidDecay_toRapidDecay}\leanok
  \uses{GaussianField-SmoothMap-Circle-toRapidDecay, GaussianField-SmoothMap-Circle-fromRapidDecay, GaussianField-SmoothMap-Circle, GaussianField-SmoothMap-Circle-hasSum-fourierBasis}
  \texttt{GaussianField.SmoothMap\_Circle.fromRapidDecay\_toRapidDecay}
\end{theorem}
\begin{definition}[\texttt{toRapidDecayLM}]\label{GaussianField-SmoothMap-Circle-toRapidDecayLM}
  \lean{GaussianField.SmoothMap_Circle.toRapidDecayLM}\leanok
  \uses{GaussianField-SmoothMap-Circle-toRapidDecay, GaussianField-RapidDecaySeq, GaussianField-SmoothMap-Circle-instAddCommGroup, GaussianField-RapidDecaySeq-instModule, GaussianField-SmoothMap-Circle-instModule, GaussianField-RapidDecaySeq-instAddCommGroup, GaussianField-SmoothMap-Circle}
  \texttt{GaussianField.SmoothMap\_Circle.toRapidDecayLM}
\end{definition}
\begin{theorem}[\texttt{contDiff\_fourierSeries}]\label{GaussianField-SmoothMap-Circle-contDiff-fourierSeries}
  \lean{GaussianField.SmoothMap_Circle.contDiff_fourierSeries}\leanok
  \uses{GaussianField-RapidDecaySeq, GaussianField-SmoothMap-Circle-fourierBasis, GaussianField-RapidDecaySeq-val, GaussianField-SmoothMap-Circle-instAddCommGroup, GaussianField-SmoothMap-Circle-fourierBasisFun-smooth, GaussianField-SmoothMap-Circle-fourierBasisFun, GaussianField-SmoothMap-Circle-instSMul, GaussianField-SmoothMap-Circle, GaussianField-SmoothMap-Circle-sobolevSeminorm-nonneg, GaussianField-SmoothMap-Circle-sobolevSeminorm, GaussianField-RapidDecaySeq-rapid-decay, GaussianField-SmoothMap-Circle-sobolevSeminorm-fourierBasis-le, Lattice-toSemilatticeInf}
  \texttt{GaussianField.SmoothMap\_Circle.contDiff\_fourierSeries}
\end{theorem}
\begin{theorem}[\texttt{toRapidDecay\_fromRapidDecay}]\label{GaussianField-SmoothMap-Circle-toRapidDecay-fromRapidDecay}
  \lean{GaussianField.SmoothMap_Circle.toRapidDecay_fromRapidDecay}\leanok
  \uses{GaussianField-SmoothMap-Circle-toRapidDecay, GaussianField-RapidDecaySeq, GaussianField-SmoothMap-Circle-fourierBasis, GaussianField-SmoothMap-Circle-fourierBasisFun-abs-le, GaussianField-RapidDecaySeq-val, GaussianField-SmoothMap-Circle-fromRapidDecay, GaussianField-SmoothMap-Circle-fourierBasisFun-smooth, GaussianField-SmoothMap-Circle-fourierCoeffReal-fourierBasis, GaussianField-SmoothMap-Circle-fourierCoeffReal, GaussianField-SmoothMap-Circle-fourierBasisFun, GaussianField-RapidDecaySeq-ext, GaussianField-RapidDecaySeq-rapid-decay, Lattice-toSemilatticeInf}
  \texttt{GaussianField.SmoothMap\_Circle.toRapidDecay\_fromRapidDecay}
\end{theorem}
\begin{theorem}[\texttt{fourierCoeff\_rapid\_decay}]\label{GaussianField-SmoothMap-Circle-fourierCoeff-rapid-decay}
  \lean{GaussianField.SmoothMap_Circle.fourierCoeff_rapid_decay}\leanok
  \uses{GaussianField-SmoothMap-Circle-fourierCoeffReal-decay, GaussianField-SmoothMap-Circle-instAddCommGroup, GaussianField-SmoothMap-Circle-fourierCoeffReal, GaussianField-SmoothMap-Circle-instSMul, GaussianField-SmoothMap-Circle-instModule, GaussianField-SmoothMap-Circle, GaussianField-SmoothMap-Circle-sobolevSeminorm, Lattice-toSemilatticeInf}
  \texttt{GaussianField.SmoothMap\_Circle.fourierCoeff\_rapid\_decay}
\end{theorem}
\begin{theorem}[\texttt{smoothCircle\_hasBiorthogonalBasis}]\label{GaussianField-smoothCircle-hasBiorthogonalBasis}
  \lean{GaussianField.smoothCircle_hasBiorthogonalBasis}\leanok
  \uses{GaussianField-SmoothMap-Circle-smoothCircle-withSeminorms, GaussianField-SmoothMap-Circle-smoothCircleRapidDecayEquiv, GaussianField-SmoothMap-Circle-instAddCommGroup, GaussianField-DyninMityaginSpace-ofRapidDecayEquiv-hasBiorthogonalBasis, GaussianField-SmoothMap-Circle-instModule, GaussianField-SmoothMap-Circle-instContinuousSMul, GaussianField-smoothCircle-dyninMityaginSpace, GaussianField-SmoothMap-Circle-instIsTopologicalAddGroup, GaussianField-SmoothMap-Circle, GaussianField-SmoothMap-Circle-sobolevSeminorm, GaussianField-DyninMityaginSpace-HasBiorthogonalBasis, GaussianField-SmoothMap-Circle-instTopologicalSpace}
  \texttt{GaussianField.smoothCircle\_hasBiorthogonalBasis}
\end{theorem}
\begin{theorem}[\texttt{fourierCoeffReal\_decay}]\label{GaussianField-SmoothMap-Circle-fourierCoeffReal-decay}
  \lean{GaussianField.SmoothMap_Circle.fourierCoeffReal_decay}\leanok
  \uses{GaussianField-SmoothMap-Circle-instAddCommGroup, GaussianField-SmoothMap-Circle-fourierCoeffReal, GaussianField-SmoothMap-Circle-instSMul, GaussianField-SmoothMap-Circle-instModule, GaussianField-SmoothMap-Circle, GaussianField-SmoothMap-Circle-sobolevSeminorm, GaussianField-SmoothMap-Circle-fourierCoeffReal-bound, Lattice-toSemilatticeInf}
  \texttt{GaussianField.SmoothMap\_Circle.fourierCoeffReal\_decay}
\end{theorem}
\begin{theorem}[\texttt{fromRapidDecay\_continuous}]\label{GaussianField-SmoothMap-Circle-fromRapidDecay-continuous}
  \lean{GaussianField.SmoothMap_Circle.fromRapidDecay_continuous}\leanok
  \uses{GaussianField-SmoothMap-Circle-instFunLike, GaussianField-RapidDecaySeq-rapidDecaySeminorm, GaussianField-SmoothMap-Circle-smoothCircle-withSeminorms, GaussianField-RapidDecaySeq, GaussianField-SmoothMap-Circle-fourierBasis, GaussianField-RapidDecaySeq-val, GaussianField-SmoothMap-Circle-instAddCommGroup, GaussianField-SmoothMap-Circle-fourierBasisFun-smooth, GaussianField-SmoothMap-Circle-fourierBasisFun, GaussianField-RapidDecaySeq-instSMul, GaussianField-RapidDecaySeq-instModule, GaussianField-SmoothMap-Circle-instSMul, GaussianField-SmoothMap-Circle-instModule, GaussianField-RapidDecaySeq-instAddCommGroup, GaussianField-SmoothMap-Circle, GaussianField-RapidDecaySeq-instTopologicalSpace, GaussianField-SmoothMap-Circle-sobolevSeminorm, GaussianField-RapidDecaySeq-rapidDecay-withSeminorms, GaussianField-RapidDecaySeq-rapid-decay, GaussianField-SmoothMap-Circle-fromRapidDecayLM, GaussianField-SmoothMap-Circle-sobolevSeminorm-fourierBasis-le, GaussianField-SmoothMap-Circle-instTopologicalSpace, Lattice-toSemilatticeInf}
  \texttt{GaussianField.SmoothMap\_Circle.fromRapidDecay\_continuous}
\end{theorem}
\begin{definition}[\texttt{smoothCircle\_dyninMityaginSpace}]\label{GaussianField-smoothCircle-dyninMityaginSpace}
  \lean{GaussianField.smoothCircle_dyninMityaginSpace}\leanok
  \uses{GaussianField-DyninMityaginSpace, GaussianField-DyninMityaginSpace-ofRapidDecayEquiv, GaussianField-SmoothMap-Circle-smoothCircle-withSeminorms, GaussianField-SmoothMap-Circle-smoothCircleRapidDecayEquiv, GaussianField-SmoothMap-Circle-instAddCommGroup, GaussianField-SmoothMap-Circle-instModule, GaussianField-SmoothMap-Circle-instContinuousSMul, GaussianField-SmoothMap-Circle-instIsTopologicalAddGroup, GaussianField-SmoothMap-Circle, GaussianField-SmoothMap-Circle-sobolevSeminorm, GaussianField-SmoothMap-Circle-instTopologicalSpace}
  \texttt{GaussianField.smoothCircle\_dyninMityaginSpace}
\end{definition}
\begin{theorem}[\texttt{periodic\_fourierSeries}]\label{GaussianField-SmoothMap-Circle-periodic-fourierSeries}
  \lean{GaussianField.SmoothMap_Circle.periodic_fourierSeries}\leanok
  \uses{GaussianField-RapidDecaySeq, GaussianField-RapidDecaySeq-val, GaussianField-SmoothMap-Circle-fourierBasisFun, GaussianField-SmoothMap-Circle-fourierBasisFun-periodic}
  \texttt{GaussianField.SmoothMap\_Circle.periodic\_fourierSeries}
\end{theorem}
\begin{theorem}[\texttt{toRapidDecay\_continuous}]\label{GaussianField-SmoothMap-Circle-toRapidDecay-continuous}
  \lean{GaussianField.SmoothMap_Circle.toRapidDecay_continuous}\leanok
  \uses{GaussianField-RapidDecaySeq-rapidDecaySeminorm, GaussianField-SmoothMap-Circle-smoothCircle-withSeminorms, GaussianField-SmoothMap-Circle-fourierCoeffReal-decay, GaussianField-RapidDecaySeq, GaussianField-SmoothMap-Circle-instAddCommGroup, GaussianField-SmoothMap-Circle-fourierCoeffReal, GaussianField-RapidDecaySeq-instModule, GaussianField-SmoothMap-Circle-instSMul, GaussianField-SmoothMap-Circle-toRapidDecayLM, GaussianField-SmoothMap-Circle-fourierCoeff-rapid-decay, GaussianField-SmoothMap-Circle-instModule, GaussianField-RapidDecaySeq-instAddCommGroup, GaussianField-SmoothMap-Circle, GaussianField-RapidDecaySeq-instTopologicalSpace, GaussianField-SmoothMap-Circle-sobolevSeminorm, GaussianField-RapidDecaySeq-rapidDecay-withSeminorms, GaussianField-SmoothMap-Circle-instTopologicalSpace, Lattice-toSemilatticeInf}
  \texttt{GaussianField.SmoothMap\_Circle.toRapidDecay\_continuous}
\end{theorem}
\begin{theorem}[\texttt{hasSum\_fourierBasis}]\label{GaussianField-SmoothMap-Circle-hasSum-fourierBasis}
  \lean{GaussianField.SmoothMap_Circle.hasSum_fourierBasis}\leanok
  \uses{GaussianField-SmoothMap-Circle-toRapidDecay, GaussianField-SmoothMap-Circle-toRapidDecay-fromRapidDecay, GaussianField-SmoothMap-Circle-instZero, GaussianField-RapidDecaySeq, GaussianField-RapidDecaySeq-val, GaussianField-SmoothMap-Circle-instAddCommGroup, GaussianField-SmoothMap-Circle-fromRapidDecay, GaussianField-SmoothMap-Circle-fourierCoeffReal, GaussianField-SmoothMap-Circle-fourierCoeffLM, GaussianField-SmoothMap-Circle-instSub, GaussianField-SmoothMap-Circle-instModule, GaussianField-SmoothMap-Circle}
  \texttt{GaussianField.SmoothMap\_Circle.hasSum\_fourierBasis}
\end{theorem}
\begin{definition}[\texttt{fromRapidDecay}]\label{GaussianField-SmoothMap-Circle-fromRapidDecay}
  \lean{GaussianField.SmoothMap_Circle.fromRapidDecay}\leanok
  \uses{GaussianField-RapidDecaySeq, GaussianField-RapidDecaySeq-val, GaussianField-SmoothMap-Circle-fourierBasisFun, GaussianField-SmoothMap-Circle-periodic-fourierSeries, GaussianField-SmoothMap-Circle, GaussianField-SmoothMap-Circle-contDiff-fourierSeries}
  \texttt{GaussianField.SmoothMap\_Circle.fromRapidDecay}
\end{definition}
\begin{definition}[\texttt{toRapidDecayCLM}]\label{GaussianField-SmoothMap-Circle-toRapidDecayCLM}
  \lean{GaussianField.SmoothMap_Circle.toRapidDecayCLM}\leanok
  \uses{GaussianField-RapidDecaySeq, GaussianField-SmoothMap-Circle-instAddCommGroup, GaussianField-RapidDecaySeq-instModule, GaussianField-SmoothMap-Circle-toRapidDecayLM, GaussianField-SmoothMap-Circle-instModule, GaussianField-RapidDecaySeq-instAddCommGroup, GaussianField-SmoothMap-Circle, GaussianField-RapidDecaySeq-instTopologicalSpace, GaussianField-SmoothMap-Circle-instTopologicalSpace, GaussianField-SmoothMap-Circle-toRapidDecay-continuous}
  \texttt{GaussianField.SmoothMap\_Circle.toRapidDecayCLM}
\end{definition}
\begin{definition}[\texttt{smoothCircleRapidDecayEquiv}]\label{GaussianField-SmoothMap-Circle-smoothCircleRapidDecayEquiv}
  \lean{GaussianField.SmoothMap_Circle.smoothCircleRapidDecayEquiv}\leanok
  \uses{GaussianField-SmoothMap-Circle-toRapidDecay-fromRapidDecay, GaussianField-RapidDecaySeq, GaussianField-SmoothMap-Circle-instAddCommGroup, GaussianField-SmoothMap-Circle-fromRapidDecay, GaussianField-RapidDecaySeq-instModule, GaussianField-SmoothMap-Circle-fromRapidDecay-continuous, GaussianField-SmoothMap-Circle-toRapidDecayLM, GaussianField-SmoothMap-Circle-instModule, GaussianField-RapidDecaySeq-instAddCommGroup, GaussianField-SmoothMap-Circle, GaussianField-RapidDecaySeq-instTopologicalSpace, GaussianField-SmoothMap-Circle-fromRapidDecay-toRapidDecay, GaussianField-SmoothMap-Circle-instTopologicalSpace, GaussianField-SmoothMap-Circle-toRapidDecay-continuous}
  \texttt{GaussianField.SmoothMap\_Circle.smoothCircleRapidDecayEquiv}
\end{definition}
\begin{definition}[\texttt{fromRapidDecayCLM}]\label{GaussianField-SmoothMap-Circle-fromRapidDecayCLM}
  \lean{GaussianField.SmoothMap_Circle.fromRapidDecayCLM}\leanok
  \uses{GaussianField-RapidDecaySeq, GaussianField-SmoothMap-Circle-instAddCommGroup, GaussianField-RapidDecaySeq-instModule, GaussianField-SmoothMap-Circle-fromRapidDecay-continuous, GaussianField-SmoothMap-Circle-instModule, GaussianField-RapidDecaySeq-instAddCommGroup, GaussianField-SmoothMap-Circle, GaussianField-RapidDecaySeq-instTopologicalSpace, GaussianField-SmoothMap-Circle-fromRapidDecayLM, GaussianField-SmoothMap-Circle-instTopologicalSpace}
  \texttt{GaussianField.SmoothMap\_Circle.fromRapidDecayCLM}
\end{definition}
\begin{definition}[\texttt{fromRapidDecayLM}]\label{GaussianField-SmoothMap-Circle-fromRapidDecayLM}
  \lean{GaussianField.SmoothMap_Circle.fromRapidDecayLM}\leanok
  \uses{GaussianField-RapidDecaySeq, GaussianField-SmoothMap-Circle-instAddCommGroup, GaussianField-SmoothMap-Circle-fromRapidDecay, GaussianField-RapidDecaySeq-instModule, GaussianField-SmoothMap-Circle-instModule, GaussianField-RapidDecaySeq-instAddCommGroup, GaussianField-SmoothMap-Circle}
  \texttt{GaussianField.SmoothMap\_Circle.fromRapidDecayLM}
\end{definition}
\begin{theorem}[\texttt{summable\_fourierBasis\_smul}]\label{GaussianField-SmoothMap-Circle-summable-fourierBasis-smul}
  \lean{GaussianField.SmoothMap_Circle.summable_fourierBasis_smul}\leanok
  \uses{GaussianField-RapidDecaySeq, GaussianField-SmoothMap-Circle-fourierBasisFun-abs-le, GaussianField-RapidDecaySeq-val, GaussianField-SmoothMap-Circle-fourierBasisFun, GaussianField-RapidDecaySeq-rapid-decay, Lattice-toSemilatticeInf}
  \texttt{GaussianField.SmoothMap\_Circle.summable\_fourierBasis\_smul}
\end{theorem}
\begin{definition}[\texttt{toRapidDecay}]\label{GaussianField-SmoothMap-Circle-toRapidDecay}
  \lean{GaussianField.SmoothMap_Circle.toRapidDecay}\leanok
  \uses{GaussianField-RapidDecaySeq, GaussianField-SmoothMap-Circle-fourierCoeffReal, GaussianField-SmoothMap-Circle-fourierCoeff-rapid-decay, GaussianField-SmoothMap-Circle}
  \texttt{GaussianField.SmoothMap\_Circle.toRapidDecay}
\end{definition}
\section{\texttt{Torus.AsymmetricTorus}}
\begin{definition}[\texttt{AsymLatticeField}]\label{GaussianField-AsymLatticeField}
  \lean{GaussianField.AsymLatticeField}\leanok
  \uses{GaussianField-AsymLatticeSites}
  \texttt{GaussianField.AsymLatticeField}
\end{definition}
\begin{theorem}[\texttt{congr\_simp}]\label{GaussianField-evalAsymTorusAtSite-congr-simp}
  \lean{GaussianField.evalAsymTorusAtSite.congr_simp}\leanok
  \uses{GaussianField-NuclearTensorProduct-instModuleReal, GaussianField-NuclearTensorProduct-instTopologicalSpace, GaussianField-AsymLatticeSites, GaussianField-evalAsymTorusAtSite, GaussianField-NuclearTensorProduct-instAddCommGroup, GaussianField-SmoothMap-Circle, GaussianField-AsymTorusTestFunction}
  \texttt{GaussianField.evalAsymTorusAtSite.congr\_simp}
\end{theorem}
\begin{theorem}[\texttt{asymTorusEmbedCLM\_apply}]\label{GaussianField-asymTorusEmbedCLM-apply}
  \lean{GaussianField.asymTorusEmbedCLM_apply}\leanok
  \uses{GaussianField-NuclearTensorProduct-instModuleReal, GaussianField-NuclearTensorProduct-instTopologicalSpace, GaussianField-AsymLatticeSites, GaussianField-NuclearTensorProduct-instIsTopologicalAddGroup, GaussianField-evalAsymTorusAtSite, GaussianField-NuclearTensorProduct-instAddCommGroup, GaussianField-Configuration, GaussianField-NuclearTensorProduct-instContinuousSMulReal, GaussianField-AsymLatticeField, GaussianField-SmoothMap-Circle, GaussianField-instFintypeAsymLatticeSitesOfNeZeroNat, GaussianField-AsymTorusTestFunction, GaussianField-asymTorusEmbedCLM}
  \texttt{GaussianField.asymTorusEmbedCLM\_apply}
\end{theorem}
\begin{definition}[\texttt{asymTorusSwap}]\label{GaussianField-asymTorusSwap}
  \lean{GaussianField.asymTorusSwap}\leanok
  \uses{GaussianField-NuclearTensorProduct-instModuleReal, GaussianField-NuclearTensorProduct-instTopologicalSpace, GaussianField-SmoothMap-Circle-instAddCommGroup, GaussianField-nuclearTensorProduct-swapCLM, GaussianField-NuclearTensorProduct-instAddCommGroup, GaussianField-SmoothMap-Circle-instModule, GaussianField-SmoothMap-Circle-instContinuousSMul, GaussianField-smoothCircle-dyninMityaginSpace, GaussianField-SmoothMap-Circle-instIsTopologicalAddGroup, GaussianField-SmoothMap-Circle, GaussianField-AsymTorusTestFunction, GaussianField-SmoothMap-Circle-instTopologicalSpace}
  \texttt{GaussianField.asymTorusSwap}
\end{definition}
\begin{theorem}[\texttt{asymTorus\_eq\_symmetric}]\label{GaussianField-asymTorus-eq-symmetric}
  \lean{GaussianField.asymTorus_eq_symmetric}\leanok
  \uses{GaussianField-AsymTorusTestFunction, GaussianField-TorusTestFunction}
  \texttt{GaussianField.asymTorus\_eq\_symmetric}
\end{theorem}
\begin{definition}[\texttt{asymTorusEmbedCLM}]\label{GaussianField-asymTorusEmbedCLM}
  \lean{GaussianField.asymTorusEmbedCLM}\leanok
  \uses{GaussianField-NuclearTensorProduct-instModuleReal, GaussianField-NuclearTensorProduct-instTopologicalSpace, GaussianField-AsymLatticeSites, GaussianField-NuclearTensorProduct-instIsTopologicalAddGroup, GaussianField-evalAsymTorusAtSite, GaussianField-NuclearTensorProduct-instAddCommGroup, GaussianField-Configuration, GaussianField-NuclearTensorProduct-instContinuousSMulReal, GaussianField-AsymLatticeField, GaussianField-SmoothMap-Circle, GaussianField-instFintypeAsymLatticeSitesOfNeZeroNat, GaussianField-AsymTorusTestFunction}
  \texttt{GaussianField.asymTorusEmbedCLM}
\end{definition}
\begin{definition}[\texttt{asymTorusTranslation}]\label{GaussianField-asymTorusTranslation}
  \lean{GaussianField.asymTorusTranslation}\leanok
  \uses{GaussianField-NuclearTensorProduct-instModuleReal, GaussianField-circleTranslation, GaussianField-NuclearTensorProduct-instTopologicalSpace, GaussianField-nuclearTensorProduct-mapCLM, GaussianField-SmoothMap-Circle-instAddCommGroup, GaussianField-NuclearTensorProduct-instAddCommGroup, GaussianField-SmoothMap-Circle-instModule, GaussianField-SmoothMap-Circle-instContinuousSMul, GaussianField-smoothCircle-dyninMityaginSpace, GaussianField-SmoothMap-Circle-instIsTopologicalAddGroup, GaussianField-SmoothMap-Circle, GaussianField-AsymTorusTestFunction, GaussianField-SmoothMap-Circle-instTopologicalSpace}
  \texttt{GaussianField.asymTorusTranslation}
\end{definition}
\begin{definition}[\texttt{instFintypeAsymLatticeSitesOfNeZeroNat}]\label{GaussianField-instFintypeAsymLatticeSitesOfNeZeroNat}
  \lean{GaussianField.instFintypeAsymLatticeSitesOfNeZeroNat}\leanok
  \uses{GaussianField-AsymLatticeSites}
  \texttt{GaussianField.instFintypeAsymLatticeSitesOfNeZeroNat}
\end{definition}
\begin{definition}[\texttt{AsymTorusTestFunction}]\label{GaussianField-AsymTorusTestFunction}
  \lean{GaussianField.AsymTorusTestFunction}\leanok
  \uses{GaussianField-NuclearTensorProduct, GaussianField-SmoothMap-Circle}
  \texttt{GaussianField.AsymTorusTestFunction}
\end{definition}
\begin{definition}[\texttt{AsymLatticeSites}]\label{GaussianField-AsymLatticeSites}
  \lean{GaussianField.AsymLatticeSites}\leanok
  \texttt{GaussianField.AsymLatticeSites}
\end{definition}
\begin{definition}[\texttt{asymTorusTimeReflection}]\label{GaussianField-asymTorusTimeReflection}
  \lean{GaussianField.asymTorusTimeReflection}\leanok
  \uses{GaussianField-NuclearTensorProduct-instModuleReal, GaussianField-NuclearTensorProduct-instTopologicalSpace, GaussianField-nuclearTensorProduct-mapCLM, GaussianField-SmoothMap-Circle-instAddCommGroup, GaussianField-NuclearTensorProduct-instAddCommGroup, GaussianField-circleReflection, GaussianField-SmoothMap-Circle-instModule, GaussianField-SmoothMap-Circle-instContinuousSMul, GaussianField-smoothCircle-dyninMityaginSpace, GaussianField-SmoothMap-Circle-instIsTopologicalAddGroup, GaussianField-SmoothMap-Circle, GaussianField-AsymTorusTestFunction, GaussianField-SmoothMap-Circle-instTopologicalSpace}
  \texttt{GaussianField.asymTorusTimeReflection}
\end{definition}
\begin{definition}[\texttt{evalAsymTorusAtSite}]\label{GaussianField-evalAsymTorusAtSite}
  \lean{GaussianField.evalAsymTorusAtSite}\leanok
  \uses{GaussianField-NuclearTensorProduct-instModuleReal, GaussianField-NuclearTensorProduct-instTopologicalSpace, GaussianField-AsymLatticeSites, GaussianField-SmoothMap-Circle-instAddCommGroup, GaussianField-circleRestriction, GaussianField-NuclearTensorProduct-instAddCommGroup, GaussianField-SmoothMap-Circle-instModule, GaussianField-SmoothMap-Circle-instContinuousSMul, GaussianField-smoothCircle-dyninMityaginSpace, GaussianField-SmoothMap-Circle-instIsTopologicalAddGroup, GaussianField-SmoothMap-Circle, GaussianField-AsymTorusTestFunction, GaussianField-NuclearTensorProduct-evalCLM, GaussianField-SmoothMap-Circle-instTopologicalSpace}
  \texttt{GaussianField.evalAsymTorusAtSite}
\end{definition}
\section{\texttt{Torus.Evaluation}}
\begin{theorem}[\texttt{torusEmbedCLM\_apply}]\label{GaussianField-torusEmbedCLM-apply}
  \lean{GaussianField.torusEmbedCLM_apply}\leanok
  \uses{GaussianField-torusEmbedCLM, GaussianField-NuclearTensorProduct-instModuleReal, GaussianField-NuclearTensorProduct-instTopologicalSpace, GaussianField-FinLatticeSites, GaussianField-NuclearTensorProduct-instIsTopologicalAddGroup, GaussianField-NuclearTensorProduct-instAddCommGroup, GaussianField-evalTorusAtSite, GaussianField-Configuration, GaussianField-NuclearTensorProduct-instContinuousSMulReal, GaussianField-SmoothMap-Circle, GaussianField-FinLatticeField, GaussianField-TorusTestFunction}
  \texttt{GaussianField.torusEmbedCLM\_apply}
\end{theorem}
\begin{theorem}[\texttt{torusEmbedCLMGJ\_apply}]\label{GaussianField-torusEmbedCLMGJ-apply}
  \lean{GaussianField.torusEmbedCLMGJ_apply}\leanok
  \uses{GaussianField-NuclearTensorProduct-instModuleReal, GaussianField-NuclearTensorProduct-instTopologicalSpace, GaussianField-FinLatticeSites, GaussianField-NuclearTensorProduct-instIsTopologicalAddGroup, GaussianField-NuclearTensorProduct-instAddCommGroup, GaussianField-Configuration, GaussianField-NuclearTensorProduct-instContinuousSMulReal, GaussianField-SmoothMap-Circle, GaussianField-evalTorusAtSiteGJ, GaussianField-torusEmbedCLMGJ, GaussianField-FinLatticeField, GaussianField-TorusTestFunction}
  \texttt{GaussianField.torusEmbedCLMGJ\_apply}
\end{theorem}
\begin{theorem}[\texttt{evalTorusAtSite\_latticeTranslation}]\label{GaussianField-evalTorusAtSite-latticeTranslation}
  \lean{GaussianField.evalTorusAtSite_latticeTranslation}\leanok
  \uses{GaussianField-NuclearTensorProduct-instModuleReal, GaussianField-circleTranslation, GaussianField-NuclearTensorProduct-instTopologicalSpace, GaussianField-SmoothMap-Circle-instFunLike, GaussianField-nuclearTensorProduct-mapCLM, GaussianField-SmoothMap-Circle-instAddCommGroup, GaussianField-FinLatticeSites, GaussianField-circleRestriction, GaussianField-NuclearTensorProduct, GaussianField-smoothCircle-coeff-basis, GaussianField-translateSites, GaussianField-NuclearTensorProduct-instAddCommGroup, GaussianField-evalTorusAtSite, GaussianField-SmoothMap-Circle-instModule, GaussianField-SmoothMap-Circle-instContinuousSMul, GaussianField-smoothCircle-dyninMityaginSpace, GaussianField-SmoothMap-Circle-instIsTopologicalAddGroup, GaussianField-SmoothMap-Circle, GaussianField-circleSpacing, GaussianField-TorusTestFunction, GaussianField-NuclearTensorProduct-evalCLM, GaussianField-SmoothMap-Circle-instTopologicalSpace, GaussianField-circlePoint, GaussianField-evalCLM-comp-mapCLM}
  \texttt{GaussianField.evalTorusAtSite\_latticeTranslation}
\end{theorem}
\begin{theorem}[\texttt{congr\_simp}]\label{GaussianField-evalTorusAtSiteGJ-congr-simp}
  \lean{GaussianField.evalTorusAtSiteGJ.congr_simp}\leanok
  \uses{GaussianField-NuclearTensorProduct-instModuleReal, GaussianField-NuclearTensorProduct-instTopologicalSpace, GaussianField-FinLatticeSites, GaussianField-NuclearTensorProduct-instAddCommGroup, GaussianField-SmoothMap-Circle, GaussianField-evalTorusAtSiteGJ, GaussianField-TorusTestFunction}
  \texttt{GaussianField.evalTorusAtSiteGJ.congr\_simp}
\end{theorem}
\begin{theorem}[\texttt{evalTorusAtSite\_timeReflection}]\label{GaussianField-evalTorusAtSite-timeReflection}
  \lean{GaussianField.evalTorusAtSite_timeReflection}\leanok
  \uses{GaussianField-NuclearTensorProduct-instModuleReal, GaussianField-NuclearTensorProduct-instTopologicalSpace, GaussianField-SmoothMap-Circle-instFunLike, GaussianField-nuclearTensorProduct-mapCLM, GaussianField-SmoothMap-Circle-instAddCommGroup, GaussianField-FinLatticeSites, GaussianField-circleRestriction, GaussianField-NuclearTensorProduct, GaussianField-smoothCircle-coeff-basis, GaussianField-NuclearTensorProduct-instAddCommGroup, GaussianField-circleReflection, GaussianField-evalTorusAtSite, GaussianField-SmoothMap-Circle-instModule, GaussianField-SmoothMap-Circle-instContinuousSMul, GaussianField-smoothCircle-dyninMityaginSpace, GaussianField-SmoothMap-Circle-instIsTopologicalAddGroup, GaussianField-SmoothMap-Circle, GaussianField-circleSpacing, GaussianField-TorusTestFunction, GaussianField-NuclearTensorProduct-evalCLM, GaussianField-SmoothMap-Circle-instTopologicalSpace, GaussianField-circlePoint, GaussianField-timeReflectSites, GaussianField-evalCLM-comp-mapCLM}
  \texttt{GaussianField.evalTorusAtSite\_timeReflection}
\end{theorem}
\begin{definition}[\texttt{evalTorusAtSiteGJ}]\label{GaussianField-evalTorusAtSiteGJ}
  \lean{GaussianField.evalTorusAtSiteGJ}\leanok
  \uses{GaussianField-NuclearTensorProduct-instModuleReal, GaussianField-NuclearTensorProduct-instTopologicalSpace, GaussianField-FinLatticeSites, GaussianField-NuclearTensorProduct-instAddCommGroup, GaussianField-evalTorusAtSite, GaussianField-SmoothMap-Circle, GaussianField-circleSpacing, GaussianField-TorusTestFunction}
  \texttt{GaussianField.evalTorusAtSiteGJ}
\end{definition}
\begin{definition}[\texttt{evalTorusAtSite}]\label{GaussianField-evalTorusAtSite}
  \lean{GaussianField.evalTorusAtSite}\leanok
  \uses{GaussianField-NuclearTensorProduct-instModuleReal, GaussianField-NuclearTensorProduct-instTopologicalSpace, GaussianField-SmoothMap-Circle-instAddCommGroup, GaussianField-FinLatticeSites, GaussianField-circleRestriction, GaussianField-NuclearTensorProduct-instAddCommGroup, GaussianField-SmoothMap-Circle-instModule, GaussianField-SmoothMap-Circle-instContinuousSMul, GaussianField-smoothCircle-dyninMityaginSpace, GaussianField-SmoothMap-Circle-instIsTopologicalAddGroup, GaussianField-SmoothMap-Circle, GaussianField-TorusTestFunction, GaussianField-NuclearTensorProduct-evalCLM, GaussianField-SmoothMap-Circle-instTopologicalSpace}
  \texttt{GaussianField.evalTorusAtSite}
\end{definition}
\begin{definition}[\texttt{translateSites}]\label{GaussianField-translateSites}
  \lean{GaussianField.translateSites}\leanok
  \uses{GaussianField-FinLatticeSites}
  \texttt{GaussianField.translateSites}
\end{definition}
\begin{theorem}[\texttt{evalTorusAtSiteGJ\_apply'}]\label{GaussianField-evalTorusAtSiteGJ-apply-}
  \lean{GaussianField.evalTorusAtSiteGJ_apply'}\leanok
  \uses{GaussianField-NuclearTensorProduct-instModuleReal, GaussianField-NuclearTensorProduct-instTopologicalSpace, GaussianField-FinLatticeSites, GaussianField-NuclearTensorProduct-instAddCommGroup, GaussianField-evalTorusAtSite, GaussianField-SmoothMap-Circle, GaussianField-circleSpacing, GaussianField-evalTorusAtSiteGJ, GaussianField-TorusTestFunction}
  \texttt{GaussianField.evalTorusAtSiteGJ\_apply'}
\end{theorem}
\begin{theorem}[\texttt{congr\_simp}]\label{GaussianField-evalTorusAtSite-congr-simp}
  \lean{GaussianField.evalTorusAtSite.congr_simp}\leanok
  \uses{GaussianField-NuclearTensorProduct-instModuleReal, GaussianField-NuclearTensorProduct-instTopologicalSpace, GaussianField-FinLatticeSites, GaussianField-NuclearTensorProduct-instAddCommGroup, GaussianField-evalTorusAtSite, GaussianField-SmoothMap-Circle, GaussianField-TorusTestFunction}
  \texttt{GaussianField.evalTorusAtSite.congr\_simp}
\end{theorem}
\begin{definition}[\texttt{timeReflectSites}]\label{GaussianField-timeReflectSites}
  \lean{GaussianField.timeReflectSites}\leanok
  \uses{GaussianField-FinLatticeSites}
  \texttt{GaussianField.timeReflectSites}
\end{definition}
\begin{definition}[\texttt{torusEmbedCLMGJ}]\label{GaussianField-torusEmbedCLMGJ}
  \lean{GaussianField.torusEmbedCLMGJ}\leanok
  \uses{GaussianField-NuclearTensorProduct-instModuleReal, GaussianField-NuclearTensorProduct-instTopologicalSpace, GaussianField-FinLatticeSites, GaussianField-NuclearTensorProduct-instIsTopologicalAddGroup, GaussianField-NuclearTensorProduct-instAddCommGroup, GaussianField-Configuration, GaussianField-NuclearTensorProduct-instContinuousSMulReal, GaussianField-SmoothMap-Circle, GaussianField-evalTorusAtSiteGJ, GaussianField-FinLatticeField, GaussianField-TorusTestFunction}
  \texttt{GaussianField.torusEmbedCLMGJ}
\end{definition}
\begin{definition}[\texttt{torusEmbedCLM}]\label{GaussianField-torusEmbedCLM}
  \lean{GaussianField.torusEmbedCLM}\leanok
  \uses{GaussianField-NuclearTensorProduct-instModuleReal, GaussianField-NuclearTensorProduct-instTopologicalSpace, GaussianField-FinLatticeSites, GaussianField-NuclearTensorProduct-instIsTopologicalAddGroup, GaussianField-NuclearTensorProduct-instAddCommGroup, GaussianField-evalTorusAtSite, GaussianField-Configuration, GaussianField-NuclearTensorProduct-instContinuousSMulReal, GaussianField-SmoothMap-Circle, GaussianField-FinLatticeField, GaussianField-TorusTestFunction}
  \texttt{GaussianField.torusEmbedCLM}
\end{definition}
\begin{definition}[\texttt{swapSites}]\label{GaussianField-swapSites}
  \lean{GaussianField.swapSites}\leanok
  \uses{GaussianField-FinLatticeSites}
  \texttt{GaussianField.swapSites}
\end{definition}
\begin{theorem}[\texttt{evalTorusAtSite\_swap}]\label{GaussianField-evalTorusAtSite-swap}
  \lean{GaussianField.evalTorusAtSite_swap}\leanok
  \uses{GaussianField-NuclearTensorProduct-instModuleReal, GaussianField-NuclearTensorProduct-instTopologicalSpace, GaussianField-SmoothMap-Circle-instAddCommGroup, GaussianField-FinLatticeSites, GaussianField-circleRestriction, GaussianField-nuclearTensorProduct-swapCLM, GaussianField-NuclearTensorProduct, GaussianField-NuclearTensorProduct-instAddCommGroup, GaussianField-evalTorusAtSite, GaussianField-SmoothMap-Circle-instModule, GaussianField-evalCLM-comp-swapCLM, GaussianField-SmoothMap-Circle-instContinuousSMul, GaussianField-smoothCircle-dyninMityaginSpace, GaussianField-SmoothMap-Circle-instIsTopologicalAddGroup, GaussianField-SmoothMap-Circle, GaussianField-swapSites, GaussianField-TorusTestFunction, GaussianField-NuclearTensorProduct-evalCLM, GaussianField-SmoothMap-Circle-instTopologicalSpace}
  \texttt{GaussianField.evalTorusAtSite\_swap}
\end{theorem}
\section{\texttt{Cylinder.PositiveTime}}
\begin{theorem}[\texttt{cylinderPositiveTime\_disjoint\_reflected}]\label{GaussianField-cylinderPositiveTime-disjoint-reflected}
  \lean{GaussianField.cylinderPositiveTime_disjoint_reflected}\leanok
  \uses{GaussianField-NuclearTensorProduct-instModuleReal, GaussianField-NuclearTensorProduct-instTopologicalSpace, GaussianField-cylinderPositiveTimeSubmodule, GaussianField-cylinderPositiveTime-negativeTime-disjoint, GaussianField-cylinderTimeReflection-pos-to-neg, GaussianField-CylinderTestFunction, GaussianField-CylinderPositiveTimeTestFunction, GaussianField-cylinderNegativeTimeSubmodule, GaussianField-NuclearTensorProduct-instAddCommGroup, GaussianField-cylinderTimeReflection, GaussianField-cylinderTimeReflection-involution, GaussianField-SmoothMap-Circle}
  \texttt{GaussianField.cylinderPositiveTime\_disjoint\_reflected}
\end{theorem}
\begin{definition}[\texttt{cylinderPositiveTimeSubmodule}]\label{GaussianField-cylinderPositiveTimeSubmodule}
  \lean{GaussianField.cylinderPositiveTimeSubmodule}\leanok
  \uses{GaussianField-NuclearTensorProduct-instModuleReal, GaussianField-NuclearTensorProduct-instTopologicalSpace, GaussianField-CylinderTestFunction, GaussianField-NuclearTensorProduct-instAddCommGroup, GaussianField-SmoothMap-Circle}
  \texttt{GaussianField.cylinderPositiveTimeSubmodule}
\end{definition}
\begin{definition}[\texttt{cylinderNegativeTimeSubmodule}]\label{GaussianField-cylinderNegativeTimeSubmodule}
  \lean{GaussianField.cylinderNegativeTimeSubmodule}\leanok
  \uses{GaussianField-NuclearTensorProduct-instModuleReal, GaussianField-NuclearTensorProduct-instTopologicalSpace, GaussianField-CylinderTestFunction, GaussianField-NuclearTensorProduct-instAddCommGroup, GaussianField-SmoothMap-Circle}
  \texttt{GaussianField.cylinderNegativeTimeSubmodule}
\end{definition}
\begin{theorem}[\texttt{cylinderTimeReflection\_pos\_to\_neg}]\label{GaussianField-cylinderTimeReflection-pos-to-neg}
  \lean{GaussianField.cylinderTimeReflection_pos_to_neg}\leanok
  \uses{GaussianField-NuclearTensorProduct-instModuleReal, GaussianField-NuclearTensorProduct-instTopologicalSpace, GaussianField-CylinderTestFunction, GaussianField-NuclearTensorProduct-instIsTopologicalAddGroup, GaussianField-CylinderPositiveTimeTestFunction, GaussianField-cylinderNegativeTimeSubmodule, GaussianField-NuclearTensorProduct-instAddCommGroup, GaussianField-cylinderTimeReflection, GaussianField-NuclearTensorProduct-instContinuousSMulReal, GaussianField-SmoothMap-Circle}
  \texttt{GaussianField.cylinderTimeReflection\_pos\_to\_neg}
\end{theorem}
\begin{theorem}[\texttt{cylinderPositiveTime\_spatialTranslation\_closed}]\label{GaussianField-cylinderPositiveTime-spatialTranslation-closed}
  \lean{GaussianField.cylinderPositiveTime_spatialTranslation_closed}\leanok
  \uses{GaussianField-NuclearTensorProduct-instModuleReal, GaussianField-NuclearTensorProduct-instTopologicalSpace, GaussianField-cylinderPositiveTimeSubmodule, GaussianField-cylinderSpatialTranslation, GaussianField-CylinderTestFunction, GaussianField-NuclearTensorProduct-instIsTopologicalAddGroup, GaussianField-CylinderPositiveTimeTestFunction, GaussianField-NuclearTensorProduct-instAddCommGroup, GaussianField-NuclearTensorProduct-instContinuousSMulReal, GaussianField-SmoothMap-Circle}
  \texttt{GaussianField.cylinderPositiveTime\_spatialTranslation\_closed}
\end{theorem}
\begin{theorem}[\texttt{pair\_weight\_le}]\label{GaussianField-pair-weight-le}
  \lean{GaussianField.pair_weight_le}\leanok
  \texttt{GaussianField.pair\_weight\_le}
\end{theorem}
\begin{theorem}[\texttt{ntpSliceSchwartz\_maps\_positive}]\label{GaussianField-ntpSliceSchwartz-maps-positive}
  \lean{GaussianField.ntpSliceSchwartz_maps_positive}\leanok
  \uses{GaussianField-NuclearTensorProduct-instModuleReal, GaussianField-NuclearTensorProduct-instTopologicalSpace, GaussianField-cylinderPositiveTimeSubmodule, GaussianField-DyninMityaginSpace-coeff, GaussianField-schwartz-dyninMityaginSpace-1D, GaussianField-SmoothMap-Circle-instAddCommGroup, GaussianField-CylinderTestFunction, GaussianField-NuclearTensorProduct-instIsTopologicalAddGroup, GaussianField-schwartzPositiveTimeSubmodule, GaussianField-ntpSliceSchwartz-pure, GaussianField-NuclearTensorProduct-instAddCommGroup, GaussianField-schwartzPositiveTimeSubmodule-isClosed, GaussianField-SmoothMap-Circle-instModule, GaussianField-SmoothMap-Circle-instContinuousSMul, GaussianField-smoothCircle-dyninMityaginSpace, GaussianField-SmoothMap-Circle-instIsTopologicalAddGroup, GaussianField-NuclearTensorProduct-instContinuousSMulReal, GaussianField-NuclearTensorProduct-pure, GaussianField-SmoothMap-Circle, GaussianField-ntpSliceSchwartz, GaussianField-SmoothMap-Circle-instTopologicalSpace}
  \texttt{GaussianField.ntpSliceSchwartz\_maps\_positive}
\end{theorem}
\begin{theorem}[\texttt{ntpExtractSlice\_val}]\label{GaussianField-ntpExtractSlice-val}
  \lean{GaussianField.ntpExtractSlice_val}\leanok
  \uses{GaussianField-ntpExtractSlice, GaussianField-RapidDecaySeq, GaussianField-RapidDecaySeq-val, GaussianField-RapidDecaySeq-instModule, GaussianField-RapidDecaySeq-instAddCommGroup, GaussianField-RapidDecaySeq-instTopologicalSpace}
  \texttt{GaussianField.ntpExtractSlice\_val}
\end{theorem}
\begin{definition}[\texttt{ntpExtractSlice}]\label{GaussianField-ntpExtractSlice}
  \lean{GaussianField.ntpExtractSlice}\leanok
  \uses{GaussianField-RapidDecaySeq, GaussianField-RapidDecaySeq-instModule, GaussianField-RapidDecaySeq-instAddCommGroup, GaussianField-RapidDecaySeq-instTopologicalSpace, GaussianField-ntpExtractSliceLM}
  \texttt{GaussianField.ntpExtractSlice}
\end{definition}
\begin{definition}[\texttt{ntpExtractSliceLM}]\label{GaussianField-ntpExtractSliceLM}
  \lean{GaussianField.ntpExtractSliceLM}\leanok
  \uses{GaussianField-ntpExtractSliceFun, GaussianField-RapidDecaySeq, GaussianField-RapidDecaySeq-instModule, GaussianField-RapidDecaySeq-instAddCommGroup}
  \texttt{GaussianField.ntpExtractSliceLM}
\end{definition}
\begin{theorem}[\texttt{pair\_right\_injective}]\label{GaussianField-pair-right-injective}
  \lean{GaussianField.pair_right_injective}\leanok
  \texttt{GaussianField.pair\_right\_injective}
\end{theorem}
\begin{definition}[\texttt{CylinderPositiveTimeTestFunction}]\label{GaussianField-CylinderPositiveTimeTestFunction}
  \lean{GaussianField.CylinderPositiveTimeTestFunction}\leanok
  \uses{GaussianField-NuclearTensorProduct-instModuleReal, GaussianField-cylinderPositiveTimeSubmodule, GaussianField-CylinderTestFunction, GaussianField-NuclearTensorProduct-instAddCommGroup, GaussianField-SmoothMap-Circle}
  \texttt{GaussianField.CylinderPositiveTimeTestFunction}
\end{definition}
\begin{theorem}[\texttt{ntpExtractSliceLM\_isBounded}]\label{GaussianField-ntpExtractSliceLM-isBounded}
  \lean{GaussianField.ntpExtractSliceLM_isBounded}\leanok
  \uses{GaussianField-ntpExtractSliceFun, GaussianField-RapidDecaySeq-rapidDecaySeminorm, GaussianField-RapidDecaySeq, GaussianField-RapidDecaySeq-val, GaussianField-RapidDecaySeq-instModule, GaussianField-RapidDecaySeq-instAddCommGroup, GaussianField-pair-right-injective, GaussianField-RapidDecaySeq-weight-nonneg, GaussianField-RapidDecaySeq-rapid-decay, GaussianField-pair-weight-le, GaussianField-ntpExtractSliceLM, Lattice-toSemilatticeInf}
  \texttt{GaussianField.ntpExtractSliceLM\_isBounded}
\end{theorem}
\begin{definition}[\texttt{ntpSliceSchwartz}]\label{GaussianField-ntpSliceSchwartz}
  \lean{GaussianField.ntpSliceSchwartz}\leanok
  \uses{GaussianField-NuclearTensorProduct-instModuleReal, GaussianField-NuclearTensorProduct-instTopologicalSpace, GaussianField-ntpExtractSlice, GaussianField-RapidDecaySeq, GaussianField-CylinderTestFunction, GaussianField-RapidDecaySeq-instModule, GaussianField-NuclearTensorProduct-instAddCommGroup, GaussianField-RapidDecaySeq-instAddCommGroup, GaussianField-SmoothMap-Circle, GaussianField-RapidDecaySeq-instTopologicalSpace, GaussianField-schwartzRapidDecayEquiv1D}
  \texttt{GaussianField.ntpSliceSchwartz}
\end{definition}
\begin{definition}[\texttt{ntpExtractSliceFun}]\label{GaussianField-ntpExtractSliceFun}
  \lean{GaussianField.ntpExtractSliceFun}\leanok
  \uses{GaussianField-RapidDecaySeq, GaussianField-RapidDecaySeq-val}
  \texttt{GaussianField.ntpExtractSliceFun}
\end{definition}
\begin{theorem}[\texttt{ntpSliceSchwartz\_pure}]\label{GaussianField-ntpSliceSchwartz-pure}
  \lean{GaussianField.ntpSliceSchwartz_pure}\leanok
  \uses{GaussianField-NuclearTensorProduct-instModuleReal, GaussianField-NuclearTensorProduct-instTopologicalSpace, GaussianField-DyninMityaginSpace-coeff, GaussianField-ntpExtractSlice, GaussianField-RapidDecaySeq, GaussianField-schwartz-dyninMityaginSpace-1D, GaussianField-RapidDecaySeq-val, GaussianField-SmoothMap-Circle-instAddCommGroup, GaussianField-CylinderTestFunction, GaussianField-RapidDecaySeq-instSMul, GaussianField-RapidDecaySeq-instModule, GaussianField-NuclearTensorProduct-instAddCommGroup, GaussianField-RapidDecaySeq-ext, GaussianField-SmoothMap-Circle-instModule, GaussianField-SmoothMap-Circle-instContinuousSMul, GaussianField-smoothCircle-dyninMityaginSpace, GaussianField-SmoothMap-Circle-instIsTopologicalAddGroup, GaussianField-NuclearTensorProduct-pure, GaussianField-RapidDecaySeq-instAddCommGroup, GaussianField-SmoothMap-Circle, GaussianField-ntpSliceSchwartz, GaussianField-RapidDecaySeq-instTopologicalSpace, GaussianField-schwartzRapidDecayEquiv1D, GaussianField-SmoothMap-Circle-instTopologicalSpace}
  \texttt{GaussianField.ntpSliceSchwartz\_pure}
\end{theorem}
\begin{theorem}[\texttt{cylinderPositiveTime\_negativeTime\_disjoint}]\label{GaussianField-cylinderPositiveTime-negativeTime-disjoint}
  \lean{GaussianField.cylinderPositiveTime_negativeTime_disjoint}\leanok
  \uses{GaussianField-NuclearTensorProduct-instModuleReal, GaussianField-NuclearTensorProduct-instTopologicalSpace, GaussianField-cylinderPositiveTimeSubmodule, GaussianField-RapidDecaySeq-instZero, GaussianField-ntpExtractSlice, GaussianField-RapidDecaySeq, GaussianField-RapidDecaySeq-val, GaussianField-CylinderTestFunction, GaussianField-RapidDecaySeq-instModule, GaussianField-schwartzPositiveTimeSubmodule, GaussianField-cylinderNegativeTimeSubmodule, GaussianField-RapidDecaySeq-zero-val, GaussianField-NuclearTensorProduct-instAddCommGroup, GaussianField-RapidDecaySeq-ext, GaussianField-RapidDecaySeq-instAddCommGroup, GaussianField-SmoothMap-Circle, GaussianField-ntpSliceSchwartz, GaussianField-RapidDecaySeq-instTopologicalSpace, GaussianField-schwartzNegativeTimeSubmodule, GaussianField-schwartzRapidDecayEquiv1D, GaussianField-ntpSliceSchwartz-maps-positive}
  \texttt{GaussianField.cylinderPositiveTime\_negativeTime\_disjoint}
\end{theorem}
\section{\texttt{SchwartzFourier.ResolventUniformBound}}
\begin{theorem}[\texttt{resolventSymbol\_mul\_quotient}]\label{GaussianField-resolventSymbol-mul-quotient}
  \lean{GaussianField.resolventSymbol_mul_quotient}\leanok
  \uses{GaussianField-resolventQuotientSymbol, GaussianField-resolventSymbol, Lattice-toSemilatticeInf}
  \texttt{GaussianField.resolventSymbol\_mul\_quotient}
\end{theorem}
\begin{theorem}[\texttt{resolventSymbol\_antitone}]\label{GaussianField-resolventSymbol-antitone}
  \lean{GaussianField.resolventSymbol_antitone}\leanok
  \uses{GaussianField-resolventSymbol, Lattice-toSemilatticeInf}
  \texttt{GaussianField.resolventSymbol\_antitone}
\end{theorem}
\begin{theorem}[\texttt{resolventSymbol\_uniform\_deriv\_bound}]\label{GaussianField-resolventSymbol-uniform-deriv-bound}
  \lean{GaussianField.resolventSymbol_uniform_deriv_bound}\leanok
  \uses{GaussianField-resolventSymbol-hasTemperateGrowth, GaussianField-resolventSymbol, Lattice-toSemilatticeInf}
  \texttt{GaussianField.resolventSymbol\_uniform\_deriv\_bound}
\end{theorem}
\begin{theorem}[\texttt{resolventQuotientSymbol\_pos}]\label{GaussianField-resolventQuotientSymbol-pos}
  \lean{GaussianField.resolventQuotientSymbol_pos}\leanok
  \uses{GaussianField-resolventQuotientSymbol, Lattice-toSemilatticeInf}
  \texttt{GaussianField.resolventQuotientSymbol\_pos}
\end{theorem}
\begin{theorem}[\texttt{fourierMultiplier\_schwartz\_bound} (axiom)]\label{GaussianField-fourierMultiplier-schwartz-bound}
  \lean{GaussianField.fourierMultiplier_schwartz_bound}\leanok
  \uses{GaussianField-realFourierMultiplierCLM}
  \texttt{GaussianField.fourierMultiplier\_schwartz\_bound}
\end{theorem}
\begin{theorem}[\texttt{resolventSchwartz\_uniformBound}]\label{GaussianField-resolventSchwartz-uniformBound}
  \lean{GaussianField.resolventSchwartz_uniformBound}\leanok
  \uses{GaussianField-fourierMultiplier-schwartz-bound, GaussianField-resolventSymbol-hasTemperateGrowth, GaussianField-realFourierMultiplierCLM, GaussianField-resolventSymbol-uniform-deriv-bound, GaussianField-resolventMultiplierCLM, GaussianField-resolventSymbol}
  \texttt{GaussianField.resolventSchwartz\_uniformBound}
\end{theorem}
\begin{theorem}[\texttt{resolventQuotientSymbol\_le\_one}]\label{GaussianField-resolventQuotientSymbol-le-one}
  \lean{GaussianField.resolventQuotientSymbol_le_one}\leanok
  \uses{GaussianField-resolventQuotientSymbol, Lattice-toSemilatticeInf}
  \texttt{GaussianField.resolventQuotientSymbol\_le\_one}
\end{theorem}
\begin{definition}[\texttt{resolventQuotientSymbol}]\label{GaussianField-resolventQuotientSymbol}
  \lean{GaussianField.resolventQuotientSymbol}\leanok
  \texttt{GaussianField.resolventQuotientSymbol}
\end{definition}
\begin{theorem}[\texttt{resolventQuotientSymbol\_even}]\label{GaussianField-resolventQuotientSymbol-even}
  \lean{GaussianField.resolventQuotientSymbol_even}\leanok
  \uses{GaussianField-resolventQuotientSymbol}
  \texttt{GaussianField.resolventQuotientSymbol\_even}
\end{theorem}
\section{\texttt{GaussianField.Construction}}
\begin{theorem}[\texttt{configuration\_eval\_measurable}]\label{GaussianField-configuration-eval-measurable}
  \lean{GaussianField.configuration_eval_measurable}\leanok
  \uses{GaussianField-Configuration, GaussianField-instMeasurableSpaceConfiguration}
  \texttt{GaussianField.configuration\_eval\_measurable}
\end{theorem}
\begin{definition}[\texttt{covariance}]\label{GaussianField-covariance}
  \lean{GaussianField.covariance}\leanok
  Lean signature ``\texttt{lean def covariance (T : E →L[ℝ] H) (f g : E) : ℝ }``  Informal: The covariance bilinear form $C(f,g) = \langle T(f), T(g) \rangle_H$.
\end{definition}
\begin{definition}[\texttt{instMeasurableSpaceConfiguration}]\label{GaussianField-instMeasurableSpaceConfiguration}
  \lean{GaussianField.instMeasurableSpaceConfiguration}\leanok
  \uses{GaussianField-Configuration}
  Lean signature ``\texttt{lean instance instMeasurableSpaceConfiguration : MeasurableSpace (Configuration E) }``  Informal: The cylindrical $\sigma$-algebra on $\mathrm{Configuration}(E)$, generated by evaluation maps $\omega \mapsto \omega(f)$.
\end{definition}
\begin{definition}[\texttt{Configuration}]\label{GaussianField-Configuration}
  \lean{GaussianField.Configuration}\leanok
  Lean signature ``\texttt{lean abbrev Configuration (E : Type*) [AddCommGroup E] [Module ℝ E] [TopologicalSpace E] [IsTopologicalAddGroup E] [ContinuousSMul ℝ E] := WeakDual ℝ E }``  Informal: The configuration space $E' = \mathrm{WeakDual}(\mathbb{R}, E)$ of continuous linear functionals on $E$.
\end{definition}
\begin{definition}[\texttt{NoiseSpace}]\label{GaussianField-NoiseSpace}
  \lean{GaussianField.NoiseSpace}\leanok
  Lean signature ``\texttt{lean abbrev NoiseSpace := ℕ → ℝ }``  Informal: The space of noise sequences $\xi : \mathbb{N} \to \mathbb{R}$.
\end{definition}
\begin{theorem}[\texttt{congr\_simp}]\label{GaussianField-DyninMityaginSpace-congr-simp}
  \lean{GaussianField.DyninMityaginSpace.congr_simp}\leanok
  \uses{GaussianField-DyninMityaginSpace}
  \texttt{GaussianField.DyninMityaginSpace.congr\_simp}
\end{theorem}
\begin{theorem}[\texttt{measure\_isProbability}]\label{GaussianField-measure-isProbability}
  \lean{GaussianField.measure_isProbability}\leanok
  \uses{GaussianField-DyninMityaginSpace, GaussianField-measure, GaussianField-ell2-separableSpace, GaussianField-ell2-not-finiteDimensional, GaussianField-Configuration, GaussianField-instMeasurableSpaceConfiguration, GaussianField-NoiseSpace, GaussianField-ell2-}
  \texttt{GaussianField.measure\_isProbability}
\end{theorem}
\begin{definition}[\texttt{measure}]\label{GaussianField-measure}
  \lean{GaussianField.measure}\leanok
  \uses{GaussianField-DyninMityaginSpace, GaussianField-ell2-separableSpace, GaussianField-ell2-not-finiteDimensional, GaussianField-Configuration, GaussianField-instMeasurableSpaceConfiguration, GaussianField-NoiseSpace, GaussianField-ell2-}
  Lean signature ``\texttt{lean def measure (T : E →L[ℝ] H) : @Measure (Configuration E) instMeasurableSpaceConfiguration }``  Informal: The Gaussian probability measure on $E'$ constructed from $T$. For infinite-dimensional $H$, this is the pushforward of the noise measure through the series-limit map. For finite-dimensional $H$, an isometric embedding into $\ell^2$ is used first.
\end{definition}
\begin{theorem}[\texttt{charFun}]\label{GaussianField-charFun}
  \lean{GaussianField.charFun}\leanok
  \uses{GaussianField-DyninMityaginSpace, GaussianField-measure, GaussianField-ell2-separableSpace, GaussianField-ell2-not-finiteDimensional, GaussianField-configuration-eval-measurable, GaussianField-Configuration, GaussianField-instMeasurableSpaceConfiguration, GaussianField-NoiseSpace, GaussianField-ell2-}
  \texttt{GaussianField.charFun}
\end{theorem}
\begin{theorem}[\texttt{configuration\_measurable\_of\_eval\_measurable}]\label{GaussianField-configuration-measurable-of-eval-measurable}
  \lean{GaussianField.configuration_measurable_of_eval_measurable}\leanok
  \uses{GaussianField-Configuration, GaussianField-instMeasurableSpaceConfiguration}
  \texttt{GaussianField.configuration\_measurable\_of\_eval\_measurable}
\end{theorem}
\section{\texttt{Cylinder.MassOperatorConstruction}}
\begin{theorem}[\texttt{cylinderMassOperator\_formula}]\label{GaussianField-cylinderMassOperator-formula}
  \lean{GaussianField.cylinderMassOperator_formula}\leanok
  \uses{GaussianField-NuclearTensorProduct-instModuleReal, GaussianField-NuclearTensorProduct-instTopologicalSpace, GaussianField-DyninMityaginSpace-coeff, GaussianField-schwartz-dyninMityaginSpace-1D, GaussianField-CylinderTestFunction, GaussianField-NuclearTensorProduct-instAddCommGroup, GaussianField-massOperatorCoord, GaussianField-cylinderMassOperator-coord, GaussianField-resolventFreq-pos, GaussianField-resolventFreq, GaussianField-cylinderMassOperator, GaussianField-SmoothMap-Circle, GaussianField-ntpSliceSchwartz, GaussianField-resolventMultiplierCLM, GaussianField-ell2-}
  \texttt{GaussianField.cylinderMassOperator\_formula}
\end{theorem}
\begin{theorem}[\texttt{massOperatorCoord\_apply}]\label{GaussianField-massOperatorCoord-apply}
  \lean{GaussianField.massOperatorCoord_apply}\leanok
  \uses{GaussianField-NuclearTensorProduct-instModuleReal, GaussianField-NuclearTensorProduct-instTopologicalSpace, GaussianField-DyninMityaginSpace-coeff, GaussianField-schwartz-dyninMityaginSpace-1D, GaussianField-CylinderTestFunction, GaussianField-NuclearTensorProduct-instAddCommGroup, GaussianField-massOperatorCoord, GaussianField-resolventFreq-pos, GaussianField-resolventFreq, GaussianField-SmoothMap-Circle, GaussianField-ntpSliceSchwartz, GaussianField-resolventMultiplierCLM}
  \texttt{GaussianField.massOperatorCoord\_apply}
\end{theorem}
\begin{theorem}[\texttt{resolventFreq\_zero\_mode}]\label{GaussianField-resolventFreq-zero-mode}
  \lean{GaussianField.resolventFreq_zero_mode}\leanok
  \uses{GaussianField-resolventFreq, GaussianField-SmoothMap-Circle-fourierFreq}
  \texttt{GaussianField.resolventFreq\_zero\_mode}
\end{theorem}
\begin{definition}[\texttt{cylinderMassOperator}]\label{GaussianField-cylinderMassOperator}
  \lean{GaussianField.cylinderMassOperator}\leanok
  \uses{GaussianField-NuclearTensorProduct-instModuleReal, GaussianField-NuclearTensorProduct-instTopologicalSpace, GaussianField-CylinderTestFunction, GaussianField-NuclearTensorProduct-instAddCommGroup, GaussianField-massOperatorCoord, GaussianField-SmoothMap-Circle, GaussianField-ell2-}
  \texttt{GaussianField.cylinderMassOperator}
\end{definition}
\begin{theorem}[\texttt{cylinderMassOperator\_coord}]\label{GaussianField-cylinderMassOperator-coord}
  \lean{GaussianField.cylinderMassOperator_coord}\leanok
  \uses{GaussianField-NuclearTensorProduct-instModuleReal, GaussianField-NuclearTensorProduct-instTopologicalSpace, GaussianField-CylinderTestFunction, GaussianField-NuclearTensorProduct-instAddCommGroup, GaussianField-massOperatorCoord, GaussianField-cylinderMassOperator, GaussianField-SmoothMap-Circle, GaussianField-ell2-}
  \texttt{GaussianField.cylinderMassOperator\_coord}
\end{theorem}
\begin{theorem}[\texttt{ntpExtractSlice\_a\_decay}]\label{GaussianField-ntpExtractSlice-a-decay}
  \lean{GaussianField.ntpExtractSlice_a_decay}\leanok
  \uses{GaussianField-RapidDecaySeq-rapidDecaySeminorm, GaussianField-ntpExtractSlice, GaussianField-RapidDecaySeq, GaussianField-RapidDecaySeq-val, GaussianField-RapidDecaySeq-instSMul, GaussianField-RapidDecaySeq-instModule, GaussianField-RapidDecaySeq-instAddCommGroup, GaussianField-RapidDecaySeq-instTopologicalSpace, GaussianField-NuclearTensorProduct-summable-inv-one-add-sq, GaussianField-RapidDecaySeq-weight-nonneg, GaussianField-RapidDecaySeq-rapid-decay, GaussianField-pair-weight-le, Lattice-toSemilatticeInf}
  \texttt{GaussianField.ntpExtractSlice\_a\_decay}
\end{theorem}
\begin{theorem}[\texttt{massOperatorCoord\_decay}]\label{GaussianField-massOperatorCoord-decay}
  \lean{GaussianField.massOperatorCoord_decay}\leanok
  \uses{GaussianField-NuclearTensorProduct-instModuleReal, GaussianField-NuclearTensorProduct-instTopologicalSpace, GaussianField-RapidDecaySeq-rapidDecaySeminorm, GaussianField-ntpExtractSlice, GaussianField-RapidDecaySeq, GaussianField-RapidDecaySeq-val, GaussianField-NuclearTensorProduct-dyninMityaginSpace, GaussianField-CylinderTestFunction, GaussianField-NuclearTensorProduct-instIsTopologicalAddGroup, GaussianField-RapidDecaySeq-instSMul, GaussianField-RapidDecaySeq-instModule, GaussianField-DyninMityaginSpace--, GaussianField-NuclearTensorProduct-instAddCommGroup, GaussianField-massOperatorCoord, GaussianField-NuclearTensorProduct-instContinuousSMulReal, GaussianField-resolventFreq-pos, GaussianField-resolventFreq, GaussianField-RapidDecaySeq-instAddCommGroup, GaussianField-SmoothMap-Circle, GaussianField-ntpSliceSchwartz, GaussianField-RapidDecaySeq-instTopologicalSpace, GaussianField-resolventMultiplierCLM, GaussianField-NuclearTensorProduct-summable-inv-one-add-sq, GaussianField-RapidDecaySeq-weight-nonneg, GaussianField-DyninMityaginSpace-p, GaussianField-RapidDecaySeq-rapid-decay, GaussianField-resolventRDS-uniformBound, GaussianField-schwartzRapidDecayEquiv1D, Lattice-toSemilatticeInf, GaussianField-ntpExtractSlice-a-decay}
  \texttt{GaussianField.massOperatorCoord\_decay}
\end{theorem}
\begin{definition}[\texttt{massOperatorCoord}]\label{GaussianField-massOperatorCoord}
  \lean{GaussianField.massOperatorCoord}\leanok
  \uses{GaussianField-NuclearTensorProduct-instModuleReal, GaussianField-NuclearTensorProduct-instTopologicalSpace, GaussianField-DyninMityaginSpace-coeff, GaussianField-schwartz-dyninMityaginSpace-1D, GaussianField-CylinderTestFunction, GaussianField-NuclearTensorProduct-instAddCommGroup, GaussianField-resolventFreq-pos, GaussianField-resolventFreq, GaussianField-SmoothMap-Circle, GaussianField-ntpSliceSchwartz, GaussianField-resolventMultiplierCLM}
  \texttt{GaussianField.massOperatorCoord}
\end{definition}
\begin{theorem}[\texttt{resolventRDS\_uniformBound}]\label{GaussianField-resolventRDS-uniformBound}
  \lean{GaussianField.resolventRDS_uniformBound}\leanok
  \uses{GaussianField-RapidDecaySeq-rapidDecaySeminorm, GaussianField-RapidDecaySeq, GaussianField-resolventFreq-mass-le, GaussianField-RapidDecaySeq-instSMul, GaussianField-RapidDecaySeq-instModule, GaussianField-resolventFreq-pos, GaussianField-resolventFreq, GaussianField-RapidDecaySeq-instAddCommGroup, GaussianField-RapidDecaySeq-instTopologicalSpace, GaussianField-resolventMultiplierCLM, GaussianField-resolventSchwartz-uniformBound, GaussianField-RapidDecaySeq-rapidDecay-withSeminorms, GaussianField-schwartzRapidDecayEquiv1D, Lattice-toSemilatticeInf}
  \texttt{GaussianField.resolventRDS\_uniformBound}
\end{theorem}
\begin{theorem}[\texttt{congr\_simp}]\label{GaussianField-resolventMultiplierCLM-congr-simp}
  \lean{GaussianField.resolventMultiplierCLM.congr_simp}\leanok
  \uses{GaussianField-resolventMultiplierCLM}
  \texttt{GaussianField.resolventMultiplierCLM.congr\_simp}
\end{theorem}
\begin{theorem}[\texttt{resolventFreq\_sq}]\label{GaussianField-resolventFreq-sq}
  \lean{GaussianField.resolventFreq_sq}\leanok
  \uses{GaussianField-resolventFreq, GaussianField-SmoothMap-Circle-fourierFreq, Lattice-toSemilatticeInf}
  \texttt{GaussianField.resolventFreq\_sq}
\end{theorem}
\begin{definition}[\texttt{resolventFreq}]\label{GaussianField-resolventFreq}
  \lean{GaussianField.resolventFreq}\leanok
  \uses{GaussianField-SmoothMap-Circle-fourierFreq}
  \texttt{GaussianField.resolventFreq}
\end{definition}
\begin{theorem}[\texttt{resolventFreq\_nonneg}]\label{GaussianField-resolventFreq-nonneg}
  \lean{GaussianField.resolventFreq_nonneg}\leanok
  \uses{GaussianField-resolventFreq, GaussianField-SmoothMap-Circle-fourierFreq}
  \texttt{GaussianField.resolventFreq\_nonneg}
\end{theorem}
\begin{theorem}[\texttt{resolventFreq\_mass\_le}]\label{GaussianField-resolventFreq-mass-le}
  \lean{GaussianField.resolventFreq_mass_le}\leanok
  \uses{GaussianField-resolventFreq, GaussianField-SmoothMap-Circle-fourierFreq, Lattice-toSemilatticeInf}
  \texttt{GaussianField.resolventFreq\_mass\_le}
\end{theorem}
\begin{theorem}[\texttt{resolventFreq\_pos}]\label{GaussianField-resolventFreq-pos}
  \lean{GaussianField.resolventFreq_pos}\leanok
  \uses{GaussianField-resolventFreq, GaussianField-SmoothMap-Circle-fourierFreq, Lattice-toSemilatticeInf}
  \texttt{GaussianField.resolventFreq\_pos}
\end{theorem}
\section{\texttt{GaussianField.TargetFactorization}}
\begin{theorem}[\texttt{nuclear\_ell2\_embedding\_from\_decay}]\label{GaussianField-nuclear-ell2-embedding-from-decay}
  \lean{GaussianField.nuclear_ell2_embedding_from_decay}\leanok
  \uses{GaussianField-DyninMityaginSpace, GaussianField-DyninMityaginSpace--, GaussianField-DyninMityaginSpace-h-with, GaussianField-DyninMityaginSpace-p, GaussianField-ell2-, Lattice-toSemilatticeInf}
  \texttt{GaussianField.nuclear\_ell2\_embedding\_from\_decay}
\end{theorem}
\begin{theorem}[\texttt{summable\_smul\_ell2}]\label{GaussianField-summable-smul-ell2}
  \lean{GaussianField.summable_smul_ell2}\leanok
  \uses{GaussianField-ell2-}
  \texttt{GaussianField.summable\_smul\_ell2}
\end{theorem}
\begin{theorem}[\texttt{nuclear\_clm\_target\_factorization}]\label{GaussianField-nuclear-clm-target-factorization}
  \lean{GaussianField.nuclear_clm_target_factorization}\leanok
  \uses{GaussianField-DyninMityaginSpace, GaussianField-DyninMityaginSpace-expansion-H, GaussianField-DyninMityaginSpace-coeff, GaussianField-DyninMityaginSpace-coeff-decay, GaussianField-DyninMityaginSpace--, GaussianField-summable-one-add-rpow-neg-two, GaussianField-DyninMityaginSpace-basis, GaussianField-nuclear-sequence-svd, GaussianField-DyninMityaginSpace-p, GaussianField-ell2-, GaussianField-clm-image-growth, GaussianField-nuclear-ell2-embedding-from-decay, Lattice-toSemilatticeInf}
  \texttt{GaussianField.nuclear\_clm\_target\_factorization}
\end{theorem}
\begin{theorem}[\texttt{ell2\_inner\_eq\_tsum}]\label{GaussianField-ell2-inner-eq-tsum}
  \lean{GaussianField.ell2_inner_eq_tsum}\leanok
  \uses{GaussianField-ell2-}
  \texttt{GaussianField.ell2\_inner\_eq\_tsum}
\end{theorem}
\section{\texttt{HeatKernel.GreenInvariance}}
\begin{theorem}[\texttt{greenFunctionBilinear\_translation\_pure}]\label{GaussianField-greenFunctionBilinear-translation-pure}
  \lean{GaussianField.greenFunctionBilinear_translation_pure}\leanok
  \uses{GaussianField-NuclearTensorProduct-instModuleReal, GaussianField-circleTranslation, GaussianField-NuclearTensorProduct-instTopologicalSpace, GaussianField-circleHasLaplacianEigenvalues, GaussianField-NuclearTensorProduct-dyninMityaginSpace, GaussianField-SmoothMap-Circle-instAddCommGroup, GaussianField-NuclearTensorProduct-instIsTopologicalAddGroup, GaussianField-tensorProductHasLaplacianEigenvalues, GaussianField-NuclearTensorProduct, GaussianField-NuclearTensorProduct-instAddCommGroup, GaussianField-SmoothMap-Circle-instModule, GaussianField-SmoothMap-Circle-instContinuousSMul, GaussianField-smoothCircle-dyninMityaginSpace, GaussianField-SmoothMap-Circle-instIsTopologicalAddGroup, GaussianField-NuclearTensorProduct-instContinuousSMulReal, GaussianField-NuclearTensorProduct-pure, GaussianField-SmoothMap-Circle, GaussianField-greenFunctionBilinear-swap-pure, GaussianField-SmoothMap-Circle-instTopologicalSpace, GaussianField-greenFunctionBilinear}
  \texttt{GaussianField.greenFunctionBilinear\_translation\_pure}
\end{theorem}
\begin{theorem}[\texttt{coeff\_eq\_fourierCoeffReal}]\label{GaussianField-coeff-eq-fourierCoeffReal}
  \lean{GaussianField.coeff_eq_fourierCoeffReal}\leanok
  \uses{GaussianField-DyninMityaginSpace-coeff, GaussianField-SmoothMap-Circle-instAddCommGroup, GaussianField-SmoothMap-Circle-fourierCoeffReal, GaussianField-SmoothMap-Circle-instModule, GaussianField-SmoothMap-Circle-instContinuousSMul, GaussianField-smoothCircle-dyninMityaginSpace, GaussianField-SmoothMap-Circle-instIsTopologicalAddGroup, GaussianField-SmoothMap-Circle, GaussianField-SmoothMap-Circle-instTopologicalSpace}
  \texttt{GaussianField.coeff\_eq\_fourierCoeffReal}
\end{theorem}
\begin{theorem}[\texttt{modePartner\_cos}]\label{GaussianField-modePartner-cos}
  \lean{GaussianField.modePartner_cos}\leanok
  \uses{GaussianField-modePartner}
  \texttt{GaussianField.modePartner\_cos}
\end{theorem}
\begin{definition}[\texttt{modePartnerEquiv}]\label{GaussianField-modePartnerEquiv}
  \lean{GaussianField.modePartnerEquiv}\leanok
  \uses{GaussianField-modePartner-involutive, GaussianField-modePartner}
  \texttt{GaussianField.modePartnerEquiv}
\end{definition}
\begin{theorem}[\texttt{greenFunctionBilinear\_translation\_invariant}]\label{GaussianField-greenFunctionBilinear-translation-invariant}
  \lean{GaussianField.greenFunctionBilinear_translation_invariant}\leanok
  \uses{GaussianField-NuclearTensorProduct-instModuleReal, GaussianField-circleTranslation, GaussianField-NuclearTensorProduct-instTopologicalSpace, GaussianField-nuclearTensorProduct-mapCLM, GaussianField-circleHasLaplacianEigenvalues, GaussianField-NuclearTensorProduct-dyninMityaginSpace, GaussianField-SmoothMap-Circle-instAddCommGroup, GaussianField-NuclearTensorProduct-instIsTopologicalAddGroup, GaussianField-nuclearTensorProduct-mapCLM-pure, GaussianField-greenFunctionBilinear-invariant-of-pure, GaussianField-tensorProductHasLaplacianEigenvalues, GaussianField-NuclearTensorProduct, GaussianField-smoothCircle-coeff-basis, GaussianField-greenFunctionBilinear-translation-pure, GaussianField-NuclearTensorProduct-instAddCommGroup, GaussianField-SmoothMap-Circle-instModule, GaussianField-SmoothMap-Circle-instContinuousSMul, GaussianField-smoothCircle-dyninMityaginSpace, GaussianField-SmoothMap-Circle-instIsTopologicalAddGroup, GaussianField-NuclearTensorProduct-instContinuousSMulReal, GaussianField-NuclearTensorProduct-pure, GaussianField-SmoothMap-Circle, GaussianField-SmoothMap-Circle-instTopologicalSpace, GaussianField-greenFunctionBilinear}
  \texttt{GaussianField.greenFunctionBilinear\_translation\_invariant}
\end{theorem}
\begin{theorem}[\texttt{modePartner\_sin}]\label{GaussianField-modePartner-sin}
  \lean{GaussianField.modePartner_sin}\leanok
  \uses{GaussianField-modePartner}
  \texttt{GaussianField.modePartner\_sin}
\end{theorem}
\begin{theorem}[\texttt{tsum\_eq\_of\_paired\_involution}]\label{GaussianField-tsum-eq-of-paired-involution}
  \lean{GaussianField.tsum_eq_of_paired_involution}\leanok
  \uses{Lattice-toSemilatticeInf}
  \texttt{GaussianField.tsum\_eq\_of\_paired\_involution}
\end{theorem}
\begin{theorem}[\texttt{smoothCircle\_coeff\_basis}]\label{GaussianField-smoothCircle-coeff-basis}
  \lean{GaussianField.smoothCircle_coeff_basis}\leanok
  \uses{GaussianField-DyninMityaginSpace-coeff, GaussianField-RapidDecaySeq, GaussianField-RapidDecaySeq-val, GaussianField-SmoothMap-Circle-smoothCircleRapidDecayEquiv, GaussianField-SmoothMap-Circle-instAddCommGroup, GaussianField-RapidDecaySeq-instModule, GaussianField-RapidDecaySeq-basisVec, GaussianField-RapidDecaySeq-coeffLM, GaussianField-SmoothMap-Circle-instModule, GaussianField-SmoothMap-Circle-instContinuousSMul, GaussianField-smoothCircle-dyninMityaginSpace, GaussianField-SmoothMap-Circle-instIsTopologicalAddGroup, GaussianField-DyninMityaginSpace-basis, GaussianField-RapidDecaySeq-instAddCommGroup, GaussianField-SmoothMap-Circle, GaussianField-RapidDecaySeq-instTopologicalSpace, GaussianField-RapidDecaySeq-coeffCLM, GaussianField-SmoothMap-Circle-instTopologicalSpace}
  \texttt{GaussianField.smoothCircle\_coeff\_basis}
\end{theorem}
\begin{theorem}[\texttt{greenFunctionBilinear\_swap\_invariant}]\label{GaussianField-greenFunctionBilinear-swap-invariant}
  \lean{GaussianField.greenFunctionBilinear_swap_invariant}\leanok
  \uses{GaussianField-NuclearTensorProduct-instModuleReal, GaussianField-NuclearTensorProduct-instTopologicalSpace, GaussianField-circleHasLaplacianEigenvalues, GaussianField-NuclearTensorProduct-dyninMityaginSpace, GaussianField-SmoothMap-Circle-instAddCommGroup, GaussianField-NuclearTensorProduct-instIsTopologicalAddGroup, GaussianField-greenFunctionBilinear-invariant-of-pure, GaussianField-tensorProductHasLaplacianEigenvalues, GaussianField-nuclearTensorProduct-swapCLM, GaussianField-NuclearTensorProduct, GaussianField-smoothCircle-coeff-basis, GaussianField-nuclearTensorProduct-swapCLM-pure, GaussianField-NuclearTensorProduct-instAddCommGroup, GaussianField-SmoothMap-Circle-instModule, GaussianField-SmoothMap-Circle-instContinuousSMul, GaussianField-smoothCircle-dyninMityaginSpace, GaussianField-SmoothMap-Circle-instIsTopologicalAddGroup, GaussianField-NuclearTensorProduct-instContinuousSMulReal, GaussianField-NuclearTensorProduct-pure, GaussianField-SmoothMap-Circle, GaussianField-greenFunctionBilinear-swap-pure, GaussianField-SmoothMap-Circle-instTopologicalSpace, GaussianField-greenFunctionBilinear}
  \texttt{GaussianField.greenFunctionBilinear\_swap\_invariant}
\end{theorem}
\begin{definition}[\texttt{modePartner}]\label{GaussianField-modePartner}
  \lean{GaussianField.modePartner}\leanok
  \texttt{GaussianField.modePartner}
\end{definition}
\begin{theorem}[\texttt{greenFunctionBilinear\_timeReflection\_invariant}]\label{GaussianField-greenFunctionBilinear-timeReflection-invariant}
  \lean{GaussianField.greenFunctionBilinear_timeReflection_invariant}\leanok
  \uses{GaussianField-NuclearTensorProduct-instModuleReal, GaussianField-NuclearTensorProduct-instTopologicalSpace, GaussianField-nuclearTensorProduct-mapCLM, GaussianField-circleHasLaplacianEigenvalues, GaussianField-NuclearTensorProduct-dyninMityaginSpace, GaussianField-SmoothMap-Circle-instAddCommGroup, GaussianField-NuclearTensorProduct-instIsTopologicalAddGroup, GaussianField-nuclearTensorProduct-mapCLM-pure, GaussianField-greenFunctionBilinear-invariant-of-pure, GaussianField-tensorProductHasLaplacianEigenvalues, GaussianField-NuclearTensorProduct, GaussianField-smoothCircle-coeff-basis, GaussianField-NuclearTensorProduct-instAddCommGroup, GaussianField-greenFunctionBilinear-reflection-pure, GaussianField-circleReflection, GaussianField-SmoothMap-Circle-instModule, GaussianField-SmoothMap-Circle-instContinuousSMul, GaussianField-smoothCircle-dyninMityaginSpace, GaussianField-SmoothMap-Circle-instIsTopologicalAddGroup, GaussianField-NuclearTensorProduct-instContinuousSMulReal, GaussianField-NuclearTensorProduct-pure, GaussianField-SmoothMap-Circle, GaussianField-SmoothMap-Circle-instTopologicalSpace, GaussianField-greenFunctionBilinear}
  \texttt{GaussianField.greenFunctionBilinear\_timeReflection\_invariant}
\end{theorem}
\begin{theorem}[\texttt{coeff\_product\_paired\_translation}]\label{GaussianField-coeff-product-paired-translation}
  \lean{GaussianField.coeff_product_paired_translation}\leanok
  \uses{GaussianField-circleTranslation, GaussianField-modePartner-sin, GaussianField-SmoothMap-Circle-instAddCommGroup, GaussianField-SmoothMap-Circle-fourierCoeffReal, GaussianField-paired-coeff-product-circleTranslation, GaussianField-fourierCoeffReal-circleTranslation-zero, GaussianField-SmoothMap-Circle-instModule, GaussianField-SmoothMap-Circle, GaussianField-modePartner-cos, GaussianField-SmoothMap-Circle-instTopologicalSpace, GaussianField-modePartner, Lattice-toSemilatticeInf}
  \texttt{GaussianField.coeff\_product\_paired\_translation}
\end{theorem}
\begin{theorem}[\texttt{greenFunctionBilinear\_reflection\_pure}]\label{GaussianField-greenFunctionBilinear-reflection-pure}
  \lean{GaussianField.greenFunctionBilinear_reflection_pure}\leanok
  \uses{GaussianField-NuclearTensorProduct-instModuleReal, GaussianField-NuclearTensorProduct-instTopologicalSpace, GaussianField-DyninMityaginSpace-coeff, GaussianField-circleHasLaplacianEigenvalues, GaussianField-RapidDecaySeq-val, GaussianField-NuclearTensorProduct-dyninMityaginSpace, GaussianField-SmoothMap-Circle-instAddCommGroup, GaussianField-NuclearTensorProduct-instIsTopologicalAddGroup, GaussianField-tensorProductHasLaplacianEigenvalues, GaussianField-SmoothMap-Circle-fourierCoeffReal, GaussianField-NuclearTensorProduct, GaussianField-NuclearTensorProduct-instAddCommGroup, GaussianField-HasLaplacianEigenvalues-eigenvalue, GaussianField-circleReflection, GaussianField-SmoothMap-Circle-instModule, GaussianField-SmoothMap-Circle-instContinuousSMul, GaussianField-smoothCircle-dyninMityaginSpace, GaussianField-SmoothMap-Circle-instIsTopologicalAddGroup, GaussianField-NuclearTensorProduct-instContinuousSMulReal, GaussianField-NuclearTensorProduct-pure, GaussianField-SmoothMap-Circle, GaussianField-coeff-eq-fourierCoeffReal, GaussianField-SmoothMap-Circle-instTopologicalSpace, GaussianField-greenFunctionBilinear}
  \texttt{GaussianField.greenFunctionBilinear\_reflection\_pure}
\end{theorem}
\begin{theorem}[\texttt{modePartner\_involutive}]\label{GaussianField-modePartner-involutive}
  \lean{GaussianField.modePartner_involutive}\leanok
  \uses{GaussianField-modePartner}
  \texttt{GaussianField.modePartner\_involutive}
\end{theorem}
\begin{theorem}[\texttt{greenFunctionBilinear\_swap\_pure}]\label{GaussianField-greenFunctionBilinear-swap-pure}
  \lean{GaussianField.greenFunctionBilinear_swap_pure}\leanok
  \uses{GaussianField-NuclearTensorProduct-instModuleReal, GaussianField-NuclearTensorProduct-instTopologicalSpace, GaussianField-DyninMityaginSpace-coeff, GaussianField-circleHasLaplacianEigenvalues, GaussianField-NuclearTensorProduct-dyninMityaginSpace, GaussianField-SmoothMap-Circle-instAddCommGroup, GaussianField-NuclearTensorProduct-instIsTopologicalAddGroup, GaussianField-tensorProductHasLaplacianEigenvalues, GaussianField-NuclearTensorProduct, GaussianField-NuclearTensorProduct-instAddCommGroup, GaussianField-HasLaplacianEigenvalues-eigenvalue, GaussianField-SmoothMap-Circle-instModule, GaussianField-SmoothMap-Circle-instContinuousSMul, GaussianField-smoothCircle-dyninMityaginSpace, GaussianField-SmoothMap-Circle-instIsTopologicalAddGroup, GaussianField-NuclearTensorProduct-instContinuousSMulReal, GaussianField-NuclearTensorProduct-pure, GaussianField-SmoothMap-Circle, GaussianField-NuclearTensorProduct-pure-val, GaussianField-SmoothMap-Circle-instTopologicalSpace, GaussianField-greenFunctionBilinear}
  \texttt{GaussianField.greenFunctionBilinear\_swap\_pure}
\end{theorem}
\begin{theorem}[\texttt{greenFunctionBilinear\_invariant\_of\_pure}]\label{GaussianField-greenFunctionBilinear-invariant-of-pure}
  \lean{GaussianField.greenFunctionBilinear_invariant_of_pure}\leanok
  \uses{GaussianField-DyninMityaginSpace, GaussianField-NuclearTensorProduct-instModuleReal, GaussianField-HasLaplacianEigenvalues, GaussianField-NuclearTensorProduct-instTopologicalSpace, GaussianField-DyninMityaginSpace-coeff, GaussianField-NuclearTensorProduct-dyninMityaginSpace, GaussianField-NuclearTensorProduct-instIsTopologicalAddGroup, GaussianField-tensorProductHasLaplacianEigenvalues, GaussianField-DyninMityaginSpace-expansion, GaussianField-NuclearTensorProduct, GaussianField-NuclearTensorProduct-instAddCommGroup, GaussianField-DyninMityaginSpace-basis, GaussianField-NuclearTensorProduct-instContinuousSMulReal, GaussianField-NuclearTensorProduct-pure, GaussianField-greenFunctionBilinear, GaussianField-greenFunctionBilinear-symm}
  \texttt{GaussianField.greenFunctionBilinear\_invariant\_of\_pure}
\end{theorem}
\section{\texttt{GaussianField.StandardGaussianBridge}}
\begin{theorem}[\texttt{gffOrthonormalProj\_measurable}]\label{GaussianField-gffOrthonormalProj-measurable}
  \lean{GaussianField.gffOrthonormalProj_measurable}\leanok
  \uses{GaussianField-FinLatticeSites, GaussianField-gffOrthonormalProj, GaussianField-configuration-eval-measurable, GaussianField-Configuration, GaussianField-massEigenvalues, GaussianField-massEigenvectorBasis, GaussianField-instMeasurableSpaceConfiguration, GaussianField-FinLatticeField}
  \texttt{GaussianField.gffOrthonormalProj\_measurable}
\end{theorem}
\begin{theorem}[\texttt{latticeCovarianceGJ\_eigenvector\_inner\_off}]\label{GaussianField-latticeCovarianceGJ-eigenvector-inner-off}
  \lean{GaussianField.latticeCovarianceGJ_eigenvector_inner_off}\leanok
  \uses{GaussianField-covariance, GaussianField-FinLatticeSites, GaussianField-latticeCovarianceGJ, GaussianField-lattice-covariance-GJ-eq-spectral, GaussianField-massEigenvalues, GaussianField-massEigenvectorBasis, GaussianField-ell2-, GaussianField-FinLatticeField}
  \texttt{GaussianField.latticeCovarianceGJ\_eigenvector\_inner\_off}
\end{theorem}
\begin{definition}[\texttt{gffOrthonormalCoord}]\label{GaussianField-gffOrthonormalCoord}
  \lean{GaussianField.gffOrthonormalCoord}\leanok
  \uses{GaussianField-FinLatticeSites, GaussianField-Configuration, GaussianField-massEigenvalues, GaussianField-massEigenvectorBasis, GaussianField-FinLatticeField}
  \texttt{GaussianField.gffOrthonormalCoord}
\end{definition}
\begin{theorem}[\texttt{gffOrthonormalCoord\_normal}]\label{GaussianField-gffOrthonormalCoord-normal}
  \lean{GaussianField.gffOrthonormalCoord_normal}\leanok
  \uses{GaussianField-latticeCovarianceGJ-eigenvector-inner-self, GaussianField-finLatticeField-dyninMityaginSpace, GaussianField-measure, GaussianField-ell2-separableSpace, GaussianField-FinLatticeSites, GaussianField-latticeCovarianceGJ, GaussianField-gffOrthonormalCoord, GaussianField-pairing-is-gaussian, GaussianField-latticeGaussianMeasure, GaussianField-configuration-eval-measurable, GaussianField-Configuration, GaussianField-massEigenvalues, GaussianField-massEigenvectorBasis, GaussianField-instMeasurableSpaceConfiguration, GaussianField-massOperatorMatrix-eigenvalues-pos, GaussianField-ell2-, GaussianField-FinLatticeField, Lattice-toSemilatticeInf}
  \texttt{GaussianField.gffOrthonormalCoord\_normal}
\end{theorem}
\begin{theorem}[\texttt{gffOrthonormalProj\_charFun}]\label{GaussianField-gffOrthonormalProj-charFun}
  \lean{GaussianField.gffOrthonormalProj_charFun}\leanok
  \uses{GaussianField-FinLatticeSites, GaussianField-gffOrthonormalProj-pushforward-eq-stdGaussian, GaussianField-gffOrthonormalProj, GaussianField-latticeGaussianMeasure, GaussianField-Configuration, GaussianField-instMeasurableSpaceConfiguration, GaussianField-FinLatticeField}
  \texttt{GaussianField.gffOrthonormalProj\_charFun}
\end{theorem}
\begin{definition}[\texttt{gffOrthonormalProj}]\label{GaussianField-gffOrthonormalProj}
  \lean{GaussianField.gffOrthonormalProj}\leanok
  \uses{GaussianField-FinLatticeSites, GaussianField-gffOrthonormalCoord, GaussianField-Configuration, GaussianField-FinLatticeField}
  \texttt{GaussianField.gffOrthonormalProj}
\end{definition}
\begin{theorem}[\texttt{gffOrthonormalProj\_pushforward\_eq\_stdGaussian}]\label{GaussianField-gffOrthonormalProj-pushforward-eq-stdGaussian}
  \lean{GaussianField.gffOrthonormalProj_pushforward_eq_stdGaussian}\leanok
  \uses{GaussianField-gffOrthonormalCoord-independent, GaussianField-FinLatticeSites, GaussianField-latticeGaussianMeasure-isProbability, GaussianField-gffOrthonormalProj, GaussianField-gffOrthonormalCoord, GaussianField-latticeGaussianMeasure, GaussianField-configuration-eval-measurable, GaussianField-Configuration, GaussianField-massEigenvalues, GaussianField-massEigenvectorBasis, GaussianField-instMeasurableSpaceConfiguration, GaussianField-FinLatticeField, GaussianField-gffOrthonormalCoord-normal}
  \texttt{GaussianField.gffOrthonormalProj\_pushforward\_eq\_stdGaussian}
\end{theorem}
\begin{theorem}[\texttt{gffOrthonormalCoord\_independent}]\label{GaussianField-gffOrthonormalCoord-independent}
  \lean{GaussianField.gffOrthonormalCoord_independent}\leanok
  \uses{GaussianField-cross-moment-eq-covariance, GaussianField-finLatticeField-dyninMityaginSpace, GaussianField-measure, GaussianField-ell2-separableSpace, GaussianField-FinLatticeSites, GaussianField-latticeCovarianceGJ, GaussianField-pairing-memLp, GaussianField-measure-centered, GaussianField-latticeCovarianceGJ-eigenvector-inner-off, GaussianField-latticeGaussianMeasure-isProbability, GaussianField-gffOrthonormalCoord, GaussianField-latticeGaussianMeasure, GaussianField-Configuration, GaussianField-massEigenvalues, GaussianField-massEigenvectorBasis, GaussianField-instMeasurableSpaceConfiguration, GaussianField-ell2-, GaussianField-FinLatticeField}
  \texttt{GaussianField.gffOrthonormalCoord\_independent}
\end{theorem}
\begin{theorem}[\texttt{latticeGaussianMeasure\_isGaussian}]\label{GaussianField-latticeGaussianMeasure-isGaussian}
  \lean{GaussianField.latticeGaussianMeasure_isGaussian}\leanok
  \uses{GaussianField-finLatticeField-dyninMityaginSpace, GaussianField-ell2-separableSpace, GaussianField-FinLatticeSites, GaussianField-latticeCovarianceGJ, GaussianField-latticeGaussianMeasure, GaussianField-measure-isGaussian, GaussianField-Configuration, GaussianField-instMeasurableSpaceConfiguration, GaussianField-ell2-, GaussianField-FinLatticeField}
  \texttt{GaussianField.latticeGaussianMeasure\_isGaussian}
\end{theorem}
\begin{theorem}[\texttt{latticeCovarianceGJ\_eigenvector\_inner\_self}]\label{GaussianField-latticeCovarianceGJ-eigenvector-inner-self}
  \lean{GaussianField.latticeCovarianceGJ_eigenvector_inner_self}\leanok
  \uses{GaussianField-covariance, GaussianField-FinLatticeSites, GaussianField-latticeCovarianceGJ, GaussianField-lattice-covariance-GJ-eq-spectral, GaussianField-massEigenvalues, GaussianField-massEigenvectorBasis, GaussianField-ell2-, GaussianField-FinLatticeField}
  \texttt{GaussianField.latticeCovarianceGJ\_eigenvector\_inner\_self}
\end{theorem}
\section{\texttt{Nuclear.TensorProductFunctorAxioms}}
\begin{definition}[\texttt{mapImage}]\label{GaussianField-mapImage}
  \lean{GaussianField.mapImage}\leanok
  \uses{GaussianField-DyninMityaginSpace, GaussianField-NuclearTensorProduct, GaussianField-DyninMityaginSpace-basis, GaussianField-NuclearTensorProduct-pure}
  \texttt{GaussianField.mapImage}
\end{definition}
\begin{theorem}[\texttt{nuclearTensorProduct\_mapCLM\_id}]\label{GaussianField-nuclearTensorProduct-mapCLM-id}
  \lean{GaussianField.nuclearTensorProduct_mapCLM_id}\leanok
  \uses{GaussianField-DyninMityaginSpace, GaussianField-NuclearTensorProduct-ext, GaussianField-NuclearTensorProduct-instModuleReal, GaussianField-NuclearTensorProduct-instTopologicalSpace, GaussianField-DyninMityaginSpace-coeff, GaussianField-nuclearTensorProduct-mapCLM, GaussianField-RapidDecaySeq-val, GaussianField-NuclearTensorProduct, GaussianField-DyninMityaginSpace-HasBiorthogonalBasis-coeff-basis, GaussianField-NuclearTensorProduct-instAddCommGroup, GaussianField-DyninMityaginSpace-basis, GaussianField-DyninMityaginSpace-HasBiorthogonalBasis}
  \texttt{GaussianField.nuclearTensorProduct\_mapCLM\_id}
\end{theorem}
\begin{theorem}[\texttt{nuclearTensorProduct\_mapCLM\_pure}]\label{GaussianField-nuclearTensorProduct-mapCLM-pure}
  \lean{GaussianField.nuclearTensorProduct_mapCLM_pure}\leanok
  \uses{GaussianField-DyninMityaginSpace, GaussianField-NuclearTensorProduct-ext, GaussianField-NuclearTensorProduct-instModuleReal, GaussianField-NuclearTensorProduct-instTopologicalSpace, GaussianField-DyninMityaginSpace-coeff, GaussianField-nuclearTensorProduct-mapCLM, GaussianField-RapidDecaySeq-val, GaussianField-DyninMityaginSpace-expansion, GaussianField-NuclearTensorProduct, GaussianField-NuclearTensorProduct-instAddCommGroup, GaussianField-DyninMityaginSpace-basis, GaussianField-NuclearTensorProduct-pure}
  \texttt{GaussianField.nuclearTensorProduct\_mapCLM\_pure}
\end{theorem}
\begin{definition}[\texttt{nuclearTensorProduct\_swapCLM}]\label{GaussianField-nuclearTensorProduct-swapCLM}
  \lean{GaussianField.nuclearTensorProduct_swapCLM}\leanok
  \uses{GaussianField-DyninMityaginSpace, GaussianField-NuclearTensorProduct-instModuleReal, GaussianField-NuclearTensorProduct-instTopologicalSpace, GaussianField-NuclearTensorProduct, GaussianField-NuclearTensorProduct-instAddCommGroup}
  \texttt{GaussianField.nuclearTensorProduct\_swapCLM}
\end{definition}
\begin{definition}[\texttt{nuclearTensorProduct\_mapCLM}]\label{GaussianField-nuclearTensorProduct-mapCLM}
  \lean{GaussianField.nuclearTensorProduct_mapCLM}\leanok
  \uses{GaussianField-DyninMityaginSpace, GaussianField-NuclearTensorProduct-instModuleReal, GaussianField-NuclearTensorProduct-instTopologicalSpace, GaussianField-RapidDecaySeq, GaussianField-RapidDecaySeq-instModule, GaussianField-NuclearTensorProduct, GaussianField-NuclearTensorProduct-instAddCommGroup, GaussianField-RapidDecaySeq-instAddCommGroup}
  \texttt{GaussianField.nuclearTensorProduct\_mapCLM}
\end{definition}
\begin{theorem}[\texttt{mapImage\_seminorm\_bound}]\label{GaussianField-mapImage-seminorm-bound}
  \lean{GaussianField.mapImage_seminorm_bound}\leanok
  \uses{GaussianField-DyninMityaginSpace, GaussianField-mapImage, GaussianField-RapidDecaySeq-rapidDecaySeminorm, GaussianField-RapidDecaySeq, GaussianField-RapidDecaySeq-instSMul, GaussianField-DyninMityaginSpace--, GaussianField-finset-sup-poly-bound, GaussianField-DyninMityaginSpace-h-with, GaussianField-DyninMityaginSpace-basis, GaussianField-NuclearTensorProduct-pure, GaussianField-RapidDecaySeq-instAddCommGroup, GaussianField-DyninMityaginSpace-basis-growth, GaussianField-DyninMityaginSpace-p, GaussianField-NuclearTensorProduct-pure-seminorm-bound}
  \texttt{GaussianField.mapImage\_seminorm\_bound}
\end{theorem}
\begin{theorem}[\texttt{evalCLM\_comp\_swapCLM}]\label{GaussianField-evalCLM-comp-swapCLM}
  \lean{GaussianField.evalCLM_comp_swapCLM}\leanok
  \uses{GaussianField-DyninMityaginSpace, GaussianField-NuclearTensorProduct-instModuleReal, GaussianField-NuclearTensorProduct-instTopologicalSpace, GaussianField-RapidDecaySeq-val, GaussianField-nuclearTensorProduct-swapCLM, GaussianField-NuclearTensorProduct, GaussianField-NuclearTensorProduct-instAddCommGroup, GaussianField-DyninMityaginSpace-basis, GaussianField-NuclearTensorProduct-evalCLM}
  \texttt{GaussianField.evalCLM\_comp\_swapCLM}
\end{theorem}
\begin{theorem}[\texttt{val\_le\_seminorm}]\label{GaussianField-val-le-seminorm}
  \lean{GaussianField.val_le_seminorm}\leanok
  \uses{GaussianField-RapidDecaySeq-rapidDecaySeminorm, GaussianField-RapidDecaySeq, GaussianField-RapidDecaySeq-val, GaussianField-RapidDecaySeq-instSMul, GaussianField-RapidDecaySeq-instAddCommGroup, GaussianField-RapidDecaySeq-rapid-decay, Lattice-toSemilatticeInf}
  \texttt{GaussianField.val\_le\_seminorm}
\end{theorem}
\begin{theorem}[\texttt{nuclearTensorProduct\_swapCLM\_pure}]\label{GaussianField-nuclearTensorProduct-swapCLM-pure}
  \lean{GaussianField.nuclearTensorProduct_swapCLM_pure}\leanok
  \uses{GaussianField-DyninMityaginSpace, GaussianField-NuclearTensorProduct-ext, GaussianField-NuclearTensorProduct-instModuleReal, GaussianField-NuclearTensorProduct-instTopologicalSpace, GaussianField-DyninMityaginSpace-coeff, GaussianField-RapidDecaySeq-val, GaussianField-nuclearTensorProduct-swapCLM, GaussianField-NuclearTensorProduct, GaussianField-NuclearTensorProduct-instAddCommGroup, GaussianField-NuclearTensorProduct-pure}
  \texttt{GaussianField.nuclearTensorProduct\_swapCLM\_pure}
\end{theorem}
\begin{theorem}[\texttt{evalCLM\_comp\_mapCLM}]\label{GaussianField-evalCLM-comp-mapCLM}
  \lean{GaussianField.evalCLM_comp_mapCLM}\leanok
  \uses{GaussianField-DyninMityaginSpace, GaussianField-NuclearTensorProduct-instModuleReal, GaussianField-NuclearTensorProduct-instTopologicalSpace, GaussianField-DyninMityaginSpace-coeff, GaussianField-RapidDecaySeq, GaussianField-nuclearTensorProduct-mapCLM, GaussianField-RapidDecaySeq-val, GaussianField-nuclearTensorProduct-mapCLM-pure, GaussianField-RapidDecaySeq-instModule, GaussianField-NuclearTensorProduct, GaussianField-RapidDecaySeq-basisVec, GaussianField-NuclearTensorProduct-instAddCommGroup, GaussianField-RapidDecaySeq-rapidDecay-expansion, GaussianField-DyninMityaginSpace-basis, GaussianField-NuclearTensorProduct-pure, GaussianField-RapidDecaySeq-instAddCommGroup, GaussianField-RapidDecaySeq-instTopologicalSpace, GaussianField-NuclearTensorProduct-basisVec-eq-pure, GaussianField-NuclearTensorProduct-evalCLM-pure, GaussianField-NuclearTensorProduct-evalCLM}
  \texttt{GaussianField.evalCLM\_comp\_mapCLM}
\end{theorem}
\begin{theorem}[\texttt{nuclearTensorProduct\_mapCLM\_comp}]\label{GaussianField-nuclearTensorProduct-mapCLM-comp}
  \lean{GaussianField.nuclearTensorProduct_mapCLM_comp}\leanok
  \uses{GaussianField-DyninMityaginSpace, GaussianField-NuclearTensorProduct-ext, GaussianField-mapImage, GaussianField-NuclearTensorProduct-instModuleReal, GaussianField-NuclearTensorProduct-instTopologicalSpace, GaussianField-RapidDecaySeq, GaussianField-nuclearTensorProduct-mapCLM, GaussianField-RapidDecaySeq-val, GaussianField-nuclearTensorProduct-mapCLM-pure, GaussianField-RapidDecaySeq-instModule, GaussianField-NuclearTensorProduct, GaussianField-RapidDecaySeq-basisVec, GaussianField-NuclearTensorProduct-instAddCommGroup, GaussianField-RapidDecaySeq-rapidDecay-expansion, GaussianField-DyninMityaginSpace-basis, GaussianField-NuclearTensorProduct-pure, GaussianField-RapidDecaySeq-instAddCommGroup, GaussianField-RapidDecaySeq-instTopologicalSpace, GaussianField-DyninMityaginSpace-HasBiorthogonalBasis, GaussianField-RapidDecaySeq-coeffCLM}
  \texttt{GaussianField.nuclearTensorProduct\_mapCLM\_comp}
\end{theorem}
\section{\texttt{SmoothCircle.HeatSemigroup}}
\begin{theorem}[\texttt{heatWeight\_zero}]\label{GaussianField-heatWeight-zero}
  \lean{GaussianField.heatWeight_zero}\leanok
  \uses{GaussianField-circleHasLaplacianEigenvalues, GaussianField-SmoothMap-Circle-instAddCommGroup, GaussianField-heatWeight, GaussianField-HasLaplacianEigenvalues-eigenvalue, GaussianField-SmoothMap-Circle-instModule, GaussianField-SmoothMap-Circle-instContinuousSMul, GaussianField-smoothCircle-dyninMityaginSpace, GaussianField-SmoothMap-Circle-instIsTopologicalAddGroup, GaussianField-SmoothMap-Circle, GaussianField-SmoothMap-Circle-instTopologicalSpace}
  \texttt{GaussianField.heatWeight\_zero}
\end{theorem}
\begin{definition}[\texttt{heatWeight}]\label{GaussianField-heatWeight}
  \lean{GaussianField.heatWeight}\leanok
  \uses{GaussianField-circleHasLaplacianEigenvalues, GaussianField-SmoothMap-Circle-instAddCommGroup, GaussianField-HasLaplacianEigenvalues-eigenvalue, GaussianField-SmoothMap-Circle-instModule, GaussianField-SmoothMap-Circle-instContinuousSMul, GaussianField-smoothCircle-dyninMityaginSpace, GaussianField-SmoothMap-Circle-instIsTopologicalAddGroup, GaussianField-SmoothMap-Circle, GaussianField-SmoothMap-Circle-instTopologicalSpace}
  \texttt{GaussianField.heatWeight}
\end{definition}
\begin{theorem}[\texttt{congr\_simp}]\label{GaussianField-heatWeight-congr-simp}
  \lean{GaussianField.heatWeight.congr_simp}\leanok
  \uses{GaussianField-heatWeight}
  \texttt{GaussianField.heatWeight.congr\_simp}
\end{theorem}
\begin{theorem}[\texttt{heatWeight\_nonneg}]\label{GaussianField-heatWeight-nonneg}
  \lean{GaussianField.heatWeight_nonneg}\leanok
  \uses{GaussianField-circleHasLaplacianEigenvalues, GaussianField-SmoothMap-Circle-instAddCommGroup, GaussianField-heatWeight, GaussianField-HasLaplacianEigenvalues-eigenvalue, GaussianField-SmoothMap-Circle-instModule, GaussianField-SmoothMap-Circle-instContinuousSMul, GaussianField-smoothCircle-dyninMityaginSpace, GaussianField-SmoothMap-Circle-instIsTopologicalAddGroup, GaussianField-SmoothMap-Circle, GaussianField-SmoothMap-Circle-instTopologicalSpace}
  \texttt{GaussianField.heatWeight\_nonneg}
\end{theorem}
\begin{theorem}[\texttt{heatWeight\_add}]\label{GaussianField-heatWeight-add}
  \lean{GaussianField.heatWeight_add}\leanok
  \uses{GaussianField-circleHasLaplacianEigenvalues, GaussianField-SmoothMap-Circle-instAddCommGroup, GaussianField-heatWeight, GaussianField-HasLaplacianEigenvalues-eigenvalue, GaussianField-SmoothMap-Circle-instModule, GaussianField-SmoothMap-Circle-instContinuousSMul, GaussianField-smoothCircle-dyninMityaginSpace, GaussianField-SmoothMap-Circle-instIsTopologicalAddGroup, GaussianField-SmoothMap-Circle, GaussianField-SmoothMap-Circle-instTopologicalSpace}
  \texttt{GaussianField.heatWeight\_add}
\end{theorem}
\begin{theorem}[\texttt{circleHeatSemigroup\_fourierCoeff}]\label{GaussianField-circleHeatSemigroup-fourierCoeff}
  \lean{GaussianField.circleHeatSemigroup_fourierCoeff}\leanok
  \uses{GaussianField-RapidDecaySeq, GaussianField-RapidDecaySeq-val, GaussianField-SmoothMap-Circle-smoothCircleRapidDecayEquiv, GaussianField-SmoothMap-Circle-instAddCommGroup, GaussianField-heatWeight, GaussianField-RapidDecaySeq-instModule, GaussianField-SmoothMap-Circle-instModule, GaussianField-RapidDecaySeq-instAddCommGroup, GaussianField-SmoothMap-Circle, GaussianField-RapidDecaySeq-instTopologicalSpace, GaussianField-circleHeatSemigroup, GaussianField-heatWeight-abs-le-one, GaussianField-SmoothMap-Circle-instTopologicalSpace}
  \texttt{GaussianField.circleHeatSemigroup\_fourierCoeff}
\end{theorem}
\begin{theorem}[\texttt{circleHeatSemigroup\_zero}]\label{GaussianField-circleHeatSemigroup-zero}
  \lean{GaussianField.circleHeatSemigroup_zero}\leanok
  \uses{GaussianField-SmoothMap-Circle-instFunLike, GaussianField-SmoothMap-Circle-ext, GaussianField-RapidDecaySeq, GaussianField-circleHasLaplacianEigenvalues, GaussianField-RapidDecaySeq-val, GaussianField-SmoothMap-Circle-smoothCircleRapidDecayEquiv, GaussianField-SmoothMap-Circle-instAddCommGroup, GaussianField-heatWeight, GaussianField-RapidDecaySeq-instModule, GaussianField-HasLaplacianEigenvalues-eigenvalue, GaussianField-RapidDecaySeq-ext, GaussianField-SmoothMap-Circle-instModule, GaussianField-SmoothMap-Circle-instContinuousSMul, GaussianField-smoothCircle-dyninMityaginSpace, GaussianField-SmoothMap-Circle-instIsTopologicalAddGroup, GaussianField-RapidDecaySeq-instAddCommGroup, GaussianField-SmoothMap-Circle, GaussianField-RapidDecaySeq-instTopologicalSpace, GaussianField-circleHeatSemigroup, GaussianField-heatWeight-abs-le-one, GaussianField-SmoothMap-Circle-instTopologicalSpace}
  \texttt{GaussianField.circleHeatSemigroup\_zero}
\end{theorem}
\begin{theorem}[\texttt{heatWeight\_le\_one}]\label{GaussianField-heatWeight-le-one}
  \lean{GaussianField.heatWeight_le_one}\leanok
  \uses{GaussianField-circleHasLaplacianEigenvalues, GaussianField-SmoothMap-Circle-instAddCommGroup, GaussianField-heatWeight, GaussianField-HasLaplacianEigenvalues-eigenvalue, GaussianField-SmoothMap-Circle-instModule, GaussianField-SmoothMap-Circle-instContinuousSMul, GaussianField-HasLaplacianEigenvalues-eigenvalue-nonneg, GaussianField-smoothCircle-dyninMityaginSpace, GaussianField-SmoothMap-Circle-instIsTopologicalAddGroup, GaussianField-SmoothMap-Circle, GaussianField-SmoothMap-Circle-instTopologicalSpace}
  \texttt{GaussianField.heatWeight\_le\_one}
\end{theorem}
\begin{theorem}[\texttt{circleHeatSemigroup\_fourierBasis}]\label{GaussianField-circleHeatSemigroup-fourierBasis}
  \lean{GaussianField.circleHeatSemigroup_fourierBasis}\leanok
  \uses{GaussianField-RapidDecaySeq, GaussianField-SmoothMap-Circle-fourierBasis, GaussianField-RapidDecaySeq-val, GaussianField-SmoothMap-Circle-smoothCircleRapidDecayEquiv, GaussianField-SmoothMap-Circle-instAddCommGroup, GaussianField-heatWeight, GaussianField-SmoothMap-Circle-fourierCoeffReal-fourierBasis, GaussianField-RapidDecaySeq-instSMul, GaussianField-RapidDecaySeq-instModule, GaussianField-SmoothMap-Circle-instSMul, GaussianField-SmoothMap-Circle-toRapidDecayLM, GaussianField-RapidDecaySeq-ext, GaussianField-SmoothMap-Circle-instModule, GaussianField-RapidDecaySeq-instAddCommGroup, GaussianField-SmoothMap-Circle, GaussianField-RapidDecaySeq-instTopologicalSpace, GaussianField-circleHeatSemigroup, GaussianField-heatWeight-abs-le-one, GaussianField-SmoothMap-Circle-instTopologicalSpace}
  \texttt{GaussianField.circleHeatSemigroup\_fourierBasis}
\end{theorem}
\begin{definition}[\texttt{circleHeatSemigroup}]\label{GaussianField-circleHeatSemigroup}
  \lean{GaussianField.circleHeatSemigroup}\leanok
  \uses{GaussianField-RapidDecaySeq, GaussianField-SmoothMap-Circle-smoothCircleRapidDecayEquiv, GaussianField-SmoothMap-Circle-instAddCommGroup, GaussianField-heatWeight, GaussianField-RapidDecaySeq-instModule, GaussianField-SmoothMap-Circle-instModule, GaussianField-RapidDecaySeq-instAddCommGroup, GaussianField-SmoothMap-Circle, GaussianField-RapidDecaySeq-instTopologicalSpace, GaussianField-heatWeight-abs-le-one, GaussianField-SmoothMap-Circle-instTopologicalSpace}
  \texttt{GaussianField.circleHeatSemigroup}
\end{definition}
\begin{theorem}[\texttt{heatWeight\_abs\_le\_one}]\label{GaussianField-heatWeight-abs-le-one}
  \lean{GaussianField.heatWeight_abs_le_one}\leanok
  \uses{GaussianField-heatWeight, GaussianField-heatWeight-nonneg, GaussianField-heatWeight-le-one, Lattice-toSemilatticeInf}
  \texttt{GaussianField.heatWeight\_abs\_le\_one}
\end{theorem}
\section{\texttt{Lattice.LatticeFourierIndexing}}
\begin{theorem}[\texttt{sum\_latticeEigenvalue\_two\_eq\_sum\_latticeEigenvalue1d\_pair}]\label{GaussianField-sum-latticeEigenvalue-two-eq-sum-latticeEigenvalue1d-pair}
  \lean{GaussianField.sum_latticeEigenvalue_two_eq_sum_latticeEigenvalue1d_pair}\leanok
  \uses{GaussianField-latticeEigenvalue1d, GaussianField-FinLatticeSites, GaussianField-latticeLaplacianEigenvalue, GaussianField-sum-rawFrequency-pair-eq-sum-latticeEigenvalue1d-pair, GaussianField-latticeEigenvalue}
  \texttt{GaussianField.sum\_latticeEigenvalue\_two\_eq\_sum\_latticeEigenvalue1d\_pair}
\end{theorem}
\begin{theorem}[\texttt{sum\_rawFrequency\_pair\_eq\_sum\_latticeEigenvalue1d\_pair}]\label{GaussianField-sum-rawFrequency-pair-eq-sum-latticeEigenvalue1d-pair}
  \lean{GaussianField.sum_rawFrequency_pair_eq_sum_latticeEigenvalue1d_pair}\leanok
  \uses{GaussianField-latticeEigenvalue1d, GaussianField-sum-rawFrequency-eq-sum-latticeEigenvalue1d}
  \texttt{GaussianField.sum\_rawFrequency\_pair\_eq\_sum\_latticeEigenvalue1d\_pair}
\end{theorem}
\begin{theorem}[\texttt{sum\_rawFrequency\_eq\_sum\_latticeEigenvalue1d}]\label{GaussianField-sum-rawFrequency-eq-sum-latticeEigenvalue1d}
  \lean{GaussianField.sum_rawFrequency_eq_sum_latticeEigenvalue1d}\leanok
  \uses{GaussianField-latticeEigenvalue1d}
  \texttt{GaussianField.sum\_rawFrequency\_eq\_sum\_latticeEigenvalue1d}
\end{theorem}
\begin{theorem}[\texttt{sum\_latticeEigenvalue\_eq\_sum\_latticeEigenvalue1d\_family}]\label{GaussianField-sum-latticeEigenvalue-eq-sum-latticeEigenvalue1d-family}
  \lean{GaussianField.sum_latticeEigenvalue_eq_sum_latticeEigenvalue1d_family}\leanok
  \uses{GaussianField-latticeEigenvalue1d, GaussianField-FinLatticeSites, GaussianField-latticeLaplacianEigenvalue, GaussianField-latticeEigenvalue}
  \texttt{GaussianField.sum\_latticeEigenvalue\_eq\_sum\_latticeEigenvalue1d\_family}
\end{theorem}
\section{\texttt{Nuclear.PointEval}}
\begin{definition}[\texttt{eval}]\label{GaussianField-HasPointEval-eval}
  \lean{GaussianField.HasPointEval.eval}\leanok
  \uses{GaussianField-HasPointEval, GaussianField-HasPointEval-pointEval}
  \texttt{GaussianField.HasPointEval.eval}
\end{definition}
\begin{definition}[\texttt{smoothCircle\_hasPointEval}]\label{GaussianField-smoothCircle-hasPointEval}
  \lean{GaussianField.smoothCircle_hasPointEval}\leanok
  \uses{GaussianField-SmoothMap-Circle-toFun, GaussianField-HasPointEval, GaussianField-SmoothMap-Circle}
  \texttt{GaussianField.smoothCircle\_hasPointEval}
\end{definition}
\begin{definition}[\texttt{ctorIdx}]\label{GaussianField-HasPointEval-ctorIdx}
  \lean{GaussianField.HasPointEval.ctorIdx}\leanok
  \uses{GaussianField-HasPointEval}
  \texttt{GaussianField.HasPointEval.ctorIdx}
\end{definition}
\begin{definition}[\texttt{nuclearTensorProduct\_hasPointEval}]\label{GaussianField-nuclearTensorProduct-hasPointEval}
  \lean{GaussianField.nuclearTensorProduct_hasPointEval}\leanok
  \uses{GaussianField-RapidDecaySeq-val, GaussianField-NuclearTensorProduct, GaussianField-HasPointEval}
  \texttt{GaussianField.nuclearTensorProduct\_hasPointEval}
\end{definition}
\begin{definition}[\texttt{pointEval}]\label{GaussianField-HasPointEval-pointEval}
  \lean{GaussianField.HasPointEval.pointEval}\leanok
  \uses{GaussianField-HasPointEval}
  \texttt{GaussianField.HasPointEval.pointEval}
\end{definition}
\begin{definition}[\texttt{HasPointEval}]\label{GaussianField-HasPointEval}
  \lean{GaussianField.HasPointEval}\leanok
  \texttt{GaussianField.HasPointEval}
\end{definition}
\begin{definition}[\texttt{schwartz\_hasPointEval}]\label{GaussianField-schwartz-hasPointEval}
  \lean{GaussianField.schwartz_hasPointEval}\leanok
  \uses{GaussianField-HasPointEval}
  \texttt{GaussianField.schwartz\_hasPointEval}
\end{definition}
\begin{definition}[\texttt{fin\_hasPointEval}]\label{GaussianField-fin-hasPointEval}
  \lean{GaussianField.fin_hasPointEval}\leanok
  \uses{GaussianField-HasPointEval}
  \texttt{GaussianField.fin\_hasPointEval}
\end{definition}
\section{\texttt{SchwartzNuclear.Periodization}}
\begin{theorem}[\texttt{periodizeCLM\_comp\_schwartzReflection}]\label{GaussianField-periodizeCLM-comp-schwartzReflection}
  \lean{GaussianField.periodizeCLM_comp_schwartzReflection}\leanok
  \uses{GaussianField-SmoothMap-Circle-instFunLike, GaussianField-SmoothMap-Circle-ext, GaussianField-periodizeCLM-apply, GaussianField-SmoothMap-Circle-instAddCommGroup, GaussianField-periodizeCLM, GaussianField-SmoothMap-Circle-toFun, GaussianField-circleReflection, GaussianField-SmoothMap-Circle-instModule, GaussianField-SmoothMap-Circle, GaussianField-schwartzReflection, GaussianField-SmoothMap-Circle-instTopologicalSpace}
  \texttt{GaussianField.periodizeCLM\_comp\_schwartzReflection}
\end{theorem}
\begin{theorem}[\texttt{periodizeFun\_periodic}]\label{GaussianField-periodizeFun-periodic}
  \lean{GaussianField.periodizeFun_periodic}\leanok
  \uses{GaussianField-periodizeFun}
  \texttt{GaussianField.periodizeFun\_periodic}
\end{theorem}
\begin{definition}[\texttt{periodizeLM}]\label{GaussianField-periodizeLM}
  \lean{GaussianField.periodizeLM}\leanok
  \uses{GaussianField-periodizeSmoothCircle-add, GaussianField-SmoothMap-Circle-instAddCommGroup, GaussianField-SmoothMap-Circle-instModule, GaussianField-SmoothMap-Circle, GaussianField-periodizeSmoothCircle}
  \texttt{GaussianField.periodizeLM}
\end{definition}
\begin{theorem}[\texttt{periodize\_summable}]\label{GaussianField-periodize-summable}
  \lean{GaussianField.periodize_summable}\leanok
  \uses{Lattice-toSemilatticeInf}
  \texttt{GaussianField.periodize\_summable}
\end{theorem}
\begin{definition}[\texttt{periodizeFun}]\label{GaussianField-periodizeFun}
  \lean{GaussianField.periodizeFun}\leanok
  \texttt{GaussianField.periodizeFun}
\end{definition}
\begin{theorem}[\texttt{periodizeCLM\_comp\_schwartzTranslation}]\label{GaussianField-periodizeCLM-comp-schwartzTranslation}
  \lean{GaussianField.periodizeCLM_comp_schwartzTranslation}\leanok
  \uses{GaussianField-circleTranslation, GaussianField-SmoothMap-Circle-instFunLike, GaussianField-SmoothMap-Circle-ext, GaussianField-periodizeCLM-apply, GaussianField-SmoothMap-Circle-instAddCommGroup, GaussianField-periodizeCLM, GaussianField-SmoothMap-Circle-toFun, GaussianField-SmoothMap-Circle-instModule, GaussianField-SmoothMap-Circle, GaussianField-schwartzTranslation, GaussianField-SmoothMap-Circle-instTopologicalSpace}
  \texttt{GaussianField.periodizeCLM\_comp\_schwartzTranslation}
\end{theorem}
\begin{theorem}[\texttt{periodizeCLM\_apply}]\label{GaussianField-periodizeCLM-apply}
  \lean{GaussianField.periodizeCLM_apply}\leanok
  \uses{GaussianField-SmoothMap-Circle-instAddCommGroup, GaussianField-periodizeCLM, GaussianField-SmoothMap-Circle-toFun, GaussianField-SmoothMap-Circle-instModule, GaussianField-SmoothMap-Circle, GaussianField-SmoothMap-Circle-instTopologicalSpace}
  \texttt{GaussianField.periodizeCLM\_apply}
\end{theorem}
\begin{theorem}[\texttt{periodizeSmoothCircle\_toFun}]\label{GaussianField-periodizeSmoothCircle-toFun}
  \lean{GaussianField.periodizeSmoothCircle_toFun}\leanok
  \uses{GaussianField-SmoothMap-Circle-toFun, GaussianField-periodizeSmoothCircle}
  \texttt{GaussianField.periodizeSmoothCircle\_toFun}
\end{theorem}
\begin{definition}[\texttt{periodizeSmoothCircle}]\label{GaussianField-periodizeSmoothCircle}
  \lean{GaussianField.periodizeSmoothCircle}\leanok
  \uses{GaussianField-periodizeFun, GaussianField-periodize-smooth, GaussianField-periodizeFun-periodic, GaussianField-SmoothMap-Circle}
  \texttt{GaussianField.periodizeSmoothCircle}
\end{definition}
\begin{definition}[\texttt{periodizeCLM}]\label{GaussianField-periodizeCLM}
  \lean{GaussianField.periodizeCLM}\leanok
  \uses{GaussianField-periodizeLM, GaussianField-SmoothMap-Circle-instAddCommGroup, GaussianField-SmoothMap-Circle-instModule, GaussianField-SmoothMap-Circle, GaussianField-SmoothMap-Circle-instTopologicalSpace}
  \texttt{GaussianField.periodizeCLM}
\end{definition}
\begin{theorem}[\texttt{periodize\_sobolevSeminorm\_le}]\label{GaussianField-periodize-sobolevSeminorm-le}
  \lean{GaussianField.periodize_sobolevSeminorm_le}\leanok
  \uses{GaussianField-periodizeFun, GaussianField-SmoothMap-Circle-instFunLike, GaussianField-periodize-summable, GaussianField-SmoothMap-Circle-Icc-nonempty, GaussianField-SmoothMap-Circle-instAddCommGroup, GaussianField-SmoothMap-Circle-instSMul, GaussianField-SmoothMap-Circle, GaussianField-periodizeSmoothCircle, GaussianField-SmoothMap-Circle-sobolevSeminorm, Lattice-toSemilatticeInf}
  \texttt{GaussianField.periodize\_sobolevSeminorm\_le}
\end{theorem}
\begin{theorem}[\texttt{periodize\_smooth}]\label{GaussianField-periodize-smooth}
  \lean{GaussianField.periodize_smooth}\leanok
  \uses{GaussianField-periodizeFun}
  \texttt{GaussianField.periodize\_smooth}
\end{theorem}
\begin{theorem}[\texttt{periodizeSmoothCircle\_smul}]\label{GaussianField-periodizeSmoothCircle-smul}
  \lean{GaussianField.periodizeSmoothCircle_smul}\leanok
  \uses{GaussianField-periodizeFun, GaussianField-SmoothMap-Circle-instFunLike, GaussianField-SmoothMap-Circle-ext, GaussianField-periodize-smooth, GaussianField-SmoothMap-Circle-instSMul, GaussianField-periodizeFun-periodic, GaussianField-SmoothMap-Circle, GaussianField-periodizeSmoothCircle}
  \texttt{GaussianField.periodizeSmoothCircle\_smul}
\end{theorem}
\begin{theorem}[\texttt{periodizeSmoothCircle\_add}]\label{GaussianField-periodizeSmoothCircle-add}
  \lean{GaussianField.periodizeSmoothCircle_add}\leanok
  \uses{GaussianField-periodizeFun, GaussianField-SmoothMap-Circle-instFunLike, GaussianField-SmoothMap-Circle-ext, GaussianField-SmoothMap-Circle-instAdd, GaussianField-periodize-summable, GaussianField-periodize-smooth, GaussianField-periodizeFun-periodic, GaussianField-SmoothMap-Circle, GaussianField-periodizeSmoothCircle}
  \texttt{GaussianField.periodizeSmoothCircle\_add}
\end{theorem}
\begin{theorem}[\texttt{periodizeCLM\_eq\_on\_large\_period}]\label{GaussianField-periodizeCLM-eq-on-large-period}
  \lean{GaussianField.periodizeCLM_eq_on_large_period}\leanok
  \uses{GaussianField-periodizeCLM-apply, GaussianField-SmoothMap-Circle-instAddCommGroup, GaussianField-periodizeCLM, GaussianField-SmoothMap-Circle-toFun, GaussianField-SmoothMap-Circle-instModule, GaussianField-SmoothMap-Circle, GaussianField-SmoothMap-Circle-instTopologicalSpace}
  \texttt{GaussianField.periodizeCLM\_eq\_on\_large\_period}
\end{theorem}
\section{\texttt{SchwartzNuclear.SchwartzSlicing}}
\begin{definition}[\texttt{euclideanSnoc}]\label{GaussianField-euclideanSnoc}
  \lean{GaussianField.euclideanSnoc}\leanok
  \texttt{GaussianField.euclideanSnoc}
\end{definition}
\begin{theorem}[\texttt{schwartz\_slice\_partial\_seminorm\_bound}]\label{GaussianField-schwartz-slice-partial-seminorm-bound}
  \lean{GaussianField.schwartz_slice_partial_seminorm_bound}\leanok
  \uses{GaussianField-schwartz-slice-partial, Lattice-toSemilatticeInf}
  \texttt{GaussianField.schwartz\_slice\_partial\_seminorm\_bound}
\end{theorem}
\begin{theorem}[\texttt{integral\_euclidean\_snoc}]\label{GaussianField-integral-euclidean-snoc}
  \lean{GaussianField.integral_euclidean_snoc}\leanok
  \uses{GaussianField-euclideanSnoc}
  \texttt{GaussianField.integral\_euclidean\_snoc}
\end{theorem}
\begin{theorem}[\texttt{schwartz\_partial\_hermiteCoeff\_iteratedFDeriv}]\label{GaussianField-schwartz-partial-hermiteCoeff-iteratedFDeriv}
  \lean{GaussianField.schwartz_partial_hermiteCoeff_iteratedFDeriv}\leanok
  \uses{GaussianField-schwartz-slice-partial, GaussianField-schwartz-partial-hermiteCoeff, GaussianField-euclideanSnoc}
  \texttt{GaussianField.schwartz\_partial\_hermiteCoeff\_iteratedFDeriv}
\end{theorem}
\begin{theorem}[\texttt{euclideanSnoc\_init\_last}]\label{GaussianField-euclideanSnoc-init-last}
  \lean{GaussianField.euclideanSnoc_init_last}\leanok
  \uses{GaussianField-euclideanSnoc, GaussianField-euclideanInit}
  \texttt{GaussianField.euclideanSnoc\_init\_last}
\end{theorem}
\begin{definition}[\texttt{schwartz\_slice\_partial}]\label{GaussianField-schwartz-slice-partial}
  \lean{GaussianField.schwartz_slice_partial}\leanok
  \uses{GaussianField-euclideanSnoc}
  \texttt{GaussianField.schwartz\_slice\_partial}
\end{definition}
\begin{definition}[\texttt{schwartz\_slice}]\label{GaussianField-schwartz-slice}
  \lean{GaussianField.schwartz_slice}\leanok
  \uses{GaussianField-euclideanSnoc}
  \texttt{GaussianField.schwartz\_slice}
\end{definition}
\begin{definition}[\texttt{euclideanInit}]\label{GaussianField-euclideanInit}
  \lean{GaussianField.euclideanInit}\leanok
  \texttt{GaussianField.euclideanInit}
\end{definition}
\begin{definition}[\texttt{schwartz\_partial\_hermiteCoeff}]\label{GaussianField-schwartz-partial-hermiteCoeff}
  \lean{GaussianField.schwartz_partial_hermiteCoeff}\leanok
  \uses{GaussianField-schwartz-slice}
  \texttt{GaussianField.schwartz\_partial\_hermiteCoeff}
\end{definition}
\begin{theorem}[\texttt{schwartz\_partial\_hermiteCoeff\_eq\_1D}]\label{GaussianField-schwartz-partial-hermiteCoeff-eq-1D}
  \lean{GaussianField.schwartz_partial_hermiteCoeff_eq_1D}\leanok
  \uses{GaussianField-schwartz-slice, GaussianField-schwartz-partial-hermiteCoeff}
  \texttt{GaussianField.schwartz\_partial\_hermiteCoeff\_eq\_1D}
\end{theorem}
\begin{theorem}[\texttt{schwartz\_slice\_eq}]\label{GaussianField-schwartz-slice-eq}
  \lean{GaussianField.schwartz_slice_eq}\leanok
  \uses{GaussianField-schwartz-slice, GaussianField-euclideanSnoc}
  \texttt{GaussianField.schwartz\_slice\_eq}
\end{theorem}
\section{\texttt{GaussianField.Wick}}
\begin{theorem}[\texttt{wick\_bound}]\label{GaussianField-wick-bound}
  \lean{GaussianField.wick_bound}\leanok
  \uses{GaussianField-DyninMityaginSpace, GaussianField-measure, GaussianField-Configuration, GaussianField-instMeasurableSpaceConfiguration}
  \texttt{GaussianField.wick\_bound}
\end{theorem}
\begin{theorem}[\texttt{odd\_moment\_vanish}]\label{GaussianField-odd-moment-vanish}
  \lean{GaussianField.odd_moment_vanish}\leanok
  \uses{GaussianField-DyninMityaginSpace, GaussianField-measure, GaussianField-measure-centered, GaussianField-Configuration, GaussianField-instMeasurableSpaceConfiguration, GaussianField-wick-recursive}
  \texttt{GaussianField.odd\_moment\_vanish}
\end{theorem}
\begin{theorem}[\texttt{gaussian\_ibp\_general}]\label{GaussianField-gaussian-ibp-general}
  \lean{GaussianField.gaussian_ibp_general}\leanok
  \uses{GaussianField-DyninMityaginSpace, GaussianField-measure, GaussianField-pairing-integrable, GaussianField-gaussian-ibp, GaussianField-configuration-eval-measurable, GaussianField-Configuration, GaussianField-product-integrable, GaussianField-instMeasurableSpaceConfiguration}
  \texttt{GaussianField.gaussian\_ibp\_general}
\end{theorem}
\begin{theorem}[\texttt{wick\_recursive}]\label{GaussianField-wick-recursive}
  \lean{GaussianField.wick_recursive}\leanok
  \uses{GaussianField-DyninMityaginSpace, GaussianField-cross-moment-eq-covariance, GaussianField-measure, GaussianField-gaussian-ibp-general, GaussianField-Configuration, GaussianField-product-integrable, GaussianField-instMeasurableSpaceConfiguration, GaussianField-measure-isProbability}
  \texttt{GaussianField.wick\_recursive}
\end{theorem}
\begin{theorem}[\texttt{gaussian\_ibp}]\label{GaussianField-gaussian-ibp}
  \lean{GaussianField.gaussian_ibp}\leanok
  \uses{GaussianField-DyninMityaginSpace, GaussianField-measure, GaussianField-Configuration, GaussianField-instMeasurableSpaceConfiguration}
  \texttt{GaussianField.gaussian\_ibp}
\end{theorem}
\begin{theorem}[\texttt{wick\_bound\_factorial}]\label{GaussianField-wick-bound-factorial}
  \lean{GaussianField.wick_bound_factorial}\leanok
  \uses{GaussianField-DyninMityaginSpace, GaussianField-wick-bound, GaussianField-measure, GaussianField-Configuration, GaussianField-instMeasurableSpaceConfiguration}
  \texttt{GaussianField.wick\_bound\_factorial}
\end{theorem}
\begin{theorem}[\texttt{product\_integrable}]\label{GaussianField-product-integrable}
  \lean{GaussianField.product_integrable}\leanok
  \uses{GaussianField-DyninMityaginSpace, GaussianField-measure, GaussianField-Configuration, GaussianField-instMeasurableSpaceConfiguration}
  \texttt{GaussianField.product\_integrable}
\end{theorem}
\section{\texttt{Cylinder.MassOperatorSpatialEquivariance}}
\begin{theorem}[\texttt{cylinderMassOperator\_spatialTranslation\_norm\_eq\_proved}]\label{GaussianField-cylinderMassOperator-spatialTranslation-norm-eq-proved}
  \lean{GaussianField.cylinderMassOperator_spatialTranslation_norm_eq_proved}\leanok
  \uses{GaussianField-NuclearTensorProduct-instModuleReal, GaussianField-NuclearTensorProduct-instTopologicalSpace, GaussianField-cylinderSpatialTranslation, GaussianField-CylinderTestFunction, GaussianField-NuclearTensorProduct-instAddCommGroup, GaussianField-cylinderMassOperator-spatialTranslation-normSq-eq-proved, GaussianField-cylinderMassOperator, GaussianField-SmoothMap-Circle, GaussianField-ell2-, Lattice-toSemilatticeInf}
  \texttt{GaussianField.cylinderMassOperator\_spatialTranslation\_norm\_eq\_proved}
\end{theorem}
\begin{theorem}[\texttt{cylinderMassOperator\_spatialTranslation\_normSq\_eq\_proved}]\label{GaussianField-cylinderMassOperator-spatialTranslation-normSq-eq-proved}
  \lean{GaussianField.cylinderMassOperator_spatialTranslation_normSq_eq_proved}\leanok
  \uses{GaussianField-NuclearTensorProduct-instModuleReal, GaussianField-NuclearTensorProduct-instTopologicalSpace, GaussianField-cylinderSpatialTranslation, GaussianField-CylinderTestFunction, GaussianField-NuclearTensorProduct-instAddCommGroup, GaussianField-cylinderMassOperator, GaussianField-SmoothMap-Circle, GaussianField-ell2-}
  \texttt{GaussianField.cylinderMassOperator\_spatialTranslation\_normSq\_eq\_proved}
\end{theorem}
\section{\texttt{Pphi2.NelsonEstimate.FieldDecomposition}}
\begin{theorem}[\texttt{congr\_simp}]\label{GaussianField-latticeHeatKernelMatrix-congr-simp}
  \lean{GaussianField.latticeHeatKernelMatrix.congr_simp}\leanok
  \uses{GaussianField-FinLatticeSites, GaussianField-latticeHeatKernelMatrix}
  \texttt{GaussianField.latticeHeatKernelMatrix.congr\_simp}
\end{theorem}
\section{\texttt{Cylinder.MassOperatorRangeDense}}
\begin{theorem}[\texttt{congr\_simp}]\label{GaussianField-resolventMultiplierCLE-congr-simp}
  \lean{GaussianField.resolventMultiplierCLE.congr_simp}\leanok
  \uses{GaussianField-resolventMultiplierCLE}
  \texttt{GaussianField.resolventMultiplierCLE.congr\_simp}
\end{theorem}
\begin{theorem}[\texttt{cylinderMassOperator\_range\_dense}]\label{GaussianField-cylinderMassOperator-range-dense}
  \lean{GaussianField.cylinderMassOperator_range_dense}\leanok
  \uses{GaussianField-NuclearTensorProduct-instModuleReal, GaussianField-NuclearTensorProduct-instTopologicalSpace, GaussianField-CylinderTestFunction, GaussianField-NuclearTensorProduct-instAddCommGroup, GaussianField-cylinderMassOperator, GaussianField-SmoothMap-Circle, GaussianField-ell2-}
  \texttt{GaussianField.cylinderMassOperator\_range\_dense}
\end{theorem}
\section{\texttt{GaussianField.HypercontractiveNat}}
\begin{theorem}[\texttt{gaussian\_hypercontractive}]\label{GaussianField-gaussian-hypercontractive}
  \lean{GaussianField.gaussian_hypercontractive}\leanok
  \uses{GaussianField-DyninMityaginSpace, GaussianField-measure, GaussianField-hypercontractive-gaussianReal, GaussianField-pairing-is-gaussian, GaussianField-configuration-eval-measurable, GaussianField-Configuration, GaussianField-instMeasurableSpaceConfiguration}
  \texttt{GaussianField.gaussian\_hypercontractive}
\end{theorem}
\begin{theorem}[\texttt{gaussian\_even\_moment\_eq\_doubleFactorial}]\label{GaussianField-gaussian-even-moment-eq-doubleFactorial}
  \lean{GaussianField.gaussian_even_moment_eq_doubleFactorial}\leanok
  \uses{GaussianField-gaussian-abs-moment-std}
  \texttt{GaussianField.gaussian\_even\_moment\_eq\_doubleFactorial}
\end{theorem}
\begin{theorem}[\texttt{hypercontractive\_1d\_even}]\label{GaussianField-hypercontractive-1d-even}
  \lean{GaussianField.hypercontractive_1d_even}\leanok
  \uses{GaussianField-doubleFactorial-hypercontractive-bound, GaussianField-gaussian-even-moment-pow-eq-doubleFactorial}
  \texttt{GaussianField.hypercontractive\_1d\_even}
\end{theorem}
\begin{theorem}[\texttt{doubleFactorial\_hypercontractive\_bound}]\label{GaussianField-doubleFactorial-hypercontractive-bound}
  \lean{GaussianField.doubleFactorial_hypercontractive_bound}\leanok
  \texttt{GaussianField.doubleFactorial\_hypercontractive\_bound}
\end{theorem}
\begin{theorem}[\texttt{hypercontractive\_gaussianReal}]\label{GaussianField-hypercontractive-gaussianReal}
  \lean{GaussianField.hypercontractive_gaussianReal}\leanok
  \uses{GaussianField-hypercontractive-1d, Lattice-toSemilatticeInf}
  \texttt{GaussianField.hypercontractive\_gaussianReal}
\end{theorem}
\begin{theorem}[\texttt{hypercontractive\_1d\_p4}]\label{GaussianField-hypercontractive-1d-p4}
  \lean{GaussianField.hypercontractive_1d_p4}\leanok
  \uses{GaussianField-hypercontractive-1d-even}
  \texttt{GaussianField.hypercontractive\_1d\_p4}
\end{theorem}
\begin{theorem}[\texttt{gaussian\_even\_moment\_pow\_eq\_doubleFactorial}]\label{GaussianField-gaussian-even-moment-pow-eq-doubleFactorial}
  \lean{GaussianField.gaussian_even_moment_pow_eq_doubleFactorial}\leanok
  \uses{GaussianField-gaussian-even-moment-eq-doubleFactorial}
  \texttt{GaussianField.gaussian\_even\_moment\_pow\_eq\_doubleFactorial}
\end{theorem}
\begin{theorem}[\texttt{hypercontractive\_1d}]\label{GaussianField-hypercontractive-1d}
  \lean{GaussianField.hypercontractive_1d}\leanok
  \uses{GaussianField-hypercontractive-1d-even}
  \texttt{GaussianField.hypercontractive\_1d}
\end{theorem}
\section{\texttt{Lattice.Convergence}}
\begin{theorem}[\texttt{lattice\_green\_tendsto\_continuum}]\label{GaussianField-lattice-green-tendsto-continuum}
  \lean{GaussianField.lattice_green_tendsto_continuum}\leanok
  \uses{GaussianField-NuclearTensorProduct-instModuleReal, GaussianField-NuclearTensorProduct-instTopologicalSpace, GaussianField-DyninMityaginSpace-coeff, GaussianField-DyninMityaginSpace-coeff-decay, GaussianField-covariance, GaussianField-latticeCovariance, GaussianField-circleHasLaplacianEigenvalues, GaussianField-NuclearTensorProduct-dyninMityaginSpace, GaussianField-SmoothMap-Circle-instAddCommGroup, GaussianField-FinLatticeSites, GaussianField-NuclearTensorProduct-instIsTopologicalAddGroup, GaussianField-tensorProductHasLaplacianEigenvalues, GaussianField-DyninMityaginSpace-expansion, GaussianField-DyninMityaginSpace--, GaussianField-NuclearTensorProduct, GaussianField-NuclearTensorProduct-instAddCommGroup, GaussianField-evalTorusAtSite, GaussianField-SmoothMap-Circle-instModule, GaussianField-SmoothMap-Circle-instContinuousSMul, GaussianField-smoothCircle-dyninMityaginSpace, GaussianField-SmoothMap-Circle-instIsTopologicalAddGroup, GaussianField-DyninMityaginSpace-basis, GaussianField-NuclearTensorProduct-instContinuousSMulReal, GaussianField-NuclearTensorProduct-pure, GaussianField-SmoothMap-Circle, GaussianField-circleSpacing, GaussianField-latticeCovariance-congr-simp, GaussianField-DyninMityaginSpace-p, GaussianField-ell2-, GaussianField-FinLatticeField, GaussianField-TorusTestFunction, GaussianField-SmoothMap-Circle-instTopologicalSpace, GaussianField-greenFunctionBilinear, Lattice-toSemilatticeInf, GaussianField-circleSpacing-pos, GaussianField-greenFunctionBilinear-symm}
  \texttt{GaussianField.lattice\_green\_tendsto\_continuum}
\end{theorem}
\begin{theorem}[\texttt{latticeDFTCoeff1d\_quadratic\_bound}]\label{GaussianField-latticeDFTCoeff1d-quadratic-bound}
  \lean{GaussianField.latticeDFTCoeff1d_quadratic_bound}\leanok
  \uses{GaussianField-latticeDFTCoeff1d-quadratic-bound-, GaussianField-SmoothMap-Circle, GaussianField-latticeDFTCoeff1d}
  \texttt{GaussianField.latticeDFTCoeff1d\_quadratic\_bound}
\end{theorem}
\begin{theorem}[\texttt{lattice\_green\_tendsto\_continuum\_pure}]\label{GaussianField-lattice-green-tendsto-continuum-pure}
  \lean{GaussianField.lattice_green_tendsto_continuum_pure}\leanok
  \uses{GaussianField-NuclearTensorProduct-instModuleReal, GaussianField-NuclearTensorProduct-instTopologicalSpace, GaussianField-covariance, GaussianField-latticeCovariance, GaussianField-circleHasLaplacianEigenvalues, GaussianField-NuclearTensorProduct-dyninMityaginSpace, GaussianField-SmoothMap-Circle-instAddCommGroup, GaussianField-FinLatticeSites, GaussianField-NuclearTensorProduct-instIsTopologicalAddGroup, GaussianField-tensorProductHasLaplacianEigenvalues, GaussianField-latticeDFTCoeff1d-quadratic-bound, GaussianField-NuclearTensorProduct-instAddCommGroup, GaussianField-evalTorusAtSite, GaussianField-SmoothMap-Circle-instModule, GaussianField-SmoothMap-Circle-instContinuousSMul, GaussianField-smoothCircle-dyninMityaginSpace, GaussianField-SmoothMap-Circle-instIsTopologicalAddGroup, GaussianField-NuclearTensorProduct-instContinuousSMulReal, GaussianField-NuclearTensorProduct-pure, GaussianField-SmoothMap-Circle, GaussianField-latticeDFTCoeff1d, GaussianField-circleSpacing, GaussianField-ell2-, GaussianField-FinLatticeField, GaussianField-TorusTestFunction, GaussianField-SmoothMap-Circle-instTopologicalSpace, GaussianField-greenFunctionBilinear, GaussianField-circleSpacing-pos}
  \texttt{GaussianField.lattice\_green\_tendsto\_continuum\_pure}
\end{theorem}
\begin{theorem}[\texttt{lattice\_covariance\_pure\_eq\_2d\_spectral}]\label{GaussianField-lattice-covariance-pure-eq-2d-spectral}
  \lean{GaussianField.lattice_covariance_pure_eq_2d_spectral}\leanok
  \uses{GaussianField-NuclearTensorProduct-instModuleReal, GaussianField-latticeEigenvalue1d, GaussianField-NuclearTensorProduct-instTopologicalSpace, GaussianField-SmoothMap-Circle-instFunLike, GaussianField-latticeFourierBasisFun, GaussianField-covariance, GaussianField-latticeCovariance, GaussianField-SmoothMap-Circle-instAddCommGroup, GaussianField-FinLatticeSites, GaussianField-circleRestriction, GaussianField-NuclearTensorProduct, GaussianField-sum-factor, GaussianField-NuclearTensorProduct-instAddCommGroup, GaussianField-latticeFourierNormSq, GaussianField-evalTorusAtSite, GaussianField-SmoothMap-Circle-instModule, GaussianField-SmoothMap-Circle-instContinuousSMul, GaussianField-smoothCircle-dyninMityaginSpace, GaussianField-SmoothMap-Circle-instIsTopologicalAddGroup, GaussianField-NuclearTensorProduct-pure, GaussianField-SmoothMap-Circle, GaussianField-latticeDFTCoeff1d, GaussianField-abstract-spectral-eq-dft-spectral-2d, GaussianField-circleSpacing, GaussianField-NuclearTensorProduct-evalCLM-pure, GaussianField-ell2-, GaussianField-FinLatticeField, GaussianField-TorusTestFunction, GaussianField-NuclearTensorProduct-evalCLM, GaussianField-SmoothMap-Circle-instTopologicalSpace, GaussianField-circlePoint, GaussianField-circleSpacing-pos}
  \texttt{GaussianField.lattice\_covariance\_pure\_eq\_2d\_spectral}
\end{theorem}
\section{\texttt{Lattice.FiniteField}}
\begin{theorem}[\texttt{congr\_simp}]\label{GaussianField-finLatticeCoeffCLM-congr-simp}
  \lean{GaussianField.finLatticeCoeffCLM.congr_simp}\leanok
  \uses{GaussianField-FinLatticeSites, GaussianField-finLatticeCoeffCLM, GaussianField-FinLatticeField}
  \texttt{GaussianField.finLatticeCoeffCLM.congr\_simp}
\end{theorem}
\begin{definition}[\texttt{finLatticeField\_dyninMityaginSpace}]\label{GaussianField-finLatticeField-dyninMityaginSpace}
  \lean{GaussianField.finLatticeField_dyninMityaginSpace}\leanok
  \uses{GaussianField-DyninMityaginSpace, GaussianField-finLatticeBasisVec, GaussianField-FinLatticeSites, GaussianField-finLatticeCoeffCLM, GaussianField-finLatticeSeminorm, GaussianField-FinLatticeField}
  \texttt{GaussianField.finLatticeField\_dyninMityaginSpace}
\end{definition}
\begin{definition}[\texttt{FinLatticeField}]\label{GaussianField-FinLatticeField}
  \lean{GaussianField.FinLatticeField}\leanok
  \uses{GaussianField-FinLatticeSites}
  \texttt{GaussianField.FinLatticeField}
\end{definition}
\begin{definition}[\texttt{finLatticeCoeffCLM}]\label{GaussianField-finLatticeCoeffCLM}
  \lean{GaussianField.finLatticeCoeffCLM}\leanok
  \uses{GaussianField-FinLatticeSites, GaussianField-FinLatticeField}
  \texttt{GaussianField.finLatticeCoeffCLM}
\end{definition}
\begin{theorem}[\texttt{congr\_simp}]\label{GaussianField-finLatticeBasisVec-congr-simp}
  \lean{GaussianField.finLatticeBasisVec.congr_simp}\leanok
  \uses{GaussianField-finLatticeBasisVec, GaussianField-FinLatticeSites}
  \texttt{GaussianField.finLatticeBasisVec.congr\_simp}
\end{theorem}
\begin{theorem}[\texttt{congr\_simp}]\label{GaussianField-finLatticeSeminorm-congr-simp}
  \lean{GaussianField.finLatticeSeminorm.congr_simp}\leanok
  \uses{GaussianField-FinLatticeSites, GaussianField-finLatticeSeminorm, GaussianField-FinLatticeField}
  \texttt{GaussianField.finLatticeSeminorm.congr\_simp}
\end{theorem}
\begin{definition}[\texttt{finLatticeField\_hasPointEval}]\label{GaussianField-finLatticeField-hasPointEval}
  \lean{GaussianField.finLatticeField_hasPointEval}\leanok
  \uses{GaussianField-FinLatticeSites, GaussianField-HasPointEval, GaussianField-FinLatticeField}
  \texttt{GaussianField.finLatticeField\_hasPointEval}
\end{definition}
\begin{definition}[\texttt{finLatticeDelta}]\label{GaussianField-finLatticeDelta}
  \lean{GaussianField.finLatticeDelta}\leanok
  \uses{GaussianField-FinLatticeSites, GaussianField-FinLatticeField}
  \texttt{GaussianField.finLatticeDelta}
\end{definition}
\begin{definition}[\texttt{finLatticeBasisVec}]\label{GaussianField-finLatticeBasisVec}
  \lean{GaussianField.finLatticeBasisVec}\leanok
  \uses{GaussianField-finLatticeDelta, GaussianField-FinLatticeSites, GaussianField-FinLatticeField}
  \texttt{GaussianField.finLatticeBasisVec}
\end{definition}
\begin{definition}[\texttt{finLatticeSeminorm}]\label{GaussianField-finLatticeSeminorm}
  \lean{GaussianField.finLatticeSeminorm}\leanok
  \uses{GaussianField-FinLatticeSites, GaussianField-FinLatticeField}
  \texttt{GaussianField.finLatticeSeminorm}
\end{definition}
\section{\texttt{Lattice.Laplacian}}
\begin{theorem}[\texttt{latticeLaplacianEigenvalue\_nonneg}]\label{GaussianField-latticeLaplacianEigenvalue-nonneg}
  \lean{GaussianField.latticeLaplacianEigenvalue_nonneg}\leanok
  \uses{GaussianField-FinLatticeSites, GaussianField-latticeLaplacianEigenvalue, Lattice-toSemilatticeInf}
  \texttt{GaussianField.latticeLaplacianEigenvalue\_nonneg}
\end{theorem}
\begin{theorem}[\texttt{latticeEigenvalue\_pos}]\label{GaussianField-latticeEigenvalue-pos}
  \lean{GaussianField.latticeEigenvalue_pos}\leanok
  \uses{GaussianField-latticeEigenvalue-eq, GaussianField-latticeLaplacianEigenvalue, GaussianField-latticeLaplacianEigenvalue-nonneg, GaussianField-latticeEigenvalue}
  \texttt{GaussianField.latticeEigenvalue\_pos}
\end{theorem}
\begin{theorem}[\texttt{latticeEigenvalue\_eq}]\label{GaussianField-latticeEigenvalue-eq}
  \lean{GaussianField.latticeEigenvalue_eq}\leanok
  \uses{GaussianField-latticeLaplacianEigenvalue, GaussianField-latticeEigenvalue}
  \texttt{GaussianField.latticeEigenvalue\_eq}
\end{theorem}
\begin{theorem}[\texttt{finiteLaplacian\_neg\_semidefinite}]\label{GaussianField-finiteLaplacian-neg-semidefinite}
  \lean{GaussianField.finiteLaplacian_neg_semidefinite}\leanok
  \uses{GaussianField-FinLatticeSites, GaussianField-finiteLaplacianFun, GaussianField-finiteLaplacian, GaussianField-FinLatticeField, Lattice-toSemilatticeInf}
  \texttt{GaussianField.finiteLaplacian\_neg\_semidefinite}
\end{theorem}
\begin{definition}[\texttt{finiteLaplacianFun}]\label{GaussianField-finiteLaplacianFun}
  \lean{GaussianField.finiteLaplacianFun}\leanok
  \uses{GaussianField-FinLatticeSites, GaussianField-FinLatticeField}
  \texttt{GaussianField.finiteLaplacianFun}
\end{definition}
\begin{definition}[\texttt{massOperator}]\label{GaussianField-massOperator}
  \lean{GaussianField.massOperator}\leanok
  \uses{GaussianField-FinLatticeSites, GaussianField-finiteLaplacian, GaussianField-FinLatticeField}
  \texttt{GaussianField.massOperator}
\end{definition}
\begin{definition}[\texttt{finiteLaplacianLM}]\label{GaussianField-finiteLaplacianLM}
  \lean{GaussianField.finiteLaplacianLM}\leanok
  \uses{GaussianField-FinLatticeSites, GaussianField-finiteLaplacianFun, GaussianField-FinLatticeField}
  \texttt{GaussianField.finiteLaplacianLM}
\end{definition}
\begin{definition}[\texttt{latticeEigenvalue}]\label{GaussianField-latticeEigenvalue}
  \lean{GaussianField.latticeEigenvalue}\leanok
  \uses{GaussianField-latticeLaplacianEigenvalue}
  \texttt{GaussianField.latticeEigenvalue}
\end{definition}
\begin{theorem}[\texttt{finiteLaplacian\_selfAdjoint}]\label{GaussianField-finiteLaplacian-selfAdjoint}
  \lean{GaussianField.finiteLaplacian_selfAdjoint}\leanok
  \uses{GaussianField-FinLatticeSites, GaussianField-finiteLaplacianFun, GaussianField-finiteLaplacian, GaussianField-FinLatticeField}
  \texttt{GaussianField.finiteLaplacian\_selfAdjoint}
\end{theorem}
\begin{theorem}[\texttt{infiniteLaplacian\_apply}]\label{GaussianField-infiniteLaplacian-apply}
  \lean{GaussianField.infiniteLaplacian_apply}\leanok
  \uses{GaussianField-RapidDecayLattice-instAddCommGroup, GaussianField-stdBasisInt, GaussianField-RapidDecayLattice, GaussianField-RapidDecayLattice-val, GaussianField-RapidDecayLattice-instTopologicalSpace, GaussianField-infiniteLaplacian, GaussianField-RapidDecayLattice-instModule}
  \texttt{GaussianField.infiniteLaplacian\_apply}
\end{theorem}
\begin{theorem}[\texttt{massOperator\_pos\_def}]\label{GaussianField-massOperator-pos-def}
  \lean{GaussianField.massOperator_pos_def}\leanok
  \uses{GaussianField-FinLatticeSites, GaussianField-massOperator, GaussianField-finiteLaplacian-neg-semidefinite, GaussianField-finiteLaplacian, GaussianField-FinLatticeField, Lattice-toSemilatticeInf}
  \texttt{GaussianField.massOperator\_pos\_def}
\end{theorem}
\begin{definition}[\texttt{infiniteLaplacian}]\label{GaussianField-infiniteLaplacian}
  \lean{GaussianField.infiniteLaplacian}\leanok
  \uses{GaussianField-RapidDecayLattice-instAddCommGroup, GaussianField-RapidDecayLattice, GaussianField-RapidDecayLattice-instTopologicalSpace, GaussianField-RapidDecayLattice-instModule}
  \texttt{GaussianField.infiniteLaplacian}
\end{definition}
\begin{theorem}[\texttt{congr\_simp}]\label{GaussianField-finiteLaplacianFun-congr-simp}
  \lean{GaussianField.finiteLaplacianFun.congr_simp}\leanok
  \uses{GaussianField-FinLatticeSites, GaussianField-finiteLaplacianFun, GaussianField-FinLatticeField}
  \texttt{GaussianField.finiteLaplacianFun.congr\_simp}
\end{theorem}
\begin{theorem}[\texttt{congr\_simp}]\label{GaussianField-finiteLaplacian-congr-simp}
  \lean{GaussianField.finiteLaplacian.congr_simp}\leanok
  \uses{GaussianField-FinLatticeSites, GaussianField-finiteLaplacian, GaussianField-FinLatticeField}
  \texttt{GaussianField.finiteLaplacian.congr\_simp}
\end{theorem}
\begin{definition}[\texttt{latticeLaplacianEigenvalue}]\label{GaussianField-latticeLaplacianEigenvalue}
  \lean{GaussianField.latticeLaplacianEigenvalue}\leanok
  \uses{GaussianField-FinLatticeSites}
  \texttt{GaussianField.latticeLaplacianEigenvalue}
\end{definition}
\begin{definition}[\texttt{finiteLaplacian}]\label{GaussianField-finiteLaplacian}
  \lean{GaussianField.finiteLaplacian}\leanok
  \uses{GaussianField-FinLatticeSites, GaussianField-FinLatticeField, GaussianField-finiteLaplacianLM}
  \texttt{GaussianField.finiteLaplacian}
\end{definition}
\section{\texttt{GaussianField.NuclearFactorization}}
\begin{theorem}[\texttt{nuclear\_clm\_representation}]\label{GaussianField-nuclear-clm-representation}
  \lean{GaussianField.nuclear_clm_representation}\leanok
  \uses{GaussianField-DyninMityaginSpace, GaussianField-DyninMityaginSpace-expansion-H, GaussianField-DyninMityaginSpace-coeff, GaussianField-DyninMityaginSpace-coeff-decay, GaussianField-DyninMityaginSpace--, GaussianField-summable-one-add-rpow-neg-two, GaussianField-DyninMityaginSpace-basis, GaussianField-DyninMityaginSpace-p, GaussianField-clm-image-growth}
  \texttt{GaussianField.nuclear\_clm\_representation}
\end{theorem}
\begin{theorem}[\texttt{summable\_one\_add\_rpow\_neg\_two}]\label{GaussianField-summable-one-add-rpow-neg-two}
  \lean{GaussianField.summable_one_add_rpow_neg_two}\leanok
  \texttt{GaussianField.summable\_one\_add\_rpow\_neg\_two}
\end{theorem}
\begin{theorem}[\texttt{clm\_image\_growth}]\label{GaussianField-clm-image-growth}
  \lean{GaussianField.clm_image_growth}\leanok
  \uses{GaussianField-DyninMityaginSpace, GaussianField-DyninMityaginSpace--, GaussianField-DyninMityaginSpace-h-with, GaussianField-DyninMityaginSpace-basis, GaussianField-DyninMityaginSpace-basis-growth, GaussianField-DyninMityaginSpace-p}
  \texttt{GaussianField.clm\_image\_growth}
\end{theorem}
\section{\texttt{Cylinder.MassOperatorEquivariance}}
\begin{definition}[\texttt{perModeL2SqNorm}]\label{GaussianField-perModeL2SqNorm}
  \lean{GaussianField.perModeL2SqNorm}\leanok
  \uses{GaussianField-NuclearTensorProduct-instModuleReal, GaussianField-NuclearTensorProduct-instTopologicalSpace, GaussianField-CylinderTestFunction, GaussianField-NuclearTensorProduct-instAddCommGroup, GaussianField-resolventFreq-pos, GaussianField-resolventFreq, GaussianField-SmoothMap-Circle, GaussianField-ntpSliceSchwartz, GaussianField-resolventMultiplierCLM}
  \texttt{GaussianField.perModeL2SqNorm}
\end{definition}
\begin{theorem}[\texttt{cylinderMassOperator\_norm\_eq\_of\_temporal\_action}]\label{GaussianField-cylinderMassOperator-norm-eq-of-temporal-action}
  \lean{GaussianField.cylinderMassOperator_norm_eq_of_temporal_action}\leanok
  \uses{GaussianField-NuclearTensorProduct-instModuleReal, GaussianField-NuclearTensorProduct-instTopologicalSpace, GaussianField-nuclearTensorProduct-mapCLM, GaussianField-schwartz-dyninMityaginSpace-1D, GaussianField-SmoothMap-Circle-instAddCommGroup, GaussianField-CylinderTestFunction, GaussianField-NuclearTensorProduct, GaussianField-NuclearTensorProduct-instAddCommGroup, GaussianField-SmoothMap-Circle-instModule, GaussianField-SmoothMap-Circle-instContinuousSMul, GaussianField-smoothCircle-dyninMityaginSpace, GaussianField-SmoothMap-Circle-instIsTopologicalAddGroup, GaussianField-resolventFreq-pos, GaussianField-resolventFreq, GaussianField-cylinderMassOperator, GaussianField-SmoothMap-Circle, GaussianField-ntpSliceSchwartz, GaussianField-resolventMultiplierCLM, GaussianField-ell2-, GaussianField-perModeL2SqNorm, GaussianField-SmoothMap-Circle-instTopologicalSpace, Lattice-toSemilatticeInf, GaussianField-cylinderMassOperator-normSq-eq-sum-perMode}
  \texttt{GaussianField.cylinderMassOperator\_norm\_eq\_of\_temporal\_action}
\end{theorem}
\begin{theorem}[\texttt{congr\_simp}]\label{GaussianField-perModeL2SqNorm-congr-simp}
  \lean{GaussianField.perModeL2SqNorm.congr_simp}\leanok
  \uses{GaussianField-CylinderTestFunction, GaussianField-perModeL2SqNorm}
  \texttt{GaussianField.perModeL2SqNorm.congr\_simp}
\end{theorem}
\begin{theorem}[\texttt{congr\_simp}]\label{GaussianField-cylinderMassOperator-congr-simp}
  \lean{GaussianField.cylinderMassOperator.congr_simp}\leanok
  \uses{GaussianField-NuclearTensorProduct-instModuleReal, GaussianField-NuclearTensorProduct-instTopologicalSpace, GaussianField-CylinderTestFunction, GaussianField-NuclearTensorProduct-instAddCommGroup, GaussianField-cylinderMassOperator, GaussianField-SmoothMap-Circle, GaussianField-ell2-}
  \texttt{GaussianField.cylinderMassOperator.congr\_simp}
\end{theorem}
\begin{theorem}[\texttt{cylinderMassOperator\_timeReflection\_norm\_eq\_proved}]\label{GaussianField-cylinderMassOperator-timeReflection-norm-eq-proved}
  \lean{GaussianField.cylinderMassOperator_timeReflection_norm_eq_proved}\leanok
  \uses{GaussianField-NuclearTensorProduct-instModuleReal, GaussianField-NuclearTensorProduct-instTopologicalSpace, GaussianField-resolventSymbol-even, GaussianField-CylinderTestFunction, GaussianField-resolventSymbol-hasTemperateGrowth, GaussianField-NuclearTensorProduct-instAddCommGroup, GaussianField-cylinderTimeReflection, GaussianField-cylinderMassOperator, GaussianField-SmoothMap-Circle, GaussianField-realFourierMultiplierCLM-even-reflection-comm, GaussianField-resolventSymbol, GaussianField-schwartzReflection, GaussianField-ell2-, GaussianField-cylinderMassOperator-norm-eq-of-temporal-action}
  \texttt{GaussianField.cylinderMassOperator\_timeReflection\_norm\_eq\_proved}
\end{theorem}
\begin{theorem}[\texttt{cylinderMassOperator\_timeTranslation\_norm\_eq\_proved}]\label{GaussianField-cylinderMassOperator-timeTranslation-norm-eq-proved}
  \lean{GaussianField.cylinderMassOperator_timeTranslation_norm_eq_proved}\leanok
  \uses{GaussianField-NuclearTensorProduct-instModuleReal, GaussianField-NuclearTensorProduct-instTopologicalSpace, GaussianField-realFourierMultiplierCLM-translation-comm, GaussianField-CylinderTestFunction, GaussianField-resolventSymbol-hasTemperateGrowth, GaussianField-NuclearTensorProduct-instAddCommGroup, GaussianField-cylinderMassOperator, GaussianField-SmoothMap-Circle, GaussianField-schwartzTranslation, GaussianField-resolventSymbol, GaussianField-ell2-, GaussianField-cylinderTimeTranslation, GaussianField-cylinderMassOperator-norm-eq-of-temporal-action}
  \texttt{GaussianField.cylinderMassOperator\_timeTranslation\_norm\_eq\_proved}
\end{theorem}
\begin{theorem}[\texttt{cylinderMassOperator\_normSq\_eq\_sum\_perMode}]\label{GaussianField-cylinderMassOperator-normSq-eq-sum-perMode}
  \lean{GaussianField.cylinderMassOperator_normSq_eq_sum_perMode}\leanok
  \uses{GaussianField-NuclearTensorProduct-instModuleReal, GaussianField-NuclearTensorProduct-instTopologicalSpace, GaussianField-cylinderMassOperator-formula, GaussianField-DyninMityaginSpace-coeff, GaussianField-schwartz-dyninMityaginSpace-1D, GaussianField-CylinderTestFunction, GaussianField-NuclearTensorProduct-instAddCommGroup, GaussianField-resolventFreq-pos, GaussianField-resolventFreq, GaussianField-cylinderMassOperator, GaussianField-SmoothMap-Circle, GaussianField-ntpSliceSchwartz, GaussianField-resolventMultiplierCLM, GaussianField-ell2-, GaussianField-perModeL2SqNorm}
  \texttt{GaussianField.cylinderMassOperator\_normSq\_eq\_sum\_perMode}
\end{theorem}
\section{\texttt{Lattice.HeatKernelConvergence1d}}
\begin{theorem}[\texttt{latticeEigenvalue1d\_tendsto\_continuum}]\label{GaussianField-latticeEigenvalue1d-tendsto-continuum}
  \lean{GaussianField.latticeEigenvalue1d_tendsto_continuum}\leanok
  \uses{GaussianField-latticeEigenvalue1d, GaussianField-circleHasLaplacianEigenvalues, GaussianField-SmoothMap-Circle-instAddCommGroup, GaussianField-latticeEigenvalue1d-tendsto, GaussianField-HasLaplacianEigenvalues-eigenvalue, GaussianField-SmoothMap-Circle-instModule, GaussianField-SmoothMap-Circle-instContinuousSMul, GaussianField-smoothCircle-dyninMityaginSpace, GaussianField-SmoothMap-Circle-instIsTopologicalAddGroup, GaussianField-SmoothMap-Circle, GaussianField-circleSpacing, GaussianField-SmoothMap-Circle-instTopologicalSpace, GaussianField-SmoothMap-Circle-fourierFreq}
  \texttt{GaussianField.latticeEigenvalue1d\_tendsto\_continuum}
\end{theorem}
\begin{definition}[\texttt{latticeFourierBasisFun}]\label{GaussianField-latticeFourierBasisFun}
  \lean{GaussianField.latticeFourierBasisFun}\leanok
  \texttt{GaussianField.latticeFourierBasisFun}
\end{definition}
\begin{definition}[\texttt{latticeDFTCoeff1d}]\label{GaussianField-latticeDFTCoeff1d}
  \lean{GaussianField.latticeDFTCoeff1d}\leanok
  \uses{GaussianField-latticeFourierBasisFun, GaussianField-SmoothMap-Circle-instAddCommGroup, GaussianField-circleRestriction, GaussianField-SmoothMap-Circle-instModule, GaussianField-SmoothMap-Circle, GaussianField-SmoothMap-Circle-instTopologicalSpace}
  \texttt{GaussianField.latticeDFTCoeff1d}
\end{definition}
\begin{theorem}[\texttt{latticeHeatKernelBilinear1d\_eq\_mathlib\_spectral}]\label{GaussianField-latticeHeatKernelBilinear1d-eq-mathlib-spectral}
  \lean{GaussianField.latticeHeatKernelBilinear1d_eq_mathlib_spectral}\leanok
  \uses{GaussianField-SmoothMap-Circle-instAddCommGroup, GaussianField-FinLatticeSites, GaussianField-circleRestriction, GaussianField-SmoothMap-Circle-instModule, GaussianField-Matrix-IsHermitian-bilinear-exp-eq-spectral, GaussianField-SmoothMap-Circle, GaussianField-circleSpacing, GaussianField-latticeHeatKernelBilinear1d, GaussianField-FinLatticeField, GaussianField-SmoothMap-Circle-instTopologicalSpace, GaussianField-latticeHeatKernelMatrix, GaussianField-negLaplacianMatrix-isHermitian, GaussianField-negLaplacianMatrix}
  \texttt{GaussianField.latticeHeatKernelBilinear1d\_eq\_mathlib\_spectral}
\end{theorem}
\begin{definition}[\texttt{latticeHeatKernelBilinear1d}]\label{GaussianField-latticeHeatKernelBilinear1d}
  \lean{GaussianField.latticeHeatKernelBilinear1d}\leanok
  \uses{GaussianField-SmoothMap-Circle-instAddCommGroup, GaussianField-FinLatticeSites, GaussianField-circleRestriction, GaussianField-SmoothMap-Circle-instModule, GaussianField-SmoothMap-Circle, GaussianField-circleSpacing, GaussianField-FinLatticeField, GaussianField-SmoothMap-Circle-instTopologicalSpace, GaussianField-latticeHeatKernelMatrix}
  \texttt{GaussianField.latticeHeatKernelBilinear1d}
\end{definition}
\begin{theorem}[\texttt{latticeDFTCoeff1d\_tendsto}]\label{GaussianField-latticeDFTCoeff1d-tendsto}
  \lean{GaussianField.latticeDFTCoeff1d_tendsto}\leanok
  \uses{GaussianField-DyninMityaginSpace-coeff, GaussianField-SmoothMap-Circle-instFunLike, GaussianField-SmoothMap-Circle-continuous, GaussianField-SmoothMap-Circle-periodic, GaussianField-SmoothMap-Circle-instAddCommGroup, GaussianField-SmoothMap-Circle-fourierBasisFun-smooth, GaussianField-SmoothMap-Circle-fourierBasisFun, GaussianField-riemann-sum-periodic-tendsto, GaussianField-SmoothMap-Circle-instModule, GaussianField-SmoothMap-Circle-instContinuousSMul, GaussianField-smoothCircle-dyninMityaginSpace, GaussianField-SmoothMap-Circle-instIsTopologicalAddGroup, GaussianField-SmoothMap-Circle, GaussianField-latticeDFTCoeff1d, GaussianField-SmoothMap-Circle-fourierBasisFun-periodic, GaussianField-SmoothMap-Circle-instTopologicalSpace, GaussianField-circlePoint}
  \texttt{GaussianField.latticeDFTCoeff1d\_tendsto}
\end{theorem}
\begin{theorem}[\texttt{latticeEigenvalue1d\_zero}]\label{GaussianField-latticeEigenvalue1d-zero}
  \lean{GaussianField.latticeEigenvalue1d_zero}\leanok
  \uses{GaussianField-latticeEigenvalue1d, GaussianField-SmoothMap-Circle-fourierFreq}
  \texttt{GaussianField.latticeEigenvalue1d\_zero}
\end{theorem}
\begin{theorem}[\texttt{latticeDFTCoeff1d\_zero\_of\_ge}]\label{GaussianField-latticeDFTCoeff1d-zero-of-ge}
  \lean{GaussianField.latticeDFTCoeff1d_zero_of_ge}\leanok
  \uses{GaussianField-latticeFourierBasisFun, GaussianField-SmoothMap-Circle-instAddCommGroup, GaussianField-circleRestriction, GaussianField-SmoothMap-Circle-instModule, GaussianField-SmoothMap-Circle, GaussianField-latticeDFTCoeff1d, GaussianField-SmoothMap-Circle-instTopologicalSpace}
  \texttt{GaussianField.latticeDFTCoeff1d\_zero\_of\_ge}
\end{theorem}
\begin{theorem}[\texttt{latticeEigenvalue1d\_nonneg}]\label{GaussianField-latticeEigenvalue1d-nonneg}
  \lean{GaussianField.latticeEigenvalue1d_nonneg}\leanok
  \uses{GaussianField-latticeEigenvalue1d, GaussianField-SmoothMap-Circle-fourierFreq, Lattice-toSemilatticeInf}
  \texttt{GaussianField.latticeEigenvalue1d\_nonneg}
\end{theorem}
\begin{theorem}[\texttt{latticeDFTCoeff1d\_flat\_bound}]\label{GaussianField-latticeDFTCoeff1d-flat-bound}
  \lean{GaussianField.latticeDFTCoeff1d_flat_bound}\leanok
  \uses{GaussianField-latticeDFTCoeff1d-zero-of-ge, GaussianField-SmoothMap-Circle-instFunLike, GaussianField-latticeFourierBasisFun, GaussianField-SmoothMap-Circle-instAddCommGroup, GaussianField-SmoothMap-Circle-norm-iteratedDeriv-le-sobolevSeminorm-, GaussianField-circleRestriction, GaussianField-SmoothMap-Circle-instSMul, GaussianField-circleRestriction-apply, GaussianField-SmoothMap-Circle-instModule, GaussianField-SmoothMap-Circle, GaussianField-latticeDFTCoeff1d, GaussianField-circleSpacing, GaussianField-SmoothMap-Circle-sobolevSeminorm-nonneg, GaussianField-SmoothMap-Circle-sobolevSeminorm, GaussianField-SmoothMap-Circle-instTopologicalSpace, GaussianField-circlePoint, Lattice-toSemilatticeInf}
  \texttt{GaussianField.latticeDFTCoeff1d\_flat\_bound}
\end{theorem}
\begin{definition}[\texttt{latticeEigenvalue1d}]\label{GaussianField-latticeEigenvalue1d}
  \lean{GaussianField.latticeEigenvalue1d}\leanok
  \uses{GaussianField-SmoothMap-Circle-fourierFreq}
  \texttt{GaussianField.latticeEigenvalue1d}
\end{definition}
\begin{theorem}[\texttt{latticeEigenvalue1d\_ge\_8m}]\label{GaussianField-latticeEigenvalue1d-ge-8m}
  \lean{GaussianField.latticeEigenvalue1d_ge_8m}\leanok
  \uses{GaussianField-latticeEigenvalue1d, GaussianField-circleSpacing, GaussianField-SmoothMap-Circle-fourierFreq, Lattice-toSemilatticeInf}
  \texttt{GaussianField.latticeEigenvalue1d\_ge\_8m}
\end{theorem}
\begin{theorem}[\texttt{latticeEigenvalue1d\_tendsto}]\label{GaussianField-latticeEigenvalue1d-tendsto}
  \lean{GaussianField.latticeEigenvalue1d_tendsto}\leanok
  \texttt{GaussianField.latticeEigenvalue1d\_tendsto}
\end{theorem}
\begin{theorem}[\texttt{riemann\_sum\_periodic\_tendsto}]\label{GaussianField-riemann-sum-periodic-tendsto}
  \lean{GaussianField.riemann_sum_periodic_tendsto}\leanok
  \uses{Lattice-toSemilatticeSup, Lattice-toSemilatticeInf}
  \texttt{GaussianField.riemann\_sum\_periodic\_tendsto}
\end{theorem}
\section{\texttt{Cylinder.MethodOfImages}}
\begin{theorem}[\texttt{l2InnerProduct\_swap}]\label{GaussianField-l2InnerProduct-swap}
  \lean{GaussianField.l2InnerProduct_swap}\leanok
  \uses{GaussianField-DyninMityaginSpace, GaussianField-l2InnerProduct, GaussianField-NuclearTensorProduct-instModuleReal, GaussianField-NuclearTensorProduct-instTopologicalSpace, GaussianField-RapidDecaySeq-val, GaussianField-NuclearTensorProduct-dyninMityaginSpace, GaussianField-NuclearTensorProduct-instIsTopologicalAddGroup, GaussianField-nuclearTensorProduct-swapCLM, GaussianField-NuclearTensorProduct, GaussianField-NuclearTensorProduct-instAddCommGroup, GaussianField-NuclearTensorProduct-instContinuousSMulReal}
  \texttt{GaussianField.l2InnerProduct\_swap}
\end{theorem}
\begin{theorem}[\texttt{torusGreen\_uniform\_bound}]\label{GaussianField-torusGreen-uniform-bound}
  \lean{GaussianField.torusGreen_uniform_bound}\leanok
  \uses{GaussianField-greenFunctionBilinear-le, GaussianField-l2InnerProduct, GaussianField-NuclearTensorProduct-instModuleReal, GaussianField-NuclearTensorProduct-instTopologicalSpace, GaussianField-asymTorusContinuumGreen, GaussianField-embed-l2-uniform-bound, GaussianField-circleHasLaplacianEigenvalues, GaussianField-NuclearTensorProduct-dyninMityaginSpace, GaussianField-SmoothMap-Circle-instAddCommGroup, GaussianField-CylinderTestFunction, GaussianField-NuclearTensorProduct-instIsTopologicalAddGroup, GaussianField-tensorProductHasLaplacianEigenvalues, GaussianField-NuclearTensorProduct-instAddCommGroup, GaussianField-cylinderToTorusEmbed, GaussianField-SmoothMap-Circle-instModule, GaussianField-SmoothMap-Circle-instContinuousSMul, GaussianField-smoothCircle-dyninMityaginSpace, GaussianField-SmoothMap-Circle-instIsTopologicalAddGroup, GaussianField-NuclearTensorProduct-instContinuousSMulReal, GaussianField-SmoothMap-Circle, GaussianField-AsymTorusTestFunction, GaussianField-SmoothMap-Circle-instTopologicalSpace, GaussianField-greenFunctionBilinear}
  \texttt{GaussianField.torusGreen\_uniform\_bound}
\end{theorem}
\begin{theorem}[\texttt{l2InnerProduct\_pure}]\label{GaussianField-l2InnerProduct-pure}
  \lean{GaussianField.l2InnerProduct_pure}\leanok
  \uses{GaussianField-DyninMityaginSpace, GaussianField-l2InnerProduct, GaussianField-NuclearTensorProduct-instModuleReal, GaussianField-NuclearTensorProduct-instTopologicalSpace, GaussianField-DyninMityaginSpace-coeff, GaussianField-RapidDecaySeq-val, GaussianField-NuclearTensorProduct-dyninMityaginSpace, GaussianField-NuclearTensorProduct-instIsTopologicalAddGroup, GaussianField-NuclearTensorProduct, GaussianField-NuclearTensorProduct-instAddCommGroup, GaussianField-NuclearTensorProduct-instContinuousSMulReal, GaussianField-NuclearTensorProduct-pure, GaussianField-l2InnerProduct-summable}
  \texttt{GaussianField.l2InnerProduct\_pure}
\end{theorem}
\begin{theorem}[\texttt{embed\_l2\_uniform\_bound} (axiom)]\label{GaussianField-embed-l2-uniform-bound}
  \lean{GaussianField.embed_l2_uniform_bound}\leanok
  \uses{GaussianField-l2InnerProduct, GaussianField-NuclearTensorProduct-instModuleReal, GaussianField-NuclearTensorProduct-instTopologicalSpace, GaussianField-NuclearTensorProduct-dyninMityaginSpace, GaussianField-CylinderTestFunction, GaussianField-NuclearTensorProduct-instIsTopologicalAddGroup, GaussianField-NuclearTensorProduct-instAddCommGroup, GaussianField-cylinderToTorusEmbed, GaussianField-NuclearTensorProduct-instContinuousSMulReal, GaussianField-SmoothMap-Circle, GaussianField-AsymTorusTestFunction}
  \texttt{GaussianField.embed\_l2\_uniform\_bound}
\end{theorem}
\begin{definition}[\texttt{asymTorusContinuumGreen}]\label{GaussianField-asymTorusContinuumGreen}
  \lean{GaussianField.asymTorusContinuumGreen}\leanok
  \uses{GaussianField-NuclearTensorProduct-instModuleReal, GaussianField-NuclearTensorProduct-instTopologicalSpace, GaussianField-circleHasLaplacianEigenvalues, GaussianField-NuclearTensorProduct-dyninMityaginSpace, GaussianField-SmoothMap-Circle-instAddCommGroup, GaussianField-tensorProductHasLaplacianEigenvalues, GaussianField-NuclearTensorProduct-instAddCommGroup, GaussianField-SmoothMap-Circle-instModule, GaussianField-SmoothMap-Circle-instContinuousSMul, GaussianField-smoothCircle-dyninMityaginSpace, GaussianField-SmoothMap-Circle-instIsTopologicalAddGroup, GaussianField-SmoothMap-Circle, GaussianField-AsymTorusTestFunction, GaussianField-SmoothMap-Circle-instTopologicalSpace, GaussianField-greenFunctionBilinear}
  ``\texttt{lean noncomputable def asymTorusContinuumGreen (mass : ℝ) (hmass : 0 < mass) (f g : AsymTorusTestFunction Lt Ls) : ℝ }``  The continuum Green's function on the asymmetric torus: $G_{L_t,L_s}(f,g) = \sum_{n_1,n_2} \hat{c}_{n_1}(f_1)\, \hat{c}_{n_2}(f_2) / ((2\pi n_1/L_t)^2 + (2\pi n_2/L_s)^2 + m^2)$.
\end{definition}
\begin{theorem}[\texttt{cylinderGreen\_continuous\_seminorm\_bound}]\label{GaussianField-cylinderGreen-continuous-seminorm-bound}
  \lean{GaussianField.cylinderGreen_continuous_seminorm_bound}\leanok
  \uses{GaussianField-NuclearTensorProduct-instModuleReal, GaussianField-NuclearTensorProduct-instTopologicalSpace, GaussianField-CylinderTestFunction, GaussianField-cylinderGreen, GaussianField-NuclearTensorProduct-instAddCommGroup, GaussianField-cylinderMassOperator, GaussianField-SmoothMap-Circle, GaussianField-ell2-}
  \texttt{GaussianField.cylinderGreen\_continuous\_seminorm\_bound}
\end{theorem}
\begin{definition}[\texttt{cylinderToTorusEmbed}]\label{GaussianField-cylinderToTorusEmbed}
  \lean{GaussianField.cylinderToTorusEmbed}\leanok
  \uses{GaussianField-NuclearTensorProduct-instModuleReal, GaussianField-NuclearTensorProduct-instTopologicalSpace, GaussianField-schwartz-dyninMityaginSpace-1D, GaussianField-nuclearTensorProduct-mapCLM-general, GaussianField-SmoothMap-Circle-instAddCommGroup, GaussianField-CylinderTestFunction, GaussianField-nuclearTensorProduct-swapCLM, GaussianField-NuclearTensorProduct, GaussianField-NuclearTensorProduct-instAddCommGroup, GaussianField-periodizeCLM, GaussianField-SmoothMap-Circle-instModule, GaussianField-SmoothMap-Circle-instContinuousSMul, GaussianField-smoothCircle-dyninMityaginSpace, GaussianField-SmoothMap-Circle-instIsTopologicalAddGroup, GaussianField-SmoothMap-Circle, GaussianField-AsymTorusTestFunction, GaussianField-SmoothMap-Circle-instTopologicalSpace}
  ``\texttt{lean def cylinderToTorusEmbed : CylinderTestFunction Ls →L[ℝ] AsymTorusTestFunction Lt Ls }``  Embeds cylinder test functions into asymmetric torus test functions by applying $\mathrm{id} \otimes \mathrm{periodize}$ followed by the NTP swap map.
\end{definition}
\section{\texttt{Pphi2.OSProofs.OS2\_WardIdentity}}
\begin{theorem}[\texttt{congr\_simp}]\label{GaussianField-evalMapMeasurableEquiv-congr-simp}
  \lean{GaussianField.evalMapMeasurableEquiv.congr_simp}\leanok
  \uses{GaussianField-FinLatticeSites, GaussianField-evalMapMeasurableEquiv, GaussianField-Configuration, GaussianField-instMeasurableSpaceConfiguration, GaussianField-FinLatticeField}
  \texttt{GaussianField.evalMapMeasurableEquiv.congr\_simp}
\end{theorem}
\section{\texttt{Lattice.Symmetry}}
\begin{theorem}[\texttt{latticeTranslationMatrix\_commute\_negLaplacian}]\label{GaussianField-latticeTranslationMatrix-commute-negLaplacian}
  \lean{GaussianField.latticeTranslationMatrix_commute_negLaplacian}\leanok
  \uses{GaussianField-FinLatticeSites, GaussianField-negLaplacianMatrix, GaussianField-latticeTranslationMatrix, GaussianField-negLaplacianMatrix-toeplitz}
  \texttt{GaussianField.latticeTranslationMatrix\_commute\_negLaplacian}
\end{theorem}
\begin{theorem}[\texttt{latticeFullReflection\_commute\_heatKernel}]\label{GaussianField-latticeFullReflection-commute-heatKernel}
  \lean{GaussianField.latticeFullReflection_commute_heatKernel}\leanok
  \uses{GaussianField-latticeFullReflectionMatrix, GaussianField-FinLatticeSites, GaussianField-latticeHeatKernelMatrix-commute, GaussianField-latticeFullReflectionMatrix-commute-negLaplacian, GaussianField-latticeHeatKernelMatrix}
  \texttt{GaussianField.latticeFullReflection\_commute\_heatKernel}
\end{theorem}
\begin{definition}[\texttt{latticeFullReflectionFun}]\label{GaussianField-latticeFullReflectionFun}
  \lean{GaussianField.latticeFullReflectionFun}\leanok
  \uses{GaussianField-FinLatticeSites, GaussianField-FinLatticeField}
  \texttt{GaussianField.latticeFullReflectionFun}
\end{definition}
\begin{definition}[\texttt{latticeTranslation}]\label{GaussianField-latticeTranslation}
  \lean{GaussianField.latticeTranslation}\leanok
  \uses{GaussianField-FinLatticeSites, GaussianField-FinLatticeField, GaussianField-latticeTranslationLM}
  \texttt{GaussianField.latticeTranslation}
\end{definition}
\begin{theorem}[\texttt{latticeFullReflectionMatrix\_commute\_negLaplacian}]\label{GaussianField-latticeFullReflectionMatrix-commute-negLaplacian}
  \lean{GaussianField.latticeFullReflectionMatrix_commute_negLaplacian}\leanok
  \uses{GaussianField-latticeFullReflectionMatrix, GaussianField-FinLatticeSites, GaussianField-negLaplacianMatrix-congr-simp, GaussianField-negLaplacianMatrix-isHermitian, GaussianField-negLaplacianMatrix, GaussianField-negLaplacianMatrix-toeplitz}
  \texttt{GaussianField.latticeFullReflectionMatrix\_commute\_negLaplacian}
\end{theorem}
\begin{definition}[\texttt{latticeTranslationFun}]\label{GaussianField-latticeTranslationFun}
  \lean{GaussianField.latticeTranslationFun}\leanok
  \uses{GaussianField-FinLatticeSites, GaussianField-FinLatticeField}
  \texttt{GaussianField.latticeTranslationFun}
\end{definition}
\begin{theorem}[\texttt{negLaplacianMatrix\_toeplitz}]\label{GaussianField-negLaplacianMatrix-toeplitz}
  \lean{GaussianField.negLaplacianMatrix_toeplitz}\leanok
  \uses{GaussianField-finLatticeDelta, GaussianField-FinLatticeSites, GaussianField-massOperator, GaussianField-finiteLaplacian-translation-commute, GaussianField-latticeTranslationFun, GaussianField-finiteLaplacian, GaussianField-FinLatticeField, GaussianField-negLaplacianMatrix}
  \texttt{GaussianField.negLaplacianMatrix\_toeplitz}
\end{theorem}
\begin{definition}[\texttt{latticeTranslationLM}]\label{GaussianField-latticeTranslationLM}
  \lean{GaussianField.latticeTranslationLM}\leanok
  \uses{GaussianField-FinLatticeSites, GaussianField-latticeTranslationFun, GaussianField-FinLatticeField}
  \texttt{GaussianField.latticeTranslationLM}
\end{definition}
\begin{theorem}[\texttt{latticeTranslation\_commute\_heatKernel}]\label{GaussianField-latticeTranslation-commute-heatKernel}
  \lean{GaussianField.latticeTranslation_commute_heatKernel}\leanok
  \uses{GaussianField-FinLatticeSites, GaussianField-latticeTranslationMatrix-commute-negLaplacian, GaussianField-latticeHeatKernelMatrix-commute, GaussianField-latticeHeatKernelMatrix, GaussianField-latticeTranslationMatrix}
  \texttt{GaussianField.latticeTranslation\_commute\_heatKernel}
\end{theorem}
\begin{theorem}[\texttt{finiteLaplacian\_reflection\_commute}]\label{GaussianField-finiteLaplacian-reflection-commute}
  \lean{GaussianField.finiteLaplacian_reflection_commute}\leanok
  \uses{GaussianField-latticeFullReflectionFun, GaussianField-FinLatticeSites, GaussianField-finiteLaplacianFun, GaussianField-finiteLaplacian, GaussianField-FinLatticeField}
  \texttt{GaussianField.finiteLaplacian\_reflection\_commute}
\end{theorem}
\begin{theorem}[\texttt{negLaplacianMatrix\_neg\_invariant}]\label{GaussianField-negLaplacianMatrix-neg-invariant}
  \lean{GaussianField.negLaplacianMatrix_neg_invariant}\leanok
  \uses{GaussianField-latticeFullReflectionFun, GaussianField-finLatticeDelta, GaussianField-FinLatticeSites, GaussianField-massOperator, GaussianField-finiteLaplacian-reflection-commute, GaussianField-finiteLaplacian, GaussianField-FinLatticeField, GaussianField-negLaplacianMatrix}
  \texttt{GaussianField.negLaplacianMatrix\_neg\_invariant}
\end{theorem}
\begin{definition}[\texttt{latticeFullReflectionMatrix}]\label{GaussianField-latticeFullReflectionMatrix}
  \lean{GaussianField.latticeFullReflectionMatrix}\leanok
  \uses{GaussianField-FinLatticeSites}
  \texttt{GaussianField.latticeFullReflectionMatrix}
\end{definition}
\begin{theorem}[\texttt{finiteLaplacian\_translation\_commute}]\label{GaussianField-finiteLaplacian-translation-commute}
  \lean{GaussianField.finiteLaplacian_translation_commute}\leanok
  \uses{GaussianField-FinLatticeSites, GaussianField-latticeTranslationFun, GaussianField-finiteLaplacianFun, GaussianField-finiteLaplacian, GaussianField-FinLatticeField}
  \texttt{GaussianField.finiteLaplacian\_translation\_commute}
\end{theorem}
\begin{definition}[\texttt{latticeTranslationMatrix}]\label{GaussianField-latticeTranslationMatrix}
  \lean{GaussianField.latticeTranslationMatrix}\leanok
  \uses{GaussianField-FinLatticeSites}
  \texttt{GaussianField.latticeTranslationMatrix}
\end{definition}
\section{\texttt{Lattice.HeatKernel}}
\begin{theorem}[\texttt{mulVec\_exp\_of\_eigenvector}]\label{GaussianField-Matrix-IsHermitian-mulVec-exp-of-eigenvector}
  \lean{GaussianField.Matrix.IsHermitian.mulVec_exp_of_eigenvector}\leanok
  \uses{GaussianField-Matrix-IsHermitian-dotProduct-eigenvectors-eq-zero, GaussianField-Matrix-IsHermitian-bilinear-exp-eq-spectral}
  \texttt{GaussianField.Matrix.IsHermitian.mulVec\_exp\_of\_eigenvector}
\end{theorem}
\begin{definition}[\texttt{latticeHeatKernelMatrix}]\label{GaussianField-latticeHeatKernelMatrix}
  \lean{GaussianField.latticeHeatKernelMatrix}\leanok
  \uses{GaussianField-FinLatticeSites, GaussianField-negLaplacianMatrix}
  \texttt{GaussianField.latticeHeatKernelMatrix}
\end{definition}
\begin{definition}[\texttt{negLaplacianMatrix}]\label{GaussianField-negLaplacianMatrix}
  \lean{GaussianField.negLaplacianMatrix}\leanok
  \uses{GaussianField-massOperatorMatrix, GaussianField-FinLatticeSites}
  \texttt{GaussianField.negLaplacianMatrix}
\end{definition}
\begin{theorem}[\texttt{bilinear\_exp\_eq\_spectral}]\label{GaussianField-Matrix-IsHermitian-bilinear-exp-eq-spectral}
  \lean{GaussianField.Matrix.IsHermitian.bilinear_exp_eq_spectral}\leanok
  \uses{GaussianField-Matrix-IsHermitian-eigenCoeff-exp-neg-smul}
  \texttt{GaussianField.Matrix.IsHermitian.bilinear\_exp\_eq\_spectral}
\end{theorem}
\begin{theorem}[\texttt{negLaplacianMatrix\_isHermitian}]\label{GaussianField-negLaplacianMatrix-isHermitian}
  \lean{GaussianField.negLaplacianMatrix_isHermitian}\leanok
  \uses{GaussianField-massOperatorMatrix-isHermitian, GaussianField-FinLatticeSites, GaussianField-negLaplacianMatrix}
  \texttt{GaussianField.negLaplacianMatrix\_isHermitian}
\end{theorem}
\begin{theorem}[\texttt{mulVec\_exp\_neg\_smul}]\label{GaussianField-Matrix-IsHermitian-mulVec-exp-neg-smul}
  \lean{GaussianField.Matrix.IsHermitian.mulVec_exp_neg_smul}\leanok
  \texttt{GaussianField.Matrix.IsHermitian.mulVec\_exp\_neg\_smul}
\end{theorem}
\begin{theorem}[\texttt{latticeHeatKernelMatrix\_zero}]\label{GaussianField-latticeHeatKernelMatrix-zero}
  \lean{GaussianField.latticeHeatKernelMatrix_zero}\leanok
  \uses{GaussianField-FinLatticeSites, GaussianField-latticeHeatKernelMatrix, GaussianField-negLaplacianMatrix}
  \texttt{GaussianField.latticeHeatKernelMatrix\_zero}
\end{theorem}
\begin{theorem}[\texttt{dotProduct\_mulVec\_comm}]\label{GaussianField-Matrix-IsHermitian-dotProduct-mulVec-comm}
  \lean{GaussianField.Matrix.IsHermitian.dotProduct_mulVec_comm}\leanok
  \texttt{GaussianField.Matrix.IsHermitian.dotProduct\_mulVec\_comm}
\end{theorem}
\begin{theorem}[\texttt{latticeHeatKernelMatrix\_commute}]\label{GaussianField-latticeHeatKernelMatrix-commute}
  \lean{GaussianField.latticeHeatKernelMatrix_commute}\leanok
  \uses{GaussianField-FinLatticeSites, GaussianField-latticeHeatKernelMatrix, GaussianField-negLaplacianMatrix}
  \texttt{GaussianField.latticeHeatKernelMatrix\_commute}
\end{theorem}
\begin{theorem}[\texttt{congr\_simp}]\label{GaussianField-negLaplacianMatrix-congr-simp}
  \lean{GaussianField.negLaplacianMatrix.congr_simp}\leanok
  \uses{GaussianField-FinLatticeSites, GaussianField-negLaplacianMatrix}
  \texttt{GaussianField.negLaplacianMatrix.congr\_simp}
\end{theorem}
\begin{theorem}[\texttt{dotProduct\_eigenvectors\_eq\_zero}]\label{GaussianField-Matrix-IsHermitian-dotProduct-eigenvectors-eq-zero}
  \lean{GaussianField.Matrix.IsHermitian.dotProduct_eigenvectors_eq_zero}\leanok
  \uses{GaussianField-Matrix-IsHermitian-dotProduct-mulVec-comm, Lattice-toSemilatticeInf}
  \texttt{GaussianField.Matrix.IsHermitian.dotProduct\_eigenvectors\_eq\_zero}
\end{theorem}
\begin{theorem}[\texttt{eigenCoeff\_exp\_neg\_smul}]\label{GaussianField-Matrix-IsHermitian-eigenCoeff-exp-neg-smul}
  \lean{GaussianField.Matrix.IsHermitian.eigenCoeff_exp_neg_smul}\leanok
  \uses{GaussianField-Matrix-IsHermitian-mulVec-exp-neg-smul}
  \texttt{GaussianField.Matrix.IsHermitian.eigenCoeff\_exp\_neg\_smul}
\end{theorem}
\begin{theorem}[\texttt{latticeHeatKernelMatrix\_isHermitian}]\label{GaussianField-latticeHeatKernelMatrix-isHermitian}
  \lean{GaussianField.latticeHeatKernelMatrix_isHermitian}\leanok
  \uses{GaussianField-FinLatticeSites, GaussianField-latticeHeatKernelMatrix, GaussianField-negLaplacianMatrix-isHermitian, GaussianField-negLaplacianMatrix}
  \texttt{GaussianField.latticeHeatKernelMatrix\_isHermitian}
\end{theorem}
\begin{theorem}[\texttt{latticeHeatKernelMatrix\_semigroup}]\label{GaussianField-latticeHeatKernelMatrix-semigroup}
  \lean{GaussianField.latticeHeatKernelMatrix_semigroup}\leanok
  \uses{GaussianField-FinLatticeSites, GaussianField-latticeHeatKernelMatrix, GaussianField-negLaplacianMatrix}
  \texttt{GaussianField.latticeHeatKernelMatrix\_semigroup}
\end{theorem}
\section{\texttt{SmoothCircle.Eigenvalues}}
\begin{definition}[\texttt{circleHasLaplacianEigenvalues}]\label{GaussianField-circleHasLaplacianEigenvalues}
  \lean{GaussianField.circleHasLaplacianEigenvalues}\leanok
  \uses{GaussianField-HasLaplacianEigenvalues, GaussianField-SmoothMap-Circle-instAddCommGroup, GaussianField-SmoothMap-Circle-instModule, GaussianField-SmoothMap-Circle-instContinuousSMul, GaussianField-smoothCircle-dyninMityaginSpace, GaussianField-SmoothMap-Circle-instIsTopologicalAddGroup, GaussianField-SmoothMap-Circle, GaussianField-SmoothMap-Circle-instTopologicalSpace, GaussianField-SmoothMap-Circle-fourierFreq}
  \texttt{GaussianField.circleHasLaplacianEigenvalues}
\end{definition}
\section{\texttt{Pphi2.InteractingMeasure.LatticeMeasure}}
\begin{theorem}[\texttt{congr\_simp}]\label{GaussianField-latticeGaussianMeasure-congr-simp}
  \lean{GaussianField.latticeGaussianMeasure.congr_simp}\leanok
  \uses{GaussianField-FinLatticeSites, GaussianField-latticeGaussianMeasure, GaussianField-Configuration, GaussianField-instMeasurableSpaceConfiguration, GaussianField-FinLatticeField}
  \texttt{GaussianField.latticeGaussianMeasure.congr\_simp}
\end{theorem}
\section{\texttt{Pphi2.InteractingMeasure.LatticeEuclideanTime}}
\begin{theorem}[\texttt{congr\_simp}]\label{GaussianField-latticeTranslation-congr-simp}
  \lean{GaussianField.latticeTranslation.congr_simp}\leanok
  \uses{GaussianField-FinLatticeSites, GaussianField-latticeTranslation, GaussianField-FinLatticeField}
  \texttt{GaussianField.latticeTranslation.congr\_simp}
\end{theorem}
\section{\texttt{GaussianField.IsGaussian}}
\begin{theorem}[\texttt{measure\_isGaussian}]\label{GaussianField-measure-isGaussian}
  \lean{GaussianField.measure_isGaussian}\leanok
  \uses{GaussianField-DyninMityaginSpace, GaussianField-measure, GaussianField-second-moment-eq-covariance, GaussianField-weakDual-clm-eq-eval, GaussianField-measure-centered, GaussianField-pairing-is-gaussian, GaussianField-configuration-eval-measurable, GaussianField-Configuration, GaussianField-instMeasurableSpaceConfiguration}
  \texttt{GaussianField.measure\_isGaussian}
\end{theorem}
\begin{theorem}[\texttt{weakDual\_clm\_eq\_eval}]\label{GaussianField-weakDual-clm-eq-eval}
  \lean{GaussianField.weakDual_clm_eq_eval}\leanok
  \uses{Lattice-toSemilatticeInf}
  \texttt{GaussianField.weakDual\_clm\_eq\_eval}
\end{theorem}
\section{\texttt{Pphi2.GeneralResults.DyninMityaginBilinear}}
\begin{theorem}[\texttt{summable\_coeff\_mul\_pow}]\label{GaussianField-summable-coeff-mul-pow}
  \lean{GaussianField.summable_coeff_mul_pow}\leanok
  \uses{GaussianField-DyninMityaginSpace, GaussianField-DyninMityaginSpace-coeff, GaussianField-DyninMityaginSpace-coeff-decay, GaussianField-DyninMityaginSpace--, GaussianField-DyninMityaginSpace-p, Lattice-toSemilatticeInf}
  \texttt{GaussianField.summable\_coeff\_mul\_pow}
\end{theorem}
\begin{theorem}[\texttt{tendsto\_of\_symmetric\_basis\_tendsto}]\label{GaussianField-tendsto-of-symmetric-basis-tendsto}
  \lean{GaussianField.tendsto_of_symmetric_basis_tendsto}\leanok
  \uses{GaussianField-DyninMityaginSpace, Lattice, GaussianField-DyninMityaginSpace-coeff, GaussianField-DyninMityaginSpace-expansion, GaussianField-DyninMityaginSpace-basis, GaussianField-summable-coeff-mul-pow, Lattice-toSemilatticeInf}
  \texttt{GaussianField.tendsto\_of\_symmetric\_basis\_tendsto}
\end{theorem}
\section{\texttt{Nuclear.GeneralMapCLM}}
\begin{theorem}[\texttt{nuclearTensorProduct\_mapCLM\_general\_pure}]\label{GaussianField-nuclearTensorProduct-mapCLM-general-pure}
  \lean{GaussianField.nuclearTensorProduct_mapCLM_general_pure}\leanok
  \uses{GaussianField-DyninMityaginSpace, GaussianField-NuclearTensorProduct-ext, GaussianField-NuclearTensorProduct-instModuleReal, GaussianField-NuclearTensorProduct-instTopologicalSpace, GaussianField-DyninMityaginSpace-coeff, GaussianField-nuclearTensorProduct-mapCLM-general, GaussianField-RapidDecaySeq-val, GaussianField-DyninMityaginSpace-expansion, GaussianField-NuclearTensorProduct, GaussianField-NuclearTensorProduct-instAddCommGroup, GaussianField-DyninMityaginSpace-basis, GaussianField-NuclearTensorProduct-pure}
  \texttt{GaussianField.nuclearTensorProduct\_mapCLM\_general\_pure}
\end{theorem}
\begin{definition}[\texttt{nuclearTensorProduct\_mapCLM\_general}]\label{GaussianField-nuclearTensorProduct-mapCLM-general}
  \lean{GaussianField.nuclearTensorProduct_mapCLM_general}\leanok
  \uses{GaussianField-DyninMityaginSpace, GaussianField-NuclearTensorProduct-instModuleReal, GaussianField-NuclearTensorProduct-instTopologicalSpace, GaussianField-NuclearTensorProduct, GaussianField-NuclearTensorProduct-instAddCommGroup}
  \texttt{GaussianField.nuclearTensorProduct\_mapCLM\_general}
\end{definition}
\section{\texttt{Pphi2.GeneralResults.WeakLimitMoment}}
\begin{theorem}[\texttt{weakLimit\_exponential\_moment}]\label{GaussianField-weakLimit-exponential-moment}
  \lean{GaussianField.weakLimit_exponential_moment}\leanok
  \uses{Lattice-toSemilatticeSup, GaussianField-configuration-eval-measurable, GaussianField-Configuration, GaussianField-instMeasurableSpaceConfiguration, Lattice-toSemilatticeInf}
  \texttt{GaussianField.weakLimit\_exponential\_moment}
\end{theorem}
\section{\texttt{GaussianField.Tightness}}
\begin{theorem}[\texttt{instBaireSpace}]\label{GaussianField-DyninMityaginSpace-instBaireSpace}
  \lean{GaussianField.DyninMityaginSpace.instBaireSpace}\leanok
  \uses{GaussianField-DyninMityaginSpace, GaussianField-DyninMityaginSpace-h-countable, GaussianField-DyninMityaginSpace-toT1Space, GaussianField-DyninMityaginSpace--, GaussianField-DyninMityaginSpace-h-with, GaussianField-DyninMityaginSpace-h-completeSpace, GaussianField-DyninMityaginSpace-p}
  \texttt{GaussianField.DyninMityaginSpace.instBaireSpace}
\end{theorem}
\begin{theorem}[\texttt{configuration\_tight\_of\_uniform\_second\_moments}]\label{GaussianField-configuration-tight-of-uniform-second-moments}
  \lean{GaussianField.configuration_tight_of_uniform_second_moments}\leanok
  \uses{GaussianField-DyninMityaginSpace, GaussianField-DyninMityaginSpace--, GaussianField-Configuration, GaussianField-DyninMityaginSpace-basis, GaussianField-instMeasurableSpaceConfiguration, GaussianField-DyninMityaginSpace-p}
  \texttt{GaussianField.configuration\_tight\_of\_uniform\_second\_moments}
\end{theorem}
\section{\texttt{Pphi2.AsymTorus.AsymTorusEmbeddingIso}}
\begin{theorem}[\texttt{congr\_simp}]\label{GaussianField-latticeCovarianceAsymGJ-congr-simp}
  \lean{GaussianField.latticeCovarianceAsymGJ.congr_simp}\leanok
  \uses{GaussianField-AsymLatticeSites, GaussianField-AsymLatticeField, GaussianField-latticeCovarianceAsymGJ, GaussianField-ell2-}
  \texttt{GaussianField.latticeCovarianceAsymGJ.congr\_simp}
\end{theorem}
\section{\texttt{GaussianField.SpectralTheorem}}
\begin{theorem}[\texttt{compact\_selfAdjoint\_orthogonalComplement\_iSup\_eigenspaces\_eq\_bot}]\label{GaussianField-compact-selfAdjoint-orthogonalComplement-iSup-eigenspaces-eq-bot}
  \lean{GaussianField.compact_selfAdjoint_orthogonalComplement_iSup_eigenspaces_eq_bot}\leanok
  \texttt{GaussianField.compact\_selfAdjoint\_orthogonalComplement\_iSup\_eigenspaces\_eq\_bot}
\end{theorem}
\begin{theorem}[\texttt{compact\_selfAdjoint\_hasEigenvector}]\label{GaussianField-compact-selfAdjoint-hasEigenvector}
  \lean{GaussianField.compact_selfAdjoint_hasEigenvector}\leanok
  \texttt{GaussianField.compact\_selfAdjoint\_hasEigenvector}
\end{theorem}
\begin{theorem}[\texttt{compact\_selfAdjoint\_spectral}]\label{GaussianField-compact-selfAdjoint-spectral}
  \lean{GaussianField.compact_selfAdjoint_spectral}\leanok
  \texttt{GaussianField.compact\_selfAdjoint\_spectral}
\end{theorem}
\section{\texttt{Cylinder.Basic}}
\begin{definition}[\texttt{CylinderTestFunction}]\label{GaussianField-CylinderTestFunction}
  \lean{GaussianField.CylinderTestFunction}\leanok
  \uses{GaussianField-NuclearTensorProduct, GaussianField-SmoothMap-Circle}
  \texttt{GaussianField.CylinderTestFunction}
\end{definition}
\section{\texttt{Pphi2.ContinuumLimit.Hypercontractivity}}
\begin{theorem}[\texttt{congr\_simp}]\label{GaussianField-spectralLatticeCovariance-congr-simp}
  \lean{GaussianField.spectralLatticeCovariance.congr_simp}\leanok
  \uses{GaussianField-FinLatticeSites, GaussianField-spectralLatticeCovariance, GaussianField-ell2-, GaussianField-FinLatticeField}
  \texttt{GaussianField.spectralLatticeCovariance.congr\_simp}
\end{theorem}
\section{\texttt{Torus.Restriction}}
\begin{definition}[\texttt{TorusTestFunction}]\label{GaussianField-TorusTestFunction}
  \lean{GaussianField.TorusTestFunction}\leanok
  \uses{GaussianField-NuclearTensorProduct, GaussianField-SmoothMap-Circle}
  \texttt{GaussianField.TorusTestFunction}
\end{definition}
\section{\texttt{GaussianField.CramerWold}}
\begin{theorem}[\texttt{cramerWold}]\label{GaussianField-cramerWold}
  \lean{GaussianField.cramerWold}\leanok
  \texttt{GaussianField.cramerWold}
\end{theorem}
\section{\texttt{Pphi2.GeneralResults.GaussianExpMoment}}
\begin{theorem}[\texttt{gaussian\_exp\_abs\_moment}]\label{GaussianField-gaussian-exp-abs-moment}
  \lean{GaussianField.gaussian_exp_abs_moment}\leanok
  \uses{GaussianField-DyninMityaginSpace, Lattice-toSemilatticeSup, GaussianField-measure, GaussianField-second-moment-eq-covariance, GaussianField-pairing-is-gaussian, GaussianField-configuration-eval-measurable, GaussianField-Configuration, GaussianField-instMeasurableSpaceConfiguration, Lattice-toSemilatticeInf}
  For the interacting measure, $\int e^{c|\omega(f)|}\, d\mu_P \le$ explicit Gaussian bound.
\end{theorem}
\chapter{pphi2: P(Φ)₂ construction}
\section{\texttt{Pphi2.NelsonEstimate.FieldDecomposition}}
\begin{theorem}[\texttt{canonicalRoughFieldFunction\_pointwise\_measurable}]\label{Pphi2-canonicalRoughFieldFunction-pointwise-measurable}
  \lean{Pphi2.canonicalRoughFieldFunction_pointwise_measurable}\leanok
  \uses{GaussianField-FinLatticeSites, Pphi2-canonicalRoughFieldFunction, Pphi2-CanonicalJoint}
  \texttt{Pphi2.canonicalRoughFieldFunction\_pointwise\_measurable}
\end{theorem}
\begin{theorem}[\texttt{canonicalSumFieldLaw\_eq\_map\_canonicalSumFieldFunction}]\label{Pphi2-canonicalSumFieldLaw-eq-map-canonicalSumFieldFunction}
  \lean{Pphi2.canonicalSumFieldLaw_eq_map_canonicalSumFieldFunction}\leanok
  \uses{Pphi2-canonicalJointMeasure-map-stdGaussianEuclidean, Pphi2-canonicalJointMeasure, GaussianField-FinLatticeSites, Pphi2-canonicalStdToFieldCLM, Pphi2-canonicalStdFieldFunction-eq-canonicalSumFieldFunction, Pphi2-canonicalSumFieldFunction, Pphi2-CanonicalJointSumIndex, Pphi2-canonicalSumFieldLaw, GaussianField-FinLatticeField, Pphi2-CanonicalJoint}
  \texttt{Pphi2.canonicalSumFieldLaw\_eq\_map\_canonicalSumFieldFunction}
\end{theorem}
\begin{theorem}[\texttt{canonicalRoughFieldFunction\_covariance\_eq}]\label{Pphi2-canonicalRoughFieldFunction-covariance-eq}
  \lean{Pphi2.canonicalRoughFieldFunction_covariance_eq}\leanok
  \uses{Pphi2-canonicalJointMeasure, Pphi2-canonicalRoughCovariance, GaussianField-FinLatticeSites, Pphi2-canonicalRoughFieldFunction, Pphi2-CanonicalJoint}
  \texttt{Pphi2.canonicalRoughFieldFunction\_covariance\_eq}
\end{theorem}
\begin{theorem}[\texttt{wickPower\_two\_site\_pi\_gaussianReal\_eq\_diag}]\label{Pphi2-wickPower-two-site-pi-gaussianReal-eq-diag}
  \lean{Pphi2.wickPower_two_site_pi_gaussianReal_eq_diag}\leanok
  \uses{Pphi2-wickMonomial, GaussianField-multiWickMonomial-pi-gaussianReal-inner}
  \texttt{Pphi2.wickPower\_two\_site\_pi\_gaussianReal\_eq\_diag}
\end{theorem}
\begin{theorem}[\texttt{wickPower\_two\_site\_pi\_gaussianReal\_eq\_zero\_of\_ne}]\label{Pphi2-wickPower-two-site-pi-gaussianReal-eq-zero-of-ne}
  \lean{Pphi2.wickPower_two_site_pi_gaussianReal_eq_zero_of_ne}\leanok
  \uses{Pphi2-wickMonomial, GaussianField-multiWickMonomial-pi-gaussianReal-inner}
  \texttt{Pphi2.wickPower\_two\_site\_pi\_gaussianReal\_eq\_zero\_of\_ne}
\end{theorem}
\begin{definition}[\texttt{canonicalSmoothGamma}]\label{Pphi2-canonicalSmoothGamma}
  \lean{Pphi2.canonicalSmoothGamma}\leanok
  \uses{GaussianField-FinLatticeSites}
  \texttt{Pphi2.canonicalSmoothGamma}
\end{definition}
\begin{theorem}[\texttt{canonicalSumConfig\_apply\_delta}]\label{Pphi2-canonicalSumConfig-apply-delta}
  \lean{Pphi2.canonicalSumConfig_apply_delta}\leanok
  \uses{Pphi2-canonicalSumConfig, GaussianField-finLatticeDelta, GaussianField-FinLatticeSites, GaussianField-Configuration, Pphi2-canonicalSumFieldFunction, Pphi2-latticeFieldToConfig, GaussianField-FinLatticeField, Pphi2-latticeFieldToConfig-apply, Pphi2-CanonicalJoint}
  \texttt{Pphi2.canonicalSumConfig\_apply\_delta}
\end{theorem}
\begin{theorem}[\texttt{congr\_simp}]\label{Pphi2-canonicalSmoothWeight-congr-simp}
  \lean{Pphi2.canonicalSmoothWeight.congr_simp}\leanok
  \uses{Pphi2-canonicalSmoothWeight}
  \texttt{Pphi2.canonicalSmoothWeight.congr\_simp}
\end{theorem}
\begin{theorem}[\texttt{canonicalRoughWickPower\_two\_site\_marginal\_eq\_diag}]\label{Pphi2-canonicalRoughWickPower-two-site-marginal-eq-diag}
  \lean{Pphi2.canonicalRoughWickPower_two_site_marginal_eq_diag}\leanok
  \uses{Pphi2-canonicalRoughFieldFunctionOfSnd-eq-sum-gamma, Pphi2-canonicalRoughFieldFunctionOfSnd, Pphi2-canonicalRoughCovariance, Pphi2-canonicalRoughCovariance-eq-sum-gamma-mul-gamma, GaussianField-FinLatticeSites, Pphi2-wickPower-two-site-pi-gaussianReal-eq-diag, Pphi2-wickMonomial, Pphi2-roughWickConstant, Pphi2-canonicalRoughGamma}
  \texttt{Pphi2.canonicalRoughWickPower\_two\_site\_marginal\_eq\_diag}
\end{theorem}
\begin{theorem}[\texttt{canonicalRoughFieldFunction\_self\_moment\_eq\_roughWickConstant}]\label{Pphi2-canonicalRoughFieldFunction-self-moment-eq-roughWickConstant}
  \lean{Pphi2.canonicalRoughFieldFunction_self_moment_eq_roughWickConstant}\leanok
  \uses{Pphi2-canonicalJointMeasure, GaussianField-FinLatticeSites, Pphi2-canonicalRoughFieldFunction-self-moment-diag, Pphi2-canonicalRoughWeight, Pphi2-canonicalRoughFieldFunction, GaussianField-latticeFourierProductBasisFun, Pphi2-roughWickConstant, GaussianField-latticeFourierProductNormSq, Pphi2-CanonicalJoint}
  \texttt{Pphi2.canonicalRoughFieldFunction\_self\_moment\_eq\_roughWickConstant}
\end{theorem}
\begin{theorem}[\texttt{canonicalRoughFieldFunction\_add}]\label{Pphi2-canonicalRoughFieldFunction-add}
  \lean{Pphi2.canonicalRoughFieldFunction_add}\leanok
  \uses{GaussianField-FinLatticeSites, Pphi2-canonicalRoughFieldFunction, Pphi2-CanonicalJoint}
  \texttt{Pphi2.canonicalRoughFieldFunction\_add}
\end{theorem}
\begin{theorem}[\texttt{canonicalJointMeasure\_isProbability}]\label{Pphi2-canonicalJointMeasure-isProbability}
  \lean{Pphi2.canonicalJointMeasure_isProbability}\leanok
  \uses{Pphi2-canonicalJointMeasure, Pphi2-CanonicalJoint}
  \texttt{Pphi2.canonicalJointMeasure\_isProbability}
\end{theorem}
\begin{theorem}[\texttt{canonicalSumFieldFunction\_covariance\_eq\_GJ}]\label{Pphi2-canonicalSumFieldFunction-covariance-eq-GJ}
  \lean{Pphi2.canonicalSumFieldFunction_covariance_eq_GJ}\leanok
  \uses{Pphi2-canonicalJointMeasure, GaussianField-latticeEigenvalue1d, GaussianField-latticeCovariance-GJ-eq-inv-smul-bare, GaussianField-covariance, GaussianField-latticeCovariance, GaussianField-abstract-spectral-eq-dft-spectral-family, GaussianField-FinLatticeSites, GaussianField-latticeCovarianceGJ, Pphi2-canonicalSumFieldFunction-covariance, Pphi2-canonicalEigenvalue, Pphi2-canonicalSumFieldFunction, GaussianField-latticeFourierProductBasisFun, GaussianField-latticeFourierProductCoeff, GaussianField-ell2-, GaussianField-FinLatticeField, GaussianField-latticeFourierProductNormSq, Pphi2-CanonicalJoint}
  \texttt{Pphi2.canonicalSumFieldFunction\_covariance\_eq\_GJ}
\end{theorem}
\begin{theorem}[\texttt{latticeFieldToConfig\_apply}]\label{Pphi2-latticeFieldToConfig-apply}
  \lean{Pphi2.latticeFieldToConfig_apply}\leanok
  \uses{GaussianField-FinLatticeSites, GaussianField-Configuration, Pphi2-latticeFieldToConfig, GaussianField-FinLatticeField}
  \texttt{Pphi2.latticeFieldToConfig\_apply}
\end{theorem}
\begin{definition}[\texttt{Joint}]\label{Pphi2-FieldDecomposition-Joint}
  \lean{Pphi2.FieldDecomposition.Joint}\leanok
  \uses{Pphi2-FieldDecomposition}
  \texttt{Pphi2.FieldDecomposition.Joint}
\end{definition}
\begin{theorem}[\texttt{abs\_canonicalSmoothCovariance\_le\_smoothWickConstant}]\label{Pphi2-abs-canonicalSmoothCovariance-le-smoothWickConstant}
  \lean{Pphi2.abs_canonicalSmoothCovariance_le_smoothWickConstant}\leanok
  \uses{Pphi2-smoothWickConstant, Pphi2-canonicalSmoothCovariance, GaussianField-FinLatticeSites, Pphi2-canonicalSmoothGamma, Lattice-toSemilatticeInf}
  \texttt{Pphi2.abs\_canonicalSmoothCovariance\_le\_smoothWickConstant}
\end{theorem}
\begin{theorem}[\texttt{canonicalRoughFieldFunction\_eq\_of\_snd}]\label{Pphi2-canonicalRoughFieldFunction-eq-of-snd}
  \lean{Pphi2.canonicalRoughFieldFunction_eq_of_snd}\leanok
  \uses{Pphi2-canonicalRoughFieldFunctionOfSnd, GaussianField-FinLatticeSites, Pphi2-canonicalRoughFieldFunction}
  \texttt{Pphi2.canonicalRoughFieldFunction\_eq\_of\_snd}
\end{theorem}
\begin{theorem}[\texttt{pushforward\_eq\_GFF}]\label{Pphi2-FieldDecomposition-pushforward-eq-GFF}
  \lean{Pphi2.FieldDecomposition.pushforward_eq_GFF}\leanok
  \uses{GaussianField-FinLatticeSites, Pphi2-FieldDecomposition---joint, Pphi2-FieldDecomposition---S, Pphi2-FieldDecomposition, Pphi2-FieldDecomposition---R, Pphi2-FieldDecomposition-Joint, GaussianField-latticeGaussianMeasure, GaussianField-Configuration, GaussianField-instMeasurableSpaceConfiguration, Pphi2-FieldDecomposition-jointMeasurable, GaussianField-FinLatticeField}
  \texttt{Pphi2.FieldDecomposition.pushforward\_eq\_GFF}
\end{theorem}
\begin{theorem}[\texttt{canonicalRoughFieldFunction\_memLp\_two}]\label{Pphi2-canonicalRoughFieldFunction-memLp-two}
  \lean{Pphi2.canonicalRoughFieldFunction_memLp_two}\leanok
  \uses{Pphi2-canonicalJointMeasure, GaussianField-FinLatticeSites, Pphi2-canonicalRoughFieldFunction, Pphi2-CanonicalJoint}
  \texttt{Pphi2.canonicalRoughFieldFunction\_memLp\_two}
\end{theorem}
\begin{theorem}[\texttt{canonicalSumFieldFunction\_smul}]\label{Pphi2-canonicalSumFieldFunction-smul}
  \lean{Pphi2.canonicalSumFieldFunction_smul}\leanok
  \uses{GaussianField-FinLatticeSites, Pphi2-canonicalRoughFieldFunction-smul, Pphi2-canonicalSumFieldFunction, Pphi2-canonicalRoughFieldFunction, Pphi2-canonicalSumFieldFunction-eq-smooth-plus-rough, Pphi2-canonicalSmoothFieldFunction, Pphi2-canonicalSmoothFieldFunction-smul, Pphi2-CanonicalJoint}
  \texttt{Pphi2.canonicalSumFieldFunction\_smul}
\end{theorem}
\begin{theorem}[\texttt{heatKernel\_row\_pairing\_eq\_diagonal}]\label{Pphi2-heatKernel-row-pairing-eq-diagonal}
  \lean{Pphi2.heatKernel_row_pairing_eq_diagonal}\leanok
  \uses{GaussianField-latticeHeatKernelMatrix-isHermitian, GaussianField-FinLatticeSites, GaussianField-latticeHeatKernelMatrix-semigroup, GaussianField-latticeHeatKernelMatrix}
  \texttt{Pphi2.heatKernel\_row\_pairing\_eq\_diagonal}
\end{theorem}
\begin{theorem}[\texttt{canonicalSmoothFieldFunction\_memLp\_two}]\label{Pphi2-canonicalSmoothFieldFunction-memLp-two}
  \lean{Pphi2.canonicalSmoothFieldFunction_memLp_two}\leanok
  \uses{Pphi2-canonicalJointMeasure, GaussianField-FinLatticeSites, Pphi2-canonicalSmoothFieldFunction, Pphi2-CanonicalJoint}
  \texttt{Pphi2.canonicalSmoothFieldFunction\_memLp\_two}
\end{theorem}
\begin{theorem}[\texttt{canonicalSmoothFieldFunction\_covariance\_eq}]\label{Pphi2-canonicalSmoothFieldFunction-covariance-eq}
  \lean{Pphi2.canonicalSmoothFieldFunction_covariance_eq}\leanok
  \uses{Pphi2-canonicalJointMeasure, Pphi2-canonicalSmoothCovariance, GaussianField-FinLatticeSites, Pphi2-canonicalSmoothFieldFunction, Pphi2-CanonicalJoint}
  \texttt{Pphi2.canonicalSmoothFieldFunction\_covariance\_eq}
\end{theorem}
\begin{theorem}[\texttt{canonicalRoughFieldFunction\_measurable}]\label{Pphi2-canonicalRoughFieldFunction-measurable}
  \lean{Pphi2.canonicalRoughFieldFunction_measurable}\leanok
  \uses{GaussianField-FinLatticeSites, Pphi2-canonicalRoughFieldFunction-pointwise-measurable, Pphi2-canonicalRoughFieldFunction, Pphi2-CanonicalJoint}
  \texttt{Pphi2.canonicalRoughFieldFunction\_measurable}
\end{theorem}
\begin{theorem}[\texttt{wickPower\_two\_site\_pi\_gaussianReal\_integrable}]\label{Pphi2-wickPower-two-site-pi-gaussianReal-integrable}
  \lean{Pphi2.wickPower_two_site_pi_gaussianReal_integrable}\leanok
  \uses{Pphi2-wickMonomial}
  \texttt{Pphi2.wickPower\_two\_site\_pi\_gaussianReal\_integrable}
\end{theorem}
\begin{definition}[\texttt{canonicalRoughWeight}]\label{Pphi2-canonicalRoughWeight}
  \lean{Pphi2.canonicalRoughWeight}\leanok
  \uses{Pphi2-canonicalEigenvalue}
  \texttt{Pphi2.canonicalRoughWeight}
\end{definition}
\begin{theorem}[\texttt{canonicalRoughFieldFunctionOfSnd\_eq\_sum\_gamma}]\label{Pphi2-canonicalRoughFieldFunctionOfSnd-eq-sum-gamma}
  \lean{Pphi2.canonicalRoughFieldFunctionOfSnd_eq_sum_gamma}\leanok
  \uses{Pphi2-canonicalRoughFieldFunctionOfSnd, GaussianField-FinLatticeSites, Pphi2-canonicalRoughGamma}
  \texttt{Pphi2.canonicalRoughFieldFunctionOfSnd\_eq\_sum\_gamma}
\end{theorem}
\begin{theorem}[\texttt{canonicalJointMeasure\_map\_sum\_pi\_gaussian}]\label{Pphi2-canonicalJointMeasure-map-sum-pi-gaussian}
  \lean{Pphi2.canonicalJointMeasure_map_sum_pi_gaussian}\leanok
  \uses{Pphi2-canonicalJointMeasure, Pphi2-CanonicalJointSumIndex}
  \texttt{Pphi2.canonicalJointMeasure\_map\_sum\_pi\_gaussian}
\end{theorem}
\begin{definition}[\texttt{CanonicalJointSumIndex}]\label{Pphi2-CanonicalJointSumIndex}
  \lean{Pphi2.CanonicalJointSumIndex}\leanok
  \texttt{Pphi2.CanonicalJointSumIndex}
\end{definition}
\begin{theorem}[\texttt{congr\_simp}]\label{Pphi2-canonicalSmoothFieldFunctionOfFst-congr-simp}
  \lean{Pphi2.canonicalSmoothFieldFunctionOfFst.congr_simp}\leanok
  \uses{GaussianField-FinLatticeSites, Pphi2-canonicalSmoothFieldFunctionOfFst}
  \texttt{Pphi2.canonicalSmoothFieldFunctionOfFst.congr\_simp}
\end{theorem}
\begin{definition}[\texttt{canonicalRoughCovariance}]\label{Pphi2-canonicalRoughCovariance}
  \lean{Pphi2.canonicalRoughCovariance}\leanok
  \uses{GaussianField-FinLatticeSites, Pphi2-canonicalRoughWeight, GaussianField-latticeFourierProductBasisFun, GaussianField-latticeFourierProductNormSq}
  \texttt{Pphi2.canonicalRoughCovariance}
\end{definition}
\begin{definition}[\texttt{FieldDecomposition}]\label{Pphi2-FieldDecomposition}
  \lean{Pphi2.FieldDecomposition}\leanok
  \texttt{Pphi2.FieldDecomposition}
\end{definition}
\begin{theorem}[\texttt{canonicalSmoothConfig\_add\_canonicalRoughConfig\_eq\_lift\_sum}]\label{Pphi2-canonicalSmoothConfig-add-canonicalRoughConfig-eq-lift-sum}
  \lean{Pphi2.canonicalSmoothConfig_add_canonicalRoughConfig_eq_lift_sum}\leanok
  \uses{Pphi2-canonicalSmoothConfig, GaussianField-FinLatticeSites, Pphi2-canonicalRoughConfig, GaussianField-Configuration, Pphi2-canonicalSumFieldFunction, Pphi2-canonicalRoughFieldFunction, Pphi2-latticeFieldToConfig, Pphi2-canonicalSmoothFieldFunction, GaussianField-FinLatticeField, Pphi2-CanonicalJoint}
  \texttt{Pphi2.canonicalSmoothConfig\_add\_canonicalRoughConfig\_eq\_lift\_sum}
\end{theorem}
\begin{definition}[\texttt{canonicalJointStdGaussianMeasurableEquiv}]\label{Pphi2-canonicalJointStdGaussianMeasurableEquiv}
  \lean{Pphi2.canonicalJointStdGaussianMeasurableEquiv}\leanok
  \uses{Pphi2-CanonicalJointSumIndex, Pphi2-CanonicalJoint}
  \texttt{Pphi2.canonicalJointStdGaussianMeasurableEquiv}
\end{definition}
\begin{theorem}[\texttt{canonicalJointMeasure\_map\_stdGaussian}]\label{Pphi2-canonicalJointMeasure-map-stdGaussian}
  \lean{Pphi2.canonicalJointMeasure_map_stdGaussian}\leanok
  \uses{Pphi2-canonicalJointMeasure, Pphi2-canonicalJointStdGaussianMeasurableEquiv, Pphi2-canonicalJointMeasure-map-sum-pi-gaussian, GaussianHilbert-stdGaussianFin, Pphi2-CanonicalJointSumIndex, Pphi2-CanonicalJoint}
  \texttt{Pphi2.canonicalJointMeasure\_map\_stdGaussian}
\end{theorem}
\begin{theorem}[\texttt{latticeHeatKernel\_entry\_eq\_eigenvalue\_average}]\label{Pphi2-latticeHeatKernel-entry-eq-eigenvalue-average}
  \lean{Pphi2.latticeHeatKernel_entry_eq_eigenvalue_average}\leanok
  \uses{GaussianField-latticeEigenvalue1d, GaussianField-FinLatticeSites, GaussianField-latticeFourierProductBasisFun, GaussianField-latticeFourierProductNormSq, GaussianField-latticeHeatKernelMatrix}
  \texttt{Pphi2.latticeHeatKernel\_entry\_eq\_eigenvalue\_average}
\end{theorem}
\begin{theorem}[\texttt{canonicalRoughCovariance\_eq\_integral\_Icc\_heatKernel}]\label{Pphi2-canonicalRoughCovariance-eq-integral-Icc-heatKernel}
  \lean{Pphi2.canonicalRoughCovariance_eq_integral_Icc_heatKernel}\leanok
  \uses{GaussianField-latticeEigenvalue1d, Pphi2-canonicalRoughCovariance, GaussianField-FinLatticeSites, Pphi2-canonicalRoughWeight, GaussianField-latticeFourierProductBasisFun, GaussianField-latticeFourierProductNormSq, GaussianField-latticeHeatKernelMatrix}
  \texttt{Pphi2.canonicalRoughCovariance\_eq\_integral\_Icc\_heatKernel}
\end{theorem}
\begin{theorem}[\texttt{canonicalSmoothFieldFunctionOfFst\_eq\_sum\_gamma}]\label{Pphi2-canonicalSmoothFieldFunctionOfFst-eq-sum-gamma}
  \lean{Pphi2.canonicalSmoothFieldFunctionOfFst_eq_sum_gamma}\leanok
  \uses{GaussianField-FinLatticeSites, Pphi2-canonicalSmoothGamma, Pphi2-canonicalSmoothFieldFunctionOfFst}
  \texttt{Pphi2.canonicalSmoothFieldFunctionOfFst\_eq\_sum\_gamma}
\end{theorem}
\begin{definition}[\texttt{canonicalSmoothConfig}]\label{Pphi2-canonicalSmoothConfig}
  \lean{Pphi2.canonicalSmoothConfig}\leanok
  \uses{GaussianField-FinLatticeSites, GaussianField-Configuration, Pphi2-latticeFieldToConfig, Pphi2-canonicalSmoothFieldFunction, GaussianField-FinLatticeField, Pphi2-CanonicalJoint}
  \texttt{Pphi2.canonicalSmoothConfig}
\end{definition}
\begin{theorem}[\texttt{heatKernel\_row\_sum\_eq\_one}]\label{Pphi2-heatKernel-row-sum-eq-one}
  \lean{Pphi2.heatKernel_row_sum_eq_one}\leanok
  \uses{GaussianField-latticeEigenvalue1d, GaussianField-massOperatorMatrix, GaussianField-massOperator-product-eigenvector-family, GaussianField-FinLatticeSites, GaussianField-massOperator, GaussianField-Matrix-IsHermitian-mulVec-exp-of-eigenvector, GaussianField-massOperator-eq-matrix-mulVec, GaussianField-latticeEigenvalue1d-zero, GaussianField-latticeFourierProductBasisFun, GaussianField-FinLatticeField, GaussianField-latticeHeatKernelMatrix, GaussianField-negLaplacianMatrix-isHermitian, GaussianField-negLaplacianMatrix}
  \texttt{Pphi2.heatKernel\_row\_sum\_eq\_one}
\end{theorem}
\begin{theorem}[\texttt{heatKernel\_diagonal\_mass\_weighted\_eq\_eigenvalue\_average}]\label{Pphi2-heatKernel-diagonal-mass-weighted-eq-eigenvalue-average}
  \lean{Pphi2.heatKernel_diagonal_mass_weighted_eq_eigenvalue_average}\leanok
  \uses{GaussianField-latticeEigenvalue1d, GaussianField-FinLatticeSites, Pphi2-canonicalEigenvalue, GaussianField-latticeHeatKernelMatrix}
  \texttt{Pphi2.heatKernel\_diagonal\_mass\_weighted\_eq\_eigenvalue\_average}
\end{theorem}
\begin{theorem}[\texttt{canonicalSmoothFieldFunction\_add}]\label{Pphi2-canonicalSmoothFieldFunction-add}
  \lean{Pphi2.canonicalSmoothFieldFunction_add}\leanok
  \uses{GaussianField-FinLatticeSites, Pphi2-canonicalSmoothFieldFunction, Pphi2-CanonicalJoint}
  \texttt{Pphi2.canonicalSmoothFieldFunction\_add}
\end{theorem}
\begin{definition}[\texttt{canonicalSmoothCovariance}]\label{Pphi2-canonicalSmoothCovariance}
  \lean{Pphi2.canonicalSmoothCovariance}\leanok
  \uses{Pphi2-canonicalSmoothWeight, GaussianField-FinLatticeSites, GaussianField-latticeFourierProductBasisFun, GaussianField-latticeFourierProductNormSq}
  \texttt{Pphi2.canonicalSmoothCovariance}
\end{definition}
\begin{definition}[\texttt{canonicalRoughGamma}]\label{Pphi2-canonicalRoughGamma}
  \lean{Pphi2.canonicalRoughGamma}\leanok
  \uses{GaussianField-FinLatticeSites}
  \texttt{Pphi2.canonicalRoughGamma}
\end{definition}
\begin{theorem}[\texttt{canonicalSumFieldFunction\_memLp\_two}]\label{Pphi2-canonicalSumFieldFunction-memLp-two}
  \lean{Pphi2.canonicalSumFieldFunction_memLp_two}\leanok
  \uses{Pphi2-canonicalJointMeasure, Pphi2-canonicalRoughFieldFunction-memLp-two, GaussianField-FinLatticeSites, Pphi2-canonicalSmoothFieldFunction-memLp-two, Pphi2-canonicalSumFieldFunction, Pphi2-canonicalRoughFieldFunction, Pphi2-canonicalSmoothFieldFunction, Pphi2-CanonicalJoint}
  \texttt{Pphi2.canonicalSumFieldFunction\_memLp\_two}
\end{theorem}
\begin{theorem}[\texttt{canonicalSmoothFieldFunction\_eq\_of\_fst}]\label{Pphi2-canonicalSmoothFieldFunction-eq-of-fst}
  \lean{Pphi2.canonicalSmoothFieldFunction_eq_of_fst}\leanok
  \uses{GaussianField-FinLatticeSites, Pphi2-canonicalSmoothFieldFunctionOfFst, Pphi2-canonicalSmoothFieldFunction}
  \texttt{Pphi2.canonicalSmoothFieldFunction\_eq\_of\_fst}
\end{theorem}
\begin{theorem}[\texttt{canonicalSmoothWickPower\_two\_site\_marginal\_eq\_zero\_of\_ne}]\label{Pphi2-canonicalSmoothWickPower-two-site-marginal-eq-zero-of-ne}
  \lean{Pphi2.canonicalSmoothWickPower_two_site_marginal_eq_zero_of_ne}\leanok
  \uses{Pphi2-smoothWickConstant, GaussianField-FinLatticeSites, Pphi2-canonicalSmoothGamma, Pphi2-wickMonomial, Pphi2-canonicalSmoothFieldFunctionOfFst, Pphi2-wickPower-two-site-pi-gaussianReal-eq-zero-of-ne, Pphi2-canonicalSmoothFieldFunctionOfFst-eq-sum-gamma}
  \texttt{Pphi2.canonicalSmoothWickPower\_two\_site\_marginal\_eq\_zero\_of\_ne}
\end{theorem}
\begin{definition}[\texttt{canonicalRoughFieldFunction\_self\_moment\_diag}]\label{Pphi2-canonicalRoughFieldFunction-self-moment-diag}
  \lean{Pphi2.canonicalRoughFieldFunction_self_moment_diag}\leanok
  \uses{Pphi2-canonicalJointMeasure, GaussianField-FinLatticeSites, Pphi2-canonicalRoughFieldFunction, Pphi2-CanonicalJoint}
  \texttt{Pphi2.canonicalRoughFieldFunction\_self\_moment\_diag}
\end{definition}
\begin{theorem}[\texttt{canonicalSumFieldFunction\_add}]\label{Pphi2-canonicalSumFieldFunction-add}
  \lean{Pphi2.canonicalSumFieldFunction_add}\leanok
  \uses{GaussianField-FinLatticeSites, Pphi2-canonicalSmoothFieldFunction-add, Pphi2-canonicalSumFieldFunction, Pphi2-canonicalRoughFieldFunction, Pphi2-canonicalSumFieldFunction-eq-smooth-plus-rough, Pphi2-canonicalRoughFieldFunction-add, Pphi2-canonicalSmoothFieldFunction, Pphi2-CanonicalJoint}
  \texttt{Pphi2.canonicalSumFieldFunction\_add}
\end{theorem}
\begin{theorem}[\texttt{canonicalSmoothFieldFunction\_measurable}]\label{Pphi2-canonicalSmoothFieldFunction-measurable}
  \lean{Pphi2.canonicalSmoothFieldFunction_measurable}\leanok
  \uses{GaussianField-FinLatticeSites, Pphi2-canonicalSmoothFieldFunction-pointwise-measurable, Pphi2-canonicalSmoothFieldFunction, Pphi2-CanonicalJoint}
  \texttt{Pphi2.canonicalSmoothFieldFunction\_measurable}
\end{theorem}
\begin{theorem}[\texttt{canonicalSmoothRough\_cross\_moment\_zero}]\label{Pphi2-canonicalSmoothRough-cross-moment-zero}
  \lean{Pphi2.canonicalSmoothRough_cross_moment_zero}\leanok
  \uses{Pphi2-canonicalJointMeasure, GaussianField-FinLatticeSites, Pphi2-canonicalRoughFieldFunction, Pphi2-canonicalSmoothFieldFunction, Pphi2-CanonicalJoint}
  \texttt{Pphi2.canonicalSmoothRough\_cross\_moment\_zero}
\end{theorem}
\begin{definition}[\texttt{canonicalSumFieldLaw}]\label{Pphi2-canonicalSumFieldLaw}
  \lean{Pphi2.canonicalSumFieldLaw}\leanok
  \uses{GaussianField-FinLatticeSites, Pphi2-canonicalStdToFieldCLM, Pphi2-CanonicalJointSumIndex, GaussianField-FinLatticeField}
  \texttt{Pphi2.canonicalSumFieldLaw}
\end{definition}
\begin{definition}[\texttt{μ\_joint}]\label{Pphi2-FieldDecomposition---joint}
  \lean{Pphi2.FieldDecomposition.μ_joint}\leanok
  \uses{Pphi2-FieldDecomposition, Pphi2-FieldDecomposition-Joint, Pphi2-FieldDecomposition-jointMeasurable}
  \texttt{Pphi2.FieldDecomposition.μ\_joint}
\end{definition}
\begin{theorem}[\texttt{canonicalSumFieldFunction\_eq\_smooth\_plus\_rough}]\label{Pphi2-canonicalSumFieldFunction-eq-smooth-plus-rough}
  \lean{Pphi2.canonicalSumFieldFunction_eq_smooth_plus_rough}\leanok
  \uses{GaussianField-FinLatticeSites, Pphi2-canonicalSumFieldFunction, Pphi2-canonicalRoughFieldFunction, Pphi2-canonicalSmoothFieldFunction, Pphi2-CanonicalJoint}
  \texttt{Pphi2.canonicalSumFieldFunction\_eq\_smooth\_plus\_rough}
\end{theorem}
\begin{definition}[\texttt{canonicalSmoothWeight}]\label{Pphi2-canonicalSmoothWeight}
  \lean{Pphi2.canonicalSmoothWeight}\leanok
  \uses{Pphi2-canonicalEigenvalue}
  \texttt{Pphi2.canonicalSmoothWeight}
\end{definition}
\begin{definition}[\texttt{canonicalRoughFieldFunction}]\label{Pphi2-canonicalRoughFieldFunction}
  \lean{Pphi2.canonicalRoughFieldFunction}\leanok
  \uses{GaussianField-FinLatticeSites, Pphi2-CanonicalJoint}
  \texttt{Pphi2.canonicalRoughFieldFunction}
\end{definition}
\begin{theorem}[\texttt{canonicalSumFieldFunction\_measurable}]\label{Pphi2-canonicalSumFieldFunction-measurable}
  \lean{Pphi2.canonicalSumFieldFunction_measurable}\leanok
  \uses{GaussianField-FinLatticeSites, Pphi2-canonicalSumFieldFunction, Pphi2-canonicalSumFieldFunction-pointwise-measurable, Pphi2-CanonicalJoint}
  \texttt{Pphi2.canonicalSumFieldFunction\_measurable}
\end{theorem}
\begin{definition}[\texttt{ctorIdx}]\label{Pphi2-FieldDecomposition-ctorIdx}
  \lean{Pphi2.FieldDecomposition.ctorIdx}\leanok
  \uses{Pphi2-FieldDecomposition}
  \texttt{Pphi2.FieldDecomposition.ctorIdx}
\end{definition}
\begin{theorem}[\texttt{canonicalRoughFieldFunction\_smul}]\label{Pphi2-canonicalRoughFieldFunction-smul}
  \lean{Pphi2.canonicalRoughFieldFunction_smul}\leanok
  \uses{GaussianField-FinLatticeSites, Pphi2-canonicalRoughFieldFunction, Pphi2-CanonicalJoint}
  \texttt{Pphi2.canonicalRoughFieldFunction\_smul}
\end{theorem}
\begin{theorem}[\texttt{congr\_simp}]\label{Pphi2-canonicalRoughFieldFunctionOfSnd-congr-simp}
  \lean{Pphi2.canonicalRoughFieldFunctionOfSnd.congr_simp}\leanok
  \uses{Pphi2-canonicalRoughFieldFunctionOfSnd, GaussianField-FinLatticeSites}
  \texttt{Pphi2.canonicalRoughFieldFunctionOfSnd.congr\_simp}
\end{theorem}
\begin{definition}[\texttt{sitePairingCLM}]\label{Pphi2-sitePairingCLM}
  \lean{Pphi2.sitePairingCLM}\leanok
  \uses{GaussianField-FinLatticeSites, GaussianField-FinLatticeField}
  \texttt{Pphi2.sitePairingCLM}
\end{definition}
\begin{theorem}[\texttt{canonicalSumFieldLaw\_isGaussian}]\label{Pphi2-canonicalSumFieldLaw-isGaussian}
  \lean{Pphi2.canonicalSumFieldLaw_isGaussian}\leanok
  \uses{GaussianField-FinLatticeSites, Pphi2-canonicalStdToFieldCLM, Pphi2-CanonicalJointSumIndex, Pphi2-canonicalSumFieldLaw, GaussianField-FinLatticeField}
  \texttt{Pphi2.canonicalSumFieldLaw\_isGaussian}
\end{theorem}
\begin{theorem}[\texttt{integral\_comp\_canonicalSumConfig}]\label{Pphi2-integral-comp-canonicalSumConfig}
  \lean{Pphi2.integral_comp_canonicalSumConfig}\leanok
  \uses{Pphi2-canonicalJointMeasure, Pphi2-canonicalSumConfig, GaussianField-FinLatticeSites, Pphi2-canonicalSumConfig-measurable, GaussianField-latticeGaussianMeasure, GaussianField-Configuration, GaussianField-instMeasurableSpaceConfiguration, Pphi2-canonicalJointMeasure-map-canonicalSumConfig, GaussianField-FinLatticeField, Pphi2-CanonicalJoint}
  \texttt{Pphi2.integral\_comp\_canonicalSumConfig}
\end{theorem}
\begin{theorem}[\texttt{φ\_S\_measurable}]\label{Pphi2-FieldDecomposition---S-measurable}
  \lean{Pphi2.FieldDecomposition.φ_S_measurable}\leanok
  \uses{GaussianField-FinLatticeSites, Pphi2-FieldDecomposition---S, Pphi2-FieldDecomposition, Pphi2-FieldDecomposition-Joint, GaussianField-Configuration, GaussianField-instMeasurableSpaceConfiguration, Pphi2-FieldDecomposition-jointMeasurable, GaussianField-FinLatticeField}
  \texttt{Pphi2.FieldDecomposition.φ\_S\_measurable}
\end{theorem}
\begin{definition}[\texttt{canonicalSmoothFieldFunction}]\label{Pphi2-canonicalSmoothFieldFunction}
  \lean{Pphi2.canonicalSmoothFieldFunction}\leanok
  \uses{GaussianField-FinLatticeSites, Pphi2-CanonicalJoint}
  \texttt{Pphi2.canonicalSmoothFieldFunction}
\end{definition}
\begin{theorem}[\texttt{canonicalSumConfig\_measurable}]\label{Pphi2-canonicalSumConfig-measurable}
  \lean{Pphi2.canonicalSumConfig_measurable}\leanok
  \uses{Pphi2-canonicalSumConfig, GaussianField-FinLatticeSites, Pphi2-canonicalSumFieldFunction-measurable, GaussianField-evalMapInv, GaussianField-Configuration, Pphi2-canonicalSumFieldFunction, GaussianField-evalMapInv-measurable, Pphi2-latticeFieldToConfig, GaussianField-instMeasurableSpaceConfiguration, Pphi2-latticeFieldToConfig-eq-evalMapInv, GaussianField-FinLatticeField, Pphi2-CanonicalJoint}
  \texttt{Pphi2.canonicalSumConfig\_measurable}
\end{theorem}
\begin{theorem}[\texttt{canonicalSmoothFieldFunction\_self\_moment\_const}]\label{Pphi2-canonicalSmoothFieldFunction-self-moment-const}
  \lean{Pphi2.canonicalSmoothFieldFunction_self_moment_const}\leanok
  \uses{Pphi2-smoothWickConstant, Pphi2-canonicalSmoothFieldFunction-self-moment-eq-smoothWickConstant, GaussianField-FinLatticeSites, Pphi2-canonicalSmoothFieldFunction-self-moment-diag}
  \texttt{Pphi2.canonicalSmoothFieldFunction\_self\_moment\_const}
\end{theorem}
\begin{definition}[\texttt{latticeFieldToConfig}]\label{Pphi2-latticeFieldToConfig}
  \lean{Pphi2.latticeFieldToConfig}\leanok
  \uses{GaussianField-FinLatticeSites, GaussianField-Configuration, GaussianField-FinLatticeField}
  \texttt{Pphi2.latticeFieldToConfig}
\end{definition}
\begin{theorem}[\texttt{canonicalSmoothFieldFunction\_smul}]\label{Pphi2-canonicalSmoothFieldFunction-smul}
  \lean{Pphi2.canonicalSmoothFieldFunction_smul}\leanok
  \uses{GaussianField-FinLatticeSites, Pphi2-canonicalSmoothFieldFunction, Pphi2-CanonicalJoint}
  \texttt{Pphi2.canonicalSmoothFieldFunction\_smul}
\end{theorem}
\begin{definition}[\texttt{canonicalStdToFieldLinearMap}]\label{Pphi2-canonicalStdToFieldLinearMap}
  \lean{Pphi2.canonicalStdToFieldLinearMap}\leanok
  \uses{GaussianField-FinLatticeSites, Pphi2-CanonicalJointSumIndex, GaussianField-FinLatticeField, Pphi2-canonicalStdFieldFunction}
  \texttt{Pphi2.canonicalStdToFieldLinearMap}
\end{definition}
\begin{theorem}[\texttt{congr\_simp}]\label{Pphi2-canonicalSmoothCovariance-congr-simp}
  \lean{Pphi2.canonicalSmoothCovariance.congr_simp}\leanok
  \uses{Pphi2-canonicalSmoothCovariance, GaussianField-FinLatticeSites}
  \texttt{Pphi2.canonicalSmoothCovariance.congr\_simp}
\end{theorem}
\begin{definition}[\texttt{canonicalJointMeasure}]\label{Pphi2-canonicalJointMeasure}
  \lean{Pphi2.canonicalJointMeasure}\leanok
  \uses{Pphi2-CanonicalJoint}
  \texttt{Pphi2.canonicalJointMeasure}
\end{definition}
\begin{theorem}[\texttt{congr\_simp}]\label{Pphi2-CanonicalJoint-congr-simp}
  \lean{Pphi2.CanonicalJoint.congr_simp}\leanok
  \uses{Pphi2-CanonicalJoint}
  \texttt{Pphi2.CanonicalJoint.congr\_simp}
\end{theorem}
\begin{theorem}[\texttt{canonicalSmoothFieldFunction\_pointwise\_measurable}]\label{Pphi2-canonicalSmoothFieldFunction-pointwise-measurable}
  \lean{Pphi2.canonicalSmoothFieldFunction_pointwise_measurable}\leanok
  \uses{GaussianField-FinLatticeSites, Pphi2-canonicalSmoothFieldFunction, Pphi2-CanonicalJoint}
  \texttt{Pphi2.canonicalSmoothFieldFunction\_pointwise\_measurable}
\end{theorem}
\begin{theorem}[\texttt{congr\_simp}]\label{Pphi2-canonicalRoughWeight-congr-simp}
  \lean{Pphi2.canonicalRoughWeight.congr_simp}\leanok
  \uses{Pphi2-canonicalRoughWeight}
  \texttt{Pphi2.canonicalRoughWeight.congr\_simp}
\end{theorem}
\begin{theorem}[\texttt{congr\_simp}]\label{Pphi2-roughWickConstant-congr-simp}
  \lean{Pphi2.roughWickConstant.congr_simp}\leanok
  \uses{Pphi2-roughWickConstant}
  \texttt{Pphi2.roughWickConstant.congr\_simp}
\end{theorem}
\begin{theorem}[\texttt{congr\_simp}]\label{Pphi2-canonicalSmoothGamma-congr-simp}
  \lean{Pphi2.canonicalSmoothGamma.congr_simp}\leanok
  \uses{GaussianField-FinLatticeSites, Pphi2-canonicalSmoothGamma}
  \texttt{Pphi2.canonicalSmoothGamma.congr\_simp}
\end{theorem}
\begin{theorem}[\texttt{jointProbability}]\label{Pphi2-FieldDecomposition-jointProbability}
  \lean{Pphi2.FieldDecomposition.jointProbability}\leanok
  \uses{Pphi2-FieldDecomposition---joint, Pphi2-FieldDecomposition, Pphi2-FieldDecomposition-Joint, Pphi2-FieldDecomposition-jointMeasurable}
  \texttt{Pphi2.FieldDecomposition.jointProbability}
\end{theorem}
\begin{theorem}[\texttt{canonicalSumFieldLaw\_eq\_latticeGaussianFieldLaw}]\label{Pphi2-canonicalSumFieldLaw-eq-latticeGaussianFieldLaw}
  \lean{Pphi2.canonicalSumFieldLaw_eq_latticeGaussianFieldLaw}\leanok
  \uses{GaussianField-measurable-evalMap, GaussianField-latticeGaussianFieldLaw-pairing-is-gaussian, GaussianField-evalMap, GaussianField-FinLatticeSites, GaussianField-latticeCovarianceGJ, Pphi2-canonicalSumFieldLaw-isGaussian, GaussianField-latticeGaussianFieldLaw, GaussianField-latticeGaussianMeasure-isProbability, GaussianField-latticeGaussianMeasure, GaussianField-cramerWold, GaussianField-Configuration, GaussianField-instMeasurableSpaceConfiguration, Pphi2-canonicalSumFieldLaw, GaussianField-ell2-, GaussianField-FinLatticeField}
  \texttt{Pphi2.canonicalSumFieldLaw\_eq\_latticeGaussianFieldLaw}
\end{theorem}
\begin{theorem}[\texttt{canonicalSmoothWickPower\_two\_site\_marginal\_eq\_diag}]\label{Pphi2-canonicalSmoothWickPower-two-site-marginal-eq-diag}
  \lean{Pphi2.canonicalSmoothWickPower_two_site_marginal_eq_diag}\leanok
  \uses{Pphi2-smoothWickConstant, Pphi2-canonicalSmoothCovariance, GaussianField-FinLatticeSites, Pphi2-canonicalSmoothGamma, Pphi2-wickPower-two-site-pi-gaussianReal-eq-diag, Pphi2-wickMonomial, Pphi2-canonicalSmoothFieldFunctionOfFst, Pphi2-canonicalSmoothFieldFunctionOfFst-eq-sum-gamma}
  \texttt{Pphi2.canonicalSmoothWickPower\_two\_site\_marginal\_eq\_diag}
\end{theorem}
\begin{definition}[\texttt{φ\_S}]\label{Pphi2-FieldDecomposition---S}
  \lean{Pphi2.FieldDecomposition.φ_S}\leanok
  \uses{GaussianField-FinLatticeSites, Pphi2-FieldDecomposition, Pphi2-FieldDecomposition-Joint, GaussianField-Configuration, GaussianField-FinLatticeField}
  \texttt{Pphi2.FieldDecomposition.φ\_S}
\end{definition}
\begin{definition}[\texttt{canonicalStdFieldFunction}]\label{Pphi2-canonicalStdFieldFunction}
  \lean{Pphi2.canonicalStdFieldFunction}\leanok
  \uses{GaussianField-FinLatticeSites, Pphi2-CanonicalJointSumIndex}
  \texttt{Pphi2.canonicalStdFieldFunction}
\end{definition}
\begin{theorem}[\texttt{canonicalSmoothFieldFunction\_self\_moment\_eq\_smoothWickConstant}]\label{Pphi2-canonicalSmoothFieldFunction-self-moment-eq-smoothWickConstant}
  \lean{Pphi2.canonicalSmoothFieldFunction_self_moment_eq_smoothWickConstant}\leanok
  \uses{Pphi2-smoothWickConstant, Pphi2-canonicalJointMeasure, Pphi2-canonicalSmoothWeight, GaussianField-FinLatticeSites, GaussianField-latticeFourierProductBasisFun, Pphi2-canonicalSmoothFieldFunction-self-moment-diag, Pphi2-canonicalSmoothFieldFunction, GaussianField-latticeFourierProductNormSq, Pphi2-CanonicalJoint}
  \texttt{Pphi2.canonicalSmoothFieldFunction\_self\_moment\_eq\_smoothWickConstant}
\end{theorem}
\begin{theorem}[\texttt{canonicalSumFieldFunction\_covariance}]\label{Pphi2-canonicalSumFieldFunction-covariance}
  \lean{Pphi2.canonicalSumFieldFunction_covariance}\leanok
  \uses{Pphi2-canonicalJointMeasure, Pphi2-canonicalSmoothWeight, Pphi2-canonicalRoughFieldFunction-memLp-two, GaussianField-FinLatticeSites, Pphi2-canonicalSmoothFieldFunction-memLp-two, Pphi2-canonicalEigenvalue, Pphi2-canonicalRoughWeight, Pphi2-canonicalSumFieldFunction, Pphi2-canonicalRoughFieldFunction, Pphi2-canonicalSmoothRough-cross-moment-zero, GaussianField-latticeFourierProductBasisFun, Pphi2-canonicalSmoothFieldFunction, GaussianField-latticeFourierProductNormSq-pos, GaussianField-latticeFourierProductNormSq, Pphi2-CanonicalJoint}
  \texttt{Pphi2.canonicalSumFieldFunction\_covariance}
\end{theorem}
\begin{theorem}[\texttt{canonicalSmoothWickPower\_two\_site\_marginal\_integrable}]\label{Pphi2-canonicalSmoothWickPower-two-site-marginal-integrable}
  \lean{Pphi2.canonicalSmoothWickPower_two_site_marginal_integrable}\leanok
  \uses{Pphi2-smoothWickConstant, GaussianField-FinLatticeSites, Pphi2-canonicalSmoothGamma, Pphi2-wickPower-two-site-pi-gaussianReal-integrable, Pphi2-wickMonomial, Pphi2-canonicalSmoothFieldFunctionOfFst, Pphi2-canonicalSmoothFieldFunctionOfFst-eq-sum-gamma}
  \texttt{Pphi2.canonicalSmoothWickPower\_two\_site\_marginal\_integrable}
\end{theorem}
\begin{theorem}[\texttt{congr\_simp}]\label{Pphi2-canonicalEigenvalue-congr-simp}
  \lean{Pphi2.canonicalEigenvalue.congr_simp}\leanok
  \uses{Pphi2-canonicalEigenvalue}
  \texttt{Pphi2.canonicalEigenvalue.congr\_simp}
\end{theorem}
\begin{definition}[\texttt{φ\_R}]\label{Pphi2-FieldDecomposition---R}
  \lean{Pphi2.FieldDecomposition.φ_R}\leanok
  \uses{GaussianField-FinLatticeSites, Pphi2-FieldDecomposition, Pphi2-FieldDecomposition-Joint, GaussianField-Configuration, GaussianField-FinLatticeField}
  \texttt{Pphi2.FieldDecomposition.φ\_R}
\end{definition}
\begin{theorem}[\texttt{canonicalRoughCovariance\_eq\_sum\_gamma\_mul\_gamma}]\label{Pphi2-canonicalRoughCovariance-eq-sum-gamma-mul-gamma}
  \lean{Pphi2.canonicalRoughCovariance_eq_sum_gamma_mul_gamma}\leanok
  \uses{Pphi2-canonicalRoughCovariance, GaussianField-FinLatticeSites, Pphi2-canonicalRoughWeight, GaussianField-latticeFourierProductBasisFun, GaussianField-latticeFourierProductNormSq-pos, Pphi2-canonicalRoughGamma, GaussianField-latticeFourierProductNormSq}
  \texttt{Pphi2.canonicalRoughCovariance\_eq\_sum\_gamma\_mul\_gamma}
\end{theorem}
\begin{theorem}[\texttt{canonicalJointMeasure\_map\_stdGaussianEuclidean}]\label{Pphi2-canonicalJointMeasure-map-stdGaussianEuclidean}
  \lean{Pphi2.canonicalJointMeasure_map_stdGaussianEuclidean}\leanok
  \uses{Pphi2-canonicalJointMeasure, Pphi2-canonicalJointMeasure-map-sum-pi-gaussian, Pphi2-CanonicalJointSumIndex, Pphi2-CanonicalJoint}
  \texttt{Pphi2.canonicalJointMeasure\_map\_stdGaussianEuclidean}
\end{theorem}
\begin{theorem}[\texttt{canonicalStdFieldFunction\_eq\_canonicalSumFieldFunction}]\label{Pphi2-canonicalStdFieldFunction-eq-canonicalSumFieldFunction}
  \lean{Pphi2.canonicalStdFieldFunction_eq_canonicalSumFieldFunction}\leanok
  \uses{GaussianField-FinLatticeSites, Pphi2-canonicalSumFieldFunction, Pphi2-CanonicalJointSumIndex, Pphi2-canonicalRoughFieldFunction, Pphi2-canonicalSmoothFieldFunction, Pphi2-canonicalStdFieldFunction, Pphi2-CanonicalJoint}
  \texttt{Pphi2.canonicalStdFieldFunction\_eq\_canonicalSumFieldFunction}
\end{theorem}
\begin{definition}[\texttt{canonicalRoughConfig}]\label{Pphi2-canonicalRoughConfig}
  \lean{Pphi2.canonicalRoughConfig}\leanok
  \uses{GaussianField-FinLatticeSites, GaussianField-Configuration, Pphi2-canonicalRoughFieldFunction, Pphi2-latticeFieldToConfig, GaussianField-FinLatticeField, Pphi2-CanonicalJoint}
  \texttt{Pphi2.canonicalRoughConfig}
\end{definition}
\begin{theorem}[\texttt{congr\_simp}]\label{Pphi2-canonicalRoughCovariance-congr-simp}
  \lean{Pphi2.canonicalRoughCovariance.congr_simp}\leanok
  \uses{Pphi2-canonicalRoughCovariance, GaussianField-FinLatticeSites}
  \texttt{Pphi2.canonicalRoughCovariance.congr\_simp}
\end{theorem}
\begin{theorem}[\texttt{canonicalRoughCovariance\_sq\_row\_sum\_eq\_double\_integral}]\label{Pphi2-canonicalRoughCovariance-sq-row-sum-eq-double-integral}
  \lean{Pphi2.canonicalRoughCovariance_sq_row_sum_eq_double_integral}\leanok
  \uses{GaussianField-latticeEigenvalue1d, Pphi2-canonicalRoughCovariance, GaussianField-FinLatticeSites, Pphi2-canonicalRoughCovariance-eq-integral-Icc-heatKernel, Pphi2-heatKernel-row-pairing-eq-diagonal, GaussianField-latticeFourierProductBasisFun, GaussianField-latticeFourierProductNormSq, GaussianField-latticeHeatKernelMatrix}
  \texttt{Pphi2.canonicalRoughCovariance\_sq\_row\_sum\_eq\_double\_integral}
\end{theorem}
\begin{theorem}[\texttt{canonicalSumFieldFunction\_pointwise\_measurable}]\label{Pphi2-canonicalSumFieldFunction-pointwise-measurable}
  \lean{Pphi2.canonicalSumFieldFunction_pointwise_measurable}\leanok
  \uses{GaussianField-FinLatticeSites, Pphi2-canonicalRoughFieldFunction-pointwise-measurable, Pphi2-canonicalSumFieldFunction, Pphi2-canonicalRoughFieldFunction, Pphi2-canonicalSmoothFieldFunction-pointwise-measurable, Pphi2-canonicalSmoothFieldFunction, Pphi2-CanonicalJoint}
  \texttt{Pphi2.canonicalSumFieldFunction\_pointwise\_measurable}
\end{theorem}
\begin{theorem}[\texttt{canonicalRoughFieldFunction\_self\_moment\_const}]\label{Pphi2-canonicalRoughFieldFunction-self-moment-const}
  \lean{Pphi2.canonicalRoughFieldFunction_self_moment_const}\leanok
  \uses{GaussianField-FinLatticeSites, Pphi2-canonicalRoughFieldFunction-self-moment-eq-roughWickConstant, Pphi2-canonicalRoughFieldFunction-self-moment-diag, Pphi2-roughWickConstant}
  \texttt{Pphi2.canonicalRoughFieldFunction\_self\_moment\_const}
\end{theorem}
\begin{theorem}[\texttt{canonicalJointMeasure\_map\_canonicalSumConfig}]\label{Pphi2-canonicalJointMeasure-map-canonicalSumConfig}
  \lean{Pphi2.canonicalJointMeasure_map_canonicalSumConfig}\leanok
  \uses{Pphi2-canonicalJointMeasure, GaussianField-measurable-evalMap, Pphi2-canonicalSumConfig, GaussianField-evalMap, GaussianField-FinLatticeSites, Pphi2-canonicalSumFieldFunction-measurable, GaussianField-evalMapInv, GaussianField-latticeGaussianFieldLaw, GaussianField-latticeGaussianMeasure, GaussianField-Configuration, Pphi2-canonicalSumFieldFunction, GaussianField-evalMapInv-evalMap, GaussianField-evalMapInv-measurable, Pphi2-canonicalSumFieldLaw-eq-map-canonicalSumFieldFunction, Pphi2-latticeFieldToConfig, GaussianField-instMeasurableSpaceConfiguration, Pphi2-canonicalSumFieldLaw-eq-latticeGaussianFieldLaw, Pphi2-latticeFieldToConfig-eq-evalMapInv, Pphi2-canonicalSumFieldLaw, GaussianField-FinLatticeField, Pphi2-CanonicalJoint}
  \texttt{Pphi2.canonicalJointMeasure\_map\_canonicalSumConfig}
\end{theorem}
\begin{theorem}[\texttt{φ\_R\_measurable}]\label{Pphi2-FieldDecomposition---R-measurable}
  \lean{Pphi2.FieldDecomposition.φ_R_measurable}\leanok
  \uses{GaussianField-FinLatticeSites, Pphi2-FieldDecomposition, Pphi2-FieldDecomposition---R, Pphi2-FieldDecomposition-Joint, GaussianField-Configuration, GaussianField-instMeasurableSpaceConfiguration, Pphi2-FieldDecomposition-jointMeasurable, GaussianField-FinLatticeField}
  \texttt{Pphi2.FieldDecomposition.φ\_R\_measurable}
\end{theorem}
\begin{definition}[\texttt{canonicalEigenvalue}]\label{Pphi2-canonicalEigenvalue}
  \lean{Pphi2.canonicalEigenvalue}\leanok
  \uses{GaussianField-latticeEigenvalue1d}
  \texttt{Pphi2.canonicalEigenvalue}
\end{definition}
\begin{theorem}[\texttt{congr\_simp}]\label{Pphi2-canonicalRoughGamma-congr-simp}
  \lean{Pphi2.canonicalRoughGamma.congr_simp}\leanok
  \uses{GaussianField-FinLatticeSites, Pphi2-canonicalRoughGamma}
  \texttt{Pphi2.canonicalRoughGamma.congr\_simp}
\end{theorem}
\begin{definition}[\texttt{canonicalSmoothFieldFunction\_self\_moment\_diag}]\label{Pphi2-canonicalSmoothFieldFunction-self-moment-diag}
  \lean{Pphi2.canonicalSmoothFieldFunction_self_moment_diag}\leanok
  \uses{Pphi2-canonicalJointMeasure, GaussianField-FinLatticeSites, Pphi2-canonicalSmoothFieldFunction, Pphi2-CanonicalJoint}
  \texttt{Pphi2.canonicalSmoothFieldFunction\_self\_moment\_diag}
\end{definition}
\begin{definition}[\texttt{canonicalRoughFieldFunctionOfSnd}]\label{Pphi2-canonicalRoughFieldFunctionOfSnd}
  \lean{Pphi2.canonicalRoughFieldFunctionOfSnd}\leanok
  \uses{GaussianField-FinLatticeSites}
  \texttt{Pphi2.canonicalRoughFieldFunctionOfSnd}
\end{definition}
\begin{definition}[\texttt{jointMeasurable}]\label{Pphi2-FieldDecomposition-jointMeasurable}
  \lean{Pphi2.FieldDecomposition.jointMeasurable}\leanok
  \uses{Pphi2-FieldDecomposition, Pphi2-FieldDecomposition-Joint}
  \texttt{Pphi2.FieldDecomposition.jointMeasurable}
\end{definition}
\begin{definition}[\texttt{canonicalSumConfig}]\label{Pphi2-canonicalSumConfig}
  \lean{Pphi2.canonicalSumConfig}\leanok
  \uses{GaussianField-FinLatticeSites, GaussianField-Configuration, Pphi2-canonicalSumFieldFunction, Pphi2-latticeFieldToConfig, GaussianField-FinLatticeField, Pphi2-CanonicalJoint}
  \texttt{Pphi2.canonicalSumConfig}
\end{definition}
\begin{definition}[\texttt{CanonicalJoint}]\label{Pphi2-CanonicalJoint}
  \lean{Pphi2.CanonicalJoint}\leanok
  \texttt{Pphi2.CanonicalJoint}
\end{definition}
\begin{definition}[\texttt{canonicalStdToFieldCLM}]\label{Pphi2-canonicalStdToFieldCLM}
  \lean{Pphi2.canonicalStdToFieldCLM}\leanok
  \uses{GaussianField-FinLatticeSites, Pphi2-CanonicalJointSumIndex, Pphi2-canonicalStdToFieldLinearMap, GaussianField-FinLatticeField}
  \texttt{Pphi2.canonicalStdToFieldCLM}
\end{definition}
\begin{theorem}[\texttt{latticeFieldToConfig\_eq\_evalMapInv}]\label{Pphi2-latticeFieldToConfig-eq-evalMapInv}
  \lean{Pphi2.latticeFieldToConfig_eq_evalMapInv}\leanok
  \uses{GaussianField-evalMapInv-apply, GaussianField-FinLatticeSites, GaussianField-evalMapInv, GaussianField-Configuration, Pphi2-latticeFieldToConfig, GaussianField-FinLatticeField}
  \texttt{Pphi2.latticeFieldToConfig\_eq\_evalMapInv}
\end{theorem}
\begin{theorem}[\texttt{canonicalRoughWickPower\_two\_site\_marginal\_eq\_zero\_of\_ne}]\label{Pphi2-canonicalRoughWickPower-two-site-marginal-eq-zero-of-ne}
  \lean{Pphi2.canonicalRoughWickPower_two_site_marginal_eq_zero_of_ne}\leanok
  \uses{Pphi2-canonicalRoughFieldFunctionOfSnd-eq-sum-gamma, Pphi2-canonicalRoughFieldFunctionOfSnd, GaussianField-FinLatticeSites, Pphi2-wickMonomial, Pphi2-roughWickConstant, Pphi2-canonicalRoughGamma, Pphi2-wickPower-two-site-pi-gaussianReal-eq-zero-of-ne}
  \texttt{Pphi2.canonicalRoughWickPower\_two\_site\_marginal\_eq\_zero\_of\_ne}
\end{theorem}
\begin{theorem}[\texttt{pi\_gaussian\_bilinear\_moment}]\label{Pphi2-pi-gaussian-bilinear-moment}
  \lean{Pphi2.pi_gaussian_bilinear_moment}\leanok
  \texttt{Pphi2.pi\_gaussian\_bilinear\_moment}
\end{theorem}
\begin{theorem}[\texttt{canonicalRoughWickPower\_two\_site\_marginal\_integrable}]\label{Pphi2-canonicalRoughWickPower-two-site-marginal-integrable}
  \lean{Pphi2.canonicalRoughWickPower_two_site_marginal_integrable}\leanok
  \uses{Pphi2-canonicalRoughFieldFunctionOfSnd-eq-sum-gamma, Pphi2-canonicalRoughFieldFunctionOfSnd, GaussianField-FinLatticeSites, Pphi2-wickPower-two-site-pi-gaussianReal-integrable, Pphi2-wickMonomial, Pphi2-roughWickConstant, Pphi2-canonicalRoughGamma}
  \texttt{Pphi2.canonicalRoughWickPower\_two\_site\_marginal\_integrable}
\end{theorem}
\begin{definition}[\texttt{canonicalSumFieldFunction}]\label{Pphi2-canonicalSumFieldFunction}
  \lean{Pphi2.canonicalSumFieldFunction}\leanok
  \uses{GaussianField-FinLatticeSites, Pphi2-canonicalRoughFieldFunction, Pphi2-canonicalSmoothFieldFunction, Pphi2-CanonicalJoint}
  \texttt{Pphi2.canonicalSumFieldFunction}
\end{definition}
\begin{definition}[\texttt{canonicalSmoothFieldFunctionOfFst}]\label{Pphi2-canonicalSmoothFieldFunctionOfFst}
  \lean{Pphi2.canonicalSmoothFieldFunctionOfFst}\leanok
  \uses{GaussianField-FinLatticeSites}
  \texttt{Pphi2.canonicalSmoothFieldFunctionOfFst}
\end{definition}
\section{\texttt{Pphi2.NelsonEstimate.LatticeSetup}}
\begin{theorem}[\texttt{latticeRoughError\_pure\_quartic\_summed}]\label{Pphi2-LatticeSetup-latticeRoughError-pure-quartic-summed}
  \lean{Pphi2.LatticeSetup.latticeRoughError_pure_quartic_summed}\leanok
  \uses{Pphi2-smoothWickConstant, GaussianField-finLatticeDelta, GaussianField-FinLatticeSites, Pphi2-LatticeSetup-latticeRoughError, Pphi2-wickMonomial, GaussianField-Configuration, Pphi2-wickConstant, Pphi2-LatticeSetup-wickPolynomial-pure-quartic-diff, Pphi2-LatticeSetup-latticeRoughError-eq-sum-diff, GaussianField-FinLatticeField, Pphi2-wickPolynomial}
  \texttt{Pphi2.LatticeSetup.latticeRoughError\_pure\_quartic\_summed}
\end{theorem}
\begin{theorem}[\texttt{latticeSmoothInteraction\_measurable}]\label{Pphi2-LatticeSetup-latticeSmoothInteraction-measurable}
  \lean{Pphi2.LatticeSetup.latticeSmoothInteraction_measurable}\leanok
  \uses{Pphi2-LatticeSetup-latticeSmoothInteraction, Pphi2-smoothWickConstant, GaussianField-finLatticeDelta, GaussianField-FinLatticeSites, Pphi2-wickPolynomial-continuous-, GaussianField-configuration-eval-measurable, GaussianField-Configuration, GaussianField-instMeasurableSpaceConfiguration, GaussianField-FinLatticeField, Pphi2-wickPolynomial}
  \texttt{Pphi2.LatticeSetup.latticeSmoothInteraction\_measurable}
\end{theorem}
\begin{theorem}[\texttt{wickPolynomial\_pure\_quartic\_diff}]\label{Pphi2-LatticeSetup-wickPolynomial-pure-quartic-diff}
  \lean{Pphi2.LatticeSetup.wickPolynomial_pure_quartic_diff}\leanok
  \uses{Pphi2-smoothWickConstant, Pphi2-LatticeSetup-wickMonomial-four-diff, GaussianField-finLatticeDelta, GaussianField-FinLatticeSites, Pphi2-LatticeSetup-wickPolynomial-of-pure, Pphi2-wickMonomial, GaussianField-Configuration, Pphi2-wickConstant, GaussianField-FinLatticeField, Lattice-toSemilatticeInf, Pphi2-wickPolynomial}
  \texttt{Pphi2.LatticeSetup.wickPolynomial\_pure\_quartic\_diff}
\end{theorem}
\begin{definition}[\texttt{latticeWickSquareSum}]\label{Pphi2-LatticeSetup-latticeWickSquareSum}
  \lean{Pphi2.LatticeSetup.latticeWickSquareSum}\leanok
  \uses{GaussianField-finLatticeDelta, GaussianField-FinLatticeSites, Pphi2-wickMonomial, GaussianField-Configuration, Pphi2-wickConstant, GaussianField-FinLatticeField}
  \texttt{Pphi2.LatticeSetup.latticeWickSquareSum}
\end{definition}
\begin{theorem}[\texttt{wickMonomial\_two\_shift}]\label{Pphi2-LatticeSetup-wickMonomial-two-shift}
  \lean{Pphi2.LatticeSetup.wickMonomial_two_shift}\leanok
  \uses{Pphi2-wickMonomial-two, Pphi2-wickMonomial}
  \texttt{Pphi2.LatticeSetup.wickMonomial\_two\_shift}
\end{theorem}
\begin{theorem}[\texttt{wickPolynomial\_of\_pure}]\label{Pphi2-LatticeSetup-wickPolynomial-of-pure}
  \lean{Pphi2.LatticeSetup.wickPolynomial_of_pure}\leanok
  \uses{Pphi2-wickMonomial, Pphi2-wickPolynomial}
  \texttt{Pphi2.LatticeSetup.wickPolynomial\_of\_pure}
\end{theorem}
\begin{definition}[\texttt{latticeRoughError}]\label{Pphi2-LatticeSetup-latticeRoughError}
  \lean{Pphi2.LatticeSetup.latticeRoughError}\leanok
  \uses{Pphi2-LatticeSetup-latticeSmoothInteraction, GaussianField-FinLatticeSites, GaussianField-Configuration, GaussianField-FinLatticeField, Pphi2-interactionFunctional}
  \texttt{Pphi2.LatticeSetup.latticeRoughError}
\end{definition}
\begin{theorem}[\texttt{latticeWickSquareSum\_measurable}]\label{Pphi2-LatticeSetup-latticeWickSquareSum-measurable}
  \lean{Pphi2.LatticeSetup.latticeWickSquareSum_measurable}\leanok
  \uses{GaussianField-finLatticeDelta, GaussianField-FinLatticeSites, Pphi2-LatticeSetup-latticeWickSquareSum, Pphi2-wickMonomial, GaussianField-configuration-eval-measurable, GaussianField-Configuration, Pphi2-wickConstant, GaussianField-instMeasurableSpaceConfiguration, Pphi2-wickMonomial-continuous-, GaussianField-FinLatticeField}
  \texttt{Pphi2.LatticeSetup.latticeWickSquareSum\_measurable}
\end{theorem}
\begin{definition}[\texttt{latticeSmoothInteraction}]\label{Pphi2-LatticeSetup-latticeSmoothInteraction}
  \lean{Pphi2.LatticeSetup.latticeSmoothInteraction}\leanok
  \uses{Pphi2-smoothWickConstant, GaussianField-finLatticeDelta, GaussianField-FinLatticeSites, GaussianField-Configuration, GaussianField-FinLatticeField, Pphi2-wickPolynomial}
  \texttt{Pphi2.LatticeSetup.latticeSmoothInteraction}
\end{definition}
\begin{theorem}[\texttt{interactionFunctional\_eq\_smooth\_plus\_rough}]\label{Pphi2-LatticeSetup-interactionFunctional-eq-smooth-plus-rough}
  \lean{Pphi2.LatticeSetup.interactionFunctional_eq_smooth_plus_rough}\leanok
  \uses{Pphi2-LatticeSetup-latticeSmoothInteraction, GaussianField-FinLatticeSites, Pphi2-LatticeSetup-latticeRoughError, GaussianField-Configuration, GaussianField-FinLatticeField, Pphi2-interactionFunctional}
  \texttt{Pphi2.LatticeSetup.interactionFunctional\_eq\_smooth\_plus\_rough}
\end{theorem}
\begin{theorem}[\texttt{wickMonomial\_four\_diff}]\label{Pphi2-LatticeSetup-wickMonomial-four-diff}
  \lean{Pphi2.LatticeSetup.wickMonomial_four_diff}\leanok
  \uses{Pphi2-wickMonomial-two, Pphi2-wickMonomial, Pphi2-wickMonomial-four}
  \texttt{Pphi2.LatticeSetup.wickMonomial\_four\_diff}
\end{theorem}
\begin{theorem}[\texttt{latticeRoughError\_eq\_sum\_diff}]\label{Pphi2-LatticeSetup-latticeRoughError-eq-sum-diff}
  \lean{Pphi2.LatticeSetup.latticeRoughError_eq_sum_diff}\leanok
  \uses{Pphi2-LatticeSetup-latticeSmoothInteraction, Pphi2-smoothWickConstant, GaussianField-finLatticeDelta, GaussianField-FinLatticeSites, Pphi2-LatticeSetup-latticeRoughError, GaussianField-Configuration, Pphi2-wickConstant, GaussianField-FinLatticeField, Pphi2-interactionFunctional, Pphi2-wickPolynomial}
  \texttt{Pphi2.LatticeSetup.latticeRoughError\_eq\_sum\_diff}
\end{theorem}
\begin{theorem}[\texttt{latticeSmoothInteraction\_lower\_bound\_at\_cutoff\_quartic}]\label{Pphi2-LatticeSetup-latticeSmoothInteraction-lower-bound-at-cutoff-quartic}
  \lean{Pphi2.LatticeSetup.latticeSmoothInteraction_lower_bound_at_cutoff_quartic}\leanok
  \uses{Pphi2-LatticeSetup-latticeSmoothInteraction, Pphi2-smoothWickConstant, Pphi2-DynamicalCutoff-dynamicalCutoffScale-log-sq-le, GaussianField-finLatticeDelta, GaussianField-FinLatticeSites, Pphi2-smooth-interaction-lower-bound-log-uniform, Pphi2-DynamicalCutoff-dynamicalCutoffScale-pos, Pphi2-wickMonomial, Pphi2-smoothBoundConstant, GaussianField-Configuration, Pphi2-DynamicalCutoff-dynamicalCutoffScale, Pphi2-LatticeSetup-latticeSmoothInteraction-of-pure, Pphi2-smoothBoundConstant-pos, GaussianField-FinLatticeField}
  \texttt{Pphi2.LatticeSetup.latticeSmoothInteraction\_lower\_bound\_at\_cutoff\_quartic}
\end{theorem}
\begin{theorem}[\texttt{latticeSmoothInteraction\_of\_pure}]\label{Pphi2-LatticeSetup-latticeSmoothInteraction-of-pure}
  \lean{Pphi2.LatticeSetup.latticeSmoothInteraction_of_pure}\leanok
  \uses{Pphi2-LatticeSetup-latticeSmoothInteraction, Pphi2-smoothWickConstant, GaussianField-finLatticeDelta, GaussianField-FinLatticeSites, Pphi2-LatticeSetup-wickPolynomial-of-pure, Pphi2-wickMonomial, GaussianField-Configuration, GaussianField-FinLatticeField, Pphi2-wickPolynomial}
  \texttt{Pphi2.LatticeSetup.latticeSmoothInteraction\_of\_pure}
\end{theorem}
\begin{theorem}[\texttt{latticeRoughError\_pure\_quartic\_chaos2\_form}]\label{Pphi2-LatticeSetup-latticeRoughError-pure-quartic-chaos2-form}
  \lean{Pphi2.LatticeSetup.latticeRoughError_pure_quartic_chaos2_form}\leanok
  \uses{Pphi2-smoothWickConstant, GaussianField-finLatticeDelta, GaussianField-FinLatticeSites, Pphi2-LatticeSetup-latticeRoughError, Pphi2-wickMonomial, Pphi2-LatticeSetup-wickMonomial-two-shift, GaussianField-Configuration, Pphi2-wickConstant, Pphi2-LatticeSetup-latticeRoughError-pure-quartic-summed, GaussianField-FinLatticeField}
  \texttt{Pphi2.LatticeSetup.latticeRoughError\_pure\_quartic\_chaos2\_form}
\end{theorem}
\begin{theorem}[\texttt{latticeRoughError\_measurable}]\label{Pphi2-LatticeSetup-latticeRoughError-measurable}
  \lean{Pphi2.LatticeSetup.latticeRoughError_measurable}\leanok
  \uses{Pphi2-LatticeSetup-latticeSmoothInteraction, Pphi2-LatticeSetup-latticeSmoothInteraction-measurable, GaussianField-FinLatticeSites, Pphi2-interactionFunctional-measurable, Pphi2-LatticeSetup-latticeRoughError, GaussianField-Configuration, GaussianField-instMeasurableSpaceConfiguration, GaussianField-FinLatticeField, Pphi2-interactionFunctional}
  \texttt{Pphi2.LatticeSetup.latticeRoughError\_measurable}
\end{theorem}
\section{\texttt{Pphi2.NelsonEstimate.CovarianceSplit}}
\begin{theorem}[\texttt{smoothVariance\_le\_log}]\label{Pphi2-smoothVariance-le-log}
  \lean{Pphi2.smoothVariance_le_log}\leanok
  \uses{Pphi2-smoothWickConstant, GaussianField-FinLatticeSites, GaussianField-latticeEigenvalue-pos, Pphi2-wickConstant, Pphi2-smoothCovEigenvalue, Pphi2-wickConstant-le-inv-mass-sq, GaussianField-latticeEigenvalue, Lattice-toSemilatticeInf}
  $c_S \le C \cdot (1 + |\log T|)$ for $d = 2$, since each $C_S(k) \le 1/\lambda_k \le 1/m^2$.
\end{theorem}
\begin{theorem}[\texttt{covariance\_split}]\label{Pphi2-covariance-split}
  \lean{Pphi2.covariance_split}\leanok
  \uses{Pphi2-roughCovEigenvalue, Pphi2-smoothCovEigenvalue, GaussianField-latticeEigenvalue}
  $\lambda_m^{-1} = C_S(T,m) + C_R(T,m)$.
\end{theorem}
\begin{definition}[\texttt{roughCovEigenvalue}]\label{Pphi2-roughCovEigenvalue}
  \lean{Pphi2.roughCovEigenvalue}\leanok
  \uses{GaussianField-latticeEigenvalue}
  \texttt{Pphi2.roughCovEigenvalue}
\end{definition}
\begin{theorem}[\texttt{roughCovEigenvalue\_pos}]\label{Pphi2-roughCovEigenvalue-pos}
  \lean{Pphi2.roughCovEigenvalue_pos}\leanok
  \uses{Pphi2-roughCovEigenvalue, GaussianField-latticeEigenvalue-pos, GaussianField-latticeEigenvalue}
  \texttt{Pphi2.roughCovEigenvalue\_pos}
\end{theorem}
\begin{definition}[\texttt{smoothWickConstant}]\label{Pphi2-smoothWickConstant}
  \lean{Pphi2.smoothWickConstant}\leanok
  \uses{GaussianField-FinLatticeSites, Pphi2-smoothCovEigenvalue}
  Averages of $C_S$ and $C_R$ over Fourier modes.
\end{definition}
\begin{theorem}[\texttt{roughCovEigenvalue\_le\_inv}]\label{Pphi2-roughCovEigenvalue-le-inv}
  \lean{Pphi2.roughCovEigenvalue_le_inv}\leanok
  \uses{Pphi2-roughCovEigenvalue, GaussianField-latticeEigenvalue-pos, GaussianField-latticeEigenvalue}
  \texttt{Pphi2.roughCovEigenvalue\_le\_inv}
\end{theorem}
\begin{definition}[\texttt{smoothCovEigenvalue}]\label{Pphi2-smoothCovEigenvalue}
  \lean{Pphi2.smoothCovEigenvalue}\leanok
  \uses{GaussianField-latticeEigenvalue}
  ``\texttt{lean def smoothCovEigenvalue (T : ℝ) (m : ℕ) : ℝ def roughCovEigenvalue (T : ℝ) (m : ℕ) : ℝ }`` $C_S(T,m) = e^{-T\lambda_m}/\lambda_m$ and $C_R(T,m) = (1 - e^{-T\lambda_m})/\lambda_m$.
\end{definition}
\begin{theorem}[\texttt{smoothVarianceConstant\_pos}]\label{Pphi2-smoothVarianceConstant-pos}
  \lean{Pphi2.smoothVarianceConstant_pos}\leanok
  \uses{Pphi2-smoothVarianceConstant}
  \texttt{Pphi2.smoothVarianceConstant\_pos}
\end{theorem}
\begin{theorem}[\texttt{roughCovEigenvalue\_le\_T}]\label{Pphi2-roughCovEigenvalue-le-T}
  \lean{Pphi2.roughCovEigenvalue_le_T}\leanok
  \uses{Pphi2-roughCovEigenvalue, GaussianField-latticeEigenvalue-pos, GaussianField-latticeEigenvalue}
  $C_R(T, m) \le T$ from the elementary bound $(1 - e^{-x})/x \le 1$.
\end{theorem}
\begin{theorem}[\texttt{smoothVariance\_le\_log\_uniform}]\label{Pphi2-smoothVariance-le-log-uniform}
  \lean{Pphi2.smoothVariance_le_log_uniform}\leanok
  \uses{Pphi2-smoothWickConstant, GaussianField-FinLatticeSites, Pphi2-smoothVarianceConstant, GaussianField-latticeEigenvalue-pos, Pphi2-wickConstant, Pphi2-smoothCovEigenvalue, Pphi2-wickConstant-le-inv-mass-sq, GaussianField-latticeEigenvalue, Lattice-toSemilatticeInf}
  \texttt{Pphi2.smoothVariance\_le\_log\_uniform}
\end{theorem}
\begin{theorem}[\texttt{smoothCovEigenvalue\_pos}]\label{Pphi2-smoothCovEigenvalue-pos}
  \lean{Pphi2.smoothCovEigenvalue_pos}\leanok
  \uses{GaussianField-latticeEigenvalue-pos, Pphi2-smoothCovEigenvalue, GaussianField-latticeEigenvalue}
  Both parts are strictly positive for $T > 0$.
\end{theorem}
\begin{definition}[\texttt{smoothVarianceConstant}]\label{Pphi2-smoothVarianceConstant}
  \lean{Pphi2.smoothVarianceConstant}\leanok
  \texttt{Pphi2.smoothVarianceConstant}
\end{definition}
\begin{theorem}[\texttt{roughCovariance\_sq\_summable}]\label{Pphi2-roughCovariance-sq-summable}
  \lean{Pphi2.roughCovariance_sq_summable}\leanok
  \uses{Pphi2-roughCovEigenvalue, GaussianField-FinLatticeSites, GaussianField-latticeEigenvalue-pos, Pphi2-wickConstant, Pphi2-roughCovEigenvalue-le-T, GaussianField-latticeEigenvalue, Pphi2-roughCovEigenvalue-le-inv}
  $\frac{1}{|\Lambda|}\sum C_R(m)^2 \le T \cdot c_a$, since $C_R^2 \le T \cdot C_R \le T/\lambda_m$.
\end{theorem}
\begin{theorem}[\texttt{congr\_simp}]\label{Pphi2-wickConstant-congr-simp}
  \lean{Pphi2.wickConstant.congr_simp}\leanok
  \uses{Pphi2-wickConstant}
  \texttt{Pphi2.wickConstant.congr\_simp}
\end{theorem}
\begin{theorem}[\texttt{wickConstant\_split}]\label{Pphi2-wickConstant-split}
  \lean{Pphi2.wickConstant_split}\leanok
  \uses{Pphi2-smoothWickConstant, Pphi2-roughCovEigenvalue, GaussianField-FinLatticeSites, Pphi2-wickConstant, Pphi2-smoothCovEigenvalue, Pphi2-covariance-split, Pphi2-roughWickConstant, GaussianField-latticeEigenvalue}
  $c_a = c_S + c_R$.
\end{theorem}
\begin{definition}[\texttt{roughWickConstant}]\label{Pphi2-roughWickConstant}
  \lean{Pphi2.roughWickConstant}\leanok
  \uses{Pphi2-roughCovEigenvalue, GaussianField-FinLatticeSites}
  \texttt{Pphi2.roughWickConstant}
\end{definition}
\section{\texttt{Pphi2.NelsonEstimate.PolynomialChaosBridge}}
\begin{theorem}[\texttt{nelson\_exponential\_estimate\_master\_bounded} (axiom)]\label{Pphi2-nelson-exponential-estimate-master-bounded}
  \lean{Pphi2.nelson_exponential_estimate_master_bounded}\leanok
  \uses{GaussianField-FinLatticeSites, GaussianField-latticeGaussianMeasure, GaussianField-Configuration, GaussianField-instMeasurableSpaceConfiguration, GaussianField-FinLatticeField, Pphi2-interactionFunctional}
  \texttt{Pphi2.nelson\_exponential\_estimate\_master\_bounded}
\end{theorem}
\begin{theorem}[\texttt{canonicalRoughError\_cutoff\_tail\_quartic\_uniform}]\label{Pphi2-canonicalRoughError-cutoff-tail-quartic-uniform}
  \lean{Pphi2.canonicalRoughError_cutoff_tail_quartic_uniform}\leanok
  \uses{Pphi2-canonicalJointMeasure, Pphi2-DynamicalCutoff-dynamicalCutoffScale-pos, Pphi2-canonicalRoughError-neg-tail-uniform-in-aN, Pphi2-DynamicalCutoff-dynamicalCutoffScale, Pphi2-ChaosTailBridge-chaosTailConstant, Pphi2-canonicalRoughError, Pphi2-CanonicalJoint}
  \texttt{Pphi2.canonicalRoughError\_cutoff\_tail\_quartic\_uniform}
\end{theorem}
\begin{definition}[\texttt{degreePiecewiseTail}]\label{Pphi2-degreePiecewiseTail}
  \lean{Pphi2.degreePiecewiseTail}\leanok
  \uses{Pphi2-degreeCutoffTail}
  \texttt{Pphi2.degreePiecewiseTail}
\end{definition}
\begin{theorem}[\texttt{polynomial\_chaos\_exp\_moment\_bridge}]\label{Pphi2-polynomial-chaos-exp-moment-bridge}
  \lean{Pphi2.polynomial_chaos_exp_moment_bridge}\leanok
  \uses{GaussianField-FinLatticeSites, Pphi2-polynomial-chaos-exp-moment-bridge-bounded, GaussianField-latticeGaussianMeasure, GaussianField-Configuration, GaussianField-instMeasurableSpaceConfiguration, GaussianField-FinLatticeField, Pphi2-interactionFunctional}
  \texttt{Pphi2.polynomial\_chaos\_exp\_moment\_bridge}
\end{theorem}
\begin{theorem}[\texttt{canonicalRoughError\_cutoff\_tail\_uniform}]\label{Pphi2-canonicalRoughError-cutoff-tail-uniform}
  \lean{Pphi2.canonicalRoughError_cutoff_tail_uniform}\leanok
  \uses{Pphi2-degreeCutoffTail, Pphi2-canonicalJointMeasure, Pphi2-two-div-degree-eq-inv-cutoffPower, Pphi2-DynamicalCutoff-degreeDynamicalCutoffScale-pos, Pphi2-DynamicalCutoff-degreeDynamicalCutoffScale, Pphi2-canonicalRoughError-neg-tail-uniform-in-aN, Pphi2-ChaosTailBridge-chaosTailConstant, Pphi2-DynamicalCutoff-degreeCutoffPower, Pphi2-canonicalRoughError, Pphi2-CanonicalJoint}
  \texttt{Pphi2.canonicalRoughError\_cutoff\_tail\_uniform}
\end{theorem}
\begin{theorem}[\texttt{polynomial\_chaos\_exp\_moment\_bridge\_bounded\_of\_cutoffTail}]\label{Pphi2-polynomial-chaos-exp-moment-bridge-bounded-of-cutoffTail}
  \lean{Pphi2.polynomial_chaos_exp_moment_bridge_bounded_of_cutoffTail}\leanok
  \uses{Pphi2-degreeCutoffTail, Pphi2-canonicalJointMeasure, GaussianField-FinLatticeSites, Pphi2-canonicalSmoothInteraction, GaussianField-latticeGaussianMeasure, Pphi2-DynamicalCutoff-degreeDynamicalCutoffScale, GaussianField-Configuration, GaussianField-instMeasurableSpaceConfiguration, GaussianField-FinLatticeField, Pphi2-canonicalRoughError, Pphi2-polynomial-chaos-exp-moment-bridge-bounded-of-tail, Pphi2-interactionFunctional, Pphi2-degreePiecewiseTail, Pphi2-CanonicalJoint}
  \texttt{Pphi2.polynomial\_chaos\_exp\_moment\_bridge\_bounded\_of\_cutoffTail}
\end{theorem}
\begin{theorem}[\texttt{polynomial\_chaos\_exp\_moment\_bridge\_bounded}]\label{Pphi2-polynomial-chaos-exp-moment-bridge-bounded}
  \lean{Pphi2.polynomial_chaos_exp_moment_bridge_bounded}\leanok
  \uses{Pphi2-degreeCutoffTail, Pphi2-canonicalJointMeasure, GaussianField-FinLatticeSites, Pphi2-canonicalRoughError-cutoff-tail-uniform, Pphi2-canonicalSmoothInteraction, GaussianField-latticeGaussianMeasure, Pphi2-DynamicalCutoff-degreeDynamicalCutoffScale, GaussianField-Configuration, GaussianField-instMeasurableSpaceConfiguration, Pphi2-canonicalSmoothInteraction-lower-bound-at-cutoff-uniform, Pphi2-polynomial-chaos-exp-moment-bridge-bounded-of-cutoffTail, GaussianField-FinLatticeField, Pphi2-canonicalRoughError, Pphi2-interactionFunctional, Pphi2-CanonicalJoint, Pphi2-degreePiecewiseTail-layerCake-lt-top}
  \texttt{Pphi2.polynomial\_chaos\_exp\_moment\_bridge\_bounded}
\end{theorem}
\begin{theorem}[\texttt{two\_div\_degree\_eq\_inv\_cutoffPower}]\label{Pphi2-two-div-degree-eq-inv-cutoffPower}
  \lean{Pphi2.two_div_degree_eq_inv_cutoffPower}\leanok
  \uses{Pphi2-DynamicalCutoff-degreeCutoffPower}
  \texttt{Pphi2.two\_div\_degree\_eq\_inv\_cutoffPower}
\end{theorem}
\begin{definition}[\texttt{degreeCutoffTail}]\label{Pphi2-degreeCutoffTail}
  \lean{Pphi2.degreeCutoffTail}\leanok
  \uses{Pphi2-DynamicalCutoff-degreeDynamicalCutoffScale, Pphi2-ChaosTailBridge-chaosTailConstant, Pphi2-DynamicalCutoff-degreeCutoffPower}
  \texttt{Pphi2.degreeCutoffTail}
\end{definition}
\begin{theorem}[\texttt{degreePiecewiseTail\_layerCake\_lt\_top}]\label{Pphi2-degreePiecewiseTail-layerCake-lt-top}
  \lean{Pphi2.degreePiecewiseTail_layerCake_lt_top}\leanok
  \uses{Pphi2-degreeCutoffTail, Pphi2-DynamicalCutoff-one-add-abs-log-degreeDynamicalCutoff-eq-rpow, Pphi2-DynamicalCutoff-degreeDynamicalCutoffScale-eq-exp, Pphi2-DynamicalCutoff-degreeCutoffPower-pos, Pphi2-IntegrabilityHelpers-lintegral-layer-cake-lt-top-of-eventual-decay, Pphi2-DynamicalCutoff-degreeDynamicalCutoffScale-pos, Pphi2-ChaosTailBridge-chaosTailConstant-pos, Pphi2-DynamicalCutoff-degreeDynamicalCutoffScale, Pphi2-ChaosTailBridge-chaosTailConstant, Pphi2-DynamicalCutoff-degreeCutoffPower, Pphi2-degreePiecewiseTail, Lattice-toSemilatticeInf}
  \texttt{Pphi2.degreePiecewiseTail\_layerCake\_lt\_top}
\end{theorem}
\begin{theorem}[\texttt{quarticPiecewiseTail\_layerCake\_lt\_top}]\label{Pphi2-quarticPiecewiseTail-layerCake-lt-top}
  \lean{Pphi2.quarticPiecewiseTail_layerCake_lt_top}\leanok
  \uses{Pphi2-DynamicalCutoff-dynamicalCutoffScale-eq-exp, Pphi2-IntegrabilityHelpers-lintegral-layer-cake-lt-top-of-eventual-decay, Pphi2-ChaosTailBridge-chaosTailConstant-pos, Pphi2-DynamicalCutoff-dynamicalCutoffScale-pos, Pphi2-DynamicalCutoff-dynamicalCutoffScale, Pphi2-ChaosTailBridge-chaosTailConstant, Lattice-toSemilatticeInf}
  \texttt{Pphi2.quarticPiecewiseTail\_layerCake\_lt\_top}
\end{theorem}
\begin{theorem}[\texttt{nelson\_exponential\_estimate\_master}]\label{Pphi2-nelson-exponential-estimate-master}
  \lean{Pphi2.nelson_exponential_estimate_master}\leanok
  \uses{GaussianField-FinLatticeSites, Pphi2-polynomial-chaos-exp-moment-bridge, GaussianField-latticeGaussianMeasure, GaussianField-Configuration, GaussianField-instMeasurableSpaceConfiguration, GaussianField-FinLatticeField, Pphi2-interactionFunctional}
  \texttt{Pphi2.nelson\_exponential\_estimate\_master}
\end{theorem}
\begin{theorem}[\texttt{polynomial\_chaos\_exp\_moment\_bridge\_bounded\_of\_tail}]\label{Pphi2-polynomial-chaos-exp-moment-bridge-bounded-of-tail}
  \lean{Pphi2.polynomial_chaos_exp_moment_bridge_bounded_of_tail}\leanok
  \uses{Pphi2-canonicalJointMeasure, GaussianField-FinLatticeSites, Pphi2-canonicalSmoothInteraction, GaussianField-latticeGaussianMeasure, Pphi2-expMoment-bound-of-cutoff-degree-tail, Pphi2-DynamicalCutoff-degreeDynamicalCutoffScale, GaussianField-Configuration, GaussianField-instMeasurableSpaceConfiguration, GaussianField-FinLatticeField, Pphi2-canonicalRoughError, Pphi2-interactionFunctional, Pphi2-CanonicalJoint}
  \texttt{Pphi2.polynomial\_chaos\_exp\_moment\_bridge\_bounded\_of\_tail}
\end{theorem}
\section{\texttt{Pphi2.TorusContinuumLimit.TorusInteractingOS}}
\begin{theorem}[\texttt{torusInteractingMeasure\_exponentialMomentBound\_cutoff}]\label{Pphi2-torusInteractingMeasure-exponentialMomentBound-cutoff}
  \lean{Pphi2.torusInteractingMeasure_exponentialMomentBound_cutoff}\leanok
  \uses{Pphi2-torusEmbeddedTwoPoint-eq-lattice-second-moment, GaussianField-NuclearTensorProduct-instModuleReal, Pphi2-gaussian-exp-abs-moment, Pphi2-partitionFunction-ge-one, Lattice-toSemilatticeSup, GaussianField-NuclearTensorProduct-instTopologicalSpace, Pphi2-nelson-exponential-estimate, Pphi2-torusEmbeddedTwoPoint, Pphi2-interactionFunctional-bounded-below, GaussianField-FinLatticeSites, GaussianField-NuclearTensorProduct-instIsTopologicalAddGroup, Pphi2-torusEmbedLift, Pphi2-partitionFunction, Pphi2-interactionFunctional-measurable, Pphi2-torusEmbedLift-measurable, Pphi2-interactingLatticeMeasure, GaussianField-NuclearTensorProduct-instAddCommGroup, GaussianField-latticeGaussianMeasure, GaussianField-configuration-eval-measurable, GaussianField-Configuration, Pphi2-density-transfer-bound, GaussianField-NuclearTensorProduct-instContinuousSMulReal, GaussianField-SmoothMap-Circle, GaussianField-circleSpacing, GaussianField-instMeasurableSpaceConfiguration, Pphi2-boltzmannWeight, Pphi2-latticeTestFn, Pphi2-torusEmbedLift-eval-eq, Pphi2-torusInteractingMeasure, GaussianField-FinLatticeField, Pphi2-interactionFunctional, GaussianField-TorusTestFunction, Pphi2-partitionFunction-pos, GaussianField-circleSpacing-pos, Pphi2-boltzmannWeight-pos}
  Cutoff-level exponential moment bound: $\|Z_{\mathbb{C},N}[f_{\text{re}}, f_{\text{im}}]\| \le \exp(c \cdot (p(f_{\text{re}}) + p(f_{\text{im}})))$.
\end{theorem}
\begin{theorem}[\texttt{torusEmbeddedTwoPoint\_le\_seminorm}]\label{Pphi2-torusEmbeddedTwoPoint-le-seminorm}
  \lean{Pphi2.torusEmbeddedTwoPoint_le_seminorm}\leanok
  \uses{GaussianField-NuclearTensorProduct-instTopologicalSpace, GaussianField-RapidDecaySeq-instZero, GaussianField-RapidDecaySeq-rapidDecaySeminorm, GaussianField-RapidDecaySeq, Pphi2-torusEmbeddedTwoPoint, GaussianField-RapidDecaySeq-instSMul, GaussianField-RapidDecaySeq-instModule, GaussianField-NuclearTensorProduct-instAddCommGroup, GaussianField-RapidDecaySeq-instAddCommGroup, GaussianField-SmoothMap-Circle, GaussianField-RapidDecaySeq-instTopologicalSpace, GaussianField-RapidDecaySeq-rapidDecay-withSeminorms, Pphi2-torusEmbeddedTwoPoint-le-seminorm-tight, GaussianField-TorusTestFunction}
  $G_N(f,f) \le p(f)^2$ for a continuous seminorm $p$, uniformly in $N$.
\end{theorem}
\begin{theorem}[\texttt{torusInteracting\_satisfies\_OS}]\label{Pphi2-torusInteracting-satisfies-OS}
  \lean{Pphi2.torusInteracting_satisfies_OS}\leanok
  \uses{GaussianField-NuclearTensorProduct-instModuleReal, GaussianField-NuclearTensorProduct-instTopologicalSpace, Pphi2-torusInteracting-os1, GaussianField-NuclearTensorProduct-instIsTopologicalAddGroup, GaussianField-NuclearTensorProduct-instAddCommGroup, Pphi2-SatisfiesTorusOS, Pphi2-torusInteracting-os2-D4, GaussianField-Configuration, GaussianField-NuclearTensorProduct-instContinuousSMulReal, GaussianField-SmoothMap-Circle, Pphi2-torusInteracting-os2-translation, GaussianField-instMeasurableSpaceConfiguration, Pphi2-torusInteracting-os0, Pphi2-torusInteractingMeasure, GaussianField-TorusTestFunction}
  Bundled: the interacting continuum limit satisfies \texttt{SatisfiesTorusOS} (OS0 + OS1 + OS2).
\end{theorem}
\begin{theorem}[\texttt{interactingLatticeMeasure\_symmetry\_invariant}]\label{Pphi2-interactingLatticeMeasure-symmetry-invariant}
  \lean{Pphi2.interactingLatticeMeasure_symmetry_invariant}\leanok
  \uses{GaussianField-finLatticeDelta, GaussianField-evalMap, GaussianField-FinLatticeSites, Pphi2-partitionFunction, GaussianField-gaussianDensityMeasure, GaussianField-evalMapMeasurableEquiv, Pphi2-interactionFunctional-measurable, GaussianField-gaussianDensityNormConst, Pphi2-interactingLatticeMeasure, GaussianField-latticeGaussianFieldLaw, GaussianField-gaussianDensity-congr-simp, GaussianField-latticeGaussianFieldLaw-eq-normalizedGaussianDensityMeasure, GaussianField-latticeGaussianMeasure, GaussianField-gaussianDensityWeight-measurable, GaussianField-Configuration, GaussianField-gaussianDensityWeight, Pphi2-wickConstant, GaussianField-circleSpacing, GaussianField-instMeasurableSpaceConfiguration, GaussianField-normalizedGaussianDensityMeasure, Pphi2-boltzmannWeight, GaussianField-gaussianDensity, GaussianField-FinLatticeField, Pphi2-interactionFunctional, Pphi2-wickPolynomial}
  For a bijective site permutation $\sigma$ preserving the Gaussian density, $\int F(\omega \circ \sigma)\, d\mu_{\text{int}} = \int F(\omega)\, d\mu_{\text{int}}$. Proved via MeasurePreserving decomposition: evalME $\to$ piCongrLeft $\to$ evalME$^{-1}$.
\end{theorem}
\begin{theorem}[\texttt{torusInteracting\_os1}]\label{Pphi2-torusInteracting-os1}
  \lean{Pphi2.torusInteracting_os1}\leanok
  \uses{Pphi2-torusContinuumGreen, GaussianField-NuclearTensorProduct-instModuleReal, GaussianField-NuclearTensorProduct-instTopologicalSpace, Pphi2-torusContinuumGreen-continuous-diag, GaussianField-NuclearTensorProduct-instIsTopologicalAddGroup, Pphi2-torusContinuumGreen-nonneg, Pphi2-TorusOS1-Regularity, GaussianField-NuclearTensorProduct-instAddCommGroup, GaussianField-Configuration, Pphi2-torusInteracting-exponentialMomentBound, GaussianField-NuclearTensorProduct-instContinuousSMulReal, GaussianField-SmoothMap-Circle, GaussianField-instMeasurableSpaceConfiguration, Pphi2-torusGeneratingFunctional-, Pphi2-torusInteractingMeasure, GaussianField-TorusTestFunction, Lattice-toSemilatticeInf}
  OS1 regularity: $\|Z_{\mathbb{C}}[f_{\text{re}}, f_{\text{im}}]\| \le \exp(c(q(f_{\text{re}}) + q(f_{\text{im}})))$.
\end{theorem}
\begin{theorem}[\texttt{torusInteractingMeasure\_gf\_latticeTranslation\_invariant}]\label{Pphi2-torusInteractingMeasure-gf-latticeTranslation-invariant}
  \lean{Pphi2.torusInteractingMeasure_gf_latticeTranslation_invariant}\leanok
  \uses{GaussianField-NuclearTensorProduct-instModuleReal, GaussianField-circleTranslation, GaussianField-NuclearTensorProduct-instTopologicalSpace, GaussianField-nuclearTensorProduct-mapCLM, GaussianField-evalTorusAtSite-latticeTranslation, GaussianField-SmoothMap-Circle-instAddCommGroup, GaussianField-FinLatticeSites, GaussianField-NuclearTensorProduct-instIsTopologicalAddGroup, Pphi2-torusEmbedLift, Pphi2-torusGeneratingFunctional, Pphi2-torusEmbedLift-measurable, Pphi2-interactingLatticeMeasure, GaussianField-translateSites, GaussianField-NuclearTensorProduct-instAddCommGroup, GaussianField-torusTranslation, GaussianField-evalTorusAtSiteGJ-apply-, GaussianField-evalTorusAtSite, GaussianField-SmoothMap-Circle-instModule, GaussianField-configuration-eval-measurable, GaussianField-SmoothMap-Circle-instContinuousSMul, GaussianField-Configuration, GaussianField-smoothCircle-dyninMityaginSpace, GaussianField-SmoothMap-Circle-instIsTopologicalAddGroup, GaussianField-NuclearTensorProduct-instContinuousSMulReal, GaussianField-SmoothMap-Circle, GaussianField-circleSpacing, GaussianField-instMeasurableSpaceConfiguration, Pphi2-latticeTestFn, Pphi2-torusEmbedLift-eval-eq, Pphi2-torusInteractingMeasure-isProbability, Pphi2-torusInteractingMeasure, GaussianField-FinLatticeField, GaussianField-TorusTestFunction, GaussianField-SmoothMap-Circle-instTopologicalSpace, GaussianField-circleSpacing-pos}
  The generating functional of the interacting measure is invariant under lattice translations.
\end{theorem}
\begin{theorem}[\texttt{torusTranslation\_continuous\_in\_v}]\label{Pphi2-torusTranslation-continuous-in-v}
  \lean{Pphi2.torusTranslation_continuous_in_v}\leanok
  \uses{GaussianField-sobolevSeminorm-affine-precomp-le, GaussianField-mapImage, GaussianField-NuclearTensorProduct-instModuleReal, GaussianField-circleTranslation, Lattice-toSemilatticeSup, GaussianField-NuclearTensorProduct-instTopologicalSpace, GaussianField-SmoothMap-Circle-instFunLike, GaussianField-RapidDecaySeq-rapidDecaySeminorm, GaussianField-RapidDecaySeq, GaussianField-nuclearTensorProduct-mapCLM, GaussianField-RapidDecaySeq-val, GaussianField-RapidDecaySeq-hasSum-basisVec, GaussianField-SmoothMap-Circle-instAddCommGroup, GaussianField-NuclearTensorProduct-instIsTopologicalAddGroup, GaussianField-nuclearTensorProduct-mapCLM-pure, GaussianField-RapidDecaySeq-instSMul, GaussianField-RapidDecaySeq-instModule, GaussianField-DyninMityaginSpace--, GaussianField-NuclearTensorProduct, GaussianField-DyninMityaginSpace-HasBiorthogonalBasis-coeff-basis, GaussianField-RapidDecaySeq-basisVec, GaussianField-NuclearTensorProduct-instAddCommGroup, GaussianField-NuclearTensorProduct-pure-continuous, GaussianField-finset-sup-poly-bound, GaussianField-smoothCircle-hasBiorthogonalBasis, GaussianField-torusTranslation, GaussianField-SmoothMap-Circle-instModule, GaussianField-SmoothMap-Circle-instContinuousSMul, GaussianField-smoothCircle-dyninMityaginSpace, GaussianField-SmoothMap-Circle-instIsTopologicalAddGroup, GaussianField-DyninMityaginSpace-basis, GaussianField-NuclearTensorProduct-pure, GaussianField-RapidDecaySeq-instAddCommGroup, GaussianField-SmoothMap-Circle, GaussianField-DyninMityaginSpace-basis-growth, GaussianField-RapidDecaySeq-instTopologicalSpace, GaussianField-NuclearTensorProduct-basisVec-eq-pure, GaussianField-DyninMityaginSpace-p, GaussianField-RapidDecaySeq-rapidDecay-withSeminorms, GaussianField-RapidDecaySeq-rapid-decay, GaussianField-NuclearTensorProduct-pure-seminorm-bound, GaussianField-TorusTestFunction, GaussianField-SmoothMap-Circle-instTopologicalSpace, Lattice-toSemilatticeInf}
  \texttt{Pphi2.torusTranslation\_continuous\_in\_v}
\end{theorem}
\begin{theorem}[\texttt{torusInteracting\_os2\_translation}]\label{Pphi2-torusInteracting-os2-translation}
  \lean{Pphi2.torusInteracting_os2_translation}\leanok
  \uses{GaussianField-NuclearTensorProduct-instModuleReal, GaussianField-NuclearTensorProduct-instTopologicalSpace, GaussianField-NuclearTensorProduct-instIsTopologicalAddGroup, GaussianField-NuclearTensorProduct-instAddCommGroup, GaussianField-Configuration, GaussianField-NuclearTensorProduct-instContinuousSMulReal, GaussianField-SmoothMap-Circle, GaussianField-instMeasurableSpaceConfiguration, Pphi2-torusInteractingMeasure, GaussianField-TorusTestFunction, Pphi2-TorusOS2-TranslationInvariance, Pphi2-torusInteractingLimit-translation-invariant}
  OS2 translation invariance for the interacting limit.
\end{theorem}
\begin{theorem}[\texttt{gaussian\_exp\_abs\_moment}]\label{Pphi2-gaussian-exp-abs-moment}
  \lean{Pphi2.gaussian_exp_abs_moment}\leanok
  \uses{Lattice-toSemilatticeSup, GaussianField-finLatticeField-dyninMityaginSpace, GaussianField-measure, GaussianField-ell2-separableSpace, GaussianField-second-moment-eq-covariance, GaussianField-FinLatticeSites, GaussianField-latticeCovarianceGJ, GaussianField-pairing-is-gaussian, GaussianField-latticeGaussianMeasure, GaussianField-configuration-eval-measurable, GaussianField-Configuration, GaussianField-circleSpacing, GaussianField-instMeasurableSpaceConfiguration, GaussianField-ell2-, GaussianField-FinLatticeField, GaussianField-circleSpacing-pos}
  For the interacting measure, $\int e^{c|\omega(f)|}\, d\mu_P \le$ explicit Gaussian bound.
\end{theorem}
\begin{theorem}[\texttt{torusInteractingLimit\_translation\_invariant}]\label{Pphi2-torusInteractingLimit-translation-invariant}
  \lean{Pphi2.torusInteractingLimit_translation_invariant}\leanok
  \uses{GaussianField-NuclearTensorProduct-instModuleReal, GaussianField-NuclearTensorProduct-instTopologicalSpace, Pphi2-torusGF-latticeApproximation-error-vanishes, GaussianField-NuclearTensorProduct-instIsTopologicalAddGroup, Pphi2-torusGeneratingFunctional, GaussianField-NuclearTensorProduct-instAddCommGroup, GaussianField-torusTranslation, GaussianField-Configuration, GaussianField-NuclearTensorProduct-instContinuousSMulReal, GaussianField-SmoothMap-Circle, GaussianField-instMeasurableSpaceConfiguration, Pphi2-torusInteractingMeasure-isProbability, Pphi2-torusInteractingMeasure, GaussianField-TorusTestFunction}
  The generating functional of the interacting continuum limit is translation invariant.
\end{theorem}
\begin{theorem}[\texttt{torusInteractingMeasure\_gf\_swap\_invariant}]\label{Pphi2-torusInteractingMeasure-gf-swap-invariant}
  \lean{Pphi2.torusInteractingMeasure_gf_swap_invariant}\leanok
  \uses{GaussianField-evalTorusAtSite-swap, GaussianField-NuclearTensorProduct-instModuleReal, GaussianField-NuclearTensorProduct-instTopologicalSpace, GaussianField-SmoothMap-Circle-instAddCommGroup, GaussianField-FinLatticeSites, GaussianField-NuclearTensorProduct-instIsTopologicalAddGroup, Pphi2-torusEmbedLift, Pphi2-torusGeneratingFunctional, GaussianField-nuclearTensorProduct-swapCLM, GaussianField-torusSwap, Pphi2-torusEmbedLift-measurable, Pphi2-interactingLatticeMeasure, GaussianField-NuclearTensorProduct-instAddCommGroup, GaussianField-evalTorusAtSiteGJ-apply-, GaussianField-evalTorusAtSite, GaussianField-SmoothMap-Circle-instModule, GaussianField-configuration-eval-measurable, GaussianField-SmoothMap-Circle-instContinuousSMul, GaussianField-Configuration, GaussianField-smoothCircle-dyninMityaginSpace, GaussianField-SmoothMap-Circle-instIsTopologicalAddGroup, GaussianField-NuclearTensorProduct-instContinuousSMulReal, GaussianField-SmoothMap-Circle, GaussianField-circleSpacing, GaussianField-instMeasurableSpaceConfiguration, Pphi2-latticeTestFn, Pphi2-torusEmbedLift-eval-eq, GaussianField-swapSites, Pphi2-torusInteractingMeasure-isProbability, Pphi2-torusInteractingMeasure, GaussianField-FinLatticeField, GaussianField-TorusTestFunction, GaussianField-SmoothMap-Circle-instTopologicalSpace, GaussianField-circleSpacing-pos}
  \texttt{Pphi2.torusInteractingMeasure\_gf\_swap\_invariant}
\end{theorem}
\begin{theorem}[\texttt{congr\_simp}]\label{Pphi2-torusGeneratingFunctional-congr-simp}
  \lean{Pphi2.torusGeneratingFunctional.congr_simp}\leanok
  \uses{GaussianField-NuclearTensorProduct-instModuleReal, GaussianField-NuclearTensorProduct-instTopologicalSpace, GaussianField-NuclearTensorProduct-instIsTopologicalAddGroup, Pphi2-torusGeneratingFunctional, GaussianField-NuclearTensorProduct-instAddCommGroup, GaussianField-Configuration, GaussianField-NuclearTensorProduct-instContinuousSMulReal, GaussianField-SmoothMap-Circle, GaussianField-instMeasurableSpaceConfiguration, GaussianField-TorusTestFunction}
  \texttt{Pphi2.torusGeneratingFunctional.congr\_simp}
\end{theorem}
\begin{theorem}[\texttt{torusInteracting\_os0}]\label{Pphi2-torusInteracting-os0}
  \lean{Pphi2.torusInteracting_os0}\leanok
  \uses{Pphi2-torusContinuumGreen, GaussianField-NuclearTensorProduct-instModuleReal, GaussianField-NuclearTensorProduct-instTopologicalSpace, GaussianField-NuclearTensorProduct-instIsTopologicalAddGroup, Pphi2-TorusOS0-Analyticity, GaussianField-NuclearTensorProduct-instAddCommGroup, GaussianField-configuration-eval-measurable, GaussianField-Configuration, Pphi2-torusInteracting-exponentialMomentBound, GaussianField-NuclearTensorProduct-instContinuousSMulReal, GaussianField-SmoothMap-Circle, GaussianField-instMeasurableSpaceConfiguration, Pphi2-torusGeneratingFunctional-, Pphi2-torusInteractingMeasure, GaussianField-TorusTestFunction, Lattice-toSemilatticeInf}
  OS0 analyticity for the interacting continuum limit: $z \mapsto Z_{\mathbb{C}}[\sum z_i J_i]$ is entire.
\end{theorem}
\begin{theorem}[\texttt{torusInteracting\_os2\_D4}]\label{Pphi2-torusInteracting-os2-D4}
  \lean{Pphi2.torusInteracting_os2_D4}\leanok
  \uses{GaussianField-torusTimeReflection, GaussianField-NuclearTensorProduct-instModuleReal, GaussianField-NuclearTensorProduct-instTopologicalSpace, Pphi2-torusInteractingMeasure-gf-swap-invariant, GaussianField-NuclearTensorProduct-instIsTopologicalAddGroup, Pphi2-torusGeneratingFunctional, GaussianField-torusSwap, GaussianField-NuclearTensorProduct-instAddCommGroup, GaussianField-Configuration, GaussianField-NuclearTensorProduct-instContinuousSMulReal, Pphi2-TorusOS2-D4Invariance, GaussianField-SmoothMap-Circle, Pphi2-torusInteractingMeasure-gf-timeReflection-invariant, GaussianField-instMeasurableSpaceConfiguration, Pphi2-torusInteractingMeasure-isProbability, Pphi2-torusInteractingMeasure, GaussianField-TorusTestFunction}
  OS2 D4 point group invariance (swap + time reflection) for the interacting limit.
\end{theorem}
\begin{theorem}[\texttt{torus\_interacting\_second\_moment\_continuous}]\label{Pphi2-torus-interacting-second-moment-continuous}
  \lean{Pphi2.torus_interacting_second_moment_continuous}\leanok
  \uses{Pphi2-torusEmbeddedTwoPoint-eq-lattice-second-moment, GaussianField-NuclearTensorProduct-instModuleReal, Pphi2-partitionFunction-ge-one, GaussianField-NuclearTensorProduct-instTopologicalSpace, GaussianField-finLatticeField-dyninMityaginSpace, Pphi2-nelson-exponential-estimate, GaussianField-measure, GaussianField-ell2-separableSpace, Pphi2-torusEmbeddedTwoPoint, GaussianField-FinLatticeSites, GaussianField-NuclearTensorProduct-instIsTopologicalAddGroup, Pphi2-torusEmbedLift, Pphi2-partitionFunction, GaussianField-latticeCovarianceGJ, GaussianField-pairing-memLp, Pphi2-torusEmbedLift-measurable, Pphi2-interactingLatticeMeasure, GaussianField-NuclearTensorProduct-instAddCommGroup, GaussianField-latticeGaussianMeasure, GaussianField-configuration-eval-measurable, GaussianField-Configuration, Pphi2-density-transfer-bound, GaussianField-NuclearTensorProduct-instContinuousSMulReal, GaussianField-gaussian-hypercontractive, GaussianField-SmoothMap-Circle, GaussianField-circleSpacing, GaussianField-instMeasurableSpaceConfiguration, Pphi2-latticeTestFn, Pphi2-torusEmbedLift-eval-eq, GaussianField-ell2-, Pphi2-torusInteractingMeasure, GaussianField-FinLatticeField, Pphi2-interactionFunctional, GaussianField-TorusTestFunction, Lattice-toSemilatticeInf, GaussianField-circleSpacing-pos}
  $E_P[(\omega f)^2] \le C \cdot G_N(f,f)$ with $C = 3\sqrt{K}$ universal.
\end{theorem}
\begin{theorem}[\texttt{torusInteractingMeasure\_gf\_timeReflection\_invariant}]\label{Pphi2-torusInteractingMeasure-gf-timeReflection-invariant}
  \lean{Pphi2.torusInteractingMeasure_gf_timeReflection_invariant}\leanok
  \uses{GaussianField-torusTimeReflection, GaussianField-NuclearTensorProduct-instModuleReal, GaussianField-NuclearTensorProduct-instTopologicalSpace, GaussianField-evalTorusAtSite-timeReflection, GaussianField-nuclearTensorProduct-mapCLM, GaussianField-SmoothMap-Circle-instAddCommGroup, GaussianField-FinLatticeSites, GaussianField-NuclearTensorProduct-instIsTopologicalAddGroup, Pphi2-torusEmbedLift, Pphi2-torusGeneratingFunctional, Pphi2-torusEmbedLift-measurable, Pphi2-interactingLatticeMeasure, GaussianField-NuclearTensorProduct-instAddCommGroup, GaussianField-evalTorusAtSiteGJ-apply-, GaussianField-circleReflection, GaussianField-evalTorusAtSite, GaussianField-SmoothMap-Circle-instModule, GaussianField-configuration-eval-measurable, GaussianField-SmoothMap-Circle-instContinuousSMul, GaussianField-Configuration, GaussianField-smoothCircle-dyninMityaginSpace, GaussianField-SmoothMap-Circle-instIsTopologicalAddGroup, GaussianField-NuclearTensorProduct-instContinuousSMulReal, GaussianField-SmoothMap-Circle, GaussianField-circleSpacing, GaussianField-instMeasurableSpaceConfiguration, Pphi2-latticeTestFn, Pphi2-torusEmbedLift-eval-eq, Pphi2-torusInteractingMeasure-isProbability, Pphi2-torusInteractingMeasure, GaussianField-FinLatticeField, GaussianField-TorusTestFunction, GaussianField-SmoothMap-Circle-instTopologicalSpace, GaussianField-circleSpacing-pos, GaussianField-timeReflectSites}
  \texttt{Pphi2.torusInteractingMeasure\_gf\_timeReflection\_invariant}
\end{theorem}
\begin{theorem}[\texttt{torusGF\_latticeApproximation\_error\_vanishes}]\label{Pphi2-torusGF-latticeApproximation-error-vanishes}
  \lean{Pphi2.torusGF_latticeApproximation_error_vanishes}\leanok
  \uses{GaussianField-NuclearTensorProduct-instModuleReal, GaussianField-NuclearTensorProduct-instTopologicalSpace, GaussianField-NuclearTensorProduct-instIsTopologicalAddGroup, Pphi2-torusGeneratingFunctional, GaussianField-NuclearTensorProduct-instAddCommGroup, GaussianField-torusTranslation, GaussianField-SmoothMap-Circle, GaussianField-circleSpacing, Pphi2-torusTranslation-continuous-in-v, Pphi2-torusInteractingMeasure-gf-latticeTranslation-invariant, Pphi2-torusInteractingMeasure-isProbability, Pphi2-torusInteractingMeasure, GaussianField-TorusTestFunction, Lattice-toSemilatticeInf, GaussianField-circleSpacing-pos}
  The cutoff GF approximation error $|Z_N[T_v f] - Z_N[T_{w_n} f]| \to 0$ along the subsequence, where $w_n$ are nearest lattice vectors to $v$.
\end{theorem}
\begin{theorem}[\texttt{torusInteracting\_exponentialMomentBound}]\label{Pphi2-torusInteracting-exponentialMomentBound}
  \lean{Pphi2.torusInteracting_exponentialMomentBound}\leanok
  \uses{Pphi2-torus-propagator-convergence, Pphi2-torusContinuumGreen, GaussianField-NuclearTensorProduct-instModuleReal, Lattice-toSemilatticeSup, GaussianField-NuclearTensorProduct-instTopologicalSpace, Pphi2-torusEmbeddedTwoPoint, GaussianField-NuclearTensorProduct-instIsTopologicalAddGroup, GaussianField-NuclearTensorProduct-instAddCommGroup, Pphi2-torusInteractingMeasure-exponentialMomentBound-cutoff, GaussianField-configuration-eval-measurable, GaussianField-Configuration, GaussianField-NuclearTensorProduct-instContinuousSMulReal, GaussianField-SmoothMap-Circle, GaussianField-instMeasurableSpaceConfiguration, Pphi2-torusInteractingMeasure, GaussianField-TorusTestFunction, Lattice-toSemilatticeInf}
  The exponential moment bound passes to the continuum limit.
\end{theorem}
\begin{theorem}[\texttt{congr\_simp}]\label{Pphi2-torusInteractingMeasure-congr-simp}
  \lean{Pphi2.torusInteractingMeasure.congr_simp}\leanok
  \uses{GaussianField-NuclearTensorProduct-instModuleReal, GaussianField-NuclearTensorProduct-instTopologicalSpace, GaussianField-NuclearTensorProduct-instIsTopologicalAddGroup, GaussianField-NuclearTensorProduct-instAddCommGroup, GaussianField-Configuration, GaussianField-NuclearTensorProduct-instContinuousSMulReal, GaussianField-SmoothMap-Circle, GaussianField-instMeasurableSpaceConfiguration, Pphi2-torusInteractingMeasure, GaussianField-TorusTestFunction}
  \texttt{Pphi2.torusInteractingMeasure.congr\_simp}
\end{theorem}
\section{\texttt{Pphi2.InteractingMeasure.WickConstantBridge}}
\begin{theorem}[\texttt{wickConstant\_eq\_gffPositionCovariance}]\label{Pphi2-wickConstant-eq-gffPositionCovariance}
  \lean{Pphi2.wickConstant_eq_gffPositionCovariance}\leanok
  \uses{GaussianField-covariance, GaussianField-finLatticeDelta, GaussianField-FinLatticeSites, Pphi2-finLatticeDelta-eq-single, GaussianField-latticeCovarianceGJ, GaussianField-gffPositionCovariance, Pphi2-wickConstant-eq-variance, GaussianField-gffPositionCovariance-eq-covarianceGJ, Pphi2-wickConstant, GaussianField-ell2-, GaussianField-FinLatticeField}
  \texttt{Pphi2.wickConstant\_eq\_gffPositionCovariance}
\end{theorem}
\begin{theorem}[\texttt{gffSmearedCovariance\_single\_single\_eq\_wickConstant}]\label{Pphi2-gffSmearedCovariance-single-single-eq-wickConstant}
  \lean{Pphi2.gffSmearedCovariance_single_single_eq_wickConstant}\leanok
  \uses{GaussianField-FinLatticeSites, Pphi2-gffSmearedCovariance-single-single, GaussianField-gffPositionCovariance, Pphi2-wickConstant, Pphi2-wickConstant-eq-gffPositionCovariance, GaussianField-gffSmearedCovariance}
  \texttt{Pphi2.gffSmearedCovariance\_single\_single\_eq\_wickConstant}
\end{theorem}
\begin{theorem}[\texttt{finLatticeDelta\_eq\_single}]\label{Pphi2-finLatticeDelta-eq-single}
  \lean{Pphi2.finLatticeDelta_eq_single}\leanok
  \uses{GaussianField-finLatticeDelta, GaussianField-FinLatticeSites, GaussianField-FinLatticeField}
  \texttt{Pphi2.finLatticeDelta\_eq\_single}
\end{theorem}
\begin{theorem}[\texttt{gffSmearedCovariance\_single\_single}]\label{Pphi2-gffSmearedCovariance-single-single}
  \lean{Pphi2.gffSmearedCovariance_single_single}\leanok
  \uses{GaussianField-FinLatticeSites, GaussianField-gffPositionCovariance, GaussianField-gffSmearedCovariance, GaussianField-gffSmearedCovariance-single-right}
  \texttt{Pphi2.gffSmearedCovariance\_single\_single}
\end{theorem}
\section{\texttt{Pphi2.InteractingMeasure.MomentIntegrability}}
\begin{theorem}[\texttt{integrable\_powMul\_wickMonomial}]\label{Pphi2-integrable-powMul-wickMonomial}
  \lean{Pphi2.integrable_powMul_wickMonomial}\leanok
  \uses{GaussianField-FinLatticeSites, Pphi2-wickMonomial, GaussianField-latticeGaussianMeasure, GaussianField-Configuration, GaussianField-instMeasurableSpaceConfiguration, GaussianField-FinLatticeField}
  \texttt{Pphi2.integrable\_powMul\_wickMonomial}
\end{theorem}
\begin{theorem}[\texttt{integrable\_powMul\_wickMonomial\_mul}]\label{Pphi2-integrable-powMul-wickMonomial-mul}
  \lean{Pphi2.integrable_powMul_wickMonomial_mul}\leanok
  \uses{GaussianField-FinLatticeSites, Pphi2-wickMonomial, GaussianField-latticeGaussianMeasure, GaussianField-Configuration, GaussianField-instMeasurableSpaceConfiguration, GaussianField-FinLatticeField}
  \texttt{Pphi2.integrable\_powMul\_wickMonomial\_mul}
\end{theorem}
\begin{theorem}[\texttt{integrable\_pow\_pairing}]\label{Pphi2-integrable-pow-pairing}
  \lean{Pphi2.integrable_pow_pairing}\leanok
  \uses{GaussianField-FinLatticeSites, GaussianField-latticeGaussianMeasure-isProbability, Pphi2-pairing-memLp-lattice, GaussianField-latticeGaussianMeasure, GaussianField-configuration-eval-measurable, GaussianField-Configuration, GaussianField-instMeasurableSpaceConfiguration, GaussianField-FinLatticeField}
  \texttt{Pphi2.integrable\_pow\_pairing}
\end{theorem}
\begin{theorem}[\texttt{integrable\_powMul\_wickPolynomial\_sq}]\label{Pphi2-integrable-powMul-wickPolynomial-sq}
  \lean{Pphi2.integrable_powMul_wickPolynomial_sq}\leanok
  \uses{GaussianField-FinLatticeSites, GaussianField-latticeGaussianMeasure, GaussianField-Configuration, GaussianField-instMeasurableSpaceConfiguration, Pphi2-integrable-powMul-wickPolynomial-mul-self, GaussianField-FinLatticeField, Pphi2-wickPolynomial}
  \texttt{Pphi2.integrable\_powMul\_wickPolynomial\_sq}
\end{theorem}
\begin{theorem}[\texttt{integrable\_powMul\_interaction}]\label{Pphi2-integrable-powMul-interaction}
  \lean{Pphi2.integrable_powMul_interaction}\leanok
  \uses{Pphi2-integrable-powMul-wickPolynomial, GaussianField-FinLatticeSites, GaussianField-latticeGaussianMeasure, GaussianField-Configuration, Pphi2-wickConstant, GaussianField-instMeasurableSpaceConfiguration, GaussianField-FinLatticeField, Pphi2-wickPolynomial}
  \texttt{Pphi2.integrable\_powMul\_interaction}
\end{theorem}
\begin{theorem}[\texttt{integrable\_powMul\_wickPolynomial\_mul}]\label{Pphi2-integrable-powMul-wickPolynomial-mul}
  \lean{Pphi2.integrable_powMul_wickPolynomial_mul}\leanok
  \uses{GaussianField-FinLatticeSites, Pphi2-integrable-powMul-wickPolynomial-sq, GaussianField-latticeGaussianMeasure, Pphi2-wickPolynomial-continuous-, GaussianField-configuration-eval-measurable, GaussianField-Configuration, GaussianField-instMeasurableSpaceConfiguration, GaussianField-FinLatticeField, Lattice-toSemilatticeInf, Pphi2-wickPolynomial}
  \texttt{Pphi2.integrable\_powMul\_wickPolynomial\_mul}
\end{theorem}
\begin{theorem}[\texttt{integrable\_powMul\_wickPolynomial\_mul\_self}]\label{Pphi2-integrable-powMul-wickPolynomial-mul-self}
  \lean{Pphi2.integrable_powMul_wickPolynomial_mul_self}\leanok
  \uses{Pphi2-integrable-powMul-wickMonomial-mulWickPoly, GaussianField-FinLatticeSites, Pphi2-wickMonomial, GaussianField-latticeGaussianMeasure, GaussianField-Configuration, GaussianField-instMeasurableSpaceConfiguration, GaussianField-FinLatticeField, Pphi2-wickPolynomial}
  \texttt{Pphi2.integrable\_powMul\_wickPolynomial\_mul\_self}
\end{theorem}
\begin{theorem}[\texttt{integrable\_powMul\_wickMonomial\_mulWickPoly}]\label{Pphi2-integrable-powMul-wickMonomial-mulWickPoly}
  \lean{Pphi2.integrable_powMul_wickMonomial_mulWickPoly}\leanok
  \uses{GaussianField-FinLatticeSites, Pphi2-wickMonomial, GaussianField-latticeGaussianMeasure, GaussianField-Configuration, GaussianField-instMeasurableSpaceConfiguration, Pphi2-integrable-powMul-wickMonomial-mul, GaussianField-FinLatticeField, Pphi2-wickPolynomial}
  \texttt{Pphi2.integrable\_powMul\_wickMonomial\_mulWickPoly}
\end{theorem}
\begin{theorem}[\texttt{integrable\_powMul\_wickPolynomial}]\label{Pphi2-integrable-powMul-wickPolynomial}
  \lean{Pphi2.integrable_powMul_wickPolynomial}\leanok
  \uses{GaussianField-FinLatticeSites, Pphi2-wickMonomial, GaussianField-latticeGaussianMeasure, GaussianField-Configuration, Pphi2-integrable-powMul-wickMonomial-, GaussianField-instMeasurableSpaceConfiguration, GaussianField-FinLatticeField, Pphi2-wickPolynomial}
  \texttt{Pphi2.integrable\_powMul\_wickPolynomial}
\end{theorem}
\begin{theorem}[\texttt{integrable\_powMul\_wickMonomial'}]\label{Pphi2-integrable-powMul-wickMonomial-}
  \lean{Pphi2.integrable_powMul_wickMonomial'}\leanok
  \uses{GaussianField-FinLatticeSites, Pphi2-wickMonomial, GaussianField-latticeGaussianMeasure, GaussianField-Configuration, GaussianField-instMeasurableSpaceConfiguration, Pphi2-integrable-powMul-wickMonomial, GaussianField-FinLatticeField}
  \texttt{Pphi2.integrable\_powMul\_wickMonomial'}
\end{theorem}
\begin{theorem}[\texttt{integrable\_pow\_pairing\_mul}]\label{Pphi2-integrable-pow-pairing-mul}
  \lean{Pphi2.integrable_pow_pairing_mul}\leanok
  \uses{GaussianField-FinLatticeSites, Pphi2-integrable-pow-pairing, GaussianField-latticeGaussianMeasure, GaussianField-configuration-eval-measurable, GaussianField-Configuration, GaussianField-instMeasurableSpaceConfiguration, GaussianField-FinLatticeField}
  \texttt{Pphi2.integrable\_pow\_pairing\_mul}
\end{theorem}
\begin{theorem}[\texttt{pairing\_memLp\_lattice}]\label{Pphi2-pairing-memLp-lattice}
  \lean{Pphi2.pairing_memLp_lattice}\leanok
  \uses{GaussianField-finLatticeField-dyninMityaginSpace, GaussianField-ell2-separableSpace, GaussianField-FinLatticeSites, GaussianField-latticeCovarianceGJ, GaussianField-pairing-memLp, GaussianField-latticeGaussianMeasure, GaussianField-Configuration, GaussianField-instMeasurableSpaceConfiguration, GaussianField-ell2-, GaussianField-FinLatticeField}
  \texttt{Pphi2.pairing\_memLp\_lattice}
\end{theorem}
\section{\texttt{Pphi2.NelsonEstimate.HeatKernelBound}}
\begin{theorem}[\texttt{heat\_kernel\_1d\_bound}]\label{Pphi2-heat-kernel-1d-bound}
  \lean{Pphi2.heat_kernel_1d_bound}\leanok
  \uses{Lattice-toSemilatticeInf}
  1D: $\sum_k e^{-t \cdot 4\sin^2(\pi k/N)/a^2} \le C(1 + 1/\sqrt{t})$. Full: $\sum_m e^{-t\lambda_m} \le C/t$ with $C = |\Lambda|/m^2$.
\end{theorem}
\begin{theorem}[\texttt{heat\_kernel\_trace\_bound}]\label{Pphi2-heat-kernel-trace-bound}
  \lean{Pphi2.heat_kernel_trace_bound}\leanok
  \uses{GaussianField-latticeEigenvalue-eq, GaussianField-FinLatticeSites, GaussianField-latticeLaplacianEigenvalue, GaussianField-latticeLaplacianEigenvalue-nonneg, GaussianField-latticeEigenvalue-pos, GaussianField-latticeEigenvalue}
  \texttt{Pphi2.heat\_kernel\_trace\_bound}
\end{theorem}
\begin{theorem}[\texttt{heat\_kernel\_1d\_bound\_uniform}]\label{Pphi2-heat-kernel-1d-bound-uniform}
  \lean{Pphi2.heat_kernel_1d_bound_uniform}\leanok
  \uses{Lattice-toSemilatticeInf}
  \texttt{Pphi2.heat\_kernel\_1d\_bound\_uniform}
\end{theorem}
\begin{theorem}[\texttt{gaussian\_sum\_bound}]\label{Pphi2-gaussian-sum-bound}
  \lean{Pphi2.gaussian_sum_bound}\leanok
  $\sum_{k=-M}^M e^{-\alpha k^2} \le 1 + \sqrt{\pi/\alpha}$. Proved by splitting into $k=0$, positive, and negative parts, using the Gaussian integral $\int_0^\infty e^{-\alpha x^2}\,dx = \sqrt{\pi/\alpha}/2$.
\end{theorem}
\begin{theorem}[\texttt{schwinger\_smooth}]\label{Pphi2-schwinger-smooth}
  \lean{Pphi2.schwinger_smooth}\leanok
  \uses{Pphi2-schwinger-smooth-Ioi}
  Schwinger parametrization: $e^{-T\lambda}/\lambda = \int_T^\infty e^{-t\lambda}\,dt$ and $(1 - e^{-T\lambda})/\lambda = \int_0^T e^{-t\lambda}\,dt$. Proved via FTC with the antiderivative $-e^{-t\lambda}/\lambda$.
\end{theorem}
\begin{theorem}[\texttt{roughVariance\_from\_heat\_kernel}]\label{Pphi2-roughVariance-from-heat-kernel}
  \lean{Pphi2.roughVariance_from_heat_kernel}\leanok
  \uses{GaussianField-latticeEigenvalue-eq, Pphi2-roughCovEigenvalue, GaussianField-FinLatticeSites, GaussianField-latticeLaplacianEigenvalue, GaussianField-latticeLaplacianEigenvalue-nonneg, Pphi2-roughCovEigenvalue-pos, Pphi2-roughCovEigenvalue-le-T, GaussianField-latticeEigenvalue, Pphi2-roughCovEigenvalue-le-inv}
  \texttt{Pphi2.roughVariance\_from\_heat\_kernel}
\end{theorem}
\begin{theorem}[\texttt{smoothVariance\_from\_heat\_kernel}]\label{Pphi2-smoothVariance-from-heat-kernel}
  \lean{Pphi2.smoothVariance_from_heat_kernel}\leanok
  \uses{GaussianField-latticeEigenvalue-eq, GaussianField-FinLatticeSites, GaussianField-latticeLaplacianEigenvalue, GaussianField-latticeLaplacianEigenvalue-nonneg, GaussianField-latticeEigenvalue-pos, Pphi2-smoothCovEigenvalue, GaussianField-latticeEigenvalue, Lattice-toSemilatticeInf}
  $\sum C_S(m) \le C(1 + |\log T|)$ and $\sum C_R(m)^2 \le CT$ derived from the heat kernel bound.
\end{theorem}
\begin{theorem}[\texttt{schwinger\_smooth\_Ioi}]\label{Pphi2-schwinger-smooth-Ioi}
  \lean{Pphi2.schwinger_smooth_Ioi}\leanok
  \uses{Lattice-toSemilatticeInf}
  \texttt{Pphi2.schwinger\_smooth\_Ioi}
\end{theorem}
\begin{theorem}[\texttt{heat\_kernel\_trace\_bound\_uniform}]\label{Pphi2-heat-kernel-trace-bound-uniform}
  \lean{Pphi2.heat_kernel_trace_bound_uniform}\leanok
  \uses{GaussianField-latticeEigenvalue-eq, GaussianField-FinLatticeSites, GaussianField-latticeLaplacianEigenvalue, Pphi2-heat-kernel-1d-bound-uniform, GaussianField-latticeEigenvalue}
  \texttt{Pphi2.heat\_kernel\_trace\_bound\_uniform}
\end{theorem}
\begin{theorem}[\texttt{schwinger\_rough\_interval}]\label{Pphi2-schwinger-rough-interval}
  \lean{Pphi2.schwinger_rough_interval}\leanok
  \texttt{Pphi2.schwinger\_rough\_interval}
\end{theorem}
\begin{theorem}[\texttt{smoothWickConstant\_le\_log\_uniform\_in\_aN\_proved}]\label{Pphi2-smoothWickConstant-le-log-uniform-in-aN-proved}
  \lean{Pphi2.smoothWickConstant_le_log_uniform_in_aN_proved}\leanok
  \uses{Pphi2-smoothWickConstant, GaussianField-FinLatticeSites, Pphi2-schwinger-smooth-Ioi, Pphi2-heat-kernel-trace-bound-uniform, GaussianField-latticeEigenvalue-pos, Pphi2-smoothCovEigenvalue, GaussianField-latticeEigenvalue, Lattice-toSemilatticeInf}
  \texttt{Pphi2.smoothWickConstant\_le\_log\_uniform\_in\_aN\_proved}
\end{theorem}
\begin{theorem}[\texttt{sin\_sq\_lower\_bound}]\label{Pphi2-sin-sq-lower-bound}
  \lean{Pphi2.sin_sq_lower_bound}\leanok
  \uses{Lattice-toSemilatticeInf}
  $(2/\pi)^2 x^2 \le \sin^2 x$ for $|x| \le \pi/2$, from Mathlib's Jordan inequality.
\end{theorem}
\begin{theorem}[\texttt{schwinger\_rough}]\label{Pphi2-schwinger-rough}
  \lean{Pphi2.schwinger_rough}\leanok
  \uses{Pphi2-schwinger-rough-interval}
  \texttt{Pphi2.schwinger\_rough}
\end{theorem}
\section{\texttt{Pphi2.NelsonEstimate.AsymFieldDecomposition}}
\begin{theorem}[\texttt{congr\_simp}]\label{Pphi2-asymCanonicalSumFieldFunction-congr-simp}
  \lean{Pphi2.asymCanonicalSumFieldFunction.congr_simp}\leanok
  \uses{GaussianField-AsymLatticeSites, Pphi2-AsymCanonicalJoint, Pphi2-asymCanonicalSumFieldFunction}
  \texttt{Pphi2.asymCanonicalSumFieldFunction.congr\_simp}
\end{theorem}
\begin{theorem}[\texttt{asymCanonicalSmoothFieldFunction\_smul}]\label{Pphi2-asymCanonicalSmoothFieldFunction-smul}
  \lean{Pphi2.asymCanonicalSmoothFieldFunction_smul}\leanok
  \uses{GaussianField-AsymLatticeSites, Pphi2-AsymCanonicalJoint, Pphi2-asymCanonicalSmoothModeCoeff, Pphi2-asymCanonicalSmoothFieldFunction, Pphi2-AsymCanonicalMode}
  \texttt{Pphi2.asymCanonicalSmoothFieldFunction\_smul}
\end{theorem}
\begin{theorem}[\texttt{congr\_simp}]\label{Pphi2-asymEvalMapInv-congr-simp}
  \lean{Pphi2.asymEvalMapInv.congr_simp}\leanok
  \uses{GaussianField-AsymLatticeSites, GaussianField-Configuration, GaussianField-AsymLatticeField, Pphi2-asymEvalMapInv}
  \texttt{Pphi2.asymEvalMapInv.congr\_simp}
\end{theorem}
\begin{theorem}[\texttt{asymField\_basis\_decomposition}]\label{Pphi2-asymField-basis-decomposition}
  \lean{Pphi2.asymField_basis_decomposition}\leanok
  \uses{GaussianField-AsymLatticeSites, GaussianField-AsymLatticeField, GaussianField-instFintypeAsymLatticeSitesOfNeZeroNat, GaussianField-asymLatticeDelta}
  \texttt{Pphi2.asymField\_basis\_decomposition}
\end{theorem}
\begin{definition}[\texttt{asymCanonicalSumFieldLaw}]\label{Pphi2-asymCanonicalSumFieldLaw}
  \lean{Pphi2.asymCanonicalSumFieldLaw}\leanok
  \uses{GaussianField-AsymLatticeSites, Pphi2-AsymCanonicalJointSumIndex, GaussianField-AsymLatticeField, Pphi2-asymCanonicalStdToFieldCLM, Pphi2-AsymCanonicalMode}
  \texttt{Pphi2.asymCanonicalSumFieldLaw}
\end{definition}
\begin{definition}[\texttt{AsymCanonicalJoint}]\label{Pphi2-AsymCanonicalJoint}
  \lean{Pphi2.AsymCanonicalJoint}\leanok
  \uses{Pphi2-AsymCanonicalMode}
  \texttt{Pphi2.AsymCanonicalJoint}
\end{definition}
\begin{theorem}[\texttt{asymCanonicalSumFieldFunction\_memLp\_two}]\label{Pphi2-asymCanonicalSumFieldFunction-memLp-two}
  \lean{Pphi2.asymCanonicalSumFieldFunction_memLp_two}\leanok
  \uses{GaussianField-AsymLatticeSites, Pphi2-AsymCanonicalJoint, Pphi2-asymCanonicalSumFieldFunction, Pphi2-asymCanonicalRoughFieldFunction, Pphi2-asymCanonicalSmoothFieldFunction-memLp-two, Pphi2-asymCanonicalRoughFieldFunction-memLp-two, Pphi2-asymCanonicalSmoothFieldFunction, Pphi2-asymCanonicalJointMeasure, Pphi2-AsymCanonicalMode}
  \texttt{Pphi2.asymCanonicalSumFieldFunction\_memLp\_two}
\end{theorem}
\begin{theorem}[\texttt{asymCanonicalJointMeasure\_map\_sumConfig}]\label{Pphi2-asymCanonicalJointMeasure-map-sumConfig}
  \lean{Pphi2.asymCanonicalJointMeasure_map_sumConfig}\leanok
  \uses{Pphi2-asymLatticeFieldToConfig-eq-evalMapInv, GaussianField-AsymLatticeSites, Pphi2-asymCanonicalSumFieldLaw, Pphi2-asymCanonicalSumFieldLaw-eq-asymLatticeGaussianFieldLaw, Pphi2-AsymCanonicalJoint, Pphi2-asymCanonicalSumFieldLaw-eq-map-asymCanonicalSumFieldFunction, Pphi2-asymEvalMapInv-measurable, Pphi2-asymEvalMap-measurable, Pphi2-asymCanonicalSumFieldFunction, Pphi2-asymCanonicalSumConfig, GaussianField-Configuration, GaussianField-AsymLatticeField, Pphi2-latticeGaussianMeasureAsym, GaussianField-instMeasurableSpaceConfiguration, Pphi2-asymEvalMap, Pphi2-asymLatticeFieldToConfig, Pphi2-asymLatticeGaussianFieldLaw, Pphi2-asymCanonicalSumFieldFunction-measurable, Pphi2-asymCanonicalJointMeasure, Pphi2-asymEvalMapInv-evalMap, Pphi2-asymEvalMapInv, Pphi2-AsymCanonicalMode}
  \texttt{Pphi2.asymCanonicalJointMeasure\_map\_sumConfig}
\end{theorem}
\begin{definition}[\texttt{asymLatticeFieldToConfig}]\label{Pphi2-asymLatticeFieldToConfig}
  \lean{Pphi2.asymLatticeFieldToConfig}\leanok
  \uses{GaussianField-AsymLatticeSites, GaussianField-Configuration, GaussianField-AsymLatticeField, GaussianField-instFintypeAsymLatticeSitesOfNeZeroNat}
  \texttt{Pphi2.asymLatticeFieldToConfig}
\end{definition}
\begin{definition}[\texttt{asymCanonicalBasis}]\label{Pphi2-asymCanonicalBasis}
  \lean{Pphi2.asymCanonicalBasis}\leanok
  \uses{GaussianField-AsymLatticeSites, GaussianField-latticeFourierBasisFun, Pphi2-AsymCanonicalMode}
  \texttt{Pphi2.asymCanonicalBasis}
\end{definition}
\begin{theorem}[\texttt{asymCanonicalSumFieldFunction\_add}]\label{Pphi2-asymCanonicalSumFieldFunction-add}
  \lean{Pphi2.asymCanonicalSumFieldFunction_add}\leanok
  \uses{GaussianField-AsymLatticeSites, Pphi2-AsymCanonicalJoint, Pphi2-asymCanonicalSumFieldFunction, Pphi2-asymCanonicalRoughFieldFunction, Pphi2-asymCanonicalRoughFieldFunction-add, Pphi2-asymCanonicalSmoothFieldFunction-add, Pphi2-asymCanonicalSmoothFieldFunction, Pphi2-AsymCanonicalMode}
  \texttt{Pphi2.asymCanonicalSumFieldFunction\_add}
\end{theorem}
\begin{theorem}[\texttt{congr\_simp}]\label{Pphi2-asymCanonicalSmoothModeCoeff-congr-simp}
  \lean{Pphi2.asymCanonicalSmoothModeCoeff.congr_simp}\leanok
  \uses{GaussianField-AsymLatticeSites, Pphi2-asymCanonicalSmoothModeCoeff, Pphi2-AsymCanonicalMode}
  \texttt{Pphi2.asymCanonicalSmoothModeCoeff.congr\_simp}
\end{theorem}
\begin{definition}[\texttt{asymSitePairingCLM}]\label{Pphi2-asymSitePairingCLM}
  \lean{Pphi2.asymSitePairingCLM}\leanok
  \uses{GaussianField-AsymLatticeSites, GaussianField-AsymLatticeField, GaussianField-instFintypeAsymLatticeSitesOfNeZeroNat}
  \texttt{Pphi2.asymSitePairingCLM}
\end{definition}
\begin{theorem}[\texttt{asymCanonicalSmoothWeight\_add\_roughWeight}]\label{Pphi2-asymCanonicalSmoothWeight-add-roughWeight}
  \lean{Pphi2.asymCanonicalSmoothWeight_add_roughWeight}\leanok
  \uses{Pphi2-asymCanonicalEigenvalue, Pphi2-asymCanonicalRoughWeight, Pphi2-asymCanonicalSmoothWeight, Pphi2-AsymCanonicalMode}
  \texttt{Pphi2.asymCanonicalSmoothWeight\_add\_roughWeight}
\end{theorem}
\begin{theorem}[\texttt{asymCanonicalSumFieldLaw\_isGaussian}]\label{Pphi2-asymCanonicalSumFieldLaw-isGaussian}
  \lean{Pphi2.asymCanonicalSumFieldLaw_isGaussian}\leanok
  \uses{GaussianField-AsymLatticeSites, Pphi2-asymCanonicalSumFieldLaw, Pphi2-AsymCanonicalJointSumIndex, GaussianField-AsymLatticeField, GaussianField-instFintypeAsymLatticeSitesOfNeZeroNat, Pphi2-asymCanonicalStdToFieldCLM, Pphi2-AsymCanonicalMode}
  \texttt{Pphi2.asymCanonicalSumFieldLaw\_isGaussian}
\end{theorem}
\begin{definition}[\texttt{asymCanonicalRoughWeight}]\label{Pphi2-asymCanonicalRoughWeight}
  \lean{Pphi2.asymCanonicalRoughWeight}\leanok
  \uses{Pphi2-asymCanonicalEigenvalue, Pphi2-AsymCanonicalMode}
  \texttt{Pphi2.asymCanonicalRoughWeight}
\end{definition}
\begin{theorem}[\texttt{congr\_simp}]\label{Pphi2-asymCanonicalNormSq-congr-simp}
  \lean{Pphi2.asymCanonicalNormSq.congr_simp}\leanok
  \uses{Pphi2-AsymCanonicalMode, Pphi2-asymCanonicalNormSq}
  \texttt{Pphi2.asymCanonicalNormSq.congr\_simp}
\end{theorem}
\begin{theorem}[\texttt{asymCanonicalRoughFieldFunction\_memLp\_two}]\label{Pphi2-asymCanonicalRoughFieldFunction-memLp-two}
  \lean{Pphi2.asymCanonicalRoughFieldFunction_memLp_two}\leanok
  \uses{GaussianField-AsymLatticeSites, Pphi2-asymCanonicalRoughModeCoeff, Pphi2-AsymCanonicalJoint, Pphi2-asymCanonicalRoughFieldFunction, Pphi2-asymCanonicalJointMeasure, Pphi2-AsymCanonicalMode}
  \texttt{Pphi2.asymCanonicalRoughFieldFunction\_memLp\_two}
\end{theorem}
\begin{theorem}[\texttt{asymCanonicalStdFieldFunction\_eq\_asymCanonicalSumFieldFunction}]\label{Pphi2-asymCanonicalStdFieldFunction-eq-asymCanonicalSumFieldFunction}
  \lean{Pphi2.asymCanonicalStdFieldFunction_eq_asymCanonicalSumFieldFunction}\leanok
  \uses{GaussianField-AsymLatticeSites, Pphi2-asymCanonicalRoughModeCoeff, Pphi2-AsymCanonicalJointSumIndex, Pphi2-AsymCanonicalJoint, Pphi2-asymCanonicalSmoothModeCoeff, Pphi2-asymCanonicalSumFieldFunction, Pphi2-asymCanonicalRoughFieldFunction, GaussianField-AsymLatticeField, Pphi2-asymCanonicalStdFieldFunction, Pphi2-asymCanonicalSmoothFieldFunction, Pphi2-AsymCanonicalMode}
  \texttt{Pphi2.asymCanonicalStdFieldFunction\_eq\_asymCanonicalSumFieldFunction}
\end{theorem}
\begin{theorem}[\texttt{asymConfig\_apply\_eq\_sum\_delta}]\label{Pphi2-asymConfig-apply-eq-sum-delta}
  \lean{Pphi2.asymConfig_apply_eq_sum_delta}\leanok
  \uses{GaussianField-AsymLatticeSites, Pphi2-asymField-basis-decomposition, GaussianField-Configuration, GaussianField-AsymLatticeField, GaussianField-instFintypeAsymLatticeSitesOfNeZeroNat, GaussianField-asymLatticeDelta}
  \texttt{Pphi2.asymConfig\_apply\_eq\_sum\_delta}
\end{theorem}
\begin{definition}[\texttt{AsymCanonicalJointSumIndex}]\label{Pphi2-AsymCanonicalJointSumIndex}
  \lean{Pphi2.AsymCanonicalJointSumIndex}\leanok
  \uses{Pphi2-AsymCanonicalMode}
  \texttt{Pphi2.AsymCanonicalJointSumIndex}
\end{definition}
\begin{theorem}[\texttt{asymCanonicalSmoothFieldFunction\_measurable}]\label{Pphi2-asymCanonicalSmoothFieldFunction-measurable}
  \lean{Pphi2.asymCanonicalSmoothFieldFunction_measurable}\leanok
  \uses{Pphi2-asymCanonicalSmoothFieldFunction-pointwise-measurable, GaussianField-AsymLatticeSites, Pphi2-AsymCanonicalJoint, GaussianField-AsymLatticeField, Pphi2-asymCanonicalSmoothFieldFunction, Pphi2-AsymCanonicalMode}
  \texttt{Pphi2.asymCanonicalSmoothFieldFunction\_measurable}
\end{theorem}
\begin{definition}[\texttt{asymCanonicalEigenvalue}]\label{Pphi2-asymCanonicalEigenvalue}
  \lean{Pphi2.asymCanonicalEigenvalue}\leanok
  \uses{GaussianField-latticeEigenvalue1d, Pphi2-AsymCanonicalMode}
  \texttt{Pphi2.asymCanonicalEigenvalue}
\end{definition}
\begin{theorem}[\texttt{asymEvalMap\_evalMapInv}]\label{Pphi2-asymEvalMap-evalMapInv}
  \lean{Pphi2.asymEvalMap_evalMapInv}\leanok
  \uses{GaussianField-AsymLatticeSites, GaussianField-Configuration, GaussianField-AsymLatticeField, Pphi2-asymEvalMapInv-apply, GaussianField-instFintypeAsymLatticeSitesOfNeZeroNat, Pphi2-asymEvalMap, GaussianField-asymLatticeDelta, Pphi2-asymEvalMapInv}
  \texttt{Pphi2.asymEvalMap\_evalMapInv}
\end{theorem}
\begin{definition}[\texttt{asymEvalMapInv}]\label{Pphi2-asymEvalMapInv}
  \lean{Pphi2.asymEvalMapInv}\leanok
  \uses{GaussianField-AsymLatticeSites, GaussianField-Configuration, GaussianField-AsymLatticeField, GaussianField-instFintypeAsymLatticeSitesOfNeZeroNat}
  \texttt{Pphi2.asymEvalMapInv}
\end{definition}
\begin{theorem}[\texttt{asymCanonicalSumFieldLaw\_eq\_map\_asymCanonicalSumFieldFunction}]\label{Pphi2-asymCanonicalSumFieldLaw-eq-map-asymCanonicalSumFieldFunction}
  \lean{Pphi2.asymCanonicalSumFieldLaw_eq_map_asymCanonicalSumFieldFunction}\leanok
  \uses{GaussianField-AsymLatticeSites, Pphi2-asymCanonicalSumFieldLaw, Pphi2-AsymCanonicalJointSumIndex, Pphi2-AsymCanonicalJoint, Pphi2-asymCanonicalSumFieldFunction, Pphi2-asymCanonicalStdFieldFunction-eq-asymCanonicalSumFieldFunction, GaussianField-AsymLatticeField, GaussianField-instFintypeAsymLatticeSitesOfNeZeroNat, Pphi2-asymCanonicalJointMeasure-map-stdGaussianEuclidean, Pphi2-asymCanonicalStdToFieldCLM, Pphi2-asymCanonicalJointMeasure, Pphi2-AsymCanonicalMode}
  \texttt{Pphi2.asymCanonicalSumFieldLaw\_eq\_map\_asymCanonicalSumFieldFunction}
\end{theorem}
\begin{theorem}[\texttt{asymCanonicalRoughFieldFunction\_add}]\label{Pphi2-asymCanonicalRoughFieldFunction-add}
  \lean{Pphi2.asymCanonicalRoughFieldFunction_add}\leanok
  \uses{GaussianField-AsymLatticeSites, Pphi2-asymCanonicalRoughModeCoeff, Pphi2-AsymCanonicalJoint, Pphi2-asymCanonicalRoughFieldFunction, Pphi2-AsymCanonicalMode}
  \texttt{Pphi2.asymCanonicalRoughFieldFunction\_add}
\end{theorem}
\begin{theorem}[\texttt{asymCanonicalSmoothFieldFunction\_memLp\_two}]\label{Pphi2-asymCanonicalSmoothFieldFunction-memLp-two}
  \lean{Pphi2.asymCanonicalSmoothFieldFunction_memLp_two}\leanok
  \uses{GaussianField-AsymLatticeSites, Pphi2-AsymCanonicalJoint, Pphi2-asymCanonicalSmoothModeCoeff, Pphi2-asymCanonicalSmoothFieldFunction, Pphi2-asymCanonicalJointMeasure, Pphi2-AsymCanonicalMode}
  \texttt{Pphi2.asymCanonicalSmoothFieldFunction\_memLp\_two}
\end{theorem}
\begin{theorem}[\texttt{asymEvalMapInv\_evalMap}]\label{Pphi2-asymEvalMapInv-evalMap}
  \lean{Pphi2.asymEvalMapInv_evalMap}\leanok
  \uses{GaussianField-AsymLatticeSites, GaussianField-Configuration, GaussianField-AsymLatticeField, GaussianField-instFintypeAsymLatticeSitesOfNeZeroNat, Pphi2-asymEvalMap, Pphi2-asymConfig-apply-eq-sum-evalMap, Pphi2-asymEvalMapInv}
  \texttt{Pphi2.asymEvalMapInv\_evalMap}
\end{theorem}
\begin{theorem}[\texttt{congr\_simp}]\label{Pphi2-asymCanonicalSmoothFieldFunction-congr-simp}
  \lean{Pphi2.asymCanonicalSmoothFieldFunction.congr_simp}\leanok
  \uses{GaussianField-AsymLatticeSites, Pphi2-AsymCanonicalJoint, Pphi2-asymCanonicalSmoothFieldFunction}
  \texttt{Pphi2.asymCanonicalSmoothFieldFunction.congr\_simp}
\end{theorem}
\begin{theorem}[\texttt{asymLatticeFieldToConfig\_eq\_evalMapInv}]\label{Pphi2-asymLatticeFieldToConfig-eq-evalMapInv}
  \lean{Pphi2.asymLatticeFieldToConfig_eq_evalMapInv}\leanok
  \uses{GaussianField-AsymLatticeSites, GaussianField-Configuration, GaussianField-AsymLatticeField, Pphi2-asymEvalMapInv-apply, GaussianField-instFintypeAsymLatticeSitesOfNeZeroNat, Pphi2-asymLatticeFieldToConfig, Pphi2-asymEvalMapInv}
  \texttt{Pphi2.asymLatticeFieldToConfig\_eq\_evalMapInv}
\end{theorem}
\begin{theorem}[\texttt{asymCanonicalRoughFieldFunction\_smul}]\label{Pphi2-asymCanonicalRoughFieldFunction-smul}
  \lean{Pphi2.asymCanonicalRoughFieldFunction_smul}\leanok
  \uses{GaussianField-AsymLatticeSites, Pphi2-asymCanonicalRoughModeCoeff, Pphi2-AsymCanonicalJoint, Pphi2-asymCanonicalRoughFieldFunction, Pphi2-AsymCanonicalMode}
  \texttt{Pphi2.asymCanonicalRoughFieldFunction\_smul}
\end{theorem}
\begin{theorem}[\texttt{asymEvalMap\_measurable}]\label{Pphi2-asymEvalMap-measurable}
  \lean{Pphi2.asymEvalMap_measurable}\leanok
  \uses{GaussianField-AsymLatticeSites, GaussianField-configuration-eval-measurable, GaussianField-Configuration, GaussianField-AsymLatticeField, GaussianField-instMeasurableSpaceConfiguration, Pphi2-asymEvalMap, GaussianField-asymLatticeDelta}
  \texttt{Pphi2.asymEvalMap\_measurable}
\end{theorem}
\begin{theorem}[\texttt{asymConfig\_apply\_eq\_sum\_evalMap}]\label{Pphi2-asymConfig-apply-eq-sum-evalMap}
  \lean{Pphi2.asymConfig_apply_eq_sum_evalMap}\leanok
  \uses{GaussianField-AsymLatticeSites, Pphi2-asymConfig-apply-eq-sum-delta, GaussianField-Configuration, GaussianField-AsymLatticeField, GaussianField-instFintypeAsymLatticeSitesOfNeZeroNat, Pphi2-asymEvalMap, GaussianField-asymLatticeDelta}
  \texttt{Pphi2.asymConfig\_apply\_eq\_sum\_evalMap}
\end{theorem}
\begin{definition}[\texttt{asymCanonicalRoughFieldFunction}]\label{Pphi2-asymCanonicalRoughFieldFunction}
  \lean{Pphi2.asymCanonicalRoughFieldFunction}\leanok
  \uses{GaussianField-AsymLatticeSites, Pphi2-asymCanonicalRoughModeCoeff, Pphi2-AsymCanonicalJoint, GaussianField-AsymLatticeField, Pphi2-AsymCanonicalMode}
  \texttt{Pphi2.asymCanonicalRoughFieldFunction}
\end{definition}
\begin{theorem}[\texttt{asymCanonicalSumFieldLaw\_eq\_asymLatticeGaussianFieldLaw}]\label{Pphi2-asymCanonicalSumFieldLaw-eq-asymLatticeGaussianFieldLaw}
  \lean{Pphi2.asymCanonicalSumFieldLaw_eq_asymLatticeGaussianFieldLaw}\leanok
  \uses{GaussianField-AsymLatticeSites, Pphi2-asymCanonicalSumFieldLaw, Pphi2-asymEvalMap-measurable, GaussianField-cramerWold, GaussianField-Configuration, GaussianField-AsymLatticeField, Pphi2-latticeGaussianMeasureAsym, GaussianField-instMeasurableSpaceConfiguration, GaussianField-instFintypeAsymLatticeSitesOfNeZeroNat, Pphi2-asymEvalMap, GaussianField-latticeCovarianceAsymGJ, Pphi2-asymLatticeGaussianFieldLaw, GaussianField-ell2-, Pphi2-latticeGaussianMeasureAsym-isProbability, Pphi2-asymCanonicalSumFieldLaw-isGaussian}
  \texttt{Pphi2.asymCanonicalSumFieldLaw\_eq\_asymLatticeGaussianFieldLaw}
\end{theorem}
\begin{definition}[\texttt{asymCanonicalSumFieldFunction}]\label{Pphi2-asymCanonicalSumFieldFunction}
  \lean{Pphi2.asymCanonicalSumFieldFunction}\leanok
  \uses{GaussianField-AsymLatticeSites, Pphi2-AsymCanonicalJoint, Pphi2-asymCanonicalRoughFieldFunction, GaussianField-AsymLatticeField, Pphi2-asymCanonicalSmoothFieldFunction}
  \texttt{Pphi2.asymCanonicalSumFieldFunction}
\end{definition}
\begin{theorem}[\texttt{asymCanonicalSumFieldFunction\_covariance}]\label{Pphi2-asymCanonicalSumFieldFunction-covariance}
  \lean{Pphi2.asymCanonicalSumFieldFunction_covariance}\leanok
  \uses{Pphi2-asymCanonicalSmoothRough-cross-moment-zero, GaussianField-AsymLatticeSites, Pphi2-asymCanonicalBasis, Pphi2-AsymCanonicalJoint, Pphi2-asymCanonicalSumFieldFunction, Pphi2-asymCanonicalRoughFieldFunction, Pphi2-asymCanonicalEigenvalue, Pphi2-asymCanonicalSmoothWeight-add-roughWeight, Pphi2-asymCanonicalSmoothFieldFunction-memLp-two, Pphi2-asymCanonicalRoughFieldFunction-memLp-two, Pphi2-asymCanonicalRoughWeight, Pphi2-asymCanonicalSmoothFieldFunction, Pphi2-asymCanonicalJointMeasure, Pphi2-asymCanonicalSmoothWeight, Pphi2-AsymCanonicalMode, Pphi2-asymCanonicalNormSq}
  \texttt{Pphi2.asymCanonicalSumFieldFunction\_covariance}
\end{theorem}
\begin{definition}[\texttt{asymCanonicalSmoothFieldFunctionOfFst}]\label{Pphi2-asymCanonicalSmoothFieldFunctionOfFst}
  \lean{Pphi2.asymCanonicalSmoothFieldFunctionOfFst}\leanok
  \uses{GaussianField-AsymLatticeSites, Pphi2-asymCanonicalSmoothModeCoeff, GaussianField-AsymLatticeField, Pphi2-AsymCanonicalMode}
  \texttt{Pphi2.asymCanonicalSmoothFieldFunctionOfFst}
\end{definition}
\begin{definition}[\texttt{asymCanonicalJointMeasure}]\label{Pphi2-asymCanonicalJointMeasure}
  \lean{Pphi2.asymCanonicalJointMeasure}\leanok
  \uses{Pphi2-AsymCanonicalJoint, Pphi2-AsymCanonicalMode}
  \texttt{Pphi2.asymCanonicalJointMeasure}
\end{definition}
\begin{theorem}[\texttt{asymCanonicalRoughFieldFunction\_eq\_of\_snd}]\label{Pphi2-asymCanonicalRoughFieldFunction-eq-of-snd}
  \lean{Pphi2.asymCanonicalRoughFieldFunction_eq_of_snd}\leanok
  \uses{Pphi2-asymCanonicalRoughFieldFunctionOfSnd, GaussianField-AsymLatticeSites, Pphi2-asymCanonicalRoughFieldFunction, Pphi2-AsymCanonicalMode}
  \texttt{Pphi2.asymCanonicalRoughFieldFunction\_eq\_of\_snd}
\end{theorem}
\begin{theorem}[\texttt{asymCanonicalSmoothRough\_cross\_moment\_zero}]\label{Pphi2-asymCanonicalSmoothRough-cross-moment-zero}
  \lean{Pphi2.asymCanonicalSmoothRough_cross_moment_zero}\leanok
  \uses{GaussianField-AsymLatticeSites, Pphi2-asymCanonicalRoughModeCoeff, Pphi2-AsymCanonicalJoint, Pphi2-asymCanonicalSmoothModeCoeff, Pphi2-asymCanonicalRoughFieldFunction, Pphi2-asymCanonicalSmoothFieldFunction, Pphi2-asymCanonicalJointMeasure, Pphi2-AsymCanonicalMode}
  \texttt{Pphi2.asymCanonicalSmoothRough\_cross\_moment\_zero}
\end{theorem}
\begin{definition}[\texttt{asymCanonicalSumConfig}]\label{Pphi2-asymCanonicalSumConfig}
  \lean{Pphi2.asymCanonicalSumConfig}\leanok
  \uses{GaussianField-AsymLatticeSites, Pphi2-AsymCanonicalJoint, Pphi2-asymCanonicalSumFieldFunction, GaussianField-Configuration, GaussianField-AsymLatticeField, Pphi2-asymLatticeFieldToConfig}
  \texttt{Pphi2.asymCanonicalSumConfig}
\end{definition}
\begin{definition}[\texttt{asymCanonicalNormSq}]\label{Pphi2-asymCanonicalNormSq}
  \lean{Pphi2.asymCanonicalNormSq}\leanok
  \uses{GaussianField-latticeFourierNormSq, Pphi2-AsymCanonicalMode}
  \texttt{Pphi2.asymCanonicalNormSq}
\end{definition}
\begin{theorem}[\texttt{asymCanonicalSumConfig\_apply\_delta}]\label{Pphi2-asymCanonicalSumConfig-apply-delta}
  \lean{Pphi2.asymCanonicalSumConfig_apply_delta}\leanok
  \uses{Pphi2-asymLatticeFieldToConfig-apply, GaussianField-AsymLatticeSites, Pphi2-AsymCanonicalJoint, Pphi2-asymCanonicalSumFieldFunction, Pphi2-asymCanonicalSumConfig, GaussianField-Configuration, GaussianField-AsymLatticeField, GaussianField-instFintypeAsymLatticeSitesOfNeZeroNat, Pphi2-asymLatticeFieldToConfig, GaussianField-asymLatticeDelta}
  \texttt{Pphi2.asymCanonicalSumConfig\_apply\_delta}
\end{theorem}
\begin{definition}[\texttt{asymLatticeGaussianFieldLaw}]\label{Pphi2-asymLatticeGaussianFieldLaw}
  \lean{Pphi2.asymLatticeGaussianFieldLaw}\leanok
  \uses{GaussianField-AsymLatticeSites, GaussianField-Configuration, GaussianField-AsymLatticeField, Pphi2-latticeGaussianMeasureAsym, GaussianField-instMeasurableSpaceConfiguration, Pphi2-asymEvalMap}
  \texttt{Pphi2.asymLatticeGaussianFieldLaw}
\end{definition}
\begin{theorem}[\texttt{asymCanonicalSmoothFieldFunction\_pointwise\_measurable}]\label{Pphi2-asymCanonicalSmoothFieldFunction-pointwise-measurable}
  \lean{Pphi2.asymCanonicalSmoothFieldFunction_pointwise_measurable}\leanok
  \uses{GaussianField-AsymLatticeSites, Pphi2-AsymCanonicalJoint, Pphi2-asymCanonicalSmoothModeCoeff, Pphi2-asymCanonicalSmoothFieldFunction, Pphi2-AsymCanonicalMode}
  \texttt{Pphi2.asymCanonicalSmoothFieldFunction\_pointwise\_measurable}
\end{theorem}
\begin{theorem}[\texttt{congr\_simp}]\label{Pphi2-asymCanonicalStdToFieldCLM-congr-simp}
  \lean{Pphi2.asymCanonicalStdToFieldCLM.congr_simp}\leanok
  \uses{GaussianField-AsymLatticeSites, Pphi2-AsymCanonicalJointSumIndex, GaussianField-AsymLatticeField, Pphi2-asymCanonicalStdToFieldCLM}
  \texttt{Pphi2.asymCanonicalStdToFieldCLM.congr\_simp}
\end{theorem}
\begin{definition}[\texttt{asymCanonicalRoughFieldFunctionOfSnd}]\label{Pphi2-asymCanonicalRoughFieldFunctionOfSnd}
  \lean{Pphi2.asymCanonicalRoughFieldFunctionOfSnd}\leanok
  \uses{GaussianField-AsymLatticeSites, Pphi2-asymCanonicalRoughModeCoeff, GaussianField-AsymLatticeField, Pphi2-AsymCanonicalMode}
  \texttt{Pphi2.asymCanonicalRoughFieldFunctionOfSnd}
\end{definition}
\begin{theorem}[\texttt{asymEvalMapInv\_apply}]\label{Pphi2-asymEvalMapInv-apply}
  \lean{Pphi2.asymEvalMapInv_apply}\leanok
  \uses{GaussianField-AsymLatticeSites, GaussianField-Configuration, GaussianField-AsymLatticeField, GaussianField-instFintypeAsymLatticeSitesOfNeZeroNat, Pphi2-asymEvalMapInv}
  \texttt{Pphi2.asymEvalMapInv\_apply}
\end{theorem}
\begin{theorem}[\texttt{asymCanonicalJointMeasure\_map\_stdGaussianEuclidean}]\label{Pphi2-asymCanonicalJointMeasure-map-stdGaussianEuclidean}
  \lean{Pphi2.asymCanonicalJointMeasure_map_stdGaussianEuclidean}\leanok
  \uses{Pphi2-AsymCanonicalJointSumIndex, Pphi2-AsymCanonicalJoint, Pphi2-asymCanonicalJointMeasure, Pphi2-AsymCanonicalMode}
  \texttt{Pphi2.asymCanonicalJointMeasure\_map\_stdGaussianEuclidean}
\end{theorem}
\begin{definition}[\texttt{asymCanonicalSmoothFieldFunction}]\label{Pphi2-asymCanonicalSmoothFieldFunction}
  \lean{Pphi2.asymCanonicalSmoothFieldFunction}\leanok
  \uses{GaussianField-AsymLatticeSites, Pphi2-AsymCanonicalJoint, Pphi2-asymCanonicalSmoothModeCoeff, GaussianField-AsymLatticeField, Pphi2-AsymCanonicalMode}
  \texttt{Pphi2.asymCanonicalSmoothFieldFunction}
\end{definition}
\begin{theorem}[\texttt{asymCanonicalRoughFieldFunction\_pointwise\_measurable}]\label{Pphi2-asymCanonicalRoughFieldFunction-pointwise-measurable}
  \lean{Pphi2.asymCanonicalRoughFieldFunction_pointwise_measurable}\leanok
  \uses{GaussianField-AsymLatticeSites, Pphi2-asymCanonicalRoughModeCoeff, Pphi2-AsymCanonicalJoint, Pphi2-asymCanonicalRoughFieldFunction, Pphi2-AsymCanonicalMode}
  \texttt{Pphi2.asymCanonicalRoughFieldFunction\_pointwise\_measurable}
\end{theorem}
\begin{theorem}[\texttt{congr\_simp}]\label{Pphi2-asymCanonicalRoughModeCoeff-congr-simp}
  \lean{Pphi2.asymCanonicalRoughModeCoeff.congr_simp}\leanok
  \uses{GaussianField-AsymLatticeSites, Pphi2-asymCanonicalRoughModeCoeff, Pphi2-AsymCanonicalMode}
  \texttt{Pphi2.asymCanonicalRoughModeCoeff.congr\_simp}
\end{theorem}
\begin{theorem}[\texttt{asymCanonicalSmoothFieldFunction\_add}]\label{Pphi2-asymCanonicalSmoothFieldFunction-add}
  \lean{Pphi2.asymCanonicalSmoothFieldFunction_add}\leanok
  \uses{GaussianField-AsymLatticeSites, Pphi2-AsymCanonicalJoint, Pphi2-asymCanonicalSmoothModeCoeff, Pphi2-asymCanonicalSmoothFieldFunction, Pphi2-AsymCanonicalMode}
  \texttt{Pphi2.asymCanonicalSmoothFieldFunction\_add}
\end{theorem}
\begin{theorem}[\texttt{asymCanonicalSmoothFieldFunction\_eq\_of\_fst}]\label{Pphi2-asymCanonicalSmoothFieldFunction-eq-of-fst}
  \lean{Pphi2.asymCanonicalSmoothFieldFunction_eq_of_fst}\leanok
  \uses{GaussianField-AsymLatticeSites, Pphi2-asymCanonicalSmoothFieldFunctionOfFst, Pphi2-asymCanonicalSmoothFieldFunction, Pphi2-AsymCanonicalMode}
  \texttt{Pphi2.asymCanonicalSmoothFieldFunction\_eq\_of\_fst}
\end{theorem}
\begin{definition}[\texttt{asymCanonicalStdFieldFunction}]\label{Pphi2-asymCanonicalStdFieldFunction}
  \lean{Pphi2.asymCanonicalStdFieldFunction}\leanok
  \uses{GaussianField-AsymLatticeSites, Pphi2-asymCanonicalRoughModeCoeff, Pphi2-AsymCanonicalJointSumIndex, Pphi2-asymCanonicalSmoothModeCoeff, GaussianField-AsymLatticeField, Pphi2-AsymCanonicalMode}
  \texttt{Pphi2.asymCanonicalStdFieldFunction}
\end{definition}
\begin{theorem}[\texttt{asymCanonicalJointMeasure\_isProbability}]\label{Pphi2-asymCanonicalJointMeasure-isProbability}
  \lean{Pphi2.asymCanonicalJointMeasure_isProbability}\leanok
  \uses{Pphi2-AsymCanonicalJoint, Pphi2-asymCanonicalJointMeasure, Pphi2-AsymCanonicalMode}
  \texttt{Pphi2.asymCanonicalJointMeasure\_isProbability}
\end{theorem}
\begin{definition}[\texttt{asymCanonicalRoughModeCoeff}]\label{Pphi2-asymCanonicalRoughModeCoeff}
  \lean{Pphi2.asymCanonicalRoughModeCoeff}\leanok
  \uses{GaussianField-AsymLatticeSites, Pphi2-asymCanonicalBasis, Pphi2-asymCanonicalRoughWeight, Pphi2-AsymCanonicalMode, Pphi2-asymCanonicalNormSq}
  \texttt{Pphi2.asymCanonicalRoughModeCoeff}
\end{definition}
\begin{theorem}[\texttt{asymCanonicalSumFieldFunction\_covariance\_eq\_GJ}]\label{Pphi2-asymCanonicalSumFieldFunction-covariance-eq-GJ}
  \lean{Pphi2.asymCanonicalSumFieldFunction_covariance_eq_GJ}\leanok
  \uses{GaussianField-spectralLatticeCovarianceAsym, GaussianField-latticeEigenvalue1d, GaussianField-AsymLatticeSites, Pphi2-asymCanonicalBasis, GaussianField-latticeFourierBasisFun, Pphi2-asymCanonicalSumFieldFunction-covariance, GaussianField-covariance, Pphi2-AsymCanonicalJoint, GaussianField-latticeFourierNormSq, Pphi2-asymCanonicalSumFieldFunction, GaussianField-AsymLatticeField, Pphi2-asymCanonicalEigenvalue, GaussianField-instFintypeAsymLatticeSitesOfNeZeroNat, GaussianField-abstract-spectral-eq-dft-spectral-2d-asym, GaussianField-latticeCovarianceAsymGJ, GaussianField-ell2-, Pphi2-asymCanonicalJointMeasure, Pphi2-AsymCanonicalMode, Pphi2-asymCanonicalNormSq}
  \texttt{Pphi2.asymCanonicalSumFieldFunction\_covariance\_eq\_GJ}
\end{theorem}
\begin{definition}[\texttt{asymCanonicalStdToFieldLinearMap}]\label{Pphi2-asymCanonicalStdToFieldLinearMap}
  \lean{Pphi2.asymCanonicalStdToFieldLinearMap}\leanok
  \uses{GaussianField-AsymLatticeSites, Pphi2-AsymCanonicalJointSumIndex, GaussianField-AsymLatticeField, Pphi2-asymCanonicalStdFieldFunction}
  \texttt{Pphi2.asymCanonicalStdToFieldLinearMap}
\end{definition}
\begin{theorem}[\texttt{asymCanonicalSumFieldFunction\_smul}]\label{Pphi2-asymCanonicalSumFieldFunction-smul}
  \lean{Pphi2.asymCanonicalSumFieldFunction_smul}\leanok
  \uses{GaussianField-AsymLatticeSites, Pphi2-asymCanonicalSmoothFieldFunction-smul, Pphi2-AsymCanonicalJoint, Pphi2-asymCanonicalSumFieldFunction, Pphi2-asymCanonicalRoughFieldFunction, Pphi2-asymCanonicalRoughFieldFunction-smul, Pphi2-asymCanonicalSmoothFieldFunction, Pphi2-AsymCanonicalMode}
  \texttt{Pphi2.asymCanonicalSumFieldFunction\_smul}
\end{theorem}
\begin{definition}[\texttt{asymEvalMap}]\label{Pphi2-asymEvalMap}
  \lean{Pphi2.asymEvalMap}\leanok
  \uses{GaussianField-AsymLatticeSites, GaussianField-Configuration, GaussianField-AsymLatticeField, GaussianField-asymLatticeDelta}
  \texttt{Pphi2.asymEvalMap}
\end{definition}
\begin{theorem}[\texttt{asymCanonicalSumFieldFunction\_pointwise\_measurable}]\label{Pphi2-asymCanonicalSumFieldFunction-pointwise-measurable}
  \lean{Pphi2.asymCanonicalSumFieldFunction_pointwise_measurable}\leanok
  \uses{Pphi2-asymCanonicalSmoothFieldFunction-pointwise-measurable, GaussianField-AsymLatticeSites, Pphi2-AsymCanonicalJoint, Pphi2-asymCanonicalRoughFieldFunction-pointwise-measurable, Pphi2-asymCanonicalSumFieldFunction, Pphi2-asymCanonicalRoughFieldFunction, Pphi2-asymCanonicalSmoothFieldFunction, Pphi2-AsymCanonicalMode}
  \texttt{Pphi2.asymCanonicalSumFieldFunction\_pointwise\_measurable}
\end{theorem}
\begin{definition}[\texttt{asymCanonicalSmoothModeCoeff}]\label{Pphi2-asymCanonicalSmoothModeCoeff}
  \lean{Pphi2.asymCanonicalSmoothModeCoeff}\leanok
  \uses{GaussianField-AsymLatticeSites, Pphi2-asymCanonicalBasis, Pphi2-asymCanonicalSmoothWeight, Pphi2-AsymCanonicalMode, Pphi2-asymCanonicalNormSq}
  \texttt{Pphi2.asymCanonicalSmoothModeCoeff}
\end{definition}
\begin{theorem}[\texttt{asymEvalMapInv\_measurable}]\label{Pphi2-asymEvalMapInv-measurable}
  \lean{Pphi2.asymEvalMapInv_measurable}\leanok
  \uses{GaussianField-AsymLatticeSites, GaussianField-configuration-measurable-of-eval-measurable, GaussianField-Configuration, GaussianField-AsymLatticeField, Pphi2-asymEvalMapInv-apply, GaussianField-instMeasurableSpaceConfiguration, GaussianField-instFintypeAsymLatticeSitesOfNeZeroNat, Pphi2-asymEvalMapInv}
  \texttt{Pphi2.asymEvalMapInv\_measurable}
\end{theorem}
\begin{theorem}[\texttt{asymCanonicalSumConfig\_measurable}]\label{Pphi2-asymCanonicalSumConfig-measurable}
  \lean{Pphi2.asymCanonicalSumConfig_measurable}\leanok
  \uses{Pphi2-asymLatticeFieldToConfig-eq-evalMapInv, GaussianField-AsymLatticeSites, Pphi2-AsymCanonicalJoint, Pphi2-asymEvalMapInv-measurable, Pphi2-asymCanonicalSumFieldFunction, Pphi2-asymCanonicalSumConfig, GaussianField-Configuration, GaussianField-AsymLatticeField, GaussianField-instMeasurableSpaceConfiguration, Pphi2-asymLatticeFieldToConfig, Pphi2-asymCanonicalSumFieldFunction-measurable, Pphi2-asymEvalMapInv, Pphi2-AsymCanonicalMode}
  \texttt{Pphi2.asymCanonicalSumConfig\_measurable}
\end{theorem}
\begin{definition}[\texttt{asymCanonicalStdToFieldCLM}]\label{Pphi2-asymCanonicalStdToFieldCLM}
  \lean{Pphi2.asymCanonicalStdToFieldCLM}\leanok
  \uses{GaussianField-AsymLatticeSites, Pphi2-AsymCanonicalJointSumIndex, GaussianField-AsymLatticeField, Pphi2-asymCanonicalStdToFieldLinearMap}
  \texttt{Pphi2.asymCanonicalStdToFieldCLM}
\end{definition}
\begin{theorem}[\texttt{congr\_simp}]\label{Pphi2-asymCanonicalRoughFieldFunction-congr-simp}
  \lean{Pphi2.asymCanonicalRoughFieldFunction.congr_simp}\leanok
  \uses{GaussianField-AsymLatticeSites, Pphi2-AsymCanonicalJoint, Pphi2-asymCanonicalRoughFieldFunction}
  \texttt{Pphi2.asymCanonicalRoughFieldFunction.congr\_simp}
\end{theorem}
\begin{theorem}[\texttt{asymLatticeFieldToConfig\_apply}]\label{Pphi2-asymLatticeFieldToConfig-apply}
  \lean{Pphi2.asymLatticeFieldToConfig_apply}\leanok
  \uses{GaussianField-AsymLatticeSites, GaussianField-Configuration, GaussianField-AsymLatticeField, GaussianField-instFintypeAsymLatticeSitesOfNeZeroNat, Pphi2-asymLatticeFieldToConfig}
  \texttt{Pphi2.asymLatticeFieldToConfig\_apply}
\end{theorem}
\begin{definition}[\texttt{AsymCanonicalMode}]\label{Pphi2-AsymCanonicalMode}
  \lean{Pphi2.AsymCanonicalMode}\leanok
  \texttt{Pphi2.AsymCanonicalMode}
\end{definition}
\begin{theorem}[\texttt{asymCanonicalSumFieldFunction\_measurable}]\label{Pphi2-asymCanonicalSumFieldFunction-measurable}
  \lean{Pphi2.asymCanonicalSumFieldFunction_measurable}\leanok
  \uses{GaussianField-AsymLatticeSites, Pphi2-asymCanonicalSumFieldFunction-pointwise-measurable, Pphi2-AsymCanonicalJoint, Pphi2-asymCanonicalSumFieldFunction, GaussianField-AsymLatticeField, Pphi2-AsymCanonicalMode}
  \texttt{Pphi2.asymCanonicalSumFieldFunction\_measurable}
\end{theorem}
\begin{definition}[\texttt{asymCanonicalSmoothWeight}]\label{Pphi2-asymCanonicalSmoothWeight}
  \lean{Pphi2.asymCanonicalSmoothWeight}\leanok
  \uses{Pphi2-asymCanonicalEigenvalue, Pphi2-AsymCanonicalMode}
  \texttt{Pphi2.asymCanonicalSmoothWeight}
\end{definition}
\begin{theorem}[\texttt{asymCanonicalRoughFieldFunction\_measurable}]\label{Pphi2-asymCanonicalRoughFieldFunction-measurable}
  \lean{Pphi2.asymCanonicalRoughFieldFunction_measurable}\leanok
  \uses{GaussianField-AsymLatticeSites, Pphi2-AsymCanonicalJoint, Pphi2-asymCanonicalRoughFieldFunction-pointwise-measurable, Pphi2-asymCanonicalRoughFieldFunction, GaussianField-AsymLatticeField, Pphi2-AsymCanonicalMode}
  \texttt{Pphi2.asymCanonicalRoughFieldFunction\_measurable}
\end{theorem}
\section{\texttt{Pphi2.NelsonEstimate.RoughErrorBound}}
\begin{theorem}[\texttt{integrable\_sq\_real\_sum\_of\_pairwise}]\label{Pphi2-integrable-sq-real-sum-of-pairwise}
  \lean{Pphi2.integrable_sq_real_sum_of_pairwise}\leanok
  \uses{Pphi2-integrable-sum-mul-sum-of-pairwise}
  \texttt{Pphi2.integrable\_sq\_real\_sum\_of\_pairwise}
\end{theorem}
\begin{theorem}[\texttt{wienerChaosLE\_mono}]\label{Pphi2-wienerChaosLE-mono}
  \lean{Pphi2.wienerChaosLE_mono}\leanok
  \uses{GaussianHilbert-stdGaussianFin, GaussianHilbert-wienerChaos, GaussianHilbert-wienerChaosLE}
  \texttt{Pphi2.wienerChaosLE\_mono}
\end{theorem}
\begin{theorem}[\texttt{wickMonomial\_eq\_root\_local}]\label{Pphi2-wickMonomial-eq-root-local}
  \lean{Pphi2.wickMonomial_eq_root_local}\leanok
  \uses{Pphi2-wickMonomial}
  \texttt{Pphi2.wickMonomial\_eq\_root\_local}
\end{theorem}
\begin{theorem}[\texttt{congr\_simp}]\label{Pphi2-latticeFieldToConfig-congr-simp}
  \lean{Pphi2.latticeFieldToConfig.congr_simp}\leanok
  \uses{GaussianField-FinLatticeSites, GaussianField-Configuration, Pphi2-latticeFieldToConfig, GaussianField-FinLatticeField}
  \texttt{Pphi2.latticeFieldToConfig.congr\_simp}
\end{theorem}
\begin{theorem}[\texttt{congr\_simp}]\label{Pphi2-CanonicalJointSumIndex-congr-simp}
  \lean{Pphi2.CanonicalJointSumIndex.congr_simp}\leanok
  \uses{Pphi2-CanonicalJointSumIndex}
  \texttt{Pphi2.CanonicalJointSumIndex.congr\_simp}
\end{theorem}
\begin{theorem}[\texttt{rough\_error\_variance}]\label{Pphi2-rough-error-variance}
  \lean{Pphi2.rough_error_variance}\leanok
  \uses{Pphi2-integrable-sum-mul-sum-of-pairwise, Pphi2-canonicalJointMeasure, Pphi2-canonicalCrossTerm-inner-eq-zero, Pphi2-canonicalRoughError-l2-sq-eq, Pphi2-integral-sum-mul-sum-eq-zero-of-orth, Pphi2-integrable-sq-real-sum-of-pairwise, Pphi2-canonicalCrossTerm-pair-integrable, Pphi2-canonicalCrossTerm-l2-sq-le, Pphi2-canonicalRoughError, Pphi2-canonicalCrossTerm, Lattice-toSemilatticeInf, Pphi2-CanonicalJoint}
  $\exists C > 0$ and ... (placeholder \texttt{True}). Intended: $\mathbb{E}[E_R^2] \le C \cdot T^{1/2}$.
\end{theorem}
\begin{theorem}[\texttt{canonicalRoughError\_neg\_tail\_uniform\_in\_aN}]\label{Pphi2-canonicalRoughError-neg-tail-uniform-in-aN}
  \lean{Pphi2.canonicalRoughError_neg_tail_uniform_in_aN}\leanok
  \uses{Pphi2-rough-error-variance, Pphi2-canonicalJointMeasure, Pphi2-canonicalJointStdGaussianMeasurableEquiv, GaussianHilbert-stdGaussianFin, GaussianHilbert-wienerChaosLE, Pphi2-CanonicalJointSumIndex, Pphi2-canonicalJointMeasure-map-stdGaussian, Pphi2-ChaosTailBridge-chaosTailConstant, Pphi2-canonicalRoughError-neg-tail-of-stdGaussian-explicit-ae, Pphi2-canonicalRoughError, Lattice-toSemilatticeInf, Pphi2-CanonicalJoint}
  \texttt{Pphi2.canonicalRoughError\_neg\_tail\_uniform\_in\_aN}
\end{theorem}
\begin{theorem}[\texttt{canonicalCrossTerm\_l2\_sq\_eq\_covSum}]\label{Pphi2-canonicalCrossTerm-l2-sq-eq-covSum}
  \lean{Pphi2.canonicalCrossTerm_l2_sq_eq_covSum}\leanok
  \uses{Pphi2-smoothWickConstant, Pphi2-canonicalJointMeasure, Pphi2-canonicalRoughFieldFunctionOfSnd, Pphi2-canonicalRoughCovariance, Pphi2-canonicalRoughWickPower-two-site-marginal-integrable, Pphi2-canonicalSmoothCovariance, GaussianField-FinLatticeSites, Pphi2-wickMonomial, Pphi2-canonicalSmoothFieldFunctionOfFst, Pphi2-canonicalSmoothWickPower-two-site-marginal-integrable, Pphi2-canonicalRoughFieldFunction, Pphi2-canonicalRoughWickPower-two-site-marginal-eq-diag, Pphi2-canonicalSmoothFieldFunction, Pphi2-roughWickConstant, Pphi2-canonicalCrossTerm, Pphi2-canonicalSmoothWickPower-two-site-marginal-eq-diag, Pphi2-CanonicalJoint}
  \texttt{Pphi2.canonicalCrossTerm\_l2\_sq\_eq\_covSum}
\end{theorem}
\begin{theorem}[\texttt{integral\_sum\_mul\_sum\_eq\_zero\_of\_orth}]\label{Pphi2-integral-sum-mul-sum-eq-zero-of-orth}
  \lean{Pphi2.integral_sum_mul_sum_eq_zero_of_orth}\leanok
  \texttt{Pphi2.integral\_sum\_mul\_sum\_eq\_zero\_of\_orth}
\end{theorem}
\begin{theorem}[\texttt{canonicalRoughError\_neg\_tail\_of\_stdGaussian\_explicit\_ae}]\label{Pphi2-canonicalRoughError-neg-tail-of-stdGaussian-explicit-ae}
  \lean{Pphi2.canonicalRoughError_neg_tail_of_stdGaussian_explicit_ae}\leanok
  \uses{Pphi2-canonicalJointMeasure, Pphi2-canonicalJointStdGaussianMeasurableEquiv, GaussianHilbert-stdGaussianFin, GaussianHilbert-wienerChaosLE, Pphi2-CanonicalJointSumIndex, Pphi2-ChaosTailBridge-chaos-neg-tail-bound-explicit, Pphi2-canonicalJointMeasure-map-stdGaussian, Pphi2-ChaosTailBridge-chaosTailConstant, Pphi2-canonicalRoughError, Pphi2-CanonicalJoint}
  \texttt{Pphi2.canonicalRoughError\_neg\_tail\_of\_stdGaussian\_explicit\_ae}
\end{theorem}
\begin{theorem}[\texttt{integral\_sq\_add\_of\_inner\_eq\_zero}]\label{Pphi2-integral-sq-add-of-inner-eq-zero}
  \lean{Pphi2.integral_sq_add_of_inner_eq_zero}\leanok
  \texttt{Pphi2.integral\_sq\_add\_of\_inner\_eq\_zero}
\end{theorem}
\begin{theorem}[\texttt{finite\_indexed\_wick\_sum\_mem\_wienerChaosLE}]\label{Pphi2-finite-indexed-wick-sum-mem-wienerChaosLE}
  \lean{Pphi2.finite_indexed_wick_sum_mem_wienerChaosLE}\leanok
  \uses{GaussianHilbert-MultiIndex-totalDegree, Pphi2-multiWickMonomial-eq-hermiteMultiEval, GaussianHilbert-hermiteMultiLp-coeFn, Pphi2-wickMonomial, GaussianHilbert-stdGaussianFin, GaussianHilbert-wienerChaos, GaussianHilbert-wienerChaosLE, GaussianHilbert-hermiteMultiLp, GaussianHilbert-hermiteMultiEval, GaussianHilbert-hermiteMultiLp-mem-wienerChaos}
  \texttt{Pphi2.finite\_indexed\_wick\_sum\_mem\_wienerChaosLE}
\end{theorem}
\begin{theorem}[\texttt{canonicalFullInteractionJoint\_moment\_le}]\label{Pphi2-canonicalFullInteractionJoint-moment-le}
  \lean{Pphi2.canonicalFullInteractionJoint_moment_le}\leanok
  \uses{Pphi2-canonicalJointMeasure, Pphi2-canonicalJointStdGaussianMeasurableEquiv, Pphi2-canonicalFullInteractionJoint, GaussianHilbert-stdGaussianFin, GaussianHilbert-wienerChaosLE, Pphi2-CanonicalJointSumIndex, Pphi2-canonicalJointMeasure-map-stdGaussian, Pphi2-chaosLE-moment-le, Pphi2-CanonicalJoint}
  \texttt{Pphi2.canonicalFullInteractionJoint\_moment\_le}
\end{theorem}
\begin{definition}[\texttt{canonicalFullInteractionJoint}]\label{Pphi2-canonicalFullInteractionJoint}
  \lean{Pphi2.canonicalFullInteractionJoint}\leanok
  \uses{GaussianField-FinLatticeSites, Pphi2-canonicalSumFieldFunction, Pphi2-wickConstant, Pphi2-wickPolynomial, Pphi2-CanonicalJoint}
  \texttt{Pphi2.canonicalFullInteractionJoint}
\end{definition}
\begin{theorem}[\texttt{canonicalFullInteractionJoint\_eq\_interactionFunctional}]\label{Pphi2-canonicalFullInteractionJoint-eq-interactionFunctional}
  \lean{Pphi2.canonicalFullInteractionJoint_eq_interactionFunctional}\leanok
  \uses{Pphi2-canonicalSumConfig, GaussianField-finLatticeDelta, GaussianField-FinLatticeSites, Pphi2-canonicalSumConfig-apply-delta, Pphi2-canonicalFullInteractionJoint, GaussianField-Configuration, Pphi2-canonicalSumFieldFunction, Pphi2-wickConstant, GaussianField-FinLatticeField, Pphi2-interactionFunctional, Pphi2-wickPolynomial, Pphi2-CanonicalJoint}
  \texttt{Pphi2.canonicalFullInteractionJoint\_eq\_interactionFunctional}
\end{theorem}
\begin{theorem}[\texttt{canonicalRoughError\_eq\_sum\_over\_cross\_terms}]\label{Pphi2-canonicalRoughError-eq-sum-over-cross-terms}
  \lean{Pphi2.canonicalRoughError_eq_sum_over_cross_terms}\leanok
  \uses{Pphi2-smoothWickConstant, GaussianField-FinLatticeSites, Pphi2-canonicalRoughError-pointwise-decomposition, Pphi2-wickMonomial, Pphi2-canonicalRoughFieldFunction, Pphi2-canonicalSmoothFieldFunction, Pphi2-roughWickConstant, Pphi2-canonicalRoughError, Pphi2-canonicalCrossTerm, Pphi2-CanonicalJoint}
  \texttt{Pphi2.canonicalRoughError\_eq\_sum\_over\_cross\_terms}
\end{theorem}
\begin{theorem}[\texttt{canonicalCrossTerm\_inner\_eq\_zero}]\label{Pphi2-canonicalCrossTerm-inner-eq-zero}
  \lean{Pphi2.canonicalCrossTerm_inner_eq_zero}\leanok
  \uses{Pphi2-smoothWickConstant, Pphi2-canonicalJointMeasure, Pphi2-canonicalRoughFieldFunctionOfSnd, Pphi2-canonicalRoughWickPower-two-site-marginal-integrable, GaussianField-FinLatticeSites, Pphi2-wickMonomial, Pphi2-canonicalSmoothFieldFunctionOfFst, Pphi2-canonicalSmoothWickPower-two-site-marginal-integrable, Pphi2-canonicalRoughFieldFunction, Pphi2-canonicalSmoothWickPower-two-site-marginal-eq-zero-of-ne, Pphi2-canonicalSmoothFieldFunction, Pphi2-canonicalRoughWickPower-two-site-marginal-eq-zero-of-ne, Pphi2-roughWickConstant, Pphi2-canonicalCrossTerm, Pphi2-CanonicalJoint}
  \texttt{Pphi2.canonicalCrossTerm\_inner\_eq\_zero}
\end{theorem}
\begin{theorem}[\texttt{canonicalRoughError\_neg\_tail\_of\_stdGaussian\_explicit}]\label{Pphi2-canonicalRoughError-neg-tail-of-stdGaussian-explicit}
  \lean{Pphi2.canonicalRoughError_neg_tail_of_stdGaussian_explicit}\leanok
  \uses{Pphi2-canonicalJointMeasure, Pphi2-canonicalJointStdGaussianMeasurableEquiv, GaussianHilbert-stdGaussianFin, GaussianHilbert-wienerChaosLE, Pphi2-CanonicalJointSumIndex, Pphi2-ChaosTailBridge-chaos-neg-tail-bound-explicit, Pphi2-canonicalJointMeasure-map-stdGaussian, Pphi2-ChaosTailBridge-chaosTailConstant, Pphi2-canonicalRoughError, Pphi2-CanonicalJoint}
  \texttt{Pphi2.canonicalRoughError\_neg\_tail\_of\_stdGaussian\_explicit}
\end{theorem}
\begin{theorem}[\texttt{canonicalSmoothInteraction\_lower\_bound\_at\_cutoff\_uniform}]\label{Pphi2-canonicalSmoothInteraction-lower-bound-at-cutoff-uniform}
  \lean{Pphi2.canonicalSmoothInteraction_lower_bound_at_cutoff_uniform}\leanok
  \uses{Pphi2-DynamicalCutoff-degreeDynamicalCutoffScale-pos, Pphi2-canonicalSmoothInteraction, Pphi2-DynamicalCutoff-degreeDynamicalCutoffScale, Pphi2-DynamicalCutoff-degreeCutoffPower, Pphi2-DynamicalCutoff-degreeDynamicalCutoffScale-log-pow-le, Pphi2-canonicalSmoothInteraction-lower-bound-log-uniform-in-aN, Pphi2-CanonicalJoint}
  \texttt{Pphi2.canonicalSmoothInteraction\_lower\_bound\_at\_cutoff\_uniform}
\end{theorem}
\begin{theorem}[\texttt{canonicalRoughError\_l2\_sq\_eq\_lead\_plus\_perCoef\_sq}]\label{Pphi2-canonicalRoughError-l2-sq-eq-lead-plus-perCoef-sq}
  \lean{Pphi2.canonicalRoughError_l2_sq_eq_lead_plus_perCoef_sq}\leanok
  \uses{Pphi2-canonicalJointMeasure, Pphi2-canonicalRoughError-eq-sum-over-cross-terms, Pphi2-integral-sq-add-of-inner-eq-zero, Pphi2-canonicalLeading-perCoeff-inner-eq-zero, Pphi2-canonicalRoughError, Pphi2-canonicalCrossTerm, Pphi2-CanonicalJoint}
  \texttt{Pphi2.canonicalRoughError\_l2\_sq\_eq\_lead\_plus\_perCoef\_sq}
\end{theorem}
\begin{definition}[\texttt{canonicalSmoothInteraction}]\label{Pphi2-canonicalSmoothInteraction}
  \lean{Pphi2.canonicalSmoothInteraction}\leanok
  \uses{Pphi2-smoothWickConstant, GaussianField-FinLatticeSites, Pphi2-canonicalSmoothFieldFunction, Pphi2-wickPolynomial, Pphi2-CanonicalJoint}
  \texttt{Pphi2.canonicalSmoothInteraction}
\end{definition}
\begin{theorem}[\texttt{canonicalSmoothInteraction\_moment\_le}]\label{Pphi2-canonicalSmoothInteraction-moment-le}
  \lean{Pphi2.canonicalSmoothInteraction_moment_le}\leanok
  \uses{Pphi2-canonicalJointMeasure, Pphi2-canonicalJointStdGaussianMeasurableEquiv, GaussianHilbert-stdGaussianFin, Pphi2-canonicalSmoothInteraction, GaussianHilbert-wienerChaosLE, Pphi2-CanonicalJointSumIndex, Pphi2-canonicalJointMeasure-map-stdGaussian, Pphi2-chaosLE-moment-le, Pphi2-CanonicalJoint}
  \texttt{Pphi2.canonicalSmoothInteraction\_moment\_le}
\end{theorem}
\begin{theorem}[\texttt{multiWickMonomial\_eq\_hermiteMultiEval}]\label{Pphi2-multiWickMonomial-eq-hermiteMultiEval}
  \lean{Pphi2.multiWickMonomial_eq_hermiteMultiEval}\leanok
  \uses{Pphi2-wickMonomial, GaussianHilbert-hermiteEval, GaussianHilbert-hermiteMultiEval}
  \texttt{Pphi2.multiWickMonomial\_eq\_hermiteMultiEval}
\end{theorem}
\begin{theorem}[\texttt{canonicalCrossTerm\_scaled\_inner\_sum\_l2\_sq}]\label{Pphi2-canonicalCrossTerm-scaled-inner-sum-l2-sq}
  \lean{Pphi2.canonicalCrossTerm_scaled_inner_sum_l2_sq}\leanok
  \uses{Pphi2-canonicalJointMeasure, Pphi2-canonicalCrossTerm-inner-eq-zero, Pphi2-integral-sq-real-sum-of-pairwise-orthogonal, Pphi2-canonicalCrossTerm, Pphi2-CanonicalJoint}
  \texttt{Pphi2.canonicalCrossTerm\_scaled\_inner\_sum\_l2\_sq}
\end{theorem}
\begin{theorem}[\texttt{integral\_exp\_neg\_interaction\_sq\_eq\_canonicalJoint}]\label{Pphi2-integral-exp-neg-interaction-sq-eq-canonicalJoint}
  \lean{Pphi2.integral_exp_neg_interaction_sq_eq_canonicalJoint}\leanok
  \uses{Pphi2-canonicalJointMeasure, Pphi2-canonicalSumConfig, Pphi2-canonicalFullInteractionJoint-eq-interactionFunctional, GaussianField-FinLatticeSites, Pphi2-canonicalFullInteractionJoint, Pphi2-interactionFunctional-measurable, GaussianField-latticeGaussianMeasure, GaussianField-Configuration, GaussianField-instMeasurableSpaceConfiguration, Pphi2-integral-comp-canonicalSumConfig, GaussianField-FinLatticeField, Pphi2-interactionFunctional, Pphi2-CanonicalJoint}
  \texttt{Pphi2.integral\_exp\_neg\_interaction\_sq\_eq\_canonicalJoint}
\end{theorem}
\begin{theorem}[\texttt{canonicalFullInteractionJoint\_memLp}]\label{Pphi2-canonicalFullInteractionJoint-memLp}
  \lean{Pphi2.canonicalFullInteractionJoint_memLp}\leanok
  \uses{Pphi2-canonicalJointMeasure, Pphi2-canonicalJointStdGaussianMeasurableEquiv, Pphi2-canonicalFullInteractionJoint, GaussianHilbert-stdGaussianFin, GaussianHilbert-wienerChaosLE, Pphi2-CanonicalJointSumIndex, Pphi2-canonicalJointMeasure-map-stdGaussian, GaussianHilbert-bonami-nelson-chaosLE, Pphi2-CanonicalJoint}
  \texttt{Pphi2.canonicalFullInteractionJoint\_memLp}
\end{theorem}
\begin{theorem}[\texttt{canonicalRoughError\_leading\_l2\_sq}]\label{Pphi2-canonicalRoughError-leading-l2-sq}
  \lean{Pphi2.canonicalRoughError_leading_l2_sq}\leanok
  \uses{Pphi2-canonicalJointMeasure, Pphi2-canonicalCrossTerm-scaled-inner-sum-l2-sq, Pphi2-canonicalCrossTerm, Pphi2-CanonicalJoint}
  \texttt{Pphi2.canonicalRoughError\_leading\_l2\_sq}
\end{theorem}
\begin{theorem}[\texttt{canonicalCrossTerm\_pair\_integrable}]\label{Pphi2-canonicalCrossTerm-pair-integrable}
  \lean{Pphi2.canonicalCrossTerm_pair_integrable}\leanok
  \uses{Pphi2-smoothWickConstant, Pphi2-canonicalJointMeasure, Pphi2-canonicalRoughFieldFunctionOfSnd, Pphi2-canonicalRoughWickPower-two-site-marginal-integrable, GaussianField-FinLatticeSites, Pphi2-wickMonomial, Pphi2-canonicalSmoothFieldFunctionOfFst, Pphi2-canonicalSmoothWickPower-two-site-marginal-integrable, Pphi2-canonicalRoughFieldFunction, Pphi2-canonicalSmoothFieldFunction, Pphi2-roughWickConstant, Pphi2-canonicalCrossTerm, Pphi2-CanonicalJoint}
  \texttt{Pphi2.canonicalCrossTerm\_pair\_integrable}
\end{theorem}
\begin{definition}[\texttt{canonicalRoughError}]\label{Pphi2-canonicalRoughError}
  \lean{Pphi2.canonicalRoughError}\leanok
  \uses{Pphi2-canonicalFullInteractionJoint, Pphi2-canonicalSmoothInteraction, Pphi2-CanonicalJoint}
  \texttt{Pphi2.canonicalRoughError}
\end{definition}
\begin{theorem}[\texttt{canonicalRoughError\_perCoef\_outer\_l2\_sq}]\label{Pphi2-canonicalRoughError-perCoef-outer-l2-sq}
  \lean{Pphi2.canonicalRoughError_perCoef_outer_l2_sq}\leanok
  \uses{Pphi2-canonicalJointMeasure, Pphi2-integral-sq-real-sum-of-pairwise-orthogonal, Pphi2-canonicalCrossTerm, Pphi2-CanonicalJoint}
  \texttt{Pphi2.canonicalRoughError\_perCoef\_outer\_l2\_sq}
\end{theorem}
\begin{theorem}[\texttt{canonicalRoughError\_pointwise\_decomposition}]\label{Pphi2-canonicalRoughError-pointwise-decomposition}
  \lean{Pphi2.canonicalRoughError_pointwise_decomposition}\leanok
  \uses{Pphi2-smoothWickConstant, Pphi2-wickConstant-split, GaussianField-FinLatticeSites, Pphi2-wickMonomial, Pphi2-canonicalRoughError-eq-sum-diff, Pphi2-canonicalSumFieldFunction, Pphi2-canonicalRoughFieldFunction, Pphi2-canonicalSumFieldFunction-eq-smooth-plus-rough, Pphi2-wickConstant, Pphi2-canonicalSmoothFieldFunction, Pphi2-roughWickConstant, Pphi2-canonicalRoughError, Pphi2-wickPolynomial-add-sub-self, Pphi2-wickPolynomial, Pphi2-CanonicalJoint}
  \texttt{Pphi2.canonicalRoughError\_pointwise\_decomposition}
\end{theorem}
\begin{definition}[\texttt{wickExpansionCoeff}]\label{Pphi2-wickExpansionCoeff}
  \lean{Pphi2.wickExpansionCoeff}\leanok
  \texttt{Pphi2.wickExpansionCoeff}
\end{definition}
\begin{theorem}[\texttt{canonicalRoughError\_moment\_le}]\label{Pphi2-canonicalRoughError-moment-le}
  \lean{Pphi2.canonicalRoughError_moment_le}\leanok
  \uses{Pphi2-canonicalJointMeasure, Pphi2-canonicalJointStdGaussianMeasurableEquiv, GaussianHilbert-stdGaussianFin, GaussianHilbert-wienerChaosLE, Pphi2-CanonicalJointSumIndex, Pphi2-canonicalJointMeasure-map-stdGaussian, Pphi2-canonicalRoughError, Pphi2-chaosLE-moment-le, Pphi2-CanonicalJoint}
  \texttt{Pphi2.canonicalRoughError\_moment\_le}
\end{theorem}
\begin{theorem}[\texttt{canonicalRoughError\_l2\_sq\_eq}]\label{Pphi2-canonicalRoughError-l2-sq-eq}
  \lean{Pphi2.canonicalRoughError_l2_sq_eq}\leanok
  \uses{Pphi2-canonicalRoughError-l2-sq-eq-lead-plus-perCoef-sq, Pphi2-canonicalJointMeasure, Pphi2-canonicalRoughError-perCoef-outer-l2-sq, Pphi2-canonicalRoughError-perCoeff-l2-sq, Pphi2-canonicalRoughError-leading-l2-sq, Pphi2-canonicalRoughError, Pphi2-canonicalCrossTerm, Pphi2-CanonicalJoint}
  \texttt{Pphi2.canonicalRoughError\_l2\_sq\_eq}
\end{theorem}
\begin{theorem}[\texttt{congr\_simp}]\label{Pphi2-LatticeSetup-latticeSmoothInteraction-congr-simp}
  \lean{Pphi2.LatticeSetup.latticeSmoothInteraction.congr_simp}\leanok
  \uses{Pphi2-LatticeSetup-latticeSmoothInteraction, GaussianField-FinLatticeSites, GaussianField-Configuration, GaussianField-FinLatticeField}
  \texttt{Pphi2.LatticeSetup.latticeSmoothInteraction.congr\_simp}
\end{theorem}
\begin{theorem}[\texttt{canonicalRoughError\_eq\_sum\_diff}]\label{Pphi2-canonicalRoughError-eq-sum-diff}
  \lean{Pphi2.canonicalRoughError_eq_sum_diff}\leanok
  \uses{Pphi2-smoothWickConstant, GaussianField-FinLatticeSites, Pphi2-canonicalFullInteractionJoint, Pphi2-canonicalSmoothInteraction, Pphi2-canonicalSumFieldFunction, Pphi2-wickConstant, Pphi2-canonicalSmoothFieldFunction, Pphi2-canonicalRoughError, Pphi2-wickPolynomial, Pphi2-CanonicalJoint}
  \texttt{Pphi2.canonicalRoughError\_eq\_sum\_diff}
\end{theorem}
\begin{theorem}[\texttt{canonicalRoughError\_sq\_integrable}]\label{Pphi2-canonicalRoughError-sq-integrable}
  \lean{Pphi2.canonicalRoughError_sq_integrable}\leanok
  \uses{Pphi2-canonicalJointMeasure, Pphi2-canonicalJointStdGaussianMeasurableEquiv, GaussianHilbert-stdGaussianFin, GaussianHilbert-wienerChaosLE, Pphi2-CanonicalJointSumIndex, Pphi2-canonicalJointMeasure-map-stdGaussian, Pphi2-canonicalRoughError, Pphi2-CanonicalJoint}
  \texttt{Pphi2.canonicalRoughError\_sq\_integrable}
\end{theorem}
\begin{theorem}[\texttt{canonicalRoughError\_neg\_tail\_of\_stdGaussian}]\label{Pphi2-canonicalRoughError-neg-tail-of-stdGaussian}
  \lean{Pphi2.canonicalRoughError_neg_tail_of_stdGaussian}\leanok
  \uses{Pphi2-canonicalJointMeasure, Pphi2-ChaosTailBridge-chaos-neg-tail-bound, Pphi2-canonicalJointStdGaussianMeasurableEquiv, GaussianHilbert-stdGaussianFin, GaussianHilbert-wienerChaosLE, Pphi2-CanonicalJointSumIndex, Pphi2-canonicalJointMeasure-map-stdGaussian, Pphi2-canonicalRoughError, Pphi2-CanonicalJoint}
  \texttt{Pphi2.canonicalRoughError\_neg\_tail\_of\_stdGaussian}
\end{theorem}
\begin{theorem}[\texttt{integrable\_sum\_mul\_sum\_of\_pairwise}]\label{Pphi2-integrable-sum-mul-sum-of-pairwise}
  \lean{Pphi2.integrable_sum_mul_sum_of_pairwise}\leanok
  \texttt{Pphi2.integrable\_sum\_mul\_sum\_of\_pairwise}
\end{theorem}
\begin{theorem}[\texttt{canonicalCrossTerm\_l2\_sq\_le}]\label{Pphi2-canonicalCrossTerm-l2-sq-le}
  \lean{Pphi2.canonicalCrossTerm_l2_sq_le}\leanok
  \uses{Pphi2-smoothWickConstant, Pphi2-canonicalJointMeasure, Pphi2-canonicalCrossTerm-l2-sq-eq-covSum, Pphi2-canonicalRoughCovariance, Pphi2-canonicalSmoothCovariance, Pphi2-smoothWickConstant-le-log-uniform-in-aN, GaussianField-FinLatticeSites, Pphi2-canonicalRoughCovariance-pow-sum-le-uniform-in-aN, GaussianField-latticeEigenvalue-pos, Pphi2-abs-canonicalSmoothCovariance-le-smoothWickConstant, Pphi2-smoothCovEigenvalue, GaussianField-latticeEigenvalue, Pphi2-canonicalCrossTerm, Lattice-toSemilatticeInf, Pphi2-CanonicalJoint}
  \texttt{Pphi2.canonicalCrossTerm\_l2\_sq\_le}
\end{theorem}
\begin{theorem}[\texttt{canonicalLeading\_perCoeff\_inner\_eq\_zero}]\label{Pphi2-canonicalLeading-perCoeff-inner-eq-zero}
  \lean{Pphi2.canonicalLeading_perCoeff_inner_eq_zero}\leanok
  \uses{Pphi2-canonicalJointMeasure, Pphi2-integral-sum-mul-sum-eq-zero-of-orth, Pphi2-canonicalCrossTerm, Pphi2-CanonicalJoint}
  \texttt{Pphi2.canonicalLeading\_perCoeff\_inner\_eq\_zero}
\end{theorem}
\begin{theorem}[\texttt{expMoment\_bound\_of\_cutoff\_quartic\_tail}]\label{Pphi2-expMoment-bound-of-cutoff-quartic-tail}
  \lean{Pphi2.expMoment_bound_of_cutoff_quartic_tail}\leanok
  \uses{Pphi2-LatticeSetup-latticeSmoothInteraction, Pphi2-canonicalJointMeasure, Pphi2-canonicalSumConfig, Pphi2-canonicalSmoothConfig, Pphi2-canonicalFullInteractionJoint-eq-interactionFunctional, GaussianField-FinLatticeSites, Pphi2-canonicalFullInteractionJoint, Pphi2-canonicalSumConfig-measurable, Pphi2-DynamicalCutoff-dynamicalCutoffScale-pos, Pphi2-canonicalSmoothInteraction-eq-latticeSmoothInteraction, Pphi2-canonicalSmoothInteraction, GaussianField-latticeGaussianMeasure, GaussianField-Configuration, Pphi2-DynamicalCutoff-dynamicalCutoffScale, GaussianField-instMeasurableSpaceConfiguration, Pphi2-canonicalJointMeasure-map-canonicalSumConfig, Pphi2-LatticeBridge-bridgeAxiom-of-varying-coupled-setup-real, Pphi2-canonicalJointMeasure-isProbability, GaussianField-FinLatticeField, Pphi2-canonicalRoughError, Pphi2-interactionFunctional, Pphi2-CanonicalJoint}
  \texttt{Pphi2.expMoment\_bound\_of\_cutoff\_quartic\_tail}
\end{theorem}
\begin{theorem}[\texttt{expMoment\_bound\_of\_cutoff\_degree\_tail}]\label{Pphi2-expMoment-bound-of-cutoff-degree-tail}
  \lean{Pphi2.expMoment_bound_of_cutoff_degree_tail}\leanok
  \uses{Pphi2-LatticeSetup-latticeSmoothInteraction, Pphi2-canonicalJointMeasure, Pphi2-canonicalSumConfig, Pphi2-canonicalSmoothConfig, Pphi2-canonicalFullInteractionJoint-eq-interactionFunctional, GaussianField-FinLatticeSites, Pphi2-canonicalFullInteractionJoint, Pphi2-DynamicalCutoff-degreeDynamicalCutoffScale-pos, Pphi2-canonicalSumConfig-measurable, Pphi2-canonicalSmoothInteraction-eq-latticeSmoothInteraction, Pphi2-canonicalSmoothInteraction, GaussianField-latticeGaussianMeasure, Pphi2-DynamicalCutoff-degreeDynamicalCutoffScale, GaussianField-Configuration, GaussianField-instMeasurableSpaceConfiguration, Pphi2-canonicalJointMeasure-map-canonicalSumConfig, Pphi2-LatticeBridge-bridgeAxiom-of-varying-coupled-setup-real, Pphi2-canonicalJointMeasure-isProbability, GaussianField-FinLatticeField, Pphi2-canonicalRoughError, Pphi2-interactionFunctional, Pphi2-CanonicalJoint}
  \texttt{Pphi2.expMoment\_bound\_of\_cutoff\_degree\_tail}
\end{theorem}
\begin{theorem}[\texttt{canonicalSmoothInteraction\_eq\_latticeSmoothInteraction}]\label{Pphi2-canonicalSmoothInteraction-eq-latticeSmoothInteraction}
  \lean{Pphi2.canonicalSmoothInteraction_eq_latticeSmoothInteraction}\leanok
  \uses{Pphi2-LatticeSetup-latticeSmoothInteraction, Pphi2-smoothWickConstant, Pphi2-canonicalSmoothConfig, GaussianField-finLatticeDelta, GaussianField-FinLatticeSites, Pphi2-canonicalSmoothInteraction, GaussianField-Configuration, Pphi2-canonicalSmoothFieldFunction, GaussianField-FinLatticeField, Pphi2-wickPolynomial, Pphi2-CanonicalJoint}
  \texttt{Pphi2.canonicalSmoothInteraction\_eq\_latticeSmoothInteraction}
\end{theorem}
\begin{theorem}[\texttt{integral\_sq\_real\_sum\_of\_pairwise\_orthogonal}]\label{Pphi2-integral-sq-real-sum-of-pairwise-orthogonal}
  \lean{Pphi2.integral_sq_real_sum_of_pairwise_orthogonal}\leanok
  \texttt{Pphi2.integral\_sq\_real\_sum\_of\_pairwise\_orthogonal}
\end{theorem}
\begin{theorem}[\texttt{canonicalRoughError\_perCoeff\_l2\_sq}]\label{Pphi2-canonicalRoughError-perCoeff-l2-sq}
  \lean{Pphi2.canonicalRoughError_perCoeff_l2_sq}\leanok
  \uses{Pphi2-canonicalJointMeasure, Pphi2-canonicalCrossTerm-scaled-inner-sum-l2-sq, Pphi2-canonicalCrossTerm, Pphi2-CanonicalJoint}
  \texttt{Pphi2.canonicalRoughError\_perCoeff\_l2\_sq}
\end{theorem}
\begin{definition}[\texttt{canonicalCrossTerm}]\label{Pphi2-canonicalCrossTerm}
  \lean{Pphi2.canonicalCrossTerm}\leanok
  \uses{Pphi2-smoothWickConstant, GaussianField-FinLatticeSites, Pphi2-wickMonomial, Pphi2-canonicalRoughFieldFunction, Pphi2-canonicalSmoothFieldFunction, Pphi2-roughWickConstant, Pphi2-CanonicalJoint}
  \texttt{Pphi2.canonicalCrossTerm}
\end{definition}
\begin{theorem}[\texttt{canonicalSmoothInteraction\_lower\_bound\_at\_cutoff\_quartic}]\label{Pphi2-canonicalSmoothInteraction-lower-bound-at-cutoff-quartic}
  \lean{Pphi2.canonicalSmoothInteraction_lower_bound_at_cutoff_quartic}\leanok
  \uses{Pphi2-LatticeSetup-latticeSmoothInteraction, Pphi2-canonicalSmoothConfig, Pphi2-canonicalSmoothInteraction-eq-latticeSmoothInteraction, Pphi2-smoothBoundConstant, Pphi2-canonicalSmoothInteraction, Pphi2-LatticeSetup-latticeSmoothInteraction-lower-bound-at-cutoff-quartic, Pphi2-DynamicalCutoff-dynamicalCutoffScale, Pphi2-CanonicalJoint}
  \texttt{Pphi2.canonicalSmoothInteraction\_lower\_bound\_at\_cutoff\_quartic}
\end{theorem}
\begin{theorem}[\texttt{canonicalSmoothInteraction\_lower\_bound\_log\_uniform\_in\_aN}]\label{Pphi2-canonicalSmoothInteraction-lower-bound-log-uniform-in-aN}
  \lean{Pphi2.canonicalSmoothInteraction_lower_bound_log_uniform_in_aN}\leanok
  \uses{Pphi2-wickPolynomial-lower-bound-general, Pphi2-LatticeSetup-latticeSmoothInteraction, Pphi2-smoothWickConstant, Pphi2-smoothCovEigenvalue-pos, Pphi2-canonicalSmoothConfig, Pphi2-smoothWickConstant-le-log-uniform-in-aN, GaussianField-finLatticeDelta, GaussianField-FinLatticeSites, Pphi2-canonicalSmoothInteraction-eq-latticeSmoothInteraction, Pphi2-canonicalSmoothInteraction, GaussianField-Configuration, Pphi2-smoothCovEigenvalue, Pphi2-DynamicalCutoff-degreeCutoffPower, GaussianField-FinLatticeField, Lattice-toSemilatticeInf, Pphi2-wickPolynomial, Pphi2-CanonicalJoint}
  \texttt{Pphi2.canonicalSmoothInteraction\_lower\_bound\_log\_uniform\_in\_aN}
\end{theorem}
\section{\texttt{Pphi2.WickOrdering.Counterterm}}
\begin{definition}[\texttt{wickConstant}]\label{Pphi2-wickConstant}
  \lean{Pphi2.wickConstant}\leanok
  \uses{GaussianField-FinLatticeSites, GaussianField-latticeEigenvalue}
  \texttt{Pphi2.wickConstant}
\end{definition}
\begin{theorem}[\texttt{wickConstant\_eq\_latticeEigenvalue1d\_family\_average}]\label{Pphi2-wickConstant-eq-latticeEigenvalue1d-family-average}
  \lean{Pphi2.wickConstant_eq_latticeEigenvalue1d_family_average}\leanok
  \uses{GaussianField-latticeEigenvalue1d, GaussianField-FinLatticeSites, Pphi2-wickConstant, GaussianField-sum-latticeEigenvalue-eq-sum-latticeEigenvalue1d-family, GaussianField-latticeEigenvalue}
  \texttt{Pphi2.wickConstant\_eq\_latticeEigenvalue1d\_family\_average}
\end{theorem}
\begin{theorem}[\texttt{wickConstant\_pos}]\label{Pphi2-wickConstant-pos}
  \lean{Pphi2.wickConstant_pos}\leanok
  \uses{GaussianField-FinLatticeSites, GaussianField-latticeEigenvalue-pos, Pphi2-wickConstant, GaussianField-latticeEigenvalue}
  $c_a > 0$ when $a > 0$ and $\text{mass} > 0$. Each eigenvalue $\lambda_m > 0$, so each inverse is positive, and the average of positive terms is positive.
\end{theorem}
\begin{theorem}[\texttt{wickConstant\_le\_inv\_mass\_sq}]\label{Pphi2-wickConstant-le-inv-mass-sq}
  \lean{Pphi2.wickConstant_le_inv_mass_sq}\leanok
  \uses{GaussianField-FinLatticeSites, GaussianField-latticeLaplacianEigenvalue, Pphi2-wickConstant, GaussianField-latticeEigenvalue, Lattice-toSemilatticeInf}
  $c_a \le m^{-2}$ since each $\lambda_m \ge m^2$, so each $1/\lambda_m \le 1/m^2$.
\end{theorem}
\begin{theorem}[\texttt{wickConstant\_antitone\_mass}]\label{Pphi2-wickConstant-antitone-mass}
  \lean{Pphi2.wickConstant_antitone_mass}\leanok
  \uses{GaussianField-FinLatticeSites, GaussianField-latticeLaplacianEigenvalue, GaussianField-latticeEigenvalue-pos, Pphi2-wickConstant, GaussianField-latticeEigenvalue}
  $c_a$ is monotone decreasing in mass: $m_1 \le m_2 \implies c_a(m_2) \le c_a(m_1)$, since larger mass increases eigenvalues.
\end{theorem}
\begin{theorem}[\texttt{wickConstant\_two\_eq\_dft\_eigenvalue\_average}]\label{Pphi2-wickConstant-two-eq-dft-eigenvalue-average}
  \lean{Pphi2.wickConstant_two_eq_dft_eigenvalue_average}\leanok
  \uses{GaussianField-latticeEigenvalue1d, GaussianField-sum-latticeEigenvalue-two-eq-sum-latticeEigenvalue1d-pair, GaussianField-FinLatticeSites, Pphi2-wickConstant, GaussianField-latticeEigenvalue}
  \texttt{Pphi2.wickConstant\_two\_eq\_dft\_eigenvalue\_average}
\end{theorem}
\section{\texttt{Pphi2.Backgrounds.EuclideanPlane}}
\begin{definition}[\texttt{continuumEuclideanAction}]\label{Pphi2-continuumEuclideanAction}
  \lean{Pphi2.continuumEuclideanAction}\leanok
  \uses{Pphi2-EuclideanPlaneBackground-SpaceTime, Pphi2-EuclideanPlaneBackground-action, Pphi2-planeBackground, Pphi2-EuclideanPlaneBackground-dim, Pphi2-ContinuumMotion, Pphi2-ContinuumTestFunction}
  \texttt{Pphi2.continuumEuclideanAction}
\end{definition}
\begin{definition}[\texttt{ContinuumComplexTestFunction}]\label{Pphi2-ContinuumComplexTestFunction}
  \lean{Pphi2.ContinuumComplexTestFunction}\leanok
  \uses{Pphi2-planeBackground, Pphi2-EuclideanPlaneBackground-ComplexTestFunction}
  \texttt{Pphi2.ContinuumComplexTestFunction}
\end{definition}
\begin{definition}[\texttt{t}]\label{Pphi2-EuclideanPlaneBackground-Motion-t}
  \lean{Pphi2.EuclideanPlaneBackground.Motion.t}\leanok
  \uses{Pphi2-EuclideanPlaneBackground-SpaceTime, Pphi2-EuclideanPlaneBackground-Motion, Pphi2-EuclideanPlaneBackground}
  \texttt{Pphi2.EuclideanPlaneBackground.Motion.t}
\end{definition}
\begin{theorem}[\texttt{translate\_apply}]\label{Pphi2-EuclideanPlaneBackground-translate-apply}
  \lean{Pphi2.EuclideanPlaneBackground.translate_apply}\leanok
  \uses{Pphi2-EuclideanPlaneBackground-SpaceTime, Pphi2-EuclideanPlaneBackground-action, Pphi2-EuclideanPlaneBackground-RealTestFunction, Pphi2-EuclideanPlaneBackground-translate, Pphi2-EuclideanPlaneBackground-inversePointAction, Pphi2-EuclideanPlaneBackground-dim, Pphi2-EuclideanPlaneBackground-translationMotion, Pphi2-EuclideanPlaneBackground, Pphi2-EuclideanPlaneBackground-action-apply, Pphi2-EuclideanPlaneBackground-inversePointAction-translationMotion}
  \texttt{Pphi2.EuclideanPlaneBackground.translate\_apply}
\end{theorem}
\begin{definition}[\texttt{ContinuumTestFunction}]\label{Pphi2-ContinuumTestFunction}
  \lean{Pphi2.ContinuumTestFunction}\leanok
  \uses{Pphi2-EuclideanPlaneBackground-RealTestFunction, Pphi2-planeBackground}
  \texttt{Pphi2.ContinuumTestFunction}
\end{definition}
\begin{theorem}[\texttt{schwartzTranslate\_apply}]\label{Pphi2-schwartzTranslate-apply}
  \lean{Pphi2.schwartzTranslate_apply}\leanok
  \uses{Pphi2-EuclideanPlaneBackground-SpaceTime, Pphi2-planeBackground, Pphi2-ContinuumSpaceTime, Pphi2-EuclideanPlaneBackground-dim, Pphi2-ContinuumTestFunction, Pphi2-EuclideanPlaneBackground-translate-apply, Pphi2-schwartzTranslate}
  \texttt{Pphi2.schwartzTranslate\_apply}
\end{theorem}
\begin{definition}[\texttt{Motion}]\label{Pphi2-EuclideanPlaneBackground-Motion}
  \lean{Pphi2.EuclideanPlaneBackground.Motion}\leanok
  \uses{Pphi2-EuclideanPlaneBackground}
  \texttt{Pphi2.EuclideanPlaneBackground.Motion}
\end{definition}
\begin{theorem}[\texttt{actionComplex\_apply}]\label{Pphi2-EuclideanPlaneBackground-actionComplex-apply}
  \lean{Pphi2.EuclideanPlaneBackground.actionComplex_apply}\leanok
  \uses{Pphi2-EuclideanPlaneBackground-SpaceTime, Pphi2-EuclideanPlaneBackground-inversePointAction, Pphi2-EuclideanPlaneBackground-dim, Pphi2-EuclideanPlaneBackground-Motion, Pphi2-EuclideanPlaneBackground, Pphi2-EuclideanPlaneBackground-actionComplex, Pphi2-EuclideanPlaneBackground-ComplexTestFunction}
  \texttt{Pphi2.EuclideanPlaneBackground.actionComplex\_apply}
\end{theorem}
\begin{definition}[\texttt{R}]\label{Pphi2-EuclideanPlaneBackground-Motion-R}
  \lean{Pphi2.EuclideanPlaneBackground.Motion.R}\leanok
  \uses{Pphi2-EuclideanPlaneBackground-OrthogonalGroup, Pphi2-EuclideanPlaneBackground-Motion, Pphi2-EuclideanPlaneBackground}
  \texttt{Pphi2.EuclideanPlaneBackground.Motion.R}
\end{definition}
\begin{definition}[\texttt{dim}]\label{Pphi2-EuclideanPlaneBackground-dim}
  \lean{Pphi2.EuclideanPlaneBackground.dim}\leanok
  \uses{Pphi2-EuclideanPlaneBackground}
  \texttt{Pphi2.EuclideanPlaneBackground.dim}
\end{definition}
\begin{definition}[\texttt{ComplexTestFunction}]\label{Pphi2-EuclideanPlaneBackground-ComplexTestFunction}
  \lean{Pphi2.EuclideanPlaneBackground.ComplexTestFunction}\leanok
  \uses{Pphi2-EuclideanPlaneBackground-SpaceTime, Pphi2-EuclideanPlaneBackground-dim, Pphi2-EuclideanPlaneBackground}
  \texttt{Pphi2.EuclideanPlaneBackground.ComplexTestFunction}
\end{definition}
\begin{theorem}[\texttt{ContinuumFieldConfig\_eq}]\label{Pphi2-ContinuumFieldConfig-eq}
  \lean{Pphi2.ContinuumFieldConfig_eq}\leanok
  \uses{GaussianField-Configuration, Pphi2-ContinuumFieldConfig}
  \texttt{Pphi2.ContinuumFieldConfig\_eq}
\end{theorem}
\begin{theorem}[\texttt{continuumEuclideanAction\_apply}]\label{Pphi2-continuumEuclideanAction-apply}
  \lean{Pphi2.continuumEuclideanAction_apply}\leanok
  \uses{Pphi2-EuclideanPlaneBackground-SpaceTime, Pphi2-planeBackground, Pphi2-EuclideanPlaneBackground-inversePointAction, Pphi2-ContinuumSpaceTime, Pphi2-EuclideanPlaneBackground-dim, Pphi2-EuclideanPlaneBackground-action-apply, Pphi2-continuumEuclideanAction, Pphi2-ContinuumMotion, Pphi2-ContinuumTestFunction}
  \texttt{Pphi2.continuumEuclideanAction\_apply}
\end{theorem}
\begin{definition}[\texttt{translateComplex}]\label{Pphi2-EuclideanPlaneBackground-translateComplex}
  \lean{Pphi2.EuclideanPlaneBackground.translateComplex}\leanok
  \uses{Pphi2-EuclideanPlaneBackground-SpaceTime, Pphi2-EuclideanPlaneBackground-dim, Pphi2-EuclideanPlaneBackground-translationMotion, Pphi2-EuclideanPlaneBackground, Pphi2-EuclideanPlaneBackground-actionComplex, Pphi2-EuclideanPlaneBackground-ComplexTestFunction}
  \texttt{Pphi2.EuclideanPlaneBackground.translateComplex}
\end{definition}
\begin{theorem}[\texttt{ContinuumOrthogonalGroup\_eq}]\label{Pphi2-ContinuumOrthogonalGroup-eq}
  \lean{Pphi2.ContinuumOrthogonalGroup_eq}\leanok
  \uses{Pphi2-ContinuumOrthogonalGroup}
  \texttt{Pphi2.ContinuumOrthogonalGroup\_eq}
\end{theorem}
\begin{definition}[\texttt{SpaceTime}]\label{Pphi2-EuclideanPlaneBackground-SpaceTime}
  \lean{Pphi2.EuclideanPlaneBackground.SpaceTime}\leanok
  \uses{Pphi2-EuclideanPlaneBackground-dim, Pphi2-EuclideanPlaneBackground}
  \texttt{Pphi2.EuclideanPlaneBackground.SpaceTime}
\end{definition}
\begin{definition}[\texttt{EuclideanPlaneBackground}]\label{Pphi2-EuclideanPlaneBackground}
  \lean{Pphi2.EuclideanPlaneBackground}\leanok
  \texttt{Pphi2.EuclideanPlaneBackground}
\end{definition}
\begin{definition}[\texttt{actionComplex}]\label{Pphi2-EuclideanPlaneBackground-actionComplex}
  \lean{Pphi2.EuclideanPlaneBackground.actionComplex}\leanok
  \uses{Pphi2-EuclideanPlaneBackground-SpaceTime, Pphi2-EuclideanPlaneBackground-inversePointAction, Pphi2-EuclideanPlaneBackground-dim, Pphi2-EuclideanPlaneBackground-Motion, Pphi2-EuclideanPlaneBackground, Pphi2-EuclideanPlaneBackground-ComplexTestFunction}
  \texttt{Pphi2.EuclideanPlaneBackground.actionComplex}
\end{definition}
\begin{theorem}[\texttt{ContinuumTestFunction\_eq}]\label{Pphi2-ContinuumTestFunction-eq}
  \lean{Pphi2.ContinuumTestFunction_eq}\leanok
  \uses{Pphi2-ContinuumTestFunction}
  \texttt{Pphi2.ContinuumTestFunction\_eq}
\end{theorem}
\begin{definition}[\texttt{Distribution}]\label{Pphi2-EuclideanPlaneBackground-Distribution}
  \lean{Pphi2.EuclideanPlaneBackground.Distribution}\leanok
  \uses{Pphi2-EuclideanPlaneBackground-SpaceTime, Pphi2-EuclideanPlaneBackground-RealTestFunction, Pphi2-EuclideanPlaneBackground-dim, Pphi2-EuclideanPlaneBackground, GaussianField-Configuration}
  \texttt{Pphi2.EuclideanPlaneBackground.Distribution}
\end{definition}
\begin{theorem}[\texttt{planeBackground\_dim}]\label{Pphi2-planeBackground-dim}
  \lean{Pphi2.planeBackground_dim}\leanok
  \uses{Pphi2-planeBackground, Pphi2-EuclideanPlaneBackground-dim}
  \texttt{Pphi2.planeBackground\_dim}
\end{theorem}
\begin{definition}[\texttt{translate}]\label{Pphi2-EuclideanPlaneBackground-translate}
  \lean{Pphi2.EuclideanPlaneBackground.translate}\leanok
  \uses{Pphi2-EuclideanPlaneBackground-SpaceTime, Pphi2-EuclideanPlaneBackground-action, Pphi2-EuclideanPlaneBackground-RealTestFunction, Pphi2-EuclideanPlaneBackground-dim, Pphi2-EuclideanPlaneBackground-translationMotion, Pphi2-EuclideanPlaneBackground}
  \texttt{Pphi2.EuclideanPlaneBackground.translate}
\end{definition}
\begin{theorem}[\texttt{ContinuumComplexTestFunction\_eq}]\label{Pphi2-ContinuumComplexTestFunction-eq}
  \lean{Pphi2.ContinuumComplexTestFunction_eq}\leanok
  \uses{Pphi2-ContinuumComplexTestFunction}
  \texttt{Pphi2.ContinuumComplexTestFunction\_eq}
\end{theorem}
\begin{definition}[\texttt{ContinuumOrthogonalGroup}]\label{Pphi2-ContinuumOrthogonalGroup}
  \lean{Pphi2.ContinuumOrthogonalGroup}\leanok
  \uses{Pphi2-EuclideanPlaneBackground-OrthogonalGroup, Pphi2-planeBackground}
  \texttt{Pphi2.ContinuumOrthogonalGroup}
\end{definition}
\begin{definition}[\texttt{inversePointAction}]\label{Pphi2-EuclideanPlaneBackground-inversePointAction}
  \lean{Pphi2.EuclideanPlaneBackground.inversePointAction}\leanok
  \uses{Pphi2-EuclideanPlaneBackground-SpaceTime, Pphi2-EuclideanPlaneBackground-Motion-t, Pphi2-EuclideanPlaneBackground-Motion-R, Pphi2-EuclideanPlaneBackground-dim, Pphi2-EuclideanPlaneBackground-Motion, Pphi2-EuclideanPlaneBackground}
  \texttt{Pphi2.EuclideanPlaneBackground.inversePointAction}
\end{definition}
\begin{definition}[\texttt{schwartzTranslate}]\label{Pphi2-schwartzTranslate}
  \lean{Pphi2.schwartzTranslate}\leanok
  \uses{Pphi2-EuclideanPlaneBackground-SpaceTime, Pphi2-planeBackground, Pphi2-EuclideanPlaneBackground-translate, Pphi2-ContinuumSpaceTime, Pphi2-EuclideanPlaneBackground-dim, Pphi2-ContinuumTestFunction}
  \texttt{Pphi2.schwartzTranslate}
\end{definition}
\begin{theorem}[\texttt{translateComplex\_apply}]\label{Pphi2-EuclideanPlaneBackground-translateComplex-apply}
  \lean{Pphi2.EuclideanPlaneBackground.translateComplex_apply}\leanok
  \uses{Pphi2-EuclideanPlaneBackground-SpaceTime, Pphi2-EuclideanPlaneBackground-inversePointAction, Pphi2-EuclideanPlaneBackground-translateComplex, Pphi2-EuclideanPlaneBackground-dim, Pphi2-EuclideanPlaneBackground-translationMotion, Pphi2-EuclideanPlaneBackground-actionComplex-apply, Pphi2-EuclideanPlaneBackground, Pphi2-EuclideanPlaneBackground-actionComplex, Pphi2-EuclideanPlaneBackground-ComplexTestFunction, Pphi2-EuclideanPlaneBackground-inversePointAction-translationMotion}
  \texttt{Pphi2.EuclideanPlaneBackground.translateComplex\_apply}
\end{theorem}
\begin{definition}[\texttt{planeBackground}]\label{Pphi2-planeBackground}
  \lean{Pphi2.planeBackground}\leanok
  \uses{Pphi2-EuclideanPlaneBackground}
  \texttt{Pphi2.planeBackground}
\end{definition}
\begin{definition}[\texttt{OrthogonalGroup}]\label{Pphi2-EuclideanPlaneBackground-OrthogonalGroup}
  \lean{Pphi2.EuclideanPlaneBackground.OrthogonalGroup}\leanok
  \uses{Pphi2-EuclideanPlaneBackground-SpaceTime, Pphi2-EuclideanPlaneBackground-dim, Pphi2-EuclideanPlaneBackground}
  \texttt{Pphi2.EuclideanPlaneBackground.OrthogonalGroup}
\end{definition}
\begin{definition}[\texttt{translationMotion}]\label{Pphi2-EuclideanPlaneBackground-translationMotion}
  \lean{Pphi2.EuclideanPlaneBackground.translationMotion}\leanok
  \uses{Pphi2-EuclideanPlaneBackground-SpaceTime, Pphi2-EuclideanPlaneBackground-dim, Pphi2-EuclideanPlaneBackground-Motion, Pphi2-EuclideanPlaneBackground}
  \texttt{Pphi2.EuclideanPlaneBackground.translationMotion}
\end{definition}
\begin{definition}[\texttt{ctorIdx}]\label{Pphi2-EuclideanPlaneBackground-Motion-ctorIdx}
  \lean{Pphi2.EuclideanPlaneBackground.Motion.ctorIdx}\leanok
  \uses{Pphi2-EuclideanPlaneBackground-Motion, Pphi2-EuclideanPlaneBackground}
  \texttt{Pphi2.EuclideanPlaneBackground.Motion.ctorIdx}
\end{definition}
\begin{theorem}[\texttt{action\_apply}]\label{Pphi2-EuclideanPlaneBackground-action-apply}
  \lean{Pphi2.EuclideanPlaneBackground.action_apply}\leanok
  \uses{Pphi2-EuclideanPlaneBackground-SpaceTime, Pphi2-EuclideanPlaneBackground-action, Pphi2-EuclideanPlaneBackground-RealTestFunction, Pphi2-EuclideanPlaneBackground-inversePointAction, Pphi2-EuclideanPlaneBackground-dim, Pphi2-EuclideanPlaneBackground-Motion, Pphi2-EuclideanPlaneBackground}
  \texttt{Pphi2.EuclideanPlaneBackground.action\_apply}
\end{theorem}
\begin{definition}[\texttt{ContinuumFieldConfig}]\label{Pphi2-ContinuumFieldConfig}
  \lean{Pphi2.ContinuumFieldConfig}\leanok
  \uses{Pphi2-planeBackground, Pphi2-EuclideanPlaneBackground-Distribution}
  \texttt{Pphi2.ContinuumFieldConfig}
\end{definition}
\begin{theorem}[\texttt{continuumEuclideanActionComplex\_apply}]\label{Pphi2-continuumEuclideanActionComplex-apply}
  \lean{Pphi2.continuumEuclideanActionComplex_apply}\leanok
  \uses{Pphi2-EuclideanPlaneBackground-SpaceTime, Pphi2-planeBackground, Pphi2-EuclideanPlaneBackground-inversePointAction, Pphi2-ContinuumSpaceTime, Pphi2-EuclideanPlaneBackground-dim, Pphi2-continuumEuclideanActionComplex, Pphi2-ContinuumComplexTestFunction, Pphi2-EuclideanPlaneBackground-actionComplex-apply, Pphi2-ContinuumMotion}
  \texttt{Pphi2.continuumEuclideanActionComplex\_apply}
\end{theorem}
\begin{definition}[\texttt{action}]\label{Pphi2-EuclideanPlaneBackground-action}
  \lean{Pphi2.EuclideanPlaneBackground.action}\leanok
  \uses{Pphi2-EuclideanPlaneBackground-SpaceTime, Pphi2-EuclideanPlaneBackground-RealTestFunction, Pphi2-EuclideanPlaneBackground-inversePointAction, Pphi2-EuclideanPlaneBackground-dim, Pphi2-EuclideanPlaneBackground-Motion, Pphi2-EuclideanPlaneBackground}
  \texttt{Pphi2.EuclideanPlaneBackground.action}
\end{definition}
\begin{definition}[\texttt{ctorIdx}]\label{Pphi2-EuclideanPlaneBackground-ctorIdx}
  \lean{Pphi2.EuclideanPlaneBackground.ctorIdx}\leanok
  \uses{Pphi2-EuclideanPlaneBackground}
  \texttt{Pphi2.EuclideanPlaneBackground.ctorIdx}
\end{definition}
\begin{theorem}[\texttt{ContinuumSpaceTime\_eq}]\label{Pphi2-ContinuumSpaceTime-eq}
  \lean{Pphi2.ContinuumSpaceTime_eq}\leanok
  \uses{Pphi2-ContinuumSpaceTime}
  \texttt{Pphi2.ContinuumSpaceTime\_eq}
\end{theorem}
\begin{definition}[\texttt{ContinuumSpaceTime}]\label{Pphi2-ContinuumSpaceTime}
  \lean{Pphi2.ContinuumSpaceTime}\leanok
  \uses{Pphi2-EuclideanPlaneBackground-SpaceTime, Pphi2-planeBackground}
  \texttt{Pphi2.ContinuumSpaceTime}
\end{definition}
\begin{theorem}[\texttt{inversePointAction\_translationMotion}]\label{Pphi2-EuclideanPlaneBackground-inversePointAction-translationMotion}
  \lean{Pphi2.EuclideanPlaneBackground.inversePointAction_translationMotion}\leanok
  \uses{Pphi2-EuclideanPlaneBackground-SpaceTime, Pphi2-EuclideanPlaneBackground-inversePointAction, Pphi2-EuclideanPlaneBackground-dim, Pphi2-EuclideanPlaneBackground-translationMotion, Pphi2-EuclideanPlaneBackground}
  \texttt{Pphi2.EuclideanPlaneBackground.inversePointAction\_translationMotion}
\end{theorem}
\begin{definition}[\texttt{RealTestFunction}]\label{Pphi2-EuclideanPlaneBackground-RealTestFunction}
  \lean{Pphi2.EuclideanPlaneBackground.RealTestFunction}\leanok
  \uses{Pphi2-EuclideanPlaneBackground-SpaceTime, Pphi2-EuclideanPlaneBackground-dim, Pphi2-EuclideanPlaneBackground}
  \texttt{Pphi2.EuclideanPlaneBackground.RealTestFunction}
\end{definition}
\begin{definition}[\texttt{continuumEuclideanActionComplex}]\label{Pphi2-continuumEuclideanActionComplex}
  \lean{Pphi2.continuumEuclideanActionComplex}\leanok
  \uses{Pphi2-EuclideanPlaneBackground-SpaceTime, Pphi2-planeBackground, Pphi2-EuclideanPlaneBackground-dim, Pphi2-ContinuumComplexTestFunction, Pphi2-EuclideanPlaneBackground-actionComplex, Pphi2-ContinuumMotion}
  \texttt{Pphi2.continuumEuclideanActionComplex}
\end{definition}
\begin{definition}[\texttt{ContinuumMotion}]\label{Pphi2-ContinuumMotion}
  \lean{Pphi2.ContinuumMotion}\leanok
  \uses{Pphi2-planeBackground, Pphi2-EuclideanPlaneBackground-Motion}
  \texttt{Pphi2.ContinuumMotion}
\end{definition}
\section{\texttt{Pphi2.InteractingMeasure.FieldRedefinition}}
\begin{theorem}[\texttt{interactingMeasure\_zero}]\label{Pphi2-interactingMeasure-zero}
  \lean{Pphi2.interactingMeasure_zero}\leanok
  \uses{Pphi2-interactingPartitionFunction, Pphi2-interactingMeasure, Pphi2-interactingBoltzmannWeight, GaussianField-Configuration, GaussianField-instMeasurableSpaceConfiguration}
  \texttt{Pphi2.interactingMeasure\_zero}
\end{theorem}
\begin{theorem}[\texttt{connectedFourPoint\_map\_const\_smul}]\label{Pphi2-connectedFourPoint-map-const-smul}
  \lean{Pphi2.connectedFourPoint_map_const_smul}\leanok
  \uses{Pphi2-connectedFourPoint, GaussianField-Configuration, Pphi2-integral-pow-map-const-smul, GaussianField-instMeasurableSpaceConfiguration}
  \texttt{Pphi2.connectedFourPoint\_map\_const\_smul}
\end{theorem}
\begin{theorem}[\texttt{fieldRescaleEquiv\_apply}]\label{Pphi2-fieldRescaleEquiv-apply}
  \lean{Pphi2.fieldRescaleEquiv_apply}\leanok
  \uses{Pphi2-fieldRescaleEquiv, GaussianField-Configuration, GaussianField-instMeasurableSpaceConfiguration}
  \texttt{Pphi2.fieldRescaleEquiv\_apply}
\end{theorem}
\begin{definition}[\texttt{fieldRescaleEquiv}]\label{Pphi2-fieldRescaleEquiv}
  \lean{Pphi2.fieldRescaleEquiv}\leanok
  \uses{Pphi2-measurable-configuration-const-smul, GaussianField-Configuration, GaussianField-instMeasurableSpaceConfiguration}
  \texttt{Pphi2.fieldRescaleEquiv}
\end{definition}
\begin{theorem}[\texttt{interactingPartitionFunction\_map\_measurableEquiv}]\label{Pphi2-interactingPartitionFunction-map-measurableEquiv}
  \lean{Pphi2.interactingPartitionFunction_map_measurableEquiv}\leanok
  \uses{Pphi2-interactingPartitionFunction, Pphi2-interactingBoltzmannWeight, GaussianField-Configuration, GaussianField-instMeasurableSpaceConfiguration}
  \texttt{Pphi2.interactingPartitionFunction\_map\_measurableEquiv}
\end{theorem}
\begin{theorem}[\texttt{interactingMeasure\_map\_measurableEquiv}]\label{Pphi2-interactingMeasure-map-measurableEquiv}
  \lean{Pphi2.interactingMeasure_map_measurableEquiv}\leanok
  \uses{Pphi2-interactingBoltzmannWeight-ennreal-measurable, Pphi2-interactingPartitionFunction, Pphi2-interactingMeasure, Pphi2-interactingPartitionFunction-map-measurableEquiv, Pphi2-interactingBoltzmannWeight, GaussianField-Configuration, GaussianField-instMeasurableSpaceConfiguration}
  \texttt{Pphi2.interactingMeasure\_map\_measurableEquiv}
\end{theorem}
\begin{theorem}[\texttt{integral\_pow\_map\_const\_smul}]\label{Pphi2-integral-pow-map-const-smul}
  \lean{Pphi2.integral_pow_map_const_smul}\leanok
  \uses{Pphi2-measurable-configuration-const-smul, GaussianField-configuration-eval-measurable, GaussianField-Configuration, GaussianField-instMeasurableSpaceConfiguration}
  \texttt{Pphi2.integral\_pow\_map\_const\_smul}
\end{theorem}
\begin{theorem}[\texttt{interactingMeasure\_zero\_smul}]\label{Pphi2-interactingMeasure-zero-smul}
  \lean{Pphi2.interactingMeasure_zero_smul}\leanok
  \uses{Pphi2-interactingMeasure, GaussianField-Configuration, Pphi2-interactingMeasure-zero, GaussianField-instMeasurableSpaceConfiguration}
  \texttt{Pphi2.interactingMeasure\_zero\_smul}
\end{theorem}
\begin{theorem}[\texttt{measurable\_configuration\_const\_smul}]\label{Pphi2-measurable-configuration-const-smul}
  \lean{Pphi2.measurable_configuration_const_smul}\leanok
  \uses{GaussianField-configuration-measurable-of-eval-measurable, GaussianField-configuration-eval-measurable, GaussianField-Configuration, GaussianField-instMeasurableSpaceConfiguration}
  \texttt{Pphi2.measurable\_configuration\_const\_smul}
\end{theorem}
\begin{theorem}[\texttt{connectedFourPoint\_gaussianMeasure\_eq\_zero}]\label{Pphi2-connectedFourPoint-gaussianMeasure-eq-zero}
  \lean{Pphi2.connectedFourPoint_gaussianMeasure_eq_zero}\leanok
  \uses{GaussianField-DyninMityaginSpace, GaussianField-measure, Pphi2-connectedFourPoint, GaussianField-integral-pow4-gaussianReal, GaussianField-pairing-is-gaussian, GaussianField-configuration-eval-measurable, GaussianField-Configuration, GaussianField-integral-sq-gaussianReal, GaussianField-instMeasurableSpaceConfiguration}
  \texttt{Pphi2.connectedFourPoint\_gaussianMeasure\_eq\_zero}
\end{theorem}
\begin{theorem}[\texttt{connectedFourPoint\_interactingMeasure\_field\_rescale}]\label{Pphi2-connectedFourPoint-interactingMeasure-field-rescale}
  \lean{Pphi2.connectedFourPoint_interactingMeasure_field_rescale}\leanok
  \uses{Pphi2-interactingMeasure, Pphi2-connectedFourPoint, Pphi2-fieldRescaleEquiv, Pphi2-interactingMeasure-map-measurableEquiv, GaussianField-Configuration, GaussianField-instMeasurableSpaceConfiguration, Pphi2-connectedFourPoint-map-const-smul}
  \texttt{Pphi2.connectedFourPoint\_interactingMeasure\_field\_rescale}
\end{theorem}
\begin{definition}[\texttt{connectedFourPoint}]\label{Pphi2-connectedFourPoint}
  \lean{Pphi2.connectedFourPoint}\leanok
  \uses{GaussianField-Configuration, GaussianField-instMeasurableSpaceConfiguration}
  \texttt{Pphi2.connectedFourPoint}
\end{definition}
\begin{theorem}[\texttt{fieldRescaleEquiv\_symm\_apply}]\label{Pphi2-fieldRescaleEquiv-symm-apply}
  \lean{Pphi2.fieldRescaleEquiv_symm_apply}\leanok
  \uses{Pphi2-fieldRescaleEquiv, GaussianField-Configuration, GaussianField-instMeasurableSpaceConfiguration}
  \texttt{Pphi2.fieldRescaleEquiv\_symm\_apply}
\end{theorem}
\begin{theorem}[\texttt{connectedFourPoint\_interactingMeasure\_zero\_smul}]\label{Pphi2-connectedFourPoint-interactingMeasure-zero-smul}
  \lean{Pphi2.connectedFourPoint_interactingMeasure_zero_smul}\leanok
  \uses{GaussianField-DyninMityaginSpace, Pphi2-interactingMeasure, GaussianField-measure, Pphi2-connectedFourPoint, Pphi2-interactingMeasure-zero-smul, GaussianField-Configuration, GaussianField-instMeasurableSpaceConfiguration, Pphi2-connectedFourPoint-gaussianMeasure-eq-zero, GaussianField-measure-isProbability}
  \texttt{Pphi2.connectedFourPoint\_interactingMeasure\_zero\_smul}
\end{theorem}
\section{\texttt{Pphi2.OSProofs.OS3\_RP\_Lattice}}
\begin{theorem}[\texttt{rp\_from\_transfer\_positivity}]\label{Pphi2-rp-from-transfer-positivity}
  \lean{Pphi2.rp_from_transfer_positivity}\leanok
  \uses{Pphi2-L2SpatialField, Pphi2-transferOperator-inner-nonneg, Pphi2-SpatialField, Pphi2-transferOperator-isSelfAdjoint, Pphi2-transferOperatorCLM}
  \texttt{Pphi2.rp\_from\_transfer\_positivity}
\end{theorem}
\begin{theorem}[\texttt{gaussian\_rp\_with\_boundary\_weight}]\label{Pphi2-gaussian-rp-with-boundary-weight}
  \lean{Pphi2.gaussian_rp_with_boundary_weight}\leanok
  \uses{Pphi2-BoundarySupported, GaussianField-gaussianDensity-integral-pos, Pphi2-fieldReflection2D, Pphi2-evalField2D, GaussianField-finLatticeDelta, GaussianField-evalMap, GaussianField-FinLatticeSites, GaussianField-latticeGaussianMeasure, GaussianField-Configuration, GaussianField-latticeGaussianMeasure-density-integral-, GaussianField-instMeasurableSpaceConfiguration, GaussianField-gaussianDensity, Pphi2-fieldFromSites, GaussianField-FinLatticeField, Pphi2-PositiveTimeSupported, Pphi2-gaussian-density-rp}
  \texttt{Pphi2.gaussian\_rp\_with\_boundary\_weight}
\end{theorem}
\begin{theorem}[\texttt{action\_decomposition}]\label{Pphi2-action-decomposition}
  \lean{Pphi2.action_decomposition}\leanok
  \uses{Pphi2-siteEquiv, Pphi2-fieldReflection2D, GaussianField-FinLatticeSites, Pphi2-fieldToSites, Pphi2-wickConstant, Pphi2-timeReflection2D-involution, Pphi2-timeReflection2D, Lattice-toSemilatticeInf, Pphi2-wickPolynomial, Pphi2-latticeInteraction}
  \texttt{Pphi2.action\_decomposition}
\end{theorem}
\begin{definition}[\texttt{Fsin}]\label{Pphi2-Fsin}
  \lean{Pphi2.Fsin}\leanok
  \uses{Pphi2-pairing2D}
  \texttt{Pphi2.Fsin}
\end{definition}
\begin{definition}[\texttt{Fcos}]\label{Pphi2-Fcos}
  \lean{Pphi2.Fcos}\leanok
  \uses{Pphi2-pairing2D}
  \texttt{Pphi2.Fcos}
\end{definition}
\begin{theorem}[\texttt{lattice\_rp}]\label{Pphi2-lattice-rp}
  \lean{Pphi2.lattice_rp}\leanok
  \uses{Pphi2-fieldReflection2D, Pphi2-evalField2D, GaussianField-FinLatticeSites, Pphi2-partitionFunction, Pphi2-interactionFunctional-measurable, Pphi2-interactingLatticeMeasure, GaussianField-latticeGaussianMeasure, Pphi2-gaussian-rp-with-boundary-weight, GaussianField-Configuration, GaussianField-instMeasurableSpaceConfiguration, Pphi2-boltzmannWeight, Pphi2-fieldFromSites, GaussianField-FinLatticeField, Pphi2-PositiveTimeSupported, Pphi2-interactionFunctional, Pphi2-boltzmannWeight-pos}
  \texttt{Pphi2.lattice\_rp}
\end{theorem}
\begin{definition}[\texttt{evalField2D}]\label{Pphi2-evalField2D}
  \lean{Pphi2.evalField2D}\leanok
  \uses{GaussianField-FinLatticeSites, GaussianField-Configuration, Pphi2-fieldFromSites, GaussianField-FinLatticeField}
  \texttt{Pphi2.evalField2D}
\end{definition}
\begin{theorem}[\texttt{lattice\_rp\_matrix\_reduction}]\label{Pphi2-lattice-rp-matrix-reduction}
  \lean{Pphi2.lattice_rp_matrix_reduction}\leanok
  \uses{Pphi2-fieldReflection2D, Pphi2-lattice-rp, Pphi2-evalField2D, GaussianField-FinLatticeSites, Pphi2-positiveTimeSupported-Fcos, Pphi2-positiveTimeSupported-Fsin, Pphi2-interactingLatticeMeasure, Pphi2-pairing2D, GaussianField-Configuration, GaussianField-instMeasurableSpaceConfiguration, Pphi2-timeReflection2D, Pphi2-Fsin, GaussianField-FinLatticeField, Pphi2-Fcos}
  \texttt{Pphi2.lattice\_rp\_matrix\_reduction}
\end{theorem}
\begin{theorem}[\texttt{massOperator\_cross\_block\_zero\_symm}]\label{Pphi2-massOperator-cross-block-zero-symm}
  \lean{Pphi2.massOperator_cross_block_zero_symm}\leanok
  \uses{GaussianField-massOperatorMatrix-isHermitian, GaussianField-massOperatorMatrix, GaussianField-FinLatticeSites, GaussianField-massOperatorEntry, Lattice-toSemilatticeInf, Pphi2-massOperator-cross-block-zero}
  \texttt{Pphi2.massOperator\_cross\_block\_zero\_symm}
\end{theorem}
\begin{definition}[\texttt{timeReflection2D}]\label{Pphi2-timeReflection2D}
  \lean{Pphi2.timeReflection2D}\leanok
  \texttt{Pphi2.timeReflection2D}
\end{definition}
\begin{theorem}[\texttt{congr\_simp}]\label{Pphi2-Fcos-congr-simp}
  \lean{Pphi2.Fcos.congr_simp}\leanok
  \uses{Pphi2-Fcos}
  \texttt{Pphi2.Fcos.congr\_simp}
\end{theorem}
\begin{theorem}[\texttt{pairing\_reflection\_reindex}]\label{Pphi2-pairing-reflection-reindex}
  \lean{Pphi2.pairing_reflection_reindex}\leanok
  \uses{Pphi2-fieldReflection2D, Pphi2-pairing2D, Pphi2-timeReflection2D-involution, Pphi2-timeReflection2D}
  \texttt{Pphi2.pairing\_reflection\_reindex}
\end{theorem}
\begin{theorem}[\texttt{positiveTimeSupported\_Fsin}]\label{Pphi2-positiveTimeSupported-Fsin}
  \lean{Pphi2.positiveTimeSupported_Fsin}\leanok
  \uses{Pphi2-pairing2D, Pphi2-Fsin, Pphi2-PositiveTimeSupported}
  \texttt{Pphi2.positiveTimeSupported\_Fsin}
\end{theorem}
\begin{theorem}[\texttt{congr\_simp}]\label{Pphi2-latticeInteraction-congr-simp}
  \lean{Pphi2.latticeInteraction.congr_simp}\leanok
  \uses{GaussianField-FinLatticeField, Pphi2-latticeInteraction}
  \texttt{Pphi2.latticeInteraction.congr\_simp}
\end{theorem}
\begin{theorem}[\texttt{congr\_simp}]\label{Pphi2-pairing2D-congr-simp}
  \lean{Pphi2.pairing2D.congr_simp}\leanok
  \uses{Pphi2-pairing2D}
  \texttt{Pphi2.pairing2D.congr\_simp}
\end{theorem}
\begin{theorem}[\texttt{gaussian\_density\_rp}]\label{Pphi2-gaussian-density-rp}
  \lean{Pphi2.gaussian_density_rp}\leanok
  \uses{GaussianField-massOperator-selfAdjoint, Pphi2-siteEquiv, Pphi2-BoundarySupported, Pphi2-fieldReflection2D, GaussianField-FinLatticeSites, GaussianField-massOperator, GaussianField-gaussianDensity, Pphi2-timeReflection2D, Pphi2-fieldFromSites, GaussianField-FinLatticeField, Pphi2-PositiveTimeSupported, Lattice-toSemilatticeInf}
  \texttt{Pphi2.gaussian\_density\_rp}
\end{theorem}
\begin{theorem}[\texttt{timeReflection2D\_involution}]\label{Pphi2-timeReflection2D-involution}
  \lean{Pphi2.timeReflection2D_involution}\leanok
  \uses{Pphi2-timeReflection2D}
  \texttt{Pphi2.timeReflection2D\_involution}
\end{theorem}
\begin{theorem}[\texttt{lattice\_rp\_matrix}]\label{Pphi2-lattice-rp-matrix}
  \lean{Pphi2.lattice_rp_matrix}\leanok
  \uses{Pphi2-siteEquiv, Pphi2-bounded-integrable-interacting, Pphi2-fieldReflection2D, Pphi2-lattice-rp-matrix-reduction, Pphi2-evalField2D, GaussianField-FinLatticeSites, Pphi2-interactingLatticeMeasure, Pphi2-pairing2D, GaussianField-configuration-eval-measurable, GaussianField-Configuration, GaussianField-instMeasurableSpaceConfiguration, Pphi2-timeReflection2D, Pphi2-fieldFromSites, Pphi2-Fsin, GaussianField-FinLatticeField, Pphi2-Fcos, Lattice-toSemilatticeInf}
  \texttt{Pphi2.lattice\_rp\_matrix}
\end{theorem}
\begin{theorem}[\texttt{positiveTimeSupported\_Fcos}]\label{Pphi2-positiveTimeSupported-Fcos}
  \lean{Pphi2.positiveTimeSupported_Fcos}\leanok
  \uses{Pphi2-pairing2D, Pphi2-PositiveTimeSupported, Pphi2-Fcos}
  \texttt{Pphi2.positiveTimeSupported\_Fcos}
\end{theorem}
\begin{definition}[\texttt{fieldToSites}]\label{Pphi2-fieldToSites}
  \lean{Pphi2.fieldToSites}\leanok
  \uses{Pphi2-siteEquiv, GaussianField-FinLatticeSites, GaussianField-FinLatticeField}
  \texttt{Pphi2.fieldToSites}
\end{definition}
\begin{theorem}[\texttt{congr\_simp}]\label{Pphi2-Fsin-congr-simp}
  \lean{Pphi2.Fsin.congr_simp}\leanok
  \uses{Pphi2-Fsin}
  \texttt{Pphi2.Fsin.congr\_simp}
\end{theorem}
\begin{definition}[\texttt{fieldReflection2D}]\label{Pphi2-fieldReflection2D}
  \lean{Pphi2.fieldReflection2D}\leanok
  \uses{Pphi2-timeReflection2D}
  \texttt{Pphi2.fieldReflection2D}
\end{definition}
\begin{definition}[\texttt{fieldFromSites}]\label{Pphi2-fieldFromSites}
  \lean{Pphi2.fieldFromSites}\leanok
  \uses{Pphi2-siteEquiv, GaussianField-FinLatticeSites, GaussianField-FinLatticeField}
  \texttt{Pphi2.fieldFromSites}
\end{definition}
\begin{definition}[\texttt{pairing2D}]\label{Pphi2-pairing2D}
  \lean{Pphi2.pairing2D}\leanok
  \texttt{Pphi2.pairing2D}
\end{definition}
\begin{definition}[\texttt{PositiveTimeSupported}]\label{Pphi2-PositiveTimeSupported}
  \lean{Pphi2.PositiveTimeSupported}\leanok
  \texttt{Pphi2.PositiveTimeSupported}
\end{definition}
\begin{definition}[\texttt{BoundarySupported}]\label{Pphi2-BoundarySupported}
  \lean{Pphi2.BoundarySupported}\leanok
  \texttt{Pphi2.BoundarySupported}
\end{definition}
\begin{definition}[\texttt{siteEquiv}]\label{Pphi2-siteEquiv}
  \lean{Pphi2.siteEquiv}\leanok
  \uses{GaussianField-FinLatticeSites}
  \texttt{Pphi2.siteEquiv}
\end{definition}
\begin{theorem}[\texttt{massOperator\_cross\_block\_zero}]\label{Pphi2-massOperator-cross-block-zero}
  \lean{Pphi2.massOperator_cross_block_zero}\leanok
  \uses{GaussianField-finLatticeDelta, GaussianField-FinLatticeSites, GaussianField-finiteLaplacianFun, GaussianField-finiteLaplacian, GaussianField-FinLatticeField, GaussianField-massOperatorEntry}
  \texttt{Pphi2.massOperator\_cross\_block\_zero}
\end{theorem}
\section{\texttt{Pphi2.InteractingMeasure.LatticeMeasure}}
\begin{definition}[\texttt{fieldAtSite}]\label{Pphi2-fieldAtSite}
  \lean{Pphi2.fieldAtSite}\leanok
  \uses{GaussianField-finLatticeDelta, GaussianField-FinLatticeSites, GaussianField-Configuration, GaussianField-FinLatticeField}
  \texttt{Pphi2.fieldAtSite}
\end{definition}
\begin{theorem}[\texttt{congr\_simp}]\label{Pphi2-partitionFunction-congr-simp}
  \lean{Pphi2.partitionFunction.congr_simp}\leanok
  \uses{Pphi2-partitionFunction}
  \texttt{Pphi2.partitionFunction.congr\_simp}
\end{theorem}
\begin{theorem}[\texttt{interactionFunctional\_bounded\_below}]\label{Pphi2-interactionFunctional-bounded-below}
  \lean{Pphi2.interactionFunctional_bounded_below}\leanok
  \uses{GaussianField-finLatticeDelta, GaussianField-FinLatticeSites, GaussianField-Configuration, Pphi2-wickConstant, GaussianField-FinLatticeField, Pphi2-interactionFunctional, Pphi2-wickPolynomial-bounded-below, Pphi2-wickPolynomial}
  \texttt{Pphi2.interactionFunctional\_bounded\_below}
\end{theorem}
\begin{definition}[\texttt{boltzmannWeight}]\label{Pphi2-boltzmannWeight}
  \lean{Pphi2.boltzmannWeight}\leanok
  \uses{GaussianField-FinLatticeSites, GaussianField-Configuration, GaussianField-FinLatticeField, Pphi2-interactionFunctional}
  \texttt{Pphi2.boltzmannWeight}
\end{definition}
\begin{theorem}[\texttt{partitionFunction\_pos}]\label{Pphi2-partitionFunction-pos}
  \lean{Pphi2.partitionFunction_pos}\leanok
  \uses{GaussianField-FinLatticeSites, Pphi2-partitionFunction, Pphi2-boltzmannWeight-integrable, GaussianField-latticeGaussianMeasure-isProbability, GaussianField-latticeGaussianMeasure, GaussianField-Configuration, GaussianField-instMeasurableSpaceConfiguration, Pphi2-boltzmannWeight, GaussianField-FinLatticeField, Pphi2-boltzmannWeight-pos}
  $Z_a > 0$ since $e^{-V_a} > 0$ everywhere.
\end{theorem}
\begin{theorem}[\texttt{interactionFunctional\_single\_site}]\label{Pphi2-interactionFunctional-single-site}
  \lean{Pphi2.interactionFunctional_single_site}\leanok
  \uses{GaussianField-finLatticeDelta, GaussianField-FinLatticeSites, GaussianField-Configuration, Pphi2-wickConstant, GaussianField-FinLatticeField, Pphi2-interactionFunctional, Pphi2-wickPolynomial}
  \texttt{Pphi2.interactionFunctional\_single\_site}
\end{theorem}
\begin{theorem}[\texttt{boltzmannWeight\_pos}]\label{Pphi2-boltzmannWeight-pos}
  \lean{Pphi2.boltzmannWeight_pos}\leanok
  \uses{GaussianField-FinLatticeSites, GaussianField-Configuration, Pphi2-boltzmannWeight, GaussianField-FinLatticeField, Pphi2-interactionFunctional}
  \texttt{Pphi2.boltzmannWeight\_pos}
\end{theorem}
\begin{definition}[\texttt{latticeSchwinger2}]\label{Pphi2-latticeSchwinger2}
  \lean{Pphi2.latticeSchwinger2}\leanok
  \uses{GaussianField-finLatticeDelta, GaussianField-FinLatticeSites, Pphi2-interactingLatticeMeasure, GaussianField-Configuration, GaussianField-instMeasurableSpaceConfiguration, GaussianField-FinLatticeField}
  \texttt{Pphi2.latticeSchwinger2}
\end{definition}
\begin{definition}[\texttt{interactingLatticeMeasure}]\label{Pphi2-interactingLatticeMeasure}
  \lean{Pphi2.interactingLatticeMeasure}\leanok
  \uses{GaussianField-FinLatticeSites, Pphi2-partitionFunction, GaussianField-latticeGaussianMeasure, GaussianField-Configuration, GaussianField-instMeasurableSpaceConfiguration, Pphi2-boltzmannWeight, GaussianField-FinLatticeField}
  \texttt{Pphi2.interactingLatticeMeasure}
\end{definition}
\begin{theorem}[\texttt{boltzmannWeight\_integrable}]\label{Pphi2-boltzmannWeight-integrable}
  \lean{Pphi2.boltzmannWeight_integrable}\leanok
  \uses{Pphi2-interactionFunctional-bounded-below, GaussianField-FinLatticeSites, Pphi2-interactionFunctional-measurable, GaussianField-latticeGaussianMeasure-isProbability, GaussianField-latticeGaussianMeasure, GaussianField-Configuration, GaussianField-instMeasurableSpaceConfiguration, Pphi2-boltzmannWeight, GaussianField-FinLatticeField, Pphi2-interactionFunctional}
  \texttt{Pphi2.boltzmannWeight\_integrable}
\end{theorem}
\begin{theorem}[\texttt{interactingLatticeMeasure\_isProbability}]\label{Pphi2-interactingLatticeMeasure-isProbability}
  \lean{Pphi2.interactingLatticeMeasure_isProbability}\leanok
  \uses{GaussianField-FinLatticeSites, Pphi2-partitionFunction, Pphi2-boltzmannWeight-integrable, Pphi2-interactingLatticeMeasure, GaussianField-latticeGaussianMeasure, GaussianField-Configuration, GaussianField-instMeasurableSpaceConfiguration, Pphi2-boltzmannWeight, GaussianField-FinLatticeField, Pphi2-partitionFunction-pos, Pphi2-boltzmannWeight-pos}
  $\mu_a$ is a probability measure.
\end{theorem}
\begin{theorem}[\texttt{congr\_simp}]\label{Pphi2-interactionFunctional-congr-simp}
  \lean{Pphi2.interactionFunctional.congr_simp}\leanok
  \uses{GaussianField-FinLatticeSites, GaussianField-Configuration, GaussianField-FinLatticeField, Pphi2-interactionFunctional}
  \texttt{Pphi2.interactionFunctional.congr\_simp}
\end{theorem}
\begin{theorem}[\texttt{interactionFunctional\_measurable}]\label{Pphi2-interactionFunctional-measurable}
  \lean{Pphi2.interactionFunctional_measurable}\leanok
  \uses{GaussianField-finLatticeDelta, GaussianField-FinLatticeSites, Pphi2-wickMonomial, GaussianField-configuration-eval-measurable, GaussianField-Configuration, Pphi2-wickConstant, GaussianField-instMeasurableSpaceConfiguration, GaussianField-FinLatticeField, Pphi2-interactionFunctional, Pphi2-wickPolynomial}
  $V_a$ is measurable (each $\omega \mapsto \omega(\delta_x)$ is measurable by the cylindrical $\sigma$-algebra).
\end{theorem}
\begin{definition}[\texttt{partitionFunction}]\label{Pphi2-partitionFunction}
  \lean{Pphi2.partitionFunction}\leanok
  \uses{GaussianField-FinLatticeSites, GaussianField-latticeGaussianMeasure, GaussianField-Configuration, GaussianField-instMeasurableSpaceConfiguration, Pphi2-boltzmannWeight, GaussianField-FinLatticeField}
  \texttt{Pphi2.partitionFunction}
\end{definition}
\begin{theorem}[\texttt{congr\_simp}]\label{Pphi2-boltzmannWeight-congr-simp}
  \lean{Pphi2.boltzmannWeight.congr_simp}\leanok
  \uses{GaussianField-FinLatticeSites, GaussianField-Configuration, Pphi2-boltzmannWeight, GaussianField-FinLatticeField}
  \texttt{Pphi2.boltzmannWeight.congr\_simp}
\end{theorem}
\begin{definition}[\texttt{interactionFunctional}]\label{Pphi2-interactionFunctional}
  \lean{Pphi2.interactionFunctional}\leanok
  \uses{GaussianField-finLatticeDelta, GaussianField-FinLatticeSites, GaussianField-Configuration, Pphi2-wickConstant, GaussianField-FinLatticeField, Pphi2-wickPolynomial}
  \texttt{Pphi2.interactionFunctional}
\end{definition}
\section{\texttt{Pphi2.ContinuumLimit.Convergence}}
\begin{theorem}[\texttt{continuumLimit}]\label{Pphi2-continuumLimit}
  \lean{Pphi2.continuumLimit}\leanok
  \uses{Pphi2-EuclideanPlaneBackground-SpaceTime, Pphi2-planeBackground, Pphi2-continuumMeasures-tight, Pphi2-continuumMeasure-isProbability, Pphi2-EuclideanPlaneBackground-dim, Pphi2-prokhorov-configuration-sequential, Pphi2-continuumMeasure, GaussianField-Configuration, GaussianField-instMeasurableSpaceConfiguration, Pphi2-ContinuumTestFunction}
  For any sequence of lattice spacings $a_n \to 0$, there exists a subsequence $a_{n_k}$ and a probability measure $\mu$ on $\mathcal{S}'(\mathbb{R}^d)$ such that $\nu_{a_{n_k}} \rightharpoonup \mu$ weakly.
\end{theorem}
\begin{theorem}[\texttt{prokhorov\_sequential}]\label{Pphi2-prokhorov-sequential}
  \lean{Pphi2.prokhorov_sequential}\leanok
  On a Polish space $X$, if a sequence of probability measures $\{\mu_n\}$ is tight, then it has a weakly convergent subsequence. Fully proved from Mathlib.
\end{theorem}
\begin{theorem}[\texttt{continuumLimit\_nonGaussian} (axiom)]\label{Pphi2-continuumLimit-nonGaussian}
  \lean{Pphi2.continuumLimit_nonGaussian}\leanok
  \uses{Pphi2-EuclideanPlaneBackground-SpaceTime, Pphi2-planeBackground, Pphi2-EuclideanPlaneBackground-dim, GaussianField-Configuration, GaussianField-instMeasurableSpaceConfiguration, Pphi2-ContinuumTestFunction}
  The continuum limit is non-Gaussian for nontrivial $P$: there exists $f$ with $S_4(f,f,f,f) - 3 S_2(f,f)^2 \ne 0$ (nonzero connected four-point function).
\end{theorem}
\begin{theorem}[\texttt{prokhorov\_configuration\_sequential}]\label{Pphi2-prokhorov-configuration-sequential}
  \lean{Pphi2.prokhorov_configuration_sequential}\leanok
  \uses{Pphi2-EuclideanPlaneBackground-SpaceTime, GaussianField-schwartz-dyninMityaginSpace, Pphi2-planeBackground, Pphi2-EuclideanPlaneBackground-dim, GaussianField-prokhorov-configuration, GaussianField-Configuration, GaussianField-instMeasurableSpaceConfiguration, Pphi2-ContinuumTestFunction}
  Sequential Prokhorov extraction on configuration space, derived from \texttt{prokhorov\_sequential} using the two topological axioms.
\end{theorem}
\begin{theorem}[\texttt{pphi2\_limit\_exists}]\label{Pphi2-pphi2-limit-exists}
  \lean{Pphi2.pphi2_limit_exists}\leanok
  \uses{Pphi2-EuclideanPlaneBackground-SpaceTime, Pphi2-FieldConfig2, Pphi2-EuclideanPlaneBackground-RealTestFunction, Pphi2-planeBackground, Pphi2-EuclideanPlaneBackground-dim, Pphi2-IsPphi2Limit, Pphi2-SpaceTime2, GaussianField-configuration-measurable-of-eval-measurable, Pphi2-PositiveTimeTestFunction2, GaussianField-configuration-eval-measurable, GaussianField-Configuration, GaussianField-instMeasurableSpaceConfiguration, Pphi2-TestFunction2, Pphi2-ContinuumTestFunction, Pphi2-compTimeReflection2, Pphi2-schwartzTranslate}
  There exists a probability measure $\mu$ on $\mathcal{S}'(\mathbb{R}^2)$ satisfying \texttt{IsPphi2Limit}. Proved by constructing the Dirac measure at $0$ as a trivial witness.
\end{theorem}
\section{\texttt{Pphi2.ContinuumLimit.Embedding}}
\begin{definition}[\texttt{IsPphi2Limit}]\label{Pphi2-IsPphi2Limit}
  \lean{Pphi2.IsPphi2Limit}\leanok
  \uses{Pphi2-EuclideanPlaneBackground-SpaceTime, Pphi2-FieldConfig2, Pphi2-EuclideanPlaneBackground-RealTestFunction, Pphi2-planeBackground, Pphi2-EuclideanPlaneBackground-dim, Pphi2-SpaceTime2, Pphi2-PositiveTimeTestFunction2, GaussianField-instMeasurableSpaceConfiguration, Pphi2-TestFunction2, Pphi2-ContinuumTestFunction, Pphi2-compTimeReflection2, Pphi2-schwartzTranslate}
  ``\texttt{lean def IsPphi2Limit \{d : ℕ\} (μ : Measure (Configuration (ContinuumTestFunction d))) (\_P : InteractionPolynomial) (\_mass : ℝ) : Prop }`` Marker predicate: $\mu$ is a $P(\varphi)_2$ continuum limit if there exist lattice spacings $a_k \to 0$ and measures $\nu_k$ with convergent Schwinger functions, $\mathbb{Z}_2$ symmetry, characteristic functional convergence, and lattice translation invariance.
\end{definition}
\begin{definition}[\texttt{latticeEmbed}]\label{Pphi2-latticeEmbed}
  \lean{Pphi2.latticeEmbed}\leanok
  \uses{Pphi2-EuclideanPlaneBackground-SpaceTime, Pphi2-planeBackground, Pphi2-latticeEmbedEval, Pphi2-EuclideanPlaneBackground-dim, GaussianField-Configuration, Pphi2-ContinuumTestFunction, GaussianField-FinLatticeField}
  \texttt{Pphi2.latticeEmbed}
\end{definition}
\begin{definition}[\texttt{latticeTestField}]\label{Pphi2-latticeTestField}
  \lean{Pphi2.latticeTestField}\leanok
  \uses{GaussianField-FinLatticeSites, Pphi2-evalAtSite, Pphi2-ContinuumTestFunction, GaussianField-FinLatticeField}
  \texttt{Pphi2.latticeTestField}
\end{definition}
\begin{theorem}[\texttt{latticeEmbedEval\_linear\_f}]\label{Pphi2-latticeEmbedEval-linear-f}
  \lean{Pphi2.latticeEmbedEval_linear_f}\leanok
  \uses{Pphi2-EuclideanPlaneBackground-SpaceTime, Pphi2-planeBackground, GaussianField-FinLatticeSites, Pphi2-evalAtSite, Pphi2-latticeEmbedEval, Pphi2-EuclideanPlaneBackground-dim, Pphi2-physicalPosition, Pphi2-ContinuumTestFunction, GaussianField-FinLatticeField}
  The embedding is linear in $f$: $\iota_a(\varphi)(c_1 f_1 + c_2 f_2) = c_1 \iota_a(\varphi)(f_1) + c_2 \iota_a(\varphi)(f_2)$. Proved by linearity of Schwartz evaluation.
\end{theorem}
\begin{definition}[\texttt{continuumMeasure}]\label{Pphi2-continuumMeasure}
  \lean{Pphi2.continuumMeasure}\leanok
  \uses{Pphi2-EuclideanPlaneBackground-SpaceTime, Pphi2-planeBackground, GaussianField-FinLatticeSites, Pphi2-EuclideanPlaneBackground-dim, Pphi2-interactingLatticeMeasure, GaussianField-Configuration, GaussianField-instMeasurableSpaceConfiguration, Pphi2-latticeEmbedLift, Pphi2-ContinuumTestFunction, GaussianField-FinLatticeField}
  \texttt{Pphi2.continuumMeasure}
\end{definition}
\begin{theorem}[\texttt{signedVal\_neg}]\label{Pphi2-signedVal-neg}
  \lean{Pphi2.signedVal_neg}\leanok
  \uses{Pphi2-signedVal}
  \texttt{Pphi2.signedVal\_neg}
\end{theorem}
\begin{definition}[\texttt{signedVal}]\label{Pphi2-signedVal}
  \lean{Pphi2.signedVal}\leanok
  \texttt{Pphi2.signedVal}
\end{definition}
\begin{theorem}[\texttt{continuumMeasure\_isProbability}]\label{Pphi2-continuumMeasure-isProbability}
  \lean{Pphi2.continuumMeasure_isProbability}\leanok
  \uses{Pphi2-EuclideanPlaneBackground-SpaceTime, Pphi2-planeBackground, GaussianField-FinLatticeSites, Pphi2-latticeEmbedLift-measurable, Pphi2-EuclideanPlaneBackground-dim, Pphi2-interactingLatticeMeasure-isProbability, Pphi2-interactingLatticeMeasure, Pphi2-continuumMeasure, GaussianField-Configuration, GaussianField-instMeasurableSpaceConfiguration, Pphi2-latticeEmbedLift, Pphi2-ContinuumTestFunction, GaussianField-FinLatticeField}
  The pushforward is a probability measure. Proved from \texttt{interactingLatticeMeasure\_isProbability}.
\end{theorem}
\begin{theorem}[\texttt{latticeEmbedLift\_measurable}]\label{Pphi2-latticeEmbedLift-measurable}
  \lean{Pphi2.latticeEmbedLift_measurable}\leanok
  \uses{Pphi2-EuclideanPlaneBackground-SpaceTime, Pphi2-planeBackground, GaussianField-FinLatticeSites, Pphi2-evalAtSite, Pphi2-EuclideanPlaneBackground-dim, GaussianField-configuration-measurable-of-eval-measurable, GaussianField-configuration-eval-measurable, GaussianField-Configuration, GaussianField-instMeasurableSpaceConfiguration, Pphi2-latticeEmbedLift, Pphi2-ContinuumTestFunction, GaussianField-FinLatticeField}
  The lifted embedding is measurable. Proved.
\end{theorem}
\begin{definition}[\texttt{latticeEmbedEval}]\label{Pphi2-latticeEmbedEval}
  \lean{Pphi2.latticeEmbedEval}\leanok
  \uses{GaussianField-FinLatticeSites, Pphi2-evalAtSite, Pphi2-ContinuumTestFunction, GaussianField-FinLatticeField}
  \texttt{Pphi2.latticeEmbedEval}
\end{definition}
\begin{theorem}[\texttt{latticeTestField\_expand}]\label{Pphi2-latticeTestField-expand}
  \lean{Pphi2.latticeTestField_expand}\leanok
  \uses{Pphi2-latticeTestField, GaussianField-FinLatticeSites, Pphi2-evalAtSite, Pphi2-ContinuumTestFunction, GaussianField-FinLatticeField}
  \texttt{Pphi2.latticeTestField\_expand}
\end{theorem}
\begin{definition}[\texttt{evalAtSite}]\label{Pphi2-evalAtSite}
  \lean{Pphi2.evalAtSite}\leanok
  \uses{Pphi2-EuclideanPlaneBackground-SpaceTime, Pphi2-planeBackground, GaussianField-FinLatticeSites, Pphi2-EuclideanPlaneBackground-dim, Pphi2-physicalPosition, Pphi2-ContinuumTestFunction}
  \texttt{Pphi2.evalAtSite}
\end{definition}
\begin{definition}[\texttt{physicalPosition}]\label{Pphi2-physicalPosition}
  \lean{Pphi2.physicalPosition}\leanok
  \uses{GaussianField-FinLatticeSites, Pphi2-signedVal}
  \texttt{Pphi2.physicalPosition}
\end{definition}
\begin{theorem}[\texttt{latticeEmbedEval\_linear\_phi}]\label{Pphi2-latticeEmbedEval-linear-phi}
  \lean{Pphi2.latticeEmbedEval_linear_phi}\leanok
  \uses{GaussianField-FinLatticeSites, Pphi2-evalAtSite, Pphi2-latticeEmbedEval, Pphi2-ContinuumTestFunction, GaussianField-FinLatticeField}
  The embedding is bilinear: $\iota_a(c_1\varphi_1 + c_2\varphi_2)(f) = c_1 \iota_a(\varphi_1)(f) + c_2 \iota_a(\varphi_2)(f)$. Proved by expanding the sum and factoring.
\end{theorem}
\begin{theorem}[\texttt{latticeEmbed\_measurable}]\label{Pphi2-latticeEmbed-measurable}
  \lean{Pphi2.latticeEmbed_measurable}\leanok
  \uses{Pphi2-EuclideanPlaneBackground-SpaceTime, Pphi2-latticeEmbed, Pphi2-planeBackground, GaussianField-FinLatticeSites, Pphi2-evalAtSite, Pphi2-EuclideanPlaneBackground-dim, GaussianField-configuration-measurable-of-eval-measurable, GaussianField-Configuration, GaussianField-instMeasurableSpaceConfiguration, Pphi2-ContinuumTestFunction, GaussianField-FinLatticeField}
  The embedding is measurable (each evaluation $\varphi \mapsto (\iota_a \varphi)(f)$ is continuous on the finite-dimensional space). Proved.
\end{theorem}
\begin{theorem}[\texttt{congr\_simp}]\label{Pphi2-latticeEmbed-congr-simp}
  \lean{Pphi2.latticeEmbed.congr_simp}\leanok
  \uses{Pphi2-EuclideanPlaneBackground-SpaceTime, Pphi2-latticeEmbed, Pphi2-planeBackground, Pphi2-EuclideanPlaneBackground-dim, GaussianField-Configuration, Pphi2-ContinuumTestFunction, GaussianField-FinLatticeField}
  \texttt{Pphi2.latticeEmbed.congr\_simp}
\end{theorem}
\begin{theorem}[\texttt{latticeEmbedLift\_eval\_eq}]\label{Pphi2-latticeEmbedLift-eval-eq}
  \lean{Pphi2.latticeEmbedLift_eval_eq}\leanok
  \uses{Pphi2-EuclideanPlaneBackground-SpaceTime, Pphi2-latticeTestField, Pphi2-planeBackground, GaussianField-FinLatticeSites, Pphi2-evalAtSite, Pphi2-EuclideanPlaneBackground-dim, GaussianField-Configuration, Pphi2-latticeEmbedLift, Pphi2-ContinuumTestFunction, Pphi2-latticeTestField-expand, GaussianField-FinLatticeField}
  \texttt{Pphi2.latticeEmbedLift\_eval\_eq}
\end{theorem}
\begin{theorem}[\texttt{congr\_simp}]\label{Pphi2-latticeEmbedEval-congr-simp}
  \lean{Pphi2.latticeEmbedEval.congr_simp}\leanok
  \uses{Pphi2-latticeEmbedEval, Pphi2-ContinuumTestFunction, GaussianField-FinLatticeField}
  \texttt{Pphi2.latticeEmbedEval.congr\_simp}
\end{theorem}
\begin{theorem}[\texttt{latticeEmbed\_eval}]\label{Pphi2-latticeEmbed-eval}
  \lean{Pphi2.latticeEmbed_eval}\leanok
  \uses{Pphi2-EuclideanPlaneBackground-SpaceTime, Pphi2-latticeEmbed, Pphi2-planeBackground, Pphi2-latticeEmbedEval, Pphi2-EuclideanPlaneBackground-dim, GaussianField-Configuration, Pphi2-ContinuumTestFunction, GaussianField-FinLatticeField}
  Agreement: $(\texttt{latticeEmbed}\ \varphi)(f) = \texttt{latticeEmbedEval}\ \varphi\ f$. Proved by \texttt{rfl}.
\end{theorem}
\begin{definition}[\texttt{latticeEmbedLift}]\label{Pphi2-latticeEmbedLift}
  \lean{Pphi2.latticeEmbedLift}\leanok
  \uses{Pphi2-EuclideanPlaneBackground-SpaceTime, Pphi2-latticeEmbed, Pphi2-planeBackground, GaussianField-FinLatticeSites, Pphi2-EuclideanPlaneBackground-dim, GaussianField-Configuration, Pphi2-ContinuumTestFunction, GaussianField-FinLatticeField}
  \texttt{Pphi2.latticeEmbedLift}
\end{definition}
\section{\texttt{Pphi2.NelsonEstimate.SmoothLowerBound}}
\begin{theorem}[\texttt{smoothBoundConstant\_pos}]\label{Pphi2-smoothBoundConstant-pos}
  \lean{Pphi2.smoothBoundConstant_pos}\leanok
  \uses{Pphi2-smoothVarianceConstant, Pphi2-smoothBoundConstant, Pphi2-smoothVarianceConstant-pos}
  \texttt{Pphi2.smoothBoundConstant\_pos}
\end{theorem}
\begin{theorem}[\texttt{smooth\_interaction\_lower\_bound}]\label{Pphi2-smooth-interaction-lower-bound}
  \lean{Pphi2.smooth_interaction_lower_bound}\leanok
  \uses{Pphi2-wickMonomial-four-lower-bound, GaussianField-FinLatticeSites, Pphi2-wickMonomial}
  Summed over all sites: $V_S = a^d \sum_x :(\varphi_S(x))^4:_{c_S} \ge -6 a^d |\Lambda| c_S^2$.
\end{theorem}
\begin{theorem}[\texttt{smooth\_interaction\_lower\_bound\_log}]\label{Pphi2-smooth-interaction-lower-bound-log}
  \lean{Pphi2.smooth_interaction_lower_bound_log}\leanok
  \uses{Pphi2-smoothWickConstant, Pphi2-smoothCovEigenvalue-pos, GaussianField-FinLatticeSites, Pphi2-smoothVariance-le-log, Pphi2-wickMonomial, Pphi2-smooth-interaction-lower-bound-volume, Pphi2-smoothCovEigenvalue}
  Combined with \texttt{smoothVariance\_le\_log}: $V_S \ge -C(1 + |\log T|)^2$ where $C$ depends on $L$ and $m$ but NOT on $N$.  Depends on: \texttt{wickMonomial\_four\_lower\_bound}, \texttt{smoothVariance\_le\_log}, \texttt{smoothCovEigenvalue\_pos}.
\end{theorem}
\begin{theorem}[\texttt{smooth\_interaction\_lower\_bound\_log\_uniform}]\label{Pphi2-smooth-interaction-lower-bound-log-uniform}
  \lean{Pphi2.smooth_interaction_lower_bound_log_uniform}\leanok
  \uses{Pphi2-smoothWickConstant, Pphi2-smoothCovEigenvalue-pos, Pphi2-smoothVariance-le-log-uniform, GaussianField-FinLatticeSites, Pphi2-smoothVarianceConstant, Pphi2-wickMonomial, Pphi2-smoothBoundConstant, Pphi2-smooth-interaction-lower-bound-volume, Pphi2-smoothVarianceConstant-pos, Pphi2-smoothCovEigenvalue}
  \texttt{Pphi2.smooth\_interaction\_lower\_bound\_log\_uniform}
\end{theorem}
\begin{definition}[\texttt{smoothBoundConstant}]\label{Pphi2-smoothBoundConstant}
  \lean{Pphi2.smoothBoundConstant}\leanok
  \uses{Pphi2-smoothVarianceConstant}
  \texttt{Pphi2.smoothBoundConstant}
\end{definition}
\begin{theorem}[\texttt{site\_interaction\_lower\_bound}]\label{Pphi2-site-interaction-lower-bound}
  \lean{Pphi2.site_interaction_lower_bound}\leanok
  \uses{Pphi2-wickMonomial-four-lower-bound, Pphi2-wickMonomial}
  $a^d \cdot :(\varphi(x))^4:_c \ge -6 a^d c^2$ at a single lattice site.
\end{theorem}
\begin{theorem}[\texttt{smooth\_interaction\_lower\_bound\_volume}]\label{Pphi2-smooth-interaction-lower-bound-volume}
  \lean{Pphi2.smooth_interaction_lower_bound_volume}\leanok
  \uses{GaussianField-FinLatticeSites, Pphi2-wickMonomial, Pphi2-smooth-interaction-lower-bound, Lattice-toSemilatticeInf}
  Volume-independent version: with $a = L/N$, we have $a^d |\Lambda| = L^d$, so $V_S \ge -6L^d c_S^2$.
\end{theorem}
\section{\texttt{Pphi2.AsymTorus.AsymWickMean}}
\begin{theorem}[\texttt{congr\_simp}]\label{Pphi2-wickConstantAsym-congr-simp}
  \lean{Pphi2.wickConstantAsym.congr_simp}\leanok
  \uses{Pphi2-wickConstantAsym}
  \texttt{Pphi2.wickConstantAsym.congr\_simp}
\end{theorem}
\begin{theorem}[\texttt{wickConstantAsym\_eq\_variance}]\label{Pphi2-wickConstantAsym-eq-variance}
  \lean{Pphi2.wickConstantAsym_eq_variance}\leanok
  \uses{GaussianField-spectralLatticeCovarianceAsym, GaussianField-AsymLatticeSites, Pphi2-latticeCovarianceAsymGJ-inner-eq-inv-a-sq-spectral, Pphi2-wickConstantAsym, GaussianField-massEigenvaluesAsym, GaussianField-AsymLatticeField, GaussianField-instFintypeAsymLatticeSitesOfNeZeroNat, GaussianField-latticeCovarianceAsymGJ, GaussianField-asymLatticeDelta, Pphi2-spectralVarianceAsym-inner-eq-massEigenvalue-average, GaussianField-ell2-}
  \texttt{Pphi2.wickConstantAsym\_eq\_variance}
\end{theorem}
\begin{theorem}[\texttt{wickMonomialAsym\_latticeGaussian}]\label{Pphi2-wickMonomialAsym-latticeGaussian}
  \lean{Pphi2.wickMonomialAsym_latticeGaussian}\leanok
  \uses{Pphi2-wickConstantAsym-pos, GaussianField-AsymLatticeSites, GaussianField-measure, GaussianField-ell2-separableSpace, Pphi2-wickConstantAsym, GaussianField-asymLatticeField-dyninMityaginSpace, Pphi2-gaussian-hermite-zero-mean, Pphi2-wickConstantAsym-eq-variance, Pphi2-wickMonomial, GaussianField-pairing-is-gaussian, GaussianField-configuration-eval-measurable, GaussianField-Configuration, GaussianField-AsymLatticeField, Pphi2-latticeGaussianMeasureAsym, GaussianField-instMeasurableSpaceConfiguration, GaussianField-latticeCovarianceAsymGJ, GaussianField-asymLatticeDelta, GaussianField-ell2-}
  \texttt{Pphi2.wickMonomialAsym\_latticeGaussian}
\end{theorem}
\begin{theorem}[\texttt{density\_transfer\_bound\_iso}]\label{Pphi2-density-transfer-bound-iso}
  \lean{Pphi2.density_transfer_bound_iso}\leanok
  \uses{GaussianField-AsymLatticeSites, Pphi2-partitionFunctionAsym-pos, Pphi2-boltzmannWeightAsym, Pphi2-boltzmannWeightAsym-pos, Pphi2-interactionFunctionalAsym-bounded-below, Pphi2-partitionFunctionAsym, GaussianField-Configuration, GaussianField-AsymLatticeField, Pphi2-latticeGaussianMeasureAsym, Pphi2-interactingLatticeMeasureAsym, Pphi2-interactionFunctionalAsym, GaussianField-instMeasurableSpaceConfiguration, Pphi2-interactionFunctionalAsym-measurable, Pphi2-latticeGaussianMeasureAsym-isProbability, Lattice-toSemilatticeInf}
  \texttt{Pphi2.density\_transfer\_bound\_iso}
\end{theorem}
\begin{theorem}[\texttt{partitionFunctionAsym\_ge\_one}]\label{Pphi2-partitionFunctionAsym-ge-one}
  \lean{Pphi2.partitionFunctionAsym_ge_one}\leanok
  \uses{Pphi2-boltzmannWeightAsym-integrable, GaussianField-AsymLatticeSites, Pphi2-partitionFunctionAsym, GaussianField-Configuration, GaussianField-AsymLatticeField, Pphi2-latticeGaussianMeasureAsym, Pphi2-interactionFunctionalAsym, GaussianField-instMeasurableSpaceConfiguration, Pphi2-latticeGaussianMeasureAsym-isProbability, Pphi2-interactionFunctionalAsym-mean-nonpos}
  \texttt{Pphi2.partitionFunctionAsym\_ge\_one}
\end{theorem}
\begin{theorem}[\texttt{interactionFunctionalAsym\_mean\_nonpos}]\label{Pphi2-interactionFunctionalAsym-mean-nonpos}
  \lean{Pphi2.interactionFunctionalAsym_mean_nonpos}\leanok
  \uses{GaussianField-AsymLatticeSites, Pphi2-wickConstantAsym, Pphi2-wickPolynomialAsym-integral-eq-coeff-zero, GaussianField-Configuration, GaussianField-AsymLatticeField, Pphi2-latticeGaussianMeasureAsym, Pphi2-interactionFunctionalAsym, GaussianField-instMeasurableSpaceConfiguration, GaussianField-instFintypeAsymLatticeSitesOfNeZeroNat, GaussianField-asymLatticeDelta, Pphi2-wickPolynomial}
  \texttt{Pphi2.interactionFunctionalAsym\_mean\_nonpos}
\end{theorem}
\begin{theorem}[\texttt{wickPolynomialAsym\_integral\_eq\_coeff\_zero}]\label{Pphi2-wickPolynomialAsym-integral-eq-coeff-zero}
  \lean{Pphi2.wickPolynomialAsym_integral_eq_coeff_zero}\leanok
  \uses{GaussianField-AsymLatticeSites, Pphi2-wickConstantAsym, Pphi2-wickMonomialAsym-latticeGaussian, Pphi2-wickMonomial, GaussianField-Configuration, GaussianField-AsymLatticeField, Pphi2-latticeGaussianMeasureAsym, GaussianField-instMeasurableSpaceConfiguration, GaussianField-asymLatticeDelta, Pphi2-latticeGaussianMeasureAsym-isProbability, Pphi2-wickPolynomial}
  \texttt{Pphi2.wickPolynomialAsym\_integral\_eq\_coeff\_zero}
\end{theorem}
\section{\texttt{Pphi2.AsymTorus.AsymB5bSingleSlice}}
\begin{theorem}[\texttt{congr\_simp}]\label{Pphi2-singleSliceLatticeField-congr-simp}
  \lean{Pphi2.singleSliceLatticeField.congr_simp}\leanok
  \uses{GaussianField-AsymLatticeSites, Pphi2-SpatialField, Pphi2-singleSliceLatticeField}
  \texttt{Pphi2.singleSliceLatticeField.congr\_simp}
\end{theorem}
\begin{definition}[\texttt{groundSliceVarianceSum}]\label{Pphi2-groundSliceVarianceSum}
  \lean{Pphi2.groundSliceVarianceSum}\leanok
  \uses{Pphi2-groundSliceVariance, Pphi2-SpatialField}
  \texttt{Pphi2.groundSliceVarianceSum}
\end{definition}
\begin{theorem}[\texttt{interacting\_second\_moment\_bound\_to\_lattice\_free\_covariance}]\label{Pphi2-interacting-second-moment-bound-to-lattice-free-covariance}
  \lean{Pphi2.interacting_second_moment_bound_to_lattice_free_covariance}\leanok
  \uses{Pphi2-freeSingleSliceCovarianceSum, GaussianField-AsymLatticeSites, Pphi2-groundVarianceFreeCovarianceConstant, Pphi2-SpatialField, Pphi2-asymSliceEquiv, Pphi2-asymTransferSystem, Pphi2-asymSliceFamilyLinear, GaussianField-Configuration, Pphi2-interacting-second-moment-bound-to-freeCovariance-sum, GaussianField-AsymLatticeField, Pphi2-latticeGaussianMeasureAsym, Pphi2-interactingLatticeMeasureAsym, GaussianField-instMeasurableSpaceConfiguration, Pphi2-groundSliceVarianceSum}
  \texttt{Pphi2.interacting\_second\_moment\_bound\_to\_lattice\_free\_covariance}
\end{theorem}
\begin{theorem}[\texttt{groundVariance\_le\_freeCovariance\_with\_constant}]\label{Pphi2-groundVariance-le-freeCovariance-with-constant}
  \lean{Pphi2.groundVariance_le_freeCovariance_with_constant}\leanok
  \uses{Pphi2-groundSliceVariance, Pphi2-groundVarianceFreeCovarianceConstant, Pphi2-SpatialField, Pphi2-freeSingleSliceCovariance, Pphi2-groundVariance-le-freeCovariance}
  \texttt{Pphi2.groundVariance\_le\_freeCovariance\_with\_constant}
\end{theorem}
\begin{theorem}[\texttt{interacting\_second\_moment\_bound\_to\_freeCovariance\_sum}]\label{Pphi2-interacting-second-moment-bound-to-freeCovariance-sum}
  \lean{Pphi2.interacting_second_moment_bound_to_freeCovariance_sum}\leanok
  \uses{Pphi2-freeSingleSliceCovarianceSum, Pphi2-asym-interacting-second-moment-eq-pathMeasure-slicePairing, GaussianField-AsymLatticeSites, Pphi2-groundVarianceFreeCovarianceConstant, Pphi2-SpatialField, Pphi2-asymSliceEquiv, Pphi2-asymTransferSystem, Pphi2-asymSliceFamilyLinear, GaussianField-Configuration, GaussianField-AsymLatticeField, Pphi2-piece3-pathMeasure-bound-to-freeCovariance-sum, Pphi2-slicePairing, Pphi2-slicePairing-congr-simp, Pphi2-interactingLatticeMeasureAsym, GaussianField-instMeasurableSpaceConfiguration, Pphi2-asymSliceFamilyLinear-eq-slicePairing, Pphi2-groundSliceVarianceSum}
  \texttt{Pphi2.interacting\_second\_moment\_bound\_to\_freeCovariance\_sum}
\end{theorem}
\begin{theorem}[\texttt{freeSingleSliceCovariance\_nonneg}]\label{Pphi2-freeSingleSliceCovariance-nonneg}
  \lean{Pphi2.freeSingleSliceCovariance_nonneg}\leanok
  \uses{GaussianField-AsymLatticeSites, GaussianField-covariance, Pphi2-SpatialField, Pphi2-freeSingleSliceCovariance, GaussianField-AsymLatticeField, GaussianField-latticeCovarianceAsymGJ, Pphi2-singleSliceLatticeField, GaussianField-ell2-}
  \texttt{Pphi2.freeSingleSliceCovariance\_nonneg}
\end{theorem}
\begin{theorem}[\texttt{groundVariance\_le\_freeCovariance} (axiom)]\label{Pphi2-groundVariance-le-freeCovariance}
  \lean{Pphi2.groundVariance_le_freeCovariance}\leanok
  \uses{Pphi2-groundSliceVariance, Pphi2-SpatialField, Pphi2-freeSingleSliceCovariance}
  \texttt{Pphi2.groundVariance\_le\_freeCovariance}
\end{theorem}
\begin{theorem}[\texttt{groundVariance\_sum\_le\_freeCovariance\_sum}]\label{Pphi2-groundVariance-sum-le-freeCovariance-sum}
  \lean{Pphi2.groundVariance_sum_le_freeCovariance_sum}\leanok
  \uses{Pphi2-freeSingleSliceCovarianceSum, Pphi2-groundSliceVariance, Pphi2-groundVarianceFreeCovarianceConstant, Pphi2-SpatialField, Pphi2-freeSingleSliceCovariance, Pphi2-groundVariance-le-freeCovariance-with-constant, Pphi2-groundSliceVarianceSum}
  \texttt{Pphi2.groundVariance\_sum\_le\_freeCovariance\_sum}
\end{theorem}
\begin{definition}[\texttt{freeSingleSliceCovarianceSum}]\label{Pphi2-freeSingleSliceCovarianceSum}
  \lean{Pphi2.freeSingleSliceCovarianceSum}\leanok
  \uses{Pphi2-SpatialField, Pphi2-freeSingleSliceCovariance}
  \texttt{Pphi2.freeSingleSliceCovarianceSum}
\end{definition}
\begin{theorem}[\texttt{freeSingleSliceCovariance\_smul}]\label{Pphi2-freeSingleSliceCovariance-smul}
  \lean{Pphi2.freeSingleSliceCovariance_smul}\leanok
  \uses{GaussianField-AsymLatticeSites, GaussianField-covariance, Pphi2-SpatialField, Pphi2-freeSingleSliceCovariance, GaussianField-AsymLatticeField, GaussianField-latticeCovarianceAsymGJ, Pphi2-singleSliceLatticeField, Pphi2-singleSliceLatticeField-smul, GaussianField-ell2-}
  \texttt{Pphi2.freeSingleSliceCovariance\_smul}
\end{theorem}
\begin{theorem}[\texttt{singleSliceLatticeField\_smul}]\label{Pphi2-singleSliceLatticeField-smul}
  \lean{Pphi2.singleSliceLatticeField_smul}\leanok
  \uses{Pphi2-singleSliceFamily-smul, GaussianField-AsymLatticeSites, Pphi2-SpatialField, Pphi2-asymSliceEquiv, Pphi2-singleSliceFamily, GaussianField-AsymLatticeField, Pphi2-singleSliceLatticeField}
  \texttt{Pphi2.singleSliceLatticeField\_smul}
\end{theorem}
\begin{definition}[\texttt{freeSingleSliceCovariance}]\label{Pphi2-freeSingleSliceCovariance}
  \lean{Pphi2.freeSingleSliceCovariance}\leanok
  \uses{GaussianField-AsymLatticeSites, GaussianField-covariance, Pphi2-SpatialField, GaussianField-AsymLatticeField, GaussianField-latticeCovarianceAsymGJ, Pphi2-singleSliceLatticeField, GaussianField-ell2-}
  \texttt{Pphi2.freeSingleSliceCovariance}
\end{definition}
\begin{theorem}[\texttt{freeSingleSliceCovarianceSum\_nonneg}]\label{Pphi2-freeSingleSliceCovarianceSum-nonneg}
  \lean{Pphi2.freeSingleSliceCovarianceSum_nonneg}\leanok
  \uses{Pphi2-freeSingleSliceCovarianceSum, Pphi2-freeSingleSliceCovariance-nonneg, Pphi2-SpatialField, Pphi2-freeSingleSliceCovariance}
  \texttt{Pphi2.freeSingleSliceCovarianceSum\_nonneg}
\end{theorem}
\begin{theorem}[\texttt{congr\_simp}]\label{Pphi2-freeSingleSliceCovarianceSum-congr-simp}
  \lean{Pphi2.freeSingleSliceCovarianceSum.congr_simp}\leanok
  \uses{Pphi2-freeSingleSliceCovarianceSum, Pphi2-SpatialField}
  \texttt{Pphi2.freeSingleSliceCovarianceSum.congr\_simp}
\end{theorem}
\begin{definition}[\texttt{groundSliceVariance}]\label{Pphi2-groundSliceVariance}
  \lean{Pphi2.groundSliceVariance}\leanok
  \uses{Pphi2-asymSliceObsLinear, Pphi2-SpatialField, Pphi2-asymGroundVector}
  \texttt{Pphi2.groundSliceVariance}
\end{definition}
\begin{theorem}[\texttt{singleSliceFamily\_smul}]\label{Pphi2-singleSliceFamily-smul}
  \lean{Pphi2.singleSliceFamily_smul}\leanok
  \uses{Pphi2-SpatialField, Pphi2-singleSliceFamily}
  \texttt{Pphi2.singleSliceFamily\_smul}
\end{theorem}
\begin{theorem}[\texttt{piece3\_pathMeasure\_bound\_to\_freeCovariance\_sum}]\label{Pphi2-piece3-pathMeasure-bound-to-freeCovariance-sum}
  \lean{Pphi2.piece3_pathMeasure_bound_to_freeCovariance_sum}\leanok
  \uses{Pphi2-freeSingleSliceCovarianceSum, Pphi2-groundVarianceFreeCovarianceConstant, Pphi2-SpatialField, Pphi2-asymTransferSystem, Pphi2-asymSliceFamilyLinear, Pphi2-groundVariance-sum-le-freeCovariance-sum, Pphi2-groundSliceVarianceSum}
  \texttt{Pphi2.piece3\_pathMeasure\_bound\_to\_freeCovariance\_sum}
\end{theorem}
\begin{theorem}[\texttt{groundVarianceFreeCovarianceConstant\_pos}]\label{Pphi2-groundVarianceFreeCovarianceConstant-pos}
  \lean{Pphi2.groundVarianceFreeCovarianceConstant_pos}\leanok
  \uses{Pphi2-groundSliceVariance, Pphi2-groundVarianceFreeCovarianceConstant, Pphi2-SpatialField, Pphi2-freeSingleSliceCovariance, Pphi2-groundVariance-le-freeCovariance}
  \texttt{Pphi2.groundVarianceFreeCovarianceConstant\_pos}
\end{theorem}
\begin{definition}[\texttt{groundVarianceFreeCovarianceConstant}]\label{Pphi2-groundVarianceFreeCovarianceConstant}
  \lean{Pphi2.groundVarianceFreeCovarianceConstant}\leanok
  \uses{Pphi2-groundSliceVariance, Pphi2-SpatialField, Pphi2-freeSingleSliceCovariance, Pphi2-groundVariance-le-freeCovariance}
  \texttt{Pphi2.groundVarianceFreeCovarianceConstant}
\end{definition}
\begin{definition}[\texttt{singleSliceLatticeField}]\label{Pphi2-singleSliceLatticeField}
  \lean{Pphi2.singleSliceLatticeField}\leanok
  \uses{GaussianField-AsymLatticeSites, Pphi2-SpatialField, Pphi2-asymSliceEquiv, Pphi2-singleSliceFamily, GaussianField-AsymLatticeField}
  \texttt{Pphi2.singleSliceLatticeField}
\end{definition}
\begin{definition}[\texttt{singleSliceFamily}]\label{Pphi2-singleSliceFamily}
  \lean{Pphi2.singleSliceFamily}\leanok
  \uses{Pphi2-SpatialField}
  \texttt{Pphi2.singleSliceFamily}
\end{definition}
\begin{theorem}[\texttt{congr\_simp}]\label{Pphi2-groundSliceVarianceSum-congr-simp}
  \lean{Pphi2.groundSliceVarianceSum.congr_simp}\leanok
  \uses{Pphi2-SpatialField, Pphi2-groundSliceVarianceSum}
  \texttt{Pphi2.groundSliceVarianceSum.congr\_simp}
\end{theorem}
\begin{theorem}[\texttt{congr\_simp}]\label{Pphi2-groundVarianceFreeCovarianceConstant-congr-simp}
  \lean{Pphi2.groundVarianceFreeCovarianceConstant.congr_simp}\leanok
  \uses{Pphi2-groundVarianceFreeCovarianceConstant}
  \texttt{Pphi2.groundVarianceFreeCovarianceConstant.congr\_simp}
\end{theorem}
\section{\texttt{Pphi2.TransferMatrix.Positivity}}
\begin{theorem}[\texttt{massGap\_pos}]\label{Pphi2-massGap-pos}
  \lean{Pphi2.massGap_pos}\leanok
  \uses{Pphi2-groundEnergyLevel, Pphi2-firstExcitedEnergyLevel, Pphi2-energyLevel-gap, Pphi2-massGap}
  $m_{\mathrm{phys}} > 0$. Immediate from \texttt{energyLevel\_gap}. Fully proved.
\end{theorem}
\begin{definition}[\texttt{transferGroundExcitedData}]\label{Pphi2-transferGroundExcitedData}
  \lean{Pphi2.transferGroundExcitedData}\leanok
  \uses{Pphi2-TransferGroundExcitedData}
  Noncomputable choice of \texttt{TransferGroundExcitedData} from \texttt{transferOperator\_ground\_simple\_spectral}.
\end{definition}
\begin{definition}[\texttt{transferFirstExcitedEigenvalue}]\label{Pphi2-transferFirstExcitedEigenvalue}
  \lean{Pphi2.transferFirstExcitedEigenvalue}\leanok
  \uses{Pphi2-TransferGroundExcitedData, Pphi2-TransferGroundExcitedData-eigenval, Pphi2-TransferGroundExcitedData-i-, Pphi2-transferGroundExcitedData}
  \texttt{Pphi2.transferFirstExcitedEigenvalue}
\end{definition}
\begin{definition}[\texttt{TransferGroundExcitedData}]\label{Pphi2-TransferGroundExcitedData}
  \lean{Pphi2.TransferGroundExcitedData}\leanok
  ``\texttt{lean structure TransferGroundExcitedData (P : InteractionPolynomial) (a mass : R) (ha : 0 < a) (hmass : 0 < mass) }`` Packaged spectral decomposition with distinguished ground ($i_0$) and first-excited ($i_1$) indices.
\end{definition}
\begin{theorem}[\texttt{energyLevel\_gap}]\label{Pphi2-energyLevel-gap}
  \lean{Pphi2.energyLevel_gap}\leanok
  \uses{Pphi2-TransferGroundExcitedData-i-, Pphi2-transferOperator-eigenvalues-pos, Pphi2-TransferGroundExcitedData, Pphi2-groundEnergyLevel, Pphi2-TransferGroundExcitedData-eigenval, Pphi2-firstExcitedEnergyLevel, Pphi2-TransferGroundExcitedData-b, Pphi2-TransferGroundExcitedData--, Pphi2-TransferGroundExcitedData-hlt, Pphi2-TransferGroundExcitedData-h-eigen, Pphi2-transferFirstExcitedEigenvalue, Pphi2-TransferGroundExcitedData-i-, Pphi2-transferGroundEigenvalue, Pphi2-transferGroundExcitedData}
  $E_0 < E_1$. From $\lambda_0 > \lambda_1 > 0$ and $-\frac{1}{a} < 0$ reversing the log inequality. Fully proved.
\end{theorem}
\begin{theorem}[\texttt{h\_eigen}]\label{Pphi2-TransferGroundExcitedData-h-eigen}
  \lean{Pphi2.TransferGroundExcitedData.h_eigen}\leanok
  \uses{Pphi2-L2SpatialField, Pphi2-SpatialField, Pphi2-TransferGroundExcitedData, Pphi2-TransferGroundExcitedData-eigenval, Pphi2-TransferGroundExcitedData-b, Pphi2-transferOperatorCLM, Pphi2-TransferGroundExcitedData--}
  \texttt{Pphi2.TransferGroundExcitedData.h\_eigen}
\end{theorem}
\begin{definition}[\texttt{ι}]\label{Pphi2-TransferGroundExcitedData--}
  \lean{Pphi2.TransferGroundExcitedData.ι}\leanok
  \uses{Pphi2-TransferGroundExcitedData}
  \texttt{Pphi2.TransferGroundExcitedData.ι}
\end{definition}
\begin{definition}[\texttt{b}]\label{Pphi2-TransferGroundExcitedData-b}
  \lean{Pphi2.TransferGroundExcitedData.b}\leanok
  \uses{Pphi2-L2SpatialField, Pphi2-SpatialField, Pphi2-TransferGroundExcitedData, Pphi2-TransferGroundExcitedData--}
  \texttt{Pphi2.TransferGroundExcitedData.b}
\end{definition}
\begin{definition}[\texttt{groundEnergyLevel}]\label{Pphi2-groundEnergyLevel}
  \lean{Pphi2.groundEnergyLevel}\leanok
  \uses{Pphi2-transferGroundEigenvalue}
  $E_0 = -\frac{1}{a}\log\lambda_0$ and $E_1 = -\frac{1}{a}\log\lambda_1$.
\end{definition}
\begin{definition}[\texttt{i₀}]\label{Pphi2-TransferGroundExcitedData-i-}
  \lean{Pphi2.TransferGroundExcitedData.i₀}\leanok
  \uses{Pphi2-TransferGroundExcitedData, Pphi2-TransferGroundExcitedData--}
  \texttt{Pphi2.TransferGroundExcitedData.i₀}
\end{definition}
\begin{definition}[\texttt{i₁}]\label{Pphi2-TransferGroundExcitedData-i-}
  \lean{Pphi2.TransferGroundExcitedData.i₁}\leanok
  \uses{Pphi2-TransferGroundExcitedData, Pphi2-TransferGroundExcitedData--}
  \texttt{Pphi2.TransferGroundExcitedData.i₁}
\end{definition}
\begin{definition}[\texttt{eigenval}]\label{Pphi2-TransferGroundExcitedData-eigenval}
  \lean{Pphi2.TransferGroundExcitedData.eigenval}\leanok
  \uses{Pphi2-TransferGroundExcitedData, Pphi2-TransferGroundExcitedData--}
  \texttt{Pphi2.TransferGroundExcitedData.eigenval}
\end{definition}
\begin{theorem}[\texttt{h\_sum}]\label{Pphi2-TransferGroundExcitedData-h-sum}
  \lean{Pphi2.TransferGroundExcitedData.h_sum}\leanok
  \uses{Pphi2-L2SpatialField, Pphi2-SpatialField, Pphi2-TransferGroundExcitedData, Pphi2-TransferGroundExcitedData-eigenval, Pphi2-TransferGroundExcitedData-b, Pphi2-transferOperatorCLM, Pphi2-TransferGroundExcitedData--}
  \texttt{Pphi2.TransferGroundExcitedData.h\_sum}
\end{theorem}
\begin{definition}[\texttt{massGap}]\label{Pphi2-massGap}
  \lean{Pphi2.massGap}\leanok
  \uses{Pphi2-groundEnergyLevel, Pphi2-firstExcitedEnergyLevel}
  ``\texttt{lean def massGap (P : InteractionPolynomial) (a mass : R) (ha : 0 < a) (hmass : 0 < mass) : R }`` Physical mass gap: $m_{\mathrm{phys}} = E_1 - E_0$.
\end{definition}
\begin{definition}[\texttt{ctorIdx}]\label{Pphi2-TransferGroundExcitedData-ctorIdx}
  \lean{Pphi2.TransferGroundExcitedData.ctorIdx}\leanok
  \uses{Pphi2-TransferGroundExcitedData}
  \texttt{Pphi2.TransferGroundExcitedData.ctorIdx}
\end{definition}
\begin{theorem}[\texttt{hlt}]\label{Pphi2-TransferGroundExcitedData-hlt}
  \lean{Pphi2.TransferGroundExcitedData.hlt}\leanok
  \uses{Pphi2-TransferGroundExcitedData-i-, Pphi2-TransferGroundExcitedData, Pphi2-TransferGroundExcitedData-eigenval, Pphi2-TransferGroundExcitedData-i-}
  \texttt{Pphi2.TransferGroundExcitedData.hlt}
\end{theorem}
\begin{definition}[\texttt{transferGroundEigenvalue}]\label{Pphi2-transferGroundEigenvalue}
  \lean{Pphi2.transferGroundEigenvalue}\leanok
  \uses{Pphi2-TransferGroundExcitedData-i-, Pphi2-TransferGroundExcitedData, Pphi2-TransferGroundExcitedData-eigenval, Pphi2-transferGroundExcitedData}
  $\lambda_0$ and $\lambda_1$ extracted from spectral data.
\end{definition}
\begin{theorem}[\texttt{hi\_ne}]\label{Pphi2-TransferGroundExcitedData-hi-ne}
  \lean{Pphi2.TransferGroundExcitedData.hi_ne}\leanok
  \uses{Pphi2-TransferGroundExcitedData-i-, Pphi2-TransferGroundExcitedData, Pphi2-TransferGroundExcitedData--, Pphi2-TransferGroundExcitedData-i-}
  \texttt{Pphi2.TransferGroundExcitedData.hi\_ne}
\end{theorem}
\begin{definition}[\texttt{firstExcitedEnergyLevel}]\label{Pphi2-firstExcitedEnergyLevel}
  \lean{Pphi2.firstExcitedEnergyLevel}\leanok
  \uses{Pphi2-transferFirstExcitedEigenvalue}
  \texttt{Pphi2.firstExcitedEnergyLevel}
\end{definition}
\section{\texttt{Pphi2.AsymTorus.AsymJentzsch}}
\begin{theorem}[\texttt{asymTransferOperator\_spectral}]\label{Pphi2-asymTransferOperator-spectral}
  \lean{Pphi2.asymTransferOperator_spectral}\leanok
  \uses{Pphi2-L2SpatialField, Pphi2-SpatialField, Pphi2-asymTransferOperator-isCompact, GaussianField-compact-selfAdjoint-spectral, Pphi2-asymTransferOperator-isSelfAdjoint, Pphi2-asymTransferOperatorCLM}
  \texttt{Pphi2.asymTransferOperator\_spectral}
\end{theorem}
\begin{theorem}[\texttt{congr\_simp}]\label{Pphi2-asymTransferOperatorCLM-congr-simp}
  \lean{Pphi2.asymTransferOperatorCLM.congr_simp}\leanok
  \uses{Pphi2-L2SpatialField, Pphi2-SpatialField, Pphi2-asymTransferOperatorCLM}
  \texttt{Pphi2.asymTransferOperatorCLM.congr\_simp}
\end{theorem}
\begin{theorem}[\texttt{asymTransferOperator\_ground\_simple\_spectral}]\label{Pphi2-asymTransferOperator-ground-simple-spectral}
  \lean{Pphi2.asymTransferOperator_ground_simple_spectral}\leanok
  \uses{Pphi2-L2SpatialField, Pphi2-SpatialField, Pphi2-asymTransferOperator-spectral, Pphi2-asymTransferOperator-ground-simple, Pphi2-asymTransferOperatorCLM}
  \texttt{Pphi2.asymTransferOperator\_ground\_simple\_spectral}
\end{theorem}
\begin{theorem}[\texttt{asymTransferOperator\_inner\_nonneg}]\label{Pphi2-asymTransferOperator-inner-nonneg}
  \lean{Pphi2.asymTransferOperator_inner_nonneg}\leanok
  \uses{Pphi2-asymTransferOperator-strictly-positive-definite, Pphi2-L2SpatialField, Pphi2-SpatialField, Pphi2-asymTransferOperatorCLM}
  \texttt{Pphi2.asymTransferOperator\_inner\_nonneg}
\end{theorem}
\begin{theorem}[\texttt{asymTransferOperator\_ground\_simple}]\label{Pphi2-asymTransferOperator-ground-simple}
  \lean{Pphi2.asymTransferOperator_ground_simple}\leanok
  \uses{Pphi2-L2SpatialField, Pphi2-asymTransferOperator-eigenvalues-pos, Pphi2-SpatialField, Pphi2-asymTransferOperator-isCompact, Pphi2-l2SpatialField-hilbertBasis-nontrivial, Pphi2-asymTransferOperator-positivityImproving, Pphi2-asymTransferOperator-isSelfAdjoint, Pphi2-asymTransferOperatorCLM}
  \texttt{Pphi2.asymTransferOperator\_ground\_simple}
\end{theorem}
\begin{theorem}[\texttt{asymTransferOperator\_strictly\_positive\_definite}]\label{Pphi2-asymTransferOperator-strictly-positive-definite}
  \lean{Pphi2.asymTransferOperator_strictly_positive_definite}\leanok
  \uses{Pphi2-asymTransferWeight-measurable, Pphi2-L2SpatialField, Pphi2-SpatialField, Pphi2-asymTransferWeight, Pphi2-transferGaussian-memLp, Pphi2-asymTransferWeight-pos, Pphi2-convolution-gaussian-strictly-positive-definite, Pphi2-asymTransferWeight-bound, Pphi2-asymTransferOperatorCLM, Pphi2-transferGaussian}
  \texttt{Pphi2.asymTransferOperator\_strictly\_positive\_definite}
\end{theorem}
\begin{theorem}[\texttt{asymTransferOperator\_eigenvalues\_pos}]\label{Pphi2-asymTransferOperator-eigenvalues-pos}
  \lean{Pphi2.asymTransferOperator_eigenvalues_pos}\leanok
  \uses{Pphi2-asymTransferOperator-strictly-positive-definite, Pphi2-L2SpatialField, Pphi2-SpatialField, Pphi2-asymTransferOperatorCLM}
  \texttt{Pphi2.asymTransferOperator\_eigenvalues\_pos}
\end{theorem}
\begin{theorem}[\texttt{asymTransferOperator\_positivityImproving}]\label{Pphi2-asymTransferOperator-positivityImproving}
  \lean{Pphi2.asymTransferOperator_positivityImproving}\leanok
  \uses{Pphi2-asymTransferWeight-measurable, Pphi2-L2SpatialField, Pphi2-SpatialField, Pphi2-asymTransferWeight, Pphi2-transferGaussian-pos, Pphi2-transferGaussian-memLp, Pphi2-transferGaussian-memLp-two, Pphi2-asymTransferWeight-pos, Pphi2-asymTransferWeight-bound, Pphi2-asymTransferOperatorCLM, Pphi2-transferGaussian}
  \texttt{Pphi2.asymTransferOperator\_positivityImproving}
\end{theorem}
\section{\texttt{Pphi2.OSAxioms}}
\begin{definition}[\texttt{schwartzIm}]\label{Pphi2-schwartzIm}
  \lean{Pphi2.schwartzIm}\leanok
  \uses{Pphi2-EuclideanOS-schwartzIm, Pphi2-TestFunction2-, Pphi2-plane2Background, Pphi2-TestFunction2}
  \texttt{Pphi2.schwartzIm}
\end{definition}
\begin{definition}[\texttt{OS4\_Ergodicity}]\label{Pphi2-OS4-Ergodicity}
  \lean{Pphi2.OS4_Ergodicity}\leanok
  \uses{Pphi2-EuclideanPlaneBackground-SpaceTime, Pphi2-FieldConfig2, Pphi2-EuclideanPlaneBackground-RealTestFunction, Pphi2-planeBackground, Pphi2-EuclideanPlaneBackground-dim, Pphi2-plane2Background, GaussianField-instMeasurableSpaceConfiguration, Pphi2-EuclideanOS-OS4-Ergodicity, Pphi2-plane2TimeStructure}
  \texttt{Pphi2.OS4\_Ergodicity}
\end{definition}
\begin{definition}[\texttt{generatingFunctional}]\label{Pphi2-generatingFunctional}
  \lean{Pphi2.generatingFunctional}\leanok
  \uses{Pphi2-EuclideanPlaneBackground-SpaceTime, Pphi2-FieldConfig2, Pphi2-EuclideanPlaneBackground-RealTestFunction, Pphi2-planeBackground, Pphi2-EuclideanPlaneBackground-dim, Pphi2-plane2Background, GaussianField-instMeasurableSpaceConfiguration, Pphi2-TestFunction2, Pphi2-EuclideanOS-generatingFunctional}
  $Z[f] = \int e^{i\omega(f)}\,d\mu$ (real) and $Z[J] = \int e^{i\langle\omega,J\rangle_\mathbb{C}}\,d\mu$ (complex).
\end{definition}
\begin{definition}[\texttt{SatisfiesFullOS}]\label{Pphi2-SatisfiesFullOS}
  \lean{Pphi2.SatisfiesFullOS}\leanok
  \uses{Pphi2-EuclideanPlaneBackground-SpaceTime, Pphi2-FieldConfig2, Pphi2-EuclideanPlaneBackground-RealTestFunction, Pphi2-planeBackground, Pphi2-EuclideanOS-SatisfiesFullOS, Pphi2-EuclideanPlaneBackground-dim, Pphi2-plane2Background, GaussianField-instMeasurableSpaceConfiguration, Pphi2-plane2TimeStructure}
  Structure bundling all five axioms.
\end{definition}
\begin{definition}[\texttt{OS0\_Analyticity}]\label{Pphi2-OS0-Analyticity}
  \lean{Pphi2.OS0_Analyticity}\leanok
  \uses{Pphi2-EuclideanPlaneBackground-SpaceTime, Pphi2-FieldConfig2, Pphi2-EuclideanPlaneBackground-RealTestFunction, Pphi2-planeBackground, Pphi2-EuclideanPlaneBackground-dim, Pphi2-plane2Background, GaussianField-instMeasurableSpaceConfiguration, Pphi2-EuclideanOS-OS0-Analyticity}
  \texttt{Pphi2.OS0\_Analyticity}
\end{definition}
\begin{definition}[\texttt{distribPairing}]\label{Pphi2-distribPairing}
  \lean{Pphi2.distribPairing}\leanok
  \uses{Pphi2-FieldConfig2, Pphi2-EuclideanOS-distribPairing, Pphi2-plane2Background, Pphi2-TestFunction2}
  \texttt{Pphi2.distribPairing}
\end{definition}
\begin{definition}[\texttt{OS3\_ReflectionPositivity}]\label{Pphi2-OS3-ReflectionPositivity}
  \lean{Pphi2.OS3_ReflectionPositivity}\leanok
  \uses{Pphi2-EuclideanPlaneBackground-SpaceTime, Pphi2-FieldConfig2, Pphi2-EuclideanPlaneBackground-RealTestFunction, Pphi2-planeBackground, Pphi2-EuclideanPlaneBackground-dim, Pphi2-plane2Background, Pphi2-EuclideanOS-OS3-ReflectionPositivity, GaussianField-instMeasurableSpaceConfiguration, Pphi2-plane2TimeStructure}
  \texttt{Pphi2.OS3\_ReflectionPositivity}
\end{definition}
\begin{definition}[\texttt{OS1\_Regularity}]\label{Pphi2-OS1-Regularity}
  \lean{Pphi2.OS1_Regularity}\leanok
  \uses{Pphi2-EuclideanPlaneBackground-SpaceTime, Pphi2-FieldConfig2, Pphi2-EuclideanPlaneBackground-RealTestFunction, Pphi2-planeBackground, Pphi2-EuclideanPlaneBackground-dim, Pphi2-plane2Background, GaussianField-instMeasurableSpaceConfiguration, Pphi2-EuclideanOS-OS1-Regularity}
  \texttt{Pphi2.OS1\_Regularity}
\end{definition}
\begin{definition}[\texttt{generatingFunctionalℂ}]\label{Pphi2-generatingFunctional-}
  \lean{Pphi2.generatingFunctionalℂ}\leanok
  \uses{Pphi2-EuclideanPlaneBackground-SpaceTime, Pphi2-FieldConfig2, Pphi2-EuclideanPlaneBackground-RealTestFunction, Pphi2-planeBackground, Pphi2-EuclideanPlaneBackground-dim, Pphi2-TestFunction2-, Pphi2-plane2Background, Pphi2-EuclideanOS-generatingFunctional-, GaussianField-instMeasurableSpaceConfiguration}
  \texttt{Pphi2.generatingFunctionalℂ}
\end{definition}
\begin{definition}[\texttt{OS2\_EuclideanInvariance}]\label{Pphi2-OS2-EuclideanInvariance}
  \lean{Pphi2.OS2_EuclideanInvariance}\leanok
  \uses{Pphi2-EuclideanPlaneBackground-SpaceTime, Pphi2-FieldConfig2, Pphi2-EuclideanOS-OS2-EuclideanInvariance, Pphi2-EuclideanPlaneBackground-RealTestFunction, Pphi2-planeBackground, Pphi2-EuclideanPlaneBackground-dim, Pphi2-plane2Background, GaussianField-instMeasurableSpaceConfiguration}
  \texttt{Pphi2.OS2\_EuclideanInvariance}
\end{definition}
\begin{definition}[\texttt{schwartzRe}]\label{Pphi2-schwartzRe}
  \lean{Pphi2.schwartzRe}\leanok
  \uses{Pphi2-EuclideanOS-schwartzRe, Pphi2-TestFunction2-, Pphi2-plane2Background, Pphi2-TestFunction2}
  Extract real/imaginary parts of complex Schwartz functions as real Schwartz functions.
\end{definition}
\begin{definition}[\texttt{OS4\_Clustering}]\label{Pphi2-OS4-Clustering}
  \lean{Pphi2.OS4_Clustering}\leanok
  \uses{Pphi2-EuclideanPlaneBackground-SpaceTime, Pphi2-FieldConfig2, Pphi2-EuclideanPlaneBackground-RealTestFunction, Pphi2-planeBackground, Pphi2-EuclideanOS-OS4-Clustering, Pphi2-EuclideanPlaneBackground-dim, Pphi2-plane2Background, GaussianField-instMeasurableSpaceConfiguration}
  \texttt{Pphi2.OS4\_Clustering}
\end{definition}
\section{\texttt{Pphi2.Backgrounds.EuclideanTime}}
\begin{definition}[\texttt{timeTranslateOnRealTestFunctions}]\label{Pphi2-EuclideanTimeStructure-timeTranslateOnRealTestFunctions}
  \lean{Pphi2.EuclideanTimeStructure.timeTranslateOnRealTestFunctions}\leanok
  \uses{Pphi2-EuclideanPlaneBackground-SpaceTime, Pphi2-EuclideanTimeStructure-timeTranslation, Pphi2-EuclideanPlaneBackground-RealTestFunction, Pphi2-EuclideanPlaneBackground-translate, Pphi2-EuclideanPlaneBackground-dim, Pphi2-EuclideanPlaneBackground, Pphi2-EuclideanTimeStructure}
  \texttt{Pphi2.EuclideanTimeStructure.timeTranslateOnRealTestFunctions}
\end{definition}
\begin{definition}[\texttt{reflectCLE}]\label{Pphi2-EuclideanTimeStructure-reflectCLE}
  \lean{Pphi2.EuclideanTimeStructure.reflectCLE}\leanok
  \uses{Pphi2-EuclideanPlaneBackground-SpaceTime, Pphi2-EuclideanTimeStructure-reflect-involutive, Pphi2-EuclideanPlaneBackground-dim, Pphi2-EuclideanPlaneBackground, Pphi2-EuclideanTimeStructure, Pphi2-EuclideanTimeStructure-reflect}
  \texttt{Pphi2.EuclideanTimeStructure.reflectCLE}
\end{definition}
\begin{theorem}[\texttt{timeTranslateOnComplexTestFunctions\_apply}]\label{Pphi2-EuclideanTimeStructure-timeTranslateOnComplexTestFunctions-apply}
  \lean{Pphi2.EuclideanTimeStructure.timeTranslateOnComplexTestFunctions_apply}\leanok
  \uses{Pphi2-EuclideanPlaneBackground-SpaceTime, Pphi2-EuclideanTimeStructure-timeTranslateOnComplexTestFunctions, Pphi2-EuclideanTimeStructure-timeTranslation, Pphi2-EuclideanTimeStructure-timeShift, Pphi2-EuclideanPlaneBackground-translateComplex, Pphi2-EuclideanPlaneBackground-dim, Pphi2-EuclideanPlaneBackground-translateComplex-apply, Pphi2-EuclideanPlaneBackground, Pphi2-EuclideanTimeStructure, Pphi2-EuclideanPlaneBackground-ComplexTestFunction}
  \texttt{Pphi2.EuclideanTimeStructure.timeTranslateOnComplexTestFunctions\_apply}
\end{theorem}
\begin{definition}[\texttt{positiveTimeSet}]\label{Pphi2-EuclideanTimeStructure-positiveTimeSet}
  \lean{Pphi2.EuclideanTimeStructure.positiveTimeSet}\leanok
  \uses{Pphi2-EuclideanPlaneBackground-SpaceTime, Pphi2-EuclideanPlaneBackground, Pphi2-EuclideanTimeStructure}
  \texttt{Pphi2.EuclideanTimeStructure.positiveTimeSet}
\end{definition}
\begin{definition}[\texttt{timeShift}]\label{Pphi2-EuclideanTimeStructure-timeShift}
  \lean{Pphi2.EuclideanTimeStructure.timeShift}\leanok
  \uses{Pphi2-EuclideanPlaneBackground-SpaceTime, Pphi2-EuclideanTimeStructure-timeTranslation, Pphi2-EuclideanPlaneBackground-dim, Pphi2-EuclideanPlaneBackground, Pphi2-EuclideanTimeStructure}
  \texttt{Pphi2.EuclideanTimeStructure.timeShift}
\end{definition}
\begin{theorem}[\texttt{timeShift\_add}]\label{Pphi2-EuclideanTimeStructure-timeShift-add}
  \lean{Pphi2.EuclideanTimeStructure.timeShift_add}\leanok
  \uses{Pphi2-EuclideanPlaneBackground-SpaceTime, Pphi2-EuclideanTimeStructure-timeTranslation, Pphi2-EuclideanTimeStructure-timeShift, Pphi2-EuclideanTimeStructure-timeAxis, Pphi2-EuclideanPlaneBackground-dim, Pphi2-EuclideanPlaneBackground, Pphi2-EuclideanTimeStructure}
  \texttt{Pphi2.EuclideanTimeStructure.timeShift\_add}
\end{theorem}
\begin{theorem}[\texttt{timeTranslateOnRealTestFunctions\_apply}]\label{Pphi2-EuclideanTimeStructure-timeTranslateOnRealTestFunctions-apply}
  \lean{Pphi2.EuclideanTimeStructure.timeTranslateOnRealTestFunctions_apply}\leanok
  \uses{Pphi2-EuclideanPlaneBackground-SpaceTime, Pphi2-EuclideanTimeStructure-timeTranslation, Pphi2-EuclideanTimeStructure-timeShift, Pphi2-EuclideanPlaneBackground-RealTestFunction, Pphi2-EuclideanPlaneBackground-translate, Pphi2-EuclideanTimeStructure-timeTranslateOnRealTestFunctions, Pphi2-EuclideanPlaneBackground-dim, Pphi2-EuclideanPlaneBackground, Pphi2-EuclideanTimeStructure, Pphi2-EuclideanPlaneBackground-translate-apply}
  \texttt{Pphi2.EuclideanTimeStructure.timeTranslateOnRealTestFunctions\_apply}
\end{theorem}
\begin{definition}[\texttt{reflect}]\label{Pphi2-EuclideanTimeStructure-reflect}
  \lean{Pphi2.EuclideanTimeStructure.reflect}\leanok
  \uses{Pphi2-EuclideanPlaneBackground-SpaceTime, Pphi2-EuclideanPlaneBackground-dim, Pphi2-EuclideanPlaneBackground, Pphi2-EuclideanTimeStructure}
  \texttt{Pphi2.EuclideanTimeStructure.reflect}
\end{definition}
\begin{definition}[\texttt{positiveTimeSubmodule}]\label{Pphi2-EuclideanTimeStructure-positiveTimeSubmodule}
  \lean{Pphi2.EuclideanTimeStructure.positiveTimeSubmodule}\leanok
  \uses{Pphi2-EuclideanPlaneBackground-SpaceTime, Pphi2-EuclideanTimeStructure-positiveTimeSet, Pphi2-EuclideanPlaneBackground-RealTestFunction, Pphi2-EuclideanPlaneBackground-dim, Pphi2-EuclideanPlaneBackground, Pphi2-EuclideanTimeStructure}
  \texttt{Pphi2.EuclideanTimeStructure.positiveTimeSubmodule}
\end{definition}
\begin{definition}[\texttt{EuclideanTimeStructure}]\label{Pphi2-EuclideanTimeStructure}
  \lean{Pphi2.EuclideanTimeStructure}\leanok
  \uses{Pphi2-EuclideanPlaneBackground}
  \texttt{Pphi2.EuclideanTimeStructure}
\end{definition}
\begin{definition}[\texttt{timeAxis}]\label{Pphi2-EuclideanTimeStructure-timeAxis}
  \lean{Pphi2.EuclideanTimeStructure.timeAxis}\leanok
  \uses{Pphi2-EuclideanPlaneBackground-SpaceTime, Pphi2-EuclideanPlaneBackground-dim, Pphi2-EuclideanPlaneBackground, Pphi2-EuclideanTimeStructure}
  \texttt{Pphi2.EuclideanTimeStructure.timeAxis}
\end{definition}
\begin{definition}[\texttt{timeTranslateOnComplexTestFunctions}]\label{Pphi2-EuclideanTimeStructure-timeTranslateOnComplexTestFunctions}
  \lean{Pphi2.EuclideanTimeStructure.timeTranslateOnComplexTestFunctions}\leanok
  \uses{Pphi2-EuclideanPlaneBackground-SpaceTime, Pphi2-EuclideanTimeStructure-timeTranslation, Pphi2-EuclideanPlaneBackground-translateComplex, Pphi2-EuclideanPlaneBackground-dim, Pphi2-EuclideanPlaneBackground, Pphi2-EuclideanTimeStructure, Pphi2-EuclideanPlaneBackground-ComplexTestFunction}
  \texttt{Pphi2.EuclideanTimeStructure.timeTranslateOnComplexTestFunctions}
\end{definition}
\begin{definition}[\texttt{reflectOnComplexTestFunctions}]\label{Pphi2-EuclideanTimeStructure-reflectOnComplexTestFunctions}
  \lean{Pphi2.EuclideanTimeStructure.reflectOnComplexTestFunctions}\leanok
  \uses{Pphi2-EuclideanPlaneBackground-SpaceTime, Pphi2-EuclideanPlaneBackground-dim, Pphi2-EuclideanPlaneBackground, Pphi2-EuclideanTimeStructure, Pphi2-EuclideanTimeStructure-reflectCLE, Pphi2-EuclideanPlaneBackground-ComplexTestFunction}
  \texttt{Pphi2.EuclideanTimeStructure.reflectOnComplexTestFunctions}
\end{definition}
\begin{definition}[\texttt{reflectOnRealTestFunctions}]\label{Pphi2-EuclideanTimeStructure-reflectOnRealTestFunctions}
  \lean{Pphi2.EuclideanTimeStructure.reflectOnRealTestFunctions}\leanok
  \uses{Pphi2-EuclideanPlaneBackground-SpaceTime, Pphi2-EuclideanPlaneBackground-RealTestFunction, Pphi2-EuclideanPlaneBackground-dim, Pphi2-EuclideanPlaneBackground, Pphi2-EuclideanTimeStructure, Pphi2-EuclideanTimeStructure-reflectCLE}
  \texttt{Pphi2.EuclideanTimeStructure.reflectOnRealTestFunctions}
\end{definition}
\begin{theorem}[\texttt{timeShift\_apply}]\label{Pphi2-EuclideanTimeStructure-timeShift-apply}
  \lean{Pphi2.EuclideanTimeStructure.timeShift_apply}\leanok
  \uses{Pphi2-EuclideanPlaneBackground-SpaceTime, Pphi2-EuclideanTimeStructure-timeTranslation, Pphi2-EuclideanTimeStructure-timeShift, Pphi2-EuclideanPlaneBackground-dim, Pphi2-EuclideanPlaneBackground, Pphi2-EuclideanTimeStructure}
  \texttt{Pphi2.EuclideanTimeStructure.timeShift\_apply}
\end{theorem}
\begin{theorem}[\texttt{reflect\_involutive}]\label{Pphi2-EuclideanTimeStructure-reflect-involutive}
  \lean{Pphi2.EuclideanTimeStructure.reflect_involutive}\leanok
  \uses{Pphi2-EuclideanPlaneBackground-SpaceTime, Pphi2-EuclideanPlaneBackground-dim, Pphi2-EuclideanPlaneBackground, Pphi2-EuclideanTimeStructure, Pphi2-EuclideanTimeStructure-reflect}
  \texttt{Pphi2.EuclideanTimeStructure.reflect\_involutive}
\end{theorem}
\begin{definition}[\texttt{ctorIdx}]\label{Pphi2-EuclideanTimeStructure-ctorIdx}
  \lean{Pphi2.EuclideanTimeStructure.ctorIdx}\leanok
  \uses{Pphi2-EuclideanPlaneBackground, Pphi2-EuclideanTimeStructure}
  \texttt{Pphi2.EuclideanTimeStructure.ctorIdx}
\end{definition}
\begin{definition}[\texttt{timeTranslation}]\label{Pphi2-EuclideanTimeStructure-timeTranslation}
  \lean{Pphi2.EuclideanTimeStructure.timeTranslation}\leanok
  \uses{Pphi2-EuclideanPlaneBackground-SpaceTime, Pphi2-EuclideanTimeStructure-timeAxis, Pphi2-EuclideanPlaneBackground-dim, Pphi2-EuclideanPlaneBackground, Pphi2-EuclideanTimeStructure}
  \texttt{Pphi2.EuclideanTimeStructure.timeTranslation}
\end{definition}
\begin{theorem}[\texttt{timeShift\_zero}]\label{Pphi2-EuclideanTimeStructure-timeShift-zero}
  \lean{Pphi2.EuclideanTimeStructure.timeShift_zero}\leanok
  \uses{Pphi2-EuclideanPlaneBackground-SpaceTime, Pphi2-EuclideanTimeStructure-timeTranslation, Pphi2-EuclideanTimeStructure-timeShift, Pphi2-EuclideanTimeStructure-timeAxis, Pphi2-EuclideanPlaneBackground-dim, Pphi2-EuclideanPlaneBackground, Pphi2-EuclideanTimeStructure}
  \texttt{Pphi2.EuclideanTimeStructure.timeShift\_zero}
\end{theorem}
\section{\texttt{Pphi2.TransferMatrix.L2Operator}}
\begin{theorem}[\texttt{transferOperator\_isSelfAdjoint}]\label{Pphi2-transferOperator-isSelfAdjoint}
  \lean{Pphi2.transferOperator_isSelfAdjoint}\leanok
  \uses{Pphi2-L2SpatialField, Pphi2-SpatialField, Pphi2-transferWeight, Pphi2-transferGaussian-memLp, Pphi2-transferWeight-measurable, Pphi2-transferOperatorCLM, Pphi2-transferGaussian, Pphi2-transferGaussian-neg, Pphi2-transferWeight-bound}
  $T$ is self-adjoint, from self-adjointness of $M_w$ and $\mathrm{Conv}_G$. Fully proved.
\end{theorem}
\begin{theorem}[\texttt{congr\_simp}]\label{Pphi2-spatialKinetic-congr-simp}
  \lean{Pphi2.spatialKinetic.congr_simp}\leanok
  \uses{Pphi2-SpatialField, Pphi2-spatialKinetic}
  \texttt{Pphi2.spatialKinetic.congr\_simp}
\end{theorem}
\begin{theorem}[\texttt{transferGaussian\_memLp}]\label{Pphi2-transferGaussian-memLp}
  \lean{Pphi2.transferGaussian_memLp}\leanok
  \uses{Pphi2-SpatialField, Pphi2-transferGaussian, Lattice-toSemilatticeInf}
  \texttt{Pphi2.transferGaussian\_memLp}
\end{theorem}
\begin{theorem}[\texttt{continuous\_transferGaussian}]\label{Pphi2-continuous-transferGaussian}
  \lean{Pphi2.continuous_transferGaussian}\leanok
  \uses{Pphi2-SpatialField, Pphi2-timeCoupling, Pphi2-transferGaussian}
  \texttt{Pphi2.continuous\_transferGaussian}
\end{theorem}
\begin{theorem}[\texttt{transferGaussian\_pos}]\label{Pphi2-transferGaussian-pos}
  \lean{Pphi2.transferGaussian_pos}\leanok
  \uses{Pphi2-SpatialField, Pphi2-timeCoupling, Pphi2-transferGaussian}
  \texttt{Pphi2.transferGaussian\_pos}
\end{theorem}
\begin{theorem}[\texttt{transferGaussian\_norm\_le\_one}]\label{Pphi2-transferGaussian-norm-le-one}
  \lean{Pphi2.transferGaussian_norm_le_one}\leanok
  \uses{Pphi2-timeCoupling-nonneg, Pphi2-SpatialField, Pphi2-transferGaussian-pos, Pphi2-timeCoupling, Pphi2-transferGaussian}
  \texttt{Pphi2.transferGaussian\_norm\_le\_one}
\end{theorem}
\begin{theorem}[\texttt{transferOperator\_isCompact}]\label{Pphi2-transferOperator-isCompact}
  \lean{Pphi2.transferOperator_isCompact}\leanok
  \uses{Pphi2-L2SpatialField, Pphi2-hilbert-schmidt-isCompact, Pphi2-transferGaussian-norm-le-one, Pphi2-SpatialField, Pphi2-transferWeight, Pphi2-transferGaussian-memLp, Pphi2-transferWeight-measurable, Pphi2-transferOperatorCLM, Pphi2-transferWeight-memLp-two, Pphi2-transferGaussian, Pphi2-transferWeight-bound}
  $T$ is compact on $L^2$. Proved from \texttt{integral\_operator\_l2\_kernel\_compact} via the tensor kernel $K(x,t) = w(x) G(t)$. Fully proved.
\end{theorem}
\begin{theorem}[\texttt{transferWeight\_pos}]\label{Pphi2-transferWeight-pos}
  \lean{Pphi2.transferWeight_pos}\leanok
  \uses{Pphi2-spatialAction, Pphi2-SpatialField, Pphi2-transferWeight, Pphi2-wickConstant}
  \texttt{Pphi2.transferWeight\_pos}
\end{theorem}
\begin{theorem}[\texttt{transferOperator\_spectral}]\label{Pphi2-transferOperator-spectral}
  \lean{Pphi2.transferOperator_spectral}\leanok
  \uses{Pphi2-L2SpatialField, Pphi2-transferOperator-isCompact, Pphi2-SpatialField, Pphi2-transferOperator-isSelfAdjoint, GaussianField-compact-selfAdjoint-spectral, Pphi2-transferOperatorCLM}
  Spectral decomposition: $\exists$ Hilbert basis $\{e_i\}$ and eigenvalues $\{\lambda_i\}$ with $Te_i = \lambda_i e_i$. From \texttt{compact\_selfAdjoint\_spectral}. Fully proved.
\end{theorem}
\begin{theorem}[\texttt{transferWeight\_gaussian\_decay}]\label{Pphi2-transferWeight-gaussian-decay}
  \lean{Pphi2.transferWeight_gaussian_decay}\leanok
  \uses{Pphi2-spatialAction, Pphi2-SpatialField, Pphi2-transferWeight, Pphi2-wickConstant}
  \texttt{Pphi2.transferWeight\_gaussian\_decay}
\end{theorem}
\begin{definition}[\texttt{transferWeight}]\label{Pphi2-transferWeight}
  \lean{Pphi2.transferWeight}\leanok
  \uses{Pphi2-spatialAction, Pphi2-SpatialField, Pphi2-wickConstant}
  ``\texttt{lean def transferWeight (P : InteractionPolynomial) (a mass : R) : SpatialField Ns -> R }`` Weight function $w(\psi) = \exp(-(a/2) \cdot h_a(\psi))$.
\end{definition}
\begin{theorem}[\texttt{transferWeight\_memLp\_two}]\label{Pphi2-transferWeight-memLp-two}
  \lean{Pphi2.transferWeight_memLp_two}\leanok
  \uses{Pphi2-transferWeight-gaussian-decay, Pphi2-SpatialField, Pphi2-transferWeight, Pphi2-transferWeight-measurable, Pphi2-transferWeight-pos, Lattice-toSemilatticeInf}
  \texttt{Pphi2.transferWeight\_memLp\_two}
\end{theorem}
\begin{theorem}[\texttt{continuous\_transferWeight}]\label{Pphi2-continuous-transferWeight}
  \lean{Pphi2.continuous_transferWeight}\leanok
  \uses{Pphi2-spatialAction, Pphi2-SpatialField, Pphi2-transferWeight, Pphi2-wickConstant}
  \texttt{Pphi2.continuous\_transferWeight}
\end{theorem}
\begin{definition}[\texttt{transferOperatorCLM}]\label{Pphi2-transferOperatorCLM}
  \lean{Pphi2.transferOperatorCLM}\leanok
  \uses{Pphi2-L2SpatialField, Pphi2-SpatialField, Pphi2-transferWeight, Pphi2-transferGaussian-memLp, Pphi2-transferWeight-measurable, Pphi2-transferGaussian, Pphi2-transferWeight-bound}
  ``\texttt{lean def transferOperatorCLM (P : InteractionPolynomial) (a mass : R) (ha : 0 < a) (hmass : 0 < mass) : L2SpatialField Ns ->L[R] L2SpatialField Ns }`` The transfer matrix $T = M_w \circ \mathrm{Conv}_G \circ M_w$ as a CLM on $L^2$.
\end{definition}
\begin{definition}[\texttt{L2SpatialField}]\label{Pphi2-L2SpatialField}
  \lean{Pphi2.L2SpatialField}\leanok
  \uses{Pphi2-SpatialField}
  ``\texttt{lean abbrev L2SpatialField := MeasureTheory.Lp (alpha := SpatialField Ns) R 2 volume }`` The Hilbert space $L^2(\mathbb{R}^{N_s})$ on which the transfer operator acts.
\end{definition}
\begin{definition}[\texttt{transferGaussian}]\label{Pphi2-transferGaussian}
  \lean{Pphi2.transferGaussian}\leanok
  \uses{Pphi2-SpatialField, Pphi2-timeCoupling}
  ``\texttt{lean def transferGaussian : SpatialField Ns -> R }`` Gaussian kernel $G(\psi) = \exp(-\frac{1}{2}\|\psi\|^2)$.  \#\#\# Key proved theorems: - \texttt{transferWeight\_measurable} / \texttt{transferWeight\_bound}: $w$ is measurable and essentially bounded. - \texttt{transferWeight\_gaussian\_decay}: $w(\psi) \le C \exp(-\alpha \sum \psi(x)^2)$ (Gaussian decay). - \texttt{transferGaussian\_memLp}: $G \in L^1$ via $\int e^{-\frac{1}{2}\|x\|^2}\,dx = (2\pi)^{N_s/2}$. - \texttt{transferGaussian\_memLp\_two}: $G \in L^2$ (product of 1D Gaussians). - \texttt{transferWeight\_memLp\_two}: $w \in L^2$ (from Gaussian decay domination).
\end{definition}
\begin{theorem}[\texttt{transferWeight\_bound}]\label{Pphi2-transferWeight-bound}
  \lean{Pphi2.transferWeight_bound}\leanok
  \uses{Pphi2-spatialAction, Pphi2-SpatialField, Pphi2-transferWeight, Pphi2-wickConstant}
  \texttt{Pphi2.transferWeight\_bound}
\end{theorem}
\begin{theorem}[\texttt{transferGaussian\_sub\_comm}]\label{Pphi2-transferGaussian-sub-comm}
  \lean{Pphi2.transferGaussian_sub_comm}\leanok
  \uses{Pphi2-SpatialField, Pphi2-transferGaussian, Pphi2-transferGaussian-neg}
  \texttt{Pphi2.transferGaussian\_sub\_comm}
\end{theorem}
\begin{theorem}[\texttt{hilbert\_schmidt\_isCompact}]\label{Pphi2-hilbert-schmidt-isCompact}
  \lean{Pphi2.hilbert_schmidt_isCompact}\leanok
  \uses{Pphi2-GeneralResults-integral-operator-l2-kernel-compact}
  \texttt{Pphi2.hilbert\_schmidt\_isCompact}
\end{theorem}
\begin{theorem}[\texttt{transferWeight\_measurable}]\label{Pphi2-transferWeight-measurable}
  \lean{Pphi2.transferWeight_measurable}\leanok
  \uses{Pphi2-SpatialField, Pphi2-transferWeight, Pphi2-continuous-transferWeight}
  \texttt{Pphi2.transferWeight\_measurable}
\end{theorem}
\begin{theorem}[\texttt{transferGaussian\_memLp\_two}]\label{Pphi2-transferGaussian-memLp-two}
  \lean{Pphi2.transferGaussian_memLp_two}\leanok
  \uses{Pphi2-SpatialField, Pphi2-transferGaussian-memLp, Pphi2-transferGaussian}
  \texttt{Pphi2.transferGaussian\_memLp\_two}
\end{theorem}
\begin{theorem}[\texttt{transferGaussian\_neg}]\label{Pphi2-transferGaussian-neg}
  \lean{Pphi2.transferGaussian_neg}\leanok
  \uses{Pphi2-SpatialField, Pphi2-timeCoupling, Pphi2-transferGaussian}
  \texttt{Pphi2.transferGaussian\_neg}
\end{theorem}
\section{\texttt{Pphi2.FormulationAdapter}}
\begin{definition}[\texttt{pphi2MeasureModel}]\label{Pphi2-pphi2MeasureModel}
  \lean{Pphi2.pphi2MeasureModel}\leanok
  \uses{Pphi2-EuclideanPlaneBackground-SpaceTime, Pphi2-FieldConfig2, Pphi2-generatingFunctional, Pphi2-EuclideanPlaneBackground-RealTestFunction, Pphi2-planeBackground, Pphi2-distribPairing, Common-QFT-Euclidean-MeasureModel, Pphi2-EuclideanPlaneBackground-dim, Pphi2-pphi2Formulation, GaussianField-instMeasurableSpaceConfiguration, Pphi2-generatingFunctional-}
  \texttt{Pphi2.pphi2MeasureModel}
\end{definition}
\begin{definition}[\texttt{pphi2PairingMeasureModel}]\label{Pphi2-pphi2PairingMeasureModel}
  \lean{Pphi2.pphi2PairingMeasureModel}\leanok
  \uses{Pphi2-EuclideanPlaneBackground-SpaceTime, Pphi2-FieldConfig2, Common-QFT-Euclidean-MeasureModel-toPairingMeasureModel, Pphi2-EuclideanPlaneBackground-RealTestFunction, Pphi2-planeBackground, Common-QFT-Euclidean-PairingMeasureModel, Pphi2-EuclideanPlaneBackground-dim, Pphi2-pphi2Formulation, GaussianField-instMeasurableSpaceConfiguration, Pphi2-pphi2MeasureModel}
  \texttt{Pphi2.pphi2PairingMeasureModel}
\end{definition}
\begin{definition}[\texttt{pphi2TensorSchwingerModel}]\label{Pphi2-pphi2TensorSchwingerModel}
  \lean{Pphi2.pphi2TensorSchwingerModel}\leanok
  \uses{Pphi2-EuclideanPlaneBackground-SpaceTime, Pphi2-FieldConfig2, Pphi2-EuclideanPlaneBackground-RealTestFunction, Pphi2-planeBackground, Common-QFT-Euclidean-Formulation-TestFunction, Common-QFT-Euclidean-TensorSchwingerModel, Pphi2-EuclideanPlaneBackground-dim, Pphi2-pphi2Formulation, GaussianField-instMeasurableSpaceConfiguration}
  \texttt{Pphi2.pphi2TensorSchwingerModel}
\end{definition}
\begin{definition}[\texttt{pphi2MeasureToTensorBridge}]\label{Pphi2-pphi2MeasureToTensorBridge}
  \lean{Pphi2.pphi2MeasureToTensorBridge}\leanok
  \uses{Pphi2-EuclideanPlaneBackground-SpaceTime, Pphi2-FieldConfig2, Pphi2-EuclideanPlaneBackground-RealTestFunction, Pphi2-planeBackground, Pphi2-EuclideanPlaneBackground-dim, Pphi2-pphi2Formulation, Pphi2-pphi2PairingMeasureModel, Common-QFT-Euclidean-MeasureToTensorBridge, GaussianField-instMeasurableSpaceConfiguration, Pphi2-pphi2TensorSchwingerModel}
  \texttt{Pphi2.pphi2MeasureToTensorBridge}
\end{definition}
\begin{definition}[\texttt{pphi2Formulation}]\label{Pphi2-pphi2Formulation}
  \lean{Pphi2.pphi2Formulation}\leanok
  \uses{Pphi2-EuclideanPlaneBackground-SpaceTime, Pphi2-planeBackground, Pphi2-EuclideanPlaneBackground-dim, Pphi2-SpaceTime2, Common-QFT-Euclidean-Formulation, Pphi2-TestFunction2-, Pphi2-positiveTimeSubmodule2, Pphi2-TestFunction2}
  \texttt{Pphi2.pphi2Formulation}
\end{definition}
\section{\texttt{Pphi2.AsymTorus.AsymMeasureFactorization}}
\begin{theorem}[\texttt{boltzmannWeightAsym\_eq\_comp\_evalMapAsym}]\label{Pphi2-boltzmannWeightAsym-eq-comp-evalMapAsym}
  \lean{Pphi2.boltzmannWeightAsym_eq_comp_evalMapAsym}\leanok
  \uses{GaussianField-AsymLatticeSites, Pphi2-wickConstantAsym, Pphi2-boltzmannWeightAsym, GaussianField-Configuration, GaussianField-AsymLatticeField, GaussianField-evalMapAsym, GaussianField-instFintypeAsymLatticeSitesOfNeZeroNat, Pphi2-wickPolynomial}
  \texttt{Pphi2.boltzmannWeightAsym\_eq\_comp\_evalMapAsym}
\end{theorem}
\begin{theorem}[\texttt{measurePreserving\_asymSliceEquiv}]\label{Pphi2-measurePreserving-asymSliceEquiv}
  \lean{Pphi2.measurePreserving_asymSliceEquiv}\leanok
  \uses{Pphi2-measurePreserving-sliceReindex, GaussianField-AsymLatticeSites, Pphi2-asymSliceMeasurableEquiv, Pphi2-sliceReindexMeasurableEquiv, Pphi2-SpatialField, Pphi2-asymSliceEquiv, Pphi2-measurePreserving-curry, GaussianField-AsymLatticeField, GaussianField-instFintypeAsymLatticeSitesOfNeZeroNat, Pphi2-asymSliceMeasurableEquiv-coe}
  \texttt{Pphi2.measurePreserving\_asymSliceEquiv}
\end{theorem}
\begin{theorem}[\texttt{asymSliceMeasurableEquiv\_coe}]\label{Pphi2-asymSliceMeasurableEquiv-coe}
  \lean{Pphi2.asymSliceMeasurableEquiv_coe}\leanok
  \uses{GaussianField-AsymLatticeSites, Pphi2-asymSliceMeasurableEquiv, Pphi2-SpatialField, Pphi2-asymSliceEquiv, Pphi2-asymSliceEquiv-apply, GaussianField-AsymLatticeField}
  \texttt{Pphi2.asymSliceMeasurableEquiv\_coe}
\end{theorem}
\begin{theorem}[\texttt{interactionFunctionalAsym\_evalMapAsymInv}]\label{Pphi2-interactionFunctionalAsym-evalMapAsymInv}
  \lean{Pphi2.interactionFunctionalAsym_evalMapAsymInv}\leanok
  \uses{GaussianField-AsymLatticeSites, Pphi2-wickConstantAsym, GaussianField-Configuration, GaussianField-AsymLatticeField, GaussianField-evalMapAsym, Pphi2-interactionFunctionalAsym, GaussianField-instFintypeAsymLatticeSitesOfNeZeroNat, GaussianField-evalMap-evalMapInvAsym, GaussianField-evalMapAsymInv, GaussianField-asymLatticeDelta, Pphi2-wickPolynomial}
  \texttt{Pphi2.interactionFunctionalAsym\_evalMapAsymInv}
\end{theorem}
\begin{theorem}[\texttt{measurePreserving\_curry}]\label{Pphi2-measurePreserving-curry}
  \lean{Pphi2.measurePreserving_curry}\leanok
  \texttt{Pphi2.measurePreserving\_curry}
\end{theorem}
\begin{theorem}[\texttt{withDensity\_comp\_map\_measurableEquiv}]\label{Pphi2-withDensity-comp-map-measurableEquiv}
  \lean{Pphi2.withDensity_comp_map_measurableEquiv}\leanok
  \texttt{Pphi2.withDensity\_comp\_map\_measurableEquiv}
\end{theorem}
\begin{theorem}[\texttt{measurePreserving\_sliceReindex}]\label{Pphi2-measurePreserving-sliceReindex}
  \lean{Pphi2.measurePreserving_sliceReindex}\leanok
  \uses{Pphi2-sliceReindexMeasurableEquiv, Pphi2-SpatialField}
  \texttt{Pphi2.measurePreserving\_sliceReindex}
\end{theorem}
\begin{theorem}[\texttt{gaussianDensity\_mul\_boltzmann\_eq\_pathDensity}]\label{Pphi2-gaussianDensity-mul-boltzmann-eq-pathDensity}
  \lean{Pphi2.gaussianDensity_mul_boltzmann_eq_pathDensity}\leanok
  \uses{GaussianField-AsymLatticeSites, Pphi2-wickConstantAsym, Pphi2-SpatialField, Pphi2-asymTransferKernel, Pphi2-asymSliceEquiv, Pphi2-asymTransferSystem, GaussianField-gaussianDensityAsym, GaussianField-AsymLatticeField, Pphi2-periodicPathDensity-asymTransferKernel-eq-exp, GaussianField-instFintypeAsymLatticeSitesOfNeZeroNat, GaussianField-massOperatorAsym, Pphi2-wickPolynomial}
  \texttt{Pphi2.gaussianDensity\_mul\_boltzmann\_eq\_pathDensity}
\end{theorem}
\begin{definition}[\texttt{sliceReindexMeasurableEquiv}]\label{Pphi2-sliceReindexMeasurableEquiv}
  \lean{Pphi2.sliceReindexMeasurableEquiv}\leanok
  \uses{Pphi2-SpatialField}
  \texttt{Pphi2.sliceReindexMeasurableEquiv}
\end{definition}
\begin{theorem}[\texttt{interactingLatticeMeasureAsym\_slice\_pushforward\_eq\_pathMeasure}]\label{Pphi2-interactingLatticeMeasureAsym-slice-pushforward-eq-pathMeasure}
  \lean{Pphi2.interactingLatticeMeasureAsym_slice_pushforward_eq_pathMeasure}\leanok
  \uses{Pphi2-withDensity-comp-map-measurableEquiv, GaussianField-normalizedQuadraticGaussianMeasure, GaussianField-AsymLatticeSites, GaussianField-measure, GaussianField-ell2-separableSpace, Pphi2-wickConstantAsym, GaussianField-gaussianDensityWeightAsym, GaussianField-latticeGaussianFieldLawAsym-eq-normalizedQuadraticGaussianMeasure, GaussianField-gaussianDensityNormConstAsym, Pphi2-asymSliceMeasurableEquiv, GaussianField-normalizedGaussianDensityMeasureAsym, GaussianField-asymLatticeField-dyninMityaginSpace, GaussianField-measurable-evalMapAsym, Pphi2-SpatialField, Pphi2-asymSliceEquiv, GaussianField-evalMapAsymMeasurableEquiv, Pphi2-boltzmannWeightAsym, Pphi2-interactingLatticeMeasureAsym-isProbability, Pphi2-asymTransferSystem, GaussianField-latticeGaussianFieldLawAsym, GaussianField-gaussianDensityAsym, GaussianField-gaussianDensityAsym-measurable, GaussianField-massEigenvaluesAsym, GaussianField-normalizedGaussianDensityMeasureAsym-eq-normalizedQuadraticGaussianMeasure, Pphi2-partitionFunctionAsym, Pphi2-wickPolynomial-continuous-, GaussianField-Configuration, GaussianField-AsymLatticeField, Pphi2-latticeGaussianMeasureAsym, GaussianField-evalMapAsym, Pphi2-interactingLatticeMeasureAsym, Pphi2-measurePreserving-asymSliceEquiv, GaussianField-instMeasurableSpaceConfiguration, GaussianField-instFintypeAsymLatticeSitesOfNeZeroNat, GaussianField-massOperatorAsym, GaussianField-latticeCovarianceAsymGJ, GaussianField-gaussianDensityAsym-nonneg, Pphi2-boltzmannWeightAsym-eq-comp-evalMapAsym, Pphi2-asymSliceMeasurableEquiv-coe, GaussianField-ell2-, Pphi2-wickPolynomial, Pphi2-gaussianDensity-mul-boltzmann-eq-pathDensity}
  \texttt{Pphi2.interactingLatticeMeasureAsym\_slice\_pushforward\_eq\_pathMeasure}
\end{theorem}
\begin{definition}[\texttt{asymSliceMeasurableEquiv}]\label{Pphi2-asymSliceMeasurableEquiv}
  \lean{Pphi2.asymSliceMeasurableEquiv}\leanok
  \uses{GaussianField-AsymLatticeSites, Pphi2-sliceReindexMeasurableEquiv, Pphi2-SpatialField, GaussianField-AsymLatticeField}
  \texttt{Pphi2.asymSliceMeasurableEquiv}
\end{definition}
\section{\texttt{Pphi2.NelsonEstimate.AsymRoughErrorChaosStd}}
\begin{theorem}[\texttt{asymCanonicalRoughError\_neg\_tail\_uniform}]\label{Pphi2-asymCanonicalRoughError-neg-tail-uniform}
  \lean{Pphi2.asymCanonicalRoughError_neg_tail_uniform}\leanok
  \uses{Pphi2-asymCanonicalRoughErrorStd-eq, Pphi2-asymCanonicalRoughErrorStd, Pphi2-asymCanonicalRoughError-neg-tail-of-stdGaussian-explicit-ae, Pphi2-AsymCanonicalJointSumIndex, Pphi2-AsymCanonicalJoint, Pphi2-asymCanonicalRoughError, GaussianHilbert-stdGaussianFin, Pphi2-asymCanonicalJointStdGaussianMeasurableEquiv, GaussianHilbert-wienerChaosLE, Pphi2-asymRoughError-variance, Pphi2-asymCanonicalJointMeasure-map-stdGaussian, Pphi2-ChaosTailBridge-chaosTailConstant, Pphi2-asymCanonicalJointMeasure, Lattice-toSemilatticeInf, Pphi2-asymCanonicalRoughErrorStd-mem-wienerChaosLE, Pphi2-AsymCanonicalMode}
  \texttt{Pphi2.asymCanonicalRoughError\_neg\_tail\_uniform}
\end{theorem}
\begin{definition}[\texttt{asymCanonicalCrossTermStd}]\label{Pphi2-asymCanonicalCrossTermStd}
  \lean{Pphi2.asymCanonicalCrossTermStd}\leanok
  \uses{GaussianField-AsymLatticeSites, Pphi2-AsymCanonicalJointSumIndex, GaussianField-instFintypeAsymLatticeSitesOfNeZeroNat, Pphi2-AsymCanonicalMode}
  \texttt{Pphi2.asymCanonicalCrossTermStd}
\end{definition}
\begin{definition}[\texttt{asymCanonicalRoughErrorStd}]\label{Pphi2-asymCanonicalRoughErrorStd}
  \lean{Pphi2.asymCanonicalRoughErrorStd}\leanok
  \uses{Pphi2-AsymCanonicalJointSumIndex, Pphi2-asymCanonicalCrossTermStd, Pphi2-AsymCanonicalMode}
  \texttt{Pphi2.asymCanonicalRoughErrorStd}
\end{definition}
\begin{theorem}[\texttt{congr\_simp}]\label{Pphi2-asymCanonicalRoughErrorStd-congr-simp}
  \lean{Pphi2.asymCanonicalRoughErrorStd.congr_simp}\leanok
  \uses{Pphi2-asymCanonicalRoughErrorStd, Pphi2-AsymCanonicalJointSumIndex, Pphi2-AsymCanonicalMode}
  \texttt{Pphi2.asymCanonicalRoughErrorStd.congr\_simp}
\end{theorem}
\begin{theorem}[\texttt{asymCanonicalRoughError\_neg\_tail\_of\_stdGaussian\_explicit\_ae}]\label{Pphi2-asymCanonicalRoughError-neg-tail-of-stdGaussian-explicit-ae}
  \lean{Pphi2.asymCanonicalRoughError_neg_tail_of_stdGaussian_explicit_ae}\leanok
  \uses{Pphi2-AsymCanonicalJointSumIndex, Pphi2-AsymCanonicalJoint, Pphi2-asymCanonicalRoughError, GaussianHilbert-stdGaussianFin, Pphi2-asymCanonicalJointStdGaussianMeasurableEquiv, GaussianHilbert-wienerChaosLE, Pphi2-ChaosTailBridge-chaos-neg-tail-bound-explicit, Pphi2-asymCanonicalJointMeasure-map-stdGaussian, Pphi2-ChaosTailBridge-chaosTailConstant, Pphi2-asymCanonicalJointMeasure, Pphi2-AsymCanonicalMode}
  \texttt{Pphi2.asymCanonicalRoughError\_neg\_tail\_of\_stdGaussian\_explicit\_ae}
\end{theorem}
\begin{theorem}[\texttt{asymCanonicalCrossTermStd\_eq}]\label{Pphi2-asymCanonicalCrossTermStd-eq}
  \lean{Pphi2.asymCanonicalCrossTermStd_eq}\leanok
  \uses{GaussianField-AsymLatticeSites, Pphi2-asymCanonicalRoughWickConstant, Pphi2-AsymCanonicalJointSumIndex, Pphi2-asymCanonicalCrossTermStd, Pphi2-asymCanonicalCrossTerm, Pphi2-AsymCanonicalJoint, Pphi2-asymCanonicalRoughFieldFunction-congr-simp, Pphi2-wickMonomial, Pphi2-asymCanonicalJointStdGaussianMeasurableEquiv, Pphi2-asymCanonicalRoughFieldFunction, Pphi2-asymSmoothWickConstant, GaussianField-instFintypeAsymLatticeSitesOfNeZeroNat, Pphi2-asymCanonicalSmoothFieldFunction, Pphi2-asymCanonicalSmoothFieldFunction-congr-simp, Pphi2-AsymCanonicalMode}
  \texttt{Pphi2.asymCanonicalCrossTermStd\_eq}
\end{theorem}
\begin{theorem}[\texttt{asymCanonicalCrossTermLinearCombo\_mem\_wienerChaosLE}]\label{Pphi2-asymCanonicalCrossTermLinearCombo-mem-wienerChaosLE}
  \lean{Pphi2.asymCanonicalCrossTermLinearCombo_mem_wienerChaosLE}\leanok
  \uses{Pphi2-AsymCanonicalJointSumIndex, Pphi2-asymCanonicalCrossTermStd, Pphi2-wienerChaosLE-mono, GaussianHilbert-stdGaussianFin, GaussianHilbert-wienerChaosLE, Pphi2-asymCanonicalCrossTermStd-mem-wienerChaosLE, Pphi2-AsymCanonicalMode}
  \texttt{Pphi2.asymCanonicalCrossTermLinearCombo\_mem\_wienerChaosLE}
\end{theorem}
\begin{theorem}[\texttt{asymCanonicalJointMeasure\_map\_sum\_pi\_gaussian}]\label{Pphi2-asymCanonicalJointMeasure-map-sum-pi-gaussian}
  \lean{Pphi2.asymCanonicalJointMeasure_map_sum_pi_gaussian}\leanok
  \uses{Pphi2-AsymCanonicalJointSumIndex, Pphi2-asymCanonicalJointMeasure, Pphi2-AsymCanonicalMode}
  \texttt{Pphi2.asymCanonicalJointMeasure\_map\_sum\_pi\_gaussian}
\end{theorem}
\begin{theorem}[\texttt{congr\_simp}]\label{Pphi2-asymCanonicalJointStdGaussianMeasurableEquiv-congr-simp}
  \lean{Pphi2.asymCanonicalJointStdGaussianMeasurableEquiv.congr_simp}\leanok
  \uses{Pphi2-AsymCanonicalJointSumIndex, Pphi2-AsymCanonicalJoint, Pphi2-asymCanonicalJointStdGaussianMeasurableEquiv, Pphi2-AsymCanonicalMode}
  \texttt{Pphi2.asymCanonicalJointStdGaussianMeasurableEquiv.congr\_simp}
\end{theorem}
\begin{definition}[\texttt{asymCanonicalJointStdGaussianMeasurableEquiv}]\label{Pphi2-asymCanonicalJointStdGaussianMeasurableEquiv}
  \lean{Pphi2.asymCanonicalJointStdGaussianMeasurableEquiv}\leanok
  \uses{Pphi2-AsymCanonicalJointSumIndex, Pphi2-AsymCanonicalJoint, Pphi2-AsymCanonicalMode}
  \texttt{Pphi2.asymCanonicalJointStdGaussianMeasurableEquiv}
\end{definition}
\begin{theorem}[\texttt{congr\_simp}]\label{Pphi2-asymCanonicalRoughError-congr-simp}
  \lean{Pphi2.asymCanonicalRoughError.congr_simp}\leanok
  \uses{Pphi2-AsymCanonicalJoint, Pphi2-asymCanonicalRoughError}
  \texttt{Pphi2.asymCanonicalRoughError.congr\_simp}
\end{theorem}
\begin{theorem}[\texttt{congr\_simp}]\label{Pphi2-asymCanonicalCrossTermStd-congr-simp}
  \lean{Pphi2.asymCanonicalCrossTermStd.congr_simp}\leanok
  \uses{Pphi2-AsymCanonicalJointSumIndex, Pphi2-asymCanonicalCrossTermStd, Pphi2-AsymCanonicalMode}
  \texttt{Pphi2.asymCanonicalCrossTermStd.congr\_simp}
\end{theorem}
\begin{theorem}[\texttt{asymCanonicalRoughErrorStd\_mem\_wienerChaosLE}]\label{Pphi2-asymCanonicalRoughErrorStd-mem-wienerChaosLE}
  \lean{Pphi2.asymCanonicalRoughErrorStd_mem_wienerChaosLE}\leanok
  \uses{Pphi2-asymCanonicalRoughErrorStd, Pphi2-asymCanonicalCrossTermLinearCombo-mem-wienerChaosLE, Pphi2-AsymCanonicalJointSumIndex, Pphi2-asymCanonicalCrossTermStd, GaussianHilbert-stdGaussianFin, GaussianHilbert-wienerChaosLE, Pphi2-AsymCanonicalMode}
  \texttt{Pphi2.asymCanonicalRoughErrorStd\_mem\_wienerChaosLE}
\end{theorem}
\begin{theorem}[\texttt{asymCanonicalJointMeasure\_map\_stdGaussian}]\label{Pphi2-asymCanonicalJointMeasure-map-stdGaussian}
  \lean{Pphi2.asymCanonicalJointMeasure_map_stdGaussian}\leanok
  \uses{Pphi2-AsymCanonicalJointSumIndex, Pphi2-AsymCanonicalJoint, Pphi2-asymCanonicalJointMeasure-map-sum-pi-gaussian, GaussianHilbert-stdGaussianFin, Pphi2-asymCanonicalJointStdGaussianMeasurableEquiv, Pphi2-asymCanonicalJointMeasure, Pphi2-AsymCanonicalMode}
  \texttt{Pphi2.asymCanonicalJointMeasure\_map\_stdGaussian}
\end{theorem}
\begin{theorem}[\texttt{asymCanonicalRoughErrorStd\_eq}]\label{Pphi2-asymCanonicalRoughErrorStd-eq}
  \lean{Pphi2.asymCanonicalRoughErrorStd_eq}\leanok
  \uses{Pphi2-asymCanonicalRoughErrorStd, Pphi2-asymCanonicalRoughError-eq-sum-over-cross-terms, Pphi2-asymCanonicalCrossTermStd-eq, Pphi2-AsymCanonicalJointSumIndex, Pphi2-asymCanonicalCrossTermStd, Pphi2-asymCanonicalCrossTerm, Pphi2-AsymCanonicalJoint, Pphi2-asymCanonicalRoughError, Pphi2-asymCanonicalJointStdGaussianMeasurableEquiv, Pphi2-AsymCanonicalMode}
  \texttt{Pphi2.asymCanonicalRoughErrorStd\_eq}
\end{theorem}
\begin{theorem}[\texttt{asymCanonicalCrossTermStd\_mem\_wienerChaosLE}]\label{Pphi2-asymCanonicalCrossTermStd-mem-wienerChaosLE}
  \lean{Pphi2.asymCanonicalCrossTermStd_mem_wienerChaosLE}\leanok
  \uses{GaussianField-AsymLatticeSites, Pphi2-AsymCanonicalJointSumIndex, Pphi2-asymCanonicalCrossTermStd, GaussianHilbert-stdGaussianFin, GaussianHilbert-wienerChaosLE, GaussianField-instFintypeAsymLatticeSitesOfNeZeroNat, Pphi2-AsymCanonicalMode}
  \texttt{Pphi2.asymCanonicalCrossTermStd\_mem\_wienerChaosLE}
\end{theorem}
\section{\texttt{Pphi2.OSProofs.OS4\_Ergodicity}}
\begin{theorem}[\texttt{clustering\_implies\_ergodicity}]\label{Pphi2-clustering-implies-ergodicity}
  \lean{Pphi2.clustering_implies_ergodicity}\leanok
  \uses{Lattice-toSemilatticeInf}
  If $\mu$ satisfies exponential clustering w.r.t. translations $T_R$, then $\mu$ is ergodic: $T_R$-invariant sets have measure $0$ or $1$. Proved by setting $F = G = \mathbf{1}_A$, computing $|\mu(A) - \mu(A)^2| \le C e^{-mR} \to 0$, then $\mu(A)(1-\mu(A)) = 0$. Fully proved.
\end{theorem}
\begin{theorem}[\texttt{os4\_lattice\_from\_gap}]\label{Pphi2-os4-lattice-from-gap}
  \lean{Pphi2.os4_lattice_from_gap}\leanok
  \uses{Pphi2-massGap-pos, Pphi2-massGap}
  Same as \texttt{os4\_lattice}, stated to make the connection to the OS axiom framework explicit.
\end{theorem}
\begin{theorem}[\texttt{unique\_vacuum}]\label{Pphi2-unique-vacuum}
  \lean{Pphi2.unique_vacuum}\leanok
  \uses{Pphi2-L2SpatialField, Pphi2-SpatialField, Pphi2-transferOperatorCLM, Pphi2-transferOperator-ground-simple-spectral}
  On the lattice: uniqueness of the ground state from \texttt{transferOperator\_ground\_simple\_spectral}. Fully proved.
\end{theorem}
\begin{theorem}[\texttt{os4\_lattice}]\label{Pphi2-os4-lattice}
  \lean{Pphi2.os4_lattice}\leanok
  \uses{Pphi2-massGap-pos, Pphi2-massGap}
  The lattice interacting measure satisfies OS4: $0 < m_{\mathrm{phys}}$. Delegates to \texttt{massGap\_pos}. Fully proved.
\end{theorem}
\begin{theorem}[\texttt{exponential\_mixing}]\label{Pphi2-exponential-mixing}
  \lean{Pphi2.exponential_mixing}\leanok
  \uses{Pphi2-massGap-pos, Pphi2-massGap}
  $\exists m > 0$ with $m \le m_{\mathrm{phys}}$, the measure is exponentially mixing. Immediate from \texttt{massGap\_pos}. Fully proved.
\end{theorem}
\section{\texttt{Pphi2.EuclideanComplex}}
\begin{theorem}[\texttt{schwartzOfReal\_real\_smul}]\label{Pphi2-schwartzOfReal-real-smul}
  \lean{Pphi2.schwartzOfReal_real_smul}\leanok
  \uses{Pphi2-EuclideanPlaneBackground-SpaceTime, Pphi2-EuclideanPlaneBackground-RealTestFunction, Pphi2-schwartzOfReal, Pphi2-EuclideanPlaneBackground-dim, Pphi2-EuclideanPlaneBackground, Pphi2-EuclideanPlaneBackground-ComplexTestFunction}
  \texttt{Pphi2.schwartzOfReal\_real\_smul}
\end{theorem}
\begin{theorem}[\texttt{schwartzOfReal\_add}]\label{Pphi2-schwartzOfReal-add}
  \lean{Pphi2.schwartzOfReal_add}\leanok
  \uses{Pphi2-EuclideanPlaneBackground-SpaceTime, Pphi2-EuclideanPlaneBackground-RealTestFunction, Pphi2-schwartzOfReal, Pphi2-EuclideanPlaneBackground-dim, Pphi2-EuclideanPlaneBackground, Pphi2-EuclideanPlaneBackground-ComplexTestFunction}
  \texttt{Pphi2.schwartzOfReal\_add}
\end{theorem}
\begin{theorem}[\texttt{schwartzRe\_sum\_complex\_smul}]\label{Pphi2-schwartzRe-sum-complex-smul}
  \lean{Pphi2.schwartzRe_sum_complex_smul}\leanok
  \uses{Pphi2-EuclideanPlaneBackground-SpaceTime, Pphi2-EuclideanPlaneBackground-RealTestFunction, Pphi2-EuclideanPlaneBackground-dim, Pphi2-EuclideanOS-schwartzIm, Pphi2-EuclideanOS-schwartzRe, Pphi2-EuclideanPlaneBackground, Pphi2-EuclideanPlaneBackground-ComplexTestFunction}
  \texttt{Pphi2.schwartzRe\_sum\_complex\_smul}
\end{theorem}
\begin{theorem}[\texttt{actionComplex\_schwartzOfReal}]\label{Pphi2-actionComplex-schwartzOfReal}
  \lean{Pphi2.actionComplex_schwartzOfReal}\leanok
  \uses{Pphi2-EuclideanPlaneBackground-SpaceTime, Pphi2-EuclideanPlaneBackground-action, Pphi2-EuclideanPlaneBackground-RealTestFunction, Pphi2-schwartzOfReal, Pphi2-EuclideanPlaneBackground-dim, Pphi2-EuclideanPlaneBackground-Motion, Pphi2-EuclideanPlaneBackground, Pphi2-EuclideanPlaneBackground-actionComplex, Pphi2-EuclideanPlaneBackground-ComplexTestFunction}
  \texttt{Pphi2.actionComplex\_schwartzOfReal}
\end{theorem}
\begin{theorem}[\texttt{schwartzIm\_schwartzOfReal}]\label{Pphi2-schwartzIm-schwartzOfReal}
  \lean{Pphi2.schwartzIm_schwartzOfReal}\leanok
  \uses{Pphi2-EuclideanPlaneBackground-SpaceTime, Pphi2-EuclideanPlaneBackground-RealTestFunction, Pphi2-schwartzOfReal, Pphi2-EuclideanPlaneBackground-dim, Pphi2-EuclideanOS-schwartzIm, Pphi2-EuclideanPlaneBackground}
  \texttt{Pphi2.schwartzIm\_schwartzOfReal}
\end{theorem}
\begin{theorem}[\texttt{schwartzOfReal\_apply}]\label{Pphi2-schwartzOfReal-apply}
  \lean{Pphi2.schwartzOfReal_apply}\leanok
  \uses{Pphi2-EuclideanPlaneBackground-SpaceTime, Pphi2-EuclideanPlaneBackground-RealTestFunction, Pphi2-schwartzOfReal, Pphi2-EuclideanPlaneBackground-dim, Pphi2-EuclideanPlaneBackground, Pphi2-EuclideanPlaneBackground-ComplexTestFunction}
  $(schwartzOfReal\; f)(x) = (f(x) : \mathbb{C})$.
\end{theorem}
\begin{theorem}[\texttt{generatingFunctionalℂ\_ofReal\_add\_real\_smul}]\label{Pphi2-generatingFunctional--ofReal-add-real-smul}
  \lean{Pphi2.generatingFunctionalℂ_ofReal_add_real_smul}\leanok
  \uses{Pphi2-EuclideanPlaneBackground-SpaceTime, Pphi2-EuclideanPlaneBackground-RealTestFunction, Pphi2-schwartzOfReal, Pphi2-EuclideanPlaneBackground-dim, Pphi2-EuclideanOS-schwartzIm, Pphi2-EuclideanOS-schwartzRe, Pphi2-schwartzRe-ofReal-add-real-smul, Pphi2-EuclideanPlaneBackground, Pphi2-EuclideanOS-generatingFunctional-, GaussianField-instMeasurableSpaceConfiguration, Pphi2-EuclideanOS-generatingFunctional, Pphi2-EuclideanPlaneBackground-ComplexTestFunction, Pphi2-schwartzIm-ofReal-add-real-smul, Pphi2-EuclideanPlaneBackground-Distribution}
  $Z_\mathbb{C}[\text{ofReal}(f) + r \cdot \text{ofReal}(h)] = Z[f + rh]$ (complex generating functional at real arguments equals real one).
\end{theorem}
\begin{theorem}[\texttt{schwartzRe\_actionComplex}]\label{Pphi2-schwartzRe-actionComplex}
  \lean{Pphi2.schwartzRe_actionComplex}\leanok
  \uses{Pphi2-EuclideanPlaneBackground-SpaceTime, Pphi2-EuclideanPlaneBackground-action, Pphi2-EuclideanPlaneBackground-RealTestFunction, Pphi2-EuclideanPlaneBackground-dim, Pphi2-EuclideanOS-schwartzRe, Pphi2-EuclideanPlaneBackground-Motion, Pphi2-EuclideanPlaneBackground, Pphi2-EuclideanPlaneBackground-actionComplex, Pphi2-EuclideanPlaneBackground-ComplexTestFunction}
  \texttt{Pphi2.schwartzRe\_actionComplex}
\end{theorem}
\begin{theorem}[\texttt{schwartz\_decompose\_actionComplex}]\label{Pphi2-schwartz-decompose-actionComplex}
  \lean{Pphi2.schwartz_decompose_actionComplex}\leanok
  \uses{Pphi2-EuclideanPlaneBackground-SpaceTime, Pphi2-EuclideanPlaneBackground-action, Pphi2-EuclideanPlaneBackground-RealTestFunction, Pphi2-schwartzOfReal, Pphi2-EuclideanPlaneBackground-dim, Pphi2-EuclideanOS-schwartzIm, Pphi2-schwartz-decompose, Pphi2-EuclideanOS-schwartzRe, Pphi2-EuclideanPlaneBackground-Motion, Pphi2-EuclideanPlaneBackground, Pphi2-schwartzRe-actionComplex, Pphi2-EuclideanPlaneBackground-actionComplex, Pphi2-schwartzIm-actionComplex, Pphi2-EuclideanPlaneBackground-ComplexTestFunction}
  \texttt{Pphi2.schwartz\_decompose\_actionComplex}
\end{theorem}
\begin{theorem}[\texttt{schwartzRe\_complex\_smul}]\label{Pphi2-schwartzRe-complex-smul}
  \lean{Pphi2.schwartzRe_complex_smul}\leanok
  \uses{Pphi2-EuclideanPlaneBackground-SpaceTime, Pphi2-EuclideanPlaneBackground-RealTestFunction, Pphi2-EuclideanPlaneBackground-dim, Pphi2-EuclideanOS-schwartzIm, Pphi2-EuclideanOS-schwartzRe, Pphi2-EuclideanPlaneBackground, Pphi2-EuclideanPlaneBackground-ComplexTestFunction}
  \texttt{Pphi2.schwartzRe\_complex\_smul}
\end{theorem}
\begin{theorem}[\texttt{schwartzIm\_complex\_smul}]\label{Pphi2-schwartzIm-complex-smul}
  \lean{Pphi2.schwartzIm_complex_smul}\leanok
  \uses{Pphi2-EuclideanPlaneBackground-SpaceTime, Pphi2-EuclideanPlaneBackground-RealTestFunction, Pphi2-EuclideanPlaneBackground-dim, Pphi2-EuclideanOS-schwartzIm, Pphi2-EuclideanOS-schwartzRe, Pphi2-EuclideanPlaneBackground, Pphi2-EuclideanPlaneBackground-ComplexTestFunction}
  \texttt{Pphi2.schwartzIm\_complex\_smul}
\end{theorem}
\begin{definition}[\texttt{schwartzOfReal}]\label{Pphi2-schwartzOfReal}
  \lean{Pphi2.schwartzOfReal}\leanok
  \uses{Pphi2-EuclideanPlaneBackground-SpaceTime, Pphi2-EuclideanPlaneBackground-RealTestFunction, Pphi2-EuclideanPlaneBackground-dim, Pphi2-EuclideanPlaneBackground, Pphi2-EuclideanPlaneBackground-ComplexTestFunction}
  ``\texttt{lean def schwartzOfReal (f : TestFunction2) : TestFunction2ℂ }`` Embeds a real Schwartz function as a complex-valued one via $\text{Complex.ofReal}$. Smoothness from \texttt{Complex.ofRealCLM.contDiff.comp}; decay from composing with a bounded linear map.
\end{definition}
\begin{theorem}[\texttt{schwartzIm\_sum\_complex\_smul}]\label{Pphi2-schwartzIm-sum-complex-smul}
  \lean{Pphi2.schwartzIm_sum_complex_smul}\leanok
  \uses{Pphi2-EuclideanPlaneBackground-SpaceTime, Pphi2-EuclideanPlaneBackground-RealTestFunction, Pphi2-EuclideanPlaneBackground-dim, Pphi2-EuclideanOS-schwartzIm, Pphi2-EuclideanOS-schwartzRe, Pphi2-EuclideanPlaneBackground, Pphi2-EuclideanPlaneBackground-ComplexTestFunction}
  \texttt{Pphi2.schwartzIm\_sum\_complex\_smul}
\end{theorem}
\begin{theorem}[\texttt{schwartzIm\_actionComplex}]\label{Pphi2-schwartzIm-actionComplex}
  \lean{Pphi2.schwartzIm_actionComplex}\leanok
  \uses{Pphi2-EuclideanPlaneBackground-SpaceTime, Pphi2-EuclideanPlaneBackground-action, Pphi2-EuclideanPlaneBackground-RealTestFunction, Pphi2-EuclideanPlaneBackground-dim, Pphi2-EuclideanOS-schwartzIm, Pphi2-EuclideanPlaneBackground-Motion, Pphi2-EuclideanPlaneBackground, Pphi2-EuclideanPlaneBackground-actionComplex, Pphi2-EuclideanPlaneBackground-ComplexTestFunction}
  \texttt{Pphi2.schwartzIm\_actionComplex}
\end{theorem}
\begin{theorem}[\texttt{schwartzRe\_ofReal\_add\_real\_smul}]\label{Pphi2-schwartzRe-ofReal-add-real-smul}
  \lean{Pphi2.schwartzRe_ofReal_add_real_smul}\leanok
  \uses{Pphi2-EuclideanPlaneBackground-SpaceTime, Pphi2-EuclideanPlaneBackground-RealTestFunction, Pphi2-schwartzOfReal, Pphi2-EuclideanPlaneBackground-dim, Pphi2-EuclideanOS-schwartzRe, Pphi2-EuclideanPlaneBackground, Pphi2-EuclideanPlaneBackground-ComplexTestFunction}
  \texttt{Pphi2.schwartzRe\_ofReal\_add\_real\_smul}
\end{theorem}
\begin{theorem}[\texttt{schwartz\_decompose\_continuumEuclideanActionComplex}]\label{Pphi2-schwartz-decompose-continuumEuclideanActionComplex}
  \lean{Pphi2.schwartz_decompose_continuumEuclideanActionComplex}\leanok
  \uses{Pphi2-EuclideanPlaneBackground-SpaceTime, Pphi2-schwartz-decompose-actionComplex, Pphi2-planeBackground, Pphi2-schwartzOfReal, Pphi2-EuclideanPlaneBackground-dim, Pphi2-continuumEuclideanActionComplex, Pphi2-EuclideanOS-schwartzIm, Pphi2-EuclideanOS-schwartzRe, Pphi2-ContinuumComplexTestFunction, Pphi2-continuumEuclideanAction, Pphi2-ContinuumMotion, Pphi2-ContinuumTestFunction, Pphi2-EuclideanPlaneBackground-ComplexTestFunction}
  \texttt{Pphi2.schwartz\_decompose\_continuumEuclideanActionComplex}
\end{theorem}
\begin{theorem}[\texttt{pairing\_sum\_complex\_smul}]\label{Pphi2-pairing-sum-complex-smul}
  \lean{Pphi2.pairing_sum_complex_smul}\leanok
  \uses{Pphi2-EuclideanPlaneBackground-SpaceTime, Pphi2-EuclideanPlaneBackground-RealTestFunction, Pphi2-schwartzIm-sum-complex-smul, Pphi2-EuclideanPlaneBackground-dim, Pphi2-EuclideanOS-schwartzIm, Pphi2-EuclideanOS-schwartzRe, Pphi2-EuclideanPlaneBackground, Pphi2-schwartzRe-sum-complex-smul, Pphi2-EuclideanPlaneBackground-ComplexTestFunction, Pphi2-EuclideanPlaneBackground-Distribution}
  \texttt{Pphi2.pairing\_sum\_complex\_smul}
\end{theorem}
\begin{theorem}[\texttt{schwartz\_decompose}]\label{Pphi2-schwartz-decompose}
  \lean{Pphi2.schwartz_decompose}\leanok
  \uses{Pphi2-EuclideanPlaneBackground-SpaceTime, Pphi2-schwartzOfReal, Pphi2-EuclideanPlaneBackground-dim, Pphi2-EuclideanOS-schwartzIm, Pphi2-EuclideanOS-schwartzRe, Pphi2-EuclideanPlaneBackground, Pphi2-EuclideanPlaneBackground-ComplexTestFunction}
  $J = \text{ofReal}(\text{Re}\;J) + i \cdot \text{ofReal}(\text{Im}\;J)$ for any complex Schwartz function $J$.
\end{theorem}
\begin{theorem}[\texttt{schwartzIm\_ofReal\_add\_real\_smul}]\label{Pphi2-schwartzIm-ofReal-add-real-smul}
  \lean{Pphi2.schwartzIm_ofReal_add_real_smul}\leanok
  \uses{Pphi2-EuclideanPlaneBackground-SpaceTime, Pphi2-EuclideanPlaneBackground-RealTestFunction, Pphi2-schwartzOfReal, Pphi2-EuclideanPlaneBackground-dim, Pphi2-EuclideanOS-schwartzIm, Pphi2-EuclideanPlaneBackground, Pphi2-EuclideanPlaneBackground-ComplexTestFunction}
  \texttt{Pphi2.schwartzIm\_ofReal\_add\_real\_smul}
\end{theorem}
\begin{theorem}[\texttt{schwartzRe\_schwartzOfReal}]\label{Pphi2-schwartzRe-schwartzOfReal}
  \lean{Pphi2.schwartzRe_schwartzOfReal}\leanok
  \uses{Pphi2-EuclideanPlaneBackground-SpaceTime, Pphi2-EuclideanPlaneBackground-RealTestFunction, Pphi2-schwartzOfReal, Pphi2-EuclideanPlaneBackground-dim, Pphi2-EuclideanOS-schwartzRe, Pphi2-EuclideanPlaneBackground}
  $\text{Re}(\text{ofReal}(f)) = f$ and $\text{Im}(\text{ofReal}(f)) = 0$.
\end{theorem}
\section{\texttt{Pphi2.ContinuumLimit.AxiomInheritance}}
\begin{theorem}[\texttt{os1\_for\_continuum\_limit}]\label{Pphi2-os1-for-continuum-limit}
  \lean{Pphi2.os1_for_continuum_limit}\leanok
  \uses{Pphi2-EuclideanPlaneBackground-SpaceTime, Pphi2-FieldConfig2, Pphi2-EuclideanPlaneBackground-RealTestFunction, Pphi2-planeBackground, Pphi2-EuclideanPlaneBackground-dim, Pphi2-IsPphi2Limit, Pphi2-SpaceTime2, Pphi2-exponential-moment-schwartz-bound, Pphi2-EuclideanOS-schwartzIm, Pphi2-EuclideanOS-schwartzRe, Pphi2-TestFunction2-, Pphi2-plane2Background, GaussianField-Configuration, Pphi2-continuum-exponential-moments, Pphi2-EuclideanOS-generatingFunctional-, GaussianField-instMeasurableSpaceConfiguration, Pphi2-TestFunction2, Pphi2-EuclideanOS-OS1-Regularity, Pphi2-EuclideanPlaneBackground-ComplexTestFunction, Lattice-toSemilatticeInf, Pphi2-EuclideanPlaneBackground-Distribution}
  \texttt{Pphi2.os1\_for\_continuum\_limit}
\end{theorem}
\begin{theorem}[\texttt{continuum\_exponential\_clustering} (axiom)]\label{Pphi2-continuum-exponential-clustering}
  \lean{Pphi2.continuum_exponential_clustering}\leanok
  \uses{Pphi2-EuclideanPlaneBackground-SpaceTime, Pphi2-FieldConfig2, Pphi2-EuclideanPlaneBackground-RealTestFunction, Pphi2-planeBackground, Pphi2-SchwartzMap-translate, Pphi2-EuclideanPlaneBackground-dim, Pphi2-IsPphi2Limit, Pphi2-SpaceTime2, Pphi2-plane2Background, GaussianField-instMeasurableSpaceConfiguration, Pphi2-TestFunction2, Pphi2-EuclideanOS-generatingFunctional}
  \texttt{Pphi2.continuum\_exponential\_clustering}
\end{theorem}
\begin{theorem}[\texttt{continuum\_exponential\_moment\_bound} (axiom)]\label{Pphi2-continuum-exponential-moment-bound}
  \lean{Pphi2.continuum_exponential_moment_bound}\leanok
  \uses{Pphi2-EuclideanPlaneBackground-SpaceTime, Pphi2-FieldConfig2, Pphi2-EuclideanPlaneBackground-RealTestFunction, Pphi2-planeBackground, Pphi2-EuclideanPlaneBackground-dim, Pphi2-IsPphi2Limit, GaussianField-instMeasurableSpaceConfiguration, Pphi2-continuumGreenBilinear, Pphi2-TestFunction2}
  \texttt{Pphi2.continuum\_exponential\_moment\_bound}
\end{theorem}
\begin{theorem}[\texttt{os0\_for\_continuum\_limit}]\label{Pphi2-os0-for-continuum-limit}
  \lean{Pphi2.os0_for_continuum_limit}\leanok
  \uses{Pphi2-EuclideanPlaneBackground-SpaceTime, Pphi2-FieldConfig2, Pphi2-EuclideanPlaneBackground-RealTestFunction, Pphi2-planeBackground, Pphi2-EuclideanPlaneBackground-dim, Pphi2-IsPphi2Limit, Pphi2-plane2Background, Pphi2-analyticOn-generatingFunctionalC, Pphi2-continuum-exponential-moments, GaussianField-instMeasurableSpaceConfiguration, Pphi2-TestFunction2, Pphi2-EuclideanOS-OS0-Analyticity, Pphi2-EuclideanPlaneBackground-ComplexTestFunction}
  \texttt{Pphi2.os0\_for\_continuum\_limit}
\end{theorem}
\begin{theorem}[\texttt{os1\_inheritance}]\label{Pphi2-os1-inheritance}
  \lean{Pphi2.os1_inheritance}\leanok
  \uses{Pphi2-EuclideanPlaneBackground-SpaceTime, Pphi2-planeBackground, Pphi2-EuclideanPlaneBackground-dim, GaussianField-Configuration, GaussianField-instMeasurableSpaceConfiguration, Pphi2-ContinuumTestFunction, Lattice-toSemilatticeInf}
  For any probability measure $\mu$, $|\int \cos(\omega(f))\, d\mu| \le 1$. Proved by $|\cos| \le 1$ and triangle inequality.
\end{theorem}
\begin{theorem}[\texttt{os4\_for\_continuum\_limit}]\label{Pphi2-os4-for-continuum-limit}
  \lean{Pphi2.os4_for_continuum_limit}\leanok
  \uses{Pphi2-EuclideanPlaneBackground-SpaceTime, Pphi2-EuclideanPlaneBackground-RealTestFunction, Pphi2-planeBackground, Pphi2-EuclideanPlaneBackground-translate, Pphi2-SchwartzMap-translate, Pphi2-EuclideanOS-OS4-Clustering, Pphi2-EuclideanPlaneBackground-dim, Pphi2-IsPphi2Limit, Pphi2-SpaceTime2, Pphi2-plane2Background, GaussianField-Configuration, GaussianField-instMeasurableSpaceConfiguration, Pphi2-continuum-exponential-clustering, Pphi2-TestFunction2, Pphi2-ContinuumTestFunction, Pphi2-EuclideanOS-generatingFunctional}
  \texttt{Pphi2.os4\_for\_continuum\_limit}
\end{theorem}
\begin{theorem}[\texttt{continuum\_exponential\_moments}]\label{Pphi2-continuum-exponential-moments}
  \lean{Pphi2.continuum_exponential_moments}\leanok
  \uses{Pphi2-EuclideanPlaneBackground-SpaceTime, Pphi2-FieldConfig2, Pphi2-EuclideanPlaneBackground-RealTestFunction, Pphi2-planeBackground, Pphi2-EuclideanPlaneBackground-dim, Pphi2-IsPphi2Limit, GaussianField-Configuration, Pphi2-continuum-exponential-moment-bound, GaussianField-instMeasurableSpaceConfiguration, Pphi2-continuumGreenBilinear, Pphi2-TestFunction2}
  \texttt{Pphi2.continuum\_exponential\_moments}
\end{theorem}
\begin{theorem}[\texttt{canonical\_continuumMeasure\_cf\_tendsto} (axiom)]\label{Pphi2-canonical-continuumMeasure-cf-tendsto}
  \lean{Pphi2.canonical_continuumMeasure_cf_tendsto}\leanok
  \uses{Pphi2-EuclideanPlaneBackground-SpaceTime, Pphi2-FieldConfig2, Pphi2-EuclideanPlaneBackground-RealTestFunction, Pphi2-planeBackground, Pphi2-EuclideanPlaneBackground-dim, Pphi2-IsPphi2Limit, Pphi2-continuumMeasure, Pphi2-plane2Background, GaussianField-Configuration, GaussianField-instMeasurableSpaceConfiguration, Pphi2-TestFunction2, Pphi2-ContinuumTestFunction, Pphi2-EuclideanOS-generatingFunctional}
  \texttt{Pphi2.canonical\_continuumMeasure\_cf\_tendsto}
\end{theorem}
\begin{theorem}[\texttt{exponential\_moment\_schwartz\_bound}]\label{Pphi2-exponential-moment-schwartz-bound}
  \lean{Pphi2.exponential_moment_schwartz_bound}\leanok
  \uses{Pphi2-EuclideanPlaneBackground-SpaceTime, Pphi2-FieldConfig2, Pphi2-EuclideanPlaneBackground-RealTestFunction, Pphi2-planeBackground, Pphi2-EuclideanPlaneBackground-dim, Pphi2-IsPphi2Limit, Pphi2-SpaceTime2, Pphi2-continuumGreenBilinear-le-mass-inv-sq, Pphi2-continuum-exponential-moment-bound, GaussianField-instMeasurableSpaceConfiguration, Pphi2-continuumGreenBilinear, Pphi2-TestFunction2, Lattice-toSemilatticeInf}
  \texttt{Pphi2.exponential\_moment\_schwartz\_bound}
\end{theorem}
\section{\texttt{Pphi2.InteractingMeasure.General}}
\begin{theorem}[\texttt{schwingerN\_perm}]\label{Pphi2-schwingerN-perm}
  \lean{Pphi2.schwingerN_perm}\leanok
  \uses{Pphi2-interactingMeasure, GaussianField-Configuration, GaussianField-instMeasurableSpaceConfiguration, Pphi2-schwingerN}
  $n$-point functions are invariant under permutation of test functions.
\end{theorem}
\begin{theorem}[\texttt{interactingPartitionFunction\_pos}]\label{Pphi2-interactingPartitionFunction-pos}
  \lean{Pphi2.interactingPartitionFunction_pos}\leanok
  \uses{Pphi2-interactingPartitionFunction, Pphi2-interactingBoltzmannWeight-integrable, Pphi2-interactingBoltzmannWeight, GaussianField-Configuration, Pphi2-interactingBoltzmannWeight-pos, GaussianField-instMeasurableSpaceConfiguration}
  $Z > 0$ since $e^{-V} > 0$ everywhere and $\mu$ is a probability measure.
\end{theorem}
\begin{theorem}[\texttt{interactingBoltzmannWeight\_antitone}]\label{Pphi2-interactingBoltzmannWeight-antitone}
  \lean{Pphi2.interactingBoltzmannWeight_antitone}\leanok
  \uses{Pphi2-interactingBoltzmannWeight, GaussianField-Configuration}
  If $V_1 \le V_2$ pointwise then $e^{-V_2} \le e^{-V_1}$ pointwise.
\end{theorem}
\begin{theorem}[\texttt{interactingBoltzmannWeight\_pos}]\label{Pphi2-interactingBoltzmannWeight-pos}
  \lean{Pphi2.interactingBoltzmannWeight_pos}\leanok
  \uses{Pphi2-interactingBoltzmannWeight, GaussianField-Configuration}
  $e^{-V(\omega)} > 0$ for all $\omega$.
\end{theorem}
\begin{theorem}[\texttt{schwinger2\_eq\_free\_expectation}]\label{Pphi2-schwinger2-eq-free-expectation}
  \lean{Pphi2.schwinger2_eq_free_expectation}\leanok
  \uses{Pphi2-interactingBoltzmannWeight-ennreal-measurable, Pphi2-interactingPartitionFunction, Pphi2-interactingBoltzmannWeight, GaussianField-Configuration, Pphi2-interactingBoltzmannWeight-pos, GaussianField-instMeasurableSpaceConfiguration, Pphi2-schwinger2, Pphi2-interactingPartitionFunction-pos}
  \texttt{Pphi2.schwinger2\_eq\_free\_expectation}
\end{theorem}
\begin{theorem}[\texttt{schwinger2\_symm}]\label{Pphi2-schwinger2-symm}
  \lean{Pphi2.schwinger2_symm}\leanok
  \uses{Pphi2-interactingMeasure, GaussianField-Configuration, GaussianField-instMeasurableSpaceConfiguration, Pphi2-schwinger2}
  \texttt{Pphi2.schwinger2\_symm}
\end{theorem}
\begin{definition}[\texttt{schwinger2}]\label{Pphi2-schwinger2}
  \lean{Pphi2.schwinger2}\leanok
  \uses{Pphi2-interactingMeasure, GaussianField-Configuration, GaussianField-instMeasurableSpaceConfiguration}
  Two-point, one-point, and $n$-point Schwinger functions under the interacting measure.
\end{definition}
\begin{theorem}[\texttt{interactingWithDensity\_mass}]\label{Pphi2-interactingWithDensity-mass}
  \lean{Pphi2.interactingWithDensity_mass}\leanok
  \uses{Pphi2-interactingPartitionFunction, Pphi2-interactingBoltzmannWeight-integrable, Pphi2-interactingBoltzmannWeight, GaussianField-Configuration, Pphi2-interactingBoltzmannWeight-pos, GaussianField-instMeasurableSpaceConfiguration}
  \texttt{Pphi2.interactingWithDensity\_mass}
\end{theorem}
\begin{definition}[\texttt{interactingGeneratingFunctional}]\label{Pphi2-interactingGeneratingFunctional}
  \lean{Pphi2.interactingGeneratingFunctional}\leanok
  \uses{Pphi2-interactingMeasure, GaussianField-Configuration, GaussianField-instMeasurableSpaceConfiguration}
  \texttt{Pphi2.interactingGeneratingFunctional}
\end{definition}
\begin{theorem}[\texttt{schwingerN\_eq\_free\_expectation}]\label{Pphi2-schwingerN-eq-free-expectation}
  \lean{Pphi2.schwingerN_eq_free_expectation}\leanok
  \uses{Pphi2-interactingBoltzmannWeight-ennreal-measurable, Pphi2-interactingPartitionFunction, Pphi2-interactingBoltzmannWeight, GaussianField-Configuration, Pphi2-interactingBoltzmannWeight-pos, GaussianField-instMeasurableSpaceConfiguration, Pphi2-schwingerN, Pphi2-interactingPartitionFunction-pos}
  \texttt{Pphi2.schwingerN\_eq\_free\_expectation}
\end{theorem}
\begin{definition}[\texttt{schwingerN}]\label{Pphi2-schwingerN}
  \lean{Pphi2.schwingerN}\leanok
  \uses{Pphi2-interactingMeasure, GaussianField-Configuration, GaussianField-instMeasurableSpaceConfiguration}
  \texttt{Pphi2.schwingerN}
\end{definition}
\begin{theorem}[\texttt{interactingMeasure\_integrable\_of\_bounded}]\label{Pphi2-interactingMeasure-integrable-of-bounded}
  \lean{Pphi2.interactingMeasure_integrable_of_bounded}\leanok
  \uses{Pphi2-interactingMeasure, GaussianField-Configuration, GaussianField-instMeasurableSpaceConfiguration, Pphi2-interactingMeasure-isProbabilityMeasure}
  \texttt{Pphi2.interactingMeasure\_integrable\_of\_bounded}
\end{theorem}
\begin{definition}[\texttt{interactingMeasure}]\label{Pphi2-interactingMeasure}
  \lean{Pphi2.interactingMeasure}\leanok
  \uses{Pphi2-interactingPartitionFunction, Pphi2-interactingBoltzmannWeight, GaussianField-Configuration, GaussianField-instMeasurableSpaceConfiguration}
  ``\texttt{lean def interactingMeasure (V : Configuration E → ℝ) (μ : Measure (Configuration E)) : Measure (Configuration E) }`` $d\mu_V = \frac{1}{Z} e^{-V}\,d\mu$.
\end{definition}
\begin{theorem}[\texttt{schwingerN\_zero}]\label{Pphi2-schwingerN-zero}
  \lean{Pphi2.schwingerN_zero}\leanok
  \uses{Pphi2-interactingMeasure, GaussianField-Configuration, GaussianField-instMeasurableSpaceConfiguration, Pphi2-schwingerN, Pphi2-interactingMeasure-isProbabilityMeasure}
  \texttt{Pphi2.schwingerN\_zero}
\end{theorem}
\begin{theorem}[\texttt{interactingMeasure\_absolutelyContinuous}]\label{Pphi2-interactingMeasure-absolutelyContinuous}
  \lean{Pphi2.interactingMeasure_absolutelyContinuous}\leanok
  \uses{Pphi2-interactingPartitionFunction, Pphi2-interactingMeasure, Pphi2-interactingBoltzmannWeight, GaussianField-Configuration, GaussianField-instMeasurableSpaceConfiguration}
  \texttt{Pphi2.interactingMeasure\_absolutelyContinuous}
\end{theorem}
\begin{definition}[\texttt{interactingPartitionFunction}]\label{Pphi2-interactingPartitionFunction}
  \lean{Pphi2.interactingPartitionFunction}\leanok
  \uses{Pphi2-interactingBoltzmannWeight, GaussianField-Configuration, GaussianField-instMeasurableSpaceConfiguration}
  ``\texttt{lean def interactingPartitionFunction (V : Configuration E → ℝ) (μ : Measure (Configuration E)) : ℝ }`` $Z = \int e^{-V(\omega)}\,d\mu(\omega)$.
\end{definition}
\begin{theorem}[\texttt{schwinger1\_eq\_free\_expectation}]\label{Pphi2-schwinger1-eq-free-expectation}
  \lean{Pphi2.schwinger1_eq_free_expectation}\leanok
  \uses{Pphi2-interactingBoltzmannWeight-ennreal-measurable, Pphi2-interactingPartitionFunction, Pphi2-schwinger1, Pphi2-interactingBoltzmannWeight, GaussianField-Configuration, Pphi2-interactingBoltzmannWeight-pos, GaussianField-instMeasurableSpaceConfiguration, Pphi2-interactingPartitionFunction-pos}
  \texttt{Pphi2.schwinger1\_eq\_free\_expectation}
\end{theorem}
\begin{theorem}[\texttt{interactingBoltzmannWeight\_ennreal\_measurable}]\label{Pphi2-interactingBoltzmannWeight-ennreal-measurable}
  \lean{Pphi2.interactingBoltzmannWeight_ennreal_measurable}\leanok
  \uses{Pphi2-interactingBoltzmannWeight, GaussianField-Configuration, GaussianField-instMeasurableSpaceConfiguration}
  \texttt{Pphi2.interactingBoltzmannWeight\_ennreal\_measurable}
\end{theorem}
\begin{definition}[\texttt{interactingBoltzmannWeight}]\label{Pphi2-interactingBoltzmannWeight}
  \lean{Pphi2.interactingBoltzmannWeight}\leanok
  \uses{GaussianField-Configuration}
  ``\texttt{lean def interactingBoltzmannWeight (V : Configuration E → ℝ) : Configuration E → ℝ }`` $\omega \mapsto e^{-V(\omega)}$.
\end{definition}
\begin{definition}[\texttt{schwinger1}]\label{Pphi2-schwinger1}
  \lean{Pphi2.schwinger1}\leanok
  \uses{Pphi2-interactingMeasure, GaussianField-Configuration, GaussianField-instMeasurableSpaceConfiguration}
  \texttt{Pphi2.schwinger1}
\end{definition}
\begin{theorem}[\texttt{interactingBoltzmannWeight\_integrable}]\label{Pphi2-interactingBoltzmannWeight-integrable}
  \lean{Pphi2.interactingBoltzmannWeight_integrable}\leanok
  \uses{Pphi2-interactingBoltzmannWeight, GaussianField-Configuration, Pphi2-interactingBoltzmannWeight-pos, GaussianField-instMeasurableSpaceConfiguration}
  $e^{-V}$ is integrable w.r.t. any probability measure when $V$ is measurable and bounded below.
\end{theorem}
\begin{theorem}[\texttt{interactingMeasure\_isProbabilityMeasure}]\label{Pphi2-interactingMeasure-isProbabilityMeasure}
  \lean{Pphi2.interactingMeasure_isProbabilityMeasure}\leanok
  \uses{Pphi2-interactingWithDensity-mass, Pphi2-interactingPartitionFunction, Pphi2-interactingMeasure, Pphi2-interactingBoltzmannWeight, GaussianField-Configuration, GaussianField-instMeasurableSpaceConfiguration, Pphi2-interactingPartitionFunction-pos}
  $\mu_V$ is a probability measure: $\mu_V(\text{univ}) = \frac{1}{Z} \cdot Z = 1$.
\end{theorem}
\begin{theorem}[\texttt{schwinger2\_nonneg}]\label{Pphi2-schwinger2-nonneg}
  \lean{Pphi2.schwinger2_nonneg}\leanok
  \uses{Pphi2-interactingMeasure, GaussianField-Configuration, GaussianField-instMeasurableSpaceConfiguration, Pphi2-schwinger2, Lattice-toSemilatticeInf}
  $S_2(f,f) = \int \omega(f)^2\,d\mu_V \ge 0$.
\end{theorem}
\section{\texttt{Pphi2.ContinuumLimit.CharacteristicFunctional}}
\begin{theorem}[\texttt{pphi2\_measure\_neg\_invariant}]\label{Pphi2-pphi2-measure-neg-invariant}
  \lean{Pphi2.pphi2_measure_neg_invariant}\leanok
  \uses{Pphi2-EuclideanPlaneBackground-SpaceTime, Pphi2-FieldConfig2, Pphi2-EuclideanPlaneBackground-RealTestFunction, Pphi2-planeBackground, Pphi2-EuclideanPlaneBackground-dim, Pphi2-IsPphi2Limit, Pphi2-SpaceTime2, Pphi2-PositiveTimeTestFunction2, GaussianField-instMeasurableSpaceConfiguration, Pphi2-TestFunction2, Pphi2-ContinuumTestFunction, Pphi2-compTimeReflection2, Pphi2-schwartzTranslate}
  \texttt{Pphi2.pphi2\_measure\_neg\_invariant}
\end{theorem}
\begin{theorem}[\texttt{complex\_gf\_invariant\_of\_real\_gf\_invariant}]\label{Pphi2-complex-gf-invariant-of-real-gf-invariant}
  \lean{Pphi2.complex_gf_invariant_of_real_gf_invariant}\leanok
  \uses{Pphi2-EuclideanPlaneBackground-SpaceTime, Pphi2-generatingFunctional--ofReal-add-real-smul, Pphi2-FieldConfig2, Pphi2-E2, Pphi2-EuclideanPlaneBackground-RealTestFunction, Pphi2-planeBackground, Pphi2-schwartzOfReal, Pphi2-EuclideanPlaneBackground-dim, Pphi2-schwartz-decompose-continuumEuclideanActionComplex, Pphi2-EuclideanOS-schwartzIm, Pphi2-schwartz-decompose, Pphi2-EuclideanOS-schwartzRe, Pphi2-TestFunction2-, Pphi2-plane2Background, Pphi2-analyticOn-generatingFunctionalC, Pphi2-EuclideanOS-generatingFunctional-, GaussianField-instMeasurableSpaceConfiguration, Pphi2-euclideanAction2-, Pphi2-TestFunction2, Pphi2-euclideanAction2, Pphi2-EuclideanOS-generatingFunctional, Pphi2-EuclideanOS-generatingFunctional--congr-simp, Pphi2-EuclideanPlaneBackground-ComplexTestFunction, Pphi2-EuclideanOS-generatingFunctional-congr-simp, Pphi2-EuclideanPlaneBackground-Distribution}
  \texttt{Pphi2.complex\_gf\_invariant\_of\_real\_gf\_invariant}
\end{theorem}
\begin{theorem}[\texttt{configuration\_neg\_measurable}]\label{Pphi2-configuration-neg-measurable}
  \lean{Pphi2.configuration_neg_measurable}\leanok
  \uses{Pphi2-EuclideanPlaneBackground-SpaceTime, Pphi2-FieldConfig2, Pphi2-EuclideanPlaneBackground-RealTestFunction, Pphi2-planeBackground, Pphi2-EuclideanPlaneBackground-dim, GaussianField-configuration-measurable-of-eval-measurable, GaussianField-configuration-eval-measurable, GaussianField-Configuration, GaussianField-instMeasurableSpaceConfiguration}
  \texttt{Pphi2.configuration\_neg\_measurable}
\end{theorem}
\begin{theorem}[\texttt{norm\_generatingFunctional\_le\_one}]\label{Pphi2-norm-generatingFunctional-le-one}
  \lean{Pphi2.norm_generatingFunctional_le_one}\leanok
  \uses{Pphi2-EuclideanPlaneBackground-SpaceTime, Pphi2-FieldConfig2, Pphi2-EuclideanPlaneBackground-RealTestFunction, Pphi2-planeBackground, Pphi2-EuclideanPlaneBackground-dim, Pphi2-plane2Background, GaussianField-instMeasurableSpaceConfiguration, Pphi2-TestFunction2, Pphi2-EuclideanOS-generatingFunctional}
  \texttt{Pphi2.norm\_generatingFunctional\_le\_one}
\end{theorem}
\begin{theorem}[\texttt{congr\_simp}]\label{Pphi2-EuclideanOS-generatingFunctional--congr-simp}
  \lean{Pphi2.EuclideanOS.generatingFunctionalℂ.congr_simp}\leanok
  \uses{Pphi2-EuclideanPlaneBackground-SpaceTime, Pphi2-EuclideanPlaneBackground-RealTestFunction, Pphi2-EuclideanPlaneBackground-dim, Pphi2-EuclideanPlaneBackground, Pphi2-EuclideanOS-generatingFunctional-, GaussianField-instMeasurableSpaceConfiguration, Pphi2-EuclideanPlaneBackground-ComplexTestFunction, Pphi2-EuclideanPlaneBackground-Distribution}
  \texttt{Pphi2.EuclideanOS.generatingFunctionalℂ.congr\_simp}
\end{theorem}
\begin{theorem}[\texttt{analyticOn\_generatingFunctionalC}]\label{Pphi2-analyticOn-generatingFunctionalC}
  \lean{Pphi2.analyticOn_generatingFunctionalC}\leanok
  \uses{Pphi2-EuclideanPlaneBackground-SpaceTime, Pphi2-FieldConfig2, Pphi2-EuclideanPlaneBackground-RealTestFunction, Pphi2-planeBackground, Pphi2-EuclideanPlaneBackground-dim, Pphi2-EuclideanOS-schwartzIm, Pphi2-EuclideanOS-schwartzRe, Pphi2-TestFunction2-, Pphi2-plane2Background, GaussianField-configuration-eval-measurable, GaussianField-Configuration, Pphi2-EuclideanOS-generatingFunctional-, GaussianField-instMeasurableSpaceConfiguration, Pphi2-TestFunction2, Pphi2-EuclideanPlaneBackground-ComplexTestFunction, Pphi2-EuclideanPlaneBackground-Distribution}
  \texttt{Pphi2.analyticOn\_generatingFunctionalC}
\end{theorem}
\begin{theorem}[\texttt{pphi2\_generating\_functional\_real}]\label{Pphi2-pphi2-generating-functional-real}
  \lean{Pphi2.pphi2_generating_functional_real}\leanok
  \uses{Pphi2-EuclideanPlaneBackground-SpaceTime, Pphi2-FieldConfig2, Pphi2-pphi2-measure-neg-invariant, Pphi2-EuclideanPlaneBackground-RealTestFunction, Pphi2-planeBackground, Pphi2-EuclideanPlaneBackground-dim, Pphi2-IsPphi2Limit, Pphi2-plane2Background, GaussianField-configuration-eval-measurable, GaussianField-Configuration, GaussianField-instMeasurableSpaceConfiguration, Pphi2-configuration-neg-measurable, Pphi2-TestFunction2, Pphi2-EuclideanOS-generatingFunctional, Pphi2-EuclideanPlaneBackground-Distribution}
  \texttt{Pphi2.pphi2\_generating\_functional\_real}
\end{theorem}
\begin{theorem}[\texttt{congr\_simp}]\label{Pphi2-EuclideanOS-generatingFunctional-congr-simp}
  \lean{Pphi2.EuclideanOS.generatingFunctional.congr_simp}\leanok
  \uses{Pphi2-EuclideanPlaneBackground-SpaceTime, Pphi2-EuclideanPlaneBackground-RealTestFunction, Pphi2-EuclideanPlaneBackground-dim, Pphi2-EuclideanPlaneBackground, GaussianField-instMeasurableSpaceConfiguration, Pphi2-EuclideanOS-generatingFunctional, Pphi2-EuclideanPlaneBackground-Distribution}
  \texttt{Pphi2.EuclideanOS.generatingFunctional.congr\_simp}
\end{theorem}
\begin{theorem}[\texttt{os4\_clustering\_implies\_ergodicity}]\label{Pphi2-os4-clustering-implies-ergodicity}
  \lean{Pphi2.os4_clustering_implies_ergodicity}\leanok
  \uses{Pphi2-EuclideanPlaneBackground-SpaceTime, Pphi2-FieldConfig2, Pphi2-EuclideanTimeStructure-timeTranslation, Pphi2-EuclideanPlaneBackground-RealTestFunction, Pphi2-planeBackground, Pphi2-EuclideanPlaneBackground-translate, Pphi2-SchwartzMap-translate, Pphi2-EuclideanOS-OS4-Clustering, Pphi2-timeTranslationVector2, Pphi2-EuclideanPlaneBackground-dim, Pphi2-IsPphi2Limit, Pphi2-SpaceTime2, Pphi2-norm-generatingFunctional-le-one, Pphi2-plane2Background, Pphi2-timeTranslationVector2-ofLp-zero, GaussianField-instMeasurableSpaceConfiguration, Pphi2-TestFunction2, Pphi2-EuclideanOS-OS4-Ergodicity, Pphi2-generatingFunctional-translate-continuous, Pphi2-EuclideanOS-generatingFunctional, Pphi2-pphi2-generating-functional-real, Pphi2-plane2TimeStructure}
  \texttt{Pphi2.os4\_clustering\_implies\_ergodicity}
\end{theorem}
\begin{theorem}[\texttt{generatingFunctional\_translate\_continuous}]\label{Pphi2-generatingFunctional-translate-continuous}
  \lean{Pphi2.generatingFunctional_translate_continuous}\leanok
  \uses{Pphi2-EuclideanPlaneBackground-SpaceTime, Pphi2-FieldConfig2, Pphi2-translate-timeTranslationVector2-eq-timeTranslationSchwartz-neg, Pphi2-EuclideanPlaneBackground-RealTestFunction, Pphi2-planeBackground, Pphi2-SchwartzMap-translate, Pphi2-timeTranslationVector2, Pphi2-EuclideanPlaneBackground-dim, Pphi2-plane2Background, GaussianField-configuration-eval-measurable, GaussianField-Configuration, GaussianField-instMeasurableSpaceConfiguration, Pphi2-TestFunction2, Pphi2-EuclideanOS-generatingFunctional}
  \texttt{Pphi2.generatingFunctional\_translate\_continuous}
\end{theorem}
\section{\texttt{Pphi2.NelsonEstimate.ChaosTailBridge}}
\begin{theorem}[\texttt{chaos\_neg\_tail\_bound\_explicit}]\label{Pphi2-ChaosTailBridge-chaos-neg-tail-bound-explicit}
  \lean{Pphi2.ChaosTailBridge.chaos_neg_tail_bound_explicit}\leanok
  \uses{Pphi2-ChaosTailBridge-chaosTailConstant-pos, GaussianHilbert-stdGaussianFin, Pphi2-ChaosTailBridge-polynomial-chaos-concentration-explicit, GaussianHilbert-wienerChaosLE, Pphi2-ChaosTailBridge-chaosTailConstant}
  \texttt{Pphi2.ChaosTailBridge.chaos\_neg\_tail\_bound\_explicit}
\end{theorem}
\begin{theorem}[\texttt{polynomial\_chaos\_concentration\_explicit}]\label{Pphi2-ChaosTailBridge-polynomial-chaos-concentration-explicit}
  \lean{Pphi2.ChaosTailBridge.polynomial_chaos_concentration_explicit}\leanok
  \uses{GaussianHilbert-stdGaussianFin, GaussianHilbert-wienerChaosLE, Pphi2-ChaosTailBridge-chaosTailConstant, Lattice-toSemilatticeInf}
  \texttt{Pphi2.ChaosTailBridge.polynomial\_chaos\_concentration\_explicit}
\end{theorem}
\begin{theorem}[\texttt{chaosTailConstant\_pos}]\label{Pphi2-ChaosTailBridge-chaosTailConstant-pos}
  \lean{Pphi2.ChaosTailBridge.chaosTailConstant_pos}\leanok
  \uses{Pphi2-ChaosTailBridge-chaosTailConstant}
  \texttt{Pphi2.ChaosTailBridge.chaosTailConstant\_pos}
\end{theorem}
\begin{definition}[\texttt{chaosTailConstant}]\label{Pphi2-ChaosTailBridge-chaosTailConstant}
  \lean{Pphi2.ChaosTailBridge.chaosTailConstant}\leanok
  \texttt{Pphi2.ChaosTailBridge.chaosTailConstant}
\end{definition}
\begin{theorem}[\texttt{chaos\_neg\_tail\_bound}]\label{Pphi2-ChaosTailBridge-chaos-neg-tail-bound}
  \lean{Pphi2.ChaosTailBridge.chaos_neg_tail_bound}\leanok
  \uses{GaussianHilbert-polynomial-chaos-concentration, GaussianHilbert-stdGaussianFin, GaussianHilbert-wienerChaosLE}
  \texttt{Pphi2.ChaosTailBridge.chaos\_neg\_tail\_bound}
\end{theorem}
\section{\texttt{Pphi2.NelsonEstimate.LatticeBridge}}
\begin{theorem}[\texttt{hintegral}]\label{Pphi2-LatticeBridge-LatticeRoughErrorSetup-hintegral}
  \lean{Pphi2.LatticeBridge.LatticeRoughErrorSetup.hintegral}\leanok
  \uses{Pphi2-LatticeBridge-LatticeRoughErrorSetup--, Pphi2-LatticeBridge-LatticeRoughErrorSetup}
  \texttt{Pphi2.LatticeBridge.LatticeRoughErrorSetup.hintegral}
\end{theorem}
\begin{definition}[\texttt{ψ}]\label{Pphi2-LatticeBridge-LatticeRoughErrorSetup--}
  \lean{Pphi2.LatticeBridge.LatticeRoughErrorSetup.ψ}\leanok
  \uses{Pphi2-LatticeBridge-LatticeRoughErrorSetup}
  \texttt{Pphi2.LatticeBridge.LatticeRoughErrorSetup.ψ}
\end{definition}
\begin{theorem}[\texttt{hdecomp}]\label{Pphi2-LatticeBridge-LatticeRoughErrorSetup-hdecomp}
  \lean{Pphi2.LatticeBridge.LatticeRoughErrorSetup.hdecomp}\leanok
  \uses{GaussianField-FinLatticeSites, Pphi2-LatticeBridge-LatticeRoughErrorSetup, GaussianField-Configuration, GaussianField-FinLatticeField, Pphi2-interactionFunctional, Pphi2-LatticeBridge-LatticeRoughErrorSetup-V-S, Pphi2-LatticeBridge-LatticeRoughErrorSetup-E-R}
  \texttt{Pphi2.LatticeBridge.LatticeRoughErrorSetup.hdecomp}
\end{theorem}
\begin{theorem}[\texttt{htail}]\label{Pphi2-LatticeBridge-LatticeRoughErrorSetup-htail}
  \lean{Pphi2.LatticeBridge.LatticeRoughErrorSetup.htail}\leanok
  \uses{GaussianField-FinLatticeSites, Pphi2-LatticeBridge-LatticeRoughErrorSetup--, Pphi2-LatticeBridge-LatticeRoughErrorSetup, GaussianField-latticeGaussianMeasure, GaussianField-Configuration, GaussianField-instMeasurableSpaceConfiguration, GaussianField-FinLatticeField, Pphi2-LatticeBridge-LatticeRoughErrorSetup-E-R}
  \texttt{Pphi2.LatticeBridge.LatticeRoughErrorSetup.htail}
\end{theorem}
\begin{definition}[\texttt{V\_S}]\label{Pphi2-LatticeBridge-LatticeRoughErrorSetup-V-S}
  \lean{Pphi2.LatticeBridge.LatticeRoughErrorSetup.V_S}\leanok
  \uses{GaussianField-FinLatticeSites, Pphi2-LatticeBridge-LatticeRoughErrorSetup, GaussianField-Configuration, GaussianField-FinLatticeField}
  \texttt{Pphi2.LatticeBridge.LatticeRoughErrorSetup.V\_S}
\end{definition}
\begin{definition}[\texttt{LatticeRoughErrorSetup}]\label{Pphi2-LatticeBridge-LatticeRoughErrorSetup}
  \lean{Pphi2.LatticeBridge.LatticeRoughErrorSetup}\leanok
  \texttt{Pphi2.LatticeBridge.LatticeRoughErrorSetup}
\end{definition}
\begin{theorem}[\texttt{bridgeAxiom\_of\_varying\_coupled\_setup\_real}]\label{Pphi2-LatticeBridge-bridgeAxiom-of-varying-coupled-setup-real}
  \lean{Pphi2.LatticeBridge.bridgeAxiom_of_varying_coupled_setup_real}\leanok
  \uses{GaussianField-FinLatticeSites, Pphi2-interactionFunctional-measurable, Pphi2-LayerCake-lintegral-expSq-neg-le-of-tail, GaussianField-latticeGaussianMeasure-isProbability, GaussianField-latticeGaussianMeasure, GaussianField-Configuration, GaussianField-instMeasurableSpaceConfiguration, GaussianField-FinLatticeField, Pphi2-interactionFunctional, Lattice-toSemilatticeInf}
  \texttt{Pphi2.LatticeBridge.bridgeAxiom\_of\_varying\_coupled\_setup\_real}
\end{theorem}
\begin{theorem}[\texttt{bridgeAxiom\_of\_setup\_real}]\label{Pphi2-LatticeBridge-bridgeAxiom-of-setup-real}
  \lean{Pphi2.LatticeBridge.bridgeAxiom_of_setup_real}\leanok
  \uses{GaussianField-FinLatticeSites, Pphi2-LatticeBridge-LatticeRoughErrorSetup--, Pphi2-LatticeBridge-LatticeRoughErrorSetup, GaussianField-latticeGaussianMeasure, Pphi2-LatticeBridge-bridgeAxiom-of-setup, GaussianField-Configuration, GaussianField-instMeasurableSpaceConfiguration, Pphi2-LatticeBridge-LatticeRoughErrorSetup-hV-meas, Pphi2-LatticeBridge-LatticeRoughErrorSetup-hintegral, GaussianField-FinLatticeField, Pphi2-interactionFunctional, Lattice-toSemilatticeInf}
  \texttt{Pphi2.LatticeBridge.bridgeAxiom\_of\_setup\_real}
\end{theorem}
\begin{definition}[\texttt{E\_R}]\label{Pphi2-LatticeBridge-LatticeRoughErrorSetup-E-R}
  \lean{Pphi2.LatticeBridge.LatticeRoughErrorSetup.E_R}\leanok
  \uses{GaussianField-FinLatticeSites, Pphi2-LatticeBridge-LatticeRoughErrorSetup, GaussianField-Configuration, GaussianField-FinLatticeField}
  \texttt{Pphi2.LatticeBridge.LatticeRoughErrorSetup.E\_R}
\end{definition}
\begin{theorem}[\texttt{hsmooth}]\label{Pphi2-LatticeBridge-LatticeRoughErrorSetup-hsmooth}
  \lean{Pphi2.LatticeBridge.LatticeRoughErrorSetup.hsmooth}\leanok
  \uses{GaussianField-FinLatticeSites, Pphi2-LatticeBridge-LatticeRoughErrorSetup, GaussianField-Configuration, GaussianField-FinLatticeField, Pphi2-LatticeBridge-LatticeRoughErrorSetup-V-S}
  \texttt{Pphi2.LatticeBridge.LatticeRoughErrorSetup.hsmooth}
\end{theorem}
\begin{theorem}[\texttt{bridgeAxiom\_of\_coupled\_setup\_real}]\label{Pphi2-LatticeBridge-bridgeAxiom-of-coupled-setup-real}
  \lean{Pphi2.LatticeBridge.bridgeAxiom_of_coupled_setup_real}\leanok
  \uses{GaussianField-FinLatticeSites, Pphi2-interactionFunctional-measurable, GaussianField-latticeGaussianMeasure, GaussianField-Configuration, Pphi2-BridgeFromTail-expSqNeg-lintegral-le-of-dynamical-cutoff, GaussianField-instMeasurableSpaceConfiguration, GaussianField-FinLatticeField, Pphi2-interactionFunctional, Lattice-toSemilatticeInf}
  \texttt{Pphi2.LatticeBridge.bridgeAxiom\_of\_coupled\_setup\_real}
\end{theorem}
\begin{theorem}[\texttt{bridgeAxiom\_of\_setup}]\label{Pphi2-LatticeBridge-bridgeAxiom-of-setup}
  \lean{Pphi2.LatticeBridge.bridgeAxiom_of_setup}\leanok
  \uses{Pphi2-LatticeBridge-LatticeRoughErrorSetup-hdecomp, GaussianField-FinLatticeSites, Pphi2-LatticeBridge-LatticeRoughErrorSetup--, Pphi2-LatticeBridge-LatticeRoughErrorSetup, GaussianField-latticeGaussianMeasure-isProbability, GaussianField-latticeGaussianMeasure, GaussianField-Configuration, Pphi2-BridgeFromTail-expSqNeg-lintegral-le-of-dynamical-cutoff, GaussianField-instMeasurableSpaceConfiguration, Pphi2-LatticeBridge-LatticeRoughErrorSetup-hV-meas, GaussianField-FinLatticeField, Pphi2-interactionFunctional, Pphi2-LatticeBridge-LatticeRoughErrorSetup-hsmooth, Pphi2-LatticeBridge-LatticeRoughErrorSetup-V-S, Pphi2-LatticeBridge-LatticeRoughErrorSetup-htail, Pphi2-LatticeBridge-LatticeRoughErrorSetup-E-R}
  \texttt{Pphi2.LatticeBridge.bridgeAxiom\_of\_setup}
\end{theorem}
\begin{definition}[\texttt{ctorIdx}]\label{Pphi2-LatticeBridge-LatticeRoughErrorSetup-ctorIdx}
  \lean{Pphi2.LatticeBridge.LatticeRoughErrorSetup.ctorIdx}\leanok
  \uses{Pphi2-LatticeBridge-LatticeRoughErrorSetup}
  \texttt{Pphi2.LatticeBridge.LatticeRoughErrorSetup.ctorIdx}
\end{definition}
\begin{theorem}[\texttt{hV\_meas}]\label{Pphi2-LatticeBridge-LatticeRoughErrorSetup-hV-meas}
  \lean{Pphi2.LatticeBridge.LatticeRoughErrorSetup.hV_meas}\leanok
  \uses{GaussianField-FinLatticeSites, Pphi2-LatticeBridge-LatticeRoughErrorSetup, GaussianField-Configuration, GaussianField-instMeasurableSpaceConfiguration, GaussianField-FinLatticeField, Pphi2-interactionFunctional}
  \texttt{Pphi2.LatticeBridge.LatticeRoughErrorSetup.hV\_meas}
\end{theorem}
\section{\texttt{Pphi2.OSProofs.OS2\_WardIdentity}}
\begin{definition}[\texttt{rotationCFPointwiseDefect}]\label{Pphi2-rotationCFPointwiseDefect}
  \lean{Pphi2.rotationCFPointwiseDefect}\leanok
  \uses{Pphi2-EuclideanPlaneBackground-SpaceTime, Pphi2-FieldConfig2, Pphi2-EuclideanPlaneBackground-RealTestFunction, Pphi2-planeBackground, Pphi2-EuclideanPlaneBackground-dim, Pphi2-TestFunction2, Pphi2-euclideanAction2, Pphi2-O2}
  \texttt{Pphi2.rotationCFPointwiseDefect}
\end{definition}
\begin{definition}[\texttt{rotationCFPointwiseDefectIntegrand}]\label{Pphi2-rotationCFPointwiseDefectIntegrand}
  \lean{Pphi2.rotationCFPointwiseDefectIntegrand}\leanok
  \uses{Pphi2-EuclideanPlaneBackground-SpaceTime, Pphi2-FieldConfig2, Pphi2-EuclideanPlaneBackground-RealTestFunction, Pphi2-planeBackground, Pphi2-EuclideanPlaneBackground-dim, Pphi2-TestFunction2, Pphi2-euclideanAction2, Pphi2-O2}
  \texttt{Pphi2.rotationCFPointwiseDefectIntegrand}
\end{definition}
\begin{theorem}[\texttt{os3\_for\_continuum\_limit}]\label{Pphi2-os3-for-continuum-limit}
  \lean{Pphi2.os3_for_continuum_limit}\leanok
  \uses{Pphi2-EuclideanPlaneBackground-SpaceTime, Pphi2-FieldConfig2, Pphi2-EuclideanPlaneBackground-RealTestFunction, Pphi2-planeBackground, Pphi2-EuclideanPlaneBackground-dim, Pphi2-IsPphi2Limit, Pphi2-SpaceTime2, Pphi2-EuclideanTimeStructure-positiveTimeSubmodule, Pphi2-PositiveTimeTestFunction2, Pphi2-OS3-ReflectionPositivity, Pphi2-plane2Background, GaussianField-Configuration, GaussianField-instMeasurableSpaceConfiguration, Pphi2-EuclideanTimeStructure-reflectOnRealTestFunctions, Pphi2-TestFunction2, Pphi2-ContinuumTestFunction, Pphi2-compTimeReflection2, Pphi2-EuclideanOS-generatingFunctional, Pphi2-plane2TimeStructure, Pphi2-schwartzTranslate, Pphi2-EuclideanPlaneBackground-Distribution}
  \texttt{Pphi2.os3\_for\_continuum\_limit}
\end{theorem}
\begin{theorem}[\texttt{anomaly\_vanishes}]\label{Pphi2-anomaly-vanishes}
  \lean{Pphi2.anomaly_vanishes}\leanok
  \uses{Pphi2-EuclideanPlaneBackground-SpaceTime, Pphi2-generatingFunctional, Pphi2-planeBackground, Pphi2-continuumMeasure-isProbability, Pphi2-EuclideanPlaneBackground-dim, Pphi2-continuumMeasure, Pphi2-TestFunction2, Pphi2-anomaly-bound-from-superrenormalizability, Pphi2-euclideanAction2, Pphi2-O2}
  \texttt{Pphi2.anomaly\_vanishes}
\end{theorem}
\begin{theorem}[\texttt{rotation\_cf\_defect\_polylog\_bound} (axiom)]\label{Pphi2-rotation-cf-defect-polylog-bound}
  \lean{Pphi2.rotation_cf_defect_polylog_bound}\leanok
  \uses{Pphi2-rotationCFDefect, Pphi2-TestFunction2, Pphi2-O2}
  \texttt{Pphi2.rotation\_cf\_defect\_polylog\_bound}
\end{theorem}
\begin{theorem}[\texttt{rotationCFDefect\_le\_pointwiseDefectIntegral}]\label{Pphi2-rotationCFDefect-le-pointwiseDefectIntegral}
  \lean{Pphi2.rotationCFDefect_le_pointwiseDefectIntegral}\leanok
  \uses{Pphi2-EuclideanPlaneBackground-SpaceTime, Pphi2-FieldConfig2, Pphi2-EuclideanPlaneBackground-RealTestFunction, Pphi2-planeBackground, Pphi2-continuumMeasure-isProbability, Pphi2-EuclideanPlaneBackground-dim, Pphi2-rotationCFDefect, Pphi2-continuumMeasure, Pphi2-plane2Background, GaussianField-Configuration, GaussianField-instMeasurableSpaceConfiguration, Pphi2-TestFunction2, Pphi2-ContinuumTestFunction, Pphi2-euclideanAction2, Pphi2-rotationCFPointwiseDefect, Pphi2-O2}
  \texttt{Pphi2.rotationCFDefect\_le\_pointwiseDefectIntegral}
\end{theorem}
\begin{theorem}[\texttt{anomaly\_scaling\_dimension}]\label{Pphi2-anomaly-scaling-dimension}
  \lean{Pphi2.anomaly_scaling_dimension}\leanok
  \uses{Lattice-toSemilatticeInf}
  \texttt{Pphi2.anomaly\_scaling\_dimension}
\end{theorem}
\begin{theorem}[\texttt{interactingLatticeMeasure\_expectation\_configEuclideanTimeShift}]\label{Pphi2-interactingLatticeMeasure-expectation-configEuclideanTimeShift}
  \lean{Pphi2.interactingLatticeMeasure_expectation_configEuclideanTimeShift}\leanok
  \uses{Pphi2-latticeTranslation, GaussianField-FinLatticeSites, GaussianField-latticeTranslation, Pphi2-latticeConfigEuclideanTimeShift, Pphi2-interactingLatticeMeasure, GaussianField-Configuration, Pphi2-latticeMeasure-translation-invariant, Pphi2-latticeEuclideanTimeShift, GaussianField-instMeasurableSpaceConfiguration, GaussianField-FinLatticeField}
  \texttt{Pphi2.interactingLatticeMeasure\_expectation\_configEuclideanTimeShift}
\end{theorem}
\begin{definition}[\texttt{rotationCFDefect}]\label{Pphi2-rotationCFDefect}
  \lean{Pphi2.rotationCFDefect}\leanok
  \uses{Pphi2-EuclideanPlaneBackground-SpaceTime, Pphi2-generatingFunctional, Pphi2-planeBackground, Pphi2-EuclideanPlaneBackground-dim, Pphi2-continuumMeasure, Pphi2-TestFunction2, Pphi2-euclideanAction2, Pphi2-O2}
  \texttt{Pphi2.rotationCFDefect}
\end{definition}
\begin{theorem}[\texttt{translation\_invariance\_continuum}]\label{Pphi2-translation-invariance-continuum}
  \lean{Pphi2.translation_invariance_continuum}\leanok
  \uses{Pphi2-EuclideanPlaneBackground-SpaceTime, Pphi2-FieldConfig2, Pphi2-generatingFunctional, Pphi2-EuclideanPlaneBackground-RealTestFunction, Pphi2-planeBackground, Pphi2-SchwartzMap-translate, Pphi2-EuclideanPlaneBackground-dim, Pphi2-IsPphi2Limit, Pphi2-SpaceTime2, Pphi2-PositiveTimeTestFunction2, GaussianField-Configuration, GaussianField-instMeasurableSpaceConfiguration, Pphi2-TestFunction2, Pphi2-ContinuumTestFunction, Pphi2-compTimeReflection2, Pphi2-schwartzTranslate}
  \texttt{Pphi2.translation\_invariance\_continuum}
\end{theorem}
\begin{theorem}[\texttt{massOperator\_translation\_commute}]\label{Pphi2-massOperator-translation-commute}
  \lean{Pphi2.massOperator_translation_commute}\leanok
  \uses{Pphi2-latticeTranslation, GaussianField-FinLatticeSites, GaussianField-massOperator, GaussianField-finiteLaplacian-translation-commute, GaussianField-latticeTranslationFun, GaussianField-finiteLaplacian, GaussianField-FinLatticeField, Lattice-toSemilatticeInf}
  \texttt{Pphi2.massOperator\_translation\_commute}
\end{theorem}
\begin{definition}[\texttt{latticeRotation90}]\label{Pphi2-latticeRotation90}
  \lean{Pphi2.latticeRotation90}\leanok
  \uses{GaussianField-FinLatticeSites}
  \texttt{Pphi2.latticeRotation90}
\end{definition}
\begin{theorem}[\texttt{latticeMeasure\_translation\_invariant}]\label{Pphi2-latticeMeasure-translation-invariant}
  \lean{Pphi2.latticeMeasure_translation_invariant}\leanok
  \uses{Pphi2-latticeTranslation, GaussianField-finLatticeDelta, GaussianField-evalMap, GaussianField-FinLatticeSites, Pphi2-partitionFunction, Pphi2-gaussianDensity-translation-invariant, GaussianField-evalMapMeasurableEquiv, Pphi2-interactionFunctional-measurable, Pphi2-interactingLatticeMeasure, GaussianField-latticeGaussianMeasure, GaussianField-Configuration, GaussianField-latticeGaussianMeasure-density-integral-, Pphi2-wickConstant, GaussianField-instMeasurableSpaceConfiguration, Pphi2-boltzmannWeight, GaussianField-gaussianDensity, GaussianField-FinLatticeField, Pphi2-interactionFunctional, Pphi2-wickPolynomial, Pphi2-boltzmannWeight-pos}
  \texttt{Pphi2.latticeMeasure\_translation\_invariant}
\end{theorem}
\begin{theorem}[\texttt{pphi2\_exists}]\label{Pphi2-pphi2-exists}
  \lean{Pphi2.pphi2_exists}\leanok
  \uses{Pphi2-EuclideanPlaneBackground-SpaceTime, Pphi2-SatisfiesFullOS, Pphi2-planeBackground, Pphi2-pphi2-limit-exists, Pphi2-EuclideanPlaneBackground-dim, Pphi2-IsPphi2Limit, Pphi2-continuumLimit-satisfies-fullOS, GaussianField-Configuration, GaussianField-instMeasurableSpaceConfiguration, Pphi2-ContinuumTestFunction}
  \texttt{Pphi2.pphi2\_exists}
\end{theorem}
\begin{definition}[\texttt{so2Generator}]\label{Pphi2-so2Generator}
  \lean{Pphi2.so2Generator}\leanok
  \uses{Pphi2-ContinuumTestFunction}
  \texttt{Pphi2.so2Generator}
\end{definition}
\begin{theorem}[\texttt{continuumLimit\_satisfies\_fullOS}]\label{Pphi2-continuumLimit-satisfies-fullOS}
  \lean{Pphi2.continuumLimit_satisfies_fullOS}\leanok
  \uses{Pphi2-EuclideanPlaneBackground-SpaceTime, Pphi2-SatisfiesFullOS, Pphi2-planeBackground, Pphi2-EuclideanPlaneBackground-dim, Pphi2-IsPphi2Limit, Pphi2-os0-for-continuum-limit, Pphi2-os4-for-continuum-limit, Pphi2-plane2Background, GaussianField-Configuration, Pphi2-os3-for-continuum-limit, GaussianField-instMeasurableSpaceConfiguration, Pphi2-ContinuumTestFunction, Pphi2-plane2TimeStructure, Pphi2-os4-clustering-implies-ergodicity, Pphi2-os1-for-continuum-limit, Pphi2-os2-for-continuum-limit}
  \texttt{Pphi2.continuumLimit\_satisfies\_fullOS}
\end{theorem}
\begin{theorem}[\texttt{congr\_simp}]\label{Pphi2-generatingFunctional-congr-simp}
  \lean{Pphi2.generatingFunctional.congr_simp}\leanok
  \uses{Pphi2-EuclideanPlaneBackground-SpaceTime, Pphi2-FieldConfig2, Pphi2-generatingFunctional, Pphi2-EuclideanPlaneBackground-RealTestFunction, Pphi2-planeBackground, Pphi2-EuclideanPlaneBackground-dim, GaussianField-instMeasurableSpaceConfiguration, Pphi2-TestFunction2}
  \texttt{Pphi2.generatingFunctional.congr\_simp}
\end{theorem}
\begin{theorem}[\texttt{rotation\_invariance\_continuum}]\label{Pphi2-rotation-invariance-continuum}
  \lean{Pphi2.rotation_invariance_continuum}\leanok
  \uses{Pphi2-EuclideanPlaneBackground-SpaceTime, Pphi2-FieldConfig2, Pphi2-generatingFunctional, Pphi2-EuclideanPlaneBackground-RealTestFunction, Pphi2-planeBackground, Pphi2-EuclideanPlaneBackground-dim, Pphi2-rotationCFDefect, Pphi2-IsPphi2Limit, Pphi2-continuumMeasure, GaussianField-Configuration, Pphi2-rotation-cf-defect-polylog-bound, GaussianField-instMeasurableSpaceConfiguration, Pphi2-TestFunction2, Pphi2-canonical-continuumMeasure-cf-tendsto, Pphi2-ContinuumTestFunction, Pphi2-euclideanAction2, Pphi2-O2}
  \texttt{Pphi2.rotation\_invariance\_continuum}
\end{theorem}
\begin{theorem}[\texttt{anomaly\_bound\_from\_superrenormalizability}]\label{Pphi2-anomaly-bound-from-superrenormalizability}
  \lean{Pphi2.anomaly_bound_from_superrenormalizability}\leanok
  \uses{Pphi2-EuclideanPlaneBackground-SpaceTime, Pphi2-generatingFunctional, Pphi2-planeBackground, Pphi2-continuumMeasure-isProbability, Pphi2-EuclideanPlaneBackground-dim, Pphi2-rotationCFDefect, Pphi2-continuumMeasure, Pphi2-rotation-cf-defect-polylog-bound, Pphi2-TestFunction2, Pphi2-euclideanAction2, Pphi2-O2}
  \texttt{Pphi2.anomaly\_bound\_from\_superrenormalizability}
\end{theorem}
\begin{theorem}[\texttt{gaussianDensity\_translation\_invariant}]\label{Pphi2-gaussianDensity-translation-invariant}
  \lean{Pphi2.gaussianDensity_translation_invariant}\leanok
  \uses{Pphi2-latticeTranslation, GaussianField-FinLatticeSites, GaussianField-massOperator, Pphi2-massOperator-translation-commute, GaussianField-gaussianDensity, GaussianField-FinLatticeField}
  \texttt{Pphi2.gaussianDensity\_translation\_invariant}
\end{theorem}
\begin{theorem}[\texttt{latticeAction\_translation\_invariant}]\label{Pphi2-latticeAction-translation-invariant}
  \lean{Pphi2.latticeAction_translation_invariant}\leanok
  \uses{Pphi2-latticeTranslation, GaussianField-FinLatticeSites, Pphi2-wickConstant, GaussianField-FinLatticeField, Pphi2-wickPolynomial, Pphi2-latticeInteraction}
  \texttt{Pphi2.latticeAction\_translation\_invariant}
\end{theorem}
\begin{definition}[\texttt{latticeTranslation}]\label{Pphi2-latticeTranslation}
  \lean{Pphi2.latticeTranslation}\leanok
  \uses{GaussianField-FinLatticeSites, GaussianField-FinLatticeField}
  \texttt{Pphi2.latticeTranslation}
\end{definition}
\begin{theorem}[\texttt{ward\_identity\_lattice}]\label{Pphi2-ward-identity-lattice}
  \lean{Pphi2.ward_identity_lattice}\leanok
  \uses{GaussianField-FinLatticeSites, Pphi2-interactingLatticeMeasure, GaussianField-Configuration, GaussianField-instMeasurableSpaceConfiguration, GaussianField-FinLatticeField, Lattice-toSemilatticeInf}
  \texttt{Pphi2.ward\_identity\_lattice}
\end{theorem}
\begin{theorem}[\texttt{os2\_for\_continuum\_limit}]\label{Pphi2-os2-for-continuum-limit}
  \lean{Pphi2.os2_for_continuum_limit}\leanok
  \uses{Pphi2-EuclideanPlaneBackground-SpaceTime, Pphi2-FieldConfig2, Pphi2-E2, Pphi2-generatingFunctional, Pphi2-EuclideanPlaneBackground-Motion-t, Pphi2-EuclideanPlaneBackground-Motion-R, Pphi2-OS2-EuclideanInvariance, Pphi2-EuclideanPlaneBackground-RealTestFunction, Pphi2-planeBackground, Pphi2-SchwartzMap-translate, Pphi2-EuclideanPlaneBackground-dim, Pphi2-IsPphi2Limit, Pphi2-EuclideanPlaneBackground-Motion, Pphi2-plane2Background, GaussianField-Configuration, Pphi2-continuum-exponential-moments, GaussianField-instMeasurableSpaceConfiguration, Pphi2-TestFunction2, Pphi2-rotation-invariance-continuum, Pphi2-ContinuumTestFunction, Pphi2-translation-invariance-continuum, Pphi2-euclideanAction2, Pphi2-EuclideanPlaneBackground-ComplexTestFunction, Pphi2-complex-gf-invariant-of-real-gf-invariant, Pphi2-O2}
  \texttt{Pphi2.os2\_for\_continuum\_limit}
\end{theorem}
\section{\texttt{Pphi2.Backgrounds.EuclideanPlane2D}}
\begin{definition}[\texttt{plane2Background}]\label{Pphi2-plane2Background}
  \lean{Pphi2.plane2Background}\leanok
  \uses{Pphi2-planeBackground, Pphi2-EuclideanPlaneBackground}
  \texttt{Pphi2.plane2Background}
\end{definition}
\begin{definition}[\texttt{compTimeReflection2ℂ}]\label{Pphi2-compTimeReflection2-}
  \lean{Pphi2.compTimeReflection2ℂ}\leanok
  \uses{Pphi2-EuclideanPlaneBackground-SpaceTime, Pphi2-EuclideanTimeStructure-reflectOnComplexTestFunctions, Pphi2-planeBackground, Pphi2-EuclideanPlaneBackground-dim, Pphi2-TestFunction2-, Pphi2-plane2Background, Pphi2-plane2TimeStructure}
  \texttt{Pphi2.compTimeReflection2ℂ}
\end{definition}
\begin{theorem}[\texttt{timeReflection2\_involution}]\label{Pphi2-timeReflection2-involution}
  \lean{Pphi2.timeReflection2_involution}\leanok
  \uses{Pphi2-planeBackground, Pphi2-EuclideanPlaneBackground-dim, Pphi2-SpaceTime2, Pphi2-timeReflection2}
  \texttt{Pphi2.timeReflection2\_involution}
\end{theorem}
\begin{theorem}[\texttt{compTimeReflection2\_apply}]\label{Pphi2-compTimeReflection2-apply}
  \lean{Pphi2.compTimeReflection2_apply}\leanok
  \uses{Pphi2-EuclideanPlaneBackground-SpaceTime, Pphi2-planeBackground, Pphi2-EuclideanPlaneBackground-dim, Pphi2-SpaceTime2, Pphi2-timeReflection2, Pphi2-TestFunction2, Pphi2-compTimeReflection2}
  \texttt{Pphi2.compTimeReflection2\_apply}
\end{theorem}
\begin{definition}[\texttt{euclideanAction2}]\label{Pphi2-euclideanAction2}
  \lean{Pphi2.euclideanAction2}\leanok
  \uses{Pphi2-EuclideanPlaneBackground-SpaceTime, Pphi2-E2, Pphi2-planeBackground, Pphi2-EuclideanPlaneBackground-dim, Pphi2-continuumEuclideanAction, Pphi2-TestFunction2}
  \texttt{Pphi2.euclideanAction2}
\end{definition}
\begin{definition}[\texttt{plane2TimeStructure}]\label{Pphi2-plane2TimeStructure}
  \lean{Pphi2.plane2TimeStructure}\leanok
  \uses{Pphi2-EuclideanPlaneBackground-SpaceTime, Pphi2-timeReflectionLinear, Pphi2-EuclideanPlaneBackground-dim, Pphi2-timeAxis2, Pphi2-plane2Background, Pphi2-EuclideanTimeStructure, Pphi2-timeReflection2-involution}
  \texttt{Pphi2.plane2TimeStructure}
\end{definition}
\begin{theorem}[\texttt{translate\_timeTranslationVector2\_eq\_timeTranslationSchwartz\_neg}]\label{Pphi2-translate-timeTranslationVector2-eq-timeTranslationSchwartz-neg}
  \lean{Pphi2.translate_timeTranslationVector2_eq_timeTranslationSchwartz_neg}\leanok
  \uses{Pphi2-EuclideanPlaneBackground-SpaceTime, Pphi2-planeBackground, Pphi2-SchwartzMap-translate, Pphi2-sub-timeTranslationVector2-eq-timeShift-neg, Pphi2-timeTranslationVector2, Pphi2-EuclideanPlaneBackground-dim, Pphi2-SpaceTime2, Pphi2-SchwartzMap-translate-apply, Pphi2-TestFunction2}
  \texttt{Pphi2.translate\_timeTranslationVector2\_eq\_timeTranslationSchwartz\_neg}
\end{theorem}
\begin{definition}[\texttt{compTimeReflection2}]\label{Pphi2-compTimeReflection2}
  \lean{Pphi2.compTimeReflection2}\leanok
  \uses{Pphi2-EuclideanPlaneBackground-SpaceTime, Pphi2-planeBackground, Pphi2-EuclideanPlaneBackground-dim, Pphi2-plane2Background, Pphi2-EuclideanTimeStructure-reflectOnRealTestFunctions, Pphi2-TestFunction2, Pphi2-plane2TimeStructure}
  \texttt{Pphi2.compTimeReflection2}
\end{definition}
\begin{definition}[\texttt{PositiveTimeTestFunction2}]\label{Pphi2-PositiveTimeTestFunction2}
  \lean{Pphi2.PositiveTimeTestFunction2}\leanok
  \uses{Pphi2-EuclideanPlaneBackground-SpaceTime, Pphi2-planeBackground, Pphi2-EuclideanPlaneBackground-dim, Pphi2-positiveTimeSubmodule2, Pphi2-TestFunction2}
  \texttt{Pphi2.PositiveTimeTestFunction2}
\end{definition}
\begin{theorem}[\texttt{timeTranslationVector2\_ofLp\_zero}]\label{Pphi2-timeTranslationVector2-ofLp-zero}
  \lean{Pphi2.timeTranslationVector2_ofLp_zero}\leanok
  \uses{Pphi2-EuclideanPlaneBackground-SpaceTime, Pphi2-timeReflectionLinear, Pphi2-EuclideanTimeStructure-mk-congr-simp, Pphi2-planeBackground, Pphi2-EuclideanTimeStructure-timeAxis, Pphi2-timeTranslationVector2, Pphi2-EuclideanPlaneBackground-dim, Pphi2-SpaceTime2, Pphi2-timeAxis2, Pphi2-plane2Background, Pphi2-EuclideanTimeStructure, Pphi2-timeReflection2-involution, Pphi2-plane2TimeStructure}
  \texttt{Pphi2.timeTranslationVector2\_ofLp\_zero}
\end{theorem}
\begin{definition}[\texttt{timeAxis2}]\label{Pphi2-timeAxis2}
  \lean{Pphi2.timeAxis2}\leanok
  \uses{Pphi2-planeBackground, Pphi2-EuclideanPlaneBackground-dim, Pphi2-SpaceTime2}
  \texttt{Pphi2.timeAxis2}
\end{definition}
\begin{definition}[\texttt{FieldConfig2}]\label{Pphi2-FieldConfig2}
  \lean{Pphi2.FieldConfig2}\leanok
  \uses{Pphi2-ContinuumFieldConfig}
  \texttt{Pphi2.FieldConfig2}
\end{definition}
\begin{definition}[\texttt{TestFunction2}]\label{Pphi2-TestFunction2}
  \lean{Pphi2.TestFunction2}\leanok
  \uses{Pphi2-ContinuumTestFunction}
  \texttt{Pphi2.TestFunction2}
\end{definition}
\begin{definition}[\texttt{TestFunction2ℂ}]\label{Pphi2-TestFunction2-}
  \lean{Pphi2.TestFunction2ℂ}\leanok
  \uses{Pphi2-ContinuumComplexTestFunction}
  \texttt{Pphi2.TestFunction2ℂ}
\end{definition}
\begin{definition}[\texttt{hasPositiveTime2}]\label{Pphi2-hasPositiveTime2}
  \lean{Pphi2.hasPositiveTime2}\leanok
  \uses{Pphi2-EuclideanPlaneBackground-SpaceTime, Pphi2-EuclideanTimeStructure-positiveTimeSet, Pphi2-SpaceTime2, Pphi2-plane2Background, Pphi2-plane2TimeStructure}
  \texttt{Pphi2.hasPositiveTime2}
\end{definition}
\begin{theorem}[\texttt{timeTranslationVector2\_ofLp\_one}]\label{Pphi2-timeTranslationVector2-ofLp-one}
  \lean{Pphi2.timeTranslationVector2_ofLp_one}\leanok
  \uses{Pphi2-EuclideanPlaneBackground-SpaceTime, Pphi2-timeReflectionLinear, Pphi2-EuclideanTimeStructure-mk-congr-simp, Pphi2-planeBackground, Pphi2-EuclideanTimeStructure-timeAxis, Pphi2-timeTranslationVector2, Pphi2-EuclideanPlaneBackground-dim, Pphi2-SpaceTime2, Pphi2-timeAxis2, Pphi2-plane2Background, Pphi2-EuclideanTimeStructure, Pphi2-timeReflection2-involution, Pphi2-plane2TimeStructure}
  \texttt{Pphi2.timeTranslationVector2\_ofLp\_one}
\end{theorem}
\begin{definition}[\texttt{timeReflection2}]\label{Pphi2-timeReflection2}
  \lean{Pphi2.timeReflection2}\leanok
  \uses{Pphi2-planeBackground, Pphi2-EuclideanPlaneBackground-dim, Pphi2-SpaceTime2}
  Time reflection $(t, x) \mapsto (-t, x)$ and its pullback on Schwartz space.
\end{definition}
\begin{definition}[\texttt{O2}]\label{Pphi2-O2}
  \lean{Pphi2.O2}\leanok
  \uses{Pphi2-ContinuumOrthogonalGroup}
  \texttt{Pphi2.O2}
\end{definition}
\begin{definition}[\texttt{translate}]\label{Pphi2-SchwartzMap-translate}
  \lean{Pphi2.SchwartzMap.translate}\leanok
  \uses{Pphi2-EuclideanPlaneBackground-SpaceTime, Pphi2-planeBackground, Pphi2-EuclideanPlaneBackground-translate, Pphi2-EuclideanPlaneBackground-dim, Pphi2-SpaceTime2, Pphi2-plane2Background, Pphi2-TestFunction2}
  \texttt{Pphi2.SchwartzMap.translate}
\end{definition}
\begin{definition}[\texttt{E2}]\label{Pphi2-E2}
  \lean{Pphi2.E2}\leanok
  \uses{Pphi2-ContinuumMotion}
  Euclidean motion $(R, t)$ with action $g \cdot x = R(x) + t$; pullback on test functions via \texttt{compCLMOfAntilipschitz}.
\end{definition}
\begin{definition}[\texttt{timeReflectionLinear}]\label{Pphi2-timeReflectionLinear}
  \lean{Pphi2.timeReflectionLinear}\leanok
  \uses{Pphi2-planeBackground, Pphi2-EuclideanPlaneBackground-dim, Pphi2-SpaceTime2, Pphi2-timeReflection2}
  \texttt{Pphi2.timeReflectionLinear}
\end{definition}
\begin{theorem}[\texttt{congr\_simp}]\label{Pphi2-EuclideanTimeStructure-mk-congr-simp}
  \lean{Pphi2.EuclideanTimeStructure.mk.congr_simp}\leanok
  \uses{Pphi2-EuclideanPlaneBackground-SpaceTime, Pphi2-EuclideanPlaneBackground-dim, Pphi2-EuclideanPlaneBackground, Pphi2-EuclideanTimeStructure}
  \texttt{Pphi2.EuclideanTimeStructure.mk.congr\_simp}
\end{theorem}
\begin{theorem}[\texttt{sub\_timeTranslationVector2\_eq\_timeShift\_neg}]\label{Pphi2-sub-timeTranslationVector2-eq-timeShift-neg}
  \lean{Pphi2.sub_timeTranslationVector2_eq_timeShift_neg}\leanok
  \uses{Pphi2-planeBackground, Pphi2-timeTranslationVector2, Pphi2-EuclideanPlaneBackground-dim, Pphi2-SpaceTime2, Pphi2-timeTranslationVector2-ofLp-one, Pphi2-timeTranslationVector2-ofLp-zero}
  \texttt{Pphi2.sub\_timeTranslationVector2\_eq\_timeShift\_neg}
\end{theorem}
\begin{definition}[\texttt{SpaceTime2}]\label{Pphi2-SpaceTime2}
  \lean{Pphi2.SpaceTime2}\leanok
  \uses{Pphi2-ContinuumSpaceTime}
  \texttt{Pphi2.SpaceTime2}
\end{definition}
\begin{definition}[\texttt{euclideanAction2ℂ}]\label{Pphi2-euclideanAction2-}
  \lean{Pphi2.euclideanAction2ℂ}\leanok
  \uses{Pphi2-EuclideanPlaneBackground-SpaceTime, Pphi2-E2, Pphi2-planeBackground, Pphi2-EuclideanPlaneBackground-dim, Pphi2-continuumEuclideanActionComplex, Pphi2-TestFunction2-}
  \texttt{Pphi2.euclideanAction2ℂ}
\end{definition}
\begin{definition}[\texttt{act}]\label{Pphi2-E2-act}
  \lean{Pphi2.E2.act}\leanok
  \uses{Pphi2-EuclideanPlaneBackground-SpaceTime, Pphi2-E2, Pphi2-EuclideanPlaneBackground-Motion-t, Pphi2-EuclideanPlaneBackground-Motion-R, Pphi2-planeBackground, Pphi2-EuclideanPlaneBackground-dim, Pphi2-SpaceTime2}
  \texttt{Pphi2.E2.act}
\end{definition}
\begin{definition}[\texttt{timeTranslationVector2}]\label{Pphi2-timeTranslationVector2}
  \lean{Pphi2.timeTranslationVector2}\leanok
  \uses{Pphi2-EuclideanTimeStructure-timeTranslation, Pphi2-SpaceTime2, Pphi2-plane2Background, Pphi2-plane2TimeStructure}
  \texttt{Pphi2.timeTranslationVector2}
\end{definition}
\begin{definition}[\texttt{positiveTimeSubmodule2}]\label{Pphi2-positiveTimeSubmodule2}
  \lean{Pphi2.positiveTimeSubmodule2}\leanok
  \uses{Pphi2-EuclideanPlaneBackground-SpaceTime, Pphi2-planeBackground, Pphi2-EuclideanPlaneBackground-dim, Pphi2-EuclideanTimeStructure-positiveTimeSubmodule, Pphi2-plane2Background, Pphi2-TestFunction2, Pphi2-plane2TimeStructure}
  \texttt{Pphi2.positiveTimeSubmodule2}
\end{definition}
\begin{theorem}[\texttt{translate\_apply}]\label{Pphi2-SchwartzMap-translate-apply}
  \lean{Pphi2.SchwartzMap.translate_apply}\leanok
  \uses{Pphi2-EuclideanPlaneBackground-SpaceTime, Pphi2-planeBackground, Pphi2-SchwartzMap-translate, Pphi2-EuclideanPlaneBackground-dim, Pphi2-SpaceTime2, Pphi2-plane2Background, Pphi2-TestFunction2, Pphi2-EuclideanPlaneBackground-translate-apply}
  \texttt{Pphi2.SchwartzMap.translate\_apply}
\end{theorem}
\begin{definition}[\texttt{timeReflectionCLE}]\label{Pphi2-timeReflectionCLE}
  \lean{Pphi2.timeReflectionCLE}\leanok
  \uses{Pphi2-planeBackground, Pphi2-EuclideanPlaneBackground-dim, Pphi2-SpaceTime2, Pphi2-plane2Background, Pphi2-EuclideanTimeStructure-reflectCLE, Pphi2-plane2TimeStructure}
  \texttt{Pphi2.timeReflectionCLE}
\end{definition}
\section{\texttt{Pphi2.InteractingMeasure.UniformBounds}}
\begin{theorem}[\texttt{torus\_volume\_eq}]\label{Pphi2-torus-volume-eq}
  \lean{Pphi2.torus_volume_eq}\leanok
  \uses{Pphi2-lattice-volume-eq, GaussianField-FinLatticeSites, GaussianField-circleSpacing-eq, GaussianField-circleSpacing}
  \texttt{Pphi2.torus\_volume\_eq}
\end{theorem}
\begin{theorem}[\texttt{boltzmann\_cauchySchwarz}]\label{Pphi2-boltzmann-cauchySchwarz}
  \lean{Pphi2.boltzmann_cauchySchwarz}\leanok
  \uses{Pphi2-interactionFunctional-bounded-below, GaussianField-FinLatticeSites, Pphi2-interactionFunctional-measurable, GaussianField-latticeGaussianMeasure-isProbability, GaussianField-latticeGaussianMeasure, GaussianField-Configuration, GaussianField-instMeasurableSpaceConfiguration, GaussianField-FinLatticeField, Pphi2-interactionFunctional, Lattice-toSemilatticeInf}
  \texttt{Pphi2.boltzmann\_cauchySchwarz}
\end{theorem}
\begin{theorem}[\texttt{exists\_uniform\_neg\_of\_uniform\_affine\_bound'}]\label{Pphi2-exists-uniform-neg-of-uniform-affine-bound-}
  \lean{Pphi2.exists_uniform_neg_of_uniform_affine_bound'}\leanok
  \uses{Pphi2-deriv-affine-bound-neg-of-continuousOn}
  \texttt{Pphi2.exists\_uniform\_neg\_of\_uniform\_affine\_bound'}
\end{theorem}
\begin{theorem}[\texttt{exists\_uniform\_neg\_of\_uniform\_affine\_bound}]\label{Pphi2-exists-uniform-neg-of-uniform-affine-bound}
  \lean{Pphi2.exists_uniform_neg_of_uniform_affine_bound}\leanok
  \uses{Pphi2-deriv-affine-bound-neg}
  \texttt{Pphi2.exists\_uniform\_neg\_of\_uniform\_affine\_bound}
\end{theorem}
\begin{theorem}[\texttt{deriv\_affine\_bound\_neg\_of\_continuousOn}]\label{Pphi2-deriv-affine-bound-neg-of-continuousOn}
  \lean{Pphi2.deriv_affine_bound_neg_of_continuousOn}\leanok
  \texttt{Pphi2.deriv\_affine\_bound\_neg\_of\_continuousOn}
\end{theorem}
\begin{theorem}[\texttt{interacting\_moment\_le\_L2\_of\_expBound}]\label{Pphi2-interacting-moment-le-L2-of-expBound}
  \lean{Pphi2.interacting_moment_le_L2_of_expBound}\leanok
  \uses{Pphi2-partitionFn-ge-one, GaussianField-FinLatticeSites, GaussianField-latticeGaussianMeasure, GaussianField-Configuration, Pphi2-expMoment-le-rpow, GaussianField-instMeasurableSpaceConfiguration, Pphi2-boltzmann-cauchySchwarz, GaussianField-FinLatticeField, Pphi2-interactionFunctional, Lattice-toSemilatticeInf}
  \texttt{Pphi2.interacting\_moment\_le\_L2\_of\_expBound}
\end{theorem}
\begin{theorem}[\texttt{latticeSites\_card}]\label{Pphi2-latticeSites-card}
  \lean{Pphi2.latticeSites_card}\leanok
  \uses{GaussianField-FinLatticeSites}
  \texttt{Pphi2.latticeSites\_card}
\end{theorem}
\begin{theorem}[\texttt{partitionFn\_ge\_one}]\label{Pphi2-partitionFn-ge-one}
  \lean{Pphi2.partitionFn_ge_one}\leanok
  \uses{Pphi2-interactionFunctional-bounded-below, GaussianField-finLatticeDelta, GaussianField-FinLatticeSites, Pphi2-interactionFunctional-measurable, GaussianField-latticeGaussianMeasure-isProbability, GaussianField-latticeGaussianMeasure, GaussianField-Configuration, Pphi2-wickConstant, GaussianField-instMeasurableSpaceConfiguration, Pphi2-interactionFunctional-integral-nonpos, GaussianField-FinLatticeField, Pphi2-interactionFunctional, Pphi2-wickPolynomial-integral-eq-coeff-zero, Pphi2-wickPolynomial}
  \texttt{Pphi2.partitionFn\_ge\_one}
\end{theorem}
\begin{theorem}[\texttt{expMoment\_le\_rpow}]\label{Pphi2-expMoment-le-rpow}
  \lean{Pphi2.expMoment_le_rpow}\leanok
  \uses{Pphi2-interactionFunctional-bounded-below, GaussianField-FinLatticeSites, Pphi2-interactionFunctional-measurable, GaussianField-latticeGaussianMeasure-isProbability, GaussianField-latticeGaussianMeasure, GaussianField-Configuration, GaussianField-instMeasurableSpaceConfiguration, GaussianField-FinLatticeField, Pphi2-interactionFunctional}
  \texttt{Pphi2.expMoment\_le\_rpow}
\end{theorem}
\begin{theorem}[\texttt{lattice\_volume\_eq}]\label{Pphi2-lattice-volume-eq}
  \lean{Pphi2.lattice_volume_eq}\leanok
  \uses{Pphi2-latticeSites-card, GaussianField-FinLatticeSites}
  \texttt{Pphi2.lattice\_volume\_eq}
\end{theorem}
\begin{theorem}[\texttt{interactionFunctional\_integral\_nonpos}]\label{Pphi2-interactionFunctional-integral-nonpos}
  \lean{Pphi2.interactionFunctional_integral_nonpos}\leanok
  \uses{GaussianField-finLatticeDelta, GaussianField-FinLatticeSites, GaussianField-latticeGaussianMeasure, GaussianField-Configuration, Pphi2-wickConstant, GaussianField-instMeasurableSpaceConfiguration, GaussianField-FinLatticeField, Pphi2-interactionFunctional, Pphi2-wickPolynomial-integral-eq-coeff-zero, Pphi2-wickPolynomial}
  \texttt{Pphi2.interactionFunctional\_integral\_nonpos}
\end{theorem}
\section{\texttt{Pphi2.NelsonEstimate.AsymRoughCovarianceHigherP}}
\begin{theorem}[\texttt{asymCanonicalRoughCovariance\_pow\_sum\_le\_uniform\_of\_three\_le}]\label{Pphi2-asymCanonicalRoughCovariance-pow-sum-le-uniform-of-three-le}
  \lean{Pphi2.asymCanonicalRoughCovariance_pow_sum_le_uniform_of_three_le}\leanok
  \uses{GaussianField-AsymLatticeSites, GaussianField-instFintypeAsymLatticeSitesOfNeZeroNat, Lattice-toSemilatticeInf, Pphi2-asymCanonicalRoughCovariance}
  \texttt{Pphi2.asymCanonicalRoughCovariance\_pow\_sum\_le\_uniform\_of\_three\_le}
\end{theorem}
\section{\texttt{Pphi2.IRLimit.CylinderEmbedding}}
\begin{definition}[\texttt{cylinderPullback}]\label{Pphi2-cylinderPullback}
  \lean{Pphi2.cylinderPullback}\leanok
  \uses{GaussianField-NuclearTensorProduct-instModuleReal, GaussianField-NuclearTensorProduct-instTopologicalSpace, GaussianField-CylinderTestFunction, GaussianField-NuclearTensorProduct-instIsTopologicalAddGroup, Pphi2-cylinderToTorusEmbed, GaussianField-NuclearTensorProduct-instAddCommGroup, GaussianField-Configuration, GaussianField-NuclearTensorProduct-instContinuousSMulReal, GaussianField-SmoothMap-Circle, GaussianField-AsymTorusTestFunction}
  ``\texttt{lean def cylinderPullback (ω : Configuration (AsymTorusTestFunction Lt Ls)) : Configuration (CylinderTestFunction Ls) }``  Pullback on configurations: given a torus configuration $\omega$, produces a cylinder configuration by precomposing with the embedding.
\end{definition}
\begin{theorem}[\texttt{cylinderToTorusEmbed\_pure}]\label{Pphi2-cylinderToTorusEmbed-pure}
  \lean{Pphi2.cylinderToTorusEmbed_pure}\leanok
  \uses{GaussianField-NuclearTensorProduct-instModuleReal, GaussianField-NuclearTensorProduct-instTopologicalSpace, GaussianField-nuclearTensorProduct-mapCLM-general-pure, GaussianField-schwartz-dyninMityaginSpace-1D, GaussianField-nuclearTensorProduct-mapCLM-general, GaussianField-SmoothMap-Circle-instAddCommGroup, GaussianField-CylinderTestFunction, GaussianField-nuclearTensorProduct-swapCLM, GaussianField-NuclearTensorProduct, Pphi2-cylinderToTorusEmbed, GaussianField-nuclearTensorProduct-swapCLM-pure, GaussianField-NuclearTensorProduct-instAddCommGroup, GaussianField-periodizeCLM, GaussianField-SmoothMap-Circle-instModule, GaussianField-SmoothMap-Circle-instContinuousSMul, GaussianField-smoothCircle-dyninMityaginSpace, GaussianField-SmoothMap-Circle-instIsTopologicalAddGroup, GaussianField-NuclearTensorProduct-pure, GaussianField-SmoothMap-Circle, GaussianField-AsymTorusTestFunction, GaussianField-SmoothMap-Circle-instTopologicalSpace}
  \texttt{Pphi2.cylinderToTorusEmbed\_pure}
\end{theorem}
\begin{theorem}[\texttt{cylinderToTorusEmbed\_comp\_timeReflection}]\label{Pphi2-cylinderToTorusEmbed-comp-timeReflection}
  \lean{Pphi2.cylinderToTorusEmbed_comp_timeReflection}\leanok
  \uses{GaussianField-NuclearTensorProduct-instModuleReal, GaussianField-NuclearTensorProduct-instTopologicalSpace, GaussianField-nuclearTensorProduct-mapCLM-general-pure, GaussianField-nuclearTensorProduct-mapCLM, GaussianField-schwartz-dyninMityaginSpace-1D, GaussianField-nuclearTensorProduct-mapCLM-general, GaussianField-SmoothMap-Circle-instAddCommGroup, GaussianField-CylinderTestFunction, GaussianField-nuclearTensorProduct-mapCLM-pure, GaussianField-nuclearTensorProduct-swapCLM, GaussianField-NuclearTensorProduct, Pphi2-cylinderToTorusEmbed, GaussianField-nuclearTensorProduct-swapCLM-pure, GaussianField-periodizeCLM-comp-schwartzReflection, GaussianField-NuclearTensorProduct-instAddCommGroup, GaussianField-cylinderTimeReflection, GaussianField-periodizeCLM, GaussianField-circleReflection, GaussianField-SmoothMap-Circle-instModule, GaussianField-SmoothMap-Circle-instContinuousSMul, GaussianField-smoothCircle-dyninMityaginSpace, GaussianField-SmoothMap-Circle-instIsTopologicalAddGroup, GaussianField-NuclearTensorProduct-pure, GaussianField-SmoothMap-Circle, GaussianField-AsymTorusTestFunction, GaussianField-schwartzReflection, GaussianField-SmoothMap-Circle-instTopologicalSpace, GaussianField-asymTorusTimeReflection}
  ``\texttt{lean theorem cylinderToTorusEmbed\_comp\_timeReflection (f : CylinderTestFunction Ls) : cylinderToTorusEmbed Lt Ls (cylinderTimeReflection Ls f) = asymTorusTimeReflection Lt Ls (cylinderToTorusEmbed Lt Ls f) }``  Informal statement: The embedding intertwines cylinder time reflection with torus time reflection: $\mathrm{embed}(\Theta f) = \Theta_{\mathrm{torus}}(\mathrm{embed}\, f)$.  Proof: Fully proved. Same strategy as time translation, using \texttt{periodizeCLM\_comp\_schwartzReflection}.  Dependencies: \texttt{cylinderToTorus\_clm\_ext\_of\_pure}, \texttt{nuclearTensorProduct\_mapCLM\_pure}, \texttt{nuclearTensorProduct\_swapCLM\_pure}, \texttt{nuclearTensorProduct\_mapCLM\_general\_pure}, \texttt{periodizeCLM\_comp\_schwartzReflection}.
\end{theorem}
\begin{theorem}[\texttt{cylinderToTorusEmbed\_comp\_timeTranslation}]\label{Pphi2-cylinderToTorusEmbed-comp-timeTranslation}
  \lean{Pphi2.cylinderToTorusEmbed_comp_timeTranslation}\leanok
  \uses{GaussianField-NuclearTensorProduct-instModuleReal, GaussianField-circleTranslation, GaussianField-NuclearTensorProduct-instTopologicalSpace, GaussianField-nuclearTensorProduct-mapCLM-general-pure, GaussianField-nuclearTensorProduct-mapCLM, GaussianField-schwartz-dyninMityaginSpace-1D, GaussianField-nuclearTensorProduct-mapCLM-general, GaussianField-SmoothMap-Circle-instAddCommGroup, GaussianField-CylinderTestFunction, GaussianField-nuclearTensorProduct-mapCLM-pure, GaussianField-nuclearTensorProduct-swapCLM, GaussianField-asymTorusTranslation, GaussianField-NuclearTensorProduct, Pphi2-cylinderToTorusEmbed, GaussianField-nuclearTensorProduct-swapCLM-pure, GaussianField-NuclearTensorProduct-instAddCommGroup, GaussianField-periodizeCLM, GaussianField-SmoothMap-Circle-instModule, GaussianField-SmoothMap-Circle-instContinuousSMul, GaussianField-smoothCircle-dyninMityaginSpace, GaussianField-SmoothMap-Circle-instIsTopologicalAddGroup, GaussianField-NuclearTensorProduct-pure, GaussianField-SmoothMap-Circle, GaussianField-periodizeCLM-comp-schwartzTranslation, GaussianField-schwartzTranslation, GaussianField-AsymTorusTestFunction, GaussianField-cylinderTranslation, GaussianField-SmoothMap-Circle-instTopologicalSpace}
  ``\texttt{lean theorem cylinderToTorusEmbed\_comp\_timeTranslation (τ : ℝ) (f : CylinderTestFunction Ls) : cylinderToTorusEmbed Lt Ls (cylinderTranslation Ls 0 τ f) = asymTorusTranslation Lt Ls (τ, 0) (cylinderToTorusEmbed Lt Ls f) }``  Informal statement: The embedding intertwines cylinder time translation with torus time translation: $\mathrm{embed}(T_{0,\tau} f) = T_{\tau,0}(\mathrm{embed}\, f)$.  Proof: Fully proved. Reduces to pure tensors via \texttt{cylinderToTorus\_clm\_ext\_of\_pure}, then uses \texttt{periodizeCLM\_comp\_schwartzTranslation} (periodization commutes with Schwartz translation).  Dependencies: \texttt{cylinderToTorus\_clm\_ext\_of\_pure}, \texttt{nuclearTensorProduct\_mapCLM\_pure}, \texttt{nuclearTensorProduct\_swapCLM\_pure}, \texttt{nuclearTensorProduct\_mapCLM\_general\_pure}, \texttt{periodizeCLM\_comp\_schwartzTranslation}.
\end{theorem}
\begin{definition}[\texttt{cylinderPullbackMeasure}]\label{Pphi2-cylinderPullbackMeasure}
  \lean{Pphi2.cylinderPullbackMeasure}\leanok
  \uses{GaussianField-NuclearTensorProduct-instModuleReal, GaussianField-NuclearTensorProduct-instTopologicalSpace, Pphi2-cylinderPullback, GaussianField-CylinderTestFunction, GaussianField-NuclearTensorProduct-instIsTopologicalAddGroup, GaussianField-NuclearTensorProduct-instAddCommGroup, GaussianField-Configuration, GaussianField-NuclearTensorProduct-instContinuousSMulReal, GaussianField-SmoothMap-Circle, GaussianField-instMeasurableSpaceConfiguration, GaussianField-AsymTorusTestFunction}
  ``\texttt{lean def cylinderPullbackMeasure (μ : Measure (Configuration (AsymTorusTestFunction Lt Ls))) : Measure (Configuration (CylinderTestFunction Ls)) }``  The pushforward (pullback on test functions) measure on cylinder configurations: $\nu = \mu \circ (\mathrm{cylinderPullback})^{-1}$.
\end{definition}
\begin{theorem}[\texttt{cylinderPullback\_continuous}]\label{Pphi2-cylinderPullback-continuous}
  \lean{Pphi2.cylinderPullback_continuous}\leanok
  \uses{GaussianField-NuclearTensorProduct-instModuleReal, GaussianField-NuclearTensorProduct-instTopologicalSpace, Pphi2-cylinderPullback, GaussianField-CylinderTestFunction, GaussianField-NuclearTensorProduct-instIsTopologicalAddGroup, Pphi2-cylinderToTorusEmbed, GaussianField-NuclearTensorProduct-instAddCommGroup, GaussianField-Configuration, GaussianField-NuclearTensorProduct-instContinuousSMulReal, GaussianField-SmoothMap-Circle, GaussianField-AsymTorusTestFunction}
  ``\texttt{lean theorem cylinderPullback\_continuous : Continuous (cylinderPullback Lt Ls) }`\texttt{  Informal statement: The pullback map is continuous (hence measurable).  Proof: Fully proved via }WeakDual.continuous\_of\_continuous\_eval`.
\end{theorem}
\begin{theorem}[\texttt{cylinderPullback\_eval}]\label{Pphi2-cylinderPullback-eval}
  \lean{Pphi2.cylinderPullback_eval}\leanok
  \uses{GaussianField-NuclearTensorProduct-instModuleReal, GaussianField-NuclearTensorProduct-instTopologicalSpace, Pphi2-cylinderPullback, GaussianField-CylinderTestFunction, GaussianField-NuclearTensorProduct-instIsTopologicalAddGroup, Pphi2-cylinderToTorusEmbed, GaussianField-NuclearTensorProduct-instAddCommGroup, GaussianField-Configuration, GaussianField-NuclearTensorProduct-instContinuousSMulReal, GaussianField-SmoothMap-Circle, GaussianField-AsymTorusTestFunction}
  ``\texttt{lean @[simp] theorem cylinderPullback\_eval (ω : Configuration (AsymTorusTestFunction Lt Ls)) (f : CylinderTestFunction Ls) : (cylinderPullback Lt Ls ω) f = ω (cylinderToTorusEmbed Lt Ls f) }``  Informal statement: Evaluation of the pullback configuration on $f$ equals evaluation of $\omega$ on the embedded $f$. Definitional (\texttt{rfl}).
\end{theorem}
\begin{theorem}[\texttt{cylinderToTorusEmbed\_comp\_spatialTranslation}]\label{Pphi2-cylinderToTorusEmbed-comp-spatialTranslation}
  \lean{Pphi2.cylinderToTorusEmbed_comp_spatialTranslation}\leanok
  \uses{GaussianField-DyninMityaginSpace, GaussianField-NuclearTensorProduct-instModuleReal, GaussianField-circleTranslation, GaussianField-NuclearTensorProduct-instTopologicalSpace, GaussianField-nuclearTensorProduct-mapCLM-general-pure, GaussianField-cylinderSpatialTranslation, GaussianField-nuclearTensorProduct-mapCLM, GaussianField-schwartz-dyninMityaginSpace-1D, GaussianField-nuclearTensorProduct-mapCLM-general, GaussianField-SmoothMap-Circle-instAddCommGroup, GaussianField-CylinderTestFunction, GaussianField-nuclearTensorProduct-mapCLM-pure, GaussianField-nuclearTensorProduct-swapCLM, GaussianField-asymTorusTranslation, GaussianField-NuclearTensorProduct, Pphi2-cylinderToTorusEmbed, GaussianField-nuclearTensorProduct-swapCLM-pure, GaussianField-NuclearTensorProduct-instAddCommGroup, GaussianField-periodizeCLM, GaussianField-SmoothMap-Circle-instModule, GaussianField-SmoothMap-Circle-instContinuousSMul, GaussianField-smoothCircle-dyninMityaginSpace, GaussianField-SmoothMap-Circle-instIsTopologicalAddGroup, GaussianField-NuclearTensorProduct-pure, GaussianField-SmoothMap-Circle, GaussianField-circleTranslation-zero, GaussianField-AsymTorusTestFunction, GaussianField-SmoothMap-Circle-instTopologicalSpace}
  \texttt{Pphi2.cylinderToTorusEmbed\_comp\_spatialTranslation}
\end{theorem}
\begin{theorem}[\texttt{cylinderPullbackMeasure\_integral\_sq}]\label{Pphi2-cylinderPullbackMeasure-integral-sq}
  \lean{Pphi2.cylinderPullbackMeasure_integral_sq}\leanok
  \uses{GaussianField-NuclearTensorProduct-instModuleReal, GaussianField-NuclearTensorProduct-instTopologicalSpace, Pphi2-cylinderPullback, GaussianField-CylinderTestFunction, GaussianField-NuclearTensorProduct-instIsTopologicalAddGroup, Pphi2-cylinderToTorusEmbed, GaussianField-configuration-measurable-of-eval-measurable, GaussianField-NuclearTensorProduct-instAddCommGroup, Pphi2-cylinderPullbackMeasure, GaussianField-configuration-eval-measurable, GaussianField-Configuration, GaussianField-NuclearTensorProduct-instContinuousSMulReal, GaussianField-SmoothMap-Circle, GaussianField-instMeasurableSpaceConfiguration, GaussianField-AsymTorusTestFunction}
  \texttt{Pphi2.cylinderPullbackMeasure\_integral\_sq}
\end{theorem}
\begin{definition}[\texttt{cylinderToTorusEmbed}]\label{Pphi2-cylinderToTorusEmbed}
  \lean{Pphi2.cylinderToTorusEmbed}\leanok
  \uses{GaussianField-NuclearTensorProduct-instModuleReal, GaussianField-NuclearTensorProduct-instTopologicalSpace, GaussianField-schwartz-dyninMityaginSpace-1D, GaussianField-nuclearTensorProduct-mapCLM-general, GaussianField-SmoothMap-Circle-instAddCommGroup, GaussianField-CylinderTestFunction, GaussianField-nuclearTensorProduct-swapCLM, GaussianField-NuclearTensorProduct, GaussianField-NuclearTensorProduct-instAddCommGroup, GaussianField-periodizeCLM, GaussianField-SmoothMap-Circle-instModule, GaussianField-SmoothMap-Circle-instContinuousSMul, GaussianField-smoothCircle-dyninMityaginSpace, GaussianField-SmoothMap-Circle-instIsTopologicalAddGroup, GaussianField-SmoothMap-Circle, GaussianField-AsymTorusTestFunction, GaussianField-SmoothMap-Circle-instTopologicalSpace}
  ``\texttt{lean def cylinderToTorusEmbed : CylinderTestFunction Ls →L[ℝ] AsymTorusTestFunction Lt Ls }``  Embeds cylinder test functions into asymmetric torus test functions by applying $\mathrm{id} \otimes \mathrm{periodize}$ followed by the NTP swap map.
\end{definition}
\section{\texttt{Pphi2.GaussianContinuumLimit.PropagatorConvergence}}
\begin{theorem}[\texttt{continuumGreenBilinear\_le\_mass\_inv\_sq}]\label{Pphi2-continuumGreenBilinear-le-mass-inv-sq}
  \lean{Pphi2.continuumGreenBilinear_le_mass_inv_sq}\leanok
  \uses{Pphi2-EuclideanPlaneBackground-SpaceTime, Pphi2-planeBackground, Pphi2-EuclideanPlaneBackground-dim, Pphi2-continuumGreenBilinear, Pphi2-ContinuumTestFunction, Lattice-toSemilatticeInf}
  \texttt{Pphi2.continuumGreenBilinear\_le\_mass\_inv\_sq}
\end{theorem}
\begin{theorem}[\texttt{latticeGreenBilinear\_basis\_tendsto\_continuum} (axiom)]\label{Pphi2-latticeGreenBilinear-basis-tendsto-continuum}
  \lean{Pphi2.latticeGreenBilinear_basis_tendsto_continuum}\leanok
  \uses{Pphi2-EuclideanPlaneBackground-SpaceTime, Pphi2-planeBackground, Pphi2-EuclideanPlaneBackground-dim, GaussianField-DyninMityaginSpace-basis, Pphi2-latticeGreenBilinear, Pphi2-continuumGreenBilinear, Pphi2-ContinuumTestFunction}
  \texttt{Pphi2.latticeGreenBilinear\_basis\_tendsto\_continuum}
\end{theorem}
\begin{theorem}[\texttt{propagator\_convergence} (axiom)]\label{Pphi2-propagator-convergence}
  \lean{Pphi2.propagator_convergence}\leanok
  \uses{Pphi2-latticeGreenBilinear-tendsto-continuum, Pphi2-embeddedTwoPoint-eq-latticeGreenBilinear, Pphi2-embeddedTwoPoint, Pphi2-latticeGreenBilinear, Pphi2-continuumGreenBilinear, Pphi2-ContinuumTestFunction}
  For $a_n \to 0$ with $N_n a_n \to \infty$: $\text{embeddedTwoPoint}(f,g) \to G(f,g)$. The lattice Riemann sum converges to the Fourier integral.
\end{theorem}
\begin{theorem}[\texttt{continuumGreenBilinear\_pos}]\label{Pphi2-continuumGreenBilinear-pos}
  \lean{Pphi2.continuumGreenBilinear_pos}\leanok
  \uses{Pphi2-EuclideanPlaneBackground-SpaceTime, Pphi2-planeBackground, Pphi2-EuclideanPlaneBackground-dim, Pphi2-continuumGreenBilinear, Pphi2-ContinuumTestFunction, Lattice-toSemilatticeInf}
  $G(f,f) > 0$ for $f \ne 0$. Proved since the integrand $|f(k)|^2/(|k|^2+m^2)$ is nonneg, and strictly positive on a set of positive measure (injectivity of Fourier transform on $\mathcal{S}$).
\end{theorem}
\begin{theorem}[\texttt{latticeGreenBilinear\_tendsto\_continuum}]\label{Pphi2-latticeGreenBilinear-tendsto-continuum}
  \lean{Pphi2.latticeGreenBilinear_tendsto_continuum}\leanok
  \uses{Pphi2-EuclideanPlaneBackground-SpaceTime, Pphi2-latticeGreenBilinear-basis-tendsto-continuum, Pphi2-planeBackground, Pphi2-EuclideanPlaneBackground-dim, GaussianField-tendsto-of-symmetric-basis-tendsto, GaussianField-DyninMityaginSpace-basis, Pphi2-latticeGreenBilinear, Pphi2-continuumGreenBilinear, Pphi2-ContinuumTestFunction}
  \texttt{Pphi2.latticeGreenBilinear\_tendsto\_continuum}
\end{theorem}
\begin{theorem}[\texttt{embeddedTwoPoint\_uniform\_bound}]\label{Pphi2-embeddedTwoPoint-uniform-bound}
  \lean{Pphi2.embeddedTwoPoint_uniform_bound}\leanok
  \uses{Pphi2-EuclideanPlaneBackground-SpaceTime, GaussianField-finLatticeField-dyninMityaginSpace, Pphi2-latticeEmbed, GaussianField-measure, GaussianField-ell2-separableSpace, GaussianField-second-moment-eq-covariance, Pphi2-planeBackground, GaussianField-covariance, GaussianField-FinLatticeSites, Pphi2-evalAtSite, GaussianField-latticeCovarianceGJ, Pphi2-EuclideanPlaneBackground-dim, Pphi2-embeddedTwoPoint, GaussianField-latticeGaussianMeasure, GaussianField-Configuration, GaussianField-instMeasurableSpaceConfiguration, Pphi2-ContinuumTestFunction, GaussianField-ell2-, GaussianField-FinLatticeField, Pphi2-embeddedTwoPoint-eq-covariance}
  $E[\Phi_a(f)^2] \le m^{-2} C_f$ uniformly in $a$. Proved by combining \texttt{covariance\_le\_mass\_inv\_sq\_norm} with \texttt{schwartz\_riemann\_sum\_bound}.
\end{theorem}
\section{\texttt{Pphi2.NelsonEstimate.AsymCrossTermL2Identity}}
\begin{theorem}[\texttt{congr\_simp}]\label{Pphi2-asymCanonicalSmoothFieldFunctionOfFst-congr-simp}
  \lean{Pphi2.asymCanonicalSmoothFieldFunctionOfFst.congr_simp}\leanok
  \uses{GaussianField-AsymLatticeSites, Pphi2-asymCanonicalSmoothFieldFunctionOfFst, Pphi2-AsymCanonicalMode}
  \texttt{Pphi2.asymCanonicalSmoothFieldFunctionOfFst.congr\_simp}
\end{theorem}
\begin{definition}[\texttt{asymCanonicalRoughGamma}]\label{Pphi2-asymCanonicalRoughGamma}
  \lean{Pphi2.asymCanonicalRoughGamma}\leanok
  \uses{GaussianField-AsymLatticeSites, Pphi2-asymCanonicalRoughModeCoeff, Pphi2-AsymCanonicalMode}
  \texttt{Pphi2.asymCanonicalRoughGamma}
\end{definition}
\begin{theorem}[\texttt{congr\_simp}]\label{Pphi2-asymCanonicalSmoothCovariance-congr-simp}
  \lean{Pphi2.asymCanonicalSmoothCovariance.congr_simp}\leanok
  \uses{Pphi2-asymCanonicalSmoothCovariance, GaussianField-AsymLatticeSites}
  \texttt{Pphi2.asymCanonicalSmoothCovariance.congr\_simp}
\end{theorem}
\begin{definition}[\texttt{asymCanonicalCrossTerm}]\label{Pphi2-asymCanonicalCrossTerm}
  \lean{Pphi2.asymCanonicalCrossTerm}\leanok
  \uses{GaussianField-AsymLatticeSites, Pphi2-asymCanonicalRoughWickConstant, Pphi2-AsymCanonicalJoint, Pphi2-wickMonomial, Pphi2-asymCanonicalRoughFieldFunction, Pphi2-asymSmoothWickConstant, GaussianField-instFintypeAsymLatticeSitesOfNeZeroNat, Pphi2-asymCanonicalSmoothFieldFunction}
  \texttt{Pphi2.asymCanonicalCrossTerm}
\end{definition}
\begin{theorem}[\texttt{asymCanonicalRoughWickPower\_two\_site\_marginal\_eq\_zero\_of\_ne}]\label{Pphi2-asymCanonicalRoughWickPower-two-site-marginal-eq-zero-of-ne}
  \lean{Pphi2.asymCanonicalRoughWickPower_two_site_marginal_eq_zero_of_ne}\leanok
  \uses{Pphi2-asymCanonicalRoughFieldFunctionOfSnd, GaussianField-AsymLatticeSites, Pphi2-asymCanonicalRoughWickConstant, Pphi2-asymCanonicalRoughGamma, Pphi2-asymCanonicalRoughFieldFunctionOfSnd-eq-sum-gamma, Pphi2-wickMonomial, Pphi2-asymCanonicalRoughGamma-sq-sum-eq-asymCanonicalRoughWickConstant, Pphi2-AsymCanonicalMode, Pphi2-wickPower-two-site-pi-gaussianReal-eq-zero-of-ne}
  \texttt{Pphi2.asymCanonicalRoughWickPower\_two\_site\_marginal\_eq\_zero\_of\_ne}
\end{theorem}
\begin{theorem}[\texttt{asymCanonicalRoughWickPower\_two\_site\_marginal\_integrable}]\label{Pphi2-asymCanonicalRoughWickPower-two-site-marginal-integrable}
  \lean{Pphi2.asymCanonicalRoughWickPower_two_site_marginal_integrable}\leanok
  \uses{Pphi2-asymCanonicalRoughFieldFunctionOfSnd, GaussianField-AsymLatticeSites, Pphi2-asymCanonicalRoughWickConstant, Pphi2-asymCanonicalRoughGamma, Pphi2-wickPower-two-site-pi-gaussianReal-integrable, Pphi2-asymCanonicalRoughFieldFunctionOfSnd-eq-sum-gamma, Pphi2-wickMonomial, Pphi2-asymCanonicalRoughGamma-sq-sum-eq-asymCanonicalRoughWickConstant, Pphi2-AsymCanonicalMode}
  \texttt{Pphi2.asymCanonicalRoughWickPower\_two\_site\_marginal\_integrable}
\end{theorem}
\begin{definition}[\texttt{asymCanonicalSmoothGamma}]\label{Pphi2-asymCanonicalSmoothGamma}
  \lean{Pphi2.asymCanonicalSmoothGamma}\leanok
  \uses{GaussianField-AsymLatticeSites, Pphi2-asymCanonicalSmoothModeCoeff, Pphi2-AsymCanonicalMode}
  \texttt{Pphi2.asymCanonicalSmoothGamma}
\end{definition}
\begin{theorem}[\texttt{asymCanonicalRoughWickPower\_two\_site\_marginal\_eq\_diag}]\label{Pphi2-asymCanonicalRoughWickPower-two-site-marginal-eq-diag}
  \lean{Pphi2.asymCanonicalRoughWickPower_two_site_marginal_eq_diag}\leanok
  \uses{Pphi2-asymCanonicalRoughFieldFunctionOfSnd, GaussianField-AsymLatticeSites, Pphi2-asymCanonicalRoughWickConstant, Pphi2-asymCanonicalRoughGamma, Pphi2-wickPower-two-site-pi-gaussianReal-eq-diag, Pphi2-asymCanonicalRoughFieldFunctionOfSnd-eq-sum-gamma, Pphi2-wickMonomial, Pphi2-asymCanonicalRoughCovariance-eq-sum-gamma-mul-gamma, Pphi2-asymCanonicalRoughGamma-sq-sum-eq-asymCanonicalRoughWickConstant, Pphi2-AsymCanonicalMode, Pphi2-asymCanonicalRoughCovariance}
  \texttt{Pphi2.asymCanonicalRoughWickPower\_two\_site\_marginal\_eq\_diag}
\end{theorem}
\begin{theorem}[\texttt{asymCanonicalRoughGamma\_sq\_sum\_eq\_asymCanonicalRoughWickConstant}]\label{Pphi2-asymCanonicalRoughGamma-sq-sum-eq-asymCanonicalRoughWickConstant}
  \lean{Pphi2.asymCanonicalRoughGamma_sq_sum_eq_asymCanonicalRoughWickConstant}\leanok
  \uses{GaussianField-AsymLatticeSites, Pphi2-asymCanonicalRoughWickConstant, Pphi2-asymCanonicalRoughGamma, Pphi2-asymCanonicalRoughCovariance-eq-sum-gamma-mul-gamma, Pphi2-AsymCanonicalMode, Pphi2-asymCanonicalRoughCovariance}
  \texttt{Pphi2.asymCanonicalRoughGamma\_sq\_sum\_eq\_asymCanonicalRoughWickConstant}
\end{theorem}
\begin{theorem}[\texttt{asymCanonicalSmoothCovariance\_eq\_sum\_gamma\_mul\_gamma}]\label{Pphi2-asymCanonicalSmoothCovariance-eq-sum-gamma-mul-gamma}
  \lean{Pphi2.asymCanonicalSmoothCovariance_eq_sum_gamma_mul_gamma}\leanok
  \uses{Pphi2-asymCanonicalSmoothCovariance, GaussianField-AsymLatticeSites, Pphi2-asymCanonicalBasis, Pphi2-asymCanonicalSmoothModeCoeff, GaussianField-latticeFourierNormSq, GaussianField-latticeFourierNormSq-pos, Pphi2-asymCanonicalSmoothGamma, Pphi2-asymCanonicalSmoothWeight, Pphi2-AsymCanonicalMode, Pphi2-asymCanonicalNormSq}
  \texttt{Pphi2.asymCanonicalSmoothCovariance\_eq\_sum\_gamma\_mul\_gamma}
\end{theorem}
\begin{theorem}[\texttt{congr\_simp}]\label{Pphi2-asymCanonicalSmoothGamma-congr-simp}
  \lean{Pphi2.asymCanonicalSmoothGamma.congr_simp}\leanok
  \uses{GaussianField-AsymLatticeSites, Pphi2-asymCanonicalSmoothGamma, Pphi2-AsymCanonicalMode}
  \texttt{Pphi2.asymCanonicalSmoothGamma.congr\_simp}
\end{theorem}
\begin{theorem}[\texttt{congr\_simp}]\label{Pphi2-asymCanonicalRoughFieldFunctionOfSnd-congr-simp}
  \lean{Pphi2.asymCanonicalRoughFieldFunctionOfSnd.congr_simp}\leanok
  \uses{Pphi2-asymCanonicalRoughFieldFunctionOfSnd, GaussianField-AsymLatticeSites, Pphi2-AsymCanonicalMode}
  \texttt{Pphi2.asymCanonicalRoughFieldFunctionOfSnd.congr\_simp}
\end{theorem}
\begin{theorem}[\texttt{asymCanonicalSmoothWickPower\_two\_site\_marginal\_eq\_zero\_of\_ne}]\label{Pphi2-asymCanonicalSmoothWickPower-two-site-marginal-eq-zero-of-ne}
  \lean{Pphi2.asymCanonicalSmoothWickPower_two_site_marginal_eq_zero_of_ne}\leanok
  \uses{GaussianField-AsymLatticeSites, Pphi2-asymCanonicalSmoothFieldFunctionOfFst-eq-sum-gamma, Pphi2-wickMonomial, Pphi2-asymCanonicalSmoothGamma-sq-sum-eq-asymSmoothWickConstant, Pphi2-asymCanonicalSmoothFieldFunctionOfFst, Pphi2-asymCanonicalSmoothGamma, Pphi2-asymSmoothWickConstant, Pphi2-AsymCanonicalMode, Pphi2-wickPower-two-site-pi-gaussianReal-eq-zero-of-ne}
  \texttt{Pphi2.asymCanonicalSmoothWickPower\_two\_site\_marginal\_eq\_zero\_of\_ne}
\end{theorem}
\begin{theorem}[\texttt{congr\_simp}]\label{Pphi2-asymCanonicalRoughGamma-congr-simp}
  \lean{Pphi2.asymCanonicalRoughGamma.congr_simp}\leanok
  \uses{GaussianField-AsymLatticeSites, Pphi2-asymCanonicalRoughGamma, Pphi2-AsymCanonicalMode}
  \texttt{Pphi2.asymCanonicalRoughGamma.congr\_simp}
\end{theorem}
\begin{theorem}[\texttt{asymCanonicalCrossTerm\_l2\_sq\_le}]\label{Pphi2-asymCanonicalCrossTerm-l2-sq-le}
  \lean{Pphi2.asymCanonicalCrossTerm_l2_sq_le}\leanok
  \uses{Pphi2-asymCanonicalSmoothFieldFunction-pointwise-measurable, Pphi2-asymCanonicalSmoothCovariance, Pphi2-asymCanonicalCrossTerm-sq-integrable, GaussianField-AsymLatticeSites, Pphi2-asymCanonicalRoughWickConstant, Pphi2-abs-asymCanonicalSmoothCovariance-le-asymSmoothWickConstant, Pphi2-asymSmoothWickConstant-le-log-uniform, Pphi2-asymCanonicalCrossTerm, Pphi2-AsymCanonicalJoint, Pphi2-asymCanonicalRoughFieldFunction-pointwise-measurable, Pphi2-wickMonomial, Pphi2-wickMonomial-continuous, Pphi2-asymCanonicalRoughFieldFunction, Pphi2-asymSmoothWickConstant, GaussianField-instFintypeAsymLatticeSitesOfNeZeroNat, Pphi2-asymCanonicalRoughCovariance-pow-sum-le-uniform, Pphi2-asymCanonicalSmoothFieldFunction, Pphi2-asymCanonicalJointMeasure, Pphi2-asymCanonicalCrossTerm-l2-sq-eq-covSum, Pphi2-asymCanonicalSmoothWeight, Lattice-toSemilatticeInf, Pphi2-AsymCanonicalMode, Pphi2-asymCanonicalRoughCovariance}
  \texttt{Pphi2.asymCanonicalCrossTerm\_l2\_sq\_le}
\end{theorem}
\begin{theorem}[\texttt{asymCanonicalSmoothFieldFunctionOfFst\_eq\_sum\_gamma}]\label{Pphi2-asymCanonicalSmoothFieldFunctionOfFst-eq-sum-gamma}
  \lean{Pphi2.asymCanonicalSmoothFieldFunctionOfFst_eq_sum_gamma}\leanok
  \uses{GaussianField-AsymLatticeSites, Pphi2-asymCanonicalSmoothModeCoeff, Pphi2-asymCanonicalSmoothFieldFunctionOfFst, Pphi2-asymCanonicalSmoothGamma, Pphi2-AsymCanonicalMode}
  \texttt{Pphi2.asymCanonicalSmoothFieldFunctionOfFst\_eq\_sum\_gamma}
\end{theorem}
\begin{theorem}[\texttt{asymCanonicalSmoothWickPower\_two\_site\_marginal\_eq\_diag}]\label{Pphi2-asymCanonicalSmoothWickPower-two-site-marginal-eq-diag}
  \lean{Pphi2.asymCanonicalSmoothWickPower_two_site_marginal_eq_diag}\leanok
  \uses{Pphi2-asymCanonicalSmoothCovariance, GaussianField-AsymLatticeSites, Pphi2-asymCanonicalSmoothFieldFunctionOfFst-eq-sum-gamma, Pphi2-wickPower-two-site-pi-gaussianReal-eq-diag, Pphi2-wickMonomial, Pphi2-asymCanonicalSmoothGamma-sq-sum-eq-asymSmoothWickConstant, Pphi2-asymCanonicalSmoothFieldFunctionOfFst, Pphi2-asymCanonicalSmoothGamma, Pphi2-asymSmoothWickConstant, Pphi2-asymCanonicalSmoothCovariance-eq-sum-gamma-mul-gamma, Pphi2-AsymCanonicalMode}
  \texttt{Pphi2.asymCanonicalSmoothWickPower\_two\_site\_marginal\_eq\_diag}
\end{theorem}
\begin{theorem}[\texttt{abs\_asymCanonicalSmoothCovariance\_le\_asymSmoothWickConstant}]\label{Pphi2-abs-asymCanonicalSmoothCovariance-le-asymSmoothWickConstant}
  \lean{Pphi2.abs_asymCanonicalSmoothCovariance_le_asymSmoothWickConstant}\leanok
  \uses{Pphi2-asymCanonicalSmoothCovariance, GaussianField-AsymLatticeSites, Pphi2-asymCanonicalSmoothGamma-sq-sum-eq-asymSmoothWickConstant, Pphi2-asymCanonicalSmoothGamma, Pphi2-asymSmoothWickConstant, Pphi2-asymCanonicalSmoothCovariance-eq-sum-gamma-mul-gamma, Lattice-toSemilatticeInf, Pphi2-AsymCanonicalMode}
  \texttt{Pphi2.abs\_asymCanonicalSmoothCovariance\_le\_asymSmoothWickConstant}
\end{theorem}
\begin{definition}[\texttt{asymCanonicalSmoothCovariance}]\label{Pphi2-asymCanonicalSmoothCovariance}
  \lean{Pphi2.asymCanonicalSmoothCovariance}\leanok
  \uses{GaussianField-AsymLatticeSites, Pphi2-asymCanonicalBasis, Pphi2-asymCanonicalSmoothWeight, Pphi2-AsymCanonicalMode, Pphi2-asymCanonicalNormSq}
  \texttt{Pphi2.asymCanonicalSmoothCovariance}
\end{definition}
\begin{theorem}[\texttt{asymCanonicalCrossTerm\_sq\_integrable}]\label{Pphi2-asymCanonicalCrossTerm-sq-integrable}
  \lean{Pphi2.asymCanonicalCrossTerm_sq_integrable}\leanok
  \uses{Pphi2-asymCanonicalRoughFieldFunctionOfSnd, GaussianField-AsymLatticeSites, Pphi2-asymCanonicalRoughWickConstant, Pphi2-asymCanonicalCrossTerm, Pphi2-AsymCanonicalJoint, Pphi2-asymCanonicalRoughWickPower-two-site-marginal-integrable, Pphi2-wickMonomial, Pphi2-asymCanonicalSmoothWickPower-two-site-marginal-integrable, Pphi2-asymCanonicalRoughFieldFunction, Pphi2-asymCanonicalSmoothFieldFunctionOfFst, Pphi2-asymSmoothWickConstant, GaussianField-instFintypeAsymLatticeSitesOfNeZeroNat, Pphi2-asymCanonicalSmoothFieldFunction, Pphi2-asymCanonicalJointMeasure, Pphi2-AsymCanonicalMode}
  \texttt{Pphi2.asymCanonicalCrossTerm\_sq\_integrable}
\end{theorem}
\begin{theorem}[\texttt{asymCanonicalRoughFieldFunctionOfSnd\_eq\_sum\_gamma}]\label{Pphi2-asymCanonicalRoughFieldFunctionOfSnd-eq-sum-gamma}
  \lean{Pphi2.asymCanonicalRoughFieldFunctionOfSnd_eq_sum_gamma}\leanok
  \uses{Pphi2-asymCanonicalRoughFieldFunctionOfSnd, GaussianField-AsymLatticeSites, Pphi2-asymCanonicalRoughModeCoeff, Pphi2-asymCanonicalRoughGamma, Pphi2-AsymCanonicalMode}
  \texttt{Pphi2.asymCanonicalRoughFieldFunctionOfSnd\_eq\_sum\_gamma}
\end{theorem}
\begin{theorem}[\texttt{asymCanonicalSmoothWickPower\_two\_site\_marginal\_integrable}]\label{Pphi2-asymCanonicalSmoothWickPower-two-site-marginal-integrable}
  \lean{Pphi2.asymCanonicalSmoothWickPower_two_site_marginal_integrable}\leanok
  \uses{GaussianField-AsymLatticeSites, Pphi2-asymCanonicalSmoothFieldFunctionOfFst-eq-sum-gamma, Pphi2-wickPower-two-site-pi-gaussianReal-integrable, Pphi2-wickMonomial, Pphi2-asymCanonicalSmoothGamma-sq-sum-eq-asymSmoothWickConstant, Pphi2-asymCanonicalSmoothFieldFunctionOfFst, Pphi2-asymCanonicalSmoothGamma, Pphi2-asymSmoothWickConstant, Pphi2-AsymCanonicalMode}
  \texttt{Pphi2.asymCanonicalSmoothWickPower\_two\_site\_marginal\_integrable}
\end{theorem}
\begin{definition}[\texttt{asymCanonicalRoughWickConstant}]\label{Pphi2-asymCanonicalRoughWickConstant}
  \lean{Pphi2.asymCanonicalRoughWickConstant}\leanok
  \uses{Pphi2-wickConstantAsym, Pphi2-asymSmoothWickConstant}
  \texttt{Pphi2.asymCanonicalRoughWickConstant}
\end{definition}
\begin{theorem}[\texttt{asymCanonicalRoughCovariance\_eq\_sum\_gamma\_mul\_gamma}]\label{Pphi2-asymCanonicalRoughCovariance-eq-sum-gamma-mul-gamma}
  \lean{Pphi2.asymCanonicalRoughCovariance_eq_sum_gamma_mul_gamma}\leanok
  \uses{GaussianField-AsymLatticeSites, Pphi2-asymCanonicalBasis, Pphi2-asymCanonicalRoughModeCoeff, Pphi2-asymCanonicalRoughGamma, Pphi2-asymCanonicalRoughGamma-congr-simp, GaussianField-latticeFourierNormSq, GaussianField-latticeFourierNormSq-pos, Pphi2-asymCanonicalRoughWeight, Pphi2-AsymCanonicalMode, Pphi2-asymCanonicalRoughCovariance, Pphi2-asymCanonicalNormSq}
  \texttt{Pphi2.asymCanonicalRoughCovariance\_eq\_sum\_gamma\_mul\_gamma}
\end{theorem}
\begin{theorem}[\texttt{asymCanonicalSmoothGamma\_sq\_sum\_eq\_asymSmoothWickConstant}]\label{Pphi2-asymCanonicalSmoothGamma-sq-sum-eq-asymSmoothWickConstant}
  \lean{Pphi2.asymCanonicalSmoothGamma_sq_sum_eq_asymSmoothWickConstant}\leanok
  \uses{GaussianField-AsymLatticeSites, Pphi2-asymCanonicalBasis, Pphi2-asymCanonicalSmoothModeCoeff, GaussianField-latticeFourierNormSq, GaussianField-latticeFourierNormSq-pos, Pphi2-asymCanonicalSmoothGamma, Pphi2-asymSmoothWickConstant, Pphi2-asymCanonicalSmoothWeight, Pphi2-AsymCanonicalMode, Pphi2-asymCanonicalNormSq}
  \texttt{Pphi2.asymCanonicalSmoothGamma\_sq\_sum\_eq\_asymSmoothWickConstant}
\end{theorem}
\begin{theorem}[\texttt{asymCanonicalCrossTerm\_l2\_sq\_eq\_covSum}]\label{Pphi2-asymCanonicalCrossTerm-l2-sq-eq-covSum}
  \lean{Pphi2.asymCanonicalCrossTerm_l2_sq_eq_covSum}\leanok
  \uses{Pphi2-asymCanonicalSmoothCovariance, Pphi2-asymCanonicalRoughFieldFunctionOfSnd, GaussianField-AsymLatticeSites, Pphi2-asymCanonicalRoughWickConstant, Pphi2-asymCanonicalCrossTerm, Pphi2-AsymCanonicalJoint, Pphi2-asymCanonicalRoughWickPower-two-site-marginal-integrable, Pphi2-wickMonomial, Pphi2-asymCanonicalSmoothWickPower-two-site-marginal-integrable, Pphi2-asymCanonicalRoughFieldFunction, Pphi2-asymCanonicalSmoothFieldFunctionOfFst, Pphi2-asymSmoothWickConstant, GaussianField-instFintypeAsymLatticeSitesOfNeZeroNat, Pphi2-asymCanonicalSmoothWickPower-two-site-marginal-eq-diag, Pphi2-asymCanonicalRoughWickPower-two-site-marginal-eq-diag, Pphi2-asymCanonicalSmoothFieldFunction, Pphi2-asymCanonicalJointMeasure, Pphi2-AsymCanonicalMode, Pphi2-asymCanonicalRoughCovariance}
  \texttt{Pphi2.asymCanonicalCrossTerm\_l2\_sq\_eq\_covSum}
\end{theorem}
\begin{theorem}[\texttt{congr\_simp}]\label{Pphi2-asymCanonicalRoughWickConstant-congr-simp}
  \lean{Pphi2.asymCanonicalRoughWickConstant.congr_simp}\leanok
  \uses{Pphi2-asymCanonicalRoughWickConstant}
  \texttt{Pphi2.asymCanonicalRoughWickConstant.congr\_simp}
\end{theorem}
\begin{theorem}[\texttt{congr\_simp}]\label{Pphi2-asymCanonicalCrossTerm-congr-simp}
  \lean{Pphi2.asymCanonicalCrossTerm.congr_simp}\leanok
  \uses{Pphi2-asymCanonicalCrossTerm, Pphi2-AsymCanonicalJoint}
  \texttt{Pphi2.asymCanonicalCrossTerm.congr\_simp}
\end{theorem}
\begin{theorem}[\texttt{asymCanonicalRoughCovariance\_pow\_sum\_le\_uniform}]\label{Pphi2-asymCanonicalRoughCovariance-pow-sum-le-uniform}
  \lean{Pphi2.asymCanonicalRoughCovariance_pow_sum_le_uniform}\leanok
  \uses{GaussianField-AsymLatticeSites, Pphi2-asymCanonicalRoughCovariance-pow-two-sum-le-uniform, GaussianField-instFintypeAsymLatticeSitesOfNeZeroNat, Pphi2-asymCanonicalRoughCovariance-pow-one-sum-le-uniform, Pphi2-asymCanonicalRoughCovariance-pow-sum-le-uniform-of-three-le, Pphi2-asymCanonicalRoughCovariance}
  \texttt{Pphi2.asymCanonicalRoughCovariance\_pow\_sum\_le\_uniform}
\end{theorem}
\section{\texttt{Pphi2.EuclideanOS}}
\begin{theorem}[\texttt{generatingFunctional\_re\_eq\_integral\_cos}]\label{Pphi2-EuclideanOS-generatingFunctional-re-eq-integral-cos}
  \lean{Pphi2.EuclideanOS.generatingFunctional_re_eq_integral_cos}\leanok
  \uses{Pphi2-EuclideanPlaneBackground-SpaceTime, Pphi2-EuclideanPlaneBackground-RealTestFunction, Pphi2-EuclideanPlaneBackground-dim, Pphi2-EuclideanPlaneBackground, GaussianField-Configuration, GaussianField-instMeasurableSpaceConfiguration, Pphi2-EuclideanOS-generatingFunctional, Pphi2-EuclideanPlaneBackground-Distribution}
  \texttt{Pphi2.EuclideanOS.generatingFunctional\_re\_eq\_integral\_cos}
\end{theorem}
\begin{theorem}[\texttt{os4\_ergodicity}]\label{Pphi2-EuclideanOS-SatisfiesFullOS-os4-ergodicity}
  \lean{Pphi2.EuclideanOS.SatisfiesFullOS.os4_ergodicity}\leanok
  \uses{Pphi2-EuclideanPlaneBackground-SpaceTime, Pphi2-EuclideanPlaneBackground-RealTestFunction, Pphi2-EuclideanOS-SatisfiesFullOS, Pphi2-EuclideanPlaneBackground-dim, Pphi2-EuclideanPlaneBackground, Pphi2-EuclideanTimeStructure, GaussianField-instMeasurableSpaceConfiguration, Pphi2-EuclideanOS-OS4-Ergodicity, Pphi2-EuclideanPlaneBackground-Distribution}
  \texttt{Pphi2.EuclideanOS.SatisfiesFullOS.os4\_ergodicity}
\end{theorem}
\begin{definition}[\texttt{generatingFunctional}]\label{Pphi2-EuclideanOS-generatingFunctional}
  \lean{Pphi2.EuclideanOS.generatingFunctional}\leanok
  \uses{Pphi2-EuclideanPlaneBackground-SpaceTime, Pphi2-EuclideanPlaneBackground-RealTestFunction, Pphi2-EuclideanPlaneBackground-dim, Pphi2-EuclideanPlaneBackground, GaussianField-instMeasurableSpaceConfiguration, Pphi2-EuclideanPlaneBackground-Distribution}
  $Z[f] = \int e^{i\omega(f)}\,d\mu$ (real) and $Z[J] = \int e^{i\langle\omega,J\rangle_\mathbb{C}}\,d\mu$ (complex).
\end{definition}
\begin{definition}[\texttt{distribPairing}]\label{Pphi2-EuclideanOS-distribPairing}
  \lean{Pphi2.EuclideanOS.distribPairing}\leanok
  \uses{Pphi2-EuclideanPlaneBackground-SpaceTime, Pphi2-EuclideanPlaneBackground-RealTestFunction, Pphi2-EuclideanPlaneBackground-dim, Pphi2-EuclideanPlaneBackground, Pphi2-EuclideanPlaneBackground-Distribution}
  \texttt{Pphi2.EuclideanOS.distribPairing}
\end{definition}
\begin{definition}[\texttt{generatingFunctionalℂ}]\label{Pphi2-EuclideanOS-generatingFunctional-}
  \lean{Pphi2.EuclideanOS.generatingFunctionalℂ}\leanok
  \uses{Pphi2-EuclideanPlaneBackground-SpaceTime, Pphi2-EuclideanPlaneBackground-RealTestFunction, Pphi2-EuclideanPlaneBackground-dim, Pphi2-EuclideanOS-schwartzIm, Pphi2-EuclideanOS-schwartzRe, Pphi2-EuclideanPlaneBackground, GaussianField-instMeasurableSpaceConfiguration, Pphi2-EuclideanPlaneBackground-ComplexTestFunction, Pphi2-EuclideanPlaneBackground-Distribution}
  \texttt{Pphi2.EuclideanOS.generatingFunctionalℂ}
\end{definition}
\begin{theorem}[\texttt{os2}]\label{Pphi2-EuclideanOS-SatisfiesFullOS-os2}
  \lean{Pphi2.EuclideanOS.SatisfiesFullOS.os2}\leanok
  \uses{Pphi2-EuclideanPlaneBackground-SpaceTime, Pphi2-EuclideanOS-OS2-EuclideanInvariance, Pphi2-EuclideanPlaneBackground-RealTestFunction, Pphi2-EuclideanOS-SatisfiesFullOS, Pphi2-EuclideanPlaneBackground-dim, Pphi2-EuclideanPlaneBackground, Pphi2-EuclideanTimeStructure, GaussianField-instMeasurableSpaceConfiguration, Pphi2-EuclideanPlaneBackground-Distribution}
  \texttt{Pphi2.EuclideanOS.SatisfiesFullOS.os2}
\end{theorem}
\begin{definition}[\texttt{OS4\_Ergodicity}]\label{Pphi2-EuclideanOS-OS4-Ergodicity}
  \lean{Pphi2.EuclideanOS.OS4_Ergodicity}\leanok
  \uses{Pphi2-EuclideanPlaneBackground-SpaceTime, Pphi2-EuclideanTimeStructure-timeTranslation, Pphi2-EuclideanPlaneBackground-RealTestFunction, Pphi2-EuclideanPlaneBackground-translate, Pphi2-EuclideanPlaneBackground-dim, Pphi2-EuclideanPlaneBackground, Pphi2-EuclideanTimeStructure, GaussianField-instMeasurableSpaceConfiguration, Pphi2-EuclideanOS-generatingFunctional, Pphi2-EuclideanPlaneBackground-Distribution}
  \texttt{Pphi2.EuclideanOS.OS4\_Ergodicity}
\end{definition}
\begin{theorem}[\texttt{os4\_clustering}]\label{Pphi2-EuclideanOS-SatisfiesFullOS-os4-clustering}
  \lean{Pphi2.EuclideanOS.SatisfiesFullOS.os4_clustering}\leanok
  \uses{Pphi2-EuclideanPlaneBackground-SpaceTime, Pphi2-EuclideanPlaneBackground-RealTestFunction, Pphi2-EuclideanOS-SatisfiesFullOS, Pphi2-EuclideanOS-OS4-Clustering, Pphi2-EuclideanPlaneBackground-dim, Pphi2-EuclideanPlaneBackground, Pphi2-EuclideanTimeStructure, GaussianField-instMeasurableSpaceConfiguration, Pphi2-EuclideanPlaneBackground-Distribution}
  \texttt{Pphi2.EuclideanOS.SatisfiesFullOS.os4\_clustering}
\end{theorem}
\begin{theorem}[\texttt{os0}]\label{Pphi2-EuclideanOS-SatisfiesFullOS-os0}
  \lean{Pphi2.EuclideanOS.SatisfiesFullOS.os0}\leanok
  \uses{Pphi2-EuclideanPlaneBackground-SpaceTime, Pphi2-EuclideanPlaneBackground-RealTestFunction, Pphi2-EuclideanOS-SatisfiesFullOS, Pphi2-EuclideanPlaneBackground-dim, Pphi2-EuclideanPlaneBackground, Pphi2-EuclideanTimeStructure, GaussianField-instMeasurableSpaceConfiguration, Pphi2-EuclideanOS-OS0-Analyticity, Pphi2-EuclideanPlaneBackground-Distribution}
  \texttt{Pphi2.EuclideanOS.SatisfiesFullOS.os0}
\end{theorem}
\begin{definition}[\texttt{OS1\_Regularity}]\label{Pphi2-EuclideanOS-OS1-Regularity}
  \lean{Pphi2.EuclideanOS.OS1_Regularity}\leanok
  \uses{Pphi2-EuclideanPlaneBackground-SpaceTime, Pphi2-EuclideanPlaneBackground-RealTestFunction, Pphi2-EuclideanPlaneBackground-dim, Pphi2-EuclideanPlaneBackground, Pphi2-EuclideanOS-generatingFunctional-, GaussianField-instMeasurableSpaceConfiguration, Pphi2-EuclideanPlaneBackground-ComplexTestFunction, Pphi2-EuclideanPlaneBackground-Distribution}
  \texttt{Pphi2.EuclideanOS.OS1\_Regularity}
\end{definition}
\begin{definition}[\texttt{schwartzIm}]\label{Pphi2-EuclideanOS-schwartzIm}
  \lean{Pphi2.EuclideanOS.schwartzIm}\leanok
  \uses{Pphi2-EuclideanPlaneBackground-SpaceTime, Pphi2-EuclideanPlaneBackground-RealTestFunction, Pphi2-EuclideanPlaneBackground-dim, Pphi2-EuclideanPlaneBackground, Pphi2-EuclideanPlaneBackground-ComplexTestFunction}
  \texttt{Pphi2.EuclideanOS.schwartzIm}
\end{definition}
\begin{theorem}[\texttt{os3}]\label{Pphi2-EuclideanOS-SatisfiesFullOS-os3}
  \lean{Pphi2.EuclideanOS.SatisfiesFullOS.os3}\leanok
  \uses{Pphi2-EuclideanPlaneBackground-SpaceTime, Pphi2-EuclideanPlaneBackground-RealTestFunction, Pphi2-EuclideanOS-SatisfiesFullOS, Pphi2-EuclideanPlaneBackground-dim, Pphi2-EuclideanPlaneBackground, Pphi2-EuclideanTimeStructure, Pphi2-EuclideanOS-OS3-ReflectionPositivity, GaussianField-instMeasurableSpaceConfiguration, Pphi2-EuclideanPlaneBackground-Distribution}
  \texttt{Pphi2.EuclideanOS.SatisfiesFullOS.os3}
\end{theorem}
\begin{definition}[\texttt{SatisfiesFullOS}]\label{Pphi2-EuclideanOS-SatisfiesFullOS}
  \lean{Pphi2.EuclideanOS.SatisfiesFullOS}\leanok
  \uses{Pphi2-EuclideanPlaneBackground-SpaceTime, Pphi2-EuclideanPlaneBackground-RealTestFunction, Pphi2-EuclideanPlaneBackground-dim, Pphi2-EuclideanPlaneBackground, Pphi2-EuclideanTimeStructure, GaussianField-instMeasurableSpaceConfiguration, Pphi2-EuclideanPlaneBackground-Distribution}
  Structure bundling all five axioms.
\end{definition}
\begin{theorem}[\texttt{generatingFunctional\_im\_eq\_integral\_sin}]\label{Pphi2-EuclideanOS-generatingFunctional-im-eq-integral-sin}
  \lean{Pphi2.EuclideanOS.generatingFunctional_im_eq_integral_sin}\leanok
  \uses{Pphi2-EuclideanPlaneBackground-SpaceTime, Pphi2-EuclideanPlaneBackground-RealTestFunction, Pphi2-EuclideanPlaneBackground-dim, Pphi2-EuclideanPlaneBackground, GaussianField-Configuration, GaussianField-instMeasurableSpaceConfiguration, Pphi2-EuclideanOS-generatingFunctional, Pphi2-EuclideanPlaneBackground-Distribution}
  \texttt{Pphi2.EuclideanOS.generatingFunctional\_im\_eq\_integral\_sin}
\end{theorem}
\begin{theorem}[\texttt{os1}]\label{Pphi2-EuclideanOS-SatisfiesFullOS-os1}
  \lean{Pphi2.EuclideanOS.SatisfiesFullOS.os1}\leanok
  \uses{Pphi2-EuclideanPlaneBackground-SpaceTime, Pphi2-EuclideanPlaneBackground-RealTestFunction, Pphi2-EuclideanOS-SatisfiesFullOS, Pphi2-EuclideanPlaneBackground-dim, Pphi2-EuclideanPlaneBackground, Pphi2-EuclideanTimeStructure, GaussianField-instMeasurableSpaceConfiguration, Pphi2-EuclideanOS-OS1-Regularity, Pphi2-EuclideanPlaneBackground-Distribution}
  \texttt{Pphi2.EuclideanOS.SatisfiesFullOS.os1}
\end{theorem}
\begin{definition}[\texttt{schwartzRe}]\label{Pphi2-EuclideanOS-schwartzRe}
  \lean{Pphi2.EuclideanOS.schwartzRe}\leanok
  \uses{Pphi2-EuclideanPlaneBackground-SpaceTime, Pphi2-EuclideanPlaneBackground-RealTestFunction, Pphi2-EuclideanPlaneBackground-dim, Pphi2-EuclideanPlaneBackground, Pphi2-EuclideanPlaneBackground-ComplexTestFunction}
  Extract real/imaginary parts of complex Schwartz functions as real Schwartz functions.
\end{definition}
\begin{definition}[\texttt{OS3\_ReflectionPositivity}]\label{Pphi2-EuclideanOS-OS3-ReflectionPositivity}
  \lean{Pphi2.EuclideanOS.OS3_ReflectionPositivity}\leanok
  \uses{Pphi2-EuclideanPlaneBackground-SpaceTime, Pphi2-EuclideanPlaneBackground-RealTestFunction, Pphi2-EuclideanPlaneBackground-dim, Pphi2-EuclideanTimeStructure-positiveTimeSubmodule, Pphi2-EuclideanPlaneBackground, Pphi2-EuclideanTimeStructure, GaussianField-instMeasurableSpaceConfiguration, Pphi2-EuclideanTimeStructure-reflectOnRealTestFunctions, Pphi2-EuclideanOS-generatingFunctional, Pphi2-EuclideanPlaneBackground-Distribution}
  \texttt{Pphi2.EuclideanOS.OS3\_ReflectionPositivity}
\end{definition}
\begin{definition}[\texttt{OS4\_Clustering}]\label{Pphi2-EuclideanOS-OS4-Clustering}
  \lean{Pphi2.EuclideanOS.OS4_Clustering}\leanok
  \uses{Pphi2-EuclideanPlaneBackground-SpaceTime, Pphi2-EuclideanPlaneBackground-RealTestFunction, Pphi2-EuclideanPlaneBackground-translate, Pphi2-EuclideanPlaneBackground-dim, Pphi2-EuclideanPlaneBackground, GaussianField-instMeasurableSpaceConfiguration, Pphi2-EuclideanOS-generatingFunctional, Pphi2-EuclideanPlaneBackground-Distribution}
  \texttt{Pphi2.EuclideanOS.OS4\_Clustering}
\end{definition}
\begin{definition}[\texttt{OS0\_Analyticity}]\label{Pphi2-EuclideanOS-OS0-Analyticity}
  \lean{Pphi2.EuclideanOS.OS0_Analyticity}\leanok
  \uses{Pphi2-EuclideanPlaneBackground-SpaceTime, Pphi2-EuclideanPlaneBackground-RealTestFunction, Pphi2-EuclideanPlaneBackground-dim, Pphi2-EuclideanPlaneBackground, Pphi2-EuclideanOS-generatingFunctional-, GaussianField-instMeasurableSpaceConfiguration, Pphi2-EuclideanPlaneBackground-ComplexTestFunction, Pphi2-EuclideanPlaneBackground-Distribution}
  \texttt{Pphi2.EuclideanOS.OS0\_Analyticity}
\end{definition}
\begin{definition}[\texttt{OS2\_EuclideanInvariance}]\label{Pphi2-EuclideanOS-OS2-EuclideanInvariance}
  \lean{Pphi2.EuclideanOS.OS2_EuclideanInvariance}\leanok
  \uses{Pphi2-EuclideanPlaneBackground-SpaceTime, Pphi2-EuclideanPlaneBackground-RealTestFunction, Pphi2-EuclideanPlaneBackground-dim, Pphi2-EuclideanPlaneBackground-Motion, Pphi2-EuclideanPlaneBackground, Pphi2-EuclideanPlaneBackground-actionComplex, Pphi2-EuclideanOS-generatingFunctional-, GaussianField-instMeasurableSpaceConfiguration, Pphi2-EuclideanPlaneBackground-ComplexTestFunction, Pphi2-EuclideanPlaneBackground-Distribution}
  \texttt{Pphi2.EuclideanOS.OS2\_EuclideanInvariance}
\end{definition}
\section{\texttt{Pphi2.TorusContinuumLimit.TorusGaussianLimit}}
\begin{theorem}[\texttt{torusGaussianLimit\_isGaussian}]\label{Pphi2-torusGaussianLimit-isGaussian}
  \lean{Pphi2.torusGaussianLimit_isGaussian}\leanok
  \uses{GaussianField-NuclearTensorProduct-instModuleReal, GaussianField-NuclearTensorProduct-instTopologicalSpace, GaussianField-NuclearTensorProduct-instIsTopologicalAddGroup, GaussianField-NuclearTensorProduct-instAddCommGroup, GaussianField-configuration-eval-measurable, GaussianField-Configuration, GaussianField-NuclearTensorProduct-instContinuousSMulReal, GaussianField-SmoothMap-Circle, GaussianField-instMeasurableSpaceConfiguration, GaussianField-TorusTestFunction, Pphi2-weakLimit-centered-gaussianReal}
  Weak limits of torus lattice Gaussians are Gaussian: $\int e^{\omega(f)}\, d\mu = \exp(\frac{1}{2}\int (\omega f)^2\, d\mu)$ for all $f$. Proved from the 1D Gaussian identification + MGF matching.
\end{theorem}
\begin{theorem}[\texttt{gaussian\_measure\_unique\_of\_covariance}]\label{Pphi2-gaussian-measure-unique-of-covariance}
  \lean{Pphi2.gaussian_measure_unique_of_covariance}\leanok
  \uses{GaussianField-NuclearTensorProduct-instModuleReal, GaussianField-NuclearTensorProduct-instTopologicalSpace, GaussianField-NuclearTensorProduct-dyninMityaginSpace, GaussianField-NuclearTensorProduct-instIsTopologicalAddGroup, GaussianField-NuclearTensorProduct-instAddCommGroup, GaussianField-Configuration, GaussianField-gaussian-measure-unique-of-covariance, GaussianField-NuclearTensorProduct-instContinuousSMulReal, GaussianField-SmoothMap-Circle, GaussianField-instMeasurableSpaceConfiguration, GaussianField-TorusTestFunction}
  Main theorem: Two centered Gaussian probability measures on Configuration($E$) with the same covariance are equal. Proved by pulling back from $\mathbb{R}^{\mathbb{N}}$ using \texttt{instMeasurableSpaceConfiguration\_eq\_comap}.
\end{theorem}
\begin{theorem}[\texttt{isGaussian}]\label{Pphi2-IsTorusGaussianContinuumLimit-isGaussian}
  \lean{Pphi2.IsTorusGaussianContinuumLimit.isGaussian}\leanok
  \uses{Pphi2-IsTorusGaussianContinuumLimit, GaussianField-NuclearTensorProduct-instModuleReal, GaussianField-NuclearTensorProduct-instTopologicalSpace, GaussianField-NuclearTensorProduct-instIsTopologicalAddGroup, GaussianField-NuclearTensorProduct-instAddCommGroup, GaussianField-Configuration, GaussianField-NuclearTensorProduct-instContinuousSMulReal, GaussianField-SmoothMap-Circle, GaussianField-instMeasurableSpaceConfiguration, GaussianField-TorusTestFunction}
  \texttt{Pphi2.IsTorusGaussianContinuumLimit.isGaussian}
\end{theorem}
\begin{theorem}[\texttt{configuration\_neg\_apply}]\label{Pphi2-configuration-neg-apply}
  \lean{Pphi2.configuration_neg_apply}\leanok
  \uses{GaussianField-NuclearTensorProduct-instModuleReal, GaussianField-NuclearTensorProduct-instTopologicalSpace, GaussianField-NuclearTensorProduct-instIsTopologicalAddGroup, GaussianField-NuclearTensorProduct-instAddCommGroup, GaussianField-Configuration, GaussianField-NuclearTensorProduct-instContinuousSMulReal, GaussianField-SmoothMap-Circle, GaussianField-TorusTestFunction}
  Negation in configuration space: $(-\omega)(f) = -\omega(f)$.
\end{theorem}
\begin{theorem}[\texttt{z2\_symmetric}]\label{Pphi2-IsTorusGaussianContinuumLimit-z2-symmetric}
  \lean{Pphi2.IsTorusGaussianContinuumLimit.z2_symmetric}\leanok
  \uses{Pphi2-IsTorusGaussianContinuumLimit, GaussianField-NuclearTensorProduct-instModuleReal, GaussianField-NuclearTensorProduct-instTopologicalSpace, GaussianField-NuclearTensorProduct-instIsTopologicalAddGroup, GaussianField-NuclearTensorProduct-instAddCommGroup, GaussianField-Configuration, GaussianField-NuclearTensorProduct-instContinuousSMulReal, GaussianField-SmoothMap-Circle, GaussianField-instMeasurableSpaceConfiguration, GaussianField-TorusTestFunction}
  \texttt{Pphi2.IsTorusGaussianContinuumLimit.z2\_symmetric}
\end{theorem}
\begin{theorem}[\texttt{torusGaussianLimit\_nontrivial}]\label{Pphi2-torusGaussianLimit-nontrivial}
  \lean{Pphi2.torusGaussianLimit_nontrivial}\leanok
  \uses{Pphi2-torusContinuumGreen, GaussianField-NuclearTensorProduct-instModuleReal, GaussianField-NuclearTensorProduct-instTopologicalSpace, GaussianField-NuclearTensorProduct-instIsTopologicalAddGroup, GaussianField-NuclearTensorProduct-instAddCommGroup, Pphi2-torusContinuumGreen-pos, GaussianField-Configuration, GaussianField-NuclearTensorProduct-instContinuousSMulReal, GaussianField-SmoothMap-Circle, GaussianField-instMeasurableSpaceConfiguration, GaussianField-TorusTestFunction}
  $\int (\omega f)^2\, d\mu > 0$ for $f \ne 0$.
\end{theorem}
\begin{theorem}[\texttt{congr\_simp}]\label{Pphi2-torusContinuumMeasure-congr-simp}
  \lean{Pphi2.torusContinuumMeasure.congr_simp}\leanok
  \uses{GaussianField-NuclearTensorProduct-instModuleReal, GaussianField-NuclearTensorProduct-instTopologicalSpace, GaussianField-NuclearTensorProduct-instIsTopologicalAddGroup, Pphi2-torusContinuumMeasure, GaussianField-NuclearTensorProduct-instAddCommGroup, GaussianField-Configuration, GaussianField-NuclearTensorProduct-instContinuousSMulReal, GaussianField-SmoothMap-Circle, GaussianField-instMeasurableSpaceConfiguration, GaussianField-TorusTestFunction}
  \texttt{Pphi2.torusContinuumMeasure.congr\_simp}
\end{theorem}
\begin{theorem}[\texttt{torusLimit\_covariance\_eq}]\label{Pphi2-torusLimit-covariance-eq}
  \lean{Pphi2.torusLimit_covariance_eq}\leanok
  \uses{Pphi2-torus-propagator-convergence, Pphi2-torusContinuumGreen, GaussianField-NuclearTensorProduct-instModuleReal, GaussianField-NuclearTensorProduct-instTopologicalSpace, Pphi2-torusEmbeddedTwoPoint, GaussianField-NuclearTensorProduct-instIsTopologicalAddGroup, Pphi2-torusContinuumMeasure, GaussianField-NuclearTensorProduct-instAddCommGroup, GaussianField-configuration-eval-measurable, GaussianField-Configuration, Pphi2-torusGaussianLimit-isGaussian, GaussianField-NuclearTensorProduct-instContinuousSMulReal, Pphi2-torusContinuumMeasure-isProbability, GaussianField-SmoothMap-Circle, GaussianField-instMeasurableSpaceConfiguration, GaussianField-TorusTestFunction, Lattice-toSemilatticeInf, Pphi2-weakLimit-centered-gaussianReal, Pphi2-torusGaussianMeasure-isGaussian}
  The covariance of the weak limit equals the continuum Green's function: $\int (\omega f)^2\, d\mu = G_L(f,f)$.
\end{theorem}
\begin{theorem}[\texttt{covariance\_eq}]\label{Pphi2-IsTorusGaussianContinuumLimit-covariance-eq}
  \lean{Pphi2.IsTorusGaussianContinuumLimit.covariance_eq}\leanok
  \uses{Pphi2-torusContinuumGreen, Pphi2-IsTorusGaussianContinuumLimit, GaussianField-NuclearTensorProduct-instModuleReal, GaussianField-NuclearTensorProduct-instTopologicalSpace, GaussianField-NuclearTensorProduct-instIsTopologicalAddGroup, GaussianField-NuclearTensorProduct-instAddCommGroup, GaussianField-Configuration, GaussianField-NuclearTensorProduct-instContinuousSMulReal, GaussianField-SmoothMap-Circle, GaussianField-instMeasurableSpaceConfiguration, GaussianField-TorusTestFunction}
  \texttt{Pphi2.IsTorusGaussianContinuumLimit.covariance\_eq}
\end{theorem}
\begin{theorem}[\texttt{torusGaussianMeasure\_isGaussian}]\label{Pphi2-torusGaussianMeasure-isGaussian}
  \lean{Pphi2.torusGaussianMeasure_isGaussian}\leanok
  \uses{Pphi2-torusEmbeddedTwoPoint-eq-lattice-second-moment, GaussianField-NuclearTensorProduct-instModuleReal, GaussianField-NuclearTensorProduct-instTopologicalSpace, GaussianField-finLatticeField-dyninMityaginSpace, Pphi2-lattice-second-moment-eq-inner, GaussianField-ell2-separableSpace, Pphi2-torusEmbeddedTwoPoint, GaussianField-FinLatticeSites, GaussianField-NuclearTensorProduct-instIsTopologicalAddGroup, Pphi2-torusEmbedLift, GaussianField-latticeCovarianceGJ, Pphi2-torusContinuumMeasure, Pphi2-torusEmbedLift-measurable, GaussianField-NuclearTensorProduct-instAddCommGroup, GaussianField-pairing-is-gaussian, GaussianField-latticeGaussianMeasure, GaussianField-configuration-eval-measurable, GaussianField-Configuration, GaussianField-NuclearTensorProduct-instContinuousSMulReal, GaussianField-SmoothMap-Circle, GaussianField-circleSpacing, GaussianField-instMeasurableSpaceConfiguration, Pphi2-latticeTestFn, Pphi2-torusEmbedLift-eval-eq, GaussianField-ell2-, GaussianField-FinLatticeField, GaussianField-TorusTestFunction, GaussianField-circleSpacing-pos}
  Each lattice Gaussian measure satisfies the Gaussian MGF identity.
\end{theorem}
\begin{theorem}[\texttt{torusGaussianLimit\_converges}]\label{Pphi2-torusGaussianLimit-converges}
  \lean{Pphi2.torusGaussianLimit_converges}\leanok
  \uses{Pphi2-torusContinuumGreen, Pphi2-IsTorusGaussianContinuumLimit, GaussianField-NuclearTensorProduct-instModuleReal, GaussianField-NuclearTensorProduct-instTopologicalSpace, Pphi2-torusEmbeddedTwoPoint, Pphi2-torusGaussianMeasure-z2-symmetric, GaussianField-NuclearTensorProduct-instIsTopologicalAddGroup, Pphi2-torusContinuumMeasure, GaussianField-NuclearTensorProduct-instAddCommGroup, Pphi2-torusGaussianLimit-fullConvergence, GaussianField-Configuration, Pphi2-torusGaussianLimit-isGaussian, GaussianField-NuclearTensorProduct-instContinuousSMulReal, Pphi2-torusContinuumMeasure-isProbability, GaussianField-SmoothMap-Circle, GaussianField-instMeasurableSpaceConfiguration, Pphi2-torusLimit-covariance-eq, Pphi2-torusGaussianLimit-exists, Pphi2-z2-symmetric-of-weakLimit, GaussianField-TorusTestFunction, Pphi2-torusGaussianMeasure-isGaussian}
  Master convergence theorem: produces a subsequence, the limit measure, its \texttt{IsTorusGaussianContinuumLimit} predicate, and the weak convergence statement.
\end{theorem}
\begin{definition}[\texttt{IsTorusGaussianContinuumLimit}]\label{Pphi2-IsTorusGaussianContinuumLimit}
  \lean{Pphi2.IsTorusGaussianContinuumLimit}\leanok
  \uses{GaussianField-NuclearTensorProduct-instModuleReal, GaussianField-NuclearTensorProduct-instTopologicalSpace, GaussianField-NuclearTensorProduct-instIsTopologicalAddGroup, GaussianField-NuclearTensorProduct-instAddCommGroup, GaussianField-Configuration, GaussianField-NuclearTensorProduct-instContinuousSMulReal, GaussianField-SmoothMap-Circle, GaussianField-instMeasurableSpaceConfiguration, GaussianField-TorusTestFunction}
  ``\texttt{lean structure IsTorusGaussianContinuumLimit (μ : ...) (mass : ℝ) (hmass : 0 < mass) : Prop }`` Predicate: $\mu$ is a probability measure, Gaussian ($\int e^{\omega f}\, d\mu = \exp(\frac{1}{2}\int (\omega f)^2\, d\mu)$), has covariance $= G_L(f,f)$, and is $\mathbb{Z}_2$-symmetric.
\end{definition}
\begin{theorem}[\texttt{isProbability}]\label{Pphi2-IsTorusGaussianContinuumLimit-isProbability}
  \lean{Pphi2.IsTorusGaussianContinuumLimit.isProbability}\leanok
  \uses{Pphi2-IsTorusGaussianContinuumLimit, GaussianField-NuclearTensorProduct-instModuleReal, GaussianField-NuclearTensorProduct-instTopologicalSpace, GaussianField-NuclearTensorProduct-instIsTopologicalAddGroup, GaussianField-NuclearTensorProduct-instAddCommGroup, GaussianField-Configuration, GaussianField-NuclearTensorProduct-instContinuousSMulReal, GaussianField-SmoothMap-Circle, GaussianField-instMeasurableSpaceConfiguration, GaussianField-TorusTestFunction}
  \texttt{Pphi2.IsTorusGaussianContinuumLimit.isProbability}
\end{theorem}
\begin{theorem}[\texttt{weakLimit\_centered\_gaussianReal}]\label{Pphi2-weakLimit-centered-gaussianReal}
  \lean{Pphi2.weakLimit_centered_gaussianReal}\leanok
  \uses{Lattice-toSemilatticeSup, Lattice-toSemilatticeInf}
  The 1D marginals of the weak limit equal $N(0, G_L(f,f))$ for each test function $f$.
\end{theorem}
\begin{theorem}[\texttt{torusGaussianMeasure\_z2\_symmetric}]\label{Pphi2-torusGaussianMeasure-z2-symmetric}
  \lean{Pphi2.torusGaussianMeasure_z2_symmetric}\leanok
  \uses{GaussianField-NuclearTensorProduct-instModuleReal, GaussianField-NuclearTensorProduct-instTopologicalSpace, Pphi2-torusEmbeddedTwoPoint, GaussianField-NuclearTensorProduct-instIsTopologicalAddGroup, Pphi2-torusContinuumMeasure, GaussianField-NuclearTensorProduct-instAddCommGroup, GaussianField-configuration-eval-measurable, GaussianField-Configuration, GaussianField-NuclearTensorProduct-instContinuousSMulReal, Pphi2-torusContinuumMeasure-isProbability, GaussianField-SmoothMap-Circle, GaussianField-instMeasurableSpaceConfiguration, Pphi2-gaussian-measure-unique-of-covariance, GaussianField-TorusTestFunction, Pphi2-torusGaussianMeasure-isGaussian}
  Each lattice Gaussian measure is $\mathbb{Z}_2$-symmetric: $\mu \circ (-)^{-1} = \mu$.
\end{theorem}
\begin{theorem}[\texttt{z2\_symmetric\_of\_weakLimit}]\label{Pphi2-z2-symmetric-of-weakLimit}
  \lean{Pphi2.z2_symmetric_of_weakLimit}\leanok
  \uses{GaussianField-NuclearTensorProduct-instModuleReal, GaussianField-NuclearTensorProduct-instTopologicalSpace, GaussianField-NuclearTensorProduct-dyninMityaginSpace, GaussianField-NuclearTensorProduct-instIsTopologicalAddGroup, GaussianField-NuclearTensorProduct-instAddCommGroup, GaussianField-configBasisEval-measurable, GaussianField-Configuration, GaussianField-DyninMityaginSpace-basis, GaussianField-NuclearTensorProduct-instContinuousSMulReal, GaussianField-SmoothMap-Circle, GaussianField-instMeasurableSpaceConfiguration, GaussianField-instMeasurableSpaceConfiguration-eq-comap, GaussianField-configBasisEval, GaussianField-configBasisEval-continuous, GaussianField-TorusTestFunction}
  $\mathbb{Z}_2$ symmetry is inherited by the weak limit.
\end{theorem}
\begin{theorem}[\texttt{torusGaussianLimit\_fullConvergence}]\label{Pphi2-torusGaussianLimit-fullConvergence}
  \lean{Pphi2.torusGaussianLimit_fullConvergence}\leanok
  \uses{Pphi2-torusContinuumGreen, GaussianField-NuclearTensorProduct-instModuleReal, GaussianField-NuclearTensorProduct-instTopologicalSpace, Pphi2-torusEmbeddedTwoPoint, GaussianField-NuclearTensorProduct-dyninMityaginSpace, GaussianField-NuclearTensorProduct-instIsTopologicalAddGroup, Pphi2-torusContinuumMeasure, GaussianField-NuclearTensorProduct-instAddCommGroup, GaussianField-prokhorov-configuration, GaussianField-Configuration, Pphi2-torusGaussianLimit-isGaussian, GaussianField-NuclearTensorProduct-instContinuousSMulReal, Pphi2-torusContinuumMeasure-isProbability, GaussianField-SmoothMap-Circle, GaussianField-instMeasurableSpaceConfiguration, Pphi2-torusLimit-covariance-eq, Pphi2-gaussian-measure-unique-of-covariance, Pphi2-torusContinuumMeasures-tight, GaussianField-TorusTestFunction, Pphi2-torusGaussianMeasure-isGaussian}
  Full convergence: there exist a subsequence and a unique Gaussian limit measure satisfying \texttt{IsTorusGaussianContinuumLimit}.
\end{theorem}
\section{\texttt{Pphi2.WickOrdering.WickPolynomial}}
\begin{theorem}[\texttt{wickMonomial\_eq\_hermite}]\label{Pphi2-wickMonomial-eq-hermite}
  \lean{Pphi2.wickMonomial_eq_hermite}\leanok
  \uses{Pphi2-wickMonomial}
  For $c > 0$: $:x^n:_c = c^{n/2} \cdot He_n(x/\sqrt{c})$ where $He_n$ is the probabilist's Hermite polynomial. Proved via the bridge to \texttt{gaussian-field}'s \texttt{wick\_eq\_hermiteR\_rpow}.
\end{theorem}
\begin{theorem}[\texttt{wickMonomial\_two}]\label{Pphi2-wickMonomial-two}
  \lean{Pphi2.wickMonomial_two}\leanok
  \uses{Pphi2-wickMonomial}
  Explicit formulas: $:x^2:_c = x^2 - c$, $:x^3:_c = x^3 - 3cx$, $:x^4:_c = x^4 - 6cx^2 + 3c^2$. Proved by \texttt{simp} + \texttt{ring}.
\end{theorem}
\begin{theorem}[\texttt{wickPolynomial\_zero\_variance}]\label{Pphi2-wickPolynomial-zero-variance}
  \lean{Pphi2.wickPolynomial_zero_variance}\leanok
  \uses{Pphi2-wickMonomial, Pphi2-wickMonomial-zero-variance, Pphi2-wickPolynomial}
  $:P(x):_0 = P(x)$ (Wick ordering at zero variance recovers the original polynomial).
\end{theorem}
\begin{theorem}[\texttt{wickMonomial\_zero}]\label{Pphi2-wickMonomial-zero}
  \lean{Pphi2.wickMonomial_zero}\leanok
  \uses{Pphi2-wickMonomial}
  Simp lemmas for the recursive cases.
\end{theorem}
\begin{definition}[\texttt{wickMonomial}]\label{Pphi2-wickMonomial}
  \lean{Pphi2.wickMonomial}\leanok
  \texttt{Pphi2.wickMonomial}
\end{definition}
\begin{theorem}[\texttt{wickPolynomial\_bounded\_below}]\label{Pphi2-wickPolynomial-bounded-below}
  \lean{Pphi2.wickPolynomial_bounded_below}\leanok
  \uses{Pphi2-poly-even-degree-bounded-below, Pphi2-wickPolynomial}
  $\exists A > 0$ such that $:P(x):_c \ge -A$ for all $x \in \mathbb{R}$. Proved by converting to a formal \texttt{Polynomial} and verifying even degree $\ge 4$ with positive leading coefficient $1/n$.
\end{theorem}
\begin{theorem}[\texttt{wickMonomial\_zero\_variance}]\label{Pphi2-wickMonomial-zero-variance}
  \lean{Pphi2.wickMonomial_zero_variance}\leanok
  \uses{Pphi2-wickMonomial}
  $:x^n:_0 = x^n$ for all $n$. Proved by induction.
\end{theorem}
\begin{theorem}[\texttt{wickPolynomial\_lower\_bound\_general}]\label{Pphi2-wickPolynomial-lower-bound-general}
  \lean{Pphi2.wickPolynomial_lower_bound_general}\leanok
  \uses{Pphi2-wickPolynomial-uniform-bounded-below, Lattice-toSemilatticeInf, Pphi2-wickPolynomial}
  \texttt{Pphi2.wickPolynomial\_lower\_bound\_general}
\end{theorem}
\begin{theorem}[\texttt{wickMonomial\_succ\_succ}]\label{Pphi2-wickMonomial-succ-succ}
  \lean{Pphi2.wickMonomial_succ_succ}\leanok
  \uses{Pphi2-wickMonomial}
  \texttt{Pphi2.wickMonomial\_succ\_succ}
\end{theorem}
\begin{definition}[\texttt{wickPolynomial}]\label{Pphi2-wickPolynomial}
  \lean{Pphi2.wickPolynomial}\leanok
  \uses{Pphi2-wickMonomial}
  \texttt{Pphi2.wickPolynomial}
\end{definition}
\begin{theorem}[\texttt{wickMonomial\_three}]\label{Pphi2-wickMonomial-three}
  \lean{Pphi2.wickMonomial_three}\leanok
  \uses{Pphi2-wickMonomial}
  \texttt{Pphi2.wickMonomial\_three}
\end{theorem}
\begin{theorem}[\texttt{wickMonomial\_one}]\label{Pphi2-wickMonomial-one}
  \lean{Pphi2.wickMonomial_one}\leanok
  \uses{Pphi2-wickMonomial}
  \texttt{Pphi2.wickMonomial\_one}
\end{theorem}
\begin{theorem}[\texttt{poly\_even\_degree\_bounded\_below}]\label{Pphi2-poly-even-degree-bounded-below}
  \lean{Pphi2.poly_even_degree_bounded_below}\leanok
  \uses{Lattice-toSemilatticeInf}
  A real polynomial with even degree $\ge 2$ and positive leading coefficient is bounded below: $\exists A > 0,\; \forall x,\; p(x) \ge -A$. Proved via \texttt{tendsto\_norm\_atTop} + continuity minimum.
\end{theorem}
\begin{theorem}[\texttt{wickPolynomial\_continuous₂}]\label{Pphi2-wickPolynomial-continuous-}
  \lean{Pphi2.wickPolynomial_continuous₂}\leanok
  \uses{Pphi2-wickMonomial, Pphi2-wickMonomial-continuous-, Pphi2-wickPolynomial}
  \texttt{Pphi2.wickPolynomial\_continuous₂}
\end{theorem}
\begin{theorem}[\texttt{wickMonomial\_four}]\label{Pphi2-wickMonomial-four}
  \lean{Pphi2.wickMonomial_four}\leanok
  \uses{Pphi2-wickMonomial}
  \texttt{Pphi2.wickMonomial\_four}
\end{theorem}
\begin{theorem}[\texttt{wickMonomial\_continuous₂}]\label{Pphi2-wickMonomial-continuous-}
  \lean{Pphi2.wickMonomial_continuous₂}\leanok
  \uses{Pphi2-wickMonomial}
  Joint continuity of $(c, x) \mapsto :x^n:_c$ and $(c, x) \mapsto :P(x):_c$. Proved by induction on the recursion.
\end{theorem}
\begin{theorem}[\texttt{wickPolynomial\_uniform\_bounded\_below}]\label{Pphi2-wickPolynomial-uniform-bounded-below}
  \lean{Pphi2.wickPolynomial_uniform_bounded_below}\leanok
  \uses{Pphi2-wickPolynomial-continuous-, Lattice-toSemilatticeInf, Pphi2-wickPolynomial}
  For $c \in [0, C]$, the bound $:P(x):_c \ge -A$ holds uniformly: the constant $A$ depends on $P$ and $C$ but not on $c$ individually. Proved by coefficient continuity + compactness of $[0,C] \times [-R, R]$.
\end{theorem}
\section{\texttt{Pphi2.AsymTorus.AsymTransferKernelOperator}}
\begin{theorem}[\texttt{asymTransferKernel\_kPow\_apply}]\label{Pphi2-asymTransferKernel-kPow-apply}
  \lean{Pphi2.asymTransferKernel_kPow_apply}\leanok
  \uses{Pphi2-L2SpatialField, Pphi2-SpatialField, Pphi2-asymTransferKernel, Pphi2-asymTransferSystem, Pphi2-asymTransferOperatorCLM-apply, Pphi2-asymTransferOperatorCLM}
  \texttt{Pphi2.asymTransferKernel\_kPow\_apply}
\end{theorem}
\begin{theorem}[\texttt{asymTransferOperatorCLM\_apply}]\label{Pphi2-asymTransferOperatorCLM-apply}
  \lean{Pphi2.asymTransferOperatorCLM_apply}\leanok
  \uses{Pphi2-asymTransferWeight-measurable, Pphi2-L2SpatialField, Pphi2-SpatialField, Pphi2-asymTransferKernel, Pphi2-asymTransferWeight, Pphi2-transferGaussian-memLp, Pphi2-asymTransferWeight-bound, Pphi2-asymTransferOperatorCLM, Pphi2-transferGaussian}
  \texttt{Pphi2.asymTransferOperatorCLM\_apply}
\end{theorem}
\begin{theorem}[\texttt{congr\_simp}]\label{Pphi2-asymTransferSystem-congr-simp}
  \lean{Pphi2.asymTransferSystem.congr_simp}\leanok
  \uses{Pphi2-SpatialField, Pphi2-asymTransferSystem}
  \texttt{Pphi2.asymTransferSystem.congr\_simp}
\end{theorem}
\section{\texttt{Pphi2.AsymTorus.AsymTorusEmbeddingIso}}
\begin{theorem}[\texttt{asymTorusEmbedLiftIso\_measurable}]\label{Pphi2-asymTorusEmbedLiftIso-measurable}
  \lean{Pphi2.asymTorusEmbedLiftIso_measurable}\leanok
  \uses{GaussianField-NuclearTensorProduct-instModuleReal, GaussianField-NuclearTensorProduct-instTopologicalSpace, GaussianField-AsymLatticeSites, Pphi2-evalAsymTorusAtSiteGJ, GaussianField-NuclearTensorProduct-instIsTopologicalAddGroup, Pphi2-asymTorusEmbedLiftIso, GaussianField-configuration-measurable-of-eval-measurable, GaussianField-NuclearTensorProduct-instAddCommGroup, GaussianField-configuration-eval-measurable, GaussianField-Configuration, GaussianField-NuclearTensorProduct-instContinuousSMulReal, GaussianField-AsymLatticeField, GaussianField-SmoothMap-Circle, GaussianField-instMeasurableSpaceConfiguration, GaussianField-instFintypeAsymLatticeSitesOfNeZeroNat, GaussianField-asymLatticeDelta, GaussianField-AsymTorusTestFunction}
  \texttt{Pphi2.asymTorusEmbedLiftIso\_measurable}
\end{theorem}
\begin{definition}[\texttt{evalAsymTorusAtSiteGJ}]\label{Pphi2-evalAsymTorusAtSiteGJ}
  \lean{Pphi2.evalAsymTorusAtSiteGJ}\leanok
  \uses{GaussianField-NuclearTensorProduct-instModuleReal, GaussianField-NuclearTensorProduct-instTopologicalSpace, GaussianField-AsymLatticeSites, GaussianField-evalAsymTorusAtSite, GaussianField-NuclearTensorProduct-instAddCommGroup, GaussianField-SmoothMap-Circle, GaussianField-AsymTorusTestFunction}
  \texttt{Pphi2.evalAsymTorusAtSiteGJ}
\end{definition}
\begin{theorem}[\texttt{second\_moment\_asym\_eq\_spectral}]\label{Pphi2-second-moment-asym-eq-spectral}
  \lean{Pphi2.second_moment_asym_eq_spectral}\leanok
  \uses{GaussianField-spectralLatticeCovarianceAsym, GaussianField-NuclearTensorProduct-instModuleReal, GaussianField-NuclearTensorProduct-instTopologicalSpace, GaussianField-AsymLatticeSites, GaussianField-covariance, Pphi2-asymLatticeTestFnIso, GaussianField-evalAsymTorusAtSite, GaussianField-NuclearTensorProduct-instAddCommGroup, GaussianField-AsymLatticeField, GaussianField-SmoothMap-Circle, GaussianField-latticeCovarianceAsymGJ, Pphi2-evalAsymTorusAtSiteGJ-apply, GaussianField-AsymTorusTestFunction, GaussianField-ell2-}
  \texttt{Pphi2.second\_moment\_asym\_eq\_spectral}
\end{theorem}
\begin{theorem}[\texttt{asymLatticeTestFnIso\_expand}]\label{Pphi2-asymLatticeTestFnIso-expand}
  \lean{Pphi2.asymLatticeTestFnIso_expand}\leanok
  \uses{GaussianField-AsymLatticeSites, Pphi2-asymLatticeTestFnIso, GaussianField-AsymLatticeField, GaussianField-instFintypeAsymLatticeSitesOfNeZeroNat, GaussianField-asymLatticeDelta, GaussianField-AsymTorusTestFunction}
  \texttt{Pphi2.asymLatticeTestFnIso\_expand}
\end{theorem}
\begin{theorem}[\texttt{asymTorusEmbedLiftIso\_eval\_eq}]\label{Pphi2-asymTorusEmbedLiftIso-eval-eq}
  \lean{Pphi2.asymTorusEmbedLiftIso_eval_eq}\leanok
  \uses{GaussianField-NuclearTensorProduct-instModuleReal, GaussianField-NuclearTensorProduct-instTopologicalSpace, GaussianField-AsymLatticeSites, Pphi2-evalAsymTorusAtSiteGJ, GaussianField-NuclearTensorProduct-instIsTopologicalAddGroup, Pphi2-asymLatticeTestFnIso, Pphi2-asymTorusEmbedLiftIso, GaussianField-NuclearTensorProduct-instAddCommGroup, GaussianField-Configuration, GaussianField-NuclearTensorProduct-instContinuousSMulReal, GaussianField-AsymLatticeField, GaussianField-SmoothMap-Circle, GaussianField-instFintypeAsymLatticeSitesOfNeZeroNat, Pphi2-asymLatticeTestFnIso-expand, GaussianField-asymLatticeDelta, GaussianField-AsymTorusTestFunction}
  \texttt{Pphi2.asymTorusEmbedLiftIso\_eval\_eq}
\end{theorem}
\begin{theorem}[\texttt{second\_moment\_asym\_tendsto}]\label{Pphi2-second-moment-asym-tendsto}
  \lean{Pphi2.second_moment_asym_tendsto}\leanok
  \uses{GaussianField-spectralLatticeCovarianceAsym, GaussianField-NuclearTensorProduct-instModuleReal, GaussianField-NuclearTensorProduct-instTopologicalSpace, GaussianField-AsymLatticeSites, Pphi2-asymTorusContinuumGreen, GaussianField-covariance, Pphi2-asymLatticeTestFnIso, GaussianField-evalAsymTorusAtSite, GaussianField-NuclearTensorProduct-instAddCommGroup, GaussianField-lattice-green-tendsto-continuum-asym, GaussianField-AsymLatticeField, GaussianField-SmoothMap-Circle, GaussianField-latticeCovarianceAsymGJ, GaussianField-AsymTorusTestFunction, GaussianField-ell2-, Pphi2-second-moment-asym-eq-spectral}
  \texttt{Pphi2.second\_moment\_asym\_tendsto}
\end{theorem}
\begin{definition}[\texttt{asymLatticeTestFnIso}]\label{Pphi2-asymLatticeTestFnIso}
  \lean{Pphi2.asymLatticeTestFnIso}\leanok
  \uses{GaussianField-NuclearTensorProduct-instModuleReal, GaussianField-NuclearTensorProduct-instTopologicalSpace, GaussianField-AsymLatticeSites, Pphi2-evalAsymTorusAtSiteGJ, GaussianField-NuclearTensorProduct-instAddCommGroup, GaussianField-AsymLatticeField, GaussianField-SmoothMap-Circle, GaussianField-AsymTorusTestFunction}
  \texttt{Pphi2.asymLatticeTestFnIso}
\end{definition}
\begin{definition}[\texttt{asymTorusEmbedLiftIso}]\label{Pphi2-asymTorusEmbedLiftIso}
  \lean{Pphi2.asymTorusEmbedLiftIso}\leanok
  \uses{GaussianField-NuclearTensorProduct-instModuleReal, GaussianField-NuclearTensorProduct-instTopologicalSpace, GaussianField-AsymLatticeSites, Pphi2-evalAsymTorusAtSiteGJ, GaussianField-NuclearTensorProduct-instIsTopologicalAddGroup, GaussianField-NuclearTensorProduct-instAddCommGroup, GaussianField-Configuration, GaussianField-NuclearTensorProduct-instContinuousSMulReal, GaussianField-AsymLatticeField, GaussianField-SmoothMap-Circle, GaussianField-instFintypeAsymLatticeSitesOfNeZeroNat, GaussianField-asymLatticeDelta, GaussianField-AsymTorusTestFunction}
  \texttt{Pphi2.asymTorusEmbedLiftIso}
\end{definition}
\begin{theorem}[\texttt{congr\_simp}]\label{Pphi2-asymTorusContinuumGreen-congr-simp}
  \lean{Pphi2.asymTorusContinuumGreen.congr_simp}\leanok
  \uses{Pphi2-asymTorusContinuumGreen, GaussianField-AsymTorusTestFunction}
  \texttt{Pphi2.asymTorusContinuumGreen.congr\_simp}
\end{theorem}
\begin{theorem}[\texttt{congr\_simp}]\label{Pphi2-asymLatticeTestFnIso-congr-simp}
  \lean{Pphi2.asymLatticeTestFnIso.congr_simp}\leanok
  \uses{GaussianField-AsymLatticeSites, Pphi2-asymLatticeTestFnIso, GaussianField-AsymTorusTestFunction}
  \texttt{Pphi2.asymLatticeTestFnIso.congr\_simp}
\end{theorem}
\begin{definition}[\texttt{asymTorusInteractingMeasureIso}]\label{Pphi2-asymTorusInteractingMeasureIso}
  \lean{Pphi2.asymTorusInteractingMeasureIso}\leanok
  \uses{GaussianField-NuclearTensorProduct-instModuleReal, GaussianField-NuclearTensorProduct-instTopologicalSpace, GaussianField-AsymLatticeSites, GaussianField-NuclearTensorProduct-instIsTopologicalAddGroup, Pphi2-asymTorusEmbedLiftIso, GaussianField-NuclearTensorProduct-instAddCommGroup, GaussianField-Configuration, GaussianField-NuclearTensorProduct-instContinuousSMulReal, GaussianField-AsymLatticeField, GaussianField-SmoothMap-Circle, Pphi2-interactingLatticeMeasureAsym, GaussianField-instMeasurableSpaceConfiguration, GaussianField-AsymTorusTestFunction}
  \texttt{Pphi2.asymTorusInteractingMeasureIso}
\end{definition}
\begin{theorem}[\texttt{congr\_simp}]\label{Pphi2-evalAsymTorusAtSiteGJ-congr-simp}
  \lean{Pphi2.evalAsymTorusAtSiteGJ.congr_simp}\leanok
  \uses{GaussianField-NuclearTensorProduct-instModuleReal, GaussianField-NuclearTensorProduct-instTopologicalSpace, GaussianField-AsymLatticeSites, Pphi2-evalAsymTorusAtSiteGJ, GaussianField-NuclearTensorProduct-instAddCommGroup, GaussianField-SmoothMap-Circle, GaussianField-AsymTorusTestFunction}
  \texttt{Pphi2.evalAsymTorusAtSiteGJ.congr\_simp}
\end{theorem}
\begin{theorem}[\texttt{evalAsymTorusAtSiteGJ\_apply}]\label{Pphi2-evalAsymTorusAtSiteGJ-apply}
  \lean{Pphi2.evalAsymTorusAtSiteGJ_apply}\leanok
  \uses{GaussianField-NuclearTensorProduct-instModuleReal, GaussianField-NuclearTensorProduct-instTopologicalSpace, GaussianField-AsymLatticeSites, Pphi2-evalAsymTorusAtSiteGJ, GaussianField-evalAsymTorusAtSite, GaussianField-NuclearTensorProduct-instAddCommGroup, GaussianField-SmoothMap-Circle, GaussianField-AsymTorusTestFunction}
  \texttt{Pphi2.evalAsymTorusAtSiteGJ\_apply}
\end{theorem}
\section{\texttt{Pphi2.InteractingMeasure.U4Derivative}}
\begin{theorem}[\texttt{exists\_pos\_lt\_zero\_of\_hasDerivWithinAt\_neg}]\label{Pphi2-exists-pos-lt-zero-of-hasDerivWithinAt-neg}
  \lean{Pphi2.exists_pos_lt_zero_of_hasDerivWithinAt_neg}\leanok
  \uses{Lattice-toSemilatticeSup}
  \texttt{Pphi2.exists\_pos\_lt\_zero\_of\_hasDerivWithinAt\_neg}
\end{theorem}
\begin{theorem}[\texttt{deriv\_affine\_bound\_neg}]\label{Pphi2-deriv-affine-bound-neg}
  \lean{Pphi2.deriv_affine_bound_neg}\leanok
  \texttt{Pphi2.deriv\_affine\_bound\_neg}
\end{theorem}
\begin{theorem}[\texttt{u4\_at\_zero}]\label{Pphi2-u4-at-zero}
  \lean{Pphi2.u4_at_zero}\leanok
  \uses{Pphi2-partitionFn-zero, GaussianField-FinLatticeSites, Pphi2-latticeFourthMoment-eq, GaussianField-latticeGaussianMeasure, GaussianField-Configuration, GaussianField-instMeasurableSpaceConfiguration, GaussianField-FinLatticeField, Pphi2-interactionFunctional}
  \texttt{Pphi2.u4\_at\_zero}
\end{theorem}
\begin{theorem}[\texttt{quartic\_kernel\_sum\_nonneg}]\label{Pphi2-quartic-kernel-sum-nonneg}
  \lean{Pphi2.quartic_kernel_sum_nonneg}\leanok
  \uses{GaussianField-FinLatticeSites, GaussianField-gffPositionCovariance, GaussianField-FinLatticeField, Lattice-toSemilatticeInf}
  \texttt{Pphi2.quartic\_kernel\_sum\_nonneg}
\end{theorem}
\begin{theorem}[\texttt{wickMonomial\_four\_eq}]\label{Pphi2-wickMonomial-four-eq}
  \lean{Pphi2.wickMonomial_four_eq}\leanok
  \texttt{Pphi2.wickMonomial\_four\_eq}
\end{theorem}
\begin{theorem}[\texttt{exists\_pos\_u4\_neg}]\label{Pphi2-exists-pos-u4-neg}
  \lean{Pphi2.exists_pos_u4_neg}\leanok
  \uses{Pphi2-u4-hasDerivWithinAt, GaussianField-FinLatticeSites, Pphi2-u4-at-zero, GaussianField-gffPositionCovariance, Pphi2-u4-slope-neg, GaussianField-latticeGaussianMeasure, GaussianField-Configuration, Pphi2-exists-pos-lt-zero-of-hasDerivWithinAt-neg, GaussianField-instMeasurableSpaceConfiguration, GaussianField-FinLatticeField, Pphi2-interactionFunctional}
  \texttt{Pphi2.exists\_pos\_u4\_neg}
\end{theorem}
\begin{theorem}[\texttt{latticeFourthMoment\_eq}]\label{Pphi2-latticeFourthMoment-eq}
  \lean{Pphi2.latticeFourthMoment_eq}\leanok
  \uses{GaussianField-finLatticeField-dyninMityaginSpace, GaussianField-measure, GaussianField-ell2-separableSpace, GaussianField-FinLatticeSites, Pphi2-connectedFourPoint, GaussianField-latticeCovarianceGJ, GaussianField-latticeGaussianMeasure, GaussianField-Configuration, GaussianField-instMeasurableSpaceConfiguration, Pphi2-connectedFourPoint-gaussianMeasure-eq-zero, GaussianField-ell2-, GaussianField-FinLatticeField}
  \texttt{Pphi2.latticeFourthMoment\_eq}
\end{theorem}
\begin{theorem}[\texttt{latticeSecondMoment\_eq\_smearedCovariance}]\label{Pphi2-latticeSecondMoment-eq-smearedCovariance}
  \lean{Pphi2.latticeSecondMoment_eq_smearedCovariance}\leanok
  \uses{GaussianField-FinLatticeSites, GaussianField-gffSmearedCoeff, GaussianField-latticeGaussianMeasure, GaussianField-Configuration, GaussianField-gff-wickPower-two-smeared-inner, GaussianField-instMeasurableSpaceConfiguration, GaussianField-FinLatticeField, GaussianField-gffSmearedCovariance}
  \texttt{Pphi2.latticeSecondMoment\_eq\_smearedCovariance}
\end{theorem}
\begin{theorem}[\texttt{exists\_u4\_slope\_neg}]\label{Pphi2-exists-u4-slope-neg}
  \lean{Pphi2.exists_u4_slope_neg}\leanok
  \uses{GaussianField-FinLatticeSites, GaussianField-gffPositionCovariance, Pphi2-u4-slope-neg, Pphi2-wickConstant-pos, Pphi2-wickConstant, GaussianField-FinLatticeField, Pphi2-wickConstant-eq-gffPositionCovariance}
  \texttt{Pphi2.exists\_u4\_slope\_neg}
\end{theorem}
\begin{theorem}[\texttt{u4\_hasDerivWithinAt}]\label{Pphi2-u4-hasDerivWithinAt}
  \lean{Pphi2.u4_hasDerivWithinAt}\leanok
  \uses{Pphi2-wickFourth-interaction-inner-quartic, Pphi2-partitionFn-zero, Pphi2-interactionFunctional-eq-single, Pphi2-normalizedMoment-hasDerivWithinAt, GaussianField-FinLatticeSites, Pphi2-integrable-powMul-interaction, Pphi2-latticeFourthMoment-eq, GaussianField-gffPositionCovariance, GaussianField-latticeGaussianMeasure, GaussianField-Configuration, Pphi2-wickConstant, GaussianField-instMeasurableSpaceConfiguration, Pphi2-latticeSecondMoment-eq-smearedCovariance, Pphi2-wickMonomial-four-eq, GaussianField-FinLatticeField, Pphi2-interactionFunctional, GaussianField-gffSmearedCovariance, Pphi2-wickPolynomial}
  \texttt{Pphi2.u4\_hasDerivWithinAt}
\end{theorem}
\begin{theorem}[\texttt{u4\_slope\_neg}]\label{Pphi2-u4-slope-neg}
  \lean{Pphi2.u4_slope_neg}\leanok
  \uses{Pphi2-quartic-kernel-sum-pos, GaussianField-FinLatticeSites, GaussianField-gffPositionCovariance, GaussianField-FinLatticeField}
  \texttt{Pphi2.u4\_slope\_neg}
\end{theorem}
\begin{theorem}[\texttt{quartic\_kernel\_sum\_pos}]\label{Pphi2-quartic-kernel-sum-pos}
  \lean{Pphi2.quartic_kernel_sum_pos}\leanok
  \uses{GaussianField-FinLatticeSites, GaussianField-gffPositionCovariance, GaussianField-FinLatticeField, Lattice-toSemilatticeInf}
  \texttt{Pphi2.quartic\_kernel\_sum\_pos}
\end{theorem}
\section{\texttt{Pphi2.NelsonEstimate.AsymCovarianceBoundsGJ}}
\begin{definition}[\texttt{asymCanonicalFullInteractionJoint}]\label{Pphi2-asymCanonicalFullInteractionJoint}
  \lean{Pphi2.asymCanonicalFullInteractionJoint}\leanok
  \uses{GaussianField-AsymLatticeSites, Pphi2-wickConstantAsym, Pphi2-AsymCanonicalJoint, Pphi2-asymCanonicalSumFieldFunction, GaussianField-instFintypeAsymLatticeSitesOfNeZeroNat, Pphi2-wickPolynomial}
  \texttt{Pphi2.asymCanonicalFullInteractionJoint}
\end{definition}
\begin{theorem}[\texttt{congr\_simp}]\label{Pphi2-asymCanonicalRoughCovariance-congr-simp}
  \lean{Pphi2.asymCanonicalRoughCovariance.congr_simp}\leanok
  \uses{GaussianField-AsymLatticeSites, Pphi2-asymCanonicalRoughCovariance}
  \texttt{Pphi2.asymCanonicalRoughCovariance.congr\_simp}
\end{theorem}
\begin{definition}[\texttt{asymCanonicalRoughError}]\label{Pphi2-asymCanonicalRoughError}
  \lean{Pphi2.asymCanonicalRoughError}\leanok
  \uses{Pphi2-asymCanonicalFullInteractionJoint, Pphi2-AsymCanonicalJoint, Pphi2-asymCanonicalSmoothInteraction}
  \texttt{Pphi2.asymCanonicalRoughError}
\end{definition}
\begin{theorem}[\texttt{asym\_heat\_kernel\_trace\_bound}]\label{Pphi2-asym-heat-kernel-trace-bound}
  \lean{Pphi2.asym_heat_kernel_trace_bound}\leanok
  \uses{Pphi2-heat-kernel-1d-bound, GaussianField-latticeEigenvalue1d, Pphi2-asym-heat-kernel-trace-factorization, Pphi2-asymCanonicalEigenvalue}
  \texttt{Pphi2.asym\_heat\_kernel\_trace\_bound}
\end{theorem}
\begin{definition}[\texttt{asymSmoothWickConstant}]\label{Pphi2-asymSmoothWickConstant}
  \lean{Pphi2.asymSmoothWickConstant}\leanok
  \uses{GaussianField-AsymLatticeSites, GaussianField-instFintypeAsymLatticeSitesOfNeZeroNat, Pphi2-asymCanonicalSmoothWeight}
  \texttt{Pphi2.asymSmoothWickConstant}
\end{definition}
\begin{theorem}[\texttt{asymCanonicalRoughCovariance\_pow\_two\_sum\_le\_uniform}]\label{Pphi2-asymCanonicalRoughCovariance-pow-two-sum-le-uniform}
  \lean{Pphi2.asymCanonicalRoughCovariance_pow_two_sum_le_uniform}\leanok
  \uses{GaussianField-AsymLatticeSites, Pphi2-asymCanonicalEigenvalue, GaussianField-instFintypeAsymLatticeSitesOfNeZeroNat, Lattice-toSemilatticeInf, Pphi2-asymCanonicalRoughCovariance}
  \texttt{Pphi2.asymCanonicalRoughCovariance\_pow\_two\_sum\_le\_uniform}
\end{theorem}
\begin{definition}[\texttt{asymCanonicalSmoothInteraction}]\label{Pphi2-asymCanonicalSmoothInteraction}
  \lean{Pphi2.asymCanonicalSmoothInteraction}\leanok
  \uses{GaussianField-AsymLatticeSites, Pphi2-AsymCanonicalJoint, Pphi2-asymSmoothWickConstant, GaussianField-instFintypeAsymLatticeSitesOfNeZeroNat, Pphi2-asymCanonicalSmoothFieldFunction, Pphi2-wickPolynomial}
  \texttt{Pphi2.asymCanonicalSmoothInteraction}
\end{definition}
\begin{definition}[\texttt{asymCanonicalRoughCovariance}]\label{Pphi2-asymCanonicalRoughCovariance}
  \lean{Pphi2.asymCanonicalRoughCovariance}\leanok
  \uses{GaussianField-AsymLatticeSites, Pphi2-asymCanonicalBasis, Pphi2-asymCanonicalRoughWeight, Pphi2-asymCanonicalNormSq}
  \texttt{Pphi2.asymCanonicalRoughCovariance}
\end{definition}
\begin{theorem}[\texttt{asym\_heat\_kernel\_trace\_factorization}]\label{Pphi2-asym-heat-kernel-trace-factorization}
  \lean{Pphi2.asym_heat_kernel_trace_factorization}\leanok
  \uses{GaussianField-latticeEigenvalue1d, Pphi2-asymCanonicalEigenvalue}
  \texttt{Pphi2.asym\_heat\_kernel\_trace\_factorization}
\end{theorem}
\begin{theorem}[\texttt{congr\_simp}]\label{Pphi2-asymSmoothWickConstant-congr-simp}
  \lean{Pphi2.asymSmoothWickConstant.congr_simp}\leanok
  \uses{Pphi2-asymSmoothWickConstant}
  \texttt{Pphi2.asymSmoothWickConstant.congr\_simp}
\end{theorem}
\begin{theorem}[\texttt{asym\_heat\_kernel\_trace\_bound\_uniform}]\label{Pphi2-asym-heat-kernel-trace-bound-uniform}
  \lean{Pphi2.asym_heat_kernel_trace_bound_uniform}\leanok
  \uses{GaussianField-latticeEigenvalue1d, Pphi2-asym-heat-kernel-trace-factorization, Pphi2-heat-kernel-1d-bound-uniform, Pphi2-asymCanonicalEigenvalue}
  \texttt{Pphi2.asym\_heat\_kernel\_trace\_bound\_uniform}
\end{theorem}
\begin{theorem}[\texttt{asymCanonicalRoughCovariance\_pow\_one\_sum\_le\_uniform}]\label{Pphi2-asymCanonicalRoughCovariance-pow-one-sum-le-uniform}
  \lean{Pphi2.asymCanonicalRoughCovariance_pow_one_sum_le_uniform}\leanok
  \uses{GaussianField-latticeEigenvalue1d, GaussianField-AsymLatticeSites, Pphi2-asymCanonicalBasis, Pphi2-schwinger-rough, GaussianField-instFintypeAsymLatticeSitesOfNeZeroNat, Pphi2-asymCanonicalRoughCovariance, Pphi2-asymCanonicalNormSq}
  \texttt{Pphi2.asymCanonicalRoughCovariance\_pow\_one\_sum\_le\_uniform}
\end{theorem}
\begin{theorem}[\texttt{asymSmoothWickConstant\_le\_log\_uniform}]\label{Pphi2-asymSmoothWickConstant-le-log-uniform}
  \lean{Pphi2.asymSmoothWickConstant_le_log_uniform}\leanok
  \uses{GaussianField-latticeEigenvalue1d, GaussianField-AsymLatticeSites, Pphi2-asym-heat-kernel-trace-bound-uniform, Pphi2-schwinger-smooth-Ioi, GaussianField-latticeEigenvalue1d-nonneg, Pphi2-asymCanonicalEigenvalue, Pphi2-asymSmoothWickConstant, GaussianField-instFintypeAsymLatticeSitesOfNeZeroNat, Pphi2-asymCanonicalSmoothWeight, Lattice-toSemilatticeInf}
  \texttt{Pphi2.asymSmoothWickConstant\_le\_log\_uniform}
\end{theorem}
\section{\texttt{Pphi2.AsymTorus.AsymEnergyFactorization}}
\begin{theorem}[\texttt{timeCoupling\_sum\_slice}]\label{Pphi2-timeCoupling-sum-slice}
  \lean{Pphi2.timeCoupling_sum_slice}\leanok
  \uses{GaussianField-AsymLatticeSites, Pphi2-SpatialField, Pphi2-asymSliceEquiv, GaussianField-AsymLatticeField, Pphi2-timeCoupling}
  \texttt{Pphi2.timeCoupling\_sum\_slice}
\end{theorem}
\begin{theorem}[\texttt{asym\_timeBond\_sum\_slice}]\label{Pphi2-asym-timeBond-sum-slice}
  \lean{Pphi2.asym_timeBond_sum_slice}\leanok
  \uses{GaussianField-AsymLatticeSites, Pphi2-SpatialField, Pphi2-asymSliceEquiv, Pphi2-asymSliceEquiv-apply, GaussianField-AsymLatticeField, GaussianField-instFintypeAsymLatticeSitesOfNeZeroNat}
  \texttt{Pphi2.asym\_timeBond\_sum\_slice}
\end{theorem}
\begin{theorem}[\texttt{spatialKinetic\_sum\_slice}]\label{Pphi2-spatialKinetic-sum-slice}
  \lean{Pphi2.spatialKinetic_sum_slice}\leanok
  \uses{GaussianField-AsymLatticeSites, Pphi2-SpatialField, Pphi2-asymSliceEquiv, Pphi2-spatialKinetic, GaussianField-AsymLatticeField}
  \texttt{Pphi2.spatialKinetic\_sum\_slice}
\end{theorem}
\begin{theorem}[\texttt{asym\_spaceBond\_sum\_slice}]\label{Pphi2-asym-spaceBond-sum-slice}
  \lean{Pphi2.asym_spaceBond_sum_slice}\leanok
  \uses{GaussianField-AsymLatticeSites, Pphi2-SpatialField, Pphi2-asymSliceEquiv, Pphi2-asymSliceEquiv-apply, GaussianField-AsymLatticeField, GaussianField-instFintypeAsymLatticeSitesOfNeZeroNat}
  \texttt{Pphi2.asym\_spaceBond\_sum\_slice}
\end{theorem}
\begin{theorem}[\texttt{asym\_sbp\_direction}]\label{Pphi2-asym-sbp-direction}
  \lean{Pphi2.asym_sbp_direction}\leanok
  \uses{GaussianField-AsymLatticeSites, GaussianField-AsymLatticeField, GaussianField-instFintypeAsymLatticeSitesOfNeZeroNat}
  \texttt{Pphi2.asym\_sbp\_direction}
\end{theorem}
\begin{theorem}[\texttt{congr\_simp}]\label{Pphi2-asymSliceEquiv-congr-simp}
  \lean{Pphi2.asymSliceEquiv.congr_simp}\leanok
  \uses{GaussianField-AsymLatticeSites, Pphi2-SpatialField, Pphi2-asymSliceEquiv, GaussianField-AsymLatticeField}
  \texttt{Pphi2.asymSliceEquiv.congr\_simp}
\end{theorem}
\begin{theorem}[\texttt{energy\_exponent\_factorization}]\label{Pphi2-energy-exponent-factorization}
  \lean{Pphi2.energy_exponent_factorization}\leanok
  \uses{Pphi2-asym-onsite-sum-slice, Pphi2-spatialAction, GaussianField-AsymLatticeSites, Pphi2-wickConstantAsym, Pphi2-spatialKinetic-sum-slice, Pphi2-SpatialField, Pphi2-asymSliceEquiv, Pphi2-massOperatorAsym-quadratic-form-slice, Pphi2-spatialKinetic, Pphi2-spatialPotential, Pphi2-timeCoupling-sum-slice, GaussianField-AsymLatticeField, GaussianField-instFintypeAsymLatticeSitesOfNeZeroNat, GaussianField-massOperatorAsym, Pphi2-timeCoupling, Pphi2-wickPolynomial}
  \texttt{Pphi2.energy\_exponent\_factorization}
\end{theorem}
\begin{theorem}[\texttt{periodicPathDensity\_asymTransferKernel\_eq\_exp}]\label{Pphi2-periodicPathDensity-asymTransferKernel-eq-exp}
  \lean{Pphi2.periodicPathDensity_asymTransferKernel_eq_exp}\leanok
  \uses{Pphi2-spatialAction, GaussianField-AsymLatticeSites, Pphi2-wickConstantAsym, Pphi2-SpatialField, Pphi2-asymTransferKernel, Pphi2-asymSliceEquiv, GaussianField-AsymLatticeField, GaussianField-instFintypeAsymLatticeSitesOfNeZeroNat, GaussianField-massOperatorAsym, Pphi2-energy-exponent-factorization, Pphi2-timeCoupling, Pphi2-wickPolynomial}
  \texttt{Pphi2.periodicPathDensity\_asymTransferKernel\_eq\_exp}
\end{theorem}
\begin{theorem}[\texttt{asym\_onsite\_sum\_slice}]\label{Pphi2-asym-onsite-sum-slice}
  \lean{Pphi2.asym_onsite_sum_slice}\leanok
  \uses{GaussianField-AsymLatticeSites, Pphi2-SpatialField, Pphi2-asymSliceEquiv, Pphi2-asymSliceEquiv-apply, GaussianField-AsymLatticeField, GaussianField-instFintypeAsymLatticeSitesOfNeZeroNat}
  \texttt{Pphi2.asym\_onsite\_sum\_slice}
\end{theorem}
\begin{theorem}[\texttt{massOperatorAsym\_quadratic\_form\_slice}]\label{Pphi2-massOperatorAsym-quadratic-form-slice}
  \lean{Pphi2.massOperatorAsym_quadratic_form_slice}\leanok
  \uses{GaussianField-AsymLatticeSites, Pphi2-asym-sq-sum-slice, Pphi2-SpatialField, Pphi2-asymSliceEquiv, Pphi2-massOperatorAsym-quadratic-form-bonds, GaussianField-AsymLatticeField, GaussianField-instFintypeAsymLatticeSitesOfNeZeroNat, GaussianField-massOperatorAsym, Pphi2-asym-timeBond-sum-slice, Pphi2-asym-spaceBond-sum-slice}
  \texttt{Pphi2.massOperatorAsym\_quadratic\_form\_slice}
\end{theorem}
\begin{theorem}[\texttt{asym\_sq\_sum\_slice}]\label{Pphi2-asym-sq-sum-slice}
  \lean{Pphi2.asym_sq_sum_slice}\leanok
  \uses{GaussianField-AsymLatticeSites, Pphi2-SpatialField, Pphi2-asymSliceEquiv, Pphi2-asymSliceEquiv-apply, GaussianField-AsymLatticeField, GaussianField-instFintypeAsymLatticeSitesOfNeZeroNat}
  \texttt{Pphi2.asym\_sq\_sum\_slice}
\end{theorem}
\begin{theorem}[\texttt{massOperatorAsym\_quadratic\_form\_bonds}]\label{Pphi2-massOperatorAsym-quadratic-form-bonds}
  \lean{Pphi2.massOperatorAsym_quadratic_form_bonds}\leanok
  \uses{Pphi2-asym-sbp-direction, GaussianField-AsymLatticeSites, GaussianField-finiteLaplacianAsym, GaussianField-finiteLaplacianAsymFun, GaussianField-AsymLatticeField, GaussianField-instFintypeAsymLatticeSitesOfNeZeroNat, GaussianField-massOperatorAsym}
  \texttt{Pphi2.massOperatorAsym\_quadratic\_form\_bonds}
\end{theorem}
\section{\texttt{Pphi2.AsymTorus.AsymObsTrunc}}
\begin{theorem}[\texttt{norm\_sq\_proj\_obsTrunc\_omega\_le}]\label{Pphi2-norm-sq-proj-obsTrunc-omega-le}
  \lean{Pphi2.norm_sq_proj_obsTrunc_omega_le}\leanok
  \uses{Pphi2-asymSliceObsLinear, Pphi2-asymSliceObsTrunc-abs-le-linear, Pphi2-L2SpatialField, Pphi2-asymSliceObsTrunc-measurable, Pphi2-SpatialField, Pphi2-asymGroundVector, Pphi2-norm-sub-groundProj-le, Pphi2-asymSliceObsTruncMulCLM, Pphi2-asymSliceObsTrunc, Pphi2-asymSliceObsTruncMulCLM-coeFn, Lattice-toSemilatticeInf}
  \texttt{Pphi2.norm\_sq\_proj\_obsTrunc\_omega\_le}
\end{theorem}
\begin{definition}[\texttt{asymSliceObsTrunc}]\label{Pphi2-asymSliceObsTrunc}
  \lean{Pphi2.asymSliceObsTrunc}\leanok
  \uses{Pphi2-asymSliceObsLinear, Pphi2-SpatialField}
  \texttt{Pphi2.asymSliceObsTrunc}
\end{definition}
\begin{theorem}[\texttt{asymSliceObsTruncMulCLM\_coeFn}]\label{Pphi2-asymSliceObsTruncMulCLM-coeFn}
  \lean{Pphi2.asymSliceObsTruncMulCLM_coeFn}\leanok
  \uses{Pphi2-L2SpatialField, Pphi2-asymSliceObsTrunc-measurable, Pphi2-SpatialField, Pphi2-asymSliceObsTruncMulCLM, Pphi2-asymSliceObsTrunc, Pphi2-asymSliceObsTrunc-abs-le-bound}
  \texttt{Pphi2.asymSliceObsTruncMulCLM\_coeFn}
\end{theorem}
\begin{definition}[\texttt{asymSliceObsTruncMulCLM}]\label{Pphi2-asymSliceObsTruncMulCLM}
  \lean{Pphi2.asymSliceObsTruncMulCLM}\leanok
  \uses{Pphi2-L2SpatialField, Pphi2-asymSliceObsTrunc-measurable, Pphi2-SpatialField, Pphi2-asymSliceObsTrunc}
  \texttt{Pphi2.asymSliceObsTruncMulCLM}
\end{definition}
\begin{theorem}[\texttt{asymSliceObsTruncMulCLM\_isSelfAdjoint}]\label{Pphi2-asymSliceObsTruncMulCLM-isSelfAdjoint}
  \lean{Pphi2.asymSliceObsTruncMulCLM_isSelfAdjoint}\leanok
  \uses{Pphi2-L2SpatialField, Pphi2-asymSliceObsTrunc-measurable, Pphi2-SpatialField, Pphi2-asymSliceObsTruncMulCLM, Pphi2-asymSliceObsTrunc, Pphi2-asymSliceObsTrunc-abs-le-bound}
  \texttt{Pphi2.asymSliceObsTruncMulCLM\_isSelfAdjoint}
\end{theorem}
\begin{definition}[\texttt{asymSliceObsLinear}]\label{Pphi2-asymSliceObsLinear}
  \lean{Pphi2.asymSliceObsLinear}\leanok
  \uses{Pphi2-SpatialField}
  \texttt{Pphi2.asymSliceObsLinear}
\end{definition}
\begin{theorem}[\texttt{asymSliceObsTrunc\_measurable}]\label{Pphi2-asymSliceObsTrunc-measurable}
  \lean{Pphi2.asymSliceObsTrunc_measurable}\leanok
  \uses{Pphi2-asymSliceObsLinear, Pphi2-SpatialField, Pphi2-asymSliceObsLinear-measurable, Pphi2-asymSliceObsTrunc}
  \texttt{Pphi2.asymSliceObsTrunc\_measurable}
\end{theorem}
\begin{theorem}[\texttt{asymSliceObsTrunc\_abs\_le\_linear}]\label{Pphi2-asymSliceObsTrunc-abs-le-linear}
  \lean{Pphi2.asymSliceObsTrunc_abs_le_linear}\leanok
  \uses{Pphi2-asymSliceObsLinear, Pphi2-SpatialField, Pphi2-asymSliceObsTrunc, Lattice-toSemilatticeInf}
  \texttt{Pphi2.asymSliceObsTrunc\_abs\_le\_linear}
\end{theorem}
\begin{definition}[\texttt{asymSliceObsTruncContract}]\label{Pphi2-asymSliceObsTruncContract}
  \lean{Pphi2.asymSliceObsTruncContract}\leanok
  \uses{Pphi2-asymSliceObsTrunc-measurable, Pphi2-asymSliceObsTruncMulCLM-isSelfAdjoint, Pphi2-SpatialField, Pphi2-asymSliceObsTruncMulCLM, Pphi2-asymSliceObsTrunc, Pphi2-asymSliceObsTruncMulCLM-coeFn}
  \texttt{Pphi2.asymSliceObsTruncContract}
\end{definition}
\begin{theorem}[\texttt{asymSliceObsLinear\_measurable}]\label{Pphi2-asymSliceObsLinear-measurable}
  \lean{Pphi2.asymSliceObsLinear_measurable}\leanok
  \uses{Pphi2-asymSliceObsLinear, Pphi2-SpatialField}
  \texttt{Pphi2.asymSliceObsLinear\_measurable}
\end{theorem}
\begin{theorem}[\texttt{asymSliceObsTrunc\_abs\_le\_bound}]\label{Pphi2-asymSliceObsTrunc-abs-le-bound}
  \lean{Pphi2.asymSliceObsTrunc_abs_le_bound}\leanok
  \uses{Pphi2-asymSliceObsLinear, Pphi2-SpatialField, Pphi2-asymSliceObsTrunc, Lattice-toSemilatticeInf}
  \texttt{Pphi2.asymSliceObsTrunc\_abs\_le\_bound}
\end{theorem}
\section{\texttt{Pphi2.AsymTorus.AsymWickVariance}}
\begin{theorem}[\texttt{congr\_simp}]\label{Pphi2-asymShift-congr-simp}
  \lean{Pphi2.asymShift.congr_simp}\leanok
  \uses{GaussianField-AsymLatticeSites, Pphi2-asymShift, GaussianField-AsymLatticeField}
  \texttt{Pphi2.asymShift.congr\_simp}
\end{theorem}
\begin{theorem}[\texttt{spectralVarianceAsym\_delta\_eq}]\label{Pphi2-spectralVarianceAsym-delta-eq}
  \lean{Pphi2.spectralVarianceAsym_delta_eq}\leanok
  \uses{GaussianField-spectralLatticeCovarianceAsym, GaussianField-AsymLatticeSites, Pphi2-covariance-spectralLatticeCovarianceAsym-translation-invariant, GaussianField-covariance, Pphi2-asymShift-delta, Pphi2-asymShift, GaussianField-AsymLatticeField, GaussianField-asymLatticeDelta, GaussianField-ell2-}
  \texttt{Pphi2.spectralVarianceAsym\_delta\_eq}
\end{theorem}
\begin{theorem}[\texttt{spectralCovAsym\_massOperator\_eq}]\label{Pphi2-spectralCovAsym-massOperator-eq}
  \lean{Pphi2.spectralCovAsym_massOperator_eq}\leanok
  \uses{GaussianField-spectralLatticeCovarianceAsym, GaussianField-AsymLatticeSites, GaussianField-covariance, GaussianField-massEigenvectorBasisAsym, GaussianField-massOperatorAsym-eigenCoeff-eq-eigenvalues-mul-eigenCoeff, GaussianField-massEigenbasisAsym-sum-mul-sum-eq-site-inner, GaussianField-massEigenvaluesAsym, GaussianField-covariance-spectralLatticeCovarianceAsym-eq, GaussianField-massOperatorMatrixAsym-eigenvalues-pos, GaussianField-AsymLatticeField, GaussianField-instFintypeAsymLatticeSitesOfNeZeroNat, GaussianField-massOperatorAsym, GaussianField-ell2-}
  \texttt{Pphi2.spectralCovAsym\_massOperator\_eq}
\end{theorem}
\begin{theorem}[\texttt{spectralVarianceAsym\_eq\_massEigenvalue\_average}]\label{Pphi2-spectralVarianceAsym-eq-massEigenvalue-average}
  \lean{Pphi2.spectralVarianceAsym_eq_massEigenvalue_average}\leanok
  \uses{GaussianField-spectralLatticeCovarianceAsym, GaussianField-AsymLatticeSites, GaussianField-covariance, Pphi2-spectralVarianceAsym-delta-eq, GaussianField-massEigenvaluesAsym, Pphi2-spectralVarianceAsym-sum-eq-massEigenvalue-trace, GaussianField-AsymLatticeField, GaussianField-instFintypeAsymLatticeSitesOfNeZeroNat, GaussianField-asymLatticeDelta, GaussianField-ell2-}
  \texttt{Pphi2.spectralVarianceAsym\_eq\_massEigenvalue\_average}
\end{theorem}
\begin{theorem}[\texttt{asymShift\_sum\_mul\_eq}]\label{Pphi2-asymShift-sum-mul-eq}
  \lean{Pphi2.asymShift_sum_mul_eq}\leanok
  \uses{GaussianField-AsymLatticeSites, Pphi2-asymShift, GaussianField-AsymLatticeField, GaussianField-instFintypeAsymLatticeSitesOfNeZeroNat}
  \texttt{Pphi2.asymShift\_sum\_mul\_eq}
\end{theorem}
\begin{theorem}[\texttt{finiteLaplacianAsym\_translation\_commute}]\label{Pphi2-finiteLaplacianAsym-translation-commute}
  \lean{Pphi2.finiteLaplacianAsym_translation_commute}\leanok
  \uses{GaussianField-AsymLatticeSites, GaussianField-finiteLaplacianAsym, Pphi2-asymShift, GaussianField-finiteLaplacianAsymFun, GaussianField-AsymLatticeField}
  \texttt{Pphi2.finiteLaplacianAsym\_translation\_commute}
\end{theorem}
\begin{theorem}[\texttt{massEigenvectorBasisAsym\_norm\_sq\_eq\_one}]\label{Pphi2-massEigenvectorBasisAsym-norm-sq-eq-one}
  \lean{Pphi2.massEigenvectorBasisAsym_norm_sq_eq_one}\leanok
  \uses{GaussianField-AsymLatticeSites, GaussianField-massEigenvectorBasisAsym, GaussianField-instFintypeAsymLatticeSitesOfNeZeroNat, Lattice-toSemilatticeInf}
  \texttt{Pphi2.massEigenvectorBasisAsym\_norm\_sq\_eq\_one}
\end{theorem}
\begin{theorem}[\texttt{spectralVarianceAsym\_sum\_eq\_massEigenvalue\_trace}]\label{Pphi2-spectralVarianceAsym-sum-eq-massEigenvalue-trace}
  \lean{Pphi2.spectralVarianceAsym_sum_eq_massEigenvalue_trace}\leanok
  \uses{GaussianField-spectralLatticeCovarianceAsym, Pphi2-massEigenvectorBasisAsym-norm-sq-eq-one, GaussianField-AsymLatticeSites, GaussianField-covariance, GaussianField-massEigenvectorBasisAsym, GaussianField-massEigenvaluesAsym, GaussianField-covariance-spectralLatticeCovarianceAsym-eq, GaussianField-AsymLatticeField, Pphi2-delta-massEigenvectorCoeff-asym, GaussianField-instFintypeAsymLatticeSitesOfNeZeroNat, GaussianField-asymLatticeDelta, GaussianField-ell2-}
  \texttt{Pphi2.spectralVarianceAsym\_sum\_eq\_massEigenvalue\_trace}
\end{theorem}
\begin{theorem}[\texttt{spectralVarianceAsym\_inner\_eq\_massEigenvalue\_average}]\label{Pphi2-spectralVarianceAsym-inner-eq-massEigenvalue-average}
  \lean{Pphi2.spectralVarianceAsym_inner_eq_massEigenvalue_average}\leanok
  \uses{GaussianField-spectralLatticeCovarianceAsym, Pphi2-spectralVarianceAsym-eq-massEigenvalue-average, GaussianField-AsymLatticeSites, GaussianField-covariance, GaussianField-massEigenvaluesAsym, GaussianField-AsymLatticeField, GaussianField-instFintypeAsymLatticeSitesOfNeZeroNat, GaussianField-asymLatticeDelta, GaussianField-ell2-}
  \texttt{Pphi2.spectralVarianceAsym\_inner\_eq\_massEigenvalue\_average}
\end{theorem}
\begin{theorem}[\texttt{covariance\_spectralLatticeCovarianceAsym\_translation\_invariant}]\label{Pphi2-covariance-spectralLatticeCovarianceAsym-translation-invariant}
  \lean{Pphi2.covariance_spectralLatticeCovarianceAsym_translation_invariant}\leanok
  \uses{GaussianField-spectralLatticeCovarianceAsym, Pphi2-asymShift-sum-mul-eq, Pphi2-spectralCovAsym-massOperator-eq, GaussianField-AsymLatticeSites, GaussianField-covariance, Pphi2-asymShift, GaussianField-AsymLatticeField, GaussianField-instFintypeAsymLatticeSitesOfNeZeroNat, GaussianField-massOperatorAsym-surjective, GaussianField-massOperatorAsym, GaussianField-ell2-, Pphi2-massOperatorAsym-translation-commute}
  \texttt{Pphi2.covariance\_spectralLatticeCovarianceAsym\_translation\_invariant}
\end{theorem}
\begin{theorem}[\texttt{asymShift\_delta}]\label{Pphi2-asymShift-delta}
  \lean{Pphi2.asymShift_delta}\leanok
  \uses{GaussianField-AsymLatticeSites, Pphi2-asymShift, GaussianField-AsymLatticeField, GaussianField-asymLatticeDelta}
  \texttt{Pphi2.asymShift\_delta}
\end{theorem}
\begin{theorem}[\texttt{latticeCovarianceAsymGJ\_inner\_eq\_inv\_a\_sq\_spectral}]\label{Pphi2-latticeCovarianceAsymGJ-inner-eq-inv-a-sq-spectral}
  \lean{Pphi2.latticeCovarianceAsymGJ_inner_eq_inv_a_sq_spectral}\leanok
  \uses{GaussianField-spectralLatticeCovarianceAsym, GaussianField-AsymLatticeSites, GaussianField-AsymLatticeField, GaussianField-latticeCovarianceAsymGJ, GaussianField-ell2-}
  \texttt{Pphi2.latticeCovarianceAsymGJ\_inner\_eq\_inv\_a\_sq\_spectral}
\end{theorem}
\begin{definition}[\texttt{asymShift}]\label{Pphi2-asymShift}
  \lean{Pphi2.asymShift}\leanok
  \uses{GaussianField-AsymLatticeSites, GaussianField-AsymLatticeField}
  \texttt{Pphi2.asymShift}
\end{definition}
\begin{theorem}[\texttt{massOperatorAsym\_translation\_commute}]\label{Pphi2-massOperatorAsym-translation-commute}
  \lean{Pphi2.massOperatorAsym_translation_commute}\leanok
  \uses{Pphi2-finiteLaplacianAsym-translation-commute, GaussianField-AsymLatticeSites, GaussianField-finiteLaplacianAsym, Pphi2-asymShift, GaussianField-AsymLatticeField, GaussianField-massOperatorAsym}
  \texttt{Pphi2.massOperatorAsym\_translation\_commute}
\end{theorem}
\begin{theorem}[\texttt{delta\_massEigenvectorCoeff\_asym}]\label{Pphi2-delta-massEigenvectorCoeff-asym}
  \lean{Pphi2.delta_massEigenvectorCoeff_asym}\leanok
  \uses{GaussianField-AsymLatticeSites, GaussianField-massEigenvectorBasisAsym, GaussianField-instFintypeAsymLatticeSitesOfNeZeroNat, GaussianField-asymLatticeDelta}
  \texttt{Pphi2.delta\_massEigenvectorCoeff\_asym}
\end{theorem}
\section{\texttt{Pphi2.TorusContinuumLimit.TorusOSAxioms}}
\begin{theorem}[\texttt{os1}]\label{Pphi2-SatisfiesTorusOS-os1}
  \lean{Pphi2.SatisfiesTorusOS.os1}\leanok
  \uses{GaussianField-NuclearTensorProduct-instModuleReal, GaussianField-NuclearTensorProduct-instTopologicalSpace, GaussianField-NuclearTensorProduct-instIsTopologicalAddGroup, Pphi2-TorusOS1-Regularity, GaussianField-NuclearTensorProduct-instAddCommGroup, Pphi2-SatisfiesTorusOS, GaussianField-Configuration, GaussianField-NuclearTensorProduct-instContinuousSMulReal, GaussianField-SmoothMap-Circle, GaussianField-instMeasurableSpaceConfiguration, GaussianField-TorusTestFunction}
  \texttt{Pphi2.SatisfiesTorusOS.os1}
\end{theorem}
\begin{theorem}[\texttt{os0}]\label{Pphi2-SatisfiesTorusOS-os0}
  \lean{Pphi2.SatisfiesTorusOS.os0}\leanok
  \uses{GaussianField-NuclearTensorProduct-instModuleReal, GaussianField-NuclearTensorProduct-instTopologicalSpace, GaussianField-NuclearTensorProduct-instIsTopologicalAddGroup, Pphi2-TorusOS0-Analyticity, GaussianField-NuclearTensorProduct-instAddCommGroup, Pphi2-SatisfiesTorusOS, GaussianField-Configuration, GaussianField-NuclearTensorProduct-instContinuousSMulReal, GaussianField-SmoothMap-Circle, GaussianField-instMeasurableSpaceConfiguration, GaussianField-TorusTestFunction}
  \texttt{Pphi2.SatisfiesTorusOS.os0}
\end{theorem}
\begin{theorem}[\texttt{torusContinuumGreen\_pointGroup\_invariant}]\label{Pphi2-torusContinuumGreen-pointGroup-invariant}
  \lean{Pphi2.torusContinuumGreen_pointGroup_invariant}\leanok
  \uses{GaussianField-greenFunctionBilinear-timeReflection-invariant, GaussianField-torusTimeReflection, Pphi2-torusContinuumGreen, GaussianField-NuclearTensorProduct-instModuleReal, GaussianField-NuclearTensorProduct-instTopologicalSpace, GaussianField-torusSwap, GaussianField-NuclearTensorProduct-instAddCommGroup, GaussianField-SmoothMap-Circle, GaussianField-greenFunctionBilinear-swap-invariant, GaussianField-TorusTestFunction}
  \texttt{Pphi2.torusContinuumGreen\_pointGroup\_invariant}
\end{theorem}
\begin{theorem}[\texttt{torusGaussianLimit\_os2\_translation}]\label{Pphi2-torusGaussianLimit-os2-translation}
  \lean{Pphi2.torusGaussianLimit_os2_translation}\leanok
  \uses{Pphi2-torusContinuumGreen, Pphi2-IsTorusGaussianContinuumLimit, GaussianField-NuclearTensorProduct-instModuleReal, GaussianField-NuclearTensorProduct-instTopologicalSpace, Pphi2-torusContinuumGreen-translation-invariant, GaussianField-NuclearTensorProduct-instIsTopologicalAddGroup, Pphi2-torusGaussianLimit-characteristic-functional, Pphi2-torusGeneratingFunctional, GaussianField-NuclearTensorProduct-instAddCommGroup, GaussianField-torusTranslation, GaussianField-Configuration, GaussianField-NuclearTensorProduct-instContinuousSMulReal, GaussianField-SmoothMap-Circle, GaussianField-instMeasurableSpaceConfiguration, GaussianField-TorusTestFunction, Pphi2-TorusOS2-TranslationInvariance}
  \texttt{Pphi2.torusGaussianLimit\_os2\_translation}
\end{theorem}
\begin{definition}[\texttt{TorusOS0\_Analyticity}]\label{Pphi2-TorusOS0-Analyticity}
  \lean{Pphi2.TorusOS0_Analyticity}\leanok
  \uses{GaussianField-NuclearTensorProduct-instModuleReal, GaussianField-NuclearTensorProduct-instTopologicalSpace, GaussianField-NuclearTensorProduct-instIsTopologicalAddGroup, GaussianField-NuclearTensorProduct-instAddCommGroup, GaussianField-Configuration, GaussianField-NuclearTensorProduct-instContinuousSMulReal, GaussianField-SmoothMap-Circle, GaussianField-instMeasurableSpaceConfiguration, Pphi2-torusGeneratingFunctional-, GaussianField-TorusTestFunction}
  Prop definitions for each OS axiom on the torus.
\end{definition}
\begin{theorem}[\texttt{torusGaussianLimit\_satisfies\_OS}]\label{Pphi2-torusGaussianLimit-satisfies-OS}
  \lean{Pphi2.torusGaussianLimit_satisfies_OS}\leanok
  \uses{Pphi2-IsTorusGaussianContinuumLimit, GaussianField-NuclearTensorProduct-instModuleReal, GaussianField-NuclearTensorProduct-instTopologicalSpace, GaussianField-NuclearTensorProduct-instIsTopologicalAddGroup, Pphi2-torusGaussianLimit-os2-D4, GaussianField-NuclearTensorProduct-instAddCommGroup, Pphi2-SatisfiesTorusOS, GaussianField-Configuration, GaussianField-NuclearTensorProduct-instContinuousSMulReal, GaussianField-SmoothMap-Circle, Pphi2-torusGaussianLimit-os0, Pphi2-torusGaussianLimit-os1, GaussianField-instMeasurableSpaceConfiguration, GaussianField-TorusTestFunction, Pphi2-torusGaussianLimit-os2-translation}
  The torus Gaussian continuum limit satisfies all torus OS axioms.
\end{theorem}
\begin{theorem}[\texttt{torusGeneratingFunctionalℂ\_analyticOnNhd}]\label{Pphi2-torusGeneratingFunctional--analyticOnNhd}
  \lean{Pphi2.torusGeneratingFunctionalℂ_analyticOnNhd}\leanok
  \uses{Pphi2-IsTorusGaussianContinuumLimit, GaussianField-NuclearTensorProduct-instModuleReal, GaussianField-NuclearTensorProduct-instTopologicalSpace, GaussianField-NuclearTensorProduct-instIsTopologicalAddGroup, GaussianField-NuclearTensorProduct-instAddCommGroup, GaussianField-configuration-eval-measurable, GaussianField-Configuration, GaussianField-NuclearTensorProduct-instContinuousSMulReal, GaussianField-SmoothMap-Circle, GaussianField-instMeasurableSpaceConfiguration, Pphi2-torusGeneratingFunctional-, GaussianField-TorusTestFunction, Lattice-toSemilatticeInf}
  $z \mapsto Z_{\mathbb{C}}[\sum \text{Re}(z_i) J_i, \sum \text{Im}(z_i) J_i]$ is entire analytic. Proved via \texttt{analyticOnNhd\_integral} with Gaussian exponential moment domination.
\end{theorem}
\begin{definition}[\texttt{TorusOS2\_D4Invariance}]\label{Pphi2-TorusOS2-D4Invariance}
  \lean{Pphi2.TorusOS2_D4Invariance}\leanok
  \uses{GaussianField-torusTimeReflection, GaussianField-NuclearTensorProduct-instModuleReal, GaussianField-NuclearTensorProduct-instTopologicalSpace, GaussianField-NuclearTensorProduct-instIsTopologicalAddGroup, Pphi2-torusGeneratingFunctional, GaussianField-torusSwap, GaussianField-NuclearTensorProduct-instAddCommGroup, GaussianField-Configuration, GaussianField-NuclearTensorProduct-instContinuousSMulReal, GaussianField-SmoothMap-Circle, GaussianField-instMeasurableSpaceConfiguration, GaussianField-TorusTestFunction}
  \texttt{Pphi2.TorusOS2\_D4Invariance}
\end{definition}
\begin{theorem}[\texttt{torusGaussianLimit\_complex\_cf\_quadratic}]\label{Pphi2-torusGaussianLimit-complex-cf-quadratic}
  \lean{Pphi2.torusGaussianLimit_complex_cf_quadratic}\leanok
  \uses{Pphi2-torusContinuumGreen, Pphi2-IsTorusGaussianContinuumLimit, GaussianField-NuclearTensorProduct-instModuleReal, GaussianField-NuclearTensorProduct-instTopologicalSpace, GaussianField-circleHasLaplacianEigenvalues, GaussianField-NuclearTensorProduct-dyninMityaginSpace, GaussianField-SmoothMap-Circle-instAddCommGroup, GaussianField-NuclearTensorProduct-instIsTopologicalAddGroup, Pphi2-torusGaussianLimit-characteristic-functional, Pphi2-torusGeneratingFunctional--congr-simp, GaussianField-tensorProductHasLaplacianEigenvalues, Pphi2-torusGeneratingFunctional, GaussianField-NuclearTensorProduct-instAddCommGroup, GaussianField-greenFunctionBilinear-finset-sum, GaussianField-SmoothMap-Circle-instModule, GaussianField-SmoothMap-Circle-instContinuousSMul, GaussianField-Configuration, GaussianField-smoothCircle-dyninMityaginSpace, GaussianField-SmoothMap-Circle-instIsTopologicalAddGroup, GaussianField-NuclearTensorProduct-instContinuousSMulReal, GaussianField-SmoothMap-Circle, GaussianField-instMeasurableSpaceConfiguration, Pphi2-torusGeneratingFunctional-, Pphi2-torusGeneratingFunctional--analyticOnNhd, GaussianField-TorusTestFunction, GaussianField-SmoothMap-Circle-instTopologicalSpace, GaussianField-greenFunctionBilinear}
  $Z_{\mathbb{C}}[z] = \exp(-\frac{1}{2}\sum_{i,j} z_i z_j G_L(J_i, J_j))$. Proved by the identity theorem: both sides are entire and agree on $\mathbb{R}^n$.
\end{theorem}
\begin{theorem}[\texttt{congr\_simp}]\label{Pphi2-torusGeneratingFunctional--congr-simp}
  \lean{Pphi2.torusGeneratingFunctionalℂ.congr_simp}\leanok
  \uses{GaussianField-NuclearTensorProduct-instModuleReal, GaussianField-NuclearTensorProduct-instTopologicalSpace, GaussianField-NuclearTensorProduct-instIsTopologicalAddGroup, GaussianField-NuclearTensorProduct-instAddCommGroup, GaussianField-Configuration, GaussianField-NuclearTensorProduct-instContinuousSMulReal, GaussianField-SmoothMap-Circle, GaussianField-instMeasurableSpaceConfiguration, Pphi2-torusGeneratingFunctional-, GaussianField-TorusTestFunction}
  \texttt{Pphi2.torusGeneratingFunctionalℂ.congr\_simp}
\end{theorem}
\begin{definition}[\texttt{torusGeneratingFunctional}]\label{Pphi2-torusGeneratingFunctional}
  \lean{Pphi2.torusGeneratingFunctional}\leanok
  \uses{GaussianField-NuclearTensorProduct-instModuleReal, GaussianField-NuclearTensorProduct-instTopologicalSpace, GaussianField-NuclearTensorProduct-instIsTopologicalAddGroup, GaussianField-NuclearTensorProduct-instAddCommGroup, GaussianField-Configuration, GaussianField-NuclearTensorProduct-instContinuousSMulReal, GaussianField-SmoothMap-Circle, GaussianField-instMeasurableSpaceConfiguration, GaussianField-TorusTestFunction}
  ``\texttt{lean def torusGeneratingFunctional (μ : ...) [IsProbabilityMeasure μ] (f : TorusTestFunction L) : ℂ }`` $Z_\mu(f) = \int e^{i\omega(f)}\, d\mu(\omega)$.
\end{definition}
\begin{theorem}[\texttt{torusGaussianLimit\_os1}]\label{Pphi2-torusGaussianLimit-os1}
  \lean{Pphi2.torusGaussianLimit_os1}\leanok
  \uses{Pphi2-torusContinuumGreen, Pphi2-IsTorusGaussianContinuumLimit, GaussianField-NuclearTensorProduct-instModuleReal, GaussianField-NuclearTensorProduct-instTopologicalSpace, Pphi2-torusContinuumGreen-continuous-diag, GaussianField-NuclearTensorProduct-instIsTopologicalAddGroup, Pphi2-torusContinuumGreen-nonneg, Pphi2-TorusOS1-Regularity, Pphi2-IsTorusGaussianContinuumLimit-isGaussian, GaussianField-NuclearTensorProduct-instAddCommGroup, GaussianField-Configuration, GaussianField-NuclearTensorProduct-instContinuousSMulReal, GaussianField-SmoothMap-Circle, Pphi2-IsTorusGaussianContinuumLimit-covariance-eq, GaussianField-instMeasurableSpaceConfiguration, Pphi2-torusGeneratingFunctional-, GaussianField-TorusTestFunction, Lattice-toSemilatticeInf}
  \texttt{Pphi2.torusGaussianLimit\_os1}
\end{theorem}
\begin{definition}[\texttt{TorusOS1\_Regularity}]\label{Pphi2-TorusOS1-Regularity}
  \lean{Pphi2.TorusOS1_Regularity}\leanok
  \uses{GaussianField-NuclearTensorProduct-instModuleReal, GaussianField-NuclearTensorProduct-instTopologicalSpace, GaussianField-NuclearTensorProduct-instIsTopologicalAddGroup, GaussianField-NuclearTensorProduct-instAddCommGroup, GaussianField-Configuration, GaussianField-NuclearTensorProduct-instContinuousSMulReal, GaussianField-SmoothMap-Circle, GaussianField-instMeasurableSpaceConfiguration, Pphi2-torusGeneratingFunctional-, GaussianField-TorusTestFunction}
  \texttt{Pphi2.TorusOS1\_Regularity}
\end{definition}
\begin{theorem}[\texttt{torusGaussianLimit\_characteristic\_functional}]\label{Pphi2-torusGaussianLimit-characteristic-functional}
  \lean{Pphi2.torusGaussianLimit_characteristic_functional}\leanok
  \uses{Pphi2-torusContinuumGreen, Pphi2-IsTorusGaussianContinuumLimit, GaussianField-NuclearTensorProduct-instModuleReal, GaussianField-NuclearTensorProduct-instTopologicalSpace, GaussianField-NuclearTensorProduct-instIsTopologicalAddGroup, Pphi2-torusGeneratingFunctional, Pphi2-torusContinuumGreen-nonneg, Pphi2-IsTorusGaussianContinuumLimit-isGaussian, GaussianField-NuclearTensorProduct-instAddCommGroup, GaussianField-Configuration, GaussianField-NuclearTensorProduct-instContinuousSMulReal, GaussianField-SmoothMap-Circle, Pphi2-IsTorusGaussianContinuumLimit-covariance-eq, GaussianField-instMeasurableSpaceConfiguration, GaussianField-TorusTestFunction}
  $E[e^{i\omega(f)}] = \exp(-\frac{1}{2}G_L(f,f))$. Proved by matching the real MGF to $N(0, G_L(f,f))$ via \texttt{mgf\_id\_gaussianReal}, then analytically continuing to $z = i$ using \texttt{eqOn\_complexMGF\_of\_mgf}.
\end{theorem}
\begin{theorem}[\texttt{torusGaussianLimit\_os2\_D4}]\label{Pphi2-torusGaussianLimit-os2-D4}
  \lean{Pphi2.torusGaussianLimit_os2_D4}\leanok
  \uses{GaussianField-torusTimeReflection, Pphi2-torusContinuumGreen, Pphi2-IsTorusGaussianContinuumLimit, GaussianField-NuclearTensorProduct-instModuleReal, GaussianField-NuclearTensorProduct-instTopologicalSpace, GaussianField-NuclearTensorProduct-instIsTopologicalAddGroup, Pphi2-torusGaussianLimit-characteristic-functional, Pphi2-torusGeneratingFunctional, GaussianField-torusSwap, GaussianField-NuclearTensorProduct-instAddCommGroup, GaussianField-Configuration, GaussianField-NuclearTensorProduct-instContinuousSMulReal, Pphi2-TorusOS2-D4Invariance, GaussianField-SmoothMap-Circle, GaussianField-instMeasurableSpaceConfiguration, Pphi2-torusContinuumGreen-pointGroup-invariant, GaussianField-TorusTestFunction}
  \texttt{Pphi2.torusGaussianLimit\_os2\_D4}
\end{theorem}
\begin{theorem}[\texttt{torusGaussianLimit\_os0}]\label{Pphi2-torusGaussianLimit-os0}
  \lean{Pphi2.torusGaussianLimit_os0}\leanok
  \uses{Pphi2-torusContinuumGreen, Pphi2-IsTorusGaussianContinuumLimit, GaussianField-NuclearTensorProduct-instModuleReal, GaussianField-NuclearTensorProduct-instTopologicalSpace, GaussianField-NuclearTensorProduct-instIsTopologicalAddGroup, Pphi2-TorusOS0-Analyticity, GaussianField-NuclearTensorProduct-instAddCommGroup, GaussianField-Configuration, GaussianField-NuclearTensorProduct-instContinuousSMulReal, GaussianField-SmoothMap-Circle, GaussianField-instMeasurableSpaceConfiguration, Pphi2-torusGeneratingFunctional-, Pphi2-torusGaussianLimit-complex-cf-quadratic, GaussianField-TorusTestFunction}
  Each OS axiom for the Gaussian limit.
\end{theorem}
\begin{definition}[\texttt{TorusOS2\_TranslationInvariance}]\label{Pphi2-TorusOS2-TranslationInvariance}
  \lean{Pphi2.TorusOS2_TranslationInvariance}\leanok
  \uses{GaussianField-NuclearTensorProduct-instModuleReal, GaussianField-NuclearTensorProduct-instTopologicalSpace, GaussianField-NuclearTensorProduct-instIsTopologicalAddGroup, Pphi2-torusGeneratingFunctional, GaussianField-NuclearTensorProduct-instAddCommGroup, GaussianField-torusTranslation, GaussianField-Configuration, GaussianField-NuclearTensorProduct-instContinuousSMulReal, GaussianField-SmoothMap-Circle, GaussianField-instMeasurableSpaceConfiguration, GaussianField-TorusTestFunction}
  \texttt{Pphi2.TorusOS2\_TranslationInvariance}
\end{definition}
\begin{theorem}[\texttt{torusGaussianOS\_exists}]\label{Pphi2-torusGaussianOS-exists}
  \lean{Pphi2.torusGaussianOS_exists}\leanok
  \uses{Pphi2-IsTorusGaussianContinuumLimit, GaussianField-NuclearTensorProduct-instModuleReal, GaussianField-NuclearTensorProduct-instTopologicalSpace, Pphi2-torusGaussianLimit-satisfies-OS, Pphi2-torusGaussianLimit-converges, GaussianField-NuclearTensorProduct-instIsTopologicalAddGroup, Pphi2-torusContinuumMeasure, GaussianField-NuclearTensorProduct-instAddCommGroup, Pphi2-SatisfiesTorusOS, GaussianField-Configuration, GaussianField-NuclearTensorProduct-instContinuousSMulReal, GaussianField-SmoothMap-Circle, GaussianField-instMeasurableSpaceConfiguration, GaussianField-TorusTestFunction}
  Existence: for any $L > 0$, $m > 0$, there exists a probability measure satisfying \texttt{SatisfiesTorusOS}.
\end{theorem}
\begin{definition}[\texttt{torusGeneratingFunctionalℂ}]\label{Pphi2-torusGeneratingFunctional-}
  \lean{Pphi2.torusGeneratingFunctionalℂ}\leanok
  \uses{GaussianField-NuclearTensorProduct-instModuleReal, GaussianField-NuclearTensorProduct-instTopologicalSpace, GaussianField-NuclearTensorProduct-instIsTopologicalAddGroup, GaussianField-NuclearTensorProduct-instAddCommGroup, GaussianField-Configuration, GaussianField-NuclearTensorProduct-instContinuousSMulReal, GaussianField-SmoothMap-Circle, GaussianField-instMeasurableSpaceConfiguration, GaussianField-TorusTestFunction}
  ``\texttt{lean def torusGeneratingFunctionalℂ (μ : ...) [IsProbabilityMeasure μ] (f\_\{\textbackslash text\{re\}\}\textbackslash  f\_\{\textbackslash text\{im\}\} : TorusTestFunction L) : ℂ }`` Complex generating functional $Z[J] = \int e^{i\omega(f_{\text{re}}) - \omega(f_{\text{im}})}\, d\mu$.
\end{definition}
\begin{theorem}[\texttt{torusContinuumGreen\_translation\_invariant}]\label{Pphi2-torusContinuumGreen-translation-invariant}
  \lean{Pphi2.torusContinuumGreen_translation_invariant}\leanok
  \uses{Pphi2-torusContinuumGreen, GaussianField-NuclearTensorProduct-instModuleReal, GaussianField-NuclearTensorProduct-instTopologicalSpace, GaussianField-greenFunctionBilinear-translation-invariant, GaussianField-NuclearTensorProduct-instAddCommGroup, GaussianField-torusTranslation, GaussianField-SmoothMap-Circle, GaussianField-TorusTestFunction}
  \texttt{Pphi2.torusContinuumGreen\_translation\_invariant}
\end{theorem}
\begin{theorem}[\texttt{os2\_translation}]\label{Pphi2-SatisfiesTorusOS-os2-translation}
  \lean{Pphi2.SatisfiesTorusOS.os2_translation}\leanok
  \uses{GaussianField-NuclearTensorProduct-instModuleReal, GaussianField-NuclearTensorProduct-instTopologicalSpace, GaussianField-NuclearTensorProduct-instIsTopologicalAddGroup, GaussianField-NuclearTensorProduct-instAddCommGroup, Pphi2-SatisfiesTorusOS, GaussianField-Configuration, GaussianField-NuclearTensorProduct-instContinuousSMulReal, GaussianField-SmoothMap-Circle, GaussianField-instMeasurableSpaceConfiguration, GaussianField-TorusTestFunction, Pphi2-TorusOS2-TranslationInvariance}
  \texttt{Pphi2.SatisfiesTorusOS.os2\_translation}
\end{theorem}
\begin{definition}[\texttt{SatisfiesTorusOS}]\label{Pphi2-SatisfiesTorusOS}
  \lean{Pphi2.SatisfiesTorusOS}\leanok
  \uses{GaussianField-NuclearTensorProduct-instModuleReal, GaussianField-NuclearTensorProduct-instTopologicalSpace, GaussianField-NuclearTensorProduct-instIsTopologicalAddGroup, GaussianField-NuclearTensorProduct-instAddCommGroup, GaussianField-Configuration, GaussianField-NuclearTensorProduct-instContinuousSMulReal, GaussianField-SmoothMap-Circle, GaussianField-instMeasurableSpaceConfiguration, GaussianField-TorusTestFunction}
  Bundled torus OS axioms: OS0 + OS1 + OS2 (translation + D4).
\end{definition}
\begin{theorem}[\texttt{os2\_D4}]\label{Pphi2-SatisfiesTorusOS-os2-D4}
  \lean{Pphi2.SatisfiesTorusOS.os2_D4}\leanok
  \uses{GaussianField-NuclearTensorProduct-instModuleReal, GaussianField-NuclearTensorProduct-instTopologicalSpace, GaussianField-NuclearTensorProduct-instIsTopologicalAddGroup, GaussianField-NuclearTensorProduct-instAddCommGroup, Pphi2-SatisfiesTorusOS, GaussianField-Configuration, GaussianField-NuclearTensorProduct-instContinuousSMulReal, Pphi2-TorusOS2-D4Invariance, GaussianField-SmoothMap-Circle, GaussianField-instMeasurableSpaceConfiguration, GaussianField-TorusTestFunction}
  \texttt{Pphi2.SatisfiesTorusOS.os2\_D4}
\end{theorem}
\begin{theorem}[\texttt{torusContinuumGreen\_continuous\_diag}]\label{Pphi2-torusContinuumGreen-continuous-diag}
  \lean{Pphi2.torusContinuumGreen_continuous_diag}\leanok
  \uses{Pphi2-torusContinuumGreen, GaussianField-NuclearTensorProduct-instModuleReal, GaussianField-NuclearTensorProduct-instTopologicalSpace, GaussianField-circleHasLaplacianEigenvalues, GaussianField-NuclearTensorProduct-dyninMityaginSpace, GaussianField-SmoothMap-Circle-instAddCommGroup, GaussianField-tensorProductHasLaplacianEigenvalues, GaussianField-NuclearTensorProduct-instAddCommGroup, GaussianField-SmoothMap-Circle-instModule, GaussianField-SmoothMap-Circle-instContinuousSMul, GaussianField-smoothCircle-dyninMityaginSpace, GaussianField-SmoothMap-Circle-instIsTopologicalAddGroup, GaussianField-SmoothMap-Circle, GaussianField-greenFunctionBilinear-continuous-diag, GaussianField-TorusTestFunction, GaussianField-SmoothMap-Circle-instTopologicalSpace}
  \texttt{Pphi2.torusContinuumGreen\_continuous\_diag}
\end{theorem}
\section{\texttt{Pphi2.InteractingMeasure.LatticeAction}}
\begin{theorem}[\texttt{latticeInteraction\_bounded\_below}]\label{Pphi2-latticeInteraction-bounded-below}
  \lean{Pphi2.latticeInteraction_bounded_below}\leanok
  \uses{GaussianField-FinLatticeSites, Pphi2-wickConstant, GaussianField-FinLatticeField, Pphi2-wickPolynomial-bounded-below, Pphi2-wickPolynomial, Pphi2-latticeInteraction}
  $\exists B,\;\forall \varphi,\; V_a(\varphi) \ge -B$. Since $:P(\varphi(x)):_{c_a} \ge -A$ for each site, the sum of $|\Lambda|$ terms gives $V_a \ge -a^d |\Lambda| A$. Depends on \texttt{wickPolynomial\_bounded\_below}.
\end{theorem}
\begin{definition}[\texttt{latticeInteraction}]\label{Pphi2-latticeInteraction}
  \lean{Pphi2.latticeInteraction}\leanok
  \uses{GaussianField-FinLatticeSites, Pphi2-wickConstant, GaussianField-FinLatticeField, Pphi2-wickPolynomial}
  \texttt{Pphi2.latticeInteraction}
\end{definition}
\begin{theorem}[\texttt{latticeInteraction\_continuous}]\label{Pphi2-latticeInteraction-continuous}
  \lean{Pphi2.latticeInteraction_continuous}\leanok
  \uses{GaussianField-FinLatticeSites, Pphi2-wickMonomial, Pphi2-wickMonomial-continuous, Pphi2-wickConstant, GaussianField-FinLatticeField, Pphi2-wickPolynomial, Pphi2-latticeInteraction}
  $V_a$ is continuous in $\varphi$, since each $:P(\varphi(x)):$ is continuous in $\varphi(x)$ and the sum is finite.
\end{theorem}
\begin{theorem}[\texttt{latticeInteraction\_single\_site}]\label{Pphi2-latticeInteraction-single-site}
  \lean{Pphi2.latticeInteraction_single_site}\leanok
  \uses{GaussianField-FinLatticeSites, Pphi2-wickConstant, GaussianField-FinLatticeField, Pphi2-wickPolynomial, Pphi2-latticeInteraction}
  $V_a(\varphi) = a^d \sum_x v_x(\varphi(x))$ where $v_x(\tau) = :P(\tau):_{c_a}$ depends on $\varphi$ only through the single-site value $\varphi(x)$. This single-site decomposition enables the FKG inequality.
\end{theorem}
\begin{theorem}[\texttt{wickMonomial\_continuous}]\label{Pphi2-wickMonomial-continuous}
  \lean{Pphi2.wickMonomial_continuous}\leanok
  \uses{Pphi2-wickMonomial}
  \texttt{Pphi2.wickMonomial\_continuous}
\end{theorem}
\section{\texttt{Pphi2.TransferMatrix.Jentzsch}}
\begin{theorem}[\texttt{transferOperator\_strictly\_positive\_definite}]\label{Pphi2-transferOperator-strictly-positive-definite}
  \lean{Pphi2.transferOperator_strictly_positive_definite}\leanok
  \uses{Pphi2-L2SpatialField, Pphi2-SpatialField, Pphi2-transferWeight, Pphi2-transferGaussian-memLp, Pphi2-transferWeight-measurable, Pphi2-transferOperatorCLM, Pphi2-convolution-gaussian-strictly-positive-definite, Pphi2-transferWeight-pos, Pphi2-transferGaussian, Pphi2-transferWeight-bound}
  $\langle f, Tf\rangle > 0$ for $f \ne 0$. From $\langle f, Tf\rangle = \langle M_w f, \mathrm{Conv}_G(M_w f)\rangle$ (self-adjointness of $M_w$), injectivity of $M_w$ ($w > 0$), and Gaussian strict PD.
\end{theorem}
\begin{theorem}[\texttt{transferOperator\_ground\_simple}]\label{Pphi2-transferOperator-ground-simple}
  \lean{Pphi2.transferOperator_ground_simple}\leanok
  \uses{Pphi2-transferOperator-positivityImproving, Pphi2-L2SpatialField, Pphi2-transferOperator-isCompact, Pphi2-SpatialField, Pphi2-transferOperator-eigenvalues-pos, Pphi2-transferOperator-isSelfAdjoint, Pphi2-l2SpatialField-hilbertBasis-nontrivial, Pphi2-transferOperatorCLM}
  $\exists i_0, i_1$ with $i_1 \ne i_0$ and $\lambda(i_1) < \lambda(i_0)$. Combines Jentzsch + nontriviality + eigenvalue positivity.
\end{theorem}
\begin{theorem}[\texttt{transferOperator\_positivityImproving}]\label{Pphi2-transferOperator-positivityImproving}
  \lean{Pphi2.transferOperator_positivityImproving}\leanok
  \uses{Pphi2-L2SpatialField, Pphi2-SpatialField, Pphi2-transferWeight, Pphi2-transferGaussian-pos, Pphi2-transferGaussian-memLp, Pphi2-transferGaussian-memLp-two, Pphi2-transferWeight-measurable, Pphi2-transferOperatorCLM, Pphi2-transferWeight-pos, Pphi2-transferGaussian, Pphi2-transferWeight-bound}
  $T$ is positivity-improving: if $f \ge 0$ a.e. and $f \not\equiv 0$, then $(Tf)(x) > 0$ a.e. Proved via the factorization $T = M_w \circ \mathrm{Conv}_G \circ M_w$ with $w > 0$, $G > 0$, and translation invariance of Lebesgue measure.
\end{theorem}
\begin{theorem}[\texttt{l2SpatialField\_hilbertBasis\_nontrivial}]\label{Pphi2-l2SpatialField-hilbertBasis-nontrivial}
  \lean{Pphi2.l2SpatialField_hilbertBasis_nontrivial}\leanok
  \uses{Pphi2-L2SpatialField, Pphi2-SpatialField, Lattice-toSemilatticeInf}
  Any Hilbert basis of $L^2(\mathbb{R}^{N_s})$ has $\ge 2$ elements. Proved by constructing orthogonal indicator functions on disjoint balls.
\end{theorem}
\begin{theorem}[\texttt{transferOperator\_inner\_nonneg}]\label{Pphi2-transferOperator-inner-nonneg}
  \lean{Pphi2.transferOperator_inner_nonneg}\leanok
  \uses{Pphi2-L2SpatialField, Pphi2-SpatialField, Pphi2-transferOperatorCLM, Pphi2-transferOperator-strictly-positive-definite}
  $\langle f, Tf\rangle \ge 0$ for all $f$. Immediate from strict PD ($> 0$ for $f \ne 0$, $= 0$ for $f = 0$).
\end{theorem}
\begin{theorem}[\texttt{convolution\_gaussian\_strictly\_positive\_definite}]\label{Pphi2-convolution-gaussian-strictly-positive-definite}
  \lean{Pphi2.convolution_gaussian_strictly_positive_definite}\leanok
  \uses{Pphi2-L2SpatialField, Pphi2-SpatialField, Pphi2-transferGaussian-memLp, Pphi2-gaussian-conv-strictlyPD, Pphi2-transferGaussian}
  $\langle g, \mathrm{Conv}_G g\rangle > 0$ for $g \ne 0$. Delegates to \texttt{gaussian\_conv\_strictlyPD}.
\end{theorem}
\begin{theorem}[\texttt{congr\_simp}]\label{Pphi2-transferOperatorCLM-congr-simp}
  \lean{Pphi2.transferOperatorCLM.congr_simp}\leanok
  \uses{Pphi2-L2SpatialField, Pphi2-SpatialField, Pphi2-transferOperatorCLM}
  \texttt{Pphi2.transferOperatorCLM.congr\_simp}
\end{theorem}
\begin{theorem}[\texttt{transferOperator\_ground\_simple\_spectral}]\label{Pphi2-transferOperator-ground-simple-spectral}
  \lean{Pphi2.transferOperator_ground_simple_spectral}\leanok
  \uses{Pphi2-L2SpatialField, Pphi2-SpatialField, Pphi2-transferOperator-ground-simple, Pphi2-transferOperatorCLM, Pphi2-transferOperator-spectral}
  Packaged spectral data with distinguished ground and first excited levels.
\end{theorem}
\begin{theorem}[\texttt{transferOperator\_eigenvalues\_pos}]\label{Pphi2-transferOperator-eigenvalues-pos}
  \lean{Pphi2.transferOperator_eigenvalues_pos}\leanok
  \uses{Pphi2-L2SpatialField, Pphi2-SpatialField, Pphi2-transferOperatorCLM, Pphi2-transferOperator-strictly-positive-definite}
  All eigenvalues $\lambda_i > 0$. From $\langle b_i, Tb_i\rangle = \lambda_i \|b_i\|^2 = \lambda_i > 0$.
\end{theorem}
\section{\texttt{Pphi2.AsymTorus.AsymTorusEmbedding}}
\begin{theorem}[\texttt{congr\_simp}]\label{Pphi2-asymGeomSpacing-congr-simp}
  \lean{Pphi2.asymGeomSpacing.congr_simp}\leanok
  \uses{Pphi2-asymGeomSpacing}
  \texttt{Pphi2.asymGeomSpacing.congr\_simp}
\end{theorem}
\begin{definition}[\texttt{evalAsymAtFinSiteGJ}]\label{Pphi2-evalAsymAtFinSiteGJ}
  \lean{Pphi2.evalAsymAtFinSiteGJ}\leanok
  \uses{GaussianField-NuclearTensorProduct-instModuleReal, GaussianField-NuclearTensorProduct-instTopologicalSpace, GaussianField-FinLatticeSites, GaussianField-NuclearTensorProduct-instAddCommGroup, Pphi2-evalAsymAtFinSite, GaussianField-SmoothMap-Circle, GaussianField-AsymTorusTestFunction, Pphi2-asymGeomSpacing}
  \texttt{Pphi2.evalAsymAtFinSiteGJ}
\end{definition}
\begin{theorem}[\texttt{asymLatticeTestFn\_expand}]\label{Pphi2-asymLatticeTestFn-expand}
  \lean{Pphi2.asymLatticeTestFn_expand}\leanok
  \uses{GaussianField-FinLatticeSites, Pphi2-asymLatticeTestFn, GaussianField-AsymTorusTestFunction, GaussianField-FinLatticeField}
  ``\texttt{lean theorem asymLatticeTestFn\_expand (N : ℕ) [NeZero N] (f : AsymTorusTestFunction Lt Ls) : asymLatticeTestFn Lt Ls N f = ∑ x : FinLatticeSites 2 N, (asymLatticeTestFn Lt Ls N f) x • Pi.single x (1 : ℝ) }`\texttt{  Informal statement: The lattice test function equals its expansion as a sum of delta functions weighted by site values.  Proof: Fully proved by }funext` and simplification.
\end{theorem}
\begin{theorem}[\texttt{asymTorusEmbedLift\_eval\_eq}]\label{Pphi2-asymTorusEmbedLift-eval-eq}
  \lean{Pphi2.asymTorusEmbedLift_eval_eq}\leanok
  \uses{GaussianField-NuclearTensorProduct-instModuleReal, GaussianField-NuclearTensorProduct-instTopologicalSpace, GaussianField-FinLatticeSites, GaussianField-NuclearTensorProduct-instIsTopologicalAddGroup, Pphi2-asymLatticeTestFn-expand, GaussianField-NuclearTensorProduct-instAddCommGroup, Pphi2-evalAsymAtFinSiteGJ, GaussianField-Configuration, GaussianField-NuclearTensorProduct-instContinuousSMulReal, GaussianField-SmoothMap-Circle, Pphi2-asymLatticeTestFn, Pphi2-asymTorusEmbedLift, GaussianField-AsymTorusTestFunction, GaussianField-FinLatticeField}
  ``\texttt{lean theorem asymTorusEmbedLift\_eval\_eq (N : ℕ) [NeZero N] (f : AsymTorusTestFunction Lt Ls) (ω : Configuration (FinLatticeField 2 N)) : (asymTorusEmbedLift Lt Ls N ω) f = ω (asymLatticeTestFn Lt Ls N f) }``  Informal statement: The embedding preserves evaluation: $(\iota\, \omega)(f) = \omega(g_f)$ where $g_f$ is the lattice test function.  Proof: Fully proved using \texttt{asymLatticeTestFn\_expand} and \texttt{map\_sum}.
\end{theorem}
\begin{theorem}[\texttt{congr\_simp}]\label{Pphi2-evalAsymAtFinSite-congr-simp}
  \lean{Pphi2.evalAsymAtFinSite.congr_simp}\leanok
  \uses{GaussianField-NuclearTensorProduct-instModuleReal, GaussianField-NuclearTensorProduct-instTopologicalSpace, GaussianField-FinLatticeSites, GaussianField-NuclearTensorProduct-instAddCommGroup, Pphi2-evalAsymAtFinSite, GaussianField-SmoothMap-Circle, GaussianField-AsymTorusTestFunction}
  \texttt{Pphi2.evalAsymAtFinSite.congr\_simp}
\end{theorem}
\begin{definition}[\texttt{evalAsymAtFinSite}]\label{Pphi2-evalAsymAtFinSite}
  \lean{Pphi2.evalAsymAtFinSite}\leanok
  \uses{GaussianField-NuclearTensorProduct-instModuleReal, GaussianField-NuclearTensorProduct-instTopologicalSpace, GaussianField-FinLatticeSites, GaussianField-evalAsymTorusAtSite, GaussianField-NuclearTensorProduct-instAddCommGroup, GaussianField-SmoothMap-Circle, GaussianField-AsymTorusTestFunction}
  ``\texttt{lean def evalAsymAtFinSite (N : ℕ) [NeZero N] (x : FinLatticeSites 2 N) : AsymTorusTestFunction Lt Ls →L[ℝ] ℝ }`\texttt{  Evaluates an asymmetric torus test function at a lattice site given in }FinLatticeSites 2 N\texttt{ form, by converting to }AsymLatticeSites\texttt{ and applying }evalAsymTorusAtSite`.
\end{definition}
\begin{theorem}[\texttt{evalAsymAtFinSiteGJ\_apply}]\label{Pphi2-evalAsymAtFinSiteGJ-apply}
  \lean{Pphi2.evalAsymAtFinSiteGJ_apply}\leanok
  \uses{GaussianField-NuclearTensorProduct-instModuleReal, GaussianField-NuclearTensorProduct-instTopologicalSpace, GaussianField-FinLatticeSites, GaussianField-NuclearTensorProduct-instAddCommGroup, Pphi2-evalAsymAtFinSite, Pphi2-evalAsymAtFinSiteGJ, GaussianField-SmoothMap-Circle, GaussianField-AsymTorusTestFunction, Pphi2-asymGeomSpacing}
  \texttt{Pphi2.evalAsymAtFinSiteGJ\_apply}
\end{theorem}
\begin{definition}[\texttt{asymLatticeTestFn}]\label{Pphi2-asymLatticeTestFn}
  \lean{Pphi2.asymLatticeTestFn}\leanok
  \uses{GaussianField-NuclearTensorProduct-instModuleReal, GaussianField-NuclearTensorProduct-instTopologicalSpace, GaussianField-FinLatticeSites, GaussianField-NuclearTensorProduct-instAddCommGroup, Pphi2-evalAsymAtFinSiteGJ, GaussianField-SmoothMap-Circle, GaussianField-AsymTorusTestFunction, GaussianField-FinLatticeField}
  ``\texttt{lean def asymLatticeTestFn (N : ℕ) [NeZero N] (f : AsymTorusTestFunction Lt Ls) : FinLatticeField 2 N }``  The asymmetric lattice test function: evaluates a torus test function $f$ at each lattice site.
\end{definition}
\begin{definition}[\texttt{asymSpaceSpacing}]\label{Pphi2-asymSpaceSpacing}
  \lean{Pphi2.asymSpaceSpacing}\leanok
  ``\texttt{lean def asymSpaceSpacing (Ns : ℕ) [NeZero Ns] : ℝ := Ls / Ns }``  Space lattice spacing $a_s = L_s / N_s$.
\end{definition}
\begin{definition}[\texttt{asymTorusContinuumGreen}]\label{Pphi2-asymTorusContinuumGreen}
  \lean{Pphi2.asymTorusContinuumGreen}\leanok
  \uses{GaussianField-NuclearTensorProduct-instModuleReal, GaussianField-NuclearTensorProduct-instTopologicalSpace, GaussianField-circleHasLaplacianEigenvalues, GaussianField-NuclearTensorProduct-dyninMityaginSpace, GaussianField-SmoothMap-Circle-instAddCommGroup, GaussianField-tensorProductHasLaplacianEigenvalues, GaussianField-NuclearTensorProduct-instAddCommGroup, GaussianField-SmoothMap-Circle-instModule, GaussianField-SmoothMap-Circle-instContinuousSMul, GaussianField-smoothCircle-dyninMityaginSpace, GaussianField-SmoothMap-Circle-instIsTopologicalAddGroup, GaussianField-SmoothMap-Circle, GaussianField-AsymTorusTestFunction, GaussianField-SmoothMap-Circle-instTopologicalSpace, GaussianField-greenFunctionBilinear}
  ``\texttt{lean noncomputable def asymTorusContinuumGreen (mass : ℝ) (hmass : 0 < mass) (f g : AsymTorusTestFunction Lt Ls) : ℝ }``  The continuum Green's function on the asymmetric torus: $G_{L_t,L_s}(f,g) = \sum_{n_1,n_2} \hat{c}_{n_1}(f_1)\, \hat{c}_{n_2}(f_2) / ((2\pi n_1/L_t)^2 + (2\pi n_2/L_s)^2 + m^2)$.
\end{definition}
\begin{theorem}[\texttt{congr\_simp}]\label{Pphi2-asymLatticeTestFn-congr-simp}
  \lean{Pphi2.asymLatticeTestFn.congr_simp}\leanok
  \uses{GaussianField-FinLatticeSites, Pphi2-asymLatticeTestFn, GaussianField-AsymTorusTestFunction}
  \texttt{Pphi2.asymLatticeTestFn.congr\_simp}
\end{theorem}
\begin{theorem}[\texttt{asymSpaceSpacing\_pos}]\label{Pphi2-asymSpaceSpacing-pos}
  \lean{Pphi2.asymSpaceSpacing_pos}\leanok
  \uses{Pphi2-asymSpaceSpacing}
  ``\texttt{lean theorem asymSpaceSpacing\_pos (Ns : ℕ) [NeZero Ns] : 0 < asymSpaceSpacing Ls Ns }``  Informal statement: $a_s > 0$.  Proof: Analogous to \texttt{asymTimeSpacing\_pos}.
\end{theorem}
\begin{theorem}[\texttt{asymTorusEmbedLift\_measurable}]\label{Pphi2-asymTorusEmbedLift-measurable}
  \lean{Pphi2.asymTorusEmbedLift_measurable}\leanok
  \uses{GaussianField-NuclearTensorProduct-instModuleReal, GaussianField-NuclearTensorProduct-instTopologicalSpace, GaussianField-FinLatticeSites, GaussianField-NuclearTensorProduct-instIsTopologicalAddGroup, GaussianField-configuration-measurable-of-eval-measurable, GaussianField-NuclearTensorProduct-instAddCommGroup, Pphi2-evalAsymAtFinSiteGJ, GaussianField-configuration-eval-measurable, GaussianField-Configuration, GaussianField-NuclearTensorProduct-instContinuousSMulReal, GaussianField-SmoothMap-Circle, GaussianField-instMeasurableSpaceConfiguration, Pphi2-asymTorusEmbedLift, GaussianField-AsymTorusTestFunction, GaussianField-FinLatticeField}
  ``\texttt{lean theorem asymTorusEmbedLift\_measurable (N : ℕ) [NeZero N] : Measurable (asymTorusEmbedLift Lt Ls N) }`\texttt{  Informal statement: The embedding lift is measurable.  Proof: Fully proved via }configuration\_measurable\_of\_eval\_measurable\texttt{ and }Finset.measurable\_sum`.
\end{theorem}
\begin{definition}[\texttt{asymTimeSpacing}]\label{Pphi2-asymTimeSpacing}
  \lean{Pphi2.asymTimeSpacing}\leanok
  ``\texttt{lean def asymTimeSpacing (Nt : ℕ) [NeZero Nt] : ℝ := Lt / Nt }``  Time lattice spacing $a_t = L_t / N_t$.
\end{definition}
\begin{theorem}[\texttt{congr\_simp}]\label{Pphi2-evalAsymAtFinSiteGJ-congr-simp}
  \lean{Pphi2.evalAsymAtFinSiteGJ.congr_simp}\leanok
  \uses{GaussianField-NuclearTensorProduct-instModuleReal, GaussianField-NuclearTensorProduct-instTopologicalSpace, GaussianField-FinLatticeSites, GaussianField-NuclearTensorProduct-instAddCommGroup, Pphi2-evalAsymAtFinSiteGJ, GaussianField-SmoothMap-Circle, GaussianField-AsymTorusTestFunction}
  \texttt{Pphi2.evalAsymAtFinSiteGJ.congr\_simp}
\end{theorem}
\begin{theorem}[\texttt{asymTimeSpacing\_pos}]\label{Pphi2-asymTimeSpacing-pos}
  \lean{Pphi2.asymTimeSpacing_pos}\leanok
  \uses{Pphi2-asymTimeSpacing}
  ``\texttt{lean theorem asymTimeSpacing\_pos (Nt : ℕ) [NeZero Nt] : 0 < asymTimeSpacing Lt Nt }``  Informal statement: $a_t > 0$.  Proof: Direct from \texttt{div\_pos} and positivity of $L_t$ and $N_t$.
\end{theorem}
\begin{definition}[\texttt{finToAsymSite}]\label{Pphi2-finToAsymSite}
  \lean{Pphi2.finToAsymSite}\leanok
  \uses{GaussianField-AsymLatticeSites, GaussianField-FinLatticeSites}
  ``\texttt{lean def finToAsymSite (N : ℕ) (x : FinLatticeSites 2 N) : AsymLatticeSites N N }`\texttt{  Converts }FinLatticeSites 2 N = Fin 2 -> ZMod N\texttt{ to }AsymLatticeSites N N = ZMod N x ZMod N`.
\end{definition}
\begin{definition}[\texttt{asymTorusInteractingMeasure}]\label{Pphi2-asymTorusInteractingMeasure}
  \lean{Pphi2.asymTorusInteractingMeasure}\leanok
  \uses{GaussianField-NuclearTensorProduct-instModuleReal, GaussianField-NuclearTensorProduct-instTopologicalSpace, GaussianField-FinLatticeSites, GaussianField-NuclearTensorProduct-instIsTopologicalAddGroup, Pphi2-interactingLatticeMeasure, Pphi2-asymGeomSpacing-pos, GaussianField-NuclearTensorProduct-instAddCommGroup, GaussianField-Configuration, GaussianField-NuclearTensorProduct-instContinuousSMulReal, GaussianField-SmoothMap-Circle, GaussianField-instMeasurableSpaceConfiguration, Pphi2-asymTorusEmbedLift, GaussianField-AsymTorusTestFunction, GaussianField-FinLatticeField, Pphi2-asymGeomSpacing}
  ``\texttt{lean def asymTorusInteractingMeasure (N : ℕ) [NeZero N] (P : InteractionPolynomial) (mass : ℝ) (hmass : 0 < mass) : Measure (Configuration (AsymTorusTestFunction Lt Ls)) }``  The interacting $P(\varphi)_2$ measure on the asymmetric torus, defined as the pushforward of the lattice interacting measure through \texttt{asymTorusEmbedLift}, using the geometric mean spacing as the effective lattice spacing.
\end{definition}
\begin{definition}[\texttt{asymGeomSpacing}]\label{Pphi2-asymGeomSpacing}
  \lean{Pphi2.asymGeomSpacing}\leanok
  \uses{Pphi2-asymTimeSpacing, Pphi2-asymSpaceSpacing}
  ``\texttt{lean noncomputable def asymGeomSpacing (N : ℕ) [NeZero N] : ℝ := Real.sqrt (asymTimeSpacing Lt N * asymSpaceSpacing Ls N) }``  The geometric mean spacing $a_{\mathrm{geom}} = \sqrt{a_t \cdot a_s} = \sqrt{L_t \cdot L_s}/N$.
\end{definition}
\begin{definition}[\texttt{asymTorusEmbedLift}]\label{Pphi2-asymTorusEmbedLift}
  \lean{Pphi2.asymTorusEmbedLift}\leanok
  \uses{GaussianField-NuclearTensorProduct-instModuleReal, GaussianField-NuclearTensorProduct-instTopologicalSpace, GaussianField-FinLatticeSites, GaussianField-NuclearTensorProduct-instIsTopologicalAddGroup, GaussianField-NuclearTensorProduct-instAddCommGroup, Pphi2-evalAsymAtFinSiteGJ, GaussianField-Configuration, GaussianField-NuclearTensorProduct-instContinuousSMulReal, GaussianField-SmoothMap-Circle, GaussianField-AsymTorusTestFunction, GaussianField-FinLatticeField}
  ``\texttt{lean def asymTorusEmbedLift (N : ℕ) [NeZero N] : Configuration (FinLatticeField 2 N) → Configuration (AsymTorusTestFunction Lt Ls) }`\texttt{  Maps lattice configurations on }FinLatticeField 2 N\texttt{ to configurations on }AsymTorusTestFunction Lt Ls` via pointwise summation: $(\iota\, \omega)(f) = \sum_x \omega(\delta_x) \cdot \mathrm{eval}_x(f)$.
\end{definition}
\begin{theorem}[\texttt{asymGeomSpacing\_pos}]\label{Pphi2-asymGeomSpacing-pos}
  \lean{Pphi2.asymGeomSpacing_pos}\leanok
  \uses{Pphi2-asymTimeSpacing, Pphi2-asymSpaceSpacing-pos, Pphi2-asymSpaceSpacing, Pphi2-asymTimeSpacing-pos, Pphi2-asymGeomSpacing}
  ``\texttt{lean theorem asymGeomSpacing\_pos (N : ℕ) [NeZero N] : 0 < asymGeomSpacing Lt Ls N }``  Informal statement: $a_{\mathrm{geom}} > 0$.  Proof: From \texttt{Real.sqrt\_pos\_of\_pos} and positivity of the individual spacings.
\end{theorem}
\section{\texttt{Pphi2.InteractingMeasure.U4SecondDerivBound}}
\begin{theorem}[\texttt{congr\_simp}]\label{Pphi2-partitionFnDeriv2-congr-simp}
  \lean{Pphi2.partitionFnDeriv2.congr_simp}\leanok
  \uses{Pphi2-partitionFnDeriv2}
  \texttt{Pphi2.partitionFnDeriv2.congr\_simp}
\end{theorem}
\begin{theorem}[\texttt{congr\_simp}]\label{Pphi2-momentDeriv-congr-simp}
  \lean{Pphi2.momentDeriv.congr_simp}\leanok
  \uses{Pphi2-momentDeriv, GaussianField-FinLatticeField}
  \texttt{Pphi2.momentDeriv.congr\_simp}
\end{theorem}
\begin{theorem}[\texttt{normalizedMoment\_abs\_le}]\label{Pphi2-normalizedMoment-abs-le}
  \lean{Pphi2.normalizedMoment_abs_le}\leanok
  \uses{Pphi2-torus-normConst-field-moment-uniform, GaussianField-FinLatticeSites, Pphi2-integrable-pow-pairing, GaussianField-latticeGaussianMeasure, GaussianField-configuration-eval-measurable, GaussianField-Configuration, Pphi2-expMoment-two-le-uniform, Pphi2-normalizedMoment, GaussianField-circleSpacing, GaussianField-instMeasurableSpaceConfiguration, Pphi2-abs-interacting-moment-le, GaussianField-FinLatticeField, Pphi2-interactionFunctional, GaussianField-circleSpacing-pos}
  \texttt{Pphi2.normalizedMoment\_abs\_le}
\end{theorem}
\begin{theorem}[\texttt{congr\_simp}]\label{Pphi2-partitionFn-congr-simp}
  \lean{Pphi2.partitionFn.congr_simp}\leanok
  \uses{Pphi2-partitionFn}
  \texttt{Pphi2.partitionFn.congr\_simp}
\end{theorem}
\begin{theorem}[\texttt{u4Deriv2\_abs\_le\_uniform}]\label{Pphi2-u4Deriv2-abs-le-uniform}
  \lean{Pphi2.u4Deriv2_abs_le_uniform}\leanok
  \uses{Pphi2-normalizedMoment-abs-le, GaussianField-FinLatticeSites, Pphi2-u4Deriv2, Pphi2-normalizedMomentDeriv-abs-le, Pphi2-normalizedMomentDeriv2, Pphi2-normalizedMoment, GaussianField-circleSpacing, Pphi2-normalizedMomentDeriv2-abs-le, Pphi2-normalizedMomentDeriv, Lattice-toSemilatticeInf, GaussianField-circleSpacing-pos}
  \texttt{Pphi2.u4Deriv2\_abs\_le\_uniform}
\end{theorem}
\begin{theorem}[\texttt{congr\_simp}]\label{Pphi2-partitionFnDeriv-congr-simp}
  \lean{Pphi2.partitionFnDeriv.congr_simp}\leanok
  \uses{Pphi2-partitionFnDeriv}
  \texttt{Pphi2.partitionFnDeriv.congr\_simp}
\end{theorem}
\begin{theorem}[\texttt{ratio\_l2\_bound}]\label{Pphi2-ratio-l2-bound}
  \lean{Pphi2.ratio_l2_bound}\leanok
  \uses{GaussianField-FinLatticeSites, GaussianField-latticeGaussianMeasure, GaussianField-Configuration, GaussianField-circleSpacing, GaussianField-instMeasurableSpaceConfiguration, Pphi2-abs-interacting-moment-le, GaussianField-FinLatticeField, Pphi2-interactionFunctional, GaussianField-circleSpacing-pos}
  \texttt{Pphi2.ratio\_l2\_bound}
\end{theorem}
\begin{theorem}[\texttt{congr\_simp}]\label{Pphi2-momentDeriv2-congr-simp}
  \lean{Pphi2.momentDeriv2.congr_simp}\leanok
  \uses{Pphi2-momentDeriv2, GaussianField-FinLatticeField}
  \texttt{Pphi2.momentDeriv2.congr\_simp}
\end{theorem}
\begin{theorem}[\texttt{normalizedMomentDeriv2\_abs\_le}]\label{Pphi2-normalizedMomentDeriv2-abs-le}
  \lean{Pphi2.normalizedMomentDeriv2_abs_le}\leanok
  \uses{Pphi2-ratio-l2-bound, Pphi2-partitionFn-ge-one, Pphi2-torus-normConst-field-moment-uniform, Pphi2-partitionFnDeriv, GaussianField-FinLatticeSites, Pphi2-partitionFnDeriv2, Pphi2-gibbsMoment, Pphi2-integrable-pow-pairing, Pphi2-normalizedMomentDeriv2, Pphi2-torus-free-product-moment-uniform, Pphi2-memLp-pairing-pow-mul-interaction-pow, GaussianField-latticeGaussianMeasure, GaussianField-configuration-eval-measurable, GaussianField-Configuration, Pphi2-expMoment-two-le-uniform, GaussianField-circleSpacing, GaussianField-instMeasurableSpaceConfiguration, Pphi2-momentDeriv2, Pphi2-partitionFn, Pphi2-momentDeriv, GaussianField-FinLatticeField, Pphi2-interactionFunctional, Lattice-toSemilatticeInf, GaussianField-circleSpacing-pos}
  \texttt{Pphi2.normalizedMomentDeriv2\_abs\_le}
\end{theorem}
\begin{theorem}[\texttt{congr\_simp}]\label{Pphi2-gibbsMoment-congr-simp}
  \lean{Pphi2.gibbsMoment.congr_simp}\leanok
  \uses{Pphi2-gibbsMoment, GaussianField-FinLatticeField}
  \texttt{Pphi2.gibbsMoment.congr\_simp}
\end{theorem}
\begin{theorem}[\texttt{normalizedMomentDeriv\_abs\_le}]\label{Pphi2-normalizedMomentDeriv-abs-le}
  \lean{Pphi2.normalizedMomentDeriv_abs_le}\leanok
  \uses{Pphi2-ratio-l2-bound, Pphi2-partitionFn-ge-one, Pphi2-torus-normConst-field-moment-uniform, Pphi2-partitionFnDeriv, GaussianField-FinLatticeSites, Pphi2-gibbsMoment, Pphi2-integrable-pow-pairing, Pphi2-torus-free-product-moment-uniform, Pphi2-memLp-pairing-pow-mul-interaction-pow, GaussianField-latticeGaussianMeasure, GaussianField-configuration-eval-measurable, GaussianField-Configuration, Pphi2-expMoment-two-le-uniform, GaussianField-circleSpacing, GaussianField-instMeasurableSpaceConfiguration, Pphi2-partitionFn, Pphi2-normalizedMomentDeriv, Pphi2-momentDeriv, GaussianField-FinLatticeField, Pphi2-interactionFunctional, Lattice-toSemilatticeInf, GaussianField-circleSpacing-pos}
  \texttt{Pphi2.normalizedMomentDeriv\_abs\_le}
\end{theorem}
\section{\texttt{Pphi2.TorusContinuumLimit.TorusPropagatorConvergence}}
\begin{theorem}[\texttt{congr\_simp}]\label{Pphi2-latticeTestFn-congr-simp}
  \lean{Pphi2.latticeTestFn.congr_simp}\leanok
  \uses{GaussianField-FinLatticeSites, Pphi2-latticeTestFn, GaussianField-TorusTestFunction}
  \texttt{Pphi2.latticeTestFn.congr\_simp}
\end{theorem}
\begin{theorem}[\texttt{torusEmbeddedTwoPoint\_eq\_lattice\_second\_moment}]\label{Pphi2-torusEmbeddedTwoPoint-eq-lattice-second-moment}
  \lean{Pphi2.torusEmbeddedTwoPoint_eq_lattice_second_moment}\leanok
  \uses{GaussianField-NuclearTensorProduct-instModuleReal, GaussianField-NuclearTensorProduct-instTopologicalSpace, Pphi2-torusEmbeddedTwoPoint, GaussianField-FinLatticeSites, GaussianField-NuclearTensorProduct-instIsTopologicalAddGroup, Pphi2-torusEmbedLift, Pphi2-torusContinuumMeasure, Pphi2-torusEmbedLift-measurable, GaussianField-NuclearTensorProduct-instAddCommGroup, GaussianField-latticeGaussianMeasure, GaussianField-configuration-eval-measurable, GaussianField-Configuration, GaussianField-NuclearTensorProduct-instContinuousSMulReal, GaussianField-SmoothMap-Circle, GaussianField-circleSpacing, GaussianField-instMeasurableSpaceConfiguration, Pphi2-latticeTestFn, Pphi2-torusEmbedLift-eval-eq, GaussianField-FinLatticeField, GaussianField-TorusTestFunction, GaussianField-circleSpacing-pos}
  \texttt{Pphi2.torusEmbeddedTwoPoint\_eq\_lattice\_second\_moment}
\end{theorem}
\begin{theorem}[\texttt{torusEmbeddedTwoPoint\_le\_seminorm\_tight}]\label{Pphi2-torusEmbeddedTwoPoint-le-seminorm-tight}
  \lean{Pphi2.torusEmbeddedTwoPoint_le_seminorm_tight}\leanok
  \uses{Pphi2-torusEmbeddedTwoPoint-eq-lattice-second-moment, GaussianField-NuclearTensorProduct-instModuleReal, Pphi2-dm-basis-eq-fourierBasis, GaussianField-NuclearTensorProduct-instTopologicalSpace, GaussianField-DyninMityaginSpace-coeff, GaussianField-latticeCovariance-GJ-eq-inv-smul-bare, GaussianField-SmoothMap-Circle-instFunLike, Pphi2-lattice-second-moment-eq-inner, GaussianField-RapidDecaySeq-rapidDecaySeminorm, Pphi2-evalTorusAtSite-basisVec, GaussianField-RapidDecaySeq, GaussianField-SmoothMap-Circle-fourierBasis, Pphi2-torusEmbeddedTwoPoint, GaussianField-covariance, GaussianField-latticeCovariance, GaussianField-RapidDecaySeq-val, GaussianField-NuclearTensorProduct-dyninMityaginSpace, GaussianField-SmoothMap-Circle-instAddCommGroup, GaussianField-FinLatticeSites, GaussianField-NuclearTensorProduct-instIsTopologicalAddGroup, GaussianField-SmoothMap-Circle-norm-iteratedDeriv-le-sobolevSeminorm-, GaussianField-circleRestriction, GaussianField-DyninMityaginSpace-expansion, GaussianField-latticeCovarianceGJ, GaussianField-RapidDecaySeq-instSMul, GaussianField-SmoothMap-Circle-instSMul, GaussianField-RapidDecaySeq-basisVec, GaussianField-circleSpacing-eq, GaussianField-NuclearTensorProduct-instAddCommGroup, GaussianField-circleRestriction-apply, GaussianField-evalTorusAtSiteGJ-apply-, GaussianField-latticeGaussianMeasure, GaussianField-evalTorusAtSite, GaussianField-SmoothMap-Circle-instModule, GaussianField-SmoothMap-Circle-instContinuousSMul, GaussianField-Configuration, GaussianField-smoothCircle-dyninMityaginSpace, GaussianField-SmoothMap-Circle-instIsTopologicalAddGroup, GaussianField-DyninMityaginSpace-basis, GaussianField-NuclearTensorProduct-instContinuousSMulReal, GaussianField-RapidDecaySeq-instAddCommGroup, GaussianField-SmoothMap-Circle, GaussianField-circleSpacing, GaussianField-instMeasurableSpaceConfiguration, GaussianField-SmoothMap-Circle-sobolevSeminorm, Pphi2-latticeTestFn, Pphi2-covariance-inner-le-mass-inv-sq-norm-sq, GaussianField-ell2-, GaussianField-RapidDecaySeq-rapid-decay, GaussianField-FinLatticeField, GaussianField-SmoothMap-Circle-sobolevSeminorm-fourierBasis-le, GaussianField-TorusTestFunction, GaussianField-SmoothMap-Circle-instTopologicalSpace, GaussianField-circlePoint, Lattice-toSemilatticeInf, GaussianField-circleSpacing-pos}
  \texttt{Pphi2.torusEmbeddedTwoPoint\_le\_seminorm\_tight}
\end{theorem}
\begin{theorem}[\texttt{congr\_simp}]\label{Pphi2-torusContinuumGreen-congr-simp}
  \lean{Pphi2.torusContinuumGreen.congr_simp}\leanok
  \uses{Pphi2-torusContinuumGreen, GaussianField-TorusTestFunction}
  \texttt{Pphi2.torusContinuumGreen.congr\_simp}
\end{theorem}
\begin{theorem}[\texttt{covariance\_inner\_le\_mass\_inv\_sq\_norm\_sq}]\label{Pphi2-covariance-inner-le-mass-inv-sq-norm-sq}
  \lean{Pphi2.covariance_inner_le_mass_inv_sq_norm_sq}\leanok
  \uses{GaussianField-latticeCovariance, GaussianField-FinLatticeSites, GaussianField-massEigenbasis-sum-mul-sum-eq-site-inner, GaussianField-spectralLatticeCovariance-norm-sq, Pphi2-massEigenvalues-ge-mass-sq, GaussianField-massEigenvalues, GaussianField-spectralLatticeCovariance, GaussianField-massEigenvectorBasis, GaussianField-ell2-, GaussianField-FinLatticeField, Lattice-toSemilatticeInf}
  $\langle Tg, Tg\rangle \le m^{-2} \sum_x g(x)^2$. Proved from eigenvalue bound + Parseval.
\end{theorem}
\begin{theorem}[\texttt{lattice\_second\_moment\_le\_mass\_inv}]\label{Pphi2-lattice-second-moment-le-mass-inv}
  \lean{Pphi2.lattice_second_moment_le_mass_inv}\leanok
  \uses{GaussianField-latticeCovariance-GJ-eq-inv-smul-bare, GaussianField-finLatticeField-dyninMityaginSpace, GaussianField-measure, GaussianField-ell2-separableSpace, GaussianField-second-moment-eq-covariance, GaussianField-latticeCovariance, GaussianField-FinLatticeSites, GaussianField-latticeCovarianceGJ, GaussianField-latticeGaussianMeasure, GaussianField-Configuration, GaussianField-instMeasurableSpaceConfiguration, Pphi2-covariance-inner-le-mass-inv-sq-norm-sq, GaussianField-ell2-, GaussianField-FinLatticeField}
  \texttt{Pphi2.lattice\_second\_moment\_le\_mass\_inv}
\end{theorem}
\begin{theorem}[\texttt{massEigenvalues\_ge\_mass\_sq}]\label{Pphi2-massEigenvalues-ge-mass-sq}
  \lean{Pphi2.massEigenvalues_ge_mass_sq}\leanok
  \uses{GaussianField-massOperatorMatrix-isHermitian, GaussianField-massOperatorMatrix, GaussianField-FinLatticeSites, GaussianField-massOperator, GaussianField-massMatrixHerm, GaussianField-massOperator-eq-matrix-mulVec, GaussianField-massEigenvalues, GaussianField-finiteLaplacian-neg-semidefinite, GaussianField-finiteLaplacian, GaussianField-FinLatticeField, Lattice-toSemilatticeInf}
  $\lambda_k \ge m^2$ for all $k$. Proved by spectral decomposition: $\langle e_k, (-\Delta+m^2)e_k\rangle \ge m^2$ since $-\Delta \ge 0$.
\end{theorem}
\begin{theorem}[\texttt{evalTorusAtSite\_basisVec}]\label{Pphi2-evalTorusAtSite-basisVec}
  \lean{Pphi2.evalTorusAtSite_basisVec}\leanok
  \uses{GaussianField-NuclearTensorProduct-instModuleReal, GaussianField-NuclearTensorProduct-instTopologicalSpace, GaussianField-RapidDecaySeq, GaussianField-SmoothMap-Circle-instAddCommGroup, GaussianField-FinLatticeSites, GaussianField-circleRestriction, GaussianField-NuclearTensorProduct, GaussianField-RapidDecaySeq-basisVec, GaussianField-NuclearTensorProduct-instAddCommGroup, GaussianField-evalTorusAtSite, GaussianField-SmoothMap-Circle-instModule, GaussianField-SmoothMap-Circle-instContinuousSMul, GaussianField-smoothCircle-dyninMityaginSpace, GaussianField-SmoothMap-Circle-instIsTopologicalAddGroup, GaussianField-DyninMityaginSpace-basis, GaussianField-NuclearTensorProduct-pure, GaussianField-SmoothMap-Circle, GaussianField-NuclearTensorProduct-evalCLM-pure, GaussianField-TorusTestFunction, GaussianField-NuclearTensorProduct-evalCLM, GaussianField-SmoothMap-Circle-instTopologicalSpace}
  \texttt{Pphi2.evalTorusAtSite\_basisVec}
\end{theorem}
\begin{theorem}[\texttt{torusEmbeddedTwoPoint\_eq\_lattice\_cross\_moment}]\label{Pphi2-torusEmbeddedTwoPoint-eq-lattice-cross-moment}
  \lean{Pphi2.torusEmbeddedTwoPoint_eq_lattice_cross_moment}\leanok
  \uses{GaussianField-NuclearTensorProduct-instModuleReal, GaussianField-NuclearTensorProduct-instTopologicalSpace, Pphi2-torusEmbeddedTwoPoint, GaussianField-FinLatticeSites, GaussianField-NuclearTensorProduct-instIsTopologicalAddGroup, Pphi2-torusEmbedLift, Pphi2-torusContinuumMeasure, Pphi2-torusEmbedLift-measurable, GaussianField-NuclearTensorProduct-instAddCommGroup, GaussianField-latticeGaussianMeasure, GaussianField-configuration-eval-measurable, GaussianField-Configuration, GaussianField-NuclearTensorProduct-instContinuousSMulReal, GaussianField-SmoothMap-Circle, GaussianField-circleSpacing, GaussianField-instMeasurableSpaceConfiguration, Pphi2-latticeTestFn, Pphi2-torusEmbedLift-eval-eq, GaussianField-FinLatticeField, GaussianField-TorusTestFunction, GaussianField-circleSpacing-pos}
  \texttt{Pphi2.torusEmbeddedTwoPoint\_eq\_lattice\_cross\_moment}
\end{theorem}
\begin{theorem}[\texttt{torusEmbeddedTwoPoint\_eq\_spectral\_sum}]\label{Pphi2-torusEmbeddedTwoPoint-eq-spectral-sum}
  \lean{Pphi2.torusEmbeddedTwoPoint_eq_spectral_sum}\leanok
  \uses{Pphi2-torusEmbeddedTwoPoint, GaussianField-covariance, GaussianField-FinLatticeSites, GaussianField-latticeCovarianceGJ, GaussianField-lattice-covariance-GJ-eq-spectral, GaussianField-latticeGaussianMeasure, GaussianField-Configuration, GaussianField-massEigenvalues, GaussianField-massEigenvectorBasis, Pphi2-torusEmbeddedTwoPoint-eq-lattice-cross-moment, GaussianField-lattice-cross-moment, GaussianField-circleSpacing, GaussianField-instMeasurableSpaceConfiguration, Pphi2-latticeTestFn, GaussianField-ell2-, GaussianField-FinLatticeField, GaussianField-TorusTestFunction, GaussianField-circleSpacing-pos}
  $\langle\Phi_N(f), \Phi_N(g)\rangle = \sum_k \lambda_k^{-1} c_k(\iota^* f) c_k(\iota^* g)$.
\end{theorem}
\begin{theorem}[\texttt{latticeTestFn\_norm\_sq\_bounded}]\label{Pphi2-latticeTestFn-norm-sq-bounded}
  \lean{Pphi2.latticeTestFn_norm_sq_bounded}\leanok
  \uses{GaussianField-NuclearTensorProduct-instModuleReal, Pphi2-dm-basis-eq-fourierBasis, GaussianField-NuclearTensorProduct-instTopologicalSpace, GaussianField-DyninMityaginSpace-coeff, GaussianField-SmoothMap-Circle-instFunLike, GaussianField-RapidDecaySeq-rapidDecaySeminorm, Pphi2-evalTorusAtSite-basisVec, GaussianField-RapidDecaySeq, GaussianField-SmoothMap-Circle-fourierBasis, GaussianField-RapidDecaySeq-val, GaussianField-NuclearTensorProduct-dyninMityaginSpace, GaussianField-SmoothMap-Circle-instAddCommGroup, GaussianField-FinLatticeSites, GaussianField-NuclearTensorProduct-instIsTopologicalAddGroup, GaussianField-SmoothMap-Circle-norm-iteratedDeriv-le-sobolevSeminorm-, GaussianField-circleRestriction, GaussianField-DyninMityaginSpace-expansion, GaussianField-RapidDecaySeq-instSMul, GaussianField-SmoothMap-Circle-instSMul, GaussianField-RapidDecaySeq-basisVec, GaussianField-circleSpacing-eq, GaussianField-NuclearTensorProduct-instAddCommGroup, GaussianField-circleRestriction-apply, GaussianField-evalTorusAtSiteGJ-apply-, GaussianField-evalTorusAtSite, GaussianField-SmoothMap-Circle-instModule, GaussianField-SmoothMap-Circle-instContinuousSMul, GaussianField-smoothCircle-dyninMityaginSpace, GaussianField-SmoothMap-Circle-instIsTopologicalAddGroup, GaussianField-DyninMityaginSpace-basis, GaussianField-NuclearTensorProduct-instContinuousSMulReal, GaussianField-RapidDecaySeq-instAddCommGroup, GaussianField-SmoothMap-Circle, GaussianField-circleSpacing, GaussianField-evalTorusAtSiteGJ, GaussianField-SmoothMap-Circle-sobolevSeminorm, Pphi2-latticeTestFn, GaussianField-RapidDecaySeq-rapid-decay, GaussianField-SmoothMap-Circle-sobolevSeminorm-fourierBasis-le, GaussianField-TorusTestFunction, GaussianField-SmoothMap-Circle-instTopologicalSpace, GaussianField-circlePoint, Lattice-toSemilatticeInf}
  $\sum_x (\iota^* f)(x)^2 \le C$ uniformly in $N$. Proved via DyninMityaginSpace expansion, Fourier basis bounds, and cancellation of $N$.
\end{theorem}
\begin{definition}[\texttt{latticeTestFn}]\label{Pphi2-latticeTestFn}
  \lean{Pphi2.latticeTestFn}\leanok
  \uses{GaussianField-NuclearTensorProduct-instModuleReal, GaussianField-NuclearTensorProduct-instTopologicalSpace, GaussianField-FinLatticeSites, GaussianField-NuclearTensorProduct-instAddCommGroup, GaussianField-SmoothMap-Circle, GaussianField-evalTorusAtSiteGJ, GaussianField-FinLatticeField, GaussianField-TorusTestFunction}
  \texttt{Pphi2.latticeTestFn}
\end{definition}
\begin{theorem}[\texttt{lattice\_second\_moment\_eq\_inner}]\label{Pphi2-lattice-second-moment-eq-inner}
  \lean{Pphi2.lattice_second_moment_eq_inner}\leanok
  \uses{GaussianField-finLatticeField-dyninMityaginSpace, GaussianField-ell2-separableSpace, GaussianField-second-moment-eq-covariance, GaussianField-FinLatticeSites, GaussianField-latticeCovarianceGJ, GaussianField-latticeGaussianMeasure, GaussianField-Configuration, GaussianField-circleSpacing, GaussianField-instMeasurableSpaceConfiguration, GaussianField-ell2-, GaussianField-FinLatticeField, GaussianField-circleSpacing-pos}
  \texttt{Pphi2.lattice\_second\_moment\_eq\_inner}
\end{theorem}
\begin{theorem}[\texttt{torusContinuumGreen\_nonneg}]\label{Pphi2-torusContinuumGreen-nonneg}
  \lean{Pphi2.torusContinuumGreen_nonneg}\leanok
  \uses{Pphi2-torusContinuumGreen, GaussianField-NuclearTensorProduct-instModuleReal, GaussianField-NuclearTensorProduct-instTopologicalSpace, GaussianField-circleHasLaplacianEigenvalues, GaussianField-NuclearTensorProduct-dyninMityaginSpace, GaussianField-SmoothMap-Circle-instAddCommGroup, GaussianField-NuclearTensorProduct-instIsTopologicalAddGroup, GaussianField-greenFunctionBilinear-nonneg, GaussianField-tensorProductHasLaplacianEigenvalues, GaussianField-NuclearTensorProduct-instAddCommGroup, GaussianField-SmoothMap-Circle-instModule, GaussianField-SmoothMap-Circle-instContinuousSMul, GaussianField-smoothCircle-dyninMityaginSpace, GaussianField-SmoothMap-Circle-instIsTopologicalAddGroup, GaussianField-NuclearTensorProduct-instContinuousSMulReal, GaussianField-SmoothMap-Circle, GaussianField-TorusTestFunction, GaussianField-SmoothMap-Circle-instTopologicalSpace}
  \texttt{Pphi2.torusContinuumGreen\_nonneg}
\end{theorem}
\begin{theorem}[\texttt{torusEmbeddedTwoPoint\_uniform\_bound}]\label{Pphi2-torusEmbeddedTwoPoint-uniform-bound}
  \lean{Pphi2.torusEmbeddedTwoPoint_uniform_bound}\leanok
  \uses{GaussianField-RapidDecaySeq-rapidDecaySeminorm, GaussianField-RapidDecaySeq, Pphi2-torusEmbeddedTwoPoint, GaussianField-RapidDecaySeq-instSMul, GaussianField-RapidDecaySeq-instAddCommGroup, Pphi2-torusEmbeddedTwoPoint-le-seminorm-tight, GaussianField-TorusTestFunction, Lattice-toSemilatticeInf}
  $E[\Phi_N(f)^2] \le m^{-2} C_f$ uniformly in $N$.
\end{theorem}
\begin{theorem}[\texttt{torusContinuumGreen\_pos}]\label{Pphi2-torusContinuumGreen-pos}
  \lean{Pphi2.torusContinuumGreen_pos}\leanok
  \uses{Pphi2-torusContinuumGreen, GaussianField-NuclearTensorProduct-instModuleReal, GaussianField-greenFunctionBilinear-pos, GaussianField-NuclearTensorProduct-instTopologicalSpace, GaussianField-DyninMityaginSpace-toT1Space, GaussianField-circleHasLaplacianEigenvalues, GaussianField-NuclearTensorProduct-dyninMityaginSpace, GaussianField-SmoothMap-Circle-instAddCommGroup, GaussianField-NuclearTensorProduct-instIsTopologicalAddGroup, GaussianField-tensorProductHasLaplacianEigenvalues, GaussianField-NuclearTensorProduct-instAddCommGroup, GaussianField-SmoothMap-Circle-instModule, GaussianField-SmoothMap-Circle-instContinuousSMul, GaussianField-smoothCircle-dyninMityaginSpace, GaussianField-SmoothMap-Circle-instIsTopologicalAddGroup, GaussianField-NuclearTensorProduct-instContinuousSMulReal, GaussianField-SmoothMap-Circle, GaussianField-TorusTestFunction, GaussianField-SmoothMap-Circle-instTopologicalSpace}
  $G_L(f,f) > 0$ for $f \ne 0$; $G_L(f,f) \ge 0$ for all $f$.
\end{theorem}
\begin{theorem}[\texttt{congr\_simp}]\label{Pphi2-torusEmbedLift-congr-simp}
  \lean{Pphi2.torusEmbedLift.congr_simp}\leanok
  \uses{GaussianField-NuclearTensorProduct-instModuleReal, GaussianField-NuclearTensorProduct-instTopologicalSpace, GaussianField-FinLatticeSites, GaussianField-NuclearTensorProduct-instIsTopologicalAddGroup, Pphi2-torusEmbedLift, GaussianField-NuclearTensorProduct-instAddCommGroup, GaussianField-Configuration, GaussianField-NuclearTensorProduct-instContinuousSMulReal, GaussianField-SmoothMap-Circle, GaussianField-FinLatticeField, GaussianField-TorusTestFunction}
  \texttt{Pphi2.torusEmbedLift.congr\_simp}
\end{theorem}
\begin{theorem}[\texttt{torusEmbedLift\_eval\_eq}]\label{Pphi2-torusEmbedLift-eval-eq}
  \lean{Pphi2.torusEmbedLift_eval_eq}\leanok
  \uses{GaussianField-NuclearTensorProduct-instModuleReal, GaussianField-NuclearTensorProduct-instTopologicalSpace, GaussianField-FinLatticeSites, GaussianField-NuclearTensorProduct-instIsTopologicalAddGroup, Pphi2-torusEmbedLift, GaussianField-NuclearTensorProduct-instAddCommGroup, GaussianField-Configuration, GaussianField-NuclearTensorProduct-instContinuousSMulReal, GaussianField-SmoothMap-Circle, GaussianField-evalTorusAtSiteGJ, Pphi2-latticeTestFn-expand, Pphi2-latticeTestFn, GaussianField-FinLatticeField, GaussianField-TorusTestFunction}
  $(\tilde\iota_N \omega)(f) = \omega(\iota^* f)$ where $\iota^* f$ is the pullback lattice test function.
\end{theorem}
\begin{theorem}[\texttt{torus\_propagator\_convergence}]\label{Pphi2-torus-propagator-convergence}
  \lean{Pphi2.torus_propagator_convergence}\leanok
  \uses{Pphi2-torusContinuumGreen, GaussianField-NuclearTensorProduct-instModuleReal, GaussianField-NuclearTensorProduct-instTopologicalSpace, GaussianField-latticeCovariance-GJ-eq-inv-smul-bare, Pphi2-torusEmbeddedTwoPoint, GaussianField-lattice-green-tendsto-continuum, GaussianField-covariance, GaussianField-latticeCovariance, GaussianField-FinLatticeSites, GaussianField-latticeCovarianceGJ, GaussianField-NuclearTensorProduct-instAddCommGroup, GaussianField-evalTorusAtSiteGJ-apply-, GaussianField-latticeGaussianMeasure, GaussianField-evalTorusAtSite, GaussianField-Configuration, GaussianField-SmoothMap-Circle, Pphi2-torusEmbeddedTwoPoint-eq-lattice-cross-moment, GaussianField-lattice-cross-moment, GaussianField-circleSpacing, GaussianField-instMeasurableSpaceConfiguration, GaussianField-evalTorusAtSiteGJ, Pphi2-latticeTestFn, GaussianField-ell2-, GaussianField-FinLatticeField, GaussianField-TorusTestFunction, GaussianField-circleSpacing-pos}
  $\text{torusEmbeddedTwoPoint}(f,g) \to G_L(f,g)$ as $N \to \infty$. Proved via \texttt{lattice\_green\_tendsto\_continuum} (mode-by-mode convergence + dominated convergence).
\end{theorem}
\begin{theorem}[\texttt{latticeTestFn\_expand}]\label{Pphi2-latticeTestFn-expand}
  \lean{Pphi2.latticeTestFn_expand}\leanok
  \uses{GaussianField-FinLatticeSites, Pphi2-latticeTestFn, GaussianField-FinLatticeField, GaussianField-TorusTestFunction}
  \texttt{Pphi2.latticeTestFn\_expand}
\end{theorem}
\begin{theorem}[\texttt{dm\_basis\_eq\_fourierBasis}]\label{Pphi2-dm-basis-eq-fourierBasis}
  \lean{Pphi2.dm_basis_eq_fourierBasis}\leanok
  \uses{GaussianField-SmoothMap-Circle-instFunLike, GaussianField-SmoothMap-Circle-ext, GaussianField-SmoothMap-Circle-fourierBasis, GaussianField-RapidDecaySeq-val, GaussianField-SmoothMap-Circle-instAddCommGroup, GaussianField-SmoothMap-Circle-fromRapidDecay, GaussianField-SmoothMap-Circle-fourierBasisFun, GaussianField-RapidDecaySeq-basisVec, GaussianField-SmoothMap-Circle-instModule, GaussianField-SmoothMap-Circle-instContinuousSMul, GaussianField-smoothCircle-dyninMityaginSpace, GaussianField-SmoothMap-Circle-instIsTopologicalAddGroup, GaussianField-DyninMityaginSpace-basis, GaussianField-SmoothMap-Circle, GaussianField-SmoothMap-Circle-instTopologicalSpace}
  \texttt{Pphi2.dm\_basis\_eq\_fourierBasis}
\end{theorem}
\begin{theorem}[\texttt{congr\_simp}]\label{Pphi2-torusEmbeddedTwoPoint-congr-simp}
  \lean{Pphi2.torusEmbeddedTwoPoint.congr_simp}\leanok
  \uses{Pphi2-torusEmbeddedTwoPoint, GaussianField-TorusTestFunction}
  \texttt{Pphi2.torusEmbeddedTwoPoint.congr\_simp}
\end{theorem}
\section{\texttt{Pphi2.InteractingMeasure.LatticeEuclideanTime}}
\begin{definition}[\texttt{latticeConfigEuclideanTimeShift}]\label{Pphi2-latticeConfigEuclideanTimeShift}
  \lean{Pphi2.latticeConfigEuclideanTimeShift}\leanok
  \uses{GaussianField-FinLatticeSites, GaussianField-latticeTranslation, GaussianField-Configuration, Pphi2-latticeEuclideanTimeShift, GaussianField-FinLatticeField}
  \texttt{Pphi2.latticeConfigEuclideanTimeShift}
\end{definition}
\begin{theorem}[\texttt{latticeEuclideanTimeShift\_mod}]\label{Pphi2-latticeEuclideanTimeShift-mod}
  \lean{Pphi2.latticeEuclideanTimeShift_mod}\leanok
  \uses{GaussianField-FinLatticeSites, Pphi2-latticeEuclideanTimeShift}
  \texttt{Pphi2.latticeEuclideanTimeShift\_mod}
\end{theorem}
\begin{theorem}[\texttt{latticeConfigEuclideanTimeShift\_mod}]\label{Pphi2-latticeConfigEuclideanTimeShift-mod}
  \lean{Pphi2.latticeConfigEuclideanTimeShift_mod}\leanok
  \uses{GaussianField-FinLatticeSites, GaussianField-latticeTranslation, Pphi2-latticeConfigEuclideanTimeShift, GaussianField-Configuration, Pphi2-latticeEuclideanTimeShift, Pphi2-latticeEuclideanTimeShift-mod, GaussianField-latticeTranslation-congr-simp, GaussianField-FinLatticeField}
  \texttt{Pphi2.latticeConfigEuclideanTimeShift\_mod}
\end{theorem}
\begin{definition}[\texttt{latticeEuclideanTimeShift}]\label{Pphi2-latticeEuclideanTimeShift}
  \lean{Pphi2.latticeEuclideanTimeShift}\leanok
  \uses{GaussianField-FinLatticeSites}
  \texttt{Pphi2.latticeEuclideanTimeShift}
\end{definition}
\begin{definition}[\texttt{latticeEuclideanTimeSeparation}]\label{Pphi2-latticeEuclideanTimeSeparation}
  \lean{Pphi2.latticeEuclideanTimeSeparation}\leanok
  \texttt{Pphi2.latticeEuclideanTimeSeparation}
\end{definition}
\begin{theorem}[\texttt{latticeEuclideanTimeSeparation\_eq\_min}]\label{Pphi2-latticeEuclideanTimeSeparation-eq-min}
  \lean{Pphi2.latticeEuclideanTimeSeparation_eq_min}\leanok
  \uses{Pphi2-latticeEuclideanTimeSeparation}
  \texttt{Pphi2.latticeEuclideanTimeSeparation\_eq\_min}
\end{theorem}
\begin{theorem}[\texttt{latticeEuclideanTimeSeparation\_val}]\label{Pphi2-latticeEuclideanTimeSeparation-val}
  \lean{Pphi2.latticeEuclideanTimeSeparation_val}\leanok
  \uses{Pphi2-latticeEuclideanTimeSeparation}
  \texttt{Pphi2.latticeEuclideanTimeSeparation\_val}
\end{theorem}
\begin{theorem}[\texttt{congr\_simp}]\label{Pphi2-latticeEuclideanTimeShift-congr-simp}
  \lean{Pphi2.latticeEuclideanTimeShift.congr_simp}\leanok
  \uses{Pphi2-latticeEuclideanTimeShift}
  \texttt{Pphi2.latticeEuclideanTimeShift.congr\_simp}
\end{theorem}
\begin{theorem}[\texttt{latticeEuclideanTimeSeparation\_mod}]\label{Pphi2-latticeEuclideanTimeSeparation-mod}
  \lean{Pphi2.latticeEuclideanTimeSeparation_mod}\leanok
  \uses{Pphi2-latticeEuclideanTimeSeparation, Pphi2-latticeEuclideanTimeSeparation-eq-min}
  \texttt{Pphi2.latticeEuclideanTimeSeparation\_mod}
\end{theorem}
\section{\texttt{Pphi2.TransferMatrix.TransferMatrix}}
\begin{theorem}[\texttt{timeCoupling\_eq\_zero\_iff}]\label{Pphi2-timeCoupling-eq-zero-iff}
  \lean{Pphi2.timeCoupling_eq_zero_iff}\leanok
  \uses{Pphi2-SpatialField, Pphi2-timeCoupling, Lattice-toSemilatticeInf}
  $K(\psi,\psi') = 0 \iff \psi = \psi'$. *Proof*: From \texttt{Finset.sum\_eq\_zero\_iff\_of\_nonneg} and \texttt{sq\_eq\_zero\_iff}. Fully proved.
\end{theorem}
\begin{theorem}[\texttt{congr\_simp}]\label{Pphi2-spatialAction-congr-simp}
  \lean{Pphi2.spatialAction.congr_simp}\leanok
  \uses{Pphi2-spatialAction, Pphi2-SpatialField}
  \texttt{Pphi2.spatialAction.congr\_simp}
\end{theorem}
\begin{definition}[\texttt{SpatialField}]\label{Pphi2-SpatialField}
  \lean{Pphi2.SpatialField}\leanok
  \texttt{Pphi2.SpatialField}
\end{definition}
\begin{theorem}[\texttt{timeCoupling\_nonneg}]\label{Pphi2-timeCoupling-nonneg}
  \lean{Pphi2.timeCoupling_nonneg}\leanok
  \uses{Pphi2-SpatialField, Pphi2-timeCoupling, Lattice-toSemilatticeInf}
  $K(\psi,\psi') \ge 0$ (sum of squares). *Proof*: Direct from \texttt{mul\_nonneg} and \texttt{sq\_nonneg}. Fully proved.
\end{theorem}
\begin{definition}[\texttt{timeCoupling}]\label{Pphi2-timeCoupling}
  \lean{Pphi2.timeCoupling}\leanok
  \uses{Pphi2-SpatialField}
  ``\texttt{lean def timeCoupling (psi psi' : SpatialField Ns) : R }`` Time coupling between adjacent slices: $K(\psi,\psi') = \frac{1}{2}\sum_x (\psi(x)-\psi'(x))^2$.
\end{definition}
\begin{theorem}[\texttt{transferKernel\_symmetric}]\label{Pphi2-transferKernel-symmetric}
  \lean{Pphi2.transferKernel_symmetric}\leanok
  \uses{Pphi2-spatialAction, Pphi2-SpatialField, Pphi2-transferKernel, Pphi2-wickConstant, Pphi2-timeCoupling}
  $T(\psi,\psi') = T(\psi',\psi)$. *Proof*: From $(a-b)^2 = (b-a)^2$ and commutativity of addition. Fully proved.
\end{theorem}
\begin{definition}[\texttt{spatialKinetic}]\label{Pphi2-spatialKinetic}
  \lean{Pphi2.spatialKinetic}\leanok
  \uses{Pphi2-SpatialField}
  ``\texttt{lean def spatialKinetic (a : R) (psi : SpatialField Ns) : R }`` Spatial kinetic energy: $\frac{1}{2}\sum_x a^{-2}(\psi(x+1)-\psi(x))^2$.
\end{definition}
\begin{definition}[\texttt{transferKernel}]\label{Pphi2-transferKernel}
  \lean{Pphi2.transferKernel}\leanok
  \uses{Pphi2-spatialAction, Pphi2-SpatialField, Pphi2-wickConstant, Pphi2-timeCoupling}
  ``\texttt{lean def transferKernel (P : InteractionPolynomial) (a mass : R) : SpatialField Ns -> SpatialField Ns -> R }`` Transfer matrix kernel: $T(\psi,\psi') = \exp\bigl(-K(\psi,\psi') - \frac{a}{2}h_a(\psi) - \frac{a}{2}h_a(\psi')\bigr)$.
\end{definition}
\begin{definition}[\texttt{spatialPotential}]\label{Pphi2-spatialPotential}
  \lean{Pphi2.spatialPotential}\leanok
  \uses{Pphi2-SpatialField, Pphi2-wickPolynomial}
  ``\texttt{lean def spatialPotential (P : InteractionPolynomial) (\_a mass : R) (c : R) (psi : SpatialField Ns) : R }`` Spatial potential energy: $\sum_x \bigl(\frac{1}{2}m^2\psi(x)^2 + {:}P(\psi(x)){:}_c\bigr)$.
\end{definition}
\begin{theorem}[\texttt{transferKernel\_pos}]\label{Pphi2-transferKernel-pos}
  \lean{Pphi2.transferKernel_pos}\leanok
  \uses{Pphi2-spatialAction, Pphi2-SpatialField, Pphi2-transferKernel, Pphi2-wickConstant, Pphi2-timeCoupling}
  $T(\psi,\psi') > 0$ for all $\psi,\psi'$. *Proof*: Immediate from \texttt{Real.exp\_pos}. Fully proved.
\end{theorem}
\begin{definition}[\texttt{spatialAction}]\label{Pphi2-spatialAction}
  \lean{Pphi2.spatialAction}\leanok
  \uses{Pphi2-SpatialField, Pphi2-spatialKinetic, Pphi2-spatialPotential}
  ``\texttt{lean def spatialAction (P : InteractionPolynomial) (a mass : R) (c : R) (psi : SpatialField Ns) : R }`` Full spatial action $h_a(\psi) = \text{spatialKinetic} + \text{spatialPotential}$.
\end{definition}
\section{\texttt{Pphi2.NelsonEstimate.WickBinomial}}
\begin{theorem}[\texttt{wickMonomial\_add\_binomial}]\label{Pphi2-wickMonomial-add-binomial}
  \lean{Pphi2.wickMonomial_add_binomial}\leanok
  \uses{Pphi2-wickMonomial-succ-succ, Pphi2-wickMonomial, Lattice-toSemilatticeInf}
  For all $n$, $c_1$, $c_2$, $x$, $y$: $$:(\{x+y\})\^{}n:\_\{c\_1+c\_2\} = \textbackslash sum\_\{k=0\}\^{}n \textbackslash binom\{n\}\{k\} :x\^{}k:\_\{c\_1\} \textbackslash cdot :y\^{}\{n-k\}:\_\{c\_2\}$$ Proved by strong induction on $n$, applying the Wick recurrence to both $x$ and $y$ factors, using binomial coefficient identities $(k+1)\binom{n+1}{k+1} = (n+1)\binom{n}{k}$ and $(n+1-k)\binom{n+1}{k} = (n+1)\binom{n}{k}$, and Pascal's rule via \texttt{sum\_antidiagonal\_choose\_succ\_nsmul}.
\end{theorem}
\begin{theorem}[\texttt{wickMonomial\_add\_sub\_self}]\label{Pphi2-wickMonomial-add-sub-self}
  \lean{Pphi2.wickMonomial_add_sub_self}\leanok
  \uses{Pphi2-wickMonomial, Pphi2-wickMonomial-add-binomial}
  \texttt{Pphi2.wickMonomial\_add\_sub\_self}
\end{theorem}
\begin{theorem}[\texttt{wickMonomial\_two\_add}]\label{Pphi2-wickMonomial-two-add}
  \lean{Pphi2.wickMonomial_two_add}\leanok
  \uses{Pphi2-wickMonomial}
  $n=2$: $:(x+y)^2:_{c_1+c_2} = :x^2:_{c_1} + 2xy + :y^2:_{c_2}$.
\end{theorem}
\begin{theorem}[\texttt{wickMonomial\_four\_add}]\label{Pphi2-wickMonomial-four-add}
  \lean{Pphi2.wickMonomial_four_add}\leanok
  \uses{Pphi2-wickMonomial-two, Pphi2-wickMonomial-three, Pphi2-wickMonomial, Pphi2-wickMonomial-four}
  $n=4$: explicit 5-term expansion for the $\varphi^4$ case.
\end{theorem}
\begin{theorem}[\texttt{wickPolynomial\_add\_sub\_self}]\label{Pphi2-wickPolynomial-add-sub-self}
  \lean{Pphi2.wickPolynomial_add_sub_self}\leanok
  \uses{Pphi2-wickMonomial, Pphi2-wickMonomial-add-sub-self, Pphi2-wickPolynomial}
  \texttt{Pphi2.wickPolynomial\_add\_sub\_self}
\end{theorem}
\begin{theorem}[\texttt{wickMonomial\_four\_lower\_bound}]\label{Pphi2-wickMonomial-four-lower-bound}
  \lean{Pphi2.wickMonomial_four_lower_bound}\leanok
  \uses{Pphi2-wickMonomial, Pphi2-wickMonomial-four, Lattice-toSemilatticeInf}
  $:x^4:_c \ge -6c^2$ for all $x$ and $c \ge 0$. Proof: $x^4 - 6cx^2 + 3c^2 = (x^2 - 3c)^2 - 6c^2 \ge -6c^2$.
\end{theorem}
\section{\texttt{Pphi2.InteractingMeasure.U4DerivativeClosedForm}}
\begin{theorem}[\texttt{normalizedMoment\_hasDerivAt\_explicit}]\label{Pphi2-normalizedMoment-hasDerivAt-explicit}
  \lean{Pphi2.normalizedMoment_hasDerivAt_explicit}\leanok
  \uses{Pphi2-partitionFnDeriv, Pphi2-normalizedMoment-hasDerivAt, GaussianField-FinLatticeSites, Pphi2-gibbsMoment, GaussianField-latticeGaussianMeasure, GaussianField-Configuration, Pphi2-normalizedMoment, GaussianField-instMeasurableSpaceConfiguration, Pphi2-partitionFn, Pphi2-normalizedMomentDeriv, Pphi2-momentDeriv, GaussianField-FinLatticeField, Pphi2-interactionFunctional}
  \texttt{Pphi2.normalizedMoment\_hasDerivAt\_explicit}
\end{theorem}
\begin{theorem}[\texttt{u4\_hasDerivAt2}]\label{Pphi2-u4-hasDerivAt2}
  \lean{Pphi2.u4_hasDerivAt2}\leanok
  \uses{Pphi2-u4Deriv2, Pphi2-normalizedMomentDeriv2, Pphi2-normalizedMoment-hasDerivAt-explicit, Pphi2-normalizedMomentDeriv-hasDerivAt, Pphi2-normalizedMoment, Pphi2-normalizedMomentDeriv, GaussianField-FinLatticeField, Pphi2-u4Deriv}
  \texttt{Pphi2.u4\_hasDerivAt2}
\end{theorem}
\begin{definition}[\texttt{partitionFnDeriv}]\label{Pphi2-partitionFnDeriv}
  \lean{Pphi2.partitionFnDeriv}\leanok
  \uses{GaussianField-FinLatticeSites, GaussianField-latticeGaussianMeasure, GaussianField-Configuration, GaussianField-instMeasurableSpaceConfiguration, GaussianField-FinLatticeField, Pphi2-interactionFunctional}
  \texttt{Pphi2.partitionFnDeriv}
\end{definition}
\begin{definition}[\texttt{u4Deriv}]\label{Pphi2-u4Deriv}
  \lean{Pphi2.u4Deriv}\leanok
  \uses{Pphi2-normalizedMoment, Pphi2-normalizedMomentDeriv, GaussianField-FinLatticeField}
  \texttt{Pphi2.u4Deriv}
\end{definition}
\begin{theorem}[\texttt{normalizedMomentDeriv\_hasDerivAt}]\label{Pphi2-normalizedMomentDeriv-hasDerivAt}
  \lean{Pphi2.normalizedMomentDeriv_hasDerivAt}\leanok
  \uses{Pphi2-partitionFn-ge-one, Pphi2-partitionFn-hasDerivAt, Pphi2-partitionFnDeriv, Pphi2-partitionFn-hasDerivAt2, GaussianField-FinLatticeSites, Pphi2-partitionFnDeriv2, Pphi2-gibbsMoment, Pphi2-moment-hasDerivAt, Pphi2-normalizedMomentDeriv2, GaussianField-latticeGaussianMeasure, GaussianField-Configuration, GaussianField-instMeasurableSpaceConfiguration, Pphi2-moment-hasDerivAt2, Pphi2-momentDeriv2, Pphi2-partitionFn, Pphi2-normalizedMomentDeriv, Pphi2-momentDeriv, GaussianField-FinLatticeField, Pphi2-interactionFunctional}
  \texttt{Pphi2.normalizedMomentDeriv\_hasDerivAt}
\end{theorem}
\begin{theorem}[\texttt{u4\_hasDerivAt}]\label{Pphi2-u4-hasDerivAt}
  \lean{Pphi2.u4_hasDerivAt}\leanok
  \uses{Pphi2-normalizedMoment-hasDerivAt-explicit, Pphi2-normalizedMoment, Pphi2-normalizedMomentDeriv, Pphi2-u4, GaussianField-FinLatticeField, Pphi2-u4Deriv}
  \texttt{Pphi2.u4\_hasDerivAt}
\end{theorem}
\begin{theorem}[\texttt{congr\_simp}]\label{Pphi2-normalizedMomentDeriv2-congr-simp}
  \lean{Pphi2.normalizedMomentDeriv2.congr_simp}\leanok
  \uses{Pphi2-normalizedMomentDeriv2, GaussianField-FinLatticeField}
  \texttt{Pphi2.normalizedMomentDeriv2.congr\_simp}
\end{theorem}
\begin{definition}[\texttt{partitionFnDeriv2}]\label{Pphi2-partitionFnDeriv2}
  \lean{Pphi2.partitionFnDeriv2}\leanok
  \uses{GaussianField-FinLatticeSites, GaussianField-latticeGaussianMeasure, GaussianField-Configuration, GaussianField-instMeasurableSpaceConfiguration, GaussianField-FinLatticeField, Pphi2-interactionFunctional}
  \texttt{Pphi2.partitionFnDeriv2}
\end{definition}
\begin{definition}[\texttt{partitionFn}]\label{Pphi2-partitionFn}
  \lean{Pphi2.partitionFn}\leanok
  \uses{GaussianField-FinLatticeSites, GaussianField-latticeGaussianMeasure, GaussianField-Configuration, GaussianField-instMeasurableSpaceConfiguration, GaussianField-FinLatticeField, Pphi2-interactionFunctional}
  \texttt{Pphi2.partitionFn}
\end{definition}
\begin{definition}[\texttt{normalizedMomentDeriv}]\label{Pphi2-normalizedMomentDeriv}
  \lean{Pphi2.normalizedMomentDeriv}\leanok
  \uses{Pphi2-partitionFnDeriv, Pphi2-gibbsMoment, Pphi2-partitionFn, Pphi2-momentDeriv, GaussianField-FinLatticeField}
  \texttt{Pphi2.normalizedMomentDeriv}
\end{definition}
\begin{definition}[\texttt{u4}]\label{Pphi2-u4}
  \lean{Pphi2.u4}\leanok
  \uses{Pphi2-normalizedMoment, GaussianField-FinLatticeField}
  \texttt{Pphi2.u4}
\end{definition}
\begin{theorem}[\texttt{congr\_simp}]\label{Pphi2-normalizedMoment-congr-simp}
  \lean{Pphi2.normalizedMoment.congr_simp}\leanok
  \uses{Pphi2-normalizedMoment, GaussianField-FinLatticeField}
  \texttt{Pphi2.normalizedMoment.congr\_simp}
\end{theorem}
\begin{definition}[\texttt{normalizedMomentDeriv2}]\label{Pphi2-normalizedMomentDeriv2}
  \lean{Pphi2.normalizedMomentDeriv2}\leanok
  \uses{Pphi2-partitionFnDeriv, Pphi2-partitionFnDeriv2, Pphi2-gibbsMoment, Pphi2-momentDeriv2, Pphi2-partitionFn, Pphi2-momentDeriv, GaussianField-FinLatticeField}
  \texttt{Pphi2.normalizedMomentDeriv2}
\end{definition}
\begin{theorem}[\texttt{congr\_simp}]\label{Pphi2-normalizedMomentDeriv-congr-simp}
  \lean{Pphi2.normalizedMomentDeriv.congr_simp}\leanok
  \uses{Pphi2-normalizedMomentDeriv, GaussianField-FinLatticeField}
  \texttt{Pphi2.normalizedMomentDeriv.congr\_simp}
\end{theorem}
\begin{definition}[\texttt{momentDeriv2}]\label{Pphi2-momentDeriv2}
  \lean{Pphi2.momentDeriv2}\leanok
  \uses{GaussianField-FinLatticeSites, GaussianField-latticeGaussianMeasure, GaussianField-Configuration, GaussianField-instMeasurableSpaceConfiguration, GaussianField-FinLatticeField, Pphi2-interactionFunctional}
  \texttt{Pphi2.momentDeriv2}
\end{definition}
\begin{definition}[\texttt{momentDeriv}]\label{Pphi2-momentDeriv}
  \lean{Pphi2.momentDeriv}\leanok
  \uses{GaussianField-FinLatticeSites, GaussianField-latticeGaussianMeasure, GaussianField-Configuration, GaussianField-instMeasurableSpaceConfiguration, GaussianField-FinLatticeField, Pphi2-interactionFunctional}
  \texttt{Pphi2.momentDeriv}
\end{definition}
\begin{theorem}[\texttt{normalizedMoment\_hasDerivAt}]\label{Pphi2-normalizedMoment-hasDerivAt}
  \lean{Pphi2.normalizedMoment_hasDerivAt}\leanok
  \uses{Pphi2-partitionFn-ge-one, Pphi2-partitionFn-hasDerivAt, GaussianField-FinLatticeSites, Pphi2-moment-hasDerivAt, GaussianField-latticeGaussianMeasure, GaussianField-Configuration, GaussianField-instMeasurableSpaceConfiguration, GaussianField-FinLatticeField, Pphi2-interactionFunctional}
  \texttt{Pphi2.normalizedMoment\_hasDerivAt}
\end{theorem}
\begin{definition}[\texttt{gibbsMoment}]\label{Pphi2-gibbsMoment}
  \lean{Pphi2.gibbsMoment}\leanok
  \uses{GaussianField-FinLatticeSites, GaussianField-latticeGaussianMeasure, GaussianField-Configuration, GaussianField-instMeasurableSpaceConfiguration, GaussianField-FinLatticeField, Pphi2-interactionFunctional}
  \texttt{Pphi2.gibbsMoment}
\end{definition}
\begin{definition}[\texttt{normalizedMoment}]\label{Pphi2-normalizedMoment}
  \lean{Pphi2.normalizedMoment}\leanok
  \uses{Pphi2-gibbsMoment, Pphi2-partitionFn, GaussianField-FinLatticeField}
  \texttt{Pphi2.normalizedMoment}
\end{definition}
\begin{definition}[\texttt{u4Deriv2}]\label{Pphi2-u4Deriv2}
  \lean{Pphi2.u4Deriv2}\leanok
  \uses{Pphi2-normalizedMomentDeriv2, Pphi2-normalizedMoment, Pphi2-normalizedMomentDeriv, GaussianField-FinLatticeField}
  \texttt{Pphi2.u4Deriv2}
\end{definition}
\begin{theorem}[\texttt{partitionFn\_hasDerivAt2}]\label{Pphi2-partitionFn-hasDerivAt2}
  \lean{Pphi2.partitionFn_hasDerivAt2}\leanok
  \uses{GaussianField-FinLatticeSites, GaussianField-latticeGaussianMeasure, GaussianField-Configuration, GaussianField-instMeasurableSpaceConfiguration, Pphi2-moment-hasDerivAt2, GaussianField-FinLatticeField, Pphi2-interactionFunctional}
  \texttt{Pphi2.partitionFn\_hasDerivAt2}
\end{theorem}
\section{\texttt{Pphi2.AsymTorus.AsymPositivity}}
\begin{theorem}[\texttt{hlt}]\label{Pphi2-AsymTransferGroundExcitedData-hlt}
  \lean{Pphi2.AsymTransferGroundExcitedData.hlt}\leanok
  \uses{Pphi2-AsymTransferGroundExcitedData-eigenval, Pphi2-AsymTransferGroundExcitedData-i-, Pphi2-AsymTransferGroundExcitedData-i-, Pphi2-AsymTransferGroundExcitedData}
  \texttt{Pphi2.AsymTransferGroundExcitedData.hlt}
\end{theorem}
\begin{definition}[\texttt{eigenval}]\label{Pphi2-AsymTransferGroundExcitedData-eigenval}
  \lean{Pphi2.AsymTransferGroundExcitedData.eigenval}\leanok
  \uses{Pphi2-AsymTransferGroundExcitedData, Pphi2-AsymTransferGroundExcitedData--}
  \texttt{Pphi2.AsymTransferGroundExcitedData.eigenval}
\end{definition}
\begin{definition}[\texttt{asymFirstExcitedEnergyLevel}]\label{Pphi2-asymFirstExcitedEnergyLevel}
  \lean{Pphi2.asymFirstExcitedEnergyLevel}\leanok
  \uses{Pphi2-asymTransferFirstExcitedEigenvalue}
  \texttt{Pphi2.asymFirstExcitedEnergyLevel}
\end{definition}
\begin{theorem}[\texttt{asymSpectral\_gap\_pos}]\label{Pphi2-asymSpectral-gap-pos}
  \lean{Pphi2.asymSpectral_gap_pos}\leanok
  \uses{Pphi2-asymMassGap-pos, Pphi2-asymMassGap}
  \texttt{Pphi2.asymSpectral\_gap\_pos}
\end{theorem}
\begin{definition}[\texttt{ctorIdx}]\label{Pphi2-AsymTransferGroundExcitedData-ctorIdx}
  \lean{Pphi2.AsymTransferGroundExcitedData.ctorIdx}\leanok
  \uses{Pphi2-AsymTransferGroundExcitedData}
  \texttt{Pphi2.AsymTransferGroundExcitedData.ctorIdx}
\end{definition}
\begin{definition}[\texttt{b}]\label{Pphi2-AsymTransferGroundExcitedData-b}
  \lean{Pphi2.AsymTransferGroundExcitedData.b}\leanok
  \uses{Pphi2-L2SpatialField, Pphi2-SpatialField, Pphi2-AsymTransferGroundExcitedData, Pphi2-AsymTransferGroundExcitedData--}
  \texttt{Pphi2.AsymTransferGroundExcitedData.b}
\end{definition}
\begin{definition}[\texttt{AsymTransferGroundExcitedData}]\label{Pphi2-AsymTransferGroundExcitedData}
  \lean{Pphi2.AsymTransferGroundExcitedData}\leanok
  \texttt{Pphi2.AsymTransferGroundExcitedData}
\end{definition}
\begin{definition}[\texttt{asymGroundEnergyLevel}]\label{Pphi2-asymGroundEnergyLevel}
  \lean{Pphi2.asymGroundEnergyLevel}\leanok
  \uses{Pphi2-asymTransferGroundEigenvalue}
  \texttt{Pphi2.asymGroundEnergyLevel}
\end{definition}
\begin{theorem}[\texttt{hi\_ne}]\label{Pphi2-AsymTransferGroundExcitedData-hi-ne}
  \lean{Pphi2.AsymTransferGroundExcitedData.hi_ne}\leanok
  \uses{Pphi2-AsymTransferGroundExcitedData-i-, Pphi2-AsymTransferGroundExcitedData-i-, Pphi2-AsymTransferGroundExcitedData, Pphi2-AsymTransferGroundExcitedData--}
  \texttt{Pphi2.AsymTransferGroundExcitedData.hi\_ne}
\end{theorem}
\begin{definition}[\texttt{ι}]\label{Pphi2-AsymTransferGroundExcitedData--}
  \lean{Pphi2.AsymTransferGroundExcitedData.ι}\leanok
  \uses{Pphi2-AsymTransferGroundExcitedData}
  \texttt{Pphi2.AsymTransferGroundExcitedData.ι}
\end{definition}
\begin{definition}[\texttt{asymTransferGroundEigenvalue}]\label{Pphi2-asymTransferGroundEigenvalue}
  \lean{Pphi2.asymTransferGroundEigenvalue}\leanok
  \uses{Pphi2-AsymTransferGroundExcitedData-eigenval, Pphi2-AsymTransferGroundExcitedData-i-, Pphi2-asymTransferGroundExcitedData, Pphi2-AsymTransferGroundExcitedData}
  \texttt{Pphi2.asymTransferGroundEigenvalue}
\end{definition}
\begin{theorem}[\texttt{asymMassGap\_pos}]\label{Pphi2-asymMassGap-pos}
  \lean{Pphi2.asymMassGap_pos}\leanok
  \uses{Pphi2-asymFirstExcitedEnergyLevel, Pphi2-asymGroundEnergyLevel, Pphi2-asymMassGap, Pphi2-asymEnergyLevel-gap}
  \texttt{Pphi2.asymMassGap\_pos}
\end{theorem}
\begin{definition}[\texttt{asymMassGap}]\label{Pphi2-asymMassGap}
  \lean{Pphi2.asymMassGap}\leanok
  \uses{Pphi2-asymFirstExcitedEnergyLevel, Pphi2-asymGroundEnergyLevel}
  \texttt{Pphi2.asymMassGap}
\end{definition}
\begin{theorem}[\texttt{asymEnergyLevel\_gap}]\label{Pphi2-asymEnergyLevel-gap}
  \lean{Pphi2.asymEnergyLevel_gap}\leanok
  \uses{Pphi2-asymFirstExcitedEnergyLevel, Pphi2-AsymTransferGroundExcitedData-eigenval, Pphi2-asymTransferOperator-eigenvalues-pos, Pphi2-AsymTransferGroundExcitedData-i-, Pphi2-asymGroundEnergyLevel, Pphi2-asymTransferGroundExcitedData, Pphi2-AsymTransferGroundExcitedData-i-, Pphi2-AsymTransferGroundExcitedData, Pphi2-asymTransferGroundEigenvalue, Pphi2-AsymTransferGroundExcitedData-b, Pphi2-AsymTransferGroundExcitedData--, Pphi2-AsymTransferGroundExcitedData-hlt, Pphi2-AsymTransferGroundExcitedData-h-eigen, Pphi2-asymTransferFirstExcitedEigenvalue}
  \texttt{Pphi2.asymEnergyLevel\_gap}
\end{theorem}
\begin{theorem}[\texttt{h\_sum}]\label{Pphi2-AsymTransferGroundExcitedData-h-sum}
  \lean{Pphi2.AsymTransferGroundExcitedData.h_sum}\leanok
  \uses{Pphi2-L2SpatialField, Pphi2-AsymTransferGroundExcitedData-eigenval, Pphi2-SpatialField, Pphi2-AsymTransferGroundExcitedData, Pphi2-AsymTransferGroundExcitedData-b, Pphi2-AsymTransferGroundExcitedData--, Pphi2-asymTransferOperatorCLM}
  \texttt{Pphi2.AsymTransferGroundExcitedData.h\_sum}
\end{theorem}
\begin{theorem}[\texttt{h\_eigen}]\label{Pphi2-AsymTransferGroundExcitedData-h-eigen}
  \lean{Pphi2.AsymTransferGroundExcitedData.h_eigen}\leanok
  \uses{Pphi2-L2SpatialField, Pphi2-AsymTransferGroundExcitedData-eigenval, Pphi2-SpatialField, Pphi2-AsymTransferGroundExcitedData, Pphi2-AsymTransferGroundExcitedData-b, Pphi2-AsymTransferGroundExcitedData--, Pphi2-asymTransferOperatorCLM}
  \texttt{Pphi2.AsymTransferGroundExcitedData.h\_eigen}
\end{theorem}
\begin{definition}[\texttt{asymTransferFirstExcitedEigenvalue}]\label{Pphi2-asymTransferFirstExcitedEigenvalue}
  \lean{Pphi2.asymTransferFirstExcitedEigenvalue}\leanok
  \uses{Pphi2-AsymTransferGroundExcitedData-eigenval, Pphi2-asymTransferGroundExcitedData, Pphi2-AsymTransferGroundExcitedData-i-, Pphi2-AsymTransferGroundExcitedData}
  \texttt{Pphi2.asymTransferFirstExcitedEigenvalue}
\end{definition}
\begin{definition}[\texttt{i₁}]\label{Pphi2-AsymTransferGroundExcitedData-i-}
  \lean{Pphi2.AsymTransferGroundExcitedData.i₁}\leanok
  \uses{Pphi2-AsymTransferGroundExcitedData, Pphi2-AsymTransferGroundExcitedData--}
  \texttt{Pphi2.AsymTransferGroundExcitedData.i₁}
\end{definition}
\begin{definition}[\texttt{asymTransferGroundExcitedData}]\label{Pphi2-asymTransferGroundExcitedData}
  \lean{Pphi2.asymTransferGroundExcitedData}\leanok
  \uses{Pphi2-AsymTransferGroundExcitedData}
  \texttt{Pphi2.asymTransferGroundExcitedData}
\end{definition}
\begin{definition}[\texttt{i₀}]\label{Pphi2-AsymTransferGroundExcitedData-i-}
  \lean{Pphi2.AsymTransferGroundExcitedData.i₀}\leanok
  \uses{Pphi2-AsymTransferGroundExcitedData, Pphi2-AsymTransferGroundExcitedData--}
  \texttt{Pphi2.AsymTransferGroundExcitedData.i₀}
\end{definition}
\begin{theorem}[\texttt{htop}]\label{Pphi2-AsymTransferGroundExcitedData-htop}
  \lean{Pphi2.AsymTransferGroundExcitedData.htop}\leanok
  \uses{Pphi2-AsymTransferGroundExcitedData-eigenval, Pphi2-AsymTransferGroundExcitedData-i-, Pphi2-AsymTransferGroundExcitedData, Pphi2-AsymTransferGroundExcitedData--}
  \texttt{Pphi2.AsymTransferGroundExcitedData.htop}
\end{theorem}
\section{\texttt{Pphi2.NelsonEstimate.LayerCake}}
\begin{theorem}[\texttt{setOf\_le\_max\_eq\_setOf\_le\_neg}]\label{Pphi2-LayerCake-setOf-le-max-eq-setOf-le-neg}
  \lean{Pphi2.LayerCake.setOf_le_max_eq_setOf_le_neg}\leanok
  \uses{Lattice-toSemilatticeSup, Lattice-toSemilatticeInf}
  \texttt{Pphi2.LayerCake.setOf\_le\_max\_eq\_setOf\_le\_neg}
\end{theorem}
\begin{theorem}[\texttt{lintegral\_expSq\_neg\_le\_layer\_cake}]\label{Pphi2-LayerCake-lintegral-expSq-neg-le-layer-cake}
  \lean{Pphi2.LayerCake.lintegral_expSq_neg_le_layer_cake}\leanok
  \uses{Pphi2-LayerCake-exp-neg-sq-le-exp-two-max, Pphi2-LayerCake-setOf-le-max-eq-setOf-le-neg, Lattice-toSemilatticeInf}
  \texttt{Pphi2.LayerCake.lintegral\_expSq\_neg\_le\_layer\_cake}
\end{theorem}
\begin{theorem}[\texttt{exp\_neg\_sq\_le\_exp\_two\_max}]\label{Pphi2-LayerCake-exp-neg-sq-le-exp-two-max}
  \lean{Pphi2.LayerCake.exp_neg_sq_le_exp_two_max}\leanok
  \texttt{Pphi2.LayerCake.exp\_neg\_sq\_le\_exp\_two\_max}
\end{theorem}
\begin{theorem}[\texttt{lintegral\_expSq\_neg\_le\_of\_tail}]\label{Pphi2-LayerCake-lintegral-expSq-neg-le-of-tail}
  \lean{Pphi2.LayerCake.lintegral_expSq_neg_le_of_tail}\leanok
  \uses{Pphi2-LayerCake-lintegral-expSq-neg-le-layer-cake, Lattice-toSemilatticeInf}
  \texttt{Pphi2.LayerCake.lintegral\_expSq\_neg\_le\_of\_tail}
\end{theorem}
\section{\texttt{Pphi2.InteractingMeasure.LeadingTerm}}
\begin{theorem}[\texttt{leadingTerm\_const\_eq}]\label{Pphi2-leadingTerm-const-eq}
  \lean{Pphi2.leadingTerm_const_eq}\leanok
  \uses{GaussianField-FinLatticeSites, GaussianField-gffPositionCovariance, Pphi2-gffPositionCovariance-row-sum}
  \texttt{Pphi2.leadingTerm\_const\_eq}
\end{theorem}
\begin{theorem}[\texttt{covarianceGJ\_massOperator}]\label{Pphi2-covarianceGJ-massOperator}
  \lean{Pphi2.covarianceGJ_massOperator}\leanok
  \uses{GaussianField-latticeCovariance-GJ-eq-inv-smul-bare, GaussianField-covariance, GaussianField-latticeCovariance, GaussianField-FinLatticeSites, GaussianField-latticeCovarianceGJ, GaussianField-massOperator, GaussianField-ell2-, GaussianField-FinLatticeField, Pphi2-spectralCov-massOperator-inner}
  \texttt{Pphi2.covarianceGJ\_massOperator}
\end{theorem}
\begin{theorem}[\texttt{massOperator\_const}]\label{Pphi2-massOperator-const}
  \lean{Pphi2.massOperator_const}\leanok
  \uses{GaussianField-FinLatticeSites, GaussianField-massOperator, Pphi2-finiteLaplacian-const, GaussianField-finiteLaplacian, GaussianField-FinLatticeField}
  \texttt{Pphi2.massOperator\_const}
\end{theorem}
\begin{theorem}[\texttt{gffPositionCovariance\_row\_sum}]\label{Pphi2-gffPositionCovariance-row-sum}
  \lean{Pphi2.gffPositionCovariance_row_sum}\leanok
  \uses{Pphi2-covarianceGJ-massOperator, GaussianField-covariance, GaussianField-FinLatticeSites, GaussianField-latticeCovarianceGJ, GaussianField-massOperator, Pphi2-massOperator-const, GaussianField-gffPositionCovariance, GaussianField-gffPositionCovariance-eq-covarianceGJ, GaussianField-ell2-, GaussianField-FinLatticeField}
  \texttt{Pphi2.gffPositionCovariance\_row\_sum}
\end{theorem}
\begin{theorem}[\texttt{leadingTerm\_const\_pos}]\label{Pphi2-leadingTerm-const-pos}
  \lean{Pphi2.leadingTerm_const_pos}\leanok
  \uses{Pphi2-leadingTerm-const-eq, GaussianField-FinLatticeSites, GaussianField-gffPositionCovariance}
  \texttt{Pphi2.leadingTerm\_const\_pos}
\end{theorem}
\begin{theorem}[\texttt{spectralCov\_massOperator\_inner}]\label{Pphi2-spectralCov-massOperator-inner}
  \lean{Pphi2.spectralCov_massOperator_inner}\leanok
  \uses{GaussianField-FinLatticeSites, GaussianField-eigenbasis-completeness, GaussianField-massOperator, GaussianField-spectralLatticeCovariance-inner, GaussianField-massEigenvalues, GaussianField-spectralLatticeCovariance, GaussianField-massEigenvectorBasis, GaussianField-massOperatorMatrix-eigenvalues-pos, GaussianField-ell2-, GaussianField-FinLatticeField, GaussianField-massOperator-eigenCoeff-eq-eigenvalues-mul-eigenCoeff}
  \texttt{Pphi2.spectralCov\_massOperator\_inner}
\end{theorem}
\begin{theorem}[\texttt{finiteLaplacian\_const}]\label{Pphi2-finiteLaplacian-const}
  \lean{Pphi2.finiteLaplacian_const}\leanok
  \uses{GaussianField-FinLatticeSites, GaussianField-finiteLaplacianFun, GaussianField-finiteLaplacian, GaussianField-FinLatticeField}
  \texttt{Pphi2.finiteLaplacian\_const}
\end{theorem}
\section{\texttt{Pphi2.AsymTorus.AsymSpectralGap}}
\begin{theorem}[\texttt{asymTransferNormalized\_gap}]\label{Pphi2-asymTransferNormalized-gap}
  \lean{Pphi2.asymTransferNormalized_gap}\leanok
  \uses{Pphi2-AsymTransferGroundExcitedData-htop, Pphi2-asymTransferNormalized, Pphi2-L2SpatialField, Pphi2-AsymTransferGroundExcitedData-eigenval, Pphi2-SpatialField, Pphi2-AsymTransferGroundExcitedData-i-, Pphi2-asymGroundVector, GaussianField-compact-selfAdjoint-hasEigenvector, Pphi2-asymTransferGroundExcitedData, Pphi2-asymTransferOperator-isCompact, Pphi2-asymTransferGroundEigenvalue-pos, Pphi2-AsymTransferGroundExcitedData, Pphi2-asymTransferGroundEigenvalue, Pphi2-asymTransferOperator-isSelfAdjoint, Pphi2-AsymTransferGroundExcitedData-b, Pphi2-AsymTransferGroundExcitedData--, Pphi2-AsymTransferGroundExcitedData-h-eigen, Pphi2-asymTransferOperatorCLM}
  \texttt{Pphi2.asymTransferNormalized\_gap}
\end{theorem}
\begin{definition}[\texttt{asymGappedTransfer'}]\label{Pphi2-asymGappedTransfer-}
  \lean{Pphi2.asymGappedTransfer'}\leanok
  \uses{Pphi2-asymTransferNormalized, Pphi2-L2SpatialField, Pphi2-SpatialField, Pphi2-asymGroundVector, Pphi2-asymGappedTransfer, Pphi2-asymTransferNormalized-gap}
  \texttt{Pphi2.asymGappedTransfer'}
\end{definition}
\section{\texttt{Pphi2.NelsonEstimate.DynamicalCutoff}}
\begin{theorem}[\texttt{smooth\_lower\_bound\_at\_cutoff}]\label{Pphi2-DynamicalCutoff-smooth-lower-bound-at-cutoff}
  \lean{Pphi2.DynamicalCutoff.smooth_lower_bound_at_cutoff}\leanok
  \uses{Pphi2-smoothWickConstant, Pphi2-DynamicalCutoff-dynamicalCutoffScale-log-sq-le, GaussianField-FinLatticeSites, Pphi2-smooth-interaction-lower-bound-log-uniform, Pphi2-DynamicalCutoff-dynamicalCutoffScale-pos, Pphi2-wickMonomial, Pphi2-smoothBoundConstant, Pphi2-DynamicalCutoff-dynamicalCutoffScale, Pphi2-smoothBoundConstant-pos}
  \texttt{Pphi2.DynamicalCutoff.smooth\_lower\_bound\_at\_cutoff}
\end{theorem}
\begin{theorem}[\texttt{one\_add\_abs\_log\_degreeDynamicalCutoff\_eq\_rpow}]\label{Pphi2-DynamicalCutoff-one-add-abs-log-degreeDynamicalCutoff-eq-rpow}
  \lean{Pphi2.DynamicalCutoff.one_add_abs_log_degreeDynamicalCutoff_eq_rpow}\leanok
  \uses{Pphi2-DynamicalCutoff-degreeDynamicalCutoffScale-eq-exp, Pphi2-DynamicalCutoff-degreeDynamicalCutoffScale, Pphi2-DynamicalCutoff-degreeCutoffPower}
  \texttt{Pphi2.DynamicalCutoff.one\_add\_abs\_log\_degreeDynamicalCutoff\_eq\_rpow}
\end{theorem}
\begin{theorem}[\texttt{degreeCutoffPower\_pos}]\label{Pphi2-DynamicalCutoff-degreeCutoffPower-pos}
  \lean{Pphi2.DynamicalCutoff.degreeCutoffPower_pos}\leanok
  \uses{Pphi2-DynamicalCutoff-degreeCutoffPower}
  \texttt{Pphi2.DynamicalCutoff.degreeCutoffPower\_pos}
\end{theorem}
\begin{theorem}[\texttt{degreeDynamicalCutoffScale\_eq\_one}]\label{Pphi2-DynamicalCutoff-degreeDynamicalCutoffScale-eq-one}
  \lean{Pphi2.DynamicalCutoff.degreeDynamicalCutoffScale_eq_one}\leanok
  \uses{Pphi2-DynamicalCutoff-degreeDynamicalCutoffScale, Pphi2-DynamicalCutoff-degreeCutoffPower}
  \texttt{Pphi2.DynamicalCutoff.degreeDynamicalCutoffScale\_eq\_one}
\end{theorem}
\begin{theorem}[\texttt{dynamicalCutoffScale\_log\_sq\_le}]\label{Pphi2-DynamicalCutoff-dynamicalCutoffScale-log-sq-le}
  \lean{Pphi2.DynamicalCutoff.dynamicalCutoffScale_log_sq_le}\leanok
  \uses{Pphi2-DynamicalCutoff-dynamicalCutoffScale-eq-exp, Pphi2-DynamicalCutoff-dynamicalCutoffScale, Pphi2-DynamicalCutoff-dynamicalCutoffScale-eq-one}
  \texttt{Pphi2.DynamicalCutoff.dynamicalCutoffScale\_log\_sq\_le}
\end{theorem}
\begin{theorem}[\texttt{measure\_le\_neg\_of\_smooth\_rough\_split}]\label{Pphi2-DynamicalCutoff-measure-le-neg-of-smooth-rough-split}
  \lean{Pphi2.DynamicalCutoff.measure_le_neg_of_smooth_rough_split}\leanok
  \texttt{Pphi2.DynamicalCutoff.measure\_le\_neg\_of\_smooth\_rough\_split}
\end{theorem}
\begin{theorem}[\texttt{degreeDynamicalCutoffScale\_pos}]\label{Pphi2-DynamicalCutoff-degreeDynamicalCutoffScale-pos}
  \lean{Pphi2.DynamicalCutoff.degreeDynamicalCutoffScale_pos}\leanok
  \uses{Pphi2-DynamicalCutoff-degreeDynamicalCutoffScale, Pphi2-DynamicalCutoff-degreeCutoffPower}
  \texttt{Pphi2.DynamicalCutoff.degreeDynamicalCutoffScale\_pos}
\end{theorem}
\begin{theorem}[\texttt{congr\_simp}]\label{Pphi2-smoothWickConstant-congr-simp}
  \lean{Pphi2.smoothWickConstant.congr_simp}\leanok
  \uses{Pphi2-smoothWickConstant}
  \texttt{Pphi2.smoothWickConstant.congr\_simp}
\end{theorem}
\begin{theorem}[\texttt{dynamicalCutoffScale\_pos}]\label{Pphi2-DynamicalCutoff-dynamicalCutoffScale-pos}
  \lean{Pphi2.DynamicalCutoff.dynamicalCutoffScale_pos}\leanok
  \uses{Pphi2-DynamicalCutoff-dynamicalCutoffScale}
  \texttt{Pphi2.DynamicalCutoff.dynamicalCutoffScale\_pos}
\end{theorem}
\begin{theorem}[\texttt{degreeDynamicalCutoffScale\_eq\_exp}]\label{Pphi2-DynamicalCutoff-degreeDynamicalCutoffScale-eq-exp}
  \lean{Pphi2.DynamicalCutoff.degreeDynamicalCutoffScale_eq_exp}\leanok
  \uses{Pphi2-DynamicalCutoff-degreeDynamicalCutoffScale, Pphi2-DynamicalCutoff-degreeCutoffPower}
  \texttt{Pphi2.DynamicalCutoff.degreeDynamicalCutoffScale\_eq\_exp}
\end{theorem}
\begin{definition}[\texttt{dynamicalCutoffScale}]\label{Pphi2-DynamicalCutoff-dynamicalCutoffScale}
  \lean{Pphi2.DynamicalCutoff.dynamicalCutoffScale}\leanok
  \texttt{Pphi2.DynamicalCutoff.dynamicalCutoffScale}
\end{definition}
\begin{definition}[\texttt{degreeCutoffPower}]\label{Pphi2-DynamicalCutoff-degreeCutoffPower}
  \lean{Pphi2.DynamicalCutoff.degreeCutoffPower}\leanok
  \texttt{Pphi2.DynamicalCutoff.degreeCutoffPower}
\end{definition}
\begin{theorem}[\texttt{degreeDynamicalCutoffScale\_log\_pow\_le}]\label{Pphi2-DynamicalCutoff-degreeDynamicalCutoffScale-log-pow-le}
  \lean{Pphi2.DynamicalCutoff.degreeDynamicalCutoffScale_log_pow_le}\leanok
  \uses{Pphi2-DynamicalCutoff-degreeDynamicalCutoffScale-eq-one, Pphi2-DynamicalCutoff-one-add-abs-log-degreeDynamicalCutoff-eq-rpow, Pphi2-DynamicalCutoff-degreeDynamicalCutoffScale, Pphi2-DynamicalCutoff-degreeCutoffPower}
  \texttt{Pphi2.DynamicalCutoff.degreeDynamicalCutoffScale\_log\_pow\_le}
\end{theorem}
\begin{definition}[\texttt{degreeDynamicalCutoffScale}]\label{Pphi2-DynamicalCutoff-degreeDynamicalCutoffScale}
  \lean{Pphi2.DynamicalCutoff.degreeDynamicalCutoffScale}\leanok
  \uses{Pphi2-DynamicalCutoff-degreeCutoffPower}
  \texttt{Pphi2.DynamicalCutoff.degreeDynamicalCutoffScale}
\end{definition}
\begin{theorem}[\texttt{dynamicalCutoffScale\_eq\_exp}]\label{Pphi2-DynamicalCutoff-dynamicalCutoffScale-eq-exp}
  \lean{Pphi2.DynamicalCutoff.dynamicalCutoffScale_eq_exp}\leanok
  \uses{Pphi2-DynamicalCutoff-dynamicalCutoffScale}
  \texttt{Pphi2.DynamicalCutoff.dynamicalCutoffScale\_eq\_exp}
\end{theorem}
\begin{theorem}[\texttt{dynamicalCutoffScale\_eq\_one}]\label{Pphi2-DynamicalCutoff-dynamicalCutoffScale-eq-one}
  \lean{Pphi2.DynamicalCutoff.dynamicalCutoffScale_eq_one}\leanok
  \uses{Pphi2-DynamicalCutoff-dynamicalCutoffScale}
  \texttt{Pphi2.DynamicalCutoff.dynamicalCutoffScale\_eq\_one}
\end{theorem}
\begin{theorem}[\texttt{dynamicalCutoffScale\_le\_one}]\label{Pphi2-DynamicalCutoff-dynamicalCutoffScale-le-one}
  \lean{Pphi2.DynamicalCutoff.dynamicalCutoffScale_le_one}\leanok
  \uses{Pphi2-DynamicalCutoff-dynamicalCutoffScale-eq-exp, Pphi2-DynamicalCutoff-dynamicalCutoffScale, Pphi2-DynamicalCutoff-dynamicalCutoffScale-eq-one}
  \texttt{Pphi2.DynamicalCutoff.dynamicalCutoffScale\_le\_one}
\end{theorem}
\section{\texttt{Pphi2.TorusContinuumLimit.TorusInteractingLimit}}
\begin{theorem}[\texttt{torusInteractingMeasure\_isProbability}]\label{Pphi2-torusInteractingMeasure-isProbability}
  \lean{Pphi2.torusInteractingMeasure_isProbability}\leanok
  \uses{GaussianField-NuclearTensorProduct-instModuleReal, GaussianField-NuclearTensorProduct-instTopologicalSpace, GaussianField-FinLatticeSites, GaussianField-NuclearTensorProduct-instIsTopologicalAddGroup, Pphi2-torusEmbedLift, Pphi2-interactingLatticeMeasure-isProbability, Pphi2-torusEmbedLift-measurable, Pphi2-interactingLatticeMeasure, GaussianField-NuclearTensorProduct-instAddCommGroup, GaussianField-Configuration, GaussianField-NuclearTensorProduct-instContinuousSMulReal, GaussianField-SmoothMap-Circle, GaussianField-circleSpacing, GaussianField-instMeasurableSpaceConfiguration, Pphi2-torusInteractingMeasure, GaussianField-FinLatticeField, GaussianField-TorusTestFunction, GaussianField-circleSpacing-pos}
  \texttt{Pphi2.torusInteractingMeasure\_isProbability}
\end{theorem}
\begin{theorem}[\texttt{torus\_interacting\_second\_moment\_uniform}]\label{Pphi2-torus-interacting-second-moment-uniform}
  \lean{Pphi2.torus_interacting_second_moment_uniform}\leanok
  \uses{Pphi2-torusEmbeddedTwoPoint-eq-lattice-second-moment, GaussianField-NuclearTensorProduct-instModuleReal, Pphi2-partitionFunction-ge-one, GaussianField-NuclearTensorProduct-instTopologicalSpace, GaussianField-finLatticeField-dyninMityaginSpace, Pphi2-nelson-exponential-estimate, GaussianField-measure, GaussianField-ell2-separableSpace, Pphi2-torusEmbeddedTwoPoint, GaussianField-FinLatticeSites, GaussianField-NuclearTensorProduct-instIsTopologicalAddGroup, Pphi2-torusEmbedLift, Pphi2-partitionFunction, GaussianField-latticeCovarianceGJ, GaussianField-pairing-memLp, Pphi2-torusEmbedLift-measurable, Pphi2-interactingLatticeMeasure, GaussianField-NuclearTensorProduct-instAddCommGroup, GaussianField-latticeGaussianMeasure, GaussianField-configuration-eval-measurable, GaussianField-Configuration, Pphi2-density-transfer-bound, GaussianField-NuclearTensorProduct-instContinuousSMulReal, GaussianField-gaussian-hypercontractive, GaussianField-SmoothMap-Circle, Pphi2-torusEmbeddedTwoPoint-uniform-bound, GaussianField-circleSpacing, GaussianField-instMeasurableSpaceConfiguration, Pphi2-latticeTestFn, Pphi2-torusEmbedLift-eval-eq, GaussianField-ell2-, Pphi2-torusInteractingMeasure, GaussianField-FinLatticeField, Pphi2-interactionFunctional, GaussianField-TorusTestFunction, Lattice-toSemilatticeInf, GaussianField-circleSpacing-pos}
  $\int (\omega f)^2\, d\mu_{P,N}^{\text{torus}} \le C$ uniformly in $N$. Proved via Cauchy-Schwarz density transfer: $E_P[F] \le \sqrt{K} \sqrt{E_{\text{GFF}}[F^2]}$, then Gaussian hypercontractivity bounds $E_{\text{GFF}}[(\omega g)^4] \le 9 (E_{\text{GFF}}[(\omega g)^2])^2$, and the Gaussian second moment is uniformly bounded.
\end{theorem}
\begin{theorem}[\texttt{torus\_interacting\_tightness}]\label{Pphi2-torus-interacting-tightness}
  \lean{Pphi2.torus_interacting_tightness}\leanok
  \uses{GaussianField-NuclearTensorProduct-instModuleReal, GaussianField-NuclearTensorProduct-instTopologicalSpace, GaussianField-finLatticeField-dyninMityaginSpace, GaussianField-measure, GaussianField-ell2-separableSpace, GaussianField-NuclearTensorProduct-dyninMityaginSpace, Pphi2-interactionFunctional-bounded-below, GaussianField-FinLatticeSites, GaussianField-NuclearTensorProduct-instIsTopologicalAddGroup, Pphi2-torusEmbedLift, Pphi2-partitionFunction, GaussianField-latticeCovarianceGJ, Pphi2-interactionFunctional-measurable, GaussianField-pairing-memLp, Pphi2-torusEmbedLift-measurable, Pphi2-interactingLatticeMeasure, GaussianField-NuclearTensorProduct-instAddCommGroup, GaussianField-latticeGaussianMeasure, GaussianField-configuration-eval-measurable, GaussianField-Configuration, GaussianField-NuclearTensorProduct-instContinuousSMulReal, GaussianField-SmoothMap-Circle, GaussianField-circleSpacing, GaussianField-instMeasurableSpaceConfiguration, Pphi2-boltzmannWeight, Pphi2-latticeTestFn, Pphi2-torusEmbedLift-eval-eq, GaussianField-ell2-, Pphi2-torusInteractingMeasure-isProbability, Pphi2-torusInteractingMeasure, GaussianField-FinLatticeField, Pphi2-torus-interacting-second-moment-uniform, Pphi2-interactionFunctional, GaussianField-configuration-tight-of-uniform-second-moments, GaussianField-TorusTestFunction, Pphi2-partitionFunction-pos, GaussianField-circleSpacing-pos, Pphi2-boltzmannWeight-pos}
  The family $\{\mu_{P,N}^{\text{torus}}\}_{N \ge 1}$ is tight. Proved from uniform second moments via \texttt{configuration\_tight\_of\_uniform\_second\_moments}.
\end{theorem}
\begin{theorem}[\texttt{nelson\_exponential\_estimate}]\label{Pphi2-nelson-exponential-estimate}
  \lean{Pphi2.nelson_exponential_estimate}\leanok
  \uses{GaussianField-FinLatticeSites, GaussianField-latticeGaussianMeasure, GaussianField-Configuration, GaussianField-circleSpacing, GaussianField-instMeasurableSpaceConfiguration, Pphi2-nelson-exponential-estimate-lattice, GaussianField-FinLatticeField, Pphi2-interactionFunctional, GaussianField-circleSpacing-pos}
  $\int e^{-2V_{L/N}(\varphi)}\, d\mu_{\text{GFF}} \le K$ uniformly in $N$, where $K$ depends on $P$, $m$, $L$ but not $N$. Proved because $a^2 N^2 = L^2$ (fixed physical volume), so $V(\omega) \ge -L^2 A$ gives $e^{-2V} \le e^{2L^2 A}$.
\end{theorem}
\begin{theorem}[\texttt{torusInteractingLimit\_exists}]\label{Pphi2-torusInteractingLimit-exists}
  \lean{Pphi2.torusInteractingLimit_exists}\leanok
  \uses{GaussianField-NuclearTensorProduct-instModuleReal, GaussianField-NuclearTensorProduct-instTopologicalSpace, GaussianField-NuclearTensorProduct-dyninMityaginSpace, GaussianField-NuclearTensorProduct-instIsTopologicalAddGroup, GaussianField-NuclearTensorProduct-instAddCommGroup, GaussianField-prokhorov-configuration, GaussianField-Configuration, GaussianField-NuclearTensorProduct-instContinuousSMulReal, GaussianField-SmoothMap-Circle, GaussianField-instMeasurableSpaceConfiguration, Pphi2-torus-interacting-tightness, Pphi2-torusInteractingMeasure-isProbability, Pphi2-torusInteractingMeasure, GaussianField-TorusTestFunction}
  For $N \to \infty$, the interacting measures have a weakly convergent subsequence. Proved from tightness + \texttt{prokhorov\_configuration}.
\end{theorem}
\begin{definition}[\texttt{torusInteractingMeasure}]\label{Pphi2-torusInteractingMeasure}
  \lean{Pphi2.torusInteractingMeasure}\leanok
  \uses{GaussianField-NuclearTensorProduct-instModuleReal, GaussianField-NuclearTensorProduct-instTopologicalSpace, GaussianField-FinLatticeSites, GaussianField-NuclearTensorProduct-instIsTopologicalAddGroup, Pphi2-torusEmbedLift, Pphi2-interactingLatticeMeasure, GaussianField-NuclearTensorProduct-instAddCommGroup, GaussianField-Configuration, GaussianField-NuclearTensorProduct-instContinuousSMulReal, GaussianField-SmoothMap-Circle, GaussianField-circleSpacing, GaussianField-instMeasurableSpaceConfiguration, GaussianField-FinLatticeField, GaussianField-TorusTestFunction, GaussianField-circleSpacing-pos}
  ``\texttt{lean def torusInteractingMeasure (N : ℕ) [NeZero N] (P : InteractionPolynomial) (mass : ℝ) (hmass : 0 < mass) : Measure (Configuration (TorusTestFunction L)) }`` Pushforward $\mu_{P,N}^{\text{torus}} = (\tilde\iota_N)_* \mu_{P,N}$.
\end{definition}
\begin{theorem}[\texttt{expMoment\_two\_le\_uniform}]\label{Pphi2-expMoment-two-le-uniform}
  \lean{Pphi2.expMoment_two_le_uniform}\leanok
  \uses{Pphi2-nelson-exponential-estimate, GaussianField-FinLatticeSites, GaussianField-latticeGaussianMeasure, GaussianField-Configuration, GaussianField-circleSpacing, GaussianField-instMeasurableSpaceConfiguration, GaussianField-FinLatticeField, Pphi2-interactionFunctional, GaussianField-circleSpacing-pos}
  \texttt{Pphi2.expMoment\_two\_le\_uniform}
\end{theorem}
\section{\texttt{Pphi2.AsymTorus.AsymTransferInstantiation}}
\begin{theorem}[\texttt{asymTransferKernel\_symm}]\label{Pphi2-asymTransferKernel-symm}
  \lean{Pphi2.asymTransferKernel_symm}\leanok
  \uses{Pphi2-transferGaussian-sub-comm, Pphi2-SpatialField, Pphi2-asymTransferKernel, Pphi2-asymTransferWeight, Pphi2-transferGaussian}
  \texttt{Pphi2.asymTransferKernel\_symm}
\end{theorem}
\begin{definition}[\texttt{asymTransferKernel}]\label{Pphi2-asymTransferKernel}
  \lean{Pphi2.asymTransferKernel}\leanok
  \uses{Pphi2-SpatialField, Pphi2-asymTransferWeight, Pphi2-transferGaussian}
  \texttt{Pphi2.asymTransferKernel}
\end{definition}
\begin{theorem}[\texttt{asymTransferKernel\_nonneg}]\label{Pphi2-asymTransferKernel-nonneg}
  \lean{Pphi2.asymTransferKernel_nonneg}\leanok
  \uses{Pphi2-SpatialField, Pphi2-asymTransferKernel, Pphi2-asymTransferWeight, Pphi2-asymTransferWeight-pos, Pphi2-timeCoupling, Pphi2-transferGaussian}
  \texttt{Pphi2.asymTransferKernel\_nonneg}
\end{theorem}
\begin{theorem}[\texttt{congr\_simp}]\label{Pphi2-asymTransferWeight-congr-simp}
  \lean{Pphi2.asymTransferWeight.congr_simp}\leanok
  \uses{Pphi2-SpatialField, Pphi2-asymTransferWeight}
  \texttt{Pphi2.asymTransferWeight.congr\_simp}
\end{theorem}
\begin{definition}[\texttt{asymTransferSystem}]\label{Pphi2-asymTransferSystem}
  \lean{Pphi2.asymTransferSystem}\leanok
  \uses{Pphi2-SpatialField, Pphi2-asymTransferKernel, Pphi2-asymTransferKernel-measurable, Pphi2-asymTransferKernel-symm, Pphi2-asymTransferKernel-nonneg}
  \texttt{Pphi2.asymTransferSystem}
\end{definition}
\begin{theorem}[\texttt{asymSliceEquiv\_apply}]\label{Pphi2-asymSliceEquiv-apply}
  \lean{Pphi2.asymSliceEquiv_apply}\leanok
  \uses{GaussianField-AsymLatticeSites, Pphi2-SpatialField, Pphi2-asymSliceEquiv, GaussianField-AsymLatticeField}
  \texttt{Pphi2.asymSliceEquiv\_apply}
\end{theorem}
\begin{theorem}[\texttt{asymTransferKernel\_measurable}]\label{Pphi2-asymTransferKernel-measurable}
  \lean{Pphi2.asymTransferKernel_measurable}\leanok
  \uses{Pphi2-asymTransferWeight-measurable, Pphi2-SpatialField, Pphi2-asymTransferKernel, Pphi2-continuous-transferGaussian, Pphi2-asymTransferWeight, Pphi2-transferGaussian}
  \texttt{Pphi2.asymTransferKernel\_measurable}
\end{theorem}
\begin{theorem}[\texttt{congr\_simp}]\label{Pphi2-asymTransferKernel-congr-simp}
  \lean{Pphi2.asymTransferKernel.congr_simp}\leanok
  \uses{Pphi2-SpatialField, Pphi2-asymTransferKernel}
  \texttt{Pphi2.asymTransferKernel.congr\_simp}
\end{theorem}
\begin{definition}[\texttt{prodCurryLinearEquiv}]\label{Pphi2-prodCurryLinearEquiv}
  \lean{Pphi2.prodCurryLinearEquiv}\leanok
  \texttt{Pphi2.prodCurryLinearEquiv}
\end{definition}
\begin{definition}[\texttt{asymSliceEquiv}]\label{Pphi2-asymSliceEquiv}
  \lean{Pphi2.asymSliceEquiv}\leanok
  \uses{GaussianField-AsymLatticeSites, Pphi2-SpatialField, Pphi2-prodCurryLinearEquiv, GaussianField-AsymLatticeField}
  \texttt{Pphi2.asymSliceEquiv}
\end{definition}
\section{\texttt{Pphi2.AsymTorus.AsymNelson}}
\begin{theorem}[\texttt{asymNelson\_exponential\_estimate\_iso}]\label{Pphi2-asymNelson-exponential-estimate-iso}
  \lean{Pphi2.asymNelson_exponential_estimate_iso}\leanok
  \uses{GaussianField-AsymLatticeSites, Pphi2-asymChaosCutoffDecomposition, GaussianField-Configuration, GaussianField-AsymLatticeField, Pphi2-latticeGaussianMeasureAsym, Pphi2-interactionFunctionalAsym, GaussianField-instMeasurableSpaceConfiguration, Pphi2-BridgeFromTail-bridgeAxiom-of-setup-real-generic, Pphi2-interactionFunctionalAsym-measurable, Pphi2-latticeGaussianMeasureAsym-isProbability}
  \texttt{Pphi2.asymNelson\_exponential\_estimate\_iso}
\end{theorem}
\begin{theorem}[\texttt{interactionFunctionalAsym\_asymCanonicalSumConfig\_eq}]\label{Pphi2-interactionFunctionalAsym-asymCanonicalSumConfig-eq}
  \lean{Pphi2.interactionFunctionalAsym_asymCanonicalSumConfig_eq}\leanok
  \uses{Pphi2-asymCanonicalFullInteractionJoint, GaussianField-AsymLatticeSites, Pphi2-wickConstantAsym, Pphi2-AsymCanonicalJoint, Pphi2-asymCanonicalSumConfig-apply-delta, Pphi2-asymCanonicalSumFieldFunction, Pphi2-asymCanonicalSumConfig, GaussianField-Configuration, GaussianField-AsymLatticeField, Pphi2-interactionFunctionalAsym, GaussianField-instFintypeAsymLatticeSitesOfNeZeroNat, GaussianField-asymLatticeDelta, Pphi2-wickPolynomial}
  \texttt{Pphi2.interactionFunctionalAsym\_asymCanonicalSumConfig\_eq}
\end{theorem}
\begin{theorem}[\texttt{congr\_simp}]\label{Pphi2-asymCanonicalSumConfig-congr-simp}
  \lean{Pphi2.asymCanonicalSumConfig.congr_simp}\leanok
  \uses{GaussianField-AsymLatticeSites, Pphi2-AsymCanonicalJoint, Pphi2-asymCanonicalSumConfig, GaussianField-Configuration, GaussianField-AsymLatticeField}
  \texttt{Pphi2.asymCanonicalSumConfig.congr\_simp}
\end{theorem}
\begin{theorem}[\texttt{asymChaosCutoffDecomposition}]\label{Pphi2-asymChaosCutoffDecomposition}
  \lean{Pphi2.asymChaosCutoffDecomposition}\leanok
  \uses{Pphi2-asymCanonicalFullInteractionJoint, Pphi2-degreeCutoffTail, Pphi2-asymCanonicalJointMeasure-map-sumConfig, GaussianField-AsymLatticeSites, Pphi2-asymCanonicalRoughError-neg-tail-uniform, Pphi2-two-div-degree-eq-inv-cutoffPower, Pphi2-AsymCanonicalJoint, Pphi2-asymCanonicalSmoothInteraction-lower-bound-at-cutoff-uniform, Pphi2-DynamicalCutoff-degreeDynamicalCutoffScale-pos, Pphi2-asymCanonicalRoughError, Pphi2-interactionFunctionalAsym-asymCanonicalSumConfig-eq, Pphi2-DynamicalCutoff-degreeDynamicalCutoffScale, Pphi2-asymCanonicalSumConfig, GaussianField-Configuration, GaussianField-AsymLatticeField, Pphi2-latticeGaussianMeasureAsym, Pphi2-interactionFunctionalAsym, GaussianField-instMeasurableSpaceConfiguration, Pphi2-asymCanonicalSmoothInteraction, Pphi2-ChaosTailBridge-chaosTailConstant, Pphi2-DynamicalCutoff-degreeCutoffPower, Pphi2-interactionFunctionalAsym-measurable, Pphi2-latticeGaussianMeasureAsym-isProbability, Pphi2-asymCanonicalJointMeasure, Pphi2-asymCanonicalSumConfig-measurable, Pphi2-degreePiecewiseTail, Pphi2-AsymCanonicalMode, Pphi2-degreePiecewiseTail-layerCake-lt-top}
  \texttt{Pphi2.asymChaosCutoffDecomposition}
\end{theorem}
\begin{theorem}[\texttt{congr\_simp}]\label{Pphi2-asymCanonicalFullInteractionJoint-congr-simp}
  \lean{Pphi2.asymCanonicalFullInteractionJoint.congr_simp}\leanok
  \uses{Pphi2-asymCanonicalFullInteractionJoint, Pphi2-AsymCanonicalJoint}
  \texttt{Pphi2.asymCanonicalFullInteractionJoint.congr\_simp}
\end{theorem}
\section{\texttt{Pphi2.AsymTorus.AsymBridgeInstance}}
\begin{theorem}[\texttt{asymTransferNormalized\_contract} (axiom)]\label{Pphi2-asymTransferNormalized-contract}
  \lean{Pphi2.asymTransferNormalized_contract}\leanok
  \uses{Pphi2-L2SpatialField, Pphi2-SpatialField, Pphi2-asymTransferGroundEigenvalue, Pphi2-asymTransferOperatorCLM}
  \texttt{Pphi2.asymTransferNormalized\_contract}
\end{theorem}
\begin{theorem}[\texttt{asymSliceObsTrunc\_pathMeasure\_connected\_two\_point\_bound}]\label{Pphi2-asymSliceObsTrunc-pathMeasure-connected-two-point-bound}
  \lean{Pphi2.asymSliceObsTrunc_pathMeasure_connected_two_point_bound}\leanok
  \uses{Pphi2-asymTransferNormalized, Pphi2-L2SpatialField, Pphi2-SpatialField, Pphi2-asymGroundVector, Pphi2-asymTransferSystem, Pphi2-asymTransferGroundEigenvalue, Pphi2-asymRemainderHypothesis, Pphi2-asymSliceObsTruncContract, Pphi2-asymGroundStateRep}
  \texttt{Pphi2.asymSliceObsTrunc\_pathMeasure\_connected\_two\_point\_bound}
\end{theorem}
\begin{theorem}[\texttt{congr\_simp}]\label{Pphi2-asymGappedTransfer-congr-simp}
  \lean{Pphi2.asymGappedTransfer.congr_simp}\leanok
  \uses{Pphi2-asymTransferNormalized, Pphi2-L2SpatialField, Pphi2-SpatialField, Pphi2-asymGroundVector, Pphi2-asymGappedTransfer}
  \texttt{Pphi2.asymGappedTransfer.congr\_simp}
\end{theorem}
\begin{definition}[\texttt{asymGroundStateRep}]\label{Pphi2-asymGroundStateRep}
  \lean{Pphi2.asymGroundStateRep}\leanok
  \uses{Pphi2-SpatialField, Pphi2-asymGroundVector}
  \texttt{Pphi2.asymGroundStateRep}
\end{definition}
\begin{theorem}[\texttt{asymFinitePeriodicBridge\_remainder\_bound} (axiom)]\label{Pphi2-asymFinitePeriodicBridge-remainder-bound}
  \lean{Pphi2.asymFinitePeriodicBridge_remainder_bound}\leanok
  \uses{Pphi2-asymTransferNormalized, Pphi2-L2SpatialField, Pphi2-SpatialField, Pphi2-asymGroundVector, Pphi2-asymTransferSystem, Pphi2-asymGappedTransfer}
  \texttt{Pphi2.asymFinitePeriodicBridge\_remainder\_bound}
\end{theorem}
\begin{theorem}[\texttt{asymGroundSemigroup\_intertwines} (axiom)]\label{Pphi2-asymGroundSemigroup-intertwines}
  \lean{Pphi2.asymGroundSemigroup_intertwines}\leanok
  \uses{Pphi2-asymTransferNormalized, Pphi2-L2SpatialField, Pphi2-SpatialField, Pphi2-asymGroundVector, Pphi2-asymGroundStateRep-measurable, Pphi2-asymGroundSemigroupData, Pphi2-asymGappedTransfer, Pphi2-asymGroundStateRep}
  \texttt{Pphi2.asymGroundSemigroup\_intertwines}
\end{theorem}
\begin{theorem}[\texttt{asymGroundStateRep\_eq\_groundIsometry\_one} (axiom)]\label{Pphi2-asymGroundStateRep-eq-groundIsometry-one}
  \lean{Pphi2.asymGroundStateRep_eq_groundIsometry_one}\leanok
  \uses{Pphi2-SpatialField, Pphi2-asymGroundVector, Pphi2-asymGroundStateRep-measurable, Pphi2-asymGroundStateRep-isProbabilityMeasure, Pphi2-asymGroundStateRep}
  \texttt{Pphi2.asymGroundStateRep\_eq\_groundIsometry\_one}
\end{theorem}
\begin{theorem}[\texttt{congr\_simp}]\label{Pphi2-asymRemainderHypothesis-congr-simp}
  \lean{Pphi2.asymRemainderHypothesis.congr_simp}\leanok
  \uses{Pphi2-asymTransferNormalized, Pphi2-L2SpatialField, Pphi2-SpatialField, Pphi2-asymGroundVector, Pphi2-asymTransferSystem, Pphi2-asymTransferGroundEigenvalue, Pphi2-asymRemainderHypothesis, Pphi2-asymGroundStateRep}
  \texttt{Pphi2.asymRemainderHypothesis.congr\_simp}
\end{theorem}
\begin{definition}[\texttt{asymRemainderHypothesis}]\label{Pphi2-asymRemainderHypothesis}
  \lean{Pphi2.asymRemainderHypothesis}\leanok
  \uses{Pphi2-asymTransferNormalized, Pphi2-L2SpatialField, Pphi2-SpatialField, Pphi2-asymGroundVector, Pphi2-asymFinitePeriodicBridge-remainder-bound, Pphi2-asymTransferSystem, Pphi2-asymTransferGroundEigenvalue, Pphi2-asymGappedTransfer, Pphi2-asymGroundVector-coeFn-eq-groundStateRep, Pphi2-asymGroundStateRep, Pphi2-asymPartition-ground-bound}
  \texttt{Pphi2.asymRemainderHypothesis}
\end{definition}
\begin{theorem}[\texttt{asymSliceObsTrunc\_pathMeasure\_connected\_susceptibility\_bound\_pos}]\label{Pphi2-asymSliceObsTrunc-pathMeasure-connected-susceptibility-bound-pos}
  \lean{Pphi2.asymSliceObsTrunc_pathMeasure_connected_susceptibility_bound_pos}\leanok
  \uses{Pphi2-asymTransferNormalized, Pphi2-L2SpatialField, Pphi2-SpatialField, Pphi2-asymGroundVector, Pphi2-asymTransferSystem, Pphi2-asymTransferGroundEigenvalue, Pphi2-asymRemainderHypothesis, Pphi2-asymSliceObsTruncContract, Pphi2-asymGroundStateRep}
  \texttt{Pphi2.asymSliceObsTrunc\_pathMeasure\_connected\_susceptibility\_bound\_pos}
\end{theorem}
\begin{theorem}[\texttt{asymGroundStateRep\_isProbabilityMeasure}]\label{Pphi2-asymGroundStateRep-isProbabilityMeasure}
  \lean{Pphi2.asymGroundStateRep_isProbabilityMeasure}\leanok
  \uses{Pphi2-SpatialField, Pphi2-asymGroundStateRep-measurable, Pphi2-asymGroundStateRep-norm-integral-eq-one, Pphi2-asymGroundStateRep}
  \texttt{Pphi2.asymGroundStateRep\_isProbabilityMeasure}
\end{theorem}
\begin{theorem}[\texttt{asymPartition\_ground\_bound} (axiom)]\label{Pphi2-asymPartition-ground-bound}
  \lean{Pphi2.asymPartition_ground_bound}\leanok
  \uses{Pphi2-SpatialField, Pphi2-asymTransferSystem, Pphi2-asymTransferGroundEigenvalue}
  \texttt{Pphi2.asymPartition\_ground\_bound}
\end{theorem}
\begin{theorem}[\texttt{congr\_simp}]\label{Pphi2-asymSliceObsTruncContract-congr-simp}
  \lean{Pphi2.asymSliceObsTruncContract.congr_simp}\leanok
  \uses{Pphi2-SpatialField, Pphi2-asymSliceObsTruncContract}
  \texttt{Pphi2.asymSliceObsTruncContract.congr\_simp}
\end{theorem}
\begin{theorem}[\texttt{asymSliceObsTrunc\_pathMeasure\_connected\_susceptibility\_bound\_pos\_integral}]\label{Pphi2-asymSliceObsTrunc-pathMeasure-connected-susceptibility-bound-pos-integral}
  \lean{Pphi2.asymSliceObsTrunc_pathMeasure_connected_susceptibility_bound_pos_integral}\leanok
  \uses{Pphi2-asymSliceObsLinear, Pphi2-asymTransferNormalized, Pphi2-asymSliceObsTrunc-pathMeasure-connected-susceptibility-bound-pos, Pphi2-L2SpatialField, Pphi2-SpatialField, Pphi2-asymGroundVector, Pphi2-asymTransferSystem, Pphi2-norm-sq-proj-obsTrunc-omega-le, Pphi2-asymSliceObsTruncMulCLM, Pphi2-asymTransferGroundEigenvalue, Pphi2-asymRemainderHypothesis, Pphi2-asymSliceObsTruncContract, Pphi2-asymGroundStateRep}
  \texttt{Pphi2.asymSliceObsTrunc\_pathMeasure\_connected\_susceptibility\_bound\_pos\_integral}
\end{theorem}
\begin{theorem}[\texttt{asymGroundStateRep\_pos\_ae} (axiom)]\label{Pphi2-asymGroundStateRep-pos-ae}
  \lean{Pphi2.asymGroundStateRep_pos_ae}\leanok
  \uses{Pphi2-SpatialField, Pphi2-asymGroundStateRep}
  \texttt{Pphi2.asymGroundStateRep\_pos\_ae}
\end{theorem}
\begin{theorem}[\texttt{asym\_interacting\_second\_moment\_eq\_pathMeasure\_slicePairing}]\label{Pphi2-asym-interacting-second-moment-eq-pathMeasure-slicePairing}
  \lean{Pphi2.asym_interacting_second_moment_eq_pathMeasure_slicePairing}\leanok
  \uses{GaussianField-AsymLatticeSites, Pphi2-interacting-second-moment-eq-pathMeasure, Pphi2-SpatialField, Pphi2-asymTransferSystem, GaussianField-Configuration, GaussianField-AsymLatticeField, Pphi2-slicePairing, Pphi2-interactingLatticeMeasureAsym, GaussianField-instMeasurableSpaceConfiguration}
  \texttt{Pphi2.asym\_interacting\_second\_moment\_eq\_pathMeasure\_slicePairing}
\end{theorem}
\begin{theorem}[\texttt{asymGroundVector\_coeFn\_eq\_groundStateRep}]\label{Pphi2-asymGroundVector-coeFn-eq-groundStateRep}
  \lean{Pphi2.asymGroundVector_coeFn_eq_groundStateRep}\leanok
  \uses{Pphi2-SpatialField, Pphi2-asymGroundVector, Pphi2-asymGroundStateRep}
  \texttt{Pphi2.asymGroundVector\_coeFn\_eq\_groundStateRep}
\end{theorem}
\begin{theorem}[\texttt{asymGroundStateRep\_norm\_integral\_eq\_one}]\label{Pphi2-asymGroundStateRep-norm-integral-eq-one}
  \lean{Pphi2.asymGroundStateRep_norm_integral_eq_one}\leanok
  \uses{Pphi2-L2SpatialField, Pphi2-SpatialField, Pphi2-asymGroundVector, Pphi2-asymGroundVector-norm-eq-one, Pphi2-asymGroundVector-coeFn-eq-groundStateRep, Pphi2-asymGroundStateRep}
  \texttt{Pphi2.asymGroundStateRep\_norm\_integral\_eq\_one}
\end{theorem}
\begin{definition}[\texttt{asymGroundGapData}]\label{Pphi2-asymGroundGapData}
  \lean{Pphi2.asymGroundGapData}\leanok
  \uses{Pphi2-asymTransferNormalized, Pphi2-L2SpatialField, Pphi2-SpatialField, Pphi2-asymGroundVector, Pphi2-asymGroundStateRep-measurable, Pphi2-asymGroundSemigroupData, Pphi2-asymGroundSemigroup-intertwines, Pphi2-asymGappedTransfer, Pphi2-asymGroundStateRep-isProbabilityMeasure, Pphi2-asymGroundStateRep}
  \texttt{Pphi2.asymGroundGapData}
\end{definition}
\begin{theorem}[\texttt{asymGroundStateRep\_measurable}]\label{Pphi2-asymGroundStateRep-measurable}
  \lean{Pphi2.asymGroundStateRep_measurable}\leanok
  \uses{Pphi2-SpatialField, Pphi2-asymGroundVector, Pphi2-asymGroundStateRep}
  \texttt{Pphi2.asymGroundStateRep\_measurable}
\end{theorem}
\begin{definition}[\texttt{asymGroundSemigroupData}]\label{Pphi2-asymGroundSemigroupData}
  \lean{Pphi2.asymGroundSemigroupData}\leanok
  \uses{Pphi2-asymTransferOperatorCLM-groundVector, Pphi2-SpatialField, Pphi2-asymGroundVector, Pphi2-asymGroundStateRep-measurable, Pphi2-asymTransferGroundEigenvalue-pos, Pphi2-asymGroundStateRep-eq-groundIsometry-one, Pphi2-asymTransferGroundEigenvalue, Pphi2-asymGroundStateRep-pos-ae, Pphi2-asymTransferNormalized-contract, Pphi2-asymGroundVector-coeFn-eq-groundStateRep, Pphi2-asymGroundStateRep-norm-integral-eq-one, Pphi2-asymGroundStateRep, Pphi2-asymTransferOperatorCLM}
  \texttt{Pphi2.asymGroundSemigroupData}
\end{definition}
\section{\texttt{Pphi2.InteractingMeasure.ConnectedCorrelatorDerivative}}
\begin{theorem}[\texttt{partitionFn\_zero}]\label{Pphi2-partitionFn-zero}
  \lean{Pphi2.partitionFn_zero}\leanok
  \uses{GaussianField-FinLatticeSites, GaussianField-latticeGaussianMeasure-isProbability, GaussianField-latticeGaussianMeasure, GaussianField-Configuration, GaussianField-instMeasurableSpaceConfiguration, GaussianField-FinLatticeField, Pphi2-interactionFunctional}
  \texttt{Pphi2.partitionFn\_zero}
\end{theorem}
\begin{theorem}[\texttt{normalizedMoment\_hasDerivWithinAt}]\label{Pphi2-normalizedMoment-hasDerivWithinAt}
  \lean{Pphi2.normalizedMoment_hasDerivWithinAt}\leanok
  \uses{Pphi2-partitionFn-zero, GaussianField-FinLatticeSites, GaussianField-latticeGaussianMeasure, GaussianField-Configuration, GaussianField-instMeasurableSpaceConfiguration, Pphi2-moment-hasDerivWithinAt, Pphi2-partitionFn-hasDerivWithinAt, GaussianField-FinLatticeField, Pphi2-interactionFunctional}
  \texttt{Pphi2.normalizedMoment\_hasDerivWithinAt}
\end{theorem}
\begin{theorem}[\texttt{partitionFn\_hasDerivWithinAt}]\label{Pphi2-partitionFn-hasDerivWithinAt}
  \lean{Pphi2.partitionFn_hasDerivWithinAt}\leanok
  \uses{GaussianField-FinLatticeSites, GaussianField-latticeGaussianMeasure, GaussianField-Configuration, GaussianField-instMeasurableSpaceConfiguration, Pphi2-moment-hasDerivWithinAt, GaussianField-FinLatticeField, Pphi2-interactionFunctional}
  \texttt{Pphi2.partitionFn\_hasDerivWithinAt}
\end{theorem}
\section{\texttt{Pphi2.TorusContinuumLimit.TorusConvergence}}
\begin{theorem}[\texttt{torusGaussianLimit\_exists}]\label{Pphi2-torusGaussianLimit-exists}
  \lean{Pphi2.torusGaussianLimit_exists}\leanok
  \uses{GaussianField-NuclearTensorProduct-instModuleReal, GaussianField-NuclearTensorProduct-instTopologicalSpace, GaussianField-NuclearTensorProduct-dyninMityaginSpace, GaussianField-NuclearTensorProduct-instIsTopologicalAddGroup, Pphi2-torusContinuumMeasure, GaussianField-NuclearTensorProduct-instAddCommGroup, GaussianField-prokhorov-configuration, GaussianField-Configuration, GaussianField-NuclearTensorProduct-instContinuousSMulReal, Pphi2-torusContinuumMeasure-isProbability, GaussianField-SmoothMap-Circle, GaussianField-instMeasurableSpaceConfiguration, Pphi2-torusContinuumMeasures-tight, GaussianField-TorusTestFunction}
  For $N \to \infty$, the torus-embedded Gaussian measures $\nu_{\text{GFF},N}$ have a weakly convergent subsequence. The limit is a probability measure on Configuration(TorusTestFunction $L$).  Proved from: 1. \texttt{torusContinuumMeasures\_tight} (tightness) 2. \texttt{configuration\_torus\_polish} (Polish space) 3. \texttt{prokhorov\_configuration} (proved Prokhorov extraction)
\end{theorem}
\section{\texttt{Pphi2.AsymTorus.AsymBridgeKLimit}}
\begin{theorem}[\texttt{asymSliceFamilyTrunc\_measurable}]\label{Pphi2-asymSliceFamilyTrunc-measurable}
  \lean{Pphi2.asymSliceFamilyTrunc_measurable}\leanok
  \uses{Pphi2-asymSliceObsTrunc-measurable, Pphi2-SpatialField, Pphi2-asymSliceFamilyTrunc, Pphi2-asymSliceObsTrunc}
  \texttt{Pphi2.asymSliceFamilyTrunc\_measurable}
\end{theorem}
\begin{theorem}[\texttt{asymInteractingLattice\_pairing\_sq\_integrable}]\label{Pphi2-asymInteractingLattice-pairing-sq-integrable}
  \lean{Pphi2.asymInteractingLattice_pairing_sq_integrable}\leanok
  \uses{GaussianField-AsymLatticeSites, GaussianField-measure, GaussianField-ell2-separableSpace, Pphi2-partitionFunctionAsym-pos, GaussianField-asymLatticeField-dyninMityaginSpace, Pphi2-boltzmannWeightAsym, GaussianField-pairing-memLp, Pphi2-boltzmannWeightAsym-pos, Pphi2-interactionFunctionalAsym-bounded-below, Pphi2-partitionFunctionAsym, GaussianField-configuration-eval-measurable, GaussianField-Configuration, GaussianField-AsymLatticeField, Pphi2-latticeGaussianMeasureAsym, Pphi2-interactingLatticeMeasureAsym, Pphi2-interactionFunctionalAsym, GaussianField-instMeasurableSpaceConfiguration, GaussianField-latticeCovarianceAsymGJ, GaussianField-ell2-, Pphi2-interactionFunctionalAsym-measurable}
  \texttt{Pphi2.asymInteractingLattice\_pairing\_sq\_integrable}
\end{theorem}
\begin{theorem}[\texttt{asymSliceFamilyLinear\_sq\_integrable\_pathMeasure}]\label{Pphi2-asymSliceFamilyLinear-sq-integrable-pathMeasure}
  \lean{Pphi2.asymSliceFamilyLinear_sq_integrable_pathMeasure}\leanok
  \uses{Pphi2-asymInteractingLattice-pairing-sq-integrable, GaussianField-AsymLatticeSites, GaussianField-measurable-evalMapAsym, Pphi2-SpatialField, Pphi2-asymSliceEquiv, Pphi2-interactingLatticeMeasureAsym-slice-pushforward-eq-pathMeasure, Pphi2-asymTransferSystem, Pphi2-asymSliceFamilyLinear, Pphi2-asymSliceFamilyLinear-measurable, Pphi2-config-pairing-eq-slice, GaussianField-Configuration, GaussianField-AsymLatticeField, Pphi2-slicePairing, GaussianField-evalMapAsym, Pphi2-interactingLatticeMeasureAsym, Pphi2-measurePreserving-asymSliceEquiv, GaussianField-instMeasurableSpaceConfiguration, GaussianField-instFintypeAsymLatticeSitesOfNeZeroNat, Pphi2-asymSliceFamilyLinear-eq-slicePairing}
  \texttt{Pphi2.asymSliceFamilyLinear\_sq\_integrable\_pathMeasure}
\end{theorem}
\begin{theorem}[\texttt{asymSliceFamilyTrunc\_sq\_integral\_tendsto\_linear}]\label{Pphi2-asymSliceFamilyTrunc-sq-integral-tendsto-linear}
  \lean{Pphi2.asymSliceFamilyTrunc_sq_integral_tendsto_linear}\leanok
  \uses{Pphi2-asymSliceFamilyTrunc-measurable, Pphi2-SpatialField, Pphi2-asymTransferSystem, Pphi2-asymSliceFamilyLinear, Pphi2-asymSliceFamilyTrunc, Pphi2-asymSliceObsLinear-sq-integrable-pathMeasure, Pphi2-asymSliceFamilyTrunc-tendsto-linear, Pphi2-asymSliceFamilyAbsLinear}
  \texttt{Pphi2.asymSliceFamilyTrunc\_sq\_integral\_tendsto\_linear}
\end{theorem}
\begin{theorem}[\texttt{asymSliceFamilyLinear\_eq\_slicePairing}]\label{Pphi2-asymSliceFamilyLinear-eq-slicePairing}
  \lean{Pphi2.asymSliceFamilyLinear_eq_slicePairing}\leanok
  \uses{Pphi2-asymSliceObsLinear, GaussianField-AsymLatticeSites, Pphi2-SpatialField, Pphi2-asymSliceEquiv, Pphi2-asymSliceFamilyLinear, GaussianField-AsymLatticeField, Pphi2-slicePairing}
  \texttt{Pphi2.asymSliceFamilyLinear\_eq\_slicePairing}
\end{theorem}
\begin{theorem}[\texttt{congr\_simp}]\label{Pphi2-slicePairing-congr-simp}
  \lean{Pphi2.slicePairing.congr_simp}\leanok
  \uses{Pphi2-SpatialField, GaussianField-AsymLatticeField, Pphi2-slicePairing}
  \texttt{Pphi2.slicePairing.congr\_simp}
\end{theorem}
\begin{theorem}[\texttt{asymSliceObsLinear\_pathMeasure\_two\_point\_bound}]\label{Pphi2-asymSliceObsLinear-pathMeasure-two-point-bound}
  \lean{Pphi2.asymSliceObsLinear_pathMeasure_two_point_bound}\leanok
  \uses{Pphi2-SpatialField, Pphi2-asymTransferSystem, Pphi2-asymSliceFamilyLinear, Pphi2-asymSliceFamilyTrunc, Pphi2-asymSliceFamilyTrunc-sq-integral-tendsto-linear}
  \texttt{Pphi2.asymSliceObsLinear\_pathMeasure\_two\_point\_bound}
\end{theorem}
\begin{theorem}[\texttt{asymSliceFamilyTrunc\_tendsto\_linear}]\label{Pphi2-asymSliceFamilyTrunc-tendsto-linear}
  \lean{Pphi2.asymSliceFamilyTrunc_tendsto_linear}\leanok
  \uses{Pphi2-asymSliceObsLinear, Pphi2-asymSliceObsTrunc-tendsto-linear, Pphi2-SpatialField, Pphi2-asymSliceFamilyLinear, Pphi2-asymSliceFamilyTrunc, Pphi2-asymSliceObsTrunc}
  \texttt{Pphi2.asymSliceFamilyTrunc\_tendsto\_linear}
\end{theorem}
\begin{definition}[\texttt{asymSliceFamilyLinear}]\label{Pphi2-asymSliceFamilyLinear}
  \lean{Pphi2.asymSliceFamilyLinear}\leanok
  \uses{Pphi2-asymSliceObsLinear, Pphi2-SpatialField}
  \texttt{Pphi2.asymSliceFamilyLinear}
\end{definition}
\begin{theorem}[\texttt{asymSliceFamilyLinear\_measurable}]\label{Pphi2-asymSliceFamilyLinear-measurable}
  \lean{Pphi2.asymSliceFamilyLinear_measurable}\leanok
  \uses{GaussianField-AsymLatticeSites, Pphi2-SpatialField, Pphi2-asymSliceEquiv, Pphi2-slicePairing-measurable, Pphi2-asymSliceFamilyLinear, GaussianField-AsymLatticeField, Pphi2-slicePairing, Pphi2-asymSliceFamilyLinear-eq-slicePairing}
  \texttt{Pphi2.asymSliceFamilyLinear\_measurable}
\end{theorem}
\begin{theorem}[\texttt{asymSliceObsTrunc\_tendsto\_linear}]\label{Pphi2-asymSliceObsTrunc-tendsto-linear}
  \lean{Pphi2.asymSliceObsTrunc_tendsto_linear}\leanok
  \uses{Pphi2-asymSliceObsLinear, Pphi2-SpatialField, Pphi2-asymSliceObsTrunc}
  \texttt{Pphi2.asymSliceObsTrunc\_tendsto\_linear}
\end{theorem}
\begin{theorem}[\texttt{asymSliceFamilyLinear\_pathMeasure\_second\_moment\_le\_B1}]\label{Pphi2-asymSliceFamilyLinear-pathMeasure-second-moment-le-B1}
  \lean{Pphi2.asymSliceFamilyLinear_pathMeasure_second_moment_le_B1}\leanok
  \uses{Pphi2-asym-interacting-second-moment-eq-pathMeasure-slicePairing, GaussianField-AsymLatticeSites, Pphi2-SpatialField, Pphi2-asymSliceEquiv, Pphi2-asymTransferSystem, Pphi2-asymSliceFamilyLinear, GaussianField-Configuration, GaussianField-AsymLatticeField, Pphi2-latticeGaussianMeasureAsym, Pphi2-slicePairing, Pphi2-interactingLatticeMeasureAsym, GaussianField-instMeasurableSpaceConfiguration, Pphi2-asymSliceFamilyLinear-eq-slicePairing, Pphi2-asymInteractingVariance-le-freeVariance-lattice}
  \texttt{Pphi2.asymSliceFamilyLinear\_pathMeasure\_second\_moment\_le\_B1}
\end{theorem}
\begin{theorem}[\texttt{congr\_simp}]\label{Pphi2-asymSliceFamilyLinear-congr-simp}
  \lean{Pphi2.asymSliceFamilyLinear.congr_simp}\leanok
  \uses{Pphi2-SpatialField, Pphi2-asymSliceFamilyLinear}
  \texttt{Pphi2.asymSliceFamilyLinear.congr\_simp}
\end{theorem}
\begin{theorem}[\texttt{asymSliceObsLinear\_sq\_integrable\_pathMeasure}]\label{Pphi2-asymSliceObsLinear-sq-integrable-pathMeasure}
  \lean{Pphi2.asymSliceObsLinear_sq_integrable_pathMeasure}\leanok
  \uses{Pphi2-asymSliceObsLinear, Lattice-toSemilatticeSup, Pphi2-SpatialField, Pphi2-asymSliceObsLinear-measurable, Pphi2-asymTransferSystem, Pphi2-asymSliceFamilyLinear, Pphi2-asymSliceFamilyAbsLinear, Pphi2-asymSliceFamilyLinear-sq-integrable-pathMeasure, Lattice-toSemilatticeInf}
  \texttt{Pphi2.asymSliceObsLinear\_sq\_integrable\_pathMeasure}
\end{theorem}
\begin{definition}[\texttt{asymSliceFamilyAbsLinear}]\label{Pphi2-asymSliceFamilyAbsLinear}
  \lean{Pphi2.asymSliceFamilyAbsLinear}\leanok
  \uses{Pphi2-asymSliceObsLinear, Pphi2-SpatialField}
  \texttt{Pphi2.asymSliceFamilyAbsLinear}
\end{definition}
\begin{definition}[\texttt{asymSliceFamilyTrunc}]\label{Pphi2-asymSliceFamilyTrunc}
  \lean{Pphi2.asymSliceFamilyTrunc}\leanok
  \uses{Pphi2-SpatialField, Pphi2-asymSliceObsTrunc}
  \texttt{Pphi2.asymSliceFamilyTrunc}
\end{definition}
\section{\texttt{Pphi2.AsymTorus.AsymTorusInteractingLimit}}
\begin{theorem}[\texttt{asymGeomSpacing\_sq\_mul\_sq}]\label{Pphi2-asymGeomSpacing-sq-mul-sq}
  \lean{Pphi2.asymGeomSpacing_sq_mul_sq}\leanok
  \uses{Pphi2-asymTimeSpacing, Pphi2-asymSpaceSpacing, Pphi2-asymGeomSpacing}
  ``\texttt{lean theorem asymGeomSpacing\_sq\_mul\_sq (N : ℕ) [NeZero N] : asymGeomSpacing Lt Ls N \^{} 2 * (N : ℝ) \^{} 2 = Lt * Ls }``  Informal statement: Physical volume identity: $a_{\mathrm{geom}}^2 \cdot N^2 = L_t \cdot L_s$.  Proof: Fully proved by unfolding \texttt{asymGeomSpacing}, applying \texttt{Real.sq\_sqrt}, and \texttt{field\_simp}.
\end{theorem}
\begin{theorem}[\texttt{asymTorus\_interacting\_second\_moment\_uniform}]\label{Pphi2-asymTorus-interacting-second-moment-uniform}
  \lean{Pphi2.asymTorus_interacting_second_moment_uniform}\leanok
  \uses{Pphi2-asymTorusInteractingMeasure, GaussianField-NuclearTensorProduct-instModuleReal, GaussianField-NuclearTensorProduct-instTopologicalSpace, GaussianField-FinLatticeSites, GaussianField-NuclearTensorProduct-instIsTopologicalAddGroup, Pphi2-asymGeomSpacing-pos, GaussianField-NuclearTensorProduct-instAddCommGroup, Pphi2-asymGaussian-second-moment-uniform-bound, GaussianField-latticeGaussianMeasure, GaussianField-Configuration, GaussianField-NuclearTensorProduct-instContinuousSMulReal, GaussianField-SmoothMap-Circle, Pphi2-asymLatticeTestFn, GaussianField-instMeasurableSpaceConfiguration, GaussianField-AsymTorusTestFunction, Pphi2-asymTorus-interacting-second-moment-density-transfer, GaussianField-FinLatticeField, Pphi2-asymGeomSpacing}
  ``\texttt{lean theorem asymTorus\_interacting\_second\_moment\_uniform (P : InteractionPolynomial) (mass : ℝ) (hmass : 0 < mass) (f : AsymTorusTestFunction Lt Ls) : ∃ C : ℝ, 0 < C ∧ ∀ (N : ℕ) [NeZero N], ∫ ω, (ω f) \^{} 2 ∂(asymTorusInteractingMeasure Lt Ls N P mass hmass) ≤ C }``  Informal statement: Uniform second moment bound for the interacting measures: $\int (\omega f)^2\, d\mu_{P,N} \le 3\sqrt{K}\, C_g$ for all $N$.  Proof: Fully proved. Pushes through the torus embedding, applies \texttt{density\_transfer\_bound} (Cauchy--Schwarz between Boltzmann weight and observable), then bounds the fourth moment via \texttt{gaussian\_hypercontractive} ($L^4/L^2$ ratio $\le 3$).  Dependencies: \texttt{asymNelson\_exponential\_estimate}, \texttt{asymGaussian\_second\_moment\_uniform\_bound}, \texttt{density\_transfer\_bound}, \texttt{gaussian\_hypercontractive}, \texttt{pairing\_memLp}.
\end{theorem}
\begin{theorem}[\texttt{asymNelson\_exponential\_estimate}]\label{Pphi2-asymNelson-exponential-estimate}
  \lean{Pphi2.asymNelson_exponential_estimate}\leanok
  \uses{GaussianField-FinLatticeSites, Pphi2-asymGeomSpacing-pos, GaussianField-latticeGaussianMeasure, GaussianField-Configuration, Pphi2-asymGeomSpacing-sq-mul-sq, GaussianField-instMeasurableSpaceConfiguration, GaussianField-FinLatticeField, Pphi2-interactionFunctional, Pphi2-nelson-exponential-estimate-master, Pphi2-asymGeomSpacing}
  ``\texttt{lean theorem asymNelson\_exponential\_estimate (P : InteractionPolynomial) (mass : ℝ) (hmass : 0 < mass) : ∃ (K : ℝ), 0 < K ∧ ∀ (N : ℕ) [NeZero N], ∫ ω, (Real.exp (-interactionFunctional 2 N P (asymGeomSpacing Lt Ls N) mass ω)) \^{} 2 ∂(latticeGaussianMeasure ...) ≤ K }``  Informal statement: The $L^2$ norm of the Boltzmann weight is bounded by a constant $K = \exp(2 L_t L_s A)$ uniformly in $N$, where $A$ is the uniform lower bound on Wick polynomials. This follows from $V(\omega) \ge -(L_t L_s A)$ and integrating the pointwise bound over a probability measure.  Proof: Fully proved. Uses \texttt{wickPolynomial\_uniform\_bounded\_below}, the physical volume identity \texttt{asymGeomSpacing\_sq\_mul\_sq}, and \texttt{integral\_mono\_of\_nonneg}.  Dependencies: \texttt{wickPolynomial\_uniform\_bounded\_below}, \texttt{asymGeomSpacing\_sq\_mul\_sq}.
\end{theorem}
\begin{theorem}[\texttt{asymTorus\_interacting\_tightness}]\label{Pphi2-asymTorus-interacting-tightness}
  \lean{Pphi2.asymTorus_interacting_tightness}\leanok
  \uses{Pphi2-asymTorusInteractingMeasure, GaussianField-NuclearTensorProduct-instModuleReal, GaussianField-NuclearTensorProduct-instTopologicalSpace, GaussianField-finLatticeField-dyninMityaginSpace, GaussianField-measure, Pphi2-asymTorus-interacting-second-moment-uniform, GaussianField-ell2-separableSpace, GaussianField-NuclearTensorProduct-dyninMityaginSpace, Pphi2-interactionFunctional-bounded-below, GaussianField-FinLatticeSites, GaussianField-NuclearTensorProduct-instIsTopologicalAddGroup, Pphi2-partitionFunction, GaussianField-latticeCovarianceGJ, Pphi2-interactionFunctional-measurable, GaussianField-pairing-memLp, Pphi2-interactingLatticeMeasure, Pphi2-asymTorusEmbedLift-eval-eq, Pphi2-asymGeomSpacing-pos, GaussianField-NuclearTensorProduct-instAddCommGroup, Pphi2-asymTorusEmbedLift-measurable, GaussianField-latticeGaussianMeasure, GaussianField-configuration-eval-measurable, GaussianField-Configuration, Pphi2-asymTorusInteractingMeasure-isProbability, GaussianField-NuclearTensorProduct-instContinuousSMulReal, GaussianField-SmoothMap-Circle, Pphi2-asymLatticeTestFn, GaussianField-instMeasurableSpaceConfiguration, Pphi2-asymTorusEmbedLift, Pphi2-boltzmannWeight, GaussianField-AsymTorusTestFunction, GaussianField-ell2-, GaussianField-FinLatticeField, Pphi2-interactionFunctional, GaussianField-configuration-tight-of-uniform-second-moments, Pphi2-asymGeomSpacing, Pphi2-partitionFunction-pos, Pphi2-boltzmannWeight-pos}
  ``\texttt{lean theorem asymTorus\_interacting\_tightness (P : InteractionPolynomial) (mass : ℝ) (hmass : 0 < mass) : ∀ (ε : ℝ), 0 < ε → ∃ (K : Set ...), IsCompact K ∧ ∀ (N : ℕ) [NeZero N], ENNReal.ofReal (1 - ε) ≤ (asymTorusInteractingMeasure Lt Ls N P mass hmass) K }``  Informal statement: The family of interacting measures $\{\mu_{P,N}\}_{N \ge 1}$ is tight on the configuration space.  Proof: Fully proved via \texttt{configuration\_tight\_of\_uniform\_second\_moments} (Mitoma--Chebyshev) and \texttt{asymTorus\_interacting\_second\_moment\_uniform}.
\end{theorem}
\begin{theorem}[\texttt{asymGaussian\_second\_moment\_uniform\_bound}]\label{Pphi2-asymGaussian-second-moment-uniform-bound}
  \lean{Pphi2.asymGaussian_second_moment_uniform_bound}\leanok
  \uses{GaussianField-latticeCovariance-GJ-eq-inv-smul-bare, GaussianField-finLatticeField-dyninMityaginSpace, GaussianField-RapidDecaySeq-rapidDecaySeminorm, GaussianField-ell2-separableSpace, GaussianField-RapidDecaySeq, GaussianField-SmoothMap-Circle-fourierBasis, GaussianField-second-moment-eq-covariance, GaussianField-covariance, GaussianField-latticeCovariance, GaussianField-SmoothMap-Circle-instAddCommGroup, GaussianField-FinLatticeSites, Pphi2-asymLatticeTestFn-norm-sq-le-tight, GaussianField-latticeCovarianceGJ, GaussianField-RapidDecaySeq-instSMul, GaussianField-SmoothMap-Circle-instSMul, Pphi2-asymGeomSpacing-pos, GaussianField-latticeGaussianMeasure, GaussianField-Configuration, GaussianField-RapidDecaySeq-instAddCommGroup, GaussianField-SmoothMap-Circle, Pphi2-asymLatticeTestFn, GaussianField-instMeasurableSpaceConfiguration, GaussianField-SmoothMap-Circle-sobolevSeminorm, Pphi2-covariance-inner-le-mass-inv-sq-norm-sq, GaussianField-AsymTorusTestFunction, GaussianField-ell2-, GaussianField-FinLatticeField, GaussianField-SmoothMap-Circle-sobolevSeminorm-fourierBasis-le, Pphi2-asymGeomSpacing}
  ``\texttt{lean theorem asymGaussian\_second\_moment\_uniform\_bound (mass : ℝ) (hmass : 0 < mass) (f : AsymTorusTestFunction Lt Ls) : ∃ Cg : ℝ, 0 < Cg ∧ ∀ (N : ℕ) [NeZero N], ∫ ω, (ω (asymLatticeTestFn Lt Ls N f)) \^{} 2 ∂(latticeGaussianMeasure ...) ≤ Cg }``  Informal statement: The Gaussian second moment $\int (\omega g_f)^2\, d\mu_{\mathrm{GFF}}$ is bounded by $m^{-2} C_f$ uniformly in $N$.  Proof: Fully proved using \texttt{second\_moment\_eq\_covariance}, \texttt{covariance\_inner\_le\_mass\_inv\_sq\_norm\_sq}, and \texttt{asymLatticeTestFn\_norm\_sq\_bounded}.
\end{theorem}
\begin{theorem}[\texttt{asymLatticeTestFn\_norm\_sq\_le\_tight}]\label{Pphi2-asymLatticeTestFn-norm-sq-le-tight}
  \lean{Pphi2.asymLatticeTestFn_norm_sq_le_tight}\leanok
  \uses{GaussianField-NuclearTensorProduct-instModuleReal, Pphi2-dm-basis-eq-fourierBasis, GaussianField-NuclearTensorProduct-instTopologicalSpace, GaussianField-DyninMityaginSpace-coeff, GaussianField-SmoothMap-Circle-instFunLike, GaussianField-RapidDecaySeq-rapidDecaySeminorm, Pphi2-evalAsymAtFinSiteGJ-apply, GaussianField-RapidDecaySeq, GaussianField-SmoothMap-Circle-fourierBasis, GaussianField-RapidDecaySeq-val, GaussianField-NuclearTensorProduct-dyninMityaginSpace, GaussianField-SmoothMap-Circle-instAddCommGroup, GaussianField-FinLatticeSites, GaussianField-NuclearTensorProduct-instIsTopologicalAddGroup, GaussianField-SmoothMap-Circle-norm-iteratedDeriv-le-sobolevSeminorm-, GaussianField-circleRestriction, GaussianField-DyninMityaginSpace-expansion, GaussianField-RapidDecaySeq-instSMul, GaussianField-SmoothMap-Circle-instSMul, GaussianField-evalAsymTorusAtSite, GaussianField-NuclearTensorProduct, GaussianField-RapidDecaySeq-basisVec, GaussianField-circleSpacing-eq, Pphi2-asymGeomSpacing-pos, GaussianField-NuclearTensorProduct-instAddCommGroup, Pphi2-evalAsymAtFinSite, GaussianField-circleRestriction-apply, Pphi2-evalAsymAtFinSiteGJ, GaussianField-SmoothMap-Circle-instModule, GaussianField-SmoothMap-Circle-instContinuousSMul, GaussianField-smoothCircle-dyninMityaginSpace, GaussianField-SmoothMap-Circle-instIsTopologicalAddGroup, GaussianField-DyninMityaginSpace-basis, GaussianField-NuclearTensorProduct-instContinuousSMulReal, GaussianField-NuclearTensorProduct-pure, GaussianField-RapidDecaySeq-instAddCommGroup, GaussianField-SmoothMap-Circle, GaussianField-circleSpacing, Pphi2-asymLatticeTestFn, GaussianField-SmoothMap-Circle-sobolevSeminorm, GaussianField-AsymTorusTestFunction, GaussianField-NuclearTensorProduct-evalCLM-pure, GaussianField-RapidDecaySeq-rapid-decay, GaussianField-NuclearTensorProduct-evalCLM, GaussianField-SmoothMap-Circle-instTopologicalSpace, GaussianField-circlePoint, Pphi2-asymGeomSpacing, Lattice-toSemilatticeInf}
  \texttt{Pphi2.asymLatticeTestFn\_norm\_sq\_le\_tight}
\end{theorem}
\begin{theorem}[\texttt{asymLatticeTestFn\_norm\_sq\_bounded}]\label{Pphi2-asymLatticeTestFn-norm-sq-bounded}
  \lean{Pphi2.asymLatticeTestFn_norm_sq_bounded}\leanok
  \uses{GaussianField-RapidDecaySeq-rapidDecaySeminorm, Pphi2-asymLatticeTestFn-norm-sq-le, GaussianField-RapidDecaySeq, GaussianField-SmoothMap-Circle-fourierBasis, GaussianField-SmoothMap-Circle-instAddCommGroup, GaussianField-FinLatticeSites, GaussianField-RapidDecaySeq-instSMul, GaussianField-SmoothMap-Circle-instSMul, GaussianField-RapidDecaySeq-instAddCommGroup, GaussianField-SmoothMap-Circle, Pphi2-asymLatticeTestFn, GaussianField-SmoothMap-Circle-sobolevSeminorm, GaussianField-AsymTorusTestFunction, GaussianField-SmoothMap-Circle-sobolevSeminorm-fourierBasis-le, Lattice-toSemilatticeInf}
  ``\texttt{lean theorem asymLatticeTestFn\_norm\_sq\_bounded (f : AsymTorusTestFunction Lt Ls) : ∃ C : ℝ, 0 < C ∧ ∀ (N : ℕ) [NeZero N], ∑ x, (asymLatticeTestFn Lt Ls N f x) \^{} 2 ≤ C }`\texttt{  Informal statement: Existential wrapper for the Riemann sum bound, hiding the Sobolev constants.  Proof: Fully proved by instantiating }asymLatticeTestFn\_norm\_sq\_le\texttt{ with constants from }sobolevSeminorm\_fourierBasis\_le`.
\end{theorem}
\begin{theorem}[\texttt{asymTorusInteractingMeasure\_isProbability}]\label{Pphi2-asymTorusInteractingMeasure-isProbability}
  \lean{Pphi2.asymTorusInteractingMeasure_isProbability}\leanok
  \uses{Pphi2-asymTorusInteractingMeasure, GaussianField-NuclearTensorProduct-instModuleReal, GaussianField-NuclearTensorProduct-instTopologicalSpace, GaussianField-FinLatticeSites, GaussianField-NuclearTensorProduct-instIsTopologicalAddGroup, Pphi2-interactingLatticeMeasure-isProbability, Pphi2-interactingLatticeMeasure, Pphi2-asymGeomSpacing-pos, GaussianField-NuclearTensorProduct-instAddCommGroup, Pphi2-asymTorusEmbedLift-measurable, GaussianField-Configuration, GaussianField-NuclearTensorProduct-instContinuousSMulReal, GaussianField-SmoothMap-Circle, GaussianField-instMeasurableSpaceConfiguration, Pphi2-asymTorusEmbedLift, GaussianField-AsymTorusTestFunction, GaussianField-FinLatticeField, Pphi2-asymGeomSpacing}
  ``\texttt{lean instance asymTorusInteractingMeasure\_isProbability (N : ℕ) [NeZero N] (P : InteractionPolynomial) (mass : ℝ) (hmass : 0 < mass) : IsProbabilityMeasure (asymTorusInteractingMeasure Lt Ls N P mass hmass) }`\texttt{  Informal statement: The asymmetric torus interacting measure is a probability measure.  Proof: Fully proved via }Measure.isProbabilityMeasure\_map`.
\end{theorem}
\begin{theorem}[\texttt{asymLatticeTestFn\_norm\_sq\_le}]\label{Pphi2-asymLatticeTestFn-norm-sq-le}
  \lean{Pphi2.asymLatticeTestFn_norm_sq_le}\leanok
  \uses{GaussianField-RapidDecaySeq-rapidDecaySeminorm, GaussianField-RapidDecaySeq, GaussianField-SmoothMap-Circle-fourierBasis, GaussianField-SmoothMap-Circle-instAddCommGroup, GaussianField-FinLatticeSites, Pphi2-asymLatticeTestFn-norm-sq-le-tight, GaussianField-RapidDecaySeq-instSMul, GaussianField-SmoothMap-Circle-instSMul, GaussianField-RapidDecaySeq-instAddCommGroup, GaussianField-SmoothMap-Circle, Pphi2-asymGeomSpacing-sq-mul-sq, Pphi2-asymLatticeTestFn, GaussianField-SmoothMap-Circle-sobolevSeminorm, GaussianField-AsymTorusTestFunction, Pphi2-asymGeomSpacing}
  ``\texttt{lean theorem asymLatticeTestFn\_norm\_sq\_le (C₀t : ℝ) (hC₀t\_pos : 0 < C₀t) (hC₀t : ...) (C₀s : ℝ) (hC₀s\_pos : 0 < C₀s) (hC₀s : ...) (f : AsymTorusTestFunction Lt Ls) (N : ℕ) [NeZero N] : ∑ x : FinLatticeSites 2 N, (asymLatticeTestFn Lt Ls N f x) \^{} 2 ≤ Lt * Ls * C₀t \^{} 2 * C₀s \^{} 2 * (RapidDecaySeq.rapidDecaySeminorm 0 f) \^{} 2 + 1 }``  Informal statement: The Riemann sum $\sum_x g_f(x)^2$ is bounded by $L_t L_s C_{0t}^2 C_{0s}^2 (p_0 f)^2 + 1$ uniformly in $N$, where $C_{0t}, C_{0s}$ are Sobolev constants and $p_0$ is the rapid-decay seminorm of order 0.  Proof: Fully proved. Bounds $|\mathrm{eval}_x(\mathrm{basisVec}\, m)| \le \sqrt{L_t/N}\, C_{0t} \cdot \sqrt{L_s/N}\, C_{0s}$ via circle restriction and Sobolev embedding, extends to general $f$ via Dynin--Mityagin expansion, then sums over $N^2$ lattice sites with $N^2 \cdot B^2 = L_t L_s C_{0t}^2 C_{0s}^2$.  Dependencies: \texttt{DyninMityaginSpace.expansion}, \texttt{circleRestriction\_apply}, \texttt{SmoothMap\_Circle.norm\_iteratedDeriv\_le\_sobolevSeminorm'}, \texttt{norm\_tsum\_le\_tsum\_norm}.
\end{theorem}
\begin{theorem}[\texttt{asymTorus\_interacting\_second\_moment\_density\_transfer}]\label{Pphi2-asymTorus-interacting-second-moment-density-transfer}
  \lean{Pphi2.asymTorus_interacting_second_moment_density_transfer}\leanok
  \uses{Pphi2-asymTorusInteractingMeasure, GaussianField-NuclearTensorProduct-instModuleReal, Pphi2-partitionFunction-ge-one, GaussianField-NuclearTensorProduct-instTopologicalSpace, GaussianField-finLatticeField-dyninMityaginSpace, GaussianField-measure, GaussianField-ell2-separableSpace, GaussianField-FinLatticeSites, GaussianField-NuclearTensorProduct-instIsTopologicalAddGroup, Pphi2-partitionFunction, GaussianField-latticeCovarianceGJ, GaussianField-pairing-memLp, Pphi2-interactingLatticeMeasure, Pphi2-asymTorusEmbedLift-eval-eq, Pphi2-asymGeomSpacing-pos, GaussianField-NuclearTensorProduct-instAddCommGroup, Pphi2-asymTorusEmbedLift-measurable, GaussianField-latticeGaussianMeasure, GaussianField-configuration-eval-measurable, GaussianField-Configuration, Pphi2-density-transfer-bound, GaussianField-NuclearTensorProduct-instContinuousSMulReal, GaussianField-gaussian-hypercontractive, GaussianField-SmoothMap-Circle, Pphi2-asymLatticeTestFn, GaussianField-instMeasurableSpaceConfiguration, Pphi2-asymTorusEmbedLift, GaussianField-AsymTorusTestFunction, GaussianField-ell2-, GaussianField-FinLatticeField, Pphi2-interactionFunctional, Pphi2-asymGeomSpacing, Lattice-toSemilatticeInf, Pphi2-asymNelson-exponential-estimate}
  \texttt{Pphi2.asymTorus\_interacting\_second\_moment\_density\_transfer}
\end{theorem}
\begin{theorem}[\texttt{congr\_simp}]\label{Pphi2-asymTorusEmbedLift-congr-simp}
  \lean{Pphi2.asymTorusEmbedLift.congr_simp}\leanok
  \uses{GaussianField-NuclearTensorProduct-instModuleReal, GaussianField-NuclearTensorProduct-instTopologicalSpace, GaussianField-FinLatticeSites, GaussianField-NuclearTensorProduct-instIsTopologicalAddGroup, GaussianField-NuclearTensorProduct-instAddCommGroup, GaussianField-Configuration, GaussianField-NuclearTensorProduct-instContinuousSMulReal, GaussianField-SmoothMap-Circle, Pphi2-asymTorusEmbedLift, GaussianField-AsymTorusTestFunction, GaussianField-FinLatticeField}
  \texttt{Pphi2.asymTorusEmbedLift.congr\_simp}
\end{theorem}
\begin{theorem}[\texttt{asymTorusInteractingLimit\_exists}]\label{Pphi2-asymTorusInteractingLimit-exists}
  \lean{Pphi2.asymTorusInteractingLimit_exists}\leanok
  \uses{Pphi2-asymTorusInteractingMeasure, GaussianField-NuclearTensorProduct-instModuleReal, GaussianField-NuclearTensorProduct-instTopologicalSpace, GaussianField-NuclearTensorProduct-dyninMityaginSpace, GaussianField-NuclearTensorProduct-instIsTopologicalAddGroup, Pphi2-asymTorus-interacting-tightness, GaussianField-NuclearTensorProduct-instAddCommGroup, GaussianField-prokhorov-configuration, GaussianField-Configuration, Pphi2-asymTorusInteractingMeasure-isProbability, GaussianField-NuclearTensorProduct-instContinuousSMulReal, GaussianField-SmoothMap-Circle, GaussianField-instMeasurableSpaceConfiguration, GaussianField-AsymTorusTestFunction}
  ``\texttt{lean theorem asymTorusInteractingLimit\_exists (P : InteractionPolynomial) (mass : ℝ) (hmass : 0 < mass) : ∃ (μ : Measure ...) (\_ : IsProbabilityMeasure μ) (φ : ℕ → ℕ) (\_ : StrictMono φ), ∀ (f : Configuration ... → ℝ), Continuous f → (∃ C, ∀ x, |f x| ≤ C) → Tendsto (fun n => ∫ ω, f ω ∂(asymTorusInteractingMeasure Lt Ls (φ n + 1) P mass hmass)) atTop (nhds (∫ ω, f ω ∂μ)) }``  Informal statement: There exists a subsequence along which the interacting measures on the asymmetric torus converge weakly to a probability measure. This is the UV limit of Route B' ($N \to \infty$ with $L_t, L_s$ fixed).  Proof: Fully proved via \texttt{prokhorov\_configuration} applied with tightness from \texttt{asymTorus\_interacting\_tightness}.  Dependencies: \texttt{asymTorus\_interacting\_tightness}, \texttt{prokhorov\_configuration}.
\end{theorem}
\section{\texttt{Pphi2.AsymTorus.AsymContinuumLimit}}
\begin{theorem}[\texttt{asymTorusIso\_cylinderUniformGreenBound}]\label{Pphi2-asymTorusIso-cylinderUniformGreenBound}
  \lean{Pphi2.asymTorusIso_cylinderUniformGreenBound}\leanok
  \uses{GaussianField-NuclearTensorProduct-instModuleReal, GaussianField-NuclearTensorProduct-instTopologicalSpace, GaussianField-AsymLatticeSites, GaussianField-NuclearTensorProduct-instIsTopologicalAddGroup, Pphi2-AsymTorusSequenceHasUniformGreenMomentBound, Pphi2-asymLatticeTestFnIso, GaussianField-NuclearTensorProduct-instAddCommGroup, Pphi2-asymTorusInteractingMeasureIso, GaussianField-Configuration, GaussianField-NuclearTensorProduct-instContinuousSMulReal, GaussianField-AsymLatticeField, GaussianField-SmoothMap-Circle, Pphi2-latticeGaussianMeasureAsym, GaussianField-instMeasurableSpaceConfiguration, Pphi2-MeasureHasGreenMomentBound, Pphi2-asymTorusIso-measureHasGreenMomentBound-of-cutoff, Pphi2-AsymTorusSequenceHasUniformGreenMomentBound-of-forall, GaussianField-AsymTorusTestFunction}
  \texttt{Pphi2.asymTorusIso\_cylinderUniformGreenBound}
\end{theorem}
\begin{theorem}[\texttt{asymTorusIso\_measureHasGreenMomentBound\_of\_nelson}]\label{Pphi2-asymTorusIso-measureHasGreenMomentBound-of-nelson}
  \lean{Pphi2.asymTorusIso_measureHasGreenMomentBound_of_nelson}\leanok
  \uses{GaussianField-NuclearTensorProduct-instModuleReal, GaussianField-NuclearTensorProduct-instTopologicalSpace, GaussianField-AsymLatticeSites, GaussianField-NuclearTensorProduct-instIsTopologicalAddGroup, Pphi2-asymLatticeTestFnIso, Pphi2-asymTorusInteractingMeasureIso-exponentialMomentBound-cutoff-of-nelson, GaussianField-NuclearTensorProduct-instAddCommGroup, Pphi2-asymTorusInteractingMeasureIso, GaussianField-Configuration, GaussianField-NuclearTensorProduct-instContinuousSMulReal, GaussianField-AsymLatticeField, GaussianField-SmoothMap-Circle, Pphi2-latticeGaussianMeasureAsym, Pphi2-interactionFunctionalAsym, GaussianField-instMeasurableSpaceConfiguration, Pphi2-MeasureHasGreenMomentBound, Pphi2-asymTorusIso-measureHasGreenMomentBound-of-cutoff, GaussianField-AsymTorusTestFunction}
  \texttt{Pphi2.asymTorusIso\_measureHasGreenMomentBound\_of\_nelson}
\end{theorem}
\begin{theorem}[\texttt{congr\_simp}]\label{Pphi2-asymTorusEmbedLiftIso-congr-simp}
  \lean{Pphi2.asymTorusEmbedLiftIso.congr_simp}\leanok
  \uses{GaussianField-NuclearTensorProduct-instModuleReal, GaussianField-NuclearTensorProduct-instTopologicalSpace, GaussianField-AsymLatticeSites, GaussianField-NuclearTensorProduct-instIsTopologicalAddGroup, Pphi2-asymTorusEmbedLiftIso, GaussianField-NuclearTensorProduct-instAddCommGroup, GaussianField-Configuration, GaussianField-NuclearTensorProduct-instContinuousSMulReal, GaussianField-AsymLatticeField, GaussianField-SmoothMap-Circle, GaussianField-AsymTorusTestFunction}
  \texttt{Pphi2.asymTorusEmbedLiftIso.congr\_simp}
\end{theorem}
\begin{theorem}[\texttt{congr\_simp}]\label{Pphi2-asymTorusInteractingMeasureIso-congr-simp}
  \lean{Pphi2.asymTorusInteractingMeasureIso.congr_simp}\leanok
  \uses{GaussianField-NuclearTensorProduct-instModuleReal, GaussianField-NuclearTensorProduct-instTopologicalSpace, GaussianField-NuclearTensorProduct-instIsTopologicalAddGroup, GaussianField-NuclearTensorProduct-instAddCommGroup, Pphi2-asymTorusInteractingMeasureIso, GaussianField-Configuration, GaussianField-NuclearTensorProduct-instContinuousSMulReal, GaussianField-SmoothMap-Circle, GaussianField-instMeasurableSpaceConfiguration, GaussianField-AsymTorusTestFunction}
  \texttt{Pphi2.asymTorusInteractingMeasureIso.congr\_simp}
\end{theorem}
\begin{theorem}[\texttt{asymTorusIso\_interacting\_limit\_exists}]\label{Pphi2-asymTorusIso-interacting-limit-exists}
  \lean{Pphi2.asymTorusIso_interacting_limit_exists}\leanok
  \uses{Pphi2-asymTorusInteractingMeasureIso-sq-integrable, GaussianField-NuclearTensorProduct-instModuleReal, GaussianField-NuclearTensorProduct-instTopologicalSpace, Pphi2-asymTorusInteractingMeasureIso-isProbability, Pphi2-asymTorusIso-interacting-second-moment-uniform, GaussianField-NuclearTensorProduct-dyninMityaginSpace, GaussianField-NuclearTensorProduct-instIsTopologicalAddGroup, GaussianField-NuclearTensorProduct-instAddCommGroup, GaussianField-prokhorov-configuration, Pphi2-asymTorusInteractingMeasureIso, GaussianField-Configuration, GaussianField-NuclearTensorProduct-instContinuousSMulReal, GaussianField-SmoothMap-Circle, GaussianField-instMeasurableSpaceConfiguration, GaussianField-AsymTorusTestFunction, GaussianField-configuration-tight-of-uniform-second-moments}
  \texttt{Pphi2.asymTorusIso\_interacting\_limit\_exists}
\end{theorem}
\begin{theorem}[\texttt{asymTorusInteractingMeasureIso\_sq\_integrable}]\label{Pphi2-asymTorusInteractingMeasureIso-sq-integrable}
  \lean{Pphi2.asymTorusInteractingMeasureIso_sq_integrable}\leanok
  \uses{GaussianField-NuclearTensorProduct-instModuleReal, GaussianField-NuclearTensorProduct-instTopologicalSpace, GaussianField-AsymLatticeSites, GaussianField-measure, Pphi2-asymTorusEmbedLiftIso-measurable, GaussianField-ell2-separableSpace, Pphi2-partitionFunctionAsym-pos, GaussianField-asymLatticeField-dyninMityaginSpace, Pphi2-asymTorusEmbedLiftIso-eval-eq, GaussianField-NuclearTensorProduct-instIsTopologicalAddGroup, Pphi2-asymLatticeTestFnIso, Pphi2-boltzmannWeightAsym, GaussianField-pairing-memLp, Pphi2-boltzmannWeightAsym-pos, Pphi2-asymTorusEmbedLiftIso, Pphi2-interactionFunctionalAsym-bounded-below, GaussianField-NuclearTensorProduct-instAddCommGroup, Pphi2-partitionFunctionAsym, Pphi2-asymTorusInteractingMeasureIso, GaussianField-configuration-eval-measurable, GaussianField-Configuration, GaussianField-NuclearTensorProduct-instContinuousSMulReal, GaussianField-AsymLatticeField, GaussianField-SmoothMap-Circle, Pphi2-latticeGaussianMeasureAsym, Pphi2-interactingLatticeMeasureAsym, Pphi2-interactionFunctionalAsym, GaussianField-instMeasurableSpaceConfiguration, GaussianField-latticeCovarianceAsymGJ, GaussianField-AsymTorusTestFunction, GaussianField-ell2-, Pphi2-interactionFunctionalAsym-measurable}
  \texttt{Pphi2.asymTorusInteractingMeasureIso\_sq\_integrable}
\end{theorem}
\begin{theorem}[\texttt{asymTorusIso\_interacting\_second\_moment\_uniform}]\label{Pphi2-asymTorusIso-interacting-second-moment-uniform}
  \lean{Pphi2.asymTorusIso_interacting_second_moment_uniform}\leanok
  \uses{GaussianField-NuclearTensorProduct-instModuleReal, GaussianField-NuclearTensorProduct-instTopologicalSpace, GaussianField-AsymLatticeSites, GaussianField-ell2-separableSpace, Pphi2-asymTorusContinuumGreen, GaussianField-second-moment-eq-covariance, GaussianField-covariance, GaussianField-asymLatticeField-dyninMityaginSpace, GaussianField-NuclearTensorProduct-instIsTopologicalAddGroup, Pphi2-asymLatticeTestFnIso, GaussianField-NuclearTensorProduct-instAddCommGroup, Pphi2-asymTorusInteractingMeasureIso, GaussianField-Configuration, GaussianField-NuclearTensorProduct-instContinuousSMulReal, GaussianField-AsymLatticeField, GaussianField-SmoothMap-Circle, Pphi2-latticeGaussianMeasureAsym, Pphi2-second-moment-asym-tendsto, GaussianField-instMeasurableSpaceConfiguration, GaussianField-latticeCovarianceAsymGJ, GaussianField-AsymTorusTestFunction, GaussianField-ell2-, Pphi2-asymTorusIso-interacting-second-moment-density-transfer, Lattice-toSemilatticeInf}
  \texttt{Pphi2.asymTorusIso\_interacting\_second\_moment\_uniform}
\end{theorem}
\begin{theorem}[\texttt{asymTorusIso\_measureHasGreenMomentBound\_of\_cutoff}]\label{Pphi2-asymTorusIso-measureHasGreenMomentBound-of-cutoff}
  \lean{Pphi2.asymTorusIso_measureHasGreenMomentBound_of_cutoff}\leanok
  \uses{GaussianField-NuclearTensorProduct-instModuleReal, GaussianField-weakLimit-exponential-moment, GaussianField-NuclearTensorProduct-instTopologicalSpace, GaussianField-AsymLatticeSites, GaussianField-asymTorusContinuumGreen, GaussianField-ell2-separableSpace, Pphi2-asymTorusInteractingMeasureIso-isProbability, Pphi2-asymTorusContinuumGreen, GaussianField-second-moment-eq-covariance, GaussianField-covariance, GaussianField-asymLatticeField-dyninMityaginSpace, GaussianField-NuclearTensorProduct-instIsTopologicalAddGroup, Pphi2-asymLatticeTestFnIso, GaussianField-NuclearTensorProduct-instAddCommGroup, Pphi2-asymTorusInteractingMeasureIso, GaussianField-Configuration, GaussianField-NuclearTensorProduct-instContinuousSMulReal, GaussianField-AsymLatticeField, GaussianField-SmoothMap-Circle, Pphi2-latticeGaussianMeasureAsym, Pphi2-second-moment-asym-tendsto, GaussianField-instMeasurableSpaceConfiguration, Pphi2-asymTorusIso-interacting-limit-exists, GaussianField-latticeCovarianceAsymGJ, Pphi2-MeasureHasGreenMomentBound, GaussianField-AsymTorusTestFunction, GaussianField-ell2-}
  \texttt{Pphi2.asymTorusIso\_measureHasGreenMomentBound\_of\_cutoff}
\end{theorem}
\begin{theorem}[\texttt{asymTorusIso\_interacting\_second\_moment\_density\_transfer}]\label{Pphi2-asymTorusIso-interacting-second-moment-density-transfer}
  \lean{Pphi2.asymTorusIso_interacting_second_moment_density_transfer}\leanok
  \uses{GaussianField-NuclearTensorProduct-instModuleReal, GaussianField-NuclearTensorProduct-instTopologicalSpace, Pphi2-partitionFunctionAsym-ge-one, GaussianField-AsymLatticeSites, GaussianField-measure, Pphi2-asymTorusEmbedLiftIso-measurable, GaussianField-ell2-separableSpace, GaussianField-asymLatticeField-dyninMityaginSpace, Pphi2-asymTorusEmbedLiftIso-eval-eq, GaussianField-NuclearTensorProduct-instIsTopologicalAddGroup, Pphi2-asymLatticeTestFnIso, Pphi2-asymNelson-exponential-estimate-iso, GaussianField-pairing-memLp, Pphi2-asymTorusEmbedLiftIso, GaussianField-NuclearTensorProduct-instAddCommGroup, Pphi2-partitionFunctionAsym, Pphi2-asymTorusInteractingMeasureIso, GaussianField-configuration-eval-measurable, GaussianField-Configuration, Pphi2-density-transfer-bound-iso, GaussianField-NuclearTensorProduct-instContinuousSMulReal, GaussianField-gaussian-hypercontractive, GaussianField-AsymLatticeField, GaussianField-SmoothMap-Circle, Pphi2-latticeGaussianMeasureAsym, Pphi2-interactingLatticeMeasureAsym, Pphi2-interactionFunctionalAsym, GaussianField-instMeasurableSpaceConfiguration, GaussianField-latticeCovarianceAsymGJ, GaussianField-AsymTorusTestFunction, GaussianField-ell2-, Lattice-toSemilatticeInf}
  \texttt{Pphi2.asymTorusIso\_interacting\_second\_moment\_density\_transfer}
\end{theorem}
\begin{theorem}[\texttt{cylinderIso\_OS\_of\_RP\_OS2}]\label{Pphi2-cylinderIso-OS-of-RP-OS2}
  \lean{Pphi2.cylinderIso_OS_of_RP_OS2}\leanok
  \uses{GaussianField-NuclearTensorProduct-instModuleReal, GaussianField-NuclearTensorProduct-instTopologicalSpace, GaussianField-cylinderPositiveTimeSubmodule, Pphi2-routeBPrimeIso-cylinder-OS, GaussianField-AsymLatticeSites, GaussianField-cylinderSpatialTranslation, GaussianField-CylinderTestFunction, GaussianField-NuclearTensorProduct-instIsTopologicalAddGroup, Pphi2-asymLatticeTestFnIso, GaussianField-NuclearTensorProduct-instAddCommGroup, Pphi2-AsymTorusSequenceHasCylinderOS2Symmetry, GaussianField-cylinderTimeReflection, Pphi2-cylinderPullbackMeasure, Pphi2-asymTorusInteractingMeasureIso, GaussianField-Configuration, GaussianField-NuclearTensorProduct-instContinuousSMulReal, GaussianField-AsymLatticeField, GaussianField-SmoothMap-Circle, Pphi2-latticeGaussianMeasureAsym, GaussianField-instMeasurableSpaceConfiguration, Pphi2-asymInteracting-expMoment-volume-uniform, GaussianField-AsymTorusTestFunction, Pphi2-CylinderMeasureSequenceEventuallyReflectionPositive, GaussianField-cylinderTranslation}
  \texttt{Pphi2.cylinderIso\_OS\_of\_RP\_OS2}
\end{theorem}
\begin{theorem}[\texttt{asymTorusIso\_measureHasGreenMomentBound}]\label{Pphi2-asymTorusIso-measureHasGreenMomentBound}
  \lean{Pphi2.asymTorusIso_measureHasGreenMomentBound}\leanok
  \uses{GaussianField-NuclearTensorProduct-instModuleReal, GaussianField-NuclearTensorProduct-instTopologicalSpace, GaussianField-AsymLatticeSites, Pphi2-asymTorusIso-measureHasGreenMomentBound-of-nelson, GaussianField-NuclearTensorProduct-instIsTopologicalAddGroup, Pphi2-asymNelson-exponential-estimate-iso, GaussianField-NuclearTensorProduct-instAddCommGroup, GaussianField-Configuration, GaussianField-NuclearTensorProduct-instContinuousSMulReal, GaussianField-AsymLatticeField, GaussianField-SmoothMap-Circle, Pphi2-latticeGaussianMeasureAsym, Pphi2-interactionFunctionalAsym, GaussianField-instMeasurableSpaceConfiguration, Pphi2-MeasureHasGreenMomentBound, GaussianField-AsymTorusTestFunction}
  \texttt{Pphi2.asymTorusIso\_measureHasGreenMomentBound}
\end{theorem}
\begin{theorem}[\texttt{asymTorusInteractingMeasureIso\_isProbability}]\label{Pphi2-asymTorusInteractingMeasureIso-isProbability}
  \lean{Pphi2.asymTorusInteractingMeasureIso_isProbability}\leanok
  \uses{GaussianField-NuclearTensorProduct-instModuleReal, GaussianField-NuclearTensorProduct-instTopologicalSpace, GaussianField-AsymLatticeSites, Pphi2-asymTorusEmbedLiftIso-measurable, GaussianField-NuclearTensorProduct-instIsTopologicalAddGroup, Pphi2-interactingLatticeMeasureAsym-isProbability, Pphi2-asymTorusEmbedLiftIso, GaussianField-NuclearTensorProduct-instAddCommGroup, Pphi2-asymTorusInteractingMeasureIso, GaussianField-Configuration, GaussianField-NuclearTensorProduct-instContinuousSMulReal, GaussianField-AsymLatticeField, GaussianField-SmoothMap-Circle, Pphi2-interactingLatticeMeasureAsym, GaussianField-instMeasurableSpaceConfiguration, GaussianField-AsymTorusTestFunction}
  \texttt{Pphi2.asymTorusInteractingMeasureIso\_isProbability}
\end{theorem}
\begin{theorem}[\texttt{asymInteracting\_expMoment\_volume\_uniform} (axiom)]\label{Pphi2-asymInteracting-expMoment-volume-uniform}
  \lean{Pphi2.asymInteracting_expMoment_volume_uniform}\leanok
  \uses{GaussianField-NuclearTensorProduct-instModuleReal, GaussianField-NuclearTensorProduct-instTopologicalSpace, GaussianField-AsymLatticeSites, GaussianField-NuclearTensorProduct-instIsTopologicalAddGroup, Pphi2-asymLatticeTestFnIso, GaussianField-NuclearTensorProduct-instAddCommGroup, Pphi2-asymTorusInteractingMeasureIso, GaussianField-Configuration, GaussianField-NuclearTensorProduct-instContinuousSMulReal, GaussianField-AsymLatticeField, GaussianField-SmoothMap-Circle, Pphi2-latticeGaussianMeasureAsym, GaussianField-instMeasurableSpaceConfiguration, GaussianField-AsymTorusTestFunction}
  \texttt{Pphi2.asymInteracting\_expMoment\_volume\_uniform}
\end{theorem}
\begin{theorem}[\texttt{routeBPrimeIso\_cylinder\_OS}]\label{Pphi2-routeBPrimeIso-cylinder-OS}
  \lean{Pphi2.routeBPrimeIso_cylinder_OS}\leanok
  \uses{GaussianField-NuclearTensorProduct-instModuleReal, GaussianField-NuclearTensorProduct-instTopologicalSpace, GaussianField-cylinderPositiveTimeSubmodule, GaussianField-AsymLatticeSites, GaussianField-cylinderSpatialTranslation, GaussianField-CylinderTestFunction, GaussianField-NuclearTensorProduct-instIsTopologicalAddGroup, Pphi2-AsymTorusSequenceHasUniformGreenMomentBound, Pphi2-asymLatticeTestFnIso, GaussianField-NuclearTensorProduct-instAddCommGroup, Pphi2-AsymTorusSequenceHasCylinderOS2Symmetry, GaussianField-cylinderTimeReflection, Pphi2-cylinderPullbackMeasure, Pphi2-asymTorusInteractingMeasureIso, GaussianField-Configuration, GaussianField-NuclearTensorProduct-instContinuousSMulReal, GaussianField-AsymLatticeField, GaussianField-SmoothMap-Circle, Pphi2-latticeGaussianMeasureAsym, Pphi2-asymTorusIso-cylinderUniformGreenBound, Pphi2-routeBPrime-cylinder-OS, GaussianField-instMeasurableSpaceConfiguration, GaussianField-AsymTorusTestFunction, Pphi2-CylinderMeasureSequenceEventuallyReflectionPositive, GaussianField-cylinderTranslation}
  \texttt{Pphi2.routeBPrimeIso\_cylinder\_OS}
\end{theorem}
\section{\texttt{Pphi2.InteractingMeasure.InteractionFourPoint}}
\begin{theorem}[\texttt{integrable\_wickFourth\_wickPolynomial}]\label{Pphi2-integrable-wickFourth-wickPolynomial}
  \lean{Pphi2.integrable_wickFourth_wickPolynomial}\leanok
  \uses{GaussianField-gffSmearedCovariance-self, Pphi2-gffSmearedCovariance-single-single-eq-wickConstant, GaussianField-integrable-wickMonomial-smeared-mul, GaussianField-FinLatticeSites, GaussianField-gffSmearedCoeff, Pphi2-wickMonomial, Pphi2-wickMonomial-eq-root-local, GaussianField-latticeGaussianMeasure, GaussianField-Configuration, Pphi2-wickConstant, GaussianField-instMeasurableSpaceConfiguration, GaussianField-FinLatticeField, GaussianField-gffSmearedCovariance, Pphi2-wickPolynomial}
  \texttt{Pphi2.integrable\_wickFourth\_wickPolynomial}
\end{theorem}
\begin{theorem}[\texttt{wickFourth\_wickPolynomial\_inner\_quartic}]\label{Pphi2-wickFourth-wickPolynomial-inner-quartic}
  \lean{Pphi2.wickFourth_wickPolynomial_inner_quartic}\leanok
  \uses{GaussianField-FinLatticeSites, GaussianField-gffPositionCovariance, GaussianField-latticeGaussianMeasure, GaussianField-Configuration, Pphi2-wickConstant, GaussianField-instMeasurableSpaceConfiguration, Pphi2-wickFourth-wickPolynomial-inner, GaussianField-FinLatticeField, GaussianField-gffSmearedCovariance, Pphi2-wickPolynomial}
  \texttt{Pphi2.wickFourth\_wickPolynomial\_inner\_quartic}
\end{theorem}
\begin{theorem}[\texttt{wickFourth\_wickPolynomial\_inner}]\label{Pphi2-wickFourth-wickPolynomial-inner}
  \lean{Pphi2.wickFourth_wickPolynomial_inner}\leanok
  \uses{GaussianField-gffSmearedCovariance-self, Pphi2-gffSmearedCovariance-single-single-eq-wickConstant, Pphi2-wickFourth-smeared-site-inner, GaussianField-integrable-wickMonomial-smeared-mul, GaussianField-FinLatticeSites, GaussianField-gffSmearedCoeff, GaussianField-gffPositionCovariance, Pphi2-wickMonomial, Pphi2-wickMonomial-eq-root-local, GaussianField-latticeGaussianMeasure, GaussianField-Configuration, Pphi2-wickConstant, GaussianField-instMeasurableSpaceConfiguration, GaussianField-FinLatticeField, GaussianField-gffSmearedCovariance, Pphi2-wickPolynomial}
  \texttt{Pphi2.wickFourth\_wickPolynomial\_inner}
\end{theorem}
\begin{theorem}[\texttt{wickFourth\_interaction\_inner\_quartic}]\label{Pphi2-wickFourth-interaction-inner-quartic}
  \lean{Pphi2.wickFourth_interaction_inner_quartic}\leanok
  \uses{GaussianField-FinLatticeSites, Pphi2-integrable-wickFourth-wickPolynomial, GaussianField-gffPositionCovariance, GaussianField-latticeGaussianMeasure, GaussianField-Configuration, Pphi2-wickFourth-wickPolynomial-inner-quartic, Pphi2-wickConstant, GaussianField-instMeasurableSpaceConfiguration, GaussianField-FinLatticeField, GaussianField-gffSmearedCovariance, Pphi2-wickPolynomial}
  \texttt{Pphi2.wickFourth\_interaction\_inner\_quartic}
\end{theorem}
\section{\texttt{Pphi2.TorusContinuumLimit.TorusNontriviality}}
\begin{definition}[\texttt{TorusIsInteractingStrict}]\label{Pphi2-TorusIsInteractingStrict}
  \lean{Pphi2.TorusIsInteractingStrict}\leanok
  \uses{GaussianField-NuclearTensorProduct-instModuleReal, GaussianField-NuclearTensorProduct-instTopologicalSpace, GaussianField-NuclearTensorProduct-instIsTopologicalAddGroup, Pphi2-torusConnectedFourPoint, GaussianField-NuclearTensorProduct-instAddCommGroup, GaussianField-Configuration, GaussianField-NuclearTensorProduct-instContinuousSMulReal, GaussianField-SmoothMap-Circle, GaussianField-instMeasurableSpaceConfiguration, GaussianField-TorusTestFunction}
  \texttt{Pphi2.TorusIsInteractingStrict}
\end{definition}
\begin{definition}[\texttt{TorusIsInteracting}]\label{Pphi2-TorusIsInteracting}
  \lean{Pphi2.TorusIsInteracting}\leanok
  \uses{GaussianField-NuclearTensorProduct-instModuleReal, GaussianField-NuclearTensorProduct-instTopologicalSpace, GaussianField-NuclearTensorProduct-instIsTopologicalAddGroup, Pphi2-torusConnectedFourPoint, GaussianField-NuclearTensorProduct-instAddCommGroup, GaussianField-Configuration, GaussianField-NuclearTensorProduct-instContinuousSMulReal, GaussianField-SmoothMap-Circle, GaussianField-instMeasurableSpaceConfiguration, GaussianField-TorusTestFunction}
  \texttt{Pphi2.TorusIsInteracting}
\end{definition}
\begin{definition}[\texttt{torusTwoPoint}]\label{Pphi2-torusTwoPoint}
  \lean{Pphi2.torusTwoPoint}\leanok
  \uses{GaussianField-NuclearTensorProduct-instModuleReal, GaussianField-NuclearTensorProduct-instTopologicalSpace, GaussianField-NuclearTensorProduct-instIsTopologicalAddGroup, GaussianField-NuclearTensorProduct-instAddCommGroup, GaussianField-Configuration, GaussianField-NuclearTensorProduct-instContinuousSMulReal, GaussianField-SmoothMap-Circle, GaussianField-instMeasurableSpaceConfiguration, GaussianField-TorusTestFunction}
  \texttt{Pphi2.torusTwoPoint}
\end{definition}
\begin{definition}[\texttt{torusConnectedFourPoint}]\label{Pphi2-torusConnectedFourPoint}
  \lean{Pphi2.torusConnectedFourPoint}\leanok
  \uses{GaussianField-NuclearTensorProduct-instModuleReal, GaussianField-NuclearTensorProduct-instTopologicalSpace, GaussianField-NuclearTensorProduct-instIsTopologicalAddGroup, GaussianField-NuclearTensorProduct-instAddCommGroup, GaussianField-Configuration, GaussianField-NuclearTensorProduct-instContinuousSMulReal, GaussianField-SmoothMap-Circle, GaussianField-instMeasurableSpaceConfiguration, GaussianField-TorusTestFunction}
  \texttt{Pphi2.torusConnectedFourPoint}
\end{definition}
\begin{theorem}[\texttt{torusPphi2Limit\_exists}]\label{Pphi2-torusPphi2Limit-exists}
  \lean{Pphi2.torusPphi2Limit_exists}\leanok
  \uses{GaussianField-NuclearTensorProduct-instModuleReal, GaussianField-NuclearTensorProduct-instTopologicalSpace, GaussianField-NuclearTensorProduct-instIsTopologicalAddGroup, Pphi2-torusInteractingLimit-exists, Pphi2-IsTorusPphi2Limit, GaussianField-NuclearTensorProduct-instAddCommGroup, GaussianField-Configuration, GaussianField-NuclearTensorProduct-instContinuousSMulReal, GaussianField-SmoothMap-Circle, GaussianField-instMeasurableSpaceConfiguration, Pphi2-torusInteractingMeasure, GaussianField-TorusTestFunction}
  \texttt{Pphi2.torusPphi2Limit\_exists}
\end{theorem}
\begin{definition}[\texttt{torusUrsell4}]\label{Pphi2-torusUrsell4}
  \lean{Pphi2.torusUrsell4}\leanok
  \uses{GaussianField-NuclearTensorProduct-instModuleReal, GaussianField-NuclearTensorProduct-instTopologicalSpace, GaussianField-NuclearTensorProduct-instIsTopologicalAddGroup, GaussianField-NuclearTensorProduct-instAddCommGroup, GaussianField-Configuration, GaussianField-NuclearTensorProduct-instContinuousSMulReal, GaussianField-SmoothMap-Circle, GaussianField-instMeasurableSpaceConfiguration, GaussianField-TorusTestFunction}
  \texttt{Pphi2.torusUrsell4}
\end{definition}
\begin{theorem}[\texttt{toInteracting}]\label{Pphi2-TorusIsInteractingStrict-toInteracting}
  \lean{Pphi2.TorusIsInteractingStrict.toInteracting}\leanok
  \uses{Pphi2-TorusIsInteractingStrict, GaussianField-NuclearTensorProduct-instModuleReal, GaussianField-NuclearTensorProduct-instTopologicalSpace, GaussianField-NuclearTensorProduct-instIsTopologicalAddGroup, Pphi2-torusConnectedFourPoint, GaussianField-NuclearTensorProduct-instAddCommGroup, GaussianField-Configuration, GaussianField-NuclearTensorProduct-instContinuousSMulReal, GaussianField-SmoothMap-Circle, GaussianField-instMeasurableSpaceConfiguration, Pphi2-TorusIsInteracting, GaussianField-TorusTestFunction}
  \texttt{Pphi2.TorusIsInteractingStrict.toInteracting}
\end{theorem}
\begin{definition}[\texttt{TorusIsNondegenerate}]\label{Pphi2-TorusIsNondegenerate}
  \lean{Pphi2.TorusIsNondegenerate}\leanok
  \uses{GaussianField-NuclearTensorProduct-instModuleReal, GaussianField-NuclearTensorProduct-instTopologicalSpace, GaussianField-NuclearTensorProduct-instIsTopologicalAddGroup, GaussianField-NuclearTensorProduct-instAddCommGroup, Pphi2-torusTwoPoint, GaussianField-Configuration, GaussianField-NuclearTensorProduct-instContinuousSMulReal, GaussianField-SmoothMap-Circle, GaussianField-instMeasurableSpaceConfiguration, GaussianField-TorusTestFunction}
  \texttt{Pphi2.TorusIsNondegenerate}
\end{definition}
\begin{definition}[\texttt{TorusInteractingTheoryExists}]\label{Pphi2-TorusInteractingTheoryExists}
  \lean{Pphi2.TorusInteractingTheoryExists}\leanok
  \uses{GaussianField-NuclearTensorProduct-instModuleReal, GaussianField-NuclearTensorProduct-instTopologicalSpace, Pphi2-IsTorusPphi2Limit, GaussianField-NuclearTensorProduct-instAddCommGroup, GaussianField-Configuration, GaussianField-SmoothMap-Circle, GaussianField-instMeasurableSpaceConfiguration, Pphi2-TorusIsInteracting, Pphi2-TorusIsNondegenerate, GaussianField-TorusTestFunction}
  \texttt{Pphi2.TorusInteractingTheoryExists}
\end{definition}
\begin{theorem}[\texttt{torusUrsell4\_diag}]\label{Pphi2-torusUrsell4-diag}
  \lean{Pphi2.torusUrsell4_diag}\leanok
  \uses{GaussianField-NuclearTensorProduct-instModuleReal, GaussianField-NuclearTensorProduct-instTopologicalSpace, GaussianField-NuclearTensorProduct-instIsTopologicalAddGroup, Pphi2-torusConnectedFourPoint, Pphi2-torusUrsell4, GaussianField-NuclearTensorProduct-instAddCommGroup, GaussianField-Configuration, GaussianField-NuclearTensorProduct-instContinuousSMulReal, GaussianField-SmoothMap-Circle, GaussianField-instMeasurableSpaceConfiguration, GaussianField-TorusTestFunction}
  \texttt{Pphi2.torusUrsell4\_diag}
\end{theorem}
\begin{definition}[\texttt{IsTorusPphi2Limit}]\label{Pphi2-IsTorusPphi2Limit}
  \lean{Pphi2.IsTorusPphi2Limit}\leanok
  \uses{GaussianField-NuclearTensorProduct-instModuleReal, GaussianField-NuclearTensorProduct-instTopologicalSpace, GaussianField-NuclearTensorProduct-instIsTopologicalAddGroup, GaussianField-NuclearTensorProduct-instAddCommGroup, GaussianField-Configuration, GaussianField-NuclearTensorProduct-instContinuousSMulReal, GaussianField-SmoothMap-Circle, GaussianField-instMeasurableSpaceConfiguration, Pphi2-torusInteractingMeasure, GaussianField-TorusTestFunction}
  \texttt{Pphi2.IsTorusPphi2Limit}
\end{definition}
\section{\texttt{Pphi2.InteractingMeasure.InteractionVariance}}
\begin{theorem}[\texttt{canonicalSmoothInteraction\_sq\_integrable}]\label{Pphi2-canonicalSmoothInteraction-sq-integrable}
  \lean{Pphi2.canonicalSmoothInteraction_sq_integrable}\leanok
  \uses{Pphi2-smoothWickConstant, Pphi2-canonicalJointMeasure, GaussianField-FinLatticeSites, Pphi2-canonicalSmoothInteraction-eq-sum-crossTerm-diag, Pphi2-canonicalSmoothInteraction, Pphi2-wickPolynomial-continuous-, Pphi2-canonicalSmoothFieldFunction-pointwise-measurable, Pphi2-canonicalCrossTerm-pair-integrable, Pphi2-smoothCovEigenvalue, Pphi2-canonicalSmoothFieldFunction, Pphi2-canonicalCrossTerm, Pphi2-wickPolynomial, Pphi2-CanonicalJoint}
  \texttt{Pphi2.canonicalSmoothInteraction\_sq\_integrable}
\end{theorem}
\begin{theorem}[\texttt{interaction\_memLp}]\label{Pphi2-interaction-memLp}
  \lean{Pphi2.interaction_memLp}\leanok
  \uses{Pphi2-canonicalJointMeasure, Pphi2-canonicalSumConfig, Pphi2-canonicalFullInteractionJoint-eq-interactionFunctional, GaussianField-FinLatticeSites, Pphi2-canonicalFullInteractionJoint, Pphi2-interactionFunctional-measurable, Pphi2-canonicalSumConfig-measurable, GaussianField-latticeGaussianMeasure, GaussianField-Configuration, GaussianField-instMeasurableSpaceConfiguration, Pphi2-canonicalJointMeasure-map-canonicalSumConfig, Pphi2-canonicalFullInteractionJoint-memLp, GaussianField-FinLatticeField, Pphi2-interactionFunctional, Pphi2-CanonicalJoint}
  \texttt{Pphi2.interaction\_memLp}
\end{theorem}
\begin{theorem}[\texttt{field\_pow\_le\_second\_pow}]\label{Pphi2-field-pow-le-second-pow}
  \lean{Pphi2.field_pow_le_second_pow}\leanok
  \uses{GaussianField-finLatticeField-dyninMityaginSpace, GaussianField-measure, GaussianField-ell2-separableSpace, GaussianField-FinLatticeSites, GaussianField-latticeCovarianceGJ, GaussianField-latticeGaussianMeasure, GaussianField-Configuration, GaussianField-gaussian-hypercontractive, GaussianField-instMeasurableSpaceConfiguration, GaussianField-ell2-, GaussianField-FinLatticeField}
  \texttt{Pphi2.field\_pow\_le\_second\_pow}
\end{theorem}
\begin{theorem}[\texttt{interactionFunctional\_pow\_integrable}]\label{Pphi2-interactionFunctional-pow-integrable}
  \lean{Pphi2.interactionFunctional_pow_integrable}\leanok
  \uses{GaussianField-FinLatticeSites, Pphi2-interaction-memLp, Pphi2-interactionFunctional-measurable, GaussianField-latticeGaussianMeasure, GaussianField-Configuration, GaussianField-instMeasurableSpaceConfiguration, GaussianField-FinLatticeField, Pphi2-interactionFunctional}
  \texttt{Pphi2.interactionFunctional\_pow\_integrable}
\end{theorem}
\begin{theorem}[\texttt{interaction\_variance\_le}]\label{Pphi2-interaction-variance-le}
  \lean{Pphi2.interaction_variance_le}\leanok
  \uses{Pphi2-rough-error-variance, Pphi2-canonicalJointMeasure, Pphi2-canonicalSmoothInteraction-variance-le, Pphi2-canonicalSmoothInteraction-sq-integrable, Pphi2-canonicalRoughError-sq-integrable, GaussianField-FinLatticeSites, Pphi2-canonicalFullInteractionJoint, GaussianField-circleSpacing-eq, Pphi2-canonicalSmoothInteraction, GaussianField-latticeGaussianMeasure, GaussianField-Configuration, GaussianField-circleSpacing, GaussianField-instMeasurableSpaceConfiguration, Pphi2-integral-interaction-sq-eq-canonicalJoint, GaussianField-FinLatticeField, Pphi2-canonicalRoughError, Pphi2-interactionFunctional, GaussianField-circleSpacing-pos, Pphi2-CanonicalJoint}
  \texttt{Pphi2.interaction\_variance\_le}
\end{theorem}
\begin{theorem}[\texttt{canonicalSmoothCovariance\_pow\_double\_sum\_le}]\label{Pphi2-canonicalSmoothCovariance-pow-double-sum-le}
  \lean{Pphi2.canonicalSmoothCovariance_pow_double_sum_le}\leanok
  \uses{Pphi2-smoothWickConstant, Pphi2-canonicalSmoothCovariance, Pphi2-smoothWickConstant-le-log-uniform-in-aN, Pphi2-lattice-volume-eq, GaussianField-FinLatticeSites, Pphi2-abs-canonicalSmoothCovariance-le-smoothWickConstant, Lattice-toSemilatticeInf}
  \texttt{Pphi2.canonicalSmoothCovariance\_pow\_double\_sum\_le}
\end{theorem}
\begin{theorem}[\texttt{canonicalSmoothInteraction\_variance\_le}]\label{Pphi2-canonicalSmoothInteraction-variance-le}
  \lean{Pphi2.canonicalSmoothInteraction_variance_le}\leanok
  \uses{Pphi2-canonicalJointMeasure, Pphi2-canonicalSmoothInteraction-eq-sum-crossTerm-diag, Pphi2-canonicalSmoothInteraction, Pphi2-canonicalCrossTerm-diag-l2-sq-le, Pphi2-canonicalCrossTerm-pair-integrable, Pphi2-canonicalCrossTerm, Lattice-toSemilatticeInf, Pphi2-CanonicalJoint}
  \texttt{Pphi2.canonicalSmoothInteraction\_variance\_le}
\end{theorem}
\begin{theorem}[\texttt{free\_product\_moment\_cauchy\_schwarz}]\label{Pphi2-free-product-moment-cauchy-schwarz}
  \lean{Pphi2.free_product_moment_cauchy_schwarz}\leanok
  \uses{GaussianField-FinLatticeSites, Pphi2-interactionFunctional-measurable, Pphi2-integrable-pow-pairing, GaussianField-latticeGaussianMeasure, GaussianField-configuration-eval-measurable, GaussianField-Configuration, GaussianField-instMeasurableSpaceConfiguration, Pphi2-interactionFunctional-pow-integrable, GaussianField-FinLatticeField, Pphi2-interactionFunctional}
  \texttt{Pphi2.free\_product\_moment\_cauchy\_schwarz}
\end{theorem}
\begin{theorem}[\texttt{canonicalSmoothInteraction\_eq\_sum\_crossTerm\_diag}]\label{Pphi2-canonicalSmoothInteraction-eq-sum-crossTerm-diag}
  \lean{Pphi2.canonicalSmoothInteraction_eq_sum_crossTerm_diag}\leanok
  \uses{Pphi2-smoothWickConstant, GaussianField-FinLatticeSites, Pphi2-wickMonomial, Pphi2-canonicalSmoothInteraction, Pphi2-canonicalRoughFieldFunction, Pphi2-canonicalSmoothFieldFunction, Pphi2-roughWickConstant, Pphi2-canonicalCrossTerm, Pphi2-wickPolynomial, Pphi2-CanonicalJoint}
  \texttt{Pphi2.canonicalSmoothInteraction\_eq\_sum\_crossTerm\_diag}
\end{theorem}
\begin{theorem}[\texttt{interaction\_moment\_le}]\label{Pphi2-interaction-moment-le}
  \lean{Pphi2.interaction_moment_le}\leanok
  \uses{Pphi2-canonicalJointMeasure, Pphi2-canonicalSumConfig, Pphi2-canonicalFullInteractionJoint-eq-interactionFunctional, GaussianField-FinLatticeSites, Pphi2-canonicalFullInteractionJoint, Pphi2-interactionFunctional-measurable, GaussianField-latticeGaussianMeasure, Pphi2-interaction-variance-le, GaussianField-Configuration, GaussianField-circleSpacing, GaussianField-instMeasurableSpaceConfiguration, Pphi2-integral-interaction-sq-eq-canonicalJoint, Pphi2-canonicalFullInteractionJoint-moment-le, Pphi2-integral-comp-canonicalSumConfig, GaussianField-FinLatticeField, Pphi2-interactionFunctional, Lattice-toSemilatticeInf, GaussianField-circleSpacing-pos, Pphi2-CanonicalJoint}
  \texttt{Pphi2.interaction\_moment\_le}
\end{theorem}
\begin{theorem}[\texttt{canonicalCrossTerm\_diag\_l2\_sq\_le}]\label{Pphi2-canonicalCrossTerm-diag-l2-sq-le}
  \lean{Pphi2.canonicalCrossTerm_diag_l2_sq_le}\leanok
  \uses{Pphi2-canonicalJointMeasure, Pphi2-canonicalSmoothCovariance-pow-double-sum-le, Pphi2-canonicalCrossTerm-l2-sq-eq-covSum, Pphi2-canonicalRoughCovariance, Pphi2-canonicalSmoothCovariance, GaussianField-FinLatticeSites, Pphi2-canonicalCrossTerm, Pphi2-CanonicalJoint}
  \texttt{Pphi2.canonicalCrossTerm\_diag\_l2\_sq\_le}
\end{theorem}
\section{\texttt{Pphi2.InteractingMeasure.Normalization}}
\begin{theorem}[\texttt{field\_second\_moment\_finite}]\label{Pphi2-field-second-moment-finite}
  \lean{Pphi2.field_second_moment_finite}\leanok
  \uses{GaussianField-finLatticeField-dyninMityaginSpace, GaussianField-measure, GaussianField-ell2-separableSpace, Pphi2-interactionFunctional-bounded-below, GaussianField-finLatticeDelta, GaussianField-FinLatticeSites, Pphi2-partitionFunction, GaussianField-latticeCovarianceGJ, Pphi2-interactionFunctional-measurable, Pphi2-interactingLatticeMeasure, GaussianField-latticeGaussianMeasure, GaussianField-configuration-eval-measurable, GaussianField-Configuration, GaussianField-instMeasurableSpaceConfiguration, Pphi2-boltzmannWeight, GaussianField-ell2-, GaussianField-FinLatticeField, Pphi2-interactionFunctional, Pphi2-partitionFunction-pos, Lattice-toSemilatticeInf, Pphi2-boltzmannWeight-pos, GaussianField-pairing-product-integrable}
  $\int |\omega(\delta_x)|^2\,d\mu_a < \infty$. Proved by dominating $|\omega(\delta_x)|^2 e^{-V}$ by $|\omega(\delta_x)|^2 e^B$ and using Gaussian integrability from \texttt{pairing\_product\_integrable}.
\end{theorem}
\begin{theorem}[\texttt{fkg\_interacting}]\label{Pphi2-fkg-interacting}
  \lean{Pphi2.fkg_interacting}\leanok
  \uses{GaussianField-fkg-perturbed, Pphi2-interactionFunctional-single-site, GaussianField-finLatticeDelta, GaussianField-FinLatticeSites, Pphi2-partitionFunction, Pphi2-boltzmannWeight-integrable, Pphi2-interactionFunctional-measurable, Pphi2-interactingLatticeMeasure, GaussianField-latticeGaussianMeasure, GaussianField-IsFieldMonotone, GaussianField-Configuration, GaussianField-instMeasurableSpaceConfiguration, Pphi2-boltzmannWeight, GaussianField-FinLatticeField, Pphi2-interactionFunctional, Pphi2-partitionFunction-pos, Pphi2-boltzmannWeight-pos}
  For monotone functions $F, G$: $\mathbb{E}_{\mu_a}[FG] \ge \mathbb{E}_{\mu_a}[F] \cdot \mathbb{E}_{\mu_a}[G]$. Proved by converting integrability from $\mu_a$ to Gaussian + weight, applying \texttt{fkg\_perturbed} (unnormalized FKG), then dividing by $Z^2$.
\end{theorem}
\begin{theorem}[\texttt{field\_all\_moments\_finite}]\label{Pphi2-field-all-moments-finite}
  \lean{Pphi2.field_all_moments_finite}\leanok
  \uses{GaussianField-finLatticeField-dyninMityaginSpace, GaussianField-measure, GaussianField-ell2-separableSpace, Pphi2-interactionFunctional-bounded-below, GaussianField-finLatticeDelta, GaussianField-FinLatticeSites, Pphi2-partitionFunction, GaussianField-latticeCovarianceGJ, Pphi2-interactionFunctional-measurable, Pphi2-interactingLatticeMeasure, GaussianField-pairing-is-gaussian, GaussianField-latticeGaussianMeasure, GaussianField-configuration-eval-measurable, GaussianField-Configuration, GaussianField-instMeasurableSpaceConfiguration, Pphi2-boltzmannWeight, GaussianField-ell2-, GaussianField-FinLatticeField, Pphi2-interactionFunctional, Pphi2-partitionFunction-pos, Lattice-toSemilatticeInf, Pphi2-boltzmannWeight-pos}
  $\int |\omega(\delta_x)|^p\,d\mu_a < \infty$ for all $p \in \mathbb{N}$. Uses \texttt{pairing\_is\_gaussian} + \texttt{integrable\_pow\_of\_mem\_interior\_integrableExpSet} for Gaussian moment finiteness.
\end{theorem}
\begin{theorem}[\texttt{bounded\_integrable\_interacting}]\label{Pphi2-bounded-integrable-interacting}
  \lean{Pphi2.bounded_integrable_interacting}\leanok
  \uses{GaussianField-FinLatticeSites, Pphi2-interactingLatticeMeasure-isProbability, Pphi2-interactingLatticeMeasure, GaussianField-Configuration, GaussianField-instMeasurableSpaceConfiguration, GaussianField-FinLatticeField}
  Bounded measurable functions are integrable under $\mu_a$ (since $\mu_a$ is a probability measure).
\end{theorem}
\begin{theorem}[\texttt{generating\_functional\_bounded}]\label{Pphi2-generating-functional-bounded}
  \lean{Pphi2.generating_functional_bounded}\leanok
  \uses{GaussianField-FinLatticeSites, Pphi2-interactingLatticeMeasure-isProbability, Pphi2-interactingLatticeMeasure, GaussianField-Configuration, GaussianField-instMeasurableSpaceConfiguration, GaussianField-FinLatticeField, Lattice-toSemilatticeInf}
  $\left|\int \cos(\omega(f))\,d\mu_a\right| \le 1$ since $|\cos| \le 1$ and $\mu_a$ is a probability measure.
\end{theorem}
\section{\texttt{Pphi2.AsymTorus.AsymExpMomentDischarge}}
\begin{theorem}[\texttt{asymInteractingVariance\_le\_freeVariance\_lattice\_Lt\_uniform} (axiom)]\label{Pphi2-asymInteractingVariance-le-freeVariance-lattice-Lt-uniform}
  \lean{Pphi2.asymInteractingVariance_le_freeVariance_lattice_Lt_uniform}\leanok
  \uses{GaussianField-AsymLatticeSites, GaussianField-Configuration, GaussianField-AsymLatticeField, Pphi2-latticeGaussianMeasureAsym, Pphi2-interactingLatticeMeasureAsym, GaussianField-instMeasurableSpaceConfiguration}
  \texttt{Pphi2.asymInteractingVariance\_le\_freeVariance\_lattice\_Lt\_uniform}
\end{theorem}
\begin{theorem}[\texttt{congr\_simp}]\label{Pphi2-interactingLatticeMeasureAsym-congr-simp}
  \lean{Pphi2.interactingLatticeMeasureAsym.congr_simp}\leanok
  \uses{GaussianField-AsymLatticeSites, GaussianField-Configuration, GaussianField-AsymLatticeField, Pphi2-interactingLatticeMeasureAsym, GaussianField-instMeasurableSpaceConfiguration}
  \texttt{Pphi2.interactingLatticeMeasureAsym.congr\_simp}
\end{theorem}
\begin{theorem}[\texttt{asymInteractingVariance\_le\_freeVariance\_Lt\_uniform}]\label{Pphi2-asymInteractingVariance-le-freeVariance-Lt-uniform}
  \lean{Pphi2.asymInteractingVariance_le_freeVariance_Lt_uniform}\leanok
  \uses{GaussianField-NuclearTensorProduct-instModuleReal, GaussianField-NuclearTensorProduct-instTopologicalSpace, GaussianField-AsymLatticeSites, Pphi2-asymTorusEmbedLiftIso-measurable, Pphi2-asymTorusEmbedLiftIso-eval-eq, GaussianField-NuclearTensorProduct-instIsTopologicalAddGroup, Pphi2-asymLatticeTestFnIso, Pphi2-asymTorusEmbedLiftIso, GaussianField-NuclearTensorProduct-instAddCommGroup, Pphi2-asymTorusInteractingMeasureIso, GaussianField-configuration-eval-measurable, GaussianField-Configuration, GaussianField-NuclearTensorProduct-instContinuousSMulReal, GaussianField-AsymLatticeField, GaussianField-SmoothMap-Circle, Pphi2-latticeGaussianMeasureAsym, Pphi2-interactingLatticeMeasureAsym, GaussianField-instMeasurableSpaceConfiguration, Pphi2-asymInteractingVariance-le-freeVariance-lattice-Lt-uniform, GaussianField-AsymTorusTestFunction}
  \texttt{Pphi2.asymInteractingVariance\_le\_freeVariance\_Lt\_uniform}
\end{theorem}
\begin{theorem}[\texttt{asymInteracting\_mgf\_gaussianDominated} (axiom)]\label{Pphi2-asymInteracting-mgf-gaussianDominated}
  \lean{Pphi2.asymInteracting_mgf_gaussianDominated}
  \uses{GaussianField-AsymLatticeSites, GaussianField-Configuration, GaussianField-AsymLatticeField, Pphi2-interactingLatticeMeasureAsym, GaussianField-instMeasurableSpaceConfiguration}
  \texttt{Pphi2.asymInteracting\_mgf\_gaussianDominated}
\end{theorem}
\begin{theorem}[\texttt{asymInteracting\_expMoment\_volume\_uniform\_proof}]\label{Pphi2-asymInteracting-expMoment-volume-uniform-proof}
  \lean{Pphi2.asymInteracting_expMoment_volume_uniform_proof}\leanok
  \uses{GaussianField-NuclearTensorProduct-instModuleReal, Lattice-toSemilatticeSup, GaussianField-NuclearTensorProduct-instTopologicalSpace, GaussianField-AsymLatticeSites, Pphi2-asymTorusEmbedLiftIso-measurable, Pphi2-asymInteractingVariance-le-freeVariance-Lt-uniform, Pphi2-asymTorusEmbedLiftIso-eval-eq, GaussianField-NuclearTensorProduct-instIsTopologicalAddGroup, Pphi2-asymInteracting-mgf-gaussianDominated, Pphi2-asymLatticeTestFnIso, Pphi2-asymTorusEmbedLiftIso, GaussianField-NuclearTensorProduct-instAddCommGroup, Pphi2-asymTorusInteractingMeasureIso, GaussianField-configuration-eval-measurable, GaussianField-Configuration, GaussianField-NuclearTensorProduct-instContinuousSMulReal, GaussianField-AsymLatticeField, GaussianField-SmoothMap-Circle, Pphi2-latticeGaussianMeasureAsym, Pphi2-interactingLatticeMeasureAsym, GaussianField-instMeasurableSpaceConfiguration, GaussianField-AsymTorusTestFunction}
  \texttt{Pphi2.asymInteracting\_expMoment\_volume\_uniform\_proof}
\end{theorem}
\section{\texttt{Pphi2.NelsonEstimate.CovarianceBoundsGJ}}
\begin{theorem}[\texttt{canonicalRoughCovariance\_pow\_sum\_le\_uniform\_in\_aN\_of\_three\_le}]\label{Pphi2-canonicalRoughCovariance-pow-sum-le-uniform-in-aN-of-three-le}
  \lean{Pphi2.canonicalRoughCovariance_pow_sum_le_uniform_in_aN_of_three_le}\leanok
  \uses{Pphi2-canonicalRoughCovariance, GaussianField-FinLatticeSites, Pphi2-canonicalRoughCovariance-eq-integral-Icc-heatKernel, GaussianField-latticeHeatKernelMatrix, Lattice-toSemilatticeInf}
  \texttt{Pphi2.canonicalRoughCovariance\_pow\_sum\_le\_uniform\_in\_aN\_of\_three\_le}
\end{theorem}
\begin{theorem}[\texttt{smoothWickConstant\_le\_log\_uniform\_in\_aN}]\label{Pphi2-smoothWickConstant-le-log-uniform-in-aN}
  \lean{Pphi2.smoothWickConstant_le_log_uniform_in_aN}\leanok
  \uses{Pphi2-smoothWickConstant, Pphi2-smoothWickConstant-le-log-uniform-in-aN-proved}
  \texttt{Pphi2.smoothWickConstant\_le\_log\_uniform\_in\_aN}
\end{theorem}
\begin{theorem}[\texttt{canonicalRoughCovariance\_pow\_one\_sum\_le\_uniform\_in\_aN\_proved}]\label{Pphi2-canonicalRoughCovariance-pow-one-sum-le-uniform-in-aN-proved}
  \lean{Pphi2.canonicalRoughCovariance_pow_one_sum_le_uniform_in_aN_proved}\leanok
  \uses{GaussianField-latticeEigenvalue1d, Pphi2-canonicalRoughCovariance, GaussianField-FinLatticeSites, Pphi2-canonicalRoughCovariance-eq-integral-Icc-heatKernel, Pphi2-schwinger-rough, GaussianField-latticeFourierProductBasisFun, Pphi2-latticeHeatKernel-entry-eq-eigenvalue-average, Pphi2-heatKernel-row-sum-eq-one, GaussianField-latticeFourierProductNormSq, GaussianField-latticeHeatKernelMatrix}
  \texttt{Pphi2.canonicalRoughCovariance\_pow\_one\_sum\_le\_uniform\_in\_aN\_proved}
\end{theorem}
\begin{theorem}[\texttt{canonicalRoughCovariance\_pow\_two\_sum\_le\_uniform\_in\_aN\_proved}]\label{Pphi2-canonicalRoughCovariance-pow-two-sum-le-uniform-in-aN-proved}
  \lean{Pphi2.canonicalRoughCovariance_pow_two_sum_le_uniform_in_aN_proved}\leanok
  \uses{Pphi2-canonicalRoughCovariance, GaussianField-FinLatticeSites}
  \texttt{Pphi2.canonicalRoughCovariance\_pow\_two\_sum\_le\_uniform\_in\_aN\_proved}
\end{theorem}
\begin{theorem}[\texttt{canonicalRoughCovariance\_pow\_sum\_le\_uniform\_in\_aN}]\label{Pphi2-canonicalRoughCovariance-pow-sum-le-uniform-in-aN}
  \lean{Pphi2.canonicalRoughCovariance_pow_sum_le_uniform_in_aN}\leanok
  \uses{Pphi2-canonicalRoughCovariance, GaussianField-FinLatticeSites, Pphi2-canonicalRoughCovariance-pow-two-sum-le-uniform-in-aN-proved, Pphi2-canonicalRoughCovariance-pow-one-sum-le-uniform-in-aN-proved, Pphi2-canonicalRoughCovariance-pow-sum-le-uniform-in-aN-of-three-le}
  \texttt{Pphi2.canonicalRoughCovariance\_pow\_sum\_le\_uniform\_in\_aN}
\end{theorem}
\section{\texttt{Pphi2.ContinuumLimit.Hypercontractivity}}
\begin{theorem}[\texttt{wickConstant\_eq\_variance\_two\_dim}]\label{Pphi2-wickConstant-eq-variance-two-dim}
  \lean{Pphi2.wickConstant_eq_variance_two_dim}\leanok
  \uses{GaussianField-latticeEigenvalue1d, GaussianField-latticeCovariance-GJ-eq-inv-smul-bare, GaussianField-latticeCovariance, GaussianField-finLatticeDelta, GaussianField-FinLatticeSites, GaussianField-latticeCovarianceGJ, Pphi2-wickConstant-two-eq-dft-eigenvalue-average, Pphi2-wickConstant, GaussianField-ell2-, GaussianField-FinLatticeField}
  \texttt{Pphi2.wickConstant\_eq\_variance\_two\_dim}
\end{theorem}
\begin{theorem}[\texttt{wickMonomial\_latticeGaussian}]\label{Pphi2-wickMonomial-latticeGaussian}
  \lean{Pphi2.wickMonomial_latticeGaussian}\leanok
  \uses{GaussianField-finLatticeField-dyninMityaginSpace, GaussianField-measure, GaussianField-ell2-separableSpace, GaussianField-finLatticeDelta, GaussianField-FinLatticeSites, GaussianField-latticeCovarianceGJ, Pphi2-gaussian-hermite-zero-mean, Pphi2-wickMonomial, GaussianField-pairing-is-gaussian, GaussianField-latticeGaussianMeasure, Pphi2-wickConstant-pos, GaussianField-configuration-eval-measurable, GaussianField-Configuration, Pphi2-wickConstant-eq-variance, Pphi2-wickConstant, GaussianField-instMeasurableSpaceConfiguration, GaussianField-ell2-, GaussianField-FinLatticeField}
  Wick monomials $:x^n:_c$ of order $n \ge 1$ have zero Gaussian mean. Proved from \texttt{gaussian\_hermite\_zero\_mean} + marginal Gaussianity + pushforward.
\end{theorem}
\begin{theorem}[\texttt{wickConstant\_eq\_variance} (axiom)]\label{Pphi2-wickConstant-eq-variance}
  \lean{Pphi2.wickConstant_eq_variance}\leanok
  \uses{GaussianField-latticeEigenvalue1d, GaussianField-latticeCovariance-GJ-eq-inv-smul-bare, GaussianField-covariance, GaussianField-latticeCovariance, GaussianField-finLatticeDelta, GaussianField-FinLatticeSites, GaussianField-latticeCovarianceGJ, Pphi2-wickConstant-eq-latticeEigenvalue1d-family-average, Pphi2-wickConstant, GaussianField-ell2-, GaussianField-FinLatticeField}
  The Wick constant $c_a = \langle T(\delta_x), T(\delta_x) \rangle = E[(\omega(\delta_x))^2]$ equals the GFF variance at a site.
\end{theorem}
\begin{theorem}[\texttt{gaussian\_hypercontractivity\_continuum}]\label{Pphi2-gaussian-hypercontractivity-continuum}
  \lean{Pphi2.gaussian_hypercontractivity_continuum}\leanok
  \uses{Pphi2-EuclideanPlaneBackground-SpaceTime, GaussianField-finLatticeField-dyninMityaginSpace, GaussianField-measure, GaussianField-ell2-separableSpace, Pphi2-latticeTestField, Pphi2-planeBackground, GaussianField-FinLatticeSites, Pphi2-latticeEmbedLift-measurable, GaussianField-latticeCovarianceGJ, Pphi2-EuclideanPlaneBackground-dim, Pphi2-latticeEmbedLift-eval-eq, GaussianField-latticeGaussianMeasure, GaussianField-configuration-eval-measurable, GaussianField-Configuration, GaussianField-gaussian-hypercontractive, GaussianField-instMeasurableSpaceConfiguration, Pphi2-latticeEmbedLift, Pphi2-ContinuumTestFunction, GaussianField-ell2-, GaussianField-FinLatticeField}
  Gaussian hypercontractivity in continuum-embedded form: $\int |\omega(f)|^{pn}\, d(\iota_a)_* \mu_{\text{GFF}} \le (p-1)^{pn/2} \left(\int |\omega(f)|^{2n}\, d(\iota_a)_* \mu_{\text{GFF}}\right)^{p/2}$. Proved by pulling back through \texttt{integral\_map} and applying \texttt{gaussian\_hypercontractive} from gaussian-field.
\end{theorem}
\begin{theorem}[\texttt{partitionFunction\_ge\_one}]\label{Pphi2-partitionFunction-ge-one}
  \lean{Pphi2.partitionFunction_ge_one}\leanok
  \uses{GaussianField-FinLatticeSites, Pphi2-partitionFunction, Pphi2-boltzmannWeight-integrable, GaussianField-latticeGaussianMeasure-isProbability, Pphi2-interactionFunctional-mean-nonpos, GaussianField-latticeGaussianMeasure, GaussianField-Configuration, GaussianField-instMeasurableSpaceConfiguration, GaussianField-FinLatticeField, Pphi2-interactionFunctional}
  $Z_a \ge 1$. Proved by Jensen's inequality applied to the convex function $\exp$.
\end{theorem}
\begin{theorem}[\texttt{congr\_simp}]\label{Pphi2-latticeEmbedLift-congr-simp}
  \lean{Pphi2.latticeEmbedLift.congr_simp}\leanok
  \uses{Pphi2-EuclideanPlaneBackground-SpaceTime, Pphi2-planeBackground, GaussianField-FinLatticeSites, Pphi2-EuclideanPlaneBackground-dim, GaussianField-Configuration, Pphi2-latticeEmbedLift, Pphi2-ContinuumTestFunction, GaussianField-FinLatticeField}
  \texttt{Pphi2.latticeEmbedLift.congr\_simp}
\end{theorem}
\begin{theorem}[\texttt{interactionFunctional\_mean\_nonpos}]\label{Pphi2-interactionFunctional-mean-nonpos}
  \lean{Pphi2.interactionFunctional_mean_nonpos}\leanok
  \uses{GaussianField-finLatticeDelta, GaussianField-FinLatticeSites, GaussianField-latticeGaussianMeasure, GaussianField-Configuration, Pphi2-wickConstant, GaussianField-instMeasurableSpaceConfiguration, GaussianField-FinLatticeField, Pphi2-interactionFunctional, Pphi2-wickPolynomial-integral-eq-coeff-zero, Pphi2-wickPolynomial}
  $\int V_a\, d\mu_{\text{GFF}} \le 0$. Proved from Hermite orthogonality and $P.\text{coeff}_0 \le 0$.
\end{theorem}
\begin{theorem}[\texttt{exponential\_moment\_bound}]\label{Pphi2-exponential-moment-bound}
  \lean{Pphi2.exponential_moment_bound}\leanok
  \uses{GaussianField-FinLatticeSites, GaussianField-latticeGaussianMeasure, Pphi2-nelson-exponential-estimate-master-bounded, GaussianField-Configuration, GaussianField-instMeasurableSpaceConfiguration, GaussianField-FinLatticeField, Pphi2-interactionFunctional}
  $\int e^{-2V_a(\varphi)}\, d\mu_{\text{GFF}} \le K$ uniformly in $a \in (0,1]$. Proved via \texttt{wickPolynomial\_uniform\_bounded\_below} and pointwise exponentiation.
\end{theorem}
\begin{theorem}[\texttt{interacting\_moment\_bound}]\label{Pphi2-interacting-moment-bound}
  \lean{Pphi2.interacting_moment_bound}\leanok
  \uses{Pphi2-EuclideanPlaneBackground-SpaceTime, Pphi2-partitionFunction-ge-one, GaussianField-finLatticeField-dyninMityaginSpace, GaussianField-measure, GaussianField-ell2-separableSpace, Pphi2-latticeTestField, Pphi2-planeBackground, GaussianField-FinLatticeSites, Pphi2-latticeEmbedLift-measurable, Pphi2-partitionFunction, Pphi2-gaussian-hypercontractivity-continuum, GaussianField-latticeCovarianceGJ, Pphi2-EuclideanPlaneBackground-dim, GaussianField-pairing-memLp, Pphi2-interactingLatticeMeasure, Pphi2-exponential-moment-bound, Pphi2-latticeEmbedLift-eval-eq, Pphi2-continuumMeasure, GaussianField-latticeGaussianMeasure, GaussianField-configuration-eval-measurable, GaussianField-Configuration, Pphi2-density-transfer-bound, GaussianField-instMeasurableSpaceConfiguration, Pphi2-latticeEmbedLift, Pphi2-boltzmannWeight, Pphi2-ContinuumTestFunction, GaussianField-ell2-, GaussianField-FinLatticeField, GaussianField-measure-isProbability, Pphi2-interactionFunctional, Pphi2-partitionFunction-pos, Lattice-toSemilatticeInf}
  $\int |\omega(f)|^{pn}\, d\mu_a \le C (2p-1)^{pn/2} (\int |\omega(f)|^{2n}\, d\mu_{\text{GFF}})^{p/2}$ with $C = K^{1/2}$ uniform in $a$. Proved by chaining density transfer + hypercontractivity.
\end{theorem}
\begin{theorem}[\texttt{wickPolynomial\_integral\_eq\_coeff\_zero}]\label{Pphi2-wickPolynomial-integral-eq-coeff-zero}
  \lean{Pphi2.wickPolynomial_integral_eq_coeff_zero}\leanok
  \uses{GaussianField-finLatticeDelta, GaussianField-FinLatticeSites, Pphi2-wickMonomial-latticeGaussian, GaussianField-latticeGaussianMeasure-isProbability, Pphi2-wickMonomial, GaussianField-latticeGaussianMeasure, GaussianField-Configuration, Pphi2-wickConstant, GaussianField-instMeasurableSpaceConfiguration, GaussianField-FinLatticeField, Pphi2-wickPolynomial}
  \texttt{Pphi2.wickPolynomial\_integral\_eq\_coeff\_zero}
\end{theorem}
\begin{theorem}[\texttt{density\_transfer\_bound}]\label{Pphi2-density-transfer-bound}
  \lean{Pphi2.density_transfer_bound}\leanok
  \uses{Pphi2-interactionFunctional-bounded-below, GaussianField-FinLatticeSites, Pphi2-partitionFunction, Pphi2-interactionFunctional-measurable, Pphi2-interactingLatticeMeasure, GaussianField-latticeGaussianMeasure-isProbability, GaussianField-latticeGaussianMeasure, GaussianField-Configuration, GaussianField-instMeasurableSpaceConfiguration, Pphi2-boltzmannWeight, GaussianField-FinLatticeField, Pphi2-interactionFunctional, Pphi2-partitionFunction-pos, Pphi2-boltzmannWeight-pos}
  For nonneg $F$: $\int F\, d\mu_{\text{int}} \le K^{1/2} (\int F^2\, d\mu_{\text{GFF}})^{1/2}$. Proved via Cauchy-Schwarz (Holder $p=q=2$).
\end{theorem}
\begin{theorem}[\texttt{congr\_simp}]\label{Pphi2-interactingLatticeMeasure-congr-simp}
  \lean{Pphi2.interactingLatticeMeasure.congr_simp}\leanok
  \uses{GaussianField-FinLatticeSites, Pphi2-interactingLatticeMeasure, GaussianField-Configuration, GaussianField-instMeasurableSpaceConfiguration, GaussianField-FinLatticeField}
  \texttt{Pphi2.interactingLatticeMeasure.congr\_simp}
\end{theorem}
\section{\texttt{Pphi2.InteractingMeasure.U4AffineBound}}
\begin{theorem}[\texttt{congr\_simp}]\label{Pphi2-u4Deriv2-congr-simp}
  \lean{Pphi2.u4Deriv2.congr_simp}\leanok
  \uses{Pphi2-u4Deriv2, GaussianField-FinLatticeField}
  \texttt{Pphi2.u4Deriv2.congr\_simp}
\end{theorem}
\begin{theorem}[\texttt{u4Deriv\_zero\_eq}]\label{Pphi2-u4Deriv-zero-eq}
  \lean{Pphi2.u4Deriv_zero_eq}\leanok
  \uses{Pphi2-u4-hasDerivWithinAt, Pphi2-normalizedMoment-hasDerivWithinAt, GaussianField-FinLatticeSites, Pphi2-normalizedMoment-zero, GaussianField-latticeGaussianMeasure-isProbability, GaussianField-gffPositionCovariance, GaussianField-latticeGaussianMeasure, GaussianField-Configuration, Pphi2-normalizedMoment, GaussianField-instMeasurableSpaceConfiguration, Pphi2-normalizedMomentDeriv, Pphi2-u4, GaussianField-FinLatticeField, Pphi2-normalizedMomentDeriv-zero, Pphi2-interactionFunctional, Pphi2-u4Deriv}
  \texttt{Pphi2.u4Deriv\_zero\_eq}
\end{theorem}
\begin{theorem}[\texttt{congr\_simp}]\label{Pphi2-u4-congr-simp}
  \lean{Pphi2.u4.congr_simp}\leanok
  \uses{Pphi2-u4, GaussianField-FinLatticeField}
  \texttt{Pphi2.u4.congr\_simp}
\end{theorem}
\begin{theorem}[\texttt{normalizedMoment\_zero}]\label{Pphi2-normalizedMoment-zero}
  \lean{Pphi2.normalizedMoment_zero}\leanok
  \uses{GaussianField-FinLatticeSites, Pphi2-gibbsMoment, GaussianField-latticeGaussianMeasure-isProbability, GaussianField-latticeGaussianMeasure, GaussianField-Configuration, Pphi2-normalizedMoment, GaussianField-instMeasurableSpaceConfiguration, Pphi2-partitionFn, GaussianField-FinLatticeField, Pphi2-interactionFunctional}
  \texttt{Pphi2.normalizedMoment\_zero}
\end{theorem}
\begin{theorem}[\texttt{congr\_simp}]\label{Pphi2-u4Deriv-congr-simp}
  \lean{Pphi2.u4Deriv.congr_simp}\leanok
  \uses{GaussianField-FinLatticeField, Pphi2-u4Deriv}
  \texttt{Pphi2.u4Deriv.congr\_simp}
\end{theorem}
\begin{theorem}[\texttt{continuousWithinAt\_u4Deriv}]\label{Pphi2-continuousWithinAt-u4Deriv}
  \lean{Pphi2.continuousWithinAt_u4Deriv}\leanok
  \uses{Pphi2-interactionFunctional-eq-single, Pphi2-partitionFn-ge-one, Pphi2-partitionFnDeriv, GaussianField-FinLatticeSites, Pphi2-interactionFunctional-measurable, Pphi2-integrable-powMul-interaction, Pphi2-gibbsMoment, GaussianField-latticeGaussianMeasure-isProbability, Pphi2-integrable-pow-pairing, GaussianField-latticeGaussianMeasure, GaussianField-configuration-eval-measurable, GaussianField-Configuration, Pphi2-continuousWithinAt-weighted-exp, Pphi2-normalizedMoment, Pphi2-wickConstant, GaussianField-instMeasurableSpaceConfiguration, Pphi2-partitionFn, Pphi2-normalizedMomentDeriv, Pphi2-momentDeriv, GaussianField-FinLatticeField, Pphi2-interactionFunctional, Pphi2-u4Deriv, Pphi2-wickPolynomial}
  \texttt{Pphi2.continuousWithinAt\_u4Deriv}
\end{theorem}
\begin{theorem}[\texttt{u4Deriv\_le\_affine}]\label{Pphi2-u4Deriv-le-affine}
  \lean{Pphi2.u4Deriv_le_affine}\leanok
  \uses{Pphi2-u4Deriv2, Pphi2-continuousWithinAt-u4Deriv, Pphi2-u4-hasDerivAt2, GaussianField-FinLatticeField, Pphi2-u4Deriv}
  \texttt{Pphi2.u4Deriv\_le\_affine}
\end{theorem}
\begin{theorem}[\texttt{normalizedMomentDeriv\_zero}]\label{Pphi2-normalizedMomentDeriv-zero}
  \lean{Pphi2.normalizedMomentDeriv_zero}\leanok
  \uses{Pphi2-partitionFnDeriv, GaussianField-FinLatticeSites, Pphi2-gibbsMoment, GaussianField-latticeGaussianMeasure-isProbability, GaussianField-latticeGaussianMeasure, GaussianField-Configuration, GaussianField-instMeasurableSpaceConfiguration, Pphi2-partitionFn, Pphi2-normalizedMomentDeriv, Pphi2-momentDeriv, GaussianField-FinLatticeField, Pphi2-interactionFunctional}
  \texttt{Pphi2.normalizedMomentDeriv\_zero}
\end{theorem}
\begin{theorem}[\texttt{continuousWithinAt\_weighted\_exp}]\label{Pphi2-continuousWithinAt-weighted-exp}
  \lean{Pphi2.continuousWithinAt_weighted_exp}\leanok
  \uses{Lattice-toSemilatticeSup, Pphi2-interactionFunctional-bounded-below, GaussianField-FinLatticeSites, Pphi2-interactionFunctional-measurable, GaussianField-latticeGaussianMeasure, GaussianField-Configuration, GaussianField-instMeasurableSpaceConfiguration, GaussianField-FinLatticeField, Pphi2-interactionFunctional, Lattice-toSemilatticeInf}
  \texttt{Pphi2.continuousWithinAt\_weighted\_exp}
\end{theorem}
\begin{theorem}[\texttt{connectedFourPoint\_interactingLatticeMeasure\_eq\_u4\_one}]\label{Pphi2-connectedFourPoint-interactingLatticeMeasure-eq-u4-one}
  \lean{Pphi2.connectedFourPoint_interactingLatticeMeasure_eq_u4_one}\leanok
  \uses{GaussianField-FinLatticeSites, Pphi2-connectedFourPoint, Pphi2-integral-pow-interactingLatticeMeasure, Pphi2-interactingLatticeMeasure, GaussianField-Configuration, Pphi2-normalizedMoment, GaussianField-instMeasurableSpaceConfiguration, Pphi2-u4, GaussianField-FinLatticeField}
  \texttt{Pphi2.connectedFourPoint\_interactingLatticeMeasure\_eq\_u4\_one}
\end{theorem}
\begin{theorem}[\texttt{integral\_pow\_interactingLatticeMeasure}]\label{Pphi2-integral-pow-interactingLatticeMeasure}
  \lean{Pphi2.integral_pow_interactingLatticeMeasure}\leanok
  \uses{GaussianField-FinLatticeSites, Pphi2-partitionFunction, Pphi2-interactionFunctional-measurable, Pphi2-interactingLatticeMeasure, Pphi2-gibbsMoment, GaussianField-latticeGaussianMeasure, GaussianField-Configuration, Pphi2-normalizedMoment, GaussianField-instMeasurableSpaceConfiguration, Pphi2-boltzmannWeight, Pphi2-partitionFn, GaussianField-FinLatticeField, Pphi2-interactionFunctional, Pphi2-partitionFunction-pos, Pphi2-boltzmannWeight-pos}
  \texttt{Pphi2.integral\_pow\_interactingLatticeMeasure}
\end{theorem}
\begin{theorem}[\texttt{lattice\_u4\_neg\_uniform}]\label{Pphi2-lattice-u4-neg-uniform}
  \lean{Pphi2.lattice_u4_neg_uniform}\leanok
  \uses{Pphi2-torus-volume-eq, Pphi2-leadingTerm-const-eq, GaussianField-FinLatticeSites, Pphi2-u4Deriv-le-affine, Pphi2-u4Deriv2, Pphi2-u4Deriv-zero-eq, Pphi2-u4-at-zero, GaussianField-gffPositionCovariance, Pphi2-u4-continuousOn, Pphi2-u4-hasDerivAt, Pphi2-u4Deriv2-abs-le-uniform, GaussianField-circleSpacing, Pphi2-deriv-affine-bound-neg-of-continuousOn, Pphi2-u4, Pphi2-u4Deriv, GaussianField-circleSpacing-pos}
  \texttt{Pphi2.lattice\_u4\_neg\_uniform}
\end{theorem}
\section{\texttt{Pphi2.AsymTorus.AsymLatticeMeasure}}
\begin{theorem}[\texttt{latticeGaussianMeasureAsym\_isProbability}]\label{Pphi2-latticeGaussianMeasureAsym-isProbability}
  \lean{Pphi2.latticeGaussianMeasureAsym_isProbability}\leanok
  \uses{GaussianField-AsymLatticeSites, GaussianField-ell2-separableSpace, GaussianField-asymLatticeField-dyninMityaginSpace, GaussianField-Configuration, GaussianField-AsymLatticeField, Pphi2-latticeGaussianMeasureAsym, GaussianField-instMeasurableSpaceConfiguration, GaussianField-latticeCovarianceAsymGJ, GaussianField-ell2-, GaussianField-measure-isProbability}
  \texttt{Pphi2.latticeGaussianMeasureAsym\_isProbability}
\end{theorem}
\begin{theorem}[\texttt{partitionFunctionAsym\_pos}]\label{Pphi2-partitionFunctionAsym-pos}
  \lean{Pphi2.partitionFunctionAsym_pos}\leanok
  \uses{Pphi2-boltzmannWeightAsym-integrable, GaussianField-AsymLatticeSites, Pphi2-boltzmannWeightAsym, Pphi2-boltzmannWeightAsym-pos, Pphi2-partitionFunctionAsym, GaussianField-Configuration, GaussianField-AsymLatticeField, Pphi2-latticeGaussianMeasureAsym, GaussianField-instMeasurableSpaceConfiguration, Pphi2-latticeGaussianMeasureAsym-isProbability}
  \texttt{Pphi2.partitionFunctionAsym\_pos}
\end{theorem}
\begin{theorem}[\texttt{boltzmannWeightAsym\_integrable}]\label{Pphi2-boltzmannWeightAsym-integrable}
  \lean{Pphi2.boltzmannWeightAsym_integrable}\leanok
  \uses{GaussianField-AsymLatticeSites, Pphi2-boltzmannWeightAsym, Pphi2-interactionFunctionalAsym-bounded-below, GaussianField-Configuration, GaussianField-AsymLatticeField, Pphi2-latticeGaussianMeasureAsym, Pphi2-interactionFunctionalAsym, GaussianField-instMeasurableSpaceConfiguration, Pphi2-interactionFunctionalAsym-measurable, Pphi2-latticeGaussianMeasureAsym-isProbability}
  \texttt{Pphi2.boltzmannWeightAsym\_integrable}
\end{theorem}
\begin{theorem}[\texttt{interactionFunctionalAsym\_bounded\_below}]\label{Pphi2-interactionFunctionalAsym-bounded-below}
  \lean{Pphi2.interactionFunctionalAsym_bounded_below}\leanok
  \uses{GaussianField-AsymLatticeSites, Pphi2-wickConstantAsym, GaussianField-Configuration, GaussianField-AsymLatticeField, Pphi2-interactionFunctionalAsym, GaussianField-instFintypeAsymLatticeSitesOfNeZeroNat, GaussianField-asymLatticeDelta, Pphi2-wickPolynomial-bounded-below, Pphi2-wickPolynomial}
  \texttt{Pphi2.interactionFunctionalAsym\_bounded\_below}
\end{theorem}
\begin{definition}[\texttt{wickConstantAsym}]\label{Pphi2-wickConstantAsym}
  \lean{Pphi2.wickConstantAsym}\leanok
  \uses{GaussianField-AsymLatticeSites, GaussianField-massEigenvaluesAsym, GaussianField-instFintypeAsymLatticeSitesOfNeZeroNat}
  \texttt{Pphi2.wickConstantAsym}
\end{definition}
\begin{theorem}[\texttt{wickConstantAsym\_pos}]\label{Pphi2-wickConstantAsym-pos}
  \lean{Pphi2.wickConstantAsym_pos}\leanok
  \uses{GaussianField-AsymLatticeSites, Pphi2-wickConstantAsym, GaussianField-massEigenvaluesAsym, GaussianField-massOperatorMatrixAsym-eigenvalues-pos, GaussianField-instFintypeAsymLatticeSitesOfNeZeroNat}
  \texttt{Pphi2.wickConstantAsym\_pos}
\end{theorem}
\begin{theorem}[\texttt{latticeGaussianMeasureAsym\_cross\_moment}]\label{Pphi2-latticeGaussianMeasureAsym-cross-moment}
  \lean{Pphi2.latticeGaussianMeasureAsym_cross_moment}\leanok
  \uses{GaussianField-cross-moment-eq-covariance, GaussianField-AsymLatticeSites, GaussianField-ell2-separableSpace, GaussianField-covariance, GaussianField-asymLatticeField-dyninMityaginSpace, GaussianField-Configuration, GaussianField-AsymLatticeField, Pphi2-latticeGaussianMeasureAsym, GaussianField-instMeasurableSpaceConfiguration, GaussianField-latticeCovarianceAsymGJ, GaussianField-ell2-}
  \texttt{Pphi2.latticeGaussianMeasureAsym\_cross\_moment}
\end{theorem}
\begin{theorem}[\texttt{congr\_simp}]\label{Pphi2-interactionFunctionalAsym-congr-simp}
  \lean{Pphi2.interactionFunctionalAsym.congr_simp}\leanok
  \uses{GaussianField-AsymLatticeSites, GaussianField-Configuration, GaussianField-AsymLatticeField, Pphi2-interactionFunctionalAsym}
  \texttt{Pphi2.interactionFunctionalAsym.congr\_simp}
\end{theorem}
\begin{definition}[\texttt{latticeGaussianMeasureAsym}]\label{Pphi2-latticeGaussianMeasureAsym}
  \lean{Pphi2.latticeGaussianMeasureAsym}\leanok
  \uses{GaussianField-AsymLatticeSites, GaussianField-measure, GaussianField-ell2-separableSpace, GaussianField-asymLatticeField-dyninMityaginSpace, GaussianField-Configuration, GaussianField-AsymLatticeField, GaussianField-instMeasurableSpaceConfiguration, GaussianField-latticeCovarianceAsymGJ, GaussianField-ell2-}
  \texttt{Pphi2.latticeGaussianMeasureAsym}
\end{definition}
\begin{theorem}[\texttt{interactingLatticeMeasureAsym\_isProbability}]\label{Pphi2-interactingLatticeMeasureAsym-isProbability}
  \lean{Pphi2.interactingLatticeMeasureAsym_isProbability}\leanok
  \uses{Pphi2-boltzmannWeightAsym-integrable, GaussianField-AsymLatticeSites, Pphi2-partitionFunctionAsym-pos, Pphi2-boltzmannWeightAsym, Pphi2-boltzmannWeightAsym-pos, Pphi2-partitionFunctionAsym, GaussianField-Configuration, GaussianField-AsymLatticeField, Pphi2-latticeGaussianMeasureAsym, Pphi2-interactingLatticeMeasureAsym, GaussianField-instMeasurableSpaceConfiguration}
  \texttt{Pphi2.interactingLatticeMeasureAsym\_isProbability}
\end{theorem}
\begin{theorem}[\texttt{interactionFunctionalAsym\_measurable}]\label{Pphi2-interactionFunctionalAsym-measurable}
  \lean{Pphi2.interactionFunctionalAsym_measurable}\leanok
  \uses{GaussianField-AsymLatticeSites, Pphi2-wickConstantAsym, Pphi2-wickMonomial, GaussianField-configuration-eval-measurable, GaussianField-Configuration, GaussianField-AsymLatticeField, Pphi2-interactionFunctionalAsym, GaussianField-instMeasurableSpaceConfiguration, GaussianField-instFintypeAsymLatticeSitesOfNeZeroNat, GaussianField-asymLatticeDelta, Pphi2-wickPolynomial}
  \texttt{Pphi2.interactionFunctionalAsym\_measurable}
\end{theorem}
\begin{theorem}[\texttt{congr\_simp}]\label{Pphi2-partitionFunctionAsym-congr-simp}
  \lean{Pphi2.partitionFunctionAsym.congr_simp}\leanok
  \uses{Pphi2-partitionFunctionAsym}
  \texttt{Pphi2.partitionFunctionAsym.congr\_simp}
\end{theorem}
\begin{theorem}[\texttt{congr\_simp}]\label{Pphi2-boltzmannWeightAsym-congr-simp}
  \lean{Pphi2.boltzmannWeightAsym.congr_simp}\leanok
  \uses{GaussianField-AsymLatticeSites, Pphi2-boltzmannWeightAsym, GaussianField-Configuration, GaussianField-AsymLatticeField}
  \texttt{Pphi2.boltzmannWeightAsym.congr\_simp}
\end{theorem}
\begin{definition}[\texttt{interactingLatticeMeasureAsym}]\label{Pphi2-interactingLatticeMeasureAsym}
  \lean{Pphi2.interactingLatticeMeasureAsym}\leanok
  \uses{GaussianField-AsymLatticeSites, Pphi2-boltzmannWeightAsym, Pphi2-partitionFunctionAsym, GaussianField-Configuration, GaussianField-AsymLatticeField, Pphi2-latticeGaussianMeasureAsym, GaussianField-instMeasurableSpaceConfiguration}
  \texttt{Pphi2.interactingLatticeMeasureAsym}
\end{definition}
\begin{definition}[\texttt{boltzmannWeightAsym}]\label{Pphi2-boltzmannWeightAsym}
  \lean{Pphi2.boltzmannWeightAsym}\leanok
  \uses{GaussianField-AsymLatticeSites, GaussianField-Configuration, GaussianField-AsymLatticeField, Pphi2-interactionFunctionalAsym}
  \texttt{Pphi2.boltzmannWeightAsym}
\end{definition}
\begin{definition}[\texttt{interactionFunctionalAsym}]\label{Pphi2-interactionFunctionalAsym}
  \lean{Pphi2.interactionFunctionalAsym}\leanok
  \uses{GaussianField-AsymLatticeSites, Pphi2-wickConstantAsym, GaussianField-Configuration, GaussianField-AsymLatticeField, GaussianField-instFintypeAsymLatticeSitesOfNeZeroNat, GaussianField-asymLatticeDelta, Pphi2-wickPolynomial}
  \texttt{Pphi2.interactionFunctionalAsym}
\end{definition}
\begin{definition}[\texttt{partitionFunctionAsym}]\label{Pphi2-partitionFunctionAsym}
  \lean{Pphi2.partitionFunctionAsym}\leanok
  \uses{GaussianField-AsymLatticeSites, Pphi2-boltzmannWeightAsym, GaussianField-Configuration, GaussianField-AsymLatticeField, Pphi2-latticeGaussianMeasureAsym, GaussianField-instMeasurableSpaceConfiguration}
  \texttt{Pphi2.partitionFunctionAsym}
\end{definition}
\begin{theorem}[\texttt{congr\_simp}]\label{Pphi2-latticeGaussianMeasureAsym-congr-simp}
  \lean{Pphi2.latticeGaussianMeasureAsym.congr_simp}\leanok
  \uses{GaussianField-AsymLatticeSites, GaussianField-Configuration, GaussianField-AsymLatticeField, Pphi2-latticeGaussianMeasureAsym, GaussianField-instMeasurableSpaceConfiguration}
  \texttt{Pphi2.latticeGaussianMeasureAsym.congr\_simp}
\end{theorem}
\begin{theorem}[\texttt{boltzmannWeightAsym\_pos}]\label{Pphi2-boltzmannWeightAsym-pos}
  \lean{Pphi2.boltzmannWeightAsym_pos}\leanok
  \uses{GaussianField-AsymLatticeSites, Pphi2-boltzmannWeightAsym, GaussianField-Configuration, GaussianField-AsymLatticeField, Pphi2-interactionFunctionalAsym}
  \texttt{Pphi2.boltzmannWeightAsym\_pos}
\end{theorem}
\section{\texttt{Pphi2.InteractingMeasure.U4DerivativeInterior}}
\begin{theorem}[\texttt{u4\_continuousOn}]\label{Pphi2-u4-continuousOn}
  \lean{Pphi2.u4_continuousOn}\leanok
  \uses{Pphi2-u4-hasDerivWithinAt, GaussianField-FinLatticeSites, GaussianField-gffPositionCovariance, GaussianField-latticeGaussianMeasure, GaussianField-Configuration, GaussianField-instMeasurableSpaceConfiguration, GaussianField-FinLatticeField, Pphi2-interactionFunctional, Pphi2-u4-differentiableAt}
  \texttt{Pphi2.u4\_continuousOn}
\end{theorem}
\begin{theorem}[\texttt{partitionFn\_hasDerivAt}]\label{Pphi2-partitionFn-hasDerivAt}
  \lean{Pphi2.partitionFn_hasDerivAt}\leanok
  \uses{GaussianField-FinLatticeSites, Pphi2-moment-hasDerivAt, GaussianField-latticeGaussianMeasure, GaussianField-Configuration, GaussianField-instMeasurableSpaceConfiguration, GaussianField-FinLatticeField, Pphi2-interactionFunctional}
  \texttt{Pphi2.partitionFn\_hasDerivAt}
\end{theorem}
\begin{theorem}[\texttt{u4\_differentiableAt}]\label{Pphi2-u4-differentiableAt}
  \lean{Pphi2.u4_differentiableAt}\leanok
  \uses{Pphi2-partitionFn-ge-one, Pphi2-partitionFn-hasDerivAt, GaussianField-FinLatticeSites, Pphi2-moment-hasDerivAt, GaussianField-latticeGaussianMeasure, GaussianField-Configuration, GaussianField-instMeasurableSpaceConfiguration, GaussianField-FinLatticeField, Pphi2-interactionFunctional}
  \texttt{Pphi2.u4\_differentiableAt}
\end{theorem}
\begin{theorem}[\texttt{moment\_hasDerivAt2}]\label{Pphi2-moment-hasDerivAt2}
  \lean{Pphi2.moment_hasDerivAt2}\leanok
  \uses{GaussianField-FinLatticeSites, Pphi2-moment-mul-interaction-hasDerivAt, GaussianField-latticeGaussianMeasure, GaussianField-Configuration, GaussianField-instMeasurableSpaceConfiguration, GaussianField-FinLatticeField, Pphi2-interactionFunctional}
  \texttt{Pphi2.moment\_hasDerivAt2}
\end{theorem}
\section{\texttt{Pphi2.ContinuumLimit.TimeReflection}}
\begin{definition}[\texttt{continuumTimeReflection}]\label{Pphi2-continuumTimeReflection}
  \lean{Pphi2.continuumTimeReflection}\leanok
  \uses{Pphi2-EuclideanPlaneBackground-SpaceTime, Pphi2-planeBackground, Pphi2-EuclideanPlaneBackground-dim, Pphi2-ContinuumTestFunction}
  \texttt{Pphi2.continuumTimeReflection}
\end{definition}
\begin{definition}[\texttt{distribTimeReflection}]\label{Pphi2-distribTimeReflection}
  \lean{Pphi2.distribTimeReflection}\leanok
  \uses{Pphi2-EuclideanPlaneBackground-SpaceTime, Pphi2-planeBackground, Pphi2-continuumTimeReflection, Pphi2-EuclideanPlaneBackground-dim, GaussianField-Configuration, Pphi2-ContinuumTestFunction}
  ``\texttt{lean def distribTimeReflection : Configuration (ContinuumTestFunction 2) → Configuration (ContinuumTestFunction 2) }`` Time reflection on distributions: $(\Theta^* \omega)(f) = \omega(\Theta f)$.
\end{definition}
\begin{theorem}[\texttt{continuumTimeReflection\_apply\_coord}]\label{Pphi2-continuumTimeReflection-apply-coord}
  \lean{Pphi2.continuumTimeReflection_apply_coord}\leanok
  \uses{Pphi2-EuclideanPlaneBackground-SpaceTime, Pphi2-planeBackground, Pphi2-continuumTimeReflection, Pphi2-EuclideanPlaneBackground-dim, Pphi2-ContinuumTestFunction}
  \texttt{Pphi2.continuumTimeReflection\_apply\_coord}
\end{theorem}
\begin{theorem}[\texttt{distribTimeReflection\_apply}]\label{Pphi2-distribTimeReflection-apply}
  \lean{Pphi2.distribTimeReflection_apply}\leanok
  \uses{Pphi2-EuclideanPlaneBackground-SpaceTime, Pphi2-planeBackground, Pphi2-continuumTimeReflection, Pphi2-EuclideanPlaneBackground-dim, Pphi2-distribTimeReflection, GaussianField-Configuration, Pphi2-ContinuumTestFunction}
  $(\Theta^* \omega)(f) = \omega(\Theta f)$.
\end{theorem}
\begin{theorem}[\texttt{distribTimeReflection\_continuous}]\label{Pphi2-distribTimeReflection-continuous}
  \lean{Pphi2.distribTimeReflection_continuous}\leanok
  \uses{Pphi2-EuclideanPlaneBackground-SpaceTime, Pphi2-planeBackground, Pphi2-continuumTimeReflection, Pphi2-EuclideanPlaneBackground-dim, Pphi2-distribTimeReflection, GaussianField-Configuration, Pphi2-ContinuumTestFunction}
  \texttt{Pphi2.distribTimeReflection\_continuous}
\end{theorem}
\section{\texttt{Pphi2.TorusContinuumLimit.TorusTightness}}
\begin{theorem}[\texttt{torus\_second\_moment\_uniform}]\label{Pphi2-torus-second-moment-uniform}
  \lean{Pphi2.torus_second_moment_uniform}\leanok
  \uses{GaussianField-NuclearTensorProduct-instModuleReal, GaussianField-NuclearTensorProduct-instTopologicalSpace, Pphi2-torusEmbeddedTwoPoint, GaussianField-NuclearTensorProduct-instIsTopologicalAddGroup, Pphi2-torusContinuumMeasure, GaussianField-NuclearTensorProduct-instAddCommGroup, GaussianField-Configuration, GaussianField-NuclearTensorProduct-instContinuousSMulReal, GaussianField-SmoothMap-Circle, Pphi2-torusEmbeddedTwoPoint-uniform-bound, GaussianField-instMeasurableSpaceConfiguration, GaussianField-TorusTestFunction}
  $\int (\omega f)^2\, d\nu_{\text{GFF},N} \le C(f,L,m)$ uniformly in $N$. Proved from \texttt{torusEmbeddedTwoPoint\_uniform\_bound}.
\end{theorem}
\begin{theorem}[\texttt{torusContinuumMeasures\_tight}]\label{Pphi2-torusContinuumMeasures-tight}
  \lean{Pphi2.torusContinuumMeasures_tight}\leanok
  \uses{GaussianField-NuclearTensorProduct-instModuleReal, GaussianField-NuclearTensorProduct-instTopologicalSpace, GaussianField-finLatticeField-dyninMityaginSpace, GaussianField-measure, GaussianField-ell2-separableSpace, GaussianField-NuclearTensorProduct-dyninMityaginSpace, GaussianField-FinLatticeSites, GaussianField-NuclearTensorProduct-instIsTopologicalAddGroup, Pphi2-torusEmbedLift, Pphi2-torus-second-moment-uniform, GaussianField-latticeCovarianceGJ, Pphi2-torusContinuumMeasure, GaussianField-pairing-memLp, Pphi2-torusEmbedLift-measurable, GaussianField-NuclearTensorProduct-instAddCommGroup, GaussianField-latticeGaussianMeasure, GaussianField-configuration-eval-measurable, GaussianField-Configuration, GaussianField-NuclearTensorProduct-instContinuousSMulReal, Pphi2-torusContinuumMeasure-isProbability, GaussianField-SmoothMap-Circle, GaussianField-circleSpacing, GaussianField-instMeasurableSpaceConfiguration, Pphi2-latticeTestFn, Pphi2-torusEmbedLift-eval-eq, GaussianField-ell2-, GaussianField-FinLatticeField, GaussianField-configuration-tight-of-uniform-second-moments, GaussianField-TorusTestFunction, GaussianField-circleSpacing-pos}
  The family $\{\nu_{\text{GFF},N}\}_{N \ge 1}$ is tight. Proved by applying \texttt{configuration\_tight\_of\_uniform\_second\_moments} with the uniform second moment bounds, nuclearity of \texttt{TorusTestFunction L}, and lattice Gaussian integrability.
\end{theorem}
\section{\texttt{Pphi2.NelsonEstimate.AsymRoughErrorVariance}}
\begin{theorem}[\texttt{asymCanonicalRoughError\_eq\_sum\_over\_cross\_terms}]\label{Pphi2-asymCanonicalRoughError-eq-sum-over-cross-terms}
  \lean{Pphi2.asymCanonicalRoughError_eq_sum_over_cross_terms}\leanok
  \uses{GaussianField-AsymLatticeSites, Pphi2-asymCanonicalRoughWickConstant, Pphi2-asymCanonicalCrossTerm, Pphi2-AsymCanonicalJoint, Pphi2-asymCanonicalRoughError, Pphi2-wickMonomial, Pphi2-asymCanonicalRoughFieldFunction, Pphi2-asymSmoothWickConstant, GaussianField-instFintypeAsymLatticeSitesOfNeZeroNat, Pphi2-asymCanonicalRoughError-pointwise-decomposition, Pphi2-asymCanonicalSmoothFieldFunction}
  \texttt{Pphi2.asymCanonicalRoughError\_eq\_sum\_over\_cross\_terms}
\end{theorem}
\begin{theorem}[\texttt{asymCanonicalRoughError\_pointwise\_decomposition}]\label{Pphi2-asymCanonicalRoughError-pointwise-decomposition}
  \lean{Pphi2.asymCanonicalRoughError_pointwise_decomposition}\leanok
  \uses{Pphi2-asymCanonicalFullInteractionJoint, GaussianField-AsymLatticeSites, Pphi2-wickConstantAsym, Pphi2-asymCanonicalRoughWickConstant, Pphi2-AsymCanonicalJoint, Pphi2-asymCanonicalRoughError, Pphi2-wickMonomial, Pphi2-asymCanonicalSumFieldFunction, Pphi2-asymCanonicalRoughFieldFunction, Pphi2-asymSmoothWickConstant, GaussianField-instFintypeAsymLatticeSitesOfNeZeroNat, Pphi2-asymCanonicalSmoothInteraction, Pphi2-wickPolynomial-add-sub-self, Pphi2-asymCanonicalSmoothFieldFunction, Pphi2-wickPolynomial}
  \texttt{Pphi2.asymCanonicalRoughError\_pointwise\_decomposition}
\end{theorem}
\begin{theorem}[\texttt{asymRoughError\_variance}]\label{Pphi2-asymRoughError-variance}
  \lean{Pphi2.asymRoughError_variance}\leanok
  \uses{Pphi2-asymCanonicalRoughError-eq-sum-over-cross-terms, Pphi2-asymCanonicalCrossTerm, Pphi2-AsymCanonicalJoint, Pphi2-asymCanonicalCrossTerm-l2-sq-le, Pphi2-asymCanonicalRoughError, Pphi2-asymCanonicalJointMeasure, Lattice-toSemilatticeInf, Pphi2-AsymCanonicalMode}
  \texttt{Pphi2.asymRoughError\_variance}
\end{theorem}
\section{\texttt{Pphi2.IRLimit.CylinderOS}}
\begin{theorem}[\texttt{of\_forall}]\label{Pphi2-CylinderMeasureSequenceEventuallyReflectionPositive-of-forall}
  \lean{Pphi2.CylinderMeasureSequenceEventuallyReflectionPositive.of_forall}\leanok
  \uses{GaussianField-NuclearTensorProduct-instModuleReal, GaussianField-NuclearTensorProduct-instTopologicalSpace, GaussianField-cylinderPositiveTimeSubmodule, GaussianField-CylinderTestFunction, GaussianField-NuclearTensorProduct-instIsTopologicalAddGroup, Pphi2-CylinderRPMatrixNonnegative, Pphi2-CylinderMeasureReflectionPositive, GaussianField-NuclearTensorProduct-instAddCommGroup, GaussianField-Configuration, GaussianField-NuclearTensorProduct-instContinuousSMulReal, GaussianField-SmoothMap-Circle, GaussianField-instMeasurableSpaceConfiguration, Pphi2-CylinderMeasureSequenceEventuallyReflectionPositive}
  \texttt{Pphi2.CylinderMeasureSequenceEventuallyReflectionPositive.of\_forall}
\end{theorem}
\begin{theorem}[\texttt{cylinderPullback\_timeReflection\_invariant}]\label{Pphi2-cylinderPullback-timeReflection-invariant}
  \lean{Pphi2.cylinderPullback_timeReflection_invariant}\leanok
  \uses{GaussianField-NuclearTensorProduct-instModuleReal, GaussianField-NuclearTensorProduct-instTopologicalSpace, Pphi2-cylinderPullback, GaussianField-CylinderTestFunction, GaussianField-NuclearTensorProduct-instIsTopologicalAddGroup, Pphi2-cylinderToTorusEmbed, GaussianField-configuration-measurable-of-eval-measurable, GaussianField-NuclearTensorProduct-instAddCommGroup, GaussianField-cylinderTimeReflection, Pphi2-cylinderPullbackMeasure, GaussianField-configuration-eval-measurable, GaussianField-Configuration, GaussianField-NuclearTensorProduct-instContinuousSMulReal, GaussianField-SmoothMap-Circle, GaussianField-instMeasurableSpaceConfiguration, GaussianField-AsymTorusTestFunction, Pphi2-cylinderToTorusEmbed-comp-timeReflection, GaussianField-asymTorusTimeReflection}
  ``\texttt{lean theorem cylinderPullback\_timeReflection\_invariant (Lt : ℝ) [Fact (0 < Lt)] (μ : Measure (Configuration (AsymTorusTestFunction Lt Ls))) [IsProbabilityMeasure μ] (hμ\_refl : ∀ g : AsymTorusTestFunction Lt Ls, ∫ ω, Complex.exp (Complex.I * ↑(ω g)) ∂μ = ∫ ω, Complex.exp (Complex.I * ↑(ω (asymTorusTimeReflection Lt Ls g))) ∂μ) (f : CylinderTestFunction Ls) : ∫ ω, Complex.exp (Complex.I * ↑(ω f)) ∂(cylinderPullbackMeasure Lt Ls μ) = ∫ ω, Complex.exp (Complex.I * ↑(ω (cylinderTimeReflection Ls f))) ∂(cylinderPullbackMeasure Lt Ls μ) }``  Informal statement: The cylinder pullback measure is exactly time-reflection invariant at every finite $L_t$, via the intertwining $\mathrm{embed} \circ \Theta = \Theta_{\mathrm{torus}} \circ \mathrm{embed}$.  Proof: Fully proved. Same structure as time translation, using \texttt{cylinderToTorusEmbed\_comp\_timeReflection}.
\end{theorem}
\begin{theorem}[\texttt{cylinderMeasureReflectionPositive\_of\_tendsto\_cf}]\label{Pphi2-cylinderMeasureReflectionPositive-of-tendsto-cf}
  \lean{Pphi2.cylinderMeasureReflectionPositive_of_tendsto_cf}\leanok
  \uses{GaussianField-NuclearTensorProduct-instModuleReal, GaussianField-NuclearTensorProduct-instTopologicalSpace, GaussianField-cylinderPositiveTimeSubmodule, GaussianField-CylinderTestFunction, GaussianField-NuclearTensorProduct-instIsTopologicalAddGroup, Pphi2-CylinderMeasureReflectionPositive, GaussianField-NuclearTensorProduct-instAddCommGroup, GaussianField-cylinderTimeReflection, GaussianField-Configuration, GaussianField-NuclearTensorProduct-instContinuousSMulReal, GaussianField-SmoothMap-Circle, GaussianField-instMeasurableSpaceConfiguration, Pphi2-CylinderMeasureSequenceEventuallyReflectionPositive}
  \texttt{Pphi2.cylinderMeasureReflectionPositive\_of\_tendsto\_cf}
\end{theorem}
\begin{theorem}[\texttt{cylinderPullback\_timeTranslation\_invariant}]\label{Pphi2-cylinderPullback-timeTranslation-invariant}
  \lean{Pphi2.cylinderPullback_timeTranslation_invariant}\leanok
  \uses{GaussianField-NuclearTensorProduct-instModuleReal, GaussianField-NuclearTensorProduct-instTopologicalSpace, Pphi2-cylinderPullback, GaussianField-CylinderTestFunction, GaussianField-NuclearTensorProduct-instIsTopologicalAddGroup, GaussianField-asymTorusTranslation, Pphi2-cylinderToTorusEmbed, GaussianField-configuration-measurable-of-eval-measurable, GaussianField-NuclearTensorProduct-instAddCommGroup, Pphi2-cylinderPullbackMeasure, GaussianField-configuration-eval-measurable, GaussianField-Configuration, GaussianField-NuclearTensorProduct-instContinuousSMulReal, GaussianField-SmoothMap-Circle, GaussianField-instMeasurableSpaceConfiguration, GaussianField-AsymTorusTestFunction, Pphi2-cylinderToTorusEmbed-comp-timeTranslation, GaussianField-cylinderTranslation}
  ``\texttt{lean theorem cylinderPullback\_timeTranslation\_invariant (Lt : ℝ) [Fact (0 < Lt)] (μ : Measure (Configuration (AsymTorusTestFunction Lt Ls))) [IsProbabilityMeasure μ] (hμ\_os2 : ∀ v f, ∫ ω, Complex.exp (Complex.I * ↑(ω f)) ∂μ = ∫ ω, Complex.exp (Complex.I * ↑(ω (asymTorusTranslation Lt Ls v f))) ∂μ) (τ : ℝ) (f : CylinderTestFunction Ls) : ∫ ω, Complex.exp (Complex.I * ↑(ω f)) ∂(cylinderPullbackMeasure Lt Ls μ) = ∫ ω, Complex.exp (Complex.I * ↑(ω (cylinderTranslation Ls 0 τ f))) ∂(cylinderPullbackMeasure Lt Ls μ) }``  Informal statement: The cylinder pullback measure is exactly time-translation invariant at every finite $L_t$: $Z_{L_t}(\tau_\tau f) = Z_{L_t}(f)$, because periodization intertwines time shifts and the torus measure is translation-invariant.  Proof: Fully proved. Unfolds \texttt{cylinderPullbackMeasure} via \texttt{integral\_map}, rewrites using \texttt{cylinderPullback\_eval} and \texttt{cylinderToTorusEmbed\_comp\_timeTranslation}, then applies the torus translation invariance hypothesis.
\end{theorem}
\begin{definition}[\texttt{CylinderRPMatrixNonnegative}]\label{Pphi2-CylinderRPMatrixNonnegative}
  \lean{Pphi2.CylinderRPMatrixNonnegative}\leanok
  \uses{GaussianField-NuclearTensorProduct-instModuleReal, GaussianField-NuclearTensorProduct-instTopologicalSpace, GaussianField-cylinderPositiveTimeSubmodule, GaussianField-CylinderTestFunction, GaussianField-NuclearTensorProduct-instIsTopologicalAddGroup, GaussianField-NuclearTensorProduct-instAddCommGroup, GaussianField-cylinderTimeReflection, GaussianField-Configuration, GaussianField-NuclearTensorProduct-instContinuousSMulReal, GaussianField-SmoothMap-Circle, GaussianField-instMeasurableSpaceConfiguration}
  \texttt{Pphi2.CylinderRPMatrixNonnegative}
\end{definition}
\begin{definition}[\texttt{AsymTorusSequenceHasCylinderOS2Symmetry}]\label{Pphi2-AsymTorusSequenceHasCylinderOS2Symmetry}
  \lean{Pphi2.AsymTorusSequenceHasCylinderOS2Symmetry}\leanok
  \uses{GaussianField-NuclearTensorProduct-instModuleReal, GaussianField-NuclearTensorProduct-instTopologicalSpace, GaussianField-NuclearTensorProduct-instIsTopologicalAddGroup, GaussianField-asymTorusTranslation, GaussianField-NuclearTensorProduct-instAddCommGroup, GaussianField-Configuration, GaussianField-NuclearTensorProduct-instContinuousSMulReal, GaussianField-SmoothMap-Circle, GaussianField-instMeasurableSpaceConfiguration, GaussianField-AsymTorusTestFunction, GaussianField-asymTorusTimeReflection}
  \texttt{Pphi2.AsymTorusSequenceHasCylinderOS2Symmetry}
\end{definition}
\begin{definition}[\texttt{CylinderMeasureReflectionPositive}]\label{Pphi2-CylinderMeasureReflectionPositive}
  \lean{Pphi2.CylinderMeasureReflectionPositive}\leanok
  \uses{GaussianField-NuclearTensorProduct-instModuleReal, GaussianField-NuclearTensorProduct-instTopologicalSpace, GaussianField-cylinderPositiveTimeSubmodule, GaussianField-CylinderTestFunction, GaussianField-NuclearTensorProduct-instIsTopologicalAddGroup, Pphi2-CylinderRPMatrixNonnegative, GaussianField-NuclearTensorProduct-instAddCommGroup, GaussianField-Configuration, GaussianField-NuclearTensorProduct-instContinuousSMulReal, GaussianField-SmoothMap-Circle, GaussianField-instMeasurableSpaceConfiguration}
  \texttt{Pphi2.CylinderMeasureReflectionPositive}
\end{definition}
\begin{theorem}[\texttt{of\_torusOS}]\label{Pphi2-AsymTorusSequenceHasCylinderOS2Symmetry-of-torusOS}
  \lean{Pphi2.AsymTorusSequenceHasCylinderOS2Symmetry.of_torusOS}\leanok
  \uses{GaussianField-NuclearTensorProduct-instModuleReal, GaussianField-NuclearTensorProduct-instTopologicalSpace, Pphi2-AsymSatisfiesTorusOS-os2-translation, Pphi2-AsymSatisfiesTorusOS, GaussianField-NuclearTensorProduct-instIsTopologicalAddGroup, GaussianField-asymTorusTranslation, GaussianField-NuclearTensorProduct-instAddCommGroup, Pphi2-AsymTorusSequenceHasCylinderOS2Symmetry, GaussianField-Configuration, GaussianField-NuclearTensorProduct-instContinuousSMulReal, GaussianField-SmoothMap-Circle, GaussianField-instMeasurableSpaceConfiguration, GaussianField-AsymTorusTestFunction, Pphi2-AsymSatisfiesTorusOS-os2-timeReflection, GaussianField-asymTorusTimeReflection}
  \texttt{Pphi2.AsymTorusSequenceHasCylinderOS2Symmetry.of\_torusOS}
\end{theorem}
\begin{theorem}[\texttt{routeBPrime\_cylinder\_OS}]\label{Pphi2-routeBPrime-cylinder-OS}
  \lean{Pphi2.routeBPrime_cylinder_OS}\leanok
  \uses{Pphi2-cylinderIRLimit-exists, GaussianField-NuclearTensorProduct-instModuleReal, GaussianField-NuclearTensorProduct-instTopologicalSpace, GaussianField-cylinderPositiveTimeSubmodule, Pphi2-cylinderPullback, GaussianField-cylinderSpatialTranslation, GaussianField-CylinderTestFunction, GaussianField-NuclearTensorProduct-instIsTopologicalAddGroup, Pphi2-AsymTorusSequenceHasUniformGreenMomentBound, GaussianField-asymTorusTranslation, Pphi2-CylinderRPMatrixNonnegative, Pphi2-cylinderToTorusEmbed, Pphi2-CylinderMeasureReflectionPositive, GaussianField-configuration-measurable-of-eval-measurable, GaussianField-NuclearTensorProduct-instAddCommGroup, Pphi2-cylinderToTorusEmbed-comp-spatialTranslation, Pphi2-AsymTorusSequenceHasCylinderOS2Symmetry, GaussianField-cylinderTimeReflection, Pphi2-cylinderPullbackMeasure, Pphi2-cylinderIR-uniform-exponential-moment, GaussianField-configuration-eval-measurable, GaussianField-Configuration, GaussianField-NuclearTensorProduct-instContinuousSMulReal, GaussianField-SmoothMap-Circle, GaussianField-instMeasurableSpaceConfiguration, Pphi2-MeasureHasGreenMomentBound, Pphi2-cylinderMeasureReflectionPositive-of-tendsto-cf, Pphi2-cylinderPullback-timeTranslation-invariant, GaussianField-AsymTorusTestFunction, Pphi2-CylinderMeasureSequenceEventuallyReflectionPositive, GaussianField-cylinderTranslation, Pphi2-cylinderPullback-timeReflection-invariant, GaussianField-asymTorusTimeReflection, Lattice-toSemilatticeInf}
  ``\texttt{lean theorem routeBPrime\_cylinder\_OS (mass : ℝ) (hmass : 0 < mass) (KG CG : ℝ) (hKG\_pos : 0 < KG) (hCG\_pos : 0 < CG) (Lt : ℕ → ℝ) (hLt : ∀ n, Fact (0 < Lt n)) (hLt\_tend : Tendsto Lt atTop atTop) (μ : ∀ n, Measure (Configuration (AsymTorusTestFunction (Lt n) Ls))) (hμ\_prob : ∀ n, IsProbabilityMeasure (μ n)) (hμ\_green : AsymTorusSequenceHasUniformGreenMomentBound ...) (hμ\_rp : CylinderMeasureSequenceEventuallyReflectionPositive ...) (hμ\_os2 : AsymTorusSequenceHasCylinderOS2Symmetry ...) : ∃ (ν : Measure (Configuration (CylinderTestFunction Ls))), IsProbabilityMeasure ν ∧ ... (OS0 ∧ OS2 ∧ OS3) }``  Informal statement: Given a sequence of asymmetric torus measures $\mu_n$ on $T_{L_{t,n}, L_s}$ with $L_{t,n} \to \infty$, eventual Green moments, eventual pullback RP, and the exact sequence-level OS2 symmetry input \texttt{AsymTorusSequenceHasCylinderOS2Symmetry}, there exists a probability measure $\nu$ on the cylinder $S^1_{L_s} \times \mathbb{R}$ satisfying OS0, OS2, and OS3.  Proof: Fully proved conditional on the explicit family inputs. Extracts a subsequential limit from \texttt{cylinderIRLimit\_exists}. OS2 is transferred through characteristic-functional convergence and exact finite-$L_t$ invariance. OS0 uses the explicit Green-moment hypothesis. OS3 is transferred from \texttt{hμ\_rp}.  Dependencies: \texttt{cylinderIRLimit\_exists}, \texttt{cylinderPullback\_timeReflection\_invariant}, \texttt{CylinderMeasureSequenceEventuallyReflectionPositive}, \texttt{AsymTorusSequenceHasUniformGreenMomentBound}, \texttt{AsymTorusSequenceHasCylinderOS2Symmetry}.
\end{theorem}
\begin{definition}[\texttt{CylinderMeasureSequenceEventuallyReflectionPositive}]\label{Pphi2-CylinderMeasureSequenceEventuallyReflectionPositive}
  \lean{Pphi2.CylinderMeasureSequenceEventuallyReflectionPositive}\leanok
  \uses{GaussianField-NuclearTensorProduct-instModuleReal, GaussianField-NuclearTensorProduct-instTopologicalSpace, GaussianField-cylinderPositiveTimeSubmodule, GaussianField-CylinderTestFunction, GaussianField-NuclearTensorProduct-instIsTopologicalAddGroup, Pphi2-CylinderRPMatrixNonnegative, GaussianField-NuclearTensorProduct-instAddCommGroup, GaussianField-Configuration, GaussianField-NuclearTensorProduct-instContinuousSMulReal, GaussianField-SmoothMap-Circle, GaussianField-instMeasurableSpaceConfiguration}
  ``\texttt{lean def CylinderRPMatrixNonnegative (ν : Measure (Configuration (CylinderTestFunction Ls))) (n : ℕ) (f : Fin n → ↥(cylinderPositiveTimeSubmodule Ls)) (c : Fin n → ℂ) : Prop  def CylinderMeasureSequenceEventuallyReflectionPositive (νseq : ℕ → Measure (Configuration (CylinderTestFunction Ls))) : Prop }`\texttt{  Informal statement (OS3 reflection positivity input): }CylinderRPMatrixNonnegative\texttt{ is the single formal RP matrix inequality. }CylinderMeasureSequenceEventuallyReflectionPositive\texttt{ says that for each fixed finite RP matrix, the cylinder pullback sequence is eventually nonnegative. This is weaker and more accurate than requiring full cylinder RP at every finite }Lt\texttt{.  Use: }routeBPrime\_cylinder\_OS\texttt{ transfers this input to the IR limit by characteristic-functional convergence and finite-sum limits. Stronger full-RP statements feed this exact input through the proved bridge theorems }CylinderMeasureSequenceEventuallyReflectionPositive.of\_forall\texttt{ and }CylinderMeasureSequenceEventuallyReflectionPositive.of\_eventually\_full`.
\end{definition}
\begin{theorem}[\texttt{of\_eventually\_full}]\label{Pphi2-CylinderMeasureSequenceEventuallyReflectionPositive-of-eventually-full}
  \lean{Pphi2.CylinderMeasureSequenceEventuallyReflectionPositive.of_eventually_full}\leanok
  \uses{GaussianField-NuclearTensorProduct-instModuleReal, GaussianField-NuclearTensorProduct-instTopologicalSpace, GaussianField-cylinderPositiveTimeSubmodule, GaussianField-CylinderTestFunction, GaussianField-NuclearTensorProduct-instIsTopologicalAddGroup, Pphi2-CylinderRPMatrixNonnegative, Pphi2-CylinderMeasureReflectionPositive, GaussianField-NuclearTensorProduct-instAddCommGroup, GaussianField-Configuration, GaussianField-NuclearTensorProduct-instContinuousSMulReal, GaussianField-SmoothMap-Circle, GaussianField-instMeasurableSpaceConfiguration, Pphi2-CylinderMeasureSequenceEventuallyReflectionPositive}
  \texttt{Pphi2.CylinderMeasureSequenceEventuallyReflectionPositive.of\_eventually\_full}
\end{theorem}
\section{\texttt{Pphi2.GaussianContinuumLimit.EmbeddedCovariance}}
\begin{definition}[\texttt{embeddedTwoPoint}]\label{Pphi2-embeddedTwoPoint}
  \lean{Pphi2.embeddedTwoPoint}\leanok
  \uses{Pphi2-EuclideanPlaneBackground-SpaceTime, Pphi2-gaussianContinuumMeasure, Pphi2-planeBackground, Pphi2-EuclideanPlaneBackground-dim, GaussianField-Configuration, GaussianField-instMeasurableSpaceConfiguration, Pphi2-ContinuumTestFunction}
  ``\texttt{lean def embeddedTwoPoint (a mass : ℝ) ... (f g : ContinuumTestFunction d) : ℝ }`` Two-point function $\int \omega(f) \omega(g)\, d\nu_{\text{GFF},a}$.
\end{definition}
\begin{definition}[\texttt{latticeGreenBilinear}]\label{Pphi2-latticeGreenBilinear}
  \lean{Pphi2.latticeGreenBilinear}\leanok
  \uses{Pphi2-latticeTestField, GaussianField-FinLatticeSites, GaussianField-massEigenvalues, GaussianField-massEigenvectorBasis, Pphi2-ContinuumTestFunction}
  \texttt{Pphi2.latticeGreenBilinear}
\end{definition}
\begin{theorem}[\texttt{congr\_simp}]\label{Pphi2-embeddedTwoPoint-congr-simp}
  \lean{Pphi2.embeddedTwoPoint.congr_simp}\leanok
  \uses{Pphi2-embeddedTwoPoint, Pphi2-ContinuumTestFunction}
  \texttt{Pphi2.embeddedTwoPoint.congr\_simp}
\end{theorem}
\begin{definition}[\texttt{continuumGreenBilinear}]\label{Pphi2-continuumGreenBilinear}
  \lean{Pphi2.continuumGreenBilinear}\leanok
  \uses{Pphi2-EuclideanPlaneBackground-SpaceTime, Pphi2-planeBackground, Pphi2-EuclideanPlaneBackground-dim, Pphi2-ContinuumTestFunction}
  ``\texttt{lean def continuumGreenBilinear (mass : ℝ) (f g : ContinuumTestFunction d) : ℝ }`` Continuum Green's function $(2\pi)^{-d} \int \hat{f}(k)\hat{g}(k)/(|k|^2 + m^2)\, dk$.
\end{definition}
\begin{theorem}[\texttt{gaussianContinuumMeasure\_isProbability}]\label{Pphi2-gaussianContinuumMeasure-isProbability}
  \lean{Pphi2.gaussianContinuumMeasure_isProbability}\leanok
  \uses{Pphi2-EuclideanPlaneBackground-SpaceTime, Pphi2-gaussianContinuumMeasure, Pphi2-planeBackground, GaussianField-FinLatticeSites, Pphi2-latticeEmbedLift-measurable, Pphi2-EuclideanPlaneBackground-dim, GaussianField-latticeGaussianMeasure-isProbability, GaussianField-latticeGaussianMeasure, GaussianField-Configuration, GaussianField-instMeasurableSpaceConfiguration, Pphi2-latticeEmbedLift, Pphi2-ContinuumTestFunction, GaussianField-FinLatticeField}
  $\nu_{\text{GFF},a}$ is a probability measure.
\end{theorem}
\begin{theorem}[\texttt{embeddedTwoPoint\_eq\_latticeGreenBilinear}]\label{Pphi2-embeddedTwoPoint-eq-latticeGreenBilinear}
  \lean{Pphi2.embeddedTwoPoint_eq_latticeGreenBilinear}\leanok
  \uses{Pphi2-latticeTestField, GaussianField-covariance, GaussianField-FinLatticeSites, GaussianField-latticeCovarianceGJ, GaussianField-lattice-covariance-GJ-eq-spectral, Pphi2-embeddedTwoPoint-eq-lattice-cross-moment, Pphi2-embeddedTwoPoint, GaussianField-latticeGaussianMeasure, GaussianField-Configuration, GaussianField-massEigenvalues, Pphi2-latticeGreenBilinear, GaussianField-massEigenvectorBasis, GaussianField-lattice-cross-moment, GaussianField-instMeasurableSpaceConfiguration, Pphi2-ContinuumTestFunction, GaussianField-ell2-, GaussianField-FinLatticeField}
  \texttt{Pphi2.embeddedTwoPoint\_eq\_latticeGreenBilinear}
\end{theorem}
\begin{theorem}[\texttt{embeddedTwoPoint\_eq\_covariance}]\label{Pphi2-embeddedTwoPoint-eq-covariance}
  \lean{Pphi2.embeddedTwoPoint_eq_covariance}\leanok
  \uses{Pphi2-EuclideanPlaneBackground-SpaceTime, Pphi2-gaussianContinuumMeasure, Pphi2-latticeEmbed, Pphi2-planeBackground, GaussianField-FinLatticeSites, Pphi2-latticeEmbedLift-measurable, Pphi2-EuclideanPlaneBackground-dim, Pphi2-embeddedTwoPoint, GaussianField-latticeGaussianMeasure, GaussianField-configuration-eval-measurable, GaussianField-Configuration, GaussianField-instMeasurableSpaceConfiguration, Pphi2-latticeEmbedLift, Pphi2-ContinuumTestFunction, GaussianField-FinLatticeField}
  Rewrites the pushforward integral as a lattice Gaussian integral.
\end{theorem}
\begin{definition}[\texttt{gaussianContinuumMeasure}]\label{Pphi2-gaussianContinuumMeasure}
  \lean{Pphi2.gaussianContinuumMeasure}\leanok
  \uses{Pphi2-EuclideanPlaneBackground-SpaceTime, Pphi2-planeBackground, GaussianField-FinLatticeSites, Pphi2-EuclideanPlaneBackground-dim, GaussianField-latticeGaussianMeasure, GaussianField-Configuration, GaussianField-instMeasurableSpaceConfiguration, Pphi2-latticeEmbedLift, Pphi2-ContinuumTestFunction, GaussianField-FinLatticeField}
  ``\texttt{lean def gaussianContinuumMeasure (a mass : ℝ) (ha : 0 < a) (hmass : 0 < mass) : Measure (Configuration (ContinuumTestFunction d)) }`` Pushforward $\nu_{\text{GFF},a} = (\tilde\iota_a)_* \mu_{\text{GFF},a}$.
\end{definition}
\begin{theorem}[\texttt{congr\_simp}]\label{Pphi2-latticeGreenBilinear-congr-simp}
  \lean{Pphi2.latticeGreenBilinear.congr_simp}\leanok
  \uses{Pphi2-latticeGreenBilinear, Pphi2-ContinuumTestFunction}
  \texttt{Pphi2.latticeGreenBilinear.congr\_simp}
\end{theorem}
\begin{theorem}[\texttt{embeddedTwoPoint\_eq\_lattice\_cross\_moment}]\label{Pphi2-embeddedTwoPoint-eq-lattice-cross-moment}
  \lean{Pphi2.embeddedTwoPoint_eq_lattice_cross_moment}\leanok
  \uses{Pphi2-EuclideanPlaneBackground-SpaceTime, Pphi2-gaussianContinuumMeasure, Pphi2-latticeTestField, Pphi2-planeBackground, GaussianField-FinLatticeSites, Pphi2-latticeEmbedLift-measurable, Pphi2-EuclideanPlaneBackground-dim, Pphi2-latticeEmbedLift-eval-eq, Pphi2-embeddedTwoPoint, GaussianField-latticeGaussianMeasure, GaussianField-configuration-eval-measurable, GaussianField-Configuration, GaussianField-instMeasurableSpaceConfiguration, Pphi2-latticeEmbedLift, Pphi2-ContinuumTestFunction, GaussianField-FinLatticeField}
  \texttt{Pphi2.embeddedTwoPoint\_eq\_lattice\_cross\_moment}
\end{theorem}
\begin{theorem}[\texttt{embeddedTwoPoint\_eq\_spectral\_sum}]\label{Pphi2-embeddedTwoPoint-eq-spectral-sum}
  \lean{Pphi2.embeddedTwoPoint_eq_spectral_sum}\leanok
  \uses{Pphi2-embeddedTwoPoint-eq-latticeGreenBilinear, Pphi2-embeddedTwoPoint, Pphi2-latticeGreenBilinear, Pphi2-ContinuumTestFunction}
  \texttt{Pphi2.embeddedTwoPoint\_eq\_spectral\_sum}
\end{theorem}
\begin{theorem}[\texttt{embeddedTwoPoint\_eq\_latticeSum}]\label{Pphi2-embeddedTwoPoint-eq-latticeSum}
  \lean{Pphi2.embeddedTwoPoint_eq_latticeSum}\leanok
  \uses{Pphi2-EuclideanPlaneBackground-SpaceTime, GaussianField-finLatticeField-dyninMityaginSpace, Pphi2-latticeEmbed, GaussianField-measure, GaussianField-ell2-separableSpace, Pphi2-planeBackground, GaussianField-covariance, GaussianField-FinLatticeSites, Pphi2-evalAtSite, GaussianField-latticeCovarianceGJ, Pphi2-EuclideanPlaneBackground-dim, Pphi2-embeddedTwoPoint, GaussianField-latticeGaussianMeasure, GaussianField-Configuration, GaussianField-lattice-cross-moment, GaussianField-instMeasurableSpaceConfiguration, Pphi2-ContinuumTestFunction, GaussianField-ell2-, GaussianField-FinLatticeField, Pphi2-embeddedTwoPoint-eq-covariance, GaussianField-pairing-product-integrable}
  $\langle\Phi_a(f),\Phi_a(g)\rangle = a^{2d} \sum_{x,y} C_a(x,y) f(ax) g(ay)$. Proved by expanding the product, swapping integral and sums, and applying the Gaussian cross-moment formula.
\end{theorem}
\section{\texttt{Pphi2.AsymTorus.AsymTorusOS}}
\begin{definition}[\texttt{AsymTorusOS2\_TranslationInvariance}]\label{Pphi2-AsymTorusOS2-TranslationInvariance}
  \lean{Pphi2.AsymTorusOS2_TranslationInvariance}\leanok
  \uses{GaussianField-NuclearTensorProduct-instModuleReal, GaussianField-NuclearTensorProduct-instTopologicalSpace, GaussianField-NuclearTensorProduct-instIsTopologicalAddGroup, GaussianField-asymTorusTranslation, GaussianField-NuclearTensorProduct-instAddCommGroup, Pphi2-asymTorusGeneratingFunctional, GaussianField-Configuration, GaussianField-NuclearTensorProduct-instContinuousSMulReal, GaussianField-SmoothMap-Circle, GaussianField-instMeasurableSpaceConfiguration, GaussianField-AsymTorusTestFunction}
  ``\texttt{lean def AsymTorusOS2\_TranslationInvariance (μ : Measure (Configuration (AsymTorusTestFunction Lt Ls))) [IsProbabilityMeasure μ] : Prop }``  OS2 translation invariance: $Z_\mu(f) = Z_\mu(T_v f)$ for all $v \in \mathbb{R}^2$.
\end{definition}
\begin{theorem}[\texttt{asymTorusInteracting\_satisfies\_OS}]\label{Pphi2-asymTorusInteracting-satisfies-OS}
  \lean{Pphi2.asymTorusInteracting_satisfies_OS}\leanok
  \uses{Pphi2-asymTorusInteractingMeasure, GaussianField-NuclearTensorProduct-instModuleReal, GaussianField-NuclearTensorProduct-instTopologicalSpace, Pphi2-AsymSatisfiesTorusOS, GaussianField-NuclearTensorProduct-instIsTopologicalAddGroup, GaussianField-NuclearTensorProduct-instAddCommGroup, GaussianField-Configuration, GaussianField-NuclearTensorProduct-instContinuousSMulReal, GaussianField-SmoothMap-Circle, GaussianField-instMeasurableSpaceConfiguration, GaussianField-AsymTorusTestFunction}
  ``\texttt{lean theorem asymTorusInteracting\_satisfies\_OS (P : InteractionPolynomial) (mass : ℝ) (hmass : 0 < mass) (μ : Measure (Configuration (AsymTorusTestFunction Lt Ls))) [IsProbabilityMeasure μ] (φ : ℕ → ℕ) (hφ : StrictMono φ) (hconv : ...) : AsymSatisfiesTorusOS Lt Ls μ }``  Informal statement: The asymmetric torus $P(\varphi)_2$ interacting continuum limit satisfies OS0--OS2 (analyticity, regularity, translation invariance, time reflection invariance).  Proof: Fully proved by assembling \texttt{asymTorusInteracting\_os0}, \texttt{asymTorusInteracting\_os1}, \texttt{asymTorusInteracting\_os2\_translation}, \texttt{asymTorusInteracting\_os2\_timeReflection}.
\end{theorem}
\begin{definition}[\texttt{AsymTorusOS0\_Analyticity}]\label{Pphi2-AsymTorusOS0-Analyticity}
  \lean{Pphi2.AsymTorusOS0_Analyticity}\leanok
  \uses{GaussianField-NuclearTensorProduct-instModuleReal, GaussianField-NuclearTensorProduct-instTopologicalSpace, GaussianField-NuclearTensorProduct-instIsTopologicalAddGroup, GaussianField-NuclearTensorProduct-instAddCommGroup, GaussianField-Configuration, GaussianField-NuclearTensorProduct-instContinuousSMulReal, GaussianField-SmoothMap-Circle, GaussianField-instMeasurableSpaceConfiguration, GaussianField-AsymTorusTestFunction}
  ``\texttt{lean def AsymTorusOS0\_Analyticity (μ : Measure (Configuration (AsymTorusTestFunction Lt Ls))) [IsProbabilityMeasure μ] : Prop }``  OS0: The multi-variable generating functional $(z_1, \ldots, z_n) \mapsto \int \exp(i\, \omega(\sum_i (\Re z_i) J_i + i\, (\Im z_i) J_i))\, d\mu$ is analytic on $\mathbb{C}^n$.
\end{definition}
\begin{theorem}[\texttt{asymGf\_sub\_norm\_le\_seminorm}]\label{Pphi2-asymGf-sub-norm-le-seminorm}
  \lean{Pphi2.asymGf_sub_norm_le_seminorm}\leanok
  \uses{Pphi2-asymTorusInteractingMeasure, GaussianField-NuclearTensorProduct-instModuleReal, GaussianField-NuclearTensorProduct-instTopologicalSpace, GaussianField-finLatticeField-dyninMityaginSpace, GaussianField-measure, GaussianField-ell2-separableSpace, Pphi2-interactionFunctional-bounded-below, GaussianField-FinLatticeSites, GaussianField-NuclearTensorProduct-instIsTopologicalAddGroup, Pphi2-partitionFunction, GaussianField-latticeCovarianceGJ, Pphi2-interactionFunctional-measurable, GaussianField-pairing-memLp, Pphi2-interactingLatticeMeasure, Pphi2-asymTorusEmbedLift-eval-eq, Pphi2-asymGeomSpacing-pos, GaussianField-NuclearTensorProduct-instAddCommGroup, Pphi2-asymTorusEmbedLift-measurable, GaussianField-latticeGaussianMeasure, Pphi2-asymTorusGeneratingFunctional, GaussianField-configuration-eval-measurable, GaussianField-Configuration, Pphi2-asymTorusInteractingMeasure-isProbability, GaussianField-NuclearTensorProduct-instContinuousSMulReal, GaussianField-SmoothMap-Circle, Pphi2-asymLatticeTestFn, GaussianField-instMeasurableSpaceConfiguration, Pphi2-asymTorusEmbedLift, Pphi2-boltzmannWeight, GaussianField-AsymTorusTestFunction, Pphi2-asymTorus-interacting-second-moment-density-transfer, GaussianField-ell2-, GaussianField-FinLatticeField, Pphi2-interactionFunctional, Pphi2-asymGeomSpacing, Pphi2-partitionFunction-pos, Lattice-toSemilatticeInf, Pphi2-boltzmannWeight-pos}
  ``\texttt{lean theorem asymGf\_sub\_norm\_le\_seminorm (P : InteractionPolynomial) (mass : ℝ) (hmass : 0 < mass) : ∃ (B : ℝ) (p : AsymTorusTestFunction Lt Ls → ℝ), Continuous p ∧ p 0 = 0 ∧ ∀ (N : ℕ) [NeZero N] (g h : AsymTorusTestFunction Lt Ls), ‖Z\_N(g) - Z\_N(h)‖ ≤ B * p (g - h) }``  Informal statement: Uniform Lipschitz bound on the cutoff generating functional: $\|Z_N(g) - Z_N(h)\| \le B \cdot p(g - h)$ with $p$ continuous and $p(0) = 0$.  Proof: Fully proved via Cauchy--Schwarz, mean value theorem for $\sin/\cos$, density transfer, hypercontractivity, and the tight Riemann sum bound \texttt{asymLatticeTestFn\_norm\_sq\_tight}.  Dependencies: \texttt{asymLatticeTestFn\_norm\_sq\_tight}, \texttt{density\_transfer\_bound}, \texttt{gaussian\_hypercontractive}, \texttt{asymNelson\_exponential\_estimate}.
\end{theorem}
\begin{definition}[\texttt{asymTorusGeneratingFunctional}]\label{Pphi2-asymTorusGeneratingFunctional}
  \lean{Pphi2.asymTorusGeneratingFunctional}\leanok
  \uses{GaussianField-NuclearTensorProduct-instModuleReal, GaussianField-NuclearTensorProduct-instTopologicalSpace, GaussianField-NuclearTensorProduct-instIsTopologicalAddGroup, GaussianField-NuclearTensorProduct-instAddCommGroup, GaussianField-Configuration, GaussianField-NuclearTensorProduct-instContinuousSMulReal, GaussianField-SmoothMap-Circle, GaussianField-instMeasurableSpaceConfiguration, GaussianField-AsymTorusTestFunction}
  ``\texttt{lean def asymTorusGeneratingFunctional (μ : Measure (Configuration (AsymTorusTestFunction Lt Ls))) [IsProbabilityMeasure μ] (f : AsymTorusTestFunction Lt Ls) : ℂ }``  The generating (characteristic) functional on the asymmetric torus: $Z_\mu(f) = \int \exp(i\, \omega(f))\, d\mu$.
\end{definition}
\begin{theorem}[\texttt{os2\_timeReflection}]\label{Pphi2-AsymSatisfiesTorusOS-os2-timeReflection}
  \lean{Pphi2.AsymSatisfiesTorusOS.os2_timeReflection}\leanok
  \uses{GaussianField-NuclearTensorProduct-instModuleReal, GaussianField-NuclearTensorProduct-instTopologicalSpace, Pphi2-AsymSatisfiesTorusOS, GaussianField-NuclearTensorProduct-instIsTopologicalAddGroup, GaussianField-NuclearTensorProduct-instAddCommGroup, Pphi2-AsymTorusOS2-TimeReflectionInvariance, GaussianField-Configuration, GaussianField-NuclearTensorProduct-instContinuousSMulReal, GaussianField-SmoothMap-Circle, GaussianField-instMeasurableSpaceConfiguration, GaussianField-AsymTorusTestFunction}
  \texttt{Pphi2.AsymSatisfiesTorusOS.os2\_timeReflection}
\end{theorem}
\begin{theorem}[\texttt{os1}]\label{Pphi2-AsymSatisfiesTorusOS-os1}
  \lean{Pphi2.AsymSatisfiesTorusOS.os1}\leanok
  \uses{GaussianField-NuclearTensorProduct-instModuleReal, GaussianField-NuclearTensorProduct-instTopologicalSpace, Pphi2-AsymSatisfiesTorusOS, GaussianField-NuclearTensorProduct-instIsTopologicalAddGroup, Pphi2-AsymTorusOS1-Regularity, GaussianField-NuclearTensorProduct-instAddCommGroup, GaussianField-Configuration, GaussianField-NuclearTensorProduct-instContinuousSMulReal, GaussianField-SmoothMap-Circle, GaussianField-instMeasurableSpaceConfiguration, GaussianField-AsymTorusTestFunction}
  \texttt{Pphi2.AsymSatisfiesTorusOS.os1}
\end{theorem}
\begin{theorem}[\texttt{asymTorusInteractingMeasure\_gf\_latticeTranslation\_invariant}]\label{Pphi2-asymTorusInteractingMeasure-gf-latticeTranslation-invariant}
  \lean{Pphi2.asymTorusInteractingMeasure_gf_latticeTranslation_invariant}\leanok
  \uses{Pphi2-asymTorusInteractingMeasure, GaussianField-NuclearTensorProduct-instModuleReal, GaussianField-NuclearTensorProduct-instTopologicalSpace, Pphi2-evalAsymAtFinSiteGJ-apply, GaussianField-FinLatticeSites, GaussianField-NuclearTensorProduct-instIsTopologicalAddGroup, GaussianField-asymTorusTranslation, Pphi2-interactingLatticeMeasure, GaussianField-translateSites, Pphi2-asymTorusEmbedLift-eval-eq, Pphi2-asymGeomSpacing-pos, GaussianField-NuclearTensorProduct-instAddCommGroup, Pphi2-evalAsymAtFinSite, Pphi2-asymTorusEmbedLift-measurable, Pphi2-evalAsymAtFinSiteGJ, Pphi2-asymTorusGeneratingFunctional, GaussianField-configuration-eval-measurable, GaussianField-Configuration, Pphi2-asymTorusInteractingMeasure-isProbability, GaussianField-NuclearTensorProduct-instContinuousSMulReal, GaussianField-SmoothMap-Circle, GaussianField-circleSpacing, Pphi2-asymLatticeTestFn, GaussianField-instMeasurableSpaceConfiguration, Pphi2-asymTorusEmbedLift, GaussianField-AsymTorusTestFunction, GaussianField-FinLatticeField, Pphi2-asymGeomSpacing}
  ``\texttt{lean theorem asymTorusInteractingMeasure\_gf\_latticeTranslation\_invariant (N : ℕ) [NeZero N] (P : InteractionPolynomial) (mass : ℝ) (hmass : 0 < mass) (j₁ j₂ : ℤ) (f : AsymTorusTestFunction Lt Ls) : asymTorusGeneratingFunctional Lt Ls (asymTorusInteractingMeasure Lt Ls N P mass hmass) f = asymTorusGeneratingFunctional Lt Ls (asymTorusInteractingMeasure Lt Ls N P mass hmass) (asymTorusTranslation Lt Ls (circleSpacing Lt N * j₁, circleSpacing Ls N * j₂) f) }``  Informal statement: The cutoff generating functional is invariant under lattice translations: $Z_N(f) = Z_N(T_{(j_1 a_t, j_2 a_s)} f)$.  Proof: Fully proved from \texttt{evalAsymAtFinSite\_latticeTranslation} and \texttt{asymInteractingLatticeMeasure\_translation\_invariant}.
\end{theorem}
\begin{definition}[\texttt{AsymTorusOS1\_Regularity}]\label{Pphi2-AsymTorusOS1-Regularity}
  \lean{Pphi2.AsymTorusOS1_Regularity}\leanok
  \uses{GaussianField-NuclearTensorProduct-instModuleReal, GaussianField-NuclearTensorProduct-instTopologicalSpace, GaussianField-NuclearTensorProduct-instIsTopologicalAddGroup, GaussianField-NuclearTensorProduct-instAddCommGroup, GaussianField-Configuration, GaussianField-NuclearTensorProduct-instContinuousSMulReal, GaussianField-SmoothMap-Circle, GaussianField-instMeasurableSpaceConfiguration, GaussianField-AsymTorusTestFunction}
  ``\texttt{lean def AsymTorusOS1\_Regularity (μ : Measure (Configuration (AsymTorusTestFunction Lt Ls))) [IsProbabilityMeasure μ] : Prop }``  OS1: The generating functional is bounded by $\|Z_\mu(f_{\mathrm{re}}, f_{\mathrm{im}})\| \le \exp(c\, (q(f_{\mathrm{re}}) + q(f_{\mathrm{im}})))$ for a continuous function $q$ and constant $c > 0$.
\end{definition}
\begin{theorem}[\texttt{os2\_translation}]\label{Pphi2-AsymSatisfiesTorusOS-os2-translation}
  \lean{Pphi2.AsymSatisfiesTorusOS.os2_translation}\leanok
  \uses{GaussianField-NuclearTensorProduct-instModuleReal, GaussianField-NuclearTensorProduct-instTopologicalSpace, Pphi2-AsymTorusOS2-TranslationInvariance, Pphi2-AsymSatisfiesTorusOS, GaussianField-NuclearTensorProduct-instIsTopologicalAddGroup, GaussianField-NuclearTensorProduct-instAddCommGroup, GaussianField-Configuration, GaussianField-NuclearTensorProduct-instContinuousSMulReal, GaussianField-SmoothMap-Circle, GaussianField-instMeasurableSpaceConfiguration, GaussianField-AsymTorusTestFunction}
  \texttt{Pphi2.AsymSatisfiesTorusOS.os2\_translation}
\end{theorem}
\begin{theorem}[\texttt{os0}]\label{Pphi2-AsymSatisfiesTorusOS-os0}
  \lean{Pphi2.AsymSatisfiesTorusOS.os0}\leanok
  \uses{GaussianField-NuclearTensorProduct-instModuleReal, GaussianField-NuclearTensorProduct-instTopologicalSpace, Pphi2-AsymSatisfiesTorusOS, GaussianField-NuclearTensorProduct-instIsTopologicalAddGroup, GaussianField-NuclearTensorProduct-instAddCommGroup, Pphi2-AsymTorusOS0-Analyticity, GaussianField-Configuration, GaussianField-NuclearTensorProduct-instContinuousSMulReal, GaussianField-SmoothMap-Circle, GaussianField-instMeasurableSpaceConfiguration, GaussianField-AsymTorusTestFunction}
  \texttt{Pphi2.AsymSatisfiesTorusOS.os0}
\end{theorem}
\begin{theorem}[\texttt{congr\_simp}]\label{Pphi2-asymTorusInteractingMeasure-congr-simp}
  \lean{Pphi2.asymTorusInteractingMeasure.congr_simp}\leanok
  \uses{Pphi2-asymTorusInteractingMeasure, GaussianField-NuclearTensorProduct-instModuleReal, GaussianField-NuclearTensorProduct-instTopologicalSpace, GaussianField-NuclearTensorProduct-instIsTopologicalAddGroup, GaussianField-NuclearTensorProduct-instAddCommGroup, GaussianField-Configuration, GaussianField-NuclearTensorProduct-instContinuousSMulReal, GaussianField-SmoothMap-Circle, GaussianField-instMeasurableSpaceConfiguration, GaussianField-AsymTorusTestFunction}
  \texttt{Pphi2.asymTorusInteractingMeasure.congr\_simp}
\end{theorem}
\begin{definition}[\texttt{AsymSatisfiesTorusOS}]\label{Pphi2-AsymSatisfiesTorusOS}
  \lean{Pphi2.AsymSatisfiesTorusOS}\leanok
  \uses{GaussianField-NuclearTensorProduct-instModuleReal, GaussianField-NuclearTensorProduct-instTopologicalSpace, GaussianField-NuclearTensorProduct-instIsTopologicalAddGroup, GaussianField-NuclearTensorProduct-instAddCommGroup, GaussianField-Configuration, GaussianField-NuclearTensorProduct-instContinuousSMulReal, GaussianField-SmoothMap-Circle, GaussianField-instMeasurableSpaceConfiguration, GaussianField-AsymTorusTestFunction}
  ``\texttt{lean structure AsymSatisfiesTorusOS (μ : Measure (Configuration (AsymTorusTestFunction Lt Ls))) [IsProbabilityMeasure μ] where os0 : AsymTorusOS0\_Analyticity Lt Ls μ os1 : AsymTorusOS1\_Regularity Lt Ls μ os2\_translation : AsymTorusOS2\_TranslationInvariance Lt Ls μ os2\_timeReflection : AsymTorusOS2\_TimeReflectionInvariance Lt Ls μ }``  Bundled OS axioms for the asymmetric torus.
\end{definition}
\begin{theorem}[\texttt{asymTorusTranslation\_continuous\_in\_v}]\label{Pphi2-asymTorusTranslation-continuous-in-v}
  \lean{Pphi2.asymTorusTranslation_continuous_in_v}\leanok
  \uses{GaussianField-sobolevSeminorm-affine-precomp-le, GaussianField-mapImage, GaussianField-NuclearTensorProduct-instModuleReal, GaussianField-circleTranslation, Lattice-toSemilatticeSup, GaussianField-NuclearTensorProduct-instTopologicalSpace, GaussianField-SmoothMap-Circle-instFunLike, GaussianField-RapidDecaySeq-rapidDecaySeminorm, GaussianField-RapidDecaySeq, GaussianField-nuclearTensorProduct-mapCLM, GaussianField-RapidDecaySeq-val, GaussianField-RapidDecaySeq-hasSum-basisVec, GaussianField-SmoothMap-Circle-instAddCommGroup, GaussianField-NuclearTensorProduct-instIsTopologicalAddGroup, GaussianField-nuclearTensorProduct-mapCLM-pure, GaussianField-RapidDecaySeq-instSMul, GaussianField-RapidDecaySeq-instModule, GaussianField-DyninMityaginSpace--, GaussianField-asymTorusTranslation, GaussianField-NuclearTensorProduct, GaussianField-DyninMityaginSpace-HasBiorthogonalBasis-coeff-basis, GaussianField-RapidDecaySeq-basisVec, GaussianField-NuclearTensorProduct-instAddCommGroup, GaussianField-NuclearTensorProduct-pure-continuous, GaussianField-finset-sup-poly-bound, GaussianField-smoothCircle-hasBiorthogonalBasis, GaussianField-SmoothMap-Circle-instModule, GaussianField-SmoothMap-Circle-instContinuousSMul, GaussianField-smoothCircle-dyninMityaginSpace, GaussianField-SmoothMap-Circle-instIsTopologicalAddGroup, GaussianField-DyninMityaginSpace-basis, GaussianField-NuclearTensorProduct-pure, GaussianField-RapidDecaySeq-instAddCommGroup, GaussianField-SmoothMap-Circle, GaussianField-DyninMityaginSpace-basis-growth, GaussianField-RapidDecaySeq-instTopologicalSpace, GaussianField-NuclearTensorProduct-basisVec-eq-pure, GaussianField-DyninMityaginSpace-p, GaussianField-AsymTorusTestFunction, GaussianField-RapidDecaySeq-rapidDecay-withSeminorms, GaussianField-RapidDecaySeq-rapid-decay, GaussianField-NuclearTensorProduct-pure-seminorm-bound, GaussianField-SmoothMap-Circle-instTopologicalSpace, Lattice-toSemilatticeInf}
  ``\texttt{lean theorem asymTorusTranslation\_continuous\_in\_v (f : AsymTorusTestFunction Lt Ls) : Continuous (fun v : ℝ × ℝ => asymTorusTranslation Lt Ls v f) }``  Informal statement: The map $v \mapsto T_v f$ is continuous in the NTP topology.  Proof: Fully proved by reducing to \texttt{WithSeminorms} characterization, bounding NTP seminorms of $T_v f - T_{v_0} f$ via pure tensor bounds, and using \texttt{circleTranslation\_continuous\_in\_s} for each factor.  Dependencies: \texttt{circleTranslation\_continuous\_in\_s}, \texttt{NuclearTensorProduct.pure\_seminorm\_bound}, \texttt{RapidDecaySeq.rapidDecay\_withSeminorms}.
\end{theorem}
\begin{definition}[\texttt{AsymTorusOS2\_TimeReflectionInvariance}]\label{Pphi2-AsymTorusOS2-TimeReflectionInvariance}
  \lean{Pphi2.AsymTorusOS2_TimeReflectionInvariance}\leanok
  \uses{GaussianField-NuclearTensorProduct-instModuleReal, GaussianField-NuclearTensorProduct-instTopologicalSpace, GaussianField-NuclearTensorProduct-instIsTopologicalAddGroup, GaussianField-NuclearTensorProduct-instAddCommGroup, Pphi2-asymTorusGeneratingFunctional, GaussianField-Configuration, GaussianField-NuclearTensorProduct-instContinuousSMulReal, GaussianField-SmoothMap-Circle, GaussianField-instMeasurableSpaceConfiguration, GaussianField-AsymTorusTestFunction, GaussianField-asymTorusTimeReflection}
  ``\texttt{lean def AsymTorusOS2\_TimeReflectionInvariance (μ : Measure (Configuration (AsymTorusTestFunction Lt Ls))) [IsProbabilityMeasure μ] : Prop }``  OS2 time reflection: $Z_\mu(f) = Z_\mu(\Theta f)$ where $\Theta: (t,x) \mapsto (-t,x)$.
\end{definition}
\begin{theorem}[\texttt{congr\_simp}]\label{Pphi2-asymTorusGeneratingFunctional-congr-simp}
  \lean{Pphi2.asymTorusGeneratingFunctional.congr_simp}\leanok
  \uses{GaussianField-NuclearTensorProduct-instModuleReal, GaussianField-NuclearTensorProduct-instTopologicalSpace, GaussianField-NuclearTensorProduct-instIsTopologicalAddGroup, GaussianField-NuclearTensorProduct-instAddCommGroup, Pphi2-asymTorusGeneratingFunctional, GaussianField-Configuration, GaussianField-NuclearTensorProduct-instContinuousSMulReal, GaussianField-SmoothMap-Circle, GaussianField-instMeasurableSpaceConfiguration, GaussianField-AsymTorusTestFunction}
  \texttt{Pphi2.asymTorusGeneratingFunctional.congr\_simp}
\end{theorem}
\begin{theorem}[\texttt{asymTorusGF\_latticeApproximation\_error\_vanishes}]\label{Pphi2-asymTorusGF-latticeApproximation-error-vanishes}
  \lean{Pphi2.asymTorusGF_latticeApproximation_error_vanishes}\leanok
  \uses{Pphi2-asymTorusInteractingMeasure, GaussianField-NuclearTensorProduct-instModuleReal, GaussianField-NuclearTensorProduct-instTopologicalSpace, Pphi2-asymTorusTranslation-continuous-in-v, GaussianField-NuclearTensorProduct-instIsTopologicalAddGroup, Pphi2-asymTorusInteractingMeasure-gf-latticeTranslation-invariant, GaussianField-asymTorusTranslation, GaussianField-NuclearTensorProduct-instAddCommGroup, Pphi2-asymTorusGeneratingFunctional, Pphi2-asymTorusInteractingMeasure-isProbability, GaussianField-SmoothMap-Circle, GaussianField-circleSpacing, Pphi2-asymGf-sub-norm-le-seminorm, GaussianField-AsymTorusTestFunction, Lattice-toSemilatticeInf, GaussianField-circleSpacing-pos}
  ``\texttt{lean theorem asymTorusGF\_latticeApproximation\_error\_vanishes (P : InteractionPolynomial) (mass : ℝ) (hmass : 0 < mass) (φ : ℕ → ℕ) (hφ : StrictMono φ) (v : ℝ × ℝ) (f : AsymTorusTestFunction Lt Ls) : Tendsto (fun n => Z\_\{φ(n)+1\}(T\_v f) - Z\_\{φ(n)+1\}(f)) atTop (nhds 0) }``  Informal statement: The difference $Z_N(T_v f) - Z_N(f) \to 0$ as $N \to \infty$, via lattice rounding and generating functional continuity.  Proof: Fully proved. Decomposes $Z_N(T_v f) - Z_N(f) = (Z_N(T_v f) - Z_N(T_{w_N} f)) + 0$ where $w_N$ is the nearest lattice point to $v$, uses \texttt{asymTorusInteractingMeasure\_gf\_latticeTranslation\_invariant} for the second term and \texttt{asymGf\_sub\_norm\_le\_seminorm} with \texttt{asymTorusTranslation\_continuous\_in\_v} for the first.  Dependencies: \texttt{asymGf\_sub\_norm\_le\_seminorm}, \texttt{asymTorusInteractingMeasure\_gf\_latticeTranslation\_invariant}, \texttt{asymTorusTranslation\_continuous\_in\_v}.
\end{theorem}
\section{\texttt{Pphi2.IRLimit.GreenFunctionComparison}}
\begin{theorem}[\texttt{cylinderPullback\_second\_moment\_density\_transfer\_cutoff}]\label{Pphi2-cylinderPullback-second-moment-density-transfer-cutoff}
  \lean{Pphi2.cylinderPullback_second_moment_density_transfer_cutoff}\leanok
  \uses{Pphi2-asymTorusInteractingMeasure, GaussianField-NuclearTensorProduct-instModuleReal, GaussianField-NuclearTensorProduct-instTopologicalSpace, GaussianField-FinLatticeSites, GaussianField-CylinderTestFunction, GaussianField-NuclearTensorProduct-instIsTopologicalAddGroup, Pphi2-cylinderToTorusEmbed, Pphi2-asymGeomSpacing-pos, GaussianField-NuclearTensorProduct-instAddCommGroup, Pphi2-cylinderPullback-second-moment-eq, Pphi2-cylinderPullbackMeasure, GaussianField-latticeGaussianMeasure, GaussianField-Configuration, Pphi2-asymTorusInteractingMeasure-isProbability, GaussianField-NuclearTensorProduct-instContinuousSMulReal, GaussianField-SmoothMap-Circle, Pphi2-asymLatticeTestFn, GaussianField-instMeasurableSpaceConfiguration, GaussianField-AsymTorusTestFunction, Pphi2-asymTorus-interacting-second-moment-density-transfer, GaussianField-FinLatticeField, Pphi2-asymGeomSpacing}
  \texttt{Pphi2.cylinderPullback\_second\_moment\_density\_transfer\_cutoff}
\end{theorem}
\begin{theorem}[\texttt{cylinderPullback\_second\_moment\_eq}]\label{Pphi2-cylinderPullback-second-moment-eq}
  \lean{Pphi2.cylinderPullback_second_moment_eq}\leanok
  \uses{GaussianField-NuclearTensorProduct-instModuleReal, GaussianField-NuclearTensorProduct-instTopologicalSpace, Pphi2-cylinderPullback, GaussianField-CylinderTestFunction, GaussianField-NuclearTensorProduct-instIsTopologicalAddGroup, Pphi2-cylinderToTorusEmbed, GaussianField-configuration-measurable-of-eval-measurable, GaussianField-NuclearTensorProduct-instAddCommGroup, Pphi2-cylinderPullbackMeasure, GaussianField-configuration-eval-measurable, GaussianField-Configuration, GaussianField-NuclearTensorProduct-instContinuousSMulReal, GaussianField-SmoothMap-Circle, GaussianField-instMeasurableSpaceConfiguration, GaussianField-AsymTorusTestFunction}
  \texttt{Pphi2.cylinderPullback\_second\_moment\_eq}
\end{theorem}
\section{\texttt{Pphi2.GeneralResults.GaussianHermiteMean}}
\begin{theorem}[\texttt{gaussian\_hermite\_zero\_mean} (axiom)]\label{Pphi2-gaussian-hermite-zero-mean}
  \lean{Pphi2.gaussian_hermite_zero_mean}\leanok
  \uses{Pphi2-wickMonomial}
  Probabilist's Hermite polynomial of order $n \ge 1$ has zero mean under matching Gaussian: $\int \text{He}_n(t; \sigma^2)\, dN(0,\sigma^2)(t) = 0$.
\end{theorem}
\section{\texttt{Pphi2.AsymTorus.AsymGappedTransfer}}
\begin{definition}[\texttt{asymTransferNormalized}]\label{Pphi2-asymTransferNormalized}
  \lean{Pphi2.asymTransferNormalized}\leanok
  \uses{Pphi2-L2SpatialField, Pphi2-SpatialField, Pphi2-asymTransferGroundEigenvalue, Pphi2-asymTransferOperatorCLM}
  \texttt{Pphi2.asymTransferNormalized}
\end{definition}
\begin{definition}[\texttt{asymGappedTransfer}]\label{Pphi2-asymGappedTransfer}
  \lean{Pphi2.asymGappedTransfer}\leanok
  \uses{Pphi2-asymTransferNormalized, Pphi2-L2SpatialField, Pphi2-SpatialField, Pphi2-asymGroundVector}
  \texttt{Pphi2.asymGappedTransfer}
\end{definition}
\begin{theorem}[\texttt{asymTransfer\_susceptibility\_le}]\label{Pphi2-asymTransfer-susceptibility-le}
  \lean{Pphi2.asymTransfer_susceptibility_le}\leanok
  \uses{Pphi2-asymTransferNormalized, Pphi2-L2SpatialField, Pphi2-SpatialField, Pphi2-asymGroundVector, Pphi2-asymGappedTransfer}
  \texttt{Pphi2.asymTransfer\_susceptibility\_le}
\end{theorem}
\begin{theorem}[\texttt{congr\_simp}]\label{Pphi2-asymTransferGroundEigenvalue-congr-simp}
  \lean{Pphi2.asymTransferGroundEigenvalue.congr_simp}\leanok
  \uses{Pphi2-asymTransferGroundEigenvalue}
  \texttt{Pphi2.asymTransferGroundEigenvalue.congr\_simp}
\end{theorem}
\begin{theorem}[\texttt{asymTransferGroundEigenvalue\_pos}]\label{Pphi2-asymTransferGroundEigenvalue-pos}
  \lean{Pphi2.asymTransferGroundEigenvalue_pos}\leanok
  \uses{Pphi2-AsymTransferGroundExcitedData-eigenval, Pphi2-asymTransferOperator-eigenvalues-pos, Pphi2-AsymTransferGroundExcitedData-i-, Pphi2-asymTransferGroundExcitedData, Pphi2-asymTransferGroundEigenvalue, Pphi2-AsymTransferGroundExcitedData-b, Pphi2-AsymTransferGroundExcitedData--, Pphi2-AsymTransferGroundExcitedData-h-eigen}
  \texttt{Pphi2.asymTransferGroundEigenvalue\_pos}
\end{theorem}
\begin{definition}[\texttt{asymGroundVector}]\label{Pphi2-asymGroundVector}
  \lean{Pphi2.asymGroundVector}\leanok
  \uses{Pphi2-L2SpatialField, Pphi2-SpatialField, Pphi2-AsymTransferGroundExcitedData-i-, Pphi2-asymTransferGroundExcitedData, Pphi2-AsymTransferGroundExcitedData-b, Pphi2-AsymTransferGroundExcitedData--}
  \texttt{Pphi2.asymGroundVector}
\end{definition}
\begin{theorem}[\texttt{congr\_simp}]\label{Pphi2-asymGroundVector-congr-simp}
  \lean{Pphi2.asymGroundVector.congr_simp}\leanok
  \uses{Pphi2-L2SpatialField, Pphi2-SpatialField, Pphi2-asymGroundVector}
  \texttt{Pphi2.asymGroundVector.congr\_simp}
\end{theorem}
\section{\texttt{Pphi2.GeneralResults.HilbertSchmidt}}
\begin{theorem}[\texttt{integral\_operator\_l2\_kernel\_compact}]\label{Pphi2-GeneralResults-integral-operator-l2-kernel-compact}
  \lean{Pphi2.GeneralResults.integral_operator_l2_kernel_compact}\leanok
  Hilbert-Schmidt compactness: if $K \in L^2(\mu \otimes \mu)$ then the corresponding integral operator is compact. Proved in \texttt{Pphi2.GeneralResults.HilbertSchmidt} and re-exported from \texttt{L2Operator.lean}. Reference: Reed-Simon I, Theorem VI.23.
\end{theorem}
\begin{theorem}[\texttt{hs\_basis\_norm\_summable}]\label{Pphi2-GeneralResults-hs-basis-norm-summable}
  \lean{Pphi2.GeneralResults.hs_basis_norm_summable}\leanok
  \uses{Lattice-toSemilatticeInf}
  \texttt{Pphi2.GeneralResults.hs\_basis\_norm\_summable}
\end{theorem}
\begin{theorem}[\texttt{isCompactOperator\_of\_basis\_norm\_summable}]\label{Pphi2-GeneralResults-isCompactOperator-of-basis-norm-summable}
  \lean{Pphi2.GeneralResults.isCompactOperator_of_basis_norm_summable}\leanok
  \texttt{Pphi2.GeneralResults.isCompactOperator\_of\_basis\_norm\_summable}
\end{theorem}
\section{\texttt{Pphi2.AsymTorus.AsymVarianceBound}}
\begin{theorem}[\texttt{asymInteractingVariance\_le\_freeVariance\_torus}]\label{Pphi2-asymInteractingVariance-le-freeVariance-torus}
  \lean{Pphi2.asymInteractingVariance_le_freeVariance_torus}\leanok
  \uses{GaussianField-NuclearTensorProduct-instModuleReal, GaussianField-NuclearTensorProduct-instTopologicalSpace, GaussianField-AsymLatticeSites, GaussianField-NuclearTensorProduct-instIsTopologicalAddGroup, Pphi2-asymLatticeTestFnIso, GaussianField-NuclearTensorProduct-instAddCommGroup, Pphi2-asymTorusInteractingMeasureIso, GaussianField-Configuration, GaussianField-NuclearTensorProduct-instContinuousSMulReal, GaussianField-AsymLatticeField, GaussianField-SmoothMap-Circle, Pphi2-latticeGaussianMeasureAsym, GaussianField-instMeasurableSpaceConfiguration, GaussianField-AsymTorusTestFunction, Pphi2-asymTorusIso-interacting-second-moment-density-transfer}
  \texttt{Pphi2.asymInteractingVariance\_le\_freeVariance\_torus}
\end{theorem}
\begin{theorem}[\texttt{asymInteractingVariance\_le\_freeVariance\_lattice}]\label{Pphi2-asymInteractingVariance-le-freeVariance-lattice}
  \lean{Pphi2.asymInteractingVariance_le_freeVariance_lattice}\leanok
  \uses{Pphi2-partitionFunctionAsym-ge-one, GaussianField-AsymLatticeSites, GaussianField-measure, GaussianField-ell2-separableSpace, GaussianField-asymLatticeField-dyninMityaginSpace, Pphi2-asymNelson-exponential-estimate-iso, GaussianField-pairing-memLp, Pphi2-partitionFunctionAsym, GaussianField-configuration-eval-measurable, GaussianField-Configuration, Pphi2-density-transfer-bound-iso, GaussianField-gaussian-hypercontractive, GaussianField-AsymLatticeField, Pphi2-latticeGaussianMeasureAsym, Pphi2-interactingLatticeMeasureAsym, Pphi2-interactionFunctionalAsym, GaussianField-instMeasurableSpaceConfiguration, GaussianField-latticeCovarianceAsymGJ, GaussianField-ell2-, Lattice-toSemilatticeInf}
  \texttt{Pphi2.asymInteractingVariance\_le\_freeVariance\_lattice}
\end{theorem}
\section{\texttt{Pphi2.Main}}
\begin{theorem}[\texttt{pphi2\_existence}]\label{Pphi2-pphi2-existence}
  \lean{Pphi2.pphi2_existence}\leanok
  \uses{Pphi2-EuclideanPlaneBackground-SpaceTime, Pphi2-SatisfiesFullOS, Pphi2-planeBackground, Pphi2-EuclideanPlaneBackground-dim, Pphi2-pphi2-exists, GaussianField-Configuration, GaussianField-instMeasurableSpaceConfiguration, Pphi2-ContinuumTestFunction}
  $\exists \mu$ probability measure on $\mathcal{S}'(\mathbb{R}^2)$ satisfying \texttt{SatisfiesFullOS}. Delegates to \texttt{pphi2\_exists}.
\end{theorem}
\begin{theorem}[\texttt{mass\_reparametrization\_invariance}]\label{Pphi2-mass-reparametrization-invariance}
  \lean{Pphi2.mass_reparametrization_invariance}\leanok
  \uses{Pphi2-EuclideanPlaneBackground-SpaceTime, Pphi2-FieldConfig2, Pphi2-EuclideanPlaneBackground-RealTestFunction, Pphi2-planeBackground, Pphi2-EuclideanPlaneBackground-dim, Pphi2-IsPphi2Limit, GaussianField-instMeasurableSpaceConfiguration}
  The continuum limit is invariant under mass reparametrization via \texttt{shiftQuadratic}.
\end{theorem}
\begin{theorem}[\texttt{bareMassParameter\_positive}]\label{Pphi2-bareMassParameter-positive}
  \lean{Pphi2.bareMassParameter_positive}\leanok
  Honest theorem name: formally only $\exists m_0 > 0$ witnessed by the bare mass hypothesis, with $m_0 = m$. No physical mass-gap identification is proved here.
\end{theorem}
\begin{theorem}[\texttt{mass\_reparametrization\_exists}]\label{Pphi2-mass-reparametrization-exists}
  \lean{Pphi2.mass_reparametrization_exists}\leanok
  \uses{Pphi2-EuclideanPlaneBackground-SpaceTime, Pphi2-FieldConfig2, Pphi2-mass-reparametrization-invariance, Pphi2-EuclideanPlaneBackground-RealTestFunction, Pphi2-planeBackground, Pphi2-pphi2-limit-exists, Pphi2-EuclideanPlaneBackground-dim, Pphi2-IsPphi2Limit, GaussianField-Configuration, GaussianField-instMeasurableSpaceConfiguration, Pphi2-ContinuumTestFunction}
  \texttt{Pphi2.mass\_reparametrization\_exists}
\end{theorem}
\begin{theorem}[\texttt{pphi2\_nontriviality} (axiom)]\label{Pphi2-pphi2-nontriviality}
  \lean{Pphi2.pphi2_nontriviality}\leanok
  \uses{Pphi2-EuclideanPlaneBackground-SpaceTime, Pphi2-planeBackground, Pphi2-EuclideanPlaneBackground-dim, GaussianField-Configuration, GaussianField-instMeasurableSpaceConfiguration, Pphi2-ContinuumTestFunction}
  $\exists \mu$ such that $S_2(f,f) = \int \Phi(f)^2\,d\mu > 0$ for all $f \ne 0$.
\end{theorem}
\begin{theorem}[\texttt{pphi2\_nonGaussianity}]\label{Pphi2-pphi2-nonGaussianity}
  \lean{Pphi2.pphi2_nonGaussianity}\leanok
  \uses{Pphi2-EuclideanPlaneBackground-SpaceTime, Pphi2-continuumLimit-nonGaussian, Pphi2-planeBackground, Pphi2-EuclideanPlaneBackground-dim, GaussianField-Configuration, GaussianField-instMeasurableSpaceConfiguration, Pphi2-ContinuumTestFunction}
  The connected four-point function is nonzero: $\exists f,\; S_4(f,f,f,f) - 3 S_2(f,f)^2 \ne 0$. Uses \texttt{continuumLimit\_nonGaussian} with lattice spacings $a_n = 1/(n+1) \to 0$.
\end{theorem}
\begin{theorem}[\texttt{os\_reconstruction}]\label{Pphi2-os-reconstruction}
  \lean{Pphi2.os_reconstruction}\leanok
  \uses{Pphi2-EuclideanPlaneBackground-SpaceTime, Pphi2-SatisfiesFullOS, Pphi2-planeBackground, Pphi2-EuclideanPlaneBackground-dim, Pphi2-massParameter-positive, GaussianField-Configuration, GaussianField-instMeasurableSpaceConfiguration, Pphi2-ContinuumTestFunction}
  Aliases for \texttt{massParameter\_positive} and \texttt{pphi2\_exists\_os\_and\_massParameter\_positive} respectively.
\end{theorem}
\begin{theorem}[\texttt{pphi2\_exists\_os\_and\_massParameter\_positive}]\label{Pphi2-pphi2-exists-os-and-massParameter-positive}
  \lean{Pphi2.pphi2_exists_os_and_massParameter_positive}\leanok
  \uses{Pphi2-EuclideanPlaneBackground-SpaceTime, Pphi2-SatisfiesFullOS, Pphi2-planeBackground, Pphi2-EuclideanPlaneBackground-dim, Pphi2-massParameter-positive, Pphi2-pphi2-exists, GaussianField-Configuration, GaussianField-instMeasurableSpaceConfiguration, Pphi2-ContinuumTestFunction}
  $\exists \mu$ satisfying \texttt{SatisfiesFullOS} and $\exists m_0 > 0$ as in \texttt{massParameter\_positive}.
\end{theorem}
\begin{theorem}[\texttt{pphi2\_mass\_gap}]\label{Pphi2-pphi2-mass-gap}
  \lean{Pphi2.pphi2_mass_gap}\leanok
  \uses{Pphi2-bareMassParameter-positive}
  Alias of \texttt{bareMassParameter\_positive}.
\end{theorem}
\begin{theorem}[\texttt{massParameter\_positive}]\label{Pphi2-massParameter-positive}
  \lean{Pphi2.massParameter_positive}\leanok
  \uses{Pphi2-EuclideanPlaneBackground-SpaceTime, Pphi2-SatisfiesFullOS, Pphi2-planeBackground, Pphi2-EuclideanPlaneBackground-dim, GaussianField-Configuration, GaussianField-instMeasurableSpaceConfiguration, Pphi2-ContinuumTestFunction}
  $\exists m_0 > 0$ witnessed by the bare mass hypothesis (not deduced from OS axioms; no Minkowski/Wightman layer).
\end{theorem}
\begin{theorem}[\texttt{pphi2\_wightman}]\label{Pphi2-pphi2-wightman}
  \lean{Pphi2.pphi2_wightman}\leanok
  \uses{Pphi2-EuclideanPlaneBackground-SpaceTime, Pphi2-SatisfiesFullOS, Pphi2-planeBackground, Pphi2-EuclideanPlaneBackground-dim, Pphi2-pphi2-exists-os-and-massParameter-positive, GaussianField-Configuration, GaussianField-instMeasurableSpaceConfiguration, Pphi2-ContinuumTestFunction}
  \texttt{Pphi2.pphi2\_wightman}
\end{theorem}
\begin{theorem}[\texttt{pphi2\_nonGaussian}]\label{Pphi2-pphi2-nonGaussian}
  \lean{Pphi2.pphi2_nonGaussian}\leanok
  \uses{Pphi2-EuclideanPlaneBackground-SpaceTime, Pphi2-planeBackground, Pphi2-EuclideanPlaneBackground-dim, GaussianField-Configuration, Pphi2-pphi2-nonGaussianity, GaussianField-instMeasurableSpaceConfiguration, Pphi2-ContinuumTestFunction}
  \texttt{Pphi2.pphi2\_nonGaussian}
\end{theorem}
\begin{theorem}[\texttt{pphi2\_nontrivial}]\label{Pphi2-pphi2-nontrivial}
  \lean{Pphi2.pphi2_nontrivial}\leanok
  \uses{Pphi2-EuclideanPlaneBackground-SpaceTime, Pphi2-pphi2-nontriviality, Pphi2-planeBackground, Pphi2-EuclideanPlaneBackground-dim, GaussianField-Configuration, GaussianField-instMeasurableSpaceConfiguration, Pphi2-ContinuumTestFunction}
  \texttt{Pphi2.pphi2\_nontrivial}
\end{theorem}
\begin{theorem}[\texttt{pphi2\_main}]\label{Pphi2-pphi2-main}
  \lean{Pphi2.pphi2_main}\leanok
  \uses{Pphi2-EuclideanPlaneBackground-SpaceTime, Pphi2-SatisfiesFullOS, Pphi2-planeBackground, Pphi2-EuclideanPlaneBackground-dim, Pphi2-IsPphi2Limit, Pphi2-continuumLimit-satisfies-fullOS, GaussianField-Configuration, GaussianField-instMeasurableSpaceConfiguration, Pphi2-ContinuumTestFunction}
  For any continuum limit measure $\mu$ (satisfying \texttt{IsPphi2Limit}), $\mu$ satisfies all five OS axioms (\texttt{SatisfiesFullOS}). Delegates to \texttt{continuumLimit\_satisfies\_fullOS}.
\end{theorem}
\section{\texttt{Pphi2.OSProofs.OS4\_MassGap}}
\begin{theorem}[\texttt{connectedTwoPoint\_nonneg\_delta}]\label{Pphi2-connectedTwoPoint-nonneg-delta}
  \lean{Pphi2.connectedTwoPoint_nonneg_delta}\leanok
  \uses{GaussianField-finLatticeDelta, GaussianField-FinLatticeSites, Pphi2-interactingLatticeMeasure-isProbability, Pphi2-interactingLatticeMeasure, Pphi2-field-second-moment-finite, Pphi2-field-all-moments-finite, GaussianField-Configuration, GaussianField-instMeasurableSpaceConfiguration, GaussianField-FinLatticeField, Lattice-toSemilatticeInf, Pphi2-connectedTwoPoint}
  $G_c(\delta_x, \delta_x) \ge 0$. This is variance nonnegativity: $\mathrm{Var}[\omega(\delta_x)] \ge 0$. Fully proved via $\int (X-c)^2 \ge 0$.
\end{theorem}
\begin{theorem}[\texttt{general\_clustering\_from\_spectral\_gap} (axiom)]\label{Pphi2-general-clustering-from-spectral-gap}
  \lean{Pphi2.general_clustering_from_spectral_gap}\leanok
  \uses{Pphi2-latticeEuclideanTimeSeparation, GaussianField-FinLatticeSites, Pphi2-latticeConfigEuclideanTimeShift, Pphi2-interactingLatticeMeasure, GaussianField-Configuration, GaussianField-instMeasurableSpaceConfiguration, Pphi2-massGap, GaussianField-FinLatticeField}
  General clustering for bounded observables with \texttt{G} evaluated on the time-shifted configuration: $|\langle F(\omega)\,G(\tau_R\omega)\rangle - \langle F\rangle\langle G\rangle| \le C_{FG} \exp(-m_{\mathrm{phys}} \, d_{\mathrm{cyc}}(R)\, a)$.
\end{theorem}
\begin{definition}[\texttt{connectedTwoPoint}]\label{Pphi2-connectedTwoPoint}
  \lean{Pphi2.connectedTwoPoint}\leanok
  \uses{GaussianField-FinLatticeSites, Pphi2-interactingLatticeMeasure, GaussianField-Configuration, GaussianField-instMeasurableSpaceConfiguration, GaussianField-FinLatticeField}
  ``\texttt{lean def connectedTwoPoint (d N : N) [NeZero N] (P : InteractionPolynomial) (a mass : R) (ha : 0 < a) (hmass : 0 < mass) (f g : FinLatticeField d N) : R }`` Connected two-point function: $G_c(f,g) = \langle\omega(f)\omega(g)\rangle - \langle\omega(f)\rangle\langle\omega(g)\rangle$.
\end{definition}
\begin{theorem}[\texttt{connectedTwoPoint\_symm}]\label{Pphi2-connectedTwoPoint-symm}
  \lean{Pphi2.connectedTwoPoint_symm}\leanok
  \uses{GaussianField-FinLatticeSites, Pphi2-interactingLatticeMeasure, GaussianField-Configuration, GaussianField-instMeasurableSpaceConfiguration, GaussianField-FinLatticeField, Pphi2-connectedTwoPoint}
  $G_c(f,g) = G_c(g,f)$. From commutativity of multiplication. Fully proved.
\end{theorem}
\begin{theorem}[\texttt{two\_point\_clustering\_lattice}]\label{Pphi2-two-point-clustering-lattice}
  \lean{Pphi2.two_point_clustering_lattice}\leanok
  \uses{Pphi2-latticeEuclideanTimeSeparation, GaussianField-finLatticeDelta, GaussianField-FinLatticeSites, Pphi2-interactingLatticeMeasure, GaussianField-Configuration, Pphi2-two-point-clustering-from-spectral-gap, GaussianField-instMeasurableSpaceConfiguration, Pphi2-massGap, GaussianField-FinLatticeField}
  Exponential clustering of the two-point function on the lattice. Delegates to the axiom. Fully proved (modulo axiom).
\end{theorem}
\begin{theorem}[\texttt{two\_point\_clustering\_from\_spectral\_gap} (axiom)]\label{Pphi2-two-point-clustering-from-spectral-gap}
  \lean{Pphi2.two_point_clustering_from_spectral_gap}\leanok
  \uses{Pphi2-latticeEuclideanTimeSeparation, GaussianField-finLatticeDelta, GaussianField-FinLatticeSites, Pphi2-interactingLatticeMeasure, GaussianField-Configuration, GaussianField-instMeasurableSpaceConfiguration, Pphi2-massGap, GaussianField-FinLatticeField}
  Two-point clustering on the periodic torus: $|\langle\phi(t,x)\phi(0,y)\rangle - \langle\phi(x)\rangle\langle\phi(y)\rangle| \le C \exp(-m_{\mathrm{phys}} \, d_{\mathrm{cyc}}(t)\, a)$. From inserting a complete set of eigenstates in the transfer matrix trace.
\end{theorem}
\begin{theorem}[\texttt{general\_clustering\_lattice}]\label{Pphi2-general-clustering-lattice}
  \lean{Pphi2.general_clustering_lattice}\leanok
  \uses{Pphi2-latticeEuclideanTimeSeparation, Pphi2-general-clustering-from-spectral-gap, GaussianField-FinLatticeSites, Pphi2-latticeConfigEuclideanTimeShift, Pphi2-interactingLatticeMeasure, GaussianField-Configuration, GaussianField-instMeasurableSpaceConfiguration, Pphi2-massGap-pos, Pphi2-massGap, GaussianField-FinLatticeField}
  Exponential clustering for general observables with rate $m = m_{\mathrm{phys}}$. Fully proved (modulo axiom).
\end{theorem}
\section{\texttt{Pphi2.OSProofs.OS3\_RP\_Inheritance}}
\begin{theorem}[\texttt{continuum\_rp\_from\_lattice}]\label{Pphi2-continuum-rp-from-lattice}
  \lean{Pphi2.continuum_rp_from_lattice}\leanok
  \uses{Pphi2-IsRP, Pphi2-rp-closed-under-weak-limit}
  Specialization: RP transfers from lattice to continuum limit. Immediate from \texttt{rp\_closed\_under\_weak\_limit}. Fully proved.
\end{theorem}
\begin{definition}[\texttt{IsRP}]\label{Pphi2-IsRP}
  \lean{Pphi2.IsRP}\leanok
  ``\texttt{lean def IsRP (mu : Measure X) (Theta : X -> X) (S : Set (X -> R)) : Prop := forall F in S, integral x, F x * F (Theta x) d mu >= 0 }`` Abstract reflection positivity: $\int F(x) F(\Theta x)\,d\mu \ge 0$ for all $F \in S$.
\end{definition}
\begin{theorem}[\texttt{rp\_closed\_under\_weak\_limit}]\label{Pphi2-rp-closed-under-weak-limit}
  \lean{Pphi2.rp_closed_under_weak_limit}\leanok
  \uses{Pphi2-IsRP, Lattice-toSemilatticeInf}
  If $\{\mu_n\}$ are RP measures converging weakly to $\mu$, then $\mu$ is RP. Proof: $x \mapsto F(x) F(\Theta x)$ is bounded continuous, so $\int F F \circ \Theta\,d\mu_n \to \int F F \circ \Theta\,d\mu$. Each term $\ge 0$, so the limit $\ge 0$ by \texttt{ge\_of\_tendsto}. Fully proved.
\end{theorem}
\section{\texttt{Pphi2.InteractingMeasure.FreeMomentBound}}
\begin{theorem}[\texttt{torus\_normConst\_field\_moment\_uniform}]\label{Pphi2-torus-normConst-field-moment-uniform}
  \lean{Pphi2.torus_normConst_field_moment_uniform}\leanok
  \uses{GaussianField-FinLatticeSites, Pphi2-torus-normConst-second-moment, GaussianField-latticeGaussianMeasure-isProbability, GaussianField-latticeGaussianMeasure, GaussianField-Configuration, GaussianField-circleSpacing, GaussianField-instMeasurableSpaceConfiguration, GaussianField-FinLatticeField, GaussianField-circleSpacing-pos, Pphi2-field-pow-le-second-pow}
  \texttt{Pphi2.torus\_normConst\_field\_moment\_uniform}
\end{theorem}
\begin{theorem}[\texttt{memLp\_pairing\_pow\_mul\_interaction\_pow}]\label{Pphi2-memLp-pairing-pow-mul-interaction-pow}
  \lean{Pphi2.memLp_pairing_pow_mul_interaction_pow}\leanok
  \uses{GaussianField-FinLatticeSites, Pphi2-interactionFunctional-measurable, Pphi2-integrable-pow-pairing, GaussianField-latticeGaussianMeasure, GaussianField-configuration-eval-measurable, GaussianField-Configuration, GaussianField-instMeasurableSpaceConfiguration, Pphi2-interactionFunctional-pow-integrable, GaussianField-FinLatticeField, Pphi2-interactionFunctional}
  \texttt{Pphi2.memLp\_pairing\_pow\_mul\_interaction\_pow}
\end{theorem}
\begin{theorem}[\texttt{normConst\_second\_moment}]\label{Pphi2-normConst-second-moment}
  \lean{Pphi2.normConst_second_moment}\leanok
  \uses{GaussianField-FinLatticeSites, GaussianField-gffPositionCovariance, Pphi2-gffPositionCovariance-row-sum, GaussianField-latticeGaussianMeasure, GaussianField-Configuration, GaussianField-instMeasurableSpaceConfiguration, Pphi2-latticeSecondMoment-eq-smearedCovariance, GaussianField-gffSmearedCovariance-eq-sum-position, GaussianField-FinLatticeField, GaussianField-gffSmearedCovariance}
  \texttt{Pphi2.normConst\_second\_moment}
\end{theorem}
\begin{theorem}[\texttt{torus\_normConst\_second\_moment}]\label{Pphi2-torus-normConst-second-moment}
  \lean{Pphi2.torus_normConst_second_moment}\leanok
  \uses{Pphi2-torus-volume-eq, GaussianField-FinLatticeSites, GaussianField-latticeGaussianMeasure, GaussianField-Configuration, Pphi2-normConst-second-moment, GaussianField-circleSpacing, GaussianField-instMeasurableSpaceConfiguration, GaussianField-FinLatticeField, GaussianField-circleSpacing-pos}
  \texttt{Pphi2.torus\_normConst\_second\_moment}
\end{theorem}
\begin{theorem}[\texttt{torus\_free\_product\_moment\_uniform}]\label{Pphi2-torus-free-product-moment-uniform}
  \lean{Pphi2.torus_free_product_moment_uniform}\leanok
  \uses{Pphi2-interaction-moment-le, Pphi2-free-product-moment-cauchy-schwarz, Pphi2-torus-normConst-field-moment-uniform, GaussianField-FinLatticeSites, GaussianField-latticeGaussianMeasure, GaussianField-Configuration, GaussianField-circleSpacing, GaussianField-instMeasurableSpaceConfiguration, GaussianField-FinLatticeField, Pphi2-interactionFunctional, GaussianField-circleSpacing-pos}
  \texttt{Pphi2.torus\_free\_product\_moment\_uniform}
\end{theorem}
\section{\texttt{Pphi2.IRLimit.Periodization}}
\begin{theorem}[\texttt{periodizeCLM\_eq\_on\_symmetric\_large\_period}]\label{Pphi2-periodizeCLM-eq-on-symmetric-large-period}
  \lean{Pphi2.periodizeCLM_eq_on_symmetric_large_period}\leanok
  \uses{GaussianField-circleTranslation, GaussianField-SmoothMap-Circle-instFunLike, GaussianField-SmoothMap-Circle-instAddCommGroup, GaussianField-periodizeCLM, GaussianField-SmoothMap-Circle-toFun, GaussianField-SmoothMap-Circle-instModule, GaussianField-SmoothMap-Circle, GaussianField-periodizeCLM-eq-on-large-period, GaussianField-periodizeCLM-comp-schwartzTranslation, GaussianField-schwartzTranslation, GaussianField-SmoothMap-Circle-instTopologicalSpace, Lattice-toSemilatticeInf}
  \texttt{Pphi2.periodizeCLM\_eq\_on\_symmetric\_large\_period}
\end{theorem}
\begin{theorem}[\texttt{periodizeCLM\_tendsto\_uniformlyOn\_symmetricCompact}]\label{Pphi2-periodizeCLM-tendsto-uniformlyOn-symmetricCompact}
  \lean{Pphi2.periodizeCLM_tendsto_uniformlyOn_symmetricCompact}\leanok
  \uses{Pphi2-periodizeCLM-eq-on-symmetric-large-period, GaussianField-SmoothMap-Circle-instAddCommGroup, GaussianField-periodizeCLM, GaussianField-SmoothMap-Circle-toFun, GaussianField-SmoothMap-Circle-instModule, GaussianField-SmoothMap-Circle, GaussianField-SmoothMap-Circle-instTopologicalSpace}
  \texttt{Pphi2.periodizeCLM\_tendsto\_uniformlyOn\_symmetricCompact}
\end{theorem}
\section{\texttt{Pphi2.NelsonEstimate.IntegrabilityHelpers}}
\begin{theorem}[\texttt{lintegral\_layer\_cake\_lt\_top\_of\_eventual\_decay}]\label{Pphi2-IntegrabilityHelpers-lintegral-layer-cake-lt-top-of-eventual-decay}
  \lean{Pphi2.IntegrabilityHelpers.lintegral_layer_cake_lt_top_of_eventual_decay}\leanok
  \uses{Pphi2-IntegrabilityHelpers-lintegral-exp-neg-lt-top, Lattice-toSemilatticeInf}
  \texttt{Pphi2.IntegrabilityHelpers.lintegral\_layer\_cake\_lt\_top\_of\_eventual\_decay}
\end{theorem}
\begin{theorem}[\texttt{lintegral\_exp\_neg\_lt\_top}]\label{Pphi2-IntegrabilityHelpers-lintegral-exp-neg-lt-top}
  \lean{Pphi2.IntegrabilityHelpers.lintegral_exp_neg_lt_top}\leanok
  \texttt{Pphi2.IntegrabilityHelpers.lintegral\_exp\_neg\_lt\_top}
\end{theorem}
\section{\texttt{Pphi2.NelsonEstimate.ChaosMoment}}
\begin{theorem}[\texttt{chaosLE\_moment\_le}]\label{Pphi2-chaosLE-moment-le}
  \lean{Pphi2.chaosLE_moment_le}\leanok
  \uses{GaussianHilbert-stdGaussianFin, GaussianHilbert-wienerChaosLE, Lattice-toSemilatticeInf, GaussianHilbert-bonami-nelson-chaosLE}
  \texttt{Pphi2.chaosLE\_moment\_le}
\end{theorem}
\section{\texttt{Pphi2.TorusContinuumLimit.TorusInteractingMoments}}
\begin{theorem}[\texttt{limit\_le\_of\_uniform\_bound}]\label{Pphi2-limit-le-of-uniform-bound}
  \lean{Pphi2.limit_le_of_uniform_bound}\leanok
  \uses{GaussianField-NuclearTensorProduct-instModuleReal, GaussianField-NuclearTensorProduct-instTopologicalSpace, GaussianField-NuclearTensorProduct-instIsTopologicalAddGroup, GaussianField-NuclearTensorProduct-instAddCommGroup, GaussianField-Configuration, GaussianField-NuclearTensorProduct-instContinuousSMulReal, GaussianField-SmoothMap-Circle, GaussianField-instMeasurableSpaceConfiguration, GaussianField-TorusTestFunction, Lattice-toSemilatticeInf}
  \texttt{Pphi2.limit\_le\_of\_uniform\_bound}
\end{theorem}
\begin{theorem}[\texttt{sub\_min\_le\_sq\_div}]\label{Pphi2-sub-min-le-sq-div}
  \lean{Pphi2.sub_min_le_sq_div}\leanok
  \uses{Lattice-toSemilatticeInf}
  \texttt{Pphi2.sub\_min\_le\_sq\_div}
\end{theorem}
\begin{theorem}[\texttt{torus\_connectedFourPoint\_tendsto}]\label{Pphi2-torus-connectedFourPoint-tendsto}
  \lean{Pphi2.torus_connectedFourPoint_tendsto}\leanok
  \uses{GaussianField-NuclearTensorProduct-instModuleReal, GaussianField-NuclearTensorProduct-instTopologicalSpace, GaussianField-NuclearTensorProduct-instIsTopologicalAddGroup, Pphi2-torus-interacting-fourth-moment-tendsto, GaussianField-NuclearTensorProduct-instAddCommGroup, GaussianField-Configuration, GaussianField-NuclearTensorProduct-instContinuousSMulReal, GaussianField-SmoothMap-Circle, GaussianField-instMeasurableSpaceConfiguration, Pphi2-torusInteractingMeasure, GaussianField-TorusTestFunction, Pphi2-torus-interacting-second-moment-tendsto}
  \texttt{Pphi2.torus\_connectedFourPoint\_tendsto}
\end{theorem}
\begin{theorem}[\texttt{torus\_interacting\_abs\_pow\_integrable}]\label{Pphi2-torus-interacting-abs-pow-integrable}
  \lean{Pphi2.torus_interacting_abs_pow_integrable}\leanok
  \uses{GaussianField-NuclearTensorProduct-instModuleReal, Lattice-toSemilatticeSup, GaussianField-NuclearTensorProduct-instTopologicalSpace, GaussianField-finLatticeField-dyninMityaginSpace, GaussianField-ell2-separableSpace, Pphi2-interactionFunctional-bounded-below, GaussianField-FinLatticeSites, GaussianField-NuclearTensorProduct-instIsTopologicalAddGroup, Pphi2-torusEmbedLift, Pphi2-partitionFunction, GaussianField-latticeCovarianceGJ, Pphi2-interactionFunctional-measurable, GaussianField-pairing-memLp, Pphi2-torusEmbedLift-measurable, Pphi2-interactingLatticeMeasure, GaussianField-NuclearTensorProduct-instAddCommGroup, GaussianField-latticeGaussianMeasure, GaussianField-configuration-eval-measurable, GaussianField-Configuration, GaussianField-NuclearTensorProduct-instContinuousSMulReal, GaussianField-SmoothMap-Circle, GaussianField-circleSpacing, GaussianField-instMeasurableSpaceConfiguration, Pphi2-boltzmannWeight, Pphi2-latticeTestFn, Pphi2-torusEmbedLift-eval-eq, GaussianField-ell2-, Pphi2-torusInteractingMeasure, GaussianField-FinLatticeField, Pphi2-interactionFunctional, GaussianField-TorusTestFunction, Pphi2-partitionFunction-pos, Lattice-toSemilatticeInf, GaussianField-circleSpacing-pos, Pphi2-boltzmannWeight-pos}
  \texttt{Pphi2.torus\_interacting\_abs\_pow\_integrable}
\end{theorem}
\begin{theorem}[\texttt{torus\_interacting\_second\_moment\_tendsto}]\label{Pphi2-torus-interacting-second-moment-tendsto}
  \lean{Pphi2.torus_interacting_second_moment_tendsto}\leanok
  \uses{GaussianField-NuclearTensorProduct-instModuleReal, GaussianField-NuclearTensorProduct-instTopologicalSpace, GaussianField-NuclearTensorProduct-instIsTopologicalAddGroup, GaussianField-NuclearTensorProduct-instAddCommGroup, Pphi2-torus-interacting-fourth-moment-uniform, GaussianField-configuration-eval-measurable, GaussianField-Configuration, Pphi2-moment-tendsto-of-uniform, GaussianField-NuclearTensorProduct-instContinuousSMulReal, GaussianField-SmoothMap-Circle, GaussianField-instMeasurableSpaceConfiguration, Pphi2-torus-interacting-abs-pow-integrable, Pphi2-limit-le-of-uniform-bound, Pphi2-torusInteractingMeasure, Pphi2-torus-interacting-second-moment-uniform, GaussianField-TorusTestFunction, Pphi2-sub-min-le-sq-div, Lattice-toSemilatticeInf}
  \texttt{Pphi2.torus\_interacting\_second\_moment\_tendsto}
\end{theorem}
\begin{theorem}[\texttt{moment\_tendsto\_of\_uniform}]\label{Pphi2-moment-tendsto-of-uniform}
  \lean{Pphi2.moment_tendsto_of_uniform}\leanok
  \uses{GaussianField-NuclearTensorProduct-instModuleReal, Lattice-toSemilatticeSup, GaussianField-NuclearTensorProduct-instTopologicalSpace, GaussianField-NuclearTensorProduct-instIsTopologicalAddGroup, GaussianField-NuclearTensorProduct-instAddCommGroup, GaussianField-Configuration, GaussianField-NuclearTensorProduct-instContinuousSMulReal, GaussianField-SmoothMap-Circle, GaussianField-instMeasurableSpaceConfiguration, GaussianField-TorusTestFunction, Lattice-toSemilatticeInf}
  \texttt{Pphi2.moment\_tendsto\_of\_uniform}
\end{theorem}
\begin{theorem}[\texttt{torus\_interacting\_fourth\_moment\_tendsto}]\label{Pphi2-torus-interacting-fourth-moment-tendsto}
  \lean{Pphi2.torus_interacting_fourth_moment_tendsto}\leanok
  \uses{GaussianField-NuclearTensorProduct-instModuleReal, GaussianField-NuclearTensorProduct-instTopologicalSpace, GaussianField-NuclearTensorProduct-instIsTopologicalAddGroup, GaussianField-NuclearTensorProduct-instAddCommGroup, Pphi2-torus-interacting-fourth-moment-uniform, GaussianField-configuration-eval-measurable, GaussianField-Configuration, Pphi2-moment-tendsto-of-uniform, GaussianField-NuclearTensorProduct-instContinuousSMulReal, GaussianField-SmoothMap-Circle, GaussianField-instMeasurableSpaceConfiguration, Pphi2-torus-interacting-abs-pow-integrable, Pphi2-limit-le-of-uniform-bound, Pphi2-torus-interacting-eighth-moment-uniform, Pphi2-torusInteractingMeasure, GaussianField-TorusTestFunction, Pphi2-sub-min-le-sq-div, Lattice-toSemilatticeInf}
  \texttt{Pphi2.torus\_interacting\_fourth\_moment\_tendsto}
\end{theorem}
\begin{theorem}[\texttt{torus\_interacting\_fourth\_moment\_uniform}]\label{Pphi2-torus-interacting-fourth-moment-uniform}
  \lean{Pphi2.torus_interacting_fourth_moment_uniform}\leanok
  \uses{Pphi2-torusEmbeddedTwoPoint-eq-lattice-second-moment, GaussianField-NuclearTensorProduct-instModuleReal, Pphi2-partitionFunction-ge-one, GaussianField-NuclearTensorProduct-instTopologicalSpace, GaussianField-finLatticeField-dyninMityaginSpace, Pphi2-nelson-exponential-estimate, GaussianField-measure, GaussianField-ell2-separableSpace, Pphi2-torusEmbeddedTwoPoint, GaussianField-FinLatticeSites, GaussianField-NuclearTensorProduct-instIsTopologicalAddGroup, Pphi2-torusEmbedLift, Pphi2-partitionFunction, GaussianField-latticeCovarianceGJ, GaussianField-pairing-memLp, Pphi2-torusEmbedLift-measurable, Pphi2-interactingLatticeMeasure, GaussianField-NuclearTensorProduct-instAddCommGroup, GaussianField-latticeGaussianMeasure, GaussianField-configuration-eval-measurable, GaussianField-Configuration, Pphi2-density-transfer-bound, GaussianField-NuclearTensorProduct-instContinuousSMulReal, GaussianField-gaussian-hypercontractive, GaussianField-SmoothMap-Circle, Pphi2-torusEmbeddedTwoPoint-uniform-bound, GaussianField-circleSpacing, GaussianField-instMeasurableSpaceConfiguration, Pphi2-latticeTestFn, Pphi2-torusEmbedLift-eval-eq, GaussianField-ell2-, Pphi2-torusInteractingMeasure, GaussianField-FinLatticeField, Pphi2-interactionFunctional, GaussianField-TorusTestFunction, Lattice-toSemilatticeInf, GaussianField-circleSpacing-pos}
  \texttt{Pphi2.torus\_interacting\_fourth\_moment\_uniform}
\end{theorem}
\begin{theorem}[\texttt{torus\_interacting\_eighth\_moment\_uniform}]\label{Pphi2-torus-interacting-eighth-moment-uniform}
  \lean{Pphi2.torus_interacting_eighth_moment_uniform}\leanok
  \uses{Pphi2-torusEmbeddedTwoPoint-eq-lattice-second-moment, GaussianField-NuclearTensorProduct-instModuleReal, Pphi2-partitionFunction-ge-one, GaussianField-NuclearTensorProduct-instTopologicalSpace, GaussianField-finLatticeField-dyninMityaginSpace, Pphi2-nelson-exponential-estimate, GaussianField-measure, GaussianField-ell2-separableSpace, Pphi2-torusEmbeddedTwoPoint, GaussianField-FinLatticeSites, GaussianField-NuclearTensorProduct-instIsTopologicalAddGroup, Pphi2-torusEmbedLift, Pphi2-partitionFunction, GaussianField-latticeCovarianceGJ, GaussianField-pairing-memLp, Pphi2-torusEmbedLift-measurable, Pphi2-interactingLatticeMeasure, GaussianField-NuclearTensorProduct-instAddCommGroup, GaussianField-latticeGaussianMeasure, GaussianField-configuration-eval-measurable, GaussianField-Configuration, Pphi2-density-transfer-bound, GaussianField-NuclearTensorProduct-instContinuousSMulReal, GaussianField-gaussian-hypercontractive, GaussianField-SmoothMap-Circle, Pphi2-torusEmbeddedTwoPoint-uniform-bound, GaussianField-circleSpacing, GaussianField-instMeasurableSpaceConfiguration, Pphi2-latticeTestFn, Pphi2-torusEmbedLift-eval-eq, GaussianField-ell2-, Pphi2-torusInteractingMeasure, GaussianField-FinLatticeField, Pphi2-interactionFunctional, GaussianField-TorusTestFunction, Lattice-toSemilatticeInf, GaussianField-circleSpacing-pos}
  \texttt{Pphi2.torus\_interacting\_eighth\_moment\_uniform}
\end{theorem}
\section{\texttt{Pphi2.TorusContinuumLimit.TorusEmbedding}}
\begin{theorem}[\texttt{torusContinuumMeasure\_isProbability}]\label{Pphi2-torusContinuumMeasure-isProbability}
  \lean{Pphi2.torusContinuumMeasure_isProbability}\leanok
  \uses{GaussianField-NuclearTensorProduct-instModuleReal, GaussianField-NuclearTensorProduct-instTopologicalSpace, GaussianField-FinLatticeSites, GaussianField-NuclearTensorProduct-instIsTopologicalAddGroup, Pphi2-torusEmbedLift, Pphi2-torusContinuumMeasure, Pphi2-torusEmbedLift-measurable, GaussianField-NuclearTensorProduct-instAddCommGroup, GaussianField-latticeGaussianMeasure-isProbability, GaussianField-latticeGaussianMeasure, GaussianField-Configuration, GaussianField-NuclearTensorProduct-instContinuousSMulReal, GaussianField-SmoothMap-Circle, GaussianField-circleSpacing, GaussianField-instMeasurableSpaceConfiguration, GaussianField-FinLatticeField, GaussianField-TorusTestFunction, GaussianField-circleSpacing-pos}
  \texttt{Pphi2.torusContinuumMeasure\_isProbability}
\end{theorem}
\begin{theorem}[\texttt{torusEmbedLift\_measurable}]\label{Pphi2-torusEmbedLift-measurable}
  \lean{Pphi2.torusEmbedLift_measurable}\leanok
  \uses{GaussianField-NuclearTensorProduct-instModuleReal, GaussianField-NuclearTensorProduct-instTopologicalSpace, GaussianField-FinLatticeSites, GaussianField-NuclearTensorProduct-instIsTopologicalAddGroup, Pphi2-torusEmbedLift, GaussianField-configuration-measurable-of-eval-measurable, GaussianField-NuclearTensorProduct-instAddCommGroup, GaussianField-configuration-eval-measurable, GaussianField-Configuration, GaussianField-NuclearTensorProduct-instContinuousSMulReal, GaussianField-SmoothMap-Circle, GaussianField-instMeasurableSpaceConfiguration, GaussianField-evalTorusAtSiteGJ, GaussianField-FinLatticeField, GaussianField-TorusTestFunction}
  The torus embedding lift is measurable. Proved by showing each evaluation is a finite sum of measurable functions.
\end{theorem}
\begin{definition}[\texttt{torusEmbeddedTwoPoint}]\label{Pphi2-torusEmbeddedTwoPoint}
  \lean{Pphi2.torusEmbeddedTwoPoint}\leanok
  \uses{GaussianField-NuclearTensorProduct-instModuleReal, GaussianField-NuclearTensorProduct-instTopologicalSpace, Pphi2-torusContinuumMeasure, GaussianField-NuclearTensorProduct-instAddCommGroup, GaussianField-Configuration, GaussianField-SmoothMap-Circle, GaussianField-instMeasurableSpaceConfiguration, GaussianField-TorusTestFunction}
  ``\texttt{lean def torusEmbeddedTwoPoint (N : ℕ) [NeZero N] (mass : ℝ) (hmass : 0 < mass) (f g : TorusTestFunction L) : ℝ }`` Two-point function $\int \omega(f) \omega(g)\, d\nu_{\text{GFF},N}$.
\end{definition}
\begin{definition}[\texttt{torusEmbedLift}]\label{Pphi2-torusEmbedLift}
  \lean{Pphi2.torusEmbedLift}\leanok
  \uses{GaussianField-NuclearTensorProduct-instModuleReal, GaussianField-NuclearTensorProduct-instTopologicalSpace, GaussianField-FinLatticeSites, GaussianField-NuclearTensorProduct-instIsTopologicalAddGroup, GaussianField-NuclearTensorProduct-instAddCommGroup, GaussianField-Configuration, GaussianField-NuclearTensorProduct-instContinuousSMulReal, GaussianField-SmoothMap-Circle, GaussianField-torusEmbedCLMGJ, GaussianField-FinLatticeField, GaussianField-TorusTestFunction}
  \texttt{Pphi2.torusEmbedLift}
\end{definition}
\begin{definition}[\texttt{torusContinuumGreen}]\label{Pphi2-torusContinuumGreen}
  \lean{Pphi2.torusContinuumGreen}\leanok
  \uses{GaussianField-NuclearTensorProduct-instModuleReal, GaussianField-NuclearTensorProduct-instTopologicalSpace, GaussianField-circleHasLaplacianEigenvalues, GaussianField-NuclearTensorProduct-dyninMityaginSpace, GaussianField-SmoothMap-Circle-instAddCommGroup, GaussianField-tensorProductHasLaplacianEigenvalues, GaussianField-NuclearTensorProduct-instAddCommGroup, GaussianField-SmoothMap-Circle-instModule, GaussianField-SmoothMap-Circle-instContinuousSMul, GaussianField-smoothCircle-dyninMityaginSpace, GaussianField-SmoothMap-Circle-instIsTopologicalAddGroup, GaussianField-SmoothMap-Circle, GaussianField-TorusTestFunction, GaussianField-SmoothMap-Circle-instTopologicalSpace, GaussianField-greenFunctionBilinear}
  ``\texttt{lean def torusContinuumGreen (mass : ℝ) (hmass : 0 < mass) (f g : TorusTestFunction L) : ℝ }`` Continuum Green's function on $T^2_L$: spectral sum $G_L(f,g) = \sum_n \hat{f}(n)\hat{g}(n)/((2\pi n/L)^2 + m^2)$.
\end{definition}
\begin{definition}[\texttt{torusContinuumMeasure}]\label{Pphi2-torusContinuumMeasure}
  \lean{Pphi2.torusContinuumMeasure}\leanok
  \uses{GaussianField-NuclearTensorProduct-instModuleReal, GaussianField-NuclearTensorProduct-instTopologicalSpace, GaussianField-FinLatticeSites, GaussianField-NuclearTensorProduct-instIsTopologicalAddGroup, Pphi2-torusEmbedLift, GaussianField-NuclearTensorProduct-instAddCommGroup, GaussianField-latticeGaussianMeasure, GaussianField-Configuration, GaussianField-NuclearTensorProduct-instContinuousSMulReal, GaussianField-SmoothMap-Circle, GaussianField-circleSpacing, GaussianField-instMeasurableSpaceConfiguration, GaussianField-FinLatticeField, GaussianField-TorusTestFunction, GaussianField-circleSpacing-pos}
  ``\texttt{lean def torusContinuumMeasure (N : ℕ) [NeZero N] (mass : ℝ) (hmass : 0 < mass) : Measure (Configuration (TorusTestFunction L)) }`` Pushforward $\nu_{\text{GFF},N} = (\tilde\iota_N)_* \mu_{\text{GFF},N}$.
\end{definition}
\section{\texttt{Pphi2.InteractingMeasure.MomentDerivative}}
\begin{theorem}[\texttt{moment\_hasDerivAt}]\label{Pphi2-moment-hasDerivAt}
  \lean{Pphi2.moment_hasDerivAt}\leanok
  \uses{Pphi2-interactionFunctional-eq-single, Pphi2-interactionFunctional-bounded-below, GaussianField-FinLatticeSites, Pphi2-interactionFunctional-measurable, Pphi2-integrable-powMul-interaction, Pphi2-integrable-pow-pairing, GaussianField-latticeGaussianMeasure, GaussianField-configuration-eval-measurable, GaussianField-Configuration, Pphi2-wickConstant, GaussianField-instMeasurableSpaceConfiguration, GaussianField-FinLatticeField, Pphi2-interactionFunctional, Lattice-toSemilatticeInf, Pphi2-wickPolynomial}
  \texttt{Pphi2.moment\_hasDerivAt}
\end{theorem}
\begin{theorem}[\texttt{moment\_mul\_interaction\_hasDerivAt}]\label{Pphi2-moment-mul-interaction-hasDerivAt}
  \lean{Pphi2.moment_mul_interaction_hasDerivAt}\leanok
  \uses{Pphi2-interactionFunctional-eq-single, Pphi2-interactionFunctional-bounded-below, GaussianField-FinLatticeSites, Pphi2-interactionFunctional-measurable, Pphi2-integrable-powMul-interaction, GaussianField-latticeGaussianMeasure, GaussianField-configuration-eval-measurable, GaussianField-Configuration, Pphi2-wickConstant, GaussianField-instMeasurableSpaceConfiguration, Pphi2-integrable-powMul-interaction-sq, GaussianField-FinLatticeField, Pphi2-interactionFunctional, Lattice-toSemilatticeInf, Pphi2-wickPolynomial}
  \texttt{Pphi2.moment\_mul\_interaction\_hasDerivAt}
\end{theorem}
\begin{theorem}[\texttt{interactionFunctional\_eq\_single}]\label{Pphi2-interactionFunctional-eq-single}
  \lean{Pphi2.interactionFunctional_eq_single}\leanok
  \uses{GaussianField-finLatticeDelta, GaussianField-FinLatticeSites, Pphi2-finLatticeDelta-eq-single, GaussianField-Configuration, Pphi2-wickConstant, GaussianField-FinLatticeField, Pphi2-interactionFunctional, Pphi2-wickPolynomial}
  \texttt{Pphi2.interactionFunctional\_eq\_single}
\end{theorem}
\begin{theorem}[\texttt{integrable\_powMul\_interaction\_sq}]\label{Pphi2-integrable-powMul-interaction-sq}
  \lean{Pphi2.integrable_powMul_interaction_sq}\leanok
  \uses{Pphi2-interactionFunctional-eq-single, GaussianField-FinLatticeSites, GaussianField-latticeGaussianMeasure, GaussianField-Configuration, Pphi2-wickConstant, GaussianField-instMeasurableSpaceConfiguration, Pphi2-integrable-powMul-wickPolynomial-mul, GaussianField-FinLatticeField, Pphi2-interactionFunctional, Pphi2-wickPolynomial}
  \texttt{Pphi2.integrable\_powMul\_interaction\_sq}
\end{theorem}
\begin{theorem}[\texttt{moment\_hasDerivWithinAt}]\label{Pphi2-moment-hasDerivWithinAt}
  \lean{Pphi2.moment_hasDerivWithinAt}\leanok
  \uses{Pphi2-interactionFunctional-eq-single, Pphi2-interactionFunctional-bounded-below, GaussianField-FinLatticeSites, Pphi2-interactionFunctional-measurable, Pphi2-integrable-powMul-interaction, Pphi2-integrable-pow-pairing, GaussianField-latticeGaussianMeasure, GaussianField-configuration-eval-measurable, GaussianField-Configuration, Pphi2-wickConstant, GaussianField-instMeasurableSpaceConfiguration, GaussianField-FinLatticeField, Pphi2-interactionFunctional, Lattice-toSemilatticeInf, Pphi2-wickPolynomial}
  \texttt{Pphi2.moment\_hasDerivWithinAt}
\end{theorem}
\section{\texttt{Pphi2.AsymTorus.AsymTraceBridge}}
\begin{theorem}[\texttt{asymTransferOperatorCLM\_pow\_norm\_le\_of\_perp}]\label{Pphi2-asymTransferOperatorCLM-pow-norm-le-of-perp}
  \lean{Pphi2.asymTransferOperatorCLM_pow_norm_le_of_perp}\leanok
  \uses{Pphi2-asymTransferNormalized, Pphi2-L2SpatialField, Pphi2-SpatialField, Pphi2-asymGroundVector, Pphi2-asymTransferGroundEigenvalue-pos, Pphi2-asymTransferGroundEigenvalue, Pphi2-asymTransferOperatorCLM-eq-smul-normalized, Pphi2-asymGappedTransfer, Pphi2-asymTransferOperatorCLM}
  \texttt{Pphi2.asymTransferOperatorCLM\_pow\_norm\_le\_of\_perp}
\end{theorem}
\begin{theorem}[\texttt{asymTransferOperatorCLM\_eq\_smul\_normalized}]\label{Pphi2-asymTransferOperatorCLM-eq-smul-normalized}
  \lean{Pphi2.asymTransferOperatorCLM_eq_smul_normalized}\leanok
  \uses{Pphi2-asymTransferNormalized, Pphi2-L2SpatialField, Pphi2-SpatialField, Pphi2-asymTransferGroundEigenvalue-pos, Pphi2-asymTransferGroundEigenvalue, Pphi2-asymTransferOperatorCLM}
  \texttt{Pphi2.asymTransferOperatorCLM\_eq\_smul\_normalized}
\end{theorem}
\begin{theorem}[\texttt{asymGroundVector\_norm\_eq\_one}]\label{Pphi2-asymGroundVector-norm-eq-one}
  \lean{Pphi2.asymGroundVector_norm_eq_one}\leanok
  \uses{Pphi2-L2SpatialField, Pphi2-SpatialField, Pphi2-AsymTransferGroundExcitedData-i-, Pphi2-asymGroundVector, Pphi2-asymTransferGroundExcitedData, Pphi2-AsymTransferGroundExcitedData-b, Pphi2-AsymTransferGroundExcitedData--}
  \texttt{Pphi2.asymGroundVector\_norm\_eq\_one}
\end{theorem}
\begin{theorem}[\texttt{asymTransferOperatorCLM\_groundVector}]\label{Pphi2-asymTransferOperatorCLM-groundVector}
  \lean{Pphi2.asymTransferOperatorCLM_groundVector}\leanok
  \uses{Pphi2-L2SpatialField, Pphi2-SpatialField, Pphi2-AsymTransferGroundExcitedData-i-, Pphi2-asymGroundVector, Pphi2-asymTransferGroundExcitedData, Pphi2-asymTransferGroundEigenvalue, Pphi2-AsymTransferGroundExcitedData-h-eigen, Pphi2-asymTransferOperatorCLM}
  \texttt{Pphi2.asymTransferOperatorCLM\_groundVector}
\end{theorem}
\begin{theorem}[\texttt{norm\_sub\_groundProj\_le}]\label{Pphi2-norm-sub-groundProj-le}
  \lean{Pphi2.norm_sub_groundProj_le}\leanok
  \uses{Pphi2-L2SpatialField, Pphi2-SpatialField, Pphi2-asymGroundVector, Pphi2-asymGroundVector-norm-eq-one, Lattice-toSemilatticeInf}
  \texttt{Pphi2.norm\_sub\_groundProj\_le}
\end{theorem}
\begin{theorem}[\texttt{asymTransferOperatorCLM\_pow\_sub\_groundProj\_norm\_le}]\label{Pphi2-asymTransferOperatorCLM-pow-sub-groundProj-norm-le}
  \lean{Pphi2.asymTransferOperatorCLM_pow_sub_groundProj_norm_le}\leanok
  \uses{Pphi2-asymTransferNormalized, Pphi2-L2SpatialField, Pphi2-SpatialField, Pphi2-asymGroundVector, Pphi2-asymGroundVector-norm-eq-one, Pphi2-norm-sub-groundProj-le, Pphi2-asymTransferGroundEigenvalue-pos, Pphi2-asymTransferOperatorCLM-pow-norm-le-of-perp, Pphi2-asymTransferGroundEigenvalue, Pphi2-asymTransferOperatorCLM-pow-groundVector, Pphi2-asymTransferOperatorCLM}
  \texttt{Pphi2.asymTransferOperatorCLM\_pow\_sub\_groundProj\_norm\_le}
\end{theorem}
\begin{theorem}[\texttt{asymTransferOperatorCLM\_pow\_groundVector}]\label{Pphi2-asymTransferOperatorCLM-pow-groundVector}
  \lean{Pphi2.asymTransferOperatorCLM_pow_groundVector}\leanok
  \uses{Pphi2-L2SpatialField, Pphi2-asymTransferOperatorCLM-groundVector, Pphi2-SpatialField, Pphi2-asymGroundVector, Pphi2-asymTransferGroundEigenvalue, Pphi2-asymTransferOperatorCLM}
  \texttt{Pphi2.asymTransferOperatorCLM\_pow\_groundVector}
\end{theorem}
\section{\texttt{Pphi2.IRLimit.IRTightness}}
\begin{theorem}[\texttt{cylinderIRLimit\_exists}]\label{Pphi2-cylinderIRLimit-exists}
  \lean{Pphi2.cylinderIRLimit_exists}\leanok
  \uses{GaussianField-NuclearTensorProduct-instModuleReal, GaussianField-NuclearTensorProduct-instTopologicalSpace, Pphi2-cylinderPullback, Pphi2-cylinderIR-uniform-second-moment, GaussianField-NuclearTensorProduct-dyninMityaginSpace, GaussianField-CylinderTestFunction, GaussianField-NuclearTensorProduct-instIsTopologicalAddGroup, Pphi2-AsymTorusSequenceHasUniformGreenMomentBound, Pphi2-cylinderToTorusEmbed, GaussianField-configuration-measurable-of-eval-measurable, GaussianField-NuclearTensorProduct-instAddCommGroup, GaussianField-prokhorov-configuration, Pphi2-cylinderPullbackMeasure, Pphi2-cylinderIR-uniform-exponential-moment, GaussianField-configuration-eval-measurable, GaussianField-Configuration, GaussianField-NuclearTensorProduct-instContinuousSMulReal, GaussianField-SmoothMap-Circle, GaussianField-instMeasurableSpaceConfiguration, Pphi2-MeasureHasGreenMomentBound, GaussianField-AsymTorusTestFunction, GaussianField-configuration-tight-of-uniform-second-moments, Lattice-toSemilatticeInf}
  ``\texttt{lean theorem cylinderIRLimit\_exists (mass : ℝ) (hmass : 0 < mass) (KG CG : ℝ) (hKG\_pos : 0 < KG) (hCG\_pos : 0 < CG) (Lt : ℕ → ℝ) (hLt : ∀ n, Fact (0 < Lt n)) (hLt\_tend : Tendsto Lt atTop atTop) (μ : ∀ n, Measure (Configuration (AsymTorusTestFunction (Lt n) Ls))) (hμ\_prob : ∀ n, IsProbabilityMeasure (μ n)) (hμ\_green : AsymTorusSequenceHasUniformGreenMomentBound Ls mass hmass KG CG Lt hLt μ) : ∃ (φ : ℕ → ℕ) (ν : Measure (Configuration (CylinderTestFunction Ls))), StrictMono φ ∧ IsProbabilityMeasure ν ∧ (∀ (f : CylinderTestFunction Ls), Tendsto (fun n => ∫ ω, Complex.exp (Complex.I * ↑(ω f)) ∂(cylinderPullbackMeasure (Lt (φ n)) Ls (μ (φ n)))) atTop (nhds (∫ ω, Complex.exp (Complex.I * ↑(ω f)) ∂ν))) }``  Informal statement: Given measures $\mu_n$ on asymmetric tori $T_{L_{t,n}, L_s}$ with $L_{t,n} \to \infty$ and an eventual Green-controlled moment bound, there exist a subsequence $\varphi$ and a probability measure $\nu$ on $\mathrm{Configuration}(\mathrm{CylinderTestFunction}\, L_s)$ such that bounded-continuous observables and characteristic functionals converge: $$\textbackslash int e\^{}\{i\textbackslash omega(f)\}\textbackslash , d\textbackslash nu\_\{L\_\{t,\textbackslash varphi(n)\}\} \textbackslash to \textbackslash int e\^{}\{i\textbackslash omega(f)\}\textbackslash , d\textbackslash nu \textbackslash quad \textbackslash text\{for all \} f.$$  Proof route: 1. Combine \texttt{AsymTorusSequenceHasUniformGreenMomentBound} with \texttt{Lt → ∞} to pass to a tail where both the Green bound and \texttt{Lt ≥ 1} hold. 2. Uniform second moment bound from \texttt{cylinderIR\_uniform\_second\_moment}. 3. Mitoma--Chebyshev tightness criterion (from gaussian-field). 4. \texttt{prokhorov\_configuration} for configuration spaces. 5. Bounded-continuous convergence gives characteristic-functional convergence by cos/sin decomposition.
\end{theorem}
\begin{theorem}[\texttt{congr\_simp}]\label{Pphi2-MeasureHasGreenMomentBound-congr-simp}
  \lean{Pphi2.MeasureHasGreenMomentBound.congr_simp}\leanok
  \uses{GaussianField-NuclearTensorProduct-instModuleReal, GaussianField-NuclearTensorProduct-instTopologicalSpace, GaussianField-NuclearTensorProduct-instIsTopologicalAddGroup, GaussianField-NuclearTensorProduct-instAddCommGroup, GaussianField-Configuration, GaussianField-NuclearTensorProduct-instContinuousSMulReal, GaussianField-SmoothMap-Circle, GaussianField-instMeasurableSpaceConfiguration, Pphi2-MeasureHasGreenMomentBound, GaussianField-AsymTorusTestFunction}
  \texttt{Pphi2.MeasureHasGreenMomentBound.congr\_simp}
\end{theorem}
\begin{theorem}[\texttt{of\_forall}]\label{Pphi2-AsymTorusSequenceHasUniformGreenMomentBound-of-forall}
  \lean{Pphi2.AsymTorusSequenceHasUniformGreenMomentBound.of_forall}\leanok
  \uses{GaussianField-NuclearTensorProduct-instModuleReal, GaussianField-NuclearTensorProduct-instTopologicalSpace, GaussianField-NuclearTensorProduct-instIsTopologicalAddGroup, Pphi2-AsymTorusSequenceHasUniformGreenMomentBound, GaussianField-NuclearTensorProduct-instAddCommGroup, GaussianField-Configuration, GaussianField-NuclearTensorProduct-instContinuousSMulReal, GaussianField-SmoothMap-Circle, GaussianField-instMeasurableSpaceConfiguration, Pphi2-MeasureHasGreenMomentBound, GaussianField-AsymTorusTestFunction}
  \texttt{Pphi2.AsymTorusSequenceHasUniformGreenMomentBound.of\_forall}
\end{theorem}
\begin{definition}[\texttt{AsymTorusSequenceHasUniformGreenMomentBound}]\label{Pphi2-AsymTorusSequenceHasUniformGreenMomentBound}
  \lean{Pphi2.AsymTorusSequenceHasUniformGreenMomentBound}\leanok
  \uses{GaussianField-NuclearTensorProduct-instModuleReal, GaussianField-NuclearTensorProduct-instTopologicalSpace, GaussianField-NuclearTensorProduct-instIsTopologicalAddGroup, GaussianField-NuclearTensorProduct-instAddCommGroup, GaussianField-Configuration, GaussianField-NuclearTensorProduct-instContinuousSMulReal, GaussianField-SmoothMap-Circle, GaussianField-instMeasurableSpaceConfiguration, Pphi2-MeasureHasGreenMomentBound, GaussianField-AsymTorusTestFunction}
  \texttt{Pphi2.AsymTorusSequenceHasUniformGreenMomentBound}
\end{definition}
\begin{theorem}[\texttt{of\_forall\_ge\_one}]\label{Pphi2-AsymTorusSequenceHasUniformGreenMomentBound-of-forall-ge-one}
  \lean{Pphi2.AsymTorusSequenceHasUniformGreenMomentBound.of_forall_ge_one}\leanok
  \uses{GaussianField-NuclearTensorProduct-instModuleReal, GaussianField-NuclearTensorProduct-instTopologicalSpace, GaussianField-NuclearTensorProduct-instIsTopologicalAddGroup, Pphi2-AsymTorusSequenceHasUniformGreenMomentBound, GaussianField-NuclearTensorProduct-instAddCommGroup, GaussianField-Configuration, GaussianField-NuclearTensorProduct-instContinuousSMulReal, GaussianField-SmoothMap-Circle, GaussianField-instMeasurableSpaceConfiguration, Pphi2-MeasureHasGreenMomentBound, GaussianField-AsymTorusTestFunction}
  \texttt{Pphi2.AsymTorusSequenceHasUniformGreenMomentBound.of\_forall\_ge\_one}
\end{theorem}
\section{\texttt{Pphi2.IRLimit.UniformExponentialMoment}}
\begin{theorem}[\texttt{cylinderIR\_uniform\_exponential\_moment}]\label{Pphi2-cylinderIR-uniform-exponential-moment}
  \lean{Pphi2.cylinderIR_uniform_exponential_moment}\leanok
  \uses{GaussianField-NuclearTensorProduct-instModuleReal, GaussianField-NuclearTensorProduct-instTopologicalSpace, Pphi2-cylinderPullback-expMoment-uniform-bound, GaussianField-CylinderTestFunction, GaussianField-NuclearTensorProduct-instIsTopologicalAddGroup, GaussianField-NuclearTensorProduct-instAddCommGroup, Pphi2-cylinderPullbackMeasure, GaussianField-Configuration, GaussianField-NuclearTensorProduct-instContinuousSMulReal, GaussianField-SmoothMap-Circle, GaussianField-instMeasurableSpaceConfiguration, Pphi2-MeasureHasGreenMomentBound, GaussianField-AsymTorusTestFunction}
  ``\texttt{lean theorem cylinderIR\_uniform\_exponential\_moment (mass : ℝ) (hmass : 0 < mass) (KG CG : ℝ) (hKG\_pos : 0 < KG) (hCG\_pos : 0 < CG) : ∃ (K C : ℝ) (q : Seminorm ℝ (CylinderTestFunction Ls)), 0 < K ∧ 0 < C ∧ Continuous q ∧ ∀ (Lt : ℝ) [Fact (0 < Lt)] (\_ : 1 ≤ Lt) (μ : Measure (Configuration (AsymTorusTestFunction Lt Ls))) [IsProbabilityMeasure μ] (\_ : MeasureHasGreenMomentBound Ls mass hmass KG CG μ) (f : CylinderTestFunction Ls), Integrable (fun ω : Configuration (CylinderTestFunction Ls) => Real.exp (|ω f|)) (cylinderPullbackMeasure Lt Ls μ) ∧ ∫ ω : Configuration (CylinderTestFunction Ls), Real.exp (|ω f|) ∂(cylinderPullbackMeasure Lt Ls μ) ≤ K * Real.exp (C * q f \^{} 2) }``  Informal statement: There exist constants $K, C > 0$ and a continuous seminorm $q$ such that for all $L_t \ge 1$ and all torus measures $\mu$ satisfying the Green-controlled moment bound with uniform constants, the cylinder pullback exponential moment is integrable and bounded: $$\textbackslash int \textbackslash exp(|\textbackslash omega(f)|)\textbackslash , d\textbackslash nu\_\{L\_t\} \textbackslash le K \textbackslash cdot \textbackslash exp\textbackslash bigl(C\textbackslash , q(f)\^{}2\textbackslash bigr)$$ uniformly in $L_t$.  Proof: Compose \texttt{MeasureHasGreenMomentBound} with \texttt{cylinderPullback\_expMoment\_uniform\_bound} and the method-of-images estimate \texttt{torusGreen\_uniform\_bound}. The abstract \texttt{AsymSatisfiesTorusOS.os1} clause is deliberately not used because its seminorm is not uniform in \texttt{Lt}.
\end{theorem}
\begin{theorem}[\texttt{cylinderIR\_uniform\_second\_moment}]\label{Pphi2-cylinderIR-uniform-second-moment}
  \lean{Pphi2.cylinderIR_uniform_second_moment}\leanok
  \uses{GaussianField-NuclearTensorProduct-instModuleReal, GaussianField-NuclearTensorProduct-instTopologicalSpace, GaussianField-CylinderTestFunction, GaussianField-NuclearTensorProduct-instIsTopologicalAddGroup, GaussianField-NuclearTensorProduct-instAddCommGroup, Pphi2-cylinderPullbackMeasure, Pphi2-cylinderIR-uniform-exponential-moment, GaussianField-configuration-eval-measurable, GaussianField-Configuration, GaussianField-NuclearTensorProduct-instContinuousSMulReal, GaussianField-SmoothMap-Circle, GaussianField-instMeasurableSpaceConfiguration, Pphi2-MeasureHasGreenMomentBound, GaussianField-AsymTorusTestFunction, Lattice-toSemilatticeInf}
  \texttt{Pphi2.cylinderIR\_uniform\_second\_moment}
\end{theorem}
\section{\texttt{Pphi2.ContinuumLimit.Tightness}}
\begin{theorem}[\texttt{continuum\_second\_moment\_uniform}]\label{Pphi2-continuum-second-moment-uniform}
  \lean{Pphi2.continuum_second_moment_uniform}\leanok
  \uses{Pphi2-EuclideanPlaneBackground-SpaceTime, Pphi2-gaussianContinuumMeasure, Pphi2-gaussian-second-moment-uniform, Pphi2-planeBackground, GaussianField-FinLatticeSites, Pphi2-EuclideanPlaneBackground-dim, Pphi2-continuumMeasure, GaussianField-latticeGaussianMeasure, GaussianField-Configuration, Pphi2-interacting-moment-bound, GaussianField-instMeasurableSpaceConfiguration, Pphi2-latticeEmbedLift, Pphi2-ContinuumTestFunction, GaussianField-FinLatticeField}
  \texttt{Pphi2.continuum\_second\_moment\_uniform}
\end{theorem}
\begin{theorem}[\texttt{congr\_simp}]\label{Pphi2-continuumMeasure-congr-simp}
  \lean{Pphi2.continuumMeasure.congr_simp}\leanok
  \uses{Pphi2-EuclideanPlaneBackground-SpaceTime, Pphi2-planeBackground, Pphi2-EuclideanPlaneBackground-dim, Pphi2-continuumMeasure, GaussianField-Configuration, GaussianField-instMeasurableSpaceConfiguration, Pphi2-ContinuumTestFunction}
  \texttt{Pphi2.continuumMeasure.congr\_simp}
\end{theorem}
\begin{theorem}[\texttt{continuumMeasures\_tight} (axiom)]\label{Pphi2-continuumMeasures-tight}
  \lean{Pphi2.continuumMeasures_tight}\leanok
  \uses{Pphi2-EuclideanPlaneBackground-SpaceTime, GaussianField-schwartz-dyninMityaginSpace, Pphi2-planeBackground, Pphi2-continuumMeasure-isProbability, Pphi2-EuclideanPlaneBackground-dim, Pphi2-continuum-second-moment-uniform, Pphi2-continuumMeasure, GaussianField-Configuration, GaussianField-instMeasurableSpaceConfiguration, Pphi2-ContinuumTestFunction, GaussianField-configuration-tight-of-uniform-second-moments}
  For every $\varepsilon > 0$, there exists a compact set $K \subset \mathcal{S}'(\mathbb{R}^d)$ such that $\nu_a(K) \ge 1 - \varepsilon$ for all $a \in (0, 1]$.  Proof outline: 1. By Mitoma's criterion, tightness of measures on $\mathcal{S}'$ reduces to tightness of 1D projections. 2. Chebyshev's inequality reduces 1D tightness to uniform second moment bounds $\sup_a \int |\Phi_a(f)|^2 \, d\nu_a < \infty$. 3. The uniform bound is provided by the hypercontractive estimate (from \texttt{Hypercontractivity.lean}).  Dependencies: \texttt{Hypercontractivity.lean}, Mitoma (1983).
\end{theorem}
\section{\texttt{Pphi2.AsymTorus.AsymVarianceDischarge}}
\begin{theorem}[\texttt{slicePairing\_measurable}]\label{Pphi2-slicePairing-measurable}
  \lean{Pphi2.slicePairing_measurable}\leanok
  \uses{GaussianField-AsymLatticeSites, Pphi2-SpatialField, Pphi2-asymSliceEquiv, GaussianField-AsymLatticeField, Pphi2-slicePairing}
  \texttt{Pphi2.slicePairing\_measurable}
\end{theorem}
\begin{theorem}[\texttt{config\_pairing\_eq\_slice}]\label{Pphi2-config-pairing-eq-slice}
  \lean{Pphi2.config_pairing_eq_slice}\leanok
  \uses{GaussianField-AsymLatticeSites, Pphi2-SpatialField, Pphi2-asym-pairing-sum-slice, Pphi2-asymSliceEquiv, GaussianField-config-apply-eq-sum-evalMapAsym, GaussianField-Configuration, GaussianField-AsymLatticeField, GaussianField-evalMapAsym, GaussianField-instFintypeAsymLatticeSitesOfNeZeroNat}
  \texttt{Pphi2.config\_pairing\_eq\_slice}
\end{theorem}
\begin{theorem}[\texttt{asym\_pairing\_sum\_slice}]\label{Pphi2-asym-pairing-sum-slice}
  \lean{Pphi2.asym_pairing_sum_slice}\leanok
  \uses{GaussianField-AsymLatticeSites, Pphi2-SpatialField, Pphi2-asymSliceEquiv, Pphi2-asymSliceEquiv-apply, GaussianField-AsymLatticeField, GaussianField-instFintypeAsymLatticeSitesOfNeZeroNat}
  \texttt{Pphi2.asym\_pairing\_sum\_slice}
\end{theorem}
\begin{definition}[\texttt{slicePairing}]\label{Pphi2-slicePairing}
  \lean{Pphi2.slicePairing}\leanok
  \uses{GaussianField-AsymLatticeSites, Pphi2-SpatialField, Pphi2-asymSliceEquiv, GaussianField-AsymLatticeField}
  \texttt{Pphi2.slicePairing}
\end{definition}
\begin{theorem}[\texttt{interacting\_second\_moment\_eq\_pathMeasure}]\label{Pphi2-interacting-second-moment-eq-pathMeasure}
  \lean{Pphi2.interacting_second_moment_eq_pathMeasure}\leanok
  \uses{GaussianField-AsymLatticeSites, Pphi2-asymSliceMeasurableEquiv, GaussianField-measurable-evalMapAsym, Pphi2-SpatialField, Pphi2-asymSliceEquiv, Pphi2-interactingLatticeMeasureAsym-slice-pushforward-eq-pathMeasure, Pphi2-asymTransferSystem, Pphi2-slicePairing-measurable, Pphi2-config-pairing-eq-slice, GaussianField-Configuration, GaussianField-AsymLatticeField, Pphi2-slicePairing, GaussianField-evalMapAsym, Pphi2-interactingLatticeMeasureAsym, GaussianField-instMeasurableSpaceConfiguration, Pphi2-asymSliceMeasurableEquiv-coe}
  \texttt{Pphi2.interacting\_second\_moment\_eq\_pathMeasure}
\end{theorem}
\section{\texttt{Pphi2.AsymTorus.MomentBoundOS1}}
\begin{theorem}[\texttt{cylinderPullback\_expMoment\_uniform\_bound}]\label{Pphi2-cylinderPullback-expMoment-uniform-bound}
  \lean{Pphi2.cylinderPullback_expMoment_uniform_bound}\leanok
  \uses{GaussianField-NuclearTensorProduct-instModuleReal, GaussianField-NuclearTensorProduct-instTopologicalSpace, GaussianField-asymTorusContinuumGreen, GaussianField-CylinderTestFunction, GaussianField-NuclearTensorProduct-instIsTopologicalAddGroup, Pphi2-cylinderToTorusEmbed, Pphi2-cylinderPullback-expMoment-le-green, GaussianField-NuclearTensorProduct-instAddCommGroup, Pphi2-cylinderPullbackMeasure, GaussianField-cylinderToTorusEmbed, GaussianField-Configuration, GaussianField-NuclearTensorProduct-instContinuousSMulReal, GaussianField-SmoothMap-Circle, GaussianField-instMeasurableSpaceConfiguration, Pphi2-MeasureHasGreenMomentBound, GaussianField-torusGreen-uniform-bound, GaussianField-AsymTorusTestFunction}
  \texttt{Pphi2.cylinderPullback\_expMoment\_uniform\_bound}
\end{theorem}
\begin{theorem}[\texttt{cylinderPullback\_expMoment\_eq}]\label{Pphi2-cylinderPullback-expMoment-eq}
  \lean{Pphi2.cylinderPullback_expMoment_eq}\leanok
  \uses{GaussianField-NuclearTensorProduct-instModuleReal, Lattice-toSemilatticeSup, GaussianField-NuclearTensorProduct-instTopologicalSpace, Pphi2-cylinderPullback, GaussianField-CylinderTestFunction, GaussianField-NuclearTensorProduct-instIsTopologicalAddGroup, Pphi2-cylinderToTorusEmbed, GaussianField-configuration-measurable-of-eval-measurable, GaussianField-NuclearTensorProduct-instAddCommGroup, Pphi2-cylinderPullbackMeasure, GaussianField-configuration-eval-measurable, GaussianField-Configuration, GaussianField-NuclearTensorProduct-instContinuousSMulReal, GaussianField-SmoothMap-Circle, GaussianField-instMeasurableSpaceConfiguration, GaussianField-AsymTorusTestFunction}
  \texttt{Pphi2.cylinderPullback\_expMoment\_eq}
\end{theorem}
\begin{theorem}[\texttt{cylinderPullback\_expMoment\_le\_green}]\label{Pphi2-cylinderPullback-expMoment-le-green}
  \lean{Pphi2.cylinderPullback_expMoment_le_green}\leanok
  \uses{GaussianField-NuclearTensorProduct-instModuleReal, GaussianField-NuclearTensorProduct-instTopologicalSpace, GaussianField-asymTorusContinuumGreen, Pphi2-cylinderPullback-expMoment-eq, GaussianField-CylinderTestFunction, GaussianField-NuclearTensorProduct-instIsTopologicalAddGroup, Pphi2-cylinderToTorusEmbed, GaussianField-NuclearTensorProduct-instAddCommGroup, Pphi2-cylinderPullbackMeasure, GaussianField-Configuration, GaussianField-NuclearTensorProduct-instContinuousSMulReal, GaussianField-SmoothMap-Circle, GaussianField-instMeasurableSpaceConfiguration, Pphi2-MeasureHasGreenMomentBound, GaussianField-AsymTorusTestFunction}
  \texttt{Pphi2.cylinderPullback\_expMoment\_le\_green}
\end{theorem}
\begin{definition}[\texttt{MeasureHasGreenMomentBound}]\label{Pphi2-MeasureHasGreenMomentBound}
  \lean{Pphi2.MeasureHasGreenMomentBound}\leanok
  \uses{GaussianField-NuclearTensorProduct-instModuleReal, GaussianField-NuclearTensorProduct-instTopologicalSpace, GaussianField-asymTorusContinuumGreen, GaussianField-NuclearTensorProduct-instIsTopologicalAddGroup, GaussianField-NuclearTensorProduct-instAddCommGroup, GaussianField-Configuration, GaussianField-NuclearTensorProduct-instContinuousSMulReal, GaussianField-SmoothMap-Circle, GaussianField-instMeasurableSpaceConfiguration, GaussianField-AsymTorusTestFunction}
  \texttt{Pphi2.MeasureHasGreenMomentBound}
\end{definition}
\section{\texttt{Pphi2.AsymTorus.AsymL2Operator}}
\begin{theorem}[\texttt{asymTransferWeight\_gaussian\_decay}]\label{Pphi2-asymTransferWeight-gaussian-decay}
  \lean{Pphi2.asymTransferWeight_gaussian_decay}\leanok
  \uses{Pphi2-spatialAction, Pphi2-wickConstantAsym, Pphi2-SpatialField, Pphi2-asymTransferWeight, Lattice-toSemilatticeInf}
  \texttt{Pphi2.asymTransferWeight\_gaussian\_decay}
\end{theorem}
\begin{theorem}[\texttt{asymTransferWeight\_pos}]\label{Pphi2-asymTransferWeight-pos}
  \lean{Pphi2.asymTransferWeight_pos}\leanok
  \uses{Pphi2-spatialAction, Pphi2-wickConstantAsym, Pphi2-SpatialField, Pphi2-asymTransferWeight}
  \texttt{Pphi2.asymTransferWeight\_pos}
\end{theorem}
\begin{theorem}[\texttt{asymTransferWeight\_measurable}]\label{Pphi2-asymTransferWeight-measurable}
  \lean{Pphi2.asymTransferWeight_measurable}\leanok
  \uses{Pphi2-continuous-asymTransferWeight, Pphi2-SpatialField, Pphi2-asymTransferWeight}
  \texttt{Pphi2.asymTransferWeight\_measurable}
\end{theorem}
\begin{definition}[\texttt{asymTransferGaussian}]\label{Pphi2-asymTransferGaussian}
  \lean{Pphi2.asymTransferGaussian}\leanok
  \uses{Pphi2-SpatialField, Pphi2-transferGaussian}
  \texttt{Pphi2.asymTransferGaussian}
\end{definition}
\begin{definition}[\texttt{asymTransferOperatorCLM}]\label{Pphi2-asymTransferOperatorCLM}
  \lean{Pphi2.asymTransferOperatorCLM}\leanok
  \uses{Pphi2-asymTransferWeight-measurable, Pphi2-L2SpatialField, Pphi2-SpatialField, Pphi2-asymTransferWeight, Pphi2-transferGaussian-memLp, Pphi2-asymTransferWeight-bound, Pphi2-transferGaussian}
  \texttt{Pphi2.asymTransferOperatorCLM}
\end{definition}
\begin{theorem}[\texttt{continuous\_asymTransferWeight}]\label{Pphi2-continuous-asymTransferWeight}
  \lean{Pphi2.continuous_asymTransferWeight}\leanok
  \uses{Pphi2-spatialAction, Pphi2-wickConstantAsym, Pphi2-SpatialField, Pphi2-asymTransferWeight, Pphi2-spatialKinetic, Pphi2-wickMonomial, Pphi2-spatialPotential, Pphi2-wickPolynomial}
  \texttt{Pphi2.continuous\_asymTransferWeight}
\end{theorem}
\begin{definition}[\texttt{asymTransferWeight}]\label{Pphi2-asymTransferWeight}
  \lean{Pphi2.asymTransferWeight}\leanok
  \uses{Pphi2-spatialAction, Pphi2-wickConstantAsym, Pphi2-SpatialField}
  \texttt{Pphi2.asymTransferWeight}
\end{definition}
\begin{theorem}[\texttt{asymTransferWeight\_memLp\_two}]\label{Pphi2-asymTransferWeight-memLp-two}
  \lean{Pphi2.asymTransferWeight_memLp_two}\leanok
  \uses{Pphi2-asymTransferWeight-measurable, Pphi2-asymTransferWeight-gaussian-decay, Pphi2-SpatialField, Pphi2-asymTransferWeight, Pphi2-asymTransferWeight-pos, Lattice-toSemilatticeInf}
  \texttt{Pphi2.asymTransferWeight\_memLp\_two}
\end{theorem}
\begin{theorem}[\texttt{asymTransferOperator\_isCompact}]\label{Pphi2-asymTransferOperator-isCompact}
  \lean{Pphi2.asymTransferOperator_isCompact}\leanok
  \uses{Pphi2-asymTransferWeight-measurable, Pphi2-L2SpatialField, Pphi2-asymTransferWeight-memLp-two, Pphi2-hilbert-schmidt-isCompact, Pphi2-transferGaussian-norm-le-one, Pphi2-SpatialField, Pphi2-asymTransferWeight, Pphi2-transferGaussian-memLp, Pphi2-asymTransferWeight-bound, Pphi2-asymTransferOperatorCLM, Pphi2-transferGaussian}
  \texttt{Pphi2.asymTransferOperator\_isCompact}
\end{theorem}
\begin{theorem}[\texttt{asymTransferOperator\_isSelfAdjoint}]\label{Pphi2-asymTransferOperator-isSelfAdjoint}
  \lean{Pphi2.asymTransferOperator_isSelfAdjoint}\leanok
  \uses{Pphi2-asymTransferWeight-measurable, Pphi2-L2SpatialField, Pphi2-SpatialField, Pphi2-asymTransferWeight, Pphi2-transferGaussian-memLp, Pphi2-asymTransferWeight-bound, Pphi2-asymTransferOperatorCLM, Pphi2-transferGaussian, Pphi2-transferGaussian-neg}
  \texttt{Pphi2.asymTransferOperator\_isSelfAdjoint}
\end{theorem}
\begin{theorem}[\texttt{asymTransferWeight\_bound}]\label{Pphi2-asymTransferWeight-bound}
  \lean{Pphi2.asymTransferWeight_bound}\leanok
  \uses{Pphi2-spatialAction, Pphi2-wickConstantAsym, Pphi2-SpatialField, Pphi2-asymTransferWeight, Lattice-toSemilatticeInf}
  \texttt{Pphi2.asymTransferWeight\_bound}
\end{theorem}
\section{\texttt{Pphi2.AsymTorus.AsymCutoffBound}}
\begin{theorem}[\texttt{asymTorusInteractingMeasureIso\_exponentialMomentBound\_cutoff\_of\_nelson}]\label{Pphi2-asymTorusInteractingMeasureIso-exponentialMomentBound-cutoff-of-nelson}
  \lean{Pphi2.asymTorusInteractingMeasureIso_exponentialMomentBound_cutoff_of_nelson}\leanok
  \uses{GaussianField-NuclearTensorProduct-instModuleReal, Lattice-toSemilatticeSup, GaussianField-NuclearTensorProduct-instTopologicalSpace, Pphi2-partitionFunctionAsym-ge-one, GaussianField-AsymLatticeSites, GaussianField-measure, Pphi2-asymTorusEmbedLiftIso-measurable, GaussianField-ell2-separableSpace, Pphi2-partitionFunctionAsym-pos, GaussianField-asymLatticeField-dyninMityaginSpace, Pphi2-asymTorusEmbedLiftIso-eval-eq, GaussianField-NuclearTensorProduct-instIsTopologicalAddGroup, Pphi2-asymLatticeTestFnIso, Pphi2-boltzmannWeightAsym, Pphi2-boltzmannWeightAsym-pos, Pphi2-asymTorusEmbedLiftIso, Pphi2-interactionFunctionalAsym-bounded-below, GaussianField-NuclearTensorProduct-instAddCommGroup, Pphi2-partitionFunctionAsym, GaussianField-gaussian-exp-abs-moment, Pphi2-asymTorusInteractingMeasureIso, GaussianField-configuration-eval-measurable, GaussianField-Configuration, Pphi2-density-transfer-bound-iso, GaussianField-NuclearTensorProduct-instContinuousSMulReal, GaussianField-AsymLatticeField, GaussianField-SmoothMap-Circle, Pphi2-latticeGaussianMeasureAsym, Pphi2-interactingLatticeMeasureAsym, Pphi2-interactionFunctionalAsym, GaussianField-instMeasurableSpaceConfiguration, GaussianField-latticeCovarianceAsymGJ, GaussianField-AsymTorusTestFunction, GaussianField-ell2-, Pphi2-interactionFunctionalAsym-measurable}
  \texttt{Pphi2.asymTorusInteractingMeasureIso\_exponentialMomentBound\_cutoff\_of\_nelson}
\end{theorem}
\begin{theorem}[\texttt{asymTorusInteractingMeasureIso\_exponentialMomentBound\_cutoff}]\label{Pphi2-asymTorusInteractingMeasureIso-exponentialMomentBound-cutoff}
  \lean{Pphi2.asymTorusInteractingMeasureIso_exponentialMomentBound_cutoff}\leanok
  \uses{GaussianField-NuclearTensorProduct-instModuleReal, GaussianField-NuclearTensorProduct-instTopologicalSpace, GaussianField-AsymLatticeSites, GaussianField-NuclearTensorProduct-instIsTopologicalAddGroup, Pphi2-asymLatticeTestFnIso, Pphi2-asymTorusInteractingMeasureIso-exponentialMomentBound-cutoff-of-nelson, Pphi2-asymNelson-exponential-estimate-iso, GaussianField-NuclearTensorProduct-instAddCommGroup, Pphi2-asymTorusInteractingMeasureIso, GaussianField-Configuration, GaussianField-NuclearTensorProduct-instContinuousSMulReal, GaussianField-AsymLatticeField, GaussianField-SmoothMap-Circle, Pphi2-latticeGaussianMeasureAsym, Pphi2-interactionFunctionalAsym, GaussianField-instMeasurableSpaceConfiguration, GaussianField-AsymTorusTestFunction}
  \texttt{Pphi2.asymTorusInteractingMeasureIso\_exponentialMomentBound\_cutoff}
\end{theorem}
\section{\texttt{Pphi2.InteractingMeasure.InteractingMomentBound}}
\begin{theorem}[\texttt{abs\_interacting\_moment\_le}]\label{Pphi2-abs-interacting-moment-le}
  \lean{Pphi2.abs_interacting_moment_le}\leanok
  \uses{Pphi2-interacting-moment-le-L2-of-expBound, Pphi2-partitionFn-ge-one, GaussianField-FinLatticeSites, GaussianField-latticeGaussianMeasure, GaussianField-Configuration, GaussianField-instMeasurableSpaceConfiguration, GaussianField-FinLatticeField, Pphi2-interactionFunctional}
  \texttt{Pphi2.abs\_interacting\_moment\_le}
\end{theorem}
\section{\texttt{Pphi2.NelsonEstimate.AsymSmoothLowerBound}}
\begin{theorem}[\texttt{asymCanonicalSmoothInteraction\_lower\_bound\_log\_uniform}]\label{Pphi2-asymCanonicalSmoothInteraction-lower-bound-log-uniform}
  \lean{Pphi2.asymCanonicalSmoothInteraction_lower_bound_log_uniform}\leanok
  \uses{Pphi2-wickPolynomial-lower-bound-general, GaussianField-AsymLatticeSites, Pphi2-asymSmoothWickConstant-le-log-uniform, Pphi2-AsymCanonicalJoint, Pphi2-asymSmoothWickConstant, GaussianField-instFintypeAsymLatticeSitesOfNeZeroNat, Pphi2-asymCanonicalSmoothInteraction, Pphi2-DynamicalCutoff-degreeCutoffPower, Pphi2-asymCanonicalSmoothFieldFunction, Lattice-toSemilatticeInf, Pphi2-wickPolynomial}
  \texttt{Pphi2.asymCanonicalSmoothInteraction\_lower\_bound\_log\_uniform}
\end{theorem}
\begin{theorem}[\texttt{congr\_simp}]\label{Pphi2-asymCanonicalSmoothInteraction-congr-simp}
  \lean{Pphi2.asymCanonicalSmoothInteraction.congr_simp}\leanok
  \uses{Pphi2-AsymCanonicalJoint, Pphi2-asymCanonicalSmoothInteraction}
  \texttt{Pphi2.asymCanonicalSmoothInteraction.congr\_simp}
\end{theorem}
\begin{theorem}[\texttt{asymCanonicalSmoothInteraction\_lower\_bound\_at\_cutoff\_uniform}]\label{Pphi2-asymCanonicalSmoothInteraction-lower-bound-at-cutoff-uniform}
  \lean{Pphi2.asymCanonicalSmoothInteraction_lower_bound_at_cutoff_uniform}\leanok
  \uses{Pphi2-AsymCanonicalJoint, Pphi2-DynamicalCutoff-degreeDynamicalCutoffScale-pos, Pphi2-DynamicalCutoff-degreeDynamicalCutoffScale, Pphi2-asymCanonicalSmoothInteraction-lower-bound-log-uniform, Pphi2-asymCanonicalSmoothInteraction, Pphi2-DynamicalCutoff-degreeCutoffPower, Pphi2-DynamicalCutoff-degreeDynamicalCutoffScale-log-pow-le}
  \texttt{Pphi2.asymCanonicalSmoothInteraction\_lower\_bound\_at\_cutoff\_uniform}
\end{theorem}
\section{\texttt{Pphi2.GaussianContinuumLimit.GaussianTightness}}
\begin{theorem}[\texttt{gaussian\_second\_moment\_uniform}]\label{Pphi2-gaussian-second-moment-uniform}
  \lean{Pphi2.gaussian_second_moment_uniform}\leanok
  \uses{Pphi2-EuclideanPlaneBackground-SpaceTime, Pphi2-gaussianContinuumMeasure, Pphi2-planeBackground, Pphi2-EuclideanPlaneBackground-dim, Pphi2-embeddedTwoPoint, GaussianField-Configuration, GaussianField-instMeasurableSpaceConfiguration, Pphi2-ContinuumTestFunction, Pphi2-embeddedTwoPoint-uniform-bound}
  $\int (\omega f)^2\, d\nu_{\text{GFF},a} \le C(f,m)$ uniformly in $a \in (0,1]$ and $N$. Proved from \texttt{embeddedTwoPoint\_uniform\_bound}.
\end{theorem}
\begin{theorem}[\texttt{gaussianContinuumMeasures\_tight} (axiom)]\label{Pphi2-gaussianContinuumMeasures-tight}
  \lean{Pphi2.gaussianContinuumMeasures_tight}\leanok
  \uses{Pphi2-EuclideanPlaneBackground-SpaceTime, Pphi2-gaussianContinuumMeasure, Pphi2-gaussianContinuumMeasure-isProbability, GaussianField-schwartz-dyninMityaginSpace, Pphi2-gaussian-second-moment-uniform, Pphi2-planeBackground, Pphi2-EuclideanPlaneBackground-dim, GaussianField-Configuration, GaussianField-instMeasurableSpaceConfiguration, Pphi2-ContinuumTestFunction, GaussianField-configuration-tight-of-uniform-second-moments}
  The family $\{\nu_{\text{GFF},a}\}_{a \in (0,1]}$ is tight on $\mathcal{S}'(\mathbb{R}^d)$. Proof outline: uniform second moments $\to$ Chebyshev $\to$ 1D tightness $\to$ Mitoma criterion.
\end{theorem}
\section{\texttt{Pphi2.NelsonEstimate.NelsonEstimate}}
\begin{theorem}[\texttt{nelson\_exponential\_estimate\_lattice}]\label{Pphi2-nelson-exponential-estimate-lattice}
  \lean{Pphi2.nelson_exponential_estimate_lattice}\leanok
  \uses{GaussianField-FinLatticeSites, GaussianField-circleSpacing-eq, GaussianField-latticeGaussianMeasure, GaussianField-Configuration, GaussianField-circleSpacing, GaussianField-instMeasurableSpaceConfiguration, GaussianField-FinLatticeField, Pphi2-interactionFunctional, Pphi2-nelson-exponential-estimate-master, GaussianField-circleSpacing-pos}
  For the $P(\varphi)_2$ lattice theory on the 2D torus of size $L$ with spacing $a = L/N$ and mass $m > 0$: $$\textbackslash int \textbackslash left(e\^{}\{-V\_a(\textbackslash omega)\}\textbackslash right)\^{}2\textbackslash ,d\textbackslash mu\_\{\textbackslash text\{GFF\}\} \textbackslash le K$$ where $K = e^{2L^2 A}$ depends on $P$, $L$, $m$ but NOT on $N$. The proof exploits $a^2 \cdot N^2 = L^2$, so the pointwise bound $V(\omega) \ge -L^2 A$ is uniform in $N$.  Depends on: \texttt{wickPolynomial\_uniform\_bounded\_below}, \texttt{wickConstant\_le\_inv\_mass\_sq}, \texttt{wickConstant\_pos}, \texttt{circleSpacing\_pos}.
\end{theorem}
\section{\texttt{Pphi2.TransferMatrix.GaussianFourier}}
\begin{theorem}[\texttt{fourier\_gaussian\_pos}]\label{Pphi2-fourier-gaussian-pos}
  \lean{Pphi2.fourier_gaussian_pos}\leanok
  For $b > 0$: $\operatorname{Re}\bigl(\widehat{e^{-b\|\cdot\|^2}}(w)\bigr) > 0$ for all $w$. Proved via \texttt{fourier\_gaussian\_innerProductSpace}. Fully proved.
\end{theorem}
\begin{theorem}[\texttt{gaussian\_conv\_strictlyPD}]\label{Pphi2-gaussian-conv-strictlyPD}
  \lean{Pphi2.gaussian_conv_strictlyPD}\leanok
  \uses{Pphi2-L2SpatialField, Pphi2-SpatialField, Pphi2-continuous-transferGaussian, Pphi2-inner-convCLM-pos-of-fourier-pos, Pphi2-transferGaussian-memLp, Pphi2-fourier-gaussian-pos, Pphi2-transferGaussian, Lattice-toSemilatticeInf}
  Gaussian convolution is strictly positive definite on $L^2$: $\langle f, \mathrm{Conv}_G f\rangle > 0$ for $f \ne 0$. Applies \texttt{inner\_convCLM\_pos\_of\_fourier\_pos} with $G(\psi) = e^{-\frac{1}{2}\|\psi\|^2}$. Fully proved.
\end{theorem}
\begin{theorem}[\texttt{inner\_convCLM\_pos\_of\_fourier\_pos}]\label{Pphi2-inner-convCLM-pos-of-fourier-pos}
  \lean{Pphi2.inner_convCLM_pos_of_fourier_pos}\leanok
  \uses{Pphi2-L2SpatialField, Pphi2-SpatialField, Lattice-toSemilatticeInf}
  If $\operatorname{Re}(\hat{g}(w)) > 0$ for all $w$, then $\langle f, \mathrm{Conv}_g f\rangle > 0$ for $f \ne 0$. Proved from the Fourier representation and Plancherel injectivity.
\end{theorem}
\section{\texttt{Pphi2.TorusContinuumLimit.TorusU4Pullback}}
\begin{theorem}[\texttt{torusConnectedFourPoint\_eq\_lattice}]\label{Pphi2-torusConnectedFourPoint-eq-lattice}
  \lean{Pphi2.torusConnectedFourPoint_eq_lattice}\leanok
  \uses{GaussianField-NuclearTensorProduct-instModuleReal, GaussianField-NuclearTensorProduct-instTopologicalSpace, GaussianField-FinLatticeSites, GaussianField-NuclearTensorProduct-instIsTopologicalAddGroup, Pphi2-connectedFourPoint, Pphi2-torusEmbedLift, Pphi2-torusConnectedFourPoint, Pphi2-torusEmbedLift-measurable, Pphi2-interactingLatticeMeasure, GaussianField-NuclearTensorProduct-instAddCommGroup, GaussianField-configuration-eval-measurable, GaussianField-Configuration, GaussianField-NuclearTensorProduct-instContinuousSMulReal, GaussianField-SmoothMap-Circle, GaussianField-circleSpacing, GaussianField-instMeasurableSpaceConfiguration, Pphi2-latticeTestFn, Pphi2-torusEmbedLift-eval-eq, Pphi2-torusInteractingMeasure, GaussianField-FinLatticeField, GaussianField-TorusTestFunction, GaussianField-circleSpacing-pos}
  \texttt{Pphi2.torusConnectedFourPoint\_eq\_lattice}
\end{theorem}
\begin{theorem}[\texttt{torusConnectedFourPoint\_eq\_u4\_one}]\label{Pphi2-torusConnectedFourPoint-eq-u4-one}
  \lean{Pphi2.torusConnectedFourPoint_eq_u4_one}\leanok
  \uses{Pphi2-connectedFourPoint-interactingLatticeMeasure-eq-u4-one, GaussianField-FinLatticeSites, Pphi2-connectedFourPoint, Pphi2-torusConnectedFourPoint, Pphi2-interactingLatticeMeasure, GaussianField-circleSpacing, Pphi2-torusConnectedFourPoint-eq-lattice, Pphi2-latticeTestFn, Pphi2-u4, Pphi2-torusInteractingMeasure, GaussianField-FinLatticeField, GaussianField-TorusTestFunction, GaussianField-circleSpacing-pos}
  \texttt{Pphi2.torusConnectedFourPoint\_eq\_u4\_one}
\end{theorem}
\section{\texttt{Pphi2.NelsonEstimate.BridgeFromTail}}
\begin{theorem}[\texttt{expSqNeg\_lintegral\_le\_of\_dynamical\_cutoff}]\label{Pphi2-BridgeFromTail-expSqNeg-lintegral-le-of-dynamical-cutoff}
  \lean{Pphi2.BridgeFromTail.expSqNeg_lintegral_le_of_dynamical_cutoff}\leanok
  \uses{Pphi2-DynamicalCutoff-measure-le-neg-of-smooth-rough-split, Pphi2-LayerCake-lintegral-expSq-neg-le-of-tail}
  \texttt{Pphi2.BridgeFromTail.expSqNeg\_lintegral\_le\_of\_dynamical\_cutoff}
\end{theorem}
\section{\texttt{Pphi2.InteractingMeasure.CovariancePointwiseBound}}
\begin{theorem}[\texttt{gffPositionCovariance\_abs\_le\_mass\_inv}]\label{Pphi2-gffPositionCovariance-abs-le-mass-inv}
  \lean{Pphi2.gffPositionCovariance_abs_le_mass_inv}\leanok
  \uses{GaussianField-finLatticeField-dyninMityaginSpace, GaussianField-measure, GaussianField-ell2-separableSpace, GaussianField-second-moment-eq-covariance, GaussianField-covariance, GaussianField-FinLatticeSites, GaussianField-latticeCovarianceGJ, GaussianField-gffPositionCovariance, GaussianField-latticeGaussianMeasure, GaussianField-Configuration, GaussianField-gffPositionCovariance-eq-covarianceGJ, GaussianField-instMeasurableSpaceConfiguration, GaussianField-ell2-, Pphi2-lattice-second-moment-le-mass-inv, GaussianField-FinLatticeField}
  \texttt{Pphi2.gffPositionCovariance\_abs\_le\_mass\_inv}
\end{theorem}
\section{\texttt{Pphi2.InteractingMeasure.InteractionL2}}
\begin{theorem}[\texttt{integral\_interaction\_sq\_eq\_canonicalJoint}]\label{Pphi2-integral-interaction-sq-eq-canonicalJoint}
  \lean{Pphi2.integral_interaction_sq_eq_canonicalJoint}\leanok
  \uses{Pphi2-canonicalJointMeasure, Pphi2-canonicalSumConfig, Pphi2-canonicalFullInteractionJoint-eq-interactionFunctional, GaussianField-FinLatticeSites, Pphi2-canonicalFullInteractionJoint, Pphi2-interactionFunctional-measurable, GaussianField-latticeGaussianMeasure, GaussianField-Configuration, GaussianField-instMeasurableSpaceConfiguration, Pphi2-integral-comp-canonicalSumConfig, GaussianField-FinLatticeField, Pphi2-interactionFunctional, Pphi2-CanonicalJoint}
  \texttt{Pphi2.integral\_interaction\_sq\_eq\_canonicalJoint}
\end{theorem}
\section{\texttt{Pphi2.TransferMatrix.SpectralGap}}
\begin{theorem}[\texttt{spectral\_gap\_pos}]\label{Pphi2-spectral-gap-pos}
  \lean{Pphi2.spectral_gap_pos}\leanok
  \uses{Pphi2-massGap-pos, Pphi2-massGap}
  $E_1 - E_0 > 0$. Delegates to \texttt{massGap\_pos}. Fully proved.
\end{theorem}
\section{\texttt{Pphi2.NelsonEstimate.GenericBridge}}
\begin{theorem}[\texttt{bridgeAxiom\_of\_setup\_real\_generic}]\label{Pphi2-BridgeFromTail-bridgeAxiom-of-setup-real-generic}
  \lean{Pphi2.BridgeFromTail.bridgeAxiom_of_setup_real_generic}\leanok
  \uses{Pphi2-BridgeFromTail-expSqNeg-lintegral-le-of-dynamical-cutoff, Lattice-toSemilatticeInf}
  \texttt{Pphi2.BridgeFromTail.bridgeAxiom\_of\_setup\_real\_generic}
\end{theorem}
\section{\texttt{Pphi2.InteractingMeasure.SmearedWickVertex}}
\begin{theorem}[\texttt{wickFourth\_smeared\_site\_inner}]\label{Pphi2-wickFourth-smeared-site-inner}
  \lean{Pphi2.wickFourth_smeared_site_inner}\leanok
  \uses{GaussianField-gffSmearedCovariance-self, GaussianField-FinLatticeSites, GaussianField-gffSmearedCoeff, GaussianField-gffPositionCovariance, GaussianField-latticeGaussianMeasure, GaussianField-Configuration, GaussianField-gff-wickPower-two-smeared-inner, GaussianField-instMeasurableSpaceConfiguration, GaussianField-FinLatticeField, GaussianField-gffSmearedCovariance, GaussianField-gffSmearedCovariance-single-right}
  \texttt{Pphi2.wickFourth\_smeared\_site\_inner}
\end{theorem}
\begin{theorem}[\texttt{wickFourth\_smeared\_vertex\_inner}]\label{Pphi2-wickFourth-smeared-vertex-inner}
  \lean{Pphi2.wickFourth_smeared_vertex_inner}\leanok
  \uses{GaussianField-gffSmearedCovariance-self, GaussianField-FinLatticeSites, GaussianField-gffSmearedCoeff, GaussianField-gffPositionCovariance, GaussianField-latticeGaussianMeasure, GaussianField-Configuration, GaussianField-gff-wickPower-two-smeared-inner, GaussianField-instMeasurableSpaceConfiguration, GaussianField-FinLatticeField, GaussianField-gffSmearedCovariance, GaussianField-gffSmearedCovariance-single-right}
  \texttt{Pphi2.wickFourth\_smeared\_vertex\_inner}
\end{theorem}
\chapter{gaussian-hilbert (upstream)}
\section{\texttt{GaussianHilbert.PolynomialDensity}}
\begin{theorem}[\texttt{zero\_mem\_interior\_integrableExpSet\_sum\_sq\_of\_subGaussian}]\label{GaussianHilbert-zero-mem-interior-integrableExpSet-sum-sq-of-subGaussian}
  \lean{GaussianHilbert.zero_mem_interior_integrableExpSet_sum_sq_of_subGaussian}\leanok
  \uses{GaussianHilbert-IsSubGaussianMeasure, Lattice-toSemilatticeInf}
  \texttt{GaussianHilbert.zero\_mem\_interior\_integrableExpSet\_sum\_sq\_of\_subGaussian}
\end{theorem}
\begin{theorem}[\texttt{polynomial\_dense\_L2\_of\_subGaussian}]\label{GaussianHilbert-polynomial-dense-L2-of-subGaussian}
  \lean{GaussianHilbert.polynomial_dense_L2_of_subGaussian}\leanok
  \uses{GaussianHilbert-withDensity-pos-eq-neg-of-mem-orthogonal-polynomialToL2, GaussianHilbert-polynomialToL2, GaussianHilbert-IsSubGaussianMeasure, GaussianHilbert-MvPolynomial-eval-memLp-two-of-subGaussian}
  \texttt{GaussianHilbert.polynomial\_dense\_L2\_of\_subGaussian}
\end{theorem}
\begin{definition}[\texttt{IsSubGaussianMeasure}]\label{GaussianHilbert-IsSubGaussianMeasure}
  \lean{GaussianHilbert.IsSubGaussianMeasure}\leanok
  \texttt{GaussianHilbert.IsSubGaussianMeasure}
\end{definition}
\begin{theorem}[\texttt{MvPolynomial\_eval\_memLp\_two\_of\_subGaussian}]\label{GaussianHilbert-MvPolynomial-eval-memLp-two-of-subGaussian}
  \lean{GaussianHilbert.MvPolynomial_eval_memLp_two_of_subGaussian}\leanok
  \uses{GaussianHilbert-memLp-monomial-eval-two-of-subGaussian, GaussianHilbert-IsSubGaussianMeasure}
  \texttt{GaussianHilbert.MvPolynomial\_eval\_memLp\_two\_of\_subGaussian}
\end{theorem}
\begin{theorem}[\texttt{withDensity\_pos\_eq\_neg\_of\_mem\_orthogonal\_polynomialToL2}]\label{GaussianHilbert-withDensity-pos-eq-neg-of-mem-orthogonal-polynomialToL2}
  \lean{GaussianHilbert.withDensity_pos_eq_neg_of_mem_orthogonal_polynomialToL2}\leanok
  \uses{Lattice-toSemilatticeSup, GaussianHilbert-polynomialToL2, GaussianHilbert-linearFunctionalPolynomial, GaussianHilbert-IsSubGaussianMeasure, GaussianHilbert-complexMGF-eq-of-moments, GaussianHilbert-memLp-exp-linear-two-of-subGaussian, GaussianHilbert-linearFunctionalPolynomial-eval, GaussianHilbert-MvPolynomial-eval-memLp-two-of-subGaussian}
  \texttt{GaussianHilbert.withDensity\_pos\_eq\_neg\_of\_mem\_orthogonal\_polynomialToL2}
\end{theorem}
\begin{theorem}[\texttt{integrable\_pow\_coord\_of\_subGaussian}]\label{GaussianHilbert-integrable-pow-coord-of-subGaussian}
  \lean{GaussianHilbert.integrable_pow_coord_of_subGaussian}\leanok
  \uses{GaussianHilbert-zero-mem-interior-integrableExpSet-coord-of-subGaussian, GaussianHilbert-IsSubGaussianMeasure}
  \texttt{GaussianHilbert.integrable\_pow\_coord\_of\_subGaussian}
\end{theorem}
\begin{theorem}[\texttt{exists\_polynomial\_near\_continuous\_on\_isCompact}]\label{GaussianHilbert-exists-polynomial-near-continuous-on-isCompact}
  \lean{GaussianHilbert.exists_polynomial_near_continuous_on_isCompact}\leanok
  \texttt{GaussianHilbert.exists\_polynomial\_near\_continuous\_on\_isCompact}
\end{theorem}
\begin{theorem}[\texttt{complexMGF\_eq\_of\_moments}]\label{GaussianHilbert-complexMGF-eq-of-moments}
  \lean{GaussianHilbert.complexMGF_eq_of_moments}\leanok
  \texttt{GaussianHilbert.complexMGF\_eq\_of\_moments}
\end{theorem}
\begin{theorem}[\texttt{integrable\_exp\_linear\_of\_subGaussian}]\label{GaussianHilbert-integrable-exp-linear-of-subGaussian}
  \lean{GaussianHilbert.integrable_exp_linear_of_subGaussian}\leanok
  \uses{GaussianHilbert-IsSubGaussianMeasure, GaussianHilbert-linearFunctional-mul-le-subGaussian-weight, Lattice-toSemilatticeInf}
  \texttt{GaussianHilbert.integrable\_exp\_linear\_of\_subGaussian}
\end{theorem}
\begin{theorem}[\texttt{linearFunctionalPolynomial\_eval}]\label{GaussianHilbert-linearFunctionalPolynomial-eval}
  \lean{GaussianHilbert.linearFunctionalPolynomial_eval}\leanok
  \uses{GaussianHilbert-linearFunctionalPolynomial}
  \texttt{GaussianHilbert.linearFunctionalPolynomial\_eval}
\end{theorem}
\begin{theorem}[\texttt{memLp\_exp\_linear\_two\_of\_subGaussian}]\label{GaussianHilbert-memLp-exp-linear-two-of-subGaussian}
  \lean{GaussianHilbert.memLp_exp_linear_two_of_subGaussian}\leanok
  \uses{GaussianHilbert-integrable-exp-linear-of-subGaussian, GaussianHilbert-IsSubGaussianMeasure}
  \texttt{GaussianHilbert.memLp\_exp\_linear\_two\_of\_subGaussian}
\end{theorem}
\begin{theorem}[\texttt{memLp\_monomial\_eval\_two\_of\_subGaussian}]\label{GaussianHilbert-memLp-monomial-eval-two-of-subGaussian}
  \lean{GaussianHilbert.memLp_monomial_eval_two_of_subGaussian}\leanok
  \uses{GaussianHilbert-IsSubGaussianMeasure, GaussianHilbert-integrable-pow-sum-sq-of-subGaussian, Lattice-toSemilatticeInf}
  \texttt{GaussianHilbert.memLp\_monomial\_eval\_two\_of\_subGaussian}
\end{theorem}
\begin{theorem}[\texttt{isSubGaussianMeasure\_pi\_gaussianReal}]\label{GaussianHilbert-isSubGaussianMeasure-pi-gaussianReal}
  \lean{GaussianHilbert.isSubGaussianMeasure_pi_gaussianReal}\leanok
  \uses{GaussianHilbert-IsSubGaussianMeasure}
  \texttt{GaussianHilbert.isSubGaussianMeasure\_pi\_gaussianReal}
\end{theorem}
\begin{definition}[\texttt{linearFunctionalPolynomial}]\label{GaussianHilbert-linearFunctionalPolynomial}
  \lean{GaussianHilbert.linearFunctionalPolynomial}\leanok
  \texttt{GaussianHilbert.linearFunctionalPolynomial}
\end{definition}
\begin{theorem}[\texttt{linearFunctional\_mul\_le\_subGaussian\_weight}]\label{GaussianHilbert-linearFunctional-mul-le-subGaussian-weight}
  \lean{GaussianHilbert.linearFunctional_mul_le_subGaussian_weight}\leanok
  \uses{Lattice-toSemilatticeInf}
  \texttt{GaussianHilbert.linearFunctional\_mul\_le\_subGaussian\_weight}
\end{theorem}
\begin{theorem}[\texttt{zero\_mem\_interior\_integrableExpSet\_coord\_of\_subGaussian}]\label{GaussianHilbert-zero-mem-interior-integrableExpSet-coord-of-subGaussian}
  \lean{GaussianHilbert.zero_mem_interior_integrableExpSet_coord_of_subGaussian}\leanok
  \uses{GaussianHilbert-IsSubGaussianMeasure}
  \texttt{GaussianHilbert.zero\_mem\_interior\_integrableExpSet\_coord\_of\_subGaussian}
\end{theorem}
\begin{theorem}[\texttt{memLp\_coord\_of\_subGaussian}]\label{GaussianHilbert-memLp-coord-of-subGaussian}
  \lean{GaussianHilbert.memLp_coord_of_subGaussian}\leanok
  \uses{GaussianHilbert-zero-mem-interior-integrableExpSet-coord-of-subGaussian, GaussianHilbert-IsSubGaussianMeasure}
  \texttt{GaussianHilbert.memLp\_coord\_of\_subGaussian}
\end{theorem}
\begin{theorem}[\texttt{integrable\_pow\_sum\_sq\_of\_subGaussian}]\label{GaussianHilbert-integrable-pow-sum-sq-of-subGaussian}
  \lean{GaussianHilbert.integrable_pow_sum_sq_of_subGaussian}\leanok
  \uses{GaussianHilbert-IsSubGaussianMeasure, GaussianHilbert-zero-mem-interior-integrableExpSet-sum-sq-of-subGaussian}
  \texttt{GaussianHilbert.integrable\_pow\_sum\_sq\_of\_subGaussian}
\end{theorem}
\begin{theorem}[\texttt{exists\_hasCompactSupport\_integral\_sq\_sub\_le}]\label{GaussianHilbert-MemLp-exists-hasCompactSupport-integral-sq-sub-le}
  \lean{GaussianHilbert.MemLp.exists_hasCompactSupport_integral_sq_sub_le}\leanok
  \texttt{GaussianHilbert.MemLp.exists\_hasCompactSupport\_integral\_sq\_sub\_le}
\end{theorem}
\begin{theorem}[\texttt{integrable\_exp\_coord\_of\_subGaussian}]\label{GaussianHilbert-integrable-exp-coord-of-subGaussian}
  \lean{GaussianHilbert.integrable_exp_coord_of_subGaussian}\leanok
  \uses{GaussianHilbert-IsSubGaussianMeasure, Lattice-toSemilatticeInf}
  \texttt{GaussianHilbert.integrable\_exp\_coord\_of\_subGaussian}
\end{theorem}
\begin{definition}[\texttt{polynomialToL2}]\label{GaussianHilbert-polynomialToL2}
  \lean{GaussianHilbert.polynomialToL2}\leanok
  \uses{GaussianHilbert-IsSubGaussianMeasure, GaussianHilbert-MvPolynomial-eval-memLp-two-of-subGaussian}
  \texttt{GaussianHilbert.polynomialToL2}
\end{definition}
\begin{theorem}[\texttt{tendsto\_integral\_indicator\_ge\_sum\_sq\_eval\_sq\_to\_zero\_of\_subGaussian}]\label{GaussianHilbert-tendsto-integral-indicator-ge-sum-sq-eval-sq-to-zero-of-subGaussian}
  \lean{GaussianHilbert.tendsto_integral_indicator_ge_sum_sq_eval_sq_to_zero_of_subGaussian}\leanok
  \uses{GaussianHilbert-tendsto-integral-indicator-ge-sum-sq-to-zero-of-integrable, GaussianHilbert-IsSubGaussianMeasure, GaussianHilbert-MvPolynomial-eval-memLp-two-of-subGaussian}
  \texttt{GaussianHilbert.tendsto\_integral\_indicator\_ge\_sum\_sq\_eval\_sq\_to\_zero\_of\_subGaussian}
\end{theorem}
\begin{theorem}[\texttt{tendsto\_integral\_indicator\_ge\_sum\_sq\_to\_zero\_of\_integrable}]\label{GaussianHilbert-tendsto-integral-indicator-ge-sum-sq-to-zero-of-integrable}
  \lean{GaussianHilbert.tendsto_integral_indicator_ge_sum_sq_to_zero_of_integrable}\leanok
  \uses{Lattice-toSemilatticeInf}
  \texttt{GaussianHilbert.tendsto\_integral\_indicator\_ge\_sum\_sq\_to\_zero\_of\_integrable}
\end{theorem}
\section{\texttt{GaussianHilbert.OUEigenfunctions}}
\begin{theorem}[\texttt{ouGenerator1D\_hermiteEval}]\label{GaussianHilbert-ouGenerator1D-hermiteEval}
  \lean{GaussianHilbert.ouGenerator1D_hermiteEval}\leanok
  \uses{GaussianHilbert-hermiteEval, GaussianHilbert-ouGenerator1D}
  \texttt{GaussianHilbert.ouGenerator1D\_hermiteEval}
\end{theorem}
\begin{theorem}[\texttt{ouScale\_sq}]\label{GaussianHilbert-ouScale-sq}
  \lean{GaussianHilbert.ouScale_sq}\leanok
  \uses{GaussianHilbert-ouScale}
  \texttt{GaussianHilbert.ouScale\_sq}
\end{theorem}
\begin{definition}[\texttt{mehlerOp}]\label{GaussianHilbert-mehlerOp}
  \lean{GaussianHilbert.mehlerOp}\leanok
  \uses{GaussianHilbert-stdGaussianFin, GaussianHilbert-mehlerLM}
  \texttt{GaussianHilbert.mehlerOp}
\end{definition}
\begin{definition}[\texttt{chaosDiagCLM}]\label{GaussianHilbert-chaosDiagCLM}
  \lean{GaussianHilbert.chaosDiagCLM}\leanok
  \uses{GaussianHilbert-stdGaussianFin, GaussianHilbert-wienerChaos}
  \texttt{GaussianHilbert.chaosDiagCLM}
\end{definition}
\begin{definition}[\texttt{chaosCoordEquiv}]\label{GaussianHilbert-chaosCoordEquiv}
  \lean{GaussianHilbert.chaosCoordEquiv}\leanok
  \uses{GaussianHilbert-wienerChaos-isHilbertSum, GaussianHilbert-stdGaussianFin, GaussianHilbert-wienerChaos}
  \texttt{GaussianHilbert.chaosCoordEquiv}
\end{definition}
\begin{theorem}[\texttt{congr\_simp}]\label{GaussianHilbert-mehlerOp-congr-simp}
  \lean{GaussianHilbert.mehlerOp.congr_simp}\leanok
  \uses{GaussianHilbert-stdGaussianFin, GaussianHilbert-mehlerOp}
  \texttt{GaussianHilbert.mehlerOp.congr\_simp}
\end{theorem}
\begin{theorem}[\texttt{mehlerFun\_zero}]\label{GaussianHilbert-mehlerFun-zero}
  \lean{GaussianHilbert.mehlerFun_zero}\leanok
  \uses{GaussianHilbert-stdGaussianFin, GaussianHilbert-mehlerFun}
  \texttt{GaussianHilbert.mehlerFun\_zero}
\end{theorem}
\begin{theorem}[\texttt{ouNoise\_nonneg}]\label{GaussianHilbert-ouNoise-nonneg}
  \lean{GaussianHilbert.ouNoise_nonneg}\leanok
  \uses{GaussianHilbert-ouNoise}
  \texttt{GaussianHilbert.ouNoise\_nonneg}
\end{theorem}
\begin{definition}[\texttt{ouScale}]\label{GaussianHilbert-ouScale}
  \lean{GaussianHilbert.ouScale}\leanok
  \texttt{GaussianHilbert.ouScale}
\end{definition}
\begin{theorem}[\texttt{mehlerOp\_eq\_ouSemigroupAct}]\label{GaussianHilbert-mehlerOp-eq-ouSemigroupAct}
  \lean{GaussianHilbert.mehlerOp_eq_ouSemigroupAct}\leanok
  \uses{GaussianHilbert-wienerChaos-iSup-topologicalClosure-eq-top, GaussianHilbert-chaosDiagCLM-apply-single, GaussianHilbert-wienerChaos-isHilbertSum, GaussianHilbert-stdGaussianFin, GaussianHilbert-wienerChaos, GaussianHilbert-mehlerOp-eq-smul-of-mem-wienerChaos, GaussianHilbert-mehlerOp, GaussianHilbert-chaosCoordEquiv, GaussianHilbert-ouSemigroupAct, GaussianHilbert-chaosDiagCLM}
  \texttt{GaussianHilbert.mehlerOp\_eq\_ouSemigroupAct}
\end{theorem}
\begin{definition}[\texttt{ouNoise}]\label{GaussianHilbert-ouNoise}
  \lean{GaussianHilbert.ouNoise}\leanok
  \texttt{GaussianHilbert.ouNoise}
\end{definition}
\begin{theorem}[\texttt{ouNoise\_sq}]\label{GaussianHilbert-ouNoise-sq}
  \lean{GaussianHilbert.ouNoise_sq}\leanok
  \uses{GaussianHilbert-ouNoise, GaussianHilbert-one-sub-exp-neg-two-nonneg}
  \texttt{GaussianHilbert.ouNoise\_sq}
\end{theorem}
\begin{theorem}[\texttt{congr\_simp}]\label{GaussianHilbert-chaosDiagCLM-congr-simp}
  \lean{GaussianHilbert.chaosDiagCLM.congr_simp}\leanok
  \uses{GaussianHilbert-stdGaussianFin, GaussianHilbert-wienerChaos, GaussianHilbert-chaosDiagCLM}
  \texttt{GaussianHilbert.chaosDiagCLM.congr\_simp}
\end{theorem}
\begin{theorem}[\texttt{mehlerFun\_apply}]\label{GaussianHilbert-mehlerFun-apply}
  \lean{GaussianHilbert.mehlerFun_apply}\leanok
  \uses{GaussianHilbert-ouAffine, GaussianHilbert-stdGaussianFin, GaussianHilbert-mehlerFun}
  \texttt{GaussianHilbert.mehlerFun\_apply}
\end{theorem}
\begin{theorem}[\texttt{mehlerFun\_memLp}]\label{GaussianHilbert-mehlerFun-memLp}
  \lean{GaussianHilbert.mehlerFun_memLp}\leanok
  \uses{GaussianHilbert-stdGaussianFin, GaussianHilbert-mehlerFun, GaussianHilbert-mehlerFun-measurable}
  \texttt{GaussianHilbert.mehlerFun\_memLp}
\end{theorem}
\begin{theorem}[\texttt{mehlerOp\_eq\_smul\_of\_mem\_wienerChaos}]\label{GaussianHilbert-mehlerOp-eq-smul-of-mem-wienerChaos}
  \lean{GaussianHilbert.mehlerOp_eq_smul_of_mem_wienerChaos}\leanok
  \uses{GaussianHilbert-MultiIndex-totalDegree, GaussianHilbert-stdGaussianFin, GaussianHilbert-wienerChaos, GaussianHilbert-hermiteMultiLp, GaussianHilbert-mehlerOp}
  \texttt{GaussianHilbert.mehlerOp\_eq\_smul\_of\_mem\_wienerChaos}
\end{theorem}
\begin{theorem}[\texttt{mehlerFun\_measurable}]\label{GaussianHilbert-mehlerFun-measurable}
  \lean{GaussianHilbert.mehlerFun_measurable}\leanok
  \uses{GaussianHilbert-ouAffine, GaussianHilbert-stdGaussianFin, GaussianHilbert-mehlerFun}
  \texttt{GaussianHilbert.mehlerFun\_measurable}
\end{theorem}
\begin{theorem}[\texttt{chaosDiagCLM\_apply\_single}]\label{GaussianHilbert-chaosDiagCLM-apply-single}
  \lean{GaussianHilbert.chaosDiagCLM_apply_single}\leanok
  \uses{GaussianHilbert-stdGaussianFin, GaussianHilbert-wienerChaos, GaussianHilbert-chaosDiagCLM}
  \texttt{GaussianHilbert.chaosDiagCLM\_apply\_single}
\end{theorem}
\begin{definition}[\texttt{mehlerLM}]\label{GaussianHilbert-mehlerLM}
  \lean{GaussianHilbert.mehlerLM}\leanok
  \uses{GaussianHilbert-stdGaussianFin, GaussianHilbert-mehlerFun}
  \texttt{GaussianHilbert.mehlerLM}
\end{definition}
\begin{definition}[\texttt{ouGenerator}]\label{GaussianHilbert-ouGenerator}
  \lean{GaussianHilbert.ouGenerator}\leanok
  \texttt{GaussianHilbert.ouGenerator}
\end{definition}
\begin{definition}[\texttt{ouAffine}]\label{GaussianHilbert-ouAffine}
  \lean{GaussianHilbert.ouAffine}\leanok
  \uses{GaussianHilbert-ouNoise, GaussianHilbert-ouScale}
  \texttt{GaussianHilbert.ouAffine}
\end{definition}
\begin{theorem}[\texttt{mehlerFun\_integral\_sq\_le}]\label{GaussianHilbert-mehlerFun-integral-sq-le}
  \lean{GaussianHilbert.mehlerFun_integral_sq_le}\leanok
  \uses{GaussianHilbert-ouAffine, GaussianHilbert-stdGaussianFin, GaussianHilbert-mehlerFun}
  \texttt{GaussianHilbert.mehlerFun\_integral\_sq\_le}
\end{theorem}
\begin{definition}[\texttt{ouSemigroupAct}]\label{GaussianHilbert-ouSemigroupAct}
  \lean{GaussianHilbert.ouSemigroupAct}\leanok
  \uses{GaussianHilbert-stdGaussianFin, GaussianHilbert-wienerChaos, GaussianHilbert-chaosCoordEquiv, GaussianHilbert-chaosDiagCLM}
  \texttt{GaussianHilbert.ouSemigroupAct}
\end{definition}
\begin{theorem}[\texttt{mehlerFun\_const}]\label{GaussianHilbert-mehlerFun-const}
  \lean{GaussianHilbert.mehlerFun_const}\leanok
  \uses{GaussianHilbert-stdGaussianFin, GaussianHilbert-mehlerFun}
  \texttt{GaussianHilbert.mehlerFun\_const}
\end{theorem}
\begin{definition}[\texttt{mehlerFun}]\label{GaussianHilbert-mehlerFun}
  \lean{GaussianHilbert.mehlerFun}\leanok
  \uses{GaussianHilbert-ouAffine, GaussianHilbert-stdGaussianFin}
  \texttt{GaussianHilbert.mehlerFun}
\end{definition}
\begin{theorem}[\texttt{ouSemigroupAct\_eq\_smul\_of\_mem\_wienerChaos}]\label{GaussianHilbert-ouSemigroupAct-eq-smul-of-mem-wienerChaos}
  \lean{GaussianHilbert.ouSemigroupAct_eq_smul_of_mem_wienerChaos}\leanok
  \uses{GaussianHilbert-chaosDiagCLM-apply-single, GaussianHilbert-wienerChaos-isHilbertSum, GaussianHilbert-stdGaussianFin, GaussianHilbert-wienerChaos, GaussianHilbert-chaosCoordEquiv, GaussianHilbert-ouSemigroupAct, GaussianHilbert-chaosDiagCLM}
  \texttt{GaussianHilbert.ouSemigroupAct\_eq\_smul\_of\_mem\_wienerChaos}
\end{theorem}
\begin{definition}[\texttt{ouGenerator1D}]\label{GaussianHilbert-ouGenerator1D}
  \lean{GaussianHilbert.ouGenerator1D}\leanok
  \texttt{GaussianHilbert.ouGenerator1D}
\end{definition}
\begin{theorem}[\texttt{mehlerOp\_apply}]\label{GaussianHilbert-mehlerOp-apply}
  \lean{GaussianHilbert.mehlerOp_apply}\leanok
  \uses{GaussianHilbert-mehlerFun-memLp, GaussianHilbert-stdGaussianFin, GaussianHilbert-mehlerFun, GaussianHilbert-mehlerOp}
  \texttt{GaussianHilbert.mehlerOp\_apply}
\end{theorem}
\begin{theorem}[\texttt{ouAffine\_apply}]\label{GaussianHilbert-ouAffine-apply}
  \lean{GaussianHilbert.ouAffine_apply}\leanok
  \uses{GaussianHilbert-ouAffine, GaussianHilbert-ouNoise, GaussianHilbert-ouScale}
  \texttt{GaussianHilbert.ouAffine\_apply}
\end{theorem}
\begin{theorem}[\texttt{one\_sub\_exp\_neg\_two\_nonneg}]\label{GaussianHilbert-one-sub-exp-neg-two-nonneg}
  \lean{GaussianHilbert.one_sub_exp_neg_two_nonneg}\leanok
  \texttt{GaussianHilbert.one\_sub\_exp\_neg\_two\_nonneg}
\end{theorem}
\begin{theorem}[\texttt{ouGenerator\_hermiteMultiEval}]\label{GaussianHilbert-ouGenerator-hermiteMultiEval}
  \lean{GaussianHilbert.ouGenerator_hermiteMultiEval}\leanok
  \uses{GaussianHilbert-MultiIndex-totalDegree, GaussianHilbert-ouGenerator1D-hermiteEval, GaussianHilbert-ouGenerator, GaussianHilbert-hermiteEval, GaussianHilbert-hermiteMultiEval, GaussianHilbert-ouGenerator1D}
  \texttt{GaussianHilbert.ouGenerator\_hermiteMultiEval}
\end{theorem}
\begin{theorem}[\texttt{congr\_simp}]\label{GaussianHilbert-mehlerLM-congr-simp}
  \lean{GaussianHilbert.mehlerLM.congr_simp}\leanok
  \uses{GaussianHilbert-stdGaussianFin, GaussianHilbert-mehlerLM}
  \texttt{GaussianHilbert.mehlerLM.congr\_simp}
\end{theorem}
\begin{theorem}[\texttt{ouScale\_sq\_add\_ouNoise\_sq}]\label{GaussianHilbert-ouScale-sq-add-ouNoise-sq}
  \lean{GaussianHilbert.ouScale_sq_add_ouNoise_sq}\leanok
  \uses{GaussianHilbert-ouNoise, GaussianHilbert-ouNoise-sq, GaussianHilbert-ouScale, GaussianHilbert-ouScale-sq}
  \texttt{GaussianHilbert.ouScale\_sq\_add\_ouNoise\_sq}
\end{theorem}
\section{\texttt{GaussianHilbert.HypercontractivityFromBE}}
\begin{theorem}[\texttt{ouSemigroupAct\_eLpNorm\_hypercontractive}]\label{GaussianHilbert-ouSemigroupAct-eLpNorm-hypercontractive}
  \lean{GaussianHilbert.ouSemigroupAct_eLpNorm_hypercontractive}\leanok
  \uses{GaussianHilbert-stdGaussianFin-dirichletMarkovSemigroup-isHypercontractive, GaussianHilbert-stdGaussianFin, GaussianHilbert-ouSemigroupAct}
  \texttt{GaussianHilbert.ouSemigroupAct\_eLpNorm\_hypercontractive}
\end{theorem}
\begin{theorem}[\texttt{stdGaussianFin\_dirichletMarkovSemigroup\_satisfiesLogSobolev}]\label{GaussianHilbert-stdGaussianFin-dirichletMarkovSemigroup-satisfiesLogSobolev}
  \lean{GaussianHilbert.stdGaussianFin_dirichletMarkovSemigroup_satisfiesLogSobolev}\leanok
  \uses{GaussianHilbert-stdGaussianFin-LSI}
  \texttt{GaussianHilbert.stdGaussianFin\_dirichletMarkovSemigroup\_satisfiesLogSobolev}
\end{theorem}
\begin{theorem}[\texttt{oneDimGaussianLSI}]\label{GaussianHilbert-oneDimGaussianLSI}
  \lean{GaussianHilbert.oneDimGaussianLSI}\leanok
  \texttt{GaussianHilbert.oneDimGaussianLSI}
\end{theorem}
\begin{theorem}[\texttt{stdGaussianFin\_dirichletMarkovSemigroup\_isHypercontractive}]\label{GaussianHilbert-stdGaussianFin-dirichletMarkovSemigroup-isHypercontractive}
  \lean{GaussianHilbert.stdGaussianFin_dirichletMarkovSemigroup_isHypercontractive}\leanok
  \texttt{GaussianHilbert.stdGaussianFin\_dirichletMarkovSemigroup\_isHypercontractive}
\end{theorem}
\begin{theorem}[\texttt{stdGaussianFin\_hypercontractive\_schema}]\label{GaussianHilbert-stdGaussianFin-hypercontractive-schema}
  \lean{GaussianHilbert.stdGaussianFin_hypercontractive_schema}\leanok
  \texttt{GaussianHilbert.stdGaussianFin\_hypercontractive\_schema}
\end{theorem}
\begin{theorem}[\texttt{stdGaussianFin\_LSI}]\label{GaussianHilbert-stdGaussianFin-LSI}
  \lean{GaussianHilbert.stdGaussianFin_LSI}\leanok
  \texttt{GaussianHilbert.stdGaussianFin\_LSI}
\end{theorem}
\section{\texttt{GaussianHilbert.WienerChaos}}
\begin{definition}[\texttt{hermiteMultiLp}]\label{GaussianHilbert-hermiteMultiLp}
  \lean{GaussianHilbert.hermiteMultiLp}\leanok
  \uses{GaussianHilbert-hermiteMultiEval-memLp, GaussianHilbert-stdGaussianFin, GaussianHilbert-hermiteMultiEval}
  \texttt{GaussianHilbert.hermiteMultiLp}
\end{definition}
\begin{theorem}[\texttt{chaosProjection\_eLpNorm\_two\_le}]\label{GaussianHilbert-chaosProjection-eLpNorm-two-le}
  \lean{GaussianHilbert.chaosProjection_eLpNorm_two_le}\leanok
  \uses{GaussianHilbert-chaosProjection, GaussianHilbert-stdGaussianFin, GaussianHilbert-wienerChaos, GaussianHilbert-wienerChaos-hasOrthogonalProjection}
  \texttt{GaussianHilbert.chaosProjection\_eLpNorm\_two\_le}
\end{theorem}
\begin{theorem}[\texttt{wienerChaos\_isClosed}]\label{GaussianHilbert-wienerChaos-isClosed}
  \lean{GaussianHilbert.wienerChaos_isClosed}\leanok
  \uses{GaussianHilbert-MultiIndex-totalDegree, GaussianHilbert-stdGaussianFin, GaussianHilbert-wienerChaos, GaussianHilbert-hermiteMultiLp}
  \texttt{GaussianHilbert.wienerChaos\_isClosed}
\end{theorem}
\begin{theorem}[\texttt{hermiteMultiEval\_continuous}]\label{GaussianHilbert-hermiteMultiEval-continuous}
  \lean{GaussianHilbert.hermiteMultiEval_continuous}\leanok
  \uses{GaussianHilbert-hermiteEval, GaussianHilbert-hermiteMultiEval}
  \texttt{GaussianHilbert.hermiteMultiEval\_continuous}
\end{theorem}
\begin{theorem}[\texttt{wienerChaos\_completeSpace}]\label{GaussianHilbert-wienerChaos-completeSpace}
  \lean{GaussianHilbert.wienerChaos_completeSpace}\leanok
  \uses{GaussianHilbert-wienerChaos-isClosed, GaussianHilbert-stdGaussianFin, GaussianHilbert-wienerChaos}
  \texttt{GaussianHilbert.wienerChaos\_completeSpace}
\end{theorem}
\begin{theorem}[\texttt{hermiteMultiLp\_coeFn}]\label{GaussianHilbert-hermiteMultiLp-coeFn}
  \lean{GaussianHilbert.hermiteMultiLp_coeFn}\leanok
  \uses{GaussianHilbert-hermiteMultiEval-memLp, GaussianHilbert-stdGaussianFin, GaussianHilbert-hermiteMultiLp, GaussianHilbert-hermiteMultiEval}
  \texttt{GaussianHilbert.hermiteMultiLp\_coeFn}
\end{theorem}
\begin{theorem}[\texttt{hermiteMultiEval\_memLp}]\label{GaussianHilbert-hermiteMultiEval-memLp}
  \lean{GaussianHilbert.hermiteMultiEval_memLp}\leanok
  \uses{GaussianHilbert-stdGaussianFin, GaussianHilbert-hermiteMultiEval-sq-integrable, GaussianHilbert-hermiteMultiEval, GaussianHilbert-hermiteMultiEval-continuous}
  \texttt{GaussianHilbert.hermiteMultiEval\_memLp}
\end{theorem}
\begin{theorem}[\texttt{wienerChaos\_orthogonal}]\label{GaussianHilbert-wienerChaos-orthogonal}
  \lean{GaussianHilbert.wienerChaos_orthogonal}\leanok
  \uses{GaussianHilbert-MultiIndex-totalDegree, GaussianHilbert-stdGaussianFin, GaussianHilbert-hermiteMultiLp-inner-eq-zero-of-totalDegree-ne, GaussianHilbert-wienerChaos, GaussianHilbert-hermiteMultiLp}
  \texttt{GaussianHilbert.wienerChaos\_orthogonal}
\end{theorem}
\begin{definition}[\texttt{wienerChaosLE}]\label{GaussianHilbert-wienerChaosLE}
  \lean{GaussianHilbert.wienerChaosLE}\leanok
  \uses{GaussianHilbert-stdGaussianFin, GaussianHilbert-wienerChaos}
  \texttt{GaussianHilbert.wienerChaosLE}
\end{definition}
\begin{theorem}[\texttt{hermiteMultiEval\_mem\_wienerChaos}]\label{GaussianHilbert-hermiteMultiEval-mem-wienerChaos}
  \lean{GaussianHilbert.hermiteMultiEval_mem_wienerChaos}\leanok
  \uses{GaussianHilbert-MultiIndex-totalDegree, GaussianHilbert-hermiteMultiEval-memLp, GaussianHilbert-stdGaussianFin, GaussianHilbert-wienerChaos, GaussianHilbert-hermiteMultiLp, GaussianHilbert-hermiteMultiEval, GaussianHilbert-hermiteMultiLp-mem-wienerChaos}
  \texttt{GaussianHilbert.hermiteMultiEval\_mem\_wienerChaos}
\end{theorem}
\begin{definition}[\texttt{hermiteFamilyOfDegree}]\label{GaussianHilbert-hermiteFamilyOfDegree}
  \lean{GaussianHilbert.hermiteFamilyOfDegree}\leanok
  \uses{GaussianHilbert-MultiIndex-totalDegree, GaussianHilbert-hermiteMultiEval}
  \texttt{GaussianHilbert.hermiteFamilyOfDegree}
\end{definition}
\begin{theorem}[\texttt{chaosProjection\_mem\_wienerChaos}]\label{GaussianHilbert-chaosProjection-mem-wienerChaos}
  \lean{GaussianHilbert.chaosProjection_mem_wienerChaos}\leanok
  \uses{GaussianHilbert-chaosProjection, GaussianHilbert-stdGaussianFin, GaussianHilbert-wienerChaos, GaussianHilbert-wienerChaos-hasOrthogonalProjection}
  \texttt{GaussianHilbert.chaosProjection\_mem\_wienerChaos}
\end{theorem}
\begin{theorem}[\texttt{wienerChaos\_hasOrthogonalProjection}]\label{GaussianHilbert-wienerChaos-hasOrthogonalProjection}
  \lean{GaussianHilbert.wienerChaos_hasOrthogonalProjection}\leanok
  \uses{GaussianHilbert-wienerChaos-completeSpace, GaussianHilbert-stdGaussianFin, GaussianHilbert-wienerChaos}
  \texttt{GaussianHilbert.wienerChaos\_hasOrthogonalProjection}
\end{theorem}
\begin{theorem}[\texttt{wienerChaos\_isHilbertSum}]\label{GaussianHilbert-wienerChaos-isHilbertSum}
  \lean{GaussianHilbert.wienerChaos_isHilbertSum}\leanok
  \uses{GaussianHilbert-wienerChaos-iSup-topologicalClosure-eq-top, GaussianHilbert-wienerChaos-completeSpace, GaussianHilbert-stdGaussianFin, GaussianHilbert-wienerChaos, GaussianHilbert-wienerChaos-orthogonal}
  \texttt{GaussianHilbert.wienerChaos\_isHilbertSum}
\end{theorem}
\begin{theorem}[\texttt{chaosProjection\_sum\_eq\_of\_mem\_wienerChaosLE}]\label{GaussianHilbert-chaosProjection-sum-eq-of-mem-wienerChaosLE}
  \lean{GaussianHilbert.chaosProjection_sum_eq_of_mem_wienerChaosLE}\leanok
  \uses{GaussianHilbert-chaosProjection, GaussianHilbert-wienerChaos-completeSpace, GaussianHilbert-stdGaussianFin, GaussianHilbert-wienerChaos, GaussianHilbert-wienerChaos-orthogonal, GaussianHilbert-wienerChaosLE, GaussianHilbert-wienerChaos-hasOrthogonalProjection}
  \texttt{GaussianHilbert.chaosProjection\_sum\_eq\_of\_mem\_wienerChaosLE}
\end{theorem}
\begin{theorem}[\texttt{wienerChaos\_iSup\_topologicalClosure\_eq\_top}]\label{GaussianHilbert-wienerChaos-iSup-topologicalClosure-eq-top}
  \lean{GaussianHilbert.wienerChaos_iSup_topologicalClosure_eq_top}\leanok
  \uses{GaussianHilbert-MultiIndex-totalDegree, GaussianHilbert-hermiteMultiEval-memLp, GaussianHilbert-hermiteMultiLp-coeFn, GaussianHilbert-stdGaussianFin, GaussianHilbert-wienerChaos, GaussianHilbert-hermiteMultiLp, GaussianHilbert-hermiteMultiEval, GaussianHilbert-hermiteMulti-dense, GaussianHilbert-hermiteMultiLp-mem-wienerChaos}
  \texttt{GaussianHilbert.wienerChaos\_iSup\_topologicalClosure\_eq\_top}
\end{theorem}
\begin{theorem}[\texttt{hermiteMultiLp\_inner\_eq\_zero\_of\_totalDegree\_ne}]\label{GaussianHilbert-hermiteMultiLp-inner-eq-zero-of-totalDegree-ne}
  \lean{GaussianHilbert.hermiteMultiLp_inner_eq_zero_of_totalDegree_ne}\leanok
  \uses{GaussianHilbert-MultiIndex-totalDegree, GaussianHilbert-hermiteMultiLp-coeFn, GaussianHilbert-stdGaussianFin, GaussianHilbert-hermiteMulti-orthogonality, GaussianHilbert-hermiteMultiLp, GaussianHilbert-hermiteMultiEval}
  \texttt{GaussianHilbert.hermiteMultiLp\_inner\_eq\_zero\_of\_totalDegree\_ne}
\end{theorem}
\begin{definition}[\texttt{chaosProjection}]\label{GaussianHilbert-chaosProjection}
  \lean{GaussianHilbert.chaosProjection}\leanok
  \uses{GaussianHilbert-stdGaussianFin, GaussianHilbert-wienerChaos, GaussianHilbert-wienerChaos-hasOrthogonalProjection}
  \texttt{GaussianHilbert.chaosProjection}
\end{definition}
\begin{definition}[\texttt{wienerChaos}]\label{GaussianHilbert-wienerChaos}
  \lean{GaussianHilbert.wienerChaos}\leanok
  \uses{GaussianHilbert-MultiIndex-totalDegree, GaussianHilbert-stdGaussianFin, GaussianHilbert-hermiteMultiLp}
  \texttt{GaussianHilbert.wienerChaos}
\end{definition}
\begin{theorem}[\texttt{hermiteMultiLp\_mem\_wienerChaos}]\label{GaussianHilbert-hermiteMultiLp-mem-wienerChaos}
  \lean{GaussianHilbert.hermiteMultiLp_mem_wienerChaos}\leanok
  \uses{GaussianHilbert-MultiIndex-totalDegree, GaussianHilbert-stdGaussianFin, GaussianHilbert-wienerChaos, GaussianHilbert-hermiteMultiLp}
  \texttt{GaussianHilbert.hermiteMultiLp\_mem\_wienerChaos}
\end{theorem}
\begin{theorem}[\texttt{hermiteMultiEval\_sq\_integrable}]\label{GaussianHilbert-hermiteMultiEval-sq-integrable}
  \lean{GaussianHilbert.hermiteMultiEval_sq_integrable}\leanok
  \uses{GaussianHilbert-stdGaussianFin, GaussianHilbert-hermiteMulti-orthogonality, GaussianHilbert-hermiteMultiEval}
  \texttt{GaussianHilbert.hermiteMultiEval\_sq\_integrable}
\end{theorem}
\section{\texttt{GaussianHilbert.HermitePolynomials}}
\begin{theorem}[\texttt{hermiteMulti\_orthogonality}]\label{GaussianHilbert-hermiteMulti-orthogonality}
  \lean{GaussianHilbert.hermiteMulti_orthogonality}\leanok
  \uses{GaussianField-wickMonomial-inner-gaussianReal-one, GaussianHilbert-stdGaussianFin, GaussianHilbert-hermiteEval, GaussianHilbert-hermiteMultiEval}
  \texttt{GaussianHilbert.hermiteMulti\_orthogonality}
\end{theorem}
\begin{theorem}[\texttt{hermiteMulti\_l2\_pos}]\label{GaussianHilbert-hermiteMulti-l2-pos}
  \lean{GaussianHilbert.hermiteMulti_l2_pos}\leanok
  \uses{GaussianHilbert-stdGaussianFin, GaussianHilbert-hermiteMulti-orthogonality, GaussianHilbert-hermiteMultiEval}
  \texttt{GaussianHilbert.hermiteMulti\_l2\_pos}
\end{theorem}
\begin{definition}[\texttt{stdGaussianFin}]\label{GaussianHilbert-stdGaussianFin}
  \lean{GaussianHilbert.stdGaussianFin}\leanok
  \texttt{GaussianHilbert.stdGaussianFin}
\end{definition}
\begin{definition}[\texttt{totalDegree}]\label{GaussianHilbert-MultiIndex-totalDegree}
  \lean{GaussianHilbert.MultiIndex.totalDegree}\leanok
  \texttt{GaussianHilbert.MultiIndex.totalDegree}
\end{definition}
\begin{definition}[\texttt{hermiteEval}]\label{GaussianHilbert-hermiteEval}
  \lean{GaussianHilbert.hermiteEval}\leanok
  \texttt{GaussianHilbert.hermiteEval}
\end{definition}
\begin{theorem}[\texttt{hermiteMulti\_dense}]\label{GaussianHilbert-hermiteMulti-dense}
  \lean{GaussianHilbert.hermiteMulti_dense}\leanok
  \uses{GaussianHilbert-isSubGaussianMeasure-pi-gaussianReal, GaussianHilbert-stdGaussianFin, GaussianHilbert-hermiteMultiEval, GaussianHilbert-IsSubGaussianMeasure, GaussianHilbert-polynomial-dense-L2-of-subGaussian}
  \texttt{GaussianHilbert.hermiteMulti\_dense}
\end{theorem}
\begin{definition}[\texttt{hermiteMultiEval}]\label{GaussianHilbert-hermiteMultiEval}
  \lean{GaussianHilbert.hermiteMultiEval}\leanok
  \uses{GaussianHilbert-hermiteEval}
  \texttt{GaussianHilbert.hermiteMultiEval}
\end{definition}
\section{\texttt{GaussianHilbert.PolynomialChaosConcentration}}
\begin{theorem}[\texttt{bonami\_nelson\_chaosLE}]\label{GaussianHilbert-bonami-nelson-chaosLE}
  \lean{GaussianHilbert.bonami_nelson_chaosLE}\leanok
  \uses{GaussianHilbert-chaosProjection, GaussianHilbert-chaosProjection-eLpNorm-two-le, GaussianHilbert-stdGaussianFin, GaussianHilbert-wienerChaos, GaussianHilbert-wienerChaosLE, GaussianHilbert-chaosProjection-sum-eq-of-mem-wienerChaosLE, GaussianHilbert-chaosProjection-mem-wienerChaos, GaussianHilbert-bonami-nelson-chaos}
  \texttt{GaussianHilbert.bonami\_nelson\_chaosLE}
\end{theorem}
\begin{theorem}[\texttt{bonami\_nelson\_chaos}]\label{GaussianHilbert-bonami-nelson-chaos}
  \lean{GaussianHilbert.bonami_nelson_chaos}\leanok
  \uses{GaussianHilbert-ouSemigroupAct-eq-smul-of-mem-wienerChaos, GaussianHilbert-stdGaussianFin, GaussianHilbert-wienerChaos, GaussianHilbert-ouSemigroupAct-eLpNorm-hypercontractive, GaussianHilbert-ouSemigroupAct}
  \texttt{GaussianHilbert.bonami\_nelson\_chaos}
\end{theorem}
\begin{theorem}[\texttt{polynomial\_chaos\_concentration}]\label{GaussianHilbert-polynomial-chaos-concentration}
  \lean{GaussianHilbert.polynomial_chaos_concentration}\leanok
  \uses{GaussianHilbert-stdGaussianFin, GaussianHilbert-wienerChaosLE, Lattice-toSemilatticeInf}
  \texttt{GaussianHilbert.polynomial\_chaos\_concentration}
\end{theorem}
\chapter{pphi2: common QFT framework}
\section{\texttt{Common.QFT.Euclidean.Formulations}}
\begin{definition}[\texttt{TestFunction}]\label{Common-QFT-Euclidean-Formulation-TestFunction}
  \lean{Common.QFT.Euclidean.Formulation.TestFunction}\leanok
  \uses{Common-QFT-Euclidean-Formulation}
  \texttt{Common.QFT.Euclidean.Formulation.TestFunction}
\end{definition}
\begin{definition}[\texttt{toPairingMeasureModel}]\label{Common-QFT-Euclidean-MeasureModel-toPairingMeasureModel}
  \lean{Common.QFT.Euclidean.MeasureModel.toPairingMeasureModel}\leanok
  \uses{Common-QFT-Euclidean-PairingMeasureModel, Common-QFT-Euclidean-MeasureModel, Common-QFT-Euclidean-Formulation}
  \texttt{Common.QFT.Euclidean.MeasureModel.toPairingMeasureModel}
\end{definition}
\begin{theorem}[\texttt{instContinuousSMulTestFunction}]\label{Common-QFT-Euclidean-SchwingerFormulation-instContinuousSMulTestFunction}
  \lean{Common.QFT.Euclidean.SchwingerFormulation.instContinuousSMulTestFunction}\leanok
  \uses{Common-QFT-Euclidean-SchwingerFormulation-instTopologicalSpaceTestFunction, Common-QFT-Euclidean-Formulation-TestFunction, Common-QFT-Euclidean-SchwingerFormulation-toFormulation, Common-QFT-Euclidean-SchwingerFormulation-instAddCommGroupTestFunction, Common-QFT-Euclidean-SchwingerFormulation-instModuleTestFunction, Common-QFT-Euclidean-SchwingerFormulation}
  \texttt{Common.QFT.Euclidean.SchwingerFormulation.instContinuousSMulTestFunction}
\end{theorem}
\begin{theorem}[\texttt{realToComplex\_add}]\label{Common-QFT-Euclidean-SchwingerFormulation-realToComplex-add}
  \lean{Common.QFT.Euclidean.SchwingerFormulation.realToComplex_add}\leanok
  \uses{Common-QFT-Euclidean-SchwingerFormulation-instAddCommGroupComplexTestFunction, Common-QFT-Euclidean-Formulation-TestFunction, Common-QFT-Euclidean-Formulation-ComplexTestFunction, Common-QFT-Euclidean-SchwingerFormulation-toFormulation, Common-QFT-Euclidean-SchwingerFormulation-realToComplex, Common-QFT-Euclidean-SchwingerFormulation-instAddCommGroupTestFunction, Common-QFT-Euclidean-SchwingerFormulation}
  \texttt{Common.QFT.Euclidean.SchwingerFormulation.realToComplex\_add}
\end{theorem}
\begin{definition}[\texttt{ctorIdx}]\label{Common-QFT-Euclidean-Formulation-ctorIdx}
  \lean{Common.QFT.Euclidean.Formulation.ctorIdx}\leanok
  \uses{Common-QFT-Euclidean-Formulation}
  \texttt{Common.QFT.Euclidean.Formulation.ctorIdx}
\end{definition}
\begin{definition}[\texttt{PairingMeasureModel}]\label{Common-QFT-Euclidean-PairingMeasureModel}
  \lean{Common.QFT.Euclidean.PairingMeasureModel}\leanok
  \uses{Common-QFT-Euclidean-Formulation}
  \texttt{Common.QFT.Euclidean.PairingMeasureModel}
\end{definition}
\begin{definition}[\texttt{complexGeneratingFunctional}]\label{Common-QFT-Euclidean-GeneratingFunctionalModel-complexGeneratingFunctional}
  \lean{Common.QFT.Euclidean.GeneratingFunctionalModel.complexGeneratingFunctional}\leanok
  \uses{Common-QFT-Euclidean-GeneratingFunctionalModel, Common-QFT-Euclidean-Formulation, Common-QFT-Euclidean-Formulation-ComplexTestFunction}
  \texttt{Common.QFT.Euclidean.GeneratingFunctionalModel.complexGeneratingFunctional}
\end{definition}
\begin{theorem}[\texttt{positiveTimeSubmodule\_eq}]\label{Common-QFT-Euclidean-SchwingerFormulation-positiveTimeSubmodule-eq}
  \lean{Common.QFT.Euclidean.SchwingerFormulation.positiveTimeSubmodule_eq}\leanok
  \uses{Common-QFT-Euclidean-Formulation-positiveTimeTestFunctions, Common-QFT-Euclidean-Formulation-TestFunction, Common-QFT-Euclidean-SchwingerFormulation-toFormulation, Common-QFT-Euclidean-SchwingerFormulation-instAddCommGroupTestFunction, Common-QFT-Euclidean-SchwingerFormulation-positiveTimeSubmodule, Common-QFT-Euclidean-SchwingerFormulation-instModuleTestFunction, Common-QFT-Euclidean-SchwingerFormulation}
  \texttt{Common.QFT.Euclidean.SchwingerFormulation.positiveTimeSubmodule\_eq}
\end{theorem}
\begin{definition}[\texttt{MeasureModel}]\label{Common-QFT-Euclidean-MeasureModel}
  \lean{Common.QFT.Euclidean.MeasureModel}\leanok
  \uses{Common-QFT-Euclidean-Formulation}
  \texttt{Common.QFT.Euclidean.MeasureModel}
\end{definition}
\begin{theorem}[\texttt{ext}]\label{Common-QFT-Euclidean-TensorSchwingerModel-ext}
  \lean{Common.QFT.Euclidean.TensorSchwingerModel.ext}\leanok
  \uses{Common-QFT-Euclidean-Formulation-TestFunction, Common-QFT-Euclidean-TensorSchwingerModel, Common-QFT-Euclidean-Formulation, Common-QFT-Euclidean-TensorSchwingerModel-schwinger}
  \texttt{Common.QFT.Euclidean.TensorSchwingerModel.ext}
\end{theorem}
\begin{definition}[\texttt{NPointTestFunction}]\label{Common-QFT-Euclidean-SchwingerFormulation-NPointTestFunction}
  \lean{Common.QFT.Euclidean.SchwingerFormulation.NPointTestFunction}\leanok
  \uses{Common-QFT-Euclidean-SchwingerFormulation}
  \texttt{Common.QFT.Euclidean.SchwingerFormulation.NPointTestFunction}
\end{definition}
\begin{definition}[\texttt{DistributionalSchwingerModel}]\label{Common-QFT-Euclidean-DistributionalSchwingerModel}
  \lean{Common.QFT.Euclidean.DistributionalSchwingerModel}\leanok
  \uses{Common-QFT-Euclidean-SchwingerFormulation}
  \texttt{Common.QFT.Euclidean.DistributionalSchwingerModel}
\end{definition}
\begin{definition}[\texttt{ctorIdx}]\label{Common-QFT-Euclidean-PairingMeasureModel-ctorIdx}
  \lean{Common.QFT.Euclidean.PairingMeasureModel.ctorIdx}\leanok
  \uses{Common-QFT-Euclidean-PairingMeasureModel, Common-QFT-Euclidean-Formulation}
  \texttt{Common.QFT.Euclidean.PairingMeasureModel.ctorIdx}
\end{definition}
\begin{definition}[\texttt{measure}]\label{Common-QFT-Euclidean-PairingMeasureModel-measure}
  \lean{Common.QFT.Euclidean.PairingMeasureModel.measure}\leanok
  \uses{Common-QFT-Euclidean-PairingMeasureModel, Common-QFT-Euclidean-Formulation, Common-QFT-Euclidean-PairingMeasureModel-instMeasurableSpaceFieldConfig}
  Lean signature ``\texttt{lean def measure (T : E →L[ℝ] H) : @Measure (Configuration E) instMeasurableSpaceConfiguration }``  Informal: The Gaussian probability measure on $E'$ constructed from $T$. For infinite-dimensional $H$, this is the pushforward of the noise measure through the series-limit map. For finite-dimensional $H$, an isometric embedding into $\ell^2$ is used first.
\end{definition}
\begin{definition}[\texttt{instAddCommGroupNPointTestFunction}]\label{Common-QFT-Euclidean-SchwingerFormulation-instAddCommGroupNPointTestFunction}
  \lean{Common.QFT.Euclidean.SchwingerFormulation.instAddCommGroupNPointTestFunction}\leanok
  \uses{Common-QFT-Euclidean-SchwingerFormulation-NPointTestFunction, Common-QFT-Euclidean-SchwingerFormulation}
  \texttt{Common.QFT.Euclidean.SchwingerFormulation.instAddCommGroupNPointTestFunction}
\end{definition}
\begin{definition}[\texttt{instModuleComplexTestFunction}]\label{Common-QFT-Euclidean-SchwingerFormulation-instModuleComplexTestFunction}
  \lean{Common.QFT.Euclidean.SchwingerFormulation.instModuleComplexTestFunction}\leanok
  \uses{Common-QFT-Euclidean-SchwingerFormulation-instAddCommGroupComplexTestFunction, Common-QFT-Euclidean-Formulation-ComplexTestFunction, Common-QFT-Euclidean-SchwingerFormulation-toFormulation, Common-QFT-Euclidean-SchwingerFormulation}
  \texttt{Common.QFT.Euclidean.SchwingerFormulation.instModuleComplexTestFunction}
\end{definition}
\begin{definition}[\texttt{schwinger}]\label{Common-QFT-Euclidean-TensorSchwingerModel-schwinger}
  \lean{Common.QFT.Euclidean.TensorSchwingerModel.schwinger}\leanok
  \uses{Common-QFT-Euclidean-Formulation-TestFunction, Common-QFT-Euclidean-TensorSchwingerModel, Common-QFT-Euclidean-Formulation}
  \texttt{Common.QFT.Euclidean.TensorSchwingerModel.schwinger}
\end{definition}
\begin{theorem}[\texttt{schwingerContinuous}]\label{Common-QFT-Euclidean-DistributionalSchwingerModel-schwingerContinuous}
  \lean{Common.QFT.Euclidean.DistributionalSchwingerModel.schwingerContinuous}\leanok
  \uses{Common-QFT-Euclidean-SchwingerFormulation-instAddCommGroupNPointTestFunction, Common-QFT-Euclidean-DistributionalSchwingerModel, Common-QFT-Euclidean-SchwingerFormulation-NPointTestFunction, Common-QFT-Euclidean-SchwingerFormulation-instTopologicalSpaceNPointTestFunction, Common-QFT-Euclidean-DistributionalSchwingerModel-schwinger, Common-QFT-Euclidean-SchwingerFormulation-instModuleNPointTestFunction, Common-QFT-Euclidean-SchwingerFormulation}
  \texttt{Common.QFT.Euclidean.DistributionalSchwingerModel.schwingerContinuous}
\end{theorem}
\begin{definition}[\texttt{instTopologicalSpaceTestFunction}]\label{Common-QFT-Euclidean-SchwingerFormulation-instTopologicalSpaceTestFunction}
  \lean{Common.QFT.Euclidean.SchwingerFormulation.instTopologicalSpaceTestFunction}\leanok
  \uses{Common-QFT-Euclidean-Formulation-TestFunction, Common-QFT-Euclidean-SchwingerFormulation-toFormulation, Common-QFT-Euclidean-SchwingerFormulation}
  \texttt{Common.QFT.Euclidean.SchwingerFormulation.instTopologicalSpaceTestFunction}
\end{definition}
\begin{definition}[\texttt{instModuleTestFunction}]\label{Common-QFT-Euclidean-SchwingerFormulation-instModuleTestFunction}
  \lean{Common.QFT.Euclidean.SchwingerFormulation.instModuleTestFunction}\leanok
  \uses{Common-QFT-Euclidean-Formulation-TestFunction, Common-QFT-Euclidean-SchwingerFormulation-toFormulation, Common-QFT-Euclidean-SchwingerFormulation-instAddCommGroupTestFunction, Common-QFT-Euclidean-SchwingerFormulation}
  \texttt{Common.QFT.Euclidean.SchwingerFormulation.instModuleTestFunction}
\end{definition}
\begin{definition}[\texttt{toGeneratingFunctionalModel}]\label{Common-QFT-Euclidean-MeasureModel-toGeneratingFunctionalModel}
  \lean{Common.QFT.Euclidean.MeasureModel.toGeneratingFunctionalModel}\leanok
  \uses{Common-QFT-Euclidean-MeasureModel, Common-QFT-Euclidean-GeneratingFunctionalModel, Common-QFT-Euclidean-Formulation}
  \texttt{Common.QFT.Euclidean.MeasureModel.toGeneratingFunctionalModel}
\end{definition}
\begin{definition}[\texttt{schwingerModel}]\label{Common-QFT-Euclidean-ReconstructionInput-schwingerModel}
  \lean{Common.QFT.Euclidean.ReconstructionInput.schwingerModel}\leanok
  \uses{Common-QFT-Euclidean-ReconstructionHypothesis, Common-QFT-Euclidean-DistributionalSchwingerModel, Common-QFT-Euclidean-ReconstructionInput, Common-QFT-Euclidean-SchwingerFormulation}
  \texttt{Common.QFT.Euclidean.ReconstructionInput.schwingerModel}
\end{definition}
\begin{theorem}[\texttt{compatible}]\label{Common-QFT-Euclidean-TensorToDistributionalBridge-compatible}
  \lean{Common.QFT.Euclidean.TensorToDistributionalBridge.compatible}\leanok
  \uses{Common-QFT-Euclidean-TensorToDistributionalBridge-distributionalSchwinger, Common-QFT-Euclidean-SchwingerFormulation-instAddCommGroupNPointTestFunction, Common-QFT-Euclidean-SchwingerFormulation-NPointTestFunction, Common-QFT-Euclidean-Formulation-TestFunction, Common-QFT-Euclidean-TensorSchwingerModel, Common-QFT-Euclidean-TensorToDistributionalBridge, Common-QFT-Euclidean-DistributionalSchwingerModel-tensorLift, Common-QFT-Euclidean-DistributionalSchwingerModel-schwinger, Common-QFT-Euclidean-SchwingerFormulation-toFormulation, Common-QFT-Euclidean-TensorSchwingerModel-schwinger, Common-QFT-Euclidean-SchwingerFormulation-instModuleNPointTestFunction, Common-QFT-Euclidean-SchwingerFormulation}
  \texttt{Common.QFT.Euclidean.TensorToDistributionalBridge.compatible}
\end{theorem}
\begin{definition}[\texttt{SpaceTime}]\label{Common-QFT-Euclidean-Formulation-SpaceTime}
  \lean{Common.QFT.Euclidean.Formulation.SpaceTime}\leanok
  \uses{Common-QFT-Euclidean-Formulation}
  \texttt{Common.QFT.Euclidean.Formulation.SpaceTime}
\end{definition}
\begin{definition}[\texttt{instTopologicalSpaceComplexTestFunction}]\label{Common-QFT-Euclidean-SchwingerFormulation-instTopologicalSpaceComplexTestFunction}
  \lean{Common.QFT.Euclidean.SchwingerFormulation.instTopologicalSpaceComplexTestFunction}\leanok
  \uses{Common-QFT-Euclidean-Formulation-ComplexTestFunction, Common-QFT-Euclidean-SchwingerFormulation-toFormulation, Common-QFT-Euclidean-SchwingerFormulation}
  \texttt{Common.QFT.Euclidean.SchwingerFormulation.instTopologicalSpaceComplexTestFunction}
\end{definition}
\begin{theorem}[\texttt{hypothesis}]\label{Common-QFT-Euclidean-ReconstructionInput-hypothesis}
  \lean{Common.QFT.Euclidean.ReconstructionInput.hypothesis}\leanok
  \uses{Common-QFT-Euclidean-ReconstructionHypothesis, Common-QFT-Euclidean-ReconstructionInput-schwingerModel, Common-QFT-Euclidean-ReconstructionInput, Common-QFT-Euclidean-SchwingerFormulation}
  \texttt{Common.QFT.Euclidean.ReconstructionInput.hypothesis}
\end{theorem}
\begin{definition}[\texttt{toFormulation}]\label{Common-QFT-Euclidean-SchwingerFormulation-toFormulation}
  \lean{Common.QFT.Euclidean.SchwingerFormulation.toFormulation}\leanok
  \uses{Common-QFT-Euclidean-Formulation, Common-QFT-Euclidean-SchwingerFormulation}
  \texttt{Common.QFT.Euclidean.SchwingerFormulation.toFormulation}
\end{definition}
\begin{definition}[\texttt{MeasureToTensorBridge}]\label{Common-QFT-Euclidean-MeasureToTensorBridge}
  \lean{Common.QFT.Euclidean.MeasureToTensorBridge}\leanok
  \uses{Common-QFT-Euclidean-PairingMeasureModel, Common-QFT-Euclidean-Formulation}
  \texttt{Common.QFT.Euclidean.MeasureToTensorBridge}
\end{definition}
\begin{definition}[\texttt{ComplexTestFunction}]\label{Common-QFT-Euclidean-Formulation-ComplexTestFunction}
  \lean{Common.QFT.Euclidean.Formulation.ComplexTestFunction}\leanok
  \uses{Common-QFT-Euclidean-Formulation}
  \texttt{Common.QFT.Euclidean.Formulation.ComplexTestFunction}
\end{definition}
\begin{theorem}[\texttt{instContinuousSMulComplexTestFunction}]\label{Common-QFT-Euclidean-SchwingerFormulation-instContinuousSMulComplexTestFunction}
  \lean{Common.QFT.Euclidean.SchwingerFormulation.instContinuousSMulComplexTestFunction}\leanok
  \uses{Common-QFT-Euclidean-SchwingerFormulation-instAddCommGroupComplexTestFunction, Common-QFT-Euclidean-SchwingerFormulation-instModuleComplexTestFunction, Common-QFT-Euclidean-Formulation-ComplexTestFunction, Common-QFT-Euclidean-SchwingerFormulation-toFormulation, Common-QFT-Euclidean-SchwingerFormulation, Common-QFT-Euclidean-SchwingerFormulation-instTopologicalSpaceComplexTestFunction}
  \texttt{Common.QFT.Euclidean.SchwingerFormulation.instContinuousSMulComplexTestFunction}
\end{theorem}
\begin{definition}[\texttt{TensorSchwingerModel}]\label{Common-QFT-Euclidean-TensorSchwingerModel}
  \lean{Common.QFT.Euclidean.TensorSchwingerModel}\leanok
  \uses{Common-QFT-Euclidean-Formulation}
  \texttt{Common.QFT.Euclidean.TensorSchwingerModel}
\end{definition}
\begin{definition}[\texttt{positiveTimeSubmodule}]\label{Common-QFT-Euclidean-SchwingerFormulation-positiveTimeSubmodule}
  \lean{Common.QFT.Euclidean.SchwingerFormulation.positiveTimeSubmodule}\leanok
  \uses{Common-QFT-Euclidean-Formulation-TestFunction, Common-QFT-Euclidean-SchwingerFormulation-toFormulation, Common-QFT-Euclidean-SchwingerFormulation-instAddCommGroupTestFunction, Common-QFT-Euclidean-SchwingerFormulation-instModuleTestFunction, Common-QFT-Euclidean-SchwingerFormulation}
  \texttt{Common.QFT.Euclidean.SchwingerFormulation.positiveTimeSubmodule}
\end{definition}
\begin{theorem}[\texttt{isProbability}]\label{Common-QFT-Euclidean-PairingMeasureModel-isProbability}
  \lean{Common.QFT.Euclidean.PairingMeasureModel.isProbability}\leanok
  \uses{Common-QFT-Euclidean-PairingMeasureModel, Common-QFT-Euclidean-Formulation, Common-QFT-Euclidean-PairingMeasureModel-measure, Common-QFT-Euclidean-PairingMeasureModel-instMeasurableSpaceFieldConfig}
  \texttt{Common.QFT.Euclidean.PairingMeasureModel.isProbability}
\end{theorem}
\begin{definition}[\texttt{ReconstructionInput}]\label{Common-QFT-Euclidean-ReconstructionInput}
  \lean{Common.QFT.Euclidean.ReconstructionInput}\leanok
  \uses{Common-QFT-Euclidean-ReconstructionHypothesis, Common-QFT-Euclidean-SchwingerFormulation}
  \texttt{Common.QFT.Euclidean.ReconstructionInput}
\end{definition}
\begin{definition}[\texttt{instModuleNPointTestFunction}]\label{Common-QFT-Euclidean-SchwingerFormulation-instModuleNPointTestFunction}
  \lean{Common.QFT.Euclidean.SchwingerFormulation.instModuleNPointTestFunction}\leanok
  \uses{Common-QFT-Euclidean-SchwingerFormulation-instAddCommGroupNPointTestFunction, Common-QFT-Euclidean-SchwingerFormulation-NPointTestFunction, Common-QFT-Euclidean-SchwingerFormulation}
  \texttt{Common.QFT.Euclidean.SchwingerFormulation.instModuleNPointTestFunction}
\end{definition}
\begin{definition}[\texttt{positiveTimeTestFunctions}]\label{Common-QFT-Euclidean-Formulation-positiveTimeTestFunctions}
  \lean{Common.QFT.Euclidean.Formulation.positiveTimeTestFunctions}\leanok
  \uses{Common-QFT-Euclidean-Formulation-TestFunction, Common-QFT-Euclidean-Formulation}
  \texttt{Common.QFT.Euclidean.Formulation.positiveTimeTestFunctions}
\end{definition}
\begin{theorem}[\texttt{realToComplex\_continuous}]\label{Common-QFT-Euclidean-SchwingerFormulation-realToComplex-continuous}
  \lean{Common.QFT.Euclidean.SchwingerFormulation.realToComplex_continuous}\leanok
  \uses{Common-QFT-Euclidean-SchwingerFormulation-instTopologicalSpaceTestFunction, Common-QFT-Euclidean-Formulation-TestFunction, Common-QFT-Euclidean-Formulation-ComplexTestFunction, Common-QFT-Euclidean-SchwingerFormulation-toFormulation, Common-QFT-Euclidean-SchwingerFormulation-realToComplex, Common-QFT-Euclidean-SchwingerFormulation, Common-QFT-Euclidean-SchwingerFormulation-instTopologicalSpaceComplexTestFunction}
  \texttt{Common.QFT.Euclidean.SchwingerFormulation.realToComplex\_continuous}
\end{theorem}
\begin{definition}[\texttt{Formulation}]\label{Common-QFT-Euclidean-Formulation}
  \lean{Common.QFT.Euclidean.Formulation}\leanok
  \texttt{Common.QFT.Euclidean.Formulation}
\end{definition}
\begin{theorem}[\texttt{instContinuousSMulNPointTestFunction}]\label{Common-QFT-Euclidean-SchwingerFormulation-instContinuousSMulNPointTestFunction}
  \lean{Common.QFT.Euclidean.SchwingerFormulation.instContinuousSMulNPointTestFunction}\leanok
  \uses{Common-QFT-Euclidean-SchwingerFormulation-instAddCommGroupNPointTestFunction, Common-QFT-Euclidean-SchwingerFormulation-NPointTestFunction, Common-QFT-Euclidean-SchwingerFormulation-instTopologicalSpaceNPointTestFunction, Common-QFT-Euclidean-SchwingerFormulation-instModuleNPointTestFunction, Common-QFT-Euclidean-SchwingerFormulation}
  \texttt{Common.QFT.Euclidean.SchwingerFormulation.instContinuousSMulNPointTestFunction}
\end{theorem}
\begin{definition}[\texttt{TensorToDistributionalBridge}]\label{Common-QFT-Euclidean-TensorToDistributionalBridge}
  \lean{Common.QFT.Euclidean.TensorToDistributionalBridge}\leanok
  \uses{Common-QFT-Euclidean-TensorSchwingerModel, Common-QFT-Euclidean-SchwingerFormulation-toFormulation, Common-QFT-Euclidean-SchwingerFormulation}
  \texttt{Common.QFT.Euclidean.TensorToDistributionalBridge}
\end{definition}
\begin{definition}[\texttt{instMeasurableSpaceFieldConfig}]\label{Common-QFT-Euclidean-PairingMeasureModel-instMeasurableSpaceFieldConfig}
  \lean{Common.QFT.Euclidean.PairingMeasureModel.instMeasurableSpaceFieldConfig}\leanok
  \uses{Common-QFT-Euclidean-PairingMeasureModel, Common-QFT-Euclidean-Formulation}
  \texttt{Common.QFT.Euclidean.PairingMeasureModel.instMeasurableSpaceFieldConfig}
\end{definition}
\begin{theorem}[\texttt{realToComplex\_smul}]\label{Common-QFT-Euclidean-SchwingerFormulation-realToComplex-smul}
  \lean{Common.QFT.Euclidean.SchwingerFormulation.realToComplex_smul}\leanok
  \uses{Common-QFT-Euclidean-SchwingerFormulation-instAddCommGroupComplexTestFunction, Common-QFT-Euclidean-SchwingerFormulation-instModuleComplexTestFunction, Common-QFT-Euclidean-Formulation-TestFunction, Common-QFT-Euclidean-Formulation-ComplexTestFunction, Common-QFT-Euclidean-SchwingerFormulation-toFormulation, Common-QFT-Euclidean-SchwingerFormulation-realToComplex, Common-QFT-Euclidean-SchwingerFormulation-instAddCommGroupTestFunction, Common-QFT-Euclidean-SchwingerFormulation-instModuleTestFunction, Common-QFT-Euclidean-SchwingerFormulation}
  \texttt{Common.QFT.Euclidean.SchwingerFormulation.realToComplex\_smul}
\end{theorem}
\begin{definition}[\texttt{ctorIdx}]\label{Common-QFT-Euclidean-ReconstructionInput-ctorIdx}
  \lean{Common.QFT.Euclidean.ReconstructionInput.ctorIdx}\leanok
  \uses{Common-QFT-Euclidean-ReconstructionHypothesis, Common-QFT-Euclidean-ReconstructionInput, Common-QFT-Euclidean-SchwingerFormulation}
  \texttt{Common.QFT.Euclidean.ReconstructionInput.ctorIdx}
\end{definition}
\begin{definition}[\texttt{instAddCommGroupComplexTestFunction}]\label{Common-QFT-Euclidean-SchwingerFormulation-instAddCommGroupComplexTestFunction}
  \lean{Common.QFT.Euclidean.SchwingerFormulation.instAddCommGroupComplexTestFunction}\leanok
  \uses{Common-QFT-Euclidean-Formulation-ComplexTestFunction, Common-QFT-Euclidean-SchwingerFormulation-toFormulation, Common-QFT-Euclidean-SchwingerFormulation}
  \texttt{Common.QFT.Euclidean.SchwingerFormulation.instAddCommGroupComplexTestFunction}
\end{definition}
\begin{definition}[\texttt{ctorIdx}]\label{Common-QFT-Euclidean-MeasureToTensorBridge-ctorIdx}
  \lean{Common.QFT.Euclidean.MeasureToTensorBridge.ctorIdx}\leanok
  \uses{Common-QFT-Euclidean-PairingMeasureModel, Common-QFT-Euclidean-Formulation, Common-QFT-Euclidean-MeasureToTensorBridge}
  \texttt{Common.QFT.Euclidean.MeasureToTensorBridge.ctorIdx}
\end{definition}
\begin{definition}[\texttt{distributionalSchwinger}]\label{Common-QFT-Euclidean-TensorToDistributionalBridge-distributionalSchwinger}
  \lean{Common.QFT.Euclidean.TensorToDistributionalBridge.distributionalSchwinger}\leanok
  \uses{Common-QFT-Euclidean-DistributionalSchwingerModel, Common-QFT-Euclidean-TensorSchwingerModel, Common-QFT-Euclidean-TensorToDistributionalBridge, Common-QFT-Euclidean-SchwingerFormulation-toFormulation, Common-QFT-Euclidean-SchwingerFormulation}
  \texttt{Common.QFT.Euclidean.TensorToDistributionalBridge.distributionalSchwinger}
\end{definition}
\begin{definition}[\texttt{tensorSchwinger}]\label{Common-QFT-Euclidean-MeasureToTensorBridge-tensorSchwinger}
  \lean{Common.QFT.Euclidean.MeasureToTensorBridge.tensorSchwinger}\leanok
  \uses{Common-QFT-Euclidean-PairingMeasureModel, Common-QFT-Euclidean-TensorSchwingerModel, Common-QFT-Euclidean-Formulation, Common-QFT-Euclidean-MeasureToTensorBridge}
  \texttt{Common.QFT.Euclidean.MeasureToTensorBridge.tensorSchwinger}
\end{definition}
\begin{theorem}[\texttt{instIsTopologicalAddGroupNPointTestFunction}]\label{Common-QFT-Euclidean-SchwingerFormulation-instIsTopologicalAddGroupNPointTestFunction}
  \lean{Common.QFT.Euclidean.SchwingerFormulation.instIsTopologicalAddGroupNPointTestFunction}\leanok
  \uses{Common-QFT-Euclidean-SchwingerFormulation-instAddCommGroupNPointTestFunction, Common-QFT-Euclidean-SchwingerFormulation-NPointTestFunction, Common-QFT-Euclidean-SchwingerFormulation-instTopologicalSpaceNPointTestFunction, Common-QFT-Euclidean-SchwingerFormulation}
  \texttt{Common.QFT.Euclidean.SchwingerFormulation.instIsTopologicalAddGroupNPointTestFunction}
\end{theorem}
\begin{theorem}[\texttt{compatible}]\label{Common-QFT-Euclidean-MeasureToTensorBridge-compatible}
  \lean{Common.QFT.Euclidean.MeasureToTensorBridge.compatible}\leanok
  \uses{Common-QFT-Euclidean-MeasureToTensorBridge-tensorSchwinger, Common-QFT-Euclidean-PairingMeasureModel, Common-QFT-Euclidean-Formulation-TestFunction, Common-QFT-Euclidean-Formulation, Common-QFT-Euclidean-MeasureToTensorBridge, Common-QFT-Euclidean-PairingMeasureModel-pairing, Common-QFT-Euclidean-TensorSchwingerModel-schwinger, Common-QFT-Euclidean-PairingMeasureModel-measure, Common-QFT-Euclidean-PairingMeasureModel-instMeasurableSpaceFieldConfig}
  \texttt{Common.QFT.Euclidean.MeasureToTensorBridge.compatible}
\end{theorem}
\begin{definition}[\texttt{instTopologicalSpaceNPointTestFunction}]\label{Common-QFT-Euclidean-SchwingerFormulation-instTopologicalSpaceNPointTestFunction}
  \lean{Common.QFT.Euclidean.SchwingerFormulation.instTopologicalSpaceNPointTestFunction}\leanok
  \uses{Common-QFT-Euclidean-SchwingerFormulation-NPointTestFunction, Common-QFT-Euclidean-SchwingerFormulation}
  \texttt{Common.QFT.Euclidean.SchwingerFormulation.instTopologicalSpaceNPointTestFunction}
\end{definition}
\begin{definition}[\texttt{ctorIdx}]\label{Common-QFT-Euclidean-DistributionalSchwingerModel-ctorIdx}
  \lean{Common.QFT.Euclidean.DistributionalSchwingerModel.ctorIdx}\leanok
  \uses{Common-QFT-Euclidean-DistributionalSchwingerModel, Common-QFT-Euclidean-SchwingerFormulation}
  \texttt{Common.QFT.Euclidean.DistributionalSchwingerModel.ctorIdx}
\end{definition}
\begin{definition}[\texttt{ctorIdx}]\label{Common-QFT-Euclidean-SchwingerFormulation-ctorIdx}
  \lean{Common.QFT.Euclidean.SchwingerFormulation.ctorIdx}\leanok
  \uses{Common-QFT-Euclidean-SchwingerFormulation}
  \texttt{Common.QFT.Euclidean.SchwingerFormulation.ctorIdx}
\end{definition}
\begin{definition}[\texttt{ctorIdx}]\label{Common-QFT-Euclidean-GeneratingFunctionalModel-ctorIdx}
  \lean{Common.QFT.Euclidean.GeneratingFunctionalModel.ctorIdx}\leanok
  \uses{Common-QFT-Euclidean-GeneratingFunctionalModel, Common-QFT-Euclidean-Formulation}
  \texttt{Common.QFT.Euclidean.GeneratingFunctionalModel.ctorIdx}
\end{definition}
\begin{definition}[\texttt{pairing}]\label{Common-QFT-Euclidean-PairingMeasureModel-pairing}
  \lean{Common.QFT.Euclidean.PairingMeasureModel.pairing}\leanok
  \uses{Common-QFT-Euclidean-PairingMeasureModel, Common-QFT-Euclidean-Formulation-TestFunction, Common-QFT-Euclidean-Formulation}
  \texttt{Common.QFT.Euclidean.PairingMeasureModel.pairing}
\end{definition}
\begin{theorem}[\texttt{instIsTopologicalAddGroupComplexTestFunction}]\label{Common-QFT-Euclidean-SchwingerFormulation-instIsTopologicalAddGroupComplexTestFunction}
  \lean{Common.QFT.Euclidean.SchwingerFormulation.instIsTopologicalAddGroupComplexTestFunction}\leanok
  \uses{Common-QFT-Euclidean-SchwingerFormulation-instAddCommGroupComplexTestFunction, Common-QFT-Euclidean-Formulation-ComplexTestFunction, Common-QFT-Euclidean-SchwingerFormulation-toFormulation, Common-QFT-Euclidean-SchwingerFormulation, Common-QFT-Euclidean-SchwingerFormulation-instTopologicalSpaceComplexTestFunction}
  \texttt{Common.QFT.Euclidean.SchwingerFormulation.instIsTopologicalAddGroupComplexTestFunction}
\end{theorem}
\begin{definition}[\texttt{PositiveTimeTestFunction}]\label{Common-QFT-Euclidean-Formulation-PositiveTimeTestFunction}
  \lean{Common.QFT.Euclidean.Formulation.PositiveTimeTestFunction}\leanok
  \uses{Common-QFT-Euclidean-Formulation-positiveTimeTestFunctions, Common-QFT-Euclidean-Formulation-TestFunction, Common-QFT-Euclidean-Formulation}
  \texttt{Common.QFT.Euclidean.Formulation.PositiveTimeTestFunction}
\end{definition}
\begin{definition}[\texttt{GeneratingFunctionalModel}]\label{Common-QFT-Euclidean-GeneratingFunctionalModel}
  \lean{Common.QFT.Euclidean.GeneratingFunctionalModel}\leanok
  \uses{Common-QFT-Euclidean-Formulation}
  \texttt{Common.QFT.Euclidean.GeneratingFunctionalModel}
\end{definition}
\begin{definition}[\texttt{ctorIdx}]\label{Common-QFT-Euclidean-TensorToDistributionalBridge-ctorIdx}
  \lean{Common.QFT.Euclidean.TensorToDistributionalBridge.ctorIdx}\leanok
  \uses{Common-QFT-Euclidean-TensorSchwingerModel, Common-QFT-Euclidean-TensorToDistributionalBridge, Common-QFT-Euclidean-SchwingerFormulation-toFormulation, Common-QFT-Euclidean-SchwingerFormulation}
  \texttt{Common.QFT.Euclidean.TensorToDistributionalBridge.ctorIdx}
\end{definition}
\begin{definition}[\texttt{ctorIdx}]\label{Common-QFT-Euclidean-TensorSchwingerModel-ctorIdx}
  \lean{Common.QFT.Euclidean.TensorSchwingerModel.ctorIdx}\leanok
  \uses{Common-QFT-Euclidean-TensorSchwingerModel, Common-QFT-Euclidean-Formulation}
  \texttt{Common.QFT.Euclidean.TensorSchwingerModel.ctorIdx}
\end{definition}
\begin{definition}[\texttt{ReconstructionHypothesis}]\label{Common-QFT-Euclidean-ReconstructionHypothesis}
  \lean{Common.QFT.Euclidean.ReconstructionHypothesis}\leanok
  \uses{Common-QFT-Euclidean-DistributionalSchwingerModel, Common-QFT-Euclidean-SchwingerFormulation}
  \texttt{Common.QFT.Euclidean.ReconstructionHypothesis}
\end{definition}
\begin{definition}[\texttt{realToComplex}]\label{Common-QFT-Euclidean-SchwingerFormulation-realToComplex}
  \lean{Common.QFT.Euclidean.SchwingerFormulation.realToComplex}\leanok
  \uses{Common-QFT-Euclidean-Formulation-TestFunction, Common-QFT-Euclidean-Formulation-ComplexTestFunction, Common-QFT-Euclidean-SchwingerFormulation-toFormulation, Common-QFT-Euclidean-SchwingerFormulation}
  \texttt{Common.QFT.Euclidean.SchwingerFormulation.realToComplex}
\end{definition}
\begin{definition}[\texttt{SchwingerFormulation}]\label{Common-QFT-Euclidean-SchwingerFormulation}
  \lean{Common.QFT.Euclidean.SchwingerFormulation}\leanok
  \texttt{Common.QFT.Euclidean.SchwingerFormulation}
\end{definition}
\begin{definition}[\texttt{schwinger}]\label{Common-QFT-Euclidean-DistributionalSchwingerModel-schwinger}
  \lean{Common.QFT.Euclidean.DistributionalSchwingerModel.schwinger}\leanok
  \uses{Common-QFT-Euclidean-SchwingerFormulation-instAddCommGroupNPointTestFunction, Common-QFT-Euclidean-DistributionalSchwingerModel, Common-QFT-Euclidean-SchwingerFormulation-NPointTestFunction, Common-QFT-Euclidean-SchwingerFormulation-instModuleNPointTestFunction, Common-QFT-Euclidean-SchwingerFormulation}
  \texttt{Common.QFT.Euclidean.DistributionalSchwingerModel.schwinger}
\end{definition}
\begin{definition}[\texttt{tensorLift}]\label{Common-QFT-Euclidean-DistributionalSchwingerModel-tensorLift}
  \lean{Common.QFT.Euclidean.DistributionalSchwingerModel.tensorLift}\leanok
  \uses{Common-QFT-Euclidean-DistributionalSchwingerModel, Common-QFT-Euclidean-SchwingerFormulation-NPointTestFunction, Common-QFT-Euclidean-Formulation-TestFunction, Common-QFT-Euclidean-SchwingerFormulation-toFormulation, Common-QFT-Euclidean-SchwingerFormulation}
  \texttt{Common.QFT.Euclidean.DistributionalSchwingerModel.tensorLift}
\end{definition}
\begin{definition}[\texttt{instAddCommGroupTestFunction}]\label{Common-QFT-Euclidean-SchwingerFormulation-instAddCommGroupTestFunction}
  \lean{Common.QFT.Euclidean.SchwingerFormulation.instAddCommGroupTestFunction}\leanok
  \uses{Common-QFT-Euclidean-Formulation-TestFunction, Common-QFT-Euclidean-SchwingerFormulation-toFormulation, Common-QFT-Euclidean-SchwingerFormulation}
  \texttt{Common.QFT.Euclidean.SchwingerFormulation.instAddCommGroupTestFunction}
\end{definition}
\begin{definition}[\texttt{generatingFunctional}]\label{Common-QFT-Euclidean-GeneratingFunctionalModel-generatingFunctional}
  \lean{Common.QFT.Euclidean.GeneratingFunctionalModel.generatingFunctional}\leanok
  \uses{Common-QFT-Euclidean-Formulation-TestFunction, Common-QFT-Euclidean-GeneratingFunctionalModel, Common-QFT-Euclidean-Formulation}
  $Z[f] = \int e^{i\omega(f)}\,d\mu$ (real) and $Z[J] = \int e^{i\langle\omega,J\rangle_\mathbb{C}}\,d\mu$ (complex).
\end{definition}
\begin{theorem}[\texttt{instIsTopologicalAddGroupTestFunction}]\label{Common-QFT-Euclidean-SchwingerFormulation-instIsTopologicalAddGroupTestFunction}
  \lean{Common.QFT.Euclidean.SchwingerFormulation.instIsTopologicalAddGroupTestFunction}\leanok
  \uses{Common-QFT-Euclidean-SchwingerFormulation-instTopologicalSpaceTestFunction, Common-QFT-Euclidean-Formulation-TestFunction, Common-QFT-Euclidean-SchwingerFormulation-toFormulation, Common-QFT-Euclidean-SchwingerFormulation-instAddCommGroupTestFunction, Common-QFT-Euclidean-SchwingerFormulation}
  \texttt{Common.QFT.Euclidean.SchwingerFormulation.instIsTopologicalAddGroupTestFunction}
\end{theorem}
\begin{theorem}[\texttt{ext\_iff}]\label{Common-QFT-Euclidean-TensorSchwingerModel-ext-iff}
  \lean{Common.QFT.Euclidean.TensorSchwingerModel.ext_iff}\leanok
  \uses{Common-QFT-Euclidean-TensorSchwingerModel-ext, Common-QFT-Euclidean-Formulation-TestFunction, Common-QFT-Euclidean-TensorSchwingerModel, Common-QFT-Euclidean-Formulation, Common-QFT-Euclidean-TensorSchwingerModel-schwinger}
  \texttt{Common.QFT.Euclidean.TensorSchwingerModel.ext\_iff}
\end{theorem}
\begin{definition}[\texttt{ctorIdx}]\label{Common-QFT-Euclidean-MeasureModel-ctorIdx}
  \lean{Common.QFT.Euclidean.MeasureModel.ctorIdx}\leanok
  \uses{Common-QFT-Euclidean-MeasureModel, Common-QFT-Euclidean-Formulation}
  \texttt{Common.QFT.Euclidean.MeasureModel.ctorIdx}
\end{definition}
\chapter{gaussian-field: lattice (upstream)}
\section{\texttt{Mathlib.Order.Lattice}}
\begin{theorem}[\texttt{le\_inf}]\label{Lattice-le-inf}
  \lean{Lattice.le_inf}\leanok
  \uses{Lattice, Lattice-toSemilatticeSup, Lattice-inf}
  \texttt{Lattice.le\_inf}
\end{theorem}
\begin{definition}[\texttt{toSemilatticeInf}]\label{Lattice-toSemilatticeInf}
  \lean{Lattice.toSemilatticeInf}\leanok
  \uses{Lattice, Lattice-toSemilatticeSup, Lattice-inf-le-left, Lattice-le-inf, Lattice-inf, Lattice-inf-le-right}
  \texttt{Lattice.toSemilatticeInf}
\end{definition}
\begin{definition}[\texttt{ctorIdx}]\label{Lattice-ctorIdx}
  \lean{Lattice.ctorIdx}\leanok
  \uses{Lattice}
  \texttt{Lattice.ctorIdx}
\end{definition}
\begin{theorem}[\texttt{inf\_le\_left}]\label{Lattice-inf-le-left}
  \lean{Lattice.inf_le_left}\leanok
  \uses{Lattice, Lattice-toSemilatticeSup, Lattice-inf}
  \texttt{Lattice.inf\_le\_left}
\end{theorem}
\begin{theorem}[\texttt{ext}]\label{Lattice-ext}
  \lean{Lattice.ext}\leanok
  \uses{Lattice, Lattice-toSemilatticeSup, Lattice-inf-le-left, Lattice-le-inf, Lattice-inf, Lattice-toSemilatticeInf, Lattice-inf-le-right}
  \texttt{Lattice.ext}
\end{theorem}
\begin{definition}[\texttt{Lattice}]\label{Lattice}
  \lean{Lattice}\leanok
  \texttt{Lattice}
\end{definition}
\begin{theorem}[\texttt{inf\_le\_right}]\label{Lattice-inf-le-right}
  \lean{Lattice.inf_le_right}\leanok
  \uses{Lattice, Lattice-toSemilatticeSup, Lattice-inf}
  \texttt{Lattice.inf\_le\_right}
\end{theorem}
\begin{definition}[\texttt{toLinearOrder}]\label{Lattice-toLinearOrder}
  \lean{Lattice.toLinearOrder}\leanok
  \uses{Lattice, Lattice-toSemilatticeSup, Lattice-toSemilatticeInf}
  \texttt{Lattice.toLinearOrder}
\end{definition}
\begin{definition}[\texttt{mk'}]\label{Lattice-mk-}
  \lean{Lattice.mk'}\leanok
  \uses{Lattice}
  \texttt{Lattice.mk'}
\end{definition}
\begin{definition}[\texttt{toSemilatticeSup}]\label{Lattice-toSemilatticeSup}
  \lean{Lattice.toSemilatticeSup}\leanok
  \uses{Lattice}
  \texttt{Lattice.toSemilatticeSup}
\end{definition}
\begin{definition}[\texttt{inf}]\label{Lattice-inf}
  \lean{Lattice.inf}\leanok
  \uses{Lattice}
  \texttt{Lattice.inf}
\end{definition}
\section{\texttt{Mathlib.Order.Copy}}
\begin{theorem}[\texttt{congr\_simp}]\label{Lattice-copy-congr-simp}
  \lean{Lattice.copy.congr_simp}\leanok
  \uses{Lattice, Lattice-toSemilatticeSup, Lattice-copy, Lattice-toSemilatticeInf}
  \texttt{Lattice.copy.congr\_simp}
\end{theorem}
\begin{definition}[\texttt{copy}]\label{Lattice-copy}
  \lean{Lattice.copy}\leanok
  \uses{Lattice, Lattice-toSemilatticeSup, Lattice-toSemilatticeInf}
  \texttt{Lattice.copy}
\end{definition}
\section{\texttt{Mathlib.Order.Atoms}}
\begin{theorem}[\texttt{isStronglyAtomic}]\label{Lattice-isStronglyAtomic}
  \lean{Lattice.isStronglyAtomic}\leanok
  \uses{Lattice, Lattice-toSemilatticeSup, Lattice-toSemilatticeInf}
  \texttt{Lattice.isStronglyAtomic}
\end{theorem}
\begin{theorem}[\texttt{isStronglyCoatomic}]\label{Lattice-isStronglyCoatomic}
  \lean{Lattice.isStronglyCoatomic}\leanok
  \uses{Lattice, Lattice-isStronglyAtomic, Lattice-toSemilatticeInf}
  \texttt{Lattice.isStronglyCoatomic}
\end{theorem}
\section{\texttt{Mathlib.Order.Bounds.Basic}}
\begin{definition}[\texttt{ofIsLUBofIsGLB}]\label{Lattice-ofIsLUBofIsGLB}
  \lean{Lattice.ofIsLUBofIsGLB}\leanok
  \uses{Lattice}
  \texttt{Lattice.ofIsLUBofIsGLB}
\end{definition}
